% ─────────────────────────────────────────────────────────────────────────────
% Errores pendientes de OpenMarkov
% Compilar con:  lualatex pending-errors.tex   (dos pasadas, por el índice)
% ─────────────────────────────────────────────────────────────────────────────

\documentclass[11pt,a4paper]{article}

\usepackage{fontspec}
\directlua{luaotfload.add_fallback("omfallback", {"DejaVuSerif:mode=node;", "DejaVuSans:mode=node;"})}
\setmainfont{Latin Modern Roman}[RawFeature={fallback=omfallback}]
\setsansfont{Latin Modern Sans}
\setmonofont{DejaVu Sans Mono}[Scale=0.84,RawFeature={fallback=omfallback}]
\usepackage[spanish,es-nodecimaldot,es-noshorthands]{babel}
\addto\captionsspanish{\renewcommand{\tablename}{Tabla}}
\usepackage[left=2.4cm,right=2.4cm,top=2.5cm,bottom=2.5cm]{geometry}
\usepackage{booktabs}
\usepackage{longtable}
\usepackage{array}
\usepackage{enumitem}
\usepackage{xcolor}
\usepackage{mdframed}
\usepackage{microtype}
\usepackage{parskip}
\usepackage{titlesec}
\usepackage{fancyhdr}
\usepackage{xurl}
\usepackage[hidelinks]{hyperref}
\usepackage{tocloft}
\usepackage{etoolbox}
\usepackage{calc}

\emergencystretch=3em
\setlist{itemsep=2pt,topsep=4pt}

% ── Código resaltado (pandoc) con cortes de línea ────────────────────────────
\usepackage{color}
\usepackage{fancyvrb}
\newcommand{\VerbBar}{|}
\newcommand{\VERB}{\Verb[commandchars=\\\{\}]}
\DefineVerbatimEnvironment{Highlighting}{Verbatim}{commandchars=\\\{\}}
% Add ',fontsize=\small' for more characters per line
\usepackage{framed}
\definecolor{shadecolor}{RGB}{248,248,248}
\newenvironment{Shaded}{\begin{snugshade}}{\end{snugshade}}
\newcommand{\AlertTok}[1]{\textcolor[rgb]{0.94,0.16,0.16}{#1}}
\newcommand{\AnnotationTok}[1]{\textcolor[rgb]{0.56,0.35,0.01}{\textbf{\textit{#1}}}}
\newcommand{\AttributeTok}[1]{\textcolor[rgb]{0.13,0.29,0.53}{#1}}
\newcommand{\BaseNTok}[1]{\textcolor[rgb]{0.00,0.00,0.81}{#1}}
\newcommand{\BuiltInTok}[1]{#1}
\newcommand{\CharTok}[1]{\textcolor[rgb]{0.31,0.60,0.02}{#1}}
\newcommand{\CommentTok}[1]{\textcolor[rgb]{0.56,0.35,0.01}{\textit{#1}}}
\newcommand{\CommentVarTok}[1]{\textcolor[rgb]{0.56,0.35,0.01}{\textbf{\textit{#1}}}}
\newcommand{\ConstantTok}[1]{\textcolor[rgb]{0.56,0.35,0.01}{#1}}
\newcommand{\ControlFlowTok}[1]{\textcolor[rgb]{0.13,0.29,0.53}{\textbf{#1}}}
\newcommand{\DataTypeTok}[1]{\textcolor[rgb]{0.13,0.29,0.53}{#1}}
\newcommand{\DecValTok}[1]{\textcolor[rgb]{0.00,0.00,0.81}{#1}}
\newcommand{\DocumentationTok}[1]{\textcolor[rgb]{0.56,0.35,0.01}{\textbf{\textit{#1}}}}
\newcommand{\ErrorTok}[1]{\textcolor[rgb]{0.64,0.00,0.00}{\textbf{#1}}}
\newcommand{\ExtensionTok}[1]{#1}
\newcommand{\FloatTok}[1]{\textcolor[rgb]{0.00,0.00,0.81}{#1}}
\newcommand{\FunctionTok}[1]{\textcolor[rgb]{0.13,0.29,0.53}{\textbf{#1}}}
\newcommand{\ImportTok}[1]{#1}
\newcommand{\InformationTok}[1]{\textcolor[rgb]{0.56,0.35,0.01}{\textbf{\textit{#1}}}}
\newcommand{\KeywordTok}[1]{\textcolor[rgb]{0.13,0.29,0.53}{\textbf{#1}}}
\newcommand{\NormalTok}[1]{#1}
\newcommand{\OperatorTok}[1]{\textcolor[rgb]{0.81,0.36,0.00}{\textbf{#1}}}
\newcommand{\OtherTok}[1]{\textcolor[rgb]{0.56,0.35,0.01}{#1}}
\newcommand{\PreprocessorTok}[1]{\textcolor[rgb]{0.56,0.35,0.01}{\textit{#1}}}
\newcommand{\RegionMarkerTok}[1]{#1}
\newcommand{\SpecialCharTok}[1]{\textcolor[rgb]{0.81,0.36,0.00}{\textbf{#1}}}
\newcommand{\SpecialStringTok}[1]{\textcolor[rgb]{0.31,0.60,0.02}{#1}}
\newcommand{\StringTok}[1]{\textcolor[rgb]{0.31,0.60,0.02}{#1}}
\newcommand{\VariableTok}[1]{\textcolor[rgb]{0.00,0.00,0.00}{#1}}
\newcommand{\VerbatimStringTok}[1]{\textcolor[rgb]{0.31,0.60,0.02}{#1}}
\newcommand{\WarningTok}[1]{\textcolor[rgb]{0.56,0.35,0.01}{\textbf{\textit{#1}}}}

\usepackage{fvextra}
\RecustomVerbatimEnvironment{Highlighting}{Verbatim}{commandchars=\\\{\},breaklines,breakanywhere,fontsize=\small}
\definecolor{shadecolor}{RGB}{246,246,246}

% pandoc emite \tightlist en las listas compactas
\providecommand{\tightlist}{\setlength{\itemsep}{0pt}\setlength{\parskip}{0pt}}

% ── Colores ──────────────────────────────────────────────────────────────────
\definecolor{fichafondo}{RGB}{245,247,250}
\definecolor{fichaborde}{RGB}{170,180,195}
\definecolor{notafondo}{RGB}{255,252,225}
\definecolor{notaborde}{RGB}{200,175,60}
\definecolor{decisionrojo}{RGB}{160,40,40}
\definecolor{corregidoverde}{RGB}{30,120,50}

\newmdenv[
  backgroundcolor=notafondo, linecolor=notaborde, linewidth=0.8pt,
  innertopmargin=8pt, innerbottommargin=8pt, innerleftmargin=10pt, innerrightmargin=10pt,
  skipabove=8pt, skipbelow=8pt
]{nota}

\usepackage[most]{tcolorbox}
% Ficha de cada error: pares «clave valor» con sangría francesa, en un recuadro que puede partirse.
\newtcolorbox{ficha}{breakable, enhanced, colback=fichafondo, colframe=fichaborde, boxrule=0.6pt,
  arc=0pt, left=8pt, right=8pt, top=4pt, bottom=4pt, before skip=6pt, after skip=8pt,
  fontupper=\small, parbox=false, before upper={\setlength{\parskip}{2.5pt}}}
\newcommand{\fichakey}[1]{\par\noindent\hangindent=3.6cm\hangafter=1\makebox[3.6cm][l]{\bfseries #1}\ignorespaces}

% ── Numeración: un número por error, corrido en todo el documento ─────────────
\counterwithout{subsection}{section}
\renewcommand{\thesubsection}{\arabic{subsection}}

\titleformat{\section}{\Large\bfseries}{\thesection\quad}{0em}{}[\titlerule]
\titleformat{\subsection}{\large\bfseries}{Error~\thesubsection\enspace·\enspace}{0em}{}
\titlespacing*{\subsection}{0pt}{2.6ex plus 1ex minus .2ex}{1.2ex plus .3ex}
\titleformat{\paragraph}[runin]{\normalsize\bfseries}{}{0em}{}[.]
\titlespacing*{\paragraph}{0pt}{1.2ex plus .5ex}{0.6em}
\setcounter{secnumdepth}{2}
\setcounter{tocdepth}{2}

\setlength{\cftsecnumwidth}{2.2em}
\setlength{\cftsubsecnumwidth}{2.6em}
\renewcommand{\cftsubsecpresnum}{}

\pagestyle{fancy}
\fancyhf{}
\fancyhead[L]{\small Errores pendientes de OpenMarkov}
\fancyhead[R]{\small\thepage}
\renewcommand{\headrulewidth}{0.4pt}

\newcommand{\clase}[1]{\texttt{#1}}
\newcommand{\err}[1]{error~\ref{#1}}


\title{\textbf{Errores pendientes de OpenMarkov}\\[0.4em]
       \large Lista unificada, comprobada contra el código y ordenada por temática}
\author{}
\date{14 de septiembre de 2026; revisado el 22 de septiembre de 2026; correcciones marcadas el 23 de septiembre de 2026}

\begin{document}
\maketitle
\thispagestyle{fancy}

\begin{nota}
\textbf{Qué es este documento.} La lista de los errores de OpenMarkov que siguen sin corregir,
reunida a partir de las revisiones de código hechas en verano de 2026 y de los defectos nuevos que
aparecieron al comprobarlas. Cada error se ha vuelto a comprobar contra el código de hoy, y se describe con lo necesario para ponerse a arreglarlo sin
haber leído nada más: qué ve el usuario, por dónde se llega, en qué clase y en qué líneas está, y
por dónde podría ir el arreglo.

\textbf{Qué no es.} No es un plan de trabajo ni un reparto de tareas, y no incluye las propuestas
de rediseño. Lo que se comprobó y no es un error por corregir (porque ya está corregido, porque el
código hace lo correcto, porque nada llega a él o porque una decisión del equipo lo deja fuera)
está en el apéndice~\ref{app:notas}, para que no haga falta volver a mirarlo.
\end{nota}

\tableofcontents

% ═════════════════════════════════════════════════════════════════════════════
\section*{Cómo leer este documento}
\addcontentsline{toc}{section}{Cómo leer este documento}

\paragraph{Qué código se ha comprobado}
El de la rama \texttt{development} tal como queda al unir los últimos commits locales (hasta
\texttt{6260a4c}, del 5 de septiembre) con los últimos de \texttt{origin/development} (hasta
\texttt{fb19d55}). El 22 de septiembre se revisó la lista contra \texttt{40a9de9}, que añade diez
commits: se volvió a comprobar cada error cuyo código o cuyo camino tocan esos commits. Uno quedó
corregido y ha salido de la lista (\hyperref[nota:184]{nota~184}); tres están ahora corregidos en
parte (\hyperref[err:29]{errores~29}, \hyperref[err:43]{43} y \hyperref[err:99]{99}). Los errores
conservan su número, así que falta el 32. En las fichas revisadas, las líneas de la fila
\emph{Ficheros} son las de \texttt{40a9de9}; las que se citan en el texto pueden ser todavía las
del código anterior. Con el tiempo se irán desplazando, así que conviene buscar por el nombre del
método.

\paragraph{Correcciones marcadas el 23 de septiembre}
Los catorce commits que van de \texttt{bbde667} a \texttt{522d926} corrigen errores de esta lista.
Esos errores no salen de ella: se quedan en su sitio con la marca
\textcolor{corregidoverde}{\textbf{Corregido}} o \textcolor{corregidoverde}{\textbf{Corregido en parte}}
en la tabla del resumen y en su título, y la fila \emph{Estado} de su ficha dice qué commit los
corrige y, si la corrección es parcial, qué sigue presente. Corregidos del todo:
\hyperref[err:1]{errores~1}, \hyperref[err:2]{2}, \hyperref[err:3]{3}, \hyperref[err:4]{4},
\hyperref[err:6]{6}, \hyperref[err:7]{7}, \hyperref[err:8]{8}, \hyperref[err:10]{10},
\hyperref[err:11]{11}, \hyperref[err:12]{12}, \hyperref[err:13]{13} y \hyperref[err:14]{14}.
Corregidos en parte: \hyperref[err:5]{errores~5}, \hyperref[err:9]{9} y \hyperref[err:157]{157}.
El resto de la ficha (qué pasa, el código, cómo arreglarlo) describe el código anterior a esos
commits. El error~1 se da por corregido sin tocar lo que quedaba de él, por lo que explica la
\hyperref[nota:185]{nota~185}; al revisarlo aparecieron el \hyperref[err:193]{error~193} y la
\hyperref[nota:186]{nota~186}.

\paragraph{Cómo se ha comprobado}
Para cada error se leyó el código de hoy, se buscaron todos los sitios que llaman a la parte
afectada y se estableció por dónde llega un usuario real (qué menú, qué diálogo, qué fichero) o un
llamador real del programa. Muchos se ejecutaron además con pequeños programas de prueba sobre el
jar de la aplicación, sin abrir la ventana. La fila \emph{Comprobado} de cada ficha dice si el
error se comprobó leyendo el código, ejecutándolo o de las dos maneras; cuando una consecuencia
solo se ha deducido leyendo, la descripción lo dice.

\paragraph{Cómo está hecha cada ficha}
Cada error tiene un número, corrido en todo el documento, y un título que dice la consecuencia.
Debajo va una ficha con:
\begin{itemize}
  \item \textbf{Estado}: presente, o presente en parte (y entonces qué parte ya está corregida).
  \item \textbf{Módulo}, \textbf{paquete} y \textbf{ficheros} con sus líneas, con las rutas
        relativas a la raíz del repositorio.
  \item \textbf{Esfuerzo} estimado: pequeño (horas), medio (días) o grande (semanas).
  \item \textbf{Decisión de equipo}: si el arreglo no se puede hacer sin que el equipo decida
        antes algo (qué debe hacer el programa, qué ficheros se aceptan…), qué hay que decidir.
        Estos errores conviene no empezarlos hasta tener la decisión.
  \item \textbf{Corregir antes}, solo si lo hay: el error que hay que corregir antes que este, o a
        la vez como mínimo, porque sin él este arreglo no basta o rompe otra cosa; y por qué.
  \item \textbf{Tapado por}, solo si lo hay: el error que hoy impide ver este, o comprobar su
        arreglo desde la aplicación, y dónde. Este sí se puede corregir sin esperar al otro.
  \item \textbf{Relacionados}: otros errores que tocan el mismo código o conviene arreglar a la
        vez, y notas del apéndice sobre lo mismo.
\end{itemize}
Después vienen tres apartados: \emph{Qué pasa} (la consecuencia, la causa en el código y lo que se
midió), \emph{El código} (extractos del código actual, cuando ayudan) y \emph{Cómo arreglarlo}
(por dónde ir, qué otros sitios deben cumplir la misma regla y qué prueba lo demostraría).

\paragraph{Decisiones de equipo ya tomadas que se han tenido en cuenta}
\begin{itemize}
  \item \textbf{La aplicación es solo en inglés.} No se añaden traducciones. Lo que solo afectaba a
        los textos en castellano, o a que un texto inglés esté escrito en el código en vez de en un
        fichero de textos, no entra en la lista; sí entra lo que se ve mal en la aplicación en inglés.
  \item \textbf{El formato Elvira está congelado}: no se arregla, no se migra y no se retira.
  \item \textbf{Dos cambios de la interfaz son intencionados}: reabrir al arrancar las redes de la
        última sesión, y crear los nodos con doble clic en el lienzo en vez de con botones.
\end{itemize}

% ═════════════════════════════════════════════════════════════════════════════
\section*{Resumen}
\addcontentsline{toc}{section}{Resumen}

La lista tiene \textbf{192 errores}; 48 de ellos son defectos nuevos, encontrados al
comprobar los demás. \textbf{45} necesitan antes una decisión del equipo. Por esfuerzo estimado: 156 pequeños, 36 medios y 0 grandes.
De los 192, 12 están ya corregidos y 12 lo están en parte.
El apéndice recoge además 186 notas con lo que se comprobó y no entra en la lista.

\begin{center}
\begin{tabular}{@{}lr@{}}
\toprule
\textbf{Tema} & \textbf{Errores} \\
\midrule
\ref{sec:potenciales}. Operaciones con potenciales & 3 \\
\ref{sec:canonicos}. Modelos canónicos y modelo de ajuste & 3 \\
\ref{sec:relaciones}. Relaciones definidas por fórmulas, regresiones, árboles y sucesos & 18 \\
\ref{sec:ediciones}. Ediciones de la red: deshacer, cancelar y pegar & 12 \\
\ref{sec:restricciones}. Restricciones de red & 3 \\
\ref{sec:inferencia}. Evidencia, inferencia y variables numéricas & 10 \\
\ref{sec:decisiones}. Decisiones: políticas, criterios, agentes y árbol de decisión & 17 \\
\ref{sec:temporal}. Modelos temporales y evolución temporal & 7 \\
\ref{sec:coste}. Coste-efectividad y análisis de sensibilidad & 21 \\
\ref{sec:des}. Simulación de sucesos discretos & 16 \\
\ref{sec:ficheros}. Abrir y guardar redes & 7 \\
\ref{sec:casos}. Bases de casos & 7 \\
\ref{sec:aprendizaje}. Aprendizaje y evaluación de redes & 25 \\
\ref{sec:editor}. Interfaz: editor de red y paneles de las relaciones & 21 \\
\ref{sec:menus}. Interfaz: menús, diálogos y herramientas & 6 \\
\ref{sec:textos}. Textos y títulos que se ven mal & 12 \\
\ref{sec:rendimiento}. Rendimiento & 4 \\
\bottomrule
\end{tabular}
\end{center}

La tabla siguiente es la lista completa, en el orden del documento. La última columna da el
esfuerzo (S, pequeño; M, medio; L, grande) y, en rojo, si hace falta una decisión de equipo.

{\small
\begin{longtable}{@{}>{\raggedleft\arraybackslash}p{0.7cm}>{\raggedright\arraybackslash}p{10.0cm}>{\raggedright\arraybackslash}p{2.8cm}>{\raggedright\arraybackslash}p{1.3cm}@{}}
\toprule
\textbf{Nº} & \textbf{Error} & \textbf{Módulo} & \textbf{Esf./Dec.} \\
\midrule
\endhead
\multicolumn{4}{@{}l}{\rule{0pt}{2.6ex}\textbf{\ref{sec:potenciales}.\ Operaciones con potenciales}}\\
\hyperref[err:1]{1} & \textcolor{corregidoverde}{\textbf{Corregido}} · Multiplicar potenciales constantes pierde el criterio de la utilidad & \texttt{core} & S \hfill No\\
\hyperref[err:2]{2} & \textcolor{corregidoverde}{\textbf{Corregido}} · Tablas de más de 2\textsuperscript{31} casillas: excepciones distintas o resultado falso & \texttt{core} & S \hfill No\\
\hyperref[err:3]{3} & \textcolor{corregidoverde}{\textbf{Corregido}} · Copiar un potencial a una red sin todas sus variables mete nulos & \texttt{core} & S \hfill No\\
\multicolumn{4}{@{}l}{\rule{0pt}{2.6ex}\textbf{\ref{sec:canonicos}.\ Modelos canónicos y modelo de ajuste}}\\
\hyperref[err:4]{4} & \textcolor{corregidoverde}{\textbf{Corregido}} · El modelo de ajuste se proyecta con otro orden y otros decimales cada vez & \texttt{core} & S \hfill No\\
\hyperref[err:5]{5} & \textcolor{corregidoverde}{\textbf{Corregido en parte}} · Parámetros que no suman uno se abren sin aviso y el muestreo discrepa & \texttt{core} & S \hfill \textcolor{decisionrojo}{Sí}\\
\hyperref[err:6]{6} & \textcolor{corregidoverde}{\textbf{Corregido}} · El modelo de ajuste se ofrece a políticas sin padres y en redes de sucesos & \texttt{core} & S \hfill No\\
\multicolumn{4}{@{}l}{\rule{0pt}{2.6ex}\textbf{\ref{sec:relaciones}.\ Relaciones definidas por fórmulas, regresiones, árboles y sucesos}}\\
\hyperref[err:7]{7} & \textcolor{corregidoverde}{\textbf{Corregido}} · Una utilidad por fórmula calcula mal si sus padres no se llaman U1, U2\ldots{} & \texttt{core} & S \hfill No\\
\hyperref[err:8]{8} & \textcolor{corregidoverde}{\textbf{Corregido}} · «Absorber padres» muere con ClassCastException en 12 redes del repositorio & \texttt{core} & S \hfill No\\
\hyperref[err:9]{9} & \textcolor{corregidoverde}{\textbf{Corregido en parte}} · «Absorb parents» de una suma falla si un padre es una distribución (formato 1.0) & \texttt{core} & M \hfill \textcolor{decisionrojo}{Sí}\\
\hyperref[err:10]{10} & \textcolor{corregidoverde}{\textbf{Corregido}} · Descontar una utilidad exponencial cambia el coeficiente equivocado & \texttt{core} & S \hfill No\\
\hyperref[err:11]{11} & \textcolor{corregidoverde}{\textbf{Corregido}} · Escalar o descontar una fórmula o un producto de utilidades da otro número & \texttt{core} & M \hfill \textcolor{decisionrojo}{Sí}\\
\hyperref[err:12]{12} & \textcolor{corregidoverde}{\textbf{Corregido}} · Una regresión no se calcula si un padre vale menos de una milésima & \texttt{core} & S \hfill No\\
\hyperref[err:13]{13} & \textcolor{corregidoverde}{\textbf{Corregido}} · Una covariable con el nombre sin llaves se acepta pero no se calcula nunca & \texttt{io} & M \hfill \textcolor{decisionrojo}{Sí}\\
\hyperref[err:14]{14} & \textcolor{corregidoverde}{\textbf{Corregido}} · Una regresión sin \texttt{Constant} o \texttt{Gamma} muere con un índice −1 sin mensaje & \texttt{core} & S \hfill No\\
\hyperref[err:15]{15} & Una utilidad de regresión no cuenta como utilidad al calcular & \texttt{core} & S \hfill No\\
\hyperref[err:16]{16} & Una tabla de fórmulas con \texttt{Complement} en vez de \texttt{\#} no se puede proyectar & \texttt{core} & S \hfill \textcolor{decisionrojo}{Sí}\\
\hyperref[err:17]{17} & Quitar una etiqueta referenciada del árbol deja la red sin poder evaluarse & \texttt{gui} & S \hfill \textcolor{decisionrojo}{Sí}\\
\hyperref[err:18]{18} & Añadir a una hoja del árbol la variable del nodo deja una red que no abre & \texttt{core} & M \hfill \textcolor{decisionrojo}{Sí}\\
\hyperref[err:19]{19} & Tras cancelar el potencial de un suceso, la simulación usa el primer padre & \texttt{core} & S \hfill No\\
\hyperref[err:20]{20} & Un nodo con tabla de transiciones copiado ya no sabe a qué estado pasa & \texttt{core} & S \hfill No\\
\hyperref[err:21]{21} & Los nodos de sucesos pegados siguen leyendo los padres del nodo original & \texttt{gui} & M \hfill No\\
\hyperref[err:22]{22} & Aceptar un árbol con sucesos deja un deshacer que borra la relación & \texttt{core} & S \hfill No\\
\hyperref[err:23]{23} & Se ofrece «Tree with Excluded Events» con solo padres suceso, y falla & \texttt{core} & S \hfill No\\
\hyperref[err:193]{193} & «Absorb parents» falla si un padre es una regresión & \texttt{core} & S \hfill No\\
\multicolumn{4}{@{}l}{\rule{0pt}{2.6ex}\textbf{\ref{sec:ediciones}.\ Ediciones de la red: deshacer, cancelar y pegar}}\\
\hyperref[err:24]{24} & Un dominio estándar de otro tamaño borra tablas que «Cancel» no devuelve & \texttt{core} & M \hfill No\\
\hyperref[err:25]{25} & Restringir un enlace cambia la tabla del hijo y deshacer no la devuelve & \texttt{gui} & S \hfill No\\
\hyperref[err:26]{26} & Enlazar a una utilidad producto la convierte en suma y cambia el resultado & \texttt{core} & S \hfill No\\
\hyperref[err:27]{27} & Rehacer un enlace y su restricción deja el enlace sin la restricción & \texttt{core} & M \hfill No\\
\hyperref[err:28]{28} & Pasar a suceso una decisión numérica sin política la deja a medias & \texttt{core} & S \hfill No\\
\hyperref[err:29]{29} & Mover un estado con un hijo sin potencial deja la tabla permutada a medias & \texttt{core} & S \hfill No\\
\hyperref[err:30]{30} & Deshacer un cambio de estados duplica tablas de los hijos y cambia resultados & \texttt{core} & S \hfill No\\
\hyperref[err:31]{31} & Pegar en una red que ya incumple una restricción se rechaza siempre & \texttt{core} & M \hfill \textcolor{decisionrojo}{Sí}\\
\hyperref[err:33]{33} & Deshacer o cancelar «mostrar el comentario al abrir» no la devuelve & \texttt{core} & S \hfill No\\
\hyperref[err:34]{34} & Deshacer o cancelar «Always append» de un suceso no devuelve su valor & \texttt{core} & S \hfill No\\
\hyperref[err:35]{35} & Una red expandida o aprendida no se marca como modificada ni pide guardar & \texttt{core} & S \hfill No\\
\hyperref[err:36]{36} & En una red expandida o aprendida, deshacer no recuerda nada al principio & \texttt{core} & S \hfill No\\
\multicolumn{4}{@{}l}{\rule{0pt}{2.6ex}\textbf{\ref{sec:restricciones}.\ Restricciones de red}}\\
\hyperref[err:37]{37} & Abrir un fichero no comprueba las restricciones de la red & \texttt{io} & S \hfill \textcolor{decisionrojo}{Sí}\\
\hyperref[err:38]{38} & Un estado puede quedarse con el nombre vacío o repetido & \texttt{core} & S \hfill \textcolor{decisionrojo}{Sí}\\
\hyperref[err:39]{39} & Ocho restricciones opcionales no llegan a ninguna red & \texttt{core} & M \hfill \textcolor{decisionrojo}{Sí}\\
\multicolumn{4}{@{}l}{\rule{0pt}{2.6ex}\textbf{\ref{sec:inferencia}.\ Evidencia, inferencia y variables numéricas}}\\
\hyperref[err:40]{40} & Una variable numérica sin precisión se discretiza a un único valor cero & \texttt{core} & S \hfill \textcolor{decisionrojo}{Sí}\\
\hyperref[err:41]{41} & Un hallazgo sobre un nodo numérico de azar hace fallar la propagación & \texttt{core} & M \hfill No\\
\hyperref[err:42]{42} & Entrar en modo inferencia falla si hay nodos de utilidad intermedios & \texttt{inference} & M \hfill \textcolor{decisionrojo}{Sí}\\
\hyperref[err:43]{43} & Un nodo numérico de azar sin potencial tumba la discretización & \texttt{core} & S \hfill No\\
\hyperref[err:44]{44} & Un hallazgo numérico con más de dos decimales se pierde al discretizar & \texttt{core} & S \hfill No\\
\hyperref[err:45]{45} & Observar un nodo numérico de azar no informa sobre sus padres & \texttt{core} & M \hfill \textcolor{decisionrojo}{Sí}\\
\hyperref[err:46]{46} & Un nodo de utilidad sin rango impide entrar en modo inferencia & \texttt{gui} & S \hfill No\\
\hyperref[err:47]{47} & «Add finding» marca siempre el primer estado y pisa el hallazgo anterior & \texttt{gui} & S \hfill No\\
\hyperref[err:48]{48} & La evidencia imposible se trata distinto según el gesto & \texttt{gui} & M \hfill \textcolor{decisionrojo}{Sí}\\
\hyperref[err:49]{49} & Construir una propagación cambia el tipo de análisis de la red del usuario & \texttt{core} & S \hfill No\\
\multicolumn{4}{@{}l}{\rule{0pt}{2.6ex}\textbf{\ref{sec:decisiones}.\ Decisiones: políticas, criterios, agentes y árbol de decisión}}\\
\hyperref[err:50]{50} & El árbol de decisión falla en casi todas las redes en formato 1.0 & \texttt{inference} & M \hfill No\\
\hyperref[err:51]{51} & «Show optimal policy» y «Show expected utility» no llegan a abrirse & \texttt{gui} & S \hfill \textcolor{decisionrojo}{Sí}\\
\hyperref[err:52]{52} & «Show expected utility» falla con utilidad constante y rellena ceros & \texttt{core} & S \hfill No\\
\hyperref[err:53]{53} & La utilidad esperada de una decisión ignora la evidencia previa & \texttt{inference} & S \hfill No\\
\hyperref[err:54]{54} & Un hallazgo que descarta la primera opción de una decisión anula el cálculo & \texttt{inference} & M \hfill No\\
\hyperref[err:55]{55} & Las ventanas de utilidad esperada y política óptima muestran la clave & \texttt{gui} & S \hfill No\\
\hyperref[err:56]{56} & Deshacer «quitar la política impuesta» pierde la tabla de la política & \texttt{core} & S \hfill No\\
\hyperref[err:57]{57} & Renombrar un criterio de decisión no se deshace ni se cancela & \texttt{core} & S \hfill No\\
\hyperref[err:58]{58} & Confirmar un criterio sin cambiar su nombre se rechaza como repetido & \texttt{core} & S \hfill \textcolor{decisionrojo}{Sí}\\
\hyperref[err:59]{59} & Evaluar una red de decisiones sin criterios falla sin explicar qué falta & \texttt{core} & S \hfill No\\
\hyperref[err:60]{60} & El agente de cada decisión se pierde al guardar y no se lee al abrir & \texttt{io} & S \hfill No\\
\hyperref[err:61]{61} & Quitar, renombrar o mover un agente descuadra las decisiones & \texttt{gui} & M \hfill No\\
\hyperref[err:62]{62} & Las propiedades de una decisión no se abren si su agente ya no figura & \texttt{gui} & S \hfill No\\
\hyperref[err:63]{63} & Evaluar DAN-arthronet falla con NullPointerException al instanciar & \texttt{inference} & S \hfill No\\
\hyperref[err:64]{64} & Instanciar una variable de una red de decisiones duplica un enlace & \texttt{inference} & S \hfill No\\
\hyperref[err:65]{65} & La estrategia óptima de una red sin decisiones lanza una excepción & \texttt{gui} & S \hfill No\\
\hyperref[err:66]{66} & El botón del árbol de decisión se enciende en redes que no admiten árbol & \texttt{gui} & S \hfill No\\
\multicolumn{4}{@{}l}{\rule{0pt}{2.6ex}\textbf{\ref{sec:temporal}.\ Modelos temporales y evolución temporal}}\\
\hyperref[err:67]{67} & Con evidencia, las curvas de la evolución temporal no suman uno & \texttt{inference} & M \hfill \textcolor{decisionrojo}{Sí}\\
\hyperref[err:68]{68} & La evolución temporal de una variable que empieza en el ciclo 1 falla & \texttt{inference} & S \hfill No\\
\hyperref[err:69]{69} & La evolución temporal de una decisión temporal muere con una excepción & \texttt{inference} & M \hfill \textcolor{decisionrojo}{Sí}\\
\hyperref[err:70]{70} & «Same as previous» revienta si el ciclo anterior no tiene la variable & \texttt{core} & S \hfill No\\
\hyperref[err:71]{71} & Al expandir un modelo temporal, las regresiones siguen nombrando el ciclo 0 & \texttt{core} & S \hfill No\\
\hyperref[err:72]{72} & Pasar una variable a atemporal puede dejar dos con el mismo nombre & \texttt{core} & S \hfill No\\
\hyperref[err:73]{73} & La duración del ciclo escrita en las propiedades de la red no se aplica & \texttt{gui} & S \hfill No\\
\multicolumn{4}{@{}l}{\rule{0pt}{2.6ex}\textbf{\ref{sec:coste}.\ Coste-efectividad y análisis de sensibilidad}}\\
\hyperref[err:74]{74} & El jar publicado mezcla dos JFreeChart y muchas gráficas no se abren & \texttt{full} & S \hfill No\\
\hyperref[err:75]{75} & La araña de coste-efectividad reescribe la red abierta con las medias & \texttt{core} & M \hfill No\\
\hyperref[err:76]{76} & Una columna con incertidumbre solo en algunas casillas se sortea mal & \texttt{core} & S \hfill \textcolor{decisionrojo}{Sí}\\
\hyperref[err:77]{77} & El análisis probabilístico de sensibilidad no se puede repetir & \texttt{core} & M \hfill No\\
\hyperref[err:78]{78} & El diálogo de incertidumbre rechaza toda Dirichlet sobre dos o más estados & \texttt{gui} & S \hfill No\\
\hyperref[err:79]{79} & Editor y muestreador reparten distinto el complemento junto a una Dirichlet & \texttt{core} & S \hfill \textcolor{decisionrojo}{Sí}\\
\hyperref[err:80]{80} & El mapa de dos parámetros elige mal la mejor opción si hay tres o más & \texttt{sensitivity\allowbreak{}Analysis} & S \hfill No\\
\hyperref[err:81]{81} & En el mapa de dos parámetros el empate y la tercera opción salen en gris & \texttt{sensitivity\allowbreak{}Analysis} & S \hfill No\\
\hyperref[err:82]{82} & El mapa de dos parámetros pinta toda la zona de datos de un solo color & \texttt{sensitivity\allowbreak{}Analysis} & S \hfill No\\
\hyperref[err:83]{83} & El mapa no se abre si gana en todo él la misma opción y no es la primera & \texttt{sensitivity\allowbreak{}Analysis} & S \hfill No\\
\hyperref[err:84]{84} & La tabla de intervenciones fronterizas nombra otra opción si dos empatan & \texttt{cost\allowbreak{}Effectiveness} & S \hfill No\\
\hyperref[err:85]{85} & El coste-efectividad falla si un hallazgo hace imposible una sola opción & \texttt{inference} & S \hfill No\\
\hyperref[err:86]{86} & Las ventanas de coste-efectividad fallan si una opción es imposible & \texttt{cost\allowbreak{}Effectiveness} & S \hfill No\\
\hyperref[err:87]{87} & El escenario del ámbito de decisión puede ser imposible y no se comprueba & \texttt{gui} & S \hfill \textcolor{decisionrojo}{Sí}\\
\hyperref[err:88]{88} & Con valores todos negativos, el tornado dibuja barras hasta cero & \texttt{sensitivity\allowbreak{}Analysis} & S \hfill No\\
\hyperref[err:89]{89} & El tornado por decisión no dibuja el rango real de la diferencia & \texttt{sensitivity\allowbreak{}Analysis} & S \hfill No\\
\hyperref[err:90]{90} & Con utilidades negativas, «Plot» y la araña aplastan y recortan las curvas & \texttt{sensitivity\allowbreak{}Analysis} & S \hfill No\\
\hyperref[err:91]{91} & El tornado con dieciséis parámetros o más deja barras fuera del recuadro & \texttt{sensitivity\allowbreak{}Analysis} & S \hfill No\\
\hyperref[err:92]{92} & El análisis de sensibilidad no se abre sobre una red sin decisiones & \texttt{sensitivity\allowbreak{}Analysis} & S \hfill No\\
\hyperref[err:93]{93} & Los parámetros inciertos salen en otro orden en cada apertura del análisis & \texttt{core} & S \hfill No\\
\hyperref[err:94]{94} & Ocultar una opción en una pestaña deja marcada su casilla en la otra & \texttt{sensitivity\allowbreak{}Analysis} & S \hfill No\\
\multicolumn{4}{@{}l}{\rule{0pt}{2.6ex}\textbf{\ref{sec:des}.\ Simulación de sucesos discretos}}\\
\hyperref[err:95]{95} & Una utilidad acumulativa vale cero hasta el primer suceso que la alcanza & \texttt{inference.\allowbreak{}des} & M \hfill \textcolor{decisionrojo}{Sí}\\
\hyperref[err:96]{96} & La simulación se cae si un nodo espera a un padre de un suceso no ocurrido & \texttt{inference.\allowbreak{}des} & S \hfill No\\
\hyperref[err:97]{97} & Un ciclo entre nodos de azar de una red de sucesos impide simular & \texttt{inference.\allowbreak{}des} & M \hfill \textcolor{decisionrojo}{Sí}\\
\hyperref[err:98]{98} & Una tabla o valor exacto con un suceso padre hace caer la simulación & \texttt{core} & M \hfill \textcolor{decisionrojo}{Sí}\\
\hyperref[err:99]{99} & Añadir o quitar un enlace deja una relación uniforme y la simulación falla & \texttt{core} & M \hfill \textcolor{decisionrojo}{Sí}\\
\hyperref[err:100]{100} & Una tabla de vida «con tasas» se simula como si fueran probabilidades & \texttt{core} & S \hfill \textcolor{decisionrojo}{Sí}\\
\hyperref[err:101]{101} & La relación «External» se ofrece, elegirla falla y no trae valores de fuera & \texttt{core} & M \hfill \textcolor{decisionrojo}{Sí}\\
\hyperref[err:102]{102} & Resultados de la simulación: la media descontada va al criterio equivocado & \texttt{inference.\allowbreak{}des} & S \hfill No\\
\hyperref[err:103]{103} & El CSV de la simulación de sucesos tiene los títulos corridos una columna & \texttt{inference.\allowbreak{}des} & S \hfill No\\
\hyperref[err:104]{104} & Simular con traza una red de sucesos sin guardar falla antes de empezar & \texttt{inference.\allowbreak{}des} & S \hfill No\\
\hyperref[err:105]{105} & Un fallo de «Tree with Excluded Events» abre una ventana desde el núcleo & \texttt{core} & S \hfill No\\
\hyperref[err:106]{106} & Un nodo de azar discretizado impide simular con un error inesperado & \texttt{inference.\allowbreak{}des} & S \hfill \textcolor{decisionrojo}{Sí}\\
\hyperref[err:107]{107} & La hoja de la propagación estocástica cambia el orden de los hallazgos & \texttt{stochastic\allowbreak{}Propagation\allowbreak{}Output} & S \hfill No\\
\hyperref[err:108]{108} & La propagación estocástica no rechaza nodos numéricos: falla sin explicarlo & \texttt{stochastic\allowbreak{}Propagation\allowbreak{}Output} & S \hfill No\\
\hyperref[err:109]{109} & Las casillas de estadísticos del panel de Monte Carlo no hacen nada & \texttt{gui} & S \hfill \textcolor{decisionrojo}{Sí}\\
\hyperref[err:110]{110} & «Sort probnet» falla en una red de sucesos con un nodo enlazado a sí mismo & \texttt{full} & S \hfill No\\
\multicolumn{4}{@{}l}{\rule{0pt}{2.6ex}\textbf{\ref{sec:ficheros}.\ Abrir y guardar redes}}\\
\hyperref[err:111]{111} & Un .pgmx mal escrito se abre sin aviso y con todas las tablas uniformes & \texttt{io} & S \hfill \textcolor{decisionrojo}{Sí}\\
\hyperref[err:112]{112} & Guardar en formato 1.0 deja la red sin poder evaluarse al reabrirla & \texttt{io} & M \hfill \textcolor{decisionrojo}{Sí}\\
\hyperref[err:113]{113} & Guardar en formato 0.2 convierte los nodos de producto en uniformes & \texttt{io} & S \hfill No\\
\hyperref[err:114]{114} & «Save» falla en tres redes con un papel de potencial mal escrito & \texttt{io} & S \hfill \textcolor{decisionrojo}{Sí}\\
\hyperref[err:115]{115} & Una distribución desconocida en una red de sucesos falla lejos de la causa & \texttt{core} & S \hfill \textcolor{decisionrojo}{Sí}\\
\hyperref[err:116]{116} & Una tabla de sucesos con incertidumbre en el fichero impide abrir la red & \texttt{io} & S \hfill \textcolor{decisionrojo}{Sí}\\
\hyperref[err:117]{117} & Abrir un .xml que no es ProbModelXML falla con una excepción de índice & \texttt{io} & S \hfill No\\
\multicolumn{4}{@{}l}{\rule{0pt}{2.6ex}\textbf{\ref{sec:casos}.\ Bases de casos}}\\
\hyperref[err:118]{118} & Un «?» en una base de Weka descuadra las filas e inventa casos sin avisar & \texttt{io} & M \hfill No\\
\hyperref[err:119]{119} & Un fichero de Weka en UTF-8 con tildes o «ñ» pierde estados y lee ceros & \texttt{io} & M \hfill No\\
\hyperref[err:120]{120} & Una columna numérica, de texto o de fecha en Weka da NullPointerException & \texttt{io} & M \hfill \textcolor{decisionrojo}{Sí}\\
\hyperref[err:121]{121} & Weka con @Relation, @Attribute o @Data en mayúsculas mezcladas no se abre & \texttt{io} & S \hfill No\\
\hyperref[err:122]{122} & Una base con una columna de números guardada en Weka no se puede reabrir & \texttt{io} & S \hfill No\\
\hyperref[err:123]{123} & Un fichero de Weka en una carpeta con tilde, «ñ» o «+» no se puede reabrir & \texttt{io} & S \hfill No\\
\hyperref[err:124]{124} & Una línea de CSV con campos de menos, de más o en blanco inventa casos & \texttt{io} & S \hfill No\\
\multicolumn{4}{@{}l}{\rule{0pt}{2.6ex}\textbf{\ref{sec:aprendizaje}.\ Aprendizaje y evaluación de redes}}\\
\hyperref[err:125]{125} & La exactitud del Bayes ingenuo selectivo usa el caso evaluado al entrenar & \texttt{learning.\allowbreak{}metric} & S \hfill No\\
\hyperref[err:126]{126} & El Bayes ingenuo selectivo elige características sin el alfa del diálogo & \texttt{learning.\allowbreak{}algorithm} & S \hfill No\\
\hyperref[err:127]{127} & El superpadre pesa sus enlaces con cuentas de casillas equivocadas & \texttt{learning.\allowbreak{}metric} & S \hfill No\\
\hyperref[err:128]{128} & El superpadre penaliza la clase verdadera: ningún enlace sube la exactitud & \texttt{learning.\allowbreak{}metric} & S \hfill No\\
\hyperref[err:129]{129} & El superpadre puntúa los enlaces candidatos con la probabilidad al revés & \texttt{learning.\allowbreak{}metric} & S \hfill No\\
\hyperref[err:130]{130} & El superpadre cuenta dos veces la hija de cada enlace al medir la exactitud & \texttt{learning.\allowbreak{}metric} & S \hfill No\\
\hyperref[err:131]{131} & La validación cruzada de dos clasificadores no prueba un tercio de casos & \texttt{learning.\allowbreak{}metric} & S \hfill No\\
\hyperref[err:132]{132} & El superpadre por defecto no aprende nada: propone un enlace a sí mismo & \texttt{learning.\allowbreak{}algorithm} & S \hfill No\\
\hyperref[err:133]{133} & El superpadre acepta el primer enlace aunque empeore al Bayes ingenuo & \texttt{learning.\allowbreak{}algorithm} & S \hfill No\\
\hyperref[err:134]{134} & Superpadre interactivo: llenar la tabla avanza el algoritmo y la vacía & \texttt{learning.\allowbreak{}algorithm} & M \hfill \textcolor{decisionrojo}{Sí}\\
\hyperref[err:135]{135} & El Bayes ingenuo selectivo hacia delante no da dos veces el mismo modelo & \texttt{learning.\allowbreak{}algorithm} & S \hfill No\\
\hyperref[err:136]{136} & El árbol aumentado (TAN) orienta sus enlaces distinto en cada aprendizaje & \texttt{learning.\allowbreak{}algorithm} & S \hfill \textcolor{decisionrojo}{Sí}\\
\hyperref[err:137]{137} & El EM deja sin aprender los nodos con potencial uniforme & \texttt{learning.\allowbreak{}algorithm} & S \hfill No\\
\hyperref[err:138]{138} & El esqueleto del algoritmo PC cambia al renombrar una variable & \texttt{learning.\allowbreak{}algorithm} & S \hfill No\\
\hyperref[err:139]{139} & El PC interactivo ofrece las dos orientaciones de un enlace, contra Meek & \texttt{learning.\allowbreak{}algorithm} & S \hfill \textcolor{decisionrojo}{Sí}\\
\hyperref[err:140]{140} & El PC interactivo adelanta la fase al llenar la tabla y descoloca botones & \texttt{learning.\allowbreak{}algorithm} & M \hfill \textcolor{decisionrojo}{Sí}\\
\hyperref[err:141]{141} & «Finish» del PC interactivo empieza de cero y descarta lo aplicado a mano & \texttt{learning.\allowbreak{}algorithm} & M \hfill \textcolor{decisionrojo}{Sí}\\
\hyperref[err:142]{142} & PC interactivo: con un borrado bloqueado, «Finish» da NullPointerException & \texttt{learning.\allowbreak{}algorithm} & S \hfill No\\
\hyperref[err:143]{143} & El bosque aumentado (FAN) falla sin características ligadas a la clase & \texttt{learning.\allowbreak{}algorithm} & S \hfill No\\
\hyperref[err:144]{144} & Discretizar por frecuencias iguales con decimales falla o cuenta mal & \texttt{learning.\allowbreak{}core} & S \hfill No\\
\hyperref[err:145]{145} & «Split Dataset» lanza siempre una excepción, elija la opción que elija & \texttt{bn\allowbreak{}Evaluation} & S \hfill No\\
\hyperref[err:146]{146} & Validación cruzada: casos por pliegue truncados y pérdida por caso inflada & \texttt{bn\allowbreak{}Evaluation} & S \hfill No\\
\hyperref[err:147]{147} & Quedarse sin memoria al aprender sale como error no previsto & \texttt{learning.\allowbreak{}gui} & S \hfill No\\
\hyperref[err:148]{148} & Una variable de la red modelo que falta en los datos da error no previsto & \texttt{learning.\allowbreak{}gui} & S \hfill No\\
\hyperref[err:149]{149} & Red modelo sin alguna columna de la base: error no previsto al cargarla & \texttt{learning.\allowbreak{}gui} & S \hfill No\\
\multicolumn{4}{@{}l}{\rule{0pt}{2.6ex}\textbf{\ref{sec:editor}.\ Interfaz: editor de red y paneles de las relaciones}}\\
\hyperref[err:150]{150} & Volver a elegir la clase de relación que ya tiene un nodo borra sus números & \texttt{gui} & S \hfill No\\
\hyperref[err:151]{151} & En una regresión, lo escrito como Cholesky se pierde al guardar la red & \texttt{gui} & S \hfill No\\
\hyperref[err:152]{152} & Añadir una fila a la tabla de vida y aceptar borra el primer tramo & \texttt{gui} & S \hfill No\\
\hyperref[err:153]{153} & El editor de una distribución univariante enseña «1» en cada parámetro & \texttt{core} & M \hfill No\\
\hyperref[err:154]{154} & En la tabla de utilidad, una restricción de enlace bloquea otras casillas & \texttt{core} & S \hfill No\\
\hyperref[err:155]{155} & La ventana de restricción de un enlace edita la tabla del padre & \texttt{gui} & S \hfill No\\
\hyperref[err:156]{156} & La tabla de utilidad enseña 922 337 203,685478 en vez de valores mayores & \texttt{core} & S \hfill No\\
\hyperref[err:157]{157} & \textcolor{corregidoverde}{\textbf{Corregido en parte}} · El panel de Weibull deja borrar «Gamma» y la relación revienta al calcular & \texttt{gui} & S \hfill No\\
\hyperref[err:158]{158} & En el panel de una regresión no se ve la columna de las covariables & \texttt{gui} & S \hfill No\\
\hyperref[err:159]{159} & «Add variables to potential» en una hoja del árbol no añade ningún padre & \texttt{gui} & S \hfill No\\
\hyperref[err:160]{160} & El árbol ofrece repetir una variable discretizada y elegirla da un error & \texttt{gui} & S \hfill No\\
\hyperref[err:161]{161} & «Split interval» con el campo vacío o sin un número da un error inesperado & \texttt{gui} & S \hfill No\\
\hyperref[err:162]{162} & La pestaña «Parents» ofrece padres que las restricciones rechazan & \texttt{gui} & S \hfill No\\
\hyperref[err:163]{163} & El límite superior del último intervalo de revelación no se puede cambiar & \texttt{gui} & S \hfill No\\
\hyperref[err:164]{164} & Un número no válido en el indicador da un error y vacía la pestaña & \texttt{gui} & S \hfill No\\
\hyperref[err:165]{165} & El doble clic en la ventanilla de expresiones mete dos veces el texto & \texttt{gui} & S \hfill No\\
\hyperref[err:166]{166} & Tras cambiar el tipo de un nodo, una copia suya no encuentra sus padres & \texttt{core} & S \hfill No\\
\hyperref[err:167]{167} & Una columna con incertidumbre fuera del primer estado se deja editar & \texttt{core} & S \hfill No\\
\hyperref[err:168]{168} & Tras alejar el zoom, los nombres de estado siguen recortados al acercarlo & \texttt{gui} & S \hfill No\\
\hyperref[err:169]{169} & Con precisión 0.0, 2.0 o 0.5, los límites de los tramos acaban en punto & \texttt{gui} & S \hfill No\\
\hyperref[err:170]{170} & Tras A, B y C, el nodo de azar nuevo se llama «E», la letra de los sucesos & \texttt{core} & S \hfill No\\
\multicolumn{4}{@{}l}{\rule{0pt}{2.6ex}\textbf{\ref{sec:menus}.\ Interfaz: menús, diálogos y herramientas}}\\
\hyperref[err:171]{171} & Abrir una red desde internet apaga «Guardar» para todas las redes & \texttt{gui} & M \hfill No\\
\hyperref[err:172]{172} & «Edits history» deja un núcleo al cien por cien y puede colgar el programa & \texttt{full} & S \hfill \textcolor{decisionrojo}{Sí}\\
\hyperref[err:173]{173} & Aceptar los dominios estándar sin marcar ninguno lanza una excepción & \texttt{gui} & S \hfill No\\
\hyperref[err:174]{174} & Aceptar los criterios estándar sin marcar nada borra los criterios de la red & \texttt{gui} & S \hfill No\\
\hyperref[err:175]{175} & «OpenMarkov configuration» edita unas preferencias que el programa no usa & \texttt{gui} & M \hfill \textcolor{decisionrojo}{Sí}\\
\hyperref[err:176]{176} & Las escalas de criterio se redondean a tres decimales y pueden quedar en 0 & \texttt{gui} & S \hfill No\\
\multicolumn{4}{@{}l}{\rule{0pt}{2.6ex}\textbf{\ref{sec:textos}.\ Textos y títulos que se ven mal}}\\
\hyperref[err:177]{177} & La ayuda de atajos da dos teclas equivocadas, y una cierra el programa & \texttt{gui} & S \hfill No\\
\hyperref[err:178]{178} & Pegar un nodo de sucesos donde no se admite da «\textgreater\textgreater\textgreater{} eventNode is null \textless\textless\textless» & \texttt{core} & S \hfill No\\
\hyperref[err:179]{179} & El selector del tipo de relación enseña nombres internos como «ProbTable» & \texttt{gui} & S \hfill \textcolor{decisionrojo}{Sí}\\
\hyperref[err:180]{180} & El diálogo de dominios estándar se abre sin título & \texttt{gui} & S \hfill No\\
\hyperref[err:181]{181} & Seis de las ventanillas del editor de árboles se abren sin título & \texttt{gui} & S \hfill No\\
\hyperref[err:182]{182} & Varios títulos de error parten las siglas: «tree a d d», «X m l invalid» & \texttt{core} & S \hfill No\\
\hyperref[err:183]{183} & El aviso de zoom fuera de rango dice «Requested zoomManager» & \texttt{gui} & S \hfill No\\
\hyperref[err:184]{184} & El rótulo del panel del indicador dice «Probability of ocurrence» & \texttt{gui} & S \hfill No\\
\hyperref[err:185]{185} & La ayuda de «Add → Utility node» dice «Create an utility node» & \texttt{gui} & S \hfill No\\
\hyperref[err:186]{186} & El mensaje de evidencia incompatible dice «old evidente» y es ilegible & \texttt{core} & S \hfill No\\
\hyperref[err:187]{187} & El error de muestras sin peso dice «Samples weigth is 0» & \texttt{core} & S \hfill No\\
\hyperref[err:188]{188} & La ayuda de «Unique superparent» dice «commont» & \texttt{learning.\allowbreak{}gui} & S \hfill No\\
\multicolumn{4}{@{}l}{\rule{0pt}{2.6ex}\textbf{\ref{sec:rendimiento}.\ Rendimiento}}\\
\hyperref[err:189]{189} & Entrar en modo inferencia con una red grande congela la ventana & \texttt{gui} & M \hfill \textcolor{decisionrojo}{Sí}\\
\hyperref[err:190]{190} & Cada resultado de una operación sobre tablas reserva una tabla que tira & \texttt{core} & S \hfill No\\
\hyperref[err:191]{191} & Cada tabla de probabilidad nueva se rellena dos veces con el mismo valor & \texttt{core} & S \hfill No\\
\hyperref[err:192]{192} & La maximización reserva la tabla de la unión de variables sin usarla & \texttt{core} & S \hfill No\\
\bottomrule
\end{longtable}
}


% ══════════════════════════════════════════════════════════════════════
\section{Operaciones con potenciales}\label{sec:potenciales}

Las operaciones aritméticas sobre potenciales (multiplicar, marginalizar, copiar) que usan todos los algoritmos. Un potencial es la función que cuantifica un nodo: una tabla de probabilidad, una utilidad, una fórmula.

\setcounter{subsection}{0}
\subsection[Multiplicar potenciales constantes pierde el criterio de la utilidad]{\textcolor{corregidoverde}{[Corregido]} Multiplicar potenciales que son todos un solo número pierde el criterio, y una utilidad sale convertida en algo que no es aditivo}\label{err:1}
\begin{ficha}
\fichakey{Estado}Corregido en el commit \texttt{bbde667} (23-09-2026): multiplicar solo potenciales constantes conserva el criterio y el árbol de estrategia, así que una utilidad constante sigue siendo aditiva. Que el producto deje fuera las variables de un solo estado no se corrige: no cambia ningún número, y conservarlas rompe dos cosas que hoy funcionan (\hyperref[nota:185]{nota~185}).
\fichakey{Módulo}core
\fichakey{Paquete}org.openmarkov.core.model.network.potential.operation
\fichakey{Ficheros}\path{core/src/main/java/org/openmarkov/core/model/network/potential/operation/TablePotentialArithmetic.java:69-99} \newline \path{core/src/main/java/org/openmarkov/core/model/network/potential/operation/TablePotentialArithmetic.java:674-678} \newline \path{core/src/main/java/org/openmarkov/core/model/network/potential/operation/AuxiliaryOperations.java:48-56} \newline \path{core/src/main/java/org/openmarkov/core/model/network/potential/operation/MaxOutVariable.java:59} \newline \path{core/src/main/java/org/openmarkov/core/model/network/potential/operation/MaxOutVariable.java:159}
\fichakey{Comprobado}Leyendo el código y ejecutándolo
\fichakey{Esfuerzo}Pequeño (horas)
\fichakey{Decisión de equipo}No
\fichakey{Relacionados}\hyperref[nota:93]{nota~93}, \hyperref[nota:185]{nota~185}
\end{ficha}
\paragraph{Qué pasa}
Quien multiplica con \texttt{Discrete\allowbreak{}Potential\allowbreak{}Operations.\allowbreak{}multiply} una lista de potenciales que son todos constantes (una sola casilla) recibe un resultado \textbf{sin criterio de decisión}. El criterio es lo único que distingue un potencial de utilidad (aditivo) de uno de probabilidad: \texttt{Potential.\allowbreak{}is\allowbreak{}Additive(\allowbreak{})} devuelve \texttt{criterion\ !=\ null}. Así, una utilidad con criterio multiplicada por la probabilidad unidad sale convertida en un potencial que ya no es aditivo.

El mismo atajo pierde también la variable de un solo estado del resultado y los árboles de estrategia de un \texttt{Strategic\allowbreak{}Table\allowbreak{}Potential} constante, porque tampoco los pasa al constructor.

Dónde está tapado hoy: \texttt{Max\allowbreak{}Out\allowbreak{}Variable} (eliminación de una decisión) multiplica la probabilidad por la utilidad en la línea 59 y, al final, vuelve a poner a mano el criterio de la utilidad de entrada (línea 159). Cuando una utilidad queda reducida a un solo número (por ejemplo, sobre una decisión con un único camino posible), el producto sale sin criterio y es esa línea la que lo arregla. No se ha encontrado un resultado visible distinto.

Hay dos llamadores de producción que no lo corrigen y pueden recibir utilidades constantes: \texttt{Product\allowbreak{}Potential.\allowbreak{}table\allowbreak{}Project} (nodo de utilidad producto de sus padres: si todos los padres se proyectan a un solo número, el resultado sale sin criterio) y \texttt{Utility\allowbreak{}Function\allowbreak{}Computer} (cálculo de rangos de utilidad para las barras de la ventana). No se ha comprobado con una red que alguno de ellos cambie lo que ve el usuario.

Por qué pasa: \texttt{Discrete\allowbreak{}Potential\allowbreak{}Operations.\allowbreak{}multiply} delega en \texttt{Table\allowbreak{}Potential\allowbreak{}Arithmetic.\allowbreak{}multiply(\allowbreak{}List,\allowbreak{}\ boolean)}. Ese método calcula al principio el criterio del resultado (\texttt{find\allowbreak{}First\allowbreak{}Non\allowbreak{}Null\allowbreak{}Criterion}, línea 81). Después aparta los potenciales «constantes» ---para \texttt{Auxiliary\allowbreak{}Operations}, cualquiera cuya tabla tenga una sola casilla, aunque tenga variables de un solo estado--- y, si no queda ninguno, devuelve \texttt{build\allowbreak{}Constant\allowbreak{}Potential(\allowbreak{}constant\allowbreak{}Factor,\allowbreak{}\ role)} (línea 98), que no recibe el criterio. El criterio ya calculado se descarta.

Medido:

\begin{Verbatim}[breaklines,breakanywhere,fontsize=\small]
utilidad constante (cost) x probabilidad unidad:            criterio=(ninguno) aditivo=false valores=[10.0]
utilidad de una casilla sobre E (cost) x probabilidad unidad: criterio=(ninguno) aditivo=false valores=[10.0]
utilidad sobre D (cost) x probabilidad unidad:                criterio=cost aditivo=true valores=[10.0, 20.0]
utilidad constante (cost) x probabilidad sobre D:             criterio=cost aditivo=true valores=[5.0, 5.0]
\end{Verbatim}
\paragraph{El código}
core/src/main/java/org/openmarkov/core/model/network/potential/operation/TablePotentialArithmetic.java:79-99, 674-678

\begin{Shaded}
\begin{Highlighting}[]
        \CommentTok{// Find out if some potential has criterion. In that case, set that criterion in}
        \CommentTok{// the resulting potential}
\NormalTok{        Criterion criterion }\OperatorTok{=} \FunctionTok{findFirstNonNullCriterion}\OperatorTok{(}\NormalTok{tablePotentials}\OperatorTok{);}
        \KeywordTok{...}
        \CommentTok{// Gets constant factor: The product of constant potentials}
        \DataTypeTok{double}\NormalTok{ constantFactor }\OperatorTok{=}\NormalTok{ AuxiliaryOperations}\OperatorTok{.}\FunctionTok{getConstantFactor}\OperatorTok{(}\NormalTok{potentials}\OperatorTok{);}

        \CommentTok{// get role}
\NormalTok{        PotentialRole role }\OperatorTok{=} \FunctionTok{getRole}\OperatorTok{(}\NormalTok{potentials}\OperatorTok{);}

\NormalTok{        potentials }\OperatorTok{=}\NormalTok{ AuxiliaryOperations}\OperatorTok{.}\FunctionTok{getNonConstantPotentials}\OperatorTok{(}\NormalTok{potentials}\OperatorTok{);}
        \ControlFlowTok{if} \OperatorTok{(}\NormalTok{potentials}\OperatorTok{.}\FunctionTok{isEmpty}\OperatorTok{())} \OperatorTok{\{}
            \ControlFlowTok{return} \FunctionTok{buildConstantPotential}\OperatorTok{(}\NormalTok{constantFactor}\OperatorTok{,}\NormalTok{ role}\OperatorTok{);}
        \OperatorTok{\}}
\KeywordTok{...}
    \KeywordTok{private} \DataTypeTok{static}\NormalTok{ TablePotential }\FunctionTok{buildConstantPotential}\OperatorTok{(}\DataTypeTok{double}\NormalTok{ constantFactor}\OperatorTok{,}\NormalTok{ PotentialRole role}\OperatorTok{)} \OperatorTok{\{}
\NormalTok{        TablePotential constantTablePotential }\OperatorTok{=} \KeywordTok{new} \FunctionTok{TablePotential}\OperatorTok{(}\KeywordTok{null}\OperatorTok{,}\NormalTok{ role}\OperatorTok{);}
\NormalTok{        constantTablePotential}\OperatorTok{.}\FunctionTok{getValues}\OperatorTok{()[}\DecValTok{0}\OperatorTok{]} \OperatorTok{=}\NormalTok{ constantFactor}\OperatorTok{;}
        \ControlFlowTok{return}\NormalTok{ constantTablePotential}\OperatorTok{;}
    \OperatorTok{\}}
\end{Highlighting}
\end{Shaded}
\paragraph{Cómo arreglarlo}
Pasar el criterio al resultado también en el atajo (\texttt{build\allowbreak{}Constant\allowbreak{}Potential} con criterio, igual que \texttt{build\allowbreak{}Result\allowbreak{}Potential}), y conservar los árboles de estrategia si el único potencial con intervenciones era constante. Es la misma regla que ya siguen la suma y la maximización sobre una colección: el criterio es el primero no nulo de la lista entera, constantes incluidas. El atajo de la línea 77 (lista vacía, devuelve la probabilidad unidad) no tiene criterio que conservar.

Otros sitios que deben cumplir la misma regla: los tres atajos de «todo son constantes» de marginalizar y maximizar (sus valores ya son correctos; la tabla de elecciones vacía que devuelve la maximización está descrita en \hyperref[nota:93]{nota~93}, donde hoy no hay nada que corregir) y \texttt{Sum\allowbreak{}Out\allowbreak{}Variable}. Con el arreglo, la línea 159 de \texttt{Max\allowbreak{}Out\allowbreak{}Variable} deja de ser necesaria, pero puede quedarse.

Prueba de regresión: multiplicar una utilidad constante con criterio por la probabilidad unidad y exigir el criterio en el resultado.

\setcounter{subsection}{1}
\subsection[Tablas de más de 2\textsuperscript{31} casillas: excepciones distintas o resultado falso]{\textcolor{corregidoverde}{[Corregido]} Una operación cuya tabla resultado pasa de 2\textsuperscript{31} casillas falla con una excepción distinta según el número, o devuelve un resultado equivocado sin avisar}\label{err:2}
\begin{ficha}
\fichakey{Estado}Corregido en el commit \texttt{76b0a41} (23-09-2026): \texttt{multiply}, \texttt{sum}, \texttt{evaluate\allowbreak{}Function\allowbreak{}Potential}, \texttt{merge} y \texttt{calculate\allowbreak{}Offsets} calculan el tamaño con \texttt{compute\allowbreak{}Table\allowbreak{}Size} antes de reservar memoria, y \texttt{multiply\allowbreak{}And\allowbreak{}Marginalize} también el número de combinaciones que elimina; toda tabla de más de 2\textsuperscript{31}~−~1 casillas lanza la misma \texttt{Invalid\allowbreak{}Argument\allowbreak{}Exception} en lugar de otra excepción o un resultado falso.
\fichakey{Módulo}core
\fichakey{Paquete}org.openmarkov.core.model.network.potential.operation
\fichakey{Ficheros}\path{core/src/main/java/org/openmarkov/core/model/network/potential/operation/TablePotentialArithmetic.java:125-127, 272-278, 600-602} \newline \path{core/src/main/java/org/openmarkov/core/model/network/potential/operation/TablePotentialMerge.java:68-71} \newline \path{core/src/main/java/org/openmarkov/core/model/network/potential/operation/TablePotentialElimination.java:115-121, 148-166} \newline \path{core/src/main/java/org/openmarkov/core/model/network/potential/AbstractIndexedPotential.java:142-169}
\fichakey{Comprobado}Leyendo el código y ejecutándolo
\fichakey{Esfuerzo}Pequeño (horas)
\fichakey{Decisión de equipo}No
\end{ficha}
\paragraph{Qué pasa}
Una tabla tiene una casilla por cada combinación de estados de sus variables, así que su tamaño es el producto de los números de estados. Cuando ese producto no cabe en un \texttt{int} (más de 2\textsuperscript{31} − 1 casillas), el llamador debería recibir siempre la misma \texttt{Invalid\allowbreak{}Argument\allowbreak{}Exception} con el mensaje «the product of their numbers of states exceeds the maximum table size, 2\textsuperscript{31} - 1». Hoy recibe una excepción distinta según el número, y en un caso ninguna: un resultado falso.

Camino de llegada: cualquier inferencia sobre una red con una camarilla muy grande. Con las redes bayesianas de \texttt{integration\allowbreak{}Tests/\allowbreak{}src/\allowbreak{}main/\allowbreak{}resources/\allowbreak{}clone\_\allowbreak{}of\_\allowbreak{}probmodelxml\_\allowbreak{}networks/\allowbreak{}bn} (20 ficheros) y con \texttt{integration\allowbreak{}Tests/\allowbreak{}src/\allowbreak{}main/\allowbreak{}resources/\allowbreak{}networks/\allowbreak{}bn/\allowbreak{}BN-cpcs360b.\allowbreak{}pgmx} e \texttt{integration\allowbreak{}Tests/\allowbreak{}src/\allowbreak{}main/\allowbreak{}resources/\allowbreak{}networks/\allowbreak{}bn/\allowbreak{}BN-cpcs422b.\allowbreak{}pgmx} (360 y 422 nodos), la eliminación de variables (lo que usa la ventana) termina bien en todas. La propagación de Hugin, que es la que usa la ventana de exportar la propagación estocástica a Excel, falla en \texttt{BN-nasonet.\allowbreak{}pgmx} y \texttt{1-0/\allowbreak{}BN-nasonet-1-0.\allowbreak{}pgmx} de esa misma carpeta, hoy con el mensaje correcto porque el desbordamiento cae en un número pequeño.

Por qué pasa: el cálculo protegido, \texttt{Abstract\allowbreak{}Indexed\allowbreak{}Potential.\allowbreak{}compute\allowbreak{}Table\allowbreak{}Size}, multiplica con \texttt{Math.\allowbreak{}multiply\allowbreak{}Exact} y lanza la excepción buena; el constructor de \texttt{Table\allowbreak{}Potential} lo usa. Pero varias operaciones no usan ese cálculo: construyen el tamaño con \texttt{dimensions{[}n-1{]}\ *\ offsets{[}n-1{]}} en un \texttt{int}, reservan \texttt{new\ double{[}table\allowbreak{}Size{]}} con ese número y solo al final construyen el potencial, que es cuando la comprobación salta. Si el producto da la vuelta a un número negativo, se lanza \texttt{Negative\allowbreak{}Array\allowbreak{}Size\allowbreak{}Exception} antes de llegar a la comprobación; si da la vuelta a un número positivo grande, se intenta reservar esa memoria y puede salir \texttt{Out\allowbreak{}Of\allowbreak{}Memory\allowbreak{}Error}; solo si da la vuelta a un número pequeño llega al constructor y se ve el mensaje correcto. Los sitios son \texttt{Table\allowbreak{}Potential\allowbreak{}Arithmetic.\allowbreak{}multiply}, \texttt{Table\allowbreak{}Potential\allowbreak{}Arithmetic.\allowbreak{}sum}, \texttt{Table\allowbreak{}Potential\allowbreak{}Arithmetic.\allowbreak{}evaluate\allowbreak{}Function\allowbreak{}Potential} y \texttt{Table\allowbreak{}Potential\allowbreak{}Merge.\allowbreak{}merge}. \texttt{calculate\allowbreak{}Offsets} también multiplica en \texttt{int} sin comprobar.

El caso peor está en \texttt{Table\allowbreak{}Potential\allowbreak{}Elimination.\allowbreak{}multiply\allowbreak{}And\allowbreak{}Marginalize} (multiplicar varios potenciales y sumar las variables que sobran). El tamaño del resultado sí usa \texttt{compute\allowbreak{}Table\allowbreak{}Size}, pero el número de combinaciones de las variables a eliminar se acumula en \texttt{int\ elimination\allowbreak{}Size} sin control. Si ese producto da la vuelta a cero o a un negativo, el bucle interior no se ejecuta y cada casilla del resultado recibe solo la primera combinación. Como el potencial resultado tiene un tamaño legal, \textbf{no hay excepción}: el resultado es simplemente falso.

No se ha encontrado un llamador de producción que elimine de una vez tantas variables como para disparar el resultado falso de \texttt{multiply\allowbreak{}And\allowbreak{}Marginalize}: la eliminación de variables elimina de una en una y Hugin elimina sobre camarillas ya construidas. El defecto de la operación sí se ha comprobado.

Medido con pequeños programas de prueba (una ejecución de cada uno; son deterministas), con potenciales de todo unos y variables binarias:

\begin{itemize}
\tightlist
\item
  \texttt{multiply} de un potencial de 16 variables por otro de 16 distintas (2\^{}32): \texttt{Invalid\allowbreak{}Argument\allowbreak{}Exception} (el \texttt{int} da cero y llega al constructor).
\item
  \texttt{multiply} y \texttt{sum} de 16 por 15 variables (2\textsuperscript{31}): \texttt{Negative\allowbreak{}Array\allowbreak{}Size\allowbreak{}Exception:\ -2147483648}.
\item
  \texttt{multiply} de 15 binarias por 15 binarias más una de cinco estados (5·2\^{}30), con 2 GB de memoria: \texttt{Out\allowbreak{}Of\allowbreak{}Memory\allowbreak{}Error:\ Java\ heap\ space}.
\item
  \texttt{multiply\allowbreak{}And\allowbreak{}Marginalize} de 16 por 16 variables conservando una (elimina 31 binarias): devuelve {[}1,0, 1,0{]}; lo correcto es {[}2\textsuperscript{31}, 2\textsuperscript{31}{]}. Con 16 por 17 (elimina 32) devuelve también {[}1,0, 1,0{]} en vez de {[}2\^{}32, 2\^{}32{]}.
\end{itemize}
\paragraph{El código}
core/src/main/java/org/openmarkov/core/model/network/potential/operation/TablePotentialArithmetic.java:125-127

\begin{Shaded}
\begin{Highlighting}[]
\DataTypeTok{int}\OperatorTok{[]}\NormalTok{ offsets }\OperatorTok{=}\NormalTok{ TablePotential}\OperatorTok{.}\FunctionTok{calculateOffsets}\OperatorTok{(}\NormalTok{resultDimension}\OperatorTok{);}
\DataTypeTok{int}\NormalTok{ tableSize }\OperatorTok{=}\NormalTok{ numVariables }\OperatorTok{\textgreater{}} \DecValTok{0} \OperatorTok{?}\NormalTok{ resultDimension}\OperatorTok{[}\NormalTok{numVariables }\OperatorTok{{-}} \DecValTok{1}\OperatorTok{]} \OperatorTok{*}\NormalTok{ offsets}\OperatorTok{[}\NormalTok{numVariables }\OperatorTok{{-}} \DecValTok{1}\OperatorTok{]} \OperatorTok{:} \DecValTok{1}\OperatorTok{;}
\DataTypeTok{double}\OperatorTok{[]}\NormalTok{ resultValues }\OperatorTok{=} \KeywordTok{new} \DataTypeTok{double}\OperatorTok{[}\NormalTok{tableSize}\OperatorTok{];}
\end{Highlighting}
\end{Shaded}

core/src/main/java/org/openmarkov/core/model/network/potential/operation/TablePotentialElimination.java:115-121

\begin{Shaded}
\begin{Highlighting}[]
\DataTypeTok{int}\NormalTok{ resultSize }\OperatorTok{=}\NormalTok{ TablePotential}\OperatorTok{.}\FunctionTok{computeTableSize}\OperatorTok{(}\NormalTok{variablesToKeep}\OperatorTok{);}
\DataTypeTok{double}\OperatorTok{[]}\NormalTok{ resultValues }\OperatorTok{=} \KeywordTok{new} \DataTypeTok{double}\OperatorTok{[}\NormalTok{resultSize}\OperatorTok{];}
\DataTypeTok{int}\NormalTok{ eliminationSize }\OperatorTok{=} \DecValTok{1}\OperatorTok{;}
\ControlFlowTok{for} \OperatorTok{(}\NormalTok{Variable variable }\OperatorTok{:}\NormalTok{ variablesToEliminate}\OperatorTok{)} \OperatorTok{\{}
\NormalTok{    eliminationSize }\OperatorTok{*=}\NormalTok{ variable}\OperatorTok{.}\FunctionTok{getNumStates}\OperatorTok{();}
\OperatorTok{\}}
\end{Highlighting}
\end{Shaded}
\paragraph{Cómo arreglarlo}
Sustituir los cuatro cálculos de \texttt{table\allowbreak{}Size} por \texttt{Table\allowbreak{}Potential.\allowbreak{}compute\allowbreak{}Table\allowbreak{}Size(\allowbreak{}result\allowbreak{}Variables)} (o \texttt{merged\allowbreak{}Variables}) antes de reservar nada, y calcular \texttt{elimination\allowbreak{}Size} con el mismo método sobre \texttt{variables\allowbreak{}To\allowbreak{}Eliminate}; \texttt{calculate\allowbreak{}Offsets} debería usar también \texttt{Math.\allowbreak{}multiply\allowbreak{}Exact}. Así toda operación que se pase del límite lanza la misma \texttt{Invalid\allowbreak{}Argument\allowbreak{}Exception} antes de reservar memoria.

Conviene buscar otros productos de números de estados en \texttt{int} en los módulos \texttt{core} e \texttt{inference} (el de \texttt{Hybrid\allowbreak{}Elimination.\allowbreak{}clique\allowbreak{}Size} ya se corrigió en el commit \texttt{075ee9f}).

Pruebas: los cuatro casos medidos convertidos en pruebas unitarias, esperando \texttt{Invalid\allowbreak{}Argument\allowbreak{}Exception} en todos, incluida \texttt{multiply\allowbreak{}And\allowbreak{}Marginalize} con 31 y 32 variables eliminadas.

\setcounter{subsection}{2}
\subsection[Copiar un potencial a una red sin todas sus variables mete nulos]{\textcolor{corregidoverde}{[Corregido]} Copiar un potencial a una red que no tiene todas sus variables mete nulos en su lista de variables sin avisar}\label{err:3}
\begin{ficha}
\fichakey{Estado}Corregido en el commit \texttt{38e9782} (23-09-2026): \texttt{Potential.\allowbreak{}deep\allowbreak{}Copy} lanza una \texttt{Invalid\allowbreak{}Argument\allowbreak{}Exception} que nombra la variable, la red y el potencial cuando la red de destino no tiene alguna de sus variables, en vez de meter un nulo; lo mismo la variable de tiempo de \texttt{Weibull\allowbreak{}Hazard\allowbreak{}Potential} y las variables de los árboles (\texttt{Tree\allowbreak{}ADDPotential}, \texttt{Tree\allowbreak{}ADDBranch}).
\fichakey{Módulo}core
\fichakey{Paquete}org.openmarkov.core.model.network.potential
\fichakey{Ficheros}\path{core/src/main/java/org/openmarkov/core/model/network/potential/Potential.java:801-821} \newline \path{core/src/main/java/org/openmarkov/core/model/network/potential/TableWithEvents.java:446-455} \newline \path{core/src/main/java/org/openmarkov/core/model/network/potential/UnivariateDistrPotential.java:401-408}
\fichakey{Comprobado}Leyendo el código y ejecutándolo
\fichakey{Esfuerzo}Pequeño (horas)
\fichakey{Decisión de equipo}No
\fichakey{Relacionados}\hyperref[err:50]{error~50}, \hyperref[err:166]{error~166}, \hyperref[nota:180]{nota~180}
\end{ficha}
\paragraph{Qué pasa}
Quien copia un potencial con \texttt{Potential.\allowbreak{}deep\allowbreak{}Copy(\allowbreak{}Prob\allowbreak{}Net\ copy\allowbreak{}Net)} a una red a la que le falta alguna de sus variables recibe, sin ningún aviso, un potencial mal formado: con un \texttt{null} en su lista de variables. En la mayoría de subclases el fallo aparece más tarde y lejos, cuando alguien pide el número de estados de esa variable. En \texttt{Table\allowbreak{}With\allowbreak{}Events} y su hija \texttt{Univariate\allowbreak{}Distr\allowbreak{}Potential} aparece en la línea siguiente, con \texttt{Null\allowbreak{}Pointer\allowbreak{}Exception}.

El camino de producción comprobado que llega aquí es \texttt{Utility\allowbreak{}Operations.\allowbreak{}transform\allowbreak{}To\allowbreak{}Unicriterion}, que copia utilidades sobre la red de una rama del árbol de decisión que ya no tiene una variable. Lo que ve el usuario en ese camino, el árbol de decisión que no se construye, está descrito en \hyperref[err:50]{error~50}. Arreglar esta copia no hace que el árbol salga, pero convierte un puntero nulo en un mensaje que nombra el potencial y la variable que faltan.

Por qué pasa: \texttt{Potential.\allowbreak{}deep\allowbreak{}Copy} crea el potencial nuevo por reflexión y le pone como variables las de \texttt{copy\allowbreak{}Net} con el mismo nombre. Si alguna no existe en \texttt{copy\allowbreak{}Net}, \texttt{get\allowbreak{}Variable} devuelve \texttt{null} y ese \texttt{null} entra en la lista de variables del potencial copiado sin ninguna comprobación. \texttt{Table\allowbreak{}With\allowbreak{}Events.\allowbreak{}deep\allowbreak{}Copy} recorre enseguida la lista recién puesta, en \texttt{set\allowbreak{}Event\allowbreak{}As\allowbreak{}States}, y pregunta a cada variable su tipo.

El defecto es de contrato: una copia que no puede resolver una variable debería decirlo, no devolver un potencial mal formado.
\paragraph{El código}
core/src/main/java/org/openmarkov/core/model/network/potential/Potential.java:811-816

\begin{Shaded}
\begin{Highlighting}[]
\BuiltInTok{List}\OperatorTok{\textless{}}\NormalTok{Variable}\OperatorTok{\textgreater{}}\NormalTok{ newReferences }\OperatorTok{=} \KeywordTok{new} \BuiltInTok{ArrayList}\OperatorTok{\textless{}\textgreater{}();}
\ControlFlowTok{for} \OperatorTok{(}\NormalTok{Variable variable }\OperatorTok{:} \KeywordTok{this}\OperatorTok{.}\FunctionTok{variables}\OperatorTok{)} \OperatorTok{\{}
\NormalTok{    newReferences}\OperatorTok{.}\FunctionTok{add}\OperatorTok{(}\NormalTok{copyNet}\OperatorTok{.}\FunctionTok{getVariable}\OperatorTok{(}\NormalTok{variable}\OperatorTok{.}\FunctionTok{getName}\OperatorTok{()));}
\OperatorTok{\}}

\NormalTok{potential}\OperatorTok{.}\FunctionTok{setVariables}\OperatorTok{(}\NormalTok{newReferences}\OperatorTok{);}
\end{Highlighting}
\end{Shaded}

core/src/main/java/org/openmarkov/core/model/network/potential/TableWithEvents.java:446-447

\begin{Shaded}
\begin{Highlighting}[]
\AttributeTok{@Override} \KeywordTok{public}\NormalTok{ Potential }\FunctionTok{deepCopy}\OperatorTok{(}\NormalTok{ProbNet copyNet}\OperatorTok{)} \OperatorTok{\{}
\NormalTok{    TableWithEvents potential }\OperatorTok{=} \OperatorTok{(}\NormalTok{TableWithEvents}\OperatorTok{)} \KeywordTok{super}\OperatorTok{.}\FunctionTok{deepCopy}\OperatorTok{(}\NormalTok{copyNet}\OperatorTok{);}
\NormalTok{    potential}\OperatorTok{.}\FunctionTok{setEventAsStates}\OperatorTok{(}\NormalTok{potential}\OperatorTok{.}\FunctionTok{variables}\OperatorTok{.}\FunctionTok{subList}\OperatorTok{(}\DecValTok{1}\OperatorTok{,}\NormalTok{ potential}\OperatorTok{.}\FunctionTok{variables}\OperatorTok{.}\FunctionTok{size}\OperatorTok{()));}
\end{Highlighting}
\end{Shaded}
\paragraph{Cómo arreglarlo}
Que \texttt{Potential.\allowbreak{}deep\allowbreak{}Copy} falle con una excepción con mensaje (qué potencial, qué variable, qué red) cuando una variable no se resuelve, en vez de insertar \texttt{null}. La misma idea forma parte de la propuesta de diseño \hyperref[nota:180]{nota~180}; allí no hay nada que corregir.

Hay que revisar las subclases que redefinen \texttt{deep\allowbreak{}Copy} y construyen listas auxiliares a partir de las variables (\texttt{Table\allowbreak{}With\allowbreak{}Events}, \texttt{Univariate\allowbreak{}Distr\allowbreak{}Potential}, \texttt{ICIPotential} y las demás que llaman a \texttt{super.\allowbreak{}deep\allowbreak{}Copy}), para que no hagan nada con la lista antes de que la base haya validado. Los llamadores de producción (\texttt{Utility\allowbreak{}Operations.\allowbreak{}transform\allowbreak{}To\allowbreak{}Unicriterion}, \texttt{apply\allowbreak{}CEUtility\allowbreak{}Scaling} y las copias de red) deben decidir si propagan la excepción o si es un error de programación.

Prueba: copiar con \texttt{deep\allowbreak{}Copy} un \texttt{Table\allowbreak{}Potential} y un \texttt{Univariate\allowbreak{}Distr\allowbreak{}Potential} a una red a la que le falte uno de sus padres y comprobar que se lanza la excepción declarada con el nombre de la variable, y que con todas las variables presentes la copia sigue siendo igual que antes.

% ══════════════════════════════════════════════════════════════════════
\section{Modelos canónicos y modelo de ajuste}\label{sec:canonicos}

Los modelos canónicos (OR, MAX y MIN ruidosos) describen la influencia de muchas causas sobre un efecto con pocos parámetros; el modelo de ajuste (tuning) es una variante de tres estados.

\setcounter{subsection}{3}
\subsection[El modelo de ajuste se proyecta con otro orden y otros decimales cada vez]{\textcolor{corregidoverde}{[Corregido]} La proyección de un modelo de ajuste (tuning) cambia de una ejecución a otra el orden de las variables de su tabla y los últimos decimales de las posteriores}\label{err:4}
\begin{ficha}
\fichakey{Estado}Corregido en el commit \texttt{ee3200d} (23-09-2026): las variables a eliminar se recorren en un \texttt{Linked\allowbreak{}Hash\allowbreak{}Set} y la tabla proyectada de todo modelo canónico sale en el orden declarado del potencial, así que la proyección del modelo de ajuste da el mismo orden y los mismos valores en cada lectura y ejecución; también \texttt{get\allowbreak{}CPT} usa \texttt{Linked\allowbreak{}Hash\allowbreak{}Set}.
\fichakey{Módulo}core
\fichakey{Paquete}org.openmarkov.core.model.network.potential.canonical
\fichakey{Ficheros}\path{core/src/main/java/org/openmarkov/core/model/network/potential/canonical/ICIPotential.java:201-231} \newline \path{core/src/main/java/org/openmarkov/core/model/network/potential/Potential.java:245-256} \newline \path{core/src/main/java/org/openmarkov/core/model/network/potential/canonical/MinMaxPotential.java:124-136}
\fichakey{Comprobado}Leyendo el código y ejecutándolo
\fichakey{Esfuerzo}Pequeño (horas)
\fichakey{Decisión de equipo}No
\fichakey{Relacionados}\hyperref[err:136]{error~136}
\end{ficha}
\paragraph{Qué pasa}
En una red con un nodo cuyo potencial es un modelo de ajuste («tuning»), la tabla proyectada de ese nodo sale con los padres en un orden distinto en cada ejecución, y las posteriores de la red no coinciden bit a bit de una ejecución a otra.

Dentro del programa las operaciones con tablas se guían por la lista de variables de cada tabla, así que el cambio de orden por sí solo no produce un número equivocado. Sí afecta a quien lea los valores por posición suponiendo el orden declarado (exportaciones, comparaciones con valores de referencia, pruebas): lee otra distribución. Y las últimas cifras cambian: una prueba fijada a valores exactos falla unas veces sí y otras no, y un resultado publicado no es exactamente reproducible.

Camino de llegada: toda propagación, evaluación o consulta de la tabla que pase por la proyección del modelo. El modelo de ajuste se elige desde el diálogo de propiedades del nodo para variables de tres estados, y se lee de ficheros \texttt{.\allowbreak{}pgmx} con \texttt{\textless{}Model\textgreater{}Tuning\textless{}/\allowbreak{}Model\textgreater{}}.

Qué es proyectar el modelo. Un modelo canónico ICI (de influencia causal independiente: OR/MAX ruidoso, AND/MIN ruidoso o modelo de ajuste) no guarda su tabla de probabilidad completa. La calcula multiplicando unas tablas pequeñas ---una por padre, una para la fuga (la parte de la probabilidad que no viene de ningún padre) y la de la función que combina--- que hablan de unas variables intermedias que el modelo inventa: las «variables auxiliares» o variables z, y la variable de la fuga. Proyectar el modelo, es decir, obtener su tabla con la evidencia aplicada, consiste en multiplicar esas tablas y sumar fuera las variables auxiliares.

Por qué pasa: \texttt{ICIPotential.\allowbreak{}table\allowbreak{}Project} reúne las variables que hay que sumar fuera en un \texttt{Hash\allowbreak{}Set\textless{}Variable\textgreater{}}, las pasa a una lista y las elimina empezando por la última. \texttt{Variable} no redefine \texttt{equals} ni \texttt{hash\allowbreak{}Code}, así que el orden de recorrido de un \texttt{Hash\allowbreak{}Set} de variables depende de las direcciones de memoria de los objetos. Ese orden cambia de una ejecución del programa a otra, e incluso entre dos lecturas del mismo fichero dentro de la misma ejecución, porque cada lectura crea objetos nuevos. Con el orden de eliminación cambian el orden en que se multiplican las tablas intermedias, el tamaño de esas tablas, el orden de las operaciones de coma flotante y el orden de las variables de la tabla que sale (salvo la primera, que \texttt{with\allowbreak{}Conditioned\allowbreak{}Variable\allowbreak{}First} fija en la variable condicionada).

Solo el modelo de ajuste usa este método: las familias MAX y MIN tienen su propia proyección (\texttt{Min\allowbreak{}Max\allowbreak{}Potential.\allowbreak{}table\allowbreak{}Project}, líneas 124-136), que recorre una lista en orden fijo, y por eso el MAX ruidoso medido abajo sale siempre igual.

Medido sobre la red \texttt{io/\allowbreak{}src/\allowbreak{}test/\allowbreak{}resources/\allowbreak{}test-ici-reading.\allowbreak{}pgmx}, que tiene un modelo de ajuste sobre \texttt{I} con padres \texttt{E}, \texttt{F} y \texttt{G}, y un MAX ruidoso sobre \texttt{H} con padres \texttt{A}, \texttt{B} y \texttt{C}.

Seis ejecuciones seguidas del mismo programa:

\begin{itemize}
\tightlist
\item
  La tabla proyectada de \texttt{I} salió con las variables en los órdenes \texttt{{[}I,\allowbreak{}\ G,\allowbreak{}\ F,\allowbreak{}\ E{]}}, \texttt{{[}I,\allowbreak{}\ F,\allowbreak{}\ E,\allowbreak{}\ G{]}}, \texttt{{[}I,\allowbreak{}\ E,\allowbreak{}\ G,\allowbreak{}\ F{]}}, \texttt{{[}I,\allowbreak{}\ F,\allowbreak{}\ G,\allowbreak{}\ E{]}}, \texttt{{[}I,\allowbreak{}\ E,\allowbreak{}\ F,\allowbreak{}\ G{]}} y \texttt{{[}I,\allowbreak{}\ G,\allowbreak{}\ E,\allowbreak{}\ F{]}}. Los números son los mismos, pero en otras posiciones.
\item
  La propagación por eliminación de variables (\texttt{VEPropagation}) de la red entera dio P(I = tercer estado) = \texttt{0.\allowbreak{}0602075} en unas ejecuciones y \texttt{0.\allowbreak{}06020749999999999} en otras. La huella de todas las posteriores tomó dos valores distintos entre ejecuciones.
\end{itemize}

Cuatro lecturas seguidas del fichero en una misma ejecución, proyectando sin evidencia \texttt{H} e \texttt{I} y comparando las tablas después de reordenarlas al orden de variables del potencial:

\begin{Verbatim}[breaklines,breakanywhere,fontsize=\small]
H MaxPotential    order: H A B C   canonical-hash=1120680360   (igual en las ocho proyecciones)
I TuningPotential order: I G F E   first6=[0.9984000000000001, 0.0016, ...]     canonical-hash=-349860942
I TuningPotential order: I F E G   first6=[0.9984, 0.0016000000000000003, ...]  canonical-hash=-630171156
I TuningPotential order: I F G E   ...                                          canonical-hash=827132973
I TuningPotential order: I G E F   ...                                          canonical-hash=1234376430
\end{Verbatim}

Aun reordenadas al mismo orden, las tablas de \texttt{I} no son idénticas bit a bit (\texttt{0.\allowbreak{}9984} frente a \texttt{0.\allowbreak{}9984000000000001}).

\texttt{Potential.\allowbreak{}get\allowbreak{}CPT} (línea 248) también reúne variables en un \texttt{Hash\allowbreak{}Set}, pero sin consecuencia: el conjunto se llena con las variables de la proyección del propio potencial y luego se le quitan las variables del potencial, así que en la práctica siempre queda vacío y no hay nada que ordenar.
\paragraph{El código}
core/src/main/java/org/openmarkov/core/model/network/potential/canonical/ICIPotential.java:201-231

\begin{Shaded}
\begin{Highlighting}[]
\KeywordTok{public} \AttributeTok{@NotNull}\NormalTok{ TablePotential }\FunctionTok{tableProject}\OperatorTok{(}\NormalTok{EvidenceCase evidenceCase}\OperatorTok{,}\NormalTok{ InferenceOptions inferenceOptions}\OperatorTok{,} \BuiltInTok{List}\OperatorTok{\textless{}}\NormalTok{TablePotential}\OperatorTok{\textgreater{}}\NormalTok{ projectedPotentials}\OperatorTok{)} \KeywordTok{throws}\NormalTok{ NonProjectablePotentialException }\OperatorTok{\{}
    \BuiltInTok{List}\OperatorTok{\textless{}}\NormalTok{TablePotential}\OperatorTok{\textgreater{}}\NormalTok{ potentials }\OperatorTok{=} \FunctionTok{internalTableProject}\OperatorTok{(}\NormalTok{evidenceCase}\OperatorTok{,}\NormalTok{ inferenceOptions}\OperatorTok{);}
    \BuiltInTok{HashSet}\OperatorTok{\textless{}}\NormalTok{Variable}\OperatorTok{\textgreater{}}\NormalTok{ variablesToEliminate }\OperatorTok{=} \KeywordTok{new} \BuiltInTok{HashSet}\OperatorTok{\textless{}\textgreater{}();}
    \CommentTok{// Fill it with variables appearing in all potentials except this}
    \ControlFlowTok{for} \OperatorTok{(}\NormalTok{TablePotential tablePotential }\OperatorTok{:}\NormalTok{ potentials}\OperatorTok{)} \OperatorTok{\{}
\NormalTok{        variablesToEliminate}\OperatorTok{.}\FunctionTok{addAll}\OperatorTok{(}\NormalTok{tablePotential}\OperatorTok{.}\FunctionTok{getVariables}\OperatorTok{());}
    \OperatorTok{\}}
\NormalTok{    variables}\OperatorTok{.}\FunctionTok{forEach}\OperatorTok{(}\NormalTok{variablesToEliminate}\OperatorTok{::}\NormalTok{remove}\OperatorTok{);}

    \BuiltInTok{List}\OperatorTok{\textless{}}\NormalTok{Variable}\OperatorTok{\textgreater{}}\NormalTok{ allVariables }\OperatorTok{=} \KeywordTok{new} \BuiltInTok{ArrayList}\OperatorTok{\textless{}\textgreater{}(}\NormalTok{variables}\OperatorTok{);}
\NormalTok{    allVariables}\OperatorTok{.}\FunctionTok{addAll}\OperatorTok{(}\NormalTok{variablesToEliminate}\OperatorTok{);}
    \ControlFlowTok{while} \OperatorTok{(}\NormalTok{allVariables}\OperatorTok{.}\FunctionTok{size}\OperatorTok{()} \OperatorTok{\textgreater{}}\NormalTok{ variables}\OperatorTok{.}\FunctionTok{size}\OperatorTok{())} \OperatorTok{\{}
\NormalTok{        Variable variableToEliminate }\OperatorTok{=}\NormalTok{ allVariables}\OperatorTok{.}\FunctionTok{getLast}\OperatorTok{();}
\NormalTok{        allVariables}\OperatorTok{.}\FunctionTok{removeLast}\OperatorTok{();}
        \KeywordTok{...}
\NormalTok{        potentials}\OperatorTok{.}\FunctionTok{addFirst}\OperatorTok{(}\NormalTok{DiscretePotentialOperations}\OperatorTok{.}\FunctionTok{multiplyAndMarginalize}\OperatorTok{(}\NormalTok{relatedPotentials}\OperatorTok{,}\NormalTok{ allVariables}\OperatorTok{));}
    \OperatorTok{\}}
    \ControlFlowTok{return} \FunctionTok{withConditionedVariableFirst}\OperatorTok{(}\NormalTok{DiscretePotentialOperations}\OperatorTok{.}\FunctionTok{multiplyAndMarginalize}\OperatorTok{(}\NormalTok{potentials}\OperatorTok{,}\NormalTok{ variables}\OperatorTok{));}
\OperatorTok{\}}
\end{Highlighting}
\end{Shaded}

core/src/main/java/org/openmarkov/core/model/network/potential/Potential.java:248

\begin{Shaded}
\begin{Highlighting}[]
\BuiltInTok{HashSet}\OperatorTok{\textless{}}\NormalTok{Variable}\OperatorTok{\textgreater{}}\NormalTok{ variablesToEliminate }\OperatorTok{=} \KeywordTok{new} \BuiltInTok{HashSet}\OperatorTok{\textless{}\textgreater{}();}
\end{Highlighting}
\end{Shaded}
\paragraph{Cómo arreglarlo}
Recorrer las variables a eliminar en un orden fijo. Hay dos maneras:

\begin{itemize}
\tightlist
\item
  Cambiar el \texttt{Hash\allowbreak{}Set} de \texttt{ICIPotential.\allowbreak{}table\allowbreak{}Project} por un \texttt{Linked\allowbreak{}Hash\allowbreak{}Set} (orden de inserción). Basta, porque \texttt{internal\allowbreak{}Table\allowbreak{}Project} devuelve los potenciales en una lista de orden fijo y cada tabla tiene su lista de variables en orden fijo, así que el orden de inserción ya es determinista.
\item
  Construir la lista de variables a eliminar directamente en el orden que la clase ya guarda: las variables auxiliares en el orden del \texttt{Linked\allowbreak{}Hash\allowbreak{}Map\ z\allowbreak{}Variables} (uno por padre, en el orden de los padres; se construye en los constructores, líneas 90-93 y 111-114) y después la variable de la fuga.
\end{itemize}

Además, la tabla que sale debería llevar las variables en el orden declarado del potencial (hijo y padres en el orden de \texttt{variables}), no solo el hijo delante: \texttt{multiply\allowbreak{}And\allowbreak{}Marginalize(\allowbreak{}potentials,\allowbreak{}\ variables)} las deja en el orden de la unión de los factores. Un \texttt{reorder(\allowbreak{}variables)} al final, en lugar de \texttt{with\allowbreak{}Conditioned\allowbreak{}Variable\allowbreak{}First}, o pasar \texttt{variables\allowbreak{}To\allowbreak{}Keep} explícitamente, lo resuelve, y así el orden tampoco depende de cómo multiplique \texttt{multiply\allowbreak{}And\allowbreak{}Marginalize}.

Cambiar también el \texttt{Hash\allowbreak{}Set} de \texttt{Potential.\allowbreak{}get\allowbreak{}CPT} (línea 248) por \texttt{Linked\allowbreak{}Hash\allowbreak{}Set}, por coherencia, sin esperar ningún efecto.

El mismo patrón («conjunto de dispersión sobre variables sin \texttt{hash\allowbreak{}Code} propio») ya se corrigió en los árboles (\texttt{Tree\allowbreak{}ADDPotential.\allowbreak{}blend\allowbreak{}Potentials} usa un \texttt{Sequenced\allowbreak{}Map} desde el commit \texttt{c9a13a5}), en la maximización, en la heurística de eliminación de Cano y Moral y en el orden parcial de las decisiones. Conviene buscar si queda otro \texttt{Hash\allowbreak{}Set} o \texttt{Hash\allowbreak{}Map} de variables que se recorra para decidir un orden, por ejemplo con \texttt{grep\ -rn\ "Hash\allowbreak{}Set\textless{}Variable\textgreater{}\textbackslash{}\textbar{}Hash\allowbreak{}Map\textless{}Variable"\ core/\allowbreak{}src/\allowbreak{}main}.

Pruebas de regresión:

\begin{itemize}
\tightlist
\item
  En \texttt{core/\allowbreak{}src/\allowbreak{}test}, leer \texttt{test-ici-reading.\allowbreak{}pgmx} varias veces (cada lectura crea objetos nuevos), proyectar el nodo \texttt{I} y comprobar que la lista de variables es siempre \texttt{I,\allowbreak{}\ E,\allowbreak{}\ F,\allowbreak{}\ G} y que los valores son idénticos entre lecturas (\texttt{assert\allowbreak{}Array\allowbreak{}Equals} sin tolerancia) e iguales bit a bit a una referencia.
\item
  Para la reproducibilidad entre ejecuciones, lanzar la propagación de \texttt{test-ici-reading.\allowbreak{}pgmx} y comparar bit a bit con valores fijados, repitiendo la prueba en varias ejecuciones de la máquina virtual o creando las variables en distinto orden para mover sus códigos de dispersión.
\end{itemize}

\setcounter{subsection}{4}
\subsection[Parámetros que no suman uno se abren sin aviso y el muestreo discrepa]{\textcolor{corregidoverde}{[Corregido en parte]} Un fichero puede traer parámetros de un modelo canónico (o una tabla) cuyas columnas no suman uno; se abre sin aviso y el muestreo y el cálculo exacto dan respuestas distintas}\label{err:5}
\begin{ficha}
\fichakey{Estado}Presente en parte. El commit \texttt{9dc49bb} (23-09-2026) optó por no rechazar: al abrir un fichero en la ventana se avisa de las columnas (de tablas y de parámetros de modelos canónicos, fuga incluida) que se apartan de uno más de 10\textsuperscript{-3}, y las que se apartan menos se normalizan en la copia de la red de cada algoritmo de inferencia sin tocar la del usuario; sigue sin normalizarse el generador de bases de casos, que muestrea la red del usuario, y fuera de la ventana (lectores, llamadores del programa) nada avisa.
\fichakey{Módulo}core
\fichakey{Paquete}org.openmarkov.core.model.network.potential.canonical
\fichakey{Ficheros}\path{core/src/main/java/org/openmarkov/core/model/network/potential/canonical/ICIPotential.java:294-304} \newline \path{core/src/main/java/org/openmarkov/core/model/network/potential/canonical/ICIPotential.java:391-401} \newline \path{core/src/main/java/org/openmarkov/core/model/network/potential/canonical/ICIPotential.java:568-598} \newline \path{core/src/main/java/org/openmarkov/core/model/network/potential/TablePotential.java:68-75} \newline \path{io/src/main/java/org/openmarkov/io/probmodel/reader/PGMXPotentialParsers.java:141-163}
\fichakey{Comprobado}Leyendo el código y ejecutándolo
\fichakey{Esfuerzo}Pequeño (horas)
\fichakey{Decisión de equipo}\textcolor{decisionrojo}{\textbf{Sí.}} Si los métodos de asignación (y el lector) deben rechazar columnas que no sumen uno, para los modelos canónicos y las tablas a la vez o para ninguno, y con qué tolerancia; cambia qué ficheros se pueden abrir
\fichakey{Relacionados}\hyperref[err:37]{error~37}
\end{ficha}
\paragraph{Qué pasa}
Un fichero \texttt{.\allowbreak{}pgmx} editado a mano o generado por otra herramienta puede traer parámetros de un modelo canónico, o una tabla, cuyas columnas no suman uno. Se abre sin ningún aviso. Después, los algoritmos de muestreo (muestreo lógico, ponderación por verosimilitud, generador de bases de casos) y la eliminación de variables contestan cosas distintas sobre la misma red, y ninguno avisa. No se ha comprobado si el editor de parámetros de la ventana deja escribir esos números.

Por qué pasa: los parámetros de un modelo canónico son, por cada padre, una tabla con una columna por estado del padre, y cada columna debería ser una distribución sobre los estados del hijo (sumar uno); la fuga es otra distribución. \texttt{ICIPotential.\allowbreak{}set\allowbreak{}Noisy\allowbreak{}Parameters} y \texttt{set\allowbreak{}Leaky\allowbreak{}Parameters} solo comprueban la longitud del array, y el lector de \texttt{.\allowbreak{}pgmx} (\texttt{PGMXPotential\allowbreak{}Parsers.\allowbreak{}get\allowbreak{}ICIPotential}) les pasa los números del fichero sin filtrar. El potencial de tabla (\texttt{Table\allowbreak{}Potential}) tampoco comprueba sumas.

El commit \texttt{afa64f9} hizo que los dos bucles de muestreo de \texttt{sample\allowbreak{}Conditioned\allowbreak{}Variable} no se salgan de la columna. Con eso ya no aparecen estados inexistentes, pero una columna que no suma uno sigue entrando, y el sorteo da al último estado de la columna toda la masa que falta.

Medido sobre un OR/MAX ruidoso \texttt{Y} (tres estados) con padre \texttt{X} (dos estados): \texttt{set\allowbreak{}Noisy\allowbreak{}Parameters(\allowbreak{}X,\allowbreak{}\ {[}0.\allowbreak{}05,\allowbreak{}\ 0.\allowbreak{}03,\allowbreak{}\ 0.\allowbreak{}02,\allowbreak{}\ 0,\allowbreak{}\ 0,\allowbreak{}\ 1{]})} y la fuga \texttt{{[}0.\allowbreak{}5,\allowbreak{}\ 0.\allowbreak{}2,\allowbreak{}\ 0.\allowbreak{}2{]}} se aceptan sin queja. Con \texttt{X} en su primer estado, 10000 sorteos dieron \texttt{{[}253,\allowbreak{}\ 328,\allowbreak{}\ 9419{]}}: el 94 \% al último estado. La tabla que usa el cálculo exacto (\texttt{get\allowbreak{}CPT(\allowbreak{})}) para esa misma columna es \texttt{{[}0.\allowbreak{}025,\allowbreak{}\ 0.\allowbreak{}031,\allowbreak{}\ 0.\allowbreak{}034{]}}, que no suma uno. \texttt{new\ Table\allowbreak{}Potential(\allowbreak{}{[}Y,\allowbreak{}\ X{]},\allowbreak{}\ condicionada,\allowbreak{}\ {[}0.\allowbreak{}05,\allowbreak{}\ 0.\allowbreak{}03,\allowbreak{}\ 0.\allowbreak{}02,\allowbreak{}\ 0.\allowbreak{}2,\allowbreak{}\ 0.\allowbreak{}2,\allowbreak{}\ 0.\allowbreak{}2{]})} también se acepta.
\paragraph{El código}
core/src/main/java/org/openmarkov/core/model/network/potential/canonical/ICIPotential.java:294-304

\begin{Shaded}
\begin{Highlighting}[]
\KeywordTok{public} \DataTypeTok{void} \FunctionTok{setNoisyParameters}\OperatorTok{(}\NormalTok{Variable parent}\OperatorTok{,} \DataTypeTok{double}\OperatorTok{[]}\NormalTok{ parameters}\OperatorTok{)} \OperatorTok{\{}
    \DataTypeTok{int}\NormalTok{ row }\OperatorTok{=} \FunctionTok{parametersRowOf}\OperatorTok{(}\NormalTok{parent}\OperatorTok{);}
    \ControlFlowTok{if} \OperatorTok{(}\NormalTok{parameters}\OperatorTok{.}\FunctionTok{length} \OperatorTok{!=}\NormalTok{ variables}\OperatorTok{.}\FunctionTok{getFirst}\OperatorTok{().}\FunctionTok{getNumStates}\OperatorTok{()} \OperatorTok{*}\NormalTok{ parent}\OperatorTok{.}\FunctionTok{getNumStates}\OperatorTok{())} \OperatorTok{\{}
        \ControlFlowTok{throw} \KeywordTok{new} \FunctionTok{UnrecoverableException}\OperatorTok{(}\KeywordTok{new} \FunctionTok{InvalidArgumentException}\OperatorTok{(}\KeywordTok{...}\OperatorTok{));}
    \OperatorTok{\}}
\NormalTok{    expandedPotential }\OperatorTok{=} \KeywordTok{null}\OperatorTok{;}
\NormalTok{    noisyParameters}\OperatorTok{[}\NormalTok{row}\OperatorTok{]} \OperatorTok{=}\NormalTok{ parameters}\OperatorTok{;}
\OperatorTok{\}}
\end{Highlighting}
\end{Shaded}
\paragraph{Cómo arreglarlo}
Hay que decidir antes si se rechazan esas columnas (ver \texttt{decision-de-equipo}).

Si se decide comprobar, los sitios son:

\begin{itemize}
\tightlist
\item
  \texttt{ICIPotential.\allowbreak{}set\allowbreak{}Noisy\allowbreak{}Parameters}, que cubre \texttt{set\allowbreak{}Noisy\allowbreak{}Potentials}, los constructores de copia, \texttt{add\allowbreak{}Variable}/\texttt{remove\allowbreak{}Variable} de las tres familias y el lector;
\item
  \texttt{ICIPotential.\allowbreak{}set\allowbreak{}Leaky\allowbreak{}Parameters};
\item
  \texttt{Tuning\allowbreak{}Potential.\allowbreak{}set\allowbreak{}Noisy\allowbreak{}Parameters}, que construye la tabla desde cuatro números;
\item
  si se decide para las tablas también, el lector de tablas del \texttt{.\allowbreak{}pgmx}, mejor que el constructor de \texttt{Table\allowbreak{}Potential}, que usan las operaciones intermedias, cuyas tablas no suman uno por diseño.
\end{itemize}

La alternativa sin rechazar es avisar al abrir. Encaja con \hyperref[err:37]{error~37}, que trata de que abrir un fichero no comprueba nada de la red al terminar de leerla.

Prueba: abrir un fichero con una columna de parámetros que suma 0,1 da el rechazo o el aviso elegido.

\setcounter{subsection}{5}
\subsection[El modelo de ajuste se ofrece a políticas sin padres y en redes de sucesos]{\textcolor{corregidoverde}{[Corregido]} El diálogo de imponer una política ofrece el modelo de ajuste para una decisión de tres estados sin padres, y en redes de simulación de eventos discretos, donde el OR y el AND sí se ocultan}\label{err:6}
\begin{ficha}
\fichakey{Estado}Corregido en el commit \texttt{f2c5239} (23-09-2026): el modelo de ajuste se ofrece solo donde se ofrecen OR/MAX y AND/MIN, y además con tres estados en todas las variables; ya no aparece para una decisión sin padres con política impuesta ni en redes de simulación de eventos discretos. Las comprobaciones comunes están en \texttt{ICIPotential.\allowbreak{}validate}.
\fichakey{Módulo}core
\fichakey{Paquete}org.openmarkov.core.model.network.potential.canonical
\fichakey{Ficheros}\path{core/src/main/java/org/openmarkov/core/model/network/potential/canonical/TuningPotential.java:83-88} \newline \path{core/src/main/java/org/openmarkov/core/model/network/potential/canonical/MaxPotential.java:60-72} \newline \path{core/src/main/java/org/openmarkov/core/model/network/potential/canonical/MinPotential.java:59-71} \newline \path{core/src/main/java/org/openmarkov/core/model/network/potential/plugin/PotentialUtils.java:93-112} \newline \path{gui/src/main/java/org/openmarkov/gui/dialog/node/PotentialEditPanel.java:662-671} \newline \path{gui/src/main/java/org/openmarkov/gui/dialog/node/ImposePolicyDialog.java:19}
\fichakey{Comprobado}Leyendo el código y ejecutándolo
\fichakey{Esfuerzo}Pequeño (horas)
\fichakey{Decisión de equipo}No
\fichakey{Relacionados}\hyperref[nota:96]{nota~96}
\end{ficha}
\paragraph{Qué pasa}
El desplegable de tipos de potencial del diálogo del nodo, y el del diálogo de imponer una política (\texttt{Impose\allowbreak{}Policy\allowbreak{}Dialog}, que usa el mismo \texttt{Potential\allowbreak{}Edit\allowbreak{}Panel}), ofrece el modelo de ajuste («tuning») en dos casos donde no corresponde: a una decisión de tres estados sin padres a la que se impone una política, y a los nodos de una red de simulación de eventos discretos. En esos mismos casos el OR/MAX y el AND/MIN sí se ocultan.

Si el usuario lo elige para la decisión sin padres, \texttt{instanciate\allowbreak{}Safely(\allowbreak{}Tuning\allowbreak{}Potential.\allowbreak{}class,\allowbreak{}\ {[}D{]},\allowbreak{}\ POLICY)} devuelve un \texttt{Tuning\allowbreak{}Potential} con una sola variable, papel \texttt{conditional\allowbreak{}Probability} (no \texttt{POLICY}), sin parámetros por padre, y cuya tabla es solo la fuga por omisión \texttt{{[}0,\allowbreak{}\ 1,\allowbreak{}\ 0{]}}: una política determinista que elige siempre el estado del medio, presentada como un modelo de ajuste. No se ejecutó el diálogo (es un \texttt{JDialog}); no se ha comprobado lo que el panel del modelo canónico enseña con cero padres.

Por qué pasa: el desplegable se llena llamando por reflexión al \texttt{validate(\allowbreak{}nodo,\allowbreak{}\ variables,\allowbreak{}\ papel)} de cada clase anotada. El del modelo de ajuste es \texttt{if\ (\allowbreak{}!ICIPotential.\allowbreak{}validate(\allowbreak{}.\allowbreak{}.\allowbreak{}.\allowbreak{})\ \&\&\ role\ ==\ Potential\allowbreak{}Role.\allowbreak{}CONDITIONAL\_\allowbreak{}PROBABILITY)\ return\ false;} y después \texttt{variables.\allowbreak{}stream(\allowbreak{}).\allowbreak{}all\allowbreak{}Match(\allowbreak{}v\ -\textgreater{}\ v.\allowbreak{}get\allowbreak{}Num\allowbreak{}States(\allowbreak{})\ ==\ 3)}.

La «y» hace que el requisito de tener algún padre (\texttt{ICIPotential.\allowbreak{}validate} pide más de una variable) solo se aplique cuando el papel es probabilidad condicionada. Con cualquier otro papel ---en particular \texttt{POLICY}, el de una política impuesta a una decisión--- una variable sola de tres estados pasa. Y a diferencia de \texttt{Max\allowbreak{}Potential.\allowbreak{}validate} y \texttt{Min\allowbreak{}Potential.\allowbreak{}validate}, no rechaza las redes de simulación de eventos discretos (\texttt{DESNetwork\allowbreak{}Type}), no limita el papel a probabilidad condicionada o política, ni exige variables de estados finitos o discretizadas.

Medido con \texttt{Potential\allowbreak{}Utils.\allowbreak{}get\allowbreak{}Filtered\allowbreak{}Potential\allowbreak{}Classes}:

\begin{itemize}
\tightlist
\item
  Decisión \texttt{D} de tres estados sin padres, con política impuesta (tabla de papel \texttt{POLICY}) en un diagrama de influencia: la lista contiene \texttt{Tuning\allowbreak{}Potential} y no contiene ni \texttt{Max\allowbreak{}Potential} ni \texttt{Min\allowbreak{}Potential}.
\item
  Nodo de azar de tres estados sin padres: no lo contiene (correcto, ahí el papel es condicionada).
\item
  Red de simulación de eventos discretos, nodo \texttt{B} de tres estados con padre \texttt{A} de tres: contiene \texttt{Tuning\allowbreak{}Potential} y no \texttt{Max\allowbreak{}Potential} ni \texttt{Min\allowbreak{}Potential}.
\end{itemize}
\paragraph{El código}
core/src/main/java/org/openmarkov/core/model/network/potential/canonical/TuningPotential.java:83-88

\begin{Shaded}
\begin{Highlighting}[]
\KeywordTok{public} \DataTypeTok{static} \DataTypeTok{boolean} \FunctionTok{validate}\OperatorTok{(}\BuiltInTok{Node}\NormalTok{ node}\OperatorTok{,} \BuiltInTok{List}\OperatorTok{\textless{}}\NormalTok{Variable}\OperatorTok{\textgreater{}}\NormalTok{ variables}\OperatorTok{,}\NormalTok{ PotentialRole role}\OperatorTok{)} \OperatorTok{\{}
    \ControlFlowTok{if} \OperatorTok{(!}\NormalTok{ICIPotential}\OperatorTok{.}\FunctionTok{validate}\OperatorTok{(}\NormalTok{node}\OperatorTok{,}\NormalTok{ variables}\OperatorTok{,}\NormalTok{ role}\OperatorTok{)} \OperatorTok{\&\&}\NormalTok{ role }\OperatorTok{==}\NormalTok{ PotentialRole}\OperatorTok{.}\FunctionTok{CONDITIONAL\_PROBABILITY}\OperatorTok{)} \OperatorTok{\{}
        \ControlFlowTok{return} \KeywordTok{false}\OperatorTok{;}
    \OperatorTok{\}}
    \ControlFlowTok{return}\NormalTok{ variables}\OperatorTok{.}\FunctionTok{stream}\OperatorTok{().}\FunctionTok{allMatch}\OperatorTok{(}\NormalTok{variable }\OperatorTok{{-}\textgreater{}}\NormalTok{ variable}\OperatorTok{.}\FunctionTok{getNumStates}\OperatorTok{()} \OperatorTok{==} \DecValTok{3}\OperatorTok{);}
\OperatorTok{\}}
\end{Highlighting}
\end{Shaded}

core/src/main/java/org/openmarkov/core/model/network/potential/canonical/MaxPotential.java:60-72

\begin{Shaded}
\begin{Highlighting}[]
\KeywordTok{public} \DataTypeTok{static} \DataTypeTok{boolean} \FunctionTok{validate}\OperatorTok{(}\BuiltInTok{Node}\NormalTok{ node}\OperatorTok{,} \BuiltInTok{List}\OperatorTok{\textless{}}\NormalTok{Variable}\OperatorTok{\textgreater{}}\NormalTok{ variables}\OperatorTok{,}\NormalTok{ PotentialRole role}\OperatorTok{)} \OperatorTok{\{}
    \ControlFlowTok{if} \OperatorTok{(}\NormalTok{node}\OperatorTok{.}\FunctionTok{getProbNet}\OperatorTok{().}\FunctionTok{getNetworkType}\OperatorTok{()} \KeywordTok{instanceof}\NormalTok{ DESNetworkType}\OperatorTok{)} \OperatorTok{\{}
        \ControlFlowTok{return} \KeywordTok{false}\OperatorTok{;}
    \OperatorTok{\}}
    \ControlFlowTok{if} \OperatorTok{(}\NormalTok{role }\OperatorTok{!=}\NormalTok{ PotentialRole}\OperatorTok{.}\FunctionTok{CONDITIONAL\_PROBABILITY} \OperatorTok{\&\&}\NormalTok{ role }\OperatorTok{!=}\NormalTok{ PotentialRole}\OperatorTok{.}\FunctionTok{POLICY}\OperatorTok{)} \OperatorTok{\{}
        \ControlFlowTok{return} \KeywordTok{false}\OperatorTok{;}
    \OperatorTok{\}}
    \ControlFlowTok{if} \OperatorTok{(!}\NormalTok{ICIPotential}\OperatorTok{.}\FunctionTok{validate}\OperatorTok{(}\NormalTok{node}\OperatorTok{,}\NormalTok{ variables}\OperatorTok{,}\NormalTok{ role}\OperatorTok{))} \OperatorTok{\{}
        \ControlFlowTok{return} \KeywordTok{false}\OperatorTok{;}
    \OperatorTok{\}}
    \ControlFlowTok{return}\NormalTok{ variables}\OperatorTok{.}\FunctionTok{stream}\OperatorTok{().}\FunctionTok{allMatch}\OperatorTok{(}\NormalTok{variable }\OperatorTok{{-}\textgreater{}}\NormalTok{ variable}\OperatorTok{.}\FunctionTok{getVariableType}\OperatorTok{()} \OperatorTok{==}\NormalTok{ VariableType}\OperatorTok{.}\FunctionTok{FINITE\_STATES}
            \OperatorTok{||}\NormalTok{ variable}\OperatorTok{.}\FunctionTok{getVariableType}\OperatorTok{()} \OperatorTok{==}\NormalTok{ VariableType}\OperatorTok{.}\FunctionTok{DISCRETIZED}\OperatorTok{);}
\OperatorTok{\}}
\end{Highlighting}
\end{Shaded}
\paragraph{Cómo arreglarlo}
Escribir \texttt{Tuning\allowbreak{}Potential.\allowbreak{}validate} con las mismas cuatro comprobaciones que \texttt{Max\allowbreak{}Potential} y \texttt{Min\allowbreak{}Potential} (red que no sea de simulación de eventos discretos, papel condicionada o política, al menos un padre, variables de estados finitos o discretizadas) más la suya de tres estados, usando la constante \texttt{NUM\_\allowbreak{}STATES} en vez del 3 escrito a mano. Mejor aún, subir las comprobaciones comunes a un método de \texttt{ICIPotential} que llamen las tres clases, para que no vuelvan a separarse. Si el equipo quiere el modelo de ajuste en redes de simulación de eventos discretos, habría que decirlo y documentarlo; hoy los otros dos modelos canónicos se ocultan ahí.

Defecto vecino, que no forma parte de este error: \texttt{Potential\allowbreak{}Utils.\allowbreak{}instanciate\allowbreak{}Safely} busca un constructor \texttt{(\allowbreak{}List,\allowbreak{}\ Potential\allowbreak{}Role)}, y como los tres modelos canónicos solo tienen \texttt{(\allowbreak{}List)}, cae en ese y el potencial nace con papel \texttt{conditional\allowbreak{}Probability} aunque se esté imponiendo una política; pasa también al elegir OR/MAX para una política.

Prueba de regresión: \texttt{get\allowbreak{}Filtered\allowbreak{}Potential\allowbreak{}Classes} no contiene \texttt{Tuning\allowbreak{}Potential} para una decisión de tres estados sin padres ni para un nodo de una red de simulación de eventos discretos, y sí para un nodo de azar de tres estados con padres de tres estados.

% ══════════════════════════════════════════════════════════════════════
\section{Relaciones definidas por fórmulas, regresiones, árboles y sucesos}\label{sec:relaciones}

Los tipos de relación que no son una tabla: fórmulas, regresiones y riesgos de Weibull, árboles (TreeADD) y las tablas que dependen de sucesos.

\setcounter{subsection}{6}
\subsection[Una utilidad por fórmula calcula mal si sus padres no se llaman U1, U2\ldots{}]{\textcolor{corregidoverde}{[Corregido]} Un nodo de utilidad definido por una fórmula calcula con los valores cambiados si sus padres no se llaman U1, U2\ldots{} en el orden de los enlaces, y revienta si se llaman de otra manera}\label{err:7}
\begin{ficha}
\fichakey{Estado}Corregido en el commit \texttt{fd093b0} (23-09-2026): cada tabla se asigna a la variable de su padre, no a \texttt{"U"\ +\ posición}, así que la fórmula da el valor correcto con cualquier nombre y cualquier orden de enlaces; una fórmula que nombra una variable que no es padre sale como \texttt{Unrecoverable\allowbreak{}Exception} con la variable que falta, en vez de \texttt{No\allowbreak{}Such\allowbreak{}Element\allowbreak{}Exception}.
\fichakey{Módulo}core
\fichakey{Paquete}org.openmarkov.core.model.network.potential.operation
\fichakey{Ficheros}\path{core/src/main/java/org/openmarkov/core/model/network/potential/operation/TablePotentialArithmetic.java:583-663} \newline \path{core/src/main/java/org/openmarkov/core/inference/BasicOperations.java:41-60} \newline \path{core/src/main/java/org/openmarkov/core/inference/BasicOperations.java:144-168} \newline \path{core/src/main/java/org/openmarkov/core/action/core/AbsorbParentsEdit.java:31-46} \newline \path{core/src/main/java/org/openmarkov/core/inference/tasks/TaskUtilities.java:119-122}
\fichakey{Comprobado}Leyendo el código y ejecutándolo
\fichakey{Esfuerzo}Pequeño (horas)
\fichakey{Decisión de equipo}No
\fichakey{Relacionados}\hyperref[err:8]{error~8}, \hyperref[err:11]{error~11}, \hyperref[err:12]{error~12}
\end{ficha}
\paragraph{Qué pasa}
Un potencial de función (\texttt{Function\allowbreak{}Potential}, tipo «Function» en el diálogo del nodo y en los ficheros \texttt{.\allowbreak{}pgmx}) define el valor de un nodo numérico ---normalmente un nodo de utilidad--- con una fórmula sobre los valores de sus padres, escrita con los nombres de los padres entre llaves, por ejemplo \texttt{abs(\allowbreak{}\{Coste\})*\{Beneficio\}}. Para calcular con él, OpenMarkov «absorbe» los padres: sustituye el nodo por una tabla con el valor de la fórmula para cada configuración.

Lo que ve el usuario: si los padres no se llaman \texttt{U1}, \texttt{U2}\ldots, la absorción muere con \texttt{No\allowbreak{}Such\allowbreak{}Element\allowbreak{}Exception} sin mensaje. Si se llaman así pero no en el orden de los enlaces, sale un número equivocado sin ningún aviso. Las redes de prueba del repositorio usan justo los nombres \texttt{U1}, \texttt{U2}\ldots{} en orden, por eso no se nota.

Camino de llegada: el menú contextual «Absorber padres» (se habilita con \texttt{Absorb\allowbreak{}Parents\allowbreak{}Validator} cuando el nodo y sus padres son numéricos y los abuelos de estados finitos) y cualquier evaluación con eliminación de variables de una red con nodos numéricos intermedios (\texttt{Variable\allowbreak{}Elimination}, \texttt{Temporal\allowbreak{}Evaluation} y \texttt{MIDTemporal\allowbreak{}Evolution} llaman al preproceso).

Por qué pasa: la absorción la hace \texttt{Table\allowbreak{}Potential\allowbreak{}Arithmetic.\allowbreak{}evaluate\allowbreak{}Function\allowbreak{}Potential}, llamado desde \texttt{Basic\allowbreak{}Operations.\allowbreak{}absorb\allowbreak{}Parent\allowbreak{}Potentials}. El método recibe la lista de tablas de los padres (\texttt{potentials}) y una lista de variables (\texttt{utility\allowbreak{}Variables}) cuyos nombres calcula y no usa: la línea que los usaría está comentada. Para dar valor a cada operando de la fórmula busca, entre las variables del potencial de función, la que se llame exactamente \texttt{"U"\ +\ (\allowbreak{}i\allowbreak{}Potential\ +\ 1)} y le asigna el número de la tabla del padre que ocupa la posición \texttt{i\allowbreak{}Potential}. Si no la encuentra, \texttt{Optional.\allowbreak{}get(\allowbreak{})} lanza \texttt{No\allowbreak{}Such\allowbreak{}Element\allowbreak{}Exception}.

Además, el llamador (\texttt{absorb\allowbreak{}Parent\allowbreak{}Potentials}) le pasa como \texttt{utility\allowbreak{}Variables} las variables de la tabla del primer padre, que no son los operandos, y solo captura las dos excepciones comprobadas que el método declara (y las envuelve como «inalcanzables»), así que la \texttt{No\allowbreak{}Such\allowbreak{}Element\allowbreak{}Exception} sale tal cual.

Medido sobre copias de \texttt{integration\allowbreak{}Tests/\allowbreak{}src/\allowbreak{}main/\allowbreak{}resources/\allowbreak{}networks/\allowbreak{}dan/\allowbreak{}DAN-two-utility-and-chance-function-sv.\allowbreak{}pgmx} (nodo \texttt{U0} con fórmula sobre dos utilidades constantes; se cambió el valor del segundo padre a 5):

\begin{itemize}
\tightlist
\item
  Control, padres \texttt{U1\ =\ -2} y \texttt{U2\ =\ 5}, fórmula \texttt{abs(\allowbreak{}\{U1\})*\{U2\}}: la absorción da \texttt{10}, correcto.
\item
  Padres renombrados a \texttt{Coste} y \texttt{Beneficio} con fórmula \texttt{abs(\allowbreak{}\{Coste\})*\{Beneficio\}}: \texttt{No\allowbreak{}Such\allowbreak{}Element\allowbreak{}Exception:\ No\ value\ present} en \texttt{Table\allowbreak{}Potential\allowbreak{}Arithmetic.\allowbreak{}java:642}, tanto por \texttt{Absorb\allowbreak{}Parents\allowbreak{}Edit} (la opción «Absorber padres» del menú contextual del nodo) como por \texttt{Task\allowbreak{}Utilities.\allowbreak{}absorb\allowbreak{}All\allowbreak{}Intermediate\allowbreak{}Numeric\allowbreak{}Nodes} (el preproceso que hacen las tareas de eliminación de variables antes de evaluar una red que no es solo de azar).
\item
  Mismos nombres \texttt{U1} y \texttt{U2}, pero con el primer enlace saliendo del padre llamado \texttt{U2} (valor -2) y el segundo del llamado \texttt{U1} (valor 5): la fórmula \texttt{abs(\allowbreak{}\{U1\})*\{U2\}} debe valer \texttt{abs(\allowbreak{}5)*(\allowbreak{}-2)\ =\ -10} y la absorción devuelve \texttt{10}. A la variable llamada \texttt{U1} se le da el valor del primer padre, se llame como se llame.
\end{itemize}

De paso, la fórmula se evalúa y su resultado de texto se convierte a número con \texttt{Double.\allowbreak{}parse\allowbreak{}Double} una vez por casilla de la tabla (línea 653), en vez de una sola vez por expresión compilada; es un coste, no un error.
\paragraph{El código}
core/src/main/java/org/openmarkov/core/model/network/potential/operation/TablePotentialArithmetic.java:620-643

\begin{Shaded}
\begin{Highlighting}[]
\BuiltInTok{List}\OperatorTok{\textless{}}\BuiltInTok{String}\OperatorTok{\textgreater{}}\NormalTok{ utilityVariablesNames }\OperatorTok{=}\NormalTok{ utilityVariables}\OperatorTok{.}\FunctionTok{stream}\OperatorTok{()}
                                                     \OperatorTok{.}\FunctionTok{map}\OperatorTok{(}\NormalTok{Variable}\OperatorTok{::}\NormalTok{getName}\OperatorTok{)}
                                                     \OperatorTok{.}\FunctionTok{toList}\OperatorTok{();}
\KeywordTok{...}
\ControlFlowTok{for} \OperatorTok{(}\DataTypeTok{int}\NormalTok{ iPotential }\OperatorTok{=} \DecValTok{0}\OperatorTok{;}\NormalTok{ iPotential }\OperatorTok{\textless{}}\NormalTok{ numPotentials}\OperatorTok{;}\NormalTok{ iPotential}\OperatorTok{++)} \OperatorTok{\{}
    \DataTypeTok{int}\NormalTok{ potentialsPositionIPotential }\OperatorTok{=}\NormalTok{ potentialsPositions}\OperatorTok{[}\NormalTok{iPotential}\OperatorTok{];}
    \BuiltInTok{String}\NormalTok{ varNameInExpressionToEvaluate }\OperatorTok{=} \StringTok{"U"} \OperatorTok{+} \OperatorTok{(}\NormalTok{iPotential }\OperatorTok{+} \DecValTok{1}\OperatorTok{);}
    \CommentTok{// String varNameInExpressionToEvaluate = utilityVariablesNames.get(iPotential);}
\NormalTok{    assignment}\OperatorTok{.}\FunctionTok{put}\OperatorTok{(}
\NormalTok{            utilityPotential}\OperatorTok{.}\FunctionTok{getVariables}\OperatorTok{()}
                            \OperatorTok{.}\FunctionTok{stream}\OperatorTok{()}
                            \OperatorTok{.}\FunctionTok{filter}\OperatorTok{(}\NormalTok{variable }\OperatorTok{{-}\textgreater{}}\NormalTok{ variable}\OperatorTok{.}\FunctionTok{getName}\OperatorTok{().}\FunctionTok{equals}\OperatorTok{(}\NormalTok{varNameInExpressionToEvaluate}\OperatorTok{))}
                            \OperatorTok{.}\FunctionTok{findFirst}\OperatorTok{()}
                            \OperatorTok{.}\FunctionTok{get}\OperatorTok{()}
            \OperatorTok{,} \StringTok{""} \OperatorTok{+}\NormalTok{ tables}\OperatorTok{[}\NormalTok{iPotential}\OperatorTok{][}\NormalTok{potentialsPositionIPotential}\OperatorTok{]);}
\end{Highlighting}
\end{Shaded}

core/src/main/java/org/openmarkov/core/inference/BasicOperations.java:51-59

\begin{Shaded}
\begin{Highlighting}[]
\OperatorTok{\}} \ControlFlowTok{else} \OperatorTok{\{}   \CommentTok{// FunctionPotential}
    \ControlFlowTok{try} \OperatorTok{\{}
        \CommentTok{// To check (probable bug): are these two exceptions really unreachable?}
\NormalTok{        newTable }\OperatorTok{=}\NormalTok{ DiscretePotentialOperations}\OperatorTok{.}\FunctionTok{evaluateFunctionPotential}\OperatorTok{(}
                \OperatorTok{(}\NormalTok{FunctionPotential}\OperatorTok{)}\NormalTok{ nodePotential}\OperatorTok{,}\NormalTok{ parentsPotentials}\OperatorTok{,}\NormalTok{ parentsPotentials}\OperatorTok{.}\FunctionTok{get}\OperatorTok{(}\DecValTok{0}\OperatorTok{).}\FunctionTok{getVariables}\OperatorTok{());}
    \OperatorTok{\}} \ControlFlowTok{catch} \OperatorTok{(}\NormalTok{NonProjectablePotentialException}\OperatorTok{.}\FunctionTok{CannotEvaluate} \OperatorTok{|}
\NormalTok{             NonProjectablePotentialException}\OperatorTok{.}\FunctionTok{CannotResolveVariable}\NormalTok{ e}\OperatorTok{)} \OperatorTok{\{}
        \ControlFlowTok{throw} \KeywordTok{new} \FunctionTok{UnreachableException}\OperatorTok{(}\NormalTok{e}\OperatorTok{);}
    \OperatorTok{\}}
\OperatorTok{\}}
\end{Highlighting}
\end{Shaded}
\paragraph{Cómo arreglarlo}
Emparejar cada tabla con su operando por la variable, no por la posición. La tabla \texttt{i} es la proyección del potencial del padre \texttt{i} de \texttt{network.\allowbreak{}get\allowbreak{}Parents(\allowbreak{}node)}, así que el operando es la variable de ese nodo padre. Ojo: la tabla proyectada de una utilidad constante no lleva esa variable dentro, así que no sirve preguntársela a la tabla.

Para ello, \texttt{absorb\allowbreak{}Parent\allowbreak{}Potentials} debe recibir y pasar la lista de variables de los padres en el mismo orden que \texttt{parents\allowbreak{}Potentials} (hoy pasa las variables de la primera tabla), y \texttt{evaluate\allowbreak{}Function\allowbreak{}Potential} asignar \texttt{assignment.\allowbreak{}put(\allowbreak{}variable\allowbreak{}Del\allowbreak{}Padre\_\allowbreak{}i,\allowbreak{}\ valor)}. Si una variable que la fórmula necesita no aparece, lanzar la excepción declarada (\texttt{Non\allowbreak{}Projectable\allowbreak{}Potential\allowbreak{}Exception.\allowbreak{}Cannot\allowbreak{}Resolve\allowbreak{}Variable}) con el nombre que falta, y que \texttt{absorb\allowbreak{}Parent\allowbreak{}Potentials} deje de declararla inalcanzable y la propague o la convierta en un error con mensaje.

Los sitios que llegan: \texttt{Basic\allowbreak{}Operations.\allowbreak{}absorb\allowbreak{}Parents} (preproceso de inferencia) y \texttt{Absorb\allowbreak{}Parents\allowbreak{}Edit.\allowbreak{}generate\allowbreak{}Edits} (menú); los dos pasan por \texttt{absorb\allowbreak{}Parent\allowbreak{}Potentials}. Opcionalmente, compilar la expresión una vez fuera del bucle.

Pruebas de regresión: con la red de dos utilidades y fórmula \texttt{abs(\allowbreak{}\{Coste\})*\{Beneficio\}}, absorber padres da \texttt{10}; con los padres \texttt{U2} (-2) y \texttt{U1} (5) enlazados en ese orden y la fórmula \texttt{abs(\allowbreak{}\{U1\})*\{U2\}}, da \texttt{-10}; con una fórmula que nombra una variable que no es padre, la excepción dice qué variable falta.

\setcounter{subsection}{7}
\subsection[«Absorber padres» muere con ClassCastException en 12 redes del repositorio]{\textcolor{corregidoverde}{[Corregido]} «Absorber padres» sobre un nodo numérico que no es suma, producto ni fórmula muere con ClassCastException (12 redes del repositorio)}\label{err:8}
\begin{ficha}
\fichakey{Estado}Corregido en el commit \texttt{e682a70} (23-09-2026): «Absorb parents» solo se ofrece, y el preproceso solo absorbe, cuando el potencial del nodo es una suma, un producto o una fórmula; con cualquier otro, \texttt{absorb\allowbreak{}Parent\allowbreak{}Potentials} lanza \texttt{Invalid\allowbreak{}Argument\allowbreak{}Exception} con mensaje en vez de \texttt{Class\allowbreak{}Cast\allowbreak{}Exception}.
\fichakey{Módulo}core
\fichakey{Paquete}org.openmarkov.core.inference
\fichakey{Ficheros}\path{core/src/main/java/org/openmarkov/core/inference/BasicOperations.java:41-60} \newline \path{core/src/main/java/org/openmarkov/core/inference/BasicOperations.java:144-219} \newline \path{core/src/main/java/org/openmarkov/core/inference/BasicOperations.java:279-299} \newline \path{core/src/main/java/org/openmarkov/core/action/core/AbsorbParentsEdit.java:32-46} \newline \path{gui/src/main/java/org/openmarkov/gui/validator/AbsorbParentsValidator.java:24-26} \newline \path{gui/src/main/java/org/openmarkov/gui/menutoolbar/menu/NodeContextualMenu.java:142}
\fichakey{Comprobado}Leyendo el código y ejecutándolo
\fichakey{Esfuerzo}Pequeño (horas)
\fichakey{Decisión de equipo}No
\fichakey{Relacionados}\hyperref[err:7]{error~7}, \hyperref[err:9]{error~9}, \hyperref[err:42]{error~42}, \hyperref[err:112]{error~112}, \hyperref[err:113]{error~113}, \hyperref[err:193]{error~193}, \hyperref[nota:28]{nota~28}
\end{ficha}
\paragraph{Qué pasa}
Absorber los padres de un nodo numérico sustituye su potencial por una tabla calculada a partir de los de sus padres. Si el potencial del nodo no es una suma, un producto ni una fórmula, la operación muere con \texttt{Class\allowbreak{}Cast\allowbreak{}Exception}, y en la aplicación sale como error inesperado.

Llegan dos caminos:

\begin{itemize}
\tightlist
\item
  El menú contextual del nodo, «Absorb parents». \texttt{Node\allowbreak{}Contextual\allowbreak{}Menu} lo enciende con \texttt{Absorb\allowbreak{}Parents\allowbreak{}Validator.\allowbreak{}validate}, que solo mira que el nodo sea numérico, que sus padres sean numéricos y sus abuelos de estados finitos (\texttt{Basic\allowbreak{}Operations.\allowbreak{}have\allowbreak{}Parents\allowbreak{}And\allowbreak{}Are\allowbreak{}All\allowbreak{}Absorbable}), no la clase del potencial. La acción construye \texttt{Absorb\allowbreak{}Parents\allowbreak{}Edit}, que llama a \texttt{absorb\allowbreak{}Parent\allowbreak{}Potentials}.
\item
  El preproceso de los algoritmos exactos, \texttt{Task\allowbreak{}Utilities.\allowbreak{}absorb\allowbreak{}All\allowbreak{}Intermediate\allowbreak{}Numeric\allowbreak{}Nodes} → \texttt{Basic\allowbreak{}Operations.\allowbreak{}absorb\allowbreak{}Parents}, que usa el mismo método. Ahí se llega menos, porque la discretización previa deja de ser numéricas la mayoría de las variables de azar (no medido).
\end{itemize}

Por qué pasa: \texttt{Basic\allowbreak{}Operations.\allowbreak{}absorb\allowbreak{}Parent\allowbreak{}Potentials} (módulo \texttt{core}) decide cómo combinar los potenciales de los padres por la clase del potencial del nodo: \texttt{Sum\allowbreak{}Potential} → suma, \texttt{Product\allowbreak{}Potential} → producto, y cualquier otra cosa → \texttt{(\allowbreak{}Function\allowbreak{}Potential)\ node\allowbreak{}Potential}, sin comprobar.

Medido ejecutando \texttt{Absorb\allowbreak{}Parents\allowbreak{}Edit} sobre cada nodo con el menú encendido en las redes de \texttt{0\_\allowbreak{}miscellaneous/\allowbreak{}networks} (sin las copias de las carpetas \texttt{1-0}): de 30 nodos así, 10 son suma y 3 producto; los otros 17 fallan.

Con \texttt{Class\allowbreak{}Cast\allowbreak{}Exception}, 12 redes: en \texttt{0\_\allowbreak{}miscellaneous/\allowbreak{}networks/\allowbreak{}mid/\allowbreak{}}, \texttt{MID-CHAP-Ryan-Griffin.\allowbreak{}pgmx}, \texttt{MID-Chancellor.\allowbreak{}pgmx}, \texttt{MID-Chancellor-corrected.\allowbreak{}pgmx}, \texttt{MID-Chancellor-new.\allowbreak{}pgmx}, \texttt{MID-Cochlear.\allowbreak{}pgmx}, \texttt{MID-HPV.\allowbreak{}pgmx}, \texttt{MID-HPV-without-supervalue.\allowbreak{}pgmx} y \texttt{MID-hip-Briggs.\allowbreak{}pgmx}; y en \texttt{0\_\allowbreak{}miscellaneous/\allowbreak{}networks/\allowbreak{}mid/\allowbreak{}DMHEE/\allowbreak{}}, \texttt{MID-dmhee-2.\allowbreak{}5.\allowbreak{}pgmx}, \texttt{MID-dmhee-3.\allowbreak{}5.\allowbreak{}pgmx}, \texttt{MID-dmhee-4.\allowbreak{}7.\allowbreak{}pgmx} y \texttt{MID-dmhee-4.\allowbreak{}8.\allowbreak{}pgmx}. Las clases que lo provocan son \texttt{Cycle\allowbreak{}Length\allowbreak{}Shift} (la relación «un ciclo más» de variables como «Age {[}1{]}» o «Time in treatment {[}1{]}»), \texttt{Linear\allowbreak{}Combination\allowbreak{}Potential} y \texttt{Tree\allowbreak{}ADDPotential} («AR QoL {[}0{]}» en \texttt{MID-HPV.\allowbreak{}pgmx}).

En otras tres redes (\texttt{0\_\allowbreak{}miscellaneous/\allowbreak{}networks/\allowbreak{}mid/\allowbreak{}MID-CHD-Walker.\allowbreak{}pgmx}, \texttt{MID-Colorectal.\allowbreak{}pgmx} y \texttt{MID-mammography.\allowbreak{}pgmx}) el fallo es anterior y distinto: \texttt{Absorb\allowbreak{}Parents\allowbreak{}Edit.\allowbreak{}generate\allowbreak{}Edits} pide la tabla del potencial de un padre y recibe \texttt{Potential\allowbreak{}Cannot\allowbreak{}Be\allowbreak{}Converted\allowbreak{}To\allowbreak{}ATable} envuelto en \texttt{Unreachable\allowbreak{}Exception}, es decir, una excepción que el código declara imposible. También sale como error inesperado. Esa causa no es la de esta ficha y alcanza también a nodos de suma y de producto: está en \hyperref[err:9]{error~9}. En estas tres redes el nodo es un \texttt{Cycle\allowbreak{}Length\allowbreak{}Shift}, así que el arreglo de aquí ya apaga su menú.
\paragraph{El código}
core/src/main/java/org/openmarkov/core/inference/BasicOperations.java:46-60

\begin{Shaded}
\begin{Highlighting}[]
\NormalTok{TablePotential newTable}\OperatorTok{;}
\ControlFlowTok{if} \OperatorTok{(}\NormalTok{nodePotential }\KeywordTok{instanceof}\NormalTok{ SumPotential}\OperatorTok{)} \OperatorTok{\{}
\NormalTok{    newTable }\OperatorTok{=}\NormalTok{ DiscretePotentialOperations}\OperatorTok{.}\FunctionTok{sum}\OperatorTok{(}\NormalTok{parentsPotentials}\OperatorTok{);}
\OperatorTok{\}} \ControlFlowTok{else} \ControlFlowTok{if} \OperatorTok{(}\NormalTok{nodePotential }\KeywordTok{instanceof}\NormalTok{ ProductPotential}\OperatorTok{)} \OperatorTok{\{}
\NormalTok{    newTable }\OperatorTok{=}\NormalTok{ DiscretePotentialOperations}\OperatorTok{.}\FunctionTok{multiply}\OperatorTok{(}\NormalTok{parentsPotentials}\OperatorTok{);}
\OperatorTok{\}} \ControlFlowTok{else} \OperatorTok{\{}   \CommentTok{// FunctionPotential}
    \ControlFlowTok{try} \OperatorTok{\{}
        \CommentTok{// To check (probable bug): are these two exceptions really unreachable?}
\NormalTok{        newTable }\OperatorTok{=}\NormalTok{ DiscretePotentialOperations}\OperatorTok{.}\FunctionTok{evaluateFunctionPotential}\OperatorTok{(}
                \OperatorTok{(}\NormalTok{FunctionPotential}\OperatorTok{)}\NormalTok{ nodePotential}\OperatorTok{,}\NormalTok{ parentsPotentials}\OperatorTok{,}\NormalTok{ parentsPotentials}\OperatorTok{.}\FunctionTok{get}\OperatorTok{(}\DecValTok{0}\OperatorTok{).}\FunctionTok{getVariables}\OperatorTok{());}
    \OperatorTok{\}} \ControlFlowTok{catch} \OperatorTok{(}\NormalTok{NonProjectablePotentialException}\OperatorTok{.}\FunctionTok{CannotEvaluate} \OperatorTok{|}
\NormalTok{             NonProjectablePotentialException}\OperatorTok{.}\FunctionTok{CannotResolveVariable}\NormalTok{ e}\OperatorTok{)} \OperatorTok{\{}
        \ControlFlowTok{throw} \KeywordTok{new} \FunctionTok{UnreachableException}\OperatorTok{(}\NormalTok{e}\OperatorTok{);}
    \OperatorTok{\}}
\OperatorTok{\}}
\end{Highlighting}
\end{Shaded}
\paragraph{Cómo arreglarlo}
Añadir la comprobación \texttt{instanceof\ Function\allowbreak{}Potential} y, para cualquier otra clase, no absorber (dejar el nodo como está) o rechazar con una excepción del dominio con mensaje.

La misma condición debe estar en los tres sitios que deciden si se puede absorber: \texttt{Absorb\allowbreak{}Parents\allowbreak{}Validator.\allowbreak{}validate} / \texttt{Basic\allowbreak{}Operations.\allowbreak{}have\allowbreak{}Parents\allowbreak{}And\allowbreak{}Are\allowbreak{}All\allowbreak{}Absorbable} (para que el menú no se ofrezca), \texttt{Basic\allowbreak{}Operations.\allowbreak{}absorb\allowbreak{}All\allowbreak{}Intermediate\allowbreak{}Numeric\allowbreak{}Nodes} (preproceso automático) y \texttt{absorb\allowbreak{}Parent\allowbreak{}Potentials}. La condición de que los potenciales de los padres se puedan convertir en tabla es \hyperref[err:9]{error~9} y debe ir en los mismos sitios.

Prueba de regresión: sobre \texttt{0\_\allowbreak{}miscellaneous/\allowbreak{}networks/\allowbreak{}mid/\allowbreak{}MID-Chancellor-new.\allowbreak{}pgmx}, el menú «Absorb parents» debe estar apagado para «Time in treatment {[}1{]}» (o la edición no debe lanzar), y el recorrido de todas las redes de \texttt{0\_\allowbreak{}miscellaneous/\allowbreak{}networks} no debe dar ninguna excepción al absorber los nodos con el menú encendido.

\setcounter{subsection}{8}
\subsection[«Absorb parents» de una suma falla si un padre es una distribución (formato 1.0)]{\textcolor{corregidoverde}{[Corregido en parte]} «Absorb parents» sobre un nodo de suma o de producto sale como error inesperado cuando un padre es una distribución, como ocurre en las redes guardadas en formato 1.0}\label{err:9}
\begin{ficha}
\fichakey{Estado}Presente en parte. El commit \texttt{551e715} (23-09-2026) hizo que una distribución «Exact» con números, como guarda el formato 1.0 las relaciones exactas, se convierta en tabla y se escale como el \texttt{Exact\allowbreak{}Distr\allowbreak{}Potential} que es, así que los trece nodos de suma y producto medidos ya se absorben y las redes 1.0 se evalúan; sigue que el menú no mira si los potenciales de los padres se pueden convertir en tabla, y \texttt{Absorb\allowbreak{}Parents\allowbreak{}Edit.\allowbreak{}generate\allowbreak{}Edits} y \texttt{Basic\allowbreak{}Operations.\allowbreak{}absorb\allowbreak{}Parents} siguen convirtiendo el fallo en \texttt{Unreachable\allowbreak{}Exception} para cualquier otra distribución de un padre (por ejemplo, una normal). Esto último se ha deducido leyendo el código.
\fichakey{Módulo}core
\fichakey{Paquete}org.openmarkov.core.action.core; org.openmarkov.core.inference
\fichakey{Ficheros}\path{core/src/main/java/org/openmarkov/core/action/core/AbsorbParentsEdit.java:36-46} \newline \path{core/src/main/java/org/openmarkov/core/inference/BasicOperations.java:144-195} \newline \path{core/src/main/java/org/openmarkov/core/inference/BasicOperations.java:279-300} \newline \path{core/src/main/java/org/openmarkov/core/model/network/potential/UnivariateDistrPotential.java:357-364} \newline \path{core/src/main/java/org/openmarkov/core/model/network/potential/UniformPotential.java:91-99} \newline \path{core/src/main/java/org/openmarkov/core/inference/tasks/TaskUtilities.java:119-121} \newline \path{gui/src/main/java/org/openmarkov/gui/validator/AbsorbParentsValidator.java:24-26} \newline \path{gui/src/main/java/org/openmarkov/gui/menutoolbar/menu/NodeContextualMenu.java:142} \newline \path{gui/src/main/java/org/openmarkov/gui/configuration/UserPreferences.java:108}
\fichakey{Comprobado}Leyendo el código y ejecutándolo
\fichakey{Esfuerzo}Medio (días)
\fichakey{Decisión de equipo}\textcolor{decisionrojo}{\textbf{Sí.}} Si «Absorb parents» debe funcionar con padres cuya relación es una distribución exacta (como las guarda el formato 1.0) o solo rechazarse con un mensaje
\fichakey{Relacionados}\hyperref[err:8]{error~8}, \hyperref[err:46]{error~46}, \hyperref[err:50]{error~50}, \hyperref[err:112]{error~112}, \hyperref[err:153]{error~153}, \hyperref[err:193]{error~193}, \hyperref[nota:28]{nota~28}
\end{ficha}
\paragraph{Qué pasa}
«Absorb parents», en el menú contextual de un nodo numérico, sustituye su relación por una tabla calculada con las de sus padres. Si el potencial de algún padre no se puede convertir en tabla, la operación termina en el diálogo de error inesperado. No es un caso raro: en el formato ProbModelXML 1.0 las utilidades y muchos valores numéricos se guardan como distribuciones (\texttt{Univariate\allowbreak{}Distr\allowbreak{}Potential}, con la distribución «exacta»), y ese potencial no se convierte nunca en tabla.

Camino del usuario:

\begin{itemize}
\tightlist
\item
  abrir \texttt{0\_\allowbreak{}miscellaneous/\allowbreak{}networks/\allowbreak{}id/\allowbreak{}1-0/\allowbreak{}ID-arthronet-1-0.\allowbreak{}pgmx}, clic derecho sobre «EVAC Total» (una suma de utilidades) y «Absorb parents»;
\item
  o abrir \texttt{0\_\allowbreak{}miscellaneous/\allowbreak{}networks/\allowbreak{}id/\allowbreak{}ID-arthronet.\allowbreak{}pgmx} (formato 0.2, donde la misma operación funciona), guardarla con «Save network as\ldots» eligiendo el formato «OpenMarkov.1.0» («Save» la guardaría en 0.2, el formato del fichero abierto), reabrirla y repetir: ya falla (\hyperref[err:112]{error~112}).
\end{itemize}

Por qué pasa: \texttt{Absorb\allowbreak{}Parents\allowbreak{}Edit.\allowbreak{}generate\allowbreak{}Edits} (módulo \texttt{core}) pide a cada padre \texttt{get\allowbreak{}Potential(\allowbreak{}).\allowbreak{}table\allowbreak{}Project(\allowbreak{}null,\allowbreak{}\ null)} y convierte \texttt{Non\allowbreak{}Projectable\allowbreak{}Potential\allowbreak{}Exception} en \texttt{Unreachable\allowbreak{}Exception}, es decir, en un fallo que el código declara imposible. \texttt{Univariate\allowbreak{}Distr\allowbreak{}Potential.\allowbreak{}table\allowbreak{}Project} lanza siempre \texttt{Potential\allowbreak{}Cannot\allowbreak{}Be\allowbreak{}Converted\allowbreak{}To\allowbreak{}ATable}, y \texttt{Uniform\allowbreak{}Potential.\allowbreak{}table\allowbreak{}Project} también cuando su variable es numérica. El menú se enciende con \texttt{Absorb\allowbreak{}Parents\allowbreak{}Validator.\allowbreak{}validate} → \texttt{Basic\allowbreak{}Operations.\allowbreak{}have\allowbreak{}Parents\allowbreak{}And\allowbreak{}Are\allowbreak{}All\allowbreak{}Absorbable}, que solo mira los tipos de variable del nodo, de sus padres y de sus abuelos, no sus potenciales.

La misma conversión, con el mismo \texttt{Unreachable\allowbreak{}Exception}, está en \texttt{Basic\allowbreak{}Operations.\allowbreak{}absorb\allowbreak{}Parents}, que usa el preproceso de los algoritmos exactos (\texttt{Task\allowbreak{}Utilities.\allowbreak{}absorb\allowbreak{}All\allowbreak{}Intermediate\allowbreak{}Numeric\allowbreak{}Nodes}): evaluar con eliminación de variables \texttt{0\_\allowbreak{}miscellaneous/\allowbreak{}networks/\allowbreak{}id/\allowbreak{}1-0/\allowbreak{}ID-arthronet-1-0.\allowbreak{}pgmx}, \texttt{ID-mediastinet-1-0.\allowbreak{}pgmx} o \texttt{ID-used-car-buyer-1-0.\allowbreak{}pgmx} muere ahí. Arreglar esta ficha no basta para evaluar esas redes: las 15 de \texttt{0\_\allowbreak{}miscellaneous/\allowbreak{}networks/\allowbreak{}id/\allowbreak{}1-0/\allowbreak{}} fallan al evaluarlas, la mayoría más adelante, al convertir en tabla esas mismas distribuciones (por ejemplo \texttt{ID-decide-test-1-0.\allowbreak{}pgmx}, que no tiene nodos de suma), mientras que sus versiones en formato 0.2 se evalúan; esa causa es \hyperref[err:112]{error~112}.

Medido ejecutando \texttt{Absorb\allowbreak{}Parents\allowbreak{}Edit} sobre cada nodo con el menú encendido en todas las redes de \texttt{0\_\allowbreak{}miscellaneous/\allowbreak{}networks}, incluidas las carpetas \texttt{1-0} (148 se abren):

\begin{itemize}
\tightlist
\item
  13 nodos de suma o de producto fallan con \texttt{Unreachable\allowbreak{}Exception} causada por \texttt{Potential\allowbreak{}Cannot\allowbreak{}Be\allowbreak{}Converted\allowbreak{}To\allowbreak{}ATable}, en 9 redes: en \texttt{0\_\allowbreak{}miscellaneous/\allowbreak{}networks/\allowbreak{}id/\allowbreak{}1-0/\allowbreak{}}, «EVAC Total» y «Coste total» de \texttt{ID-arthronet-1-0.\allowbreak{}pgmx}, «Net\_QALE» y «Total\_Economic\_Cost» de \texttt{ID-mediastinet-1-0.\allowbreak{}pgmx} y «Total» de \texttt{ID-used-car-buyer-1-0.\allowbreak{}pgmx}; en \texttt{0\_\allowbreak{}miscellaneous/\allowbreak{}networks/\allowbreak{}dan/\allowbreak{}1-0/\allowbreak{}}, \texttt{DAN-arthronet-1-0.\allowbreak{}pgmx} (dos), \texttt{DAN-economic-mediastinet-1-0.\allowbreak{}pgmx}, \texttt{DAN-mediastinet-1-0.\allowbreak{}pgmx} (dos), \texttt{DAN-qale-mediastinet-1-0.\allowbreak{}pgmx} y \texttt{DAN-used-car-buyer-1-0.\allowbreak{}pgmx}; y «Treatment cost {[}1{]}» de \texttt{0\_\allowbreak{}miscellaneous/\allowbreak{}networks/\allowbreak{}mid/\allowbreak{}1-0/\allowbreak{}MID-mammography-1-0.\allowbreak{}pgmx}. En los doce primeros los padres son \texttt{Univariate\allowbreak{}Distr\allowbreak{}Potential}; en el último, un árbol cuyas hojas son distribuciones;
\item
  los mismos 13 nodos en las versiones en formato 0.2 de esas redes (por ejemplo \texttt{0\_\allowbreak{}miscellaneous/\allowbreak{}networks/\allowbreak{}id/\allowbreak{}ID-arthronet.\allowbreak{}pgmx}) se absorben sin error, y son las únicas sumas y productos con el menú encendido fuera de las carpetas \texttt{1-0};
\item
  los otros 6 casos con \texttt{Unreachable\allowbreak{}Exception} son «Age {[}1{]}» de \texttt{0\_\allowbreak{}miscellaneous/\allowbreak{}networks/\allowbreak{}mid/\allowbreak{}MID-CHD-Walker.\allowbreak{}pgmx}, \texttt{MID-Colorectal.\allowbreak{}pgmx}, \texttt{MID-mammography.\allowbreak{}pgmx} y sus copias en \texttt{1-0}: un \texttt{Cycle\allowbreak{}Length\allowbreak{}Shift} cuyo padre «Age {[}0{]}» tiene un \texttt{Uniform\allowbreak{}Potential} sobre la variable numérica. Ahí el nodo tampoco se podría absorber aunque el padre se convirtiera en tabla (\hyperref[err:8]{error~8});
\item
  guardar \texttt{0\_\allowbreak{}miscellaneous/\allowbreak{}networks/\allowbreak{}id/\allowbreak{}ID-arthronet.\allowbreak{}pgmx} con \texttt{PGMXWriter\_\allowbreak{}1\_\allowbreak{}0} y reabrirla convierte el \texttt{Exact\allowbreak{}Distr\allowbreak{}Potential} de «EVAC Implante» en \texttt{Univariate\allowbreak{}Distr\allowbreak{}Potential}, y «Absorb parents» sobre «EVAC Total» pasa de funcionar a lanzar \texttt{Unreachable\allowbreak{}Exception:\ Potential\ P(\allowbreak{}EVAC\ Implante\ \textbar{}\ Realizar\ Implante)\ cannot\ be\ converted\ to\ a\ table}.
\end{itemize}
\paragraph{El código}
core/src/main/java/org/openmarkov/core/action/core/AbsorbParentsEdit.java:36-46

\begin{Shaded}
\begin{Highlighting}[]
\BuiltInTok{List}\OperatorTok{\textless{}}\BuiltInTok{Node}\OperatorTok{\textgreater{}}\NormalTok{ parents }\OperatorTok{=}\NormalTok{ probNet}\OperatorTok{.}\FunctionTok{getParents}\OperatorTok{(}\NormalTok{node}\OperatorTok{);}
\BuiltInTok{ArrayList}\OperatorTok{\textless{}}\NormalTok{TablePotential}\OperatorTok{\textgreater{}}\NormalTok{ parentsPotential }\OperatorTok{=} \KeywordTok{new} \BuiltInTok{ArrayList}\OperatorTok{\textless{}\textgreater{}();}
\ControlFlowTok{for} \OperatorTok{(}\BuiltInTok{Node}\NormalTok{ node }\OperatorTok{:}\NormalTok{ parents}\OperatorTok{)} \OperatorTok{\{}
    \ControlFlowTok{try} \OperatorTok{\{}
\NormalTok{        parentsPotential}\OperatorTok{.}\FunctionTok{add}\OperatorTok{(}\NormalTok{node}\OperatorTok{.}\FunctionTok{getPotential}\OperatorTok{().}\FunctionTok{tableProject}\OperatorTok{(}\KeywordTok{null}\OperatorTok{,} \KeywordTok{null}\OperatorTok{));}
    \OperatorTok{\}} \ControlFlowTok{catch} \OperatorTok{(}\NormalTok{NonProjectablePotentialException e}\OperatorTok{)} \OperatorTok{\{}
        \ControlFlowTok{throw} \KeywordTok{new} \FunctionTok{UnreachableException}\OperatorTok{(}\NormalTok{e}\OperatorTok{);}
    \OperatorTok{\}}
\OperatorTok{\}}

\NormalTok{Potential potential }\OperatorTok{=}\NormalTok{ BasicOperations}\OperatorTok{.}\FunctionTok{absorbParentPotentials}\OperatorTok{(}\NormalTok{nodeVariable}\OperatorTok{,}\NormalTok{node}\OperatorTok{.}\FunctionTok{getPotential}\OperatorTok{(),}\NormalTok{parentsPotential}\OperatorTok{,}\KeywordTok{null}\OperatorTok{);}
\end{Highlighting}
\end{Shaded}

core/src/main/java/org/openmarkov/core/model/network/potential/UnivariateDistrPotential.java:357-359

\begin{Shaded}
\begin{Highlighting}[]
\KeywordTok{public} \AttributeTok{@NotNull}\NormalTok{ TablePotential }\FunctionTok{tableProject}\OperatorTok{(}\NormalTok{EvidenceCase evidenceCase}\OperatorTok{,}\NormalTok{ InferenceOptions inferenceOptions}\OperatorTok{)} \KeywordTok{throws}\NormalTok{ NonProjectablePotentialException}\OperatorTok{.}\FunctionTok{PotentialCannotBeConvertedToATable} \OperatorTok{\{}
    \ControlFlowTok{throw} \KeywordTok{new}\NormalTok{ NonProjectablePotentialException}\OperatorTok{.}\FunctionTok{PotentialCannotBeConvertedToATable}\OperatorTok{(}\KeywordTok{this}\OperatorTok{);}
\OperatorTok{\}}
\end{Highlighting}
\end{Shaded}
\paragraph{Cómo arreglarlo}
Hay que decidir qué debe pasar con un padre cuyo potencial no es una tabla:

\begin{itemize}
\tightlist
\item
  rechazarlo: añadir a la condición del menú que el potencial de cada padre se pueda convertir en tabla, y en \texttt{generate\allowbreak{}Edits} convertir el fallo en una excepción del dominio con un mensaje, en vez de \texttt{Unreachable\allowbreak{}Exception};
\item
  aceptarlo cuando tenga sentido: una distribución exacta sobre un número, la que escribe el formato 1.0, dice lo mismo que un \texttt{Exact\allowbreak{}Distr\allowbreak{}Potential}, y se podría tabular antes de absorber. Es una de las opciones de arreglo de \hyperref[err:112]{error~112}; si se toma allí, los trece nodos de suma y producto medidos dejan de fallar por esta causa.
\end{itemize}

La condición tiene que ser la misma en todos los sitios que deciden si se absorbe: \texttt{Absorb\allowbreak{}Parents\allowbreak{}Validator.\allowbreak{}validate} / \texttt{Basic\allowbreak{}Operations.\allowbreak{}have\allowbreak{}Parents\allowbreak{}And\allowbreak{}Are\allowbreak{}All\allowbreak{}Absorbable} (menú), \texttt{Basic\allowbreak{}Operations.\allowbreak{}absorb\allowbreak{}All\allowbreak{}Intermediate\allowbreak{}Numeric\allowbreak{}Nodes} (preproceso) y los dos lugares que convierten los padres, \texttt{Absorb\allowbreak{}Parents\allowbreak{}Edit.\allowbreak{}generate\allowbreak{}Edits} y \texttt{Basic\allowbreak{}Operations.\allowbreak{}absorb\allowbreak{}Parents}. \hyperref[err:8]{Error~8} añade a esa misma condición la clase del potencial del nodo.

Prueba de regresión: sobre \texttt{0\_\allowbreak{}miscellaneous/\allowbreak{}networks/\allowbreak{}id/\allowbreak{}1-0/\allowbreak{}ID-arthronet-1-0.\allowbreak{}pgmx}, el menú «Absorb parents» de «EVAC Total» debe estar apagado o la edición debe terminar sin \texttt{Unreachable\allowbreak{}Exception}; y el recorrido de todas las redes de \texttt{0\_\allowbreak{}miscellaneous/\allowbreak{}networks}, incluidas las carpetas \texttt{1-0}, no debe dar \texttt{Unreachable\allowbreak{}Exception} al absorber.

\setcounter{subsection}{9}
\subsection[Descontar una utilidad exponencial cambia el coeficiente equivocado]{\textcolor{corregidoverde}{[Corregido]} Escalar o descontar una utilidad exponencial cambia el coeficiente de otra covariable cuando el término constante no va el primero}\label{err:10}
\begin{ficha}
\fichakey{Estado}Corregido en el commit \texttt{27c5226} (23-09-2026): escalar o descontar un potencial exponencial suma \texttt{ln(\allowbreak{}escala)} al coeficiente de \texttt{Constant}, esté donde esté, también en los coeficientes muestreados; sin \texttt{Constant} lanza \texttt{Invalid\allowbreak{}Argument\allowbreak{}Exception} con mensaje en vez de indexar con −1.
\fichakey{Módulo}core
\fichakey{Paquete}org.openmarkov.core.model.network.potential
\fichakey{Ficheros}\path{core/src/main/java/org/openmarkov/core/model/network/potential/ExponentialPotential.java:133-140} \newline \path{core/src/main/java/org/openmarkov/core/model/network/potential/GLMPotential.java:308-318} \newline \path{core/src/main/java/org/openmarkov/core/model/network/TemporalNetOperations.java:349} \newline \path{core/src/main/java/org/openmarkov/core/model/network/TemporalNetOperations.java:363-391} \newline \path{core/src/main/java/org/openmarkov/core/model/network/UtilityOperations.java:28-86}
\fichakey{Comprobado}Leyendo el código y ejecutándolo
\fichakey{Esfuerzo}Pequeño (horas)
\fichakey{Decisión de equipo}No
\fichakey{Relacionados}\hyperref[err:11]{error~11}, \hyperref[err:14]{error~14}, \hyperref[err:15]{error~15}, \hyperref[err:71]{error~71}, \hyperref[nota:110]{nota~110}
\end{ficha}
\paragraph{Qué pasa}
En un modelo con una utilidad exponencial cuyo término constante no va el primero de la lista de covariables, el descuento o la escala de las utilidades dan un resultado equivocado sin aviso: en vez de multiplicar la utilidad por el factor, cambian el coeficiente de otra covariable.

Un potencial exponencial (\texttt{Exponential\allowbreak{}Potential}, de la familia de regresiones \texttt{GLMPotential}) vale \texttt{exp(\allowbreak{}β\_\allowbreak{}const\ +\ Σ\ β\_\allowbreak{}j\ ·\ covariable\_\allowbreak{}j)}. El término constante es la covariable marcada con el texto \texttt{Constant}, que puede estar en cualquier posición de la lista.

Nada obliga a que vaya la primera: el lector de \texttt{.\allowbreak{}pgmx} conserva el orden del fichero, y el panel de edición del potencial (\texttt{GLMPanel}, que hereda de \texttt{Key\allowbreak{}Table\allowbreak{}Panel} los botones de subir y bajar filas y que inserta las covariables nuevas detrás de la fila seleccionada) permite cambiarlo (esto último comprobado leyendo). Y hay redes del repositorio con \texttt{Constant} en último lugar: los potenciales exponenciales de los nodos de utilidad \texttt{Inpatient\ costs} y \texttt{Outpatient\ costs} de \texttt{0\_\allowbreak{}miscellaneous/\allowbreak{}networks/\allowbreak{}mid/\allowbreak{}MID-CHAP-Ryan-Griffin.\allowbreak{}pgmx} (y de su copia \texttt{0\_\allowbreak{}miscellaneous/\allowbreak{}networks/\allowbreak{}mid/\allowbreak{}1-0/\allowbreak{}MID-CHAP-Ryan-Griffin-1-0.\allowbreak{}pgmx}, y de \texttt{integration\allowbreak{}Tests/\allowbreak{}src/\allowbreak{}main/\allowbreak{}resources/\allowbreak{}networks/\allowbreak{}mid/\allowbreak{}chap.\allowbreak{}pgmx} e \texttt{integration\allowbreak{}Tests/\allowbreak{}src/\allowbreak{}main/\allowbreak{}resources/\allowbreak{}networks/\allowbreak{}mid/\allowbreak{}chap-sv.\allowbreak{}pgmx}) tienen las covariables \texttt{\{Age\ {[}0{]}\}*12,\allowbreak{}\ CD4\%\ 8–15\ {[}0{]},\allowbreak{}\ CD4\%\ \textgreater{}\ 15\ {[}0{]},\allowbreak{}\ Therapy,\allowbreak{}\ Constant}.

Camino de llegada: \texttt{scale\allowbreak{}Potential} se ejecuta en tres sitios de la preparación de toda evaluación de un modelo con utilidades:

\begin{itemize}
\tightlist
\item
  el descuento de los nodos de utilidad temporales (\texttt{Temporal\allowbreak{}Net\allowbreak{}Operations.\allowbreak{}apply\allowbreak{}Discount\allowbreak{}To\allowbreak{}Utility\allowbreak{}Nodes}, factor \texttt{1/\allowbreak{}(\allowbreak{}1+d)\^{}t}, que se aplica con cualquier tasa de descuento distinta de cero);
\item
  el paso a un único criterio (\texttt{Utility\allowbreak{}Operations.\allowbreak{}transform\allowbreak{}To\allowbreak{}Unicriterion});
\item
  la escala de coste-efectividad (\texttt{Utility\allowbreak{}Operations.\allowbreak{}apply\allowbreak{}CEUtility\allowbreak{}Scaling}).
\end{itemize}

Un árbol (\texttt{Tree\allowbreak{}ADDPotential.\allowbreak{}scale\allowbreak{}Potential}) reenvía el escalado a cada rama, así que alcanza también a los exponenciales colgados de un árbol, como los de CHAP.

Por qué pasa: la clase busca la posición de \texttt{Constant} con \texttt{GLMPotential.\allowbreak{}get\allowbreak{}Constant\allowbreak{}Index} y la usa al proyectar (\texttt{table\allowbreak{}Project}, línea 84). Multiplicar el potencial por un factor \texttt{s} equivale a sumar \texttt{ln(\allowbreak{}s)} al coeficiente del término constante, que es lo que dice el comentario de \texttt{scale\allowbreak{}Potential}. Pero el código suma \texttt{ln(\allowbreak{}s)} a \texttt{coefficients{[}0{]}}, que solo es el del término constante si \texttt{Constant} va la primera.

Medido:

\begin{itemize}
\tightlist
\item
  Potencial \texttt{exp(\allowbreak{}0,\allowbreak{}5\ +\ 1·X)} con \texttt{X\ =\ 3}: con \texttt{Constant} primero, escalar por 2 da 66,23 (el doble de 33,12, correcto); con las mismas covariables en orden \texttt{X,\allowbreak{}\ Constant}, da 264,92, porque el coeficiente de \texttt{X} pasa de 1 a 1,693.
\item
  Potencial real de \texttt{Inpatient\ costs\ {[}0{]}} de \texttt{MID-CHAP-Ryan-Griffin.\allowbreak{}pgmx}, escalado por el descuento de 5 ciclos al 3,5 \% (0,842, logaritmo −0,172): los coeficientes pasan de \texttt{{[}-0.\allowbreak{}0016163,\allowbreak{}\ -0.\allowbreak{}9416483,\allowbreak{}\ -1.\allowbreak{}226022,\allowbreak{}\ -0.\allowbreak{}8080182,\allowbreak{}\ 4.\allowbreak{}016486{]}} a \texttt{{[}-0.\allowbreak{}1736234,\allowbreak{}\ -0.\allowbreak{}9416483,\allowbreak{}\ -1.\allowbreak{}226022,\allowbreak{}\ -0.\allowbreak{}8080182,\allowbreak{}\ 4.\allowbreak{}016486{]}}. El que cambia es el de la edad en meses, no el término constante: con una edad de 40 años el coste hospitalario de ese ciclo queda multiplicado por \texttt{exp(\allowbreak{}-0,\allowbreak{}172·480)}, prácticamente cero, en vez de por 0,842.
\end{itemize}

Con la red tal como está no se puede medir el efecto en el resultado final del análisis de coste-efectividad de CHAP: \texttt{VECEAnalysis} se detiene antes por otro defecto, \hyperref[err:71]{error~71}. Reescribiendo a mano las fórmulas que ese defecto deja mal, sí se puede. Con los descuentos del fichero (0 \%) el resultado no cambia, porque el factor es 1. Con un descuento del 3,5 \% anual en los dos criterios, el coste esperado de cotrimoxazol sale 18,52; poniendo antes \texttt{Constant} la primera en cada exponencial, que equivale al arreglo, sale 162,4.
\paragraph{El código}
core/src/main/java/org/openmarkov/core/model/network/potential/ExponentialPotential.java:133-140

\begin{Shaded}
\begin{Highlighting}[]
\AttributeTok{@Override} \KeywordTok{public} \DataTypeTok{void} \FunctionTok{scalePotential}\OperatorTok{(}\DataTypeTok{double}\NormalTok{ scale}\OperatorTok{)} \OperatorTok{\{}
    \CommentTok{/*}
\CommentTok{     * Add ln(scale) to the first coefficient (constant covariate) is the same as}
\CommentTok{     * multiply all the exponential potential by the scale}
\CommentTok{     */}
\NormalTok{    coefficients}\OperatorTok{[}\DecValTok{0}\OperatorTok{]} \OperatorTok{+=} \BuiltInTok{Math}\OperatorTok{.}\FunctionTok{log}\OperatorTok{(}\NormalTok{scale}\OperatorTok{);}
    
\OperatorTok{\}}
\end{Highlighting}
\end{Shaded}

core/src/main/java/org/openmarkov/core/model/network/potential/ExponentialPotential.java:84 y 107

\begin{Shaded}
\begin{Highlighting}[]
\DataTypeTok{int}\NormalTok{ constantIndex }\OperatorTok{=} \FunctionTok{getConstantIndex}\OperatorTok{(}\NormalTok{covariates}\OperatorTok{);}
\KeywordTok{...}
\DataTypeTok{double}\NormalTok{ regression }\OperatorTok{=}\NormalTok{ coefficients}\OperatorTok{[}\NormalTok{constantIndex}\OperatorTok{];}
\end{Highlighting}
\end{Shaded}
\paragraph{Cómo arreglarlo}
Sumar \texttt{Math.\allowbreak{}log(\allowbreak{}scale)} a \texttt{coefficients{[}get\allowbreak{}Constant\allowbreak{}Index(\allowbreak{}covariates){]}}. Si \texttt{get\allowbreak{}Constant\allowbreak{}Index} devuelve −1 (el potencial no tiene \texttt{Constant}), fallar con un mensaje que lo diga en vez de indexar con −1; es el mismo problema de fondo que \hyperref[err:14]{error~14}.

Conviene revisar dos detalles del mismo método al tocarlo:

\begin{itemize}
\tightlist
\item
  \texttt{scale\allowbreak{}Potential(\allowbreak{}0)} (lo usa \texttt{Temporal\allowbreak{}Net\allowbreak{}Operations}, línea 349, para anular las utilidades del ciclo que se poda según el momento de transición) suma \texttt{ln(\allowbreak{}0)\ =\ −∞}. Con \texttt{Constant} primero da \texttt{exp(\allowbreak{}−∞)\ =\ 0}, correcto. Con \texttt{Constant} en otra posición, el −∞ cae en el coeficiente de otra covariable; medido con el potencial \texttt{X,\allowbreak{}\ Constant}, la utilidad que debía anularse vale 0 si \texttt{X\ =\ 3}, \textbf{NaN} si \texttt{X\ =\ 0} e \textbf{Infinity} si \texttt{X\ =\ −3}. El arreglo de arriba lo corrige también; aun así, conviene tratar el factor 0 aparte (devolver cero) para no dejar un coeficiente infinito en la copia.
\item
  Si el potencial ya tiene coeficientes muestreados (\texttt{sampled\allowbreak{}Coefficients}, que es lo que usa \texttt{table\allowbreak{}Project} durante un análisis probabilístico de sensibilidad), escalar solo \texttt{coefficients} no cambia lo que se proyecta. En ese análisis se muestrea antes de descontar, así que el descuento se perdería también con \texttt{Constant} en su sitio; hoy ningún análisis llega a ese punto (\hyperref[nota:110]{nota~110}), pero conviene que el arreglo escale también \texttt{sampled\allowbreak{}Coefficients}.
\end{itemize}

Los demás potenciales de regresión no tienen este defecto: \texttt{Linear\allowbreak{}Combination\allowbreak{}Potential.\allowbreak{}scale\allowbreak{}Potential} multiplica todos los coeficientes y \texttt{Weibull\allowbreak{}Hazard\allowbreak{}Potential.\allowbreak{}scale\allowbreak{}Potential} no admite escalado.

Prueba de regresión: en \texttt{core/\allowbreak{}src/\allowbreak{}test}, un \texttt{Exponential\allowbreak{}Potential} con covariables \texttt{\{X\},\allowbreak{}\ Constant} y otro con \texttt{Constant,\allowbreak{}\ \{X\}} (mismos coeficientes emparejados), escalar ambos por 2 y comprobar que la proyección con \texttt{X\ =\ 3} es el doble de la original en los dos casos.

\setcounter{subsection}{10}
\subsection[Escalar o descontar una fórmula o un producto de utilidades da otro número]{\textcolor{corregidoverde}{[Corregido]} Con una escala de criterio o un descuento, un nodo de utilidad que combina otras utilidades con una fórmula o un producto recibe el factor varias veces, y la fórmula además sin paréntesis: el resultado sale equivocado sin aviso}\label{err:11}
\begin{ficha}
\fichakey{Estado}Corregido en el commit \texttt{4e7c473} (23-09-2026): cuando un nodo de utilidad combina otras con un producto o una fórmula, la escala o el descuento se aplica una sola vez, en ese nodo, y no en sus padres; una fórmula se escala entera, entre paréntesis. Una suma sigue dejando el factor a sus padres.
\fichakey{Módulo}core
\fichakey{Paquete}org.openmarkov.core.model.network; org.openmarkov.core.model.network.potential
\fichakey{Ficheros}\path{core/src/main/java/org/openmarkov/core/model/network/potential/FunctionPotential.java:167-170} \newline \path{core/src/main/java/org/openmarkov/core/model/network/potential/FunctionPotentialOld.java:167-173} \newline \path{core/src/main/java/org/openmarkov/core/model/network/potential/ProductPotential.java:117-119} \newline \path{core/src/main/java/org/openmarkov/core/model/network/potential/SumPotential.java:146-148} \newline \path{core/src/main/java/org/openmarkov/core/model/network/UtilityOperations.java:28-60} \newline \path{core/src/main/java/org/openmarkov/core/model/network/UtilityOperations.java:67-88} \newline \path{core/src/main/java/org/openmarkov/core/model/network/TemporalNetOperations.java:346-353} \newline \path{core/src/main/java/org/openmarkov/core/model/network/TemporalNetOperations.java:363-392} \newline \path{inference/src/main/java/org/openmarkov/inference/algorithm/variableElimination/tasks/VariableElimination.java:66-116}
\fichakey{Comprobado}Leyendo el código y ejecutándolo
\fichakey{Esfuerzo}Medio (días)
\fichakey{Decisión de equipo}\textcolor{decisionrojo}{\textbf{Sí.}} A qué nodos se aplican la escala y el descuento cuando un nodo de utilidad combina otras utilidades (solo al nodo final que las combina, o solo a los padres y nunca al nodo que los combina), porque con productos y fórmulas que no son sumas las dos opciones dan resultados distintos
\fichakey{Relacionados}\hyperref[err:7]{error~7}, \hyperref[err:10]{error~10}, \hyperref[err:12]{error~12}, \hyperref[err:26]{error~26}, \hyperref[err:176]{error~176}
\end{ficha}
\paragraph{Qué pasa}
Un nodo de utilidad puede definirse combinando otros nodos de utilidad: como suma, como producto o con una fórmula (relación «Function», que el diálogo del nodo ofrece para cualquier nodo de utilidad). Antes de calcular, OpenMarkov multiplica cada utilidad por un factor en tres casos:

\begin{itemize}
\tightlist
\item
  la escala de un solo criterio (\texttt{Utility\allowbreak{}Operations.\allowbreak{}transform\allowbreak{}To\allowbreak{}Unicriterion});
\item
  la escala de coste-efectividad (\texttt{Utility\allowbreak{}Operations.\allowbreak{}apply\allowbreak{}CEUtility\allowbreak{}Scaling});
\item
  el descuento de los nodos temporales, \texttt{1/\allowbreak{}(\allowbreak{}1+d)\^{}t} (\texttt{Temporal\allowbreak{}Net\allowbreak{}Operations.\allowbreak{}apply\allowbreak{}Discount\allowbreak{}To\allowbreak{}Utility\allowbreak{}Nodes}), con cualquier tasa distinta de cero.
\end{itemize}

Las dos escalas se fijan en el diálogo de opciones de inferencia, que en una red con más de un criterio aparece al entrar en modo inferencia, al pedir el análisis de coste-efectividad, el de sensibilidad o el árbol de decisión, y desde el menú de opciones de inferencia.

Los tres métodos recorren \textbf{todos} los nodos de utilidad y escalan el potencial de cada uno, sin mirar si combina otros. Después, \texttt{Basic\allowbreak{}Operations.\allowbreak{}absorb\allowbreak{}All\allowbreak{}Intermediate\allowbreak{}Numeric\allowbreak{}Nodes} combina los padres ya escalados. Para una suma sale bien, porque \texttt{Sum\allowbreak{}Potential.\allowbreak{}scale\allowbreak{}Potential} no hace nada: la suma de padres escalados queda escalada una vez. Para los otros dos, no:

\begin{itemize}
\tightlist
\item
  \textbf{Fórmula.} \texttt{Function\allowbreak{}Potential.\allowbreak{}scale\allowbreak{}Potential} se escala a sí misma además de a sus padres: una fórmula lineal queda escalada dos veces y un producto de dos padres, tres. Y lo hace anteponiendo \texttt{factor*} al texto de la fórmula sin paréntesis, así que en \texttt{\{U1\}+100} el factor solo multiplica al primer término.
\item
  \textbf{Producto.} \texttt{Product\allowbreak{}Potential.\allowbreak{}scale\allowbreak{}Potential} no hace nada, pero sus dos padres sí se escalan: el producto queda escalado dos veces.
\end{itemize}

Lo mismo pasa con la anulación de las utilidades del ciclo que se poda según el momento de la transición (\texttt{Temporal\allowbreak{}Net\allowbreak{}Operations}, línea 349), que escala por 0: una fórmula \texttt{\{U1\}+100} quedaría en \texttt{0.\allowbreak{}0*\{U1\}+100} y no se anularía (por lectura).

El nodo producto está en una red del repositorio: \texttt{0\_\allowbreak{}miscellaneous/\allowbreak{}networks/\allowbreak{}mid/\allowbreak{}MID-HPV.\allowbreak{}pgmx}, donde «QoL {[}0{]}» es el producto de «AR QoL {[}0{]}» y «HSR QoL {[}0{]}», dos utilidades temporales de efectividad; basta con poner un descuento o una escala a la efectividad. Los nodos de fórmula sobre utilidades solo aparecen en redes de prueba (\texttt{integration\allowbreak{}Tests/\allowbreak{}src/\allowbreak{}main/\allowbreak{}resources/\allowbreak{}networks/\allowbreak{}dan/\allowbreak{}}), así que se llega construyendo el modelo en el editor. La fórmula tiene que nombrar a sus padres \texttt{U1}, \texttt{U2}\ldots{} (\hyperref[err:7]{error~7}).

Medido:

\begin{itemize}
\tightlist
\item
  En \texttt{0\_\allowbreak{}miscellaneous/\allowbreak{}networks/\allowbreak{}id/\allowbreak{}ID-CEA-minimal.\allowbreak{}pgmx}, cambiando el coste por un nodo \texttt{U1} (0 sin terapia, 14 000 con terapia) y añadiendo el nodo «Total cost» del mismo criterio con distintas fórmulas sobre \texttt{U1}; coste de la opción «yes» con ámbito de decisión, según la escala de coste-efectividad del coste (con la escala de un solo criterio sale igual):
\end{itemize}

\begin{Verbatim}[breaklines,breakanywhere,fontsize=\small]
fórmula                 escala 1   escala 2: debería / sale   escala 0,5: debería / sale
{U1}                    14000      28000 / 56000              7000 / 3500
{U1}*2                  28000      56000 / 112000             14000 / 7000
{U1}+100                14100      28200 / 56100              7050 / 3600
abs({U1})+abs({U1})     28000      56000 / 84000              14000 / 10500
\end{Verbatim}

\begin{itemize}
\tightlist
\item
  \texttt{integration\allowbreak{}Tests/\allowbreak{}src/\allowbreak{}main/\allowbreak{}resources/\allowbreak{}networks/\allowbreak{}dan/\allowbreak{}DAN-two-utility-and-chance-function-sv.\allowbreak{}pgmx} (\texttt{U0\ =\ abs(\allowbreak{}\{U1\})*\{U2\}}), evaluación con la escala de un solo criterio: −6 con escala 1; −48 con escala 2, en vez de −12; −6·10⁻⁹ con escala 0,001, en vez de −0,006.
\item
  \texttt{0\_\allowbreak{}miscellaneous/\allowbreak{}networks/\allowbreak{}mid/\allowbreak{}MID-HPV.\allowbreak{}pgmx}, análisis de coste-efectividad global: con escala 2 en la efectividad, las efectividades pasan de 60,16 a 240,63 (×4); en \texttt{0\_\allowbreak{}miscellaneous/\allowbreak{}networks/\allowbreak{}mid/\allowbreak{}MID-HPV-without-supervalue.\allowbreak{}pgmx}, que no tiene el nodo producto, pasan de 71,82 a 143,64 (×2). Con un descuento del 1,5 \% anual en la efectividad, en \texttt{MID-HPV.\allowbreak{}pgmx} queda el 42 \% (de 60,16 a 25,39) y en la versión sin producto, el 60 \% (de 71,82 a 43,40). Las dos redes no dan la misma efectividad ni sin escala; aquí solo se compara cuánto cambia cada una.
\end{itemize}
\paragraph{El código}
core/src/main/java/org/openmarkov/core/model/network/potential/FunctionPotential.java:167-170

\begin{Shaded}
\begin{Highlighting}[]
\AttributeTok{@Override} \KeywordTok{public} \DataTypeTok{void} \FunctionTok{scalePotential}\OperatorTok{(}\DataTypeTok{double}\NormalTok{ scale}\OperatorTok{)} \OperatorTok{\{}
    \BuiltInTok{String}\NormalTok{ scaleString }\OperatorTok{=} \BuiltInTok{String}\OperatorTok{.}\FunctionTok{valueOf}\OperatorTok{(}\NormalTok{scale}\OperatorTok{);}
\NormalTok{    covariates}\OperatorTok{[}\DecValTok{0}\OperatorTok{]} \OperatorTok{=} \KeywordTok{new} \FunctionTok{VariableExpression}\OperatorTok{(}\KeywordTok{this}\OperatorTok{.}\FunctionTok{variables}\OperatorTok{,}\NormalTok{ scaleString }\OperatorTok{+} \StringTok{"*"} \OperatorTok{+} \KeywordTok{this}\OperatorTok{.}\FunctionTok{covariates}\OperatorTok{[}\DecValTok{0}\OperatorTok{].}\FunctionTok{asStringExpression}\OperatorTok{());}
\OperatorTok{\}}
\end{Highlighting}
\end{Shaded}

core/src/main/java/org/openmarkov/core/model/network/potential/ProductPotential.java:117-119 (igual en \texttt{Sum\allowbreak{}Potential.\allowbreak{}java}, líneas 146-148)

\begin{Shaded}
\begin{Highlighting}[]
\AttributeTok{@Override} \KeywordTok{public} \DataTypeTok{void} \FunctionTok{scalePotential}\OperatorTok{(}\DataTypeTok{double}\NormalTok{ scale}\OperatorTok{)} \OperatorTok{\{}

\OperatorTok{\}}
\end{Highlighting}
\end{Shaded}

core/src/main/java/org/openmarkov/core/model/network/UtilityOperations.java:67-80

\begin{Shaded}
\begin{Highlighting}[]
\KeywordTok{public} \DataTypeTok{static} \DataTypeTok{void} \FunctionTok{applyCEUtilityScaling}\OperatorTok{(}\NormalTok{ProbNet probNet}\OperatorTok{)} \OperatorTok{\{}
    \ControlFlowTok{for} \OperatorTok{(}\BuiltInTok{Node}\NormalTok{ utilityNode }\OperatorTok{:}\NormalTok{ probNet}\OperatorTok{.}\FunctionTok{getNodes}\OperatorTok{(}\NormalTok{NodeType}\OperatorTok{.}\FunctionTok{UTILITY}\OperatorTok{))} \OperatorTok{\{}
        \ControlFlowTok{if} \OperatorTok{(}\NormalTok{utilityNode}\OperatorTok{.}\FunctionTok{getVariable}\OperatorTok{().}\FunctionTok{getDecisionCriterion}\OperatorTok{()} \OperatorTok{!=} \KeywordTok{null}\OperatorTok{)} \OperatorTok{\{}
            \CommentTok{// Save the actual criterion scale}
            \DataTypeTok{double}\NormalTok{ scale }\OperatorTok{=}\NormalTok{ utilityNode}\OperatorTok{.}\FunctionTok{getVariable}\OperatorTok{().}\FunctionTok{getDecisionCriterion}\OperatorTok{().}\FunctionTok{getCeScale}\OperatorTok{();}
            \ControlFlowTok{if} \OperatorTok{(}\NormalTok{scale }\OperatorTok{!=} \DecValTok{0}\OperatorTok{)} \OperatorTok{\{}
                \CommentTok{// Transform the potential with the scale}
\NormalTok{                Potential potential }\OperatorTok{=}\NormalTok{ utilityNode}\OperatorTok{.}\FunctionTok{getPotentials}\OperatorTok{().}\FunctionTok{get}\OperatorTok{(}\DecValTok{0}\OperatorTok{).}\FunctionTok{deepCopy}\OperatorTok{(}\NormalTok{probNet}\OperatorTok{);}
                \ControlFlowTok{try} \OperatorTok{\{}
\NormalTok{                    potential}\OperatorTok{.}\FunctionTok{scalePotential}\OperatorTok{(}\NormalTok{scale}\OperatorTok{);}
\end{Highlighting}
\end{Shaded}
\paragraph{Cómo arreglarlo}
Hay dos cosas que arreglar.

La primera, local: \texttt{Function\allowbreak{}Potential.\allowbreak{}scale\allowbreak{}Potential} (y su gemela \texttt{Function\allowbreak{}Potential\allowbreak{}Old.\allowbreak{}scale\allowbreak{}Potential}) debe escribir \texttt{factor*(\allowbreak{}fórmula)}, con el número sin notación científica, que es el arreglo de \hyperref[err:12]{error~12} para esas mismas líneas.

La segunda exige decidir dónde se aplica el factor cuando un nodo combina utilidades, y es la que decide los números:

\begin{itemize}
\tightlist
\item
  \textbf{Solo en el nodo final.} Escalar y descontar solo los nodos de utilidad que no tienen hijos de utilidad, cada uno con el factor de su criterio, y no sus padres. Es correcto para sumas, productos y cualquier fórmula, pero cambia qué criterio se aplica cuando los padres tienen criterios distintos (por ejemplo, una fórmula \texttt{30000*\{U1\}-\{U2\}} que junta efectividad y coste), y el orden de \texttt{Variable\allowbreak{}Elimination.\allowbreak{}general\allowbreak{}Preprocessing}, que descuenta antes de absorber.
\item
  \textbf{Solo en los padres.} Seguir escalando los padres y que ningún nodo que combina utilidades se escale a sí mismo, como ya hace la suma. Respeta el criterio de cada padre, pero un producto o una fórmula no lineal sigue recibiendo el factor más de una vez.
\end{itemize}

Sea cual sea, la regla tiene que cumplirse a la vez en \texttt{Utility\allowbreak{}Operations.\allowbreak{}transform\allowbreak{}To\allowbreak{}Unicriterion}, \texttt{Utility\allowbreak{}Operations.\allowbreak{}apply\allowbreak{}CEUtility\allowbreak{}Scaling}, \texttt{Temporal\allowbreak{}Net\allowbreak{}Operations.\allowbreak{}apply\allowbreak{}Discount\allowbreak{}To\allowbreak{}Utility\allowbreak{}Nodes} y la anulación por factor 0 de \texttt{Temporal\allowbreak{}Net\allowbreak{}Operations} (línea 349), y en \texttt{scale\allowbreak{}Potential} de \texttt{Sum\allowbreak{}Potential}, \texttt{Product\allowbreak{}Potential} y \texttt{Function\allowbreak{}Potential}.

Prueba de regresión: con \texttt{0\_\allowbreak{}miscellaneous/\allowbreak{}networks/\allowbreak{}mid/\allowbreak{}MID-HPV.\allowbreak{}pgmx}, una escala 2 en la efectividad debe duplicar exactamente las efectividades del análisis de coste-efectividad; y con un diagrama de influencia de dos criterios donde «Total cost» sea \texttt{\{U1\}+100} sobre un \texttt{U1} de 14 000, la escala 2 del coste debe dar 28 200.

\setcounter{subsection}{11}
\subsection[Una regresión no se calcula si un padre vale menos de una milésima]{\textcolor{corregidoverde}{[Corregido]} Una regresión deja de poder calcularse cuando el valor de un padre está por debajo de una milésima}\label{err:12}
\begin{ficha}
\fichakey{Estado}Corregido en el commit \texttt{f44ac1a} (23-09-2026): los números llegan al evaluador de expresiones sin notación científica, así que una regresión o una fórmula se calcula aunque un padre, un hallazgo o un factor de escala valga menos de 0,001.
\fichakey{Módulo}core
\fichakey{Paquete}org.openmarkov.core.model.network.potential
\fichakey{Ficheros}\path{core/src/main/java/org/openmarkov/core/model/network/potential/GLMPotential.java:227-241} \newline \path{core/src/main/java/org/openmarkov/core/model/network/potential/LinearCombinationPotential.java:104-111} \newline \path{core/src/main/java/org/openmarkov/core/model/network/potential/ExponentialPotential.java:98-105} \newline \path{core/src/main/java/org/openmarkov/core/model/network/potential/WeibullHazardPotential.java:185-192} \newline \path{core/src/main/java/org/openmarkov/core/expression/ReferencedExpression.java:106-113}
\fichakey{Comprobado}Leyendo el código y ejecutándolo
\fichakey{Esfuerzo}Pequeño (horas)
\fichakey{Decisión de equipo}No
\fichakey{Relacionados}\hyperref[err:7]{error~7}, \hyperref[err:11]{error~11}, \hyperref[err:13]{error~13}, \hyperref[err:15]{error~15}, \hyperref[err:16]{error~16}, \hyperref[err:176]{error~176}, \hyperref[nota:110]{nota~110}
\end{ficha}
\paragraph{Qué pasa}
Un nodo con un potencial de regresión deja de poder calcularse en cuanto algún padre vale, por hallazgo o por nombre de estado, menos de 0,001 en valor absoluto sin ser cero. La proyección del potencial falla con una excepción \texttt{Cannot\allowbreak{}Evaluate} que habla de la fórmula y no del número. Un valor pequeño ---una probabilidad o una tasa pequeña, de lo más corriente en un modelo de coste-efectividad--- basta. Un valor negativo pequeño (\texttt{-1.\allowbreak{}0E-4}) falla igual; el cero (\texttt{0.\allowbreak{}0}) no.

Camino de llegada: cualquier inferencia (propagación, evaluación, análisis de coste-efectividad, evolución temporal) sobre una red cuyos nodos lleven uno de estos potenciales, por ejemplo las redes de \texttt{0\_\allowbreak{}miscellaneous/\allowbreak{}networks/\allowbreak{}mid/\allowbreak{}} con \texttt{Hazard\ (\allowbreak{}Exponential)}, \texttt{Hazard\ (\allowbreak{}Weibull)} o \texttt{Linear\ combination}. No se ha encontrado en las redes de ejemplo un padre con ese valor; la consecuencia se ha comprobado construyendo el potencial en un pequeño programa de prueba, no abriendo una red del repositorio.

Por qué pasa: los potenciales de regresión ---la familia \texttt{GLMPotential} (modelo lineal generalizado): combinación lineal, exponencial, riesgo de Weibull y riesgo exponencial--- calculan cada covariable (cada término de la regresión, que puede ser una fórmula como \texttt{\{Edad\}*12}) con el evaluador de expresiones jeval, una biblioteca externa que trabaja sobre texto. \texttt{Referenced\allowbreak{}Expression.\allowbreak{}processed\allowbreak{}Expression} sustituye cada referencia \texttt{\{Variable\}} por el texto que le den para esa variable, y \texttt{evaluate\allowbreak{}With} le pasa la cadena resultante a jeval.

Quien prepara ese texto lo hace con \texttt{String.\allowbreak{}value\allowbreak{}Of(\allowbreak{}double)}: \texttt{GLMPotential.\allowbreak{}table\allowbreak{}Project} para los padres con hallazgo (línea 241), y \texttt{Linear\allowbreak{}Combination\allowbreak{}Potential}, \texttt{Exponential\allowbreak{}Potential} y \texttt{Weibull\allowbreak{}Hazard\allowbreak{}Potential} para los padres sin hallazgo, cuyo valor sale del nombre del estado (líneas 111, 105 y 192). \texttt{String.\allowbreak{}value\allowbreak{}Of} de un \texttt{double} cambia a notación científica cuando el valor absoluto es menor que 0,001 (o mayor o igual que 10⁷): \texttt{String.\allowbreak{}value\allowbreak{}Of(\allowbreak{}0.\allowbreak{}0001)} es \texttt{"1.\allowbreak{}0E-4"}. Y jeval no acepta un exponente con signo: medido, \texttt{1e3*2} da 2000 y \texttt{1.\allowbreak{}0E7*2} da 2,0E7, pero \texttt{1.\allowbreak{}0E-4*2}, \texttt{1e-4*2} y también \texttt{1E+3*2} dan \texttt{Evaluation\allowbreak{}Exception:\ Expression\ is\ invalid}. El corte está exactamente donde \texttt{Double.\allowbreak{}to\allowbreak{}String} cambia de forma y no tiene relación con el modelo.

Medido con un \texttt{Linear\allowbreak{}Combination\allowbreak{}Potential} de hijo numérico \texttt{Y} y padre numérico \texttt{X}, covariable \texttt{\{X\}*2}, y el hallazgo \texttt{X\ =\ v}:

\begin{Verbatim}[breaklines,breakanywhere,fontsize=\small]
X=50.0    -> [100.0]
X=0.001   -> [0.002]
X=1.0E-4  -> NonProjectablePotentialException$CannotEvaluate
X=1.0E-5  -> NonProjectablePotentialException$CannotEvaluate
X=1.0E7   -> [2.0E7]
\end{Verbatim}

Lo mismo ocurre sin hallazgo, con un padre de estados finitos cuyos estados se llaman \texttt{0.\allowbreak{}0001} y \texttt{1}, e incluso cuando la covariable es la variable sola (\texttt{\{D\}}), porque el nombre del estado se convierte a \texttt{double} y vuelve a texto con \texttt{String.\allowbreak{}value\allowbreak{}Of}.
\paragraph{El código}
core/src/main/java/org/openmarkov/core/model/network/potential/GLMPotential.java:227-241

\begin{Shaded}
\begin{Highlighting}[]
\DataTypeTok{double}\NormalTok{ numericValue}\OperatorTok{;}
\NormalTok{Finding finding }\OperatorTok{=}\NormalTok{ evidenceCase}\OperatorTok{.}\FunctionTok{getFinding}\OperatorTok{(}\NormalTok{variable}\OperatorTok{);}
\ControlFlowTok{if} \OperatorTok{(}\NormalTok{variable}\OperatorTok{.}\FunctionTok{getVariableType}\OperatorTok{()} \OperatorTok{==}\NormalTok{ VariableType}\OperatorTok{.}\FunctionTok{NUMERIC}
        \OperatorTok{||}\NormalTok{ variable}\OperatorTok{.}\FunctionTok{getVariableType}\OperatorTok{()} \OperatorTok{==}\NormalTok{ VariableType}\OperatorTok{.}\FunctionTok{DISCRETIZED}\OperatorTok{)} \OperatorTok{\{}
\NormalTok{    numericValue }\OperatorTok{=}\NormalTok{ finding}\OperatorTok{.}\FunctionTok{getNumericalValue}\OperatorTok{();}
\OperatorTok{\}} \ControlFlowTok{else} \OperatorTok{\{}
    \KeywordTok{...}
\OperatorTok{\}}
\NormalTok{variableValues}\OperatorTok{.}\FunctionTok{put}\OperatorTok{(}\NormalTok{variable}\OperatorTok{,} \BuiltInTok{String}\OperatorTok{.}\FunctionTok{valueOf}\OperatorTok{(}\NormalTok{numericValue}\OperatorTok{));}
\end{Highlighting}
\end{Shaded}

core/src/main/java/org/openmarkov/core/expression/ReferencedExpression.java:106-113

\begin{Shaded}
\begin{Highlighting}[]
\KeywordTok{public} \BuiltInTok{String} \FunctionTok{evaluateWith}\OperatorTok{(}\BuiltInTok{Map}\OperatorTok{\textless{}}\NormalTok{T}\OperatorTok{,} \BuiltInTok{String}\OperatorTok{\textgreater{}}\NormalTok{ variablesValues}\OperatorTok{)} \KeywordTok{throws} \KeywordTok{...} \OperatorTok{\{}
    \BuiltInTok{String}\NormalTok{ processedExpression }\OperatorTok{=} \KeywordTok{this}\OperatorTok{.}\FunctionTok{processedExpression}\OperatorTok{(}\NormalTok{variablesValues}\OperatorTok{);}
    \BuiltInTok{String}\NormalTok{ result}\OperatorTok{;}
    \ControlFlowTok{try} \OperatorTok{\{}
\NormalTok{        result }\OperatorTok{=} \KeywordTok{new}\NormalTok{ net}\OperatorTok{.}\FunctionTok{sourceforge}\OperatorTok{.}\FunctionTok{jeval}\OperatorTok{.}\FunctionTok{Evaluator}\OperatorTok{().}\FunctionTok{evaluate}\OperatorTok{(}\NormalTok{processedExpression}\OperatorTok{);}
    \OperatorTok{\}} \ControlFlowTok{catch} \OperatorTok{(}\NormalTok{EvaluationException e}\OperatorTok{)} \OperatorTok{\{}
        \ControlFlowTok{throw} \KeywordTok{new}\NormalTok{ NonProjectablePotentialException}\OperatorTok{.}\FunctionTok{CannotEvaluate}\OperatorTok{(}\NormalTok{processedExpression}\OperatorTok{,}\NormalTok{ e}\OperatorTok{);}
    \OperatorTok{\}}
\end{Highlighting}
\end{Shaded}
\paragraph{Cómo arreglarlo}
Escribir el número en una forma que jeval acepte siempre, en un único sitio. La opción más contenida es una función auxiliar (por ejemplo en \texttt{Referenced\allowbreak{}Expression} o junto a \texttt{Variable\allowbreak{}Expression}) que convierta un \texttt{double} a texto sin notación científica ---\texttt{Big\allowbreak{}Decimal.\allowbreak{}value\allowbreak{}Of(\allowbreak{}v).\allowbreak{}to\allowbreak{}Plain\allowbreak{}String(\allowbreak{})} conserva todas las cifras de \texttt{Double.\allowbreak{}to\allowbreak{}String} y nunca escribe exponente--- y usarla en todos los sitios que hoy hacen \texttt{String.\allowbreak{}value\allowbreak{}Of(\allowbreak{}double)} o \texttt{""\ +\ double} antes de evaluar:

\begin{itemize}
\tightlist
\item
  \texttt{GLMPotential.\allowbreak{}table\allowbreak{}Project} (línea 241);
\item
  \texttt{Linear\allowbreak{}Combination\allowbreak{}Potential.\allowbreak{}table\allowbreak{}Project} (línea 111), \texttt{Exponential\allowbreak{}Potential.\allowbreak{}table\allowbreak{}Project} (línea 105) y \texttt{Weibull\allowbreak{}Hazard\allowbreak{}Potential.\allowbreak{}table\allowbreak{}Project} (línea 192);
\item
  \texttt{Function\allowbreak{}Potential.\allowbreak{}sample\allowbreak{}Conditioned\allowbreak{}Variable} (\texttt{core/\allowbreak{}src/\allowbreak{}main/\allowbreak{}java/\allowbreak{}org/\allowbreak{}openmarkov/\allowbreak{}core/\allowbreak{}model/\allowbreak{}network/\allowbreak{}potential/\allowbreak{}Function\allowbreak{}Potential.\allowbreak{}java:103}, simulación de eventos discretos) y su gemela \texttt{Function\allowbreak{}Potential\allowbreak{}Old} (línea 276);
\item
  \texttt{Table\allowbreak{}Potential\allowbreak{}Arithmetic.\allowbreak{}evaluate\allowbreak{}Function\allowbreak{}Potential} (\texttt{core/\allowbreak{}src/\allowbreak{}main/\allowbreak{}java/\allowbreak{}org/\allowbreak{}openmarkov/\allowbreak{}core/\allowbreak{}model/\allowbreak{}network/\allowbreak{}potential/\allowbreak{}operation/\allowbreak{}Table\allowbreak{}Potential\allowbreak{}Arithmetic.\allowbreak{}java:643}, utilidades combinadas con una fórmula);
\item
  \texttt{Function\allowbreak{}Potential.\allowbreak{}scale\allowbreak{}Potential} (\texttt{Function\allowbreak{}Potential.\allowbreak{}java:167-169}, y \texttt{Function\allowbreak{}Potential\allowbreak{}Old.\allowbreak{}java:171-172}), que antepone \texttt{String.\allowbreak{}value\allowbreak{}Of(\allowbreak{}scale)\ +\ "*"} a la fórmula (sin paréntesis, y sobre padres que ya se han escalado: eso es otro defecto, \hyperref[err:11]{error~11}): un factor de descuento o de escala menor que 0,001 (por ejemplo 1/1,035\^{}210 en un horizonte largo) deja escrita en el potencial una fórmula que jeval no acepta. Medido con \texttt{integration\allowbreak{}Tests/\allowbreak{}src/\allowbreak{}main/\allowbreak{}resources/\allowbreak{}networks/\allowbreak{}dan/\allowbreak{}DAN-two-utility-and-chance-function-sv.\allowbreak{}pgmx} (nodo \texttt{U0} con la fórmula \texttt{abs(\allowbreak{}\{U1\})*\{U2\}}) y la escala de un solo criterio a 0,0005: la evaluación termina con \texttt{Unreachable\allowbreak{}Exception} causada por \texttt{Cannot\ evaluate\ \textquotesingle{}5.\allowbreak{}0E-4*abs(\allowbreak{}-0.\allowbreak{}001)*-0.\allowbreak{}0015\textquotesingle{}}. El diálogo de opciones de inferencia redondea las escalas a tres decimales, así que un factor así solo llega desde un fichero o con el descuento de un nodo de fórmula temporal en un ciclo lejano (con un 6 \% anual y ciclos de un año, a partir del ciclo 119).
\end{itemize}

La alternativa de cambiar el lugar de la conversión ---que \texttt{Referenced\allowbreak{}Expression} reciba números y los escriba él--- evita que un llamador futuro repita el error, pero cambia la firma \texttt{Map\textless{}T,\allowbreak{}\ String\textgreater{}} que usan también la tabla aumentada (\texttt{Augmented\allowbreak{}Prob\allowbreak{}Table\allowbreak{}Potential}) y la interfaz.

Prueba de regresión: en \texttt{core/\allowbreak{}src/\allowbreak{}test}, proyectar un \texttt{Linear\allowbreak{}Combination\allowbreak{}Potential} con covariable \texttt{\{X\}*2} y hallazgos \texttt{X\ =\ 0.\allowbreak{}0001}, \texttt{X\ =\ -0.\allowbreak{}0001} y \texttt{X\ =\ 1.\allowbreak{}0E-12}, y comprobar que devuelve \texttt{2·X}; y otra con un padre de estados finitos llamados \texttt{0.\allowbreak{}0001} y \texttt{1} sin hallazgo.

\setcounter{subsection}{12}
\subsection[Una covariable con el nombre sin llaves se acepta pero no se calcula nunca]{\textcolor{corregidoverde}{[Corregido]} Una covariable escrita como fórmula con el nombre de la variable sin llaves se acepta al leer y al editar, pero no se puede calcular nunca}\label{err:13}
\begin{ficha}
\fichakey{Estado}Corregido en el commit \texttt{54a1aec} (23-09-2026): se eligió aceptar los nombres sin llaves; al construir la expresión se les ponen, así que la covariable se calcula, se guarda con llaves y lo que el diálogo acepta se puede calcular. Las redes CHAP pasan ahora de este punto y se paran en el \hyperref[err:71]{error~71}.
\fichakey{Módulo}io, core, gui
\fichakey{Paquete}org.openmarkov.io.probmodel.reader; org.openmarkov.core.expression; org.openmarkov.gui.dialog.common
\fichakey{Ficheros}\path{io/src/main/java/org/openmarkov/io/probmodel/reader/PGMXReader_0_2.java:1480-1495} \newline \path{core/src/main/java/org/openmarkov/core/expression/ReferencedExpression.java:22-43} \newline \path{io/src/main/java/org/openmarkov/io/probmodel/writer/PGMXWriter_0_2.java:1060-1068} \newline \path{gui/src/main/java/org/openmarkov/gui/dialog/common/ArithmeticExpressionDialog.java:185-204} \newline \path{gui/src/main/java/org/openmarkov/gui/dialog/common/VariableExpressionTextField.java:80-99}
\fichakey{Comprobado}Leyendo el código y ejecutándolo
\fichakey{Esfuerzo}Medio (días)
\fichakey{Decisión de equipo}\textcolor{decisionrojo}{\textbf{Sí.}} Si una fórmula puede nombrar variables sin llaves (y entonces el evaluador debe reconocerlas) o si las llaves son obligatorias (y entonces el lector y el diálogo deben rechazar o avisar)
\fichakey{Relacionados}\hyperref[err:12]{error~12}, \hyperref[err:71]{error~71}, \hyperref[nota:28]{nota~28}
\end{ficha}
\paragraph{Qué pasa}
Una covariable de regresión escrita como fórmula que nombra una variable sin llaves, por ejemplo \texttt{pow(\allowbreak{}Age\ {[}1{]},\allowbreak{}8)}, se abre desde un fichero sin aviso y el diálogo de edición la da por buena, pero cualquier inferencia que proyecte ese potencial falla con \texttt{Non\allowbreak{}Projectable\allowbreak{}Potential\allowbreak{}Exception.\allowbreak{}Cannot\allowbreak{}Evaluate}. Guardar y volver a abrir no la repara.

Camino de llegada: abrir \texttt{0\_\allowbreak{}miscellaneous/\allowbreak{}networks/\allowbreak{}mid/\allowbreak{}MID-mammography.\allowbreak{}pgmx} (o cualquier fichero con una fórmula así) y lanzar cualquier inferencia que proyecte esos potenciales; o, en cualquier red, escribir a mano en el diálogo de expresiones una covariable con el nombre de un padre sin llaves.

En \texttt{0\_\allowbreak{}miscellaneous/\allowbreak{}networks/\allowbreak{}mid/\allowbreak{}}, las covariables con fórmula de \texttt{MID-CHAP-Ryan-Griffin.\allowbreak{}pgmx} y de \texttt{1-0/\allowbreak{}MID-mammography-1-0.\allowbreak{}pgmx} llevan llaves y se calculan. Van sin llaves las de \texttt{MID-mammography.\allowbreak{}pgmx} (cuatro covariables en los potenciales de \texttt{Death\ (\allowbreak{}OC)\ {[}1{]}} e \texttt{Invasive\ cancer\ {[}0{]}}) y, entre las redes de las pruebas, las de \texttt{integration\allowbreak{}Tests/\allowbreak{}src/\allowbreak{}main/\allowbreak{}resources/\allowbreak{}networks/\allowbreak{}mid/\allowbreak{}chap.\allowbreak{}pgmx} e \texttt{integration\allowbreak{}Tests/\allowbreak{}src/\allowbreak{}main/\allowbreak{}resources/\allowbreak{}networks/\allowbreak{}mid/\allowbreak{}chap-sv.\allowbreak{}pgmx} (covariable \texttt{Age\ {[}0{]}*12} de los costes hospitalarios y ambulatorios).

Por qué pasa: en los potenciales de regresión (familia \texttt{GLMPotential}: combinación lineal, exponencial, riesgos de Weibull y exponencial) cada covariable es una expresión de la clase \texttt{Variable\allowbreak{}Expression}. Su clase madre, \texttt{Referenced\allowbreak{}Expression}, trocea el texto con la expresión regular \texttt{\textbackslash{}\{(\allowbreak{}{[}\^{}\textbackslash{}\{{]}+?)\textbackslash{}\}} y solo trata como referencia a una variable lo que va \textbf{entre llaves} y coincide con el nombre de una variable del potencial; el resto se guarda como texto sin interpretar y se le pasa tal cual al evaluador jeval. Por tanto \texttt{pow(\allowbreak{}\{Age\ {[}1{]}\},\allowbreak{}8)} se calcula, y \texttt{pow(\allowbreak{}Age\ {[}1{]},\allowbreak{}8)} llega a jeval con el nombre de la variable dentro y falla.

Hay tres sitios que no siguen esa misma regla:

\begin{itemize}
\tightlist
\item
  \textbf{El lector de \texttt{.\allowbreak{}pgmx}} (\texttt{PGMXReader\_\allowbreak{}0\_\allowbreak{}2.\allowbreak{}get\allowbreak{}Covariates}, que también usa \texttt{PGMXReader\_\allowbreak{}1\_\allowbreak{}0} a través de \texttt{PGMXPotential\allowbreak{}Parsers}) pone las llaves solo cuando el texto de la covariable es \textbf{exactamente} el nombre de una variable (\texttt{Age\ {[}0{]}} → \texttt{\{Age\ {[}0{]}\}}). Si el nombre aparece dentro de una fórmula, la covariable se construye tal cual, sin referencias, y el fichero se abre sin ningún aviso.
\item
  \textbf{El escritor} (\texttt{PGMXWriter\_\allowbreak{}0\_\allowbreak{}2.\allowbreak{}get\allowbreak{}Covariates\allowbreak{}Element}) escribe \texttt{as\allowbreak{}String\allowbreak{}Expression(\allowbreak{})}, que devuelve el texto sin llaves de lo que no se reconoció; así que guardar y volver a abrir deja la covariable igual de rota. No hay manera de repararla desde el propio programa salvo reescribirla a mano con llaves.
\item
  \textbf{El diálogo de edición de expresiones} (\texttt{Arithmetic\allowbreak{}Expression\allowbreak{}Dialog}, que se abre al pulsar «añadir» o al hacer doble clic en una covariable en el panel \texttt{GLMPanel} del potencial) valida el texto con su propio \texttt{process\allowbreak{}Expression}, que sustituye por variables de jeval \textbf{tanto} \texttt{\{Nombre\}} \textbf{como} \texttt{Nombre} a secas. Si el usuario escribe el nombre a mano, sin pulsar en la lista de variables (que sí inserta las llaves), el diálogo pinta la expresión como válida y la acepta con «Aceptar»; la inferencia falla después. \texttt{Variable\allowbreak{}Expression\allowbreak{}Text\allowbreak{}Field} repite la misma validación. Esto se ha comprobado leyendo, no ejecutando (el diálogo no se puede construir sin pantalla).
\end{itemize}

Medido con un pequeño programa que abre la red, evalúa cada covariable con fórmula dando a todas las variables el valor 50, la guarda y la vuelve a abrir:

\begin{Verbatim}[breaklines,breakanywhere,fontsize=\small]
MID-mammography.pgmx (formato 0.2)
Death (OC) [1]      | ExponentialHazardPotential | pow(Age [1],8)                  | refs=0 -> CannotEvaluate
Invasive cancer [0] | ExponentialHazardPotential | log(0.849828*(-0.0003717*Age [0]+ ...)) | refs=0 -> CannotEvaluate
(dos covariables más iguales)                                                       -> CannotEvaluate
--- tras guardar y volver a abrir: las cuatro siguen con refs=0 y CannotEvaluate
1-0/MID-mammography-1-0.pgmx (formato 1.0)
Death (OC) [1]      | ExponentialHazardPotential | pow({Age [1]},8)                | refs=1 = 3.90625E13
Invasive cancer [0] | ExponentialHazardPotential | log(0.849828*(-0.0003717*{Age [0]}+ ...)) | refs=3 = -6.3118...
\end{Verbatim}
\paragraph{El código}
io/src/main/java/org/openmarkov/io/probmodel/reader/PGMXReader\_0\_2.java:1480-1495

\begin{Shaded}
\begin{Highlighting}[]
\KeywordTok{protected} \DataTypeTok{static}\NormalTok{ VariableExpression}\OperatorTok{[]} \FunctionTok{getCovariates}\OperatorTok{(}\BuiltInTok{Element}\NormalTok{ xmlCovariates}\OperatorTok{,} \BuiltInTok{List}\OperatorTok{\textless{}}\NormalTok{Variable}\OperatorTok{\textgreater{}}\NormalTok{ variables}\OperatorTok{)} \OperatorTok{\{}
\NormalTok{    VariableExpression}\OperatorTok{[]}\NormalTok{ covariates }\OperatorTok{=} \KeywordTok{new}\NormalTok{ VariableExpression}\OperatorTok{[}\NormalTok{xmlCovariates}\OperatorTok{.}\FunctionTok{getChildren}\OperatorTok{().}\FunctionTok{size}\OperatorTok{()];}
    \DataTypeTok{int}\NormalTok{ i }\OperatorTok{=} \DecValTok{0}\OperatorTok{;}
    \ControlFlowTok{for} \OperatorTok{(}\BuiltInTok{Element}\NormalTok{ xmlCovariate }\OperatorTok{:}\NormalTok{ xmlCovariates}\OperatorTok{.}\FunctionTok{getChildren}\OperatorTok{())} \OperatorTok{\{}
        \BuiltInTok{String}\NormalTok{ xmlCovariateText }\OperatorTok{=}\NormalTok{ xmlCovariate}\OperatorTok{.}\FunctionTok{getText}\OperatorTok{();}
        \DataTypeTok{var}\NormalTok{ variableCovariate }\OperatorTok{=}\NormalTok{ variables}\OperatorTok{.}\FunctionTok{stream}\OperatorTok{()}
                                         \OperatorTok{.}\FunctionTok{filter}\OperatorTok{(}\DataTypeTok{var} \OperatorTok{{-}\textgreater{}} \DataTypeTok{var}\OperatorTok{.}\FunctionTok{getName}\OperatorTok{().}\FunctionTok{equals}\OperatorTok{(}\NormalTok{xmlCovariateText}\OperatorTok{))}
                                         \OperatorTok{.}\FunctionTok{findFirst}\OperatorTok{();}
        \ControlFlowTok{if} \OperatorTok{(}\NormalTok{variableCovariate}\OperatorTok{.}\FunctionTok{isPresent}\OperatorTok{())} \OperatorTok{\{}
\NormalTok{            covariates}\OperatorTok{[}\NormalTok{i}\OperatorTok{++]} \OperatorTok{=}\NormalTok{ variableCovariate}\OperatorTok{.}\FunctionTok{get}\OperatorTok{().}\FunctionTok{asVariableExpression}\OperatorTok{();}
        \OperatorTok{\}} \ControlFlowTok{else} \OperatorTok{\{}
\NormalTok{            covariates}\OperatorTok{[}\NormalTok{i}\OperatorTok{++]} \OperatorTok{=} \KeywordTok{new} \FunctionTok{VariableExpression}\OperatorTok{(}\NormalTok{variables}\OperatorTok{,}\NormalTok{ xmlCovariateText}\OperatorTok{);}
        \OperatorTok{\}}
    \OperatorTok{\}}
    \ControlFlowTok{return}\NormalTok{ covariates}\OperatorTok{;}
\OperatorTok{\}}
\end{Highlighting}
\end{Shaded}

gui/src/main/java/org/openmarkov/gui/dialog/common/ArithmeticExpressionDialog.java:195-204

\begin{Shaded}
\begin{Highlighting}[]
\KeywordTok{private} \BuiltInTok{String} \FunctionTok{processExpression}\OperatorTok{(}\BuiltInTok{String}\NormalTok{ expression}\OperatorTok{)} \OperatorTok{\{}
    \BuiltInTok{String}\NormalTok{ processedExpression }\OperatorTok{=}\NormalTok{ expression}\OperatorTok{;}
    \ControlFlowTok{for} \OperatorTok{(}\DataTypeTok{int}\NormalTok{ i }\OperatorTok{=} \DecValTok{0}\OperatorTok{;}\NormalTok{ i }\OperatorTok{\textless{}}\NormalTok{ variables}\OperatorTok{.}\FunctionTok{size}\OperatorTok{();}\NormalTok{ i}\OperatorTok{++)} \OperatorTok{\{}
\NormalTok{        processedExpression }\OperatorTok{=}\NormalTok{ processedExpression}\OperatorTok{.}\FunctionTok{replace}\OperatorTok{(}\StringTok{"\{"} \OperatorTok{+}\NormalTok{ variables}\OperatorTok{.}\FunctionTok{get}\OperatorTok{(}\NormalTok{i}\OperatorTok{).}\FunctionTok{getName}\OperatorTok{()} \OperatorTok{+} \StringTok{"\}"}\OperatorTok{,} \StringTok{"\#\{v"} \OperatorTok{+}\NormalTok{ i }\OperatorTok{+} \StringTok{"\}"}\OperatorTok{);}
    \OperatorTok{\}}
    \ControlFlowTok{for} \OperatorTok{(}\DataTypeTok{int}\NormalTok{ i }\OperatorTok{=} \DecValTok{0}\OperatorTok{;}\NormalTok{ i }\OperatorTok{\textless{}}\NormalTok{ variables}\OperatorTok{.}\FunctionTok{size}\OperatorTok{();}\NormalTok{ i}\OperatorTok{++)} \OperatorTok{\{}
\NormalTok{        processedExpression }\OperatorTok{=}\NormalTok{ processedExpression}\OperatorTok{.}\FunctionTok{replace}\OperatorTok{(}\NormalTok{variables}\OperatorTok{.}\FunctionTok{get}\OperatorTok{(}\NormalTok{i}\OperatorTok{).}\FunctionTok{getName}\OperatorTok{(),} \StringTok{"\#\{v"} \OperatorTok{+}\NormalTok{ i }\OperatorTok{+} \StringTok{"\}"}\OperatorTok{);}
    \OperatorTok{\}}
    \ControlFlowTok{return}\NormalTok{ processedExpression}\OperatorTok{;}
\OperatorTok{\}}
\end{Highlighting}
\end{Shaded}
\paragraph{Cómo arreglarlo}
Hace falta decidir cuál es la regla (ver \texttt{decision-de-equipo}), y aplicarla igual en los cuatro sitios: lector, escritor, diálogo y evaluador.

\begin{itemize}
\tightlist
\item
  \textbf{Opción A --- las llaves son obligatorias.} El lector detecta una covariable con fórmula que no contiene ninguna referencia pero sí el nombre de alguna variable del potencial, y lo dice (aviso al abrir o excepción de lectura con el nombre del nodo y el texto). \texttt{Arithmetic\allowbreak{}Expression\allowbreak{}Dialog.\allowbreak{}process\allowbreak{}Expression} y \texttt{Variable\allowbreak{}Expression\allowbreak{}Text\allowbreak{}Field.\allowbreak{}process\allowbreak{}Expression} quitan el segundo bucle (el que acepta nombres sin llaves), para que el diálogo no dé por buena una expresión que la inferencia no puede calcular. \texttt{MID-mammography.\allowbreak{}pgmx} se corrige a mano poniendo las llaves.
\item
  \textbf{Opción B --- los nombres sin llaves se aceptan al leer.} El lector reescribe el texto poniendo llaves a los nombres de variables que aparezcan. Ojo: una sustitución ingenua por \texttt{String.\allowbreak{}replace} (como la que ya hace \texttt{GLMPotential.\allowbreak{}process\allowbreak{}Covariate}, líneas 278-290, hoy sin llamadores) confunde un nombre con otro que lo contiene (\texttt{Age\ {[}0{]}} dentro de \texttt{Age\ at\ state\ entry\ {[}0{]}}), y los nombres llevan espacios y corchetes; habría que sustituir primero los nombres más largos y respetar los límites de palabra. El escritor ya escribiría entonces las llaves, y el fichero se repararía solo al guardarlo.
\end{itemize}

En ambos casos, la prueba de regresión: leer un \texttt{.\allowbreak{}pgmx} mínimo con una covariable \texttt{pow(\allowbreak{}X\ {[}0{]},\allowbreak{}2)} y comprobar el comportamiento elegido (aviso o rechazo, o que se evalúa a X² y que al guardar sale \texttt{pow(\allowbreak{}\{X\ {[}0{]}\},\allowbreak{}2)}); y una prueba de la validación del diálogo con el mismo texto.

\setcounter{subsection}{13}
\subsection[Una regresión sin \texttt{Constant} o \texttt{Gamma} muere con un índice −1 sin mensaje]{\textcolor{corregidoverde}{[Corregido]} Una regresión sin su covariable \texttt{Constant} (o un riesgo de Weibull sin \texttt{Gamma}) muere al calcularse con un índice −1 que no dice qué falta}\label{err:14}
\begin{ficha}
\fichakey{Estado}Corregido en el commit \texttt{522d926} (23-09-2026): al proyectarse, una regresión sin \texttt{Constant} o un riesgo de Weibull sin \texttt{Gamma} lanza \texttt{Missing\allowbreak{}Mandatory\allowbreak{}Covariate} con el nombre de la covariable que falta, y los accesores lo dicen también; el panel ya no deja quitar \texttt{Gamma} y la red de prueba lleva un riesgo exponencial.
\fichakey{Módulo}core
\fichakey{Paquete}org.openmarkov.core.model.network.potential
\fichakey{Ficheros}\path{core/src/main/java/org/openmarkov/core/model/network/potential/GLMPotential.java:168-174} \newline \path{core/src/main/java/org/openmarkov/core/model/network/potential/GLMPotential.java:308-318} \newline \path{core/src/main/java/org/openmarkov/core/model/network/potential/WeibullHazardPotential.java:122-128} \newline \path{core/src/main/java/org/openmarkov/core/model/network/potential/WeibullHazardPotential.java:137-138} \newline \path{core/src/main/java/org/openmarkov/core/model/network/potential/WeibullHazardPotential.java:176} \newline \path{core/src/main/java/org/openmarkov/core/model/network/potential/WeibullHazardPotential.java:339-349} \newline \path{core/src/main/java/org/openmarkov/core/model/network/potential/LinearCombinationPotential.java:97-113} \newline \path{core/src/main/java/org/openmarkov/core/model/network/potential/ExponentialPotential.java:84-107}
\fichakey{Comprobado}Leyendo el código y ejecutándolo
\fichakey{Esfuerzo}Pequeño (horas)
\fichakey{Decisión de equipo}No
\fichakey{Relacionados}\hyperref[err:10]{error~10}, \hyperref[err:15]{error~15}, \hyperref[err:16]{error~16}, \hyperref[err:71]{error~71}, \hyperref[err:157]{error~157}
\end{ficha}
\paragraph{Qué pasa}
Un potencial de regresión al que le falta su covariable \texttt{Constant}, o un riesgo de Weibull al que le falta \texttt{Gamma}, muere al calcularse con \texttt{Array\allowbreak{}Index\allowbreak{}Out\allowbreak{}Of\allowbreak{}Bounds\allowbreak{}Exception:\ Index\ -1\ out\ of\ bounds\ for\ length\ N}, que no dice qué covariable falta ni en qué nodo.

Caminos de llegada:

\begin{itemize}
\tightlist
\item
  \textbf{Fichero:} el lector de \texttt{.\allowbreak{}pgmx} no comprueba que estén los marcadores obligatorios; cualquier inferencia que proyecte un potencial así falla.
\item
  \textbf{Interfaz:} en el panel de covariables (\texttt{GLMPanel}), el botón de quitar se apaga cuando la fila elegida es obligatoria, pero la lista de obligatorias se pide siempre a \texttt{GLMPotential.\allowbreak{}get\allowbreak{}Mandatory\allowbreak{}Covariates(\allowbreak{})}, que solo contiene \texttt{Constant}. Para un riesgo de Weibull se puede quitar la fila \texttt{Gamma} y la proyección falla después; es lo que describe \hyperref[err:157]{error~157} (comprobado leyendo). \texttt{Constant} sí está protegida en la interfaz.
\end{itemize}

Hay una red del propio repositorio así: \texttt{integration\allowbreak{}Tests/\allowbreak{}src/\allowbreak{}main/\allowbreak{}resources/\allowbreak{}networks\_\allowbreak{}for\_\allowbreak{}every\_\allowbreak{}potential/\allowbreak{}Dynamic-LIMID-every-potential.\allowbreak{}pgmx}, la que usan \texttt{Potentials\allowbreak{}Obey\allowbreak{}The\allowbreak{}Equality\allowbreak{}Contract\allowbreak{}Test} y \texttt{Potentials\allowbreak{}Can\allowbreak{}Be\allowbreak{}Cloned\allowbreak{}Test} para recorrer todos los tipos de potencial, tiene el nodo \texttt{Exponential\allowbreak{}Hazard\ {[}0{]}} escrito con \texttt{type="Hazard\ (\allowbreak{}Weibull)"} y \texttt{Constant} como única covariable. Esas pruebas no proyectan el potencial, así que no lo detectan.

Por qué pasa: los potenciales de regresión (familia \texttt{GLMPotential}) marcan dos covariables especiales con un texto fijo: \texttt{Constant} (el término independiente, obligatorio en todas) y \texttt{Gamma} (el logaritmo del parámetro de forma, obligatorio en el riesgo de Weibull, \texttt{Weibull\allowbreak{}Hazard\allowbreak{}Potential}). Su posición se busca con \texttt{GLMPotential.\allowbreak{}get\allowbreak{}Constant\allowbreak{}Index} y \texttt{Weibull\allowbreak{}Hazard\allowbreak{}Potential.\allowbreak{}get\allowbreak{}Gamma\allowbreak{}Index}, que devuelven \textbf{−1} si el marcador no está. Quien las llama indexa directamente con lo devuelto:

\begin{itemize}
\tightlist
\item
  \texttt{Weibull\allowbreak{}Hazard\allowbreak{}Potential.\allowbreak{}table\allowbreak{}Project}: \texttt{Math.\allowbreak{}exp(\allowbreak{}coefficients{[}gamma\allowbreak{}Index{]})} (línea 176) y \texttt{coefficients{[}constant\allowbreak{}Index{]}};
\item
  \texttt{Linear\allowbreak{}Combination\allowbreak{}Potential.\allowbreak{}table\allowbreak{}Project} (línea 113) y \texttt{Exponential\allowbreak{}Potential.\allowbreak{}table\allowbreak{}Project} (línea 107): \texttt{coefficients{[}constant\allowbreak{}Index{]}};
\item
  los accesores \texttt{GLMPotential.\allowbreak{}get\allowbreak{}Constant}/\texttt{set\allowbreak{}Constant} y \texttt{Weibull\allowbreak{}Hazard\allowbreak{}Potential.\allowbreak{}get\allowbreak{}Gamma}/\texttt{set\allowbreak{}Gamma}.
\end{itemize}

Medido (el primer caso es el nodo de la red del repositorio citada arriba):

\begin{Verbatim}[breaklines,breakanywhere,fontsize=\small]
ExponentialHazard [0]  WeibullHazardPotential [Constant]         -> ArrayIndexOutOfBoundsException: Index -1 (WeibullHazardPotential.java:176)
WeibullHazard [0]      WeibullHazardPotential [Gamma, Constant]  -> [1.0, 0.0]
LinearCombination sin Constant                                    -> ArrayIndexOutOfBoundsException: Index -1 (LinearCombinationPotential.java:113)
Exponential sin Constant                                          -> ArrayIndexOutOfBoundsException: Index -1 (ExponentialPotential.java:107)
getConstant() sin Constant                                        -> ArrayIndexOutOfBoundsException: Index -1
\end{Verbatim}
\paragraph{El código}
core/src/main/java/org/openmarkov/core/model/network/potential/GLMPotential.java:308-318

\begin{Shaded}
\begin{Highlighting}[]
\KeywordTok{protected} \DataTypeTok{static} \DataTypeTok{int} \FunctionTok{getConstantIndex}\OperatorTok{(}\NormalTok{VariableExpression}\OperatorTok{[]}\NormalTok{ covariates}\OperatorTok{)} \OperatorTok{\{}
    \DataTypeTok{int}\NormalTok{ constantIndex }\OperatorTok{=} \OperatorTok{{-}}\DecValTok{1}\OperatorTok{;}
    \DataTypeTok{int}\NormalTok{ i }\OperatorTok{=} \DecValTok{0}\OperatorTok{;}
    \ControlFlowTok{while} \OperatorTok{(}\NormalTok{i }\OperatorTok{\textless{}}\NormalTok{ covariates}\OperatorTok{.}\FunctionTok{length} \OperatorTok{\&\&}\NormalTok{ constantIndex }\OperatorTok{==} \OperatorTok{{-}}\DecValTok{1}\OperatorTok{)} \OperatorTok{\{}
        \ControlFlowTok{if} \OperatorTok{(}\NormalTok{covariates}\OperatorTok{[}\NormalTok{i}\OperatorTok{].}\FunctionTok{asStringExpression}\OperatorTok{().}\FunctionTok{equals}\OperatorTok{(}\NormalTok{VariableExpression}\OperatorTok{.}\FunctionTok{Common}\OperatorTok{.}\FunctionTok{CONSTANT}\OperatorTok{.}\FunctionTok{asStringExpression}\OperatorTok{()))} \OperatorTok{\{}
\NormalTok{            constantIndex }\OperatorTok{=}\NormalTok{ i}\OperatorTok{;}
        \OperatorTok{\}}
        \OperatorTok{++}\NormalTok{i}\OperatorTok{;}
    \OperatorTok{\}}
    \ControlFlowTok{return}\NormalTok{ constantIndex}\OperatorTok{;}
\OperatorTok{\}}
\end{Highlighting}
\end{Shaded}

core/src/main/java/org/openmarkov/core/model/network/potential/WeibullHazardPotential.java:137-138 y 176

\begin{Shaded}
\begin{Highlighting}[]
\DataTypeTok{int}\NormalTok{ gammaIndex }\OperatorTok{=} \FunctionTok{getGammaIndex}\OperatorTok{(}\NormalTok{covariates}\OperatorTok{);}
\DataTypeTok{int}\NormalTok{ constantIndex }\OperatorTok{=} \FunctionTok{getConstantIndex}\OperatorTok{(}\NormalTok{covariates}\OperatorTok{);}
\KeywordTok{...}
\DataTypeTok{double}\NormalTok{ shape }\OperatorTok{=} \BuiltInTok{Math}\OperatorTok{.}\FunctionTok{exp}\OperatorTok{(}\NormalTok{coefficients}\OperatorTok{[}\NormalTok{gammaIndex}\OperatorTok{]);}
\end{Highlighting}
\end{Shaded}
\paragraph{Cómo arreglarlo}
Comprobar los marcadores en un único punto por el que pasan todas las proyecciones: \texttt{GLMPotential.\allowbreak{}table\allowbreak{}Project(\allowbreak{}Evidence\allowbreak{}Case,\allowbreak{}\ Inference\allowbreak{}Options,\allowbreak{}\ List)} (líneas 203-245), que ya comprueba que haya tantos coeficientes como covariables y lanza \texttt{Non\allowbreak{}Projectable\allowbreak{}Potential\allowbreak{}Exception.\allowbreak{}Coefficients\allowbreak{}Do\allowbreak{}Not\allowbreak{}Match\allowbreak{}Covariates}. Añadir ahí una comprobación de que cada covariable obligatoria de la clase concreta está presente, lanzando una excepción del mismo tipo que diga el nodo y la covariable que falta.

Para que la lista de obligatorias dependa de la clase, hace falta un método de instancia (por ejemplo \texttt{get\allowbreak{}Mandatory\allowbreak{}Covariates\allowbreak{}Of\allowbreak{}This\allowbreak{}Potential(\allowbreak{})}), que \texttt{Weibull\allowbreak{}Hazard\allowbreak{}Potential} redefina con \texttt{MANDATORY\_\allowbreak{}COVARIATES} (líneas 40-43); \texttt{Exponential\allowbreak{}Hazard\allowbreak{}Potential} antepone \texttt{Gamma} por su cuenta al proyectar y no la necesita en su lista.

La misma regla debe cumplirse en:

\begin{itemize}
\tightlist
\item
  \texttt{GLMPanel.\allowbreak{}value\allowbreak{}Changed} (\texttt{gui/\allowbreak{}src/\allowbreak{}main/\allowbreak{}java/\allowbreak{}org/\allowbreak{}openmarkov/\allowbreak{}gui/\allowbreak{}dialog/\allowbreak{}common/\allowbreak{}GLMPanel.\allowbreak{}java:136}) y en el doble clic de edición (línea 164), que deben usar ese método de instancia en vez del estático de \texttt{GLMPotential} (es el arreglo de \hyperref[err:157]{error~157});
\item
  los accesores \texttt{get\allowbreak{}Constant}/\texttt{set\allowbreak{}Constant}/\texttt{get\allowbreak{}Gamma}/\texttt{set\allowbreak{}Gamma}, que deben lanzar la misma excepción clara;
\item
  opcionalmente, el lector de \texttt{.\allowbreak{}pgmx} (\texttt{PGMXPotential\allowbreak{}Parsers}), que podría avisar al abrir.
\end{itemize}

Además, corregir la red \texttt{Dynamic-LIMID-every-potential.\allowbreak{}pgmx}: o se cambia el tipo del nodo \texttt{Exponential\allowbreak{}Hazard\ {[}0{]}} a \texttt{Hazard\ (\allowbreak{}Exponential)}, o se le añade \texttt{Gamma}.

Prueba de regresión: en \texttt{core/\allowbreak{}src/\allowbreak{}test}, proyectar un \texttt{Weibull\allowbreak{}Hazard\allowbreak{}Potential} sin \texttt{Gamma}, un \texttt{Linear\allowbreak{}Combination\allowbreak{}Potential} y un \texttt{Exponential\allowbreak{}Potential} sin \texttt{Constant}, y comprobar que lanzan la excepción nueva con el nombre de la covariable ausente; y una prueba de integración que proyecte todos los potenciales de \texttt{Dynamic-LIMID-every-potential.\allowbreak{}pgmx}.

\setcounter{subsection}{14}
\subsection[Una utilidad de regresión no cuenta como utilidad al calcular]{Un nodo de utilidad definido como regresión no cuenta como utilidad: el coste-efectividad falla con ClassCastException o pierde ese coste sin avisar, y la evaluación puede dar utilidad 0}\label{err:15}
\begin{ficha}
\fichakey{Estado}Presente
\fichakey{Módulo}core, inference
\fichakey{Paquete}org.openmarkov.core.model.network.potential; org.openmarkov.inference.algorithm.variableElimination
\fichakey{Ficheros}\path{core/src/main/java/org/openmarkov/core/model/network/potential/LinearCombinationPotential.java:84-131} \newline \path{core/src/main/java/org/openmarkov/core/model/network/potential/ExponentialPotential.java:80-118} \newline \path{core/src/main/java/org/openmarkov/core/model/network/potential/ExactDistrPotential.java:69-75} \newline \path{core/src/main/java/org/openmarkov/core/model/network/potential/ExactDistrPotential.java:93-100} \newline \path{core/src/main/java/org/openmarkov/core/model/network/potential/treeadd/TreeADDPotential.java:669-690} \newline \path{core/src/main/java/org/openmarkov/core/model/network/potential/Potential.java:432-434} \newline \path{core/src/main/java/org/openmarkov/core/model/network/ProbNetPotentialQueries.java:101-121} \newline \path{inference/src/main/java/org/openmarkov/inference/algorithm/variableElimination/VariableEliminationCore.java:192-217} \newline \path{inference/src/main/java/org/openmarkov/inference/algorithm/variableElimination/DecisionVariableElimination.java:43-49}
\fichakey{Comprobado}Leyendo el código y ejecutándolo
\fichakey{Esfuerzo}Pequeño (horas)
\fichakey{Decisión de equipo}No
\fichakey{Tapado por}\hyperref[err:71]{error~71}: en las redes CHAP del repositorio, que fallan antes. \newline \hyperref[err:74]{error~74}: en el jar publicado la ventana del ámbito de decisión no se abre; desde el IDE se ve.
\fichakey{Relacionados}\hyperref[err:10]{error~10}, \hyperref[err:12]{error~12}, \hyperref[err:14]{error~14}, \hyperref[err:71]{error~71}, \hyperref[err:74]{error~74}, \hyperref[nota:28]{nota~28}
\end{ficha}
\paragraph{Qué pasa}
Un nodo de utilidad cuya relación es una regresión ---«Linear combination» o «Exponential», que el diálogo del nodo ofrece para cualquier nodo de utilidad--- no se trata como utilidad al calcular. Lo que ve el usuario depende de qué variables use el nodo:

\begin{itemize}
\tightlist
\item
  \textbf{Análisis de coste-efectividad con ámbito global} (menú Tools → «Cost-effectiveness analysis»), si el nodo depende de alguna decisión: aviso de error inesperado con \texttt{Class\allowbreak{}Cast\allowbreak{}Exception:\ GTable\allowbreak{}Potential\ cannot\ be\ cast\ to\ Table\allowbreak{}Potential}.
\item
  \textbf{Ámbito de decisión}, si el nodo solo depende de la decisión elegida (el coste de una prueba o de una prótesis, por ejemplo): el análisis termina y la ventana enseña los costes, o las efectividades, \textbf{sin ese nodo}, sin ningún aviso. Si depende de otra decisión, falla como el global.
\item
  \textbf{Si el nodo depende de una variable de azar}: la evaluación de un solo criterio da utilidad 0 y el coste-efectividad da efectividad (o coste) 0 en todas las opciones, sin aviso.
\end{itemize}

Lo mismo pasa con un árbol (\texttt{Tree\allowbreak{}ADDPotential}) cuyas hojas son todas regresiones. Un árbol que mezcla una hoja de regresión con una «Exact» sí funciona, porque toma el criterio de la hoja «Exact»: por eso las redes CHAP del repositorio, cuyos costes son exponenciales colgadas de un árbol con una rama «Exact» para «dead», no lo muestran (fallan antes por \hyperref[err:71]{error~71}). Ninguna otra red de \texttt{0\_\allowbreak{}miscellaneous/\allowbreak{}networks} en formato 0.2 tiene una regresión como utilidad, así que se llega construyendo el modelo en el editor. Con el jar ejecutable publicado, la ventana del ámbito de decisión no llega a abrirse (\hyperref[err:74]{error~74}); los costes equivocados se ven desde el IDE.

Por qué pasa. Durante la inferencia, un potencial cuenta como utilidad si lleva criterio de decisión (\texttt{Potential.\allowbreak{}is\allowbreak{}Additive(\allowbreak{})} es \texttt{criterion\ !=\ null}); en edición el criterio lo tiene la variable del nodo, no el potencial. \texttt{Exact\allowbreak{}Distr\allowbreak{}Potential.\allowbreak{}table\allowbreak{}Project} copia a la tabla proyectada el criterio de la variable del nodo, pero \texttt{Linear\allowbreak{}Combination\allowbreak{}Potential.\allowbreak{}table\allowbreak{}Project} y \texttt{Exponential\allowbreak{}Potential.\allowbreak{}table\allowbreak{}Project} crean la tabla con \texttt{new\ Table\allowbreak{}Potential(\allowbreak{}variables,\allowbreak{}\ role)} y no se lo ponen, y \texttt{Tree\allowbreak{}ADDPotential} toma el de la primera rama que lo tenga. Sin criterio:

\begin{itemize}
\tightlist
\item
  \texttt{Variable\allowbreak{}Elimination\allowbreak{}Core.\allowbreak{}create\allowbreak{}CEPPotential}, que junta costes y efectividades en la partición de coste-efectividad antes de eliminar la primera decisión, solo recoge los potenciales con criterio, así que esa tabla se queda fuera;
\item
  al eliminar una decisión de la que depende, la tabla llega a \texttt{Decision\allowbreak{}Variable\allowbreak{}Elimination} junto a la partición, y la suma (\texttt{Table\allowbreak{}Potential\allowbreak{}Arithmetic.\allowbreak{}sum}, línea 218, llamada desde la línea 48) hace la conversión que falla;
\item
  si depende de una variable de azar, el valor se pierde al eliminar esa variable y la utilidad sale 0 (medido; no se ha seguido paso a paso).
\end{itemize}

Medido, sustituyendo en redes del repositorio la utilidad «Exact» por la regresión que da los mismos números, creada como la crea el diálogo del nodo (\texttt{Potential\allowbreak{}Utils.\allowbreak{}instanciate\allowbreak{}Safely} con el papel del potencial anterior):

\begin{itemize}
\tightlist
\item
  \texttt{0\_\allowbreak{}miscellaneous/\allowbreak{}networks/\allowbreak{}mid/\allowbreak{}MID-hip-Briggs.\allowbreak{}pgmx}, «Prosthesis cost» como combinación lineal \texttt{394\ +\ 185·\{Prosthesis\ type\}} (proyecta a 394 y 579, como el original): ámbito global → \texttt{Class\allowbreak{}Cast\allowbreak{}Exception}; ámbito de decisión «Prosthesis type» → costes 163,6 y 43,6, cuando con el original son 557,6 y 622,6 (la prótesis NP1 pasa a parecer la barata). La evaluación de un solo criterio da lo mismo que el original, 636,23.
\item
  \texttt{0\_\allowbreak{}miscellaneous/\allowbreak{}networks/\allowbreak{}id/\allowbreak{}ID-CEA-minimal.\allowbreak{}pgmx}, «Cost» como \texttt{14000·\{Therapy\}}: global → \texttt{Class\allowbreak{}Cast\allowbreak{}Exception}; ámbito de decisión → costes 0 y 0 en vez de 0 y 14 000. La exponencial da lo mismo.
\item
  \texttt{0\_\allowbreak{}miscellaneous/\allowbreak{}networks/\allowbreak{}id/\allowbreak{}ID-CEA-test-2therapies-new-test.\allowbreak{}pgmx}, «C: Cheap test» como \texttt{18·\{Dec:\ Test\}}: global y ámbitos de «Dec: New test» y «Therapy» → \texttt{Class\allowbreak{}Cast\allowbreak{}Exception}; ámbito de «Dec: Test» → el coste de 18 de la prueba desaparece de las dos opciones.
\item
  La misma red, con «QALE» (depende de «Disease» y «Therapy») definida como «Exact» y como combinación lineal con los mismos números \texttt{5\ −\ 2·\{Disease\}\ +\ \{Therapy\}}: con «Exact», utilidad 151 600 y efectividades 4,72, 5,72 y 6,72; con la combinación lineal, utilidad 0 y efectividades 0, 0 y 0.
\item
  Un árbol con dos hojas de combinación lineal en «Cost» de \texttt{ID-CEA-minimal.\allowbreak{}pgmx} falla igual; con una hoja lineal y otra «Exact», da los costes correctos.
\end{itemize}
\paragraph{El código}
core/src/main/java/org/openmarkov/core/model/network/potential/ExactDistrPotential.java:69-75

\begin{Shaded}
\begin{Highlighting}[]
\KeywordTok{public} \AttributeTok{@NotNull}\NormalTok{ TablePotential }\FunctionTok{tableProject}\OperatorTok{(}\NormalTok{EvidenceCase evidenceCase}\OperatorTok{,}\NormalTok{ InferenceOptions inferenceOptions}\OperatorTok{)} \KeywordTok{throws}\NormalTok{ NonProjectablePotentialException }\OperatorTok{\{}
    \CommentTok{// get the projected TablePotential, which will be returned inside a list}
\NormalTok{    TablePotential projectedPotential }\OperatorTok{=}\NormalTok{ tablePotential}\OperatorTok{.}\FunctionTok{tableProject}\OperatorTok{(}\NormalTok{evidenceCase}\OperatorTok{,}\NormalTok{ inferenceOptions}\OperatorTok{);}
\NormalTok{    projectedPotential}\OperatorTok{.}\FunctionTok{setCriterion}\OperatorTok{(}\FunctionTok{getChildVariable}\OperatorTok{().}\FunctionTok{getDecisionCriterion}\OperatorTok{());}
\NormalTok{    projectedPotential}\OperatorTok{.}\FunctionTok{setPotentialRole}\OperatorTok{(}\NormalTok{PotentialRole}\OperatorTok{.}\FunctionTok{UNSPECIFIED}\OperatorTok{);}
    \ControlFlowTok{return}\NormalTok{ projectedPotential}\OperatorTok{;}
\OperatorTok{\}}
\end{Highlighting}
\end{Shaded}

core/src/main/java/org/openmarkov/core/model/network/potential/LinearCombinationPotential.java:93-95 (lo mismo en \texttt{Exponential\allowbreak{}Potential.\allowbreak{}java}, línea 88)

\begin{Shaded}
\begin{Highlighting}[]
\BuiltInTok{List}\OperatorTok{\textless{}}\NormalTok{Variable}\OperatorTok{\textgreater{}}\NormalTok{ projectedPotentialVariables }\OperatorTok{=} \KeywordTok{new} \BuiltInTok{ArrayList}\OperatorTok{\textless{}\textgreater{}(}\NormalTok{evidencelessVariables}\OperatorTok{);}
\NormalTok{projectedPotentialVariables}\OperatorTok{.}\FunctionTok{addFirst}\OperatorTok{(}\NormalTok{variables}\OperatorTok{.}\FunctionTok{getFirst}\OperatorTok{());}
\NormalTok{TablePotential projectedPotential }\OperatorTok{=} \KeywordTok{new} \FunctionTok{TablePotential}\OperatorTok{(}\NormalTok{projectedPotentialVariables}\OperatorTok{,}\NormalTok{ role}\OperatorTok{);}
\end{Highlighting}
\end{Shaded}

inference/src/main/java/org/openmarkov/inference/algorithm/variableElimination/VariableEliminationCore.java:197-204

\begin{Shaded}
\begin{Highlighting}[]
\ControlFlowTok{for} \OperatorTok{(}\NormalTok{Potential potential }\OperatorTok{:}\NormalTok{ markovDecisionNetwork}\OperatorTok{.}\FunctionTok{getAdditivePotentials}\OperatorTok{())} \OperatorTok{\{}
\NormalTok{    markovDecisionNetwork}\OperatorTok{.}\FunctionTok{removePotential}\OperatorTok{(}\NormalTok{potential}\OperatorTok{);}
    \ControlFlowTok{if} \OperatorTok{(}\NormalTok{potential}\OperatorTok{.}\FunctionTok{getCriterion}\OperatorTok{().}\FunctionTok{getCECriterion}\OperatorTok{()} \OperatorTok{==}\NormalTok{ Criterion}\OperatorTok{.}\FunctionTok{CECriterion}\OperatorTok{.}\FunctionTok{Cost}\OperatorTok{)} \OperatorTok{\{}
\NormalTok{        costPotentials}\OperatorTok{.}\FunctionTok{add}\OperatorTok{((}\NormalTok{TablePotential}\OperatorTok{)}\NormalTok{ potential}\OperatorTok{);}
    \OperatorTok{\}} \ControlFlowTok{else} \OperatorTok{\{}
\NormalTok{        effectivenessPotentials}\OperatorTok{.}\FunctionTok{add}\OperatorTok{((}\NormalTok{TablePotential}\OperatorTok{)}\NormalTok{ potential}\OperatorTok{);}
    \OperatorTok{\}}
\OperatorTok{\}}
\end{Highlighting}
\end{Shaded}
\paragraph{Cómo arreglarlo}
Que toda tabla proyectada desde el potencial de un nodo de utilidad lleve el criterio de la variable del nodo. Hay dos formas:

\begin{itemize}
\tightlist
\item
  clase por clase, como \texttt{Exact\allowbreak{}Distr\allowbreak{}Potential}: en \texttt{GLMPotential.\allowbreak{}table\allowbreak{}Project} (por donde pasan combinación lineal, exponencial, riesgos de Weibull y exponencial y función), en \texttt{Table\allowbreak{}Potential.\allowbreak{}table\allowbreak{}Project} (que tampoco lo copia: con una utilidad en forma de tabla y una política impuesta determinista, la utilidad esperada sale 0; medido sobre \texttt{integration\allowbreak{}Tests/\allowbreak{}src/\allowbreak{}main/\allowbreak{}resources/\allowbreak{}networks/\allowbreak{}id/\allowbreak{}ID-no-knowledge.\allowbreak{}pgmx} cambiando a mano el potencial de utilidad) y en \texttt{Tree\allowbreak{}ADDPotential}, que hoy toma el criterio de sus ramas;
\item
  en un solo sitio, \texttt{Prob\allowbreak{}Net\allowbreak{}Potential\allowbreak{}Queries.\allowbreak{}table\allowbreak{}Project\allowbreak{}Potentials}, que proyecta todos los potenciales de la red antes de eliminar: si la variable del nodo tiene criterio y la tabla no, ponérselo. Así ninguna clase nueva puede olvidarlo.
\end{itemize}

Además, \texttt{Variable\allowbreak{}Elimination\allowbreak{}Core.\allowbreak{}create\allowbreak{}CEPPotential} no debería dejar fuera en silencio una tabla de un nodo de utilidad sin criterio: mejor fallar con un mensaje que lo diga.

Prueba de regresión: con \texttt{0\_\allowbreak{}miscellaneous/\allowbreak{}networks/\allowbreak{}id/\allowbreak{}ID-CEA-minimal.\allowbreak{}pgmx} y «Cost» como combinación lineal \texttt{14000·\{Therapy\}}, \texttt{VECEAnalysis} debe dar costes 0 y 14 000 en ámbito global y en el de «Therapy»; y con \texttt{0\_\allowbreak{}miscellaneous/\allowbreak{}networks/\allowbreak{}id/\allowbreak{}ID-CEA-test-2therapies-new-test.\allowbreak{}pgmx} y «QALE» como \texttt{5\ −\ 2·\{Disease\}\ +\ \{Therapy\}}, \texttt{VEEvaluation} debe dar 151 600, lo mismo que la tabla «Exact» con esos números.

\setcounter{subsection}{15}
\subsection[Una tabla de fórmulas con \texttt{Complement} en vez de \texttt{\#} no se puede proyectar]{Una tabla de fórmulas que marca «lo que sobra» con la palabra \texttt{Complement} en vez de \texttt{\#} no se puede proyectar, y la red de prueba de todos los potenciales la usa}\label{err:16}
\begin{ficha}
\fichakey{Estado}Presente
\fichakey{Módulo}core, io, integrationTests
\fichakey{Paquete}org.openmarkov.core.expression; org.openmarkov.core.model.network.potential.operation; org.openmarkov.io.probmodel.reader
\fichakey{Ficheros}\path{core/src/main/java/org/openmarkov/core/expression/VariableExpression.java:22} \newline \path{core/src/main/java/org/openmarkov/core/model/network/potential/operation/AugmentedProbTableInference.java:18-54} \newline \path{io/src/main/java/org/openmarkov/io/probmodel/reader/PGMXReader_1_0.java:189-205} \newline \path{integrationTests/src/main/resources/networks_for_every_potential/Dynamic-LIMID-every-potential.pgmx:304-311}
\fichakey{Comprobado}Leyendo el código y ejecutándolo
\fichakey{Esfuerzo}Pequeño (horas)
\fichakey{Decisión de equipo}\textcolor{decisionrojo}{\textbf{Sí.}} Si la lectura debe aceptar \texttt{Complement} como sinónimo de \texttt{\#} (compatibilidad con ficheros que lo escriban así) o si basta con corregir la red de prueba
\fichakey{Relacionados}\hyperref[err:12]{error~12}, \hyperref[err:14]{error~14}
\end{ficha}
\paragraph{Qué pasa}
Una tabla de probabilidad aumentada («tabla de fórmulas») cuyas casillas de «lo que sobra» están escritas con la palabra \texttt{Complement} en vez de \texttt{\#} se abre sin aviso, pero cualquier inferencia que proyecte ese nodo falla con \texttt{Non\allowbreak{}Projectable\allowbreak{}Potential\allowbreak{}Exception\$Cannot\allowbreak{}Evaluate}. El mensaje habla de una expresión inválida y no dice que la palabra esperada es \texttt{\#}.

Camino de llegada: abrir un \texttt{.\allowbreak{}pgmx} cuya tabla de fórmulas use \texttt{Complement} y lanzar cualquier inferencia que proyecte ese nodo. Dentro del repositorio, el único fichero así es \texttt{integration\allowbreak{}Tests/\allowbreak{}src/\allowbreak{}main/\allowbreak{}resources/\allowbreak{}networks\_\allowbreak{}for\_\allowbreak{}every\_\allowbreak{}potential/\allowbreak{}Dynamic-LIMID-every-potential.\allowbreak{}pgmx}, cuyo nodo \texttt{Augmented\allowbreak{}Prob\allowbreak{}Table\ {[}0{]}} trae \texttt{\textless{}Functions\textgreater{}"1"\ "Complement"\ "1"\ "Complement"\ "1"\ "Complement"\ "1"\ "Complement"\textless{}/\allowbreak{}Functions\textgreater{}}. Las demás apariciones de \texttt{Complement} en los \texttt{.\allowbreak{}pgmx} son la distribución de incertidumbre \texttt{distribution="Complement"}, que es otra cosa y no está afectada. Las pruebas que usan esa red (\texttt{Potentials\allowbreak{}Obey\allowbreak{}The\allowbreak{}Equality\allowbreak{}Contract\allowbreak{}Test}, \texttt{Potentials\allowbreak{}Can\allowbreak{}Be\allowbreak{}Cloned\allowbreak{}Test}) no proyectan, así que no lo detectan.

\texttt{Augmented\allowbreak{}Prob\allowbreak{}Table} conserva comentada una constante \texttt{COMPLEMENT\_\allowbreak{}FUNCTION} que usaba la palabra \texttt{Complement} (\texttt{core/\allowbreak{}src/\allowbreak{}main/\allowbreak{}java/\allowbreak{}org/\allowbreak{}openmarkov/\allowbreak{}core/\allowbreak{}model/\allowbreak{}network/\allowbreak{}potential/\allowbreak{}Augmented\allowbreak{}Prob\allowbreak{}Table.\allowbreak{}java:58-66}, en un bloque comentado), así que es posible que existan ficheros escritos con esa grafía fuera del repositorio.

Cómo funciona el marcador: una tabla de probabilidad aumentada (\texttt{Augmented\allowbreak{}Prob\allowbreak{}Table\allowbreak{}Potential}) guarda en cada casilla una expresión en vez de un número. Una casilla puede decir «lo que sobre»: al proyectar, \texttt{Augmented\allowbreak{}Prob\allowbreak{}Table\allowbreak{}Inference} calcula primero las fórmulas de cada columna y reparte después lo que falta para sumar 1 entre las casillas marcadas como complemento. Una casilla se reconoce como complemento comparando su texto con \texttt{Variable\allowbreak{}Expression.\allowbreak{}Common.\allowbreak{}COMPLEMENT}, que es la cadena \texttt{\#}.

Con \texttt{\#} el mecanismo funciona también desde un fichero escrito por OpenMarkov. La interfaz usa ese mismo marcador: al editar una casilla, \texttt{Augmented\allowbreak{}Potential\allowbreak{}Value\allowbreak{}Edit} (\texttt{gui/\allowbreak{}src/\allowbreak{}main/\allowbreak{}java/\allowbreak{}org/\allowbreak{}openmarkov/\allowbreak{}gui/\allowbreak{}action/\allowbreak{}Augmented\allowbreak{}Potential\allowbreak{}Value\allowbreak{}Edit.\allowbreak{}java:158} y 199) pone \texttt{COMPLEMENT} en las demás casillas de la columna; el escritor \texttt{PGMXWriter\_\allowbreak{}1\_\allowbreak{}0} escribe el texto tal cual (\texttt{\textless{}Functions\textgreater{}"1"\ "\#"\ .\allowbreak{}.\allowbreak{}.\allowbreak{}}) y el lector \texttt{PGMXReader\_\allowbreak{}1\_\allowbreak{}0.\allowbreak{}get\allowbreak{}Augmented\allowbreak{}Prob\allowbreak{}Table} lo relee.

Lo que falla es la palabra \texttt{Complement}: el lector no la traduce, así que la casilla no se reconoce como complemento, se intenta calcular como fórmula y jeval la rechaza.

Medido con la red citada:

\begin{Verbatim}[breaklines,breakanywhere,fontsize=\small]
fichero con "Complement"  -> NonProjectablePotentialException$CannotEvaluate: Cannot evaluate 'Complement' due to Expression is invalid.
mismo fichero con "#"      -> [1.0, 0.0, 1.0, 0.0, 1.0, 0.0, 1.0, 0.0]
escrito por el programa    -> <Functions>"1" "#" "1" "#" "1" "#" "1" "#"</Functions>, y al releer [1.0, 0.0, ...]
\end{Verbatim}
\paragraph{El código}
core/src/main/java/org/openmarkov/core/expression/VariableExpression.java:22

\begin{Shaded}
\begin{Highlighting}[]
\KeywordTok{public} \DataTypeTok{static} \DataTypeTok{final}\NormalTok{ VariableExpression COMPLEMENT }\OperatorTok{=} \KeywordTok{new} \FunctionTok{VariableExpression}\OperatorTok{(}\BuiltInTok{Collections}\OperatorTok{.}\FunctionTok{emptyList}\OperatorTok{(),} \StringTok{"\#"}\OperatorTok{);}
\end{Highlighting}
\end{Shaded}

io/src/main/java/org/openmarkov/io/probmodel/reader/PGMXReader\_1\_0.java:191-203

\begin{Shaded}
\begin{Highlighting}[]
\BuiltInTok{String}\NormalTok{ functionsValue }\OperatorTok{=}\NormalTok{ xmlPotential}\OperatorTok{.}\FunctionTok{getChild}\OperatorTok{(}\NormalTok{XMLTags}\OperatorTok{.}\FunctionTok{FUNCTIONS}\OperatorTok{.}\FunctionTok{toString}\OperatorTok{()).}\FunctionTok{getValue}\OperatorTok{();}
\BuiltInTok{List}\OperatorTok{\textless{}}\BuiltInTok{String}\OperatorTok{\textgreater{}}\NormalTok{ uncertainParametersList }\OperatorTok{=} \BuiltInTok{Pattern}\OperatorTok{.}\FunctionTok{compile}\OperatorTok{(}\StringTok{"}\SpecialCharTok{\textbackslash{}"}\StringTok{(.*?)}\SpecialCharTok{\textbackslash{}"}\StringTok{"}\OperatorTok{)}
                                              \OperatorTok{.}\FunctionTok{matcher}\OperatorTok{(}\NormalTok{functionsValue}\OperatorTok{)}
                                              \OperatorTok{.}\FunctionTok{results}\OperatorTok{()}
                                              \OperatorTok{.}\FunctionTok{map}\OperatorTok{(}\NormalTok{m }\OperatorTok{{-}\textgreater{}}\NormalTok{ m}\OperatorTok{.}\FunctionTok{group}\OperatorTok{(}\DecValTok{1}\OperatorTok{))}
                                              \OperatorTok{.}\FunctionTok{toList}\OperatorTok{();}
\NormalTok{VariableExpression}\OperatorTok{[]}\NormalTok{ functionValues }\OperatorTok{=} \KeywordTok{new}\NormalTok{ VariableExpression}\OperatorTok{[}\NormalTok{uncertainParametersList}\OperatorTok{.}\FunctionTok{size}\OperatorTok{()];}
\DataTypeTok{int}\NormalTok{ i }\OperatorTok{=} \DecValTok{0}\OperatorTok{;}
\ControlFlowTok{for} \OperatorTok{(}\BuiltInTok{String}\NormalTok{ uncertainParameter }\OperatorTok{:}\NormalTok{ uncertainParametersList}\OperatorTok{)} \OperatorTok{\{}
\NormalTok{    functionValues}\OperatorTok{[}\NormalTok{i}\OperatorTok{++]} \OperatorTok{=} \KeywordTok{new} \FunctionTok{VariableExpression}\OperatorTok{(}
\NormalTok{            Stream}\OperatorTok{.}\FunctionTok{concat}\OperatorTok{(}\NormalTok{vDistributionTable}\OperatorTok{.}\FunctionTok{stream}\OperatorTok{(),}\NormalTok{ variables}\OperatorTok{.}\FunctionTok{stream}\OperatorTok{()).}\FunctionTok{distinct}\OperatorTok{().}\FunctionTok{toList}\OperatorTok{(),}
\NormalTok{            uncertainParameter}\OperatorTok{);}
\OperatorTok{\}}
\end{Highlighting}
\end{Shaded}
\paragraph{Cómo arreglarlo}
Dos opciones, a decidir (ver \texttt{decision-de-equipo}):

\begin{itemize}
\tightlist
\item
  \textbf{Aceptar las dos grafías al leer.} En \texttt{PGMXReader\_\allowbreak{}1\_\allowbreak{}0.\allowbreak{}get\allowbreak{}Augmented\allowbreak{}Prob\allowbreak{}Table}, si el texto de una casilla es \texttt{Complement} (conviene decidir si sin distinguir mayúsculas), construir \texttt{Variable\allowbreak{}Expression.\allowbreak{}Common.\allowbreak{}COMPLEMENT}. El escritor seguiría escribiendo \texttt{\#}, de modo que el fichero se normaliza al guardarlo. Es el único sitio de lectura: \texttt{get\allowbreak{}Augmented\allowbreak{}Prob\allowbreak{}Table\allowbreak{}Potential} y la lectura de \texttt{Univariate\allowbreak{}Distr\allowbreak{}Potential} (líneas 155-160) pasan los dos por \texttt{get\allowbreak{}Augmented\allowbreak{}Prob\allowbreak{}Table}.
\item
  \textbf{Solo \texttt{\#}.} Corregir la red de pruebas (\texttt{Complement} → \texttt{\#}) y, para que el usuario entienda el fallo, que el lector o \texttt{Augmented\allowbreak{}Prob\allowbreak{}Table\allowbreak{}Inference} digan que una casilla no es una fórmula válida indicando el nodo y la casilla.
\end{itemize}

En ambos casos conviene añadir a las pruebas de integración una que proyecte todos los potenciales de \texttt{Dynamic-LIMID-every-potential.\allowbreak{}pgmx} (la misma que propone \hyperref[err:14]{error~14}), y en \texttt{io/\allowbreak{}src/\allowbreak{}test} una prueba que lea una tabla de fórmulas con el marcador elegido y compruebe el reparto.

\setcounter{subsection}{16}
\subsection[Quitar una etiqueta referenciada del árbol deja la red sin poder evaluarse]{Quitar la etiqueta de una rama que otra rama referencia deja una referencia rota, y al guardar y reabrir la red no se puede evaluar}\label{err:17}
\begin{ficha}
\fichakey{Estado}Presente
\fichakey{Módulo}gui
\fichakey{Paquete}org.openmarkov.gui.dialog.treeadd
\fichakey{Ficheros}\path{gui/src/main/java/org/openmarkov/gui/dialog/treeadd/TreeADDEditorPanel.java:1260-1264} \newline \path{core/src/main/java/org/openmarkov/core/model/network/potential/treeadd/TreeADDBranch.java:196-212} \newline \path{core/src/main/java/org/openmarkov/core/model/network/potential/treeadd/TreeADDBranch.java:368-372} \newline \path{core/src/main/java/org/openmarkov/core/model/network/potential/treeadd/TreeADDPotential.java:594-606} \newline \path{io/src/main/java/org/openmarkov/io/probmodel/writer/PGMXWriter_0_2.java:1230-1240} \newline \path{io/src/main/java/org/openmarkov/io/probmodel/reader/PGMXReader_0_2.java:1214-1224} \newline \path{core/src/main/java/org/openmarkov/core/model/network/potential/treeadd/TreeADDPotential.java:358}
\fichakey{Comprobado}Leyendo el código y ejecutándolo
\fichakey{Esfuerzo}Pequeño (horas)
\fichakey{Decisión de equipo}\textcolor{decisionrojo}{\textbf{Sí.}} Qué debe pasar al quitar una etiqueta en uso (prohibirlo, avisar, o convertir cada referencia en una copia del potencial)
\end{ficha}
\paragraph{Qué pasa}
En un potencial en forma de árbol (\texttt{Tree\allowbreak{}ADDPotential}, módulo \texttt{core}), cada rama (\texttt{Tree\allowbreak{}ADDBranch}) puede llevar una \textbf{etiqueta} («label»). Otra rama puede, en vez de tener su propio potencial, llevar una \textbf{referencia} a esa etiqueta: así las dos ramas comparten la misma tabla sin escribirla dos veces. En el editor de árboles (\texttt{Tree\allowbreak{}ADDEditor\allowbreak{}Panel}, módulo \texttt{gui}; diálogo de propiedades del nodo, pestaña del potencial con relación de tipo árbol) se llega con el botón derecho sobre una rama: «Set label», «Set reference», «Remove label», «Remove reference».

Si el usuario quita con «Remove label» una etiqueta que otra rama referencia, el editor no dice nada, la referencia queda rota, y la red guardada y reabierta muere con \texttt{Null\allowbreak{}Pointer\allowbreak{}Exception} al evaluarla. La red guardada queda inservible para cualquier inferencia sin que nada lo haya avisado.

Por qué pasa: «Remove label» (\texttt{remove\allowbreak{}Label}, líneas 1260-1264) solo pone la etiqueta a \texttt{null}. No recorre el árbol para ver si alguna rama la estaba referenciando. La rama que la referenciaba se queda con su campo \texttt{reference} diciendo el nombre antiguo y con su puntero interno \texttt{referenced\allowbreak{}Branch} apuntando a la rama que ya no está etiquetada.

Medido sobre \texttt{0\_\allowbreak{}miscellaneous/\allowbreak{}networks/\allowbreak{}id/\allowbreak{}ID-Monty-Hall-spanish.\allowbreak{}pgmx} (nodo «Puerta escogida 2»): se etiquetó la rama «no» como \texttt{the\allowbreak{}One}, se hizo que la rama «sí» la referenciara, se abrió el editor y se ejecutó \texttt{remove\allowbreak{}Label} sobre «no». Resultado:

\begin{enumerate}
\def\labelenumi{\arabic{enumi}.}
\tightlist
\item
  En el editor, el árbol ya no tiene ninguna etiqueta, pero la rama «sí» sigue siendo una referencia (\texttt{is\allowbreak{}Reference(\allowbreak{})\ ==\ true}), el árbol la dibuja con el texto \texttt{the\allowbreak{}One} y su menú sigue ofreciendo «Remove reference». Mientras dura la sesión, la rama sigue usando el potencial de «no» a través del puntero interno, así que la inferencia en memoria no falla.
\item
  Al aceptar el diálogo (\texttt{Set\allowbreak{}Potential\allowbreak{}Edit}, que no copia el árbol) y guardar la red en ProbModelXML (\texttt{.\allowbreak{}pgmx}), el escritor (\texttt{PGMXWriter\_\allowbreak{}0\_\allowbreak{}2}, líneas 1230-1240) escribe \texttt{\textless{}Reference\textgreater{}the\allowbreak{}One\textless{}/\allowbreak{}Reference\textgreater{}} en la rama «sí» y ninguna \texttt{\textless{}Label\textgreater{}} en todo el árbol.
\item
  Al reabrir el fichero, el lector (\texttt{PGMXReader\_\allowbreak{}0\_\allowbreak{}2}, líneas 1214-1224) crea la rama con la referencia \texttt{the\allowbreak{}One}; \texttt{Tree\allowbreak{}ADDPotential.\allowbreak{}update\allowbreak{}References} no encuentra ninguna rama con esa etiqueta y la rama se queda \textbf{sin potencial} (\texttt{get\allowbreak{}Potential(\allowbreak{})\ ==\ null}).
\item
  Evaluar esa red (\texttt{VEEvaluation}) termina en \texttt{Null\allowbreak{}Pointer\allowbreak{}Exception} en \texttt{Tree\allowbreak{}ADDPotential.\allowbreak{}table\allowbreak{}Project} (línea 358).
\end{enumerate}

Hay además una incoherencia en la copia: \texttt{Tree\allowbreak{}ADDBranch.\allowbreak{}copy(\allowbreak{})} (líneas 196-212) llama a \texttt{set\allowbreak{}Referenced\allowbreak{}Branch}, que toma la referencia de la etiqueta actual de la rama referenciada; como esa etiqueta ya es \texttt{null}, la copia pierde la referencia pero sigue apuntando a la rama del árbol \textbf{original}. Si en algún camino se copia el árbol antes de guardar, el resultado es otro distinto (se guardaría una copia del potencial en lugar de la referencia).
\paragraph{El código}
gui/src/main/java/org/openmarkov/gui/dialog/treeadd/TreeADDEditorPanel.java:1260-1264

\begin{Shaded}
\begin{Highlighting}[]
\KeywordTok{private} \DataTypeTok{void} \FunctionTok{removeLabel}\OperatorTok{(}\BuiltInTok{ActionEvent}\NormalTok{ ae}\OperatorTok{,}\NormalTok{ TreeADDBranch branch}\OperatorTok{,} \BuiltInTok{TreePath}\NormalTok{ path}\OperatorTok{)} \OperatorTok{\{}
\NormalTok{    branch}\OperatorTok{.}\FunctionTok{setLabel}\OperatorTok{(}\KeywordTok{null}\OperatorTok{);}
\NormalTok{    TreeADDModel model }\OperatorTok{=} \OperatorTok{(}\NormalTok{TreeADDModel}\OperatorTok{)}\NormalTok{ jTree}\OperatorTok{.}\FunctionTok{getModel}\OperatorTok{();}
\NormalTok{    model}\OperatorTok{.}\FunctionTok{notifyTreeStructureChanged}\OperatorTok{(}\NormalTok{path}\OperatorTok{);}
\OperatorTok{\}}
\end{Highlighting}
\end{Shaded}

core/src/main/java/org/openmarkov/core/model/network/potential/treeadd/TreeADDBranch.java:368-372

\begin{Shaded}
\begin{Highlighting}[]
\KeywordTok{public} \DataTypeTok{void} \FunctionTok{setReferencedBranch}\OperatorTok{(}\NormalTok{TreeADDBranch treeADDBranch}\OperatorTok{)} \OperatorTok{\{}
    \KeywordTok{this}\OperatorTok{.}\FunctionTok{reference} \OperatorTok{=}\NormalTok{ treeADDBranch}\OperatorTok{.}\FunctionTok{getLabel}\OperatorTok{();}
    \KeywordTok{this}\OperatorTok{.}\FunctionTok{referencedBranch} \OperatorTok{=}\NormalTok{ treeADDBranch}\OperatorTok{;}
    \KeywordTok{this}\OperatorTok{.}\FunctionTok{potential} \OperatorTok{=} \KeywordTok{null}\OperatorTok{;}
\OperatorTok{\}}
\end{Highlighting}
\end{Shaded}
\paragraph{Cómo arreglarlo}
Hace falta decidir el comportamiento (ver \texttt{decision-de-equipo}):

\begin{itemize}
\tightlist
\item
  \begin{enumerate}
  \def\labelenumi{(\alph{enumi})}
  \tightlist
  \item
    no ofrecer «Remove label», o rechazarlo con un mensaje, mientras alguna rama referencie la etiqueta;
  \end{enumerate}
\item
  \begin{enumerate}
  \def\labelenumi{(\alph{enumi})}
  \setcounter{enumi}{1}
  \tightlist
  \item
    quitarla y convertir cada rama que la referenciaba en una rama con su propia copia del potencial (lo mismo que hace hoy «Remove reference», líneas 1289-1295);
  \end{enumerate}
\item
  \begin{enumerate}
  \def\labelenumi{(\alph{enumi})}
  \setcounter{enumi}{2}
  \tightlist
  \item
    quitarla y quitar también las referencias, avisando.
  \end{enumerate}
\end{itemize}

Sitios donde la regla tiene que cumplirse: \texttt{Tree\allowbreak{}ADDEditor\allowbreak{}Panel.\allowbreak{}remove\allowbreak{}Label}; \texttt{Tree\allowbreak{}ADDBranch.\allowbreak{}copy(\allowbreak{})} y \texttt{set\allowbreak{}Referenced\allowbreak{}Branch}, que no deberían depender de que la etiqueta siga puesta; y, como red de seguridad, el lector de \texttt{.\allowbreak{}pgmx}, que hoy acepta en silencio una referencia sin etiqueta y deja una rama sin potencial en lugar de dar un error de lectura claro.

Prueba de regresión: árbol con una rama etiquetada y otra que la referencia; quitar la etiqueta desde el panel; guardar y reabrir; comprobar que la red se evalúa (o que la operación se rechazó) y que ninguna rama queda sin potencial.

\setcounter{subsection}{17}
\subsection[Añadir a una hoja del árbol la variable del nodo deja una red que no abre]{El editor de árboles ofrece añadir a una hoja la propia variable del nodo; aceptarlo la escribe dos veces y la red guardada ya no se puede abrir}\label{err:18}
\begin{ficha}
\fichakey{Estado}Presente
\fichakey{Módulo}core
\fichakey{Paquete}org.openmarkov.core.model.network.potential
\fichakey{Ficheros}\path{core/src/main/java/org/openmarkov/core/model/network/potential/treeadd/TreeADDBranch.java:238-250} \newline \path{core/src/main/java/org/openmarkov/core/model/network/potential/Potential.java:700-717} \newline \path{core/src/main/java/org/openmarkov/core/model/network/potential/TablePotential.java:628-642} \newline \path{gui/src/main/java/org/openmarkov/gui/dialog/treeadd/TreeADDEditorPanel.java:288-323} \newline \path{core/src/main/java/org/openmarkov/core/action/base/linkEdits/AddLinkEdit.java:207-229}
\fichakey{Comprobado}Leyendo el código y ejecutándolo
\fichakey{Esfuerzo}Medio (días)
\fichakey{Decisión de equipo}\textcolor{decisionrojo}{\textbf{Sí.}} Qué contrato tiene \texttt{add\allowbreak{}Variable} con una variable que ya está (hoy depende de la clase), y cómo se distingue de él el caso legítimo de los bucles sobre sí mismo de las redes de simulación de eventos discretos
\fichakey{Relacionados}\hyperref[err:159]{error~159}, \hyperref[nota:164]{nota~164}, \hyperref[nota:182]{nota~182}
\end{ficha}
\paragraph{Qué pasa}
En el editor de árboles, el menú de una hoja (\texttt{Tree\allowbreak{}ADDEditor\allowbreak{}Panel.\allowbreak{}set\allowbreak{}Contextual\allowbreak{}Menu\allowbreak{}Potential}) y el diálogo «Add variables to potential» ofrecen siempre la propia variable del nodo, en cualquier tipo de red. Si el usuario la elige, la hoja queda como P(X \textbar{} \ldots, X), con la variable escrita dos veces. La comprobación de restricciones no dice nada y la evaluación en memoria sigue dando un número, pero la red guardada ya no se puede volver a abrir. Como se describe en \hyperref[err:159]{error~159}, la variable condicionada es además la única elección de ese diálogo que hace algo.

Los bucles sobre sí mismo solo son legítimos en las redes de simulación de eventos discretos, donde un nodo de azar o de suceso puede ser padre de sí mismo.

Por qué pasa: se combinan dos piezas.

\textbf{1. Qué se ofrece.} \texttt{Tree\allowbreak{}ADDBranch.\allowbreak{}get\allowbreak{}Addable\allowbreak{}Variables} (módulo \texttt{core}, líneas 238-250) calcula las variables que se pueden añadir a la hoja de una rama. Parte de las variables del árbol, quita las que ya tiene el potencial de la hoja y, si «hay bucle sobre sí mismo», vuelve a añadir la variable condicionada. Pero la condición \texttt{self\allowbreak{}Loop\ =\ addable\allowbreak{}Variables.\allowbreak{}contains(\allowbreak{}potential.\allowbreak{}get\allowbreak{}Conditioned\allowbreak{}Variable(\allowbreak{}))} pregunta si la lista de variables del árbol contiene la variable condicionada, y eso es verdad en \textbf{todo} árbol de probabilidad condicionada (la condicionada es siempre la primera).

\textbf{2. Qué hace al aceptarlo.} \texttt{Potential.\allowbreak{}add\allowbreak{}Variable} (líneas 700-717) añade la variable si no está; si ya está y es la primera, la \textbf{añade otra vez} (el propio código lo marca \texttt{FIXME} y cita los bucles de las redes de simulación de eventos discretos). \texttt{Table\allowbreak{}Potential.\allowbreak{}add\allowbreak{}Variable} (líneas 628-642) ni siquiera mira si ya está: siempre la añade.

Medido sobre \texttt{0\_\allowbreak{}miscellaneous/\allowbreak{}networks/\allowbreak{}id/\allowbreak{}ID-Monty-Hall-spanish.\allowbreak{}pgmx} (nodo «Puerta escogida 2», hoja «no», una \texttt{Table\allowbreak{}Potential}):

\begin{itemize}
\tightlist
\item
  \texttt{add\allowbreak{}Variable} de la variable condicionada sobre un \texttt{Uniform\allowbreak{}Potential} y sobre una \texttt{Table\allowbreak{}Potential} de una sola variable devuelve en ambos casos \texttt{{[}Puerta\ escogida\ 2,\allowbreak{}\ Puerta\ escogida\ 2{]}}.
\item
  Poniendo en la hoja la tabla duplicada \texttt{{[}Puerta\ escogida\ 2,\allowbreak{}\ Puerta\ escogida\ 1,\allowbreak{}\ Puerta\ escogida\ 2{]}} (27 casillas), la comprobación de restricciones de la red (\texttt{check\allowbreak{}Prob\allowbreak{}Net}) no dice nada y la evaluación en memoria (\texttt{VEEvaluation}) sigue dando un número.
\item
  Guardada la red con \texttt{PGMXWriter\_\allowbreak{}1\_\allowbreak{}0}, el fichero lleva la variable repetida en \texttt{\textless{}Variables\textgreater{}}; al \textbf{reabrirlo}, \texttt{PGMXReader} falla con \texttt{Invalid\allowbreak{}Argument\allowbreak{}Exception:\ Argument\ table\ (\allowbreak{}27)\ is\ invalid\ because\ its\ length\ must\ be\ 9,\allowbreak{}\ the\ product\ of\ the\ numbers\ of\ states\ of\ the\ variables.\allowbreak{}} El usuario pierde la posibilidad de abrir la red que acaba de guardar.
\end{itemize}
\paragraph{El código}
core/src/main/java/org/openmarkov/core/model/network/potential/treeadd/TreeADDBranch.java:238-250

\begin{Shaded}
\begin{Highlighting}[]
\KeywordTok{public} \BuiltInTok{List}\OperatorTok{\textless{}}\NormalTok{Variable}\OperatorTok{\textgreater{}} \FunctionTok{getAddableVariables}\OperatorTok{()} \OperatorTok{\{}
    \BuiltInTok{List}\OperatorTok{\textless{}}\NormalTok{Variable}\OperatorTok{\textgreater{}}\NormalTok{ addableVariables }\OperatorTok{=} \KeywordTok{new} \BuiltInTok{ArrayList}\OperatorTok{\textless{}\textgreater{}(}\NormalTok{parentVariables}\OperatorTok{);}
    \DataTypeTok{boolean}\NormalTok{ selfLoop }\OperatorTok{=}\NormalTok{ addableVariables}\OperatorTok{.}\FunctionTok{contains}\OperatorTok{(}\NormalTok{potential}\OperatorTok{.}\FunctionTok{getConditionedVariable}\OperatorTok{());}
\NormalTok{    addableVariables}\OperatorTok{.}\FunctionTok{removeAll}\OperatorTok{(}\NormalTok{potential}\OperatorTok{.}\FunctionTok{getVariables}\OperatorTok{());}
    \ControlFlowTok{if} \OperatorTok{(}\NormalTok{selfLoop}\OperatorTok{)} \OperatorTok{\{}
\NormalTok{        addableVariables}\OperatorTok{.}\FunctionTok{add}\OperatorTok{(}\NormalTok{potential}\OperatorTok{.}\FunctionTok{getConditionedVariable}\OperatorTok{());}
    \OperatorTok{\}}
    \KeywordTok{...}
\end{Highlighting}
\end{Shaded}

core/src/main/java/org/openmarkov/core/model/network/potential/Potential.java:703-716

\begin{Shaded}
\begin{Highlighting}[]
\ControlFlowTok{if} \OperatorTok{(!}\NormalTok{variables}\OperatorTok{.}\FunctionTok{contains}\OperatorTok{(}\NormalTok{variable}\OperatorTok{))} \OperatorTok{\{}
    \KeywordTok{...}
\OperatorTok{\}} \ControlFlowTok{else} \OperatorTok{\{}
    \CommentTok{//DESnets {-} 18/03/2023 {-} due to self{-}loop, the conditioned variable can be repeated {-} }\AlertTok{FIXME}
    \ControlFlowTok{if} \OperatorTok{(}\NormalTok{variable}\OperatorTok{.}\FunctionTok{equals}\OperatorTok{(}\NormalTok{variables}\OperatorTok{.}\FunctionTok{get}\OperatorTok{(}\DecValTok{0}\OperatorTok{)))\{}
        \BuiltInTok{List}\OperatorTok{\textless{}}\NormalTok{Variable}\OperatorTok{\textgreater{}}\NormalTok{ newVariables }\OperatorTok{=} \KeywordTok{new} \BuiltInTok{ArrayList}\OperatorTok{\textless{}}\NormalTok{Variable}\OperatorTok{\textgreater{}(}\NormalTok{variables}\OperatorTok{);}
\NormalTok{        newVariables}\OperatorTok{.}\FunctionTok{add}\OperatorTok{(}\NormalTok{variable}\OperatorTok{);}
\NormalTok{        newPotential }\OperatorTok{=} \KeywordTok{new} \FunctionTok{UniformPotential}\OperatorTok{(}\NormalTok{newVariables}\OperatorTok{,}\NormalTok{ role}\OperatorTok{);}
    \OperatorTok{\}} \ControlFlowTok{else}\OperatorTok{\{}
\NormalTok{        newPotential }\OperatorTok{=} \KeywordTok{this}\OperatorTok{;}
    \OperatorTok{\}}
\OperatorTok{\}}
\end{Highlighting}
\end{Shaded}
\paragraph{Cómo arreglarlo}
La detección del bucle sobre sí mismo tiene que mirar si de verdad lo hay (por ejemplo, que el nodo sea padre de sí mismo en la red, o que la variable aparezca dos veces entre las del árbol), no si la condicionada está en la lista. Con eso, fuera de las redes de simulación de eventos discretos la variable condicionada deja de ofrecerse.

El contrato de \texttt{add\allowbreak{}Variable} con una variable que ya está debe ser uno solo para todas las clases: hoy \texttt{Potential} la duplica solo si es la primera, \texttt{Table\allowbreak{}Potential} siempre, y \texttt{Tree\allowbreak{}ADDPotential} y \texttt{Sum\allowbreak{}Potential} mutan el receptor. Hay que decidirlo (ver \texttt{decision-de-equipo}); los contratos sin escribir de las operaciones sobre potenciales están recogidos también en \hyperref[nota:164]{nota~164}, donde no hay nada que corregir. Si se decide que duplicar solo es válido con bucle sobre sí mismo, la comprobación debe estar en \texttt{add\allowbreak{}Variable} (o en quien la llama) y no depender de cada llamador.

Sitios que llaman a \texttt{add\allowbreak{}Variable} y deben respetar la misma regla: \texttt{Tree\allowbreak{}ADDEditor\allowbreak{}Panel.\allowbreak{}add\allowbreak{}Variables\allowbreak{}To\allowbreak{}Potential} (línea 785), \texttt{Add\allowbreak{}Link\allowbreak{}Edit.\allowbreak{}do\allowbreak{}Edit} (líneas 207-229, que ya trata aparte el bucle sobre sí mismo de las redes de simulación de eventos discretos), \texttt{Invert\allowbreak{}Link\allowbreak{}Edit} (línea 124), \texttt{Tree\allowbreak{}ADDPotential.\allowbreak{}add\allowbreak{}Variable} (línea 331), \texttt{Discretized\allowbreak{}Cauchy\allowbreak{}Potential.\allowbreak{}add\allowbreak{}Variable} y \texttt{MIDTemporal\allowbreak{}Evolution} (línea 584).

Además, el lector de \texttt{.\allowbreak{}pgmx} debería dar un mensaje claro ante una variable repetida en un potencial (hoy da el error de tamaño de tabla), y la comprobación de red podría detectarla antes de guardar.

Prueba: en un diagrama de influencia, el diálogo «Add variables to potential» de una hoja no ofrece la variable del nodo; en una red de simulación de eventos discretos con bucle sobre sí mismo sí la ofrece; y guardar y reabrir la red tras añadir variables funciona.

\setcounter{subsection}{18}
\subsection[Tras cancelar el potencial de un suceso, la simulación usa el primer padre]{Tras cancelar la ventana de potencial de un suceso con dos sucesos padres, la simulación usa siempre los parámetros del primero}\label{err:19}
\begin{ficha}
\fichakey{Estado}Presente
\fichakey{Módulo}core
\fichakey{Paquete}org.openmarkov.core.model.network.potential
\fichakey{Ficheros}\path{core/src/main/java/org/openmarkov/core/model/network/potential/TableWithEvents.java:112-124} \newline \path{core/src/main/java/org/openmarkov/core/model/network/potential/TableWithEvents.java:150-162} \newline \path{core/src/main/java/org/openmarkov/core/model/network/potential/TableWithEvents.java:445-454} \newline \path{core/src/main/java/org/openmarkov/core/model/network/potential/DistributionTablePotential.java:104-112} \newline \path{core/src/main/java/org/openmarkov/core/model/network/potential/TablePotential.java:238-268} \newline \path{gui/src/main/java/org/openmarkov/gui/dialog/node/PotentialEditPanel.java:148-153, 570-579, 709-715} \newline \path{gui/src/main/java/org/openmarkov/gui/action/PasteEdit.java:116-118}
\fichakey{Comprobado}Leyendo el código y ejecutándolo
\fichakey{Esfuerzo}Pequeño (horas)
\fichakey{Decisión de equipo}No
\fichakey{Tapado por}\hyperref[err:20]{error~20}: en parte: la copia de \texttt{Transition\allowbreak{}Table\allowbreak{}Potential} por su constructor de copia no se ve porque \texttt{copy(\allowbreak{})} devuelve otra clase.
\fichakey{Relacionados}\hyperref[err:20]{error~20}, \hyperref[err:21]{error~21}
\end{ficha}
\paragraph{Qué pasa}
En una red de simulación de eventos discretos, un nodo de suceso cuyo tiempo hasta ocurrir depende de qué suceso lo desencadenó lleva un \texttt{Distribution\allowbreak{}Table\allowbreak{}Potential}: una distribución (exponencial, exacta\ldots) cuyos parámetros se leen de una tabla interna, \texttt{Table\allowbreak{}With\allowbreak{}Events}.

Si ese nodo tiene dos o más sucesos padres y el potencial se copia, la copia lee siempre los parámetros del primer suceso, cualquiera que sea el que haya ocurrido. La simulación no falla ni avisa: da tiempos hasta el suceso con los parámetros de otra columna y, por tanto, costes y efectividades distintos de los que el modelo dice. Con un solo suceso padre no se nota (la única columna es la buena).

Caminos de llegada:

\begin{enumerate}
\def\labelenumi{\arabic{enumi}.}
\tightlist
\item
  Ventana de edición del potencial de un nodo (doble clic en el nodo, pestaña de potencial, o «Editar potencial»). Al abrirse, \texttt{Potential\allowbreak{}Edit\allowbreak{}Panel} guarda \texttt{this.\allowbreak{}node.\allowbreak{}get\allowbreak{}Potential(\allowbreak{}).\allowbreak{}deep\allowbreak{}Copy(\allowbreak{}.\allowbreak{}.\allowbreak{}.\allowbreak{})} como potencial original; al pulsar Cancelar, \texttt{uncommit\allowbreak{}Changes} → \texttt{remove\allowbreak{}Potential\allowbreak{}On\allowbreak{}Close} vuelve a poner en el nodo \textbf{esa copia}, no el objeto que tenía. \texttt{Potential.\allowbreak{}deep\allowbreak{}Copy} construye la copia con el constructor de copia de la clase (por reflexión) y \texttt{Distribution\allowbreak{}Table\allowbreak{}Potential} no redefine \texttt{deep\allowbreak{}Copy}, así que la copia defectuosa se queda en el nodo. Si se acepta con cambios, la edición \texttt{Set\allowbreak{}Potential\allowbreak{}Edit} guarda esa misma copia como estado anterior y la restituye al deshacer. La ventana no se ha podido construir sin pantalla, así que este paso está comprobado leyendo \texttt{Potential\allowbreak{}Edit\allowbreak{}Panel}; la consecuencia sobre el potencial se midió ejecutando \texttt{deep\allowbreak{}Copy}.
\item
  Copiar y pegar un nodo de suceso con dos o más sucesos padres (\texttt{Paste\allowbreak{}Edit} usa \texttt{copy(\allowbreak{})}).
\end{enumerate}

Por qué pasa: la tabla interna no tiene una columna por cada suceso padre como variable separada. \texttt{set\allowbreak{}Event\allowbreak{}As\allowbreak{}States} fabrica una variable auxiliar llamada \texttt{Events} cuyos estados son los nombres de los sucesos padres, y \texttt{convert} traduce «ha ocurrido el suceso X» a un hallazgo sobre esa variable \texttt{Events}.

El constructor de copia de \texttt{Table\allowbreak{}With\allowbreak{}Events} hace dos cosas que no encajan. Primero copia la tabla interna (\texttt{new\ Table\allowbreak{}Potential(\allowbreak{}potential.\allowbreak{}get\allowbreak{}Table\allowbreak{}Potential(\allowbreak{}))}), que conserva las variables del original y, entre ellas, \textbf{la variable \texttt{Events} del original}. Después llama a \texttt{set\allowbreak{}Event\allowbreak{}As\allowbreak{}States}, que construye una variable \texttt{Events} \textbf{nueva}, y es esa la que usa \texttt{convert}. Como \texttt{Variable} compara por identidad de objeto, la tabla busca en la evidencia una variable que no está. Y \texttt{Table\allowbreak{}Potential.\allowbreak{}get\allowbreak{}Value(\allowbreak{}Evidence\allowbreak{}Case)} no protesta cuando falta una de sus variables: la salta, lo que equivale a leer su estado 0. \texttt{Distribution\allowbreak{}Table\allowbreak{}Potential(\allowbreak{}Distribution\allowbreak{}Table\allowbreak{}Potential)} copia su tabla interna con ese mismo \texttt{copy(\allowbreak{})}, así que hereda el defecto.

Medido:

\begin{itemize}
\tightlist
\item
  Un \texttt{Distribution\allowbreak{}Table\allowbreak{}Potential} exponencial (parametrización \texttt{Lambda}) con dos sucesos padres A y B y tasas 0,1 y 0,3. Ocurrido B, el original devuelve 2,3104906018664844 (la mediana con tasa 0,3, con número aleatorio 0,5) y su \texttt{copy(\allowbreak{})} devuelve 6,931471805599453, que es la de tasa 0,1. En la copia, la tabla interna contiene la variable \texttt{Events} del original y no la suya.
\item
  Camino real, en una red de sucesos de prueba en la que el suceso terminal \texttt{Death} depende del suceso \texttt{Onset} y del nodo \texttt{Ill}: se añade a \texttt{Death} un segundo suceso padre (\texttt{Checkup}), con parámetro «exacto» 3,0 si lo desencadena \texttt{Onset} y 1,0 si lo desencadena \texttt{Checkup} (con \texttt{Ill}=yes). Antes: tras \texttt{Checkup}, el tiempo hasta la muerte es 1,0. Tras sustituir el potencial por su \texttt{deep\allowbreak{}Copy(\allowbreak{}red)} ---lo que hace la ventana al cancelar---, el mismo caso devuelve 3,0.
\end{itemize}

La copia de \texttt{Transition\allowbreak{}Table\allowbreak{}Potential} por su constructor de copia tiene el mismo desajuste (su muestreo lanza \texttt{Null\allowbreak{}Pointer\allowbreak{}Exception}); hoy no se ve porque \texttt{Table\allowbreak{}With\allowbreak{}Events.\allowbreak{}deep\allowbreak{}Copy} rehace las variables internas y porque \texttt{copy(\allowbreak{})} devuelve otra clase, lo que se describe en \hyperref[err:20]{error~20}.
\paragraph{El código}
core/src/main/java/org/openmarkov/core/model/network/potential/TableWithEvents.java:112-124

\begin{Shaded}
\begin{Highlighting}[]
    \KeywordTok{public} \FunctionTok{TableWithEvents}\OperatorTok{(}\NormalTok{TableWithEvents potential}\OperatorTok{)} \OperatorTok{\{}
        \KeywordTok{super}\OperatorTok{(}\NormalTok{potential}\OperatorTok{);}
        \KeywordTok{this}\OperatorTok{.}\FunctionTok{setTablePotential}\OperatorTok{(}\KeywordTok{new} \FunctionTok{TablePotential}\OperatorTok{(}\NormalTok{potential}\OperatorTok{.}\FunctionTok{getTablePotential}\OperatorTok{()));}
        \KeywordTok{this}\OperatorTok{.}\FunctionTok{useTableWithFunctions} \OperatorTok{=}\NormalTok{ potential}\OperatorTok{.}\FunctionTok{useTableWithFunctions}\OperatorTok{;}
        \ControlFlowTok{if} \OperatorTok{(}\NormalTok{potential}\OperatorTok{.}\FunctionTok{tableWithFunctions} \OperatorTok{!=} \KeywordTok{null}\OperatorTok{)} \OperatorTok{\{}
            \KeywordTok{this}\OperatorTok{.}\FunctionTok{tableWithFunctions} \OperatorTok{=} \OperatorTok{(}\NormalTok{TableWithFunctions}\OperatorTok{)}\NormalTok{ potential}\OperatorTok{.}\FunctionTok{tableWithFunctions}\OperatorTok{.}\FunctionTok{copy}\OperatorTok{();}
        \OperatorTok{\}}
        \KeywordTok{this}\OperatorTok{.}\FunctionTok{hasImpossibleConfigurations} \OperatorTok{=}\NormalTok{ potential}\OperatorTok{.}\FunctionTok{hasImpossibleConfigurations}\OperatorTok{;}
        \KeywordTok{this}\OperatorTok{.}\FunctionTok{impossibleConfigurations} \OperatorTok{=}\NormalTok{ potential}\OperatorTok{.}\FunctionTok{impossibleConfigurations} \OperatorTok{==} \KeywordTok{null}
                \OperatorTok{?} \KeywordTok{new} \BuiltInTok{ArrayList}\OperatorTok{\textless{}\textgreater{}()} \OperatorTok{:} \KeywordTok{new} \BuiltInTok{ArrayList}\OperatorTok{\textless{}\textgreater{}(}\NormalTok{potential}\OperatorTok{.}\FunctionTok{impossibleConfigurations}\OperatorTok{);}
        \FunctionTok{setEventAsStates}\OperatorTok{(}\NormalTok{variables}\OperatorTok{.}\FunctionTok{subList}\OperatorTok{(}\DecValTok{1}\OperatorTok{,}\NormalTok{ variables}\OperatorTok{.}\FunctionTok{size}\OperatorTok{()));}
        \FunctionTok{rebuildTableVariables}\OperatorTok{();}
    \OperatorTok{\}}
\end{Highlighting}
\end{Shaded}

core/src/main/java/org/openmarkov/core/model/network/potential/TablePotential.java:238-248

\begin{Shaded}
\begin{Highlighting}[]
    \KeywordTok{public} \DataTypeTok{double} \FunctionTok{getValue}\OperatorTok{(}\BuiltInTok{List}\OperatorTok{\textless{}}\NormalTok{Variable}\OperatorTok{\textgreater{}}\NormalTok{ variables}\OperatorTok{,} \DataTypeTok{int}\OperatorTok{[]}\NormalTok{ statesIndices}\OperatorTok{)} \OperatorTok{\{}
        \DataTypeTok{int}\NormalTok{ position }\OperatorTok{=} \DecValTok{0}\OperatorTok{;}
        \ControlFlowTok{for} \OperatorTok{(}\DataTypeTok{int}\NormalTok{ i }\OperatorTok{=} \DecValTok{0}\OperatorTok{;}\NormalTok{ i }\OperatorTok{\textless{}}\NormalTok{ variables}\OperatorTok{.}\FunctionTok{size}\OperatorTok{();}\NormalTok{ i}\OperatorTok{++)} \OperatorTok{\{}
\NormalTok{            Variable variable }\OperatorTok{=}\NormalTok{ variables}\OperatorTok{.}\FunctionTok{get}\OperatorTok{(}\NormalTok{i}\OperatorTok{);}
            \DataTypeTok{int}\NormalTok{ indexVariable }\OperatorTok{=} \KeywordTok{this}\OperatorTok{.}\FunctionTok{variables}\OperatorTok{.}\FunctionTok{indexOf}\OperatorTok{(}\NormalTok{variable}\OperatorTok{);}
            \ControlFlowTok{if} \OperatorTok{(}\NormalTok{indexVariable }\OperatorTok{!=} \OperatorTok{{-}}\DecValTok{1}\OperatorTok{)} \OperatorTok{\{}
\NormalTok{                position }\OperatorTok{+=}\NormalTok{ offsets}\OperatorTok{[}\NormalTok{indexVariable}\OperatorTok{]} \OperatorTok{*}\NormalTok{ statesIndices}\OperatorTok{[}\NormalTok{i}\OperatorTok{];}
            \OperatorTok{\}}
        \OperatorTok{\}}
        \ControlFlowTok{return}\NormalTok{ values}\OperatorTok{[}\NormalTok{position}\OperatorTok{];}
    \OperatorTok{\}}
\end{Highlighting}
\end{Shaded}

gui/src/main/java/org/openmarkov/gui/dialog/node/PotentialEditPanel.java:147-152 y 697-701

\begin{Shaded}
\begin{Highlighting}[]
\NormalTok{        Potential lastPotential }\OperatorTok{=} \KeywordTok{this}\OperatorTok{.}\FunctionTok{node}\OperatorTok{.}\FunctionTok{getPotential}\OperatorTok{();}
        \ControlFlowTok{if} \OperatorTok{(}\NormalTok{lastPotential }\OperatorTok{!=} \KeywordTok{null}\OperatorTok{)} \OperatorTok{\{}
            \KeywordTok{this}\OperatorTok{.}\FunctionTok{originalPotential} \OperatorTok{=} \KeywordTok{this}\OperatorTok{.}\FunctionTok{node}\OperatorTok{.}\FunctionTok{getPotential}\OperatorTok{().}\FunctionTok{deepCopy}\OperatorTok{(}\KeywordTok{this}\OperatorTok{.}\FunctionTok{node}\OperatorTok{.}\FunctionTok{getProbNet}\OperatorTok{());}
        \KeywordTok{...}
    \KeywordTok{protected} \DataTypeTok{void} \FunctionTok{removePotentialOnClose}\OperatorTok{(}\AttributeTok{@Nullable}\NormalTok{ Potential originalPotential}\OperatorTok{)} \OperatorTok{\{}
        \ControlFlowTok{if} \OperatorTok{(}\NormalTok{originalPotential }\OperatorTok{!=} \KeywordTok{null}\OperatorTok{)} \OperatorTok{\{}
\NormalTok{            LinkRestrictionPotentialOperations}\OperatorTok{.}\FunctionTok{setPotentialWithRestrictions}\OperatorTok{(}\KeywordTok{this}\OperatorTok{.}\FunctionTok{node}\OperatorTok{,}\NormalTok{ originalPotential}\OperatorTok{);}
        \OperatorTok{\}}
    \OperatorTok{\}}
\end{Highlighting}
\end{Shaded}
\paragraph{Cómo arreglarlo}
En el constructor de copia de \texttt{Table\allowbreak{}With\allowbreak{}Events}, después de \texttt{set\allowbreak{}Event\allowbreak{}As\allowbreak{}States} y \texttt{rebuild\allowbreak{}Table\allowbreak{}Variables}, poner las variables de la tabla interna (y de \texttt{table\allowbreak{}With\allowbreak{}Functions} si existe) sobre \texttt{table\allowbreak{}Variables}, igual que ya hace \texttt{Table\allowbreak{}With\allowbreak{}Events.\allowbreak{}deep\allowbreak{}Copy} (líneas 445-454). Así la copia y su variable \texttt{Events} vuelven a hablar de lo mismo, y \texttt{deep\allowbreak{}Copy} puede dejar de repetirlo.

Sitios donde la misma regla («las variables de la tabla interna son las derivadas del potencial que la contiene») tiene que cumplirse:

\begin{itemize}
\tightlist
\item
  \texttt{Table\allowbreak{}With\allowbreak{}Events(\allowbreak{}Table\allowbreak{}With\allowbreak{}Events)} (copia) y \texttt{Table\allowbreak{}With\allowbreak{}Events.\allowbreak{}deep\allowbreak{}Copy}.
\item
  \texttt{Distribution\allowbreak{}Table\allowbreak{}Potential(\allowbreak{}Distribution\allowbreak{}Table\allowbreak{}Potential)}: comparte \texttt{distribution\allowbreak{}Variable} y copia la tabla interna con \texttt{copy(\allowbreak{})}; y no redefine \texttt{deep\allowbreak{}Copy}, de modo que una copia profunda hacia \textbf{otra} red deja la tabla interna con las variables de la red de origen (hoy el único \texttt{deep\allowbreak{}Copy} alcanzable para este potencial es el de la ventana, dentro de la misma red).
\item
  \texttt{Table\allowbreak{}With\allowbreak{}Events.\allowbreak{}set\allowbreak{}Variables} (líneas 257-260) pone en la tabla interna las variables exteriores, que no incluyen \texttt{Events}: cualquier llamador que lo use deja la tabla desajustada de otra forma.
\item
  \texttt{Table\allowbreak{}With\allowbreak{}Events.\allowbreak{}replace\allowbreak{}Variable} no está redefinido (lo usa \texttt{Paste\allowbreak{}Edit}); lo que eso provoca al pegar se describe en \hyperref[err:21]{error~21}.
\end{itemize}

Hay otra opción, que exigiría una decisión del equipo: que \texttt{Table\allowbreak{}Potential.\allowbreak{}get\allowbreak{}Value(\allowbreak{}Evidence\allowbreak{}Case)} falle cuando falta un hallazgo de una variable de la tabla, en vez de leer el estado 0. Haría ruidoso este defecto y otros parecidos, pero cambia un método que usa todo el proyecto.

Pruebas de regresión: para \texttt{Distribution\allowbreak{}Table\allowbreak{}Potential} y \texttt{Transition\allowbreak{}Table\allowbreak{}Potential} con dos sucesos padres, \texttt{copy(\allowbreak{})} y \texttt{deep\allowbreak{}Copy(\allowbreak{}misma\ red)} muestrean lo mismo que el original para cada suceso; lo mismo con un padre numérico (rama \texttt{table\allowbreak{}With\allowbreak{}Functions}).

\setcounter{subsection}{19}
\subsection[Un nodo con tabla de transiciones copiado ya no sabe a qué estado pasa]{Copiar un nodo de azar con tabla de transiciones deja un nodo que ya no sabe a qué estado pasa}\label{err:20}
\begin{ficha}
\fichakey{Estado}Presente
\fichakey{Módulo}core
\fichakey{Paquete}org.openmarkov.core.model.network.potential
\fichakey{Ficheros}\path{core/src/main/java/org/openmarkov/core/model/network/potential/TransitionTablePotential.java:27-80} \newline \path{core/src/main/java/org/openmarkov/core/model/network/potential/TableWithEvents.java:430-433} \newline \path{core/src/main/java/org/openmarkov/core/model/network/potential/TableWithEvents.java:486-492} \newline \path{gui/src/main/java/org/openmarkov/gui/action/PasteEdit.java:116-118}
\fichakey{Comprobado}Leyendo el código y ejecutándolo
\fichakey{Esfuerzo}Pequeño (horas)
\fichakey{Decisión de equipo}No
\fichakey{Corregir antes}\hyperref[err:19]{error~19}: a la vez, como mínimo: sin él, la tabla interna del nodo copiado sigue apuntando a la variable «Events» del original.
\fichakey{Relacionados}\hyperref[err:19]{error~19}, \hyperref[err:21]{error~21}, \hyperref[nota:184]{nota~184}
\end{ficha}
\paragraph{Qué pasa}
En una red de simulación de eventos discretos (en OpenMarkov, redes de tipo \texttt{DESNet}), un nodo de azar que cambia de estado cuando ocurre un suceso lleva un potencial \texttt{Transition\allowbreak{}Table\allowbreak{}Potential}: una tabla que dice, para cada suceso que puede ocurrir, con qué probabilidad pasa el nodo a cada estado. Esa clase hereda de \texttt{Table\allowbreak{}With\allowbreak{}Events} y lo único que añade es saber muestrearse (\texttt{sample\allowbreak{}Conditioned\allowbreak{}Variable}).

Lo que ve el usuario: al copiar y pegar un nodo así, el nodo pegado se queda en su primer estado para siempre durante la simulación, sin aviso.

Camino de llegada: el único llamador de producción de \texttt{copy(\allowbreak{})} sobre estos potenciales es pegar nodos. \texttt{Paste\allowbreak{}Edit} (menú Edición, Copiar y Pegar en el editor de redes) copia cada potencial con \texttt{original\allowbreak{}Potential.\allowbreak{}copy(\allowbreak{})}. En la simulación, el \texttt{Na\allowbreak{}N} («no es un número») que devuelve la copia acaba convertido a entero (\texttt{(\allowbreak{}int)\ Na\allowbreak{}N} vale 0 en Java) al leer el estado del nodo (\texttt{Chance\allowbreak{}Evaluation.\allowbreak{}compute\allowbreak{}Chance\allowbreak{}Value}, \texttt{DESLog\allowbreak{}Text\allowbreak{}Writer.\allowbreak{}log\allowbreak{}Chance\allowbreak{}Change}, \texttt{Configuration.\allowbreak{}add\allowbreak{}Finding}). Esta última parte, lo que hace la simulación con el \texttt{Na\allowbreak{}N}, está comprobada leyendo el código, no ejecutando la simulación con la red pegada.

El otro camino de copia, la copia profunda \texttt{Potential.\allowbreak{}deep\allowbreak{}Copy} (que usa la ventana de edición de potenciales para poder cancelar), no pasa por \texttt{copy(\allowbreak{})}: invoca por reflexión el constructor de copia de la clase exacta, y \texttt{Table\allowbreak{}With\allowbreak{}Events.\allowbreak{}deep\allowbreak{}Copy} rehace después las variables internas, así que para esta clase sí funciona.

Por qué pasa: \texttt{Transition\allowbreak{}Table\allowbreak{}Potential} no redefine \texttt{copy(\allowbreak{})}, así que se ejecuta el de \texttt{Table\allowbreak{}With\allowbreak{}Events}, que construye un objeto de la clase madre: la copia de una tabla de transiciones es una \texttt{Table\allowbreak{}With\allowbreak{}Events} a secas. Y el muestreo de \texttt{Table\allowbreak{}With\allowbreak{}Events} no está implementado: devuelve siempre \texttt{Double.\allowbreak{}Na\allowbreak{}N}. No hay excepción; quien usa la copia recibe \texttt{Na\allowbreak{}N}. Si alguien la moldea a su tipo original, \texttt{(\allowbreak{}Transition\allowbreak{}Table\allowbreak{}Potential)\ p.\allowbreak{}copy(\allowbreak{})} lanza \texttt{Class\allowbreak{}Cast\allowbreak{}Exception}.

Medido:

\begin{itemize}
\tightlist
\item
  Una \texttt{Transition\allowbreak{}Table\allowbreak{}Potential} sobre \texttt{{[}Ill,\allowbreak{}\ E1{]}} con valores \texttt{\{0,\allowbreak{}\ 1\}} (al ocurrir E1 el nodo pasa a «yes») contesta 1.0 al muestrearla con E1 ocurrido; su \texttt{copy(\allowbreak{})} es de clase \texttt{Table\allowbreak{}With\allowbreak{}Events} y contesta \texttt{Na\allowbreak{}N}; el moldeado a \texttt{Transition\allowbreak{}Table\allowbreak{}Potential} lanza \texttt{Class\allowbreak{}Cast\allowbreak{}Exception}.
\item
  En una red de sucesos de prueba, con un suceso \texttt{Onset} que hace pasar \texttt{Ill} de «no» a «yes», copiar \texttt{Onset} e \texttt{Ill} y pegarlos deja a \texttt{Ill\textquotesingle{}} con un potencial \texttt{Table\allowbreak{}With\allowbreak{}Events}; muestrearlo con \texttt{Onset\textquotesingle{}} ocurrido devuelve \texttt{Na\allowbreak{}N}, mientras el original devuelve 1.0.
\end{itemize}
\paragraph{El código}
core/src/main/java/org/openmarkov/core/model/network/potential/TableWithEvents.java:430-433 y 486-492

\begin{Shaded}
\begin{Highlighting}[]
    \CommentTok{//}\AlertTok{TODO}
    \AttributeTok{@Override} \KeywordTok{public}\NormalTok{ Potential }\FunctionTok{copy}\OperatorTok{()} \OperatorTok{\{}
        \ControlFlowTok{return} \KeywordTok{new} \FunctionTok{TableWithEvents}\OperatorTok{(}\KeywordTok{this}\OperatorTok{);}
    \OperatorTok{\}}
    \KeywordTok{...}
    \CommentTok{/**}
     \CommentTok{*}\NormalTok{ Not implemented for this potential}\CommentTok{:}\NormalTok{ always returns }\AnnotationTok{\{@link}\NormalTok{ Double}\AnnotationTok{\#}\NormalTok{NaN}\AnnotationTok{\}}\CommentTok{.}
     \CommentTok{*/}
    \AttributeTok{@Override}
    \KeywordTok{public} \DataTypeTok{double} \FunctionTok{sampleConditionedVariable}\OperatorTok{(}\DataTypeTok{double}\OperatorTok{[]}\NormalTok{ randomNumbers}\OperatorTok{,}\NormalTok{ EvidenceCase parents}\OperatorTok{)} \KeywordTok{throws}\NormalTok{ IncompatibleEvidenceException}\OperatorTok{.}\FunctionTok{EvidenceIsIncompatibleWithOther} \OperatorTok{\{}
        \ControlFlowTok{return} \BuiltInTok{Double}\OperatorTok{.}\FunctionTok{NaN}\OperatorTok{;}
    \OperatorTok{\}}
\end{Highlighting}
\end{Shaded}

core/src/main/java/org/openmarkov/core/model/network/potential/TransitionTablePotential.java:43-46 (tiene constructor de copia, pero no \texttt{copy(\allowbreak{})})

\begin{Shaded}
\begin{Highlighting}[]
    \CommentTok{//}\AlertTok{TODO}
    \KeywordTok{public} \FunctionTok{TransitionTablePotential}\OperatorTok{(}\NormalTok{TransitionTablePotential potential}\OperatorTok{)} \OperatorTok{\{}
        \KeywordTok{super}\OperatorTok{(}\NormalTok{potential}\OperatorTok{);}
    \OperatorTok{\}}
\end{Highlighting}
\end{Shaded}
\paragraph{Cómo arreglarlo}
Redefinir \texttt{copy(\allowbreak{})} en \texttt{Transition\allowbreak{}Table\allowbreak{}Potential} para que devuelva \texttt{new\ Transition\allowbreak{}Table\allowbreak{}Potential(\allowbreak{}this)}.

\textbf{No basta solo con eso}: el constructor de copia que hereda (\texttt{Table\allowbreak{}With\allowbreak{}Events(\allowbreak{}Table\allowbreak{}With\allowbreak{}Events)}) deja la tabla interna apuntando a la variable \texttt{Events} del original, que es el defecto de \hyperref[err:19]{error~19}. Con el \texttt{copy(\allowbreak{})} corregido, el muestreo de la copia pasa de devolver \texttt{Na\allowbreak{}N} a lanzar \texttt{Null\allowbreak{}Pointer\allowbreak{}Exception} (medido llamando al constructor de copia directamente: \texttt{Evidence\allowbreak{}Case.\allowbreak{}get\allowbreak{}Finding(\allowbreak{}.\allowbreak{}.\allowbreak{}.\allowbreak{})} devuelve \texttt{null} dentro de \texttt{Table\allowbreak{}Potential.\allowbreak{}sample\allowbreak{}Conditioned\allowbreak{}Variable}). Hay que arreglar \hyperref[err:19]{error~19} a la vez.

Conviene revisar en el mismo paso todas las subclases de \texttt{Table\allowbreak{}With\allowbreak{}Events} y de \texttt{Potential} que implementan \texttt{DESSimulable\allowbreak{}Potential}: cada una debe redefinir \texttt{copy(\allowbreak{})} (hoy \texttt{Univariate\allowbreak{}Distr\allowbreak{}Potential} y \texttt{Distribution\allowbreak{}Table\allowbreak{}Potential} lo hacen; \texttt{Transition\allowbreak{}Table\allowbreak{}Potential} no).

Aparte, y no es el mismo defecto: al pegar, las relaciones de sucesos siguen consultando por dentro las variables del nodo original, lo que se describe en \hyperref[err:21]{error~21} (medido: la tabla interna de \texttt{Ill\textquotesingle{}} sigue siendo \texttt{{[}Ill,\allowbreak{}\ Events\{Onset\}{]}}). Arreglar \texttt{copy(\allowbreak{})} no hace por sí solo que pegar nodos de sucesos funcione.

Prueba de regresión: para una \texttt{Transition\allowbreak{}Table\allowbreak{}Potential} con al menos dos sucesos padres, \texttt{copy(\allowbreak{})} devuelve la misma clase y, para cada suceso, la copia muestrea lo mismo que el original con los mismos números aleatorios.

\setcounter{subsection}{20}
\subsection[Los nodos de sucesos pegados siguen leyendo los padres del nodo original]{Pegar nodos de una red de sucesos deja sus relaciones consultando los padres del original: la simulación usa siempre los números del primer estado de cada padre, o se cae}\label{err:21}
\begin{ficha}
\fichakey{Estado}Presente
\fichakey{Módulo}gui, core
\fichakey{Paquete}org.openmarkov.gui.action; org.openmarkov.core.model.network.potential; org.openmarkov.core.model.network.potential.treeadd
\fichakey{Ficheros}\path{gui/src/main/java/org/openmarkov/gui/action/PasteEdit.java:113-141} \newline \path{core/src/main/java/org/openmarkov/core/model/network/potential/Potential.java:331-334} \newline \path{core/src/main/java/org/openmarkov/core/model/network/potential/Potential.java:726-736} \newline \path{core/src/main/java/org/openmarkov/core/model/network/potential/DistributionTablePotential.java:104-112,260-301} \newline \path{core/src/main/java/org/openmarkov/core/model/network/potential/TableWithEvents.java:112-124,172-187,445-454} \newline \path{core/src/main/java/org/openmarkov/core/model/network/potential/ExactDistrPotential.java:57-60,158-160,166-169} \newline \path{core/src/main/java/org/openmarkov/core/model/network/potential/treeadd/TreeWithEventsPotential.java:229-231} \newline \path{core/src/main/java/org/openmarkov/core/model/network/potential/treeadd/TreeWithExcludedEventsPotential.java:180-182} \newline \path{core/src/main/java/org/openmarkov/core/model/network/potential/TablePotential.java:238-248} \newline \path{core/src/main/java/org/openmarkov/core/model/network/potential/treeadd/TreeADDPotential.java:840-849}
\fichakey{Comprobado}Leyendo el código y ejecutándolo
\fichakey{Esfuerzo}Medio (días)
\fichakey{Decisión de equipo}No
\fichakey{Corregir antes}\hyperref[err:19]{error~19} y \hyperref[err:20]{error~20}: sin ellos, pegar sigue sin funcionar del todo para las tablas de transiciones y para los sucesos con dos o más sucesos padres.
\fichakey{Relacionados}\hyperref[err:19]{error~19}, \hyperref[err:20]{error~20}, \hyperref[err:99]{error~99}, \hyperref[err:105]{error~105}
\end{ficha}
\paragraph{Qué pasa}
En una red DES (simulación de eventos discretos), copiar nodos y pegarlos (menú Edición, Copiar y Pegar del editor de redes) deja los nodos pegados con relaciones que, por dentro, siguen consultando las variables de los nodos originales. El portapapeles es común a todas las redes abiertas, así que se puede pegar en otra pestaña, por ejemplo en una red de sucesos nueva.

La simulación de la red pegada no avisa de nada. Cada tiempo hasta un suceso y cada valor de una utilidad se lee como si todos sus padres de estados (la decisión, los nodos de azar) estuvieran en su primer estado. Con las relaciones de árbol con sucesos, muestrear el nodo pegado lanza una excepción (medido sobre la relación, sin lanzar la simulación).

Medido con una simulación completa, en una red de sucesos de prueba:

\begin{itemize}
\tightlist
\item
  una decisión \texttt{Therapy} (none, drug) y un nodo de azar \texttt{Sex} (male, female, al 50 \%);
\item
  un suceso \texttt{Onset} con tiempo exacto de 2 o 5 años para los hombres (none, drug) y de 3 o 6 para las mujeres;
\item
  un suceso terminal \texttt{Death} 4 años después de \texttt{Onset};
\item
  dos utilidades acumulativas: \texttt{QALY}, 1 por año, y \texttt{Monitoring}, 10 por año con none y 40 con drug;
\item
  200 individuos y horizonte de 20 años.
\end{itemize}

\begin{Verbatim}[breaklines,breakanywhere,fontsize=\small]
red original:                         none: efectividad 6.52, coste 65.2   drug: efectividad 9.52, coste 380.8
todos sus nodos y enlaces pegados
en otra red con los mismos criterios: none: efectividad 6.0,  coste 60.0   drug: efectividad 6.0,  coste 60.0
\end{Verbatim}

En la red pegada, el tratamiento y el sexo no cambian nada: \texttt{Onset} ocurre siempre a los 2 años y \texttt{Monitoring} cuesta siempre 10.

Por qué pasa: \texttt{Paste\allowbreak{}Edit} (paso final, líneas 113-141) copia cada potencial con \texttt{copy(\allowbreak{})} y después cambia sus variables por las de los nodos pegados, una a una, con \texttt{potential.\allowbreak{}replace\allowbreak{}Variable(\allowbreak{}i,\allowbreak{}\ variable)}. La versión de la clase base (\texttt{Potential.\allowbreak{}replace\allowbreak{}Variable}, líneas 331-334) solo cambia la lista exterior \texttt{variables}. Los potenciales compuestos guardan otra lista por dentro, y ninguno de los de sucesos la actualiza:

\begin{itemize}
\tightlist
\item
  \texttt{Distribution\allowbreak{}Table\allowbreak{}Potential}: su \texttt{table\allowbreak{}With\allowbreak{}Events}, copiado en el constructor de copia (líneas 104-112), sigue sobre \texttt{{[}distribution\allowbreak{}Variable,\allowbreak{}\ Therapy,\allowbreak{}\ Sex{]}} del original. Al muestrear (\texttt{sample\allowbreak{}Conditioned\allowbreak{}Variable}, líneas 260-301) busca cada parámetro con \texttt{table\allowbreak{}With\allowbreak{}Events.\allowbreak{}get\allowbreak{}Table\allowbreak{}Potential(\allowbreak{}).\allowbreak{}get\allowbreak{}Value(\allowbreak{}.\allowbreak{}.\allowbreak{}.\allowbreak{})} usando los hallazgos de los padres pegados, y \texttt{Table\allowbreak{}Potential.\allowbreak{}get\allowbreak{}Value} se salta cualquier variable que no sea suya (líneas 238-248), lo que equivale a leer su estado 0.
\item
  \texttt{Exact\allowbreak{}Distr\allowbreak{}Potential} (el valor exacto de las utilidades de una red de sucesos; el lector convierte en él las \texttt{Univariate\allowbreak{}Distr} de las utilidades): su \texttt{table\allowbreak{}Potential} sigue sobre los padres del original, y \texttt{sample\allowbreak{}Conditioned\allowbreak{}Variable} (líneas 158-160) lee \texttt{table\allowbreak{}Potential.\allowbreak{}get\allowbreak{}Value(\allowbreak{}parents)}. Su \texttt{set\allowbreak{}Variables} (líneas 166-169) sí actualiza la tabla interna, pero \texttt{Paste\allowbreak{}Edit} no lo usa.
\item
  \texttt{Table\allowbreak{}With\allowbreak{}Events} y \texttt{Transition\allowbreak{}Table\allowbreak{}Potential}: la tabla interna sigue con las variables del original. (La copia de una tabla de transiciones falla antes por otro motivo, descrito en \hyperref[err:20]{error~20}.)
\item
  \texttt{Tree\allowbreak{}With\allowbreak{}Events\allowbreak{}Potential} y \texttt{Tree\allowbreak{}With\allowbreak{}Excluded\allowbreak{}Events\allowbreak{}Potential} («Tree with Events» y «Tree with Excluded Events»): su \texttt{replace\allowbreak{}Variable} está vacío (\texttt{/\allowbreak{}/\allowbreak{}TODO}), así que ni siquiera cambian las variables exteriores.
\end{itemize}

\texttt{Tree\allowbreak{}ADDPotential.\allowbreak{}replace\allowbreak{}Variable} (líneas 840-849) sí propaga el cambio a sus ramas: es el modelo que les falta a los demás.

Medido también muestreando cada relación pegada con los mismos hallazgos que el original, en una segunda red de prueba. Es como la anterior, más un nodo de azar \texttt{Ill} que pasa a «yes» cuando ocurre \texttt{Onset}, una utilidad \texttt{Care\allowbreak{}Cost} que vale 100 o 2000 según \texttt{Ill}, y un \texttt{Death} que ocurre a los 8 o a los 3 años según \texttt{Ill}. Se pegó en la misma red (nombres con apóstrofo) y en otra red (mismos nombres), con resultados iguales en los dos casos:

\begin{itemize}
\tightlist
\item
  \texttt{Onset} (tabla de distribución sobre \texttt{Therapy} y \texttt{Sex}), con drug y mujer: el original da 6,0 y la copia 2,0.
\item
  \texttt{Care\allowbreak{}Cost} (valor exacto sobre \texttt{Ill}), con \texttt{Ill}=yes: el original da 2000 y la copia 100. \texttt{Monitoring}, con drug: 40 y 10.
\item
  \texttt{Death} (tabla de distribución sobre el suceso \texttt{Onset} y sobre \texttt{Ill}), con \texttt{Onset} ocurrido e \texttt{Ill}=yes: el original da 3,0 y la copia 8,0. Aquí se suma el desajuste de la variable interna \texttt{Events} descrito en \hyperref[err:19]{error~19}.
\item
  Un \texttt{Ill} con «Tree with Events» sobre \texttt{{[}Ill,\allowbreak{}\ Onset,\allowbreak{}\ Sex{]}}, pegado en otra red: las tres variables de la relación pegada son las del original, y muestrearla lanza \texttt{Null\allowbreak{}Pointer\allowbreak{}Exception} (\texttt{get\allowbreak{}Tree} devuelve \texttt{null}). Con «Tree with Excluded Events», sin pantalla sale \texttt{Headless\allowbreak{}Exception} desde la ventana de aviso descrita en \hyperref[err:105]{error~105}.
\end{itemize}

Dónde se ve: pegar un modelo entero en la misma red no se puede simular, porque la simulación solo admite una decisión (\texttt{DESInference.\allowbreak{}initialize}, líneas 416-420, por lectura). Por eso el caso que llega a la simulación es pegar en otra red, o pegar en la misma red una parte que no incluya la decisión. Si un padre no se copia, \texttt{Paste\allowbreak{}Edit} lo quita con \texttt{remove\allowbreak{}Variable}, cuya versión de la clase base (\texttt{Potential.\allowbreak{}remove\allowbreak{}Variable}, líneas 726-736) devuelve una relación uniforme: medido, pegar \texttt{Onset} o \texttt{Ill} sin sus padres deja un \texttt{Uniform\allowbreak{}Potential}, con la consecuencia que describe \hyperref[err:99]{error~99}.
\paragraph{El código}
gui/src/main/java/org/openmarkov/gui/action/PasteEdit.java:117-131

\begin{Shaded}
\begin{Highlighting}[]
                    \ControlFlowTok{for} \OperatorTok{(}\NormalTok{Potential originalPotential }\OperatorTok{:}\NormalTok{ originalNode}\OperatorTok{.}\FunctionTok{getPotentials}\OperatorTok{())} \OperatorTok{\{}
\NormalTok{                        Potential potential }\OperatorTok{=}\NormalTok{ originalPotential}\OperatorTok{.}\FunctionTok{copy}\OperatorTok{();}
                        \BuiltInTok{List}\OperatorTok{\textless{}}\NormalTok{Variable}\OperatorTok{\textgreater{}}\NormalTok{ externalVars }\OperatorTok{=} \KeywordTok{new} \BuiltInTok{ArrayList}\OperatorTok{\textless{}\textgreater{}();}
                        \ControlFlowTok{for} \OperatorTok{(}\DataTypeTok{int}\NormalTok{ i }\OperatorTok{=} \DecValTok{0}\OperatorTok{;}\NormalTok{ i }\OperatorTok{\textless{}}\NormalTok{ potential}\OperatorTok{.}\FunctionTok{getNumVariables}\OperatorTok{();} \OperatorTok{++}\NormalTok{i}\OperatorTok{)} \OperatorTok{\{}
                            \BuiltInTok{String}\NormalTok{ variableName }\OperatorTok{=}\NormalTok{ potential}\OperatorTok{.}\FunctionTok{getVariable}\OperatorTok{(}\NormalTok{i}\OperatorTok{).}\FunctionTok{getName}\OperatorTok{();}
                            \ControlFlowTok{if} \OperatorTok{(}\NormalTok{PasteEdit}\OperatorTok{.}\FunctionTok{this}\OperatorTok{.}\FunctionTok{newVariables}\OperatorTok{.}\FunctionTok{containsKey}\OperatorTok{(}\NormalTok{variableName}\OperatorTok{))} \OperatorTok{\{}
\NormalTok{                                Variable variable }\OperatorTok{=} \KeywordTok{this}\OperatorTok{.}\FunctionTok{probNet}\OperatorTok{.}\FunctionTok{getVariable}\OperatorTok{(}\NormalTok{PasteEdit}\OperatorTok{.}\FunctionTok{this}\OperatorTok{.}\FunctionTok{newVariables}\OperatorTok{.}\FunctionTok{get}\OperatorTok{(}\NormalTok{variableName}\OperatorTok{));}
\NormalTok{                                potential}\OperatorTok{.}\FunctionTok{replaceVariable}\OperatorTok{(}\NormalTok{i}\OperatorTok{,}\NormalTok{ variable}\OperatorTok{);}
                            \OperatorTok{\}} \ControlFlowTok{else} \OperatorTok{\{}
\NormalTok{                                externalVars}\OperatorTok{.}\FunctionTok{add}\OperatorTok{(}\NormalTok{potential}\OperatorTok{.}\FunctionTok{getVariable}\OperatorTok{(}\NormalTok{i}\OperatorTok{));}
                            \OperatorTok{\}}
                        \OperatorTok{\}}
                        \ControlFlowTok{for} \OperatorTok{(}\NormalTok{Variable extVar }\OperatorTok{:}\NormalTok{ externalVars}\OperatorTok{)} \OperatorTok{\{}
\NormalTok{                            potential }\OperatorTok{=}\NormalTok{ potential}\OperatorTok{.}\FunctionTok{removeVariable}\OperatorTok{(}\NormalTok{extVar}\OperatorTok{);}
                        \OperatorTok{\}}
\end{Highlighting}
\end{Shaded}

core/src/main/java/org/openmarkov/core/model/network/potential/Potential.java:331-334

\begin{Shaded}
\begin{Highlighting}[]
    \KeywordTok{public} \DataTypeTok{void} \FunctionTok{replaceVariable}\OperatorTok{(}\DataTypeTok{int}\NormalTok{ position}\OperatorTok{,}\NormalTok{ Variable variable}\OperatorTok{)} \OperatorTok{\{}
\NormalTok{        variables}\OperatorTok{.}\FunctionTok{remove}\OperatorTok{(}\NormalTok{position}\OperatorTok{);}
\NormalTok{        variables}\OperatorTok{.}\FunctionTok{add}\OperatorTok{(}\NormalTok{position}\OperatorTok{,}\NormalTok{ variable}\OperatorTok{);}
    \OperatorTok{\}}
\end{Highlighting}
\end{Shaded}

core/src/main/java/org/openmarkov/core/model/network/potential/treeadd/TreeWithEventsPotential.java:229-231

\begin{Shaded}
\begin{Highlighting}[]
    \AttributeTok{@Override} \KeywordTok{public} \DataTypeTok{void} \FunctionTok{replaceVariable}\OperatorTok{(}\DataTypeTok{int}\NormalTok{ position}\OperatorTok{,}\NormalTok{ Variable variable}\OperatorTok{)} \OperatorTok{\{}
        \CommentTok{//}\AlertTok{TODO}
    \OperatorTok{\}}
\end{Highlighting}
\end{Shaded}
\paragraph{Cómo arreglarlo}
Cada potencial compuesto debe propagar \texttt{replace\allowbreak{}Variable} a lo que guarda por dentro, como ya hace \texttt{Tree\allowbreak{}ADDPotential}:

\begin{itemize}
\tightlist
\item
  \texttt{Distribution\allowbreak{}Table\allowbreak{}Potential}: las variables de \texttt{table\allowbreak{}With\allowbreak{}Events} (la variable de parámetros más los padres que no son numéricos), su tabla y su \texttt{table\allowbreak{}With\allowbreak{}Functions}, y la lista \texttt{numeric\allowbreak{}Variables}.
\item
  \texttt{Table\allowbreak{}With\allowbreak{}Events}: la tabla interna y \texttt{table\allowbreak{}With\allowbreak{}Functions}. Si la variable cambiada es un suceso, hay que rehacer también la variable \texttt{events}, cuyos estados son los \textbf{nombres} de los sucesos padres: un suceso que pasa a llamarse \texttt{Onset\textquotesingle{}} tiene que cambiar el estado, o \texttt{convert} (líneas 172-187) no lo encontrará. \texttt{Table\allowbreak{}With\allowbreak{}Events.\allowbreak{}deep\allowbreak{}Copy} (líneas 445-454) ya rehace todo esto a partir de las variables exteriores y se puede reutilizar.
\item
  \texttt{Exact\allowbreak{}Distr\allowbreak{}Potential}: su \texttt{table\allowbreak{}Potential}, igual que hace ya \texttt{set\allowbreak{}Variables} (líneas 166-169).
\item
  \texttt{Tree\allowbreak{}With\allowbreak{}Events\allowbreak{}Potential} (las claves de su mapa de árboles y cada árbol) y \texttt{Tree\allowbreak{}With\allowbreak{}Excluded\allowbreak{}Events\allowbreak{}Potential} (su \texttt{no\allowbreak{}Event\allowbreak{}Tree}): escribir el método, que hoy está vacío.
\end{itemize}

Otra vía, que no quita la anterior: que \texttt{Paste\allowbreak{}Edit} construya la lista completa de variables nuevas y la ponga con \texttt{set\allowbreak{}Variables}, o copie con una correspondencia de nombres como hace \texttt{Potential.\allowbreak{}deep\allowbreak{}Copy}. \texttt{Distribution\allowbreak{}Table\allowbreak{}Potential} y los dos árboles tampoco propagan hoy \texttt{set\allowbreak{}Variables}, así que el trabajo en los potenciales hace falta igualmente.

\texttt{Paste\allowbreak{}Edit} es hoy el único llamador de producción de \texttt{replace\allowbreak{}Variable} sobre estos potenciales. Con el arreglo, pegar sigue sin funcionar del todo para las tablas de transiciones y para los sucesos con dos o más sucesos padres mientras no se corrijan \hyperref[err:20]{error~20} y \hyperref[err:19]{error~19}.

Pruebas de regresión:

\begin{itemize}
\tightlist
\item
  Para cada relación de sucesos (\texttt{Distribution\allowbreak{}Table\allowbreak{}Potential}, \texttt{Transition\allowbreak{}Table\allowbreak{}Potential}, \texttt{Exact\allowbreak{}Distr\allowbreak{}Potential}, los dos árboles), \texttt{copy(\allowbreak{})} más \texttt{replace\allowbreak{}Variable} de todas sus variables no deja ninguna referencia a las variables del original, tampoco en las listas internas.
\item
  Pegar todos los nodos de la red de prueba descrita en otra red da los mismos resultados de simulación que el original.
\end{itemize}

\setcounter{subsection}{21}
\subsection[Aceptar un árbol con sucesos deja un deshacer que borra la relación]{Abrir el editor de un «Tree with Events» o «Tree with Excluded Events» y aceptar deja un deshacer que borra la relación del nodo}\label{err:22}
\begin{ficha}
\fichakey{Estado}Presente
\fichakey{Módulo}core
\fichakey{Paquete}org.openmarkov.core.model.network.potential.treeadd
\fichakey{Ficheros}\path{core/src/main/java/org/openmarkov/core/model/network/potential/treeadd/TreeWithEventsPotential.java:217-226,256-259,275-278} \newline \path{core/src/main/java/org/openmarkov/core/model/network/potential/treeadd/TreeWithExcludedEventsPotential.java:168-177,207-210,226-229} \newline \path{gui/src/main/java/org/openmarkov/gui/dialog/node/PotentialEditPanel.java:148-153, 537-561, 563-568} \newline \path{core/src/main/java/org/openmarkov/core/action/core/SetPotentialEdit.java:69-91} \newline \path{core/src/main/java/org/openmarkov/core/model/network/ProbNetCopier.java:111} \newline \path{core/src/main/java/org/openmarkov/core/model/network/potential/Potential.java:801-820}
\fichakey{Comprobado}Leyendo el código y ejecutándolo
\fichakey{Esfuerzo}Pequeño (horas)
\fichakey{Decisión de equipo}No
\fichakey{Relacionados}\hyperref[err:23]{error~23}, \hyperref[err:105]{error~105}
\end{ficha}
\paragraph{Qué pasa}
En una red de simulación de eventos discretos, si el usuario abre las propiedades de un nodo cuya relación es «Tree with Events» o «Tree with Excluded Events» (o «Edit probability»), pulsa OK sin tocar nada y después Ctrl+Z, el nodo se queda sin relación: su lista de relaciones pasa a ser \texttt{{[}null{]}}. A partir de ahí \texttt{Prob\allowbreak{}Net.\allowbreak{}deep\allowbreak{}Copy} falla con \texttt{Null\allowbreak{}Pointer\allowbreak{}Exception} en \texttt{Prob\allowbreak{}Net\allowbreak{}Copier}, línea 111. Por lectura, el diálogo siguiente sobre ese nodo no podrá construir su selector de relación (\texttt{get\allowbreak{}Potentials(\allowbreak{}).\allowbreak{}get\allowbreak{}First(\allowbreak{}).\allowbreak{}get\allowbreak{}Class(\allowbreak{})}), y el escritor de ficheros encontrará un \texttt{null}; esto no se ha comprobado con la aplicación abierta.

Por qué pasa: \texttt{Tree\allowbreak{}With\allowbreak{}Events\allowbreak{}Potential} y \texttt{Tree\allowbreak{}With\allowbreak{}Excluded\allowbreak{}Events\allowbreak{}Potential} (relaciones de la simulación de eventos discretos que combinan árboles con sucesos) tienen varios métodos sin escribir, con \texttt{/\allowbreak{}/\allowbreak{}TODO}: \texttt{deep\allowbreak{}Copy(\allowbreak{}Prob\allowbreak{}Net)} devuelve \texttt{null} (líneas 256-259 y 207-210), \texttt{sample(\allowbreak{})} y \texttt{get\allowbreak{}Induced\allowbreak{}Findings} devuelven \texttt{null}, y \texttt{to\allowbreak{}String(\allowbreak{})} devuelve \texttt{""}. \texttt{deep\allowbreak{}Copy} anula el de la clase base (\texttt{Potential.\allowbreak{}deep\allowbreak{}Copy}, líneas 801-820), que sí copiaría la relación por su constructor de copia; bastaba con no anularlo. Medido: \texttt{deep\allowbreak{}Copy} de las dos devuelve \texttt{null} y \texttt{to\allowbreak{}String} la cadena vacía.

\texttt{deep\allowbreak{}Copy} tiene al menos dos llamadores de producción:

\begin{enumerate}
\def\labelenumi{\arabic{enumi}.}
\tightlist
\item
  \textbf{El editor de relaciones.} \texttt{Potential\allowbreak{}Edit\allowbreak{}Panel} (módulo \texttt{gui}, líneas 145-151) guarda al abrirse \texttt{original\allowbreak{}Potential\ =\ node.\allowbreak{}get\allowbreak{}Potential(\allowbreak{}).\allowbreak{}deep\allowbreak{}Copy(\allowbreak{}.\allowbreak{}.\allowbreak{}.\allowbreak{})}, que para estas dos clases es \texttt{null}. Al aceptar, \texttt{commit\allowbreak{}Changes} (líneas 537-549) compara la relación actual con \texttt{null} (\texttt{potential\allowbreak{}Has\allowbreak{}Changed}, líneas 552-557; medido \texttt{Reflection\allowbreak{}Equality.\allowbreak{}are\allowbreak{}Equals(\allowbreak{}null,\allowbreak{}\ relación)\ =\ false}), concluye que ha cambiado aunque el usuario no haya tocado nada, y ejecuta \texttt{Set\allowbreak{}Potential\allowbreak{}Edit(\allowbreak{}nodo,\allowbreak{}\ null,\allowbreak{}\ relación)}. Deshacer esa edición (\texttt{Set\allowbreak{}Potential\allowbreak{}Edit.\allowbreak{}undo}, líneas 88-91) pone como relación del nodo \texttt{null}. Medido: tras ejecutar y deshacer esa edición, la lista de relaciones del nodo es \texttt{{[}null{]}}, y \texttt{Prob\allowbreak{}Net.\allowbreak{}deep\allowbreak{}Copy} falla como se ha dicho.
\item
  \textbf{La copia profunda de una red.} \texttt{Prob\allowbreak{}Net\allowbreak{}Copier.\allowbreak{}deep\allowbreak{}Copy} (línea 111) copia cada relación con \texttt{deep\allowbreak{}Copy}; medido que la copia de una red deja el nodo con \texttt{{[}null{]}}. Hoy su único llamador de producción es la expansión temporal (\texttt{Temporal\allowbreak{}Net\allowbreak{}Operations}, línea 457), que no se usa con redes de simulación de eventos discretos, así que este camino no se alcanza hoy.
\end{enumerate}

\texttt{sample(\allowbreak{})}, \texttt{get\allowbreak{}Induced\allowbreak{}Findings} y \texttt{to\allowbreak{}String} no tienen consecuencia medida; se listan porque son del mismo tipo y el arreglo es el mismo.

Otro hallazgo sobre estos árboles, que muestrearlos vacía la configuración de padres que reciben, está en \hyperref[nota:74]{nota~74}; allí no hay nada que corregir.
\paragraph{El código}
core/src/main/java/org/openmarkov/core/model/network/potential/treeadd/TreeWithExcludedEventsPotential.java:207-210

\begin{Shaded}
\begin{Highlighting}[]
\AttributeTok{@Override} \KeywordTok{public}\NormalTok{ Potential }\FunctionTok{deepCopy}\OperatorTok{(}\NormalTok{ProbNet copyNet}\OperatorTok{)} \OperatorTok{\{}
    \CommentTok{//}\AlertTok{TODO}
    \ControlFlowTok{return} \KeywordTok{null}\OperatorTok{;}
\OperatorTok{\}}
\end{Highlighting}
\end{Shaded}

gui/src/main/java/org/openmarkov/gui/dialog/node/PotentialEditPanel.java:145-151

\begin{Shaded}
\begin{Highlighting}[]
\NormalTok{Potential lastPotential }\OperatorTok{=} \KeywordTok{this}\OperatorTok{.}\FunctionTok{node}\OperatorTok{.}\FunctionTok{getPotential}\OperatorTok{();}
\ControlFlowTok{if} \OperatorTok{(}\NormalTok{lastPotential }\OperatorTok{!=} \KeywordTok{null}\OperatorTok{)} \OperatorTok{\{}
    \KeywordTok{this}\OperatorTok{.}\FunctionTok{originalPotential} \OperatorTok{=} \KeywordTok{this}\OperatorTok{.}\FunctionTok{node}\OperatorTok{.}\FunctionTok{getPotential}\OperatorTok{().}\FunctionTok{deepCopy}\OperatorTok{(}\KeywordTok{this}\OperatorTok{.}\FunctionTok{node}\OperatorTok{.}\FunctionTok{getProbNet}\OperatorTok{());}
\OperatorTok{\}} \ControlFlowTok{else} \OperatorTok{\{}
    \KeywordTok{this}\OperatorTok{.}\FunctionTok{originalPotential} \OperatorTok{=} \KeywordTok{null}\OperatorTok{;}
\OperatorTok{\}}
\end{Highlighting}
\end{Shaded}
\paragraph{Cómo arreglarlo}
Quitar los \texttt{deep\allowbreak{}Copy} que devuelven \texttt{null} (la base ya copia por el constructor de copia) o implementarlos copiando también los árboles internos con las variables de la red destino, como hace \texttt{Tree\allowbreak{}ADDPotential.\allowbreak{}deep\allowbreak{}Copy}; \texttt{Tree\allowbreak{}With\allowbreak{}Events\allowbreak{}Potential} debe copiar su mapa de árboles por suceso y \texttt{Tree\allowbreak{}With\allowbreak{}Excluded\allowbreak{}Events\allowbreak{}Potential} su \texttt{no\allowbreak{}Event\allowbreak{}Tree}. Implementar \texttt{to\allowbreak{}String} con algo que identifique la relación; \texttt{sample(\allowbreak{})} y \texttt{get\allowbreak{}Induced\allowbreak{}Findings} deberían lanzar \texttt{Unsupported\allowbreak{}Operation\allowbreak{}Exception} como la base en vez de devolver \texttt{null}.

Aparte, \texttt{Set\allowbreak{}Potential\allowbreak{}Edit} con relación anterior \texttt{null} no debería poder dejar un \texttt{null} en la lista del nodo; eso protege a cualquier otra relación con el mismo fallo.

Prueba: las pruebas de copia (\texttt{copy}/\texttt{deep\allowbreak{}Copy}) del proyecto deberían cubrir estas dos clases: copiar una red de simulación de eventos discretos con cada una y comprobar que la relación copiada no es nula y usa las variables de la copia.

\setcounter{subsection}{22}
\subsection[Se ofrece «Tree with Excluded Events» con solo padres suceso, y falla]{El editor ofrece «Tree with Excluded Events» a un nodo cuyos padres son todos sucesos, y elegirlo lanza un error no previsto}\label{err:23}
\begin{ficha}
\fichakey{Estado}Presente
\fichakey{Módulo}core
\fichakey{Paquete}org.openmarkov.core.model.network.potential.treeadd
\fichakey{Ficheros}\path{core/src/main/java/org/openmarkov/core/model/network/potential/treeadd/TreeWithExcludedEventsPotential.java:67-81,102-111} \newline \path{core/src/main/java/org/openmarkov/core/model/network/potential/treeadd/TreeADDPotential.java:189-191} \newline \path{core/src/main/java/org/openmarkov/core/model/network/potential/plugin/PotentialUtils.java:90-109, 118-135} \newline \path{gui/src/main/java/org/openmarkov/gui/dialog/node/PotentialEditPanel.java:488-522, 675-683}
\fichakey{Comprobado}Leyendo el código y ejecutándolo
\fichakey{Esfuerzo}Pequeño (horas)
\fichakey{Decisión de equipo}No
\fichakey{Relacionados}\hyperref[err:22]{error~22}, \hyperref[err:101]{error~101}, \hyperref[err:105]{error~105}, \hyperref[err:162]{error~162}, \hyperref[nota:96]{nota~96}
\end{ficha}
\paragraph{Qué pasa}
En una red de simulación de eventos discretos, el selector de clase de relación ofrece «Tree with Excluded Events» a un nodo cuyos padres son todos sucesos, y elegirlo termina en un error no previsto. La relación del nodo no llega a cambiar, pero el selector queda mostrando la clase elegida.

Camino del usuario (la parte de la interfaz está comprobada leyendo): propiedades del nodo → pestaña de relación → selector de clase → «Tree with Excluded Events». \texttt{Potential\allowbreak{}Edit\allowbreak{}Panel.\allowbreak{}potential\allowbreak{}Type\allowbreak{}Changed} (líneas 476-507) llama a \texttt{instanciate\allowbreak{}Potential} (líneas 663-671) → \texttt{Potential\allowbreak{}Utils.\allowbreak{}instanciate\allowbreak{}Safely} (líneas 93-112), que recibe la excepción envuelta en \texttt{Invocation\allowbreak{}Target\allowbreak{}Exception}, prueba un constructor de un argumento que no existe y termina lanzando \texttt{Unreachable\allowbreak{}Exception}. El manejador global la enseña como error no previsto.

Por qué pasa: \texttt{Tree\allowbreak{}With\allowbreak{}Excluded\allowbreak{}Events\allowbreak{}Potential} es un árbol de relaciones que ramifica solo sobre los padres que \textbf{no} son sucesos. Su \texttt{validate} (líneas 102-111), que decide si el selector la ofrece, exige que la red sea de simulación de eventos discretos y que \textbf{algún} padre sea un suceso, pero no que haya algún padre que no lo sea. Su constructor (líneas 67-81) arma el árbol con la variable hija más los padres que no son sucesos y llama a \texttt{new\ Tree\allowbreak{}ADDPotential(\allowbreak{}variables,\allowbreak{}\ role)}, que ramifica sobre \texttt{variables.\allowbreak{}get(\allowbreak{}1)} (\texttt{Tree\allowbreak{}ADDPotential}, líneas 189-191). Si todos los padres son sucesos, la lista tiene un solo elemento y el constructor lanza \texttt{Index\allowbreak{}Out\allowbreak{}Of\allowbreak{}Bounds\allowbreak{}Exception}.

La hermana \texttt{Tree\allowbreak{}With\allowbreak{}Events\allowbreak{}Potential} no tiene este problema: arma un árbol por suceso con la hija y ese suceso, así que siempre tiene dos variables.

Medido (red de tipo \texttt{DESNetwork\allowbreak{}Type}, nodo de azar Z con dos padres suceso E y E2 y nada más):

\begin{Verbatim}[breaklines,breakanywhere,fontsize=\small]
validate = true
new TreeWithExcludedEventsPotential([Z, E, E2]) -> IndexOutOfBoundsException: Index 1 out of bounds for length 1
tipos ofrecidos para Z: [TransitionTablePotential, ExternalPotential, DeltaPotential, WeibullHazardPotential,
                         TreeWithExcludedEventsPotential, TreeWithEventsPotential, TreeADDPotential, ICIPotential]
\end{Verbatim}
\paragraph{El código}
core/src/main/java/org/openmarkov/core/model/network/potential/treeadd/TreeWithExcludedEventsPotential.java:102-111

\begin{Shaded}
\begin{Highlighting}[]
\KeywordTok{public} \DataTypeTok{static} \DataTypeTok{boolean} \FunctionTok{validate}\OperatorTok{(}\BuiltInTok{Node}\NormalTok{ node}\OperatorTok{,} \BuiltInTok{List}\OperatorTok{\textless{}}\NormalTok{Variable}\OperatorTok{\textgreater{}}\NormalTok{ variables}\OperatorTok{,}\NormalTok{ PotentialRole role}\OperatorTok{)} \OperatorTok{\{}
    \ControlFlowTok{if} \OperatorTok{(!(}\NormalTok{node}\OperatorTok{.}\FunctionTok{getProbNet}\OperatorTok{().}\FunctionTok{getNetworkType}\OperatorTok{()} \KeywordTok{instanceof}\NormalTok{ DESNetworkType}\OperatorTok{))} \ControlFlowTok{return} \KeywordTok{false}\OperatorTok{;}
    \KeywordTok{...}
    \DataTypeTok{boolean}\NormalTok{ validate }\OperatorTok{=}\NormalTok{ variables}\OperatorTok{.}\FunctionTok{subList}\OperatorTok{(}\DecValTok{1}\OperatorTok{,}\NormalTok{ variables}\OperatorTok{.}\FunctionTok{size}\OperatorTok{())}
                                \OperatorTok{.}\FunctionTok{stream}\OperatorTok{()}
                                \OperatorTok{.}\FunctionTok{anyMatch}\OperatorTok{(}\NormalTok{variable }\OperatorTok{{-}\textgreater{}}\NormalTok{ variable}\OperatorTok{.}\FunctionTok{getVariableType}\OperatorTok{()} \OperatorTok{==}\NormalTok{ VariableType}\OperatorTok{.}\FunctionTok{EVENT}\OperatorTok{);}
    \ControlFlowTok{return}\NormalTok{ validate}\OperatorTok{;}
\OperatorTok{\}}
\end{Highlighting}
\end{Shaded}
\paragraph{Cómo arreglarlo}
Añadir a \texttt{validate} la condición de que algún padre no sea suceso (\texttt{any\allowbreak{}Match(\allowbreak{}.\allowbreak{}.\allowbreak{}.\allowbreak{}\ !=\ Variable\allowbreak{}Type.\allowbreak{}EVENT)}), para que el selector no la ofrezca donde no se puede construir.

Como la relación también puede venir de un fichero (\texttt{PGMXReader\_\allowbreak{}1\_\allowbreak{}0.\allowbreak{}get\allowbreak{}Tree\allowbreak{}With\allowbreak{}Excluded\allowbreak{}Events\allowbreak{}Potential}) y los padres pueden cambiar después (quitar el último padre que no es suceso), el constructor debería rechazar ese caso con una excepción con mensaje en vez de un índice fuera de rango; y quitar un enlace que deje la relación inválida debería tratarse igual que en las demás relaciones con requisitos sobre sus padres.

De paso, \texttt{instanciate\allowbreak{}Safely} convierte cualquier excepción del constructor en \texttt{Unreachable\allowbreak{}Exception}; un constructor que rechaza sus argumentos no es un caso inalcanzable, así que conviene dejar pasar la causa.

Prueba: para un nodo con solo padres suceso, \texttt{validate} devuelve \texttt{false} y el nodo no la ofrece; para uno con un suceso y un padre de azar, se construye sin error.

\setcounter{subsection}{192}
\subsection[«Absorb parents» falla si un padre es una regresión]{«Absorb parents» sobre un nodo de suma, de producto o de fórmula sale como error inesperado cuando un padre está definido por una regresión, porque la tabla de la regresión trae la variable del propio padre}\label{err:193}
\begin{ficha}
\fichakey{Estado}Presente
\fichakey{Módulo}core
\fichakey{Paquete}org.openmarkov.core.action.core, org.openmarkov.core.inference
\fichakey{Ficheros}\path{core/src/main/java/org/openmarkov/core/action/core/AbsorbParentsEdit.java:33-68} \newline \path{core/src/main/java/org/openmarkov/core/inference/BasicOperations.java:43-76, 201-238} \newline \path{core/src/main/java/org/openmarkov/core/model/network/potential/LinearCombinationPotential.java:94-95} \newline \path{core/src/main/java/org/openmarkov/core/model/network/potential/ExponentialPotential.java:88-89}
\fichakey{Comprobado}Leyendo el código y ejecutándolo
\fichakey{Esfuerzo}Pequeño (horas)
\fichakey{Decisión de equipo}No
\fichakey{Relacionados}\hyperref[err:1]{error~1}, \hyperref[err:8]{error~8}, \hyperref[err:9]{error~9}, \hyperref[nota:185]{nota~185}
\end{ficha}
\paragraph{Qué pasa}
Sobre un nodo numérico cuya relación es una suma, un producto o una fórmula, «Absorb parents» sale como error inesperado cuando alguno de sus padres está definido por una regresión (combinación lineal o exponencial). El menú está encendido, porque \texttt{Basic\allowbreak{}Operations.\allowbreak{}have\allowbreak{}Parents\allowbreak{}And\allowbreak{}Are\allowbreak{}All\allowbreak{}Absorbable} solo mira que el nodo y sus padres sean numéricos y que los abuelos tengan estados finitos.

Camino del usuario (la parte de la interfaz está comprobada leyendo): en el editor, un nodo numérico con relación de suma, de producto o de fórmula que tiene como padre un nodo numérico con relación de regresión; menú contextual del nodo → «Absorb parents». \texttt{GUIUtils.\allowbreak{}execute\allowbreak{}UIAction} envuelve el fallo en \texttt{Unrecoverable\allowbreak{}Exception}, y el manejador global lo enseña como error inesperado. En el repositorio no hay ninguna red así: las únicas regresiones de nodos numéricos, las de \texttt{integration\allowbreak{}Tests/\allowbreak{}src/\allowbreak{}main/\allowbreak{}resources/\allowbreak{}networks\_\allowbreak{}for\_\allowbreak{}every\_\allowbreak{}potential/\allowbreak{}Dynamic-LIMID-every-potential.\allowbreak{}pgmx}, no tienen hijos.

Por qué pasa: la tabla de una regresión lleva delante la variable del propio nodo (\texttt{Linear\allowbreak{}Combination\allowbreak{}Potential.\allowbreak{}table\allowbreak{}Project}, línea 95; \texttt{Exponential\allowbreak{}Potential.\allowbreak{}table\allowbreak{}Project}, línea 89), una variable numérica de un solo estado. La absorción combina las tablas de los padres (\texttt{Basic\allowbreak{}Operations.\allowbreak{}absorb\allowbreak{}Parent\allowbreak{}Potentials}) y convierte cada variable del resultado en un padre del nodo, así que el padre absorbido vuelve a aparecer como padre. \texttt{Absorb\allowbreak{}Parents\allowbreak{}Edit.\allowbreak{}generate\allowbreak{}Edits} pide a la vez borrar ese padre y añadir el enlace desde él, y la comprobación de restricciones rechaza la edición porque ese enlace ya existe (\texttt{Constraint\allowbreak{}Violated\allowbreak{}Exception.\allowbreak{}Link\allowbreak{}Already\allowbreak{}Exists}). El preproceso de la inferencia (\texttt{Basic\allowbreak{}Operations.\allowbreak{}absorb\allowbreak{}Parents}) no falla, pero vuelve a enlazar el padre y, como ahora tiene un hijo, no lo borra.

Medido con \texttt{Dynamic-LIMID-every-potential.\allowbreak{}pgmx}, enlazando \texttt{Linear\allowbreak{}Combination\ [0]} y \texttt{Exponential\ [0]}, dos regresiones sin padres, con \texttt{Sum\ [0]} o con \texttt{Product\ [0]}:

\begin{Verbatim}[breaklines,breakanywhere,fontsize=\small]
                               menú «Absorb parents»   preproceso de la inferencia
suma, un padre de regresión    falla                   el padre sigue en la red
suma, dos padres de regresión  falla                   los padres siguen en la red
producto, un padre             falla                   el padre sigue en la red
producto, dos padres           funciona                funciona
\end{Verbatim}

El menú falla con \texttt{Cannot\allowbreak{}Do\allowbreak{}Edit\allowbreak{}Exception}, causada por \texttt{Link\allowbreak{}Already\allowbreak{}Exists} (con dos padres, por \texttt{Multiple\allowbreak{}Constraints\allowbreak{}Violateds}).

El producto de dos padres funciona solo porque la multiplicación deja fuera las variables de un solo estado de las tablas de una casilla (\hyperref[nota:185]{nota~185}). Con regresiones que tengan padres propios, sus tablas tienen más de una casilla, el producto conserva la variable y falla también; esto último se ha deducido leyendo.
\paragraph{El código}
core/src/main/java/org/openmarkov/core/action/core/AbsorbParentsEdit.java:49-65

\begin{Shaded}
\begin{Highlighting}[]
        \ControlFlowTok{for} \OperatorTok{(}\BuiltInTok{Node}\NormalTok{ parent }\OperatorTok{:}\NormalTok{ parents}\OperatorTok{)} \OperatorTok{\{}
\NormalTok{            PNEdit newEdit}\OperatorTok{;}
            \ControlFlowTok{if} \OperatorTok{(}\NormalTok{parent}\OperatorTok{.}\FunctionTok{getChildren}\OperatorTok{().}\FunctionTok{size}\OperatorTok{()} \OperatorTok{\textgreater{}} \DecValTok{1}\OperatorTok{)} \OperatorTok{\{}
\NormalTok{                newEdit }\OperatorTok{=} \KeywordTok{new} \FunctionTok{RemoveLinkEdit}\OperatorTok{(}\NormalTok{probNet}\OperatorTok{,}\NormalTok{ parent}\OperatorTok{.}\FunctionTok{getVariable}\OperatorTok{(),}\NormalTok{ nodeVariable}\OperatorTok{,} \KeywordTok{true}\OperatorTok{,} \KeywordTok{false}\OperatorTok{);}
            \OperatorTok{\}} \ControlFlowTok{else} \OperatorTok{\{}  \CommentTok{// parent.getChildren().size() == 1}
\NormalTok{                newEdit }\OperatorTok{=} \KeywordTok{new} \FunctionTok{CRemoveNodeEdit}\OperatorTok{(}\NormalTok{probNet}\OperatorTok{,}\NormalTok{ parent}\OperatorTok{);}
            \OperatorTok{\}}
\NormalTok{            edits}\OperatorTok{.}\FunctionTok{add}\OperatorTok{(}\NormalTok{newEdit}\OperatorTok{);}
        \OperatorTok{\}}

        \ControlFlowTok{for} \OperatorTok{(}\NormalTok{Variable variable }\OperatorTok{:}\NormalTok{ potential}\OperatorTok{.}\FunctionTok{getVariables}\OperatorTok{())} \OperatorTok{\{}
            \ControlFlowTok{if} \OperatorTok{(}\NormalTok{variable }\OperatorTok{!=}\NormalTok{ nodeVariable}\OperatorTok{)} \OperatorTok{\{}
\NormalTok{                AddLinkEdit addLinkEdit }\OperatorTok{=} \KeywordTok{new} \FunctionTok{AddLinkEdit}\OperatorTok{(}\NormalTok{probNet}\OperatorTok{,}\NormalTok{ variable}\OperatorTok{,}\NormalTok{ nodeVariable}\OperatorTok{,} \KeywordTok{true}\OperatorTok{);}
                \KeywordTok{...}
\NormalTok{                edits}\OperatorTok{.}\FunctionTok{add}\OperatorTok{(}\NormalTok{addLinkEdit}\OperatorTok{);}
            \OperatorTok{\}}
        \OperatorTok{\}}
\end{Highlighting}
\end{Shaded}
\paragraph{Cómo arreglarlo}
Al absorber, quitar de la tabla de cada padre la variable del propio padre antes de combinarlas: tiene un solo estado, así que ningún valor cambia (\texttt{Table\allowbreak{}Potential.\allowbreak{}remove\allowbreak{}Variable} la proyecta en su único estado). La regla tiene que cumplirse en los dos sitios que arman esas tablas, el menú (\texttt{Absorb\allowbreak{}Parents\allowbreak{}Edit.\allowbreak{}generate\allowbreak{}Edits}, líneas 38-45) y el preproceso de la inferencia (\texttt{Basic\allowbreak{}Operations.\allowbreak{}absorb\allowbreak{}Parents}, líneas 205-212); lo natural es que los dos llamen al mismo método.

La otra salida, que la tabla de una regresión no lleve la variable del nodo, cambia lo que reciben todos los que proyectan una regresión, la inferencia incluida.

Prueba de regresión: en \texttt{Dynamic-LIMID-every-potential.\allowbreak{}pgmx}, con \texttt{Linear\allowbreak{}Combination\ [0]} enlazado a \texttt{Sum\ [0]}, «Absorb parents» debe terminar sin error, borrar \texttt{Linear\allowbreak{}Combination\ [0]}, dejar \texttt{Sum\ [0]} sin padres y darle el valor de la regresión; lo mismo con dos padres, con un producto y con el preproceso de la inferencia.

% ══════════════════════════════════════════════════════════════════════
\section{Ediciones de la red: deshacer, cancelar y pegar}\label{sec:ediciones}

Cada cambio que el usuario hace sobre una red es una edición que se apunta en el historial y se puede deshacer. Aquí van los cambios que no se deshacen bien, los que dejan la red a medias y los que se rechazan sin motivo.

\setcounter{subsection}{23}
\subsection[Un dominio estándar de otro tamaño borra tablas que «Cancel» no devuelve]{Elegir un dominio estándar con otro número de estados borra las tablas del nodo y de sus hijos, y ni «Cancel» ni «deshacer» las devuelven}\label{err:24}
\begin{ficha}
\fichakey{Estado}Presente
\fichakey{Módulo}core
\fichakey{Paquete}org.openmarkov.core.action.core
\fichakey{Ficheros}\path{core/src/main/java/org/openmarkov/core/action/core/NodeReplaceStatesEdit.java} \newline \path{core/src/main/java/org/openmarkov/core/model/network/VariableStateOperations.java:69-110,292-320} \newline \path{core/src/main/java/org/openmarkov/core/action/core/NodeStateEdit.java:167-210} \newline \path{gui/src/main/java/org/openmarkov/gui/dialog/node/NodeDomainValuesTablePanel.java:1078-1117} \newline \path{gui/src/main/java/org/openmarkov/gui/dialog/node/NodePropertiesDialog.java:399-402}
\fichakey{Comprobado}Leyendo el código y ejecutándolo
\fichakey{Esfuerzo}Medio (días)
\fichakey{Decisión de equipo}No
\fichakey{Corregir antes}\hyperref[err:30]{error~30}: si se comparte con \texttt{Node\allowbreak{}State\allowbreak{}Edit} la fotografía de los potenciales, como propone esta ficha, primero hay que corregirla.
\fichakey{Relacionados}\hyperref[err:25]{error~25}, \hyperref[err:30]{error~30}, \hyperref[err:173]{error~173}, \hyperref[nota:5]{nota~5}, \hyperref[nota:15]{nota~15}, \hyperref[nota:150]{nota~150}
\end{ficha}
\paragraph{Qué pasa}
Abrir el diálogo de un nodo, probar en la pestaña de dominio un dominio estándar («Standard domains») con otro número de estados y cancelar borra sin aviso la tabla del nodo y la de sus hijos, que el usuario cree intactas porque ha cancelado. Lo mismo pasa al aceptar y deshacer con el menú de edición.

Por qué pasa: el botón «Standard domains» sustituye los estados del nodo (y de sus gemelos en otros ciclos) mediante \texttt{Node\allowbreak{}Replace\allowbreak{}States\allowbreak{}Edit}. Su \texttt{do\allowbreak{}Edit(\allowbreak{})} llama a \texttt{Variable.\allowbreak{}replace\allowbreak{}States} → \texttt{Variable\allowbreak{}State\allowbreak{}Operations.\allowbreak{}replace\allowbreak{}States}, que pone los estados nuevos y, si el número de estados cambia: sustituye el potencial del nodo por un \texttt{Uniform\allowbreak{}Potential}, sustituye el primer potencial de cada hijo por otro \texttt{Uniform\allowbreak{}Potential}, y \texttt{reset\allowbreak{}Link} quita las restricciones de los enlaces del nodo y las condiciones de revelación de sus enlaces a hijos. Además, si la variable es discretizada, pone siempre intervalos por defecto, aunque el número de estados no cambie.

La edición solo guarda \texttt{last\allowbreak{}States}. Su \texttt{undo(\allowbreak{})} vuelve a llamar a \texttt{replace\allowbreak{}States(\allowbreak{}node,\allowbreak{}\ last\allowbreak{}States)}, que devuelve los nombres pero, como el número de estados vuelve a cambiar, deja otra vez potenciales uniformes; nada de lo borrado vuelve. En cambio, \texttt{Node\allowbreak{}State\allowbreak{}Edit} (añadir, quitar o mover un estado) sí fotografía los potenciales, los intervalos y las restricciones, y los restaura en su \texttt{undo(\allowbreak{})}, aunque su fotografía de potenciales recoge también las tablas de los hijos (\hyperref[err:30]{error~30}).

«Cancel» llega a este \texttt{undo(\allowbreak{})} porque el diálogo del nodo deshace su sub-historia al cancelar.

Medido sobre \texttt{0\_\allowbreak{}miscellaneous/\allowbreak{}networks/\allowbreak{}bn/\allowbreak{}BN-one-disease.\allowbreak{}pgmx}, reproduciendo la secuencia del diálogo: antes, «Disease» tiene \texttt{{[}0.\allowbreak{}97,\allowbreak{}\ 0.\allowbreak{}03{]}} y «Symptom» \texttt{{[}0.\allowbreak{}98,\allowbreak{}\ 0.\allowbreak{}02,\allowbreak{}\ 0.\allowbreak{}05,\allowbreak{}\ 0.\allowbreak{}95{]}}; tras elegir \texttt{low\ -\ medium\ -\ high} los dos pasan a \texttt{Uniform\allowbreak{}Potential}; tras «Cancel» (\texttt{cancel\allowbreak{}Last\allowbreak{}Sub\allowbreak{}Edit\allowbreak{}History}) «Disease» vuelve a \texttt{absent,\allowbreak{}\ present} pero los dos siguen con \texttt{Uniform\allowbreak{}Potential}. Lo mismo tras aceptar y deshacer con el menú de edición.
\paragraph{El código}
core/src/main/java/org/openmarkov/core/action/core/NodeReplaceStatesEdit.java

\begin{Shaded}
\begin{Highlighting}[]
\KeywordTok{public} \FunctionTok{NodeReplaceStatesEdit}\OperatorTok{(}\BuiltInTok{Node}\NormalTok{ node}\OperatorTok{,} \BuiltInTok{State}\OperatorTok{[]}\NormalTok{ newStates}\OperatorTok{)} \OperatorTok{\{}
    \KeywordTok{super}\OperatorTok{(}\NormalTok{node}\OperatorTok{.}\FunctionTok{getProbNet}\OperatorTok{());}
    \KeywordTok{this}\OperatorTok{.}\FunctionTok{node} \OperatorTok{=}\NormalTok{ node}\OperatorTok{;}
    \KeywordTok{this}\OperatorTok{.}\FunctionTok{lastStates} \OperatorTok{=}\NormalTok{ node}\OperatorTok{.}\FunctionTok{getVariable}\OperatorTok{().}\FunctionTok{getStates}\OperatorTok{();}
    \KeywordTok{this}\OperatorTok{.}\FunctionTok{newStates} \OperatorTok{=}\NormalTok{ newStates}\OperatorTok{;}
\OperatorTok{\}}

\AttributeTok{@Override} \KeywordTok{protected} \DataTypeTok{void} \FunctionTok{doEdit}\OperatorTok{()} \OperatorTok{\{}
\NormalTok{    node}\OperatorTok{.}\FunctionTok{getVariable}\OperatorTok{().}\FunctionTok{replaceStates}\OperatorTok{(}\NormalTok{node}\OperatorTok{,}\NormalTok{newStates}\OperatorTok{);}
\OperatorTok{\}}

\AttributeTok{@Override} \KeywordTok{public} \DataTypeTok{void} \FunctionTok{undo}\OperatorTok{()} \OperatorTok{\{}
    \KeywordTok{super}\OperatorTok{.}\FunctionTok{undo}\OperatorTok{();}
\NormalTok{    node}\OperatorTok{.}\FunctionTok{getVariable}\OperatorTok{().}\FunctionTok{replaceStates}\OperatorTok{(}\NormalTok{node}\OperatorTok{,}\NormalTok{lastStates}\OperatorTok{);}
\OperatorTok{\}}
\end{Highlighting}
\end{Shaded}

core/src/main/java/org/openmarkov/core/model/network/VariableStateOperations.java:76-101

\begin{Shaded}
\begin{Highlighting}[]
\NormalTok{node}\OperatorTok{.}\FunctionTok{getVariable}\OperatorTok{().}\FunctionTok{setStates}\OperatorTok{(}\NormalTok{newStates}\OperatorTok{);}
\ControlFlowTok{if} \OperatorTok{(}\NormalTok{newStates}\OperatorTok{.}\FunctionTok{length} \OperatorTok{!=}\NormalTok{ lastStates}\OperatorTok{.}\FunctionTok{length}\OperatorTok{)} \OperatorTok{\{}
    \ControlFlowTok{if} \OperatorTok{(!}\NormalTok{lastPotential}\OperatorTok{.}\FunctionTok{isEmpty}\OperatorTok{())} \OperatorTok{\{}
\NormalTok{        UniformPotential newPotential }\OperatorTok{=} \KeywordTok{new} \FunctionTok{UniformPotential}\OperatorTok{(}
\NormalTok{                lastPotential}\OperatorTok{.}\FunctionTok{getFirst}\OperatorTok{().}\FunctionTok{getVariables}\OperatorTok{(),}
\NormalTok{                lastPotential}\OperatorTok{.}\FunctionTok{getFirst}\OperatorTok{().}\FunctionTok{getPotentialRole}\OperatorTok{());}
        \KeywordTok{...}
\NormalTok{        node}\OperatorTok{.}\FunctionTok{setPotentials}\OperatorTok{(}\NormalTok{newPotentials}\OperatorTok{);}
    \OperatorTok{\}}
\NormalTok{    nodes }\OperatorTok{=}\NormalTok{ probNet}\OperatorTok{.}\FunctionTok{getChildren}\OperatorTok{(}\NormalTok{node}\OperatorTok{);}
    \ControlFlowTok{for} \OperatorTok{(}\BuiltInTok{Node}\NormalTok{ child }\OperatorTok{:}\NormalTok{ nodes}\OperatorTok{)} \OperatorTok{\{}
        \KeywordTok{...}
\NormalTok{            child}\OperatorTok{.}\FunctionTok{setPotentials}\OperatorTok{(}\NormalTok{container}\OperatorTok{);}
    \OperatorTok{\}}
    \FunctionTok{resetLink}\OperatorTok{(}\NormalTok{node}\OperatorTok{);}
\OperatorTok{\}}
\end{Highlighting}
\end{Shaded}
\paragraph{Cómo arreglarlo}
Que \texttt{Node\allowbreak{}Replace\allowbreak{}States\allowbreak{}Edit} fotografíe en su constructor (o antes de \texttt{do\allowbreak{}Edit}) lo mismo que \texttt{Node\allowbreak{}State\allowbreak{}Edit}: la lista de potenciales del nodo y de cada hijo, el intervalo partido si la variable es discretizada, los potenciales de restricción de sus enlaces y las condiciones de revelación de los enlaces a sus hijos. Y que \texttt{undo(\allowbreak{})} ponga los estados antiguos directamente con \texttt{set\allowbreak{}States} y restaure todo eso, sin volver a llamar a \texttt{replace\allowbreak{}States}. Lo más limpio es compartir la fotografía y la restauración con \texttt{Node\allowbreak{}State\allowbreak{}Edit} en vez de escribirla dos veces, una vez corregida su fotografía de potenciales (\hyperref[err:30]{error~30}).

El mismo código toca el \texttt{reset\allowbreak{}Link} de \texttt{Node\allowbreak{}State\allowbreak{}Edit} y la restauración de las restricciones de enlace. Los dos mapas sin uso de \texttt{Variable\allowbreak{}State\allowbreak{}Operations.\allowbreak{}reset\allowbreak{}Link} están descritos en \hyperref[nota:5]{nota~5}; allí no hay nada que corregir.

Prueba de regresión: sobre \texttt{0\_\allowbreak{}miscellaneous/\allowbreak{}networks/\allowbreak{}bn/\allowbreak{}BN-one-disease.\allowbreak{}pgmx}, ejecutar \texttt{Node\allowbreak{}Replace\allowbreak{}States\allowbreak{}Edit} con un dominio de tres estados dentro de una sub-historia, cancelarla y exigir que los potenciales de «Disease» y «Symptom» sean los originales; repetir con un enlace con restricción y con una variable discretizada.

\setcounter{subsection}{24}
\subsection[Restringir un enlace cambia la tabla del hijo y deshacer no la devuelve]{Marcar una casilla como incompatible en la restricción de un enlace cambia la tabla del hijo fuera del historial, y ni «Cancel», ni deshacer, ni volver a pulsar la casilla la devuelven}\label{err:25}
\begin{ficha}
\fichakey{Estado}Presente
\fichakey{Módulo}gui, core
\fichakey{Paquete}org.openmarkov.gui.component; org.openmarkov.gui.action; org.openmarkov.core.model.network.potential.operation
\fichakey{Ficheros}\path{gui/src/main/java/org/openmarkov/gui/component/LinkRestrictionValuesTable.java:54-80} \newline \path{gui/src/main/java/org/openmarkov/gui/action/LinkRestrictionPotentialValueEdit.java:73-115} \newline \path{core/src/main/java/org/openmarkov/core/model/network/potential/operation/LinkRestrictionPotentialOperations.java:163-200,218-253} \newline \path{gui/src/main/java/org/openmarkov/gui/dialog/link/LinkRestrictionPanel.java:140-151} \newline \path{gui/src/main/java/org/openmarkov/gui/dialog/link/LinkRestrictionEditDialog.java:43,99-108} \newline \path{gui/src/main/java/org/openmarkov/gui/window/MainPanelListenerAssistant.java:243-255}
\fichakey{Comprobado}Leyendo el código y ejecutándolo
\fichakey{Esfuerzo}Pequeño (horas)
\fichakey{Decisión de equipo}No
\fichakey{Relacionados}\hyperref[err:24]{error~24}, \hyperref[err:27]{error~27}, \hyperref[err:154]{error~154}, \hyperref[err:155]{error~155}, \hyperref[nota:1]{nota~1}, \hyperref[nota:23]{nota~23}
\end{ficha}
\paragraph{Qué pasa}
En el menú contextual de un enlace, «Add restrictions» (o «Edit restrictions», si ya tiene) abre la ventana «Link Restriction: Link between \ldots{} and \ldots», con una casilla por cada pareja de estados del padre y del hijo. Pulsar una casilla la marca como incompatible (0) o compatible (1). Al marcarla incompatible, la ventana pone además esa casilla a cero en la tabla de probabilidad del hijo y reparte la probabilidad de la columna entre los estados que siguen permitidos. Ese cambio de la tabla no forma parte de ninguna edición:

\begin{itemize}
\tightlist
\item
  «Cancel» quita la restricción, pero la tabla del hijo se queda con el cero y la columna repartida;
\item
  lo mismo al aceptar con «OK» y deshacer después con el menú de edición;
\item
  volver a pulsar la casilla antes de aceptar la deja compatible, y al aceptar el enlace queda sin restricción, pero la tabla del hijo sigue cambiada.
\end{itemize}

Nada lo avisa: la red da otras probabilidades y se guarda así. Cada pulsación cambia además la tabla del padre, o lanza una excepción antes de llegar a la del hijo, según la casilla; eso es \hyperref[err:155]{error~155}.

Por qué pasa: la ventana (\texttt{Link\allowbreak{}Restriction\allowbreak{}Edit\allowbreak{}Dialog}) abre una sub-historia de ediciones al construirse (línea 43) y la cancela con «Cancel» (líneas 106-108). Cada pulsación llega a \texttt{Link\allowbreak{}Restriction\allowbreak{}Values\allowbreak{}Table.\allowbreak{}set\allowbreak{}Value\allowbreak{}At} (módulo \texttt{gui}), que ejecuta un \texttt{Link\allowbreak{}Restriction\allowbreak{}Potential\allowbreak{}Value\allowbreak{}Edit} y, si el valor nuevo es 0, llama a \texttt{Link\allowbreak{}Restriction\allowbreak{}Potential\allowbreak{}Operations.\allowbreak{}update\allowbreak{}Potential\allowbreak{}By\allowbreak{}Add\allowbreak{}Link\allowbreak{}Restriction} y pone el resultado en el hijo con \texttt{node2.\allowbreak{}set\allowbreak{}Potentials} (líneas 66-76). Ese método modifica en su sitio la tabla que ya tiene el hijo: pone el cero y, si el hijo es de azar, llama a \texttt{redistribute\allowbreak{}Probabilities} (líneas 244-249). \texttt{Link\allowbreak{}Restriction\allowbreak{}Potential\allowbreak{}Value\allowbreak{}Edit} solo guarda y restaura la tabla de la restricción del enlace (\texttt{last\allowbreak{}Table}, \texttt{new\allowbreak{}Table}; \texttt{undo}, líneas 107-115), así que ni cancelar la sub-historia ni deshacer tocan la tabla del hijo. Con un hijo de utilidad se pone el cero sin repartir, con el mismo problema (por lectura).

Medido con un pequeño programa de prueba que repite los pasos de la ventana con su panel (\texttt{Link\allowbreak{}Restriction\allowbreak{}Panel}) y su tabla reales, sin construir la ventana, sobre \texttt{0\_\allowbreak{}miscellaneous/\allowbreak{}networks/\allowbreak{}bn/\allowbreak{}BN-asia.\allowbreak{}pgmx} y el enlace VisitToAsia → Tuberculosis (la tabla de Tuberculosis vale \texttt{{[}0.\allowbreak{}99,\allowbreak{}\ 0.\allowbreak{}01,\allowbreak{}\ 0.\allowbreak{}95,\allowbreak{}\ 0.\allowbreak{}05{]}}: «no» y «yes» cuando VisitToAsia es «no», y lo mismo cuando es «yes»), marcando la casilla VisitToAsia = «no», Tuberculosis = «yes»:

\begin{itemize}
\tightlist
\item
  tras pulsar la casilla, la tabla de Tuberculosis es \texttt{{[}1.\allowbreak{}0,\allowbreak{}\ 0.\allowbreak{}0,\allowbreak{}\ 0.\allowbreak{}95,\allowbreak{}\ 0.\allowbreak{}05{]}};
\item
  tras «Cancel», la restricción ha desaparecido y la tabla sigue en \texttt{{[}1.\allowbreak{}0,\allowbreak{}\ 0.\allowbreak{}0,\allowbreak{}\ 0.\allowbreak{}95,\allowbreak{}\ 0.\allowbreak{}05{]}}; la probabilidad de Tuberculosis = «yes» sin hallazgos pasa de 0,0104 a 0,0005, y la de Dyspnea = «yes», de 0,4360 a 0,4324; no queda nada que deshacer;
\item
  tras «OK» y deshacer, el mismo resultado; guardada la red y reabierta, la tabla sigue en \texttt{{[}1.\allowbreak{}0,\allowbreak{}\ 0.\allowbreak{}0,\allowbreak{}\ 0.\allowbreak{}95,\allowbreak{}\ 0.\allowbreak{}05{]}};
\item
  volviendo a pulsar la casilla y aceptando, el enlace queda sin restricción y la tabla sigue en \texttt{{[}1.\allowbreak{}0,\allowbreak{}\ 0.\allowbreak{}0,\allowbreak{}\ 0.\allowbreak{}95,\allowbreak{}\ 0.\allowbreak{}05{]}}.
\end{itemize}
\paragraph{El código}
gui/src/main/java/org/openmarkov/gui/component/LinkRestrictionValuesTable.java:66-76

\begin{Shaded}
\begin{Highlighting}[]
\ControlFlowTok{if} \OperatorTok{(}\NormalTok{newNumericValue }\OperatorTok{==} \DecValTok{0}\OperatorTok{)} \OperatorTok{\{}
    \ControlFlowTok{if} \OperatorTok{(!}\NormalTok{node2}\OperatorTok{.}\FunctionTok{getPotentials}\OperatorTok{().}\FunctionTok{isEmpty}\OperatorTok{()} \OperatorTok{\&\&}\NormalTok{ node2}\OperatorTok{.}\FunctionTok{getPotentials}\OperatorTok{().}\FunctionTok{getFirst}\OperatorTok{()} \KeywordTok{instanceof}\NormalTok{ TablePotential}\OperatorTok{)} \OperatorTok{\{}
\NormalTok{        Potential potential }\OperatorTok{=}\NormalTok{ LinkRestrictionPotentialOperations}
                \OperatorTok{.}\FunctionTok{updatePotentialByAddLinkRestriction}\OperatorTok{(}\NormalTok{node2}\OperatorTok{,}
\NormalTok{                                                     link}\OperatorTok{.}\FunctionTok{getRestrictionsPotential}\OperatorTok{(),}\NormalTok{ variable1Index}\OperatorTok{,}
\NormalTok{                                                     variable2Index}\OperatorTok{);}
        \BuiltInTok{ArrayList}\OperatorTok{\textless{}}\NormalTok{Potential}\OperatorTok{\textgreater{}}\NormalTok{ potentials }\OperatorTok{=} \KeywordTok{new} \BuiltInTok{ArrayList}\OperatorTok{\textless{}\textgreater{}();}
\NormalTok{        potentials}\OperatorTok{.}\FunctionTok{add}\OperatorTok{(}\NormalTok{potential}\OperatorTok{);}
\NormalTok{        node2}\OperatorTok{.}\FunctionTok{setPotentials}\OperatorTok{(}\NormalTok{potentials}\OperatorTok{);}
    \OperatorTok{\}}
\OperatorTok{\}}
\end{Highlighting}
\end{Shaded}

gui/src/main/java/org/openmarkov/gui/action/LinkRestrictionPotentialValueEdit.java:107-115

\begin{Shaded}
\begin{Highlighting}[]
\AttributeTok{@Override} \KeywordTok{public} \DataTypeTok{void} \FunctionTok{undo}\OperatorTok{()} \OperatorTok{\{}
    \KeywordTok{super}\OperatorTok{.}\FunctionTok{undo}\OperatorTok{();}
    \ControlFlowTok{if} \OperatorTok{(!}\NormalTok{link}\OperatorTok{.}\FunctionTok{hasRestrictions}\OperatorTok{())} \OperatorTok{\{}
\NormalTok{        LinkOperations}\OperatorTok{.}\FunctionTok{initializesRestrictionsPotential}\OperatorTok{(}\NormalTok{link}\OperatorTok{);}
        \KeywordTok{this}\OperatorTok{.}\FunctionTok{tablePotential} \OperatorTok{=}\NormalTok{ link}\OperatorTok{.}\FunctionTok{getRestrictionsPotential}\OperatorTok{();}
    \OperatorTok{\}}
\NormalTok{    tablePotential}\OperatorTok{.}\FunctionTok{setValues}\OperatorTok{(}\NormalTok{lastTable}\OperatorTok{);}
    \FunctionTok{checkRestrictionPotential}\OperatorTok{(}\NormalTok{lastTable}\OperatorTok{);}
\OperatorTok{\}}
\end{Highlighting}
\end{Shaded}
\paragraph{Cómo arreglarlo}
Meter el cambio de la tabla del hijo en la edición: que \texttt{Link\allowbreak{}Restriction\allowbreak{}Potential\allowbreak{}Value\allowbreak{}Edit} (o una edición compuesta con ella) guarde una copia de la tabla del hijo antes de cambiarla, aplique \texttt{update\allowbreak{}Potential\allowbreak{}By\allowbreak{}Add\allowbreak{}Link\allowbreak{}Restriction} sobre la tabla del nodo en su \texttt{do\allowbreak{}Edit}, restaure la copia en \texttt{undo} y vuelva a poner la tabla cambiada en \texttt{redo}; y que \texttt{Link\allowbreak{}Restriction\allowbreak{}Values\allowbreak{}Table.\allowbreak{}set\allowbreak{}Value\allowbreak{}At} deje de tocar el hijo (en ese mismo método se edita también la tabla del padre: \hyperref[err:155]{error~155}). La copia es necesaria porque \texttt{update\allowbreak{}Potential\allowbreak{}By\allowbreak{}Add\allowbreak{}Link\allowbreak{}Restriction} modifica la tabla en su sitio.

Con eso, «Cancel» y deshacer devuelven la tabla, y volver a pulsar la casilla dentro de la misma ventana también, si al pulsarla otra vez se deshace la edición que la marcó en lugar de añadir otra. Una casilla que se desmarca en otra visita a la ventana no tiene tabla anterior que devolver; eso no es este fallo.

\texttt{update\allowbreak{}Potential\allowbreak{}By\allowbreak{}Add\allowbreak{}Link\allowbreak{}Restriction} solo se llama desde aquí. Al tocar esta edición conviene arreglar a la vez \hyperref[err:27]{error~27}, que es la misma clase: al rehacer escribe en un objeto enlace que puede no ser ya el de la red.

Prueba de regresión: sobre \texttt{0\_\allowbreak{}miscellaneous/\allowbreak{}networks/\allowbreak{}bn/\allowbreak{}BN-asia.\allowbreak{}pgmx}, abrir una sub-historia, marcar como incompatible (VisitToAsia = «no», Tuberculosis = «yes») con la edición y cancelar la sub-historia: la tabla de Tuberculosis debe volver a \texttt{{[}0.\allowbreak{}99,\allowbreak{}\ 0.\allowbreak{}01,\allowbreak{}\ 0.\allowbreak{}95,\allowbreak{}\ 0.\allowbreak{}05{]}}; lo mismo cerrando la sub-historia y deshaciendo; y rehacer debe dejarla en \texttt{{[}1.\allowbreak{}0,\allowbreak{}\ 0.\allowbreak{}0,\allowbreak{}\ 0.\allowbreak{}95,\allowbreak{}\ 0.\allowbreak{}05{]}}.

\setcounter{subsection}{25}
\subsection[Enlazar a una utilidad producto la convierte en suma y cambia el resultado]{Trazar un enlace a un nodo de utilidad que multiplica a sus padres lo convierte en uno que suma, y el modelo cambia de resultado}\label{err:26}
\begin{ficha}
\fichakey{Estado}Presente
\fichakey{Módulo}core
\fichakey{Paquete}org.openmarkov.core.action.base.linkEdits
\fichakey{Ficheros}\path{core/src/main/java/org/openmarkov/core/action/base/linkEdits/AddLinkEdit.java:189-233} \newline \path{core/src/main/java/org/openmarkov/core/action/base/linkEdits/RemoveLinkEdit.java:92-105} \newline \path{core/src/main/java/org/openmarkov/core/action/core/VariableTypeEdit.java:157-170} \newline \path{core/src/main/java/org/openmarkov/core/model/network/Node.java:604-617} \newline \path{core/src/main/java/org/openmarkov/core/action/base/linkEdits/MultiAddLinkEdit.java:30-51}
\fichakey{Comprobado}Leyendo el código y ejecutándolo
\fichakey{Esfuerzo}Pequeño (horas)
\fichakey{Decisión de equipo}No
\fichakey{Relacionados}\hyperref[err:11]{error~11}
\end{ficha}
\paragraph{Qué pasa}
Un nodo de utilidad cuyos padres son todos numéricos (un «supervalor») no guarda números propios: su potencial dice cómo combinar los de los padres, \texttt{Sum\allowbreak{}Potential} (suma) o \texttt{Product\allowbreak{}Potential} (producto). Si el usuario le traza un enlace a un nodo producto, el nodo pasa a sumar, y el modelo da otro resultado sin ningún aviso. Deshacer lo devuelve a producto, pero rehacer lo vuelve a sumar, y quitar después ese mismo enlace no lo devuelve a producto.

Camino de llegada: en modo edición, dibujar un enlace desde un nodo numérico hacia un nodo de utilidad producto (la ventana usa \texttt{Multi\allowbreak{}Add\allowbreak{}Link\allowbreak{}Edit}, que construye un \texttt{Add\allowbreak{}Link\allowbreak{}Edit} por par). En el repositorio hay nodos producto con padres solo numéricos en \texttt{0\_\allowbreak{}miscellaneous/\allowbreak{}networks/\allowbreak{}id/\allowbreak{}ID-mediastinet.\allowbreak{}pgmx} (\texttt{Net\_\allowbreak{}QALE}, \texttt{Weighted\_\allowbreak{}Economic\_\allowbreak{}Cost}), \texttt{0\_\allowbreak{}miscellaneous/\allowbreak{}networks/\allowbreak{}dan/\allowbreak{}DAN-qale-mediastinet.\allowbreak{}pgmx} y \texttt{0\_\allowbreak{}miscellaneous/\allowbreak{}networks/\allowbreak{}dan/\allowbreak{}DAN-mediastinet.\allowbreak{}pgmx} (\texttt{Net\ QALE}), \texttt{0\_\allowbreak{}miscellaneous/\allowbreak{}networks/\allowbreak{}id/\allowbreak{}ID-arthronet.\allowbreak{}pgmx} y \texttt{0\_\allowbreak{}miscellaneous/\allowbreak{}networks/\allowbreak{}dan/\allowbreak{}DAN-arthronet.\allowbreak{}pgmx} («Coste ajustado»), \texttt{0\_\allowbreak{}miscellaneous/\allowbreak{}networks/\allowbreak{}mid/\allowbreak{}MID-HPV.\allowbreak{}pgmx} (\texttt{Qo\allowbreak{}L\ {[}0{]}}) y varias redes de prueba de redes de análisis de decisiones.

Por qué pasa: al trazar el enlace, \texttt{Add\allowbreak{}Link\allowbreak{}Edit.\allowbreak{}do\allowbreak{}Edit(\allowbreak{})} (módulo \texttt{core}), después de añadir el enlace, comprueba \texttt{node\allowbreak{}To.\allowbreak{}get\allowbreak{}Node\allowbreak{}Type(\allowbreak{})\ ==\ UTILITY\ \&\&\ node\allowbreak{}To.\allowbreak{}only\allowbreak{}Numerical\allowbreak{}Parents(\allowbreak{})} y, para cada potencial antiguo, crea \texttt{new\ Sum\allowbreak{}Potential(\allowbreak{}variables,\allowbreak{}\ rol)}, sea cual sea la clase que tenía. Un producto pasa a ser una suma. Lo mismo ocurre con cualquier otra clase que tenga el nodo en ese momento (en las redes del repositorio hay nodos así con \texttt{Tree\allowbreak{}ADDPotential}, \texttt{Function\allowbreak{}Potential} y otros).

\texttt{Remove\allowbreak{}Link\allowbreak{}Edit.\allowbreak{}do\allowbreak{}Edit(\allowbreak{})} y \texttt{Variable\allowbreak{}Type\allowbreak{}Edit.\allowbreak{}set\allowbreak{}Potentials\allowbreak{}Node\allowbreak{}And\allowbreak{}Children} repiten el mismo supuesto (crean \texttt{Sum\allowbreak{}Potential} para un hijo de utilidad con padres numéricos), así que quitar un padre a un nodo producto, o cambiar el tipo de la variable de un padre, también lo convierte en suma (leído, no ejecutado para estas dos vías).

Medido sobre \texttt{0\_\allowbreak{}miscellaneous/\allowbreak{}networks/\allowbreak{}id/\allowbreak{}ID-mediastinet.\allowbreak{}pgmx}: \texttt{Net\_\allowbreak{}QALE} es \texttt{Product\allowbreak{}Potential} de tres padres; con \texttt{Add\allowbreak{}Link\allowbreak{}Edit(\allowbreak{}TBNA\_\allowbreak{}Morbidity\ →\ Net\_\allowbreak{}QALE)} pasa a \texttt{Sum\allowbreak{}Potential} de cuatro; deshacer lo devuelve a producto, pero rehacer lo vuelve a sumar; y quitando después ese mismo enlace con \texttt{Remove\allowbreak{}Link\allowbreak{}Edit} queda \texttt{Sum\allowbreak{}Potential} de los tres padres originales. La utilidad esperada global (evaluación por eliminación de variables) pasa de 3,984 a 6,059 sin ningún aviso.
\paragraph{El código}
core/src/main/java/org/openmarkov/core/action/base/linkEdits/AddLinkEdit.java:189-204

\begin{Shaded}
\begin{Highlighting}[]
\AttributeTok{@Override} \KeywordTok{protected} \DataTypeTok{void} \FunctionTok{doEdit}\OperatorTok{()} \OperatorTok{\{}
\NormalTok{    probNet}\OperatorTok{.}\FunctionTok{addLink}\OperatorTok{(}\NormalTok{nodeFrom}\OperatorTok{,}\NormalTok{ nodeTo}\OperatorTok{,}\NormalTok{ isDirected}\OperatorTok{);}
    \KeywordTok{this}\OperatorTok{.}\FunctionTok{link} \OperatorTok{=}\NormalTok{ probNet}\OperatorTok{.}\FunctionTok{getLink}\OperatorTok{(}\NormalTok{nodeFrom}\OperatorTok{,}\NormalTok{ nodeTo}\OperatorTok{,}\NormalTok{ isDirected}\OperatorTok{);}
    \ControlFlowTok{if} \OperatorTok{(}\NormalTok{updatePotentials}\OperatorTok{)} \OperatorTok{\{}
        \KeywordTok{this}\OperatorTok{.}\FunctionTok{oldPotentials} \OperatorTok{=}\NormalTok{ nodeTo}\OperatorTok{.}\FunctionTok{getPotentials}\OperatorTok{();}
        \CommentTok{// }\AlertTok{TODO}\CommentTok{ Check if this UTILITY label is outdated}
        \ControlFlowTok{if} \OperatorTok{(}\NormalTok{nodeTo}\OperatorTok{.}\FunctionTok{getNodeType}\OperatorTok{()} \OperatorTok{==}\NormalTok{ NodeType}\OperatorTok{.}\FunctionTok{UTILITY} \OperatorTok{\&\&}\NormalTok{ nodeTo}\OperatorTok{.}\FunctionTok{onlyNumericalParents}\OperatorTok{())} \OperatorTok{\{}
            \CommentTok{// Add a default Sum potential to utility supervalue nodes}
            \ControlFlowTok{for} \OperatorTok{(}\NormalTok{Potential oldPotential }\OperatorTok{:}\NormalTok{ oldPotentials}\OperatorTok{)} \OperatorTok{\{}
                \CommentTok{// Update potential}
                \BuiltInTok{List}\OperatorTok{\textless{}}\NormalTok{Variable}\OperatorTok{\textgreater{}}\NormalTok{ variables }\OperatorTok{=}\NormalTok{ oldPotential}\OperatorTok{.}\FunctionTok{getVariables}\OperatorTok{();}
                \ControlFlowTok{if} \OperatorTok{(!}\NormalTok{variables}\OperatorTok{.}\FunctionTok{contains}\OperatorTok{(}\NormalTok{nodeFrom}\OperatorTok{.}\FunctionTok{getVariable}\OperatorTok{()))} \OperatorTok{\{}
\NormalTok{                    variables}\OperatorTok{.}\FunctionTok{add}\OperatorTok{(}\NormalTok{nodeFrom}\OperatorTok{.}\FunctionTok{getVariable}\OperatorTok{());}
                \OperatorTok{\}}
\NormalTok{                Potential newPotential }\OperatorTok{=} \KeywordTok{new} \FunctionTok{SumPotential}\OperatorTok{(}\NormalTok{variables}\OperatorTok{,}\NormalTok{ oldPotential}\OperatorTok{.}\FunctionTok{getPotentialRole}\OperatorTok{());}
\end{Highlighting}
\end{Shaded}
\paragraph{Cómo arreglarlo}
Conservar la clase del potencial: si es \texttt{Product\allowbreak{}Potential}, crear un \texttt{Product\allowbreak{}Potential} con la variable añadida (o usar \texttt{add\allowbreak{}Variable} del propio potencial si da el resultado adecuado); usar \texttt{Sum\allowbreak{}Potential} solo cuando el nodo no tenía ya un combinador (por ejemplo, venía de una tabla o de un potencial uniforme). Hay que decidir qué hacer con clases que no son suma ni producto (\texttt{Function\allowbreak{}Potential}, \texttt{Tree\allowbreak{}ADDPotential}): mantenerlas y añadir la variable, o avisar.

La misma regla debe cumplirse en \texttt{Add\allowbreak{}Link\allowbreak{}Edit.\allowbreak{}do\allowbreak{}Edit}, \texttt{Remove\allowbreak{}Link\allowbreak{}Edit.\allowbreak{}do\allowbreak{}Edit} (quitar la variable manteniendo la clase) y \texttt{Variable\allowbreak{}Type\allowbreak{}Edit.\allowbreak{}set\allowbreak{}Potentials\allowbreak{}Node\allowbreak{}And\allowbreak{}Children}.

Prueba: en \texttt{0\_\allowbreak{}miscellaneous/\allowbreak{}networks/\allowbreak{}id/\allowbreak{}ID-mediastinet.\allowbreak{}pgmx}, añadir y luego quitar un padre numérico a \texttt{Net\_\allowbreak{}QALE} debe dejarlo como \texttt{Product\allowbreak{}Potential} de sus tres padres y con la misma utilidad esperada que al principio.

\setcounter{subsection}{26}
\subsection[Rehacer un enlace y su restricción deja el enlace sin la restricción]{Rehacer un enlace trazado y la restricción o la revelación que se le puso después deja el enlace sin ellas: se escriben en un enlace que ya no está en la red}\label{err:27}
\begin{ficha}
\fichakey{Estado}Presente
\fichakey{Módulo}core; gui
\fichakey{Paquete}org.openmarkov.core.action.base.linkEdits; org.openmarkov.core.action.core; org.openmarkov.gui.action
\fichakey{Ficheros}\path{core/src/main/java/org/openmarkov/core/action/base/linkEdits/AddLinkEdit.java:254-264} \newline \path{core/src/main/java/org/openmarkov/core/action/base/linkEdits/RemoveLinkEdit.java:171-187} \newline \path{gui/src/main/java/org/openmarkov/gui/action/LinkRestrictionPotentialValueEdit.java:73-115} \newline \path{core/src/main/java/org/openmarkov/core/action/core/RevelationStateEdit.java:28-52} \newline \path{gui/src/main/java/org/openmarkov/gui/action/RevelationIntervalEdit.java:67-69} \newline \path{core/src/main/java/org/openmarkov/core/action/core/VariableTypeEdit.java:271-307} \newline \path{gui/src/main/java/org/openmarkov/gui/window/MainPanelListenerAssistant.java:243-255} \newline \path{gui/src/main/java/org/openmarkov/gui/graphic/VisualLink.java:142}
\fichakey{Comprobado}Leyendo el código y ejecutándolo
\fichakey{Esfuerzo}Medio (días)
\fichakey{Decisión de equipo}No
\fichakey{Relacionados}\hyperref[err:25]{error~25}, \hyperref[err:155]{error~155}, \hyperref[nota:23]{nota~23}, \hyperref[nota:150]{nota~150}
\end{ficha}
\paragraph{Qué pasa}
Trazar un enlace A → B, ponerle una restricción (en su menú contextual, «Edit link restriction», marcar una casilla como incompatible y «OK»), deshacer las dos cosas con el menú de edición y rehacerlas deja el enlace A → B sin restricción. Rehacer no avisa ni falla: escribe la restricción en el objeto enlace que existía antes de deshacer, que ya no está en la red. Lo mismo ocurre con las condiciones de revelación de un enlace (qué estados del padre revelan al hijo), que se editan en otra ventanilla y con otras ediciones.

Por qué pasa: \texttt{Add\allowbreak{}Link\allowbreak{}Edit.\allowbreak{}redo} (módulo \texttt{core}) vuelve a llamar a \texttt{prob\allowbreak{}Net.\allowbreak{}add\allowbreak{}Link}, que crea un objeto \texttt{Link} nuevo, y actualiza su propio campo \texttt{link}. Pero las ediciones que se hicieron después sobre ese enlace guardan una referencia al objeto de su momento: \texttt{Link\allowbreak{}Restriction\allowbreak{}Potential\allowbreak{}Value\allowbreak{}Edit} (módulo \texttt{gui}) guarda en el constructor el \texttt{Link} y su tabla de restricciones, y \texttt{Revelation\allowbreak{}State\allowbreak{}Edit} y \texttt{Revelation\allowbreak{}Interval\allowbreak{}Edit} guardan también el \texttt{Link}. Al rehacerse, escriben en el objeto viejo.

\texttt{Remove\allowbreak{}Link\allowbreak{}Edit.\allowbreak{}undo} también crea un objeto enlace nuevo, pero le pasa las restricciones y revelaciones del viejo, que comparten la misma tabla, así que ahí el daño es menor. Medido: poner una restricción, quitar el enlace, deshacer y deshacer después la restricción deja en el enlace de la red una tabla de restricciones con todo a 1, en lugar de ninguna tabla. Con esa tabla el enlace cuenta como restringido: por lectura, se dibuja rayado (\texttt{Visual\allowbreak{}Link}, línea 142) y se guarda con restricciones.

Medido con un pequeño programa de prueba que repite las ediciones de la ventana (el \texttt{Multi\allowbreak{}Add\allowbreak{}Link\allowbreak{}Edit} que crea el trazado, la sub-historia del diálogo de restricción con un \texttt{Link\allowbreak{}Restriction\allowbreak{}Potential\allowbreak{}Value\allowbreak{}Edit}, y deshacer y rehacer con \texttt{PNESupport}) sobre una red bayesiana con dos nodos A y B de dos estados:

\begin{itemize}
\tightlist
\item
  tras poner la restricción, el enlace de la red tiene la tabla \texttt{{[}1,\allowbreak{}\ 1,\allowbreak{}\ 0,\allowbreak{}\ 1{]}};
\item
  tras deshacer dos veces no hay enlace; tras rehacer dos veces hay un enlace distinto del original y sin restricciones, y es el objeto viejo el que tiene \texttt{{[}1,\allowbreak{}\ 1,\allowbreak{}\ 0,\allowbreak{}\ 1{]}};
\item
  en una red de análisis de decisiones, con \texttt{Revelation\allowbreak{}State\allowbreak{}Edit} para que el estado «a1» de A revele B: tras deshacer y rehacer las dos ediciones, el enlace de la red no revela nada y el viejo revela «a1».
\end{itemize}
\paragraph{El código}
core/src/main/java/org/openmarkov/core/action/base/linkEdits/AddLinkEdit.java:254-264

\begin{Shaded}
\begin{Highlighting}[]
\AttributeTok{@Override} \KeywordTok{public} \DataTypeTok{void} \FunctionTok{redo}\OperatorTok{()} \OperatorTok{\{}
    \FunctionTok{setTypicalRedo}\OperatorTok{(}\KeywordTok{false}\OperatorTok{);}
    \KeywordTok{super}\OperatorTok{.}\FunctionTok{redo}\OperatorTok{();}

\NormalTok{    nodeTo }\OperatorTok{=}\NormalTok{ probNet}\OperatorTok{.}\FunctionTok{getNode}\OperatorTok{(}\NormalTok{variableTo}\OperatorTok{.}\FunctionTok{getName}\OperatorTok{());}
\NormalTok{    probNet}\OperatorTok{.}\FunctionTok{addLink}\OperatorTok{(}\NormalTok{nodeFrom}\OperatorTok{,}\NormalTok{ nodeTo}\OperatorTok{,}\NormalTok{ isDirected}\OperatorTok{);}
    \KeywordTok{this}\OperatorTok{.}\FunctionTok{link} \OperatorTok{=}\NormalTok{ probNet}\OperatorTok{.}\FunctionTok{getLink}\OperatorTok{(}\NormalTok{nodeFrom}\OperatorTok{,}\NormalTok{ nodeTo}\OperatorTok{,}\NormalTok{ isDirected}\OperatorTok{);}
    \ControlFlowTok{if} \OperatorTok{(}\NormalTok{updatePotentials}\OperatorTok{)} \OperatorTok{\{}
\NormalTok{        nodeTo}\OperatorTok{.}\FunctionTok{setPotentials}\OperatorTok{(}\NormalTok{newPotentials}\OperatorTok{);}
    \OperatorTok{\}}
\OperatorTok{\}}
\end{Highlighting}
\end{Shaded}

gui/src/main/java/org/openmarkov/gui/action/LinkRestrictionPotentialValueEdit.java:73-105

\begin{Shaded}
\begin{Highlighting}[]
\KeywordTok{public} \FunctionTok{LinkRestrictionPotentialValueEdit}\OperatorTok{(}\NormalTok{Link}\OperatorTok{\textless{}}\BuiltInTok{Node}\OperatorTok{\textgreater{}}\NormalTok{ link}\OperatorTok{,} \BuiltInTok{Integer}\NormalTok{ newValue}\OperatorTok{,} \DataTypeTok{int}\NormalTok{ row}\OperatorTok{,} \DataTypeTok{int}\NormalTok{ col}\OperatorTok{)} \OperatorTok{\{}
    \KeywordTok{super}\OperatorTok{(}\NormalTok{link}\OperatorTok{.}\FunctionTok{getFrom}\OperatorTok{().}\FunctionTok{getProbNet}\OperatorTok{());}
    \KeywordTok{this}\OperatorTok{.}\FunctionTok{link} \OperatorTok{=}\NormalTok{ link}\OperatorTok{;}
    \KeywordTok{...}
    \KeywordTok{this}\OperatorTok{.}\FunctionTok{tablePotential} \OperatorTok{=}\NormalTok{ link}\OperatorTok{.}\FunctionTok{getRestrictionsPotential}\OperatorTok{();}
    \KeywordTok{...}
\OperatorTok{\}}
\KeywordTok{...}
\AttributeTok{@Override} \KeywordTok{public} \DataTypeTok{void} \FunctionTok{redo}\OperatorTok{()} \OperatorTok{\{}
    \KeywordTok{this}\OperatorTok{.}\FunctionTok{setTypicalRedo}\OperatorTok{(}\KeywordTok{false}\OperatorTok{);}
    \KeywordTok{super}\OperatorTok{.}\FunctionTok{redo}\OperatorTok{();}
    \ControlFlowTok{if} \OperatorTok{(!}\NormalTok{link}\OperatorTok{.}\FunctionTok{hasRestrictions}\OperatorTok{())} \OperatorTok{\{}
\NormalTok{        LinkOperations}\OperatorTok{.}\FunctionTok{initializesRestrictionsPotential}\OperatorTok{(}\NormalTok{link}\OperatorTok{);}
        \KeywordTok{this}\OperatorTok{.}\FunctionTok{tablePotential} \OperatorTok{=}\NormalTok{ link}\OperatorTok{.}\FunctionTok{getRestrictionsPotential}\OperatorTok{();}
    \OperatorTok{\}}
\NormalTok{    tablePotential}\OperatorTok{.}\FunctionTok{setValues}\OperatorTok{(}\NormalTok{newTable}\OperatorTok{);}
    \FunctionTok{checkRestrictionPotential}\OperatorTok{(}\NormalTok{newTable}\OperatorTok{);}
\OperatorTok{\}}
\end{Highlighting}
\end{Shaded}
\paragraph{Cómo arreglarlo}
Dos caminos:

\begin{itemize}
\tightlist
\item
  que las ediciones que actúan sobre un enlace guarden sus extremos (y si es dirigido) en vez del objeto \texttt{Link}, y lo busquen en la red cada vez que hacen, deshacen o rehacen (\texttt{prob\allowbreak{}Net.\allowbreak{}get\allowbreak{}Link(\allowbreak{}nodo1,\allowbreak{}\ nodo2,\allowbreak{}\ dirigido)}). Sitios: \texttt{Link\allowbreak{}Restriction\allowbreak{}Potential\allowbreak{}Value\allowbreak{}Edit} (también su campo \texttt{table\allowbreak{}Potential}), \texttt{Revelation\allowbreak{}State\allowbreak{}Edit}, \texttt{Revelation\allowbreak{}Interval\allowbreak{}Edit} y las ediciones internas \texttt{Raw\allowbreak{}Set\allowbreak{}Revelation} y \texttt{Raw\allowbreak{}Set\allowbreak{}Restriction} de \texttt{Variable\allowbreak{}Type\allowbreak{}Edit};
\item
  que \texttt{Add\allowbreak{}Link\allowbreak{}Edit.\allowbreak{}redo} y \texttt{Remove\allowbreak{}Link\allowbreak{}Edit.\allowbreak{}undo} vuelvan a meter en el grafo el mismo objeto \texttt{Link} que quitaron, en vez de crear otro; esto arregla todas las ediciones a la vez, pero exige que el grafo acepte reinsertar un enlace ya construido.
\end{itemize}

Conviene revisar con el mismo criterio las demás ediciones que quitan y vuelven a crear enlaces.

Prueba de regresión: la secuencia medida (trazar, restringir, deshacer dos veces, rehacer dos veces) debe dejar en el enlace de la red la tabla \texttt{{[}1,\allowbreak{}\ 1,\allowbreak{}\ 0,\allowbreak{}\ 1{]}}; lo mismo con una revelación; y restringir, quitar el enlace, deshacer y deshacer la restricción debe dejar el enlace sin tabla de restricciones.

\setcounter{subsection}{27}
\subsection[Pasar a suceso una decisión numérica sin política la deja a medias]{Cambiar a suceso un nodo de decisión numérico sin política revienta y deja el nodo convertido a medias, sin poder deshacerlo}\label{err:28}
\begin{ficha}
\fichakey{Estado}Presente
\fichakey{Módulo}core
\fichakey{Paquete}org.openmarkov.core.action.core
\fichakey{Ficheros}\path{core/src/main/java/org/openmarkov/core/action/core/VariableTypeEdit.java:80-135} \newline \path{core/src/main/java/org/openmarkov/core/action/base/MultiStepEdit.java:36-48} \newline \path{gui/src/main/java/org/openmarkov/gui/action/ChangeNodeTypeEdit.java:47-62} \newline \path{gui/src/main/java/org/openmarkov/gui/dialog/node/NodeDefinitionPanel.java:332-342} \newline \path{core/src/main/java/org/openmarkov/core/model/network/VariableType.java:58-65}
\fichakey{Comprobado}Leyendo el código y ejecutándolo
\fichakey{Esfuerzo}Pequeño (horas)
\fichakey{Decisión de equipo}No
\fichakey{Relacionados}\hyperref[err:31]{error~31}, \hyperref[nota:184]{nota~184}, \hyperref[nota:142]{nota~142}
\end{ficha}
\paragraph{Qué pasa}
Cambiar a suceso un nodo numérico que no tiene potenciales termina en el diálogo de «error no previsto», y el nodo queda convertido a medias: tipo de nodo \texttt{event}, variable de tipo \texttt{Event} y un solo estado. El cambio no entra en el historial, así que no se puede deshacer, y cancelar el diálogo de propiedades no lo repara porque el cambio no está en su historial provisional.

Camino del usuario, en una red de simulación de eventos discretos (la única en la que el botón «Event» está habilitado): crear un nodo de decisión; en su dominio elegir «Numeric» (el tipo de variable ofrece Numeric para decisiones, y eso le pone un potencial uniforme); quitarle la política con «quitar política» (\texttt{Remove\allowbreak{}Policy\allowbreak{}Edit} vacía sus potenciales); abrir sus propiedades y pulsar el tipo «Event». También llega un nodo numérico de azar o de decisión sin potencial leído de un fichero.

Por qué pasa: \texttt{Change\allowbreak{}Node\allowbreak{}Type\allowbreak{}Edit} (el botón de tipo de nodo en la pestaña de definición del diálogo de propiedades del nodo) es una edición de varios pasos. Primero cambia el tipo de nodo (\texttt{Set\allowbreak{}Node\allowbreak{}Type\allowbreak{}Edit}) y, si el tipo de variable ya no vale para el tipo nuevo, ejecuta \texttt{Variable\allowbreak{}Type\allowbreak{}Edit(\allowbreak{}node,\allowbreak{}\ tipo,\allowbreak{}\ false)}, es decir, sin recalcular el potencial del propio nodo. Un nodo de suceso solo admite variables de tipo EVENT (\texttt{Variable\allowbreak{}Type.\allowbreak{}of}), así que pasar un nodo numérico a suceso entra en esa rama. \texttt{Variable\allowbreak{}Type\allowbreak{}Edit}, cuando el tipo original es NUMERIC, construye un potencial uniforme con el papel del primer potencial del nodo: \texttt{node.\allowbreak{}get\allowbreak{}Potentials(\allowbreak{}).\allowbreak{}get\allowbreak{}First(\allowbreak{}).\allowbreak{}get\allowbreak{}Potential\allowbreak{}Role(\allowbreak{})} (línea 127). Si el nodo no tiene potenciales, \texttt{get\allowbreak{}First(\allowbreak{})} lanza \texttt{No\allowbreak{}Such\allowbreak{}Element\allowbreak{}Exception}.

Por qué el daño es mayor que una excepción: \texttt{Multi\allowbreak{}Step\allowbreak{}Edit.\allowbreak{}do\allowbreak{}Edit(\allowbreak{})} solo revierte los pasos ya ejecutados si captura \texttt{Do\allowbreak{}Edit\allowbreak{}Exception}; una excepción de tiempo de ejecución la atraviesa sin revertir nada. El primer paso de \texttt{Variable\allowbreak{}Type\allowbreak{}Edit} (\texttt{Raw\allowbreak{}Set\allowbreak{}Type}, que pone el tipo de variable nuevo y cambia los estados) y el \texttt{Set\allowbreak{}Node\allowbreak{}Type\allowbreak{}Edit} de fuera ya se han aplicado, y como la edición no termina, no entra en el historial.

Medido ejecutando esa misma secuencia de ediciones sin ventana: salta \texttt{No\allowbreak{}Such\allowbreak{}Element\allowbreak{}Exception} en \texttt{Variable\allowbreak{}Type\allowbreak{}Edit.\allowbreak{}java:127}; después el nodo queda con tipo \texttt{event}, variable de tipo \texttt{Event} y un solo estado, y el historial tiene las mismas tres entradas que antes. En la aplicación, la excepción llega desde el botón del diálogo al manejador global (el diálogo de «error no previsto»).
\paragraph{El código}
core/src/main/java/org/openmarkov/core/action/core/VariableTypeEdit.java:114-130

\begin{Shaded}
\begin{Highlighting}[]
\ControlFlowTok{if} \OperatorTok{(}\NormalTok{originalVariableType }\OperatorTok{==}\NormalTok{ VariableType}\OperatorTok{.}\FunctionTok{NUMERIC}\OperatorTok{)} \OperatorTok{\{}
    \CommentTok{// Build the list of variables (node and its parents)}
    \BuiltInTok{List}\OperatorTok{\textless{}}\NormalTok{Variable}\OperatorTok{\textgreater{}}\NormalTok{ variables }\OperatorTok{=} \KeywordTok{new} \BuiltInTok{ArrayList}\OperatorTok{\textless{}\textgreater{}();}
    \KeywordTok{...}
    \CommentTok{// Create and assign uniform potential}
\NormalTok{    UniformPotential uniformPotential }\OperatorTok{=} \KeywordTok{new} \FunctionTok{UniformPotential}\OperatorTok{(}
\NormalTok{            variables}\OperatorTok{,}
\NormalTok{            node}\OperatorTok{.}\FunctionTok{getPotentials}\OperatorTok{().}\FunctionTok{getFirst}\OperatorTok{().}\FunctionTok{getPotentialRole}\OperatorTok{()}
    \OperatorTok{);}
\NormalTok{    stepExecuter}\OperatorTok{.}\FunctionTok{execute}\OperatorTok{(}\KeywordTok{new} \FunctionTok{RawSetPotentialEdit}\OperatorTok{(}\NormalTok{node}\OperatorTok{,}\NormalTok{ uniformPotential}\OperatorTok{));}
\end{Highlighting}
\end{Shaded}

core/src/main/java/org/openmarkov/core/action/base/MultiStepEdit.java:36-48

\begin{Shaded}
\begin{Highlighting}[]
\AttributeTok{@Override} \KeywordTok{protected} \DataTypeTok{final} \DataTypeTok{void} \FunctionTok{doEdit}\OperatorTok{()} \KeywordTok{throws}\NormalTok{ DoEditException }\OperatorTok{\{}
\NormalTok{    StepExecuter stepExecuter }\OperatorTok{=} \KeywordTok{new} \FunctionTok{StepExecuter}\OperatorTok{();}
    \ControlFlowTok{try} \OperatorTok{\{}
        \KeywordTok{this}\OperatorTok{.}\FunctionTok{edits}\OperatorTok{.}\FunctionTok{clear}\OperatorTok{();}
        \KeywordTok{this}\OperatorTok{.}\FunctionTok{doMultiStepEdit}\OperatorTok{(}\NormalTok{stepExecuter}\OperatorTok{);}
        \KeywordTok{this}\OperatorTok{.}\FunctionTok{edits}\OperatorTok{.}\FunctionTok{addAll}\OperatorTok{(}\NormalTok{stepExecuter}\OperatorTok{.}\FunctionTok{executedEdits}\OperatorTok{);}
    \OperatorTok{\}} \ControlFlowTok{catch} \OperatorTok{(}\NormalTok{DoEditException e}\OperatorTok{)} \OperatorTok{\{}
        \ControlFlowTok{for} \OperatorTok{(}\NormalTok{PNEdit editToUndo }\OperatorTok{:}\NormalTok{ stepExecuter}\OperatorTok{.}\FunctionTok{executedEdits}\OperatorTok{.}\FunctionTok{reversed}\OperatorTok{())} \OperatorTok{\{}
\NormalTok{            editToUndo}\OperatorTok{.}\FunctionTok{undo}\OperatorTok{();}
        \OperatorTok{\}}
        \ControlFlowTok{throw}\NormalTok{ e}\OperatorTok{;}
    \OperatorTok{\}}
\OperatorTok{\}}
\end{Highlighting}
\end{Shaded}
\paragraph{Cómo arreglarlo}
Hay dos piezas.

La local: el papel del potencial no debe salir de un potencial que puede no existir. Se puede deducir del tipo de nodo (como hace \texttt{Potential\allowbreak{}Operations.\allowbreak{}get\allowbreak{}Uniform\allowbreak{}Potential}, que usa POLICY para decisiones y CONDITIONAL\_PROBABILITY en otro caso) o comprobar la lista vacía. Hay que decidir también qué potencial debe tener un nodo de suceso (probablemente el que genere \texttt{Potential\allowbreak{}Utils.\allowbreak{}generate\allowbreak{}Default\allowbreak{}Potential}, que \texttt{Change\allowbreak{}Node\allowbreak{}Type\allowbreak{}Edit} ya aplica después).

La de la maquinaria: que la reversión de \texttt{Multi\allowbreak{}Step\allowbreak{}Edit} se haga ante cualquier excepción (\texttt{catch} de \texttt{Runtime\allowbreak{}Exception} además de \texttt{Do\allowbreak{}Edit\allowbreak{}Exception}, relanzando la original), y que \texttt{Step\allowbreak{}Executer.\allowbreak{}execute} apunte la edición antes de ejecutarla para que un paso que falle a medias también se revierta. La misma reversión falta en \texttt{Multi\allowbreak{}Edit.\allowbreak{}do\allowbreak{}Edit(\allowbreak{})}. El problema general de diseño, que una edición que modifica la red y después falla no se revierte, está descrito en \hyperref[nota:142]{nota~142}; allí no hay nada que corregir.

Pruebas:

\begin{itemize}
\tightlist
\item
  Red de simulación de eventos discretos, nodo de decisión numérico sin potenciales, \texttt{Change\allowbreak{}Node\allowbreak{}Type\allowbreak{}Edit} a EVENT: exigir que termine (o lance una excepción del dominio) y que, si falla, el nodo siga siendo de decisión y numérico.
\item
  De la maquinaria: un \texttt{Multi\allowbreak{}Step\allowbreak{}Edit} cuyo segundo paso lance una excepción no comprobada debe dejar la red como estaba.
\end{itemize}

\setcounter{subsection}{28}
\subsection[Mover un estado con un hijo sin potencial deja la tabla permutada a medias]{Mover un estado de un nodo cuyo hijo no tiene potencial revienta y deja la tabla del nodo permutada con los estados en el orden viejo}\label{err:29}
\begin{ficha}
\fichakey{Estado}Presente en parte. Desde \texttt{2c4fca6} (ajustado en \texttt{72d422b}), un nodo de azar o de utilidad sin potencial se salta al reordenar, así que mover un estado ya no lanza \texttt{NoSuchElementException} ni deja la tabla permutada con los estados en el orden viejo. Sigue presente el segundo efecto: solo se reordena el primer potencial del nodo y se instala con \texttt{setPotential}, que vacía la lista, de modo que un nodo con dos potenciales pierde los demás.
\fichakey{Módulo}core
\fichakey{Paquete}org.openmarkov.core.model.network
\fichakey{Ficheros}\path{core/src/main/java/org/openmarkov/core/model/network/VariableStateOperations.java:200-236} \newline \path{core/src/main/java/org/openmarkov/core/action/core/NodeStateEdit.java:154-163} \newline \path{core/src/main/java/org/openmarkov/core/model/network/Node.java:215-220}
\fichakey{Comprobado}Leyendo el código y ejecutándolo
\fichakey{Esfuerzo}Pequeño (horas)
\fichakey{Decisión de equipo}No
\fichakey{Relacionados}\hyperref[err:30]{error~30}, \hyperref[err:70]{error~70}
\end{ficha}
\paragraph{Qué pasa}
Subir o bajar un estado de un nodo cuyo hijo (o él mismo) es de azar o de utilidad y no tiene potencial termina en un error inesperado. Lo grave es cómo queda la red: la tabla del nodo ya se ha reordenado, pero sus estados siguen en el orden viejo, así que las probabilidades quedan intercambiadas en silencio. La edición no entra en el historial, así que ni deshacer ni «Cancelar» en el diálogo la revierten.

Camino de llegada de la parte ya corregida: un nodo de azar o de utilidad sin potencial. Desde los commits \texttt{710de2c} y \texttt{7776e28}, un nodo recién creado no recibe potencial hasta que se abre su relación; además, el lector de \texttt{.\allowbreak{}pgmx} acepta un nodo al que le falta su \texttt{\textless{}Potential\textgreater{}} sin añadir nada: en el repositorio lo hace \texttt{io/\allowbreak{}src/\allowbreak{}test/\allowbreak{}resources/\allowbreak{}Imposed\allowbreak{}Policy.\allowbreak{}pgmx} (nodo de utilidad U). Basta abrir un fichero así y mover un estado del nodo o de uno de sus padres en la pestaña de dominio. Lo que pasa con un nodo numérico de azar sin potencial al propagar o evaluar está en \hyperref[err:43]{error~43}.

Por qué pasa: \texttt{Variable\allowbreak{}State\allowbreak{}Operations.\allowbreak{}reorder\allowbreak{}States\allowbreak{}In\allowbreak{}Node\allowbreak{}And\allowbreak{}Children} (módulo \texttt{core}) primero pregunta a todos los potenciales si aceptan el reordenado (\texttt{refuse\allowbreak{}If\allowbreak{}Any\allowbreak{}Potential\allowbreak{}Cannot\allowbreak{}Be\allowbreak{}Reordered}: esa parte ya es transaccional), luego llama a \texttt{set\allowbreak{}Potential\allowbreak{}After\allowbreak{}Reordering\allowbreak{}First\allowbreak{}Potential} para el nodo y para cada hijo, y al final cambia los estados de la variable. Ese método, para nodos de azar o de utilidad, hace \texttt{aux\allowbreak{}Node.\allowbreak{}get\allowbreak{}Potentials(\allowbreak{}).\allowbreak{}get\allowbreak{}First(\allowbreak{})} sin comprobar que haya alguno. \texttt{Node\allowbreak{}State\allowbreak{}Edit.\allowbreak{}do\allowbreak{}Edit} solo traduce \texttt{Not\allowbreak{}Supported\allowbreak{}Operation\allowbreak{}Exception}, así que en la ventana la \texttt{No\allowbreak{}Such\allowbreak{}Element\allowbreak{}Exception} sale como error inesperado.

Medido, con A → B, A con tabla {[}0,3; 0,7{]} y B (de azar) sin potencial: \texttt{Node\allowbreak{}State\allowbreak{}Edit(\allowbreak{}A,\allowbreak{}\ DOWN,\allowbreak{}\ 0,\allowbreak{}\ "")} lanza \texttt{No\allowbreak{}Such\allowbreak{}Element\allowbreak{}Exception} desde la línea 230. La tabla de A ya se ha reordenado e instalado (queda {[}0,7; 0,3{]}), pero los estados de A siguen siendo «a0, a1» porque \texttt{set\allowbreak{}States} no llegó a ejecutarse; la excepción salta antes de \texttt{add\allowbreak{}Edit}.

Segundo efecto del mismo método: \texttt{aux\allowbreak{}Node.\allowbreak{}set\allowbreak{}Potential(\allowbreak{}new\allowbreak{}Potential)} vacía la lista antes de añadir, así que un nodo con dos potenciales se queda solo con el primero reordenado; los demás no se quedan con los números viejos, se eliminan. Medido en \texttt{0\_\allowbreak{}miscellaneous/\allowbreak{}networks/\allowbreak{}bn/\allowbreak{}BN-prostanet.\allowbreak{}pgmx} («Cáncer de próstata» pasa de 2 potenciales a 1) y en \texttt{0\_\allowbreak{}miscellaneous/\allowbreak{}networks/\allowbreak{}mid/\allowbreak{}MID-Cochlear.\allowbreak{}pgmx} («Processor 2 age {[}1{]}»). En esas dos redes el segundo potencial parece un duplicado o una tabla uniforme, así que el cambio numérico es nulo o pequeño, pero el modelo cambia sin aviso.
\paragraph{El código}
core/src/main/java/org/openmarkov/core/model/network/VariableStateOperations.java:200-236

\begin{Shaded}
\begin{Highlighting}[]
\KeywordTok{private} \DataTypeTok{static} \DataTypeTok{void} \FunctionTok{reorderStatesInNodeAndChildren}\OperatorTok{(}\BuiltInTok{Node}\NormalTok{ node}\OperatorTok{,}\NormalTok{ Variable variable}\OperatorTok{,} \BuiltInTok{State}\OperatorTok{[]}\NormalTok{ newStates}\OperatorTok{)} \OperatorTok{\{}
    \FunctionTok{refuseIfAnyPotentialCannotBeReordered}\OperatorTok{(}\NormalTok{node}\OperatorTok{,}\NormalTok{ variable}\OperatorTok{);}
    \FunctionTok{setPotentialAfterReorderingFirstPotential}\OperatorTok{(}\NormalTok{node}\OperatorTok{,}\NormalTok{ variable}\OperatorTok{,}\NormalTok{ newStates}\OperatorTok{);}
    \ControlFlowTok{for} \OperatorTok{(}\BuiltInTok{Node}\NormalTok{ child }\OperatorTok{:}\NormalTok{ node}\OperatorTok{.}\FunctionTok{getChildren}\OperatorTok{())} \OperatorTok{\{}
        \FunctionTok{setPotentialAfterReorderingFirstPotential}\OperatorTok{(}\NormalTok{child}\OperatorTok{,}\NormalTok{ variable}\OperatorTok{,}\NormalTok{ newStates}\OperatorTok{);}
    \OperatorTok{\}}
\NormalTok{    variable}\OperatorTok{.}\FunctionTok{setStates}\OperatorTok{(}\NormalTok{newStates}\OperatorTok{);}
    \FunctionTok{resetLink}\OperatorTok{(}\NormalTok{node}\OperatorTok{);}
\OperatorTok{\}}
\KeywordTok{...}
\KeywordTok{private} \DataTypeTok{static} \DataTypeTok{void} \FunctionTok{setPotentialAfterReorderingFirstPotential}\OperatorTok{(}\BuiltInTok{Node}\NormalTok{ auxNode}\OperatorTok{,}\NormalTok{ Variable variable}\OperatorTok{,} \BuiltInTok{State}\OperatorTok{[]}\NormalTok{ newStates}\OperatorTok{)} \OperatorTok{\{}
    \ControlFlowTok{if} \OperatorTok{(}\NormalTok{auxNode}\OperatorTok{.}\FunctionTok{getNodeType}\OperatorTok{()} \OperatorTok{==}\NormalTok{ NodeType}\OperatorTok{.}\FunctionTok{CHANCE} \OperatorTok{||}\NormalTok{ auxNode}\OperatorTok{.}\FunctionTok{getNodeType}\OperatorTok{()} \OperatorTok{==}\NormalTok{ NodeType}\OperatorTok{.}\FunctionTok{UTILITY}\OperatorTok{)} \OperatorTok{\{}
\NormalTok{        Potential oldPotential }\OperatorTok{=}\NormalTok{ auxNode}\OperatorTok{.}\FunctionTok{getPotentials}\OperatorTok{().}\FunctionTok{getFirst}\OperatorTok{();}
\NormalTok{        Potential newPotential }\OperatorTok{=}\NormalTok{ oldPotential}\OperatorTok{.}\FunctionTok{reorder}\OperatorTok{(}\NormalTok{variable}\OperatorTok{,}\NormalTok{ newStates}\OperatorTok{);}
        \ControlFlowTok{if} \OperatorTok{(}\NormalTok{newPotential }\OperatorTok{!=} \KeywordTok{null}\OperatorTok{)} \OperatorTok{\{}
\NormalTok{            auxNode}\OperatorTok{.}\FunctionTok{setPotential}\OperatorTok{(}\NormalTok{newPotential}\OperatorTok{);}
        \OperatorTok{\}}
    \OperatorTok{\}}
\OperatorTok{\}}
\end{Highlighting}
\end{Shaded}
\paragraph{Cómo arreglarlo}
Reordenar todos los potenciales del nodo (y de cada hijo) y no tocar un nodo sin potenciales; instalar los resultados con \texttt{set\allowbreak{}Potentials(\allowbreak{}lista)} en vez de \texttt{set\allowbreak{}Potential(\allowbreak{}uno)}. Para no dejar la red a medias, calcular primero todos los potenciales reordenados y solo después instalarlos y cambiar los estados: la misma idea que ya sigue \texttt{refuse\allowbreak{}If\allowbreak{}Any\allowbreak{}Potential\allowbreak{}Cannot\allowbreak{}Be\allowbreak{}Reordered}, que pregunta a todos antes de tocar nada.

Conviene revisar a la vez si el lector debe aceptar nodos de azar sin potencial (hoy los acepta en silencio).

Prueba: A → B con B sin potencial, mover un estado de A: no debe lanzar, y los estados y la tabla de A deben quedar permutados juntos; y un nodo con dos potenciales debe conservar los dos, ambos reordenados.

\setcounter{subsection}{29}
\subsection[Deshacer un cambio de estados duplica tablas de los hijos y cambia resultados]{Deshacer o cancelar un cambio de estados copia las tablas de los hijos en el nodo y en sus padres, y la red da otras probabilidades y se guarda así}\label{err:30}
\begin{ficha}
\fichakey{Estado}Presente
\fichakey{Módulo}core
\fichakey{Paquete}org.openmarkov.core.action.core
\fichakey{Ficheros}\path{core/src/main/java/org/openmarkov/core/action/core/NodeStateEdit.java:137-145, 203-207} \newline \path{core/src/main/java/org/openmarkov/core/model/network/ProbNet.java:576-600} \newline \path{core/src/main/java/org/openmarkov/core/model/network/Node.java:241-248} \newline \path{gui/src/main/java/org/openmarkov/gui/dialog/node/NodePropertiesDialog.java:399-402} \newline \path{gui/src/main/java/org/openmarkov/gui/dialog/common/TablePotentialPanel.java:170-171} \newline \path{io/src/main/java/org/openmarkov/io/probmodel/writer/PGMXWriter_1_0.java:140-153}
\fichakey{Comprobado}Leyendo el código y ejecutándolo
\fichakey{Esfuerzo}Pequeño (horas)
\fichakey{Decisión de equipo}No
\fichakey{Relacionados}\hyperref[err:24]{error~24}, \hyperref[err:29]{error~29}, \hyperref[nota:5]{nota~5}, \hyperref[nota:150]{nota~150}
\end{ficha}
\paragraph{Qué pasa}
En la pestaña de dominio del diálogo de propiedades de un nodo, añadir, quitar, subir o bajar un estado y después pulsar «Cancel» (o aceptar y deshacer con el menú de edición) deja la red con tablas repetidas. El nodo recupera sus estados, pero su lista de potenciales incluye ahora también las tablas de sus hijos, y cada padre recibe a su vez las tablas de los hijos de ese padre. Nada lo avisa.

Consecuencias medidas en \texttt{0\_\allowbreak{}miscellaneous/\allowbreak{}networks/\allowbreak{}bn/\allowbreak{}BN-prostanet.\allowbreak{}pgmx} con el nodo «Cáncer de próstata» (2 padres y 10 hijos), añadiendo un estado y cancelando (la misma secuencia que el diálogo: sub-historia abierta, \texttt{Node\allowbreak{}State\allowbreak{}Edit} y \texttt{cancel\allowbreak{}Last\allowbreak{}Sub\allowbreak{}Edit\allowbreak{}History}):

\begin{itemize}
\tightlist
\item
  la red pasa de 50 potenciales a 82: «Cáncer de próstata» tiene 12 (tenía 2), «Metástasis» 13, «Cistitis» 6;
\item
  la propagación sin hallazgos da otros números: la probabilidad del primer estado de «PSA aux» pasa de 0,674 a 0,958 (la mayor diferencia en toda la red es 0,28);
\item
  guardar escribe los 82 potenciales; al reabrir, cada tabla repetida se cuelga del nodo de su primera variable («Cáncer de próstata» con 6, «TAC» con 3\ldots) y las posteriores siguen mal (diferencia máxima 0,30), así que el daño queda en el fichero;
\item
  el primer potencial de «Cáncer de próstata» pasa a ser la tabla de «Biopsia». La pestaña de la relación enseña el primer potencial del nodo (\texttt{Table\allowbreak{}Potential\allowbreak{}Panel}, líneas 170-171), así que, por lectura, enseñaría la tabla de «Biopsia» como si fuera la del nodo; no se ha construido el panel.
\end{itemize}

Lo mismo al quitar un estado, al bajar uno y deshacer, e incluso con una edición que no cambia nada (bajar el primer estado, medido por programa). Con \texttt{0\_\allowbreak{}miscellaneous/\allowbreak{}networks/\allowbreak{}bn/\allowbreak{}BN-asia.\allowbreak{}pgmx} y el nodo «TuberculosisOrCancer» la red pasa de 8 potenciales a 12 y la probabilidad del estado «yes» de «X-ray» pasa de 0,110 a 0,071. Renombrar un estado no lo provoca, porque su deshacer no toca los potenciales.

Por qué pasa: el constructor de \texttt{Node\allowbreak{}State\allowbreak{}Edit} (módulo \texttt{core}) fotografía, para poder deshacer, \texttt{prob\allowbreak{}Net.\allowbreak{}get\allowbreak{}Potentials(\allowbreak{}variable)} del nodo y de cada vecino. Ese método no devuelve los potenciales guardados en el nodo, sino todos los potenciales de la red que contienen la variable, buscados en el nodo, en sus vecinos y en los otros padres de sus hijos; entre ellos está la tabla de cada hijo. \texttt{undo(\allowbreak{})} instala esas listas con \texttt{set\allowbreak{}Potentials} en el nodo y en cada vecino, de modo que las tablas de los hijos quedan en dos o tres nodos a la vez. Además, \texttt{get\allowbreak{}Potentials(\allowbreak{}variable)} recorre los vecinos antes que el propio nodo, y por eso la tabla de un hijo queda la primera de la lista.

Llegan a esta edición todos los botones de estados de la pestaña de dominio (\texttt{Discretize\allowbreak{}Table\allowbreak{}Panel}, líneas 746, 784, 830 y 848, y el arrastre de filas de la línea 183; \texttt{Prefixed\allowbreak{}Key\allowbreak{}Table\allowbreak{}Panel}, líneas 154, 174, 207 y 231), también cuando el cambio se propaga a los nodos gemelos de otros ciclos. «Cancel» llega a \texttt{undo(\allowbreak{})} porque el diálogo del nodo cancela su historial provisional (\texttt{Node\allowbreak{}Properties\allowbreak{}Dialog.\allowbreak{}do\allowbreak{}Cancel\allowbreak{}Click\allowbreak{}Before\allowbreak{}Hide}).
\paragraph{El código}
core/src/main/java/org/openmarkov/core/action/core/NodeStateEdit.java:136-145

\begin{Shaded}
\begin{Highlighting}[]
\CommentTok{// Save node potentials and neighbours}
\KeywordTok{this}\OperatorTok{.}\FunctionTok{oldPotentials} \OperatorTok{=}\NormalTok{ probNet}\OperatorTok{.}\FunctionTok{getPotentials}\OperatorTok{(}\NormalTok{variable}\OperatorTok{);}
\KeywordTok{this}\OperatorTok{.}\FunctionTok{listOldPotentials} \OperatorTok{=} \KeywordTok{new} \BuiltInTok{HashMap}\OperatorTok{\textless{}\textgreater{}();}
\ControlFlowTok{for} \OperatorTok{(}\BuiltInTok{Node}\NormalTok{ nodeNeighbour }\OperatorTok{:}\NormalTok{ probNet}\OperatorTok{.}\FunctionTok{getNeighbors}\OperatorTok{(}\NormalTok{node}\OperatorTok{))} \OperatorTok{\{}
    \KeywordTok{this}\OperatorTok{.}\FunctionTok{listOldPotentials}\OperatorTok{.}\FunctionTok{put}\OperatorTok{(}
\NormalTok{            nodeNeighbour}\OperatorTok{.}\FunctionTok{getVariable}\OperatorTok{(),}
\NormalTok{            probNet}\OperatorTok{.}\FunctionTok{getPotentials}\OperatorTok{(}\NormalTok{nodeNeighbour}\OperatorTok{.}\FunctionTok{getVariable}\OperatorTok{())}
    \OperatorTok{);}
\OperatorTok{\}}
\end{Highlighting}
\end{Shaded}

core/src/main/java/org/openmarkov/core/action/core/NodeStateEdit.java:203-207

\begin{Shaded}
\begin{Highlighting}[]
\CommentTok{// We restore the potentials}
\NormalTok{node}\OperatorTok{.}\FunctionTok{setPotentials}\OperatorTok{(}\NormalTok{oldPotentials}\OperatorTok{);}
\ControlFlowTok{for} \OperatorTok{(}\NormalTok{Variable variable }\OperatorTok{:}\NormalTok{ listOldPotentials}\OperatorTok{.}\FunctionTok{keySet}\OperatorTok{())} \OperatorTok{\{}
\NormalTok{    probNet}\OperatorTok{.}\FunctionTok{getNode}\OperatorTok{(}\NormalTok{variable}\OperatorTok{).}\FunctionTok{setPotentials}\OperatorTok{(}\NormalTok{listOldPotentials}\OperatorTok{.}\FunctionTok{get}\OperatorTok{(}\NormalTok{variable}\OperatorTok{));}
\OperatorTok{\}}
\end{Highlighting}
\end{Shaded}

core/src/main/java/org/openmarkov/core/model/network/ProbNet.java:586-599

\begin{Shaded}
\begin{Highlighting}[]
\BuiltInTok{Set}\OperatorTok{\textless{}}\BuiltInTok{Node}\OperatorTok{\textgreater{}}\NormalTok{ semiNeighbors }\OperatorTok{=} \KeywordTok{new} \BuiltInTok{LinkedHashSet}\OperatorTok{\textless{}\textgreater{}(}\FunctionTok{getNeighbors}\OperatorTok{(}\NormalTok{node}\OperatorTok{));}
\NormalTok{semiNeighbors}\OperatorTok{.}\FunctionTok{add}\OperatorTok{(}\NormalTok{node}\OperatorTok{);}
\KeywordTok{...}
\ControlFlowTok{for} \OperatorTok{(}\BuiltInTok{Node}\NormalTok{ neighbor }\OperatorTok{:}\NormalTok{ semiNeighbors}\OperatorTok{)} \OperatorTok{\{}
    \BuiltInTok{List}\OperatorTok{\textless{}}\NormalTok{Potential}\OperatorTok{\textgreater{}}\NormalTok{ nodePotentials }\OperatorTok{=}\NormalTok{ neighbor}\OperatorTok{.}\FunctionTok{getPotentials}\OperatorTok{();}
    \ControlFlowTok{for} \OperatorTok{(}\NormalTok{Potential potential }\OperatorTok{:}\NormalTok{ nodePotentials}\OperatorTok{)} \OperatorTok{\{}
        \ControlFlowTok{if} \OperatorTok{(}\NormalTok{potential}\OperatorTok{.}\FunctionTok{contains}\OperatorTok{(}\NormalTok{variable}\OperatorTok{))} \OperatorTok{\{}
\NormalTok{            potentials}\OperatorTok{.}\FunctionTok{add}\OperatorTok{(}\NormalTok{potential}\OperatorTok{);}
        \OperatorTok{\}}
    \OperatorTok{\}}
\OperatorTok{\}}
\end{Highlighting}
\end{Shaded}
\paragraph{Cómo arreglarlo}
Fotografiar la lista propia de cada nodo, copiada: \texttt{new\ Array\allowbreak{}List\textless{}\textgreater{}(\allowbreak{}node.\allowbreak{}get\allowbreak{}Potentials(\allowbreak{}))} para el nodo y lo mismo para cada nodo cuya lista cambie la edición (la copia hace falta porque \texttt{set\allowbreak{}Potentials} vacía la lista del nodo antes de rellenarla). \texttt{Variable\allowbreak{}State\allowbreak{}Operations} solo sustituye los potenciales del nodo y de sus hijos (\texttt{set\allowbreak{}Uniform\allowbreak{}Potential} al añadir o quitar, \texttt{reorder\allowbreak{}States\allowbreak{}In\allowbreak{}Node\allowbreak{}And\allowbreak{}Children} al mover), así que bastan el nodo y sus hijos; guardar también los padres con su propia lista no hace daño.

\texttt{Prob\allowbreak{}Net.\allowbreak{}get\allowbreak{}Potentials(\allowbreak{}Variable)} hace lo que dice su documentación y tiene otros llamadores que solo leen (\texttt{Prob\allowbreak{}Net\allowbreak{}Operations}, \texttt{DANOperations}, \texttt{CENetwork\allowbreak{}Builder}); no hay que cambiarlo, sino dejar de usarlo para fotografiar.

Si la fotografía se comparte con \texttt{Node\allowbreak{}Replace\allowbreak{}States\allowbreak{}Edit}, como propone \hyperref[err:24]{error~24}, hay que compartir la corregida.

Prueba de regresión: sobre \texttt{0\_\allowbreak{}miscellaneous/\allowbreak{}networks/\allowbreak{}bn/\allowbreak{}BN-prostanet.\allowbreak{}pgmx}, dentro de una sub-historia, ejecutar \texttt{Node\allowbreak{}State\allowbreak{}Edit} con \texttt{ADD}, \texttt{REMOVE}, \texttt{UP} y \texttt{DOWN} sobre «Cáncer de próstata» y cancelar: el número de potenciales de cada nodo, el primero de cada lista y las posteriores sin hallazgos deben ser los de antes.

\setcounter{subsection}{30}
\subsection[Pegar en una red que ya incumple una restricción se rechaza siempre]{Pegar en una red que ya incumple una restricción se rechaza siempre, con un mensaje sobre la violación que ya estaba}\label{err:31}
\begin{ficha}
\fichakey{Estado}Presente
\fichakey{Módulo}core
\fichakey{Paquete}org.openmarkov.core.action.base
\fichakey{Ficheros}\path{core/src/main/java/org/openmarkov/core/action/base/MultiStepEdit.java:36-37} \newline \path{gui/src/main/java/org/openmarkov/gui/action/PasteEdit.java:151-161} \newline \path{core/src/main/java/org/openmarkov/core/action/core/VariableTypeEdit.java:77} \newline \path{gui/src/main/java/org/openmarkov/gui/action/ChangeNodeTypeEdit.java:45, 77-86} \newline \path{core/src/main/java/org/openmarkov/core/action/base/linkEdits/InvertLinkEdit.java:131-134}
\fichakey{Comprobado}Leyendo el código y ejecutándolo
\fichakey{Esfuerzo}Medio (días)
\fichakey{Decisión de equipo}\textcolor{decisionrojo}{\textbf{Sí.}} Qué estrategia rige para las ediciones de varios pasos («comprobar antes» o «probar y revertir») y si una violación previa debe bloquear ediciones que no la empeoran
\fichakey{Corregir antes}\hyperref[err:37]{error~37}: si se decide que una red que incumple sus restricciones no pueda entrar al abrir un fichero, este error deja de alcanzarse.
\fichakey{Relacionados}\hyperref[err:28]{error~28}, \hyperref[nota:184]{nota~184}, \hyperref[err:37]{error~37}, \hyperref[nota:142]{nota~142}, \hyperref[nota:158]{nota~158}
\end{ficha}
\paragraph{Qué pasa}
En una red que ya incumple una restricción (por ejemplo, una red bayesiana con un ciclo leída de un fichero), todo pegado se rechaza y se revierte, y el mensaje habla de esa violación que ya estaba, no de lo pegado. Lo mismo ocurrirá al cambiar el tipo de cualquier nodo o al invertir cualquier enlace en esa red, porque las tres operaciones comprueban la red entera.

Hoy solo se llega con una red inválida, y la única puerta hacia una red inválida es abrir un fichero, que no comprueba ninguna restricción: eso es \hyperref[err:37]{error~37}, y los dos defectos van juntos.

Por qué pasa: \texttt{PNEdit.\allowbreak{}check\allowbreak{}Constraints\allowbreak{}Will\allowbreak{}Be\allowbreak{}Met} se documenta como la promesa de que ninguna restricción quedará violada tras la edición, y \texttt{Multi\allowbreak{}Edit} la implementa recorriendo sus sub-ediciones. \texttt{Multi\allowbreak{}Step\allowbreak{}Edit} la anula con un cuerpo vacío, así que sus tres descendientes (\texttt{Paste\allowbreak{}Edit}, \texttt{Variable\allowbreak{}Type\allowbreak{}Edit}, \texttt{Change\allowbreak{}Node\allowbreak{}Type\allowbreak{}Edit}) no comprueban nada por adelantado: ejecutan los pasos y revierten si uno falla. \texttt{Variable\allowbreak{}Type\allowbreak{}Edit} y \texttt{Change\allowbreak{}Node\allowbreak{}Type\allowbreak{}Edit} terminan su comprobación con \texttt{super.\allowbreak{}check\allowbreak{}Constraints\allowbreak{}Will\allowbreak{}Be\allowbreak{}Met(\allowbreak{}.\allowbreak{}.\allowbreak{}.\allowbreak{})}, que parece delegar en el recorrido de sub-ediciones y resuelve al cuerpo vacío: esas líneas no hacen nada.

Para comprobar el resultado, \texttt{Paste\allowbreak{}Edit} añade al final una edición anónima cuyo \texttt{do\allowbreak{}Edit(\allowbreak{})} no hace nada y cuya comprobación previa es \texttt{prob\allowbreak{}Net.\allowbreak{}check\allowbreak{}Constraints\allowbreak{}In(\allowbreak{}.\allowbreak{}.\allowbreak{}.\allowbreak{})}: la red entera, después de pegar. \texttt{Change\allowbreak{}Node\allowbreak{}Type\allowbreak{}Edit} hace lo mismo dentro de \texttt{Set\allowbreak{}Node\allowbreak{}Type\allowbreak{}Edit} (\texttt{prob\allowbreak{}Net.\allowbreak{}check\allowbreak{}Constraints(\allowbreak{})} tras cambiar el tipo) e \texttt{Invert\allowbreak{}Link\allowbreak{}Edit} tras invertir. Esa comprobación posterior no distingue lo que la edición ha roto de lo que ya estaba roto.

Medido: se abrió con el lector de \texttt{.\allowbreak{}pgmx} una red bayesiana con el ciclo A→B→A (se carga sin aviso) y se pegó en ella un nodo suelto Z copiado de otra red. Un \texttt{Add\allowbreak{}Node\allowbreak{}Edit} solo para Z pasa su comprobación, pero el pegado se rechaza con «Could not met the constraint: No cycles paths are allowed. - There is a cycle traversing between \ldots{} A \ldots{} and \ldots{} B \ldots» y la red queda sin Z.
\paragraph{El código}
gui/src/main/java/org/openmarkov/gui/action/PasteEdit.java:151-161

\begin{Shaded}
\begin{Highlighting}[]
\NormalTok{PNEdit verifier }\OperatorTok{=} \KeywordTok{new} \FunctionTok{PNEdit}\OperatorTok{(}\NormalTok{probNet}\OperatorTok{)} \OperatorTok{\{}
    
    \AttributeTok{@Override} \KeywordTok{public} \DataTypeTok{void} \FunctionTok{checkConstraintsWillBeMet}\OperatorTok{(}\NormalTok{ConstraintChecker constraintChecker}\OperatorTok{)} \OperatorTok{\{}
\NormalTok{        probNet}\OperatorTok{.}\FunctionTok{checkConstraintsIn}\OperatorTok{(}\NormalTok{constraintChecker}\OperatorTok{);}
    \OperatorTok{\}}
    
    \AttributeTok{@Override} \KeywordTok{protected} \DataTypeTok{void} \FunctionTok{doEdit}\OperatorTok{()} \OperatorTok{\{}
    
    \OperatorTok{\}}
\OperatorTok{\};}
\NormalTok{stepExecuter}\OperatorTok{.}\FunctionTok{execute}\OperatorTok{(}\NormalTok{verifier}\OperatorTok{);}
\end{Highlighting}
\end{Shaded}

core/src/main/java/org/openmarkov/core/action/base/MultiStepEdit.java:33-34

\begin{Shaded}
\begin{Highlighting}[]
\AttributeTok{@Override} \KeywordTok{public} \DataTypeTok{void} \FunctionTok{checkConstraintsWillBeMet}\OperatorTok{(}\NormalTok{ConstraintChecker constraintChecker}\OperatorTok{)} \OperatorTok{\{}
\OperatorTok{\}}
\end{Highlighting}
\end{Shaded}
\paragraph{Cómo arreglarlo}
Hay que decidir qué hacer con una red que ya llega inválida (ver \texttt{decision-de-equipo}). Una opción es comparar las violaciones antes y después y rechazar solo las nuevas; otra, impedir que entre inválida al abrir el fichero, que es el arreglo de \hyperref[err:37]{error~37}.

A más largo plazo, la anulación a vacío y la edición verificadora sobran cuando cada restricción sepa juzgar una edición; esa propuesta de diseño está en \hyperref[nota:153]{nota~153} y \hyperref[nota:158]{nota~158}, donde no hay nada que corregir.

Sitios con la misma comprobación posterior: \texttt{Paste\allowbreak{}Edit} (verificador), \texttt{Change\allowbreak{}Node\allowbreak{}Type\allowbreak{}Edit.\allowbreak{}Set\allowbreak{}Node\allowbreak{}Type\allowbreak{}Edit.\allowbreak{}do\allowbreak{}Edit} e \texttt{Invert\allowbreak{}Link\allowbreak{}Edit.\allowbreak{}do\allowbreak{}Edit} (líneas 131-134); y las dos llamadas \texttt{super.\allowbreak{}check\allowbreak{}Constraints\allowbreak{}Will\allowbreak{}Be\allowbreak{}Met} que no hacen nada en \texttt{Variable\allowbreak{}Type\allowbreak{}Edit} y \texttt{Change\allowbreak{}Node\allowbreak{}Type\allowbreak{}Edit}.

Prueba de regresión: abrir una red con un ciclo, pegar un nodo aislado y exigir que se pegue (o, si se decide lo contrario, que el mensaje diga que la red ya era inválida).

\setcounter{subsection}{32}
\subsection[Deshacer o cancelar «mostrar el comentario al abrir» no la devuelve]{Deshacer o cancelar la casilla «mostrar el comentario al abrir» no la devuelve a su valor}\label{err:33}
\begin{ficha}
\fichakey{Estado}Presente
\fichakey{Módulo}core
\fichakey{Paquete}org.openmarkov.core.action.core
\fichakey{Ficheros}\path{core/src/main/java/org/openmarkov/core/action/core/NetworkCommentEdit.java:34-50} \newline \path{gui/src/main/java/org/openmarkov/gui/dialog/network/NetworkDefinitionPanel.java:301-316} \newline \path{gui/src/main/java/org/openmarkov/gui/dialog/network/NetworkDefinitionPanel.java:351-358} \newline \path{gui/src/main/java/org/openmarkov/gui/dialog/network/NetworkPropertiesDialog.java:280-283} \newline \path{io/src/main/java/org/openmarkov/io/probmodel/writer/PGMXWriter_0_2.java:356}
\fichakey{Comprobado}Leyendo el código y ejecutándolo
\fichakey{Esfuerzo}Pequeño (horas)
\fichakey{Decisión de equipo}No
\end{ficha}
\paragraph{Qué pasa}
El usuario marca la casilla «mostrar el comentario al abrir» en las propiedades de la red, pulsa «Cancelar», y la red queda con la opción cambiada. Al guardar se escribe en el fichero (\texttt{PGMXWriter\_\allowbreak{}0\_\allowbreak{}2}, atributo del comentario) y al volver a abrir la red se enseña el comentario. Deshacer tampoco la devuelve.

Camino de llegada: Red → Propiedades → pestaña de definición. La casilla ejecuta, cada vez que se marca o desmarca, un \texttt{Network\allowbreak{}Comment\allowbreak{}Edit} con el comentario actual y el valor nuevo de la casilla (\texttt{Network\allowbreak{}Definition\allowbreak{}Panel.\allowbreak{}get\allowbreak{}Check\allowbreak{}Box\allowbreak{}Show\allowbreak{}Comment\allowbreak{}On\allowbreak{}Opening}); editar el comentario también lo construye, pasando el estado de la casilla. El diálogo abre un historial parcial y «Cancelar» deshace lo hecho en él.

Por qué pasa: \texttt{Network\allowbreak{}Comment\allowbreak{}Edit} (módulo \texttt{core}) cambia dos cosas en \texttt{do\allowbreak{}Edit(\allowbreak{})}: el comentario de la red y la opción \texttt{show\allowbreak{}Comment\allowbreak{}When\allowbreak{}Opening} (si el comentario se enseña al abrir el fichero). El constructor solo guarda el comentario anterior (\texttt{current\allowbreak{}Comment}), y \texttt{undo(\allowbreak{})} solo repone el comentario.

Medido: con la opción a falso, ejecutar la edición con la casilla marcada y deshacer deja la opción a \texttt{true}; dentro de un historial parcial, cancelar también la deja a \texttt{true}.
\paragraph{El código}
core/src/main/java/org/openmarkov/core/action/core/NetworkCommentEdit.java:34-50

\begin{Shaded}
\begin{Highlighting}[]
\KeywordTok{public} \FunctionTok{NetworkCommentEdit}\OperatorTok{(}\NormalTok{ProbNet probNet}\OperatorTok{,} \BuiltInTok{String}\NormalTok{ newComment}\OperatorTok{,} \DataTypeTok{boolean}\NormalTok{ showCommentWhenOpening}\OperatorTok{)} \OperatorTok{\{}
    \KeywordTok{super}\OperatorTok{(}\NormalTok{probNet}\OperatorTok{);}
    \KeywordTok{this}\OperatorTok{.}\FunctionTok{newComment} \OperatorTok{=}\NormalTok{ newComment}\OperatorTok{;}
    \KeywordTok{this}\OperatorTok{.}\FunctionTok{currentComment} \OperatorTok{=}\NormalTok{ probNet}\OperatorTok{.}\FunctionTok{getComment}\OperatorTok{();}
    \KeywordTok{this}\OperatorTok{.}\FunctionTok{showCommentWhenOpening} \OperatorTok{=}\NormalTok{ showCommentWhenOpening}\OperatorTok{;}
\OperatorTok{\}}

\AttributeTok{@Override} \KeywordTok{protected} \DataTypeTok{void} \FunctionTok{doEdit}\OperatorTok{()} \OperatorTok{\{}
\NormalTok{    probNet}\OperatorTok{.}\FunctionTok{setComment}\OperatorTok{(}\NormalTok{newComment}\OperatorTok{);}
\NormalTok{    probNet}\OperatorTok{.}\FunctionTok{setShowCommentWhenOpening}\OperatorTok{(}\NormalTok{showCommentWhenOpening}\OperatorTok{);}
\OperatorTok{\}}

\AttributeTok{@Override} \KeywordTok{public} \DataTypeTok{void} \FunctionTok{undo}\OperatorTok{()} \OperatorTok{\{}
    \KeywordTok{super}\OperatorTok{.}\FunctionTok{undo}\OperatorTok{();}
\NormalTok{    probNet}\OperatorTok{.}\FunctionTok{setComment}\OperatorTok{(}\NormalTok{currentComment}\OperatorTok{);}
\OperatorTok{\}}
\end{Highlighting}
\end{Shaded}
\paragraph{Cómo arreglarlo}
Guardar en el constructor el valor anterior de \texttt{prob\allowbreak{}Net.\allowbreak{}get\allowbreak{}Show\allowbreak{}Comment\allowbreak{}When\allowbreak{}Opening(\allowbreak{})} y reponerlo en \texttt{undo(\allowbreak{})} junto al comentario.

Prueba: con la opción a falso, ejecutar \texttt{Network\allowbreak{}Comment\allowbreak{}Edit(\allowbreak{}red,\allowbreak{}\ comentario,\allowbreak{}\ true)} y deshacer debe dejarla a falso; y lo mismo con \texttt{open\allowbreak{}New\allowbreak{}Sub\allowbreak{}Edit\allowbreak{}History} + \texttt{cancel\allowbreak{}Last\allowbreak{}Sub\allowbreak{}Edit\allowbreak{}History}.

\setcounter{subsection}{33}
\subsection[Deshacer o cancelar «Always append» de un suceso no devuelve su valor]{Deshacer o cancelar el cambio de «Always append» de un suceso no devuelve el valor que tenía}\label{err:34}
\begin{ficha}
\fichakey{Estado}Presente en parte. Copiar una red ya conserva «Always append»: \texttt{Node.\allowbreak{}copy\allowbreak{}Properties\allowbreak{}To}, por la que pasan las dos formas de copiar una red, copia el campo desde el commit \texttt{db16fd4}.
\fichakey{Módulo}core
\fichakey{Paquete}org.openmarkov.core.action.core
\fichakey{Ficheros}\path{core/src/main/java/org/openmarkov/core/action/core/EventNodeAlwaysAppendEdit.java:35-56} \newline \path{gui/src/main/java/org/openmarkov/gui/dialog/node/NodeDefinitionPanel.java:674-690, 735-744, 1001-1002} \newline \path{gui/src/main/java/org/openmarkov/gui/dialog/node/NodePropertiesDialog.java:399-402} \newline \path{core/src/main/java/org/openmarkov/core/model/network/Node.java:141-149, 733-743}
\fichakey{Comprobado}Leyendo el código y ejecutándolo
\fichakey{Esfuerzo}Pequeño (horas)
\fichakey{Decisión de equipo}No
\fichakey{Relacionados}\hyperref[nota:78]{nota~78}
\end{ficha}
\paragraph{Qué pasa}
Un nodo de suceso (los nodos de las redes de simulación de eventos discretos) tiene dos campos booleanos distintos en \texttt{Node}: \texttt{always\allowbreak{}Observed} («su valor se conoce siempre») y \texttt{always\allowbreak{}Append} («cuando otro suceso lo dispara, se encola siempre una instancia nueva» frente a «solo hay una en la cola y se le cambia la hora»). El segundo decide cómo se simula el suceso.

Lo que ve el usuario: un suceso que tenía «Always append» marcado, se desmarca y se cancela, \textbf{se queda desmarcado}, y la simulación lo trata desde entonces como un suceso que no se encola. No hay aviso.

Camino del usuario: en una red de simulación de eventos discretos, abrir las propiedades de un nodo de suceso. El diálogo enseña la casilla «Always append» en la fila de agentes/criterio/observado (\texttt{Node\allowbreak{}Definition\allowbreak{}Panel.\allowbreak{}get\allowbreak{}Agents\allowbreak{}Or\allowbreak{}Decision\allowbreak{}Criteria\allowbreak{}Or\allowbreak{}Observed}, líneas 674-689). Marcarla o desmarcarla ejecuta la edición al momento (líneas 996-998). Si el usuario pulsa \textbf{Cancelar}, el diálogo deshace todo lo hecho (\texttt{Node\allowbreak{}Properties\allowbreak{}Dialog.\allowbreak{}do\allowbreak{}Cancel\allowbreak{}Click\allowbreak{}Before\allowbreak{}Hide} → \texttt{cancel\allowbreak{}Last\allowbreak{}Sub\allowbreak{}Edit\allowbreak{}History}, líneas 389-392); si pulsa Aceptar y luego Ctrl+Z, igual. En los dos casos se ejecuta el \texttt{undo(\allowbreak{})} de la edición.

Por qué pasa: la edición que cambia \texttt{always\allowbreak{}Append}, \texttt{Event\allowbreak{}Node\allowbreak{}Always\allowbreak{}Append\allowbreak{}Edit} (módulo \texttt{core}), guarda en su constructor como valor anterior el \textbf{otro} campo, \texttt{node.\allowbreak{}is\allowbreak{}Always\allowbreak{}Observed(\allowbreak{})}, y su \texttt{undo(\allowbreak{})} escribe ese valor en \texttt{always\allowbreak{}Append}. Deshacer, por tanto, no devuelve lo que había: deja en «encolar siempre» lo que valiera «siempre observado».

Como la casilla «siempre observado» no se ofrece para nodos de suceso (\texttt{Always\allowbreak{}Observed\allowbreak{}Property\allowbreak{}Validator} solo la admite en nodos de azar), en la práctica \texttt{always\allowbreak{}Observed} vale \texttt{false} en un suceso, y de ahí sale el resultado descrito arriba.

Medido (red \texttt{DESNetwork\allowbreak{}Type}, historial de deshacer activado):

\begin{Verbatim}[breaklines,breakanywhere,fontsize=\small]
antes alwaysAppend=true  alwaysObserved=false -> tras la edición=false -> tras deshacer=false   (debía ser true)
antes alwaysAppend=false alwaysObserved=false -> tras la edición=true  -> tras deshacer=false   (correcto por casualidad)
antes alwaysAppend=true  alwaysObserved=true  -> tras la edición=false -> tras deshacer=true    (correcto por casualidad)
antes alwaysAppend=false alwaysObserved=true  -> tras la edición=true  -> tras deshacer=true    (debía ser false)
\end{Verbatim}

La prueba de ediciones del proyecto (\texttt{integration\allowbreak{}Tests/\allowbreak{}src/\allowbreak{}test/\allowbreak{}java/\allowbreak{}org/\allowbreak{}openmarkov/\allowbreak{}integration\allowbreak{}Tests/\allowbreak{}action/\allowbreak{}Edit\allowbreak{}Undo\allowbreak{}Cases.\allowbreak{}java}, línea 98) deja esta edición sin caso («needs an event node of a simulation network; pending»), por eso no se detecta.
\paragraph{El código}
core/src/main/java/org/openmarkov/core/action/core/EventNodeAlwaysAppendEdit.java:35-56

\begin{Shaded}
\begin{Highlighting}[]
\KeywordTok{public} \FunctionTok{EventNodeAlwaysAppendEdit}\OperatorTok{(}\BuiltInTok{Node}\NormalTok{ node}\OperatorTok{,} \DataTypeTok{boolean}\NormalTok{ newAlwaysAppend}\OperatorTok{)} \OperatorTok{\{}
    \KeywordTok{super}\OperatorTok{(}\NormalTok{node}\OperatorTok{.}\FunctionTok{getProbNet}\OperatorTok{());}
    \KeywordTok{this}\OperatorTok{.}\FunctionTok{lastAlwaysAppend} \OperatorTok{=}\NormalTok{ node}\OperatorTok{.}\FunctionTok{isAlwaysObserved}\OperatorTok{();}
    \KeywordTok{this}\OperatorTok{.}\FunctionTok{newAlwaysAppend} \OperatorTok{=}\NormalTok{ newAlwaysAppend}\OperatorTok{;}
    \KeywordTok{this}\OperatorTok{.}\FunctionTok{node} \OperatorTok{=}\NormalTok{ node}\OperatorTok{;}
\OperatorTok{\}}
\KeywordTok{...}
\AttributeTok{@Override} \KeywordTok{public} \DataTypeTok{void} \FunctionTok{undo}\OperatorTok{()} \OperatorTok{\{}
    \KeywordTok{super}\OperatorTok{.}\FunctionTok{undo}\OperatorTok{();}
\NormalTok{    node}\OperatorTok{.}\FunctionTok{setAlwaysAppend}\OperatorTok{(}\NormalTok{lastAlwaysAppend}\OperatorTok{);}
\OperatorTok{\}}
\end{Highlighting}
\end{Shaded}
\paragraph{Cómo arreglarlo}
Guardar \texttt{node.\allowbreak{}is\allowbreak{}Always\allowbreak{}Append(\allowbreak{})} en el constructor (línea 37). Una edición genérica de propiedades que lea y escriba el mismo campo evitaría este tipo de error de raíz, pero es un problema de diseño de las ediciones que no hace falta resolver para esto.

Por coherencia, el constructor de copia \texttt{Node(\allowbreak{}Node)} (líneas 141-149) puede copiar también \texttt{always\allowbreak{}Append} (y lo que copia \texttt{copy\allowbreak{}Properties\allowbreak{}To}). Hoy no lo copia, pero su único llamador de producción es \texttt{ICIOption\allowbreak{}Listener\allowbreak{}Assistant} (línea 110), que trabaja con nodos de azar de modelos canónicos, nunca con sucesos, así que nadie lo alcanza con un suceso.

Prueba de regresión: rellenar el caso pendiente de \texttt{Edit\allowbreak{}Undo\allowbreak{}Cases} con un nodo de suceso de una red \texttt{DESNetwork\allowbreak{}Type} que empiece con \texttt{always\allowbreak{}Append=true} y \texttt{always\allowbreak{}Observed=false}; ejecutar la edición con \texttt{false}, deshacer y comprobar que vuelve a \texttt{true} (y el caso simétrico con \texttt{always\allowbreak{}Observed=true}).

\setcounter{subsection}{34}
\subsection[Una red expandida o aprendida no se marca como modificada ni pide guardar]{La ventana de una red expandida o aprendida a partir de otra no se marca nunca como modificada, y cerrarla no pide guardar}\label{err:35}
\begin{ficha}
\fichakey{Estado}Presente
\fichakey{Módulo}core
\fichakey{Paquete}org.openmarkov.core.model.network
\fichakey{Ficheros}\path{core/src/main/java/org/openmarkov/core/model/network/ProbNetCopier.java:39-67, 119} \newline \path{core/src/main/java/org/openmarkov/core/action/base/PNESupport.java:55-62, 92-95, 286-295} \newline \path{core/src/main/java/org/openmarkov/core/model/network/TemporalNetOperations.java:101-108, 451-469} \newline \path{core/src/main/java/org/openmarkov/core/model/network/ProbNetOperations.java:431-432} \newline \path{learning.core/src/main/java/org/openmarkov/learning/core/LearningManager.java:185-186} \newline \path{gui/src/main/java/org/openmarkov/gui/window/InferenceHandler.java:161-186} \newline \path{gui/src/main/java/org/openmarkov/gui/window/edition/networkEditorPanel/NetworkEditorPanel.java:1152-1155} \newline \path{gui/src/main/java/org/openmarkov/gui/window/NetworkFileHandler.java:368-388} \newline \path{learning.gui/src/main/java/org/openmarkov/learning/gui/LearningDialog.java:899-904, 1019-1023}
\fichakey{Comprobado}Leyendo el código y ejecutándolo
\fichakey{Esfuerzo}Pequeño (horas)
\fichakey{Decisión de equipo}No
\fichakey{Relacionados}\hyperref[err:36]{error~36}, \hyperref[err:49]{error~49}, \hyperref[err:75]{error~75}, \hyperref[err:174]{error~174}, \hyperref[nota:38]{nota~38}, \hyperref[nota:177]{nota~177}
\end{ficha}
\paragraph{Qué pasa}
Cuando una red expandida o aprendida a partir de otra se abre en su propia pestaña, lo que el usuario edite después en esa pestaña no la pone nunca como modificada: el título no cambia de color y, al cerrarla (o al cerrar la aplicación), no se pregunta si se quiere guardar, así que los cambios se pierden sin aviso.

Hay dos caminos en los que una copia acaba abierta en su propia pestaña del editor:

\begin{enumerate}
\def\labelenumi{\arabic{enumi}.}
\tightlist
\item
  \textbf{Expandir una red temporal.} Clic derecho sobre el lienzo vacío de una red con nodos temporales (un MID, diagrama de influencia de Markov) → «Expand network» (\texttt{Editor\allowbreak{}Input\allowbreak{}Handler.\allowbreak{}java:667-670}, \texttt{Inference\allowbreak{}Handler.\allowbreak{}expand\allowbreak{}Network}, líneas 164-185). La red expandida sale de \texttt{Temporal\allowbreak{}Net\allowbreak{}Operations.\allowbreak{}expand\allowbreak{}Network}, que hace \texttt{deep\allowbreak{}Copy} y luego \texttt{copy(\allowbreak{})} (líneas 457 y 102), y \texttt{Prob\allowbreak{}Net\allowbreak{}Operations.\allowbreak{}convert\allowbreak{}Numerical\allowbreak{}Variables\allowbreak{}To\allowbreak{}FS} hace otra \texttt{copy(\allowbreak{})} (línea 432). Se abre con \texttt{create\allowbreak{}New\allowbreak{}Frame}.
\item
  \textbf{Aprender partiendo de una red modelo.} Herramientas de aprendizaje con la opción de usar como red modelo la red abierta (\texttt{Learning\allowbreak{}Dialog.\allowbreak{}java:1019-1023}) y «empezar desde la red modelo»: \texttt{Learning\allowbreak{}Manager.\allowbreak{}apply\allowbreak{}Model\allowbreak{}Net} hace \texttt{model\allowbreak{}Net.\allowbreak{}copy(\allowbreak{})} (línea 186) y esa red aprendida se abre en una pestaña nueva (\texttt{Learning\allowbreak{}Dialog.\allowbreak{}java:899-904}). Además, todas las ediciones que el algoritmo de aprendizaje ejecuta sobre la red aprendida avisan a la ventana de la red modelo.
\end{enumerate}

No se ha arrancado la aplicación: la cadena desde la edición hasta el indicador está leída en el código, y el comportamiento del indicador está medido.

\textbf{Qué es un oyente aquí.} Cada red (\texttt{Prob\allowbreak{}Net}) tiene un \texttt{PNESupport}, que guarda la historia de deshacer/rehacer y una lista de oyentes (\texttt{PNEdit\allowbreak{}Listener}). Cada vez que se ejecuta una edición sobre la red (\texttt{PNEdit.\allowbreak{}execute\allowbreak{}Edit}), se avisa a todos los oyentes de esa lista. Uno de esos oyentes es el propio \texttt{PNESupport} de la red: es él quien, al recibir el aviso, pone a verdadero su indicador \texttt{network\allowbreak{}Is\allowbreak{}Modified} (líneas 286-295). La ventana del editor (\texttt{Network\allowbreak{}Editor\allowbreak{}Panel.\allowbreak{}after\allowbreak{}Edit\allowbreak{}Executes}, líneas 1158-1161) lee ese indicador después de cada edición y lo copia a su \texttt{modified}; de \texttt{modified} dependen el color del título de la pestaña (\texttt{Main\allowbreak{}Panel.\allowbreak{}java:519-537}) y que al cerrar la pestaña se pregunte si se quiere guardar (\texttt{Network\allowbreak{}File\allowbreak{}Handler.\allowbreak{}network\allowbreak{}Can\allowbreak{}Be\allowbreak{}Closed}, líneas 372-392).

\textbf{Por qué pasa.} Las dos copias de red de \texttt{Prob\allowbreak{}Net\allowbreak{}Copier} (la superficial, \texttt{shallow\allowbreak{}Copy}, que es \texttt{Prob\allowbreak{}Net.\allowbreak{}copy(\allowbreak{})}, y la profunda, \texttt{deep\allowbreak{}Copy}) terminan con \texttt{dest.\allowbreak{}get\allowbreak{}PNESupport(\allowbreak{}).\allowbreak{}set\allowbreak{}Listeners(\allowbreak{}source.\allowbreak{}get\allowbreak{}PNESupport(\allowbreak{}).\allowbreak{}get\allowbreak{}Listeners(\allowbreak{}))}. \texttt{set\allowbreak{}Listeners} vacía la lista de la copia y mete en ella los oyentes del original. Resultado:

\begin{itemize}
\tightlist
\item
  la copia deja de oírse a sí misma: su \texttt{PNESupport} ya no está en su propia lista, así que su indicador de modificación no cambia nunca;
\item
  la copia avisa, en cada edición, a los oyentes del original: al \texttt{PNESupport} del original y a la ventana del original (su \texttt{Network\allowbreak{}Editor\allowbreak{}Panel}, su \texttt{Visual\allowbreak{}Network}, su \texttt{PNEdit\allowbreak{}Event\allowbreak{}Handler}), que se redibujan y reconstruyen su información visual por una edición que no ocurrió en su red.
\end{itemize}

Medido sobre \texttt{0\_\allowbreak{}miscellaneous/\allowbreak{}networks/\allowbreak{}bn/\allowbreak{}BN-asia.\allowbreak{}pgmx}: la copia no contiene su propio \texttt{PNESupport} entre sus oyentes y sí el del original y el oyente que se añadió al original; tras añadir un enlace en la copia (con deshacer activado) se avisa una vez al oyente del original, la copia sigue con \texttt{network\allowbreak{}Is\allowbreak{}Modified(\allowbreak{})\ =\ false} y el original no tiene el enlace. La misma edición sobre una red leída de fichero da \texttt{network\allowbreak{}Is\allowbreak{}Modified(\allowbreak{})\ =\ true}. Sobre \texttt{0\_\allowbreak{}miscellaneous/\allowbreak{}networks/\allowbreak{}mid/\allowbreak{}MID-Chancellor-corrected.\allowbreak{}pgmx}, la red que devuelve \texttt{Temporal\allowbreak{}Net\allowbreak{}Operations.\allowbreak{}expand\allowbreak{}Network(\allowbreak{}red,\allowbreak{}\ evidencia,\allowbreak{}\ nombre)} tampoco se oye a sí misma y oye al \texttt{PNESupport} y a la ventana del MID original.

En esas dos pestañas, además, el deshacer está apagado, porque \texttt{create\allowbreak{}New\allowbreak{}Frame} no lo enciende (sí lo hacen \texttt{open\allowbreak{}Network} y \texttt{create\allowbreak{}New\allowbreak{}Network}); eso está descrito en \hyperref[err:36]{error~36}. Aun encendiéndolo, el indicador seguiría sin cambiar, como muestra la medida. Otro defecto de las mismas copias, que comparten con el original las opciones de inferencia, está en \hyperref[err:49]{error~49}.
\paragraph{El código}
core/src/main/java/org/openmarkov/core/model/network/ProbNetCopier.java:39-66

\begin{Shaded}
\begin{Highlighting}[]
\DataTypeTok{static}\NormalTok{ ProbNet }\FunctionTok{shallowCopy}\OperatorTok{(}\NormalTok{ProbNet source}\OperatorTok{)} \OperatorTok{\{}
\NormalTok{    ProbNet dest }\OperatorTok{=} \KeywordTok{new} \FunctionTok{ProbNet}\OperatorTok{(}\NormalTok{source}\OperatorTok{.}\FunctionTok{getNetworkType}\OperatorTok{());}
    \KeywordTok{...}
    \FunctionTok{copyLinks}\OperatorTok{(}\NormalTok{source}\OperatorTok{,}\NormalTok{ dest}\OperatorTok{,} \KeywordTok{false}\OperatorTok{);}
    \FunctionTok{copyMetadata}\OperatorTok{(}\NormalTok{source}\OperatorTok{,}\NormalTok{ dest}\OperatorTok{,} \KeywordTok{false}\OperatorTok{);}
\NormalTok{    dest}\OperatorTok{.}\FunctionTok{getPNESupport}\OperatorTok{().}\FunctionTok{setListeners}\OperatorTok{(}\NormalTok{source}\OperatorTok{.}\FunctionTok{getPNESupport}\OperatorTok{().}\FunctionTok{getListeners}\OperatorTok{());}
    \KeywordTok{...}
    \CommentTok{// Shared like the potentials: a shallow copy carries the very options of the original.}
\NormalTok{    dest}\OperatorTok{.}\FunctionTok{setInferenceOptions}\OperatorTok{(}\NormalTok{source}\OperatorTok{.}\FunctionTok{getInferenceOptions}\OperatorTok{());}
    \ControlFlowTok{return}\NormalTok{ dest}\OperatorTok{;}
\OperatorTok{\}}
\end{Highlighting}
\end{Shaded}

core/src/main/java/org/openmarkov/core/action/base/PNESupport.java:92-95 y 286-295

\begin{Shaded}
\begin{Highlighting}[]
\KeywordTok{public} \DataTypeTok{void} \FunctionTok{setListeners}\OperatorTok{(}\BuiltInTok{Collection}\OperatorTok{\textless{}?} \KeywordTok{extends}\NormalTok{ PNEditListener}\OperatorTok{\textgreater{}}\NormalTok{ listeners}\OperatorTok{)} \OperatorTok{\{}
    \KeywordTok{this}\OperatorTok{.}\FunctionTok{listeners}\OperatorTok{.}\FunctionTok{clear}\OperatorTok{();}
    \KeywordTok{this}\OperatorTok{.}\FunctionTok{listeners}\OperatorTok{.}\FunctionTok{addAll}\OperatorTok{(}\NormalTok{listeners}\OperatorTok{);}
\OperatorTok{\}}
\KeywordTok{...}
\KeywordTok{private} \DataTypeTok{void} \FunctionTok{onEditChanges}\OperatorTok{(}\NormalTok{PNEdit edit}\OperatorTok{)} \OperatorTok{\{}
    \KeywordTok{...}
    \KeywordTok{this}\OperatorTok{.}\FunctionTok{networkIsModified} \OperatorTok{=} \OperatorTok{(!}\KeywordTok{this}\OperatorTok{.}\FunctionTok{withUndo} \OperatorTok{||} \KeywordTok{this}\OperatorTok{.}\FunctionTok{editsHistoryStacker}\OperatorTok{.}\FunctionTok{getCurrentUndoManager}\OperatorTok{()}
                                                  \OperatorTok{.}\FunctionTok{nextEditToUndo}\OperatorTok{()} \OperatorTok{!=}\NormalTok{ lastDoneEditWhenSave}\OperatorTok{);}
\OperatorTok{\}}
\end{Highlighting}
\end{Shaded}
\paragraph{Cómo arreglarlo}
Que ninguna de las dos copias transfiera oyentes por omisión: la copia se queda con la lista que le da su propio constructor (\texttt{PNESupport(\allowbreak{}Prob\allowbreak{}Net)} se añade a sí mismo, línea 61). Si algún llamador necesita de verdad que la copia avise a otros, que lo pida con una llamada explícita.

Hay que quitarlo en los dos sitios: \texttt{Prob\allowbreak{}Net\allowbreak{}Copier.\allowbreak{}shallow\allowbreak{}Copy} (línea 56) y \texttt{Prob\allowbreak{}Net\allowbreak{}Copier.\allowbreak{}deep\allowbreak{}Copy} (línea 119). Conviene revisar los llamadores de \texttt{copy(\allowbreak{})}/\texttt{deep\allowbreak{}Copy(\allowbreak{})} que registran oyentes después (por ejemplo \texttt{Hugin\allowbreak{}Forest}, líneas 108 y 114, que añade y quita la heurística sobre su red de Markov) para confirmar que ninguno cuenta con heredar los del original.

Pruebas de regresión (en \texttt{core}, sin ventana):

\begin{itemize}
\tightlist
\item
  Leer o construir una red, añadirle un oyente contador, copiarla con \texttt{copy(\allowbreak{})} y con \texttt{deep\allowbreak{}Copy(\allowbreak{})}; en cada copia comprobar que su propio \texttt{PNESupport} está entre sus oyentes, que el contador del original no está, y que tras ejecutar una edición sobre la copia \texttt{network\allowbreak{}Is\allowbreak{}Modified(\allowbreak{})} pasa a verdadero y el contador del original sigue a cero.
\item
  Otra sobre \texttt{Temporal\allowbreak{}Net\allowbreak{}Operations.\allowbreak{}expand\allowbreak{}Network(\allowbreak{}red,\allowbreak{}\ evidencia,\allowbreak{}\ nombre)} con un MID del repositorio, con las mismas comprobaciones.
\end{itemize}

\setcounter{subsection}{35}
\subsection[En una red expandida o aprendida, deshacer no recuerda nada al principio]{En la pestaña de una red expandida o aprendida, deshacer no recuerda nada hasta que se crea un nodo o se abre ciertos diálogos}\label{err:36}
\begin{ficha}
\fichakey{Estado}Presente en parte. En las redes que el usuario crea o abre (fichero, dirección de internet, sesión anterior, línea de órdenes), \texttt{Network\allowbreak{}File\allowbreak{}Handler} ya enciende el deshacer desde el principio.
\fichakey{Módulo}core, gui, learning.gui
\fichakey{Paquete}org.openmarkov.core.action.base; org.openmarkov.gui.window
\fichakey{Ficheros}\path{core/src/main/java/org/openmarkov/core/action/base/PNESupport.java:55-58, 196-212, 223-232} \newline \path{core/src/main/java/org/openmarkov/core/action/base/PNEdit.java:94-96} \newline \path{core/src/main/java/org/openmarkov/core/model/network/ProbNet.java:116} \newline \path{gui/src/main/java/org/openmarkov/gui/window/NetworkFileHandler.java:97, 103-118, 126-129, 178} \newline \path{gui/src/main/java/org/openmarkov/gui/window/InferenceHandler.java:161-186} \newline \path{learning.gui/src/main/java/org/openmarkov/learning/gui/LearningDialog.java:903-905} \newline \path{gui/src/main/java/org/openmarkov/gui/window/decisiontree/DecisionTreeEditor.java:270-281} \newline \path{gui/src/main/java/org/openmarkov/gui/window/edition/networkEditorPanel/EditorInputHandler.java:714}
\fichakey{Comprobado}Leyendo el código y ejecutándolo
\fichakey{Esfuerzo}Pequeño (horas)
\fichakey{Decisión de equipo}No
\fichakey{Relacionados}\hyperref[err:35]{error~35}, \hyperref[nota:177]{nota~177}
\end{ficha}
\paragraph{Qué pasa}
Hay pestañas de red que se abren con \texttt{Network\allowbreak{}File\allowbreak{}Handler.\allowbreak{}create\allowbreak{}New\allowbreak{}Frame} o \texttt{open\allowbreak{}Network(\allowbreak{}Prob\allowbreak{}Net)} sin encender el deshacer. En ellas, mover o borrar nodos, añadir o quitar enlaces, etc. no se apunta y la opción Deshacer queda deshabilitada, hasta que el usuario hace alguna de las acciones que lo encienden (por ejemplo, crear un nodo o abrir las propiedades de la red); a partir de ahí funciona el resto de la sesión. La consecuencia en la ventana (Deshacer deshabilitado) se ha deducido de la lectura y de la medida sobre \texttt{PNESupport}, no se ha visto con la aplicación.

Esas pestañas son:

\begin{itemize}
\tightlist
\item
  «Expand network» (menú contextual del lienzo en modo inferencia sobre una red temporal): \texttt{Inference\allowbreak{}Handler.\allowbreak{}expand\allowbreak{}Network} abre en una pestaña nueva el resultado de \texttt{Temporal\allowbreak{}Net\allowbreak{}Operations.\allowbreak{}expand\allowbreak{}Network}, que parte de \texttt{prob\allowbreak{}Net.\allowbreak{}copy(\allowbreak{})}.
\item
  El aprendizaje automático: \texttt{Learning\allowbreak{}Dialog} (línea 904) abre la red que devuelve \texttt{Learning\allowbreak{}Manager.\allowbreak{}get\allowbreak{}Learned\allowbreak{}Net(\allowbreak{})}, creada con \texttt{new\ Prob\allowbreak{}Net(\allowbreak{})}; solo el aprendizaje interactivo enciende el deshacer.
\item
  La opción de abrir la red asociada a un nodo del árbol de decisión (\texttt{Decision\allowbreak{}Tree\allowbreak{}Editor.\allowbreak{}open\allowbreak{}Associated\allowbreak{}Network}); no se comprobó de dónde sale esa red.
\end{itemize}

Por qué pasa: cada red (\texttt{Prob\allowbreak{}Net}) tiene su gestor de ediciones, \texttt{PNESupport} (módulo \texttt{core}), que nace con \texttt{with\allowbreak{}Undo\ =\ false}; \texttt{PNEdit.\allowbreak{}execute\allowbreak{}Edit} solo apunta una edición en el historial de deshacer si ese indicador está encendido (líneas 94-96). Como \texttt{Prob\allowbreak{}Net} crea siempre un \texttt{PNESupport} nuevo (línea 116), también una copia de una red nace con el deshacer apagado.

Hoy encienden el deshacer: \texttt{Network\allowbreak{}File\allowbreak{}Handler} al crear una red nueva (línea 105) y al abrir un fichero o una dirección de internet (URL, línea 182, que también usa la reapertura de la sesión anterior y la apertura desde la línea de órdenes); y además \texttt{Editor\allowbreak{}Input\allowbreak{}Handler.\allowbreak{}create\allowbreak{}Node} (línea 722, al crear un nodo con doble clic o desde el menú del lienzo), \texttt{Potential\allowbreak{}Edit\allowbreak{}Panel} (línea 157, al editar un potencial), \texttt{Network\allowbreak{}Properties\allowbreak{}Dialog} (línea 82), \texttt{Reorder\allowbreak{}Variables\allowbreak{}Dialog} (línea 31), \texttt{Add\allowbreak{}Finding\allowbreak{}Dialog} (línea 90) e \texttt{Interactive\allowbreak{}Learning\allowbreak{}Dialog} (línea 86).

Medido: tras encender el deshacer en una red y copiarla con \texttt{copy(\allowbreak{})}, la copia tiene \texttt{with\allowbreak{}Undo=false}; añadir un nodo a la copia deja \texttt{get\allowbreak{}Can\allowbreak{}Undo(\allowbreak{})=false}, y en la original \texttt{true}.

En las pestañas de red expandida o aprendida hay además otro defecto, el indicador de modificación que no cambia nunca, descrito en \hyperref[err:35]{error~35}.
\paragraph{El código}
gui/src/main/java/org/openmarkov/gui/window/NetworkFileHandler.java:134-137

\begin{Shaded}
\begin{Highlighting}[]
    \DataTypeTok{void} \FunctionTok{openNetwork}\OperatorTok{(}\NormalTok{ProbNet probNet}\OperatorTok{)} \OperatorTok{\{}
\NormalTok{        NetworkEditorPanel newNetworkEditorPanel }\OperatorTok{=} \KeywordTok{this}\OperatorTok{.}\FunctionTok{createNewFrame}\OperatorTok{(}\NormalTok{probNet}\OperatorTok{,} \KeywordTok{null}\OperatorTok{);}
\NormalTok{        networkPanels}\OperatorTok{.}\FunctionTok{add}\OperatorTok{(}\NormalTok{newNetworkEditorPanel}\OperatorTok{);}
    \OperatorTok{\}}
\end{Highlighting}
\end{Shaded}
\paragraph{Cómo arreglarlo}
Encender el deshacer en un único sitio por el que pasen todas las pestañas de red: \texttt{Network\allowbreak{}File\allowbreak{}Handler.\allowbreak{}create\allowbreak{}New\allowbreak{}Frame} (líneas 111-126), que usan la creación, la apertura, la expansión, el aprendizaje y el árbol de decisión. Así sobran las llamadas de las líneas 105 y 182. Si el aprendizaje automático no debe dejar deshacer sus pasos, encenderlo después de \texttt{learning\allowbreak{}Manager.\allowbreak{}learn(\allowbreak{})}.

Las llamadas repartidas por diálogos (\texttt{Potential\allowbreak{}Edit\allowbreak{}Panel}, \texttt{Network\allowbreak{}Properties\allowbreak{}Dialog}, \texttt{Reorder\allowbreak{}Variables\allowbreak{}Dialog}, \texttt{Add\allowbreak{}Finding\allowbreak{}Dialog}, \texttt{Editor\allowbreak{}Input\allowbreak{}Handler}) quedarían redundantes; quitarlas es una limpieza aparte. Si se hace, la propuesta de diseño de \hyperref[nota:177]{nota~177} (extraer de la ventana un documento de red que lleve, entre otras cosas, el estado de modificado) es el lugar natural; allí no hay nada que corregir.

Prueba: abrir en una pestaña una red copiada con \texttt{copy(\allowbreak{})}, ejecutar una edición y comprobar que \texttt{get\allowbreak{}Can\allowbreak{}Undo(\allowbreak{})} es verdadero.

% ══════════════════════════════════════════════════════════════════════
\section{Restricciones de red}\label{sec:restricciones}

Una restricción es una regla que dice qué permite cada tipo de red (por ejemplo, que no haya ciclos). Se imponen dos veces: al comprobar una red entera (PNConstraint.checkProbNet) y antes de cada edición (PNEdit.checkConstraintsWillBeMet).

\setcounter{subsection}{36}
\subsection[Abrir un fichero no comprueba las restricciones de la red]{Abrir un fichero no comprueba ninguna restricción: una red bayesiana con un ciclo se abre sin aviso}\label{err:37}
\begin{ficha}
\fichakey{Estado}Presente
\fichakey{Módulo}io
\fichakey{Paquete}org.openmarkov.io.probmodel.reader
\fichakey{Ficheros}\path{io/src/main/java/org/openmarkov/io/probmodel/reader/PGMXReader_0_2.java:134-150} \newline \path{io/src/main/java/org/openmarkov/io/probmodel/reader/PGMXReader_0_2.java:272-286} \newline \path{gui/src/main/java/org/openmarkov/gui/window/NetworkFileHandler.java:154-181} \newline \path{core/src/main/java/org/openmarkov/core/model/network/ProbNet.java:103-124} \newline \path{core/src/main/java/org/openmarkov/core/model/network/ProbNet.java:190-216} \newline \path{core/src/main/java/org/openmarkov/core/model/network/ProbNet.java:339-356}
\fichakey{Comprobado}Leyendo el código y ejecutándolo
\fichakey{Esfuerzo}Pequeño (horas)
\fichakey{Decisión de equipo}\textcolor{decisionrojo}{\textbf{Sí.}} Qué hacer al abrir una red que no cumple sus restricciones: rechazar el fichero o abrirlo avisando (abrir avisando no deja fuera ficheros antiguos)
\fichakey{Relacionados}\hyperref[err:5]{error~5}, \hyperref[err:31]{error~31}, \hyperref[err:38]{error~38}, \hyperref[err:39]{error~39}, \hyperref[err:111]{error~111}, \hyperref[nota:153]{nota~153}
\end{ficha}
\paragraph{Qué pasa}
Un fichero \texttt{.\allowbreak{}pgmx} con una red que incumple sus restricciones se abre sin ningún aviso. Medido: se escribió con \texttt{PGMXWriter\_\allowbreak{}1\_\allowbreak{}0} una red bayesiana A→B, se añadió a mano en el XML el enlace B→A, y \texttt{PGMXReader.\allowbreak{}read} la devolvió sin error, con los dos enlaces y con \texttt{get\allowbreak{}Unsatisfied\allowbreak{}Constraints(\allowbreak{})} = \texttt{{[}No\allowbreak{}Cycle{]}}.

A partir de ahí la red inválida se edita y se guarda como si nada. Además, pegar, cambiar el tipo de un nodo o invertir un enlace se rechazan con un mensaje sobre el ciclo que ya estaba (medido para pegar: es \hyperref[err:31]{error~31}). No se ha medido qué hace cada algoritmo de inferencia con una red así.

Las restricciones de una red se comprueban por dos vías: la de cada edición antes de ejecutarse y la de red completa (\texttt{PNConstraint.\allowbreak{}check\allowbreak{}Prob\allowbreak{}Net}, que se llama desde \texttt{Prob\allowbreak{}Net.\allowbreak{}check\allowbreak{}Constraints}, \texttt{check\allowbreak{}Constraints\allowbreak{}In} y \texttt{get\allowbreak{}Unsatisfied\allowbreak{}Constraints}). En producción, la de red completa solo la llaman \texttt{Invert\allowbreak{}Link\allowbreak{}Edit}, \texttt{Change\allowbreak{}Node\allowbreak{}Type\allowbreak{}Edit} y el verificador de \texttt{Paste\allowbreak{}Edit}, las tres después de modificar la red; \texttt{Inference\allowbreak{}Algorithm.\allowbreak{}check\allowbreak{}Constraints}, con las restricciones del algoritmo y no las de la red; y \texttt{Prob\allowbreak{}Net.\allowbreak{}set\allowbreak{}Network\allowbreak{}Type}, para las restricciones del tipo nuevo.

En el módulo \texttt{io} no hay ninguna llamada. \texttt{PGMXReader\_\allowbreak{}0\_\allowbreak{}2.\allowbreak{}read} construye la red con \texttt{new\ Prob\allowbreak{}Net(\allowbreak{}tipo)} (cuyo javadoc dice expresamente que instala las restricciones sin comprobarlas), añade nodos, enlaces y potenciales directamente, y la devuelve. \texttt{Network\allowbreak{}File\allowbreak{}Handler.\allowbreak{}open\allowbreak{}Network\allowbreak{}Source} la pone en una pestaña sin preguntar nada. Tampoco comprueba nada \texttt{resttemplate} (hoy es solo un esqueleto) ni el aprendizaje, que construye la red con ediciones.
\paragraph{El código}
io/src/main/java/org/openmarkov/io/probmodel/reader/PGMXReader\_0\_2.java:134-150

\begin{Shaded}
\begin{Highlighting}[]
\KeywordTok{public}\NormalTok{ PGMXReader}\OperatorTok{.}\FunctionTok{NetworkAndEvidence} \FunctionTok{read}\OperatorTok{(}\BuiltInTok{URL}\NormalTok{ networkSource}\OperatorTok{)} \KeywordTok{throws}\NormalTok{ ProbNetParserException }\OperatorTok{\{}
\NormalTok{    FormatManager formatManager }\OperatorTok{=}\NormalTok{ FormatManager}\OperatorTok{.}\FunctionTok{getInstance}\OperatorTok{();}
    \ControlFlowTok{try} \OperatorTok{\{}
\NormalTok{        formatManager}\OperatorTok{.}\FunctionTok{checkVersion}\OperatorTok{(}\NormalTok{networkSource}\OperatorTok{);}
\NormalTok{        formatManager}\OperatorTok{.}\FunctionTok{checkStructure}\OperatorTok{(}\NormalTok{networkSource}\OperatorTok{);}
    \OperatorTok{\}} \ControlFlowTok{catch} \OperatorTok{(}\BuiltInTok{SAXException} \OperatorTok{|} \BuiltInTok{IOException}\NormalTok{ e}\OperatorTok{)} \OperatorTok{\{}
        \ControlFlowTok{throw} \KeywordTok{new}\NormalTok{ ProbNetParserException}\OperatorTok{.}\FunctionTok{PGMXInvalid}\OperatorTok{(}\NormalTok{e}\OperatorTok{.}\FunctionTok{getMessage}\OperatorTok{());}
    \OperatorTok{\}}
    \BuiltInTok{Element}\NormalTok{ root }\OperatorTok{=} \FunctionTok{getRootElement}\OperatorTok{(}\NormalTok{networkSource}\OperatorTok{);}
\NormalTok{    ProbNet probNet }\OperatorTok{=} \KeywordTok{this}\OperatorTok{.}\FunctionTok{getProbNet}\OperatorTok{(}\NormalTok{root}\OperatorTok{,}\NormalTok{ networkSource}\OperatorTok{.}\FunctionTok{getFile}\OperatorTok{());}
    \KeywordTok{this}\OperatorTok{.}\FunctionTok{getInferenceOptions}\OperatorTok{(}\NormalTok{root}\OperatorTok{,}\NormalTok{ probNet}\OperatorTok{);}
    \BuiltInTok{List}\OperatorTok{\textless{}}\NormalTok{EvidenceCase}\OperatorTok{\textgreater{}}\NormalTok{ evidence }\OperatorTok{=} \KeywordTok{this}\OperatorTok{.}\FunctionTok{getEvidence}\OperatorTok{(}\NormalTok{root}\OperatorTok{,}\NormalTok{ probNet}\OperatorTok{);}
    \KeywordTok{this}\OperatorTok{.}\FunctionTok{getPolicies}\OperatorTok{(}\NormalTok{root}\OperatorTok{,}\NormalTok{ probNet}\OperatorTok{);}
    \KeywordTok{this}\OperatorTok{.}\FunctionTok{setVariableType}\OperatorTok{(}\NormalTok{root}\OperatorTok{,}\NormalTok{ probNet}\OperatorTok{);}
    \KeywordTok{this}\OperatorTok{.}\FunctionTok{setDefaultStates}\OperatorTok{(}\NormalTok{root}\OperatorTok{,}\NormalTok{ probNet}\OperatorTok{);}
    \ControlFlowTok{return} \KeywordTok{new}\NormalTok{ PGMXReader}\OperatorTok{.}\FunctionTok{NetworkAndEvidence}\OperatorTok{(}\NormalTok{probNet}\OperatorTok{,}\NormalTok{ evidence}\OperatorTok{);}
\OperatorTok{\}}
\end{Highlighting}
\end{Shaded}
\paragraph{Cómo arreglarlo}
Llamar a \texttt{prob\allowbreak{}Net.\allowbreak{}get\allowbreak{}Unsatisfied\allowbreak{}Constraints(\allowbreak{})} (o \texttt{check\allowbreak{}Constraints(\allowbreak{})}) al terminar de leer.

Hay que elegir dónde. El sitio común a todos los formatos es \texttt{Nets\allowbreak{}IO.\allowbreak{}open\allowbreak{}Network\allowbreak{}URL} o \texttt{Network\allowbreak{}File\allowbreak{}Handler.\allowbreak{}open\allowbreak{}Network\allowbreak{}Source}, pero entonces los que abren redes sin ventana quedan fuera. Si se pone en cada lector, hay que hacerlo en PGMX (las versiones 0.2 y 1.0 comparten \texttt{read}), en XMLBIF y en los demás lectores registrados (el formato Elvira no se mantiene).

Qué hacer con el resultado es la decisión de producto: lanzar una excepción de análisis con la lista de violaciones, o devolver la red con avisos que la ventana muestre. Las redes de ejemplo y de pruebas del repositorio podrían no pasar todas; conviene recorrerlas antes de elegir.

Prueba de regresión: el fichero con A→B y B→A debe dar la excepción o el aviso que se decida.

\setcounter{subsection}{37}
\subsection[Un estado puede quedarse con el nombre vacío o repetido]{La comprobación de nombres de estado no se ejecuta nunca: se puede dejar un estado con el nombre vacío, o con uno repetido por programa}\label{err:38}
\begin{ficha}
\fichakey{Estado}Presente
\fichakey{Módulo}core
\fichakey{Paquete}org.openmarkov.core.model.network.constraint
\fichakey{Ficheros}\path{core/src/main/java/org/openmarkov/core/model/network/constraint/ValidState.java:23-83} \newline \path{core/src/main/java/org/openmarkov/core/action/core/NodeStateEdit.java:123-124} \newline \path{core/src/main/java/org/openmarkov/core/action/core/NodeStateEdit.java:147-152} \newline \path{core/src/main/java/org/openmarkov/core/model/network/Variable.java:292-307} \newline \path{gui/src/main/java/org/openmarkov/gui/component/DiscretizeTablePanel.java:733-746} \newline \path{gui/src/main/java/org/openmarkov/gui/component/DiscretizeTablePanel.java:907-940}
\fichakey{Comprobado}Leyendo el código y ejecutándolo
\fichakey{Esfuerzo}Pequeño (horas)
\fichakey{Decisión de equipo}\textcolor{decisionrojo}{\textbf{Sí.}} Si «dos estados con el mismo nombre» se decide ignorando mayúsculas o no, y si la regla vive en una restricción activa o directamente en la variable y en las ediciones
\fichakey{Relacionados}\hyperref[err:37]{error~37}, \hyperref[err:39]{error~39}, \hyperref[err:58]{error~58}, \hyperref[err:72]{error~72}, \hyperref[nota:143]{nota~143}, \hyperref[nota:147]{nota~147}
\end{ficha}
\paragraph{Qué pasa}
\texttt{Valid\allowbreak{}State} es la regla «un estado no puede tener el nombre vacío ni repetir el de otro estado de la misma variable». Tiene las dos vías: \texttt{check\allowbreak{}State} para una edición (la usa \texttt{Node\allowbreak{}State\allowbreak{}Edit.\allowbreak{}check\allowbreak{}Constraints\allowbreak{}Will\allowbreak{}Be\allowbreak{}Met}) y \texttt{check\allowbreak{}Prob\allowbreak{}Net} para la red entera. Pero ninguna de las dos se ejecuta.

Desde la ventana, los estados se añaden y renombran en la pestaña de dominio de las propiedades del nodo (\texttt{Discretize\allowbreak{}Table\allowbreak{}Panel}). Añadir inventa un nombre libre («\ldots{} 1», «\ldots{} 2»), así que por ahí no entran repetidos. Renombrar ejecuta \texttt{Node\allowbreak{}State\allowbreak{}Edit(\allowbreak{}RENAME)}, que llama a \texttt{Variable.\allowbreak{}rename\allowbreak{}State}: esta rechaza un nombre exacto que ya tenga otro estado, pero \textbf{no} rechaza el nombre vacío, y compara con igualdad exacta, así que acepta «Yes» junto a «yes».

Medido con una variable de estados «yes» y «no»: renombrar «no» a ``\,'' se acepta; renombrar después «yes» a ``\,'' se rechaza por repetido; la red escrita con \texttt{PGMXWriter\_\allowbreak{}1\_\allowbreak{}0} y leída con \texttt{PGMXReader} vuelve con los estados \texttt{{[}yes{]}} y \texttt{{[}{]}}. Como tampoco se comprueba nada al abrir un fichero (\hyperref[err:37]{error~37}), un fichero con estados vacíos o repetidos también entra sin aviso.

Por programa también se alcanza el duplicado por alta, porque \texttt{Node\allowbreak{}State\allowbreak{}Edit} es pública. Medido directamente sobre \texttt{Node\allowbreak{}State\allowbreak{}Edit}: añadir un estado llamado igual que otro («a0» a una variable con «a0, a1») se acepta; renombrar un estado a la cadena vacía se acepta; la red se escribe en \texttt{.\allowbreak{}pgmx} con \texttt{\textless{}State\ name="a0"/\allowbreak{}\textgreater{}} repetido o \texttt{\textless{}State\ name=""/\allowbreak{}\textgreater{}} y se vuelve a leer igual. Con un nombre repetido, \texttt{Variable.\allowbreak{}get\allowbreak{}State\allowbreak{}Index(\allowbreak{}"a0")} devuelve 0 también para el tercer estado, de modo que toda búsqueda por nombre (lector, escritor, evidencia) encuentra el primero.

Por qué pasa: \texttt{Node\allowbreak{}State\allowbreak{}Edit.\allowbreak{}check\allowbreak{}Constraints\allowbreak{}Will\allowbreak{}Be\allowbreak{}Met} (módulo \texttt{core}) solo hace algo si \texttt{prob\allowbreak{}Net.\allowbreak{}get\allowbreak{}Constraint\allowbreak{}Of\allowbreak{}Class(\allowbreak{}Valid\allowbreak{}State.\allowbreak{}class)} devuelve una restricción. \texttt{Valid\allowbreak{}State} no lleva la anotación \texttt{@Constraint}, así que \texttt{Constraint\allowbreak{}Manager} no la descubre, ningún tipo de red la incluye y nadie hace \texttt{new\ Valid\allowbreak{}State(\allowbreak{})}. Todo el contenido de \texttt{Node\allowbreak{}State\allowbreak{}Edit.\allowbreak{}check\allowbreak{}Constraints\allowbreak{}Will\allowbreak{}Be\allowbreak{}Met} es código muerto. Comprobado ejecutando: en redes bayesianas, diagramas de influencia, MID (diagramas de influencia de Markov) y DAN (redes de análisis de decisiones) recién creadas, \texttt{get\allowbreak{}Constraint\allowbreak{}Of\allowbreak{}Class(\allowbreak{}Valid\allowbreak{}State.\allowbreak{}class)} es \texttt{null}; en cambio \texttt{Valid\allowbreak{}Criterion\allowbreak{}Name}, que sí está anotada, aparece en todas.

Las copias de la regla ya no coinciden entre sí, aunque hoy no se note porque la clase está inactiva: \texttt{Valid\allowbreak{}State.\allowbreak{}exist\allowbreak{}State} compara con \texttt{equals\allowbreak{}Ignore\allowbreak{}Case}, mientras que \texttt{Valid\allowbreak{}State.\allowbreak{}check\allowbreak{}Prob\allowbreak{}Net} y \texttt{Variable.\allowbreak{}rename\allowbreak{}State} comparan exacto. Si se activara la restricción tal como está, añadir «Yes» junto a «yes» se rechazaría por la vía de la edición y pasaría por la de la red.

Hay además una trampa para el arreglo: para \texttt{RENAME}, \texttt{Node\allowbreak{}State\allowbreak{}Edit} guarda en \texttt{new\allowbreak{}State} el estado \textbf{seleccionado} (línea 124), no uno con el nombre nuevo, y la comprobación usa \texttt{get\allowbreak{}New\allowbreak{}State(\allowbreak{}).\allowbreak{}get\allowbreak{}Name(\allowbreak{})}. Si se activara \texttt{Valid\allowbreak{}State} tal cual, cada renombrado preguntaría por el nombre actual del propio estado, que existe siempre, y \textbf{todos} los renombrados se rechazarían.
\paragraph{El código}
core/src/main/java/org/openmarkov/core/action/core/NodeStateEdit.java:123-124, 147-152

\begin{Shaded}
\begin{Highlighting}[]
\BuiltInTok{State}\NormalTok{ selectedState }\OperatorTok{=}\NormalTok{ oldStates}\OperatorTok{[}\NormalTok{selectedStateIndex}\OperatorTok{];}
\KeywordTok{this}\OperatorTok{.}\FunctionTok{newState} \OperatorTok{=} \OperatorTok{(}\NormalTok{stateAction }\OperatorTok{!=}\NormalTok{ StateAction}\OperatorTok{.}\FunctionTok{RENAME}\OperatorTok{)} \OperatorTok{?} \KeywordTok{new} \BuiltInTok{State}\OperatorTok{(}\NormalTok{newName}\OperatorTok{)} \OperatorTok{:}\NormalTok{ selectedState}\OperatorTok{;}
\KeywordTok{...}
\AttributeTok{@Override} \KeywordTok{public} \DataTypeTok{void} \FunctionTok{checkConstraintsWillBeMet}\OperatorTok{(}\NormalTok{ConstraintChecker constraintChecker}\OperatorTok{)} \OperatorTok{\{}
    \ControlFlowTok{if} \OperatorTok{(}\NormalTok{probNet}\OperatorTok{.}\FunctionTok{getConstraintOfClass}\OperatorTok{(}\NormalTok{ValidState}\OperatorTok{.}\FunctionTok{class}\OperatorTok{)} \KeywordTok{instanceof}\NormalTok{ ValidState constraint}\OperatorTok{)} \OperatorTok{\{}
\NormalTok{        constraint}\OperatorTok{.}\FunctionTok{checkState}\OperatorTok{(}\NormalTok{constraintChecker}\OperatorTok{,} \KeywordTok{this}\OperatorTok{.}\FunctionTok{getNewState}\OperatorTok{()}
                                                     \OperatorTok{.}\FunctionTok{getName}\OperatorTok{(),} \KeywordTok{this}\OperatorTok{.}\FunctionTok{getNode}\OperatorTok{(),} \KeywordTok{this}\OperatorTok{.}\FunctionTok{getStateAction}\OperatorTok{());}
    \OperatorTok{\}}
\OperatorTok{\}}
\end{Highlighting}
\end{Shaded}

core/src/main/java/org/openmarkov/core/model/network/constraint/ValidState.java:23, 57-60

\begin{Shaded}
\begin{Highlighting}[]
\KeywordTok{public} \KeywordTok{class}\NormalTok{ ValidState }\KeywordTok{extends}\NormalTok{ PNConstraint }\OperatorTok{\{}      \CommentTok{// sin @Constraint}
\KeywordTok{...}
\KeywordTok{private} \DataTypeTok{static} \DataTypeTok{boolean} \FunctionTok{existState}\OperatorTok{(}\BuiltInTok{String}\NormalTok{ state}\OperatorTok{,} \BuiltInTok{Node}\NormalTok{ node}\OperatorTok{)} \OperatorTok{\{}
    \ControlFlowTok{for} \OperatorTok{(}\BuiltInTok{State}\NormalTok{ states }\OperatorTok{:}\NormalTok{ node}\OperatorTok{.}\FunctionTok{getVariable}\OperatorTok{().}\FunctionTok{getStates}\OperatorTok{())} \OperatorTok{\{}
        \ControlFlowTok{if} \OperatorTok{(}\NormalTok{states}\OperatorTok{.}\FunctionTok{getName}\OperatorTok{().}\FunctionTok{equalsIgnoreCase}\OperatorTok{(}\NormalTok{state}\OperatorTok{))} \OperatorTok{\{}
\end{Highlighting}
\end{Shaded}

core/src/main/java/org/openmarkov/core/model/network/Variable.java:292-299

\begin{Shaded}
\begin{Highlighting}[]
\KeywordTok{public} \DataTypeTok{void} \FunctionTok{renameState}\OperatorTok{(}\BuiltInTok{State}\NormalTok{ state}\OperatorTok{,} \BuiltInTok{String}\NormalTok{ newName}\OperatorTok{)} \OperatorTok{\{}
    \ControlFlowTok{for} \OperatorTok{(}\BuiltInTok{State}\NormalTok{ s }\OperatorTok{:}\NormalTok{ states}\OperatorTok{)} \OperatorTok{\{}
        \CommentTok{// The state being renamed is skipped: renaming it to its own current name is a no{-}op,}
        \CommentTok{// not a collision.}
        \ControlFlowTok{if} \OperatorTok{((}\NormalTok{s }\OperatorTok{!=}\NormalTok{ state}\OperatorTok{)} \OperatorTok{\&\&}\NormalTok{ s}\OperatorTok{.}\FunctionTok{getName}\OperatorTok{().}\FunctionTok{contentEquals}\OperatorTok{(}\NormalTok{newName}\OperatorTok{))} \OperatorTok{\{}
            \ControlFlowTok{throw} \KeywordTok{new} \FunctionTok{InvalidArgumentException}\OperatorTok{(}\NormalTok{newName}\OperatorTok{,} \StringTok{"new name"}\OperatorTok{,} \StringTok{"that name already belongs to another state"}\OperatorTok{);}
        \OperatorTok{\}}
    \OperatorTok{\}}
\end{Highlighting}
\end{Shaded}
\paragraph{Cómo arreglarlo}
Hay dos opciones, y elegir es cosa del equipo:

\begin{itemize}
\tightlist
\item
  \begin{enumerate}
  \def\labelenumi{(\alph{enumi})}
  \tightlist
  \item
    Anotar \texttt{Valid\allowbreak{}State} como obligatoria, dejarle una sola regla y corregir \texttt{Node\allowbreak{}State\allowbreak{}Edit} para que compruebe \texttt{new\allowbreak{}Name} (no el nombre del estado seleccionado) y, al renombrar, excluya al propio estado.
  \end{enumerate}
\item
  \begin{enumerate}
  \def\labelenumi{(\alph{enumi})}
  \setcounter{enumi}{1}
  \tightlist
  \item
    Borrar \texttt{Valid\allowbreak{}State} y su consulta, y hacer que la regla viva en un único sitio del dominio: \texttt{Variable.\allowbreak{}rename\allowbreak{}State}, donde ya está la de repetidos, y el alta de estados de \texttt{Variable\allowbreak{}State\allowbreak{}Operations.\allowbreak{}add\allowbreak{}State\allowbreak{}At}, donde falta; en los dos hay que añadir el rechazo del nombre vacío.
  \end{enumerate}
\end{itemize}

En cualquier caso hay que decidir antes si «Yes» y «yes» chocan.

Sitios donde la regla debe cumplirse con el mismo criterio:

\begin{itemize}
\tightlist
\item
  la comprobación de red completa, \texttt{Valid\allowbreak{}State.\allowbreak{}check\allowbreak{}Prob\allowbreak{}Net}, y la previa de la edición, \texttt{Node\allowbreak{}State\allowbreak{}Edit.\allowbreak{}check\allowbreak{}Constraints\allowbreak{}Will\allowbreak{}Be\allowbreak{}Met} con \texttt{Valid\allowbreak{}State.\allowbreak{}check\allowbreak{}State}, si se conserva la clase;
\item
  \texttt{Variable.\allowbreak{}rename\allowbreak{}State};
\item
  el alta en \texttt{Variable\allowbreak{}State\allowbreak{}Operations.\allowbreak{}add\allowbreak{}State\allowbreak{}At};
\item
  la sustitución de todos los estados, \texttt{Variable\allowbreak{}State\allowbreak{}Operations.\allowbreak{}replace\allowbreak{}States}, que usa \texttt{Node\allowbreak{}Replace\allowbreak{}States\allowbreak{}Edit} (dominios estándar);
\item
  la propagación a los nodos gemelos de otros ciclos, \texttt{propagate\allowbreak{}Node\allowbreak{}State\allowbreak{}Edit\allowbreak{}Related\allowbreak{}Variables} en \texttt{Discretize\allowbreak{}Table\allowbreak{}Panel} y en \texttt{Prefixed\allowbreak{}Key\allowbreak{}Table\allowbreak{}Panel};
\item
  el lector de ficheros, si se decide comprobar las restricciones al abrir (\hyperref[err:37]{error~37}).
\end{itemize}

Pruebas: añadir un estado con nombre repetido y renombrar un estado a ``\,'' deben rechazarse con un mensaje, y el nombre original debe seguir en la tabla; renombrar un estado a su propio nombre y a un nombre libre debe aceptarse.

\setcounter{subsection}{38}
\subsection[Ocho restricciones opcionales no llegan a ninguna red]{Las restricciones opcionales no entran en ninguna red: ocho reglas escritas no defienden a nadie}\label{err:39}
\begin{ficha}
\fichakey{Estado}Presente
\fichakey{Módulo}core
\fichakey{Paquete}org.openmarkov.core.model.network.constraint
\fichakey{Ficheros}\path{core/src/main/java/org/openmarkov/core/model/network/constraint/ConstraintManager.java:57-106} \newline \path{core/src/main/java/org/openmarkov/core/model/network/ProbNet.java:115-124} \newline \path{core/src/main/java/org/openmarkov/core/model/network/ProbNet.java:311-331} \newline \path{core/src/test/java/org/openmarkov/core/model/network/constraint/EveryConstraintReachesANetworkTest.java:34-52} \newline \path{core/src/main/java/org/openmarkov/core/action/base/linkEdits/AddLinkEdit.java:99-130} \newline \path{inference/src/main/java/org/openmarkov/inference/algorithm/variableElimination/tasks/VariableElimination.java:127-132}
\fichakey{Comprobado}Leyendo el código
\fichakey{Esfuerzo}Medio (días)
\fichakey{Decisión de equipo}\textcolor{decisionrojo}{\textbf{Sí.}} Qué significa OPTIONAL: la activa el tipo de red; la activa el usuario y el formato la guarda; o el nivel desaparece
\fichakey{Relacionados}\hyperref[err:37]{error~37}, \hyperref[err:38]{error~38}, \hyperref[nota:89]{nota~89}, \hyperref[nota:146]{nota~146}
\end{ficha}
\paragraph{Qué pasa}
Cada restricción declara en su anotación \texttt{@Constraint} un comportamiento por omisión: YES (obligatoria), NO u OPTIONAL. Ocho restricciones marcadas OPTIONAL no entran en ninguna red: \texttt{No\allowbreak{}Loops}, \texttt{No\allowbreak{}Mixed\allowbreak{}Parents}, \texttt{No\allowbreak{}Super\allowbreak{}Value\allowbreak{}Node}, \texttt{Only\allowbreak{}Finite\allowbreak{}States\allowbreak{}Variables}, \texttt{Only\allowbreak{}Numeric\allowbreak{}Variables}, \texttt{Only\allowbreak{}One\allowbreak{}Utility\allowbreak{}Node}, \texttt{Proper\allowbreak{}Utility\allowbreak{}Potentials} y \texttt{Only\allowbreak{}Unlabeled\allowbreak{}Links}.

Las ediciones que consultan esas restricciones reciben \texttt{null} y no comprueban nada: \texttt{Add\allowbreak{}Link\allowbreak{}Edit} con \texttt{No\allowbreak{}Loops}, \texttt{No\allowbreak{}Mixed\allowbreak{}Parents} y \texttt{No\allowbreak{}Super\allowbreak{}Value\allowbreak{}Node}; \texttt{Add\allowbreak{}Node\allowbreak{}Edit} y \texttt{Change\allowbreak{}Node\allowbreak{}Type\allowbreak{}Edit} con \texttt{Only\allowbreak{}One\allowbreak{}Utility\allowbreak{}Node}; \texttt{Variable\allowbreak{}Type\allowbreak{}Edit} con \texttt{Only\allowbreak{}Finite\allowbreak{}States\allowbreak{}Variables} y \texttt{Only\allowbreak{}Numeric\allowbreak{}Variables}; \texttt{Remove\allowbreak{}Node\allowbreak{}Edit} con \texttt{Proper\allowbreak{}Utility\allowbreak{}Potentials}; \texttt{Invert\allowbreak{}Link\allowbreak{}Edit} con \texttt{No\allowbreak{}Mixed\allowbreak{}Parents}. Sus pruebas pasan porque añaden la restricción a mano.

Un caso visible: la eliminación de variables exige \texttt{No\allowbreak{}Mixed\allowbreak{}Parents} como restricción del algoritmo. El editor deja dar a un nodo de utilidad padres de utilidad y de azar a la vez sin quejarse, y el error solo aparece al evaluar (red no evaluable), como quedó documentado en el commit \texttt{098cf85}.

Hay otra vía por la que una opcional puede entrar sin que nadie lo decida: la sección \texttt{Additional\allowbreak{}Constraints} del lector PGMX instancia cualquier clase cuyo nombre venga en el fichero (ver \hyperref[nota:89]{nota~89}; allí no hay nada que corregir hoy).

Por qué pasa: \texttt{Constraint\allowbreak{}Manager.\allowbreak{}build\allowbreak{}Constraint\allowbreak{}List(\allowbreak{}tipo,\allowbreak{}\ include\allowbreak{}Optionals)} incluye las opcionales solo si se le pasa \texttt{true}, pero en producción solo se llama la versión de un argumento, que pasa \texttt{false} (desde el constructor de \texttt{Prob\allowbreak{}Net}, \texttt{set\allowbreak{}Network\allowbreak{}Type} y \texttt{get\allowbreak{}Additional\allowbreak{}Constraints}). Una restricción opcional solo entra en una red si un tipo la sube a YES o si alguien la añade con \texttt{add\allowbreak{}Constraint}.

La prueba \texttt{Every\allowbreak{}Constraint\allowbreak{}Reaches\allowbreak{}ANetwork\allowbreak{}Test} (commit \texttt{e2158e7}) fija la situación de hoy a la espera de la decisión: cuatro restricciones se añaden a mano (\texttt{Max\allowbreak{}Num\allowbreak{}Parents} y \texttt{Model\allowbreak{}Network\allowbreak{}Constraint} desde el aprendizaje; \texttt{Only\allowbreak{}Discrete\allowbreak{}Variables} y \texttt{Only\allowbreak{}Continuous\allowbreak{}Variables} desde el diálogo de propiedades de la red), dos las sube el tipo MDP (proceso de decisión de Markov), y las ocho de arriba no entran en ninguna red.
\paragraph{El código}
core/src/main/java/org/openmarkov/core/model/network/constraint/ConstraintManager.java:57-66

\begin{Shaded}
\begin{Highlighting}[]
\KeywordTok{public} \BuiltInTok{ArrayList}\OperatorTok{\textless{}}\NormalTok{PNConstraint}\OperatorTok{\textgreater{}} \FunctionTok{buildConstraintList}\OperatorTok{(}\NormalTok{NetworkType type}\OperatorTok{,} \DataTypeTok{boolean}\NormalTok{ includeOptionals}\OperatorTok{)} \OperatorTok{\{}
    \CommentTok{// Init the list with those constraints that have the default value set to YES}
    \BuiltInTok{ArrayList}\OperatorTok{\textless{}}\NormalTok{PNConstraint}\OperatorTok{\textgreater{}}\NormalTok{ constraints }\OperatorTok{=} \KeywordTok{new} \BuiltInTok{ArrayList}\OperatorTok{\textless{}\textgreater{}();}
    \ControlFlowTok{for} \OperatorTok{(}\BuiltInTok{Class}\OperatorTok{\textless{}?} \KeywordTok{extends}\NormalTok{ PNConstraint}\OperatorTok{\textgreater{}}\NormalTok{ constraintClass }\OperatorTok{:}\NormalTok{ defaultConstraintBehaviors}\OperatorTok{.}\FunctionTok{keySet}\OperatorTok{())} \OperatorTok{\{}
        \DataTypeTok{boolean}\NormalTok{ isMandatoryConstraint }\OperatorTok{=} \FunctionTok{getDefaultBehavior}\OperatorTok{(}\NormalTok{constraintClass}\OperatorTok{)} \OperatorTok{==}\NormalTok{ ConstraintBehavior}\OperatorTok{.}\FunctionTok{YES}\OperatorTok{;}
        \DataTypeTok{boolean}\NormalTok{ isOptionalConstraintToUse }\OperatorTok{=}\NormalTok{ includeOptionals }\OperatorTok{\&\&} \FunctionTok{getDefaultBehavior}\OperatorTok{(}\NormalTok{constraintClass}\OperatorTok{)} \OperatorTok{==}\NormalTok{ ConstraintBehavior}\OperatorTok{.}\FunctionTok{OPTIONAL}\OperatorTok{;}
        \ControlFlowTok{if} \OperatorTok{(}\NormalTok{isMandatoryConstraint }\OperatorTok{||}\NormalTok{ isOptionalConstraintToUse}\OperatorTok{)} \OperatorTok{\{}
            \FunctionTok{addIfInstantiable}\OperatorTok{(}\NormalTok{constraints}\OperatorTok{,}\NormalTok{ constraintClass}\OperatorTok{);}
        \OperatorTok{\}}
    \OperatorTok{\}}
\end{Highlighting}
\end{Shaded}

core/src/main/java/org/openmarkov/core/model/network/constraint/ConstraintManager.java:104-106

\begin{Shaded}
\begin{Highlighting}[]
\KeywordTok{public} \DataTypeTok{final} \BuiltInTok{ArrayList}\OperatorTok{\textless{}}\NormalTok{PNConstraint}\OperatorTok{\textgreater{}} \FunctionTok{buildConstraintList}\OperatorTok{(}\NormalTok{NetworkType type}\OperatorTok{)} \OperatorTok{\{}
    \ControlFlowTok{return} \FunctionTok{buildConstraintList}\OperatorTok{(}\NormalTok{type}\OperatorTok{,} \KeywordTok{false}\OperatorTok{);}
\OperatorTok{\}}
\end{Highlighting}
\end{Shaded}
\paragraph{Cómo arreglarlo}
Hay tres salidas, y elegir es cosa del equipo:

\begin{itemize}
\tightlist
\item
  \begin{enumerate}
  \def\labelenumi{(\arabic{enumi})}
  \tightlist
  \item
    Cada tipo de red declara YES las opcionales que quiere, y las que no quiera ningún tipo se borran.
  \end{enumerate}
\item
  \begin{enumerate}
  \def\labelenumi{(\arabic{enumi})}
  \setcounter{enumi}{1}
  \tightlist
  \item
    Las activa el usuario desde las propiedades de la red y el formato \texttt{.\allowbreak{}pgmx} las guarda. Esto exige que la sección de restricciones adicionales del fichero escriba y lea lo mismo (ver \hyperref[nota:89]{nota~89}).
  \end{enumerate}
\item
  \begin{enumerate}
  \def\labelenumi{(\arabic{enumi})}
  \setcounter{enumi}{2}
  \tightlist
  \item
    Desaparece OPTIONAL y todo es YES o NO por tipo.
  \end{enumerate}
\end{itemize}

Detalle latente que hay que resolver a la vez: las anulaciones de un tipo solo tratan YES y NO (\texttt{Constraint\allowbreak{}Manager}, líneas 69-84); un tipo que declarase OPTIONAL no conseguiría nada.

Tras decidir, vaciar la lista \texttt{WITHOUT\_\allowbreak{}A\_\allowbreak{}NETWORK} de \texttt{Every\allowbreak{}Constraint\allowbreak{}Reaches\allowbreak{}ANetwork\allowbreak{}Test}, que es la prueba de regresión. Si una restricción pasa a una red, revisar que sus dos vías (la de \texttt{check\allowbreak{}Prob\allowbreak{}Net} y la copia en cada edición) coincidan.

% ══════════════════════════════════════════════════════════════════════
\section{Evidencia, inferencia y variables numéricas}\label{sec:inferencia}

La propagación de la evidencia, la entrada en modo inferencia y la discretización de las variables numéricas que se hace antes de resolver con un algoritmo exacto.

\setcounter{subsection}{39}
\subsection[Una variable numérica sin precisión se discretiza a un único valor cero]{Una variable numérica de azar sin precisión se discretiza a un único valor cero y el modelo responde con otros números}\label{err:40}
\begin{ficha}
\fichakey{Estado}Presente
\fichakey{Módulo}core, io
\fichakey{Paquete}org.openmarkov.core.model.network; org.openmarkov.io.probmodel.reader
\fichakey{Ficheros}\path{core/src/main/java/org/openmarkov/core/model/network/Variable.java:170-175} \newline \path{core/src/main/java/org/openmarkov/core/model/network/Variable.java:660-662} \newline \path{core/src/main/java/org/openmarkov/core/model/network/ProbNetOperations.java:495} \newline \path{core/src/main/java/org/openmarkov/core/model/network/ProbNetOperations.java:528-535} \newline \path{io/src/main/java/org/openmarkov/io/probmodel/reader/PGMXReader_0_2.java:1006-1010}
\fichakey{Comprobado}Leyendo el código y ejecutándolo
\fichakey{Esfuerzo}Pequeño (horas)
\fichakey{Decisión de equipo}\textcolor{decisionrojo}{\textbf{Sí.}} Qué significa una precisión cero: que no se redondea, o que es un valor no válido que se sustituye por uno por defecto al leer o al construir la variable
\fichakey{Relacionados}\hyperref[err:44]{error~44}, \hyperref[err:45]{error~45}, \hyperref[err:169]{error~169}
\end{ficha}
\paragraph{Qué pasa}
Si se abre un \texttt{.\allowbreak{}pgmx} con una variable numérica de azar que no trae \texttt{\textless{}Precision\textgreater{}} y cuyo valor depende de sus padres, al evaluar o propagar la variable se trata como si valiera siempre cero. Lo que el modelo decía que valía se pierde sin aviso ni excepción, y los resultados cambian.

Medido con un diagrama de influencia de tres nodos escrito a mano: P (azar, dos estados al 50 \%) → N (numérica de azar, vale 3 si P=p0 y 7 si P=p1) → U (utilidad en árbol: 0 si N ≤ 5, 10 si N \textgreater{} 5).

\begin{Verbatim}[breaklines,breakanywhere,fontsize=\small]
precision de N = 0.01 -> estados de N discretizada: [3, 7] -> utilidad esperada = 5.0
precision de N = 0.0  -> estados de N discretizada: [0]    -> utilidad esperada = 0.0
\end{Verbatim}

Y directamente: con \texttt{new\ Variable(\allowbreak{}"X")}, \texttt{round(\allowbreak{}12.\allowbreak{}5)=0.\allowbreak{}0}, \texttt{round(\allowbreak{}-3.\allowbreak{}7)=-0.\allowbreak{}0}, \texttt{round(\allowbreak{}0)=0.\allowbreak{}0}; con valores 3, 7 y −2 la variable convertida sale con los estados «-0» y «0».

Si la variable numérica tiene hallazgo (por ejemplo, porque un potencial delta lo induce, como el «Time in treatment» de los modelos de Markov), no pasa por esta rama y no le afecta: quitar \texttt{\textless{}Precision\textgreater{}} de \texttt{0\_\allowbreak{}miscellaneous/\allowbreak{}networks/\allowbreak{}mid/\allowbreak{}MID-Chancellor.\allowbreak{}pgmx} no cambia su utilidad esperada (50608,78 en los dos casos).

Por qué pasa. Antes de resolver con un algoritmo exacto (eliminación de variables: propagación, evaluación, análisis de coste-efectividad, evolución temporal), \texttt{Prob\allowbreak{}Net\allowbreak{}Operations.\allowbreak{}convert\allowbreak{}Numerical\allowbreak{}Variables\allowbreak{}To\allowbreak{}FS} convierte cada variable numérica de azar que no tenga hallazgo en una variable de estados finitos. Para ello recorre todas las configuraciones de los padres, calcula el valor que toma la variable en cada una y lo redondea a la precisión de la variable con \texttt{old\allowbreak{}Variable.\allowbreak{}round(\allowbreak{}.\allowbreak{}.\allowbreak{}.\allowbreak{})} (línea 495). Cada valor distinto tras el redondeo se convierte en un estado.

\texttt{Variable.\allowbreak{}round} calcula \texttt{Math.\allowbreak{}round(\allowbreak{}value\ /\allowbreak{}\ precision)\ *\ precision}. Si la precisión es 0, la división da infinito (o \texttt{Na\allowbreak{}N} si el valor también es 0), \texttt{Math.\allowbreak{}round} devuelve el mayor (o menor) entero largo, o 0, y al multiplicar por 0 el resultado es siempre 0,0 (o −0,0 para valores negativos). Es decir: todos los valores se redondean a cero. La variable discretizada sale con un único estado «0» (o con dos, «-0» y «0», si había valores negativos y positivos, porque la lista de estados distingue 0,0 de −0,0), y la tabla que la define apunta todas las configuraciones a ese estado.

La precisión cero sí aparece en la práctica:

\begin{itemize}
\tightlist
\item
  El lector del formato propio (\texttt{PGMXReader\_\allowbreak{}0\_\allowbreak{}2.\allowbreak{}get\allowbreak{}XMLContinuous\allowbreak{}Variable}, líneas 1006-1010) pone 0,0 cuando al elemento \texttt{\textless{}Variable\textgreater{}} de una variable numérica le falta el hijo \texttt{\textless{}Precision\textgreater{}} (y también, claro, si trae \texttt{\textless{}Precision\textgreater{}0.\allowbreak{}0\textless{}/\allowbreak{}Precision\textgreater{}}). Un fichero escrito por OpenMarkov siempre lleva la etiqueta, pero uno escrito a mano o por otra herramienta puede no llevarla; ninguna de las redes de ejemplo del repositorio la omite.
\item
  El constructor \texttt{new\ Variable(\allowbreak{}String\ name)} (líneas 170-175), el de variables numéricas sobre toda la recta, deja la precisión en 0,0, mientras que el resto de constructores la dejan en 0,01. Hoy en producción solo se usa para nodos de utilidad y variables auxiliares que no pasan por la discretización, así que el camino real es el del fichero.
\item
  Desde la ventana de propiedades del nodo no se puede poner cero: el desplegable de precisión solo ofrece 1, 0,1, 0,01, 0,001 y 0,0001.
\end{itemize}
\paragraph{El código}
core/src/main/java/org/openmarkov/core/model/network/Variable.java:170-175, 660-662

\begin{Shaded}
\begin{Highlighting}[]
    \KeywordTok{public} \FunctionTok{Variable}\OperatorTok{(}\BuiltInTok{String}\NormalTok{ name}\OperatorTok{)} \OperatorTok{\{}
        
        \KeywordTok{this}\OperatorTok{(}\NormalTok{name}\OperatorTok{,} \KeywordTok{new} \BuiltInTok{State}\OperatorTok{[]\{}\KeywordTok{new} \BuiltInTok{State}\OperatorTok{(}\StringTok{"0"}\OperatorTok{)\},}
             \KeywordTok{new} \FunctionTok{PartitionedInterval}\OperatorTok{(}\KeywordTok{false}\OperatorTok{,} \BuiltInTok{Double}\OperatorTok{.}\FunctionTok{NEGATIVE\_INFINITY}\OperatorTok{,} \BuiltInTok{Double}\OperatorTok{.}\FunctionTok{POSITIVE\_INFINITY}\OperatorTok{,} \KeywordTok{false}\OperatorTok{),} \FloatTok{0.0}\OperatorTok{);}
        \KeywordTok{this}\OperatorTok{.}\FunctionTok{variableType} \OperatorTok{=}\NormalTok{ VariableType}\OperatorTok{.}\FunctionTok{NUMERIC}\OperatorTok{;}
    \OperatorTok{\}}
\KeywordTok{...}
    \KeywordTok{public} \DataTypeTok{double} \FunctionTok{round}\OperatorTok{(}\DataTypeTok{double}\NormalTok{ value}\OperatorTok{)} \OperatorTok{\{}
        \ControlFlowTok{return} \BuiltInTok{Math}\OperatorTok{.}\FunctionTok{round}\OperatorTok{(}\NormalTok{value }\OperatorTok{/}\NormalTok{ precision}\OperatorTok{)} \OperatorTok{*}\NormalTok{ precision}\OperatorTok{;}
    \OperatorTok{\}}
\end{Highlighting}
\end{Shaded}

io/src/main/java/org/openmarkov/io/probmodel/reader/PGMXReader\_0\_2.java:1006-1010

\begin{Shaded}
\begin{Highlighting}[]
        \CommentTok{// Precision}
        \BuiltInTok{Element}\NormalTok{ xMLPrecision }\OperatorTok{=}\NormalTok{ variableElement}\OperatorTok{.}\FunctionTok{getChild}\OperatorTok{(}\NormalTok{XMLTags}\OperatorTok{.}\FunctionTok{PRECISION}\OperatorTok{.}\FunctionTok{toString}\OperatorTok{());}
        \DataTypeTok{double}\NormalTok{ precision }\OperatorTok{=}\NormalTok{ xMLPrecision }\OperatorTok{!=} \KeywordTok{null} \OperatorTok{?} \BuiltInTok{Double}\OperatorTok{.}\FunctionTok{parseDouble}\OperatorTok{(}\NormalTok{xMLPrecision}\OperatorTok{.}\FunctionTok{getText}\OperatorTok{())} \OperatorTok{:} \FloatTok{0.0}\OperatorTok{;}
        
        \ControlFlowTok{return} \KeywordTok{new} \FunctionTok{Variable}\OperatorTok{(}\NormalTok{variableName}\OperatorTok{,}\NormalTok{ leftClosed}\OperatorTok{,}\NormalTok{ min}\OperatorTok{,}\NormalTok{ max}\OperatorTok{,}\NormalTok{ rightClosed}\OperatorTok{,}\NormalTok{ precision}\OperatorTok{);}
\end{Highlighting}
\end{Shaded}

core/src/main/java/org/openmarkov/core/model/network/ProbNetOperations.java:494-499

\begin{Shaded}
\begin{Highlighting}[]
\NormalTok{                        scalarValue }\OperatorTok{=}\NormalTok{ oldPotential}\OperatorTok{.}\FunctionTok{tableProject}\OperatorTok{(}\NormalTok{configuration}\OperatorTok{,}\NormalTok{ inferenceOptions}\OperatorTok{).}\FunctionTok{getValues}\OperatorTok{()[}\DecValTok{0}\OperatorTok{];}
\NormalTok{                        scalarValue }\OperatorTok{=}\NormalTok{ oldVariable}\OperatorTok{.}\FunctionTok{round}\OperatorTok{(}\NormalTok{scalarValue}\OperatorTok{);}
\NormalTok{                        projectedValues}\OperatorTok{[}\NormalTok{index}\OperatorTok{++]} \OperatorTok{=}\NormalTok{ scalarValue}\OperatorTok{;}
                        \ControlFlowTok{if} \OperatorTok{(!}\NormalTok{newStates}\OperatorTok{.}\FunctionTok{contains}\OperatorTok{(}\NormalTok{scalarValue}\OperatorTok{))} \OperatorTok{\{}
\NormalTok{                            newStates}\OperatorTok{.}\FunctionTok{add}\OperatorTok{(}\NormalTok{scalarValue}\OperatorTok{);}
                        \OperatorTok{\}}
\end{Highlighting}
\end{Shaded}
\paragraph{Cómo arreglarlo}
Hay que decidir qué significa una precisión cero. Dos opciones razonables:

\begin{itemize}
\tightlist
\item
  \begin{enumerate}
  \def\labelenumi{(\alph{enumi})}
  \tightlist
  \item
    «No redondear»: \texttt{Variable.\allowbreak{}round} devuelve el valor tal cual cuando la precisión es menor o igual que cero. Con esta opción conviene además normalizar −0,0 a 0,0 al construir la lista de estados, para no fabricar dos estados «-0» y «0».
  \end{enumerate}
\item
  \begin{enumerate}
  \def\labelenumi{(\alph{enumi})}
  \setcounter{enumi}{1}
  \tightlist
  \item
    «Valor no válido»: el lector y el constructor \texttt{Variable(\allowbreak{}String)} usan el valor por defecto de la clase (0,01) cuando no hay precisión, y \texttt{set\allowbreak{}Precision} rechaza el cero.
  \end{enumerate}
\end{itemize}

Sitios donde la regla tiene que cumplirse:

\begin{itemize}
\tightlist
\item
  \texttt{Variable.\allowbreak{}round}, usado en \texttt{Prob\allowbreak{}Net\allowbreak{}Operations.\allowbreak{}convert\allowbreak{}Numerical\allowbreak{}Variables\allowbreak{}To\allowbreak{}FS} (líneas 495 y 582) y en \texttt{Variable.\allowbreak{}get\allowbreak{}State\allowbreak{}Index(\allowbreak{}double)} (línea 343);
\item
  el valor por defecto del lector \texttt{PGMXReader\_\allowbreak{}0\_\allowbreak{}2.\allowbreak{}get\allowbreak{}XMLContinuous\allowbreak{}Variable} (línea 1008);
\item
  el constructor \texttt{Variable(\allowbreak{}String)};
\item
  los usos de la precisión como paso: \texttt{Variable.\allowbreak{}get\allowbreak{}Default\allowbreak{}Interval} (suma la precisión para construir intervalos) y el selector numérico de \texttt{Add\allowbreak{}Finding\allowbreak{}Dialog} (usa la precisión como incremento).
\end{itemize}

La forma en que se escriben los límites de los tramos con una precisión como 0,0 es \hyperref[err:169]{error~169}, y conviene decidir las dos cosas a la vez.

Prueba de regresión: una red P → N (numérica, precisión 0, valores 3 y 7) convertida con \texttt{convert\allowbreak{}Numerical\allowbreak{}Variables\allowbreak{}To\allowbreak{}FS} debe dar dos estados; y leer un \texttt{.\allowbreak{}pgmx} sin \texttt{\textless{}Precision\textgreater{}} como el descrito arriba y evaluarlo con \texttt{VEEvaluation} debe dar 5,0.

\setcounter{subsection}{40}
\subsection[Un hallazgo sobre un nodo numérico de azar hace fallar la propagación]{Poner un hallazgo, en modo inferencia, sobre un nodo numérico de azar hace fallar la propagación con un NullPointerException}\label{err:41}
\begin{ficha}
\fichakey{Estado}Presente
\fichakey{Módulo}core, inference
\fichakey{Paquete}org.openmarkov.core.model.network; org.openmarkov.inference.algorithm.variableElimination.tasks
\fichakey{Ficheros}\path{core/src/main/java/org/openmarkov/core/model/network/ProbNetOperations.java:63-70,203-225,324-333,431-459} \newline \path{inference/src/main/java/org/openmarkov/inference/algorithm/variableElimination/tasks/VEPropagation.java:117-157,228-242} \newline \path{inference/src/main/java/org/openmarkov/inference/algorithm/variableElimination/tasks/VariableElimination.java:95-100} \newline \path{gui/src/main/java/org/openmarkov/gui/window/edition/networkEditorPanel/EvidenceManager.java:581-600}
\fichakey{Comprobado}Leyendo el código y ejecutándolo
\fichakey{Esfuerzo}Medio (días)
\fichakey{Decisión de equipo}No
\fichakey{Relacionados}\hyperref[err:42]{error~42}, \hyperref[err:44]{error~44}, \hyperref[err:45]{error~45}, \hyperref[err:46]{error~46}, \hyperref[nota:55]{nota~55}
\end{ficha}
\paragraph{Qué pasa}
El usuario abre una red con nodos numéricos de azar, pasa a modo inferencia y añade un hallazgo a uno de ellos (el diálogo pide un valor numérico). La propagación falla con \texttt{Null\allowbreak{}Pointer\allowbreak{}Exception}: la excepción no es de las que \texttt{Evidence\allowbreak{}Manager.\allowbreak{}do\allowbreak{}Propagation} recoge y llega al gestor global, que muestra la ventana de error inesperado; el panel no enseña probabilidades para ese caso de evidencia.

Medido sobre \texttt{0\_\allowbreak{}miscellaneous/\allowbreak{}networks/\allowbreak{}mid/\allowbreak{}DMHEE/\allowbreak{}MID-dmhee-2.\allowbreak{}5.\allowbreak{}pgmx} (un modelo de Markov con decisiones): con un hallazgo posterior de valor 1 sobre el nodo numérico de azar «Duration of treatment {[}0{]}», \texttt{get\allowbreak{}Posterior\allowbreak{}Values} lanza \texttt{Null\allowbreak{}Pointer\allowbreak{}Exception:\ Cannot\ invoke\ "Node.\allowbreak{}get\allowbreak{}Neighbors(\allowbreak{})"\ because\ "node"\ is\ null} en \texttt{Prob\allowbreak{}Net\allowbreak{}Operations.\allowbreak{}remove\allowbreak{}Unreachable\allowbreak{}Nodes(\allowbreak{}Prob\allowbreak{}Net\allowbreak{}Operations.\allowbreak{}java:218)}, llamado desde \texttt{VEPropagation.\allowbreak{}prune\allowbreak{}Network(\allowbreak{}VEPropagation.\allowbreak{}java:241)}. Con un hallazgo sobre un nodo de estados finitos («State {[}1{]}») la propagación termina bien.

Por qué pasa. En modo inferencia, cada vez que cambia la evidencia, \texttt{Evidence\allowbreak{}Manager.\allowbreak{}do\allowbreak{}Propagation} (módulo \texttt{gui}) calcula las probabilidades a posteriori con \texttt{VEPropagation} (eliminación de variables, módulo \texttt{inference}), pasando como evidencia posterior a la resolución los hallazgos que el usuario ha puesto.

\texttt{VEPropagation.\allowbreak{}resolve} prepara la red: la expande si es temporal y luego \texttt{exact\allowbreak{}Algorithms\allowbreak{}Preprocessing} (\texttt{Variable\allowbreak{}Elimination}, líneas 95-100) llama a \texttt{Task\allowbreak{}Utilities.\allowbreak{}discretize\allowbreak{}Non\allowbreak{}Observed\allowbreak{}Numeric\allowbreak{}Variables(\allowbreak{}prob\allowbreak{}Net,\allowbreak{}\ get\allowbreak{}Pre\allowbreak{}Resolution\allowbreak{}Evidence(\allowbreak{}))}, es decir, a \texttt{Prob\allowbreak{}Net\allowbreak{}Operations.\allowbreak{}convert\allowbreak{}Numerical\allowbreak{}Variables\allowbreak{}To\allowbreak{}FS}. Esa conversión solo tiene en cuenta la evidencia \textbf{previa} a la resolución: un nodo numérico de azar sin hallazgo previo se convierte en uno de estados finitos poniéndole un objeto \texttt{Variable} nuevo con el mismo nombre (\texttt{node.\allowbreak{}set\allowbreak{}Variable(\allowbreak{}new\allowbreak{}Variable)}). El hallazgo posterior sigue apuntando a la variable numérica original.

Después, para cada variable de interés, \texttt{prune\allowbreak{}Network} (líneas 228-242) mete en la lista de variables a conservar las de todos los hallazgos y llama a \texttt{Prob\allowbreak{}Net\allowbreak{}Operations.\allowbreak{}get\allowbreak{}Pruned}. En \texttt{remove\allowbreak{}Unreachable\allowbreak{}Nodes} (líneas 210-212) se hace \texttt{nodes\allowbreak{}To\allowbreak{}Keep.\allowbreak{}add(\allowbreak{}prob\allowbreak{}Net.\allowbreak{}get\allowbreak{}Node(\allowbreak{}variable))} sin comprobar el resultado: para la variable original \texttt{get\allowbreak{}Node} devuelve \texttt{null}, el \texttt{null} entra en el conjunto y en la línea 218 \texttt{node.\allowbreak{}get\allowbreak{}Neighbors(\allowbreak{})} lanza la excepción. La función hermana \texttt{get\allowbreak{}Evidence\allowbreak{}Nodes} (líneas 324-333) sí descarta los \texttt{null}.

Comprobado aparte que, tras la expansión, \texttt{get\allowbreak{}Node} con la variable original encuentra el nodo y, tras la discretización, ya no (buscando por nombre sí).
\paragraph{El código}
core/src/main/java/org/openmarkov/core/model/network/ProbNetOperations.java:209-218

\begin{Shaded}
\begin{Highlighting}[]
        \CommentTok{// Store nodes of variablesOfInterest in nodesToKeep}
        \ControlFlowTok{for} \OperatorTok{(}\NormalTok{Variable variable }\OperatorTok{:}\NormalTok{ variablesOfInterest}\OperatorTok{)} \OperatorTok{\{}
\NormalTok{            nodesToKeep}\OperatorTok{.}\FunctionTok{add}\OperatorTok{(}\NormalTok{probNet}\OperatorTok{.}\FunctionTok{getNode}\OperatorTok{(}\NormalTok{variable}\OperatorTok{));}
        \OperatorTok{\}}
        
        \CommentTok{// Add neighbors of variablesOfInterest as nodesToKeep and store them in}
        \CommentTok{// nodesToExplore}
        \BuiltInTok{HashSet}\OperatorTok{\textless{}}\BuiltInTok{Node}\OperatorTok{\textgreater{}}\NormalTok{ nodesToKeepClon }\OperatorTok{=} \KeywordTok{new} \BuiltInTok{HashSet}\OperatorTok{\textless{}\textgreater{}(}\NormalTok{nodesToKeep}\OperatorTok{);}
        \ControlFlowTok{for} \OperatorTok{(}\BuiltInTok{Node}\NormalTok{ node }\OperatorTok{:}\NormalTok{ nodesToKeepClon}\OperatorTok{)} \OperatorTok{\{}
            \BuiltInTok{List}\OperatorTok{\textless{}}\BuiltInTok{Node}\OperatorTok{\textgreater{}}\NormalTok{ neighbors }\OperatorTok{=}\NormalTok{ node}\OperatorTok{.}\FunctionTok{getNeighbors}\OperatorTok{();}
\end{Highlighting}
\end{Shaded}
\paragraph{Cómo arreglarlo}
Saltar el \texttt{null} en \texttt{remove\allowbreak{}Unreachable\allowbreak{}Nodes} evita la excepción pero no basta: el hallazgo sobre la variable numérica se ignoraría sin decir nada, la misma pérdida silenciosa que describe \hyperref[nota:55]{nota~55} (allí no hay nada que corregir).

El arreglo completo es que la preparación de la red tenga en cuenta toda la evidencia (previa y posterior) al discretizar y al absorber nodos numéricos, o que tras la discretización los hallazgos se traduzcan a las variables nuevas (buscando por nombre) antes de podar.

Sitios afectados:

\begin{itemize}
\tightlist
\item
  \texttt{Variable\allowbreak{}Elimination.\allowbreak{}exact\allowbreak{}Algorithms\allowbreak{}Preprocessing};
\item
  \texttt{Prob\allowbreak{}Net\allowbreak{}Operations.\allowbreak{}convert\allowbreak{}Numerical\allowbreak{}Variables\allowbreak{}To\allowbreak{}FS};
\item
  \texttt{VEPropagation.\allowbreak{}prune\allowbreak{}Network} y \texttt{get\allowbreak{}All\allowbreak{}Evidence};
\item
  \texttt{remove\allowbreak{}Unreachable\allowbreak{}Nodes}, que debería tratar una variable ausente igual que \texttt{get\allowbreak{}Evidence\allowbreak{}Nodes}, o fallar con una excepción del dominio con mensaje.
\end{itemize}

Prueba: \texttt{0\_\allowbreak{}miscellaneous/\allowbreak{}networks/\allowbreak{}mid/\allowbreak{}DMHEE/\allowbreak{}MID-dmhee-2.\allowbreak{}5.\allowbreak{}pgmx} con un hallazgo posterior sobre «Duration of treatment {[}0{]}»; la propagación debe terminar y el valor observado debe influir en el resultado.

\setcounter{subsection}{41}
\subsection[Entrar en modo inferencia falla si hay nodos de utilidad intermedios]{Entrar en modo inferencia en una red con nodos de utilidad que otro nodo de utilidad suma o multiplica (ID-arthronet, ID-mediastinet, ID-used-car-buyer\ldots), o con un nodo numérico de azar sin observar, falla con NullPointerException}\label{err:42}
\begin{ficha}
\fichakey{Estado}Presente
\fichakey{Módulo}inference, core
\fichakey{Paquete}org.openmarkov.inference.algorithm.variableElimination.tasks; org.openmarkov.core.model.network
\fichakey{Ficheros}\path{inference/src/main/java/org/openmarkov/inference/algorithm/variableElimination/tasks/VEPropagation.java:141-157, 175-217, 228-242} \newline \path{inference/src/main/java/org/openmarkov/inference/algorithm/variableElimination/tasks/VariableElimination.java:95-97} \newline \path{core/src/main/java/org/openmarkov/core/inference/BasicOperations.java:146-178} \newline \path{core/src/main/java/org/openmarkov/core/model/network/ProbNetOperations.java:431-459, 1050-1054} \newline \path{gui/src/main/java/org/openmarkov/gui/window/edition/networkEditorPanel/EvidenceManager.java:584-588} \newline \path{gui/src/main/java/org/openmarkov/gui/window/edition/networkEditorPanel/InferencePresenter.java:60-69, 118-126} \newline \path{gui/src/main/java/org/openmarkov/gui/window/MainPanelListenerAssistant.java:145-163}
\fichakey{Comprobado}Leyendo el código y ejecutándolo
\fichakey{Esfuerzo}Medio (días)
\fichakey{Decisión de equipo}\textcolor{decisionrojo}{\textbf{Sí.}} Qué debe enseñar el modo inferencia en un nodo de utilidad que otro nodo de utilidad suma o multiplica: su utilidad esperada, calculada aparte, o nada
\fichakey{Tapado por}\hyperref[err:46]{error~46}: en nueve redes, entre ellas MID-HPV, el cálculo de rangos falla antes.
\fichakey{Relacionados}\hyperref[err:8]{error~8}, \hyperref[err:41]{error~41}, \hyperref[err:44]{error~44}, \hyperref[err:45]{error~45}, \hyperref[err:46]{error~46}, \hyperref[err:108]{error~108}, \hyperref[err:112]{error~112}, \hyperref[err:113]{error~113}, \hyperref[nota:109]{nota~109}
\end{ficha}
\paragraph{Qué pasa}
Con \texttt{0\_\allowbreak{}miscellaneous/\allowbreak{}networks/\allowbreak{}id/\allowbreak{}ID-arthronet.\allowbreak{}pgmx} abierto y sin evidencia, pulsar el botón de modo inferencia termina en la ventana de error no previsto con un \texttt{Null\allowbreak{}Pointer\allowbreak{}Exception}, y la red vuelve a modo edición: el modo inferencia no se puede usar en esa red. El camino es \texttt{Main\allowbreak{}Panel\allowbreak{}Listener\allowbreak{}Assistant} (líneas 155-171) → \texttt{Inference\allowbreak{}Handler.\allowbreak{}set\allowbreak{}Working\allowbreak{}Mode} → \texttt{Evidence\allowbreak{}Manager.\allowbreak{}update\allowbreak{}Individual\allowbreak{}Probabilities\allowbreak{}And\allowbreak{}Utilities} → \texttt{do\allowbreak{}Propagation}, que construye un \texttt{VEPropagation} (propagación por eliminación de variables) y le pide los resultados de \textbf{todas} las variables de la red (líneas 585-589). \texttt{Main\allowbreak{}Panel\allowbreak{}Listener\allowbreak{}Assistant} recoge la excepción, vuelve al modo edición y la relanza envuelta; \texttt{OMException\allowbreak{}Handler}, al encontrar debajo un \texttt{Null\allowbreak{}Pointer\allowbreak{}Exception}, abre el diálogo de error no previsto.

Medido con un pequeño programa de prueba que hace las mismas llamadas que \texttt{do\allowbreak{}Propagation} (primero el cálculo de rangos de las utilidades, después la propagación) sobre todas las redes que se pueden leer del repositorio. Terminan con este \texttt{Null\allowbreak{}Pointer\allowbreak{}Exception}, con el cálculo de rangos sin fallar antes:

\begin{itemize}
\tightlist
\item
  en \texttt{0\_\allowbreak{}miscellaneous/\allowbreak{}networks}: \texttt{id/\allowbreak{}ID-arthronet.\allowbreak{}pgmx}, \texttt{id/\allowbreak{}ID-mediastinet.\allowbreak{}pgmx}, \texttt{id/\allowbreak{}ID-used-car-buyer.\allowbreak{}pgmx}, \texttt{dan/\allowbreak{}DAN-mediastinet.\allowbreak{}pgmx}, \texttt{dan/\allowbreak{}DAN-qale-mediastinet.\allowbreak{}pgmx} y \texttt{dan/\allowbreak{}DAN-used-car-buyer.\allowbreak{}pgmx} (DAN: red de análisis de decisiones);
\item
  en las redes de prueba: once de \texttt{integration\allowbreak{}Tests/\allowbreak{}src/\allowbreak{}main/\allowbreak{}resources/\allowbreak{}networks} (\texttt{id/\allowbreak{}ID-decide-test-sv.\allowbreak{}pgmx} y diez de \texttt{dan}, casi todas con «-sv» en el nombre) y \texttt{core/\allowbreak{}src/\allowbreak{}test/\allowbreak{}resources/\allowbreak{}nets/\allowbreak{}Discrete\allowbreak{}Potential\allowbreak{}Operations/\allowbreak{}IDDecide\allowbreak{}Test\allowbreak{}Symptom.\allowbreak{}pgmx}.
\end{itemize}

En otras nueve (entre ellas \texttt{0\_\allowbreak{}miscellaneous/\allowbreak{}networks/\allowbreak{}mid/\allowbreak{}MID-HPV.\allowbreak{}pgmx}) el cálculo de rangos falla antes (\hyperref[err:46]{error~46}); arreglado ese, llegarían a este. Las versiones con dos criterios (\texttt{ID-arthronet-ce.\allowbreak{}pgmx}, \texttt{ID-mediastinet-ce.\allowbreak{}pgmx}) no tienen nodos de utilidad intermedios y se propagan bien.

Hay una segunda forma de llegar, sin ninguna red del repositorio que la tenga: un nodo numérico de azar sin hallazgo cuyo valor se pueda enumerar a partir de sus padres. Medido con una red bayesiana construida por programa, Disease (estados finitos) → N (numérico, relación exacta 3 o 7): pedir todas las variables lanza la misma excepción; pedir solo Disease devuelve {[}0,9, 0,1{]}. Si la red tiene además un nodo de utilidad que depende de N, el cálculo de rangos falla antes (\hyperref[err:46]{error~46}).

Por qué pasa. \texttt{VEPropagation.\allowbreak{}resolve} trabaja sobre una copia de la red que \texttt{Variable\allowbreak{}Elimination.\allowbreak{}exact\allowbreak{}Algorithms\allowbreak{}Preprocessing} (líneas 95-97) modifica de dos maneras que dejan fuera variables de la red del usuario:

\begin{itemize}
\tightlist
\item
  \texttt{Task\allowbreak{}Utilities.\allowbreak{}absorb\allowbreak{}All\allowbreak{}Intermediate\allowbreak{}Numeric\allowbreak{}Nodes} (\texttt{Basic\allowbreak{}Operations}, líneas 144-170) absorbe en sus hijos los nodos numéricos que tienen hijos, y los quita de la red. Entre ellos están los nodos de utilidad cuyo valor suma o multiplica otro nodo de utilidad (los antiguos nodos de «supervalor»). Medido: en \texttt{ID-arthronet.\allowbreak{}pgmx} desaparecen 12 de sus 13 nodos de utilidad, en \texttt{ID-mediastinet.\allowbreak{}pgmx} 19 y en \texttt{ID-used-car-buyer.\allowbreak{}pgmx} 5; solo queda el nodo final.
\item
  \texttt{Task\allowbreak{}Utilities.\allowbreak{}discretize\allowbreak{}Non\allowbreak{}Observed\allowbreak{}Numeric\allowbreak{}Variables} (\texttt{Prob\allowbreak{}Net\allowbreak{}Operations.\allowbreak{}convert\allowbreak{}Numerical\allowbreak{}Variables\allowbreak{}To\allowbreak{}FS}) sustituye cada nodo numérico de azar sin hallazgo por uno de estados finitos con un objeto \texttt{Variable} nuevo del mismo nombre, así que buscar el nodo con la variable original ya no lo encuentra.
\end{itemize}

Después, para cada variable de interés, \texttt{resolve} busca la variable por nombre, pero solo para ver si tiene hallazgo (línea 147). Luego llama a \texttt{prune\allowbreak{}Network} con el objeto que le dio el llamador, y \texttt{prune\allowbreak{}Network} hace \texttt{get\allowbreak{}Node\allowbreak{}Ancestors(\allowbreak{}preprocessed\allowbreak{}Network.\allowbreak{}get\allowbreak{}Node(\allowbreak{}variable\allowbreak{}Of\allowbreak{}Interest))} (línea 233). \texttt{get\allowbreak{}Node} devuelve \texttt{null} para una variable que ya no está, y \texttt{Prob\allowbreak{}Net\allowbreak{}Operations.\allowbreak{}get\allowbreak{}Node\allowbreak{}Ancestors} falla al meter ese \texttt{null} en un \texttt{Array\allowbreak{}Deque} (línea 1054). Si se pasara de ahí, \texttt{Invoke\allowbreak{}Variable\allowbreak{}Elimination\allowbreak{}Core} volvería a fallar en \texttt{prob\allowbreak{}Net.\allowbreak{}get\allowbreak{}Node(\allowbreak{}variable\allowbreak{}Of\allowbreak{}Interest).\allowbreak{}get\allowbreak{}Node\allowbreak{}Type(\allowbreak{})} (línea 192).

Es de la misma familia que \hyperref[err:41]{error~41} (la red preparada ya no contiene una variable de la red del usuario), pero allí la variable perdida es la de un hallazgo posterior y la excepción sale en \texttt{remove\allowbreak{}Unreachable\allowbreak{}Nodes}; arreglar uno no arregla el otro.
\paragraph{El código}
inference/src/main/java/org/openmarkov/inference/algorithm/variableElimination/tasks/VEPropagation.java:145-155, 228-235

\begin{Shaded}
\begin{Highlighting}[]
        \ControlFlowTok{if} \OperatorTok{(}\NormalTok{variablesOfInterest }\OperatorTok{!=} \KeywordTok{null}\OperatorTok{)} \OperatorTok{\{}
            \ControlFlowTok{for} \OperatorTok{(}\NormalTok{Variable variableOfInterest }\OperatorTok{:}\NormalTok{ variablesOfInterest}\OperatorTok{)} \OperatorTok{\{}
\NormalTok{                Variable variableOfInterestInProbnet }\OperatorTok{=}\NormalTok{ probNet}\OperatorTok{.}\FunctionTok{getVariable}\OperatorTok{(}\NormalTok{variableOfInterest}\OperatorTok{.}\FunctionTok{getName}\OperatorTok{());}
                \ControlFlowTok{if} \OperatorTok{(}\NormalTok{evidenceVariables}\OperatorTok{.}\FunctionTok{contains}\OperatorTok{(}\NormalTok{variableOfInterestInProbnet}\OperatorTok{))} \OperatorTok{\{}
\NormalTok{                    variablesOfInterestBelongingToEvidence}\OperatorTok{.}\FunctionTok{add}\OperatorTok{(}\NormalTok{variableOfInterestInProbnet}\OperatorTok{);}
                \OperatorTok{\}} \ControlFlowTok{else} \OperatorTok{\{}
\NormalTok{                    ProbNet preprocessedNetwork }\OperatorTok{=} \FunctionTok{pruneNetwork}\OperatorTok{(}\NormalTok{probNet}\OperatorTok{.}\FunctionTok{copy}\OperatorTok{(),}\NormalTok{ variableOfInterest}\OperatorTok{);}
                    \KeywordTok{...}
    \KeywordTok{private}\NormalTok{ ProbNet }\FunctionTok{pruneNetwork}\OperatorTok{(}\NormalTok{ProbNet preprocessedNetwork}\OperatorTok{,}\NormalTok{ Variable variableOfInterest}\OperatorTok{)}
            \KeywordTok{throws}\NormalTok{ IncompatibleEvidenceException}\OperatorTok{.}\FunctionTok{EvidenceIsIncompatibleWithOther} \OperatorTok{\{}
        \CommentTok{//Prune all the nodes except the variable of interest and its ancestors (and the corresponding findings).}
        \BuiltInTok{List}\OperatorTok{\textless{}}\NormalTok{Variable}\OperatorTok{\textgreater{}}\NormalTok{ variablesNotToBePruned }\OperatorTok{=} \KeywordTok{new} \BuiltInTok{ArrayList}\OperatorTok{\textless{}\textgreater{}();}
\NormalTok{        variablesNotToBePruned}\OperatorTok{.}\FunctionTok{add}\OperatorTok{(}\NormalTok{variableOfInterest}\OperatorTok{);}
        \ControlFlowTok{for} \OperatorTok{(}\BuiltInTok{Node}\NormalTok{ node }\OperatorTok{:}\NormalTok{ ProbNetOperations}\OperatorTok{.}\FunctionTok{getNodeAncestors}\OperatorTok{(}\NormalTok{preprocessedNetwork}\OperatorTok{.}\FunctionTok{getNode}\OperatorTok{(}\NormalTok{variableOfInterest}\OperatorTok{)))} \OperatorTok{\{}
\end{Highlighting}
\end{Shaded}

core/src/main/java/org/openmarkov/core/model/network/ProbNetOperations.java:1050-1054

\begin{Shaded}
\begin{Highlighting}[]
    \KeywordTok{public} \DataTypeTok{static} \BuiltInTok{Set}\OperatorTok{\textless{}}\BuiltInTok{Node}\OperatorTok{\textgreater{}} \FunctionTok{getNodeAncestors}\OperatorTok{(}\BuiltInTok{Node}\NormalTok{ node}\OperatorTok{)} \OperatorTok{\{}
        \BuiltInTok{Set}\OperatorTok{\textless{}}\BuiltInTok{Node}\OperatorTok{\textgreater{}}\NormalTok{ ancestors }\OperatorTok{=} \KeywordTok{new} \BuiltInTok{HashSet}\OperatorTok{\textless{}\textgreater{}();}
        
        \BuiltInTok{Deque}\OperatorTok{\textless{}}\BuiltInTok{Node}\OperatorTok{\textgreater{}}\NormalTok{ noExploredNodes }\OperatorTok{=} \KeywordTok{new} \BuiltInTok{ArrayDeque}\OperatorTok{\textless{}\textgreater{}();}
\NormalTok{        noExploredNodes}\OperatorTok{.}\FunctionTok{add}\OperatorTok{(}\NormalTok{node}\OperatorTok{);}
\end{Highlighting}
\end{Shaded}
\paragraph{Cómo arreglarlo}
\texttt{VEPropagation.\allowbreak{}resolve} tiene que traducir cada variable de interés a la red preparada antes de usarla, y decidir qué hace con las que ya no están:

\begin{itemize}
\tightlist
\item
  Nodo numérico de azar discretizado: buscarlo por nombre (como ya hace la línea 147) en \texttt{prune\allowbreak{}Network} y en \texttt{Invoke\allowbreak{}Variable\allowbreak{}Elimination\allowbreak{}Core}, y guardar el resultado con la variable del llamador como clave, que es la que busca la ventana. La ventana no usa la tabla de un nodo numérico de azar (\texttt{Inference\allowbreak{}Presenter}, líneas 118-126, pinta el valor del hallazgo o nada), así que también vale no calcularlo; lo que no puede es lanzar una excepción.
\item
  Nodo de utilidad absorbido: no queda nodo al que preguntar. Hay que decidir entre (a) calcular su utilidad esperada aparte, sobre una copia en la que ese nodo sea el último (quitando antes sus descendientes para que no se absorba), o (b) no calcularla y no pintar nada en ese nodo. La opción (b) obliga a tocar también \texttt{Inference\allowbreak{}Presenter.\allowbreak{}paint\allowbreak{}Inference\allowbreak{}Results\allowbreak{}Utility\allowbreak{}Node} (línea 66, que hace \texttt{individual\allowbreak{}Probabilities.\allowbreak{}get(\allowbreak{}variable).\allowbreak{}get\allowbreak{}Values(\allowbreak{}){[}0{]}} y daría otro \texttt{Null\allowbreak{}Pointer\allowbreak{}Exception}).
\end{itemize}

Otros sitios de la misma regla: \texttt{Prob\allowbreak{}Net\allowbreak{}Operations.\allowbreak{}get\allowbreak{}Node\allowbreak{}Ancestors} y \texttt{remove\allowbreak{}Unreachable\allowbreak{}Nodes} deberían rechazar una variable ausente con una excepción que la nombre, en vez de dejar pasar un \texttt{null}; y el arreglo de \hyperref[err:41]{error~41} (los hallazgos) debe usar la misma traducción de variables.

Prueba de regresión: \texttt{VEPropagation} sobre \texttt{0\_\allowbreak{}miscellaneous/\allowbreak{}networks/\allowbreak{}id/\allowbreak{}ID-arthronet.\allowbreak{}pgmx} y \texttt{0\_\allowbreak{}miscellaneous/\allowbreak{}networks/\allowbreak{}id/\allowbreak{}ID-mediastinet.\allowbreak{}pgmx} con todas las variables de interés y evidencia vacía debe terminar y devolver, para los nodos de utilidad intermedios, lo que se decida; y la red bayesiana Disease → N con todas las variables no debe lanzar excepción.

\setcounter{subsection}{42}
\subsection[Un nodo numérico de azar sin potencial tumba la discretización]{Propagar o evaluar una red con un nodo numérico de azar sin potencial muere con IndexOutOfBoundsException al discretizar}\label{err:43}
\begin{ficha}
\fichakey{Estado}Presente en parte. Desde el commit \texttt{83ff8ab} la ventana principal comprueba, antes de entrar en modo inferencia, evaluar, pedir la estrategia óptima o el árbol de decisión, expandir la red o pulsar el botón de coste-efectividad de la barra, que todo nodo que no sea de decisión tenga potencial, y si no avisa con un mensaje que nombra el nodo. Sigue sin corregir \texttt{ProbNetOperations.convertNumericalVariablesToFS}, que continúa suponiendo que el nodo tiene potencial: el análisis de coste-efectividad lanzado desde el menú Herramientas, el análisis de sensibilidad y cualquier llamador de la API siguen muriendo con \texttt{IndexOutOfBoundsException} sin decir qué nodo está mal.
\fichakey{Módulo}core
\fichakey{Paquete}org.openmarkov.core.model.network
\fichakey{Ficheros}\path{core/src/main/java/org/openmarkov/core/model/network/ProbNetOperations.java:431-443} \newline \path{core/src/main/java/org/openmarkov/core/model/network/ProbNetOperations.java:561-566} \newline \path{core/src/main/java/org/openmarkov/core/inference/tasks/TaskUtilities.java:88-90} \newline \path{inference/src/main/java/org/openmarkov/inference/algorithm/variableElimination/tasks/VariableElimination.java:95-97}
\fichakey{Comprobado}Leyendo el código y ejecutándolo
\fichakey{Esfuerzo}Pequeño (horas)
\fichakey{Decisión de equipo}No
\fichakey{Relacionados}\hyperref[err:70]{error~70}
\end{ficha}
\paragraph{Qué pasa}
Si en un fichero \texttt{.\allowbreak{}pgmx} a un nodo numérico de azar le falta su potencial (el lector lo acepta en silencio), el usuario ve un error inesperado al propagar o evaluar, sin indicación de qué nodo está mal. Es el mismo camino de llegada que \hyperref[err:29]{error~29}; desde la ventana no se crea un nodo así.

Comprobado ejecutando con una red bayesiana A → X, con X numérica de azar sin potencial: \texttt{Prob\allowbreak{}Net\allowbreak{}Operations.\allowbreak{}convert\allowbreak{}Numerical\allowbreak{}Variables\allowbreak{}To\allowbreak{}FS(\allowbreak{}net)} lanza \texttt{Index\allowbreak{}Out\allowbreak{}Of\allowbreak{}Bounds\allowbreak{}Exception:\ Index\ 0\ out\ of\ bounds\ for\ length\ 0} en la línea 443, y \texttt{VEPropagation} (propagación por eliminación de variables, módulo \texttt{inference}) sobre la misma red muere igual al pedir los posteriores, vía \texttt{Variable\allowbreak{}Elimination.\allowbreak{}exact\allowbreak{}Algorithms\allowbreak{}Preprocessing} → \texttt{Task\allowbreak{}Utilities.\allowbreak{}discretize\allowbreak{}Non\allowbreak{}Observed\allowbreak{}Numeric\allowbreak{}Variables}.

Por qué pasa: \texttt{Prob\allowbreak{}Net\allowbreak{}Operations.\allowbreak{}convert\allowbreak{}Numerical\allowbreak{}Variables\allowbreak{}To\allowbreak{}FS} (módulo \texttt{core}) convierte en variables de estados finitos las variables numéricas de azar antes de resolver con un algoritmo exacto. Para cada nodo numérico de azar empieza con \texttt{Potential\ old\allowbreak{}Potential\ =\ node.\allowbreak{}get\allowbreak{}Potentials(\allowbreak{}).\allowbreak{}get(\allowbreak{}0);}. Unas líneas más abajo, para los nodos no numéricos, sí comprueba \texttt{!node.\allowbreak{}get\allowbreak{}Potentials(\allowbreak{}).\allowbreak{}is\allowbreak{}Empty(\allowbreak{})}.

Llaman a este método también la evaluación temporal (\texttt{Temporal\allowbreak{}Evaluation}), la evolución temporal (\texttt{MIDTemporal\allowbreak{}Evolution}) y la expansión de redes temporales (\texttt{Temporal\allowbreak{}Net\allowbreak{}Operations}).
\paragraph{El código}
core/src/main/java/org/openmarkov/core/model/network/ProbNetOperations.java:438-443

\begin{Shaded}
\begin{Highlighting}[]
\ControlFlowTok{for} \OperatorTok{(}\BuiltInTok{Node}\NormalTok{ node }\OperatorTok{:}\NormalTok{ sortedNodes}\OperatorTok{)} \OperatorTok{\{}
\NormalTok{    Variable oldVariable }\OperatorTok{=}\NormalTok{ node}\OperatorTok{.}\FunctionTok{getVariable}\OperatorTok{();}
    \CommentTok{// Should the node be converted}
    \ControlFlowTok{if} \OperatorTok{(}\NormalTok{oldVariable}\OperatorTok{.}\FunctionTok{getVariableType}\OperatorTok{()} \OperatorTok{==}\NormalTok{ VariableType}\OperatorTok{.}\FunctionTok{NUMERIC} \OperatorTok{\&\&}\NormalTok{ node}\OperatorTok{.}\FunctionTok{getNodeType}\OperatorTok{()} \OperatorTok{==}\NormalTok{ NodeType}\OperatorTok{.}\FunctionTok{CHANCE}\OperatorTok{)} \OperatorTok{\{}
\NormalTok{        EvidenceCase configuration }\OperatorTok{=} \KeywordTok{new} \FunctionTok{EvidenceCase}\OperatorTok{(}\NormalTok{evidence}\OperatorTok{);}
\NormalTok{        Potential oldPotential }\OperatorTok{=}\NormalTok{ node}\OperatorTok{.}\FunctionTok{getPotentials}\OperatorTok{().}\FunctionTok{get}\OperatorTok{(}\DecValTok{0}\OperatorTok{);}
\end{Highlighting}
\end{Shaded}
\paragraph{Cómo arreglarlo}
Dejar sin convertir un nodo numérico de azar que no tiene potencial (no hay valores que enumerar), o rechazarlo con una excepción del dominio que nombre el nodo. Qué conviene es cosa del algoritmo, pero no un \texttt{Index\allowbreak{}Out\allowbreak{}Of\allowbreak{}Bounds\allowbreak{}Exception}.

El mismo supuesto («todo nodo de azar tiene potencial») está en \hyperref[err:29]{error~29} y en el lector, que es donde podría imponerse de una vez.

Prueba: con la red A → X anterior, \texttt{convert\allowbreak{}Numerical\allowbreak{}Variables\allowbreak{}To\allowbreak{}FS} no debe lanzar esa excepción y \texttt{VEPropagation} debe dar el posterior de A o un error con mensaje.

\setcounter{subsection}{43}
\subsection[Un hallazgo numérico con más de dos decimales se pierde al discretizar]{El hallazgo sobre una variable numérica observada con un valor más fino que una centésima desaparece de la evidencia al discretizar}\label{err:44}
\begin{ficha}
\fichakey{Estado}Presente
\fichakey{Módulo}core
\fichakey{Paquete}org.openmarkov.core.model.network
\fichakey{Ficheros}\path{core/src/main/java/org/openmarkov/core/model/network/ProbNetOperations.java:444-459} \newline \path{core/src/main/java/org/openmarkov/core/model/network/ProbNetOperations.java:575-588} \newline \path{core/src/main/java/org/openmarkov/core/model/network/Variable.java:75}
\fichakey{Comprobado}Leyendo el código y ejecutándolo
\fichakey{Esfuerzo}Pequeño (horas)
\fichakey{Decisión de equipo}No
\fichakey{Relacionados}\hyperref[err:40]{error~40}, \hyperref[err:41]{error~41}, \hyperref[err:42]{error~42}, \hyperref[err:45]{error~45}
\end{ficha}
\paragraph{Qué pasa}
Cualquier tarea exacta con evidencia previa a la resolución sobre una variable numérica de azar puede perder ese hallazgo sin aviso. La evidencia llega por el diálogo de añadir hallazgo en modo edición (cuyo selector usa la precisión de la variable como paso), desde el fichero, o inducida por potenciales delta y de desplazamiento de ciclo.

Las consecuencias visibles medidas son menores:

\begin{itemize}
\tightlist
\item
  La utilidad esperada (\texttt{VEEvaluation}) y las probabilidades de los demás nodos no cambian (N=7,0 y N=7,123 dan 10,0 en la red descrita más abajo). El valor observado ya queda recogido en la variable de un solo estado cuya tabla vale 1, y los hijos toman la rama que corresponde leyendo el nombre del estado; el hallazgo perdido sobre una variable de un solo estado no añade información. Por eso, hacia los hijos, el modelo no se resuelve como si nadie hubiera observado nada; hacia los padres sí, con cualquier valor observado (\hyperref[err:45]{error~45}).
\item
  La propagación tal como la hace la ventana (\texttt{Evidence\allowbreak{}Manager.\allowbreak{}do\allowbreak{}Propagation}, con evidencia posterior a la resolución vacía pero no nula) tampoco cambia: al extender la evidencia posterior, la tabla de un solo estado induce de nuevo el hallazgo.
\item
  Si se llama a \texttt{VEPropagation} sin evidencia posterior (como hace \texttt{VETemporal\allowbreak{}Evolution}) y se pregunta por esa misma variable numérica, con 7,0 devuelve {[}1,0{]} y con 7,123 lanza \texttt{Null\allowbreak{}Pointer\allowbreak{}Exception} en \texttt{Prob\allowbreak{}Net\allowbreak{}Operations.\allowbreak{}get\allowbreak{}Node\allowbreak{}Ancestors}, llamado desde \texttt{VEPropagation.\allowbreak{}prune\allowbreak{}Network}. Ningún llamador de producción pregunta hoy por una variable numérica de azar por ese camino (la exportación de la evolución temporal las salta).
\item
  En \texttt{MIDTemporal\allowbreak{}Evolution} sobre \texttt{0\_\allowbreak{}miscellaneous/\allowbreak{}networks/\allowbreak{}mid/\allowbreak{}MID-Chancellor.\allowbreak{}pgmx} con la duración de ciclo cambiada a 0,1, la probabilidad de «Time in treatment {[}3{]}» sale 1,0000000002241949 en lugar de 1,0: al perderse el hallazgo, el valor se calcula eliminando variables en lugar de salir directamente de la evidencia.
\end{itemize}

Por qué pasa. En \texttt{Prob\allowbreak{}Net\allowbreak{}Operations.\allowbreak{}convert\allowbreak{}Numerical\allowbreak{}Variables\allowbreak{}To\allowbreak{}FS} (la discretización previa a todo algoritmo exacto), una variable numérica de azar \textbf{con hallazgo} se sustituye por una variable de un solo estado cuyo nombre es el valor observado escrito con \texttt{String.\allowbreak{}value\allowbreak{}Of(\allowbreak{}value)} (línea 450), y su potencial pasa a ser una tabla con un único 1. Al final del método (líneas 575-588) se recorren los hallazgos para volver a apuntarlos sobre la variable nueva: se redondea el valor con la precisión de la \textbf{variable nueva} y se busca el estado cuyo nombre sea \texttt{String.\allowbreak{}value\allowbreak{}Of} de ese redondeo. Si no se encuentra, el \texttt{if} se salta la línea que vuelve a añadir el hallazgo, y el hallazgo desaparece de la evidencia de la tarea sin aviso.

Las dos cosas no encajan por dos razones:

\begin{enumerate}
\def\labelenumi{\arabic{enumi}.}
\tightlist
\item
  La variable nueva se construye con el constructor de estados, que deja la precisión por defecto, 0,01 (línea 75 de \texttt{Variable}), en vez de heredar la de la variable numérica. Cualquier valor con más de dos decimales (7,123; 0,001) se redondea a otro número y su nombre no coincide.
\item
  Incluso con dos decimales, \texttt{Math.\allowbreak{}round(\allowbreak{}v\ /\allowbreak{}\ 0.\allowbreak{}01)\ *\ 0.\allowbreak{}01} no siempre devuelve exactamente \texttt{v} en coma flotante: para 0,35 da 0,35000000000000003. Se comprobó que 129 de los 1001 valores 0,00; 0,01; \ldots; 10,00 pierden así su hallazgo. Un caso realista es la variable de tiempo de un modelo de Markov con duración de ciclo 0,1: a partir del tercer ciclo el valor inducido es 0,30000000000000004 y su hallazgo se pierde.
\end{enumerate}

Medido con un diagrama de influencia de tres nodos, P (azar, dos estados al 50 \%) → N (numérica de azar, vale 3 si P=p0 y 7 si P=p1) → U (utilidad: 0 si N ≤ 5, 10 si N \textgreater{} 5), el mismo de \hyperref[err:40]{error~40}, con un hallazgo previo a la resolución sobre N:

\begin{Verbatim}[breaklines,breakanywhere,fontsize=\small]
valor 7.0   -> estado "7.0",   round 7.0                 -> hallazgo sobre N tras convertir: true
valor 0.35  -> estado "0.35",  round 0.35000000000000003 -> hallazgo sobre N tras convertir: false
valor 7.123 -> estado "7.123", round 7.12                -> hallazgo sobre N tras convertir: false
valor 0.001 -> estado "0.001", round 0.0                 -> hallazgo sobre N tras convertir: false
\end{Verbatim}
\paragraph{El código}
core/src/main/java/org/openmarkov/core/model/network/ProbNetOperations.java:447-451, 575-588

\begin{Shaded}
\begin{Highlighting}[]
\NormalTok{                    Finding finding }\OperatorTok{=}\NormalTok{ configuration}\OperatorTok{.}\FunctionTok{removeFinding}\OperatorTok{(}\NormalTok{oldVariable}\OperatorTok{);}
                    \ControlFlowTok{if} \OperatorTok{(}\NormalTok{finding }\OperatorTok{!=} \KeywordTok{null}\OperatorTok{)} \OperatorTok{\{}
                        \DataTypeTok{double}\NormalTok{ value }\OperatorTok{=}\NormalTok{ finding}\OperatorTok{.}\FunctionTok{numericalValue}\OperatorTok{;}
\NormalTok{                        Variable newVariable }\OperatorTok{=} \KeywordTok{new} \FunctionTok{Variable}\OperatorTok{(}\NormalTok{oldVariable}\OperatorTok{.}\FunctionTok{getName}\OperatorTok{(),} \BuiltInTok{String}\OperatorTok{.}\FunctionTok{valueOf}\OperatorTok{(}\NormalTok{value}\OperatorTok{));}
\NormalTok{                        node}\OperatorTok{.}\FunctionTok{setVariable}\OperatorTok{(}\NormalTok{newVariable}\OperatorTok{);}
\KeywordTok{...}
        \CommentTok{// Update evidence, replacing references to old variables with new ones}
        \BuiltInTok{List}\OperatorTok{\textless{}}\NormalTok{Finding}\OperatorTok{\textgreater{}}\NormalTok{ findings }\OperatorTok{=}\NormalTok{ evidence}\OperatorTok{.}\FunctionTok{getFindings}\OperatorTok{();}
        \ControlFlowTok{for} \OperatorTok{(}\NormalTok{Finding finding }\OperatorTok{:}\NormalTok{ findings}\OperatorTok{)} \OperatorTok{\{}
\NormalTok{            Variable originalVariable }\OperatorTok{=}\NormalTok{ finding}\OperatorTok{.}\FunctionTok{getVariable}\OperatorTok{();}
            \ControlFlowTok{if} \OperatorTok{(}\NormalTok{convertedVariables}\OperatorTok{.}\FunctionTok{containsKey}\OperatorTok{(}\NormalTok{originalVariable}\OperatorTok{))} \OperatorTok{\{}
\NormalTok{                evidence}\OperatorTok{.}\FunctionTok{removeFinding}\OperatorTok{(}\NormalTok{originalVariable}\OperatorTok{);}
\NormalTok{                Variable convertedVariable }\OperatorTok{=}\NormalTok{ convertedVariables}\OperatorTok{.}\FunctionTok{get}\OperatorTok{(}\NormalTok{originalVariable}\OperatorTok{);}
                \DataTypeTok{double}\NormalTok{ numericalValue }\OperatorTok{=}\NormalTok{ convertedVariable}\OperatorTok{.}\FunctionTok{round}\OperatorTok{(}\NormalTok{finding}\OperatorTok{.}\FunctionTok{getNumericalValue}\OperatorTok{());}
                \DataTypeTok{int}\NormalTok{ stateIndex }\OperatorTok{=}\NormalTok{ convertedVariable}\OperatorTok{.}\FunctionTok{getStateIndex}\OperatorTok{(}\BuiltInTok{String}\OperatorTok{.}\FunctionTok{valueOf}\OperatorTok{(}\NormalTok{numericalValue}\OperatorTok{));}
                \ControlFlowTok{if} \OperatorTok{(}\NormalTok{stateIndex }\OperatorTok{!=} \OperatorTok{{-}}\DecValTok{1}\OperatorTok{)} \OperatorTok{\{}
\NormalTok{                    evidence}\OperatorTok{.}\FunctionTok{addFinding}\OperatorTok{(}\KeywordTok{new} \FunctionTok{Finding}\OperatorTok{(}\NormalTok{convertedVariable}\OperatorTok{,}\NormalTok{ stateIndex}\OperatorTok{));}
                \OperatorTok{\}}
            \OperatorTok{\}}
        \OperatorTok{\}}
\end{Highlighting}
\end{Shaded}
\paragraph{Cómo arreglarlo}
En la rama de las variables observadas el estado es único y se sabe cuál es: el hallazgo se puede volver a apuntar con el índice 0 directamente, sin redondear ni buscar por nombre (o, si se prefiere buscar, nombrar y buscar con la misma función). Si se decide redondear el valor observado, debe hacerse una sola vez, con la precisión de la variable original y antes de nombrar el estado. Conviene también que la variable nueva herede la precisión de la original.

Otros sitios de la misma regla, «el nombre de un estado discretizado y su búsqueda usan la misma escritura del número»:

\begin{itemize}
\tightlist
\item
  la rama de variables sin hallazgo del mismo método nombra con \texttt{Decimal\allowbreak{}Format(\allowbreak{}"\#.\allowbreak{}\#\#\#\#\#")} (línea 531), y \texttt{Tree\allowbreak{}ADDPotential.\allowbreak{}replace\allowbreak{}Numeric\allowbreak{}Variable} vuelve a leer esos nombres con \texttt{Double.\allowbreak{}parse\allowbreak{}Double};
\item
  \texttt{Variable.\allowbreak{}get\allowbreak{}State\allowbreak{}Index(\allowbreak{}double)} (línea 343) busca por \texttt{String.\allowbreak{}value\allowbreak{}Of(\allowbreak{}round(\allowbreak{}value))}, que tiene el mismo problema para variables de estados finitos con nombres numéricos.
\end{itemize}

El \texttt{Null\allowbreak{}Pointer\allowbreak{}Exception} de \texttt{VEPropagation.\allowbreak{}prune\allowbreak{}Network} al preguntar por una variable numérica que la discretización ha sustituido es un defecto distinto, \hyperref[err:42]{error~42}; arreglar este error no lo arregla.

Prueba de regresión: convertir la red P → N con un hallazgo N=7,123 y otro N=0,35 y comprobar que la evidencia sigue conteniendo un hallazgo sobre la N convertida.

\setcounter{subsection}{44}
\subsection[Observar un nodo numérico de azar no informa sobre sus padres]{Observar un nodo numérico de azar no cambia la probabilidad de sus padres, y la utilidad esperada se calcula como si esa observación no dijera nada de ellos}\label{err:45}
\begin{ficha}
\fichakey{Estado}Presente
\fichakey{Módulo}core
\fichakey{Paquete}org.openmarkov.core.model.network
\fichakey{Ficheros}\path{core/src/main/java/org/openmarkov/core/model/network/ProbNetOperations.java:420-459} \newline \path{core/src/main/java/org/openmarkov/core/inference/tasks/TaskUtilities.java:88-90} \newline \path{inference/src/main/java/org/openmarkov/inference/algorithm/variableElimination/tasks/VariableElimination.java:95-97} \newline \path{gui/src/main/java/org/openmarkov/gui/window/MainPanelMenuAssistant.java:688-692} \newline \path{core/src/main/java/org/openmarkov/core/model/network/modelUncertainty/ProbDensFunction.java:24-64}
\fichakey{Comprobado}Leyendo el código y ejecutándolo
\fichakey{Esfuerzo}Medio (días)
\fichakey{Decisión de equipo}\textcolor{decisionrojo}{\textbf{Sí.}} Qué hacer con un hallazgo sobre un nodo cuya relación es una distribución de probabilidad (normal, gamma\ldots): usar la densidad en el valor observado, que hoy no se puede calcular, o rechazar el hallazgo con un mensaje
\fichakey{Relacionados}\hyperref[err:40]{error~40}, \hyperref[err:41]{error~41}, \hyperref[err:42]{error~42}, \hyperref[err:44]{error~44}, \hyperref[err:48]{error~48}, \hyperref[nota:109]{nota~109}
\end{ficha}
\paragraph{Qué pasa}
En una red con un nodo numérico de azar que depende de un nodo de estados finitos (por ejemplo, una prueba cuyo resultado numérico depende de si hay enfermedad), el usuario pone en modo edición un hallazgo sobre el nodo numérico: menú contextual del nodo, «Add finding» (la opción está activa para todo nodo que no sea de utilidad, salvo los temporales con un padre del mismo nombre). Al entrar en modo inferencia, la probabilidad del padre sigue siendo la previa, como si no se hubiera observado nada. Como la causa está en la discretización que hacen todas las tareas exactas, lo mismo vale para las que reciben evidencia previa a la resolución: evaluación del diagrama de influencia, coste-efectividad, sensibilidad y evolución temporal (medido en la propagación y en la evaluación). No hay ningún aviso.

Hacia los hijos el hallazgo sí llega: el nodo queda fijado en el valor observado, y los nodos que dependen de él se calculan con ese valor (es lo que se mide en \hyperref[err:44]{error~44}). Lo que se pierde es lo que el valor observado dice sobre los padres.

Medido con redes de tres nodos construidas por programa: Disease (absent 0,9 / present 0,1) → N (numérico de azar, precisión 0,01) y, en el diagrama de influencia, además Disease → U (utilidad 0 si absent, 10 si present). Con un hallazgo previo a la resolución sobre N, propagando como la ventana (\texttt{VEPropagation} con todas las variables, igual que \texttt{Evidence\allowbreak{}Manager.\allowbreak{}do\allowbreak{}Propagation}) y evaluando con \texttt{VEEvaluation}:

\begin{Verbatim}[breaklines,breakanywhere,fontsize=\small]
N de estados finitos «3»/«7», tabla determinista, N = «7»         -> P(Disease) = [0.0, 1.0]  utilidad esperada = 10.0
N numérico, relación exacta (3 si absent, 7 si present), N = 7    -> P(Disease) = [0.9, 0.1]  utilidad esperada = 1.0
N numérico, normal (media 5 o 8, desviación 1), N = 8             -> P(Disease) = [0.9, 0.1]  utilidad esperada = 1.0
\end{Verbatim}

La primera fila es la misma red con N de estados finitos, que da la respuesta correcta. En la tercera, la respuesta correcta es P(present) ≈ 0,909 (el cociente de las dos densidades normales en 8 multiplica por 90 la razón de probabilidades previa). Sin hallazgo, la red con la normal ni siquiera se puede propagar («Potential P(N \textbar{} Disease) cannot be converted to a table»); con hallazgo se propaga sin quejarse y sin usar la distribución.

Por qué pasa. Antes de resolver con eliminación de variables, \texttt{Variable\allowbreak{}Elimination.\allowbreak{}exact\allowbreak{}Algorithms\allowbreak{}Preprocessing} llama a \texttt{Task\allowbreak{}Utilities.\allowbreak{}discretize\allowbreak{}Non\allowbreak{}Observed\allowbreak{}Numeric\allowbreak{}Variables}, que es \texttt{Prob\allowbreak{}Net\allowbreak{}Operations.\allowbreak{}convert\allowbreak{}Numerical\allowbreak{}Variables\allowbreak{}To\allowbreak{}FS}. Para cada nodo numérico de azar con hallazgo previo, ese método (líneas 444-459) pone en el nodo una variable nueva de un solo estado, cuyo nombre es el valor observado, y cambia su potencial por una tabla sobre esa única variable con el valor 1. La tabla nueva no tiene a los padres: desaparece el factor que decía cuánto de probable es el valor observado en cada configuración de los padres, y los padres se quedan con su distribución previa.

El javadoc del método solo habla de convertir las variables sin hallazgo. Tratar las observadas así es correcto cuando su valor lo fija un potencial delta y sus padres también están fijados, que es el caso de los modelos de Markov del repositorio («Time in treatment», «Age»): ahí perder el factor no cambia nada.

Camino de llegada. Ninguna red del repositorio lo alcanza tal cual: los nodos numéricos de azar de \texttt{0\_\allowbreak{}miscellaneous/\allowbreak{}networks} y de \texttt{integration\allowbreak{}Tests/\allowbreak{}src/\allowbreak{}main/\allowbreak{}resources/\allowbreak{}networks} no tienen padres de estados finitos, o son temporales con un padre del mismo nombre (y entonces «Add finding» está desactivado), o su red falla antes por otros motivos. Se llega construyendo la red en el editor: un nodo numérico de azar admite las relaciones «Exact» y «UnivariateDistr» con padres de estados finitos.
\paragraph{El código}
core/src/main/java/org/openmarkov/core/model/network/ProbNetOperations.java:441-459

\begin{Shaded}
\begin{Highlighting}[]
            \ControlFlowTok{if} \OperatorTok{(}\NormalTok{oldVariable}\OperatorTok{.}\FunctionTok{getVariableType}\OperatorTok{()} \OperatorTok{==}\NormalTok{ VariableType}\OperatorTok{.}\FunctionTok{NUMERIC} \OperatorTok{\&\&}\NormalTok{ node}\OperatorTok{.}\FunctionTok{getNodeType}\OperatorTok{()} \OperatorTok{==}\NormalTok{ NodeType}\OperatorTok{.}\FunctionTok{CHANCE}\OperatorTok{)} \OperatorTok{\{}
\NormalTok{                EvidenceCase configuration }\OperatorTok{=} \KeywordTok{new} \FunctionTok{EvidenceCase}\OperatorTok{(}\NormalTok{evidence}\OperatorTok{);}
\NormalTok{                Potential oldPotential }\OperatorTok{=}\NormalTok{ node}\OperatorTok{.}\FunctionTok{getPotentials}\OperatorTok{().}\FunctionTok{get}\OperatorTok{(}\DecValTok{0}\OperatorTok{);}
                \ControlFlowTok{if} \OperatorTok{(}\NormalTok{configuration}\OperatorTok{.}\FunctionTok{contains}\OperatorTok{(}\NormalTok{oldVariable}\OperatorTok{))} \OperatorTok{\{}
                    \CommentTok{// Convert numerical variables with evidence to one{-}state}
                    \CommentTok{// variables}
\NormalTok{                    Finding finding }\OperatorTok{=}\NormalTok{ configuration}\OperatorTok{.}\FunctionTok{removeFinding}\OperatorTok{(}\NormalTok{oldVariable}\OperatorTok{);}
                    \ControlFlowTok{if} \OperatorTok{(}\NormalTok{finding }\OperatorTok{!=} \KeywordTok{null}\OperatorTok{)} \OperatorTok{\{}
                        \DataTypeTok{double}\NormalTok{ value }\OperatorTok{=}\NormalTok{ finding}\OperatorTok{.}\FunctionTok{numericalValue}\OperatorTok{;}
\NormalTok{                        Variable newVariable }\OperatorTok{=} \KeywordTok{new} \FunctionTok{Variable}\OperatorTok{(}\NormalTok{oldVariable}\OperatorTok{.}\FunctionTok{getName}\OperatorTok{(),} \BuiltInTok{String}\OperatorTok{.}\FunctionTok{valueOf}\OperatorTok{(}\NormalTok{value}\OperatorTok{));}
\NormalTok{                        node}\OperatorTok{.}\FunctionTok{setVariable}\OperatorTok{(}\NormalTok{newVariable}\OperatorTok{);}
                        \KeywordTok{...}
\NormalTok{                        TablePotential potential }\OperatorTok{=} \KeywordTok{new} \FunctionTok{TablePotential}\OperatorTok{(}\BuiltInTok{Arrays}\OperatorTok{.}\FunctionTok{asList}\OperatorTok{(}\NormalTok{newVariable}\OperatorTok{),}
\NormalTok{                                                                      oldPotential}\OperatorTok{.}\FunctionTok{getPotentialRole}\OperatorTok{());}
\NormalTok{                        potential}\OperatorTok{.}\FunctionTok{getValues}\OperatorTok{()[}\DecValTok{0}\OperatorTok{]} \OperatorTok{=} \DecValTok{1}\OperatorTok{;}
\NormalTok{                        node}\OperatorTok{.}\FunctionTok{setPotential}\OperatorTok{(}\NormalTok{potential}\OperatorTok{);}
                    \OperatorTok{\}}
\end{Highlighting}
\end{Shaded}
\paragraph{Cómo arreglarlo}
El potencial del nodo observado tiene que conservar a sus padres: una tabla sobre la variable de un solo estado y los padres que diga, para cada configuración de los padres, cuánto de probable es haber observado ese valor.

\begin{itemize}
\tightlist
\item
  Si la relación es determinista (tabla, «Exact», función, suma o producto), basta con el recorrido de configuraciones que ya hace la rama de las variables sin hallazgo (líneas 461-526): proyectar el potencial en cada configuración, redondear con la precisión de la variable y poner 1 donde coincide con el valor observado y 0 donde no. Si ninguna configuración da ese valor, la evidencia es imposible y debe tratarse como tal (\hyperref[err:48]{error~48}).
\item
  Si la relación es una distribución (\texttt{Univariate\allowbreak{}Distr\allowbreak{}Potential}), lo que corresponde es la densidad en el valor observado para cada configuración. \texttt{Prob\allowbreak{}Dens\allowbreak{}Function} no tiene hoy ningún método que la calcule (tiene media, varianza, muestreo, intervalos\ldots). Hay que decidir si se añade ese cálculo a las distribuciones o si se rechaza el hallazgo con un mensaje que diga que no se puede usar.
\end{itemize}

Sitios de la misma regla: la rama de hallazgos de \texttt{Prob\allowbreak{}Net\allowbreak{}Operations.\allowbreak{}convert\allowbreak{}Numerical\allowbreak{}Variables\allowbreak{}To\allowbreak{}FS} es el único sitio, pero la usan todas las tareas exactas (\texttt{Variable\allowbreak{}Elimination.\allowbreak{}exact\allowbreak{}Algorithms\allowbreak{}Preprocessing}; \texttt{Temporal\allowbreak{}Evaluation}, línea 105; \texttt{MIDTemporal\allowbreak{}Evolution}, línea 339; \texttt{Temporal\allowbreak{}Net\allowbreak{}Operations}, línea 467), así que el arreglo llega a todas. Conviene hacerlo a la vez que \hyperref[err:44]{error~44}, que toca el mismo bloque (el nombre del estado nuevo y la búsqueda del hallazgo), y que \hyperref[err:40]{error~40} (precisión cero). Hugin y el muestreo no pasan por aquí; cómo tratan ellos los nodos numéricos está en \hyperref[nota:109]{nota~109}, donde no hay nada que corregir.

Prueba de regresión: la red Disease → N con relación exacta y hallazgo N = 7 debe dar P(Disease) = {[}0, 1{]} con \texttt{VEPropagation} y utilidad esperada 10 con \texttt{VEEvaluation}, igual que la misma red con N de estados finitos; y la utilidad esperada de \texttt{0\_\allowbreak{}miscellaneous/\allowbreak{}networks/\allowbreak{}mid/\allowbreak{}MID-Chancellor.\allowbreak{}pgmx} no debe cambiar.

\setcounter{subsection}{45}
\subsection[Un nodo de utilidad sin rango impide entrar en modo inferencia]{Un nodo de utilidad cuyo rango no se puede calcular impide entrar en modo inferencia aunque la propagación funcionaría}\label{err:46}
\begin{ficha}
\fichakey{Estado}Presente
\fichakey{Módulo}gui, core
\fichakey{Paquete}org.openmarkov.gui.window.edition.networkEditorPanel; org.openmarkov.core.model.network
\fichakey{Ficheros}\path{gui/src/main/java/org/openmarkov/gui/window/edition/networkEditorPanel/EvidenceManager.java:558-598} \newline \path{core/src/main/java/org/openmarkov/core/model/network/UtilityFunctionComputer.java:82-119} \newline \path{gui/src/main/java/org/openmarkov/gui/window/edition/networkEditorPanel/InferencePresenter.java:60-69} \newline \path{gui/src/main/java/org/openmarkov/gui/graphic/NumericVariableBox.java:101-123} \newline \path{gui/src/main/java/org/openmarkov/gui/window/InferenceHandler.java:70-125} \newline \path{gui/src/main/java/org/openmarkov/gui/window/MainPanelListenerAssistant.java:145-163}
\fichakey{Comprobado}Leyendo el código y ejecutándolo
\fichakey{Esfuerzo}Pequeño (horas)
\fichakey{Decisión de equipo}No
\fichakey{Relacionados}\hyperref[err:9]{error~9}, \hyperref[err:41]{error~41}, \hyperref[err:42]{error~42}, \hyperref[err:50]{error~50}, \hyperref[err:71]{error~71}, \hyperref[err:112]{error~112}, \hyperref[err:113]{error~113}
\end{ficha}
\paragraph{Qué pasa}
Al entrar en modo inferencia en una red con un nodo de utilidad cuyo rango no se puede calcular, la red vuelve a modo edición y sale la ventana de error. El usuario no puede usar el modo inferencia en esa red, aunque la propagación en sí funcionaría y aunque los demás nodos de utilidad sí tengan rango.

El camino es \texttt{Inference\allowbreak{}Handler.\allowbreak{}set\allowbreak{}Working\allowbreak{}Mode} → \texttt{update\allowbreak{}Individual\allowbreak{}Probabilities\allowbreak{}And\allowbreak{}Utilities} → \texttt{Evidence\allowbreak{}Manager.\allowbreak{}do\allowbreak{}Propagation}; la excepción llega a \texttt{Main\allowbreak{}Panel\allowbreak{}Listener\allowbreak{}Assistant} (líneas 155-171), que devuelve la red a modo edición y muestra la ventana de error.

Por qué pasa. En modo inferencia, los nodos de utilidad se dibujan con una barra cuyo largo se escala entre un mínimo y un máximo aproximados. \texttt{Evidence\allowbreak{}Manager.\allowbreak{}do\allowbreak{}Propagation} (módulo \texttt{gui}, líneas 581-600) calcula esos límites al principio, con \texttt{calculate\allowbreak{}Min\allowbreak{}And\allowbreak{}Max\allowbreak{}Utility\allowbreak{}Ranges} (líneas 559-572), que para cada nodo de utilidad llama a \texttt{Utility\allowbreak{}Function\allowbreak{}Computer.\allowbreak{}approximate\allowbreak{}Min\allowbreak{}Utility} y \texttt{approximate\allowbreak{}Max\allowbreak{}Utility} (módulo \texttt{core}).

Estas funciones convierten el potencial del nodo en tabla (\texttt{table\allowbreak{}Project}), y eso falla para varias clases de potencial: una distribución sobre un número (\texttt{Univariate\allowbreak{}Distr\allowbreak{}Potential}), un \texttt{Uniform\allowbreak{}Potential} sobre una variable numérica, un árbol (\texttt{Tree\allowbreak{}ADDPotential}) o una suma o producto (\texttt{Sum\allowbreak{}Potential}, \texttt{Product\allowbreak{}Potential}) cuyos sumandos fallan, o un potencial que depende de una variable numérica sin hallazgo (\texttt{Missing\allowbreak{}Evidence\allowbreak{}In\allowbreak{}Variable}). La excepción (\texttt{Non\allowbreak{}Projectable\allowbreak{}Potential\allowbreak{}Exception}) no se recoge en \texttt{do\allowbreak{}Propagation} (su \texttt{try} solo recoge \texttt{Out\allowbreak{}Of\allowbreak{}Memory\allowbreak{}Error}), así que no se propaga nada.

Medido con un pequeño programa de prueba sobre las redes de \texttt{0\_\allowbreak{}miscellaneous/\allowbreak{}networks} (148 se leen, 2 no): 779 nodos de utilidad, 414 sin rango, en 73 redes. Para esas 73 se lanzó además \texttt{VEPropagation} sin evidencia y sin calcular rangos:

\begin{itemize}
\tightlist
\item
  En 4 la propagación termina bien y el cálculo de rangos es lo único que bloquea el modo inferencia, con un nodo sin rango en cada una: \texttt{0\_\allowbreak{}miscellaneous/\allowbreak{}networks/\allowbreak{}mid/\allowbreak{}DMHEE/\allowbreak{}MID-dmhee-2.\allowbreak{}5.\allowbreak{}pgmx} («Cost {[}0{]}»), \texttt{0\_\allowbreak{}miscellaneous/\allowbreak{}networks/\allowbreak{}mid/\allowbreak{}DMHEE/\allowbreak{}MID-dmhee-4.\allowbreak{}7.\allowbreak{}pgmx} («Drug cost {[}0{]}», error «There is required evidence to be present in variable Numeric variable Time in treatment {[}0{]}\ldots»), \texttt{0\_\allowbreak{}miscellaneous/\allowbreak{}networks/\allowbreak{}mid/\allowbreak{}MID-Chancellor.\allowbreak{}pgmx} y \texttt{0\_\allowbreak{}miscellaneous/\allowbreak{}networks/\allowbreak{}mid/\allowbreak{}MID-Cochlear.\allowbreak{}pgmx}.
\item
  En las otras 69 la propagación falla también por sí sola (sobre todo \texttt{Potential\allowbreak{}Cannot\allowbreak{}Be\allowbreak{}Converted\allowbreak{}To\allowbreak{}ATable}, y en algunas \texttt{Unreachable\allowbreak{}Exception}, \texttt{Cannot\allowbreak{}Resolve\allowbreak{}Variable}, \texttt{Not\allowbreak{}Applicable\allowbreak{}Network} o \texttt{Null\allowbreak{}Pointer\allowbreak{}Exception}, que es \hyperref[err:42]{error~42}); 60 de ellas son las copias en formato 1.0 de las carpetas \texttt{1-0}. Arreglar este error no hace funcionar esas 69; son otros defectos (en las copias en formato 1.0, sobre todo \hyperref[err:112]{error~112}).
\end{itemize}

Es decir, casi la mitad de las redes tiene algún nodo de utilidad sin rango, pero las que quedan bloqueadas solo por esta causa son 4.
\paragraph{El código}
gui/src/main/java/org/openmarkov/gui/window/edition/networkEditorPanel/EvidenceManager.java:559-570,581-590

\begin{Shaded}
\begin{Highlighting}[]
    \KeywordTok{private} \DataTypeTok{void} \FunctionTok{calculateMinAndMaxUtilityRanges}\OperatorTok{()} \KeywordTok{throws}\NormalTok{ NonProjectablePotentialException }\OperatorTok{\{}
        \BuiltInTok{List}\OperatorTok{\textless{}}\NormalTok{Variable}\OperatorTok{\textgreater{}}\NormalTok{ utilityVariables }\OperatorTok{=} \KeywordTok{this}\OperatorTok{.}\FunctionTok{networkEditorPanel}\OperatorTok{.}\FunctionTok{getVisualNetwork}\OperatorTok{().}\FunctionTok{getProbNet}\OperatorTok{().}\FunctionTok{getVariables}\OperatorTok{(}\NormalTok{NodeType}\OperatorTok{.}\FunctionTok{UTILITY}\OperatorTok{);}
        \ControlFlowTok{for} \OperatorTok{(}\NormalTok{Variable utility }\OperatorTok{:}\NormalTok{ utilityVariables}\OperatorTok{)} \OperatorTok{\{}
\NormalTok{            ProbNet newNet }\OperatorTok{=} \KeywordTok{this}\OperatorTok{.}\FunctionTok{networkEditorPanel}\OperatorTok{.}\FunctionTok{getVisualNetwork}\OperatorTok{().}\FunctionTok{getProbNet}\OperatorTok{().}\FunctionTok{copy}\OperatorTok{();}
            \KeywordTok{...}
            \KeywordTok{this}\OperatorTok{.}\FunctionTok{minUtilityRange}\OperatorTok{.}\FunctionTok{put}\OperatorTok{(}\NormalTok{utility}\OperatorTok{,}\NormalTok{ UtilityFunctionComputer}\OperatorTok{.}\FunctionTok{approximateMinUtility}\OperatorTok{(}\NormalTok{node}\OperatorTok{));}
            \KeywordTok{this}\OperatorTok{.}\FunctionTok{maxUtilityRange}\OperatorTok{.}\FunctionTok{put}\OperatorTok{(}\NormalTok{utility}\OperatorTok{,}\NormalTok{ UtilityFunctionComputer}\OperatorTok{.}\FunctionTok{approximateMaxUtility}\OperatorTok{(}\NormalTok{node}\OperatorTok{));}
        \OperatorTok{\}}
    \OperatorTok{\}}
    \KeywordTok{...}
        \ControlFlowTok{try} \OperatorTok{\{}
            \KeywordTok{this}\OperatorTok{.}\FunctionTok{calculateMinAndMaxUtilityRanges}\OperatorTok{();}
\NormalTok{            Propagation vePosteriorValues }\OperatorTok{=} \KeywordTok{new} \FunctionTok{VEPropagation}\OperatorTok{(}\KeywordTok{this}\OperatorTok{.}\FunctionTok{networkEditorPanel}\OperatorTok{.}\FunctionTok{getVisualNetwork}\OperatorTok{().}\FunctionTok{getProbNet}\OperatorTok{());}
\end{Highlighting}
\end{Shaded}
\paragraph{Cómo arreglarlo}
Calcular el rango nodo a nodo recogiendo el fallo de cada uno: el nodo sin rango se queda sin límites (o con los del dominio de su variable, que \texttt{Numeric\allowbreak{}Variable\allowbreak{}Box} ya usa al construirse, líneas 83-84) y los demás siguen.

Hay que tocar a la vez \texttt{Inference\allowbreak{}Presenter.\allowbreak{}paint\allowbreak{}Inference\allowbreak{}Results\allowbreak{}Utility\allowbreak{}Node} (líneas 67-68): hoy pasa el \texttt{Double} del mapa a \texttt{set\allowbreak{}Min\allowbreak{}Value(\allowbreak{}double)}, y un rango ausente daría \texttt{Null\allowbreak{}Pointer\allowbreak{}Exception} al desempaquetar.

De paso, copiar la red una sola vez fuera del bucle (hoy se copia una vez por nodo de utilidad).

Prueba: entrar en modo inferencia en \texttt{0\_\allowbreak{}miscellaneous/\allowbreak{}networks/\allowbreak{}mid/\allowbreak{}DMHEE/\allowbreak{}MID-dmhee-4.\allowbreak{}7.\allowbreak{}pgmx} (o llamar a \texttt{do\allowbreak{}Propagation} sobre su panel) debe terminar con las probabilidades calculadas y el nodo «Drug cost {[}0{]}» sin barra escalada.

\setcounter{subsection}{46}
\subsection[«Add finding» marca siempre el primer estado y pisa el hallazgo anterior]{El diálogo de añadir hallazgo marca siempre el primer estado, de modo que aceptar sin tocar cambia el hallazgo que ya había, y deshacer ese cambio no devuelve el anterior}\label{err:47}
\begin{ficha}
\fichakey{Estado}Presente
\fichakey{Módulo}gui
\fichakey{Paquete}org.openmarkov.gui.dialog.node
\fichakey{Ficheros}\path{gui/src/main/java/org/openmarkov/gui/dialog/node/AddFindingDialog.java:119-141, 172-186} \newline \path{gui/src/main/java/org/openmarkov/gui/window/edition/networkEditorPanel/EvidenceManager.java:123-136, 489-509} \newline \path{gui/src/main/java/org/openmarkov/gui/action/AddFindingEdit.java:29-58}
\fichakey{Comprobado}Leyendo el código
\fichakey{Esfuerzo}Pequeño (horas)
\fichakey{Decisión de equipo}No
\end{ficha}
\paragraph{Qué pasa}
Para una variable de estados finitos, el diálogo de añadir hallazgo se abre siempre con el estado 0 marcado (el que aparece último en la lista), sea cual sea el hallazgo que ya tenga el nodo. Consecuencias, deducidas leyendo el código (no se ha abierto el diálogo):

\begin{enumerate}
\def\labelenumi{\arabic{enumi}.}
\tightlist
\item
  Si el nodo ya tiene un hallazgo en otro estado y el usuario abre el diálogo (menú contextual del nodo, «Add finding») y acepta sin tocar nada, \texttt{do\allowbreak{}Ok\allowbreak{}Click\allowbreak{}Before\allowbreak{}Hide} crea el hallazgo con el estado marcado, el 0, y \texttt{Evidence\allowbreak{}Manager.\allowbreak{}set\allowbreak{}New\allowbreak{}Finding} sustituye el que había.
\item
  En modo edición, si el nodo ya tenía hallazgo previo a la resolución, \texttt{do\allowbreak{}Ok\allowbreak{}Click\allowbreak{}Before\allowbreak{}Hide} pasa \texttt{previous\allowbreak{}Finding} a \texttt{set\allowbreak{}New\allowbreak{}Finding}, que crea \texttt{Add\allowbreak{}Finding\allowbreak{}Edit(\allowbreak{}visual\allowbreak{}Node,\allowbreak{}\ evidence\allowbreak{}Case,\allowbreak{}\ previous\allowbreak{}Finding,\allowbreak{}\ finding)}. Al deshacer, \texttt{Add\allowbreak{}Finding\allowbreak{}Edit.\allowbreak{}undo} pone \texttt{previous\allowbreak{}Finding}, es decir, el estado 0, en lugar del hallazgo real anterior.
\end{enumerate}

Las variables numéricas no están afectadas: el diálogo inicia el selector con el valor del hallazgo existente.

Por qué pasa. \texttt{Evidence\allowbreak{}Manager.\allowbreak{}add\allowbreak{}Finding} (líneas 130-137) abre \texttt{Add\allowbreak{}Finding\allowbreak{}Dialog} para el nodo seleccionado pasándole el hallazgo que ya tenga en el caso de evidencia actual (en modo inferencia) o en la evidencia previa a la resolución (en modo edición). Para una variable de estados finitos, \texttt{get\allowbreak{}Principal\allowbreak{}Panel} crea un botón de opción por estado, recorriendo los estados del último al primero, todos en el mismo grupo:

\begin{itemize}
\tightlist
\item
  si hay hallazgo, marca el botón cuyo nombre coincide con el estado del hallazgo, y además, en todas las vueltas, pone \texttt{previous\allowbreak{}Finding} al estado de esa vuelta;
\item
  en la vuelta \texttt{i\ ==\ 0} marca el botón del estado 0 y pone \texttt{previous\allowbreak{}Finding} al estado 0, sin mirar si había hallazgo.
\end{itemize}

Como marcar un botón del grupo desmarca los demás, la última vuelta pisa la marca del hallazgo existente, y \texttt{previous\allowbreak{}Finding} vale siempre el estado 0.
\paragraph{El código}
gui/src/main/java/org/openmarkov/gui/dialog/node/AddFindingDialog.java:122-138

\begin{Shaded}
\begin{Highlighting}[]
            \ControlFlowTok{for} \OperatorTok{(}\DataTypeTok{int}\NormalTok{ i }\OperatorTok{=}\NormalTok{ states}\OperatorTok{.}\FunctionTok{length} \OperatorTok{{-}} \DecValTok{1}\OperatorTok{;}\NormalTok{ i }\OperatorTok{\textgreater{}=} \DecValTok{0}\OperatorTok{;}\NormalTok{ i}\OperatorTok{{-}{-})} \OperatorTok{\{}
                \BuiltInTok{String}\NormalTok{ stateName }\OperatorTok{=}\NormalTok{ states}\OperatorTok{[}\NormalTok{i}\OperatorTok{].}\FunctionTok{getName}\OperatorTok{();}
                \BuiltInTok{JRadioButton}\NormalTok{ jRadioButton }\OperatorTok{=} \KeywordTok{new} \BuiltInTok{JRadioButton}\OperatorTok{(}\NormalTok{stateName}\OperatorTok{);}
                \ControlFlowTok{if} \OperatorTok{(}\NormalTok{finding }\OperatorTok{!=} \KeywordTok{null}\OperatorTok{)} \OperatorTok{\{}
\NormalTok{                    jRadioButton}\OperatorTok{.}\FunctionTok{setSelected}\OperatorTok{(}\NormalTok{finding}\OperatorTok{.}\FunctionTok{getState}\OperatorTok{().}\FunctionTok{equals}\OperatorTok{(}\NormalTok{stateName}\OperatorTok{));}
\NormalTok{                    previousFinding }\OperatorTok{=} \KeywordTok{new} \FunctionTok{Finding}\OperatorTok{(}\KeywordTok{this}\OperatorTok{.}\FunctionTok{variable}\OperatorTok{,} \KeywordTok{new} \BuiltInTok{State}\OperatorTok{(}\NormalTok{stateName}\OperatorTok{));}
                \OperatorTok{\}}
\NormalTok{                radioButtonsPanel}\OperatorTok{.}\FunctionTok{add}\OperatorTok{(}\NormalTok{jRadioButton}\OperatorTok{);}
\NormalTok{                jRadioButton}\OperatorTok{.}\FunctionTok{setActionCommand}\OperatorTok{(}\NormalTok{stateName}\OperatorTok{);}
                \ControlFlowTok{if} \OperatorTok{(}\NormalTok{i }\OperatorTok{==} \DecValTok{0}\OperatorTok{)} \OperatorTok{\{}
\NormalTok{                    jRadioButton}\OperatorTok{.}\FunctionTok{setSelected}\OperatorTok{(}\KeywordTok{true}\OperatorTok{);}
\NormalTok{                    previousFinding }\OperatorTok{=} \KeywordTok{new} \FunctionTok{Finding}\OperatorTok{(}\KeywordTok{this}\OperatorTok{.}\FunctionTok{variable}\OperatorTok{,} \KeywordTok{new} \BuiltInTok{State}\OperatorTok{(}\NormalTok{stateName}\OperatorTok{));}
                \OperatorTok{\}}
\end{Highlighting}
\end{Shaded}
\paragraph{Cómo arreglarlo}
Marcar el estado del hallazgo existente si lo hay, y el estado 0 solo cuando no hay hallazgo (\texttt{if\ (\allowbreak{}i\ ==\ 0\ \&\&\ finding\ ==\ null)}); y fijar \texttt{previous\allowbreak{}Finding\ =\ finding} una sola vez fuera del bucle (\texttt{null} si no hay hallazgo).

Prueba: construir el diálogo sin mostrarlo, o extraer la elección del estado inicial a un método, para una variable con tres estados y un hallazgo en el estado 2; debe quedar marcado el 2 y \texttt{previous\allowbreak{}Finding} debe ser ese hallazgo. En modo edición, cambiar un hallazgo previo de 2 a 1 y deshacer debe devolver el 2.

\setcounter{subsection}{47}
\subsection[La evidencia imposible se trata distinto según el gesto]{La evidencia imposible recibe un tratamiento distinto según el gesto, y en el más común se presenta como un error del programa}\label{err:48}
\begin{ficha}
\fichakey{Estado}Presente
\fichakey{Módulo}gui
\fichakey{Paquete}org.openmarkov.gui.window.edition.networkEditorPanel
\fichakey{Ficheros}\path{gui/src/main/java/org/openmarkov/gui/window/edition/networkEditorPanel/EditorInputHandler.java:260-266} \newline \path{gui/src/main/java/org/openmarkov/gui/window/edition/networkEditorPanel/EvidenceManager.java:183-190, 527-541, 394-401, 633-825} \newline \path{gui/src/main/java/org/openmarkov/gui/window/MainPanelListenerAssistant.java:145-163} \newline \path{gui/src/main/java/org/openmarkov/gui/dialog/PropagationOptionsDialogListener.java:67-72} \newline \path{stochasticPropagationOutput/src/main/java/org/openmarkov/stochasticPropagationOutput/StochasticPropagationOutputFrame.java:194-201} \newline \path{gui/src/main/java/org/openmarkov/gui/dialog/OMExceptionHandler.java:17-59} \newline \path{core/src/main/java/org/openmarkov/core/model/network/potential/operation/TablePotentialTransform.java:45-48} \newline \path{core/src/main/resources/core/localize/core_class_localizations_en.xml:96-97} \newline \path{io/src/main/java/org/openmarkov/io/probmodel/reader/PGMXReader_0_2.java:534-536}
\fichakey{Comprobado}Leyendo el código
\fichakey{Esfuerzo}Medio (días)
\fichakey{Decisión de equipo}\textcolor{decisionrojo}{\textbf{Sí.}} Qué respuesta única dar (una excepción del dominio con mensaje propio, dónde se lanza y dónde se captura) y si al fallar la propagación debe apagarse la propagación automática o no
\fichakey{Relacionados}\hyperref[err:45]{error~45}, \hyperref[err:54]{error~54}, \hyperref[err:67]{error~67}, \hyperref[err:85]{error~85}, \hyperref[err:86]{error~86}, \hyperref[err:108]{error~108}, \hyperref[err:147]{error~147}, \hyperref[err:186]{error~186}, \hyperref[nota:91]{nota~91}, \hyperref[nota:105]{nota~105}, \hyperref[nota:167]{nota~167}
\end{ficha}
\paragraph{Qué pasa}
La evidencia imposible es una situación normal de uso: el usuario introduce observaciones que el modelo considera imposibles. Hay dos formas:

\begin{itemize}
\tightlist
\item
  dos hallazgos contradictorios sobre la misma variable, que detecta \texttt{Evidence\allowbreak{}Case.\allowbreak{}add\allowbreak{}Finding} (\texttt{Incompatible\allowbreak{}Evidence\allowbreak{}Exception.\allowbreak{}Evidence\allowbreak{}Is\allowbreak{}Incompatible\allowbreak{}With\allowbreak{}Other}). Uno de los dos puede no haberlo puesto el usuario: la extensión de la evidencia añade el hallazgo que fija una tabla con un 1. Medido con \texttt{0\_\allowbreak{}miscellaneous/\allowbreak{}networks/\allowbreak{}mid/\allowbreak{}DMHEE/\allowbreak{}MID-dmhee-2.\allowbreak{}5.\allowbreak{}pgmx}, cuyo «State {[}0{]}» tiene probabilidad 1 en «State A»: un hallazgo en cualquier otro estado, previo o posterior a la resolución, hace fallar \texttt{VEPropagation.\allowbreak{}get\allowbreak{}All\allowbreak{}Evidence} (línea 263) con el texto «The new evidence \ldots{} is not compatible with an old evidente: \ldots», que habla de una evidencia que el usuario no ha puesto y tiene una errata («evidente» por «evidence», \texttt{core/\allowbreak{}src/\allowbreak{}main/\allowbreak{}resources/\allowbreak{}core/\allowbreak{}localize/\allowbreak{}core\_\allowbreak{}class\_\allowbreak{}localizations\_\allowbreak{}en.\allowbreak{}xml}, línea 220);
\item
  un conjunto de hallazgos con probabilidad cero en la red: la probabilidad posterior sale toda a ceros y la normalización lanza \texttt{Cannot\allowbreak{}Normalize\allowbreak{}Potential\allowbreak{}Exception} (\texttt{Table\allowbreak{}Potential\allowbreak{}Transform.\allowbreak{}normalize}, líneas 45-48); en la propagación estocástica el equivalente es \texttt{Samples\allowbreak{}Weight\allowbreak{}Is\allowbreak{}Zero}.
\end{itemize}

Todo acaba en el manejador global \texttt{OMException\allowbreak{}Handler}: lo envuelto en \texttt{Unreachable\allowbreak{}Exception} (o una excepción no comprobada) abre el diálogo de «algo que un desarrollador no previó»; lo envuelto en \texttt{Unrecoverable\allowbreak{}Exception} abre un diálogo de error normal. Con la propagación automática activa (la opción por omisión), en modo inferencia, cada gesto acaba así:

\begin{itemize}
\tightlist
\item
  Doble clic en un estado de un nodo (el gesto principal): \texttt{Editor\allowbreak{}Input\allowbreak{}Handler}, líneas 270-277, envuelve todo en \texttt{Unreachable\allowbreak{}Exception} → diálogo de error no previsto.
\item
  Quitar un hallazgo dejando una evidencia restante imposible: \texttt{Evidence\allowbreak{}Manager.\allowbreak{}remove\allowbreak{}Finding}, líneas 191-194, apaga la propagación y envuelve en \texttt{Unreachable\allowbreak{}Exception} → diálogo de error no previsto.
\item
  Añadir un hallazgo por el diálogo del menú contextual: \texttt{set\allowbreak{}New\allowbreak{}Finding} (líneas 529-540) deshace el hallazgo y relanza; el botón Aceptar de \texttt{Ok\allowbreak{}Cancel\allowbreak{}Dialog} lo envuelve en \texttt{Unrecoverable\allowbreak{}Exception} → error normal.
\item
  Entrar en modo inferencia: \texttt{Main\allowbreak{}Panel\allowbreak{}Listener\allowbreak{}Assistant}, líneas 158-171, vuelve al modo anterior y lanza \texttt{Unrecoverable\allowbreak{}Exception} → error normal.
\item
  Crear un caso de evidencia, ir al primero, anterior, siguiente o último: \texttt{Evidence\allowbreak{}Manager}, líneas 656-822, apagan la propagación (\texttt{set\allowbreak{}Propagation\allowbreak{}Active(\allowbreak{}false)}) y relanzan; \texttt{GUIUtils.\allowbreak{}execute\allowbreak{}UIAction} lo convierte en error normal.
\item
  Quitar todos los hallazgos (líneas 394-401): apaga la propagación y relanza.
\item
  Aceptar el diálogo de opciones de propagación: \texttt{Propagation\allowbreak{}Options\allowbreak{}Dialog\allowbreak{}Listener}, líneas 68-72 → error normal.
\item
  Propagación estocástica: \texttt{Stochastic\allowbreak{}Propagation\allowbreak{}Output\allowbreak{}Frame}, líneas 201-208 → error normal.
\end{itemize}

Además, el texto de \texttt{Cannot\allowbreak{}Normalize\allowbreak{}Potential\allowbreak{}Exception} es «All the values are 0.0 in potential \{potential\}», que no dice al usuario que su evidencia es imposible; y apagar la propagación cambia el estado de la ventana sin avisar. El único sitio que trata la contradicción como un problema de los datos, con contexto, es el lector PGMX (\texttt{PGMXReader\_\allowbreak{}0\_\allowbreak{}2}, líneas 534-536, \texttt{Evidence\allowbreak{}Incompatible\allowbreak{}In\allowbreak{}File}).

No se ha ejecutado la ventana; la cadena se ha seguido leyendo el código.
\paragraph{El código}
gui/src/main/java/org/openmarkov/gui/window/edition/networkEditorPanel/EditorInputHandler.java:270-277

\begin{Shaded}
\begin{Highlighting}[]
\ControlFlowTok{try} \OperatorTok{\{}
    \KeywordTok{this}\OperatorTok{.}\FunctionTok{networkEditorPanel}\OperatorTok{.}\FunctionTok{getEvidenceManager}\OperatorTok{().}\FunctionTok{toggleFinding}\OperatorTok{(}\NormalTok{visualNode}\OperatorTok{,}\NormalTok{ visualState}\OperatorTok{);}
\OperatorTok{\}} \ControlFlowTok{catch} \OperatorTok{(}\NormalTok{IncompatibleEvidenceException }\OperatorTok{|}\NormalTok{ NotEvaluableNetworkException }\OperatorTok{|}
\NormalTok{         NonProjectablePotentialException }\OperatorTok{|}
\NormalTok{         NotEnoughMemoryException }\OperatorTok{|}\NormalTok{ DoEditException }\OperatorTok{|}\NormalTok{ CannotNormalizePotentialException }\OperatorTok{|}
\NormalTok{         ConstraintViolatedException ex}\OperatorTok{)} \OperatorTok{\{}
    \ControlFlowTok{throw} \KeywordTok{new} \FunctionTok{UnreachableException}\OperatorTok{(}\NormalTok{ex}\OperatorTok{);}
\OperatorTok{\}}
\end{Highlighting}
\end{Shaded}

gui/src/main/java/org/openmarkov/gui/window/edition/networkEditorPanel/EvidenceManager.java:191-194

\begin{Shaded}
\begin{Highlighting}[]
\OperatorTok{\}} \ControlFlowTok{catch} \OperatorTok{(}\NormalTok{NotEvaluableNetworkException }\OperatorTok{|}\NormalTok{ NonProjectablePotentialException }\OperatorTok{|}\NormalTok{ NotEnoughMemoryException }\OperatorTok{|}
\NormalTok{         IncompatibleEvidenceException }\OperatorTok{|}\NormalTok{ CannotNormalizePotentialException }\OperatorTok{|}\NormalTok{ ConstraintViolatedException e}\OperatorTok{)} \OperatorTok{\{}
    \KeywordTok{this}\OperatorTok{.}\FunctionTok{networkEditorPanel}\OperatorTok{.}\FunctionTok{setPropagationActive}\OperatorTok{(}\KeywordTok{false}\OperatorTok{);}
    \ControlFlowTok{throw} \KeywordTok{new} \FunctionTok{UnreachableException}\OperatorTok{(}\NormalTok{e}\OperatorTok{);}
\end{Highlighting}
\end{Shaded}
\paragraph{Cómo arreglarlo}
Una salida posible: una excepción del dominio con mensaje («esta evidencia es imposible en esta red», «estos dos hallazgos se contradicen»), lanzada donde ya se detecta ---la normalización de la posterior en la propagación (\texttt{VEPropagation}), los pesos de las muestras (\texttt{Stochastic\allowbreak{}Propagation}) y \texttt{Evidence\allowbreak{}Case.\allowbreak{}add\allowbreak{}Finding}--- y capturada una sola vez en los embudos que ya existen (\texttt{GUIUtils.\allowbreak{}execute\allowbreak{}UIAction}, \texttt{JMenu\allowbreak{}Item\allowbreak{}Builder}, \texttt{Ok\allowbreak{}Cancel\allowbreak{}Dialog}, \texttt{OMException\allowbreak{}Handler}). Así sobran los envoltorios «inalcanzable» de \texttt{Editor\allowbreak{}Input\allowbreak{}Handler} y de \texttt{Evidence\allowbreak{}Manager.\allowbreak{}remove\allowbreak{}Finding}.

Hay que decidir a la vez qué pasa con la propagación automática al fallar (hoy unos sitios la apagan y otros no).

Cuidado: la normalización por columnas del aprendizaje y la normalización condicionada que necesita la evolución temporal (\hyperref[err:67]{error~67}) no son evidencia imposible y no deben recibir ese mensaje.

Pruebas: con una red de dos nodos y una tabla con un cero, introducir por cada gesto la evidencia de probabilidad cero y comprobar que el mensaje y el tipo de diálogo son los mismos. Como la ventana no se prueba sin pantalla, al menos una prueba del núcleo que afirme el tipo y el mensaje de la excepción de \texttt{VEPropagation}.

\setcounter{subsection}{48}
\subsection[Construir una propagación cambia el tipo de análisis de la red del usuario]{Construir una propagación o una evolución temporal cambia el tipo de análisis guardado en la red del usuario, sin pasar por deshacer ni marcarla como modificada}\label{err:49}
\begin{ficha}
\fichakey{Estado}Presente
\fichakey{Módulo}core
\fichakey{Paquete}org.openmarkov.core.model.network
\fichakey{Ficheros}\path{core/src/main/java/org/openmarkov/core/model/network/ProbNetCopier.java:64-65} \newline \path{core/src/main/java/org/openmarkov/core/model/network/ProbNetPotentialQueries.java:133-142} \newline \path{inference/src/main/java/org/openmarkov/inference/algorithm/variableElimination/tasks/VEPropagation.java:83-94} \newline \path{inference/src/main/java/org/openmarkov/inference/algorithm/temporalevaluation/tasks/MIDTemporalEvolution.java:173,545-551} \newline \path{core/src/main/java/org/openmarkov/core/inference/InferenceAlgorithm.java:35-38,73-84} \newline \path{core/src/main/java/org/openmarkov/core/action/core/MulticriteriaEdit.java:56,70} \newline \path{gui/src/main/java/org/openmarkov/gui/window/edition/networkEditorPanel/EvidenceManager.java:585} \newline \path{gui/src/main/java/org/openmarkov/gui/dialog/inference/temporalevolution/TraceTemporalEvolutionDialog.java:319,373-376}
\fichakey{Comprobado}Leyendo el código y ejecutándolo
\fichakey{Esfuerzo}Pequeño (horas)
\fichakey{Decisión de equipo}No
\fichakey{Relacionados}\hyperref[err:35]{error~35}, \hyperref[err:75]{error~75}, \hyperref[nota:6]{nota~6}, \hyperref[nota:48]{nota~48}
\end{ficha}
\paragraph{Qué pasa}
Construir una propagación por eliminación de variables (\texttt{VEPropagation}) o una evolución temporal (\texttt{MIDTemporal\allowbreak{}Evolution}) sobre una red pone el tipo de análisis de esa red a un solo criterio. El cambio no pasa por el sistema de ediciones: no se deshace y no marca la red como modificada. La otra cara del mismo defecto de las copias de red, la lista de oyentes, es \hyperref[err:35]{error~35}.

Medido: se lee \texttt{0\_\allowbreak{}miscellaneous/\allowbreak{}networks/\allowbreak{}id/\allowbreak{}ID-CEA-test-2therapies.\allowbreak{}pgmx} (dos criterios), se pone su tipo a \texttt{COST\_\allowbreak{}EFFECTIVENESS} y se construye \texttt{new\ VEPropagation(\allowbreak{}red)} sin llamar a nada más; después, la red del llamador dice \texttt{UNICRITERION}, no está marcada como modificada y no hay nada que deshacer.

En la ventana se ve menos de lo que parece. La ventana construye \texttt{VEPropagation} sobre la red del usuario cada vez que propaga evidencia en modo inferencia (\texttt{Evidence\allowbreak{}Manager.\allowbreak{}java:585}), pero para entrar en modo inferencia con una red temporal o con más de un criterio pasa antes por el diálogo de opciones con el tipo forzado a \texttt{UNICRITERION} (\texttt{Inference\allowbreak{}Handler.\allowbreak{}java:84-85}), que ya lo pone mediante \texttt{Multicriteria\allowbreak{}Edit}. En los recorridos seguidos, la escritura del constructor coincide con lo que la ventana ya había dejado, así que no se ha encontrado un gesto en el que el usuario vea un número distinto por ella. Lo comprobado es:

\begin{itemize}
\tightlist
\item
  el tipo de análisis de la red del usuario cambia fuera del sistema de ediciones: no se deshace y no marca la red como modificada; si luego se guarda por otro motivo, el fichero lleva el tipo que dejó el algoritmo;
\item
  cualquier llamador sin ventana (una prueba, el módulo \texttt{resttemplate}, un programa que use OpenMarkov como biblioteca) que construya una propagación sobre un modelo de coste-efectividad se queda con el modelo pasado a un criterio;
\item
  el orden de construcción importa: \texttt{MIDTemporal\allowbreak{}Evolution(\allowbreak{}red,\allowbreak{}\ utilidades)} lee el tipo de las opciones compartidas (líneas 189-191) justo después de que otro constructor lo haya podido pisar.
\end{itemize}

\textbf{Contexto.} Cada red guarda unas opciones de inferencia (\texttt{Inference\allowbreak{}Options}): entre ellas, el tipo de análisis multicriterio (\texttt{Multicriteria\allowbreak{}Options.\allowbreak{}Type}: \texttt{UNICRITERION}, un solo criterio, o \texttt{COST\_\allowbreak{}EFFECTIVENESS}, coste-efectividad), las opciones temporales y las de Monte Carlo. Esas opciones las elige el usuario en el diálogo de opciones de inferencia y se guardan en el fichero \texttt{.\allowbreak{}pgmx}. Cuando el diálogo las cambia, lo hace con una edición (\texttt{Multicriteria\allowbreak{}Edit}), que se puede deshacer y marca la red como modificada.

\textbf{Por qué pasa.} Todo algoritmo de inferencia (\texttt{Inference\allowbreak{}Algorithm}) documenta su campo \texttt{prob\allowbreak{}Net} como «una copia de la red recibida» y hace \texttt{network.\allowbreak{}copy(\allowbreak{})} (líneas 35-38 y 84). Pero la copia superficial \textbf{no copia} las opciones: \texttt{Prob\allowbreak{}Net\allowbreak{}Copier.\allowbreak{}shallow\allowbreak{}Copy} le pone a la copia el mismo objeto (\texttt{dest.\allowbreak{}set\allowbreak{}Inference\allowbreak{}Options(\allowbreak{}source.\allowbreak{}get\allowbreak{}Inference\allowbreak{}Options(\allowbreak{}))}, línea 65). La red de Markov que construye \texttt{Prob\allowbreak{}Net\allowbreak{}Potential\allowbreak{}Queries.\allowbreak{}build\allowbreak{}Markov\allowbreak{}Decision\allowbreak{}Network} hace lo mismo (línea 140). La copia profunda, en cambio, sí las copia (\texttt{new\ Inference\allowbreak{}Options(\allowbreak{}.\allowbreak{}.\allowbreak{}.\allowbreak{})}, línea 83).

Dos constructores escriben en las opciones de «su copia», que son las de la red del llamador:

\begin{itemize}
\tightlist
\item
  \texttt{VEPropagation} (los dos constructores, líneas 85-86 y 91-92) pone el tipo a \texttt{UNICRITERION}.
\item
  \texttt{MIDTemporal\allowbreak{}Evolution.\allowbreak{}force\allowbreak{}Unicriterion} (líneas 545-551), llamado desde el constructor de una variable (línea 173) y desde el diálogo de evolución temporal por criterio (\texttt{Trace\allowbreak{}Temporal\allowbreak{}Evolution\allowbreak{}Dialog.\allowbreak{}java:376}), hace lo mismo.
\end{itemize}
\paragraph{El código}
inference/src/main/java/org/openmarkov/inference/algorithm/variableElimination/tasks/VEPropagation.java:83-87

\begin{Shaded}
\begin{Highlighting}[]
\KeywordTok{public} \FunctionTok{VEPropagation}\OperatorTok{(}\NormalTok{ProbNet network}\OperatorTok{)} \KeywordTok{throws} \KeywordTok{...} \OperatorTok{\{}
    \KeywordTok{super}\OperatorTok{(}\NormalTok{network}\OperatorTok{);}
\NormalTok{    probNet}\OperatorTok{.}\FunctionTok{getInferenceOptions}\OperatorTok{().}\FunctionTok{getMultiCriteriaOptions}\OperatorTok{()}
           \OperatorTok{.}\FunctionTok{setMulticriteriaType}\OperatorTok{(}\NormalTok{MulticriteriaOptions}\OperatorTok{.}\FunctionTok{Type}\OperatorTok{.}\FunctionTok{UNICRITERION}\OperatorTok{);}
\OperatorTok{\}}
\end{Highlighting}
\end{Shaded}

core/src/main/java/org/openmarkov/core/model/network/ProbNetCopier.java:64-65

\begin{Shaded}
\begin{Highlighting}[]
\CommentTok{// Shared like the potentials: a shallow copy carries the very options of the original.}
\NormalTok{dest}\OperatorTok{.}\FunctionTok{setInferenceOptions}\OperatorTok{(}\NormalTok{source}\OperatorTok{.}\FunctionTok{getInferenceOptions}\OperatorTok{());}
\end{Highlighting}
\end{Shaded}
\paragraph{Cómo arreglarlo}
Que la copia superficial copie las opciones (\texttt{new\ Inference\allowbreak{}Options(\allowbreak{}source.\allowbreak{}get\allowbreak{}Inference\allowbreak{}Options(\allowbreak{}))}, como ya hace la profunda en la línea 83). Es un objeto pequeño, y compartir las tablas de números, que es lo que hace barata la copia, no cambia. El mismo cambio hace falta en \texttt{Prob\allowbreak{}Net\allowbreak{}Potential\allowbreak{}Queries.\allowbreak{}build\allowbreak{}Markov\allowbreak{}Decision\allowbreak{}Network} (línea 140), que también reparte el mismo objeto. Con eso las escrituras de \texttt{VEPropagation} y \texttt{MIDTemporal\allowbreak{}Evolution.\allowbreak{}force\allowbreak{}Unicriterion} quedan en la copia del algoritmo.

Hay que revisar a la vez si algún llamador depende hoy de leer, en la red original, lo que el algoritmo escribió. Hay que buscar lecturas de \texttt{get\allowbreak{}Multicriteria\allowbreak{}Type(\allowbreak{})} posteriores a construir un algoritmo: \texttt{MIDTemporal\allowbreak{}Evolution.\allowbreak{}java:189-191}, \texttt{Temporal\allowbreak{}Evaluation.\allowbreak{}java:81}, \texttt{Decomposition\allowbreak{}Generate\allowbreak{}Decision\allowbreak{}Tree.\allowbreak{}java:58}, \texttt{Sensitivity\allowbreak{}Analysis\allowbreak{}Controller.\allowbreak{}java:97}, \texttt{Decision\allowbreak{}Tree\allowbreak{}Editor.\allowbreak{}java:310}.

Prueba de regresión: leer un modelo con dos criterios, poner su tipo a coste-efectividad, construir \texttt{VEPropagation} y \texttt{MIDTemporal\allowbreak{}Evolution} (constructor de una variable) sobre él y comprobar que el tipo de la red sigue siendo coste-efectividad; y comprobar que \texttt{copy(\allowbreak{})} devuelve un objeto de opciones distinto con los mismos valores.

% ══════════════════════════════════════════════════════════════════════
\section{Decisiones: políticas, criterios, agentes y árbol de decisión}\label{sec:decisiones}

Todo lo que tiene que ver con los nodos de decisión: la utilidad esperada y la política óptima que se enseñan desde el menú, las políticas impuestas, los criterios de decisión, los agentes, las redes de análisis de decisiones y el árbol de decisión.

\setcounter{subsection}{49}
\subsection[El árbol de decisión falla en casi todas las redes en formato 1.0]{El árbol de decisión no se puede construir en casi ninguna red guardada en formato 1.0 (Mediastinet, decide-test\ldots): sale un error inesperado}\label{err:50}
\begin{ficha}
\fichakey{Estado}Presente
\fichakey{Módulo}inference
\fichakey{Paquete}org.openmarkov.inference.algorithm.decompositionIntoSymmetricDANs.core
\fichakey{Ficheros}\path{inference/src/main/java/org/openmarkov/inference/algorithm/decompositionIntoSymmetricDANs/core/DANOperations.java:60-171} \newline \path{inference/src/main/java/org/openmarkov/inference/algorithm/decompositionIntoSymmetricDANs/core/DANOperations.java:396-446} \newline \path{core/src/main/java/org/openmarkov/core/model/network/UtilityOperations.java:28-50} \newline \path{gui/src/main/java/org/openmarkov/gui/window/InferenceHandler.java:193-205}
\fichakey{Comprobado}Leyendo el código y ejecutándolo
\fichakey{Esfuerzo}Medio (días)
\fichakey{Decisión de equipo}No
\fichakey{Relacionados}\hyperref[err:3]{error~3}, \hyperref[err:9]{error~9}, \hyperref[err:46]{error~46}, \hyperref[err:112]{error~112}, \hyperref[err:166]{error~166}, \hyperref[nota:180]{nota~180}
\end{ficha}
\paragraph{Qué pasa}
Con una red abierta de tipo diagrama de influencia o red de análisis de decisiones, el botón «Decision tree» de la barra de herramientas llama a \texttt{Inference\allowbreak{}Handler.\allowbreak{}show\allowbreak{}Decision\allowbreak{}Tree}, que construye un \texttt{Decision\allowbreak{}Tree\allowbreak{}Editor}; este pide el árbol a \texttt{Decision\allowbreak{}Tree\allowbreak{}Manager\allowbreak{}Impl.\allowbreak{}build\allowbreak{}Decision\allowbreak{}Tree(\allowbreak{}prob\allowbreak{}Net,\allowbreak{}\ 5)}. Con \texttt{0\_\allowbreak{}miscellaneous/\allowbreak{}networks/\allowbreak{}id/\allowbreak{}1-0/\allowbreak{}ID-mediastinet-1-0.\allowbreak{}pgmx} la construcción muere con

\begin{Verbatim}[breaklines,breakanywhere,fontsize=\small]
java.lang.NullPointerException: Cannot invoke "Variable.getVariableType()" because "variable" is null
  at TableWithEvents.setEventAsStates(TableWithEvents.java:154)
  at TableWithEvents.deepCopy(TableWithEvents.java:447)
  at UnivariateDistrPotential.deepCopy(UnivariateDistrPotential.java:402)
  at UtilityOperations.transformToUnicriterion(UtilityOperations.java:41)
  at TaskUtilities.scaleUtilitiesUnicriterion(TaskUtilities.java:79)
  at VariableElimination.unicriterionPreprocess(VariableElimination.java:107)
  at VEEvaluation.resolve(VEEvaluation.java:68)
  at DANConditionalSymmetricInference.<init>(DANConditionalSymmetricInference.java:54)
  ... DANDecisionTreeInference / IDDecisionTreeInference (recursión por ramas)
\end{Verbatim}

Es una excepción de tiempo de ejecución, así que el manejador global (\texttt{OMException\allowbreak{}Handler}) la enseña en la ventana de «error inesperado» y la pestaña del árbol no se abre. La versión antigua de la misma red, \texttt{0\_\allowbreak{}miscellaneous/\allowbreak{}networks/\allowbreak{}id/\allowbreak{}ID-mediastinet.\allowbreak{}pgmx} (con tablas en vez de distribuciones), sí da árbol.

No es un caso raro de Mediastinet. Medido construyendo el árbol de todas las redes de \texttt{0\_\allowbreak{}miscellaneous/\allowbreak{}networks/\allowbreak{}id/\allowbreak{}1-0/\allowbreak{}} y \texttt{0\_\allowbreak{}miscellaneous/\allowbreak{}networks/\allowbreak{}dan/\allowbreak{}1-0/\allowbreak{}} (formato ProbModelXML 1.0, en el que las utilidades y muchas probabilidades se guardan como potenciales \texttt{Univariate\allowbreak{}Distr}, distribución «exacta» sobre un número):

\begin{itemize}
\tightlist
\item
  43 de las 50 redes dan este mismo \texttt{Null\allowbreak{}Pointer\allowbreak{}Exception}, entre ellas \texttt{ID-decide-test-1-0}, \texttt{ID-arthronet-1-0}, \texttt{ID-CEA-minimal-1-0}, \texttt{DAN-decide-test-1-0} y las cuatro de Mediastinet;
\item
  tres fallan con «Potential \ldots{} cannot be converted to a table» (\texttt{ID-Monty-Hall-spanish-1-0}, \texttt{DAN-dating-1-0}, \texttt{DAN-used-car-buyer-1-0});
\item
  una falla con otro puntero nulo en \texttt{DANOperations.\allowbreak{}instantiate} (\texttt{DAN-arthronet-1-0}, que también falla en su versión 0.2; es \hyperref[err:63]{error~63});
\item
  dos no terminan en 60 segundos.
\end{itemize}

Las mismas redes en formato 0.2 (\texttt{0\_\allowbreak{}miscellaneous/\allowbreak{}networks/\allowbreak{}id/\allowbreak{}*.\allowbreak{}pgmx}, \texttt{0\_\allowbreak{}miscellaneous/\allowbreak{}networks/\allowbreak{}dan/\allowbreak{}*.\allowbreak{}pgmx}, con tablas) dan árbol salvo \texttt{DAN-arthronet} y \texttt{DAN-king}.

Por qué ocurre. El árbol se construye rama a rama: para cada valor de una decisión o de una variable observada, \texttt{DANOperations.\allowbreak{}instantiate} copia la red y proyecta sobre ese valor los potenciales que dependen de la variable (\texttt{project\allowbreak{}Potentials}), y si la variable es una decisión la quita de la copia (línea 165). \texttt{project\allowbreak{}Potentials} sabe proyectar \texttt{Table\allowbreak{}Potential} y \texttt{Exact\allowbreak{}Distr\allowbreak{}Potential}; para el resto intenta \texttt{table\allowbreak{}Project} y, si el potencial no se puede convertir en tabla, \textbf{se traga la excepción} y deja el potencial como estaba, todavía con la variable que se va a quitar. Un \texttt{Univariate\allowbreak{}Distr\allowbreak{}Potential} no se convierte en tabla (\texttt{Potential\allowbreak{}Cannot\allowbreak{}Be\allowbreak{}Converted\allowbreak{}To\allowbreak{}ATable}), así que queda colgando de una variable que ya no está en la red.

Medido llamando a \texttt{instantiate} sobre cada decisión de \texttt{ID-mediastinet-1-0}: tras instanciar \texttt{Decision\_\allowbreak{}TBNA} los potenciales de \texttt{TBNA\_\allowbreak{}Morbidity} y \texttt{Economic\_\allowbreak{}Cost\_\allowbreak{}TBNA} siguen nombrando \texttt{Decision\_\allowbreak{}TBNA}; tras \texttt{Treatment}, los de \texttt{Survivors\_\allowbreak{}QALE}, \texttt{Inmediate\_\allowbreak{}Survival} y \texttt{Economic\_\allowbreak{}Cost\_\allowbreak{}Treatment} siguen nombrando \texttt{Treatment}, y lo mismo con \texttt{Decision\_\allowbreak{}PET}, \texttt{Decision\_\allowbreak{}MED} y \texttt{Decision\_\allowbreak{}EBUS\_\allowbreak{}EUS}.

Después, la evaluación de la rama escala las utilidades a un solo criterio (\texttt{Utility\allowbreak{}Operations.\allowbreak{}transform\allowbreak{}To\allowbreak{}Unicriterion}), que hace \texttt{potential.\allowbreak{}deep\allowbreak{}Copy(\allowbreak{}prob\allowbreak{}Net)} de cada utilidad sobre la red de la rama. \texttt{Potential.\allowbreak{}deep\allowbreak{}Copy} busca cada variable por nombre y mete \texttt{null} por la que falta (eso es \hyperref[err:3]{error~3}), y \texttt{Table\allowbreak{}With\allowbreak{}Events.\allowbreak{}deep\allowbreak{}Copy} (clase madre de \texttt{Univariate\allowbreak{}Distr\allowbreak{}Potential}) recorre en la línea siguiente esa lista y muere.

La evaluación de cada rama (\texttt{DANConditional\allowbreak{}Symmetric\allowbreak{}Inference}) solo trata como rama de valor cero la evidencia imposible; cualquier otro error sube hasta el usuario, y por eso este fallo se ve como excepción en vez de como un árbol con valores falsos. Aunque el \texttt{null} de \texttt{deep\allowbreak{}Copy} se convirtiera en una excepción con mensaje (\hyperref[err:3]{error~3}), el árbol seguiría sin salir: lo que falta es que la rama no deje potenciales apuntando a variables quitadas.
\paragraph{El código}
inference/src/main/java/org/openmarkov/inference/algorithm/decompositionIntoSymmetricDANs/core/DANOperations.java:410-425

\begin{Shaded}
\begin{Highlighting}[]
\ControlFlowTok{for} \OperatorTok{(}\NormalTok{Potential potential }\OperatorTok{:}\NormalTok{ probNet}\OperatorTok{.}\FunctionTok{getPotentials}\OperatorTok{(}\NormalTok{variable}\OperatorTok{))} \OperatorTok{\{}
    \CommentTok{// Project all potentials except the potential of the node}
    \ControlFlowTok{if} \OperatorTok{(}\NormalTok{nodePotential }\OperatorTok{==} \KeywordTok{null} \OperatorTok{||} \OperatorTok{(!}\NormalTok{nodePotential}\OperatorTok{.}\FunctionTok{equals}\OperatorTok{(}\NormalTok{potential}\OperatorTok{)))} \OperatorTok{\{}
\NormalTok{        Variable conditionedVariable }\OperatorTok{=}\NormalTok{ potential}\OperatorTok{.}\FunctionTok{getConditionedVariable}\OperatorTok{();}
        \BuiltInTok{Node}\NormalTok{ node }\OperatorTok{=}\NormalTok{ probNet}\OperatorTok{.}\FunctionTok{getNode}\OperatorTok{(}\NormalTok{conditionedVariable}\OperatorTok{);}
        \ControlFlowTok{if} \OperatorTok{(}\NormalTok{potential }\KeywordTok{instanceof}\NormalTok{ TablePotential }\OperatorTok{||}\NormalTok{ potential }\KeywordTok{instanceof}\NormalTok{ ExactDistrPotential}\OperatorTok{)} \OperatorTok{\{}
            \CommentTok{// Maintain the type of potential.}
\NormalTok{            node}\OperatorTok{.}\FunctionTok{setPotential}\OperatorTok{(}\NormalTok{potential}\OperatorTok{.}\FunctionTok{project}\OperatorTok{(}\NormalTok{decisionEvidence}\OperatorTok{));}
        \OperatorTok{\}} \ControlFlowTok{else} \OperatorTok{\{}
            \KeywordTok{...}
\NormalTok{            TablePotential potentialProjectedFromTreeADD }\OperatorTok{=} \KeywordTok{null}\OperatorTok{;}
            \ControlFlowTok{try} \OperatorTok{\{}
\NormalTok{                potentialProjectedFromTreeADD }\OperatorTok{=}\NormalTok{ potential}\OperatorTok{.}\FunctionTok{tableProject}\OperatorTok{(}\NormalTok{decisionEvidence}\OperatorTok{,} \KeywordTok{null}\OperatorTok{);}
            \OperatorTok{\}} \ControlFlowTok{catch} \OperatorTok{(}\NormalTok{NonProjectablePotentialException ignored}\OperatorTok{)} \OperatorTok{\{}
            \OperatorTok{\}}
\end{Highlighting}
\end{Shaded}

core/src/main/java/org/openmarkov/core/model/network/UtilityOperations.java:39-42

\begin{Shaded}
\begin{Highlighting}[]
\BuiltInTok{List}\OperatorTok{\textless{}}\NormalTok{Potential}\OperatorTok{\textgreater{}}\NormalTok{ scaledPotentials }\OperatorTok{=} \KeywordTok{new} \BuiltInTok{ArrayList}\OperatorTok{\textless{}\textgreater{}();}
\ControlFlowTok{for} \OperatorTok{(}\NormalTok{Potential potential }\OperatorTok{:}\NormalTok{ utilityPotentials}\OperatorTok{)} \OperatorTok{\{}
\NormalTok{        Potential scaledPotential }\OperatorTok{=}\NormalTok{ potential}\OperatorTok{.}\FunctionTok{deepCopy}\OperatorTok{(}\NormalTok{probNet}\OperatorTok{);}
\end{Highlighting}
\end{Shaded}
\paragraph{Cómo arreglarlo}
El sitio donde se decide es \texttt{DANOperations.\allowbreak{}project\allowbreak{}Potentials}: un potencial que depende de la variable instanciada no puede quedarse sin proyectar. Opciones:

\begin{itemize}
\tightlist
\item
  \begin{enumerate}
  \def\labelenumi{(\alph{enumi})}
  \tightlist
  \item
    Proyectarlo por un camino que no exija convertirlo en tabla: los \texttt{Univariate\allowbreak{}Distr\allowbreak{}Potential} y demás potenciales con eventos necesitarían su propio \texttt{project} sobre una evidencia de decisión.
  \end{enumerate}
\item
  \begin{enumerate}
  \def\labelenumi{(\alph{enumi})}
  \setcounter{enumi}{1}
  \tightlist
  \item
    Si no hay forma de proyectarlo, no tragarse \texttt{Non\allowbreak{}Projectable\allowbreak{}Potential\allowbreak{}Exception}: dejarla subir para que \texttt{build\allowbreak{}Decision\allowbreak{}Tree} falle con un mensaje comprensible («el modelo contiene un potencial de tipo X que el árbol de decisión no sabe tratar») en vez de con un puntero nulo lejos de allí.
  \end{enumerate}
\end{itemize}

La opción (a) es la que da árbol; la (b) solo cambia el mensaje.

Hay que revisar también \texttt{update\allowbreak{}Potentials\allowbreak{}With\allowbreak{}New\allowbreak{}Variable} y el borrado de nodos por restricciones de \texttt{instantiate} (líneas 100-160), que también quitan variables de la copia sin mirar si queda algún potencial que las nombre, y \texttt{prune\allowbreak{}Zero\allowbreak{}Probability\allowbreak{}Chance\allowbreak{}Nodes}.

\hyperref[err:46]{Error~46} describe la misma incapacidad (un nodo de utilidad que es una distribución no se deja convertir en tabla) en el modo inferencia. Por qué las redes en formato 1.0 tienen distribuciones donde las 0.2 tienen tablas, y cómo guardar una red en ese formato la deja así, está en \hyperref[err:112]{error~112}; si allí se decide que una distribución exacta se pueda convertir en tabla, \texttt{project\allowbreak{}Potentials} la proyectaría por su rama de \texttt{table\allowbreak{}Project} (por lectura).

Prueba: construir con \texttt{Decision\allowbreak{}Tree\allowbreak{}Manager\allowbreak{}Impl.\allowbreak{}build\allowbreak{}Decision\allowbreak{}Tree(\allowbreak{}net,\allowbreak{}\ 5)} el árbol de \texttt{0\_\allowbreak{}miscellaneous/\allowbreak{}networks/\allowbreak{}id/\allowbreak{}1-0/\allowbreak{}ID-mediastinet-1-0.\allowbreak{}pgmx}, \texttt{0\_\allowbreak{}miscellaneous/\allowbreak{}networks/\allowbreak{}id/\allowbreak{}1-0/\allowbreak{}ID-decide-test-1-0.\allowbreak{}pgmx} y \texttt{0\_\allowbreak{}miscellaneous/\allowbreak{}networks/\allowbreak{}dan/\allowbreak{}1-0/\allowbreak{}DAN-decide-test-1-0.\allowbreak{}pgmx} y comprobar que sale (o, si se elige la opción b, que falla con la excepción declarada y su mensaje, nunca con \texttt{Null\allowbreak{}Pointer\allowbreak{}Exception}); y una prueba unitaria de \texttt{DANOperations.\allowbreak{}instantiate} que compruebe que ningún potencial de la red devuelta nombra una variable ausente de ella.

\setcounter{subsection}{50}
\subsection[«Show optimal policy» y «Show expected utility» no llegan a abrirse]{«Show optimal policy» y «Show expected utility» de una decisión fallan con NullPointerException porque el panel del tipo de política nunca se construye}\label{err:51}
\begin{ficha}
\fichakey{Estado}Presente
\fichakey{Módulo}gui
\fichakey{Paquete}org.openmarkov.gui.dialog.node
\fichakey{Ficheros}\path{gui/src/main/java/org/openmarkov/gui/dialog/node/PotentialEditPanel.java:478-486, 617-632, 299-300} \newline \path{gui/src/main/java/org/openmarkov/gui/dialog/node/PotentialEditDialog.java:95-102} \newline \path{gui/src/main/java/org/openmarkov/gui/window/edition/networkEditorPanel/NetworkEditorPanel.java:560-600} \newline \path{gui/src/main/java/org/openmarkov/gui/menutoolbar/menu/NodeContextualMenu.java:191-203} \newline \path{gui/src/main/java/org/openmarkov/gui/window/MainPanelListenerAssistant.java:276-278} \newline \path{gui/src/main/java/org/openmarkov/gui/dialog/common/PolicyTypePanel.java} \newline \path{core/src/main/java/org/openmarkov/core/model/network/Node.java:95}
\fichakey{Comprobado}Leyendo el código y ejecutándolo
\fichakey{Esfuerzo}Pequeño (horas)
\fichakey{Decisión de equipo}\textcolor{decisionrojo}{\textbf{Sí.}} Si el tipo de política (óptima, determinista, probabilista) debe poder elegirse en la interfaz o si se retira \texttt{Policy\allowbreak{}Type\allowbreak{}Panel}
\fichakey{Tapado por}\hyperref[err:52]{error~52}: en algunas redes el cálculo de la utilidad esperada falla antes de abrir la ventana.
\fichakey{Relacionados}\hyperref[err:52]{error~52}, \hyperref[err:55]{error~55}, \hyperref[err:56]{error~56}, \hyperref[nota:32]{nota~32}
\end{ficha}
\paragraph{Qué pasa}
En modo inferencia, con la evidencia compilada, el menú contextual de un nodo de decisión ofrece «Show expected utility» y «Show optimal policy» (\texttt{Node\allowbreak{}Contextual\allowbreak{}Menu}, líneas 191-203). Pedir cualquiera de las dos saca el diálogo de «error inesperado» con un \texttt{Null\allowbreak{}Pointer\allowbreak{}Exception}, y la ventana de la política o de la utilidad esperada no llega a mostrarse.

El camino desde el menú, seguido en el código: las acciones (\texttt{Main\allowbreak{}Panel\allowbreak{}Listener\allowbreak{}Assistant}, líneas 284-286) llaman a \texttt{Network\allowbreak{}Editor\allowbreak{}Panel.\allowbreak{}show\allowbreak{}Expected\allowbreak{}Utility\allowbreak{}Of\allowbreak{}Node} y \texttt{show\allowbreak{}Optimal\allowbreak{}Policy\allowbreak{}Of\allowbreak{}Node} (líneas 568-609). Las dos construyen un nodo auxiliar de tipo decisión (en la segunda, con \texttt{set\allowbreak{}Policy\allowbreak{}Type(\allowbreak{}Policy\allowbreak{}Type.\allowbreak{}OPTIMAL)}; en la primera, con el tipo de política por omisión de \texttt{Node}, que también es \texttt{OPTIMAL}, línea 95 de \texttt{Node}) y abren un \texttt{Potential\allowbreak{}Edit\allowbreak{}Dialog} de solo lectura con \texttt{request\allowbreak{}Values(\allowbreak{})}.

\texttt{Potential\allowbreak{}Edit\allowbreak{}Dialog.\allowbreak{}request\allowbreak{}Values} (líneas 94-100), cuando el nodo es de decisión, de política óptima y de solo lectura, llama a \texttt{Potential\allowbreak{}Edit\allowbreak{}Panel.\allowbreak{}set\allowbreak{}Enabled\allowbreak{}Decision\allowbreak{}Options} antes de hacer visible la ventana. Ese método (líneas 609-621) termina con \texttt{this.\allowbreak{}get\allowbreak{}Politicy\allowbreak{}Type\allowbreak{}Panel(\allowbreak{}).\allowbreak{}set\allowbreak{}Enabled\allowbreak{}Decision\allowbreak{}Options(\allowbreak{}true)}. Pero \texttt{get\allowbreak{}Politicy\allowbreak{}Type\allowbreak{}Panel} (líneas 466-474) tiene el cuerpo que construiría el panel dentro de un comentario, bajo un \texttt{/\allowbreak{}/\allowbreak{}TODO}, y el campo \texttt{pnl\allowbreak{}Policy\allowbreak{}Type} no se asigna en ningún otro sitio (tampoco se añade al panel: líneas 287-288 comentadas). Devuelve siempre \texttt{null}.

Medido con un pequeño programa de prueba: sobre un \texttt{Potential\allowbreak{}Edit\allowbreak{}Panel} con un nodo de decisión de política óptima (instancia creada sin ejecutar el constructor, porque el constructor crea un editor de comentarios que es un \texttt{JDialog} y no se puede construir sin pantalla), \texttt{set\allowbreak{}Enabled\allowbreak{}Decision\allowbreak{}Options} lanza \texttt{Null\allowbreak{}Pointer\allowbreak{}Exception:\ Cannot\ invoke\ "Policy\allowbreak{}Type\allowbreak{}Panel.\allowbreak{}set\allowbreak{}Enabled\allowbreak{}Decision\allowbreak{}Options(\allowbreak{}boolean)"\ because\ the\ return\ value\ of\ "Potential\allowbreak{}Edit\allowbreak{}Panel.\allowbreak{}get\allowbreak{}Politicy\allowbreak{}Type\allowbreak{}Panel(\allowbreak{})"\ is\ null}, en la línea 620. En la aplicación, la excepción no está envuelta y llega al manejador global.

No se ha ejecutado el camino completo desde el menú (la evaluación previa y la construcción de la ventana no se pueden probar sin pantalla), así que en alguna red la petición puede fallar antes por otro motivo: por ejemplo, \hyperref[err:52]{error~52} salta durante el cálculo de la utilidad esperada, antes de abrir la ventana.

El tipo de política no se puede elegir en ninguna parte de la interfaz, pero hoy eso no tiene otra consecuencia: fuera de la interfaz, el tipo de política solo lo lee \texttt{Remove\allowbreak{}Policy\allowbreak{}Edit} (ver \hyperref[err:56]{error~56}); ni la inferencia ni el formato \texttt{.\allowbreak{}pgmx} lo usan.
\paragraph{El código}
gui/src/main/java/org/openmarkov/gui/dialog/node/PotentialEditPanel.java:466-474

\begin{Shaded}
\begin{Highlighting}[]
\KeywordTok{private}\NormalTok{ PolicyTypePanel }\FunctionTok{getPoliticyTypePanel}\OperatorTok{()} \OperatorTok{\{}
    \CommentTok{//}\AlertTok{TODO}
    \CommentTok{/*}
\CommentTok{    if (this.pnlPolicyType == null) \{}
\CommentTok{        this.pnlPolicyType = new PolicyTypePanel(this, this.node);}
\CommentTok{    \}}
\CommentTok{     */}
    \ControlFlowTok{return} \KeywordTok{this}\OperatorTok{.}\FunctionTok{pnlPolicyType}\OperatorTok{;}
\OperatorTok{\}}
\end{Highlighting}
\end{Shaded}

gui/src/main/java/org/openmarkov/gui/dialog/node/PotentialEditPanel.java:609-621

\begin{Shaded}
\begin{Highlighting}[]
\DataTypeTok{void} \FunctionTok{setEnabledDecisionOptions}\OperatorTok{()} \OperatorTok{\{}
    \ControlFlowTok{switch} \OperatorTok{(}\KeywordTok{this}\OperatorTok{.}\FunctionTok{node}\OperatorTok{.}\FunctionTok{getPolicyType}\OperatorTok{())} \OperatorTok{\{}
        \ControlFlowTok{case}\NormalTok{ OPTIMAL}\OperatorTok{,}\NormalTok{ DETERMINISTIC }\OperatorTok{{-}\textgreater{}} \KeywordTok{this}\OperatorTok{.}\FunctionTok{getPotentialTypeJCombobox}\OperatorTok{().}\FunctionTok{setEnabled}\OperatorTok{(}\KeywordTok{false}\OperatorTok{);}
        \ControlFlowTok{case}\NormalTok{ PROBABILISTIC }\OperatorTok{{-}\textgreater{}} \OperatorTok{\{}
            \KeywordTok{...}
        \OperatorTok{\}}
    \OperatorTok{\}}
    \KeywordTok{this}\OperatorTok{.}\FunctionTok{getPoliticyTypePanel}\OperatorTok{().}\FunctionTok{setEnabledDecisionOptions}\OperatorTok{(}\KeywordTok{true}\OperatorTok{);}
\OperatorTok{\}}
\end{Highlighting}
\end{Shaded}
\paragraph{Cómo arreglarlo}
Arreglo inmediato: que \texttt{set\allowbreak{}Enabled\allowbreak{}Decision\allowbreak{}Options} no dependa de un panel que no existe (quitar la última línea o comprobar \texttt{null}).

La decisión de fondo es si el selector del tipo de política se incorpora de verdad o se retiran \texttt{Policy\allowbreak{}Type\allowbreak{}Panel}, \texttt{get\allowbreak{}Politicy\allowbreak{}Type\allowbreak{}Panel} y el campo. El constructor comentado no compila tal cual: \texttt{Policy\allowbreak{}Type\allowbreak{}Panel} espera un \texttt{Potential\allowbreak{}Edit\allowbreak{}Dialog} y se le pasaba el panel. Si se incorpora, el cambio del tipo de política debería ser una edición deshacible y no una asignación directa (ver \hyperref[err:56]{error~56}), y \texttt{Impose\allowbreak{}Policy\allowbreak{}Panel.\allowbreak{}initialize\allowbreak{}Potential} dejaría de fijar \texttt{OPTIMAL} a mano.

Al corregirlo, estas dos ventanas empezarán a abrirse, y con ellas se verán \hyperref[err:53]{error~53} (la utilidad esperada ignora la evidencia) y \hyperref[err:55]{error~55} (el título enseña el nombre de la clave).

Prueba: una prueba de interfaz que, sobre \texttt{0\_\allowbreak{}miscellaneous/\allowbreak{}networks/\allowbreak{}id/\allowbreak{}ID-decide-test.\allowbreak{}pgmx} en modo inferencia, pida «Show optimal policy» y «Show expected utility» de «Therapy» y compruebe que se abre la ventana sin excepción. El panel no se puede construir sin pantalla, así que la prueba necesita una.

\setcounter{subsection}{51}
\subsection[«Show expected utility» falla con utilidad constante y rellena ceros]{«Show expected utility» sobre una decisión muere con un índice fuera de rango cuando la utilidad que resulta no tiene variables, y rellena con ceros cuando le faltan variables}\label{err:52}
\begin{ficha}
\fichakey{Estado}Presente
\fichakey{Módulo}core, inference, gui
\fichakey{Paquete}org.openmarkov.core.model.network.potential; org.openmarkov.inference.algorithm.variableElimination.tasks
\fichakey{Ficheros}\path{inference/src/main/java/org/openmarkov/inference/algorithm/variableElimination/tasks/VEExpectedUtilityDecision.java:55-66} \newline \path{core/src/main/java/org/openmarkov/core/model/network/potential/AbstractIndexedPotential.java:183-208} \newline \path{core/src/main/java/org/openmarkov/core/model/network/potential/TablePotential.java:742-771} \newline \path{core/src/main/java/org/openmarkov/core/model/network/potential/StrategicTablePotential.java:107-110} \newline \path{gui/src/main/java/org/openmarkov/gui/window/edition/networkEditorPanel/NetworkEditorPanel.java:568-582}
\fichakey{Comprobado}Leyendo el código y ejecutándolo
\fichakey{Esfuerzo}Pequeño (horas)
\fichakey{Decisión de equipo}No
\fichakey{Tapado por}\hyperref[err:51]{error~51}: la tabla con ceros no se ve porque la ventana no llega a abrirse.
\fichakey{Relacionados}\hyperref[err:51]{error~51}, \hyperref[err:53]{error~53}, \hyperref[err:67]{error~67}
\end{ficha}
\paragraph{Qué pasa}
Con una red con decisiones en modo inferencia y el caso de evidencia compilado, clic derecho sobre una decisión sin política impuesta → «Show expected utility» (\texttt{Action\allowbreak{}Commands.\allowbreak{}DECISION\_\allowbreak{}SHOW\_\allowbreak{}EXPECTED\_\allowbreak{}UTILITY} → \texttt{Network\allowbreak{}Editor\allowbreak{}Panel.\allowbreak{}show\allowbreak{}Expected\allowbreak{}Utility\allowbreak{}Of\allowbreak{}Node}). Ese método crea \texttt{VEExpected\allowbreak{}Utility\allowbreak{}Decision(\allowbreak{}red,\allowbreak{}\ decisión)} y pide \texttt{get\allowbreak{}Expected\allowbreak{}Utility(\allowbreak{})}.

\texttt{VEExpected\allowbreak{}Utility\allowbreak{}Decision.\allowbreak{}resolve} evalúa la red con \texttt{VEEvaluation}, condicionando a los predecesores informativos de la decisión, y \textbf{reordena} la utilidad resultante para poner la decisión delante: \texttt{ve\allowbreak{}Evaluation.\allowbreak{}get\allowbreak{}Utility(\allowbreak{}).\allowbreak{}reorder(\allowbreak{}ordered\allowbreak{}Variables)}. La lista \texttt{ordered\allowbreak{}Variables} es la decisión seguida de sus predecesores informativos, y no tiene por qué coincidir con las variables de la utilidad que devuelve la evaluación. \texttt{Table\allowbreak{}Potential.\allowbreak{}reorder(\allowbreak{}List\textless{}Variable\textgreater{})} supone que el orden nuevo es una permutación de las variables del potencial. Dos casos rompen esa suposición:

\begin{enumerate}
\def\labelenumi{\arabic{enumi}.}
\tightlist
\item
  \textbf{La utilidad es una constante (cero variables).} Pasa en un modelo sin nodos de utilidad o donde la decisión no influye en ninguna utilidad. El reordenamiento llama a \texttt{get\allowbreak{}Accumulated\allowbreak{}Offsets(\allowbreak{}this.\allowbreak{}variables,\allowbreak{}\ other\allowbreak{}Variables)} con \texttt{this.\allowbreak{}variables} vacía: el método crea \texttt{ordering\ =\ new\ int{[}0{]}} y lee \texttt{ordering{[}0{]}} (línea 201), y lanza \texttt{Array\allowbreak{}Index\allowbreak{}Out\allowbreak{}Of\allowbreak{}Bounds\allowbreak{}Exception}. Es una excepción de tiempo de ejecución que el manejador global enseña en la ventana de «error inesperado». Salta durante el cálculo, antes de abrir la ventana de la utilidad esperada, así que en este caso no se llega a \hyperref[err:51]{error~51}.
\item
  \textbf{La utilidad tiene menos variables que el orden pedido.} No hay excepción: el recorrido copia los valores en las casillas del primer estado de las variables que faltan y deja el resto a cero, así que la ventana enseñaría una tabla con ceros que nadie ha calculado, como si fueran utilidades. Hoy esa ventana no llega a abrirse por \hyperref[err:51]{error~51}; esta tabla con ceros aparecerá en cuanto se corrija.
\end{enumerate}

Medido ejecutando \texttt{VEExpected\allowbreak{}Utility\allowbreak{}Decision} sobre \texttt{integration\allowbreak{}Tests/\allowbreak{}src/\allowbreak{}main/\allowbreak{}resources/\allowbreak{}networks/\allowbreak{}id/\allowbreak{}ID-only-decision-no-utility.\allowbreak{}pgmx}, \texttt{integration\allowbreak{}Tests/\allowbreak{}src/\allowbreak{}main/\allowbreak{}resources/\allowbreak{}networks/\allowbreak{}id/\allowbreak{}ID-three-dec-two-util.\allowbreak{}pgmx}, \texttt{integration\allowbreak{}Tests/\allowbreak{}src/\allowbreak{}main/\allowbreak{}resources/\allowbreak{}networks/\allowbreak{}id/\allowbreak{}ID-only-imposed-decision-and-chance-no-utility.\allowbreak{}pgmx}, \texttt{integration\allowbreak{}Tests/\allowbreak{}src/\allowbreak{}main/\allowbreak{}resources/\allowbreak{}networks/\allowbreak{}dan/\allowbreak{}DAN-only-decision-no-utility.\allowbreak{}pgmx} y \texttt{DAN-dating.\allowbreak{}pgmx} (está en \texttt{integration\allowbreak{}Tests/\allowbreak{}src/\allowbreak{}main/\allowbreak{}resources/\allowbreak{}networks/\allowbreak{}dan/\allowbreak{}} y en \texttt{0\_\allowbreak{}miscellaneous/\allowbreak{}networks/\allowbreak{}dan/\allowbreak{}}):

\begin{Verbatim}[breaklines,breakanywhere,fontsize=\small]
ID-only-decision-no-utility.pgmx / D   : ArrayIndexOutOfBoundsException: Index 0 out of bounds for length 0
ID-three-dec-two-util.pgmx / D         : ArrayIndexOutOfBoundsException
ID-three-dec-two-util.pgmx / D1        : [D1, D] [4.0, 5.0, 0.0, 0.0]
ID-three-dec-two-util.pgmx / D2        : [D2, D, D1] [4.0, 1.0, 0.0, 0.0, 5.0, 2.0, 0.0, 0.0]
DAN-dating.pgmx / Ask?                 : ArrayIndexOutOfBoundsException
DAN-dating.pgmx / NClub?               : [NClub?, Ask?, Accept] [0.52, 16.82, 0.0, 0.0, 0.0, 0.0, 0.0, 0.0]
DAN-only-decision-no-utility.pgmx / D  : ArrayIndexOutOfBoundsException
\end{Verbatim}

La pila de las excepciones:

\begin{Verbatim}[breaklines,breakanywhere,fontsize=\small]
DAN-dating.pgmx, Ask?  → ArrayIndexOutOfBoundsException: Index 0 out of bounds for length 0
   at AbstractIndexedPotential.getAccumulatedOffsets(AbstractIndexedPotential.java:201)
   at AbstractIndexedPotential.getAccumulatedOffsets(AbstractIndexedPotential.java:308)
   at StrategicTablePotential.reorder(StrategicTablePotential.java:109)
ID-only-imposed-decision-and-chance-no-utility.pgmx, D → la misma excepción, vía TablePotential.reorder(TablePotential.java:744)
   y VEExpectedUtilityDecision.resolve(VEExpectedUtilityDecision.java:65)
\end{Verbatim}

Y directamente sobre \texttt{Table\allowbreak{}Potential}: una tabla sin variables reordenada a {[}D{]} lanza la misma excepción; una tabla sobre X con valores {[}1, 2{]} reordenada a {[}D, X{]} devuelve {[}1, 0, 2, 0{]}.
\paragraph{El código}
inference/src/main/java/org/openmarkov/inference/algorithm/variableElimination/tasks/VEExpectedUtilityDecision.java:55-66

\begin{Shaded}
\begin{Highlighting}[]
\KeywordTok{private} \DataTypeTok{void} \FunctionTok{resolve}\OperatorTok{()} \KeywordTok{throws} \KeywordTok{...} \OperatorTok{\{}
    \BuiltInTok{List}\OperatorTok{\textless{}}\NormalTok{Variable}\OperatorTok{\textgreater{}}\NormalTok{ informationalPredecesors }\OperatorTok{=}\NormalTok{ ProbNetOperations}\OperatorTok{.}\FunctionTok{getInformationalPredecessors}\OperatorTok{(}\NormalTok{probNet}\OperatorTok{,}\NormalTok{ decision}\OperatorTok{);}
    \BuiltInTok{List}\OperatorTok{\textless{}}\NormalTok{Variable}\OperatorTok{\textgreater{}}\NormalTok{ orderedVariables }\OperatorTok{=} \KeywordTok{new} \BuiltInTok{ArrayList}\OperatorTok{\textless{}\textgreater{}(}\NormalTok{informationalPredecesors}\OperatorTok{);}
\NormalTok{    orderedVariables}\OperatorTok{.}\FunctionTok{remove}\OperatorTok{(}\NormalTok{decision}\OperatorTok{);}
\NormalTok{    orderedVariables}\OperatorTok{.}\FunctionTok{add}\OperatorTok{(}\DecValTok{0}\OperatorTok{,}\NormalTok{ decision}\OperatorTok{);}

\NormalTok{    VEEvaluation veEvaluation }\OperatorTok{=} \KeywordTok{new} \FunctionTok{VEEvaluation}\OperatorTok{(}\NormalTok{probNet}\OperatorTok{);}
\NormalTok{    veEvaluation}\OperatorTok{.}\FunctionTok{setConditioningVariables}\OperatorTok{(}\NormalTok{informationalPredecesors}\OperatorTok{);}

\NormalTok{    result }\OperatorTok{=}\NormalTok{ veEvaluation}\OperatorTok{.}\FunctionTok{getUtility}\OperatorTok{().}\FunctionTok{reorder}\OperatorTok{(}\NormalTok{orderedVariables}\OperatorTok{);}
\OperatorTok{\}}
\end{Highlighting}
\end{Shaded}

core/src/main/java/org/openmarkov/core/model/network/potential/AbstractIndexedPotential.java:183-203

\begin{Shaded}
\begin{Highlighting}[]
\KeywordTok{public} \DataTypeTok{static} \DataTypeTok{int}\OperatorTok{[]} \FunctionTok{getAccumulatedOffsets}\OperatorTok{(}\BuiltInTok{List}\OperatorTok{\textless{}}\NormalTok{Variable}\OperatorTok{\textgreater{}}\NormalTok{ variables}\OperatorTok{,} \BuiltInTok{List}\OperatorTok{\textless{}}\NormalTok{Variable}\OperatorTok{\textgreater{}}\NormalTok{ otherVariables}\OperatorTok{)} \OperatorTok{\{}
    \DataTypeTok{int}\NormalTok{ otherSize }\OperatorTok{=}\NormalTok{ otherVariables}\OperatorTok{.}\FunctionTok{size}\OperatorTok{();}
    \DataTypeTok{int}\NormalTok{ thisSize  }\OperatorTok{=}\NormalTok{ variables}\OperatorTok{.}\FunctionTok{size}\OperatorTok{();}
    \DataTypeTok{int}\OperatorTok{[]}\NormalTok{ accOffsetXY }\OperatorTok{=} \KeywordTok{new} \DataTypeTok{int}\OperatorTok{[}\NormalTok{thisSize}\OperatorTok{];}
    \ControlFlowTok{if} \OperatorTok{(}\NormalTok{otherSize }\OperatorTok{==} \DecValTok{0}\OperatorTok{)} \OperatorTok{\{}
        \ControlFlowTok{return}\NormalTok{ accOffsetXY}\OperatorTok{;} \CommentTok{// Initialized to 0}
    \OperatorTok{\}}
    \DataTypeTok{int}\OperatorTok{[]}\NormalTok{ ordering }\OperatorTok{=} \KeywordTok{new} \DataTypeTok{int}\OperatorTok{[}\NormalTok{thisSize}\OperatorTok{];}
    \KeywordTok{...}
    \DataTypeTok{int}\NormalTok{ ordering\_0 }\OperatorTok{=}\NormalTok{ ordering}\OperatorTok{[}\DecValTok{0}\OperatorTok{];}
\end{Highlighting}
\end{Shaded}
\paragraph{Cómo arreglarlo}
Hay dos sitios, que se complementan.

En el cálculo de desplazamientos y el reordenamiento:

\begin{itemize}
\tightlist
\item
  \texttt{Abstract\allowbreak{}Indexed\allowbreak{}Potential.\allowbreak{}get\allowbreak{}Accumulated\allowbreak{}Offsets} debe devolver el array vacío también cuando \texttt{variables} está vacía (hoy solo se protege \texttt{other\allowbreak{}Size\ ==\ 0}). Los \texttt{reorder} de \texttt{Table\allowbreak{}Potential} y \texttt{Strategic\allowbreak{}Table\allowbreak{}Potential} usan el mismo cálculo y quedan cubiertos por este cambio.
\item
  \texttt{Table\allowbreak{}Potential.\allowbreak{}reorder} debe comprobar que el orden nuevo contiene exactamente las mismas variables que el potencial y lanzar una excepción con mensaje si no; y tratar el potencial sin variables (con orden vacío, devolver una copia).
\end{itemize}

En \texttt{VEExpected\allowbreak{}Utility\allowbreak{}Decision.\allowbreak{}resolve}, no pedir un reordenamiento que no es una permutación. Lo más directo es extender antes la utilidad a todas las variables de \texttt{ordered\allowbreak{}Variables}: una utilidad que no depende de una variable tiene el mismo valor para todos sus estados, así que se repiten valores y no se ponen ceros. También vale reordenar solo sobre las variables que la utilidad tiene y extender después. Qué debe enseñar la ventana cuando la utilidad no depende de la decisión no es del todo obvio: además de repetir el valor constante, cabe un aviso de que la utilidad esperada no depende de la decisión; en cualquier caso, de forma explícita.

Otros llamadores de \texttt{reorder} que pueden recibir potenciales constantes u órdenes que no son permutaciones, y deben revisarse con la misma regla: \texttt{VESens\allowbreak{}An\allowbreak{}Tornado\allowbreak{}Spider} (línea 83), \texttt{VESens\allowbreak{}An\allowbreak{}Map} (76), \texttt{VESens\allowbreak{}An\allowbreak{}Plot} (74), \texttt{VEPropagation.\allowbreak{}Invoke\allowbreak{}Variable\allowbreak{}Elimination\allowbreak{}Core} (208), \texttt{MIDTemporal\allowbreak{}Evolution.\allowbreak{}compute\allowbreak{}Posterior\allowbreak{}Probability} (610, que además puede recibir una lista con un nulo, ver \hyperref[err:67]{error~67}) y \texttt{Table\allowbreak{}Potential\allowbreak{}Maximization} (367).

De paso, \texttt{show\allowbreak{}Expected\allowbreak{}Utility\allowbreak{}Of\allowbreak{}Node} no pasa la evidencia a la tarea (\hyperref[err:53]{error~53}).

Pruebas de regresión: \texttt{new\ VEExpected\allowbreak{}Utility\allowbreak{}Decision(\allowbreak{}red,\allowbreak{}\ Ask?).\allowbreak{}get\allowbreak{}Expected\allowbreak{}Utility(\allowbreak{})} sobre \texttt{DAN-dating.\allowbreak{}pgmx} no lanza \texttt{Array\allowbreak{}Index\allowbreak{}Out\allowbreak{}Of\allowbreak{}Bounds\allowbreak{}Exception}; \texttt{VEExpected\allowbreak{}Utility\allowbreak{}Decision} sobre \texttt{ID-only-decision-no-utility.\allowbreak{}pgmx} no lanza y devuelve una tabla sobre D con ceros legítimos (no hay utilidad); sobre \texttt{ID-three-dec-two-util.\allowbreak{}pgmx} para D1, la tabla no tiene casillas a cero fabricadas; un \texttt{Table\allowbreak{}Potential} sin variables \texttt{.\allowbreak{}reorder(\allowbreak{}List.\allowbreak{}of(\allowbreak{}v))} no se sale del array; y una prueba unitaria de \texttt{reorder} con un orden que no es permutación que exija la excepción.

\setcounter{subsection}{52}
\subsection[La utilidad esperada de una decisión ignora la evidencia previa]{La utilidad esperada de una decisión ignora la evidencia previa a la resolución}\label{err:53}
\begin{ficha}
\fichakey{Estado}Presente
\fichakey{Módulo}inference, gui
\fichakey{Paquete}org.openmarkov.inference.algorithm.variableElimination.tasks; org.openmarkov.gui.window.edition.networkEditorPanel
\fichakey{Ficheros}\path{inference/src/main/java/org/openmarkov/inference/algorithm/variableElimination/tasks/VEExpectedUtilityDecision.java:55-66} \newline \path{inference/src/main/java/org/openmarkov/inference/algorithm/variableElimination/tasks/VEOptimalIntervention.java:120-126} \newline \path{gui/src/main/java/org/openmarkov/gui/window/edition/networkEditorPanel/NetworkEditorPanel.java:568-580} \newline \path{gui/src/main/java/org/openmarkov/gui/window/edition/networkEditorPanel/NetworkEditorPanel.java:586-590}
\fichakey{Comprobado}Leyendo el código y ejecutándolo
\fichakey{Esfuerzo}Pequeño (horas)
\fichakey{Decisión de equipo}No
\fichakey{Tapado por}\hyperref[err:51]{error~51}: la ventana no llega a abrirse.
\fichakey{Relacionados}\hyperref[err:52]{error~52}, \hyperref[err:54]{error~54}
\end{ficha}
\paragraph{Qué pasa}
La utilidad esperada de cada opción de una decisión se calcula como si no se hubiera observado nada, aunque haya hallazgos previos a la resolución.

Medido con \texttt{0\_\allowbreak{}miscellaneous/\allowbreak{}networks/\allowbreak{}id/\allowbreak{}ID-decide-test.\allowbreak{}pgmx}:

\begin{Verbatim}[breaklines,breakanywhere,fontsize=\small]
Do test? sin evidencia:                         [Do test?] [9.16, 9.3289]
Do test? con Disease=present en la tarea:       [Do test?] [9.16, 9.3289]
Do test? VEEvaluation con Disease=present:      [Do test?] [7.25, 7.05]
\end{Verbatim}

Con la evidencia, la mejor opción cambia (sin evidencia conviene hacer la prueba; con Disease=present, no hacerla), y la tarea no lo refleja.

Camino del usuario: en modo inferencia, el menú contextual de un nodo de decisión sin política impuesta ofrece «Show expected utility», que llama a \texttt{Network\allowbreak{}Editor\allowbreak{}Panel.\allowbreak{}show\allowbreak{}Expected\allowbreak{}Utility\allowbreak{}Of\allowbreak{}Node}. Ese método crea la tarea \textbf{sin} llamar a \texttt{set\allowbreak{}Pre\allowbreak{}Resolution\allowbreak{}Evidence}, así que aunque la clase pasara la evidencia, la ventana seguiría enseñando la utilidad sin los hallazgos previos a la resolución que el usuario haya puesto en modo edición. Lo mismo ocurre con «Show optimal policy» (\texttt{show\allowbreak{}Optimal\allowbreak{}Policy\allowbreak{}Of\allowbreak{}Node}, que crea \texttt{VEEvaluation} sin evidencia). En cambio, la propagación que pinta el modo inferencia sí usa esa evidencia para calcular las políticas (\texttt{VEPropagation.\allowbreak{}calculate\allowbreak{}Optimal\allowbreak{}Policies}), de modo que las dos ventanas pueden contradecir lo que se ve en los nodos. Esto último se ha comprobado leyendo el código; no se ha ejecutado la ventana.

Hoy esas dos ventanas no llegan a abrirse por \hyperref[err:51]{error~51}; la contradicción descrita aparecerá en cuanto se corrija.

Por qué pasa: \texttt{VEExpected\allowbreak{}Utility\allowbreak{}Decision} (la tarea que calcula la utilidad esperada de cada opción de una decisión en función de lo que se sabe al tomarla) hereda de \texttt{Variable\allowbreak{}Elimination}, así que tiene \texttt{set\allowbreak{}Pre\allowbreak{}Resolution\allowbreak{}Evidence}. Pero su \texttt{resolve(\allowbreak{})} crea un \texttt{VEEvaluation} nuevo y nunca le pasa esa evidencia.

\texttt{VEOptimal\allowbreak{}Intervention} tiene \texttt{set\allowbreak{}Pre\allowbreak{}Resolution\allowbreak{}Evidence} y \texttt{set\allowbreak{}Conditioning\allowbreak{}Variables} con el cuerpo vacío, pero su constructor sí recibe la evidencia y se la pasa a su evaluación, y el único llamador de producción (\texttt{Inference\allowbreak{}Handler}, línea 226, «Mostrar estrategia óptima») usa ese constructor. Nadie llama a los métodos vacíos, así que por ahí hoy no se llega a ningún fallo.
\paragraph{El código}
gui/src/main/java/org/openmarkov/gui/window/edition/networkEditorPanel/NetworkEditorPanel.java:568-573

\begin{Shaded}
\begin{Highlighting}[]
    \KeywordTok{public} \DataTypeTok{void} \FunctionTok{showExpectedUtilityOfNode}\OperatorTok{()} \KeywordTok{throws}\NormalTok{ IncompatibleEvidenceException}\OperatorTok{.}\FunctionTok{EvidenceIsIncompatibleWithOther}\OperatorTok{,}\NormalTok{ NonProjectablePotentialException}\OperatorTok{,}\NormalTok{ NotEvaluableNetworkException}\OperatorTok{.}\FunctionTok{NotApplicableNetwork}\OperatorTok{,}\NormalTok{ ConstraintViolatedException }\OperatorTok{\{}
\NormalTok{        VisualNode visualNode }\OperatorTok{=} \KeywordTok{this}\OperatorTok{.}\FunctionTok{visualNetwork}\OperatorTok{.}\FunctionTok{getLastSelectedNode}\OperatorTok{();}
        \BuiltInTok{Node}\NormalTok{ node }\OperatorTok{=}\NormalTok{ visualNode}\OperatorTok{.}\FunctionTok{getNode}\OperatorTok{();}
\NormalTok{        VEExpectedUtilityDecision veExpectedUtilityDecision}
                \OperatorTok{=} \KeywordTok{new} \FunctionTok{VEExpectedUtilityDecision}\OperatorTok{(}\KeywordTok{this}\OperatorTok{.}\FunctionTok{visualNetwork}\OperatorTok{.}\FunctionTok{getProbNet}\OperatorTok{(),}\NormalTok{ node}\OperatorTok{.}\FunctionTok{getVariable}\OperatorTok{());}
\NormalTok{        Potential expectedUtility }\OperatorTok{=}\NormalTok{ veExpectedUtilityDecision}\OperatorTok{.}\FunctionTok{getExpectedUtility}\OperatorTok{();}
\end{Highlighting}
\end{Shaded}

inference/src/main/java/org/openmarkov/inference/algorithm/variableElimination/tasks/VEOptimalIntervention.java:120-126

\begin{Shaded}
\begin{Highlighting}[]
    \AttributeTok{@Override} \KeywordTok{public} \DataTypeTok{void} \FunctionTok{setPreResolutionEvidence}\OperatorTok{(}\NormalTok{EvidenceCase preresolutionEvidence}\OperatorTok{)} \OperatorTok{\{}

    \OperatorTok{\}}

    \AttributeTok{@Override} \KeywordTok{public} \DataTypeTok{void} \FunctionTok{setConditioningVariables}\OperatorTok{(}\BuiltInTok{List}\OperatorTok{\textless{}}\NormalTok{Variable}\OperatorTok{\textgreater{}}\NormalTok{ conditioningVariables}\OperatorTok{)} \OperatorTok{\{}

    \OperatorTok{\}}
\end{Highlighting}
\end{Shaded}
\paragraph{Cómo arreglarlo}
En \texttt{VEExpected\allowbreak{}Utility\allowbreak{}Decision.\allowbreak{}resolve(\allowbreak{})}, pasar \texttt{get\allowbreak{}Pre\allowbreak{}Resolution\allowbreak{}Evidence(\allowbreak{})} a la \texttt{VEEvaluation} interna.

Para que el usuario lo vea, \texttt{Network\allowbreak{}Editor\allowbreak{}Panel.\allowbreak{}show\allowbreak{}Expected\allowbreak{}Utility\allowbreak{}Of\allowbreak{}Node} y \texttt{show\allowbreak{}Optimal\allowbreak{}Policy\allowbreak{}Of\allowbreak{}Node} tienen que pasar la evidencia previa a la resolución del \texttt{Evidence\allowbreak{}Manager} (la misma que usa \texttt{do\allowbreak{}Propagation}).

En \texttt{VEOptimal\allowbreak{}Intervention}, los dos métodos vacíos deben delegar en \texttt{ve\allowbreak{}Evaluation} (o lanzar \texttt{Unsupported\allowbreak{}Operation\allowbreak{}Exception}), para que un llamador futuro no reciba silencio.

Prueba de regresión: \texttt{VEExpected\allowbreak{}Utility\allowbreak{}Decision} sobre \texttt{0\_\allowbreak{}miscellaneous/\allowbreak{}networks/\allowbreak{}id/\allowbreak{}ID-decide-test.\allowbreak{}pgmx} para «Do test?» con Disease=present debe dar {[}7,25; 7,05{]}.

\setcounter{subsection}{53}
\subsection[Un hallazgo que descarta la primera opción de una decisión anula el cálculo]{Un hallazgo previo posible que la primera opción de una decisión hace imposible se toma por evidencia imposible: la evaluación da utilidad 0 y el análisis de coste-efectividad devuelve la partición cero y revienta}\label{err:54}
\begin{ficha}
\fichakey{Estado}Presente
\fichakey{Módulo}inference, core
\fichakey{Paquete}org.openmarkov.inference.algorithm.variableElimination; org.openmarkov.core.model.network.potential.operation
\fichakey{Ficheros}\path{inference/src/main/java/org/openmarkov/inference/algorithm/variableElimination/DecisionVariableElimination.java:50-57} \newline \path{core/src/main/java/org/openmarkov/core/model/network/potential/operation/TablePotentialTransform.java:99-104} \newline \path{core/src/main/java/org/openmarkov/core/model/network/potential/operation/DiscretePotentialOperations.java:386-388} \newline \path{inference/src/main/java/org/openmarkov/inference/algorithm/variableElimination/VariableEliminationCore.java:160-173} \newline \path{inference/src/main/java/org/openmarkov/inference/algorithm/variableElimination/operation/CEAlgebra.java:181-203}
\fichakey{Comprobado}Leyendo el código y ejecutándolo
\fichakey{Esfuerzo}Medio (días)
\fichakey{Decisión de equipo}No
\fichakey{Relacionados}\hyperref[err:48]{error~48}, \hyperref[err:53]{error~53}, \hyperref[err:85]{error~85}, \hyperref[err:86]{error~86}, \hyperref[err:87]{error~87}
\end{ficha}
\paragraph{Qué pasa}
Un hallazgo previo a la resolución (el que se pone en modo edición con «Add finding») sobre una variable de azar que depende de una decisión puede ser perfectamente posible y, aun así, la evaluación y el análisis de coste-efectividad lo tratan como si fuera imposible. Pasa cuando el valor observado es imposible con la \textbf{primera} opción de la decisión y posible con otra, y antes de esa decisión hay otras variables en el modelo (una decisión anterior, o algo que se conoce al tomarla). El caso típico es el resultado de una prueba: «negative» o «positive» solo pueden darse si se decidió hacerla, y «no» suele ser la primera opción.

Lo que ve el usuario:

\begin{itemize}
\tightlist
\item
  Menú Tools → «Cost-effectiveness analysis», ámbito global: la tarea devuelve la partición de probabilidad cero y la ventana \texttt{CEPDialog} revienta con \texttt{Null\allowbreak{}Pointer\allowbreak{}Exception} (\hyperref[err:86]{error~86}).
\item
  Ámbito de decisión, con una decisión anterior a la afectada: todas sus opciones salen con la partición cero y la ventana revienta igual (\hyperref[err:86]{error~86}).
\item
  Evaluación de un solo criterio: la probabilidad de la evidencia sale 0 y la utilidad no se calcula (sale 0 en las dos redes medidas). «Show optimal strategy» del modo inferencia (\texttt{Inference\allowbreak{}Handler.\allowbreak{}show\allowbreak{}Optimal\allowbreak{}Strategy}, que pasa la evidencia previa a \texttt{VEOptimal\allowbreak{}Intervention}) recibe una estrategia que sale de un empate entre ceros, o ninguna (\texttt{null}).
\item
  La propagación que pinta el modo inferencia no se ve afectada.
\end{itemize}

Camino del usuario: abrir \texttt{0\_\allowbreak{}miscellaneous/\allowbreak{}networks/\allowbreak{}id/\allowbreak{}ID-CEA-test-2therapies-new-test.\allowbreak{}pgmx} (decisiones «Dec: Test», «Dec: New test» y «Therapy», en ese orden); en modo edición, clic derecho sobre «New test» → «Add finding» → «negative»; menú Tools → «Cost-effectiveness analysis», aceptar las opciones y elegir el ámbito global.

Por qué pasa. La eliminación de variables resuelve el modelo desde la última decisión hacia atrás. Al eliminar una decisión, \texttt{Decision\allowbreak{}Variable\allowbreak{}Elimination} maximiza la utilidad y quita la decisión de la probabilidad con \texttt{Discrete\allowbreak{}Potential\allowbreak{}Operations.\allowbreak{}project\allowbreak{}Out\allowbreak{}Variable}, que la fija en su primera opción: \texttt{Table\allowbreak{}Potential\allowbreak{}Transform.\allowbreak{}project\allowbreak{}Out\allowbreak{}Variable} pone un hallazgo en el estado 0 y proyecta. Eso supone que la probabilidad no depende de la decisión, lo que es cierto sin evidencia (una decisión no cambia la probabilidad de lo que pasó antes), pero deja de serlo en cuanto se ha observado algo que depende de ella. Con la primera opción, la evidencia tiene probabilidad cero; al eliminar después las variables de azar anteriores, \texttt{CEAlgebra.\allowbreak{}multiply\allowbreak{}And\allowbreak{}Marginalize} ve probabilidades cero y pone la partición cero, y en un solo criterio la utilidad queda multiplicada por cero.

Si antes de la decisión afectada no hay ninguna variable (como «Dec: Test» en esa red), no queda nada que multiplicar por cero: la utilidad sale bien, aunque la probabilidad de la evidencia se quede en 0; en coste-efectividad ese caso falla antes, en \texttt{optimal\allowbreak{}CEP} (\hyperref[err:85]{error~85}). En \texttt{ID-arthronet-ce.\allowbreak{}pgmx}, en cambio, «Realizar Implante» es la primera decisión pero se toma conociendo «IMC», «Diabetes» y «Alergia ATB», que se eliminan después.

Medido, siguiendo la eliminación con \texttt{0\_\allowbreak{}miscellaneous/\allowbreak{}networks/\allowbreak{}id/\allowbreak{}ID-CEA-test-2therapies-new-test.\allowbreak{}pgmx} y «New test = negative»: antes de eliminar «Dec: New test», la probabilidad sobre «Dec: New test», «Test» y «Dec: Test» vale 0,8138, 0,732 y 0,0818 en tres casillas con «Dec: New test = yes» y 0 en todas las de «no»; tras eliminar la decisión, la probabilidad sobre «Test» y «Dec: Test» es cero entera, y al eliminar «Test» todas las particiones pasan a ser la partición cero.

Resultados medidos, comparados con el mismo cálculo imponiendo como política la única opción compatible con el hallazgo (lo que debería dar):

\begin{Verbatim}[breaklines,breakanywhere,fontsize=\small]
ID-CEA-test-2therapies-new-test.pgmx
  New test = negative   coste-efectividad global: partición cero (CEPDialog: NullPointerException)
                        evaluación: utilidad 0, probabilidad 0; estrategia óptima «Dec: Test = no», de un empate
                        con «Dec: New test = yes» impuesta: 1 tramo, coste 150, efectividad 9,8486; utilidad 295 308,34
  New test = positive   igual; con la política impuesta: 3 tramos; utilidad 157 075,56
  New test = not done   (compatible con la primera opción) 4 tramos; utilidad 267 513,8: sale bien
  Test = positive       («Dec: Test» no tiene nada antes) evaluación 194 271,92, igual que con la política impuesta
ID-arthronet-ce.pgmx
  Isquemia = Menor 1 hora y media   partición cero; evaluación: utilidad 0; estrategia óptima nula
                                    con «Realizar Implante = si» impuesta: utilidad 69 423,50
\end{Verbatim}

En \texttt{0\_\allowbreak{}miscellaneous/\allowbreak{}networks/\allowbreak{}id/\allowbreak{}ID-arthronet-ce.\allowbreak{}pgmx}, el análisis global devuelve también la partición cero con «Isquemia = Mayor 1 hora y media» y con los otros dos valores de «CC\_Drenaje»; en \texttt{0\_\allowbreak{}miscellaneous/\allowbreak{}networks/\allowbreak{}id/\allowbreak{}ID-mediastinet-ce.\allowbreak{}pgmx}, con «MED = negative» y «MED = positive». Para «CC\_Drenaje = menor 800 cc o mayor 1000 cc» y «MED = negative» se ha comprobado que la propagación del modo inferencia acepta el hallazgo y pinta la decisión afectada en su segunda opción con probabilidad 1; con «New test = negative», «Dec: New test = yes».
\paragraph{El código}
inference/src/main/java/org/openmarkov/inference/algorithm/variableElimination/DecisionVariableElimination.java:50-57

\begin{Shaded}
\begin{Highlighting}[]
\NormalTok{TablePotential jointProbability}\OperatorTok{;}

\ControlFlowTok{if} \OperatorTok{(!}\NormalTok{probPotentials}\OperatorTok{.}\FunctionTok{isEmpty}\OperatorTok{())} \OperatorTok{\{}
\NormalTok{    jointProbability }\OperatorTok{=}\NormalTok{ DiscretePotentialOperations}\OperatorTok{.}\FunctionTok{multiply}\OperatorTok{(}\NormalTok{probPotentials}\OperatorTok{);}
\NormalTok{    projectedProbability }\OperatorTok{=}\NormalTok{ DiscretePotentialOperations}\OperatorTok{.}\FunctionTok{projectOutVariable}\OperatorTok{(}\NormalTok{variableToDelete}\OperatorTok{,}\NormalTok{ jointProbability}\OperatorTok{);}
\OperatorTok{\}} \ControlFlowTok{else} \OperatorTok{\{}
\NormalTok{    projectedProbability }\OperatorTok{=}\NormalTok{ DiscretePotentialOperations}\OperatorTok{.}\FunctionTok{createUnityProbabilityPotential}\OperatorTok{();}
\OperatorTok{\}}
\end{Highlighting}
\end{Shaded}

core/src/main/java/org/openmarkov/core/model/network/potential/operation/TablePotentialTransform.java:99-104

\begin{Shaded}
\begin{Highlighting}[]
\DataTypeTok{static}\NormalTok{ TablePotential }\FunctionTok{projectOutVariable}\OperatorTok{(}\NormalTok{Variable variable}\OperatorTok{,}\NormalTok{ TablePotential inputPotential}\OperatorTok{)}
        \KeywordTok{throws}\NormalTok{ IncompatibleEvidenceException}\OperatorTok{.}\FunctionTok{EvidenceIsIncompatibleWithOther}\OperatorTok{,}\NormalTok{ NonProjectablePotentialException }\OperatorTok{\{}
\NormalTok{    EvidenceCase evi }\OperatorTok{=} \KeywordTok{new} \FunctionTok{EvidenceCase}\OperatorTok{();}
\NormalTok{    evi}\OperatorTok{.}\FunctionTok{addFinding}\OperatorTok{(}\KeywordTok{new} \FunctionTok{Finding}\OperatorTok{(}\NormalTok{variable}\OperatorTok{,}\NormalTok{ variable}\OperatorTok{.}\FunctionTok{getStates}\OperatorTok{()[}\DecValTok{0}\OperatorTok{]));}
    \ControlFlowTok{return}\NormalTok{ inputPotential}\OperatorTok{.}\FunctionTok{tableProject}\OperatorTok{(}\NormalTok{evi}\OperatorTok{,} \KeywordTok{null}\OperatorTok{);}
\OperatorTok{\}}
\end{Highlighting}
\end{Shaded}
\paragraph{Cómo arreglarlo}
Al eliminar una decisión, la probabilidad que queda no puede tomarse en la primera opción cuando depende de la decisión. Opciones técnicas:

\begin{itemize}
\tightlist
\item
  multiplicar la utilidad por esa probabilidad antes de maximizar, trabajando con utilidades sin normalizar, y dividir al final por la probabilidad de la evidencia, que es la forma habitual de la eliminación de variables en diagramas de influencia con evidencia; para coste-efectividad ya existe la multiplicación de una partición por una probabilidad (\texttt{CEAlgebra.\allowbreak{}multiply});
\item
  tomar la probabilidad en la opción óptima: en un solo criterio es la política que calcula \texttt{Max\allowbreak{}Out\allowbreak{}Variable}, pero en coste-efectividad la opción óptima cambia de un tramo a otro de la disposición a pagar, así que ahí no basta.
\end{itemize}

En cualquier caso, \texttt{project\allowbreak{}Out\allowbreak{}Variable} no debería proyectar en silencio una probabilidad que depende de la variable: al menos, comprobarlo.

La misma regla vale en los tres sitios que usan \texttt{Decision\allowbreak{}Variable\allowbreak{}Elimination}: \texttt{Variable\allowbreak{}Elimination\allowbreak{}Core} (evaluación, coste-efectividad y políticas de la propagación), \texttt{Decision\allowbreak{}Tree\allowbreak{}Expansion} (línea 389) y \texttt{DANInference} (línea 113). Este arreglo y el de \hyperref[err:85]{error~85} se necesitan los dos: arreglado solo \texttt{optimal\allowbreak{}CEP}, «New test = negative» seguiría dando la partición cero; arreglado solo esto, «Test = positive» seguiría fallando en \texttt{optimal\allowbreak{}CEP}.

Prueba de regresión: con \texttt{0\_\allowbreak{}miscellaneous/\allowbreak{}networks/\allowbreak{}id/\allowbreak{}ID-CEA-test-2therapies-new-test.\allowbreak{}pgmx} y «New test = negative», \texttt{VEEvaluation} debe dar 295 308,34 y \texttt{VECEAnalysis} global un tramo con coste 150 y efectividad 9,8486, lo mismo que imponiendo «Dec: New test = yes»; y la probabilidad de la evidencia debe ser la misma en los dos cálculos.

\setcounter{subsection}{54}
\subsection[Las ventanas de utilidad esperada y política óptima muestran la clave]{Las ventanas «utilidad esperada» y «política óptima» de una decisión se titulan con el nombre de la clave en bruto}\label{err:55}
\begin{ficha}
\fichakey{Estado}Presente
\fichakey{Módulo}gui
\fichakey{Paquete}org.openmarkov.gui.window.edition.networkEditorPanel
\fichakey{Ficheros}\path{gui/src/main/java/org/openmarkov/gui/window/edition/networkEditorPanel/NetworkEditorPanel.java:568-570, 594-597} \newline \path{gui/src/main/java/org/openmarkov/gui/dialog/node/PotentialEditDialog.java:76-85} \newline \path{gui/src/main/resources/gui/localize/Dialogs_en.xml:824-829}
\fichakey{Comprobado}Leyendo el código y ejecutándolo
\fichakey{Esfuerzo}Pequeño (horas)
\fichakey{Decisión de equipo}No
\fichakey{Tapado por}\hyperref[err:51]{error~51}: la ventana no llega a abrirse.
\fichakey{Relacionados}\hyperref[err:51]{error~51}, \hyperref[nota:32]{nota~32}
\end{ficha}
\paragraph{Qué pasa}
En modo inferencia, clic derecho sobre una decisión sin política impuesta → «Show expected utility» o «Show optimal policy». Las dos opciones abren un \texttt{Potential\allowbreak{}Edit\allowbreak{}Dialog} de solo lectura cuya barra de título dice \texttt{Expected\allowbreak{}Utility\allowbreak{}Dialog.\allowbreak{}Title} y \texttt{Optimal\allowbreak{}Policy\allowbreak{}Dialog.\allowbreak{}Title}, sin el nombre de la decisión.

Hoy esas dos ventanas no llegan a abrirse por \hyperref[err:51]{error~51}; el título con la clave en bruto se verá en cuanto se corrija.

Por qué pasa: el constructor de \texttt{Potential\allowbreak{}Edit\allowbreak{}Dialog} pone un título correcto sacado de la base de textos, \texttt{get\allowbreak{}Base\allowbreak{}Title(\allowbreak{})\ +\ ":\ "\ +\ nombre\allowbreak{}Del\allowbreak{}Nodo}, pero \texttt{Network\allowbreak{}Editor\allowbreak{}Panel.\allowbreak{}show\allowbreak{}Expected\allowbreak{}Utility\allowbreak{}Of\allowbreak{}Node} y \texttt{show\allowbreak{}Optimal\allowbreak{}Policy\allowbreak{}Of\allowbreak{}Node} lo sustituyen a continuación por la cadena literal de la clave.

Comprobado consultando la base de textos: \texttt{Expected\allowbreak{}Utility\allowbreak{}Dialog.\allowbreak{}Title} → «Expected utility» y \texttt{Optimal\allowbreak{}Policy\allowbreak{}Dialog.\allowbreak{}Title} → «Optimal policy» existen en \texttt{Dialogs\_\allowbreak{}en.\allowbreak{}xml}. Es un texto que se ve mal en inglés, así que la decisión de tener la aplicación solo en inglés no lo deja fuera.
\paragraph{El código}
gui/src/main/java/org/openmarkov/gui/window/edition/networkEditorPanel/NetworkEditorPanel.java:576-579

\begin{Shaded}
\begin{Highlighting}[]
\NormalTok{PotentialEditDialog expectedUtilityDialog }\OperatorTok{=} \KeywordTok{new} \FunctionTok{PotentialEditDialog}\OperatorTok{(}\NormalTok{ComponentUtilities}\OperatorTok{.}\FunctionTok{getOwner}\OperatorTok{(}\KeywordTok{this}\OperatorTok{),}\NormalTok{ dummyNode}\OperatorTok{,} \KeywordTok{true}\OperatorTok{);}
\NormalTok{expectedUtilityDialog}\OperatorTok{.}\FunctionTok{setTitle}\OperatorTok{(}\StringTok{"ExpectedUtilityDialog.Title"}\OperatorTok{);}
\NormalTok{expectedUtilityDialog}\OperatorTok{.}\FunctionTok{requestValues}\OperatorTok{();}
\end{Highlighting}
\end{Shaded}

gui/src/main/java/org/openmarkov/gui/window/edition/networkEditorPanel/NetworkEditorPanel.java:602-605

\begin{Shaded}
\begin{Highlighting}[]
\NormalTok{PotentialEditDialog optimalPolicyDialog }\OperatorTok{=}
        \KeywordTok{new} \FunctionTok{PotentialEditDialog}\OperatorTok{(}\NormalTok{ComponentUtilities}\OperatorTok{.}\FunctionTok{getOwner}\OperatorTok{(}\KeywordTok{this}\OperatorTok{),}\NormalTok{ dummyNode}\OperatorTok{,} \KeywordTok{true}\OperatorTok{);}
\NormalTok{optimalPolicyDialog}\OperatorTok{.}\FunctionTok{setTitle}\OperatorTok{(}\StringTok{"OptimalPolicyDialog.Title"}\OperatorTok{);}
\NormalTok{optimalPolicyDialog}\OperatorTok{.}\FunctionTok{requestValues}\OperatorTok{();}
\end{Highlighting}
\end{Shaded}
\paragraph{Cómo arreglarlo}
Pedir el texto a la base: \texttt{String\allowbreak{}Database.\allowbreak{}get\allowbreak{}Unique\allowbreak{}Instance(\allowbreak{}).\allowbreak{}get\allowbreak{}String(\allowbreak{}"Expected\allowbreak{}Utility\allowbreak{}Dialog.\allowbreak{}Title")\ +\ ":\ "\ +\ node.\allowbreak{}get\allowbreak{}Name(\allowbreak{})} (y lo mismo con \texttt{Optimal\allowbreak{}Policy\allowbreak{}Dialog.\allowbreak{}Title}). Otra opción es dar a \texttt{Potential\allowbreak{}Edit\allowbreak{}Dialog} el título base por parámetro, o por una subclase que redefina \texttt{get\allowbreak{}Base\allowbreak{}Title(\allowbreak{})}, como ya hace \texttt{Impose\allowbreak{}Policy\allowbreak{}Dialog}, para que el título lleve el nombre del nodo como las demás ventanas de potencial. Los dos sitios están en el mismo fichero.

Prueba: construir el diálogo por el mismo camino (sin mostrarlo) y comprobar que \texttt{get\allowbreak{}Title(\allowbreak{})} empieza por «Expected utility» o «Optimal policy» y contiene el nombre de la decisión; o, en general, una prueba que busque \texttt{set\allowbreak{}Title(\allowbreak{}"} seguido de una clave con punto en el código de la interfaz.

\setcounter{subsection}{55}
\subsection[Deshacer «quitar la política impuesta» pierde la tabla de la política]{Deshacer «quitar la política impuesta» no devuelve la política: se pierde la tabla}\label{err:56}
\begin{ficha}
\fichakey{Estado}Presente
\fichakey{Módulo}core
\fichakey{Paquete}org.openmarkov.core.action.core
\fichakey{Ficheros}\path{core/src/main/java/org/openmarkov/core/action/core/RemovePolicyEdit.java:27-49} \newline \path{core/src/main/java/org/openmarkov/core/model/network/Node.java:95} \newline \path{gui/src/main/java/org/openmarkov/gui/dialog/node/ImposePolicyPanel.java:31-43} \newline \path{gui/src/main/java/org/openmarkov/gui/window/edition/networkEditorPanel/NetworkEditorPanel.java:546-557} \newline \path{gui/src/main/java/org/openmarkov/gui/window/MainPanelMenuAssistant.java:646-652} \newline \path{gui/src/main/java/org/openmarkov/gui/dialog/node/PotentialEditPanel.java:464-473}
\fichakey{Comprobado}Leyendo el código y ejecutándolo
\fichakey{Esfuerzo}Pequeño (horas)
\fichakey{Decisión de equipo}No
\fichakey{Relacionados}\hyperref[err:51]{error~51}
\end{ficha}
\paragraph{Qué pasa}
En modo edición, con una decisión que tiene política impuesta, el usuario pide en el menú contextual del nodo «quitar política» (\texttt{Network\allowbreak{}Editor\allowbreak{}Panel.\allowbreak{}remove\allowbreak{}Policy\allowbreak{}From\allowbreak{}Node}, habilitado cuando la decisión tiene potenciales) y luego Edición → Deshacer. La política impuesta, que el usuario escribió a mano, se pierde; la entrada del menú «Deshacer» aparenta haber funcionado.

Comprobado con un pequeño programa que reproduce esos pasos (política \{0, 1\} puesta como hace \texttt{Impose\allowbreak{}Policy\allowbreak{}Edit}, \texttt{Remove\allowbreak{}Policy\allowbreak{}Edit}, deshacer): tras deshacer, la decisión tiene 0 potenciales.

Por qué pasa: \texttt{Remove\allowbreak{}Policy\allowbreak{}Edit} (módulo \texttt{core}) quita la política impuesta a una decisión. En el constructor guarda el potencial solo si \texttt{node.\allowbreak{}get\allowbreak{}Policy\allowbreak{}Type(\allowbreak{})\ !=\ Policy\allowbreak{}Type.\allowbreak{}OPTIMAL}; en \texttt{do\allowbreak{}Edit(\allowbreak{})} vacía los potenciales del nodo siempre; en \texttt{undo(\allowbreak{})}, si no guardó nada, vuelve a vaciarlos.

El tipo de política nunca es otro que \texttt{OPTIMAL} en producción:

\begin{itemize}
\tightlist
\item
  el campo \texttt{Node.\allowbreak{}policy\allowbreak{}Type} nace en \texttt{OPTIMAL};
\item
  \texttt{Impose\allowbreak{}Policy\allowbreak{}Panel.\allowbreak{}initialize\allowbreak{}Potential}, que es lo que corre al imponer una política desde la ventana, hace \texttt{node.\allowbreak{}set\allowbreak{}Policy\allowbreak{}Type(\allowbreak{}Policy\allowbreak{}Type.\allowbreak{}OPTIMAL)} antes de crear la tabla;
\item
  el lector de \texttt{.\allowbreak{}pgmx} no escribe el tipo;
\item
  el panel que permitiría elegir «determinista» o «probabilista» (\texttt{Policy\allowbreak{}Type\allowbreak{}Panel}) está desactivado: su construcción está comentada en \texttt{Potential\allowbreak{}Edit\allowbreak{}Panel} (ver \hyperref[err:51]{error~51}).
\end{itemize}

Recorriendo las redes del repositorio se encontraron decisiones con política impuesta y tipo «óptima» en \texttt{io/\allowbreak{}src/\allowbreak{}test/\allowbreak{}resources/\allowbreak{}Imposed\allowbreak{}Policy.\allowbreak{}pgmx}, \texttt{integration\allowbreak{}Tests/\allowbreak{}src/\allowbreak{}main/\allowbreak{}resources/\allowbreak{}networks/\allowbreak{}mid/\allowbreak{}21-gene.\allowbreak{}pgmx} y otras.

Además, el \texttt{get\allowbreak{}Potentials(\allowbreak{}).\allowbreak{}get(\allowbreak{}0)} del constructor lanza \texttt{Index\allowbreak{}Out\allowbreak{}Of\allowbreak{}Bounds\allowbreak{}Exception} si una decisión marcada como no óptima no tiene tabla (comprobado construyendo el caso por programa), pero no se alcanza, porque nada marca hoy una decisión como no óptima. Tampoco es un cambio registrado el tipo de política: \texttt{set\allowbreak{}Policy\allowbreak{}Type} se llama directamente, sin edición.
\paragraph{El código}
core/src/main/java/org/openmarkov/core/action/core/RemovePolicyEdit.java:27-49

\begin{Shaded}
\begin{Highlighting}[]
\KeywordTok{public} \FunctionTok{RemovePolicyEdit}\OperatorTok{(}\BuiltInTok{Node}\NormalTok{ node}\OperatorTok{)} \OperatorTok{\{}
    \KeywordTok{super}\OperatorTok{(}\NormalTok{node}\OperatorTok{.}\FunctionTok{getProbNet}\OperatorTok{());}
    \KeywordTok{this}\OperatorTok{.}\FunctionTok{node} \OperatorTok{=}\NormalTok{ node}\OperatorTok{;}
    \ControlFlowTok{if} \OperatorTok{(}\NormalTok{node}\OperatorTok{.}\FunctionTok{getNodeType}\OperatorTok{()} \OperatorTok{==}\NormalTok{ NodeType}\OperatorTok{.}\FunctionTok{DECISION} \OperatorTok{\&\&}\NormalTok{ node}\OperatorTok{.}\FunctionTok{getPolicyType}\OperatorTok{()} \OperatorTok{!=}\NormalTok{ PolicyType}\OperatorTok{.}\FunctionTok{OPTIMAL}\OperatorTok{)} \OperatorTok{\{}
\NormalTok{        oldPotential }\OperatorTok{=}\NormalTok{ node}\OperatorTok{.}\FunctionTok{getPotentials}\OperatorTok{().}\FunctionTok{get}\OperatorTok{(}\DecValTok{0}\OperatorTok{);}
    \OperatorTok{\}} \ControlFlowTok{else} \OperatorTok{\{}
\NormalTok{        oldPotential }\OperatorTok{=} \KeywordTok{null}\OperatorTok{;}
    \OperatorTok{\}}
\OperatorTok{\}}

\AttributeTok{@Override} \KeywordTok{protected} \DataTypeTok{void} \FunctionTok{doEdit}\OperatorTok{()} \OperatorTok{\{}
\NormalTok{    node}\OperatorTok{.}\FunctionTok{clearPotentials}\OperatorTok{();}
\OperatorTok{\}}

\AttributeTok{@Override} \KeywordTok{public} \DataTypeTok{void} \FunctionTok{undo}\OperatorTok{()} \OperatorTok{\{}
    \KeywordTok{super}\OperatorTok{.}\FunctionTok{undo}\OperatorTok{();}
    \ControlFlowTok{if} \OperatorTok{(}\NormalTok{oldPotential }\OperatorTok{!=} \KeywordTok{null}\OperatorTok{)} \OperatorTok{\{}
\NormalTok{        node}\OperatorTok{.}\FunctionTok{setPotential}\OperatorTok{(}\NormalTok{oldPotential}\OperatorTok{);}
    \OperatorTok{\}} \ControlFlowTok{else} \OperatorTok{\{}
\NormalTok{        node}\OperatorTok{.}\FunctionTok{clearPotentials}\OperatorTok{();}
    \OperatorTok{\}}
\OperatorTok{\}}
\end{Highlighting}
\end{Shaded}
\paragraph{Cómo arreglarlo}
Guardar lo que el nodo tenga, sea cual sea el tipo de política: una copia de la lista \texttt{node.\allowbreak{}get\allowbreak{}Potentials(\allowbreak{})} (vacía si no hay), y reponerla entera en \texttt{undo(\allowbreak{})} con \texttt{set\allowbreak{}Potentials}. Eso elimina también el \texttt{get(\allowbreak{}0)} sin guarda.

Si además se quiere que el tipo de política signifique algo, \texttt{Impose\allowbreak{}Policy\allowbreak{}Panel} no debería fijarlo en \texttt{OPTIMAL} al imponer, y cambiarlo debería ser una edición; eso es un cambio aparte.

Prueba: imponer una política \{0, 1\} con \texttt{Impose\allowbreak{}Policy\allowbreak{}Edit}, ejecutar \texttt{Remove\allowbreak{}Policy\allowbreak{}Edit} y deshacer: la decisión debe volver a tener exactamente esa tabla. Y construir \texttt{Remove\allowbreak{}Policy\allowbreak{}Edit} sobre una decisión sin tabla no debe lanzar.

\setcounter{subsection}{56}
\subsection[Renombrar un criterio de decisión no se deshace ni se cancela]{Renombrar un criterio de decisión no se deshace, ni con Deshacer ni con «Cancelar» en las propiedades de la red}\label{err:57}
\begin{ficha}
\fichakey{Estado}Presente
\fichakey{Módulo}core
\fichakey{Paquete}org.openmarkov.core.action.core
\fichakey{Ficheros}\path{core/src/main/java/org/openmarkov/core/action/core/DecisionCriteriaEdit.java:44} \newline \path{core/src/main/java/org/openmarkov/core/action/core/DecisionCriteriaEdit.java:67-83} \newline \path{core/src/main/java/org/openmarkov/core/action/core/DecisionCriteriaEdit.java:97-101} \newline \path{gui/src/main/java/org/openmarkov/gui/dialog/network/DecisionCriteriaTablePanel.java:56-70} \newline \path{gui/src/main/java/org/openmarkov/gui/dialog/network/NetworkPropertiesDialog.java:280-283}
\fichakey{Comprobado}Leyendo el código y ejecutándolo
\fichakey{Esfuerzo}Pequeño (horas)
\fichakey{Decisión de equipo}No
\fichakey{Relacionados}\hyperref[err:58]{error~58}, \hyperref[nota:36]{nota~36}, \hyperref[nota:81]{nota~81}
\end{ficha}
\paragraph{Qué pasa}
En Red → Propiedades → pestaña de criterios de decisión, el usuario edita el nombre de un criterio en la tabla (\texttt{Decision\allowbreak{}Criteria\allowbreak{}Table\allowbreak{}Panel.\allowbreak{}table\allowbreak{}Changed} construye la edición \texttt{RENAME}). El diálogo abre un historial parcial al entrar y, al pulsar «Cancelar», deshace lo hecho (\texttt{cancel\allowbreak{}Last\allowbreak{}Sub\allowbreak{}Edit\allowbreak{}History}). Pero el renombrado no se deshace: \textbf{«Cancelar» no cancela el renombrado}, y tampoco lo deshace Edición → Deshacer tras «Aceptar».

Comprobado con un pequeño programa de prueba: renombrar «cost» a «coste» y deshacer deja «coste»; lo mismo dentro de un historial parcial cancelado.

Por qué pasa: \texttt{Decision\allowbreak{}Criteria\allowbreak{}Edit} (módulo \texttt{core}) hace en el constructor \texttt{last\allowbreak{}Criteria\ =\ new\ Array\allowbreak{}List\textless{}\textgreater{}(\allowbreak{}probnet.\allowbreak{}get\allowbreak{}Decision\allowbreak{}Criteria(\allowbreak{}))}, una lista nueva con los mismos objetos \texttt{Criterion}. Para \texttt{RENAME}, \texttt{do\allowbreak{}Edit(\allowbreak{})} cambia el nombre del objeto (\texttt{modified\allowbreak{}Criterion.\allowbreak{}set\allowbreak{}Criterion\allowbreak{}Name(\allowbreak{}new\allowbreak{}Name)}), y \texttt{undo(\allowbreak{})} repone la lista guardada, que contiene ese mismo objeto ya renombrado. Deshacer no cambia nada. Añadir, quitar, subir y bajar sí se deshacen, porque ahí lo que cambia es la lista.

Los criterios se comparten por referencia: la variable de cada nodo de utilidad apunta a uno (\texttt{Variable.\allowbreak{}get\allowbreak{}Decision\allowbreak{}Criterion}), así que es el propio objeto el que hay que devolver a su nombre.
\paragraph{El código}
core/src/main/java/org/openmarkov/core/action/core/DecisionCriteriaEdit.java:44, 79, 97-101

\begin{Shaded}
\begin{Highlighting}[]
\KeywordTok{this}\OperatorTok{.}\FunctionTok{lastCriteria} \OperatorTok{=} \KeywordTok{new} \BuiltInTok{ArrayList}\OperatorTok{\textless{}\textgreater{}(}\NormalTok{probnet}\OperatorTok{.}\FunctionTok{getDecisionCriteria}\OperatorTok{());}
\KeywordTok{...}
\ControlFlowTok{case}\NormalTok{ RENAME }\OperatorTok{{-}\textgreater{}} \KeywordTok{this}\OperatorTok{.}\FunctionTok{modifiedCriterion}\OperatorTok{.}\FunctionTok{setCriterionName}\OperatorTok{(}\NormalTok{newName}\OperatorTok{);}
\KeywordTok{...}
\AttributeTok{@Override} \KeywordTok{public} \DataTypeTok{void} \FunctionTok{undo}\OperatorTok{()} \OperatorTok{\{}
    \KeywordTok{super}\OperatorTok{.}\FunctionTok{undo}\OperatorTok{();}
\NormalTok{    probNet}\OperatorTok{.}\FunctionTok{setDecisionCriteria}\OperatorTok{(}\NormalTok{lastCriteria}\OperatorTok{);}
\OperatorTok{\}}
\end{Highlighting}
\end{Shaded}
\paragraph{Cómo arreglarlo}
Guardar el nombre anterior del criterio en el constructor (o en \texttt{do\allowbreak{}Edit}) y, para \texttt{RENAME}, devolvérselo al mismo objeto en \texttt{undo(\allowbreak{})}; para las demás acciones, reponer la lista como ahora.

Prueba: renombrar un criterio y deshacer debe devolver el nombre viejo; la misma prueba dentro de \texttt{open\allowbreak{}New\allowbreak{}Sub\allowbreak{}Edit\allowbreak{}History} / \texttt{cancel\allowbreak{}Last\allowbreak{}Sub\allowbreak{}Edit\allowbreak{}History} reproduce el «Cancelar» del diálogo. Conviene comprobar también que las variables de utilidad que apuntan al criterio muestran el nombre viejo tras deshacer.

\setcounter{subsection}{57}
\subsection[Confirmar un criterio sin cambiar su nombre se rechaza como repetido]{Confirmar el nombre de un criterio sin cambiarlo, o cambiar solo sus mayúsculas, se rechaza como si chocara consigo mismo}\label{err:58}
\begin{ficha}
\fichakey{Estado}Presente
\fichakey{Módulo}core
\fichakey{Paquete}org.openmarkov.core.action.core
\fichakey{Ficheros}\path{core/src/main/java/org/openmarkov/core/action/core/DecisionCriteriaEdit.java:47-65} \newline \path{core/src/main/java/org/openmarkov/core/model/network/constraint/ValidCriterionName.java:31-54} \newline \path{gui/src/main/java/org/openmarkov/gui/dialog/network/DecisionCriteriaTablePanel.java:56-84}
\fichakey{Comprobado}Leyendo el código y ejecutándolo
\fichakey{Esfuerzo}Pequeño (horas)
\fichakey{Decisión de equipo}\textcolor{decisionrojo}{\textbf{Sí.}} Si dos criterios pueden llamarse igual salvo mayúsculas (la comprobación de la edición las ignora y la de red completa no)
\fichakey{Relacionados}\hyperref[err:38]{error~38}, \hyperref[err:57]{error~57}, \hyperref[nota:143]{nota~143}
\end{ficha}
\paragraph{Qué pasa}
En Red → Propiedades → criterios de decisión, terminar de editar la celda del nombre de un criterio sin cambiar el texto, o cambiando solo sus mayúsculas («cost» → «Cost»), se rechaza como si el nombre ya lo tuviera otro criterio. En la aplicación, el gestor global de excepciones muestra un diálogo de error.

Comprobado construyendo el panel sin pantalla: escribir en la celda el mismo nombre que ya tiene lanza \texttt{Unrecoverable\allowbreak{}Exception} con causa \texttt{Criterion\allowbreak{}Name\allowbreak{}Is\allowbreak{}Already\allowbreak{}Present}. Medido también directamente sobre la edición: «cost» → «cost» y «cost» → «Cost» se rechazan ambos.

\texttt{Decision\allowbreak{}Criteria\allowbreak{}Table\allowbreak{}Panel.\allowbreak{}table\allowbreak{}Changed} construye un \texttt{RENAME} con cada evento de actualización de la columna del nombre, y una tabla Swing emite ese evento al terminar de editar una celda aunque el texto no haya cambiado (el modelo de la tabla no filtra valores iguales).

Por qué pasa: \texttt{Decision\allowbreak{}Criteria\allowbreak{}Edit.\allowbreak{}check\allowbreak{}Constraints\allowbreak{}Will\allowbreak{}Be\allowbreak{}Met} (módulo \texttt{core}), para \texttt{ADD} y \texttt{RENAME}, pasa el nombre nuevo a minúsculas y sin espacios y lo compara con \textbf{todos} los criterios de \texttt{last\allowbreak{}Criteria}, incluido \texttt{modified\allowbreak{}Criterion}. Al renombrar un criterio a su propio nombre, o a su nombre con otras mayúsculas, encuentra el propio criterio y añade \texttt{Criterion\allowbreak{}Name\allowbreak{}Is\allowbreak{}Already\allowbreak{}Present}. La restricción \texttt{Valid\allowbreak{}Criterion\allowbreak{}Name} lleva \texttt{@Constraint(\allowbreak{}default\allowbreak{}Behavior\ =\ YES)}, así que está en todas las redes (comprobado en redes bayesianas, diagramas de influencia, MID ---diagramas de influencia de Markov--- y DAN ---redes de análisis de decisiones---), y la comprobación corre siempre.

Las dos copias de la regla difieren: la comprobación de red completa, \texttt{Valid\allowbreak{}Criterion\allowbreak{}Name.\allowbreak{}check\allowbreak{}Prob\allowbreak{}Net}, compara con igualdad exacta y sí admitiría «Cost» junto a «cost».

Subir el primer criterio o bajar el último, en la misma pestaña, es otro asunto, que la ventana no deja pedir y donde no hay nada que corregir (\hyperref[nota:36]{nota~36}).
\paragraph{El código}
core/src/main/java/org/openmarkov/core/action/core/DecisionCriteriaEdit.java:52-61

\begin{Shaded}
\begin{Highlighting}[]
\ControlFlowTok{case}\NormalTok{ ADD}\OperatorTok{,}\NormalTok{ RENAME }\OperatorTok{{-}\textgreater{}} \OperatorTok{\{}
    \BuiltInTok{String}\NormalTok{ name }\OperatorTok{=} \KeywordTok{this}\OperatorTok{.}\FunctionTok{newName}\OperatorTok{.}\FunctionTok{trim}\OperatorTok{().}\FunctionTok{toLowerCase}\OperatorTok{();}
    \ControlFlowTok{if} \OperatorTok{(}\NormalTok{name}\OperatorTok{.}\FunctionTok{isEmpty}\OperatorTok{())} \OperatorTok{\{}
\NormalTok{        constraintChecker}\OperatorTok{.}\FunctionTok{addException}\OperatorTok{(}\KeywordTok{new}\NormalTok{ ConstraintViolatedException}\OperatorTok{.}\FunctionTok{CriterionNameIsEmpty}\OperatorTok{(}\NormalTok{constraint}\OperatorTok{,}\NormalTok{ modifiedCriterion}\OperatorTok{));}
    \OperatorTok{\}}
    \ControlFlowTok{for} \OperatorTok{(}\NormalTok{Criterion criterion }\OperatorTok{:} \KeywordTok{this}\OperatorTok{.}\FunctionTok{lastCriteria}\OperatorTok{)} \OperatorTok{\{}
        \ControlFlowTok{if} \OperatorTok{(}\NormalTok{criterion}\OperatorTok{.}\FunctionTok{getCriterionName}\OperatorTok{().}\FunctionTok{trim}\OperatorTok{().}\FunctionTok{toLowerCase}\OperatorTok{().}\FunctionTok{equals}\OperatorTok{(}\NormalTok{name}\OperatorTok{))} \OperatorTok{\{}
\NormalTok{            constraintChecker}\OperatorTok{.}\FunctionTok{addException}\OperatorTok{(}\KeywordTok{new}\NormalTok{ ConstraintViolatedException}\OperatorTok{.}\FunctionTok{CriterionNameIsAlreadyPresent}\OperatorTok{(}\NormalTok{constraint}\OperatorTok{,}\NormalTok{ modifiedCriterion}\OperatorTok{,}\NormalTok{ name}\OperatorTok{));}
        \OperatorTok{\}}
    \OperatorTok{\}}
\OperatorTok{\}}
\end{Highlighting}
\end{Shaded}
\paragraph{Cómo arreglarlo}
Al renombrar, saltarse \texttt{modified\allowbreak{}Criterion} en el bucle (comparación por identidad), y en la tabla no construir la edición si el nombre no ha cambiado.

Decidir si las mayúsculas cuentan y aplicar la misma regla en los dos sitios: la comprobación previa de la edición (\texttt{Decision\allowbreak{}Criteria\allowbreak{}Edit.\allowbreak{}check\allowbreak{}Constraints\allowbreak{}Will\allowbreak{}Be\allowbreak{}Met}) y la de red completa (\texttt{Valid\allowbreak{}Criterion\allowbreak{}Name.\allowbreak{}check\allowbreak{}Prob\allowbreak{}Net}). La comprobación de la edición también la usa el diálogo de criterios estándar (\texttt{Standard\allowbreak{}Criteria\allowbreak{}Dialog}, que añade criterios con \texttt{ADD}).

Prueba: renombrar un criterio a su propio nombre no debe lanzar; renombrarlo al nombre de otro criterio sí.

\setcounter{subsection}{58}
\subsection[Evaluar una red de decisiones sin criterios falla sin explicar qué falta]{Evaluar una red de decisiones que no tiene ningún criterio de decisión muere con IndexOutOfBoundsException en vez de explicar qué falta}\label{err:59}
\begin{ficha}
\fichakey{Estado}Presente
\fichakey{Módulo}core
\fichakey{Paquete}org.openmarkov.core.model.network.potential.operation
\fichakey{Ficheros}\path{core/src/main/java/org/openmarkov/core/model/network/potential/operation/TablePotentialFactory.java:41-47} \newline \path{inference/src/main/java/org/openmarkov/inference/algorithm/decompositionIntoSymmetricDANs/core/DANInference.java:64-69} \newline \path{io/src/main/java/org/openmarkov/io/probmodel/reader/PGMXReader_0_2.java:338-362} \newline \path{io/src/main/java/org/openmarkov/io/probmodel/reader/PGMXReader_0_2.java:714-727} \newline \path{gui/src/main/java/org/openmarkov/gui/dialog/network/StandardCriteriaDialog.java:82-110} \newline \path{gui/src/main/java/org/openmarkov/gui/dialog/network/StandardCriteriaDialog.java:126-136}
\fichakey{Comprobado}Leyendo el código y ejecutándolo
\fichakey{Esfuerzo}Pequeño (horas)
\fichakey{Decisión de equipo}No
\fichakey{Relacionados}\hyperref[err:174]{error~174}
\end{ficha}
\paragraph{Qué pasa}
Un criterio de decisión es la unidad en la que se mide una utilidad (por ejemplo, coste en euros o efectividad en años de vida ajustados por calidad). Cada red guarda una lista de criterios y cada nodo de utilidad apunta a uno. En OpenMarkov, «ser una utilidad» se reconoce precisamente porque el potencial lleva criterio.

Evaluar una red de análisis de decisiones con la lista de criterios vacía termina en un error interno sin relación aparente con los criterios. Medido con copias de \texttt{0\_\allowbreak{}miscellaneous/\allowbreak{}networks/\allowbreak{}dan/\allowbreak{}DAN-dating.\allowbreak{}pgmx} y \texttt{0\_\allowbreak{}miscellaneous/\allowbreak{}networks/\allowbreak{}dan/\allowbreak{}DAN-decide-test.\allowbreak{}pgmx} a las que se dejó el elemento \texttt{\textless{}Decision\allowbreak{}Criteria\textgreater{}\textless{}/\allowbreak{}Decision\allowbreak{}Criteria\textgreater{}} vacío: el lector las abre sin queja con cero criterios, y la evaluación por descomposición en redes simétricas (\texttt{DANDecomposition\allowbreak{}Into\allowbreak{}Symmetric\allowbreak{}DANs\allowbreak{}Evaluation}, la que usa la ventana para la estrategia óptima de una red de análisis de decisiones) muere en las dos con \texttt{Index\allowbreak{}Out\allowbreak{}Of\allowbreak{}Bounds\allowbreak{}Exception:\ Index\ 0\ out\ of\ bounds\ for\ length\ 0} en \texttt{DANInference.\allowbreak{}set\allowbreak{}Utility(\allowbreak{}DANInference.\allowbreak{}java:66)}, llamado desde \texttt{maximize\allowbreak{}And\allowbreak{}Set\allowbreak{}Utility} (línea 121).

Cómo se llega a una red sin criterios:

\begin{itemize}
\tightlist
\item
  Abriendo un \texttt{.\allowbreak{}pgmx} cuyo elemento \texttt{\textless{}Decision\allowbreak{}Criteria\textgreater{}} esté vacío. El escritor de OpenMarkov no lo escribe así (si la lista está vacía omite el elemento, y al leer un fichero sin ese elemento el lector crea un criterio por omisión salvo en redes solo de azar), así que el fichero tiene que venir de fuera o estar editado a mano.
\item
  Desde la ventana, medido ejecutando el método de aceptar del diálogo sin construir la ventana (\hyperref[err:174]{error~174}): el botón de criterios estándar de las propiedades de la red abre \texttt{Standard\allowbreak{}Criteria\allowbreak{}Dialog}, que no preselecciona ninguna opción; al aceptar, primero borra todos los criterios con ediciones y después, si no se eligió ninguna opción, \texttt{get\allowbreak{}Default\allowbreak{}Criteria} hace \texttt{selected\allowbreak{}Default\allowbreak{}Criteria.\allowbreak{}equals(\allowbreak{}.\allowbreak{}.\allowbreak{}.\allowbreak{})} sobre un nulo (línea 136) y lanza \texttt{Null\allowbreak{}Pointer\allowbreak{}Exception}, dejando la red sin criterios si después se aceptan las propiedades de la red. El botón de borrar un criterio de la tabla, en cambio, solo se habilita con dos o más filas, así que por ahí no se vacía la lista.
\end{itemize}

Por qué pasa: \texttt{Table\allowbreak{}Potential\allowbreak{}Factory.\allowbreak{}create\allowbreak{}Zero\allowbreak{}Utility\allowbreak{}Potential(\allowbreak{}red)} fabrica una utilidad que vale cero y le pone \texttt{red.\allowbreak{}get\allowbreak{}Decision\allowbreak{}Criteria(\allowbreak{}).\allowbreak{}get(\allowbreak{}0)}: comprueba que la red no sea nula, pero no que la lista tenga algún criterio. La llaman cinco sitios de la evaluación de redes de análisis de decisiones y de árboles de decisión (\texttt{Decision\allowbreak{}Tree\allowbreak{}Expansion:395}, \texttt{DANInference:118}, \texttt{DANOperations:566} y \texttt{:579}, \texttt{DANConditional\allowbreak{}Symmetric\allowbreak{}Inference:44}) cuando una rama se queda sin utilidad.

El mismo supuesto aparece en \texttt{DANInference.\allowbreak{}set\allowbreak{}Utility} (línea 66), que da el primer criterio de la red a toda utilidad que llegue sin criterio; en el lector de ficheros (\texttt{PGMXReader\_\allowbreak{}0\_\allowbreak{}2}, línea 726: un nodo de utilidad sin elemento \texttt{\textless{}Criterion\textgreater{}} recibe el primero de la red); y en \texttt{Standard\allowbreak{}Criteria\allowbreak{}Dialog} (línea 104).

Con una red de ese tipo sin nodos de utilidad la evaluación contestó 0 sin fallar, así que no se consiguió hacer pasar la ejecución por \texttt{create\allowbreak{}Zero\allowbreak{}Utility\allowbreak{}Potential} con una red real; leyendo el código, falla igual (\texttt{get(\allowbreak{}0)} sobre lista vacía).
\paragraph{El código}
core/src/main/java/org/openmarkov/core/model/network/potential/operation/TablePotentialFactory.java:41-47

\begin{Shaded}
\begin{Highlighting}[]
\KeywordTok{public} \DataTypeTok{static}\NormalTok{ TablePotential }\FunctionTok{createZeroUtilityPotential}\OperatorTok{(}\NormalTok{ProbNet dan}\OperatorTok{)} \OperatorTok{\{}
\NormalTok{    TablePotential newPotential }\OperatorTok{=} \FunctionTok{createOneValuePotential}\OperatorTok{(}\NormalTok{PotentialRole}\OperatorTok{.}\FunctionTok{UNSPECIFIED}\OperatorTok{,} \FloatTok{0.0}\OperatorTok{);}
    \ControlFlowTok{if} \OperatorTok{(}\NormalTok{dan }\OperatorTok{!=} \KeywordTok{null}\OperatorTok{)} \OperatorTok{\{}
\NormalTok{        newPotential}\OperatorTok{.}\FunctionTok{setCriterion}\OperatorTok{(}\NormalTok{dan}\OperatorTok{.}\FunctionTok{getDecisionCriteria}\OperatorTok{().}\FunctionTok{get}\OperatorTok{(}\DecValTok{0}\OperatorTok{));}
    \OperatorTok{\}}
    \ControlFlowTok{return}\NormalTok{ newPotential}\OperatorTok{;}
\OperatorTok{\}}
\end{Highlighting}
\end{Shaded}

inference/src/main/java/org/openmarkov/inference/algorithm/decompositionIntoSymmetricDANs/core/DANInference.java:64-69

\begin{Shaded}
\begin{Highlighting}[]
\KeywordTok{protected} \DataTypeTok{void} \FunctionTok{setUtility}\OperatorTok{(}\NormalTok{Potential util}\OperatorTok{)} \OperatorTok{\{}
    \ControlFlowTok{if} \OperatorTok{(}\NormalTok{util}\OperatorTok{.}\FunctionTok{getCriterion}\OperatorTok{()} \OperatorTok{==} \KeywordTok{null}\OperatorTok{)} \OperatorTok{\{}
\NormalTok{        util}\OperatorTok{.}\FunctionTok{setCriterion}\OperatorTok{(}\KeywordTok{this}\OperatorTok{.}\FunctionTok{probNet}\OperatorTok{.}\FunctionTok{getDecisionCriteria}\OperatorTok{().}\FunctionTok{get}\OperatorTok{(}\DecValTok{0}\OperatorTok{));}
    \OperatorTok{\}}
    \KeywordTok{this}\OperatorTok{.}\FunctionTok{utility} \OperatorTok{=}\NormalTok{ util}\OperatorTok{;}
\OperatorTok{\}}
\end{Highlighting}
\end{Shaded}
\paragraph{Cómo arreglarlo}
Hay que decidir qué significa una red de decisiones sin criterios y aplicarlo en todos los sitios que suponen que hay uno: \texttt{Table\allowbreak{}Potential\allowbreak{}Factory.\allowbreak{}create\allowbreak{}Zero\allowbreak{}Utility\allowbreak{}Potential}, \texttt{DANInference.\allowbreak{}set\allowbreak{}Utility}, \texttt{PGMXReader\_\allowbreak{}0\_\allowbreak{}2} (línea 726) y \texttt{Standard\allowbreak{}Criteria\allowbreak{}Dialog} (línea 104). Opciones:

\begin{itemize}
\tightlist
\item
  rechazar la red al evaluarla con una excepción del dominio que diga «la red no tiene criterios de decisión» (por ejemplo \texttt{Not\allowbreak{}Evaluable\allowbreak{}Network\allowbreak{}Exception});
\item
  crear un criterio por omisión al leer un \texttt{\textless{}Decision\allowbreak{}Criteria\textgreater{}} vacío, como ya se hace cuando el elemento falta.
\end{itemize}

Que \texttt{Standard\allowbreak{}Criteria\allowbreak{}Dialog} borre los criterios antes de saber que hay una opción elegida es un defecto propio, \hyperref[err:174]{error~174}.

Prueba de regresión: leer \texttt{DAN-dating.\allowbreak{}pgmx} con \texttt{\textless{}Decision\allowbreak{}Criteria\textgreater{}\textless{}/\allowbreak{}Decision\allowbreak{}Criteria\textgreater{}} y evaluarla con \texttt{DANDecomposition\allowbreak{}Into\allowbreak{}Symmetric\allowbreak{}DANs\allowbreak{}Evaluation}: debe dar el mensaje elegido (o evaluar con el criterio por omisión), no \texttt{Index\allowbreak{}Out\allowbreak{}Of\allowbreak{}Bounds\allowbreak{}Exception}; y \texttt{create\allowbreak{}Zero\allowbreak{}Utility\allowbreak{}Potential} con una red sin criterios no debe lanzar ese error.

\setcounter{subsection}{59}
\subsection[El agente de cada decisión se pierde al guardar y no se lee al abrir]{El agente asignado a cada decisión se pierde al guardar y no se lee al abrir}\label{err:60}
\begin{ficha}
\fichakey{Estado}Presente
\fichakey{Módulo}io
\fichakey{Paquete}org.openmarkov.io.probmodel.reader; org.openmarkov.io.probmodel.writer
\fichakey{Ficheros}\path{io/src/main/java/org/openmarkov/io/probmodel/reader/PGMXReader_0_2.java:318-332,669-789} \newline \path{io/src/main/java/org/openmarkov/io/probmodel/writer/PGMXWriter_0_2.java:235-266,410-432,503-508} \newline \path{core/src/main/resources/val_v4.xsd:58-61} \newline \path{io/src/test/resources/test-decpomdp-manual.pgmx:55,68,81,93} \newline \path{0_miscellaneous/networks/pomdp/Dec-POMDP-wireless-network.pgmx:153,163} \newline \path{io/src/test/java/org/openmarkov/io/probmodel/AgentsTest.java}
\fichakey{Comprobado}Leyendo el código y ejecutándolo
\fichakey{Esfuerzo}Pequeño (horas)
\fichakey{Decisión de equipo}No
\fichakey{Corregir antes}\hyperref[err:62]{error~62}: a la vez, como mínimo: en cuanto el lector asigne agentes, abrir las propiedades de una decisión leída de un fichero puede lanzar una excepción.
\fichakey{Relacionados}\hyperref[err:61]{error~61}, \hyperref[err:62]{error~62}, \hyperref[nota:72]{nota~72}
\end{ficha}
\paragraph{Qué pasa}
Abrir un Dec-POMDP (proceso de decisión de Markov parcialmente observable descentralizado: varios decisores que no comparten lo que ven) deja todas sus decisiones sin agente, y guardar una red después de asignar agentes a mano los borra del fichero. La asignación solo vive mientras la red está abierta. Hoy ningún algoritmo de inferencia lee \texttt{Variable.\allowbreak{}get\allowbreak{}Agent(\allowbreak{})} (solo la interfaz, la restricción \texttt{Only\allowbreak{}One\allowbreak{}Agent} y \texttt{Add\allowbreak{}Node\allowbreak{}Edit}), así que la pérdida no cambia ningún resultado numérico: lo que se pierde es el modelo que el usuario escribió.

Un Dec-POMDP necesita dos datos: la lista de agentes de la red y, para cada nodo de decisión, a qué agente pertenece. En OpenMarkov la lista está en \texttt{Prob\allowbreak{}Net.\allowbreak{}get\allowbreak{}Agents(\allowbreak{})} y la pertenencia en \texttt{Variable.\allowbreak{}get\allowbreak{}Agent(\allowbreak{})} (módulo \texttt{core}). El usuario asigna el agente de una decisión en el diálogo de propiedades del nodo, con el selector «Agent» (\texttt{Node\allowbreak{}Definition\allowbreak{}Panel.\allowbreak{}get\allowbreak{}JCombo\allowbreak{}Box\allowbreak{}Network\allowbreak{}Agents}), que aparece en las redes cuyo tipo tiene agentes (por ejemplo, el tipo \texttt{DEC\_\allowbreak{}POMDP}).

El formato propio lo contempla: el esquema \texttt{core/\allowbreak{}src/\allowbreak{}main/\allowbreak{}resources/\allowbreak{}val\_\allowbreak{}v4.\allowbreak{}xsd} (líneas 58-61) admite un \texttt{\textless{}Agent\textgreater{}} dentro de cada \texttt{\textless{}Variable\textgreater{}}, junto a \texttt{\textless{}Criterion\textgreater{}}, y los dos ficheros Dec-POMDP del repositorio lo usan (\texttt{io/\allowbreak{}src/\allowbreak{}test/\allowbreak{}resources/\allowbreak{}test-decpomdp-manual.\allowbreak{}pgmx}, cuatro decisiones; \texttt{0\_\allowbreak{}miscellaneous/\allowbreak{}networks/\allowbreak{}pomdp/\allowbreak{}Dec-POMDP-wireless-network.\allowbreak{}pgmx}, dos).

Medido con el lector \texttt{PGMXReader} y el escritor \texttt{PGMXWriter\_\allowbreak{}1\_\allowbreak{}0}:

\begin{Verbatim}[breaklines,breakanywhere,fontsize=\small]
leer test-decpomdp-manual.pgmx  agentes=[Agent 1, Agent 2]  D1 [0]=none D2 [0]=none D1 [1]=none D2 [1]=none
asignar Agent 1 a D1 [0], guardar: el fichero solo contiene <Agents><Agent name="Agent 1"/><Agent name="Agent 2"/></Agents>
releer                          D1 [0]=none ...
leer Dec-POMDP-wireless-network  agentes=[A, B]  Action A [0]=none Action B [0]=none   (el fichero dice A y B)
\end{Verbatim}

Por qué pasa: el lector y el escritor (módulo \texttt{io}) solo tratan la lista de la red.

\begin{itemize}
\tightlist
\item
  \texttt{PGMXReader\_\allowbreak{}0\_\allowbreak{}2.\allowbreak{}get\allowbreak{}Agents} (líneas 318-332) lee \texttt{\textless{}Agents\textgreater{}} bajo \texttt{\textless{}Prob\allowbreak{}Net\textgreater{}}; el método que lee cada variable (\texttt{load\allowbreak{}Variable\allowbreak{}Advanced\allowbreak{}Information}, líneas 669-789) lee el criterio de los nodos de utilidad y \texttt{\textless{}Always\allowbreak{}Observed\textgreater{}}, y no mira \texttt{\textless{}Agent\textgreater{}}. \texttt{PGMXReader\_\allowbreak{}1\_\allowbreak{}0} hereda ese método.
\item
  \texttt{PGMXWriter\_\allowbreak{}0\_\allowbreak{}2.\allowbreak{}get\allowbreak{}Agents} (líneas 235-266) escribe la lista; \texttt{get\allowbreak{}Variable\allowbreak{}Children} (líneas 410-432) escribe el criterio (\texttt{get\allowbreak{}Decision\allowbreak{}Criterion}, líneas 503-508) y no el agente. \texttt{PGMXWriter\_\allowbreak{}1\_\allowbreak{}0} no lo cambia.
\end{itemize}

La prueba \texttt{io/\allowbreak{}src/\allowbreak{}test/\allowbreak{}java/\allowbreak{}org/\allowbreak{}openmarkov/\allowbreak{}io/\allowbreak{}probmodel/\allowbreak{}Agents\allowbreak{}Test.\allowbreak{}java} solo comprueba que se leen los dos nombres de la lista, y por eso no lo detecta.
\paragraph{El código}
io/src/main/java/org/openmarkov/io/probmodel/writer/PGMXWriter\_0\_2.java:426-432

\begin{Shaded}
\begin{Highlighting}[]
\FunctionTok{getAlwaysObservedAttribute}\OperatorTok{(}\NormalTok{variableElement}\OperatorTok{,}\NormalTok{ node}\OperatorTok{);}
\NormalTok{NodeType nodeType }\OperatorTok{=}\NormalTok{ node}\OperatorTok{.}\FunctionTok{getNodeType}\OperatorTok{();}
\ControlFlowTok{if} \OperatorTok{(}\NormalTok{nodeType }\OperatorTok{==}\NormalTok{ NodeType}\OperatorTok{.}\FunctionTok{UTILITY}\OperatorTok{)} \OperatorTok{\{}
    \BuiltInTok{Element}\NormalTok{ decisionCriterionElement }\OperatorTok{=} \KeywordTok{new} \BuiltInTok{Element}\OperatorTok{(}\NormalTok{XMLTags}\OperatorTok{.}\FunctionTok{CRITERION}\OperatorTok{.}\FunctionTok{toString}\OperatorTok{());}
    \FunctionTok{getDecisionCriterion}\OperatorTok{(}\NormalTok{variableElement}\OperatorTok{,}\NormalTok{ decisionCriterionElement}\OperatorTok{,}\NormalTok{ node}\OperatorTok{);}
\OperatorTok{\}}
\CommentTok{// (no hay nada equivalente para XMLTags.AGENT en una variable de decisión)}
\end{Highlighting}
\end{Shaded}
\paragraph{Cómo arreglarlo}
Escribir \texttt{\textless{}Agent\ name=".\allowbreak{}.\allowbreak{}.\allowbreak{}"/\allowbreak{}\textgreater{}} dentro de la variable cuando \texttt{node.\allowbreak{}get\allowbreak{}Variable(\allowbreak{}).\allowbreak{}get\allowbreak{}Agent(\allowbreak{})\ !=\ null}, en el orden que fija el esquema (después de \texttt{\textless{}Criterion\textgreater{}}, antes de \texttt{\textless{}Always\allowbreak{}Observed\textgreater{}}). Y leerlo en el método de lectura de variables: buscar en \texttt{prob\allowbreak{}Net.\allowbreak{}get\allowbreak{}Agents(\allowbreak{})} el agente con ese nombre (comparando texto con \texttt{equals}) y asignar \textbf{ese objeto} a la variable, igual que se hace con el criterio.

Hay que decidir qué hacer si el nombre no está en la lista: lanzar la excepción de fichero mal formado del lector, o avisar. Lo que no conviene es inventar un agente nuevo que la lista no tiene.

Hay que arreglarlo \textbf{a la vez que \hyperref[err:62]{error~62}}: en cuanto el lector asigne agentes, abrir las propiedades de una decisión de un fichero leído puede lanzar \texttt{Illegal\allowbreak{}Argument\allowbreak{}Exception} si el objeto cadena no es el mismo que el del selector. Asignar el mismo objeto de la lista, como se propone arriba, evita la excepción, pero el selector debe arreglarse igualmente.

Comprobar también que la restricción \texttt{Only\allowbreak{}One\allowbreak{}Agent} (redes de un solo agente) sigue rechazando un fichero que asigne agentes en una red que no los admite, una vez el lector los lea.

Prueba de regresión: en \texttt{Agents\allowbreak{}Test}, leer \texttt{test-decpomdp-manual.\allowbreak{}pgmx} y comprobar que \texttt{D1\ {[}0{]}} y \texttt{D1\ {[}1{]}} tienen «Agent 1» y \texttt{D2\ {[}0{]}}, \texttt{D2\ {[}1{]}} «Agent 2»; guardar con \texttt{PGMXWriter\_\allowbreak{}1\_\allowbreak{}0}, releer y comprobar lo mismo.

\setcounter{subsection}{60}
\subsection[Quitar, renombrar o mover un agente descuadra las decisiones]{Quitar, renombrar o mover un agente de la red deja las decisiones con agentes que la lista ya no tiene}\label{err:61}
\begin{ficha}
\fichakey{Estado}Presente
\fichakey{Módulo}gui
\fichakey{Paquete}org.openmarkov.gui.action; org.openmarkov.core.model.network
\fichakey{Ficheros}\path{gui/src/main/java/org/openmarkov/gui/action/NetworkAgentEdit.java:48-59} \newline \path{core/src/main/java/org/openmarkov/core/model/network/ProbNetAgentManager.java:47-86} \newline \path{gui/src/main/java/org/openmarkov/gui/dialog/network/NetworkAgentsTablePanel.java:44-147} \newline \path{gui/src/main/java/org/openmarkov/gui/dialog/network/NetworkPropertiesDialog.java:79-90,269-282}
\fichakey{Comprobado}Leyendo el código y ejecutándolo
\fichakey{Esfuerzo}Medio (días)
\fichakey{Decisión de equipo}No
\fichakey{Relacionados}\hyperref[err:60]{error~60}, \hyperref[err:62]{error~62}
\end{ficha}
\paragraph{Qué pasa}
Las redes con varios decisores (tipo Dec-POMDP, proceso de decisión de Markov parcialmente observable descentralizado) tienen una lista de agentes, editable en el menú File → «Network properties\ldots» (Ctrl+D) → pestaña de agentes (\texttt{Network\allowbreak{}Agents\allowbreak{}Table\allowbreak{}Panel}), y cada decisión puede pertenecer a uno (\texttt{Variable.\allowbreak{}get\allowbreak{}Agent(\allowbreak{})}). Cada operación de la pestaña ejecuta una \texttt{Network\allowbreak{}Agent\allowbreak{}Edit} (módulo \texttt{gui}, paquete \texttt{org.\allowbreak{}openmarkov.\allowbreak{}gui.\allowbreak{}action}), que delega en \texttt{Prob\allowbreak{}Net\allowbreak{}Agent\allowbreak{}Manager.\allowbreak{}modify\allowbreak{}Agent} (módulo \texttt{core}). Tres cosas dejan la lista y los nodos diciendo cosas distintas.

\textbf{1. Deshacer un borrado no devuelve el agente a sus decisiones.} Quitar un agente (REMOVE, líneas 58-77 de \texttt{Prob\allowbreak{}Net\allowbreak{}Agent\allowbreak{}Manager}) lo quita de la lista y pone a \texttt{null} el agente de las decisiones que eran suyas, lo cual es correcto. Pero \texttt{Network\allowbreak{}Agent\allowbreak{}Edit.\allowbreak{}undo(\allowbreak{})} (líneas 55-59) solo repone la lista; el propio código lo dice con un \texttt{/\allowbreak{}/\allowbreak{}TODO\ restore\ agents\ in\ nodes}. Esto se alcanza también sin pulsar Ctrl+Z: el botón \textbf{Cancelar} del diálogo de propiedades de la red deshace todo lo hecho en él (\texttt{do\allowbreak{}Cancel\allowbreak{}Click\allowbreak{}Before\allowbreak{}Hide} → \texttt{cancel\allowbreak{}Last\allowbreak{}Sub\allowbreak{}Edit\allowbreak{}History}, líneas 280-282). Un usuario que quita un agente por error y cancela ve la lista como estaba, pero sus decisiones se han quedado sin dueño, sin aviso.

\textbf{2. Renombrar, subir o bajar un agente borra sus propiedades.} Un agente es un \texttt{String\allowbreak{}With\allowbreak{}Properties}: un nombre y unas propiedades adicionales que el formato guarda. Para RENAME, UP y DOWN (líneas 78-84) la lista se reconstruye con \texttt{new\ String\allowbreak{}With\allowbreak{}Properties(\allowbreak{}(\allowbreak{}String)\ objects{[}0{]})}, es decir, solo con los nombres. Subir y bajar no cambian ningún nombre y aun así pierden las propiedades. Hoy la interfaz no permite escribir propiedades de agentes, así que solo se pierden las que vengan de un fichero (\texttt{PGMXReader\_\allowbreak{}0\_\allowbreak{}2.\allowbreak{}get\allowbreak{}Agents} las lee); ninguna red del repositorio las tiene. Al guardar después, se pierden del fichero.

\textbf{3. Renombrar un agente no renombra a sus decisiones.} Las decisiones que eran de «Agent 1» siguen diciendo «Agent 1» cuando la lista ya dice otra cosa. Ese nombre huérfano es el que hace que el diálogo de propiedades de esas decisiones no se abra (\hyperref[err:62]{error~62}).

Medido con una red \texttt{DECPOMDPType} con historial de deshacer, con «Agent 1» con una propiedad, D1 de «Agent 1» y D2 de «Agent 2»:

\begin{Verbatim}[breaklines,breakanywhere,fontsize=\small]
inicio                        agentes=Agent 1(1 prop.) Agent 2(0)   D1=Agent 1 D2=Agent 2
quitar Agent 1                agentes=Agent 2(0)                    D1=none    D2=Agent 2
deshacer                      agentes=Agent 1(1 prop.) Agent 2(0)   D1=none    D2=Agent 2
subir Agent 2                 agentes=Agent 2(0) Agent 1(0 prop.)   D1=Agent 1 D2=Agent 2
renombrar Agent 1 -> Doctor   agentes=Doctor(0) Agent 2(0)          D1=Agent 1 D2=Agent 2
\end{Verbatim}

El commit \texttt{8055530} corrigió que se pudieran crear agentes con el mismo nombre (issue 577) y las excepciones del panel al añadir, quitar o mover filas (issue 646); ninguno de los tres defectos de arriba cambia con eso. El deshacer del cambio de tipo de red que borraba los agentes sí está corregido y tiene prueba (\texttt{Change\allowbreak{}Network\allowbreak{}Type\allowbreak{}Edit\allowbreak{}Agents\allowbreak{}Test}); es otro camino.
\paragraph{El código}
gui/src/main/java/org/openmarkov/gui/action/NetworkAgentEdit.java:48-59

\begin{Shaded}
\begin{Highlighting}[]
\AttributeTok{@Override} \KeywordTok{protected} \DataTypeTok{void} \FunctionTok{doEdit}\OperatorTok{()} \OperatorTok{\{}
    \KeywordTok{this}\OperatorTok{.}\FunctionTok{lastAgents} \OperatorTok{=} \KeywordTok{super}\OperatorTok{.}\FunctionTok{getProbNet}\OperatorTok{().}\FunctionTok{getAgents}\OperatorTok{().}\FunctionTok{stream}\OperatorTok{()}
            \OperatorTok{.}\FunctionTok{collect}\OperatorTok{(}\NormalTok{Collectors}\OperatorTok{.}\FunctionTok{toList}\OperatorTok{());}
\NormalTok{    probNet}\OperatorTok{.}\FunctionTok{modifyAgent}\OperatorTok{(}\NormalTok{stateAction}\OperatorTok{,}\NormalTok{agentName}\OperatorTok{,}\NormalTok{dataTable}\OperatorTok{);}
\OperatorTok{\}}

\AttributeTok{@Override} \KeywordTok{public} \DataTypeTok{void} \FunctionTok{undo}\OperatorTok{()} \OperatorTok{\{}
    \KeywordTok{super}\OperatorTok{.}\FunctionTok{undo}\OperatorTok{();}
\NormalTok{    probNet}\OperatorTok{.}\FunctionTok{setAgents}\OperatorTok{(}\NormalTok{lastAgents}\OperatorTok{);}
    \CommentTok{//}\AlertTok{TODO}\CommentTok{ restore agents in nodes}
\OperatorTok{\}}
\end{Highlighting}
\end{Shaded}

core/src/main/java/org/openmarkov/core/model/network/ProbNetAgentManager.java:78-84

\begin{Shaded}
\begin{Highlighting}[]
\ControlFlowTok{case}\NormalTok{ DOWN}\OperatorTok{,}\NormalTok{ RENAME}\OperatorTok{,}\NormalTok{ UP}\OperatorTok{:}
    \BuiltInTok{ArrayList}\OperatorTok{\textless{}}\NormalTok{StringWithProperties}\OperatorTok{\textgreater{}}\NormalTok{ modifiedAgent }\OperatorTok{=} \KeywordTok{new} \BuiltInTok{ArrayList}\OperatorTok{\textless{}\textgreater{}();}
    \ControlFlowTok{for} \OperatorTok{(}\BuiltInTok{Object}\OperatorTok{[]}\NormalTok{ objects }\OperatorTok{:}\NormalTok{ dataTable}\OperatorTok{)} \OperatorTok{\{}
\NormalTok{        modifiedAgent}\OperatorTok{.}\FunctionTok{add}\OperatorTok{(}\KeywordTok{new} \FunctionTok{StringWithProperties}\OperatorTok{((}\BuiltInTok{String}\OperatorTok{)}\NormalTok{ objects}\OperatorTok{[}\DecValTok{0}\OperatorTok{]));}
    \OperatorTok{\}}
\NormalTok{    probNet}\OperatorTok{.}\FunctionTok{setAgents}\OperatorTok{(}\NormalTok{modifiedAgent}\OperatorTok{);}
    \ControlFlowTok{break}\OperatorTok{;}
\end{Highlighting}
\end{Shaded}
\paragraph{Cómo arreglarlo}
\begin{itemize}
\tightlist
\item
  Deshacer: guardar en \texttt{do\allowbreak{}Edit}, además de la lista, qué agente tenía cada decisión (o solo las que cambian), y reponerlo en \texttt{undo(\allowbreak{})}.
\item
  Mover: reordenar los objetos \texttt{String\allowbreak{}With\allowbreak{}Properties} existentes en vez de crear otros con el nombre.
\item
  Renombrar: crear el agente nuevo copiando las propiedades del viejo y cambiar el agente de las decisiones que tenían el viejo; deshacerlo debe devolver las dos cosas.
\end{itemize}

Hay que decidir de paso cómo se identifica un agente dentro de la red (por objeto o por nombre) y aplicarlo igual en \texttt{Prob\allowbreak{}Net\allowbreak{}Agent\allowbreak{}Manager}, en el selector de agente del nodo (\hyperref[err:62]{error~62}), en el lector y el escritor de ficheros (\hyperref[err:60]{error~60}) y en la copia de redes (\texttt{Variable} clona el agente con \texttt{Clone\allowbreak{}Utils.\allowbreak{}safe\allowbreak{}Clone}). Hoy \texttt{REMOVE} compara por texto y el selector por objeto.

Prueba de regresión: sobre una red \texttt{DECPOMDPType} con dos decisiones de agentes distintos y un agente con una propiedad, comprobar para cada operación (quitar, renombrar, subir, bajar) y su deshacer que la lista y el agente de cada decisión quedan coherentes y que las propiedades se conservan.

\setcounter{subsection}{61}
\subsection[Las propiedades de una decisión no se abren si su agente ya no figura]{Las propiedades de una decisión no se abren si su agente ya no coincide, como objeto, con uno de la lista}\label{err:62}
\begin{ficha}
\fichakey{Estado}Presente
\fichakey{Módulo}gui
\fichakey{Paquete}org.openmarkov.gui.dialog.node
\fichakey{Ficheros}\path{gui/src/main/java/org/openmarkov/gui/dialog/node/NodeDefinitionPanel.java:129-138,622-671,674-677} \newline \path{gui/src/main/java/org/openmarkov/gui/dialog/node/NodePropertiesDialog.java:66-81,330-338} \newline \path{core/src/main/java/org/openmarkov/core/model/network/ProbNetAgentManager.java:78-84} \newline \path{gui/src/main/java/org/openmarkov/gui/dialog/network/NetworkAgentsTablePanel.java:44-78}
\fichakey{Comprobado}Leyendo el código y ejecutándolo
\fichakey{Esfuerzo}Pequeño (horas)
\fichakey{Decisión de equipo}No
\fichakey{Relacionados}\hyperref[err:60]{error~60}, \hyperref[err:61]{error~61}, \hyperref[nota:30]{nota~30}, \hyperref[nota:72]{nota~72}
\end{ficha}
\paragraph{Qué pasa}
El diálogo de propiedades de una decisión deja de abrirse después de renombrar su agente en las propiedades de la red. Camino de hoy:

\begin{enumerate}
\def\labelenumi{\arabic{enumi}.}
\tightlist
\item
  Red de tipo Dec-POMDP (proceso de decisión de Markov parcialmente observable descentralizado; tiene «Agent 1» y «Agent 2»). Abrir las propiedades de una decisión y elegir «Agent 1». El nodo guarda la misma cadena que tenía el selector, así que de momento todo cuadra.
\item
  Abrir File → «Network properties\ldots», pestaña de agentes, y renombrar «Agent 1» a «Doctor» (\texttt{Network\allowbreak{}Agents\allowbreak{}Table\allowbreak{}Panel.\allowbreak{}table\allowbreak{}Changed} → \texttt{Network\allowbreak{}Agent\allowbreak{}Edit} RENAME). La lista pasa a decir «Doctor», pero el nodo sigue diciendo «Agent 1» (ver \hyperref[err:61]{error~61}).
\item
  Abrir otra vez las propiedades de esa decisión: el nombre «Agent 1» no está en la lista y el diálogo no se abre.
\end{enumerate}

Aunque el usuario vuelva a renombrar «Doctor» a «Agent 1», la cadena tecleada es otro objeto y la comparación sigue fallando.

Medido llamando al método real del panel sobre una red \texttt{DECPOMDPType}:

\begin{Verbatim}[breaklines,breakanywhere,fontsize=\small]
mismo objeto que en la lista                     -> se abre, marcado el puesto 1
mismo texto, otro objeto                         -> IllegalArgumentException: setSelectedIndex: 3 out of bounds
tras renombrar Agent 1 -> Doctor, decisión de Agent 1 -> IllegalArgumentException: setSelectedIndex: 3 out of bounds
tras renombrar Agent 1 -> Doctor, decisión de Agent 2 -> se abre, marcado el puesto 2
\end{Verbatim}

Por qué pasa. Cuando la red tiene lista de agentes, la fila «Agent» del diálogo es un selector con un hueco vacío y los nombres de los agentes de la red. Lo construye \texttt{Node\allowbreak{}Definition\allowbreak{}Panel.\allowbreak{}get\allowbreak{}JCombo\allowbreak{}Box\allowbreak{}Network\allowbreak{}Agents} (módulo \texttt{gui}, líneas 622-671), que busca cuál marcar comparando el nombre del agente del nodo con cada nombre de la lista mediante \texttt{==}. Entre dos \texttt{String}, \texttt{==} pregunta si son el mismo objeto en memoria, no si dicen lo mismo. Si no hay coincidencia de objeto, el bucle termina con \texttt{i} igual al número de elementos y \texttt{set\allowbreak{}Selected\allowbreak{}Index(\allowbreak{}i)} lanza \texttt{Illegal\allowbreak{}Argument\allowbreak{}Exception}.

Ese método se ejecuta dentro del constructor del panel (líneas 129-138 y, a través de \texttt{get\allowbreak{}Agents\allowbreak{}Or\allowbreak{}Decision\allowbreak{}Criteria\allowbreak{}Or\allowbreak{}Observed}, líneas 674-677), y el panel se construye dentro del constructor de \texttt{Node\allowbreak{}Properties\allowbreak{}Dialog} (líneas 66-81, \texttt{reinitialize} → \texttt{get\allowbreak{}Node\allowbreak{}Definition\allowbreak{}Panel}, líneas 330-338). Así que la excepción impide que el diálogo de propiedades se abra. Además, leyendo el código, el constructor del diálogo ya ha abierto un grupo de ediciones (\texttt{open\allowbreak{}New\allowbreak{}Sub\allowbreak{}Edit\allowbreak{}History}, línea 75) que nadie cierra al fallar; no se ha medido qué efecto tiene eso después.

Cuando se corrija \hyperref[err:60]{error~60} (leer el agente de cada decisión desde el fichero), habrá un segundo camino: si el lector crea cadenas nuevas, abrir las propiedades de cualquier decisión de un Dec-POMDP recién leído lanzará la misma excepción.
\paragraph{El código}
gui/src/main/java/org/openmarkov/gui/dialog/node/NodeDefinitionPanel.java:655-666

\begin{Shaded}
\begin{Highlighting}[]
\ControlFlowTok{if} \OperatorTok{(}\NormalTok{node}\OperatorTok{.}\FunctionTok{getVariable}\OperatorTok{().}\FunctionTok{getAgent}\OperatorTok{()} \OperatorTok{!=} \KeywordTok{null} \OperatorTok{\&\&}\NormalTok{ agents }\OperatorTok{!=} \KeywordTok{null}\OperatorTok{)} \OperatorTok{\{}
    \BuiltInTok{String}\NormalTok{ name }\OperatorTok{=}\NormalTok{ node}\OperatorTok{.}\FunctionTok{getVariable}\OperatorTok{().}\FunctionTok{getAgent}\OperatorTok{().}\FunctionTok{getString}\OperatorTok{();}
    \DataTypeTok{int}\NormalTok{ i}\OperatorTok{;}
    \ControlFlowTok{for} \OperatorTok{(}\NormalTok{i }\OperatorTok{=} \DecValTok{0}\OperatorTok{;}\NormalTok{ i }\OperatorTok{\textless{}}\NormalTok{ agentNames}\OperatorTok{.}\FunctionTok{length}\OperatorTok{;}\NormalTok{ i}\OperatorTok{++)} \OperatorTok{\{}
        \ControlFlowTok{if} \OperatorTok{(}\NormalTok{name }\OperatorTok{==}\NormalTok{ agentNames}\OperatorTok{[}\NormalTok{i}\OperatorTok{])} \OperatorTok{\{}
            \ControlFlowTok{break}\OperatorTok{;}
        \OperatorTok{\}}
    \OperatorTok{\}}
\NormalTok{    jComboBoxNetworkAgents}\OperatorTok{.}\FunctionTok{setSelectedIndex}\OperatorTok{(}\NormalTok{i}\OperatorTok{);}
\OperatorTok{\}} \ControlFlowTok{else} \OperatorTok{\{}
\NormalTok{    jComboBoxNetworkAgents}\OperatorTok{.}\FunctionTok{setSelectedIndex}\OperatorTok{(}\DecValTok{0}\OperatorTok{);}
\OperatorTok{\}}
\end{Highlighting}
\end{Shaded}
\paragraph{Cómo arreglarlo}
Comparar con \texttt{name.\allowbreak{}equals(\allowbreak{}agent\allowbreak{}Names{[}i{]})} y decidir qué se marca si el nombre no está en la lista (lo natural es el hueco vacío, puesto 0, y no una excepción). Con eso el diálogo se abre siempre; el nombre huérfano sigue siendo un defecto de la red, que corrige \hyperref[err:61]{error~61} (renombrar debe renombrar también en los nodos).

Es el mismo patrón que \hyperref[nota:30]{nota~30} (comparar nombres con \texttt{==} en el editor de árboles), donde hoy no hay nada que corregir; conviene buscar otros \texttt{==} entre cadenas en \texttt{gui}.

Prueba de regresión: construir el selector para una decisión cuyo agente sea \texttt{new\ String\allowbreak{}With\allowbreak{}Properties(\allowbreak{}new\ String(\allowbreak{}"Agent\ 1"))} en una red con «Agent 1» en la lista y comprobar que queda marcado ese agente; y con un nombre que no está en la lista, que queda marcado el hueco vacío.

\setcounter{subsection}{62}
\subsection[Evaluar DAN-arthronet falla con NullPointerException al instanciar]{Evaluar DAN-arthronet (y cualquier red de análisis de decisiones con la misma forma) termina con NullPointerException al instanciar}\label{err:63}
\begin{ficha}
\fichakey{Estado}Presente
\fichakey{Módulo}inference
\fichakey{Paquete}org.openmarkov.inference.algorithm.decompositionIntoSymmetricDANs.core
\fichakey{Ficheros}\path{inference/src/main/java/org/openmarkov/inference/algorithm/decompositionIntoSymmetricDANs/core/DANOperations.java:71-79} \newline \path{inference/src/main/java/org/openmarkov/inference/algorithm/decompositionIntoSymmetricDANs/core/DANOperations.java:100-146}
\fichakey{Comprobado}Leyendo el código y ejecutándolo
\fichakey{Esfuerzo}Pequeño (horas)
\fichakey{Decisión de equipo}No
\fichakey{Relacionados}\hyperref[err:64]{error~64}, \hyperref[err:113]{error~113}
\end{ficha}
\paragraph{Qué pasa}
El usuario abre la red y pide la estrategia óptima o la evaluación (\texttt{Inference\allowbreak{}Handler} usa \texttt{DANDecomposition\allowbreak{}Into\allowbreak{}Symmetric\allowbreak{}DANs\allowbreak{}Evaluation} para las redes de análisis de decisiones), o el árbol de decisión (\texttt{DANDecision\allowbreak{}Tree\allowbreak{}Inference} y \texttt{Decision\allowbreak{}Tree\allowbreak{}Expansion}, que también llaman a \texttt{instantiate}). La red no se puede evaluar: sale un \texttt{Null\allowbreak{}Pointer\allowbreak{}Exception}.

Medido con \texttt{0\_\allowbreak{}miscellaneous/\allowbreak{}networks/\allowbreak{}dan/\allowbreak{}DAN-arthronet.\allowbreak{}pgmx}:

\begin{Verbatim}[breaklines,breakanywhere,fontsize=\small]
instanciar «Realizar Implante» = no:
  NullPointerException: Cannot invoke "Node.getNodeType()" because "destinationNode" is null
    at DANOperations.instantiate(DANOperations.java:79)
evaluación completa (DANDecompositionIntoSymmetricDANsEvaluation): la misma excepción
\end{Verbatim}

Por qué pasa: \texttt{DANOperations.\allowbreak{}instantiate} recorre los enlaces que salen de la variable instanciada en la red \textbf{original} y, para cada uno, busca el nodo de destino en la \textbf{copia} (\texttt{instantiated\allowbreak{}Net.\allowbreak{}get\allowbreak{}Node(\allowbreak{}.\allowbreak{}.\allowbreak{}.\allowbreak{})}, línea 75) y lo usa sin mirar si es nulo (\texttt{destination\allowbreak{}Node.\allowbreak{}get\allowbreak{}Node\allowbreak{}Type(\allowbreak{})}, línea 79; y más abajo \texttt{destination\allowbreak{}Node.\allowbreak{}get\allowbreak{}Variable(\allowbreak{})}, línea 101).

El propio bucle quita nodos de la copia: cuando un enlace con restricciones prohíbe todos los estados de su destino para el estado instanciado (líneas 106-146), se eliminan el destino y sus descendientes. Si otro enlace de la misma variable, que llega más tarde en el recorrido, apuntaba a uno de esos nodos ya eliminados, \texttt{get\allowbreak{}Node} devuelve nulo y la línea 79 lanza. Si el enlace eliminado llega antes o después depende del orden en que la red original devuelve sus enlaces.
\paragraph{El código}
inference/src/main/java/org/openmarkov/inference/algorithm/decompositionIntoSymmetricDANs/core/DANOperations.java:71-79

\begin{Shaded}
\begin{Highlighting}[]
        \ControlFlowTok{for} \OperatorTok{(}\NormalTok{Link}\OperatorTok{\textless{}}\BuiltInTok{Node}\OperatorTok{\textgreater{}}\NormalTok{ link }\OperatorTok{:}\NormalTok{ originalDAN}\OperatorTok{.}\FunctionTok{getLinks}\OperatorTok{(}\NormalTok{originalNode}\OperatorTok{))} \OperatorTok{\{}
            \ControlFlowTok{if} \OperatorTok{(!}\NormalTok{link}\OperatorTok{.}\FunctionTok{getFrom}\OperatorTok{().}\FunctionTok{equals}\OperatorTok{(}\NormalTok{originalNode}\OperatorTok{))} \OperatorTok{\{}
                \ControlFlowTok{continue}\OperatorTok{;}
            \OperatorTok{\}}
            \BuiltInTok{Node}\NormalTok{ destinationNode }\OperatorTok{=}\NormalTok{ instantiatedNet}\OperatorTok{.}\FunctionTok{getNode}\OperatorTok{(}\NormalTok{link}\OperatorTok{.}\FunctionTok{getTo}\OperatorTok{().}\FunctionTok{getVariable}\OperatorTok{());}
            \CommentTok{// Remove link between restricting originalNode and restricted originalNode}
\NormalTok{            instantiatedNet}\OperatorTok{.}\FunctionTok{removeLink}\OperatorTok{(}\NormalTok{link}\OperatorTok{.}\FunctionTok{getFrom}\OperatorTok{().}\FunctionTok{getVariable}\OperatorTok{(),}\NormalTok{ link}\OperatorTok{.}\FunctionTok{getTo}\OperatorTok{().}\FunctionTok{getVariable}\OperatorTok{(),} \KeywordTok{true}\OperatorTok{);}
            
            \ControlFlowTok{if} \OperatorTok{(}\NormalTok{destinationNode}\OperatorTok{.}\FunctionTok{getNodeType}\OperatorTok{()} \OperatorTok{==}\NormalTok{ NodeType}\OperatorTok{.}\FunctionTok{CHANCE}\OperatorTok{)} \OperatorTok{\{}
\end{Highlighting}
\end{Shaded}
\paragraph{Cómo arreglarlo}
Cuando \texttt{destination\allowbreak{}Node} es nulo, pasar al enlace siguiente (\texttt{continue}): el nodo ya se quitó porque en esta rama no cabe, y no queda nada que revelar ni restringir en él. Conviene hacerlo antes de \texttt{remove\allowbreak{}Link}, que no falla con nodos ausentes pero no tiene sentido.

Revisar a la vez los demás bucles del mismo método que usan nodos de la copia después de haber podido quitarlos: la cascada de borrado (líneas 110-144) y \texttt{prune\allowbreak{}Zero\allowbreak{}Probability\allowbreak{}Chance\allowbreak{}Nodes}.

Prueba de regresión: evaluar \texttt{0\_\allowbreak{}miscellaneous/\allowbreak{}networks/\allowbreak{}dan/\allowbreak{}DAN-arthronet.\allowbreak{}pgmx} con \texttt{DANDecomposition\allowbreak{}Into\allowbreak{}Symmetric\allowbreak{}DANs\allowbreak{}Evaluation} sin excepción, e instanciar «Realizar Implante» = no.

\setcounter{subsection}{63}
\subsection[Instanciar una variable de una red de decisiones duplica un enlace]{Al instanciar una variable de una red de análisis de decisiones se crea un segundo enlace entre dos nodos ya enlazados, y la red resultante incumple su propia restricción de enlaces distintos}\label{err:64}
\begin{ficha}
\fichakey{Estado}Presente
\fichakey{Módulo}inference
\fichakey{Paquete}org.openmarkov.inference.algorithm.decompositionIntoSymmetricDANs.core
\fichakey{Ficheros}\path{inference/src/main/java/org/openmarkov/inference/algorithm/decompositionIntoSymmetricDANs/core/DANOperations.java:79-97} \newline \path{core/src/main/java/org/openmarkov/core/model/graph/Graph.java:186-205} \newline \path{core/src/main/java/org/openmarkov/core/model/network/ProbNetOperations.java:815-835} \newline \path{core/src/main/java/org/openmarkov/core/inference/BasicOperations.java:413-425}
\fichakey{Comprobado}Leyendo el código y ejecutándolo
\fichakey{Esfuerzo}Pequeño (horas)
\fichakey{Decisión de equipo}No
\fichakey{Relacionados}\hyperref[err:63]{error~63}
\end{ficha}
\paragraph{Qué pasa}
Toda evaluación de una red de análisis de decisiones (DAN, del inglés \emph{decision analysis network}) por descomposición (\texttt{DANDecomposition\allowbreak{}Into\allowbreak{}Symmetric\allowbreak{}DANs\allowbreak{}Evaluation}, que usa la ventana para «Mostrar estrategia óptima» en \texttt{Inference\allowbreak{}Handler}) y la construcción del árbol de decisión (\texttt{DANDecision\allowbreak{}Tree\allowbreak{}Inference}, \texttt{Decision\allowbreak{}Tree\allowbreak{}Expansion}) pasan por \texttt{DANOperations.\allowbreak{}instantiate}. En algunas redes, la red intermedia que devuelve queda con dos enlaces entre el mismo par de nodos e incumple una restricción de su propio tipo.

Medido instanciando cada variable en cada estado en las 34 redes que se pueden leer de \texttt{0\_\allowbreak{}miscellaneous/\allowbreak{}networks/\allowbreak{}dan}:

\begin{Verbatim}[breaklines,breakanywhere,fontsize=\small]
DAN-delayed-result-of-test.pgmx -> Therapy today=no:  Do test? -> Result of test (2 enlaces)
                                   Do test?=yes:      Therapy today -> Result of test (2 enlaces)
DAN-king.pgmx                   -> War=no / War=yes:  Marriage -> Wealth (2 enlaces)
DAN-used-car-buyer.pgmx         -> Dec: Second Test=differential: Dec: First Test -> Second result (2 enlaces)
\end{Verbatim}

Las tres redes tienen la restricción \texttt{Distinct\allowbreak{}Links}, que la copia hereda, así que la red intermedia incumple una restricción de su propio tipo.

Qué puede cambiar: quienes leen si un nodo queda revelado toman el primer enlace que encuentran entre los dos nodos. \texttt{Prob\allowbreak{}Net.\allowbreak{}get\allowbreak{}Link} (usado en \texttt{Basic\allowbreak{}Operations.\allowbreak{}get\allowbreak{}Variables\allowbreak{}Revealed\allowbreak{}Transitively\allowbreak{}By\allowbreak{}Variable}, que da el orden de eliminación de una DAN) y el bucle de \texttt{Prob\allowbreak{}Net\allowbreak{}Operations.\allowbreak{}get\allowbreak{}Observable\allowbreak{}Variables} se quedan con el primero, que es el original; el revelador añadido puede no contar, o contar con las restricciones del otro enlace. En \texttt{DAN-delayed-result-of-test.\allowbreak{}pgmx} instanciando «Therapy today=no», el primer enlace es el original (revelador solo para «yes» y con restricciones) y el añadido revela para «no» y «yes».

No se ha podido medir si esto cambia la utilidad o la estrategia final: dos de las tres redes no se evalúan enteras por otros motivos (\texttt{Potential\allowbreak{}Cannot\allowbreak{}Be\allowbreak{}Converted\allowbreak{}To\allowbreak{}ATable}), y la tercera (\texttt{DAN-used-car-buyer}, utilidad 32,96) no tiene con qué compararse sin modificar el código.

Por qué pasa. Una DAN se evalúa descomponiéndola: \texttt{DANOperations.\allowbreak{}instantiate} fija una variable en uno de sus estados y devuelve una copia de la red simplificada para esa rama. Cuando un enlace sale de la variable instanciada con estados reveladores (el destino pasa a conocerse si la variable toma ese estado) y el destino tiene decisiones predecesoras todavía por tomar, el código no lo marca como observado, sino que traza un enlace revelador desde cada una de esas decisiones hasta el destino (líneas 89-95), con todos los estados de la decisión como reveladores.

\texttt{Prob\allowbreak{}Net.\allowbreak{}add\allowbreak{}Link} delega en \texttt{Graph.\allowbreak{}add\allowbreak{}Link}, que por diseño crea siempre un enlace nuevo (su javadoc dice que impedir duplicados es trabajo de la restricción \texttt{Distinct\allowbreak{}Links}). Si la decisión predecesora ya era padre directo del destino ---lo corriente: la decisión que revela algo suele además influir en ello---, la copia queda con dos enlaces entre el mismo par de nodos: el original (sin esos estados reveladores, y a veces con restricciones) y el nuevo.
\paragraph{El código}
inference/src/main/java/org/openmarkov/inference/algorithm/decompositionIntoSymmetricDANs/core/DANOperations.java:82-95

\begin{Shaded}
\begin{Highlighting}[]
                        \BuiltInTok{List}\OperatorTok{\textless{}}\BuiltInTok{Node}\OperatorTok{\textgreater{}}\NormalTok{ predecessorDecisions }\OperatorTok{=}\NormalTok{ ProbNetOperations}
                                \OperatorTok{.}\FunctionTok{getPredecessorDecisions}\OperatorTok{(}\NormalTok{destinationNode}\OperatorTok{,}\NormalTok{ instantiatedNet}\OperatorTok{);}
                        \CommentTok{// If it has predecessor decisions, do not reveal it yet, but add revealing links}
                        \CommentTok{// from every predecessor decision to the originalNode}
                        \ControlFlowTok{if} \OperatorTok{(}\NormalTok{predecessorDecisions}\OperatorTok{.}\FunctionTok{isEmpty}\OperatorTok{())} \OperatorTok{\{}
\NormalTok{                            destinationNode}\OperatorTok{.}\FunctionTok{setAlwaysObserved}\OperatorTok{(}\KeywordTok{true}\OperatorTok{);}
                        \OperatorTok{\}} \ControlFlowTok{else} \OperatorTok{\{}
                            \ControlFlowTok{for} \OperatorTok{(}\BuiltInTok{Node}\NormalTok{ predecessorDecision }\OperatorTok{:}\NormalTok{ predecessorDecisions}\OperatorTok{)} \OperatorTok{\{}
\NormalTok{                                Link}\OperatorTok{\textless{}}\BuiltInTok{Node}\OperatorTok{\textgreater{}}\NormalTok{ revealingArc }\OperatorTok{=}\NormalTok{ instantiatedNet}
                                        \OperatorTok{.}\FunctionTok{addLink}\OperatorTok{(}\NormalTok{predecessorDecision}\OperatorTok{,}\NormalTok{ destinationNode}\OperatorTok{,} \KeywordTok{true}\OperatorTok{);}
                                \BuiltInTok{State}\OperatorTok{[]}\NormalTok{ predecessorDecisionStates }\OperatorTok{=}\NormalTok{ predecessorDecision}\OperatorTok{.}\FunctionTok{getVariable}\OperatorTok{().}\FunctionTok{getStates}\OperatorTok{();}
                                \ControlFlowTok{for} \OperatorTok{(}\BuiltInTok{State}\NormalTok{ predecessorDecisionState }\OperatorTok{:}\NormalTok{ predecessorDecisionStates}\OperatorTok{)}
\NormalTok{                                    revealingArc}\OperatorTok{.}\FunctionTok{addRevealingState}\OperatorTok{(}\NormalTok{predecessorDecisionState}\OperatorTok{);}
                            \OperatorTok{\}}
                        \OperatorTok{\}}
\end{Highlighting}
\end{Shaded}
\paragraph{Cómo arreglarlo}
Antes de añadir, buscar con \texttt{instantiated\allowbreak{}Net.\allowbreak{}get\allowbreak{}Link(\allowbreak{}predecessor\allowbreak{}Decision,\allowbreak{}\ destination\allowbreak{}Node,\allowbreak{}\ true)}: si ya existe, añadirle los estados reveladores que le falten en lugar de crear otro enlace; si no, crearlo como hoy. Hay que decidir qué hacer si el enlace existente tiene restricciones; lo natural es conservarlas, porque siguen siendo ciertas.

Otros sitios que deben respetar la misma regla:

\begin{itemize}
\tightlist
\item
  los que crean enlaces en redes intermedias de la inferencia: todos los \texttt{add\allowbreak{}Link} de \texttt{DANOperations} y de \texttt{Decision\allowbreak{}Tree\allowbreak{}Expansion};
\item
  los que leen suponiendo un único enlace por par: \texttt{Prob\allowbreak{}Net.\allowbreak{}get\allowbreak{}Link}, \texttt{Prob\allowbreak{}Net\allowbreak{}Operations.\allowbreak{}get\allowbreak{}Observable\allowbreak{}Variables}, \texttt{Basic\allowbreak{}Operations.\allowbreak{}get\allowbreak{}Variables\allowbreak{}Revealed\allowbreak{}Transitively\allowbreak{}By\allowbreak{}Variable}.
\end{itemize}

Prueba de regresión: instanciar «Therapy today=no» en \texttt{0\_\allowbreak{}miscellaneous/\allowbreak{}networks/\allowbreak{}dan/\allowbreak{}DAN-delayed-result-of-test.\allowbreak{}pgmx} y comprobar que hay un solo enlace «Do test?» → «Result of test», que revela para los dos estados, y que la copia cumple \texttt{check\allowbreak{}Constraints(\allowbreak{})}.

\setcounter{subsection}{64}
\subsection[La estrategia óptima de una red sin decisiones lanza una excepción]{Pedir la estrategia óptima de un diagrama de influencia o una red de análisis de decisiones sin decisiones lanza una excepción inesperada}\label{err:65}
\begin{ficha}
\fichakey{Estado}Presente
\fichakey{Módulo}gui
\fichakey{Paquete}org.openmarkov.gui.window
\fichakey{Ficheros}\path{gui/src/main/java/org/openmarkov/gui/window/MainPanelMenuAssistant.java:373-378,393-398} \newline \path{gui/src/main/java/org/openmarkov/gui/window/InferenceHandler.java:207-243} \newline \path{inference/src/main/java/org/openmarkov/inference/algorithm/variableElimination/tasks/VEOptimalIntervention.java:64-90} \newline \path{gui/src/main/java/org/openmarkov/gui/dialog/network/OptimalStrategyDialog.java:31-37} \newline \path{gui/src/main/java/org/openmarkov/gui/menutoolbar/toolbar/StandardToolBar.java:260-263}
\fichakey{Comprobado}Leyendo el código y ejecutándolo
\fichakey{Esfuerzo}Pequeño (horas)
\fichakey{Decisión de equipo}No
\fichakey{Relacionados}\hyperref[err:66]{error~66}, \hyperref[err:92]{error~92}
\end{ficha}
\paragraph{Qué pasa}
El botón de la barra de herramientas «Show optimal strategy» (\texttt{DECISION\_\allowbreak{}SHOW\_\allowbreak{}OPTIMAL\_\allowbreak{}STRATEGY}) está activo en los diagramas de influencia, las redes de análisis de decisiones, los MID (diagramas de influencia de Markov) y demás tipos con decisiones aunque la red no tenga ninguna decisión, incluido un diagrama de influencia recién creado y vacío. Al pulsarlo, el gestor global muestra el diálogo de «excepción inesperada» pidiendo que se envíe a los desarrolladores; el programa no se cierra. En las redes del repositorio no hay diagramas de influencia ni redes de análisis de decisiones sin decisiones, pero cualquier modelo a medio construir lo es.

Medido con una red de prueba (un nodo de azar «Disease» con dos estados y un nodo de utilidad «U» hijo suyo, escrita como diagrama de influencia y como red de análisis de decisiones): el diagrama de influencia lanza \texttt{Null\allowbreak{}Pointer\allowbreak{}Exception:\ Cannot\ load\ from\ object\ array\ because\ the\ return\ value\ of\ "Strategic\allowbreak{}Table\allowbreak{}Potential.\allowbreak{}get\allowbreak{}Strategy\allowbreak{}Trees(\allowbreak{})"\ is\ null}; un diagrama de influencia vacío construido en código lanza la misma; la red de análisis de decisiones lanza \texttt{Class\allowbreak{}Cast\allowbreak{}Exception:\ Table\allowbreak{}Potential\ cannot\ be\ cast\ to\ Strategy\allowbreak{}Carrier}. Las dos son excepciones de ejecución.

Por qué pasa. El botón se activa en \texttt{Main\allowbreak{}Panel\allowbreak{}Menu\allowbreak{}Assistant}: en modo edición con \texttt{!current\allowbreak{}Prob\allowbreak{}Net.\allowbreak{}has\allowbreak{}Constraint\allowbreak{}Of\allowbreak{}Class(\allowbreak{}Only\allowbreak{}Chance\allowbreak{}Nodes.\allowbreak{}class)} y en modo inferencia con la misma condición. Esos tipos de red no llevan esa restricción, así que el botón queda activo sin mirar si hay decisiones.

Al pulsarlo, \texttt{Inference\allowbreak{}Handler.\allowbreak{}show\allowbreak{}Optimal\allowbreak{}Strategy} sigue uno de dos caminos:

\begin{itemize}
\tightlist
\item
  Red de análisis de decisiones: \texttt{DANDecomposition\allowbreak{}Into\allowbreak{}Symmetric\allowbreak{}DANs\allowbreak{}Evaluation(\allowbreak{}.\allowbreak{}.\allowbreak{}.\allowbreak{}).\allowbreak{}get\allowbreak{}Utility(\allowbreak{})} se convierte a \texttt{Strategy\allowbreak{}Carrier} y se toma \texttt{get\allowbreak{}Strategy\allowbreak{}Trees(\allowbreak{}){[}0{]}} (línea 218). Sin decisiones la utilidad es un \texttt{Table\allowbreak{}Potential} normal y la conversión lanza \texttt{Class\allowbreak{}Cast\allowbreak{}Exception}.
\item
  Resto (diagramas de influencia, MID\ldots): \texttt{new\ Optimal\allowbreak{}Strategy\allowbreak{}Dialog(\allowbreak{}.\allowbreak{}.\allowbreak{}.\allowbreak{},\allowbreak{}\ new\ VEOptimal\allowbreak{}Intervention(\allowbreak{}.\allowbreak{}.\allowbreak{}.\allowbreak{}))}, cuyo constructor llama a \texttt{get\allowbreak{}Optimal\allowbreak{}Intervention(\allowbreak{})}. Si la utilidad no trae árboles de estrategia, el método intenta construir uno a partir de \texttt{get\allowbreak{}An\allowbreak{}Admissible\allowbreak{}Order\allowbreak{}Of\allowbreak{}Decisions}, pero ese relleno está dentro de \texttt{if\ (\allowbreak{}!decisions.\allowbreak{}is\allowbreak{}Empty(\allowbreak{}))}; sin decisiones no hace nada y la última línea, \texttt{strategic.\allowbreak{}get\allowbreak{}Strategy\allowbreak{}Trees(\allowbreak{}){[}0{]}}, lanza \texttt{Null\allowbreak{}Pointer\allowbreak{}Exception}.
\end{itemize}
\paragraph{El código}
inference/src/main/java/org/openmarkov/inference/algorithm/variableElimination/tasks/VEOptimalIntervention.java:74-89

\begin{Shaded}
\begin{Highlighting}[]
\NormalTok{StrategyTree}\OperatorTok{[]}\NormalTok{ strategyTrees }\OperatorTok{=}\NormalTok{ strategic}\OperatorTok{.}\FunctionTok{getStrategyTrees}\OperatorTok{();}
\CommentTok{// Create strategy tree if it is empty}
\ControlFlowTok{if} \OperatorTok{((}\NormalTok{strategyTrees}\OperatorTok{==}\KeywordTok{null}\OperatorTok{)||(}\NormalTok{strategyTrees}\OperatorTok{.}\FunctionTok{length}\OperatorTok{==}\DecValTok{0}\OperatorTok{))} \OperatorTok{\{}
    \BuiltInTok{List}\OperatorTok{\textless{}}\NormalTok{Variable}\OperatorTok{\textgreater{}}\NormalTok{ decisions }\OperatorTok{=}\NormalTok{ BasicOperations}\OperatorTok{.}\FunctionTok{getAnAdmissibleOrderOfDecisions}\OperatorTok{(}\NormalTok{probNet}\OperatorTok{);}
    \ControlFlowTok{if} \OperatorTok{(!}\NormalTok{decisions}\OperatorTok{.}\FunctionTok{isEmpty}\OperatorTok{())} \OperatorTok{\{}
\NormalTok{        Variable lastVar }\OperatorTok{=}\NormalTok{ decisions}\OperatorTok{.}\FunctionTok{getLast}\OperatorTok{();}
\NormalTok{        StrategyTree auxST }\OperatorTok{=} \KeywordTok{new} \FunctionTok{StrategyTree}\OperatorTok{(}\NormalTok{lastVar}\OperatorTok{,}\NormalTok{lastVar}\OperatorTok{.}\FunctionTok{getStates}\OperatorTok{());}
        \KeywordTok{...}
\NormalTok{        strategic}\OperatorTok{.}\FunctionTok{strategyTrees} \OperatorTok{=} \KeywordTok{new}\NormalTok{ StrategyTree}\OperatorTok{[}\DecValTok{1}\OperatorTok{];}
\NormalTok{        strategic}\OperatorTok{.}\FunctionTok{strategyTrees}\OperatorTok{[}\DecValTok{0}\OperatorTok{]} \OperatorTok{=}\NormalTok{ auxST}\OperatorTok{;}
    \OperatorTok{\}}
\OperatorTok{\}}
\ControlFlowTok{return}\NormalTok{ strategic}\OperatorTok{.}\FunctionTok{getStrategyTrees}\OperatorTok{()[}\DecValTok{0}\OperatorTok{];}
\end{Highlighting}
\end{Shaded}

gui/src/main/java/org/openmarkov/gui/window/MainPanelMenuAssistant.java:373-378

\begin{Shaded}
\begin{Highlighting}[]
\DataTypeTok{boolean}\NormalTok{ isOnlyChanceNodes }\OperatorTok{=}\NormalTok{ currentProbNet}\OperatorTok{.}\FunctionTok{hasConstraintOfClass}\OperatorTok{(}\NormalTok{OnlyChanceNodes}\OperatorTok{.}\FunctionTok{class}\OperatorTok{);}
\FunctionTok{setOptionEnabled}\OperatorTok{(}\NormalTok{ActionCommands}\OperatorTok{.}\FunctionTok{DECISION\_CREATION}\OperatorTok{,} \OperatorTok{!}\NormalTok{isOnlyChanceNodes}\OperatorTok{);}
\FunctionTok{setOptionEnabled}\OperatorTok{(}\NormalTok{ActionCommands}\OperatorTok{.}\FunctionTok{UTILITY\_CREATION}\OperatorTok{,} \OperatorTok{!}\NormalTok{isOnlyChanceNodes}\OperatorTok{);}
\KeywordTok{...}
\FunctionTok{setOptionEnabled}\OperatorTok{(}\NormalTok{ActionCommands}\OperatorTok{.}\FunctionTok{DECISION\_SHOW\_OPTIMAL\_STRATEGY}\OperatorTok{,} \OperatorTok{!}\NormalTok{isOnlyChanceNodes}\OperatorTok{);}
\end{Highlighting}
\end{Shaded}
\paragraph{Cómo arreglarlo}
Hay que fijar el comportamiento en dos sitios coherentes.

En la interfaz: activar el botón solo si la red tiene al menos un nodo de decisión, en los dos modos (\texttt{Main\allowbreak{}Panel\allowbreak{}Menu\allowbreak{}Assistant}, ramas \texttt{EDITION} e \texttt{INFERENCE}), y comprobar que se recalcula cuando se añade o se quita la última decisión.

En el cálculo, para que un llamador que no pase por el botón no reciba una excepción de ejecución: \texttt{VEOptimal\allowbreak{}Intervention.\allowbreak{}get\allowbreak{}Optimal\allowbreak{}Intervention} debería lanzar una excepción declarada con mensaje claro («la red no tiene decisiones») o devolver un resultado vacío documentado; y \texttt{Inference\allowbreak{}Handler.\allowbreak{}show\allowbreak{}Optimal\allowbreak{}Strategy} debería comprobar el tipo antes de convertir a \texttt{Strategy\allowbreak{}Carrier} en la rama de las redes de análisis de decisiones.

El botón del árbol de decisión tiene una regla de activación parecida (\hyperref[err:66]{error~66}).

Prueba de regresión: sobre las redes de prueba sin decisiones (diagrama de influencia y red de análisis de decisiones), comprobar que el botón queda desactivado tras cargar la red en los dos modos, y que \texttt{VEOptimal\allowbreak{}Intervention.\allowbreak{}get\allowbreak{}Optimal\allowbreak{}Intervention} falla con la excepción declarada y no con \texttt{Null\allowbreak{}Pointer\allowbreak{}Exception}.

\setcounter{subsection}{65}
\subsection[El botón del árbol de decisión se enciende en redes que no admiten árbol]{El botón del árbol de decisión se enciende en modo inferencia para redes que no admiten árbol (MID, POMDP, LIMID\ldots), y al pulsarlo sale un error}\label{err:66}
\begin{ficha}
\fichakey{Estado}Presente
\fichakey{Módulo}gui
\fichakey{Paquete}org.openmarkov.gui.window
\fichakey{Ficheros}\path{gui/src/main/java/org/openmarkov/gui/window/MainPanelMenuAssistant.java:365-398} \newline \path{inference/src/main/java/org/openmarkov/inference/decisiontree/operation/DecisionTreeManagerImpl.java:55-74} \newline \path{gui/src/main/java/org/openmarkov/gui/menutoolbar/toolbar/StandardToolBar.java:275-287}
\fichakey{Comprobado}Leyendo el código y ejecutándolo
\fichakey{Esfuerzo}Pequeño (horas)
\fichakey{Decisión de equipo}No
\fichakey{Relacionados}\hyperref[err:65]{error~65}
\end{ficha}
\paragraph{Qué pasa}
Un MID (diagrama de influencia de Markov), un POMDP (proceso de decisión de Markov parcialmente observable), un LIMID (diagrama de influencia con memoria limitada) o una red de sucesos discretos tienen el botón «Decision tree» apagado mientras se editan y encendido mientras se hace inferencia. Al pulsarlo en inferencia, el usuario recibe un error que dice que la opción no vale para esa red.

Medido sobre \texttt{0\_\allowbreak{}miscellaneous/\allowbreak{}networks/\allowbreak{}mid/\allowbreak{}MID-Chancellor-corrected.\allowbreak{}pgmx}: \texttt{Decision\allowbreak{}Tree\allowbreak{}Manager\allowbreak{}Impl.\allowbreak{}build\allowbreak{}Decision\allowbreak{}Tree} rechaza la red con \texttt{Not\allowbreak{}Evaluable\allowbreak{}Network\allowbreak{}Exception.\allowbreak{}Not\allowbreak{}Applicable\allowbreak{}Network} y el mensaje «This algorithm cannot evaluate Markov Influence Diagram network named MID-Chancellor-corrected.pgmx, as it can only evaluate networks of type: {[}Influence diagram, Decision analysis{]}». En la aplicación esa excepción comprobada llega envuelta en \texttt{Unrecoverable\allowbreak{}Exception} al manejador global, que la enseña en la ventana de error.

Por qué pasa. El árbol de decisión solo se ofrece desde el botón «Decision tree» de la barra de herramientas estándar (\texttt{Standard\allowbreak{}Tool\allowbreak{}Bar.\allowbreak{}get\allowbreak{}Decision\allowbreak{}Tree\allowbreak{}Button}). \texttt{Main\allowbreak{}Panel\allowbreak{}Menu\allowbreak{}Assistant} decide si ese botón está activo, y lo decide con dos reglas distintas según el modo de trabajo:

\begin{itemize}
\tightlist
\item
  En modo edición (línea 379): activo si la red no tiene la restricción \texttt{Only\allowbreak{}Chance\allowbreak{}Nodes} \textbf{y} es un diagrama de influencia (\texttt{Influence\allowbreak{}Diagram\allowbreak{}Type}) o una red de análisis de decisiones (\texttt{Decision\allowbreak{}Analysis\allowbreak{}Network\allowbreak{}Type}).
\item
  En modo inferencia (líneas 395-397): activo si la red no tiene la restricción \texttt{Only\allowbreak{}Chance\allowbreak{}Nodes}, sin mirar el tipo.
\end{itemize}

Ninguna de las dos reglas mira si hay decisiones, pero para el árbol eso no produce fallo: con diagramas de influencia y redes de análisis de decisiones sin ninguna decisión (\texttt{integration\allowbreak{}Tests/\allowbreak{}src/\allowbreak{}main/\allowbreak{}resources/\allowbreak{}networks/\allowbreak{}id/\allowbreak{}ID-only-utility.\allowbreak{}pgmx}, \texttt{integration\allowbreak{}Tests/\allowbreak{}src/\allowbreak{}main/\allowbreak{}resources/\allowbreak{}networks/\allowbreak{}id/\allowbreak{}ID-one-chance.\allowbreak{}pgmx}, \texttt{integration\allowbreak{}Tests/\allowbreak{}src/\allowbreak{}main/\allowbreak{}resources/\allowbreak{}networks/\allowbreak{}dan/\allowbreak{}DAN-one-chance.\allowbreak{}pgmx}) el árbol se construye (un árbol con solo nodos de azar).
\paragraph{El código}
gui/src/main/java/org/openmarkov/gui/window/MainPanelMenuAssistant.java:377-380, 391-397

\begin{Shaded}
\begin{Highlighting}[]
\FunctionTok{setOptionEnabled}\OperatorTok{(}\NormalTok{ActionCommands}\OperatorTok{.}\FunctionTok{DECISION\_SHOW\_OPTIMAL\_STRATEGY}\OperatorTok{,} \OperatorTok{!}\NormalTok{isOnlyChanceNodes}\OperatorTok{);}
\FunctionTok{setOptionEnabled}\OperatorTok{(}\NormalTok{ActionCommands}\OperatorTok{.}\FunctionTok{DECISION\_TREE}\OperatorTok{,} \OperatorTok{!}\NormalTok{isOnlyChanceNodes }\OperatorTok{\&\&} \OperatorTok{((}\NormalTok{currentProbNet}\OperatorTok{.}\FunctionTok{getNetworkType}\OperatorTok{()} \KeywordTok{instanceof}\NormalTok{ InfluenceDiagramType}\OperatorTok{)}
        \OperatorTok{||} \OperatorTok{(}\NormalTok{currentProbNet}\OperatorTok{.}\FunctionTok{getNetworkType}\OperatorTok{()} \KeywordTok{instanceof}\NormalTok{ DecisionAnalysisNetworkType}\OperatorTok{)));}
\KeywordTok{...}
\ControlFlowTok{case}\NormalTok{ INFERENCE }\OperatorTok{{-}\textgreater{}} \OperatorTok{\{}
    \KeywordTok{...}
    \ControlFlowTok{if} \OperatorTok{(!}\NormalTok{currentNetworkEditorPanel}\OperatorTok{.}\FunctionTok{getProbNet}\OperatorTok{().}\FunctionTok{hasConstraintOfClass}\OperatorTok{(}\NormalTok{OnlyChanceNodes}\OperatorTok{.}\FunctionTok{class}\OperatorTok{))} \OperatorTok{\{}
        \FunctionTok{setOptionEnabled}\OperatorTok{(}\NormalTok{ActionCommands}\OperatorTok{.}\FunctionTok{DECISION\_TREE}\OperatorTok{,} \KeywordTok{true}\OperatorTok{);}
        \FunctionTok{setOptionEnabled}\OperatorTok{(}\NormalTok{ActionCommands}\OperatorTok{.}\FunctionTok{DECISION\_SHOW\_OPTIMAL\_STRATEGY}\OperatorTok{,} \KeywordTok{true}\OperatorTok{);}
    \OperatorTok{\}}
\OperatorTok{\}}
\end{Highlighting}
\end{Shaded}
\paragraph{Cómo arreglarlo}
Usar una sola regla para los dos modos, la misma que aplica quien construye el árbol: tipo de red diagrama de influencia o red de análisis de decisiones (y sin \texttt{Only\allowbreak{}Chance\allowbreak{}Nodes}). Lo más robusto es sacar esa condición a un único método (por ejemplo en \texttt{Decision\allowbreak{}Tree\allowbreak{}Manager}, «¿admite esta red un árbol?») y usarlo en \texttt{Main\allowbreak{}Panel\allowbreak{}Menu\allowbreak{}Assistant} en las dos ramas y en \texttt{Decision\allowbreak{}Tree\allowbreak{}Manager\allowbreak{}Impl.\allowbreak{}build\allowbreak{}Decision\allowbreak{}Tree}, para que no vuelvan a divergir.

El botón de estrategia óptima de las mismas líneas tiene su propio problema (\hyperref[err:65]{error~65}), que no se arregla con esto.

Prueba: con una red MID abierta en modo inferencia, comprobar que la opción \texttt{DECISION\_\allowbreak{}TREE} queda desactivada; con un diagrama de influencia, activada en los dos modos.

% ══════════════════════════════════════════════════════════════════════
\section{Modelos temporales y evolución temporal}\label{sec:temporal}

Las redes con variables que se repiten ciclo a ciclo (por ejemplo, los diagramas de influencia de Markov) y la ventana de evolución temporal.

\setcounter{subsection}{66}
\subsection[Con evidencia, las curvas de la evolución temporal no suman uno]{Con evidencia, las curvas de la evolución temporal no suman uno: se muestran probabilidades conjuntas en vez de condicionadas}\label{err:67}
\begin{ficha}
\fichakey{Estado}Presente
\fichakey{Módulo}inference
\fichakey{Paquete}org.openmarkov.inference.algorithm.temporalevaluation.tasks; org.openmarkov.inference.algorithm.variableElimination.tasks
\fichakey{Ficheros}\path{inference/src/main/java/org/openmarkov/inference/algorithm/temporalevaluation/tasks/MIDTemporalEvolution.java:273-283} \newline \path{inference/src/main/java/org/openmarkov/inference/algorithm/temporalevaluation/tasks/MIDTemporalEvolution.java:291-295} \newline \path{inference/src/main/java/org/openmarkov/inference/algorithm/temporalevaluation/tasks/MIDTemporalEvolution.java:602-619} \newline \path{inference/src/main/java/org/openmarkov/inference/algorithm/variableElimination/tasks/VEPropagation.java:210-213} \newline \path{core/src/main/java/org/openmarkov/core/inference/InferenceAlgorithm.java:200-204} \newline \path{gui/src/main/java/org/openmarkov/gui/dialog/inference/temporalevolution/TraceTemporalEvolutionDialog.java:315-334}
\fichakey{Comprobado}Leyendo el código y ejecutándolo
\fichakey{Esfuerzo}Medio (días)
\fichakey{Decisión de equipo}\textcolor{decisionrojo}{\textbf{Sí.}} Si las curvas de la evolución temporal deben normalizarse por cada opción de la decisión (y qué significa una columna que suma cero, es decir, una alternativa imposible con esa evidencia), o si basta con decir en la ventana que se muestran valores sin normalizar
\fichakey{Tapado por}\hyperref[err:74]{error~74}: en el jar publicado; desde el IDE se ve.
\fichakey{Relacionados}\hyperref[err:48]{error~48}, \hyperref[err:52]{error~52}, \hyperref[err:69]{error~69}, \hyperref[err:74]{error~74}
\end{ficha}
\paragraph{Qué pasa}
En modo inferencia de un modelo de Markov con decisiones (MID, diagrama de influencia de Markov), la evolución temporal de un nodo de azar temporal dibuja, ciclo a ciclo, la probabilidad de cada estado de la variable para cada opción de la decisión elegida. Para llegar al fallo basta con tener hallazgos previos a la resolución (puestos en modo edición) o elegir un escenario en el diálogo de la evolución temporal (\texttt{Temporal\allowbreak{}Evolution\allowbreak{}Dialog} añade esos hallazgos a la evidencia), seleccionar un nodo de azar temporal y pedir su evolución temporal.

Con evidencia, la gráfica y la tabla muestran la probabilidad conjunta P(V, e) y no la condicionada P(V \textbar{} e). Son números que no suman uno, presentados como probabilidades y escalados de forma distinta para cada opción de la decisión, lo que distorsiona la comparación entre ellas. Sin evidencia no se nota, porque la tabla ya suma uno por columna.

Al final de \texttt{MIDTemporal\allowbreak{}Evolution.\allowbreak{}compute\allowbreak{}Posterior\allowbreak{}Probability} (líneas 612-615) la normalización está protegida por \texttt{if\ (\allowbreak{}get\allowbreak{}Conditioning\allowbreak{}Variables(\allowbreak{})\ ==\ null\ \textbar{}\textbar{}\ get\allowbreak{}Conditioning\allowbreak{}Variables(\allowbreak{}).\allowbreak{}is\allowbreak{}Empty(\allowbreak{}))}, con un «TODO - Realizar la normalización condicionada» encima.

Esa lista no está vacía nunca: \texttt{common\allowbreak{}Preprocessing} (línea 294) la fija a \texttt{Collections.\allowbreak{}singleton\allowbreak{}List(\allowbreak{}conditioning\allowbreak{}Decision)}, y \texttt{Inference\allowbreak{}Algorithm.\allowbreak{}set\allowbreak{}Conditioning\allowbreak{}Variables} guarda cualquier lista no nula. Con decisión, es una lista de un elemento; sin decisión (\texttt{conditioning\allowbreak{}Decision} nulo), es una lista de un elemento que es nulo. En los dos casos la condición es falsa y la normalización no se ejecuta.

La misma lista tiene dos consecuencias más:

\begin{itemize}
\tightlist
\item
  El brazo \texttt{Cannot\allowbreak{}Normalize\allowbreak{}Potential\allowbreak{}Exception} del \texttt{catch} de \texttt{resolve(\allowbreak{})} (líneas 278-281, declarado inalcanzable con un «TODO: Maybe this can actually happen») no puede saltar mientras la normalización no se ejecute: con evidencia imposible, las curvas saldrían todas a cero sin aviso. Si se implementa la normalización, ese \texttt{catch} convierte la excepción en \texttt{Unreachable\allowbreak{}Exception}.
\item
  En la línea 609 se construye \texttt{Arrays.\allowbreak{}as\allowbreak{}List(\allowbreak{}probability\allowbreak{}Variable,\allowbreak{}\ conditioning\allowbreak{}Decision)} para reordenar; sin decisión, se reordena sobre una lista con un nulo dentro. No se ha visto dispararse en las pruebas hechas.
\end{itemize}

La misma guarda está copiada en \texttt{VEPropagation.\allowbreak{}Invoke\allowbreak{}Variable\allowbreak{}Elimination\allowbreak{}Core} (líneas 210-213). Allí sí normaliza cuando no hay variables condicionantes (la propagación normal), pero no cuando \texttt{VETemporal\allowbreak{}Evolution} le pasa la decisión, que es lo que hace la exportación a Excel de la evolución temporal.

Medido con \texttt{0\_\allowbreak{}miscellaneous/\allowbreak{}networks/\allowbreak{}mid/\allowbreak{}MID-Chancellor.\allowbreak{}pgmx}, variable «State», decisión «Therapy type» (suma de cada columna, una por opción):

\begin{Verbatim}[breaklines,breakanywhere,fontsize=\small]
sin evidencia:                State [0..4]: 1.000000 1.000000
State [1] = State A:          State [0], [1]: 1.000000 1.000000 ; State [2..4]: 0.721453 0.858220
State [1] = State B:          State [2..4]: 0.201845 0.102739
State [1] = State C:          State [2..4]: 0.066897 0.034051
\end{Verbatim}

Los ciclos 0 y 1 suman 1 porque la variable observada se trata aparte con un potencial determinista. Llamando a la tarea sin decisión, cosa que la ventana no permite, las columnas suman 2,0: la decisión sin política se suma como si fuera una variable de azar.

Un primer intento de medida con \texttt{0\_\allowbreak{}miscellaneous/\allowbreak{}networks/\allowbreak{}mid/\allowbreak{}MID-Chancellor-new.\allowbreak{}pgmx} (horizonte 1, variable \texttt{State\ {[}1{]}}, hallazgo en \texttt{State\ {[}0{]}}) dio las mismas sumas con y sin evidencia (1,0 en el ciclo 0 y 0,0 en el ciclo 1), quizá porque el programa de prueba no preparaba la red como lo hace la ventana. Ese resultado no se ha sabido interpretar; la medida que vale es la de \texttt{MID-Chancellor.\allowbreak{}pgmx}.
\paragraph{El código}
inference/src/main/java/org/openmarkov/inference/algorithm/temporalevaluation/tasks/MIDTemporalEvolution.java:292-294

\begin{Shaded}
\begin{Highlighting}[]
    \KeywordTok{private} \DataTypeTok{void} \FunctionTok{commonPreprocessing}\OperatorTok{()} \KeywordTok{throws}\NormalTok{ NotAllNodesHavePoliciesException}\OperatorTok{,}\NormalTok{ IncompatibleEvidenceException}\OperatorTok{.}\FunctionTok{EvidenceIsIncompatibleWithOther}\OperatorTok{,}\NormalTok{ NonProjectablePotentialException }\OperatorTok{\{}
        \FunctionTok{checkDecision}\OperatorTok{(}\NormalTok{probNet}\OperatorTok{,}\NormalTok{ probNet}\OperatorTok{.}\FunctionTok{getNode}\OperatorTok{(}\NormalTok{conditioningDecision}\OperatorTok{));}
        \FunctionTok{setConditioningVariables}\OperatorTok{(}\BuiltInTok{Collections}\OperatorTok{.}\FunctionTok{singletonList}\OperatorTok{(}\NormalTok{conditioningDecision}\OperatorTok{));}
\end{Highlighting}
\end{Shaded}

inference/src/main/java/org/openmarkov/inference/algorithm/temporalevaluation/tasks/MIDTemporalEvolution.java:605-617

\begin{Shaded}
\begin{Highlighting}[]
\NormalTok{        TablePotential posteriorValue }\OperatorTok{=}\NormalTok{ variableEliminationCore}\OperatorTok{.}\FunctionTok{getProbability}\OperatorTok{();}
        \CommentTok{//copied from org.openmarkov.inference.variableElimination.tasks.VEPropagation.InvokeVariableEliminationCore lines 220{-}236}
        \ControlFlowTok{if} \OperatorTok{(}\NormalTok{posteriorValue}\OperatorTok{.}\FunctionTok{getVariables}\OperatorTok{().}\FunctionTok{getFirst}\OperatorTok{()} \OperatorTok{!=}\NormalTok{ probabilityVariable}\OperatorTok{)} \OperatorTok{\{}
            \CommentTok{// }\AlertTok{TODO}\CommentTok{ {-} Comprobar este código}
            \BuiltInTok{List}\OperatorTok{\textless{}}\NormalTok{Variable}\OperatorTok{\textgreater{}}\NormalTok{ orderedVariables }\OperatorTok{=} \KeywordTok{new} \BuiltInTok{ArrayList}\OperatorTok{\textless{}\textgreater{}(}\BuiltInTok{Arrays}\OperatorTok{.}\FunctionTok{asList}\OperatorTok{(}\NormalTok{probabilityVariable}\OperatorTok{,}\NormalTok{ conditioningDecision}\OperatorTok{));}
\NormalTok{            posteriorValue }\OperatorTok{=}\NormalTok{ posteriorValue}\OperatorTok{.}\FunctionTok{reorder}\OperatorTok{(}\NormalTok{orderedVariables}\OperatorTok{);}
        \OperatorTok{\}}
        \CommentTok{// }\AlertTok{TODO}\CommentTok{ {-} Realizar la normalización condicionada}
        \ControlFlowTok{if} \OperatorTok{(}\FunctionTok{getConditioningVariables}\OperatorTok{()} \OperatorTok{==} \KeywordTok{null} \OperatorTok{||} \FunctionTok{getConditioningVariables}\OperatorTok{().}\FunctionTok{isEmpty}\OperatorTok{())} \OperatorTok{\{}
\NormalTok{            DiscretePotentialOperations}\OperatorTok{.}\FunctionTok{normalize}\OperatorTok{(}\NormalTok{posteriorValue}\OperatorTok{);}
        \OperatorTok{\}}
        \CommentTok{//end copied}
        \ControlFlowTok{return}\NormalTok{ posteriorValue}\OperatorTok{;}
\end{Highlighting}
\end{Shaded}
\paragraph{Cómo arreglarlo}
Antes de tocar el código hay que tomar la decisión de equipo.

Si se normaliza:

\begin{itemize}
\tightlist
\item
  Sin decisión condicionante, \texttt{common\allowbreak{}Preprocessing} debe dejar la lista vacía (no una lista con un nulo) y normalizar la tabla de una variable.
\item
  Con decisión, normalizar P(V, D, e) columna a columna, una columna por opción de D.
\item
  Una columna que suma cero significa que la evidencia es imposible bajo esa alternativa, aunque no lo sea bajo otras; no puede tratarse como error. La normalización por columnas se niega hoy a dividir por una columna que suma cero y lanza una excepción (commit \texttt{87129d6}), así que esa columna hay que marcarla como alternativa imposible en vez de dejar que la excepción salte.
\item
  Ajustar a la vez el \texttt{catch} de \texttt{resolve(\allowbreak{})}, que convertiría esa excepción en \texttt{Unreachable\allowbreak{}Exception}, y el hilo del monitor de progreso de \texttt{Trace\allowbreak{}Temporal\allowbreak{}Evolution\allowbreak{}Dialog} (líneas 315-334), que espera un aviso que no llega si hay una excepción.
\end{itemize}

Si no se normaliza: la ventana y el javadoc deben decir que los valores son P(V, e) y no probabilidades condicionadas.

En cualquier caso, la línea 609 no debería construir listas con nulos.

Sitios donde la misma regla tiene que cumplirse: \texttt{MIDTemporal\allowbreak{}Evolution.\allowbreak{}compute\allowbreak{}Posterior\allowbreak{}Probability} y \texttt{VEPropagation.\allowbreak{}Invoke\allowbreak{}Variable\allowbreak{}Elimination\allowbreak{}Core} (usado por \texttt{VETemporal\allowbreak{}Evolution} desde la exportación a Excel). Los potenciales deterministas de las variables observadas (\texttt{create\allowbreak{}Deterministic\allowbreak{}Potential}) ya suman uno.

Pruebas de regresión:

\begin{itemize}
\tightlist
\item
  Evolución temporal de «State» en \texttt{0\_\allowbreak{}miscellaneous/\allowbreak{}networks/\allowbreak{}mid/\allowbreak{}MID-Chancellor.\allowbreak{}pgmx} con el hallazgo State {[}1{]} = State A: si se decide normalizar, cada columna de los ciclos 2 en adelante debe sumar 1.
\item
  Un MID pequeño con evidencia en una variable del ciclo 0: la evolución de otra variable debe sumar uno en cada ciclo o, si se decide no normalizar, sumar P(e).
\end{itemize}

\setcounter{subsection}{67}
\subsection[La evolución temporal de una variable que empieza en el ciclo 1 falla]{La evolución temporal de una variable que no empieza en el ciclo cero termina con NullPointerException}\label{err:68}
\begin{ficha}
\fichakey{Estado}Presente
\fichakey{Módulo}inference
\fichakey{Paquete}org.openmarkov.inference.algorithm.temporalevaluation.tasks
\fichakey{Ficheros}\path{inference/src/main/java/org/openmarkov/inference/algorithm/temporalevaluation/tasks/MIDTemporalEvolution.java:161-174} \newline \path{inference/src/main/java/org/openmarkov/inference/algorithm/temporalevaluation/tasks/MIDTemporalEvolution.java:440-451} \newline \path{inference/src/main/java/org/openmarkov/inference/algorithm/temporalevaluation/tasks/MIDTemporalEvolution.java:553-565}
\fichakey{Comprobado}Leyendo el código y ejecutándolo
\fichakey{Esfuerzo}Pequeño (horas)
\fichakey{Decisión de equipo}No
\fichakey{Relacionados}\hyperref[err:69]{error~69}
\end{ficha}
\paragraph{Qué pasa}
En modo inferencia de un modelo de Markov con decisiones (MID, diagrama de influencia de Markov), la evolución temporal de una variable que el modelo solo define a partir del ciclo 1 no se puede pedir: sale un error inesperado y la ventana con la gráfica no se abre. Son variables frecuentes: costes y utilidades que se cuentan desde el primer ciclo, variables de tratamiento aplicado\ldots{} No se saltan ciclos en silencio: no se obtiene nada.

Camino de llegada: abrir \texttt{0\_\allowbreak{}miscellaneous/\allowbreak{}networks/\allowbreak{}mid/\allowbreak{}MID-Chancellor-corrected.\allowbreak{}pgmx}, pasar a modo inferencia, clic derecho sobre un nodo temporal cuyo primer ciclo es el 1 (por ejemplo «Cost AZT {[}1{]}» o «Therapy applied {[}1{]}») → «Temporal Evolution» → aceptar. \texttt{Network\allowbreak{}Editor\allowbreak{}Panel} abre \texttt{Temporal\allowbreak{}Evolution\allowbreak{}Dialog}, que al aceptar crea \texttt{Trace\allowbreak{}Temporal\allowbreak{}Evolution\allowbreak{}Dialog}, y este lanza \texttt{new\ MIDTemporal\allowbreak{}Evolution(\allowbreak{}red,\allowbreak{}\ variable)} en un hilo aparte. La excepción no está entre las que ese hilo captura, así que llega al manejador general de excepciones de la aplicación, que la enseña como error inesperado, y la ventana de resultados no se abre (esto último, por lectura del código).

\texttt{MIDTemporal\allowbreak{}Evolution} calcula el valor de una variable ciclo a ciclo. Su constructor averigua en qué ciclo empieza la variable pedida con \texttt{get\allowbreak{}First\allowbreak{}Slice} y busca la variable de ese ciclo:
\texttt{first\allowbreak{}Slice\ =\ get\allowbreak{}First\allowbreak{}Slice(\allowbreak{}temporal\allowbreak{}Variable,\allowbreak{}\ prob\allowbreak{}Net);\ this.\allowbreak{}temporal\allowbreak{}Variable\ =\ prob\allowbreak{}Net.\allowbreak{}get\allowbreak{}Variable(\allowbreak{}temporal\allowbreak{}Variable.\allowbreak{}get\allowbreak{}Base\allowbreak{}Name(\allowbreak{}),\allowbreak{}\ first\allowbreak{}Slice);}

\texttt{get\allowbreak{}First\allowbreak{}Slice} busca la variable en el ciclo 0, descarta el resultado y devuelve 0 siempre. La cabecera de la clase promete contemplar variables que no empiezan en el ciclo cero, y el javadoc del propio método lleva un «FIXME» que reconoce no saber hacerlo.

Para una variable que solo existe desde el ciclo 1, \texttt{prob\allowbreak{}Net.\allowbreak{}get\allowbreak{}Variable(\allowbreak{}nombre,\allowbreak{}\ 0)} devuelve nulo, \texttt{this.\allowbreak{}temporal\allowbreak{}Variable} queda nulo y la evaluación muere en \texttt{evaluate\allowbreak{}Slice} (línea 450) al pedirle el nombre base.

Medido sobre \texttt{0\_\allowbreak{}miscellaneous/\allowbreak{}networks/\allowbreak{}mid/\allowbreak{}MID-Chancellor-corrected.\allowbreak{}pgmx} recorriendo todos sus nodos temporales, llamando a la tarea como lo hace la ventana, sin evidencia, con la decisión condicionante que pone el diálogo y también sin ella. La decisión de esa red se llama «Therapy choice» («Therapy type» es la de \texttt{MID-Chancellor.\allowbreak{}pgmx}). Siete de sus nueve variables temporales no se pueden pedir:

\begin{Verbatim}[breaklines,breakanywhere,fontsize=\small]
Community care cost, Cost AZT, Cost lamivudine, Direct medical cost, Life years (utilidad, primer ciclo 1),
Therapy applied, Transition inhibited (azar, primer ciclo 1):
  java.lang.NullPointerException: Cannot invoke "Variable.getBaseName()" because "this.temporalVariable" is null
     at MIDTemporalEvolution.evaluateSlice(MIDTemporalEvolution.java:450)
     at MIDTemporalEvolution.evaluate(MIDTemporalEvolution.java:522)
     at MIDTemporalEvolution.resolve(MIDTemporalEvolution.java:277)
State, Time in treatment (primer ciclo 0): 21 resultados, ciclos 0..20
\end{Verbatim}
\paragraph{El código}
inference/src/main/java/org/openmarkov/inference/algorithm/temporalevaluation/tasks/MIDTemporalEvolution.java:166-168

\begin{Shaded}
\begin{Highlighting}[]
        \CommentTok{//Gets the first slice in which temporalVariable sequence is defined}
\NormalTok{        firstSlice }\OperatorTok{=} \FunctionTok{getFirstSlice}\OperatorTok{(}\NormalTok{temporalVariable}\OperatorTok{,}\NormalTok{ probNet}\OperatorTok{);}
        \KeywordTok{this}\OperatorTok{.}\FunctionTok{temporalVariable} \OperatorTok{=}\NormalTok{ probNet}\OperatorTok{.}\FunctionTok{getVariable}\OperatorTok{(}\NormalTok{temporalVariable}\OperatorTok{.}\FunctionTok{getBaseName}\OperatorTok{(),}\NormalTok{ firstSlice}\OperatorTok{);}
\end{Highlighting}
\end{Shaded}

inference/src/main/java/org/openmarkov/inference/algorithm/temporalevaluation/tasks/MIDTemporalEvolution.java:561-565

\begin{Shaded}
\begin{Highlighting}[]
    \KeywordTok{private} \DataTypeTok{static} \DataTypeTok{int} \FunctionTok{getFirstSlice}\OperatorTok{(}\NormalTok{Variable temporalVariable}\OperatorTok{,}\NormalTok{ ProbNet probNet}\OperatorTok{)} \OperatorTok{\{}
        \DataTypeTok{int}\NormalTok{ timeSlice }\OperatorTok{=} \DecValTok{0}\OperatorTok{;}
\NormalTok{        probNet}\OperatorTok{.}\FunctionTok{getVariable}\OperatorTok{(}\NormalTok{temporalVariable}\OperatorTok{.}\FunctionTok{getBaseName}\OperatorTok{(),}\NormalTok{ timeSlice}\OperatorTok{);}
        \ControlFlowTok{return}\NormalTok{ timeSlice}\OperatorTok{;}
    \OperatorTok{\}}
\end{Highlighting}
\end{Shaded}

inference/src/main/java/org/openmarkov/inference/algorithm/temporalevaluation/tasks/MIDTemporalEvolution.java:449-451

\begin{Shaded}
\begin{Highlighting}[]
\OperatorTok{\}} \ControlFlowTok{else} \ControlFlowTok{if} \OperatorTok{(}\NormalTok{firstSlice }\OperatorTok{\textless{}=}\NormalTok{ slice}\OperatorTok{)} \OperatorTok{\{}
\NormalTok{    temporalVariableSlice }\OperatorTok{=}\NormalTok{ probNet}\OperatorTok{.}\FunctionTok{getVariable}\OperatorTok{(}\NormalTok{temporalVariable}\OperatorTok{.}\FunctionTok{getBaseName}\OperatorTok{(),}\NormalTok{ slice}\OperatorTok{);}
\OperatorTok{\}}
\end{Highlighting}
\end{Shaded}
\paragraph{Cómo arreglarlo}
Hacer que \texttt{get\allowbreak{}First\allowbreak{}Slice} devuelva el menor ciclo en que la red tiene una variable con ese nombre base. Dos maneras: recorrer los ciclos desde 0 hasta el horizonte hasta que \texttt{get\allowbreak{}Variable(\allowbreak{}base\allowbreak{}Name,\allowbreak{}\ k)} no sea \texttt{null}, o recorrer \texttt{prob\allowbreak{}Net.\allowbreak{}get\allowbreak{}Variables(\allowbreak{})} y quedarse con el mínimo \texttt{get\allowbreak{}Time\allowbreak{}Slice(\allowbreak{})} de las que comparten \texttt{get\allowbreak{}Base\allowbreak{}Name(\allowbreak{})}. Si no hay ninguna, lanzar una excepción con mensaje.

Con eso \texttt{this.\allowbreak{}temporal\allowbreak{}Variable} deja de ser \texttt{null} y la condición \texttt{first\allowbreak{}Slice\ \textless{}=\ slice} de \texttt{evaluate\allowbreak{}Slice} (líneas 449 y 467), que ya compara con el ciclo en curso, hace lo que promete la cabecera de la clase («contempla la posibilidad de que la secuencia no empiece en el ciclo cero»).

Hay que comprobar que la tabla de resultados (\texttt{Temporal\allowbreak{}Evolution\allowbreak{}Table\allowbreak{}Pane}) y la gráfica aceptan una evolución sin ciclo 0: el constructor de tabla por criterio ya escribe «-» antes del primer ciclo; el constructor de tabla por variable recorre ciclos desde 0 y rellena con 0,0 los que faltan.

Prueba de regresión: sobre \texttt{0\_\allowbreak{}miscellaneous/\allowbreak{}networks/\allowbreak{}mid/\allowbreak{}MID-Chancellor-corrected.\allowbreak{}pgmx}, \texttt{new\ MIDTemporal\allowbreak{}Evolution(\allowbreak{}red,\allowbreak{}\ variable).\allowbreak{}get\allowbreak{}Temporal\allowbreak{}Evolution(\allowbreak{})} debe devolver resultados para los ciclos 1 a 20, sin excepción, para cada uno de los siete nodos citados (en particular «Cost AZT {[}1{]}» y «Therapy applied {[}1{]}»).

\setcounter{subsection}{68}
\subsection[La evolución temporal de una decisión temporal muere con una excepción]{Pedir la evolución temporal de una decisión temporal muere con NoSuchElementException}\label{err:69}
\begin{ficha}
\fichakey{Estado}Presente
\fichakey{Módulo}inference
\fichakey{Paquete}org.openmarkov.inference.algorithm.temporalevaluation.tasks
\fichakey{Ficheros}\path{inference/src/main/java/org/openmarkov/inference/algorithm/temporalevaluation/tasks/MIDTemporalEvolution.java:602-619} \newline \path{inference/src/main/java/org/openmarkov/inference/algorithm/temporalevaluation/tasks/MIDTemporalEvolution.java:440-481} \newline \path{gui/src/main/java/org/openmarkov/gui/window/MainPanelMenuAssistant.java:629-643} \newline \path{gui/src/main/java/org/openmarkov/gui/menutoolbar/menu/NodeContextualMenu.java:168-169} \newline \path{gui/src/main/java/org/openmarkov/gui/dialog/inference/temporalevolution/TraceTemporalEvolutionDialog.java:311-340}
\fichakey{Comprobado}Leyendo el código y ejecutándolo
\fichakey{Esfuerzo}Medio (días)
\fichakey{Decisión de equipo}\textcolor{decisionrojo}{\textbf{Sí.}} Si la evolución temporal de una decisión tiene sentido (probabilidad de elegir cada opción en cada ciclo) o si la opción no debe ofrecerse sobre decisiones
\fichakey{Relacionados}\hyperref[err:67]{error~67}, \hyperref[err:68]{error~68}
\end{ficha}
\paragraph{Qué pasa}
En un MID (diagrama de influencia de Markov), el diálogo de propiedades de cualquier nodo ofrece el selector de ciclo temporal (\texttt{Node\allowbreak{}Definition\allowbreak{}Panel}, visible cuando \texttt{variables\allowbreak{}Could\allowbreak{}Be\allowbreak{}Temporal(\allowbreak{})}), así que una decisión puede ser temporal («Therapy {[}0{]}»). Sobre un nodo temporal, el menú contextual añade «Temporal Evolution» (\texttt{Node\allowbreak{}Contextual\allowbreak{}Menu}, líneas 168-169) y \texttt{Main\allowbreak{}Panel\allowbreak{}Menu\allowbreak{}Assistant} la activa para todo nodo temporal salvo los de azar que no son de estados finitos (líneas 629-642): decisiones incluidas. Al aceptar el diálogo de ciclos, \texttt{Trace\allowbreak{}Temporal\allowbreak{}Evolution\allowbreak{}Dialog} lanza en un hilo aparte \texttt{new\ MIDTemporal\allowbreak{}Evolution(\allowbreak{}red,\allowbreak{}\ variable).\allowbreak{}get\allowbreak{}Temporal\allowbreak{}Evolution(\allowbreak{})}.

Sobre una decisión, el cálculo muere con \texttt{No\allowbreak{}Such\allowbreak{}Element\allowbreak{}Exception}. Nadie la recoge en el hilo (el \texttt{catch} del diálogo solo recoge \texttt{Index\allowbreak{}Out\allowbreak{}Of\allowbreak{}Bounds\allowbreak{}Exception} y excepciones comprobadas), así que llega al manejador general, que la enseña como error inesperado, y la ventana de resultados no se abre.

En cada ciclo, \texttt{evaluate\allowbreak{}Slice} quita la variable de interés de la lista de variables a eliminar y pide a \texttt{Variable\allowbreak{}Elimination\allowbreak{}Core} la probabilidad. Para una decisión el resultado es un potencial \textbf{sin variables}, y \texttt{compute\allowbreak{}Posterior\allowbreak{}Probability} hace \texttt{posterior\allowbreak{}Value.\allowbreak{}get\allowbreak{}Variables(\allowbreak{}).\allowbreak{}get\allowbreak{}First(\allowbreak{})} sobre una lista vacía.

Medido con un pequeño programa de prueba sobre un MID mínimo: \texttt{Illness\ {[}0{]}} (azar, estados healthy/ill), \texttt{Therapy\ {[}0{]}} (decisión, no/yes) y \texttt{Cost\ {[}0{]}} (utilidad que depende de las dos), horizonte 5. Se probó con la decisión condicionante que pone el diálogo por omisión (la única decisión sin política, \texttt{Therapy\ {[}0{]}}) y también sin decisión condicionante:

\begin{Verbatim}[breaklines,breakanywhere,fontsize=\small]
decision Therapy [0] EXC java.util.NoSuchElementException
   at java.util.ArrayList.getFirst(ArrayList.java:440)
   at MIDTemporalEvolution.computePosteriorProbability(MIDTemporalEvolution.java:607)
   at MIDTemporalEvolution.evaluateSlice(MIDTemporalEvolution.java:475)
\end{Verbatim}

\texttt{Illness\ {[}0{]}} y \texttt{Cost\ {[}0{]}} sí se calculan.

Ninguna red de \texttt{0\_\allowbreak{}miscellaneous/\allowbreak{}networks/\allowbreak{}mid} ni de \texttt{integration\allowbreak{}Tests} tiene decisiones temporales, así que para reproducirlo hace falta construir un MID como el descrito.
\paragraph{El código}
inference/src/main/java/org/openmarkov/inference/algorithm/temporalevaluation/tasks/MIDTemporalEvolution.java:602-612

\begin{Shaded}
\begin{Highlighting}[]
\KeywordTok{private}\NormalTok{ TablePotential }\FunctionTok{computePosteriorProbability}\OperatorTok{(}\NormalTok{VariableEliminationCore variableEliminationCore}\OperatorTok{,}\NormalTok{ Variable}
\NormalTok{        probabilityVariable}\OperatorTok{)} \KeywordTok{throws}\NormalTok{ CannotNormalizePotentialException }\OperatorTok{\{}
    \CommentTok{// 13/10/2022 only one conditioning decision variable}
\NormalTok{    TablePotential posteriorValue }\OperatorTok{=}\NormalTok{ variableEliminationCore}\OperatorTok{.}\FunctionTok{getProbability}\OperatorTok{();}
    \CommentTok{//copied from org.openmarkov.inference.variableElimination.tasks.VEPropagation.InvokeVariableEliminationCore lines 220{-}236}
    \ControlFlowTok{if} \OperatorTok{(}\NormalTok{posteriorValue}\OperatorTok{.}\FunctionTok{getVariables}\OperatorTok{().}\FunctionTok{getFirst}\OperatorTok{()} \OperatorTok{!=}\NormalTok{ probabilityVariable}\OperatorTok{)} \OperatorTok{\{}
        \CommentTok{// }\AlertTok{TODO}\CommentTok{ {-} Comprobar este código}
        \BuiltInTok{List}\OperatorTok{\textless{}}\NormalTok{Variable}\OperatorTok{\textgreater{}}\NormalTok{ orderedVariables }\OperatorTok{=} \KeywordTok{new} \BuiltInTok{ArrayList}\OperatorTok{\textless{}\textgreater{}(}\BuiltInTok{Arrays}\OperatorTok{.}\FunctionTok{asList}\OperatorTok{(}\NormalTok{probabilityVariable}\OperatorTok{,}\NormalTok{ conditioningDecision}\OperatorTok{));}
\NormalTok{        posteriorValue }\OperatorTok{=}\NormalTok{ posteriorValue}\OperatorTok{.}\FunctionTok{reorder}\OperatorTok{(}\NormalTok{orderedVariables}\OperatorTok{);}
    \OperatorTok{\}}
\end{Highlighting}
\end{Shaded}

gui/src/main/java/org/openmarkov/gui/window/MainPanelMenuAssistant.java:629-642

\begin{Shaded}
\begin{Highlighting}[]
\ControlFlowTok{if} \OperatorTok{(}\NormalTok{visualNode}\OperatorTok{.}\FunctionTok{getNode}\OperatorTok{().}\FunctionTok{getVariable}\OperatorTok{().}\FunctionTok{isTemporal}\OperatorTok{())} \OperatorTok{\{}
    \KeywordTok{...}
    \ControlFlowTok{if} \OperatorTok{(!(}
\NormalTok{            visualNode}\OperatorTok{.}\FunctionTok{getNode}\OperatorTok{().}\FunctionTok{getNodeType}\OperatorTok{()} \OperatorTok{==}\NormalTok{ NodeType}\OperatorTok{.}\FunctionTok{CHANCE} \OperatorTok{\&\&}
\NormalTok{                    visualNode}\OperatorTok{.}\FunctionTok{getNode}\OperatorTok{()}
                              \OperatorTok{.}\FunctionTok{getVariable}\OperatorTok{()}
                              \OperatorTok{.}\FunctionTok{getVariableType}\OperatorTok{()} \OperatorTok{!=}\NormalTok{ VariableType}\OperatorTok{.}\FunctionTok{FINITE\_STATES}
    \OperatorTok{))} \OperatorTok{\{}
\NormalTok{        canLog }\OperatorTok{=} \KeywordTok{true}\OperatorTok{;}
\NormalTok{        canTemporalEvolution }\OperatorTok{=} \KeywordTok{true}\OperatorTok{;}
    \OperatorTok{\}}
\OperatorTok{\}}
\end{Highlighting}
\end{Shaded}
\paragraph{Cómo arreglarlo}
Hay que decidir qué significa la evolución temporal de una decisión.

\begin{itemize}
\tightlist
\item
  Si se quiere contestar (probabilidad de tomar cada opción en cada ciclo según la política óptima o impuesta), \texttt{MIDTemporal\allowbreak{}Evolution} necesita un camino propio para decisiones: hoy \texttt{evaluate\allowbreak{}Slice} las trata como variables de probabilidad y la eliminación no produce un potencial sobre ellas.
\item
  Si no se quiere, la opción no debe activarse sobre decisiones: hay que cambiar la condición de \texttt{Main\allowbreak{}Panel\allowbreak{}Menu\allowbreak{}Assistant} (líneas 634-642) y la construcción del menú en \texttt{Node\allowbreak{}Contextual\allowbreak{}Menu} (líneas 168-169). Además, \texttt{MIDTemporal\allowbreak{}Evolution} debería rechazar la variable con una excepción declarada en su constructor, como ya hace con \texttt{Variable\allowbreak{}Is\allowbreak{}Not\allowbreak{}Temporal}.
\end{itemize}

En ambos casos, \texttt{compute\allowbreak{}Posterior\allowbreak{}Probability} no debe llamar a \texttt{get\allowbreak{}First(\allowbreak{})} sin comprobar que la lista no esté vacía.

Prueba: construir el MID mínimo descrito y comprobar que \texttt{new\ MIDTemporal\allowbreak{}Evolution(\allowbreak{}red,\allowbreak{}\ Therapy{[}0{]}).\allowbreak{}get\allowbreak{}Temporal\allowbreak{}Evolution(\allowbreak{})} o bien devuelve una tabla por ciclo, o bien el constructor lanza la excepción declarada; y que la opción de menú queda desactivada si se elige no ofrecerla.

\setcounter{subsection}{69}
\subsection[«Same as previous» revienta si el ciclo anterior no tiene la variable]{«Igual que el ciclo anterior» revienta si el ciclo anterior no tiene esa variable o su nodo no tiene potencial}\label{err:70}
\begin{ficha}
\fichakey{Estado}Presente
\fichakey{Módulo}core
\fichakey{Paquete}org.openmarkov.core.model.network.potential
\fichakey{Ficheros}\path{core/src/main/java/org/openmarkov/core/model/network/potential/SameAsPrevious.java:61-63} \newline \path{core/src/main/java/org/openmarkov/core/model/network/potential/SameAsPrevious.java:85-103} \newline \path{core/src/main/java/org/openmarkov/core/model/network/TemporalNetOperations.java:66-76} \newline \path{gui/src/main/java/org/openmarkov/gui/window/MainPanelMenuAssistant.java:237-265}
\fichakey{Comprobado}Leyendo el código y ejecutándolo
\fichakey{Esfuerzo}Pequeño (horas)
\fichakey{Decisión de equipo}No
\fichakey{Relacionados}\hyperref[err:29]{error~29}, \hyperref[err:43]{error~43}, \hyperref[err:71]{error~71}
\end{ficha}
\paragraph{Qué pasa}
En un modelo temporal (red bayesiana dinámica, diagrama de influencia de Markov, etc.) cada variable existe en varios ciclos o cortes de tiempo (\texttt{Edad\ {[}0{]}}, \texttt{Edad\ {[}1{]}}\ldots). El potencial \texttt{Same\allowbreak{}As\allowbreak{}Previous} («Same as previous», igual que el anterior) de un nodo del ciclo \emph{t} no guarda números: indica que su distribución es la del mismo nodo en el ciclo anterior.

Basta con abrir un fichero en el que un nodo con «Same as previous» no tiene la variable en el ciclo de antes, o con elegir «Same as previous» para un nodo así, para que salte una excepción al refrescar los menús, y vuelva a saltar con cada edición posterior. Una variable que aparece a mitad del modelo ---un tratamiento que empieza en el ciclo 1 o en el 3--- tiene exactamente esa forma.

El camino se ha leído en el código, no ejecutado en la ventana: \texttt{Main\allowbreak{}Panel\allowbreak{}Menu\allowbreak{}Assistant.\allowbreak{}update\allowbreak{}Inference\allowbreak{}Buttons} recorre todos los potenciales de la red y llama a \texttt{get\allowbreak{}Original\allowbreak{}Potential} para cada \texttt{Same\allowbreak{}As\allowbreak{}Previous}, y se ejecuta al abrir una red (\texttt{update\allowbreak{}Options\allowbreak{}New\allowbreak{}Network\allowbreak{}Open}) y tras cada modificación (\texttt{update\allowbreak{}Options\allowbreak{}Network\allowbreak{}Modified}). En la evaluación, \texttt{Temporal\allowbreak{}Net\allowbreak{}Operations.\allowbreak{}compact\allowbreak{}Network} (usada por \texttt{expand\allowbreak{}Network}) llega al caso 3 de la medida de abajo cuando un nodo \texttt{Same\allowbreak{}As\allowbreak{}Previous} cuelga de otro \texttt{Same\allowbreak{}As\allowbreak{}Previous} sin variable en el ciclo de antes.

Nada impide llegar a esas formas. \texttt{Same\allowbreak{}As\allowbreak{}Previous.\allowbreak{}validate}, que decide si el editor de potenciales ofrece este tipo para un nodo, solo exige que la variable sea temporal y de ciclo mayor que cero; no mira si el ciclo anterior tiene la variable. El lector de \texttt{.\allowbreak{}pgmx} (\texttt{PGMXPotential\allowbreak{}Parsers.\allowbreak{}get\allowbreak{}Same\allowbreak{}As\allowbreak{}Previous}) solo exige que la variable sea temporal.

\texttt{get\allowbreak{}Original\allowbreak{}Potential} busca el potencial retrocediendo ciclo a ciclo hasta encontrar uno que no sea también \texttt{Same\allowbreak{}As\allowbreak{}Previous}. El bucle hace tres cosas seguidas sin comprobar ninguna: \texttt{prob\allowbreak{}Net.\allowbreak{}get\allowbreak{}Variable(\allowbreak{}nombre\allowbreak{}Base,\allowbreak{}\ -\/-ciclo)} (que devuelve \texttt{null} si ese ciclo no tiene la variable), \texttt{prob\allowbreak{}Net.\allowbreak{}get\allowbreak{}Node(\allowbreak{}variable)} (que devuelve \texttt{null} si recibe \texttt{null}) y \texttt{get\allowbreak{}Potentials(\allowbreak{}).\allowbreak{}get\allowbreak{}First(\allowbreak{})} (que lanza si la lista está vacía).

Además, cuando el método devuelve \texttt{null} (si el retroceso llega al ciclo 0 sin encontrar un potencial que no sea \texttt{Same\allowbreak{}As\allowbreak{}Previous}), sus dos llamadores lo usan sin comprobar: \texttt{Temporal\allowbreak{}Net\allowbreak{}Operations.\allowbreak{}compact\allowbreak{}Network} hace \texttt{.\allowbreak{}get\allowbreak{}Original\allowbreak{}Potential(\allowbreak{}prob\allowbreak{}Net).\allowbreak{}copy(\allowbreak{})} y \texttt{Main\allowbreak{}Panel\allowbreak{}Menu\allowbreak{}Assistant.\allowbreak{}update\allowbreak{}Inference\allowbreak{}Buttons} hace \texttt{original\allowbreak{}Potential.\allowbreak{}is\allowbreak{}Uncertain(\allowbreak{})}.

Medido con un pequeño programa de prueba sobre redes construidas en código:

\begin{Verbatim}[breaklines,breakanywhere,fontsize=\small]
1. solo existe Treatment [1] (SameAsPrevious)          -> NullPointerException (getNode(...) es null)
2. A [0] sin potenciales, A [1] SameAsPrevious          -> NoSuchElementException
3. B [2] y B [3] SameAsPrevious, no existe B [1]        -> NullPointerException
4. C [0] uniforme, C [1] SameAsPrevious (caso normal)   -> devuelve P(C [0]) = Uniform
\end{Verbatim}
\paragraph{El código}
core/src/main/java/org/openmarkov/core/model/network/potential/SameAsPrevious.java:85-103

\begin{Shaded}
\begin{Highlighting}[]
\KeywordTok{private} \DataTypeTok{static} \AttributeTok{@Nullable}\NormalTok{ Potential }\FunctionTok{getOriginalPotential}\OperatorTok{(}\NormalTok{ProbNet probNet}\OperatorTok{,}\NormalTok{ Variable variable}\OperatorTok{)} \OperatorTok{\{}
\NormalTok{    Potential previousPotential }\OperatorTok{=} \KeywordTok{null}\OperatorTok{;}
    \ControlFlowTok{if} \OperatorTok{(!}\NormalTok{variable}\OperatorTok{.}\FunctionTok{isTemporal}\OperatorTok{())} \OperatorTok{\{}
        \ControlFlowTok{return} \KeywordTok{null}\OperatorTok{;}
    \OperatorTok{\}}
    \DataTypeTok{int}\NormalTok{ timeSlice }\OperatorTok{=}\NormalTok{ variable}\OperatorTok{.}\FunctionTok{getTimeSlice}\OperatorTok{();}
    \AttributeTok{@Nullable}\NormalTok{ Variable previousVariable }\OperatorTok{=} \KeywordTok{null}\OperatorTok{;}
    \ControlFlowTok{while} \OperatorTok{(}\NormalTok{timeSlice }\OperatorTok{\textgreater{}} \DecValTok{0} \OperatorTok{\&\&}\NormalTok{ previousVariable }\OperatorTok{==} \KeywordTok{null}\OperatorTok{)} \OperatorTok{\{}
\NormalTok{        previousVariable }\OperatorTok{=}\NormalTok{ probNet}\OperatorTok{.}\FunctionTok{getVariable}\OperatorTok{(}\NormalTok{variable}\OperatorTok{.}\FunctionTok{getBaseName}\OperatorTok{(),} \OperatorTok{{-}{-}}\NormalTok{timeSlice}\OperatorTok{);}
\NormalTok{        previousPotential }\OperatorTok{=}\NormalTok{ probNet}\OperatorTok{.}\FunctionTok{getNode}\OperatorTok{(}\NormalTok{previousVariable}\OperatorTok{).}\FunctionTok{getPotentials}\OperatorTok{().}\FunctionTok{getFirst}\OperatorTok{();}
        \ControlFlowTok{if} \OperatorTok{(}\NormalTok{previousPotential }\KeywordTok{instanceof}\NormalTok{ SameAsPrevious}\OperatorTok{)} \OperatorTok{\{}
\NormalTok{            previousVariable }\OperatorTok{=} \KeywordTok{null}\OperatorTok{;}
        \OperatorTok{\}}
    \OperatorTok{\}}
    \ControlFlowTok{if} \OperatorTok{(}\NormalTok{previousVariable }\OperatorTok{==} \KeywordTok{null}\OperatorTok{)} \OperatorTok{\{}\CommentTok{// There is no previous variable}
        \ControlFlowTok{return} \KeywordTok{null}\OperatorTok{;}
    \OperatorTok{\}}
    \ControlFlowTok{return}\NormalTok{ previousPotential}\OperatorTok{;}
\OperatorTok{\}}
\end{Highlighting}
\end{Shaded}
\paragraph{Cómo arreglarlo}
Dentro del bucle, si \texttt{get\allowbreak{}Variable} devuelve \texttt{null} o el nodo no tiene potenciales, pasar al ciclo anterior sin tocar el nodo; al terminar sin encontrar nada, devolver \texttt{null} (ya documentado con \texttt{@Nullable}).

Después, los dos llamadores tienen que tratar ese \texttt{null}:

\begin{itemize}
\tightlist
\item
  \texttt{Temporal\allowbreak{}Net\allowbreak{}Operations.\allowbreak{}compact\allowbreak{}Network} (líneas 70-76): lanzar una excepción del dominio que diga qué nodo tiene «Same as previous» sin ciclo anterior del que copiar, en vez de un \texttt{Null\allowbreak{}Pointer\allowbreak{}Exception}.
\item
  \texttt{Main\allowbreak{}Panel\allowbreak{}Menu\allowbreak{}Assistant.\allowbreak{}update\allowbreak{}Inference\allowbreak{}Buttons} (líneas 253-258): tratarlo como potencial sin incertidumbre.
\end{itemize}

La regla puede imponerse también antes: \texttt{Same\allowbreak{}As\allowbreak{}Previous.\allowbreak{}validate} (líneas 61-63) podría exigir que exista la variable en el ciclo anterior, para que el editor no ofrezca este tipo donde no tiene sentido, y el lector podría avisar. Si se endurece \texttt{validate}, un fichero que ya tenga esa forma debe seguir abriéndose sin reventar, así que el arreglo del bucle es necesario en cualquier caso.

Prueba de regresión: en \texttt{core/\allowbreak{}src/\allowbreak{}test}, los cuatro casos de arriba (variable sin ciclo anterior, ciclo anterior sin potencial, cadena de dos \texttt{Same\allowbreak{}As\allowbreak{}Previous} sin ciclo anterior, caso normal), comprobando que los tres primeros devuelven \texttt{null} sin excepción y el cuarto devuelve el potencial del ciclo 0.

\setcounter{subsection}{70}
\subsection[Al expandir un modelo temporal, las regresiones siguen nombrando el ciclo 0]{Un modelo de Markov con regresiones que nombran variables del ciclo (CHAP) no se puede evaluar: al expandirlo, las fórmulas de los ciclos siguientes siguen nombrando las variables del ciclo 0}\label{err:71}
\begin{ficha}
\fichakey{Estado}Presente
\fichakey{Módulo}core
\fichakey{Paquete}org.openmarkov.core.model.network.potential; org.openmarkov.core.model.network
\fichakey{Ficheros}\path{core/src/main/java/org/openmarkov/core/model/network/potential/GLMPotential.java:152-158,320-322} \newline \path{core/src/main/java/org/openmarkov/core/model/network/potential/WeibullHazardPotential.java:320-325} \newline \path{core/src/main/java/org/openmarkov/core/model/network/potential/Potential.java:500-502} \newline \path{core/src/main/java/org/openmarkov/core/model/network/TemporalNetOperations.java:51-93,100-107,180-197} \newline \path{core/src/main/java/org/openmarkov/core/expression/VariableExpression.java:11-18} \newline \path{core/src/main/java/org/openmarkov/core/expression/ReferencedExpression.java:76-88} \newline \path{integrationTests/src/test/java/org/openmarkov/integrationTests/costeffectiveness/CEAGlobalAnalysisTest.java:75-78}
\fichakey{Comprobado}Leyendo el código y ejecutándolo
\fichakey{Esfuerzo}Pequeño (horas)
\fichakey{Decisión de equipo}No
\fichakey{Tapado por}\hyperref[err:46]{error~46}: en el modo inferencia, donde el cálculo de rangos falla antes.
\fichakey{Relacionados}\hyperref[err:10]{error~10}, \hyperref[err:13]{error~13}, \hyperref[err:14]{error~14}, \hyperref[err:15]{error~15}, \hyperref[err:46]{error~46}, \hyperref[err:70]{error~70}, \hyperref[err:114]{error~114}
\end{ficha}
\paragraph{Qué pasa}
Con \texttt{0\_\allowbreak{}miscellaneous/\allowbreak{}networks/\allowbreak{}mid/\allowbreak{}MID-CHAP-Ryan-Griffin.\allowbreak{}pgmx} abierto (un diagrama de influencia de Markov, MID, sobre profilaxis con cotrimoxazol), la herramienta de análisis de coste-efectividad (\texttt{Cost\allowbreak{}Effectiveness\allowbreak{}Plugin}), con ámbito global y tras aceptar sus dos diálogos de opciones, termina en un diálogo de error con el texto «An expression cannot be resolved as there is no value resolved for variable Age {[}0{]}» (\texttt{Non\allowbreak{}Projectable\allowbreak{}Potential\allowbreak{}Exception\$Cannot\allowbreak{}Resolve\allowbreak{}Variable}). El mensaje no dice qué potencial ni que el modelo esté mal; de hecho el modelo está bien. Con el mismo mensaje fallan, medido con un pequeño programa de prueba que hace las mismas llamadas:

\begin{itemize}
\tightlist
\item
  la evaluación (\texttt{VEEvaluation}) y la propagación del modo inferencia (esta, detrás del error de rangos de \hyperref[err:46]{error~46});
\item
  «Expand network» (\texttt{Temporal\allowbreak{}Net\allowbreak{}Operations.\allowbreak{}expand\allowbreak{}Network(\allowbreak{}red,\allowbreak{}\ evidencia,\allowbreak{}\ nombre)});
\item
  la evolución temporal (\texttt{MIDTemporal\allowbreak{}Evolution}), aquí envuelta en \texttt{Unreachable\allowbreak{}Exception}, que la ventana enseña como error no previsto.
\end{itemize}

Afecta a las cuatro redes CHAP del repositorio: la anterior, su copia \texttt{0\_\allowbreak{}miscellaneous/\allowbreak{}networks/\allowbreak{}mid/\allowbreak{}1-0/\allowbreak{}MID-CHAP-Ryan-Griffin-1-0.\allowbreak{}pgmx}, e \texttt{integration\allowbreak{}Tests/\allowbreak{}src/\allowbreak{}main/\allowbreak{}resources/\allowbreak{}networks/\allowbreak{}mid/\allowbreak{}chap.\allowbreak{}pgmx} y \texttt{chap-sv.\allowbreak{}pgmx}. Es la causa por la que se quitaron las pruebas \texttt{test\allowbreak{}CHAP} y \texttt{test\allowbreak{}CHAPSV} de \texttt{CEAGlobal\allowbreak{}Analysis\allowbreak{}Test} (el comentario de las líneas 75-78 lo dice). \texttt{0\_\allowbreak{}miscellaneous/\allowbreak{}networks/\allowbreak{}mid/\allowbreak{}1-0/\allowbreak{}MID-mammography-1-0.\allowbreak{}pgmx} tiene el mismo problema en las regresiones de «Invasive cancer», pero falla antes por otro motivo.

Por qué pasa. Un MID se describe con los ciclos 0 y 1, y antes de resolver se expande hasta el horizonte. \texttt{Temporal\allowbreak{}Net\allowbreak{}Operations.\allowbreak{}compact\allowbreak{}Network} crea el ciclo 1 de los nodos que solo existen en el ciclo 0, y \texttt{expand\allowbreak{}Network} (líneas 100-107) va copiando el último ciclo al siguiente. En los dos casos, \texttt{expand\allowbreak{}Potential\allowbreak{}And\allowbreak{}Links} (líneas 180-197) copia el potencial y llama a \texttt{Potential.\allowbreak{}shift}, que cambia su lista de variables por las del ciclo nuevo (líneas 500-502).

Las regresiones (la familia \texttt{GLMPotential}: combinación lineal, exponencial, función, riesgos de Weibull y exponencial) guardan cada covariable como un \texttt{Variable\allowbreak{}Expression}, una fórmula que apunta a objetos \texttt{Variable} concretos. \texttt{GLMPotential.\allowbreak{}shift} (líneas 320-322) solo llama a \texttt{super.\allowbreak{}shift}, y \texttt{Weibull\allowbreak{}Hazard\allowbreak{}Potential.\allowbreak{}shift} desplaza además la variable de tiempo, pero ninguno toca las fórmulas. Así, tras expandir CHAP:

\begin{Verbatim}[breaklines,breakanywhere,fontsize=\small]
Age at state entry [1]: variables [Age at state entry [1], Age [1], Time in state [1]]
                        fórmula   {Age [0]} + -1.0*{Time in state [0]}
\end{Verbatim}

Al calcular, \texttt{GLMPotential.\allowbreak{}table\allowbreak{}Project} da valor solo a las variables del potencial (\texttt{Age\ {[}1{]}}, \texttt{Time\ in\ state\ {[}1{]}}), y \texttt{Referenced\allowbreak{}Expression} busca \texttt{Age\ {[}0{]}} por objeto y por nombre sin encontrarlo (líneas 76-88). En CHAP expandida a 20 ciclos hay 426 referencias así: las de «Age at state entry {[}k{]}» y las de las exponenciales de los árboles de «Inpatient costs {[}k{]}» y «Outpatient costs {[}k{]}».

Medido: expandiendo CHAP y reescribiendo a mano cada fórmula para que nombre la variable del mismo nombre base de su potencial (350 fórmulas), \texttt{VECEAnalysis} sobre la red expandida termina: placebo, coste 183,2 y efectividad 2,629; cotrimoxazol, 188,8 y 4,082; el umbral entre ambas está en 3,87. Con esa reescritura se ve también el efecto de \hyperref[err:10]{error~10} sobre CHAP (allí se describe la medida).
\paragraph{El código}
core/src/main/java/org/openmarkov/core/model/network/potential/GLMPotential.java:320-322

\begin{Shaded}
\begin{Highlighting}[]
    \AttributeTok{@Override} \KeywordTok{public} \DataTypeTok{void} \FunctionTok{shift}\OperatorTok{(}\NormalTok{ProbNet probNet}\OperatorTok{,} \DataTypeTok{int}\NormalTok{ timeDifference}\OperatorTok{)} \OperatorTok{\{}
        \KeywordTok{super}\OperatorTok{.}\FunctionTok{shift}\OperatorTok{(}\NormalTok{probNet}\OperatorTok{,}\NormalTok{ timeDifference}\OperatorTok{);}
    \OperatorTok{\}}
\end{Highlighting}
\end{Shaded}

core/src/main/java/org/openmarkov/core/model/network/TemporalNetOperations.java:186-194

\begin{Shaded}
\begin{Highlighting}[]
\NormalTok{            Potential newPotential}\OperatorTok{;}
            \ControlFlowTok{if} \OperatorTok{(}\NormalTok{oldPotential }\KeywordTok{instanceof}\NormalTok{ CycleLengthShift}\OperatorTok{)} \OperatorTok{\{}
\NormalTok{                newPotential }\OperatorTok{=} \KeywordTok{new} \FunctionTok{CycleLengthShift}\OperatorTok{(}\NormalTok{oldPotential}\OperatorTok{.}\FunctionTok{getShiftedVariables}\OperatorTok{(}\NormalTok{probNet}\OperatorTok{,}\NormalTok{ timeDifference}\OperatorTok{),}
\NormalTok{                                                    probNet}\OperatorTok{.}\FunctionTok{getCycleLength}\OperatorTok{());}
            \OperatorTok{\}} \ControlFlowTok{else} \OperatorTok{\{}
\NormalTok{                newPotential }\OperatorTok{=}\NormalTok{ oldPotential}\OperatorTok{.}\FunctionTok{copy}\OperatorTok{();}
\NormalTok{                newPotential}\OperatorTok{.}\FunctionTok{shift}\OperatorTok{(}\NormalTok{probNet}\OperatorTok{,}\NormalTok{ timeDifference}\OperatorTok{);}
            \OperatorTok{\}}
\end{Highlighting}
\end{Shaded}
\paragraph{Cómo arreglarlo}
\texttt{GLMPotential.\allowbreak{}shift} tiene que desplazar también las covariables: para cada fórmula, sustituir cada variable temporal a la que apunta por \texttt{prob\allowbreak{}Net.\allowbreak{}get\allowbreak{}Shifted\allowbreak{}Variable(\allowbreak{}variable,\allowbreak{}\ time\allowbreak{}Difference)} (las atemporales se quedan) y construir un \texttt{Variable\allowbreak{}Expression} nuevo sobre la lista de variables desplazada, porque la fórmula guarda las variables por objeto y no basta con cambiar el texto. \texttt{Weibull\allowbreak{}Hazard\allowbreak{}Potential.\allowbreak{}shift} y \texttt{Exponential\allowbreak{}Hazard\allowbreak{}Potential} heredan de ahí; \texttt{Tree\allowbreak{}ADDPotential.\allowbreak{}shift} ya reenvía el desplazamiento a las ramas, así que los árboles quedan cubiertos.

Otros sitios de la misma regla, «un potencial y sus fórmulas nombran las mismas variables»:

\begin{itemize}
\tightlist
\item
  \texttt{Univariate\allowbreak{}Distr\allowbreak{}Potential} y \texttt{Augmented\allowbreak{}Prob\allowbreak{}Table}, cuyos parámetros también pueden ser fórmulas y que heredan \texttt{Potential.\allowbreak{}shift} sin tocarlas (no se ha encontrado ninguna red del repositorio con ese caso tras expandir);
\item
  \texttt{GLMPotential.\allowbreak{}deep\allowbreak{}Copy} (línea 343) y el constructor de copia (línea 84) copian el array de covariables pero no las fórmulas, que siguen apuntando a las variables de la red original; hoy funciona porque la búsqueda cae al nombre, y conviene comprobarlo al tocar esto.
\end{itemize}

Prueba de regresión: una prueba de \texttt{core} que expanda un MID de dos ciclos con un nodo \texttt{Y\ {[}0{]}} definido como combinación lineal de \texttt{\{X\ {[}0{]}\}} y compruebe que la fórmula de \texttt{Y\ {[}1{]}} nombra \texttt{X\ {[}1{]}}; y volver a poner \texttt{test\allowbreak{}CHAP} en \texttt{CEAGlobal\allowbreak{}Analysis\allowbreak{}Test}, con \texttt{VECEAnalysis} sobre \texttt{integration\allowbreak{}Tests/\allowbreak{}src/\allowbreak{}main/\allowbreak{}resources/\allowbreak{}networks/\allowbreak{}mid/\allowbreak{}chap.\allowbreak{}pgmx} terminando sin excepción.

\setcounter{subsection}{71}
\subsection[Pasar una variable a atemporal puede dejar dos con el mismo nombre]{Pasar a atemporal una variable temporal puede dejar dos variables con el mismo nombre, y la red se guarda y se reabre así}\label{err:72}
\begin{ficha}
\fichakey{Estado}Presente
\fichakey{Módulo}core
\fichakey{Paquete}org.openmarkov.core.action.core
\fichakey{Ficheros}\path{core/src/main/java/org/openmarkov/core/action/core/TimeSliceEdit.java:49-57} \newline \path{core/src/main/java/org/openmarkov/core/action/core/TimeSliceEdit.java:94-105} \newline \path{core/src/main/java/org/openmarkov/core/model/network/Variable.java:580-628} \newline \path{gui/src/main/java/org/openmarkov/gui/dialog/node/NodeDefinitionPanel.java:845-866}
\fichakey{Comprobado}Leyendo el código y ejecutándolo
\fichakey{Esfuerzo}Pequeño (horas)
\fichakey{Decisión de equipo}No
\fichakey{Relacionados}\hyperref[err:38]{error~38}
\end{ficha}
\paragraph{Qué pasa}
Una variable de un modelo temporal se llama, por ejemplo, «Therapy {[}2{]}»: nombre base e instante. Si la red ya tiene una variable atemporal llamada «Therapy» y el usuario pasa a atemporal la variable «Therapy {[}0{]}», el cambio se acepta y la red queda con dos variables llamadas «Therapy».

Camino de llegada: diálogo de propiedades del nodo, pestaña de definición, desplegable del instante → «Atemporal» (\texttt{Node\allowbreak{}Definition\allowbreak{}Panel}, que construye \texttt{Time\allowbreak{}Slice\allowbreak{}Edit(\allowbreak{}node,\allowbreak{}\ Integer.\allowbreak{}MIN\_\allowbreak{}VALUE)}), en una red donde ya existe una variable atemporal con el nombre base.

\texttt{Time\allowbreak{}Slice\allowbreak{}Edit} (módulo \texttt{core}) cambia el instante de una variable. Si la red tiene la restricción \texttt{Distinct\allowbreak{}Variable\allowbreak{}Names} (obligatoria por omisión), su comprobación previa quita de la lista de nombres \texttt{get\allowbreak{}Previous\allowbreak{}Name(\allowbreak{})} y mira si queda \texttt{get\allowbreak{}New\allowbreak{}Name(\allowbreak{})}. Los dos métodos componen siempre \texttt{nombre\allowbreak{}Base\ +\ "\ {[}"\ +\ instante\ +\ "{]}"}.

Cuando el instante nuevo es \texttt{Integer.\allowbreak{}MIN\_\allowbreak{}VALUE}, que significa «atemporal», el nombre que la variable tendrá de verdad es el nombre base («Therapy», así lo construye \texttt{Variable.\allowbreak{}rebuild\allowbreak{}Full\allowbreak{}Name}). Pero la comprobación pregunta por «Therapy {[}-2147483648{]}», que no existe nunca, y siempre pasa.

Comprobado ejecutando, en una red MID (diagrama de influencia de Markov) con «Therapy {[}0{]}» y «Therapy»:

\begin{itemize}
\tightlist
\item
  \texttt{get\allowbreak{}New\allowbreak{}Name(\allowbreak{})} devuelve «Therapy {[}-2147483648{]}»; la edición se ejecuta y la red queda con los nombres {[}Therapy, Therapy{]}.
\item
  \texttt{prob\allowbreak{}Net.\allowbreak{}get\allowbreak{}Variable(\allowbreak{}"Therapy")} devuelve la primera (estados «no, yes»), de modo que la otra queda oculta a toda búsqueda por nombre.
\item
  Deshacer devuelve «Therapy {[}0{]}» y rehacer vuelve a duplicar.
\item
  La comprobación de red completa (\texttt{prob\allowbreak{}Net.\allowbreak{}check\allowbreak{}Constraints(\allowbreak{})}) sí detecta \texttt{Variable\allowbreak{}Name\allowbreak{}Is\allowbreak{}Already\allowbreak{}Present}, pero la edición no la usa.
\item
  La red se escribe en \texttt{.\allowbreak{}pgmx} y se vuelve a leer con dos variables llamadas «Therapy»; como el fichero enlaza potenciales y enlaces por nombre, no vuelve como se guardó.
\end{itemize}

El paso contrario (de atemporal a un instante) no tiene el mismo hueco: la comprobación pregunta por el nombre correcto («Therapy {[}0{]}»).
\paragraph{El código}
core/src/main/java/org/openmarkov/core/action/core/TimeSliceEdit.java:49-57, 94-96

\begin{Shaded}
\begin{Highlighting}[]
\AttributeTok{@Override} \KeywordTok{public} \DataTypeTok{void} \FunctionTok{checkConstraintsWillBeMet}\OperatorTok{(}\NormalTok{ConstraintChecker constraintChecker}\OperatorTok{)} \OperatorTok{\{}
    \ControlFlowTok{if} \OperatorTok{(}\NormalTok{probNet}\OperatorTok{.}\FunctionTok{getConstraintOfClass}\OperatorTok{(}\NormalTok{DistinctVariableNames}\OperatorTok{.}\FunctionTok{class}\OperatorTok{)} \KeywordTok{instanceof}\NormalTok{ DistinctVariableNames constraint}\OperatorTok{)} \OperatorTok{\{}
        \BuiltInTok{List}\OperatorTok{\textless{}}\BuiltInTok{String}\OperatorTok{\textgreater{}}\NormalTok{ variablesProbNetNames }\OperatorTok{=}\NormalTok{ probNet}\OperatorTok{.}\FunctionTok{getVariablesNames}\OperatorTok{();}
\NormalTok{        variablesProbNetNames}\OperatorTok{.}\FunctionTok{remove}\OperatorTok{(}\KeywordTok{this}\OperatorTok{.}\FunctionTok{getPreviousName}\OperatorTok{());}
        \ControlFlowTok{if} \OperatorTok{(}\NormalTok{variablesProbNetNames}\OperatorTok{.}\FunctionTok{contains}\OperatorTok{(}\KeywordTok{this}\OperatorTok{.}\FunctionTok{getNewName}\OperatorTok{()))} \OperatorTok{\{}
\NormalTok{            constraintChecker}\OperatorTok{.}\FunctionTok{addException}\OperatorTok{(}\KeywordTok{new}\NormalTok{ ConstraintViolatedException}\OperatorTok{.}\FunctionTok{VariableNameIsAlreadyPresent}\OperatorTok{(}\NormalTok{constraint}\OperatorTok{,} \KeywordTok{this}\OperatorTok{.}\FunctionTok{getNewName}\OperatorTok{()));}
        \OperatorTok{\}}
    \OperatorTok{\}}
\OperatorTok{\}}
\KeywordTok{...}
\KeywordTok{public} \BuiltInTok{String} \FunctionTok{getNewName}\OperatorTok{()} \OperatorTok{\{}
    \ControlFlowTok{return}\NormalTok{ lastBaseName }\OperatorTok{+} \StringTok{" "} \OperatorTok{+} \StringTok{"["} \OperatorTok{+} \BuiltInTok{String}\OperatorTok{.}\FunctionTok{valueOf}\OperatorTok{(}\NormalTok{newTimeSlice}\OperatorTok{)} \OperatorTok{+} \StringTok{"]"}\OperatorTok{;}
\OperatorTok{\}}
\end{Highlighting}
\end{Shaded}
\paragraph{Cómo arreglarlo}
Que \texttt{get\allowbreak{}New\allowbreak{}Name(\allowbreak{})} y \texttt{get\allowbreak{}Previous\allowbreak{}Name(\allowbreak{})} devuelvan el nombre base solo cuando el instante es \texttt{Integer.\allowbreak{}MIN\_\allowbreak{}VALUE} (lo mismo que hace \texttt{Variable.\allowbreak{}rebuild\allowbreak{}Full\allowbreak{}Name}), o mejor, calcular el nombre con la propia regla de \texttt{Variable} para que no haya dos copias.

La regla «nombres distintos» vive en \texttt{Distinct\allowbreak{}Variable\allowbreak{}Names.\allowbreak{}check\allowbreak{}Prob\allowbreak{}Net} (red completa) y en cada edición que cambia un nombre: en esta, \texttt{Time\allowbreak{}Slice\allowbreak{}Edit.\allowbreak{}check\allowbreak{}Constraints\allowbreak{}Will\allowbreak{}Be\allowbreak{}Met}, y en la edición del nombre de nodo (\texttt{Node\allowbreak{}Base\allowbreak{}Name\allowbreak{}Edit}), que conviene revisar con la misma prueba.

Prueba: en una red con «Therapy {[}0{]}» y «Therapy», pasar la primera a atemporal debe rechazarse; pasar a atemporal sin conflicto debe aceptarse.

\setcounter{subsection}{72}
\subsection[La duración del ciclo escrita en las propiedades de la red no se aplica]{Cambiar la duración del ciclo en la pestaña de opciones temporales de las propiedades de la red no cambia la red}\label{err:73}
\begin{ficha}
\fichakey{Estado}Presente
\fichakey{Módulo}gui
\fichakey{Paquete}org.openmarkov.gui.dialog.network
\fichakey{Ficheros}\path{gui/src/main/java/org/openmarkov/gui/dialog/network/NetworkTemporalOptionsPanel.java:44-114} \newline \path{gui/src/main/java/org/openmarkov/gui/dialog/network/NetworkPropertiesDialog.java:130,213-218,269-272} \newline \path{gui/src/main/java/org/openmarkov/gui/dialog/network/NetworkDefinitionPanel.java:347-350} \newline \path{core/src/main/java/org/openmarkov/core/action/core/CycleLengthEdit.java:28-49}
\fichakey{Comprobado}Leyendo el código
\fichakey{Esfuerzo}Pequeño (horas)
\fichakey{Decisión de equipo}No
\fichakey{Relacionados}\hyperref[nota:66]{nota~66}
\end{ficha}
\paragraph{Qué pasa}
En una red temporal (las que no llevan la restricción \texttt{Only\allowbreak{}Atemporal\allowbreak{}Variables}: modelos de Markov con decisiones o MID ---diagramas de influencia de Markov---, POMDP ---procesos de decisión de Markov parcialmente observables---, redes bayesianas dinámicas, LIMID ---diagramas de influencia con memoria limitada--- dinámicos, DAN ---redes de análisis de decisiones--- de Markov\ldots), el diálogo «File → Network properties\ldots» tiene la pestaña «Temporal options», con un cuadro de texto para el valor de la duración del ciclo y un desplegable para la unidad (día, mes, año\ldots).

Escribir «6», elegir «months» y aceptar deja la duración del ciclo como estaba. Al reabrir el diálogo se ve el valor antiguo, al guardar se escribe el antiguo (\texttt{PGMXWriter\_\allowbreak{}0\_\allowbreak{}2.\allowbreak{}java:197-202}) y los cálculos que dependen de él (el ajuste del descuento por unidad de tiempo en \texttt{MIDTemporal\allowbreak{}Evolution.\allowbreak{}java:631-632}, el potencial \texttt{Cycle\allowbreak{}Length\allowbreak{}Shift}) usan el antiguo. No hay mensaje: el usuario cree que lo ha cambiado.

Ningún otro diálogo de la interfaz llama a \texttt{set\allowbreak{}Cycle\allowbreak{}Length}, así que la duración del ciclo de una red existente no se puede cambiar desde el programa.

\texttt{Network\allowbreak{}Temporal\allowbreak{}Options\allowbreak{}Panel} pone los dos controles y los rellena con lo que dice la red. Eso es todo lo que hace: no registra ningún oyente en ellos, los guarda en campos privados sin métodos de acceso y ningún otro código los lee.

Al pulsar «OK», \texttt{Network\allowbreak{}Properties\allowbreak{}Dialog.\allowbreak{}do\allowbreak{}Ok\allowbreak{}Click\allowbreak{}Before\allowbreak{}Hide} cierra la sub-historia de ediciones y devuelve \texttt{Network\allowbreak{}Definition\allowbreak{}Panel.\allowbreak{}check\allowbreak{}Name(\allowbreak{})}, que siempre devuelve \texttt{true}; no mira esa pestaña.

La edición necesaria existe (\texttt{Cycle\allowbreak{}Length\allowbreak{}Edit}, con su deshacer), pero no la crea ningún código de producción.

Además, si la red no tiene duración de ciclo, \texttt{initialize(\allowbreak{})} le asigna una nueva con \texttt{prob\allowbreak{}Net.\allowbreak{}set\allowbreak{}Cycle\allowbreak{}Length(\allowbreak{}.\allowbreak{}.\allowbreak{}.\allowbreak{})} directamente, fuera del sistema de ediciones, solo por abrir la pestaña.
\paragraph{El código}
gui/src/main/java/org/openmarkov/gui/dialog/network/NetworkTemporalOptionsPanel.java:55-95

\begin{Shaded}
\begin{Highlighting}[]
\KeywordTok{private} \DataTypeTok{void} \FunctionTok{initialize}\OperatorTok{()} \OperatorTok{\{}
    \FunctionTok{add}\OperatorTok{(}\FunctionTok{getUnitScale}\OperatorTok{());}
    \FunctionTok{add}\OperatorTok{(}\FunctionTok{getTemporalUnits}\OperatorTok{());}
\NormalTok{    CycleLength temporalUnit}\OperatorTok{;}
    \ControlFlowTok{if} \OperatorTok{(}\NormalTok{probNet}\OperatorTok{.}\FunctionTok{getCycleLength}\OperatorTok{()} \OperatorTok{!=} \KeywordTok{null}\OperatorTok{)} \OperatorTok{\{}
\NormalTok{        temporalUnit }\OperatorTok{=}\NormalTok{ probNet}\OperatorTok{.}\FunctionTok{getCycleLength}\OperatorTok{();}
    \OperatorTok{\}} \ControlFlowTok{else} \OperatorTok{\{}
\NormalTok{        temporalUnit }\OperatorTok{=} \KeywordTok{new} \FunctionTok{CycleLength}\OperatorTok{();}
\NormalTok{        probNet}\OperatorTok{.}\FunctionTok{setCycleLength}\OperatorTok{(}\NormalTok{temporalUnit}\OperatorTok{);}
    \OperatorTok{\}}
\NormalTok{    probNetUnit }\OperatorTok{=}\NormalTok{ temporalUnit}\OperatorTok{.}\FunctionTok{getUnit}\OperatorTok{();}
\NormalTok{    probNetScale }\OperatorTok{=}\NormalTok{ temporalUnit}\OperatorTok{.}\FunctionTok{getValue}\OperatorTok{();}
\NormalTok{    temporalUnits}\OperatorTok{.}\FunctionTok{setSelectedItem}\OperatorTok{(}
\NormalTok{            stringDatabase}\OperatorTok{.}\FunctionTok{getString}\OperatorTok{(}\StringTok{"NetworkTemporalOptionsPanel.TemporalOptions.Unit."} \OperatorTok{+}\NormalTok{ probNetUnit}\OperatorTok{.}\FunctionTok{toString}\OperatorTok{()));}
\NormalTok{    unitScale}\OperatorTok{.}\FunctionTok{setText}\OperatorTok{(}\BuiltInTok{String}\OperatorTok{.}\FunctionTok{valueOf}\OperatorTok{(}\NormalTok{probNetScale}\OperatorTok{));}
\OperatorTok{\}}
\KeywordTok{...}
\KeywordTok{private} \BuiltInTok{Component} \FunctionTok{getUnitScale}\OperatorTok{()} \OperatorTok{\{}
    \ControlFlowTok{if} \OperatorTok{(}\NormalTok{unitScale }\OperatorTok{==} \KeywordTok{null}\OperatorTok{)} \OperatorTok{\{}
\NormalTok{        unitScale }\OperatorTok{=} \KeywordTok{new} \BuiltInTok{JTextField}\OperatorTok{();}
    \OperatorTok{\}}
    \ControlFlowTok{return}\NormalTok{ unitScale}\OperatorTok{;}
\OperatorTok{\}}
\end{Highlighting}
\end{Shaded}

gui/src/main/java/org/openmarkov/gui/dialog/network/NetworkPropertiesDialog.java:269-272

\begin{Shaded}
\begin{Highlighting}[]
\AttributeTok{@Override} \KeywordTok{protected} \DataTypeTok{boolean} \FunctionTok{doOkClickBeforeHide}\OperatorTok{()} \OperatorTok{\{}
\NormalTok{    probNet}\OperatorTok{.}\FunctionTok{getPNESupport}\OperatorTok{().}\FunctionTok{closeSubEditHistory}\OperatorTok{();}
    \ControlFlowTok{return}\NormalTok{ NetworkDefinitionPanel}\OperatorTok{.}\FunctionTok{checkName}\OperatorTok{();}
\OperatorTok{\}}
\end{Highlighting}
\end{Shaded}
\paragraph{Cómo arreglarlo}
Hacer que el cambio llegue a la red mediante \texttt{Cycle\allowbreak{}Length\allowbreak{}Edit}, dentro de la sub-historia que el diálogo ya abre, para que «Cancel» lo revierta (\texttt{cancel\allowbreak{}Last\allowbreak{}Sub\allowbreak{}Edit\allowbreak{}History}) y deshacer funcione.

Dos opciones razonables:

\begin{itemize}
\tightlist
\item
  ejecutar la edición cuando cambia cualquiera de los dos controles (oyente de acción en el desplegable y de pérdida de foco o «Enter» en el cuadro de texto);
\item
  leer los dos controles en \texttt{do\allowbreak{}Ok\allowbreak{}Click\allowbreak{}Before\allowbreak{}Hide}.
\end{itemize}

En ambos casos hay que validar el número (positivo, bien escrito) y avisar si no lo es, en vez de lanzar una excepción. Conviene también que la asignación por defecto de \texttt{initialize(\allowbreak{})} no cambie la red fuera de las ediciones.

\texttt{Cycle\allowbreak{}Length\allowbreak{}Edit} es una de las ediciones de opciones de la red que ejecutan su cambio dos veces al rehacer (\hyperref[nota:66]{nota~66}); allí no hay nada que corregir mientras nadie la use, pero antes de usarla desde la interfaz hay que revisar su \texttt{redo(\allowbreak{})}.

Prueba de regresión: construir el panel sin pantalla sobre una red MID, cambiar el texto y la unidad, aceptar y comprobar \texttt{prob\allowbreak{}Net.\allowbreak{}get\allowbreak{}Cycle\allowbreak{}Length(\allowbreak{})}; y comprobar que cancelar y deshacer lo devuelven.

% ══════════════════════════════════════════════════════════════════════
\section{Coste-efectividad y análisis de sensibilidad}\label{sec:coste}

El análisis de coste-efectividad, el análisis de sensibilidad determinista (tornado, araña, mapa de dos parámetros) y el probabilístico, con sus ventanas de resultados.

\setcounter{subsection}{73}
\subsection[El jar publicado mezcla dos JFreeChart y muchas gráficas no se abren]{El jar ejecutable que se publica lleva mezcladas dos versiones de JFreeChart, y la evolución temporal, el coste-efectividad por decisión y casi todas las gráficas del análisis de sensibilidad fallan con NoSuchMethodError}\label{err:74}
\begin{ficha}
\fichakey{Estado}Presente
\fichakey{Módulo}full, inference.des
\fichakey{Ficheros}\path{inference.des/pom.xml:24-28} \newline \path{pom.xml:322-326} \newline \path{full/pom.xml:608-654} \newline \path{full/pom.xml:178-196} \newline \path{.github/workflows/maven.yml:51-57} \newline \path{sensitivityAnalysis/src/main/java/org/openmarkov/sensitivityanalysis/dialog/MapDialog.java:389-390} \newline \path{gui/src/main/java/org/openmarkov/gui/dialog/inference/temporalevolution/TraceTemporalEvolutionDialog.java:654,680,1056} \newline \path{costEffectiveness/src/main/java/org/openmarkov/costEffectiveness/CEDecisionResults.java:844,856} \newline \path{sensitivityAnalysis/src/main/java/org/openmarkov/sensitivityanalysis/dialog/TornadoSpiderDialog.java:450,458,623,631,652} \newline \path{sensitivityAnalysis/src/main/java/org/openmarkov/sensitivityanalysis/dialog/PlotDialog.java:294} \newline \path{sensitivityAnalysis/src/main/java/org/openmarkov/sensitivityanalysis/dialog/CEProbabilisticDialog.java:523,535,558,568,571} \newline \path{sensitivityAnalysis/src/main/java/org/openmarkov/sensitivityanalysis/dialog/CESpiderDialog.java:377,388,477,496}
\fichakey{Comprobado}Leyendo el código y ejecutándolo
\fichakey{Esfuerzo}Pequeño (horas)
\fichakey{Decisión de equipo}No
\fichakey{Corregir antes}\hyperref[err:83]{error~83}: arreglado solo este, el jar publicado pasa a fallar también en el mapa de ámbito de decisión.
\fichakey{Relacionados}\hyperref[err:15]{error~15}, \hyperref[err:67]{error~67}, \hyperref[err:80]{error~80}, \hyperref[err:81]{error~81}, \hyperref[err:82]{error~82}, \hyperref[err:83]{error~83}, \hyperref[err:84]{error~84}, \hyperref[err:86]{error~86}, \hyperref[err:88]{error~88}, \hyperref[err:89]{error~89}, \hyperref[err:90]{error~90}, \hyperref[err:91]{error~91}, \hyperref[nota:37]{nota~37}, \hyperref[nota:58]{nota~58}, \hyperref[nota:169]{nota~169}
\end{ficha}
\paragraph{Qué pasa}
Con el \texttt{Open\allowbreak{}Markov.\allowbreak{}jar} que se reparte, varias ventanas de resultados no llegan a abrirse. Es el jar que construye la integración continua en \texttt{.\allowbreak{}github/\allowbreak{}workflows/\allowbreak{}maven.\allowbreak{}yml} (paso «Build OpenMarkov.jar»: \texttt{mvn\ package\ -pl\ full\ -am\ -Dskip\allowbreak{}Tests\ -\/-batch-mode\ -PGenerate\allowbreak{}Full\allowbreak{}Jar}) y publica en GitHub Pages (\texttt{https:/\allowbreak{}/\allowbreak{}openmarkov.\allowbreak{}github.\allowbreak{}io/\allowbreak{}Open\allowbreak{}Markov/\allowbreak{}Open\allowbreak{}Markov.\allowbreak{}jar}). Al construir su gráfica, cada una de esas ventanas llama a un método de JFreeChart (la biblioteca de gráficas) que no existe en las clases que lleva el jar, y Java lanza \texttt{java.\allowbreak{}lang.\allowbreak{}No\allowbreak{}Such\allowbreak{}Method\allowbreak{}Error}. Como no es una \texttt{Runtime\allowbreak{}Exception}, el gestor global (\texttt{OMException\allowbreak{}Handler}) la trata como error previsto y \texttt{Exception\allowbreak{}Dialog} enseña un aviso titulado «Internal error», con «An internal error has happened:», el nombre del método y la pila de llamadas (esto último, por lectura).

Ventanas que fallan con ese jar:

\begin{itemize}
\tightlist
\item
  \textbf{Evolución temporal} (modo inferencia, clic derecho sobre un nodo → «Temporal Evolution» → aceptar, y también la de la red entera): falla siempre, por variable y por criterio (\texttt{Trace\allowbreak{}Temporal\allowbreak{}Evolution\allowbreak{}Dialog.\allowbreak{}get\allowbreak{}Charts\allowbreak{}Panel}, línea 1056, y \texttt{get\allowbreak{}Charts\allowbreak{}By\allowbreak{}Criterion\allowbreak{}Panel}, líneas 654 y 680).
\item
  \textbf{Menú Tools → «Cost-effectiveness analysis» con ámbito de decisión}: el constructor de \texttt{CEDecision\allowbreak{}Results} arma la pestaña del plano coste-efectividad, que falla en la línea 844, y la ventana no se abre. El ámbito global usa otra ventana (\texttt{CEPDialog}), que no está afectada.
\item
  \textbf{Menú Tools → «Sensitivity analysis»}: «Tornado/Spider» (la pestaña de araña falla siempre, \texttt{Tornado\allowbreak{}Spider\allowbreak{}Dialog} línea 652; la de tornado ya en la línea 450 con ámbito global), «Plot» (\texttt{Plot\allowbreak{}Dialog}, línea 294), «Map» con ámbito global (\texttt{Map\allowbreak{}Dialog}, línea 389), el plano y la curva de aceptabilidad del análisis probabilístico (\texttt{CEProbabilistic\allowbreak{}Dialog}, líneas 523 y 571) y «Spider CE» (\texttt{CESpider\allowbreak{}Dialog}, línea 377). El mapa con ámbito de decisión sí se abre.
\item
  \texttt{CEDESDialog}, de la simulación de sucesos discretos, tiene cinco llamadas de este tipo, pero solo se construye con el análisis probabilístico de la simulación, que hoy no se puede activar (\hyperref[nota:58]{nota~58}).
\end{itemize}

Mientras dure, en el jar no se llega a ver lo que describen \hyperref[err:67]{error~67} (evolución temporal), \hyperref[err:84]{error~84} (coste-efectividad por decisión), \hyperref[err:88]{error~88} y \hyperref[err:91]{error~91} (tornado), ni el mapa de ámbito global de \hyperref[err:82]{error~82}.

Por qué pasa. \texttt{inference.\allowbreak{}des/\allowbreak{}pom.\allowbreak{}xml} declara SSJ 3.3.1, la biblioteca de simulación estocástica de la que OpenMarkov solo usa los generadores de números aleatorios (\texttt{umontreal.\allowbreak{}ssj.\allowbreak{}rng}). El \texttt{pom} de SSJ declara a su vez, en ámbito de ejecución, \texttt{jfree:jfreechart:1.\allowbreak{}0.\allowbreak{}12}, que trae \texttt{jfree:jcommon:1.\allowbreak{}0.\allowbreak{}15}. El proyecto usa \texttt{org.\allowbreak{}jfree:jfreechart:1.\allowbreak{}5.\allowbreak{}6}, declarada en el \texttt{pom.\allowbreak{}xml} raíz. Entre esas versiones JFreeChart cambió de grupo (\texttt{jfree} → \texttt{org.\allowbreak{}jfree}), así que para Maven son dos bibliotecas distintas y no elige una: las dos entran en las dependencias de \texttt{full}.

El perfil \texttt{Generate\allowbreak{}Full\allowbreak{}Jar} de \texttt{full/\allowbreak{}pom.\allowbreak{}xml} usa \texttt{maven-assembly-plugin} con el descriptor predefinido \texttt{jar-with-dependencies}, que desempaqueta todas las dependencias en un solo jar y guarda una sola copia de cada fichero. Medido sobre un jar construido con ese perfil a partir del código actual y sobre el publicado (fechado el 8 de septiembre de 2026), comparando cada clase con las de los dos jar de JFreeChart:

\begin{itemize}
\tightlist
\item
  de las 692 clases de la 1.5.6, las 536 que también existen con el mismo nombre en la 1.0.12 están con la copia de la 1.0.12;
\item
  las 156 que solo tiene la 1.5.6 también están, por ejemplo \texttt{org.\allowbreak{}jfree.\allowbreak{}chart.\allowbreak{}ui.\allowbreak{}Rectangle\allowbreak{}Edge} (en la 1.0.12 esa clase era \texttt{org.\allowbreak{}jfree.\allowbreak{}ui.\allowbreak{}Rectangle\allowbreak{}Edge}, de JCommon, que también va dentro).
\end{itemize}

Como cada clase está una sola vez, no hay orden de carga que decida nada: con \texttt{java\ -jar} esas 536 clases son siempre las de la 1.0.12. El código de OpenMarkov, compilado contra la 1.5.6, encuentra las clases pero no algunos métodos: \texttt{Legend\allowbreak{}Title.\allowbreak{}set\allowbreak{}Position} existe en la 1.0.12 solo con un \texttt{org.\allowbreak{}jfree.\allowbreak{}ui.\allowbreak{}Rectangle\allowbreak{}Edge}, y \texttt{set\allowbreak{}Default\allowbreak{}Tool\allowbreak{}Tip\allowbreak{}Generator} se llamaba allí \texttt{set\allowbreak{}Base\allowbreak{}Tool\allowbreak{}Tip\allowbreak{}Generator}.

Medido:

\begin{itemize}
\tightlist
\item
  Recorriendo en las clases de OpenMarkov del jar todas las llamadas, accesos a campos y creaciones de objetos de bibliotecas de terceros (4 215), solo 8 métodos no existen dentro del jar, todos de JFreeChart, usados en 27 sitios (lista abajo). Contra la 1.5.6 sola no falta ninguno.
\item
  Construyendo las gráficas sin pantalla con los métodos reales: con las clases del jar, \texttt{Map\allowbreak{}Dialog.\allowbreak{}get\allowbreak{}Map\allowbreak{}Chart} en ámbito global lanza \texttt{No\allowbreak{}Such\allowbreak{}Method\allowbreak{}Error} en la línea 389; \texttt{Trace\allowbreak{}Temporal\allowbreak{}Evolution\allowbreak{}Dialog.\allowbreak{}get\allowbreak{}Charts\allowbreak{}Panel}, en la 1056; y \texttt{CEDecision\allowbreak{}Results}, con \texttt{0\_\allowbreak{}miscellaneous/\allowbreak{}networks/\allowbreak{}id/\allowbreak{}ID-arthronet-ce.\allowbreak{}pgmx}, construye la pestaña de análisis y falla en la del plano, línea 844. Con la 1.5.6 por delante de las clases del jar, el mapa global y la evolución temporal se construyen.
\end{itemize}

Desde el IDE, con el sistema de módulos de Java, no pasa. Los \texttt{module-info.\allowbreak{}java} piden el módulo \texttt{org.\allowbreak{}jfree.\allowbreak{}jfreechart}, que es el nombre que declara el manifiesto de la 1.5.6; nadie pide \texttt{jfreechart}, que sería el nombre de la 1.0.12. Las dos no pueden estar a la vez en la ruta de módulos: medido, al resolver los dos jar de JFreeChart junto con el de SSJ, Java se niega con \texttt{Resolution\allowbreak{}Exception:\ Modules\ org.\allowbreak{}jfree.\allowbreak{}jfreechart\ and\ jfreechart\ export\ package\ org.\allowbreak{}jfree.\allowbreak{}data.\allowbreak{}time\ to\ module\ ssj}, así que la aplicación ni arrancaría. Y si la 1.0.12 queda en la ruta de clases, un paquete que pertenece a un módulo con nombre se carga siempre del módulo. El IDE no se ha arrancado: esto sale de los descriptores y de esa medida.

Los instaladores del perfil \texttt{Installer}, que la integración continua no publica, tienen la misma mezcla por otro camino: copian todas las dependencias a \texttt{full/\allowbreak{}target/\allowbreak{}libs}, las dos versiones incluidas, y \texttt{jpackage} pone en la ruta de clases los jar de esa carpeta ordenados por nombre, con \texttt{jfreechart-1.\allowbreak{}0.\allowbreak{}12.\allowbreak{}jar} antes que \texttt{jfreechart-1.\allowbreak{}5.\allowbreak{}6.\allowbreak{}jar} (comprobado en las clases de \texttt{jpackage} de Java 25; no se ha construido ningún instalador).
\paragraph{El código}
inference.des/pom.xml:24-28

\begin{Shaded}
\begin{Highlighting}[]
\NormalTok{\textless{}}\KeywordTok{dependency}\NormalTok{\textgreater{}}
\NormalTok{    \textless{}}\KeywordTok{groupId}\NormalTok{\textgreater{}ca.umontreal.iro.simul\textless{}/}\KeywordTok{groupId}\NormalTok{\textgreater{}}
\NormalTok{    \textless{}}\KeywordTok{artifactId}\NormalTok{\textgreater{}ssj\textless{}/}\KeywordTok{artifactId}\NormalTok{\textgreater{}}
\NormalTok{    \textless{}}\KeywordTok{version}\NormalTok{\textgreater{}3.3.1\textless{}/}\KeywordTok{version}\NormalTok{\textgreater{}}
\NormalTok{\textless{}/}\KeywordTok{dependency}\NormalTok{\textgreater{}}
\end{Highlighting}
\end{Shaded}

full/pom.xml:621-626 (perfil \texttt{Generate\allowbreak{}Full\allowbreak{}Jar})

\begin{Shaded}
\begin{Highlighting}[]
\NormalTok{\textless{}}\KeywordTok{finalName}\NormalTok{\textgreater{}OpenMarkov\textless{}/}\KeywordTok{finalName}\NormalTok{\textgreater{}}
\NormalTok{\textless{}}\KeywordTok{appendAssemblyId}\NormalTok{\textgreater{}false\textless{}/}\KeywordTok{appendAssemblyId}\NormalTok{\textgreater{}}
\NormalTok{\textless{}}\KeywordTok{attach}\NormalTok{\textgreater{}false\textless{}/}\KeywordTok{attach}\NormalTok{\textgreater{}}
\NormalTok{\textless{}}\KeywordTok{descriptorRefs}\NormalTok{\textgreater{}}
\NormalTok{    \textless{}}\KeywordTok{descriptorRef}\NormalTok{\textgreater{}jar{-}with{-}dependencies\textless{}/}\KeywordTok{descriptorRef}\NormalTok{\textgreater{}}
\NormalTok{\textless{}/}\KeywordTok{descriptorRefs}\NormalTok{\textgreater{}}
\end{Highlighting}
\end{Shaded}

Métodos de JFreeChart 1.5.6 que no existen dentro de \texttt{Open\allowbreak{}Markov.\allowbreak{}jar}, y dónde se llaman:

\begin{Verbatim}[breaklines,breakanywhere,fontsize=\small]
LegendTitle.setPosition(RectangleEdge)          CEDecisionResults:844  CEProbabilisticDialog:523,571
                                                CESpiderDialog:377  TornadoSpiderDialog:652  CEDESDialog:495,545
PaintScaleLegend.setPosition(RectangleEdge)     MapDialog:389
PaintScaleLegend.setMargin(RectangleInsets)     MapDialog:390
ValueMarker.setLabelTextAnchor(TextAnchor)      PlotDialog:294  TornadoSpiderDialog:450,458,623,631
Marker.setLabelTextAnchor(TextAnchor)           CEProbabilisticDialog:558  CEDESDialog:532
XYItemRenderer.setDefaultToolTipGenerator       CEDecisionResults:856  CEProbabilisticDialog:535,568
                                                CEDESDialog:505,542
XYLineAndShapeRenderer.setDefaultToolTipGenerator
                                                TraceTemporalEvolutionDialog:654,680,1056  CESpiderDialog:388
XYSeriesCollection.getSeriesIndex(Comparable)   CESpiderDialog:477,496
\end{Verbatim}
\paragraph{Cómo arreglarlo}
Quitar la 1.0.12 del árbol de dependencias: en \texttt{inference.\allowbreak{}des/\allowbreak{}pom.\allowbreak{}xml}, añadir a la dependencia de SSJ una exclusión de \texttt{jfree:jfreechart}; con ella se va también \texttt{jfree:jcommon}, que solo llega por ahí. De SSJ, OpenMarkov solo usa \texttt{umontreal.\allowbreak{}ssj.\allowbreak{}rng.\allowbreak{}MRG31k3p} y \texttt{Random\allowbreak{}Stream} (en \texttt{DESRandom\allowbreak{}Provider} y \texttt{Not\allowbreak{}Random\allowbreak{}Stream}), y ninguna clase de ese paquete hace referencia a JFreeChart (comprobado en el jar de SSJ 3.3.1); la usa solo su paquete de gráficas, \texttt{umontreal.\allowbreak{}ssj.\allowbreak{}charts}.

Con la exclusión, la regla se cumple a la vez en los tres sitios que reúnen las dependencias: el jar del perfil \texttt{Generate\allowbreak{}Full\allowbreak{}Jar}, la carpeta \texttt{full/\allowbreak{}target/\allowbreak{}libs} de la que tiran los instaladores del perfil \texttt{Installer}, y la ruta de módulos del IDE y de las pruebas. Para que no vuelva a pasar con otra biblioteca que cambie de grupo, una opción es una comprobación en la construcción que falle cuando dos dependencias traen la misma clase (por ejemplo, la regla \texttt{ban\allowbreak{}Duplicate\allowbreak{}Classes} de \texttt{extra-enforcer-rules} para \texttt{maven-enforcer-plugin}).

Al quitar la 1.0.12, el jar pasa a comportarse como el IDE también en \hyperref[err:83]{error~83}, que hoy solo se ve desde el IDE.

Prueba: tras construir el jar, comprobar que no contiene \texttt{org/\allowbreak{}jfree/\allowbreak{}ui/\allowbreak{}} (JCommon) y que \texttt{org/\allowbreak{}jfree/\allowbreak{}chart/\allowbreak{}title/\allowbreak{}Legend\allowbreak{}Title.\allowbreak{}class} es la de la 1.5.6; y lanzar con \texttt{java\ -cp\ Open\allowbreak{}Markov.\allowbreak{}jar} un pequeño programa sin pantalla que construya la gráfica de \texttt{Trace\allowbreak{}Temporal\allowbreak{}Evolution\allowbreak{}Dialog.\allowbreak{}get\allowbreak{}Charts\allowbreak{}Panel} o la del mapa en ámbito global sin error. La integración continua puede ejecutarlo justo después del paso «Build OpenMarkov.jar».

\setcounter{subsection}{74}
\subsection[La araña de coste-efectividad reescribe la red abierta con las medias]{El diagrama de araña de coste-efectividad reescribe los números de la red abierta con la media de cada distribución}\label{err:75}
\begin{ficha}
\fichakey{Estado}Presente
\fichakey{Módulo}core
\fichakey{Paquete}org.openmarkov.core.model.network.modelUncertainty
\fichakey{Ficheros}\path{core/src/main/java/org/openmarkov/core/model/network/modelUncertainty/SystematicSampling.java:252-321,382-387} \newline \path{core/src/main/java/org/openmarkov/core/model/network/modelUncertainty/UncertainParameter.java:36-38} \newline \path{inference/src/main/java/org/openmarkov/inference/algorithm/variableElimination/tasks/VECESensAnSpider.java:59-86} \newline \path{sensitivityAnalysis/src/main/java/org/openmarkov/sensitivityanalysis/model/SensitivityAnalysisController.java:88-92} \newline \path{sensitivityAnalysis/src/main/java/org/openmarkov/sensitivityanalysis/dialog/CESpiderDialog.java:158-168}
\fichakey{Comprobado}Leyendo el código y ejecutándolo
\fichakey{Esfuerzo}Medio (días)
\fichakey{Decisión de equipo}No
\fichakey{Relacionados}\hyperref[err:35]{error~35}, \hyperref[err:49]{error~49}, \hyperref[nota:103]{nota~103}, \hyperref[nota:104]{nota~104}
\end{ficha}
\paragraph{Qué pasa}
Ejecutar el diagrama de araña de coste-efectividad cambia las tablas de la red abierta en el editor. Si una tabla no coincide con la media de la distribución de su parámetro incierto, al terminar la araña la tabla tiene la media. Si coinciden, queda un residuo de redondeo. Si el análisis lanza una excepción a mitad del barrido, la red se queda con el punto barrido, porque no hay \texttt{finally}. En todos los casos la red cambia sin pasar por una edición, así que no se puede deshacer ni se marca como modificada.

La ventana de incertidumbre escribe en la tabla las medias al aceptar (\texttt{Uncertain\allowbreak{}Values\allowbreak{}Dialog.\allowbreak{}calculate\allowbreak{}Reference\allowbreak{}Values}), así que en un modelo editado solo desde la aplicación la media y la tabla suelen coincidir y el usuario ve solo el residuo de redondeo; el cambio visible aparece con un fichero en el que no coinciden.

La red que llega al cálculo es la red abierta: \texttt{Sensitivity\allowbreak{}Analysis\allowbreak{}Controller} toma \texttt{get\allowbreak{}Current\allowbreak{}Network\allowbreak{}Editor\allowbreak{}Panel(\allowbreak{}).\allowbreak{}get\allowbreak{}Prob\allowbreak{}Net(\allowbreak{})} y \texttt{CESpider\allowbreak{}Dialog} la pasa a \texttt{new\ VECESens\allowbreak{}An\allowbreak{}Spider(\allowbreak{}prob\allowbreak{}Net)}, cuyas copias (\texttt{Inference\allowbreak{}Algorithm} y la línea 63) son superficiales.

\texttt{Systematic\allowbreak{}Sampling.\allowbreak{}network\allowbreak{}Sample(\allowbreak{}red,\allowbreak{}\ parámetro,\allowbreak{}\ min,\allowbreak{}\ max,\allowbreak{}\ num\allowbreak{}Intervalos,\allowbreak{}\ nombre,\allowbreak{}\ intervalo)} debería devolver una red en la que el parámetro incierto vale el punto \texttt{intervalo} de su barrido. Empieza con \texttt{Prob\allowbreak{}Net\ net\ =\ original\allowbreak{}Net.\allowbreak{}copy(\allowbreak{})}, pero \texttt{Prob\allowbreak{}Net.\allowbreak{}copy(\allowbreak{})} es una copia superficial que comparte los potenciales. Además, el método toma como tabla a modificar \texttt{original\allowbreak{}Sub\allowbreak{}Potential\allowbreak{}Table} (línea 282: \texttt{Table\allowbreak{}Potential\ new\allowbreak{}Sub\allowbreak{}Potential\allowbreak{}Table\ =\ original\allowbreak{}Sub\allowbreak{}Potential\allowbreak{}Table;}) y escribe la columna en su vector de valores (\texttt{copy\allowbreak{}In\allowbreak{}Array}, línea 307). Así que la red que recibe queda cambiada.

Su propio comentario (líneas 296-298) lo reconoce, y para devolver la red a su estado añade un caso especial: si \texttt{interval\ \textgreater{}\ num\allowbreak{}Intervals}, escribe \texttt{uncertain\allowbreak{}Parameter.\allowbreak{}get\allowbreak{}Base\allowbreak{}Line\allowbreak{}Value(\allowbreak{})}, que es la \textbf{media} de la distribución (\texttt{Uncertain\allowbreak{}Parameter.\allowbreak{}get\allowbreak{}Base\allowbreak{}Line\allowbreak{}Value} devuelve \texttt{get\allowbreak{}Prob\allowbreak{}Dens\allowbreak{}Function(\allowbreak{}).\allowbreak{}get\allowbreak{}Mean(\allowbreak{})}), no el número que tenía la tabla.

El único llamador de producción es \texttt{VECESens\allowbreak{}An\allowbreak{}Spider.\allowbreak{}resolve}, que recorre \texttt{for\ (\allowbreak{}int\ i\ =\ 0;\ i\ \textless{}=\ (\allowbreak{}iterations+1);\ ++i)} y descarta el resultado de la última vuelta (\texttt{if(\allowbreak{}i\textless{}=iterations)}). Esa vuelta solo sirve para reponer el valor, y cuesta un análisis de coste-efectividad completo más por parámetro.

\texttt{sample\allowbreak{}Network}, el método hermano que usan el tornado y la araña de un solo criterio, no tiene este problema: construye un potencial nuevo y lo sustituye en la copia. Comprobado sobre \texttt{0\_\allowbreak{}miscellaneous/\allowbreak{}networks/\allowbreak{}id/\allowbreak{}ID-decide-test.\allowbreak{}pgmx}: la red original queda idéntica.

Medido con pequeños programas de prueba (deterministas, una ejecución de cada uno) sobre \texttt{0\_\allowbreak{}miscellaneous/\allowbreak{}networks/\allowbreak{}id/\allowbreak{}ID-decide-test.\allowbreak{}pgmx} y \texttt{0\_\allowbreak{}miscellaneous/\allowbreak{}networks/\allowbreak{}id/\allowbreak{}ID-CEA-test-2therapies.\allowbreak{}pgmx}, con una variación de ±20 \% sobre el valor de referencia (la opción «Percentage over reference value», PORV en el código) y 4 intervalos:

\begin{itemize}
\tightlist
\item
  Una llamada a \texttt{network\allowbreak{}Sample} con \texttt{interval\ =\ 0} para «prevalence» cambia la tabla de \texttt{Disease} de la red recibida de {[}0,86, 0,14{]} a {[}0,888, 0,112{]}; para «specificity», la de \texttt{Result\ of\ test} pasa de 0,97/0,03 a 0,776/0,224.
\item
  Tras la llamada de reposición, los valores vuelven, pero con residuo de redondeo: después de la araña completa, la tabla de \texttt{Result\ of\ test} de la red recibida queda con 0,030000000000000023 y 0,08999999999999997 donde tenía 0,03 y 0,09.
\item
  Si la tabla no coincide con la media de su distribución, la araña la pisa: con la tabla de \texttt{Disease} puesta a {[}0,8, 0,2{]} y la distribución de «prevalence» Beta(14, 86) (media 0,14), tras ejecutar la araña la red recibida tiene {[}0,86, 0,14{]}.
\end{itemize}
\paragraph{El código}
core/src/main/java/org/openmarkov/core/model/network/modelUncertainty/SystematicSampling.java:282-308

\begin{Shaded}
\begin{Highlighting}[]
\NormalTok{TablePotential newSubPotentialTable }\OperatorTok{=}\NormalTok{ originalSubPotentialTable}\OperatorTok{;}
\KeywordTok{...}
\DataTypeTok{double}\NormalTok{ auxValueToAssign }\OperatorTok{=}\NormalTok{ min}\OperatorTok{+(}\NormalTok{pointsDistance}\OperatorTok{*}\NormalTok{interval}\OperatorTok{);}
\CommentTok{//Returns the net to its original potentials (after giving the last net with the last intervals potentials)}
\CommentTok{//for the next uncertainParameter. This is necessary because even if you operate on a copy of the original network}
\CommentTok{//with "net.copy()", the changes are transferred to the original network.}
\ControlFlowTok{if}\OperatorTok{(}\NormalTok{interval }\OperatorTok{\textgreater{}}\NormalTok{ numIntervals}\OperatorTok{)} \OperatorTok{\{}
\NormalTok{    auxValueToAssign }\OperatorTok{=}\NormalTok{ uncertainParameter}\OperatorTok{.}\FunctionTok{getBaseLineValue}\OperatorTok{();}
\OperatorTok{\}}
\DataTypeTok{double}\OperatorTok{[]}\NormalTok{ newSubpotentialTableValues }\OperatorTok{=}\NormalTok{ newSubPotentialTable}\OperatorTok{.}\FunctionTok{getValues}\OperatorTok{();}
\BuiltInTok{System}\OperatorTok{.}\FunctionTok{arraycopy}\OperatorTok{(}\NormalTok{sampledConfigurationValues}\OperatorTok{,} \DecValTok{0}\OperatorTok{,}\NormalTok{ auxSampledConfigurationValues}\OperatorTok{,} \DecValTok{0}\OperatorTok{,}\NormalTok{ numStates}\OperatorTok{);}
\FunctionTok{replaceValueAndRedistributeComplements}\OperatorTok{(}\NormalTok{auxSampledConfigurationValues}\OperatorTok{,}\NormalTok{ sampler}\OperatorTok{,}\NormalTok{ posUncertainInColumn}\OperatorTok{,}
\NormalTok{            auxValueToAssign}\OperatorTok{);}
\FunctionTok{copyInArray}\OperatorTok{(}\NormalTok{newSubpotentialTableValues}\OperatorTok{,}\NormalTok{ configurationBasePositionInitColumn}\OperatorTok{,}
\NormalTok{            auxSampledConfigurationValues}\OperatorTok{);}
\end{Highlighting}
\end{Shaded}
\paragraph{Cómo arreglarlo}
\texttt{network\allowbreak{}Sample} debe construir un potencial nuevo (con una copia del vector de valores) y sustituirlo en la copia de la red, como ya hace \texttt{sample\allowbreak{}Network} (líneas 242-243: \texttt{net.\allowbreak{}remove\allowbreak{}Potential(\allowbreak{}original\allowbreak{}Potential);\ net.\allowbreak{}add\allowbreak{}Potential(\allowbreak{}new\allowbreak{}Potential);}), en vez de escribir en \texttt{original\allowbreak{}Sub\allowbreak{}Potential\allowbreak{}Table}.

Con eso sobran el caso \texttt{interval\ \textgreater{}\ num\allowbreak{}Intervals} y la vuelta extra de \texttt{VECESens\allowbreak{}An\allowbreak{}Spider.\allowbreak{}resolve}: \texttt{i\ \textless{}=\ iterations+1} pasa a \texttt{i\ \textless{}=\ iterations}, y desaparece el \texttt{if(\allowbreak{}i\textless{}=iterations)}.

El arreglo tiene que estar en \texttt{network\allowbreak{}Sample}: la copia superficial de la red seguirá compartiendo las tablas.

Pruebas:

\begin{itemize}
\tightlist
\item
  Tras \texttt{VECESens\allowbreak{}An\allowbreak{}Spider.\allowbreak{}get\allowbreak{}CEPotentials(\allowbreak{})} sobre una red cuya tabla no coincide con la media de la distribución del parámetro, los vectores de valores de todos los potenciales de la red recibida deben ser idénticos (con \texttt{Arrays.\allowbreak{}equals}) a los de antes.
\item
  Una llamada suelta a \texttt{network\allowbreak{}Sample} no debe modificar la red que recibe.
\end{itemize}

\setcounter{subsection}{75}
\subsection[Una columna con incertidumbre solo en algunas casillas se sortea mal]{En el análisis probabilístico, una columna con incertidumbre en solo algunas casillas no se sortea o pierde sus números fijos}\label{err:76}
\begin{ficha}
\fichakey{Estado}Presente
\fichakey{Módulo}core
\fichakey{Paquete}org.openmarkov.core.model.network.modelUncertainty
\fichakey{Ficheros}\path{core/src/main/java/org/openmarkov/core/model/network/modelUncertainty/TablePotentialSampler.java:60-77} \newline \path{core/src/main/java/org/openmarkov/core/model/network/modelUncertainty/Sampler.java:56-82,126-179} \newline \path{io/src/main/java/org/openmarkov/io/probmodel/reader/PGMXReader_0_2.java:1324-1356} \newline \path{inference/src/main/java/org/openmarkov/inference/algorithm/variableElimination/tasks/VECEPSA.java:105-113} \newline \path{sensitivityAnalysis/src/main/java/org/openmarkov/sensitivityanalysis/dialog/CEProbabilisticDialog.java:130-142}
\fichakey{Comprobado}Leyendo el código y ejecutándolo
\fichakey{Esfuerzo}Pequeño (horas)
\fichakey{Decisión de equipo}\textcolor{decisionrojo}{\textbf{Sí.}} Qué significa una columna de probabilidad en la que unas casillas llevan distribución y otras un número fijo (prohibirla al leer y al editar, o sortear las inciertas conservando las fijas)
\fichakey{Relacionados}\hyperref[err:77]{error~77}, \hyperref[err:116]{error~116}, \hyperref[err:167]{error~167}, \hyperref[nota:58]{nota~58}, \hyperref[nota:110]{nota~110}
\end{ficha}
\paragraph{Qué pasa}
Una tabla de probabilidad con incertidumbre lleva en sus casillas, además del número, una distribución (Beta, Dirichlet, «Complement»\ldots). Si en una columna (una configuración de los padres) solo algunas casillas llevan distribución, el análisis probabilístico de sensibilidad de coste-efectividad la sortea mal, sin ningún aviso:

\begin{itemize}
\tightlist
\item
  si la casilla del primer estado no tiene distribución, la columna no se sortea nunca y las distribuciones de las demás casillas se ignoran;
\item
  si la tiene, la columna se sortea, pero las casillas sin distribución quedan a cero y su número fijo se pierde.
\end{itemize}

Camino de llegada: el análisis probabilístico de sensibilidad de coste-efectividad (\texttt{CEProbabilistic\allowbreak{}Dialog}, que construye \texttt{VECEPSA} y en cada simulación llama a \texttt{potential.\allowbreak{}sample(\allowbreak{})} sobre todos los potenciales de la red). También \texttt{Tree\allowbreak{}ADDPotential.\allowbreak{}sample} y \texttt{Exact\allowbreak{}Distr\allowbreak{}Potential.\allowbreak{}sample} delegan en este código.

La ventana de incertidumbre (\texttt{Uncertain\allowbreak{}Values\allowbreak{}Dialog.\allowbreak{}read\allowbreak{}Data\allowbreak{}From\allowbreak{}Table}) siempre escribe una distribución en todas las casillas de la columna, así que desde la edición no se produce una columna a medias. Sí se produce leyendo un \texttt{.\allowbreak{}pgmx}: el formato admite \texttt{\textless{}Value\ /\allowbreak{}\textgreater{}} vacío en cualquier casilla de \texttt{Uncertain\allowbreak{}Values}, y \texttt{PGMXReader\_\allowbreak{}0\_\allowbreak{}2.\allowbreak{}get\allowbreak{}Uncertain\allowbreak{}Value} lo convierte en \texttt{null}.

Se leyeron los 77 \texttt{.\allowbreak{}pgmx} del repositorio que contienen \texttt{Uncertain\allowbreak{}Values} (bajo \texttt{0\_\allowbreak{}miscellaneous/\allowbreak{}networks} y los \texttt{src/\allowbreak{}*/\allowbreak{}resources} de los módulos): 63 se leen, 51 tienen incertidumbre y ninguno tiene una columna a medias (las columnas son o todas nulas o todas con distribución). Para llegar al fallo hace falta un fichero hecho a mano o por otra herramienta.

\texttt{Table\allowbreak{}Potential\allowbreak{}Sampler.\allowbreak{}sample} recorre la tabla columna a columna y decide si sortearla con \texttt{has\allowbreak{}Uncertainty\ =\ uncertain\allowbreak{}Values.\allowbreak{}get(\allowbreak{}0)\ !=\ null}, es decir, mirando solo la casilla del primer estado. Si esa casilla no tiene distribución, la columna entera se copia tal cual. Si la tiene, \texttt{generate\allowbreak{}Sample} reparte las casillas entre la familia de distribuciones independientes, la de Dirichlet y la de complemento; las casillas sin distribución no están en ninguna (los métodos de \texttt{Sampler} saltan los nulos), así que quedan a cero en la columna sorteada.

Medido con un pequeño programa de prueba (5 sorteos, con el generador del propio programa, que no admite semilla), sobre una tabla P(X \textbar{} Y) con X de tres estados:

\begin{itemize}
\tightlist
\item
  Columna y0 = {[}0,2 fijo, Beta(2,5), Complement{]}: los cinco sorteos devuelven {[}0,2 0,5 0,3{]}, la columna original, sin sortear nunca.
\item
  Columna y1 = {[}Beta(2,5), 0,3 fijo, Complement{]}: los cinco sorteos devuelven la casilla fija a 0,0000 (por ejemplo {[}0,4084 0,0000 0,5916{]}); la columna suma 1 pero ya no dice que P(x1 \textbar{} y1) = 0,3.
\end{itemize}
\paragraph{El código}
core/src/main/java/org/openmarkov/core/model/network/modelUncertainty/TablePotentialSampler.java:60-76

\begin{Shaded}
\begin{Highlighting}[]
\ControlFlowTok{for} \OperatorTok{(}\DataTypeTok{int}\NormalTok{ configurationIndex }\OperatorTok{=} \DecValTok{0}\OperatorTok{;}\NormalTok{ configurationIndex }\OperatorTok{\textless{}}\NormalTok{ numConfigurations}\OperatorTok{;}\NormalTok{ configurationIndex}\OperatorTok{++)} \OperatorTok{\{}
    \DataTypeTok{int}\NormalTok{ configurationBasePosition }\OperatorTok{=}\NormalTok{ numStates }\OperatorTok{*}\NormalTok{ configurationIndex}\OperatorTok{;}
\NormalTok{    uncertainValues }\OperatorTok{=} \FunctionTok{getUncertainValuesChance}\OperatorTok{(}\NormalTok{uTable}\OperatorTok{,}\NormalTok{ configurationBasePosition}\OperatorTok{,}\NormalTok{ numStates}\OperatorTok{);}
\NormalTok{    hasUncertainty }\OperatorTok{=}\NormalTok{ uncertainValues}\OperatorTok{.}\FunctionTok{get}\OperatorTok{(}\DecValTok{0}\OperatorTok{)} \OperatorTok{!=} \KeywordTok{null}\OperatorTok{;}
    \ControlFlowTok{if} \OperatorTok{(}\NormalTok{hasUncertainty}\OperatorTok{)} \OperatorTok{\{}
\NormalTok{        sampledConfigurationValues }\OperatorTok{=} \FunctionTok{generateSample}\OperatorTok{(}\NormalTok{uncertainValues}\OperatorTok{,}\NormalTok{ numStates}\OperatorTok{,}\NormalTok{ functionTypes}\OperatorTok{);}
        \FunctionTok{copyInArray}\OperatorTok{(}\NormalTok{sampledValues}\OperatorTok{,}\NormalTok{ configurationBasePosition}\OperatorTok{,}\NormalTok{ sampledConfigurationValues}\OperatorTok{);}
    \OperatorTok{\}} \ControlFlowTok{else} \OperatorTok{\{}
        \ControlFlowTok{for} \OperatorTok{(}\DataTypeTok{int}\NormalTok{ stateIndex }\OperatorTok{=} \DecValTok{0}\OperatorTok{;}\NormalTok{ stateIndex }\OperatorTok{\textless{}}\NormalTok{ numStates}\OperatorTok{;}\NormalTok{ stateIndex}\OperatorTok{++)} \OperatorTok{\{}
\NormalTok{            sampledValues}\OperatorTok{[}\NormalTok{configurationBasePosition }\OperatorTok{+}\NormalTok{ stateIndex}\OperatorTok{]} \OperatorTok{=}\NormalTok{ originalValues}\OperatorTok{[}\NormalTok{configurationBasePosition}
                    \OperatorTok{+}\NormalTok{ stateIndex}\OperatorTok{];}
        \OperatorTok{\}}
    \OperatorTok{\}}
\OperatorTok{\}}
\end{Highlighting}
\end{Shaded}
\paragraph{Cómo arreglarlo}
La columna tiene incertidumbre si alguna de sus casillas la tiene.

Qué hacer con las casillas fijas de una columna incierta es la decisión de equipo. Una opción es rechazar la columna a medias al leer el fichero (y en \texttt{Uncertain\allowbreak{}Values\allowbreak{}Edit}); otra es tratar cada casilla fija como un valor exacto que resta masa al complemento.

La misma pregunta, «¿tiene incertidumbre esta columna?», se decide también en la tabla de la ventana (\texttt{Table\allowbreak{}Potential.\allowbreak{}has\allowbreak{}Uncertainty}, que mira solo la primera casilla: ver \hyperref[err:167]{error~167}) y en \texttt{Uncertain\allowbreak{}Values\allowbreak{}Dialog.\allowbreak{}has\allowbreak{}Uncertain\allowbreak{}Values}. Conviene que los tres sitios usen un único método.

Prueba: la tabla descrita arriba, sorteada, debe sortear la columna y0 y no poner a cero la casilla fija de y1 (o el lector debe rechazarla, según lo que se decida).

\setcounter{subsection}{76}
\subsection[El análisis probabilístico de sensibilidad no se puede repetir]{El análisis probabilístico de sensibilidad no se puede repetir: cada columna se sortea con un generador nuevo sembrado con el reloj}\label{err:77}
\begin{ficha}
\fichakey{Estado}Presente
\fichakey{Módulo}core
\fichakey{Paquete}org.openmarkov.core.model.network.modelUncertainty
\fichakey{Ficheros}\path{core/src/main/java/org/openmarkov/core/model/network/modelUncertainty/Sampler.java:162-179} \newline \path{core/src/main/java/org/openmarkov/core/model/network/modelUncertainty/TablePotentialSampler.java:85-88} \newline \path{core/src/main/java/org/openmarkov/core/model/network/modelUncertainty/XORShiftRandom.java} \newline \path{core/src/main/java/org/openmarkov/core/model/network/potential/TablePotential.java:566-575} \newline \path{inference/src/main/java/org/openmarkov/inference/algorithm/variableElimination/tasks/VECEPSA.java:105-113} \newline \path{sensitivityAnalysis/src/main/java/org/openmarkov/sensitivityanalysis/dialog/CEProbabilisticDialog.java:130-142}
\fichakey{Comprobado}Leyendo el código y ejecutándolo
\fichakey{Esfuerzo}Medio (días)
\fichakey{Decisión de equipo}No
\fichakey{Relacionados}\hyperref[err:76]{error~76}, \hyperref[err:136]{error~136}, \hyperref[nota:58]{nota~58}, \hyperref[nota:103]{nota~103}, \hyperref[nota:104]{nota~104}, \hyperref[nota:110]{nota~110}
\end{ficha}
\paragraph{Qué pasa}
Dos ejecuciones del análisis probabilístico de coste-efectividad sobre el mismo modelo nunca dan los mismos resultados, y no hay forma de volver a obtener una ejecución concreta ni de escribir una prueba de regresión con números. Además, las columnas sorteadas una tras otra salen ligeramente correlacionadas, aunque el modelo declare sus parámetros independientes.

Camino de llegada: menú de análisis de sensibilidad, análisis probabilístico sobre una red de coste-efectividad con parámetros inciertos (\texttt{CEProbabilistic\allowbreak{}Dialog} → \texttt{VECEPSA.\allowbreak{}resolve} → \texttt{sample\allowbreak{}Network\allowbreak{}Potentials} → \texttt{Table\allowbreak{}Potential.\allowbreak{}sample}).

\texttt{Sampler.\allowbreak{}generate\allowbreak{}Sample} pide el generador de números aleatorios dentro del método (\texttt{Random\ random\allowbreak{}Generator\ =\ create\allowbreak{}Random\allowbreak{}Generator(\allowbreak{});}), y \texttt{Table\allowbreak{}Potential\allowbreak{}Sampler.\allowbreak{}create\allowbreak{}Random\allowbreak{}Generator} devuelve \texttt{new\ XORShift\allowbreak{}Random(\allowbreak{})}, cuyo estado inicial es \texttt{System.\allowbreak{}nano\allowbreak{}Time(\allowbreak{})}. Es decir: cada columna de cada tabla de cada simulación usa un generador nuevo sembrado con el reloj.

\texttt{Table\allowbreak{}Potential.\allowbreak{}sample(\allowbreak{})} no recibe generador ni semilla, \texttt{VECEPSA} no tiene ningún método para fijarla y \texttt{CEProbabilistic\allowbreak{}Dialog} no la pide. La propagación estocástica (\texttt{Stochastic\allowbreak{}Propagation}) sí admite semilla, así que el proyecto ya tiene el patrón.

La correlación viene de sembrar con el reloj generadores que se crean con nanosegundos de diferencia: XORShift es lineal, y estados iniciales cercanos dan salidas emparentadas.

Medido con pequeños programas de prueba:

\begin{itemize}
\tightlist
\item
  Repetición: 3 ejecuciones de 3 sorteos de una Beta(3,7); las nueve extracciones son distintas entre ejecuciones.
\item
  Correlación: una tabla de cuatro columnas, cada una Beta(3,7) más complemento, 200.000 sorteos del potencial entero, repetido 3 veces. Las medias (0,299-0,301) y desviaciones (0,137-0,138) coinciden con las de la Beta, pero la correlación entre columnas contiguas sale entre 0,011 y 0,024 en las nueve medidas (el error típico con independencia es 0,0022, así que son de 5 a 11 errores típicos, siempre positivas), y 0,009-0,011 entre la primera y la cuarta.
\item
  Control: la misma Beta extraída de un único generador sembrado una vez da correlaciones de 0,0007, 0,0009 y -0,0031.
\end{itemize}

Es un efecto pequeño, pero introduce correlación entre parámetros que el modelo declara independientes.
\paragraph{El código}
core/src/main/java/org/openmarkov/core/model/network/modelUncertainty/Sampler.java:162-168

\begin{Shaded}
\begin{Highlighting}[]
\KeywordTok{protected} \DataTypeTok{double}\OperatorTok{[]} \FunctionTok{generateSample}\OperatorTok{(}\NormalTok{FamilyDistribution otherFamily}\OperatorTok{,}\NormalTok{ DirichletFamily dirFamily}\OperatorTok{,}
\NormalTok{        ComplementFamily complementFamily}\OperatorTok{,} \DataTypeTok{int}\OperatorTok{[]}\NormalTok{ indexesOther}\OperatorTok{,} \DataTypeTok{int}\OperatorTok{[]}\NormalTok{ indexesDirichlet}\OperatorTok{,} \DataTypeTok{int}\OperatorTok{[]}\NormalTok{ indexesComplement}\OperatorTok{,}
        \DataTypeTok{int}\NormalTok{ numStates}\OperatorTok{)} \OperatorTok{\{}
    \BuiltInTok{Random}\NormalTok{ randomGenerator }\OperatorTok{=} \FunctionTok{createRandomGenerator}\OperatorTok{();}
    \DataTypeTok{double}\OperatorTok{[]}\NormalTok{ sampledConfigurationValues }\OperatorTok{=} \KeywordTok{new} \DataTypeTok{double}\OperatorTok{[}\NormalTok{numStates}\OperatorTok{];}
\end{Highlighting}
\end{Shaded}

core/src/main/java/org/openmarkov/core/model/network/modelUncertainty/TablePotentialSampler.java:85-88

\begin{Shaded}
\begin{Highlighting}[]
\AttributeTok{@Override}
\KeywordTok{protected} \BuiltInTok{Random} \FunctionTok{createRandomGenerator}\OperatorTok{()} \OperatorTok{\{}
    \ControlFlowTok{return} \KeywordTok{new} \FunctionTok{XORShiftRandom}\OperatorTok{();}
\OperatorTok{\}}
\end{Highlighting}
\end{Shaded}
\paragraph{Cómo arreglarlo}
Un único generador por análisis, creado una vez con una semilla que se pueda fijar (y mostrar), pasado hacia abajo:

\begin{itemize}
\tightlist
\item
  \texttt{Potential.\allowbreak{}sample(\allowbreak{}Random)} en lugar de \texttt{sample(\allowbreak{})}, y sus redefiniciones en \texttt{Table\allowbreak{}Potential}, \texttt{Exact\allowbreak{}Distr\allowbreak{}Potential}, \texttt{Tree\allowbreak{}ADDPotential} y cualquier otra subclase que sortee;
\item
  \texttt{Sampler.\allowbreak{}generate\allowbreak{}Sample} recibiéndolo como parámetro.
\end{itemize}

\texttt{Systematic\allowbreak{}Sampling.\allowbreak{}create\allowbreak{}Random\allowbreak{}Generator} devuelve hoy \texttt{null} y no lo usa, así que el cambio de firma no le afecta.

Si se usan varios hilos (\texttt{VECEPSA.\allowbreak{}set\allowbreak{}Use\allowbreak{}Multithreading}), cada hilo necesita su propio generador derivado de la semilla, además de su propia red; es el problema de varios hilos que describe \hyperref[nota:103]{nota~103}, donde hoy no hay nada que corregir.

Pruebas:

\begin{itemize}
\tightlist
\item
  Dos análisis con la misma semilla devuelven exactamente los mismos resultados.
\item
  La correlación entre dos columnas independientes sorteadas muchas veces queda dentro de tres errores típicos de cero.
\end{itemize}

\setcounter{subsection}{77}
\subsection[El diálogo de incertidumbre rechaza toda Dirichlet sobre dos o más estados]{El diálogo de incertidumbre rechaza cualquier Dirichlet sobre dos o más estados, incluidas las que traen las redes del repositorio}\label{err:78}
\begin{ficha}
\fichakey{Estado}Presente
\fichakey{Módulo}gui
\fichakey{Paquete}org.openmarkov.gui.dialog.node
\fichakey{Ficheros}\path{gui/src/main/java/org/openmarkov/gui/dialog/node/UncertainValuesDialog.java:521-549, 578-583, 649-674, 676-684} \newline \path{gui/src/main/resources/gui/localize/gui_class_localizations_en.xml:53-56} \newline \path{gui/src/main/java/org/openmarkov/gui/dialog/common/TablePotentialPanel.java:701-724} \newline \path{gui/src/main/java/org/openmarkov/gui/dialog/common/OkCancelDialog.java:68-79}
\fichakey{Comprobado}Leyendo el código y ejecutándolo
\fichakey{Esfuerzo}Pequeño (horas)
\fichakey{Decisión de equipo}No
\fichakey{Relacionados}\hyperref[err:79]{error~79}
\end{ficha}
\paragraph{Qué pasa}
En OpenMarkov una distribución de Dirichlet sobre una columna de una tabla de probabilidad se escribe poniendo en cada casilla una \texttt{Dirichlet\allowbreak{}Function} con su parámetro alfa: una Dirichlet sobre tres estados son tres casillas Dirichlet. Así la escriben las redes del repositorio: \texttt{0\_\allowbreak{}miscellaneous/\allowbreak{}networks/\allowbreak{}mid/\allowbreak{}MID-Chancellor-corrected.\allowbreak{}pgmx} tiene columnas \texttt{Dirichlet\ 15,\allowbreak{}\ Dirichlet\ 512,\allowbreak{}\ Dirichlet\ 731,\allowbreak{}\ Exact\ 0} y \texttt{Dirichlet\ 17,\allowbreak{}\ 116,\allowbreak{}\ 350,\allowbreak{}\ 1251}.

El diálogo de incertidumbre no acepta ninguna de esas columnas. Camino del usuario: abrir \texttt{0\_\allowbreak{}miscellaneous/\allowbreak{}networks/\allowbreak{}mid/\allowbreak{}MID-Chancellor-corrected.\allowbreak{}pgmx}, abrir la tabla del nodo que lleva esa incertidumbre, clic derecho sobre una de esas columnas → «Edit uncertainty» y pulsar «OK» sin tocar nada. \texttt{Ok\allowbreak{}Cancel\allowbreak{}Dialog} envuelve la excepción en \texttt{Unrecoverable\allowbreak{}Exception}, el gestor global muestra el mensaje de la regla 3 y el diálogo no se cierra; la única salida es «Cancel». Tampoco se puede crear desde cero una Dirichlet sobre varios estados con «Assign uncertainty».

Al pulsar «OK», \texttt{Uncertain\allowbreak{}Values\allowbreak{}Dialog.\allowbreak{}do\allowbreak{}Ok\allowbreak{}Click\allowbreak{}Before\allowbreak{}Hide} comprueba las reglas de familia 1, 2 y 3. La regla 3 cuenta las casillas Dirichlet: con cero no hace nada; con una exige que todas las demás sean exactas con valor 0; con dos o más \textbf{lanza siempre} \texttt{Rule3Broken}.

El mensaje que ve el usuario dice lo contrario: «If one of the distributions is a Dirichlet, then: all the others must be Exact with v = 0 or Dirichlet. At least one of the others must also be a Dirichlet». Es decir, el mensaje describe justamente la columna que el código rechaza, y el código solo acepta una Dirichlet de una sola casilla (que siempre vale 1 y no tiene incertidumbre).

Medido con un pequeño programa de prueba que llama a la verificación real: 3 Dirichlet → \texttt{Rule3Broken}; 2 Dirichlet → \texttt{Rule3Broken}; 2 Dirichlet + exacta a 0 → \texttt{Rule3Broken}; 1 Dirichlet + 2 exactas a 0 → aceptada; beta + complemento → aceptada.

En la misma clase, \texttt{do\allowbreak{}Verify\allowbreak{}Rule4} (rechaza una columna formada solo por complementos) no la llama nadie, cosa que el propio código anota con \texttt{@To\allowbreak{}Check}; no se ha comprobado qué consecuencia tiene aceptar esa columna.
\paragraph{El código}
gui/src/main/java/org/openmarkov/gui/dialog/node/UncertainValuesDialog.java:649-674

\begin{Shaded}
\begin{Highlighting}[]
\KeywordTok{private} \DataTypeTok{static} \DataTypeTok{boolean} \FunctionTok{doVerifyRule3}\OperatorTok{(}\NormalTok{FamilyDistribution family}\OperatorTok{)} \KeywordTok{throws}\NormalTok{ FamilyDistributionRuleBrokenException}\OperatorTok{.}\FunctionTok{Rule3Broken} \OperatorTok{\{}
    \BuiltInTok{List}\OperatorTok{\textless{}}\NormalTok{UncertainValue}\OperatorTok{\textgreater{}}\NormalTok{ uncertainFamily }\OperatorTok{=}\NormalTok{ family}\OperatorTok{.}\FunctionTok{getFamily}\OperatorTok{();}
    \DataTypeTok{int}\NormalTok{ totalSizeFamily }\OperatorTok{=}\NormalTok{ uncertainFamily}\OperatorTok{.}\FunctionTok{size}\OperatorTok{();}
    \BuiltInTok{List}\OperatorTok{\textless{}}\NormalTok{UncertainValue}\OperatorTok{\textgreater{}}\NormalTok{ dirUncertain }\OperatorTok{=} \FunctionTok{getUncertainValuesOfClass}\OperatorTok{(}\NormalTok{uncertainFamily}\OperatorTok{,}\NormalTok{ DirichletFunction}\OperatorTok{.}\FunctionTok{class}\OperatorTok{);}
    \DataTypeTok{int}\NormalTok{ numDirichlet }\OperatorTok{=}\NormalTok{ dirUncertain}\OperatorTok{.}\FunctionTok{size}\OperatorTok{();}
    \ControlFlowTok{switch} \OperatorTok{(}\NormalTok{numDirichlet}\OperatorTok{)} \OperatorTok{\{}
        \ControlFlowTok{case} \DecValTok{0} \OperatorTok{{-}\textgreater{}} \OperatorTok{\{}
        \OperatorTok{\}}
        \ControlFlowTok{case} \DecValTok{1} \OperatorTok{{-}\textgreater{}} \OperatorTok{\{}
            \BuiltInTok{List}\OperatorTok{\textless{}}\NormalTok{UncertainValue}\OperatorTok{\textgreater{}}\NormalTok{ exactUncertain }\OperatorTok{=} \FunctionTok{getUncertainValuesOfClass}\OperatorTok{(}\NormalTok{uncertainFamily}\OperatorTok{,}\NormalTok{ ExactFunction}\OperatorTok{.}\FunctionTok{class}\OperatorTok{);}
            \DataTypeTok{int}\NormalTok{ numExact }\OperatorTok{=}\NormalTok{ exactUncertain}\OperatorTok{.}\FunctionTok{size}\OperatorTok{();}
            \DataTypeTok{boolean}\NormalTok{ verify }\OperatorTok{=} \OperatorTok{(}
                    \OperatorTok{(}\NormalTok{numExact }\OperatorTok{+}\NormalTok{ numDirichlet }\OperatorTok{==}\NormalTok{ totalSizeFamily}\OperatorTok{)} \OperatorTok{\&\&} \FunctionTok{areAllZero}\OperatorTok{(}
                            \KeywordTok{new} \FunctionTok{FamilyDistribution}\OperatorTok{(}\NormalTok{exactUncertain}\OperatorTok{).}\FunctionTok{getMean}\OperatorTok{())}
            \OperatorTok{);}
            \ControlFlowTok{if} \OperatorTok{(!}\NormalTok{verify}\OperatorTok{)} \OperatorTok{\{}
                \ControlFlowTok{throw} \KeywordTok{new}\NormalTok{ FamilyDistributionRuleBrokenException}\OperatorTok{.}\FunctionTok{Rule3Broken}\OperatorTok{(}\NormalTok{family}\OperatorTok{);}
            \OperatorTok{\}}
        \OperatorTok{\}}
        \KeywordTok{default} \OperatorTok{{-}\textgreater{}} \OperatorTok{\{}
            \ControlFlowTok{throw} \KeywordTok{new}\NormalTok{ FamilyDistributionRuleBrokenException}\OperatorTok{.}\FunctionTok{Rule3Broken}\OperatorTok{(}\NormalTok{family}\OperatorTok{);}
        \OperatorTok{\}}
    \OperatorTok{\}}
    \ControlFlowTok{return} \KeywordTok{true}\OperatorTok{;}
\OperatorTok{\}}
\end{Highlighting}
\end{Shaded}
\paragraph{Cómo arreglarlo}
Reescribir la regla 3 según su mensaje: si hay alguna casilla Dirichlet, todas las demás deben ser Dirichlet o exactas con valor 0, y debe haber al menos dos Dirichlet (\texttt{num\allowbreak{}Dirichlet\ \textgreater{}=\ 2\ \&\&\ num\allowbreak{}Dirichlet\ +\ num\allowbreak{}Exactas\allowbreak{}ACero\ ==\ total}). El caso de una sola Dirichlet pasa a ser inválido, como dice el texto. El mensaje en \texttt{gui\_\allowbreak{}class\_\allowbreak{}localizations\_\allowbreak{}en.\allowbreak{}xml:53-54} ya describe esa regla.

Si las reglas de familia se mueven a \texttt{core} para aplicarlas también al leer o al muestrear (ver \hyperref[err:79]{error~79}), la regla corregida es la que debe moverse.

Queda por ver, de paso, si \texttt{do\allowbreak{}Verify\allowbreak{}Rule4} debe llamarse o borrarse.

Pruebas de regresión:

\begin{itemize}
\tightlist
\item
  Una prueba de la verificación de familia con las columnas «3 Dirichlet» y «2 Dirichlet + exacta a 0» (aceptadas), y «1 Dirichlet + exactas a 0», «Dirichlet + complemento» y «Dirichlet + beta» (rechazadas).
\item
  Otra que abra las columnas de \texttt{0\_\allowbreak{}miscellaneous/\allowbreak{}networks/\allowbreak{}mid/\allowbreak{}MID-Chancellor-corrected.\allowbreak{}pgmx} y compruebe que pasan la verificación.
\end{itemize}

\setcounter{subsection}{78}
\subsection[Editor y muestreador reparten distinto el complemento junto a una Dirichlet]{El editor de incertidumbre y el muestreador reparten de forma distinta la probabilidad del complemento cuando una columna mezcla Dirichlet y complemento}\label{err:79}
\begin{ficha}
\fichakey{Estado}Presente
\fichakey{Módulo}core
\fichakey{Paquete}org.openmarkov.core.model.network.modelUncertainty
\fichakey{Ficheros}\path{core/src/main/java/org/openmarkov/core/model/network/modelUncertainty/Sampler.java:162-179} \newline \path{gui/src/main/java/org/openmarkov/gui/dialog/node/UncertainValuesDialog.java:253-289} \newline \path{core/src/main/java/org/openmarkov/core/model/network/modelUncertainty/TablePotentialSampler.java:38-80} \newline \path{core/src/main/java/org/openmarkov/core/model/network/modelUncertainty/SystematicSampling.java:216-217,292-293}
\fichakey{Comprobado}Leyendo el código y ejecutándolo
\fichakey{Esfuerzo}Pequeño (horas)
\fichakey{Decisión de equipo}\textcolor{decisionrojo}{\textbf{Sí.}} Si una columna que mezcla Dirichlet con complemento es válida (y entonces qué masa recibe el complemento) o se rechaza también al leer el fichero
\fichakey{Relacionados}\hyperref[err:78]{error~78}
\end{ficha}
\paragraph{Qué pasa}
Una columna de una tabla de probabilidad condicionada puede llevar incertidumbre: cada casilla tiene una distribución (exacta, beta, Dirichlet, complemento\ldots). Si una columna mezcla Dirichlet y complemento, el análisis de sensibilidad probabilístico muestrea una columna que suma más de 1, sin aviso.

El editor no deja escribir esa columna en un nodo de azar, porque su regla 3 la rechaza (ver \hyperref[err:78]{error~78}). Pero el muestreador no mira las reglas del editor ni las valida nadie al leer: un \texttt{.\allowbreak{}pgmx} escrito a mano o por otro programa con una columna así llega al análisis. En las redes del repositorio no hay ninguna columna con esa mezcla: las de \texttt{0\_\allowbreak{}miscellaneous/\allowbreak{}networks/\allowbreak{}mid/\allowbreak{}MID-Chancellor-*.\allowbreak{}pgmx} tienen Dirichlet y complemento en el mismo potencial, pero en columnas distintas.

De una columna incierta salen números de dos maneras:

\begin{itemize}
\tightlist
\item
  El editor (\texttt{Uncertain\allowbreak{}Values\allowbreak{}Dialog.\allowbreak{}calculate\allowbreak{}Reference\allowbreak{}Values}, al aceptar el diálogo de incertidumbre) calcula los valores de referencia de la tabla: cada distribución en su media, las casillas Dirichlet normalizadas entre sí, y los complementos repartiéndose \texttt{1\ -\ (\allowbreak{}suma\ de\ las\ medias\ de\ las\ «otras»\ +\ suma\ de\ las\ Dirichlet)}.
\item
  El muestreador que usan los análisis de sensibilidad (\texttt{Sampler.\allowbreak{}generate\allowbreak{}Sample}, llamado desde \texttt{Table\allowbreak{}Potential\allowbreak{}Sampler.\allowbreak{}sample} en el análisis probabilístico y desde \texttt{Systematic\allowbreak{}Sampling} en tornado y araña) reparte entre los complementos \texttt{1\ -\ (\allowbreak{}suma\ de\ las\ «otras»)}, sin descontar lo que se llevan las Dirichlet.
\end{itemize}

Solo difieren en columnas que mezclan Dirichlet y complemento. Como la regla 3 del editor rechaza esa columna, la fórmula del editor nunca se aplica a ese caso.

Medido con un pequeño programa de prueba sobre la columna {[}Dirichlet(2), Complemento(1){]}: el editor da \texttt{{[}1.\allowbreak{}0,\allowbreak{}\ 0.\allowbreak{}0{]}} (suma 1) y el muestreador da \texttt{{[}1.\allowbreak{}0,\allowbreak{}\ 1.\allowbreak{}0{]}} (suma 2). Con una columna solo de Dirichlet los dos coinciden en forma.
\paragraph{El código}
core/src/main/java/org/openmarkov/core/model/network/modelUncertainty/Sampler.java:170-176

\begin{Shaded}
\begin{Highlighting}[]
\DataTypeTok{double}\OperatorTok{[]}\NormalTok{ sampleDir }\OperatorTok{=} \FunctionTok{getSample}\OperatorTok{(}\NormalTok{dirFamily}\OperatorTok{,}\NormalTok{ randomGenerator}\OperatorTok{);}
\FunctionTok{placeInArray}\OperatorTok{(}\NormalTok{sampledConfigurationValues}\OperatorTok{,}\NormalTok{ indexesDirichlet}\OperatorTok{,}\NormalTok{ sampleDir}\OperatorTok{);}
\CommentTok{// Process complements}
\DataTypeTok{double}\NormalTok{ massForComp }\OperatorTok{=} \FloatTok{1.0} \OperatorTok{{-}} \OperatorTok{(}\NormalTok{Tools}\OperatorTok{.}\FunctionTok{sum}\OperatorTok{(}\NormalTok{sampleOther}\OperatorTok{));}
\NormalTok{complementFamily}\OperatorTok{.}\FunctionTok{setProbMass}\OperatorTok{(}\NormalTok{massForComp}\OperatorTok{);}
\DataTypeTok{double}\OperatorTok{[]}\NormalTok{ sampleComp }\OperatorTok{=}\NormalTok{ complementFamily}\OperatorTok{.}\FunctionTok{getSample}\OperatorTok{();}
\end{Highlighting}
\end{Shaded}

gui/src/main/java/org/openmarkov/gui/dialog/node/UncertainValuesDialog.java:275-281

\begin{Shaded}
\begin{Highlighting}[]
\DataTypeTok{double}\OperatorTok{[]}\NormalTok{ meanDir }\OperatorTok{=}\NormalTok{ dir}\OperatorTok{.}\FunctionTok{getMean}\OperatorTok{();}
\FunctionTok{placeInArray}\OperatorTok{(}\NormalTok{refValues}\OperatorTok{,}\NormalTok{ dirichletIndexes}\OperatorTok{,}\NormalTok{ meanDir}\OperatorTok{);}
\CommentTok{// Process complements}
\DataTypeTok{double}\NormalTok{ massForComp }\OperatorTok{=} \FloatTok{1.0} \OperatorTok{{-}} \OperatorTok{(}\NormalTok{Tools}\OperatorTok{.}\FunctionTok{sum}\OperatorTok{(}\NormalTok{meanOther}\OperatorTok{)} \OperatorTok{+}\NormalTok{ Tools}\OperatorTok{.}\FunctionTok{sum}\OperatorTok{(}\NormalTok{meanDir}\OperatorTok{));}
\NormalTok{comp}\OperatorTok{.}\FunctionTok{setProbMass}\OperatorTok{(}\NormalTok{massForComp}\OperatorTok{);}
\DataTypeTok{double}\OperatorTok{[]}\NormalTok{ meanComp }\OperatorTok{=}\NormalTok{ comp}\OperatorTok{.}\FunctionTok{getMean}\OperatorTok{();}
\end{Highlighting}
\end{Shaded}
\paragraph{Cómo arreglarlo}
Lo primero es la decisión de equipo:

\begin{itemize}
\tightlist
\item
  Si esa mezcla es inválida (que es lo que dice hoy la regla 3 del editor), hay que comprobar las mismas reglas al leer o antes de muestrear, y avisar.
\item
  Si es válida, las dos fórmulas deben ser la misma, y lo natural es que vivan en un único sitio de \texttt{core} (por ejemplo, en \texttt{Sampler} o \texttt{Family\allowbreak{}Distribution}) al que llamen tanto el editor como el muestreador, de modo que no puedan volver a divergir.
\end{itemize}

Las reglas de familia hoy viven solo en el diálogo (\texttt{gui}), así que moverlas a \texttt{core} es lo que permite aplicarlas en el muestreo y en la lectura.

Prueba de regresión: para un conjunto de columnas (exacta+complemento, beta+complemento, solo Dirichlet, Dirichlet+exactas a 0 y, según lo decidido, Dirichlet+complemento), exigir que la media que calcula el editor y la esperanza del muestreador coincidan y que ninguna columna aceptada sume más de 1.

\setcounter{subsection}{79}
\subsection[El mapa de dos parámetros elige mal la mejor opción si hay tres o más]{El mapa de dos parámetros pinta como mejor una opción que no lo es cuando la decisión tiene tres o más opciones}\label{err:80}
\begin{ficha}
\fichakey{Estado}Presente
\fichakey{Módulo}sensitivityAnalysis
\fichakey{Paquete}org.openmarkov.sensitivityanalysis.dialog
\fichakey{Ficheros}\path{sensitivityAnalysis/src/main/java/org/openmarkov/sensitivityanalysis/dialog/MapDialog.java:270-298}
\fichakey{Comprobado}Leyendo el código y ejecutándolo
\fichakey{Esfuerzo}Pequeño (horas)
\fichakey{Decisión de equipo}No
\fichakey{Tapado por}\hyperref[err:82]{error~82}: la zona de datos sale de un solo color. \newline \hyperref[err:81]{error~81}: las opciones tercera y siguientes salen en gris.
\fichakey{Relacionados}\hyperref[err:74]{error~74}, \hyperref[err:81]{error~81}, \hyperref[err:82]{error~82}, \hyperref[err:83]{error~83}
\end{ficha}
\paragraph{Qué pasa}
El análisis de sensibilidad de tipo «Map (two-way analysis)» (análisis de dos parámetros a la vez) dibuja una rejilla: cada casilla es una combinación de valores de los dos parámetros inciertos elegidos. Cuando el ámbito del análisis es una decisión («Scope: decision»), cada casilla se pinta del color de la opción de esa decisión que da mayor utilidad esperada en ese punto, o del color de «draw» (empate) si hay empate entre las mejores.

Con dos opciones el resultado es correcto. Con tres o más, gana la última opción que supere a la primera, aunque no sea la mejor, y se declara empate en cuanto cualquier opción iguala a la primera, aunque otra las supere a las dos.

Camino del usuario: con una red de decisión abierta que tenga dos parámetros inciertos con nombre, menú de herramientas → «Sensitivity analysis», aceptar las opciones de inferencia, elegir el tipo «Map (two-way analysis)», ámbito de decisión y una decisión con tres o más opciones.

En las redes de ejemplo del repositorio las decisiones con dos parámetros inciertos con nombre son de dos opciones, así que no se ha visto con una red real; la consecuencia (color de una opción equivocada) se ha comprobado sobre los datos que la gráfica recibe. Hoy la consecuencia visible queda además tapada por otros dos defectos del mismo mapa: toda la zona de datos sale de un solo color (\hyperref[err:82]{error~82}) y las opciones tercera y siguientes salen en gris (\hyperref[err:81]{error~81}).

El método privado \texttt{Map\allowbreak{}Dialog.\allowbreak{}get\allowbreak{}Map\allowbreak{}Chart(\allowbreak{})} busca la opción ganadora así: guarda en \texttt{z\allowbreak{}Axis\allowbreak{}Value} la utilidad de la primera opción y recorre las demás. Pero \texttt{z\allowbreak{}Axis\allowbreak{}Value} no se actualiza cuando aparece una opción mejor, de modo que cada opción se compara con la primera y no con la mejor hasta ese momento.

Medido con un pequeño programa de prueba que construye la gráfica del mapa (sin abrir la ventana) con una decisión de tres opciones \texttt{a}, \texttt{b}, \texttt{c} y una casilla con utilidades escritas a mano. El valor que se pinta es el índice de la opción ganadora (−1 es empate):

\begin{Verbatim}[breaklines,breakanywhere,fontsize=\small]
utilidades [1, 5, 3] -> se pinta 2 (opción c)     debería ser 1 (opción b)
utilidades [1, 5, 1] -> se pinta -1 (empate)      debería ser 1 (opción b)
utilidades [1, 1, 5] -> se pinta 2 (opción c)     correcto
utilidades [5, 5, 3] -> se pinta -1 (empate)      correcto
\end{Verbatim}
\paragraph{El código}
sensitivityAnalysis/src/main/java/org/openmarkov/sensitivityanalysis/dialog/MapDialog.java:274-297

\begin{Shaded}
\begin{Highlighting}[]
\DataTypeTok{double}\NormalTok{ zAxisValue }\OperatorTok{=}\NormalTok{ uncertainParameterPotential}\OperatorTok{.}\FunctionTok{getValues}\OperatorTok{()[}\NormalTok{dataPosition}\OperatorTok{];}
\KeywordTok{...}
\ControlFlowTok{if} \OperatorTok{(}\NormalTok{decisionVariable }\OperatorTok{!=} \KeywordTok{null}\OperatorTok{)} \OperatorTok{\{}
    \DataTypeTok{int}\NormalTok{ greatUtilityStateIndex }\OperatorTok{=} \DecValTok{0}\OperatorTok{;}
    \ControlFlowTok{for} \OperatorTok{(}\DataTypeTok{int}\NormalTok{ decisionVariableStateIndex }\OperatorTok{=} \DecValTok{1}\OperatorTok{;}
\NormalTok{         decisionVariableStateIndex }\OperatorTok{\textless{}}\NormalTok{ decisionVariableStates}\OperatorTok{;}\NormalTok{ decisionVariableStateIndex}\OperatorTok{++)} \OperatorTok{\{}
        \DataTypeTok{int}\NormalTok{ dataPositionForThatDecision }\OperatorTok{=}\NormalTok{ dataPosition }\OperatorTok{+} \OperatorTok{(}
\NormalTok{                decisionVariableStateIndex }\OperatorTok{*}\NormalTok{ totalIterations }\OperatorTok{*}\NormalTok{ totalIterations}
        \OperatorTok{);}
        \DataTypeTok{double}\NormalTok{ zAxisValueForThatDecision }\OperatorTok{=}\NormalTok{ uncertainParameterPotential}
                \OperatorTok{.}\FunctionTok{getValues}\OperatorTok{()[}\NormalTok{dataPositionForThatDecision}\OperatorTok{];}
        \ControlFlowTok{if} \OperatorTok{(}\NormalTok{zAxisValueForThatDecision }\OperatorTok{\textgreater{}}\NormalTok{ zAxisValue}\OperatorTok{)} \OperatorTok{\{}
\NormalTok{            greatUtilityStateIndex }\OperatorTok{=}\NormalTok{ decisionVariableStateIndex}\OperatorTok{;}
        \OperatorTok{\}} \ControlFlowTok{else} \ControlFlowTok{if} \OperatorTok{(}\NormalTok{zAxisValueForThatDecision }\OperatorTok{==}\NormalTok{ zAxisValue}\OperatorTok{)} \OperatorTok{\{}
\NormalTok{            greatUtilityStateIndex }\OperatorTok{=} \OperatorTok{{-}}\DecValTok{1}\OperatorTok{;}
        \OperatorTok{\}}
    \OperatorTok{\}}
\NormalTok{    zAxisValue }\OperatorTok{=}\NormalTok{ greatUtilityStateIndex}\OperatorTok{;}
\OperatorTok{\}}
\end{Highlighting}
\end{Shaded}
\paragraph{Cómo arreglarlo}
Llevar en el bucle el mejor valor visto hasta el momento (actualizarlo cuando una opción lo supera) y declarar empate solo cuando una opción iguala a ese mejor valor; si luego otra lo supera, el empate desaparece.

Conviene separar la búsqueda de la opción ganadora en un método estático que reciba el vector de utilidades de una casilla y devuelva el índice o −1, para poder probarlo sin la ventana.

Es el único sitio del módulo que decide el ganador de una casilla del mapa (los demás diálogos ---tornado, gráfica de un parámetro, araña--- no pintan ganadores por casilla). Si se decide comparar con tolerancia en vez de con \texttt{==}, debe hacerse igual en la comparación de «mayor» y en la de «igual».

Prueba de regresión: el método extraído con \texttt{\{1,\allowbreak{}5,\allowbreak{}3\}} → 1, \texttt{\{1,\allowbreak{}5,\allowbreak{}1\}} → 1, \texttt{\{1,\allowbreak{}1,\allowbreak{}5\}} → 2, \texttt{\{5,\allowbreak{}5,\allowbreak{}3\}} → −1, \texttt{\{3,\allowbreak{}5,\allowbreak{}5\}} → −1.

\setcounter{subsection}{80}
\subsection[En el mapa de dos parámetros el empate y la tercera opción salen en gris]{En el mapa de dos parámetros el empate y las opciones tercera y siguientes salen en gris, no del color que enseña la leyenda}\label{err:81}
\begin{ficha}
\fichakey{Estado}Presente
\fichakey{Módulo}sensitivityAnalysis
\fichakey{Paquete}org.openmarkov.sensitivityanalysis.dialog
\fichakey{Ficheros}\path{sensitivityAnalysis/src/main/java/org/openmarkov/sensitivityanalysis/dialog/MapDialog.java:352-381}
\fichakey{Comprobado}Leyendo el código y ejecutándolo
\fichakey{Esfuerzo}Pequeño (horas)
\fichakey{Decisión de equipo}No
\fichakey{Tapado por}\hyperref[err:82]{error~82}: la zona de datos sale de un solo color.
\fichakey{Relacionados}\hyperref[err:74]{error~74}, \hyperref[err:80]{error~80}, \hyperref[err:82]{error~82}, \hyperref[err:83]{error~83}
\end{ficha}
\paragraph{Qué pasa}
En el mapa de dos parámetros con ámbito de decisión, la leyenda dice que el empate es rojo y que la tercera opción es amarilla, pero en el mapa los dos salen del mismo gris claro, indistinguibles entre sí. El empate gris afecta también a decisiones de dos opciones, que son las de las redes de ejemplo; la tercera opción gris, solo a decisiones de tres o más.

Camino del usuario: menú de herramientas → «Sensitivity analysis», tipo «Map (two-way analysis)», ámbito de decisión. Hoy la consecuencia visible queda además tapada porque la zona de datos sale entera de un color (\hyperref[err:82]{error~82}).

Cada casilla del mapa recibe un número: el índice de la opción ganadora (0, 1, 2\ldots) o −1 si hay empate (ver \hyperref[err:80]{error~80}). Para convertir ese número en color, \texttt{Map\allowbreak{}Dialog.\allowbreak{}get\allowbreak{}Map\allowbreak{}Chart(\allowbreak{})} crea una \texttt{Lookup\allowbreak{}Paint\allowbreak{}Scale} de JFreeChart (la biblioteca de gráficas) y le añade un color por cada valor, el mismo que pone en la leyenda.

El constructor sin argumentos de \texttt{Lookup\allowbreak{}Paint\allowbreak{}Scale} fija una escala que solo cubre de 0,0 a 1,0 y usa gris claro (192,192,192) para todo lo que cae fuera. Por eso las entradas añadidas para −1 (empate) y para 2, 3\ldots{} (tercera opción en adelante) quedan fuera del rango y se pintan en gris claro.

Medido con un pequeño programa de prueba que construye la gráfica del mapa (sin abrir la ventana) con una decisión de tres opciones, y le pregunta a la escala de qué color pinta cada valor:

\begin{Verbatim}[breaklines,breakanywhere,fontsize=\small]
escala de 0.0 a 1.0, color por omisión 192,192,192
leyenda «draw»   255,85,85     escala para -1.0: 192,192,192
leyenda «D = a»   85,85,255    escala para  0.0:  85,85,255
leyenda «D = b»   85,255,85    escala para  1.0:  85,255,85
leyenda «D = c»  255,255,85    escala para  2.0: 192,192,192
\end{Verbatim}
\paragraph{El código}
sensitivityAnalysis/src/main/java/org/openmarkov/sensitivityanalysis/dialog/MapDialog.java:355-381

\begin{Shaded}
\begin{Highlighting}[]
\NormalTok{DefaultDrawingSupplier defaultDrawingSupplier }\OperatorTok{=} \KeywordTok{new} \FunctionTok{DefaultDrawingSupplier}\OperatorTok{();}
\NormalTok{LegendItemCollection chartLegend }\OperatorTok{=} \KeywordTok{new} \FunctionTok{LegendItemCollection}\OperatorTok{();}
\NormalTok{LookupPaintScale paintScale }\OperatorTok{=} \KeywordTok{new} \FunctionTok{LookupPaintScale}\OperatorTok{();}
\KeywordTok{...}
\NormalTok{chartLegend}\OperatorTok{.}\FunctionTok{add}\OperatorTok{(}\NormalTok{drawLegendItem}\OperatorTok{);}
\NormalTok{paintScale}\OperatorTok{.}\FunctionTok{add}\OperatorTok{({-}}\FloatTok{1.0}\OperatorTok{,}\NormalTok{ drawLegendItem}\OperatorTok{.}\FunctionTok{getFillPaint}\OperatorTok{());}

\ControlFlowTok{for} \OperatorTok{(}\DataTypeTok{int}\NormalTok{ stateIndex }\OperatorTok{=} \DecValTok{0}\OperatorTok{;}\NormalTok{ stateIndex }\OperatorTok{\textless{}}\NormalTok{ decisionVariableStates}\OperatorTok{;}\NormalTok{ stateIndex}\OperatorTok{++)} \OperatorTok{\{}
\NormalTok{    LegendItem legendItem }\OperatorTok{=} \KeywordTok{new} \FunctionTok{LegendItem}\OperatorTok{(}
\NormalTok{            decisionVariable}\OperatorTok{.}\FunctionTok{getName}\OperatorTok{()} \OperatorTok{+} \StringTok{" = "} \OperatorTok{+}\NormalTok{ decisionVariable}\OperatorTok{.}\FunctionTok{getStateName}\OperatorTok{(}\NormalTok{stateIndex}\OperatorTok{),} \CommentTok{// Legend label}
            \KeywordTok{...}
\NormalTok{            defaultDrawingSupplier}\OperatorTok{.}\FunctionTok{getNextPaint}\OperatorTok{());} \CommentTok{// Fill paint}
\NormalTok{    paintScale}\OperatorTok{.}\FunctionTok{add}\OperatorTok{(}\NormalTok{stateIndex}\OperatorTok{,}\NormalTok{ legendItem}\OperatorTok{.}\FunctionTok{getFillPaint}\OperatorTok{());}
\NormalTok{    chartLegend}\OperatorTok{.}\FunctionTok{add}\OperatorTok{(}\NormalTok{legendItem}\OperatorTok{);}
\OperatorTok{\}}
\NormalTok{plot}\OperatorTok{.}\FunctionTok{setFixedLegendItems}\OperatorTok{(}\NormalTok{chartLegend}\OperatorTok{);}
\NormalTok{renderer}\OperatorTok{.}\FunctionTok{setPaintScale}\OperatorTok{(}\NormalTok{paintScale}\OperatorTok{);}
\end{Highlighting}
\end{Shaded}
\paragraph{Cómo arreglarlo}
Construir la escala con sus límites reales: \texttt{new\ Lookup\allowbreak{}Paint\allowbreak{}Scale(\allowbreak{}-1.\allowbreak{}0,\allowbreak{}\ decision\allowbreak{}Variable\allowbreak{}States\ -\ 1,\allowbreak{}\ color)}, de modo que cubra desde el empate (−1) hasta el índice de la última opción (los dos límites están incluidos en la escala). El color por omisión debería ser uno que no se confunda con ninguna entrada de la leyenda.

Es el único sitio que crea esta escala (la rama de ámbito global usa \texttt{Gray\allowbreak{}Paint\allowbreak{}Scale} con los límites mínimo y máximo, que no tiene este problema). \texttt{Default\allowbreak{}Drawing\allowbreak{}Supplier} repite colores a partir de cierto número de opciones; no es parte de este fallo.

Prueba de regresión: construir la gráfica con una decisión de tres opciones (se puede hacer fabricando el diálogo sin ejecutar su constructor) y comprobar que \texttt{get\allowbreak{}Paint(\allowbreak{}z)} coincide con el color del elemento de la leyenda correspondiente para z = −1, 0, 1, 2.

\setcounter{subsection}{81}
\subsection[El mapa de dos parámetros pinta toda la zona de datos de un solo color]{El mapa de dos parámetros pinta toda la zona de datos de un solo color porque cada casilla mide una unidad}\label{err:82}
\begin{ficha}
\fichakey{Estado}Presente
\fichakey{Módulo}sensitivityAnalysis
\fichakey{Paquete}org.openmarkov.sensitivityanalysis.dialog
\fichakey{Ficheros}\path{sensitivityAnalysis/src/main/java/org/openmarkov/sensitivityanalysis/dialog/MapDialog.java:211-246} \newline \path{sensitivityAnalysis/src/main/java/org/openmarkov/sensitivityanalysis/dialog/MapDialog.java:339-341}
\fichakey{Comprobado}Leyendo el código y ejecutándolo
\fichakey{Esfuerzo}Pequeño (horas)
\fichakey{Decisión de equipo}No
\fichakey{Tapado por}\hyperref[err:74]{error~74}: en el jar publicado el mapa de ámbito global no llega a dibujarse.
\fichakey{Relacionados}\hyperref[err:74]{error~74}, \hyperref[err:80]{error~80}, \hyperref[err:81]{error~81}, \hyperref[err:83]{error~83}, \hyperref[err:90]{error~90}
\end{ficha}
\paragraph{Qué pasa}
El análisis de sensibilidad «Map (two-way analysis)» dibuja la zona de datos entera del color de una sola casilla, así que el mapa no enseña nada. Además, los ejes quedan de −0,75 a 0,75, con media unidad vacía a cada lado de los datos. Afecta a los dos ámbitos, el de decisión y el global (el global pinta una escala de grises con el mismo renderizador).

Camino del usuario: menú de herramientas → «Sensitivity analysis», tipo «Map (two-way analysis)», con cualquier red que tenga dos parámetros inciertos con nombre.

\texttt{Map\allowbreak{}Dialog.\allowbreak{}get\allowbreak{}Map\allowbreak{}Chart(\allowbreak{})} dibuja el mapa con un \texttt{XYBlock\allowbreak{}Renderer} de JFreeChart, que pinta un rectángulo («bloque») centrado en cada punto de datos. El tamaño del bloque se da en unidades de los datos y por omisión es 1 × 1. \texttt{Map\allowbreak{}Dialog} no lo cambia nunca.

Las coordenadas del mapa son la variación de cada parámetro. Para las variaciones en porcentaje o en razón (todas menos «UDIN», intervalo definido por el usuario) van de −v/100 a +v/100: con la variación por omisión del 25 \%, de −0,25 a 0,25, y con 50 puntos por parámetro (valor por omisión) las casillas están a 0,01 unas de otras.

Cada bloque mide entonces 1 × 1, cien veces más que la separación, y cualquier bloque centrado en un punto entre −0,25 y 0,25 cubre por completo el cuadrado de datos. Como los bloques se pintan en orden, el último (la esquina superior derecha) tapa a todos. JFreeChart estira los ejes para que quepa el bloque, y de ahí el rango de −0,75 a 0,75.

Con el tipo «UDIN» la separación depende de los límites que escriba el usuario; con los de omisión (0 a 1) los bloques también se solapan mucho.

Medido con un pequeño programa de prueba que construye la gráfica del mapa sin abrir la ventana, con 5 × 5 casillas (4 intervalos, variación del 25 \%) y una decisión de dos opciones: en las tres columnas de la izquierda gana «no» y en las dos de la derecha gana «yes». Pintada en una imagen de 800 × 800 y leyendo el color en el centro de cada casilla:

\begin{Verbatim}[breaklines,breakanywhere,fontsize=\small]
anchura del bloque 1.0, altura 1.0
eje horizontal de -0.75 a 0.75, eje vertical de -0.75 a 0.75
las 25 casillas, esperado «no» (azul 85,85,255) en x = -0.25, -0.125, 0 y «yes» (verde) en x = 0.125, 0.25:
todas leen 85,255,85 (verde)
\end{Verbatim}

Da lo mismo con la versión 1.5.6 de JFreeChart que declara el \texttt{pom.\allowbreak{}xml} y, en ámbito de decisión, con el jar ejecutable publicado; con ese jar el mapa de ámbito global no llega a dibujarse (\hyperref[err:74]{error~74}).
\paragraph{El código}
sensitivityAnalysis/src/main/java/org/openmarkov/sensitivityanalysis/dialog/MapDialog.java:221-222, 245-246, 339-341

\begin{Shaded}
\begin{Highlighting}[]
\CommentTok{// Size of every horizontal interval}
\DataTypeTok{double}\NormalTok{ hVariationInterval }\OperatorTok{=} \OperatorTok{(}\NormalTok{hMaxVariationValue }\OperatorTok{{-}}\NormalTok{ hMinVariationValue}\OperatorTok{)} \OperatorTok{/}\NormalTok{ iterations}\OperatorTok{;}
\KeywordTok{...}
\CommentTok{// Size of every vertical interval}
\DataTypeTok{double}\NormalTok{ vVariationInterval }\OperatorTok{=} \OperatorTok{(}\NormalTok{vMaxVariationValue }\OperatorTok{{-}}\NormalTok{ vMinVariationValue}\OperatorTok{)} \OperatorTok{/}\NormalTok{ iterations}\OperatorTok{;}
\KeywordTok{...}
\NormalTok{XYBlockRenderer renderer }\OperatorTok{=} \KeywordTok{new} \FunctionTok{XYBlockRenderer}\OperatorTok{();}

\NormalTok{XYPlot plot }\OperatorTok{=} \KeywordTok{new} \FunctionTok{XYPlot}\OperatorTok{(}\NormalTok{dataset}\OperatorTok{,}\NormalTok{ xAxis}\OperatorTok{,}\NormalTok{ yAxis}\OperatorTok{,}\NormalTok{ renderer}\OperatorTok{);}
\end{Highlighting}
\end{Shaded}
\paragraph{Cómo arreglarlo}
Tras crear el renderizador, darle el tamaño de una casilla: \texttt{renderer.\allowbreak{}set\allowbreak{}Block\allowbreak{}Width(\allowbreak{}h\allowbreak{}Variation\allowbreak{}Interval)} y \texttt{renderer.\allowbreak{}set\allowbreak{}Block\allowbreak{}Height(\allowbreak{}v\allowbreak{}Variation\allowbreak{}Interval)}; las dos variables ya están calculadas. Con eso los bloques se tocan sin solaparse y los ejes llegan medio intervalo más allá del primer y el último punto. Sirve para los dos ámbitos, porque el renderizador es común.

No hay otro sitio del módulo que use \texttt{XYBlock\allowbreak{}Renderer}.

Prueba de regresión: construir la gráfica con un patrón conocido (media rejilla gana una opción, media otra), pintarla en una \texttt{Buffered\allowbreak{}Image} sin pantalla y comprobar el color en el centro de cada casilla; comprobar también que el rango del eje horizontal es {[}mín − intervalo/2, máx + intervalo/2{]}.

\setcounter{subsection}{82}
\subsection[El mapa no se abre si gana en todo él la misma opción y no es la primera]{El mapa de dos parámetros no se abre cuando gana la misma opción en todas las casillas y no es la primera, o cuando la utilidad es la misma en todo el mapa}\label{err:83}
\begin{ficha}
\fichakey{Estado}Presente
\fichakey{Módulo}sensitivityAnalysis
\fichakey{Paquete}org.openmarkov.sensitivityanalysis.dialog
\fichakey{Ficheros}\path{sensitivityAnalysis/src/main/java/org/openmarkov/sensitivityanalysis/dialog/MapDialog.java:248-249} \newline \path{sensitivityAnalysis/src/main/java/org/openmarkov/sensitivityanalysis/dialog/MapDialog.java:300-308} \newline \path{sensitivityAnalysis/src/main/java/org/openmarkov/sensitivityanalysis/dialog/MapDialog.java:349-350} \newline \path{sensitivityAnalysis/src/main/java/org/openmarkov/sensitivityanalysis/dialog/MapDialog.java:384}
\fichakey{Comprobado}Leyendo el código y ejecutándolo
\fichakey{Esfuerzo}Pequeño (horas)
\fichakey{Decisión de equipo}No
\fichakey{Tapado por}\hyperref[err:74]{error~74}: en el jar publicado, en ámbito de decisión; desde el IDE se ve.
\fichakey{Relacionados}\hyperref[err:74]{error~74}, \hyperref[err:80]{error~80}, \hyperref[err:81]{error~81}, \hyperref[err:82]{error~82}, \hyperref[err:88]{error~88}, \hyperref[err:90]{error~90}
\end{ficha}
\paragraph{Qué pasa}
El análisis de sensibilidad «Map (two-way analysis)» con ámbito de decisión no se abre cuando en todas las casillas del mapa gana la misma opción y esa opción no es la primera de la decisión: sale el aviso de error inesperado con \texttt{Illegal\allowbreak{}Argument\allowbreak{}Exception:\ A\ positive\ range\ length\ is\ required:\ Range{[}1.\allowbreak{}0,\allowbreak{}1.\allowbreak{}0{]}}. Es un caso corriente: basta con que la mejor opción no cambie dentro de los intervalos de los dos parámetros y que no sea la primera (que suele ser «no»).

Camino del usuario: abrir \texttt{0\_\allowbreak{}miscellaneous/\allowbreak{}networks/\allowbreak{}id/\allowbreak{}ID-decide-test.\allowbreak{}pgmx}, menú Tools → «Sensitivity analysis», aceptar las opciones de inferencia, tipo «Map (two-way analysis)», parámetros «cost of test» y «cost of therapy» con la variación y los puntos por omisión (25 \% y 50 puntos), ámbito de decisión, decisión «Do test?».

En ámbito global pasa lo mismo cuando la utilidad esperada es igual y positiva en todas las casillas, por ejemplo con dos parámetros que no influyen en el resultado.

Así se comporta con la versión de JFreeChart que declara el proyecto (1.5.6), que es la que se usa desde el IDE. Con el jar ejecutable publicado, que lleva dentro las clases de JFreeChart 1.0.12 (\hyperref[err:74]{error~74}), el ámbito de decisión no falla, porque esa versión no comprueba la longitud del rango; el ámbito global falla igual, antes: con utilidades iguales y positivas, en \texttt{new\ Gray\allowbreak{}Paint\allowbreak{}Scale} (línea 384, \texttt{Requires\ lower\allowbreak{}Bound\ \textless{}\ upper\allowbreak{}Bound}), y con cualquier otra, por el \texttt{No\allowbreak{}Such\allowbreak{}Method\allowbreak{}Error} de \hyperref[err:74]{error~74}. Al arreglar \hyperref[err:74]{error~74}, el jar pasa a fallar también en ámbito de decisión.

Por qué pasa: \texttt{Map\allowbreak{}Dialog.\allowbreak{}get\allowbreak{}Map\allowbreak{}Chart} guarda en cada casilla, en ámbito de decisión, el índice de la opción ganadora (−1 si hay empate), y en ámbito global la utilidad. Mientras recorre las casillas lleva el menor y el mayor valor (\texttt{min\allowbreak{}Range\allowbreak{}Utility}, \texttt{max\allowbreak{}Range\allowbreak{}Utility}), y después crea el eje de la escala y le da ese rango con \texttt{scale\allowbreak{}Axis.\allowbreak{}set\allowbreak{}Range(\allowbreak{}min\allowbreak{}Range\allowbreak{}Utility,\allowbreak{}\ max\allowbreak{}Range\allowbreak{}Utility)} (línea 350). Lo hace también en ámbito de decisión, donde ese eje no se usa. JFreeChart 1.5.6 rechaza un rango de longitud cero.

Si en todas las casillas gana la opción de índice k ≥ 1, el mínimo y el máximo valen k. Con la primera opción (0) o con empate (−1) en todo el mapa no falla, pero solo porque el máximo arranca en \texttt{Double.\allowbreak{}MIN\_\allowbreak{}VALUE}, que es un positivo diminuto y no el número más negativo (el mismo inicio de \hyperref[err:88]{error~88} y \hyperref[err:90]{error~90}): el rango queda de 0, o de −1, a ese número. Por lo mismo, en ámbito global una utilidad constante negativa o cero no falla, y una positiva sí.

Medido construyendo la gráfica sin pantalla con los datos reales de la tarea \texttt{VESens\allowbreak{}An\allowbreak{}Map} sobre \texttt{0\_\allowbreak{}miscellaneous/\allowbreak{}networks/\allowbreak{}id/\allowbreak{}ID-decide-test.\allowbreak{}pgmx} (51 × 51 casillas, variación del 25 \%) y JFreeChart 1.5.6:

\begin{Verbatim}[breaklines,breakanywhere,fontsize=\small]
cost of test x cost of therapy, decisión Do test?: «yes» en 2601 casillas     -> IllegalArgumentException (MapDialog.java:350)
cost of test x cost of therapy, decisión Therapy:  «no» en 2601 casillas      -> se construye
prevalence x sensitivity, decisión Do test?:       «no» en 15, «yes» en 2586  -> se construye
\end{Verbatim}

Con datos escritos a mano: tres opciones con la tercera ganadora en todas las casillas → \texttt{Range{[}2.\allowbreak{}0,\allowbreak{}2.\allowbreak{}0{]}}; la primera o empate en todas → se construye; ámbito global con todas las utilidades a 7 → \texttt{Range{[}7.\allowbreak{}0,\allowbreak{}7.\allowbreak{}0{]}}; a −7 o a 0 → se construye.
\paragraph{El código}
sensitivityAnalysis/src/main/java/org/openmarkov/sensitivityanalysis/dialog/MapDialog.java:248-249, 300-308, 349-350, 384

\begin{Shaded}
\begin{Highlighting}[]
\DataTypeTok{double}\NormalTok{ minRangeUtility }\OperatorTok{=} \BuiltInTok{Double}\OperatorTok{.}\FunctionTok{MAX\_VALUE}\OperatorTok{;}
\DataTypeTok{double}\NormalTok{ maxRangeUtility }\OperatorTok{=} \BuiltInTok{Double}\OperatorTok{.}\FunctionTok{MIN\_VALUE}\OperatorTok{;}
\KeywordTok{...}
                \CommentTok{// Update minimum utility found}
                \ControlFlowTok{if} \OperatorTok{(}\NormalTok{zAxisValue }\OperatorTok{\textless{}}\NormalTok{ minRangeUtility}\OperatorTok{)} \OperatorTok{\{}
\NormalTok{                    minRangeUtility }\OperatorTok{=}\NormalTok{ zAxisValue}\OperatorTok{;}
                \OperatorTok{\}}
                \CommentTok{// Update maximum utility found}
                \ControlFlowTok{if} \OperatorTok{(}\NormalTok{zAxisValue }\OperatorTok{\textgreater{}}\NormalTok{ maxRangeUtility}\OperatorTok{)} \OperatorTok{\{}
\NormalTok{                    maxRangeUtility }\OperatorTok{=}\NormalTok{ zAxisValue}\OperatorTok{;}
                \OperatorTok{\}}
\KeywordTok{...}
\NormalTok{NumberAxis scaleAxis }\OperatorTok{=} \KeywordTok{new} \FunctionTok{NumberAxis}\OperatorTok{(}\NormalTok{stringDatabase}\OperatorTok{.}\FunctionTok{getString}\OperatorTok{(}\StringTok{"SensitivityAnalysis.General.Scale"}\OperatorTok{));}
\NormalTok{scaleAxis}\OperatorTok{.}\FunctionTok{setRange}\OperatorTok{(}\NormalTok{minRangeUtility}\OperatorTok{,}\NormalTok{ maxRangeUtility}\OperatorTok{);}
\KeywordTok{...}
\NormalTok{    PaintScale paintScale }\OperatorTok{=} \KeywordTok{new} \FunctionTok{GrayPaintScale}\OperatorTok{(}\NormalTok{minRangeUtility}\OperatorTok{,}\NormalTok{ maxRangeUtility}\OperatorTok{);}
\end{Highlighting}
\end{Shaded}
\paragraph{Cómo arreglarlo}
Crear el eje de la escala y darle rango solo en la rama global, que es donde se usa (para \texttt{Paint\allowbreak{}Scale\allowbreak{}Legend}). Allí, si el mínimo y el máximo coinciden, abrir el rango un poco a los dos lados antes de dárselo al eje y a \texttt{Gray\allowbreak{}Paint\allowbreak{}Scale}, que tampoco admite un rango vacío.

Este arreglo debe ir junto con el del inicio del máximo: \hyperref[err:88]{error~88} propone arrancarlo con \texttt{Double.\allowbreak{}NEGATIVE\_\allowbreak{}INFINITY} también en \texttt{Map\allowbreak{}Dialog} (línea 249), y con ese cambio solo, un mapa donde gana la primera opción o hay empate en todas las casillas, o una utilidad global constante negativa o cero, pasarían a fallar igual.

En ámbito de decisión la escala de colores es una \texttt{Lookup\allowbreak{}Paint\allowbreak{}Scale}, que no necesita el eje; su rango propio es el de \hyperref[err:81]{error~81}.

Prueba de regresión: construir la gráfica sin pantalla (o extraer el cálculo del rango a un método sin interfaz) con una rejilla en la que gane siempre la segunda opción, otra en la que gane siempre la primera, otra con empate en todas las casillas y, en ámbito global, con utilidad constante positiva, cero y negativa: en todos los casos debe construirse.

\setcounter{subsection}{83}
\subsection[La tabla de intervenciones fronterizas nombra otra opción si dos empatan]{La tabla de intervenciones fronterizas pone el nombre de otra opción cuando dos opciones tienen el mismo coste y la misma efectividad}\label{err:84}
\begin{ficha}
\fichakey{Estado}Presente
\fichakey{Módulo}costEffectiveness
\fichakey{Paquete}org.openmarkov.costEffectiveness
\fichakey{Ficheros}\path{costEffectiveness/src/main/java/org/openmarkov/costEffectiveness/CEDecisionResults.java:479-487,547-590} \newline \path{core/src/main/java/org/openmarkov/core/model/network/CEP.java:480-485} \newline \path{inference/src/main/java/org/openmarkov/inference/algorithm/variableElimination/operation/CEBaseOperations.java:395-409}
\fichakey{Comprobado}Leyendo el código y ejecutándolo
\fichakey{Esfuerzo}Pequeño (horas)
\fichakey{Decisión de equipo}No
\fichakey{Tapado por}\hyperref[err:74]{error~74}: en el jar publicado; desde el IDE se ve.
\fichakey{Relacionados}\hyperref[err:74]{error~74}, \hyperref[err:85]{error~85}, \hyperref[err:86]{error~86}
\end{ficha}
\paragraph{Qué pasa}
La ventana de coste-efectividad con ámbito de decisión (\texttt{CEDecision\allowbreak{}Results}) tiene una pestaña de intervenciones fronterizas: las opciones de la decisión que forman la frontera eficiente, con su coste, su efectividad y la razón coste-efectividad incremental (ICER). Cuando dos opciones dan exactamente el mismo coste y la misma efectividad, la frontera contiene la segunda y la fila se titula con el nombre de la primera. Los números son los mismos, pero la tabla nombra una opción que el análisis no eligió.

Camino del usuario: diagrama de influencia o modelo de Markov con criterio de coste-efectividad, menú de herramientas → «Cost-effectiveness analysis», ámbito de decisión, pestaña de intervenciones fronterizas, cuando dos opciones de la decisión dan exactamente el mismo coste y la misma efectividad (por ejemplo, opciones que no afectan a ninguna utilidad). Es un caso poco frecuente y la consecuencia es un nombre equivocado, no un número equivocado. No se ha buscado una red de ejemplo del repositorio que lo muestre.

\texttt{calculate\allowbreak{}Frontier\allowbreak{}Interventions} calcula la frontera con \texttt{CEBase\allowbreak{}Operations.\allowbreak{}deterministic\allowbreak{}CEA} y devuelve, por identidad (\texttt{==}), los objetos \texttt{CEP} (las particiones de coste-efectividad) de las opciones elegidas. Después, \texttt{get\allowbreak{}Frontier\allowbreak{}Interventions\allowbreak{}Table} tiene que saber de qué opción es cada uno para escribir su nombre, y lo busca con \texttt{Arrays.\allowbreak{}as\allowbreak{}List(\allowbreak{}ceps\allowbreak{}For\allowbreak{}Decision).\allowbreak{}index\allowbreak{}Of(\allowbreak{}.\allowbreak{}.\allowbreak{}.\allowbreak{})}, que compara con \texttt{equals}.

\texttt{CEP.\allowbreak{}equals} compara por valor (umbrales, costes y efectividades; es una redefinición real desde el commit \texttt{325e5c6}). Si dos opciones tienen particiones iguales, \texttt{index\allowbreak{}Of} devuelve siempre la primera. Y \texttt{deterministic\allowbreak{}CEA}, ante costes iguales, se queda con la última opción (condición \texttt{i\ \textgreater{}\ 0} de la línea 403).

Medido con un pequeño programa de prueba sobre dos particiones construidas a mano (coste 0 y efectividad 0 para «no» y para «yes»), usando los métodos reales de la ventana fabricada sin constructor:

\begin{Verbatim}[breaklines,breakanywhere,fontsize=\small]
mismo objeto: false  equals: true  indexOf(segunda): 0
deterministicCEA elige: yes (índice 1)
la partición de la frontera es el objeto de «yes»: true
tabla de la frontera: 1 fila, nombre de la fila 0: no
\end{Verbatim}
\paragraph{El código}
costEffectiveness/src/main/java/org/openmarkov/costEffectiveness/CEDecisionResults.java:479-487

\begin{Shaded}
\begin{Highlighting}[]
\KeywordTok{public} \BuiltInTok{JTable} \FunctionTok{getFrontierInterventionsTable}\OperatorTok{()} \OperatorTok{\{}

    \BuiltInTok{List}\OperatorTok{\textless{}}\NormalTok{CEP}\OperatorTok{\textgreater{}}\NormalTok{ frontierInterventions }\OperatorTok{=} \FunctionTok{calculateFrontierInterventions}\OperatorTok{(}\NormalTok{AnalysisTab}\OperatorTok{.}\FunctionTok{FRONTIER\_INTERVENTIONS}\OperatorTok{);}
    \BuiltInTok{String}\OperatorTok{[]}\NormalTok{ interventionsNames }\OperatorTok{=} \KeywordTok{new} \BuiltInTok{String}\OperatorTok{[}\NormalTok{frontierInterventions}\OperatorTok{.}\FunctionTok{size}\OperatorTok{()];}
    \BuiltInTok{List}\OperatorTok{\textless{}}\NormalTok{CEP}\OperatorTok{\textgreater{}}\NormalTok{ cepsForDecisionList }\OperatorTok{=} \BuiltInTok{Arrays}\OperatorTok{.}\FunctionTok{asList}\OperatorTok{(}\NormalTok{cepsForDecision}\OperatorTok{);}
    \ControlFlowTok{for} \OperatorTok{(}\DataTypeTok{int}\NormalTok{ i }\OperatorTok{=} \DecValTok{0}\OperatorTok{;}\NormalTok{ i }\OperatorTok{\textless{}}\NormalTok{ frontierInterventions}\OperatorTok{.}\FunctionTok{size}\OperatorTok{();}\NormalTok{ i}\OperatorTok{++)} \OperatorTok{\{}
        \DataTypeTok{int}\NormalTok{ stateIndex }\OperatorTok{=}\NormalTok{ cepsForDecisionList}\OperatorTok{.}\FunctionTok{indexOf}\OperatorTok{(}\NormalTok{frontierInterventions}\OperatorTok{.}\FunctionTok{get}\OperatorTok{(}\NormalTok{i}\OperatorTok{));}
\NormalTok{        interventionsNames}\OperatorTok{[}\NormalTok{i}\OperatorTok{]} \OperatorTok{=}\NormalTok{ decisionVariable}\OperatorTok{.}\FunctionTok{getStateName}\OperatorTok{(}\NormalTok{stateIndex}\OperatorTok{);}
    \OperatorTok{\}}
\end{Highlighting}
\end{Shaded}
\paragraph{Cómo arreglarlo}
Que \texttt{calculate\allowbreak{}Frontier\allowbreak{}Interventions} devuelva también el índice del estado de cada intervención (ya lo tiene: \texttt{selected\allowbreak{}Indices.\allowbreak{}get(\allowbreak{}j)}, línea 584), por ejemplo devolviendo índices en vez de particiones, y que la tabla use ese índice en lugar de buscarlo.

Hay otro llamador de \texttt{calculate\allowbreak{}Frontier\allowbreak{}Interventions} (línea 930, la línea de la frontera en el plano), que solo usa costes y efectividades y habría que adaptar si cambia el tipo devuelto. No hay más búsquedas por valor de particiones en el módulo.

Prueba de regresión: dos particiones iguales, como en la medida de arriba, comprobando que la fila se llama «yes».

\setcounter{subsection}{84}
\subsection[El coste-efectividad falla si un hallazgo hace imposible una sola opción]{El análisis de coste-efectividad de un diagrama de influencia falla con NullPointerException cuando un hallazgo previo hace imposible una opción de una decisión y no las demás}\label{err:85}
\begin{ficha}
\fichakey{Estado}Presente
\fichakey{Módulo}inference
\fichakey{Paquete}org.openmarkov.inference.algorithm.variableElimination.operation
\fichakey{Ficheros}\path{inference/src/main/java/org/openmarkov/inference/algorithm/variableElimination/operation/CEBaseOperations.java:64-120} \newline \path{inference/src/main/java/org/openmarkov/inference/algorithm/variableElimination/operation/CEBaseOperations.java:609-661} \newline \path{inference/src/main/java/org/openmarkov/inference/algorithm/variableElimination/operation/CEBaseOperations.java:543-556} \newline \path{inference/src/main/java/org/openmarkov/inference/algorithm/variableElimination/operation/CEAlgebra.java:195-203} \newline \path{inference/src/main/java/org/openmarkov/inference/algorithm/variableElimination/operation/CEAlgebra.java:333-358} \newline \path{core/src/main/java/org/openmarkov/core/model/network/CEP.java:146-151}
\fichakey{Comprobado}Leyendo el código y ejecutándolo
\fichakey{Esfuerzo}Pequeño (horas)
\fichakey{Decisión de equipo}No
\fichakey{Relacionados}\hyperref[err:48]{error~48}, \hyperref[err:54]{error~54}, \hyperref[err:84]{error~84}, \hyperref[err:86]{error~86}, \hyperref[err:87]{error~87}, \hyperref[nota:46]{nota~46}
\end{ficha}
\paragraph{Qué pasa}
Con un hallazgo previo a la resolución sobre una variable de azar cuyo valor depende de una decisión, el análisis de coste-efectividad no se completa: sale el aviso de error inesperado con \texttt{Null\allowbreak{}Pointer\allowbreak{}Exception:\ Cannot\ read\ the\ array\ length\ because\ "thresholds{[}i{]}"\ is\ null}. El caso típico es el resultado de una prueba, que vale «not done» cuando se decide no hacerla: anotar el resultado como ya conocido basta.

Camino del usuario: abrir \texttt{0\_\allowbreak{}miscellaneous/\allowbreak{}networks/\allowbreak{}id/\allowbreak{}ID-CEA-test-2therapies-new-test.\allowbreak{}pgmx}; en modo edición, clic derecho sobre «Test» → «Add finding» y elegir cualquier estado; menú Tools → «Cost-effectiveness analysis», aceptar las opciones de inferencia y elegir el ámbito global. Con ámbito de decisión falla cuando la decisión elegida es anterior a aquella cuyas opciones deja imposibles el hallazgo y el escenario del diálogo no fija esa otra decisión: por ejemplo, en \texttt{0\_\allowbreak{}miscellaneous/\allowbreak{}networks/\allowbreak{}id/\allowbreak{}ID-arthronet-ce.\allowbreak{}pgmx} con un hallazgo en «Ga67 Tc99» y la decisión «Realizar Implante». En \texttt{ID-CEA-test-2therapies-new-test.\allowbreak{}pgmx}, el ámbito de «Dec: New test» no falla, porque su escenario fija «Dec: Test»; y con el de «Dec: Test» el análisis termina, pero la ventana recibe la partición de probabilidad cero de una opción y revienta al pintarla (\hyperref[err:86]{error~86}).

Por qué pasa. El análisis por eliminación de variables representa el resultado de cada configuración con una partición de coste-efectividad (\texttt{CEP}: para cada tramo de la disposición a pagar, la intervención óptima con su coste y su efectividad). Cuando una configuración tiene probabilidad cero, \texttt{CEAlgebra.\allowbreak{}multiply\allowbreak{}And\allowbreak{}Marginalize} pone en su lugar la partición de probabilidad cero (\texttt{CEP.\allowbreak{}get\allowbreak{}Zero\allowbreak{}Partition(\allowbreak{})}), un objeto sin tramos cuyos umbrales, costes, efectividades e intervenciones son \texttt{null}.

Al eliminar una decisión, \texttt{CEAlgebra.\allowbreak{}ce\allowbreak{}Maximize} pasa a \texttt{CEBase\allowbreak{}Operations.\allowbreak{}optimal\allowbreak{}CEP} la partición de cada opción. \texttt{optimal\allowbreak{}CEP} solo contempla dos casos: que todas sean de probabilidad cero (devuelve la partición cero) o que no lo sea ninguna. Si alguna no es cero, llama a \texttt{get\allowbreak{}Union\allowbreak{}Thresholds(\allowbreak{}partitions,\allowbreak{}\ null)}, que lee los umbrales de todas, también los de la partición cero, y \texttt{get\allowbreak{}Min\allowbreak{}Threshold} falla al pedir \texttt{thresholds{[}i{]}.\allowbreak{}length} (líneas 649-650). Más abajo haría lo mismo con \texttt{get\allowbreak{}Intervention}, \texttt{get\allowbreak{}Cost} y \texttt{get\allowbreak{}Effectiveness}, y el umbral máximo lo toma de \texttt{partitions.\allowbreak{}get(\allowbreak{}0)}, que en la partición cero vale 0.

Un hallazgo sobre una variable que depende de una decisión deja con probabilidad cero las opciones incompatibles con él y no las demás: «Test = not done» hace imposible «Dec: Test = yes», y «Test = negative» o «positive» hacen imposible «Dec: Test = no». Esa mezcla es la que llega a \texttt{optimal\allowbreak{}CEP}.

Medido:

\begin{itemize}
\tightlist
\item
  \texttt{optimal\allowbreak{}CEP} con las particiones {[}normal, cero{]} y {[}cero, normal{]} lanza \texttt{Null\allowbreak{}Pointer\allowbreak{}Exception} en \texttt{get\allowbreak{}Min\allowbreak{}Threshold}; con {[}cero, cero{]} devuelve la partición cero.
\item
  Con \texttt{ID-CEA-test-2therapies-new-test.\allowbreak{}pgmx}, poniendo un hallazgo en cada estado de cada variable de azar y lanzando \texttt{VECEAnalysis} en ámbito global como lo hace el menú: los tres estados de «Test» fallan en \texttt{optimal\allowbreak{}CEP}; los demás hallazgos, no. En el ámbito de «Dec: Test» las particiones salen \texttt{{[}cep,\allowbreak{}\ zero{]}} («not done») o \texttt{{[}zero,\allowbreak{}\ cep{]}} («negative», «positive»). Sin hallazgos, la red se analiza bien.
\item
  Con \texttt{0\_\allowbreak{}miscellaneous/\allowbreak{}networks/\allowbreak{}id/\allowbreak{}ID-arthronet-ce.\allowbreak{}pgmx}: los tres estados de «Ga67 Tc99» (la gammagrafía, que vale «Prueba no realizada» si no se hace) fallan igual en ámbito global, y también en el ámbito de «Realizar Implante» con un escenario para sus tres predecesores informativos («IMC», «Diabetes», «Alergia ATB»), como el que añade el diálogo.
\end{itemize}
\paragraph{El código}
inference/src/main/java/org/openmarkov/inference/algorithm/variableElimination/operation/CEBaseOperations.java:64-72

\begin{Shaded}
\begin{Highlighting}[]
\KeywordTok{public} \DataTypeTok{static}\NormalTok{ CEP }\FunctionTok{optimalCEP}\OperatorTok{(}\NormalTok{Variable decisionVariable}\OperatorTok{,} \BuiltInTok{List}\OperatorTok{\textless{}}\NormalTok{CEP}\OperatorTok{\textgreater{}}\NormalTok{ partitions}\OperatorTok{)} \OperatorTok{\{}
\NormalTok{    CEP result}\OperatorTok{;}
    \DataTypeTok{int}\NormalTok{ i}\OperatorTok{;}
    \ControlFlowTok{for} \OperatorTok{(}\NormalTok{i }\OperatorTok{=} \DecValTok{0}\OperatorTok{;}\NormalTok{ i }\OperatorTok{\textless{}}\NormalTok{ partitions}\OperatorTok{.}\FunctionTok{size}\OperatorTok{()} \OperatorTok{\&\&}\NormalTok{ partitions}\OperatorTok{.}\FunctionTok{get}\OperatorTok{(}\NormalTok{i}\OperatorTok{).}\FunctionTok{isZero}\OperatorTok{();}\NormalTok{ i}\OperatorTok{++)}
        \OperatorTok{;}
    \ControlFlowTok{if} \OperatorTok{(}\NormalTok{i }\OperatorTok{==}\NormalTok{ partitions}\OperatorTok{.}\FunctionTok{size}\OperatorTok{())} \OperatorTok{\{}
\NormalTok{        result }\OperatorTok{=}\NormalTok{ CEP}\OperatorTok{.}\FunctionTok{getZeroPartition}\OperatorTok{();}
    \OperatorTok{\}} \ControlFlowTok{else} \OperatorTok{\{}
        \DataTypeTok{double}\OperatorTok{[]}\NormalTok{ allThresholds }\OperatorTok{=} \FunctionTok{getUnionThresholds}\OperatorTok{(}\NormalTok{partitions}\OperatorTok{,} \KeywordTok{null}\OperatorTok{);}
\end{Highlighting}
\end{Shaded}

inference/src/main/java/org/openmarkov/inference/algorithm/variableElimination/operation/CEBaseOperations.java:612-616, 648-650

\begin{Shaded}
\begin{Highlighting}[]
\ControlFlowTok{for} \OperatorTok{(}\DataTypeTok{int}\NormalTok{ i }\OperatorTok{=} \DecValTok{0}\OperatorTok{;}\NormalTok{ i }\OperatorTok{\textless{}}\NormalTok{ indexes}\OperatorTok{.}\FunctionTok{length}\OperatorTok{;}\NormalTok{ i}\OperatorTok{++)} \OperatorTok{\{}
    \CommentTok{// When probabilities = null, all partitions must be taken into account.}
    \ControlFlowTok{if} \OperatorTok{(}\NormalTok{probabilities }\OperatorTok{==} \KeywordTok{null} \OperatorTok{||}\NormalTok{ probabilities}\OperatorTok{[}\NormalTok{i}\OperatorTok{]} \OperatorTok{!=} \FloatTok{0.0}\OperatorTok{)} \OperatorTok{\{}
\NormalTok{        thresholds}\OperatorTok{[}\NormalTok{i}\OperatorTok{]} \OperatorTok{=}\NormalTok{ partitions}\OperatorTok{.}\FunctionTok{get}\OperatorTok{(}\NormalTok{i}\OperatorTok{).}\FunctionTok{getThresholds}\OperatorTok{();}
\KeywordTok{...}
\ControlFlowTok{for} \OperatorTok{(}\DataTypeTok{int}\NormalTok{ i }\OperatorTok{=} \DecValTok{0}\OperatorTok{;}\NormalTok{ i }\OperatorTok{\textless{}}\NormalTok{ indexes}\OperatorTok{.}\FunctionTok{length}\OperatorTok{;}\NormalTok{ i}\OperatorTok{++)} \OperatorTok{\{}
    \ControlFlowTok{if} \OperatorTok{((}\NormalTok{probabilities }\OperatorTok{==} \KeywordTok{null} \OperatorTok{||}\NormalTok{ probabilities}\OperatorTok{[}\NormalTok{i}\OperatorTok{]} \OperatorTok{!=} \FloatTok{0.0}\OperatorTok{)}
            \OperatorTok{\&\&}\NormalTok{ indexes}\OperatorTok{[}\NormalTok{i}\OperatorTok{]} \OperatorTok{\textless{}}\NormalTok{ thresholds}\OperatorTok{[}\NormalTok{i}\OperatorTok{].}\FunctionTok{length}\OperatorTok{)} \OperatorTok{\{}\CommentTok{// Threshold list not empty}
\end{Highlighting}
\end{Shaded}
\paragraph{Cómo arreglarlo}
Que las opciones cuya partición es de probabilidad cero no participen en la maximización: una opción incompatible con lo observado no se puede elegir. En \texttt{optimal\allowbreak{}CEP}, calcular los umbrales solo con las particiones que no son cero, hacer el análisis determinista de cada tramo (\texttt{deterministic\allowbreak{}CEA}) solo con esas opciones, conservando el índice de estado de cada una para construir su \texttt{Strategy\allowbreak{}Tree}, y tomar el umbral máximo de una partición que no sea cero.

La misma regla (una partición cero no aporta umbrales ni valores) debe cumplirse en los otros lectores de umbrales: \texttt{get\allowbreak{}Union\allowbreak{}Thresholds} hoy solo se salta las particiones cuya probabilidad es cero, no las particiones cero; \texttt{weighted\allowbreak{}Average}, al que llama \texttt{CEAlgebra.\allowbreak{}multiply\allowbreak{}And\allowbreak{}Marginalize}, lee los umbrales y los valores de toda partición con probabilidad distinta de cero; y \texttt{CEAlgebra.\allowbreak{}add\allowbreak{}Marginalize} (hoy sin llamadores) calcula qué particiones tener en cuenta y luego no lo usa. La ventana del ámbito de decisión tiene su propio problema con la partición cero, descrito en \hyperref[err:86]{error~86}.

Prueba de regresión: en \texttt{inference/\allowbreak{}src/\allowbreak{}test}, \texttt{optimal\allowbreak{}CEP} con {[}normal, cero{]} y con {[}cero, normal{]} debe devolver la partición normal con la intervención de su opción; y, con \texttt{0\_\allowbreak{}miscellaneous/\allowbreak{}networks/\allowbreak{}id/\allowbreak{}ID-CEA-test-2therapies-new-test.\allowbreak{}pgmx} y el hallazgo «Test = positive», \texttt{VECEAnalysis} en ámbito global debe terminar con una partición que no sea cero.

\setcounter{subsection}{85}
\subsection[Las ventanas de coste-efectividad fallan si una opción es imposible]{Las ventanas de resultados del análisis de coste-efectividad fallan con NullPointerException cuando una opción de la decisión, o la evidencia entera, tiene probabilidad cero}\label{err:86}
\begin{ficha}
\fichakey{Estado}Presente
\fichakey{Módulo}costEffectiveness, gui
\fichakey{Paquete}org.openmarkov.costEffectiveness; org.openmarkov.gui.dialog.costeffectiveness
\fichakey{Ficheros}\path{costEffectiveness/src/main/java/org/openmarkov/costEffectiveness/CEDecisionResults.java:273-291} \newline \path{costEffectiveness/src/main/java/org/openmarkov/costEffectiveness/CEDecisionResults.java:380-397} \newline \path{costEffectiveness/src/main/java/org/openmarkov/costEffectiveness/CEDecisionResults.java:547-576} \newline \path{costEffectiveness/src/main/java/org/openmarkov/costEffectiveness/CEDecisionResults.java:866-902} \newline \path{core/src/main/java/org/openmarkov/core/model/network/CEP.java:149-151} \newline \path{core/src/main/java/org/openmarkov/core/model/network/CEP.java:167-173} \newline \path{core/src/main/java/org/openmarkov/core/model/network/CEP.java:227-229} \newline \path{gui/src/main/java/org/openmarkov/gui/dialog/costeffectiveness/CEPDialog.java:204-208} \newline \path{costEffectiveness/src/main/java/org/openmarkov/costEffectiveness/CostEffectivenessPlugin.java:83-110}
\fichakey{Comprobado}Leyendo el código y ejecutándolo
\fichakey{Esfuerzo}Pequeño (horas)
\fichakey{Decisión de equipo}No
\fichakey{Corregir antes}\hyperref[err:48]{error~48}: allí se decide el mensaje para la evidencia imposible, que estas ventanas deben usar.
\fichakey{Relacionados}\hyperref[err:48]{error~48}, \hyperref[err:54]{error~54}, \hyperref[err:74]{error~74}, \hyperref[err:84]{error~84}, \hyperref[err:85]{error~85}, \hyperref[err:87]{error~87}, \hyperref[nota:46]{nota~46}
\end{ficha}
\paragraph{Qué pasa}
El análisis de coste-efectividad con ámbito de decisión no llega a abrir su ventana cuando alguna opción de la decisión es imposible con la evidencia: sale el aviso de error inesperado con \texttt{Null\allowbreak{}Pointer\allowbreak{}Exception:\ Cannot\ read\ the\ array\ length\ because\ "this.\allowbreak{}strategy\allowbreak{}Trees"\ is\ null}. La ventana del ámbito global hace lo mismo cuando es el resultado entero el que tiene probabilidad cero.

Camino del usuario: abrir \texttt{0\_\allowbreak{}miscellaneous/\allowbreak{}networks/\allowbreak{}id/\allowbreak{}ID-CEA-test-2therapies-new-test.\allowbreak{}pgmx}; en modo edición, clic derecho sobre «Test» → «Add finding» y elegir cualquier estado, por ejemplo «not done»; menú Tools → «Cost-effectiveness analysis» y aceptar las opciones de inferencia; en el diálogo de ámbito, que propone el ámbito de decisión con «Dec: New test», elegir la decisión «Dec: Test» (no tiene escenario que fijar) y aceptar.

El resultado de ese análisis es una partición de coste-efectividad por opción (\texttt{CEP}: para cada tramo de la disposición a pagar, el coste, la efectividad y la intervención óptima). A una opción incompatible con la evidencia le corresponde la partición de probabilidad cero (\texttt{CEP.\allowbreak{}get\allowbreak{}Zero\allowbreak{}Partition(\allowbreak{})}), un objeto sin tramos cuyos vectores son \texttt{null}. Con «Test = not done», «Dec: Test = yes» es imposible y las particiones son \texttt{{[}normal,\allowbreak{}\ cero{]}}; con «negative» o «positive», \texttt{{[}cero,\allowbreak{}\ normal{]}}.

Por qué pasa. El constructor de \texttt{CEDecision\allowbreak{}Results} monta primero la pestaña de análisis, y \texttt{get\allowbreak{}Intervals\allowbreak{}Panel} pide a cada partición \texttt{get\allowbreak{}Num\allowbreak{}Intervals(\allowbreak{})}, que devuelve \texttt{strategy\allowbreak{}Trees.\allowbreak{}length} sin mirar si es la partición cero. No es el único sitio: \texttt{get\allowbreak{}Analysis\allowbreak{}Table} pide después a cada opción \texttt{get\allowbreak{}Cost}, \texttt{get\allowbreak{}Effectiveness} y \texttt{get\allowbreak{}Intervention}, y \texttt{CEP.\allowbreak{}index} lee \texttt{costs.\allowbreak{}length}; lo mismo hacen \texttt{calculate\allowbreak{}Frontier\allowbreak{}Interventions} (pestaña de intervenciones fronterizas y línea del plano), \texttt{get\allowbreak{}Scatter\allowbreak{}Plot\allowbreak{}Data} y \texttt{get\allowbreak{}Line\allowbreak{}Frontier\allowbreak{}Interventions\allowbreak{}Data} (plano). Ningún método de la ventana consulta \texttt{is\allowbreak{}Zero(\allowbreak{})}.

La ventana global, \texttt{CEPDialog}, construye su tabla con \texttt{get\allowbreak{}Data\allowbreak{}From\allowbreak{}CEP}, que lee \texttt{cep.\allowbreak{}get\allowbreak{}Costs(\allowbreak{}).\allowbreak{}length}, y con la partición cero lanza \texttt{Null\allowbreak{}Pointer\allowbreak{}Exception:\ Cannot\ read\ the\ array\ length\ because\ "costs"\ is\ null}. Le llega esa partición cuando la evidencia previa a la resolución es imposible, por ejemplo en \texttt{0\_\allowbreak{}miscellaneous/\allowbreak{}networks/\allowbreak{}id/\allowbreak{}ID-arthronet-ce.\allowbreak{}pgmx} con «Isquemia = Implante no realizado» y «CC\_Drenaje = mayor 800 cc y menor 1000 cc», y también con evidencia posible por el defecto de \hyperref[err:54]{error~54}.

Otros caminos que dejan particiones cero en la ventana por decisión: un escenario imposible en el diálogo de ámbito (\hyperref[err:87]{error~87}) y \hyperref[err:54]{error~54}, que dejan todas las opciones a cero. Con ámbito global, un hallazgo de este tipo hace fallar el análisis antes, en \texttt{optimal\allowbreak{}CEP} (\hyperref[err:85]{error~85}), o le hace devolver la partición cero entera (\hyperref[err:54]{error~54}). En las redes de análisis de decisiones el selector desactiva el ámbito de decisión, así que por ahí no se llega a esta ventana.

Con el jar ejecutable publicado la ventana por decisión tampoco se abre por \hyperref[err:74]{error~74}, pero este error sale antes: la pestaña de análisis se construye antes que la del plano.

Medido, reproduciendo el camino del menú con el panel de ámbito real (\texttt{Scope\allowbreak{}Selector\allowbreak{}Panel}, construido sin pantalla), la tarea \texttt{VECEAnalysis} y los métodos de la ventana:

\begin{itemize}
\tightlist
\item
  \texttt{ID-CEA-test-2therapies-new-test.\allowbreak{}pgmx}, hallazgo «Test = not done», decisión «Dec: Test»: particiones \texttt{{[}normal,\allowbreak{}\ cero{]}}; \texttt{get\allowbreak{}Intervals\allowbreak{}Panel} lanza la excepción en \texttt{CEP.\allowbreak{}get\allowbreak{}Num\allowbreak{}Intervals} (línea 228) y \texttt{get\allowbreak{}Analysis\allowbreak{}Table} en \texttt{CEP.\allowbreak{}index} (línea 168). Con «Test = positive», \texttt{{[}cero,\allowbreak{}\ normal{]}} y lo mismo. Sin hallazgo, o con «Test = not done» y las decisiones «Dec: New test» o «Therapy», las particiones son normales y los dos métodos construyen su panel.
\item
  Recorriendo todas las decisiones que ofrece el diálogo, sin hallazgo y con cada hallazgo suelto sobre cada variable de azar, con el escenario por omisión, en las redes de coste-efectividad de \texttt{0\_\allowbreak{}miscellaneous/\allowbreak{}networks/\allowbreak{}id/\allowbreak{}} (las cuatro \texttt{ID-CEA-*}, \texttt{ID-arthronet-ce.\allowbreak{}pgmx} y las tres \texttt{ID-mediastinet*-ce.\allowbreak{}pgmx}): la ventana recibe alguna partición cero, pero no todas, en 3 casos de \texttt{ID-CEA-test-2therapies-new-test.\allowbreak{}pgmx} y en 12 de cada \texttt{ID-mediastinet*-ce.\allowbreak{}pgmx}; y todas cero en 4, 2, 27, 37 y 33 casos de \texttt{ID-CEA-test-2therapies-new-test.\allowbreak{}pgmx}, \texttt{ID-arthronet-ce.\allowbreak{}pgmx}, \texttt{ID-mediastinet-ce.\allowbreak{}pgmx}, \texttt{ID-mediastinet-first-ebus-ce.\allowbreak{}pgmx} e \texttt{ID-mediastinet-first-eus-ce.\allowbreak{}pgmx}. Sin hallazgo no hay ninguno.
\item
  \texttt{CEPDialog.\allowbreak{}get\allowbreak{}Data\allowbreak{}From\allowbreak{}CEP} con la partición global de \texttt{ID-arthronet-ce.\allowbreak{}pgmx} y la evidencia imposible de arriba: \texttt{Null\allowbreak{}Pointer\allowbreak{}Exception} en la línea 207.
\end{itemize}
\paragraph{El código}
costEffectiveness/src/main/java/org/openmarkov/costEffectiveness/CEDecisionResults.java:282-291

\begin{Shaded}
\begin{Highlighting}[]
\ControlFlowTok{for} \OperatorTok{(}\DataTypeTok{int}\NormalTok{ i }\OperatorTok{=} \DecValTok{0}\OperatorTok{;}\NormalTok{ i }\OperatorTok{\textless{}}\NormalTok{ gtablePotentialResult}\OperatorTok{.}\FunctionTok{elementTable}\OperatorTok{.}\FunctionTok{size}\OperatorTok{();}\NormalTok{ i}\OperatorTok{++)} \OperatorTok{\{}
\NormalTok{    CEP cep }\OperatorTok{=} \OperatorTok{(}\NormalTok{CEP}\OperatorTok{)}\NormalTok{ gtablePotentialResult}\OperatorTok{.}\FunctionTok{elementTable}\OperatorTok{.}\FunctionTok{get}\OperatorTok{(}\NormalTok{i}\OperatorTok{);}
\NormalTok{    cepsForDecision}\OperatorTok{[}\NormalTok{i}\OperatorTok{]} \OperatorTok{=}\NormalTok{ cep}\OperatorTok{;}
    \ControlFlowTok{if} \OperatorTok{(}\NormalTok{cep}\OperatorTok{.}\FunctionTok{getNumIntervals}\OperatorTok{()} \OperatorTok{\textgreater{}} \DecValTok{1}\OperatorTok{)} \OperatorTok{\{}
\NormalTok{        moreThanOneInterval }\OperatorTok{=} \KeywordTok{true}\OperatorTok{;}
        \ControlFlowTok{for} \OperatorTok{(}\DataTypeTok{double}\NormalTok{ threshold }\OperatorTok{:}\NormalTok{ cep}\OperatorTok{.}\FunctionTok{getThresholds}\OperatorTok{())} \OperatorTok{\{}
\NormalTok{            thresholds}\OperatorTok{.}\FunctionTok{add}\OperatorTok{(}\NormalTok{threshold}\OperatorTok{);}
        \OperatorTok{\}}
    \OperatorTok{\}}
\OperatorTok{\}}
\end{Highlighting}
\end{Shaded}

core/src/main/java/org/openmarkov/core/model/network/CEP.java:149-151, 227-229

\begin{Shaded}
\begin{Highlighting}[]
\KeywordTok{private} \FunctionTok{CEP}\OperatorTok{()} \OperatorTok{\{}
\NormalTok{    zeroProbability }\OperatorTok{=} \KeywordTok{true}\OperatorTok{;}
\OperatorTok{\}}
\KeywordTok{...}
\KeywordTok{public} \DataTypeTok{int} \FunctionTok{getNumIntervals}\OperatorTok{()} \OperatorTok{\{}
    \ControlFlowTok{return}\NormalTok{ strategyTrees}\OperatorTok{.}\FunctionTok{length}\OperatorTok{;}
\OperatorTok{\}}
\end{Highlighting}
\end{Shaded}

gui/src/main/java/org/openmarkov/gui/dialog/costeffectiveness/CEPDialog.java:204-207

\begin{Shaded}
\begin{Highlighting}[]
\KeywordTok{private} \DataTypeTok{static} \BuiltInTok{Object}\OperatorTok{[][]} \FunctionTok{getDataFromCEP}\OperatorTok{(}\NormalTok{CEP cep}\OperatorTok{)} \OperatorTok{\{}
    \DataTypeTok{double}\OperatorTok{[]}\NormalTok{ costs }\OperatorTok{=}\NormalTok{ cep}\OperatorTok{.}\FunctionTok{getCosts}\OperatorTok{();}
    \DataTypeTok{double}\OperatorTok{[]}\NormalTok{ effectiveness }\OperatorTok{=}\NormalTok{ cep}\OperatorTok{.}\FunctionTok{getEffectivities}\OperatorTok{();}
    \DataTypeTok{int}\NormalTok{ numRows }\OperatorTok{=}\NormalTok{ costs}\OperatorTok{.}\FunctionTok{length}\OperatorTok{;}
\end{Highlighting}
\end{Shaded}
\paragraph{Cómo arreglarlo}
La ventana por decisión debe tratar la partición cero como una opción imposible con la evidencia:

\begin{itemize}
\tightlist
\item
  en \texttt{get\allowbreak{}Intervals\allowbreak{}Panel}, no pedirle tramos ni umbrales (comprobar \texttt{is\allowbreak{}Zero(\allowbreak{})} antes);
\item
  en \texttt{get\allowbreak{}Analysis\allowbreak{}Table}, enseñar su fila sin coste ni efectividad, con un texto que diga que es imposible con la evidencia (el texto exacto es cosa de la interfaz);
\item
  dejarla fuera del plano (\texttt{get\allowbreak{}Scatter\allowbreak{}Plot\allowbreak{}Data}), de la frontera (\texttt{calculate\allowbreak{}Frontier\allowbreak{}Interventions}, que alimenta la pestaña de intervenciones fronterizas y \texttt{get\allowbreak{}Line\allowbreak{}Frontier\allowbreak{}Interventions\allowbreak{}Data}) y del selector «relative to» de \texttt{get\allowbreak{}Absolute\allowbreak{}Relative\allowbreak{}Panel}, porque \texttt{get\allowbreak{}Scatter\allowbreak{}Plot\allowbreak{}Data} resta el coste y la efectividad de la opción de referencia.
\end{itemize}

Si todas las particiones son cero, o la global es la partición cero, no hay nada que enseñar: \texttt{Cost\allowbreak{}Effectiveness\allowbreak{}Plugin.\allowbreak{}on\allowbreak{}Click} debería avisar, en sus dos ramas y antes de construir \texttt{CEDecision\allowbreak{}Results} o \texttt{CEPDialog}, de que la evidencia (con el escenario elegido) es imposible, con el mismo mensaje que se decida en \hyperref[err:48]{error~48}.

Otros lectores de particiones que no miran \texttt{is\allowbreak{}Zero(\allowbreak{})} y reciben la evidencia previa del editor con su escenario: \texttt{CESpider\allowbreak{}Dialog} (líneas 443 y 451) y \texttt{CEProbabilistic\allowbreak{}Dialog} (líneas 148, 488 y 498), del análisis de sensibilidad (por lectura, no medido). \texttt{CEP.\allowbreak{}clone} tampoco la mira, pero hoy no se llega (\hyperref[nota:46]{nota~46}). Que \texttt{CEP.\allowbreak{}get\allowbreak{}Num\allowbreak{}Intervals} devuelva 0 para la partición cero evitaría la primera excepción, no las de \texttt{get\allowbreak{}Cost}, así que no sustituye a lo anterior.

Prueba de regresión: sacar de la ventana el cálculo de los umbrales y de las filas de la tabla a métodos sin interfaz y probarlos con \texttt{{[}normal,\allowbreak{}\ cero{]}}, \texttt{{[}cero,\allowbreak{}\ normal{]}} y \texttt{{[}cero,\allowbreak{}\ cero{]}}; y comprobar que \texttt{VECEAnalysis} sobre \texttt{0\_\allowbreak{}miscellaneous/\allowbreak{}networks/\allowbreak{}id/\allowbreak{}ID-CEA-test-2therapies-new-test.\allowbreak{}pgmx}, con el hallazgo «Test = not done» y la decisión «Dec: Test», da \texttt{{[}normal,\allowbreak{}\ cero{]}}, que es el caso que la ventana debe cubrir.

\setcounter{subsection}{86}
\subsection[El escenario del ámbito de decisión puede ser imposible y no se comprueba]{El ámbito de decisión del análisis de coste-efectividad propone por omisión, y deja elegir, escenarios imposibles con la evidencia: la ventana revienta o enseña efectividades 0}\label{err:87}
\begin{ficha}
\fichakey{Estado}Presente
\fichakey{Módulo}gui, costEffectiveness
\fichakey{Paquete}org.openmarkov.gui.dialog.inference.common; org.openmarkov.costEffectiveness
\fichakey{Ficheros}\path{gui/src/main/java/org/openmarkov/gui/dialog/inference/common/ScopeSelectorPanel.java:227-285} \newline \path{costEffectiveness/src/main/java/org/openmarkov/costEffectiveness/CostEffectivenessPlugin.java:84-91} \newline \path{sensitivityAnalysis/src/main/java/org/openmarkov/sensitivityanalysis/dialog/ScopePanel.java:264-300} \newline \path{sensitivityAnalysis/src/main/java/org/openmarkov/sensitivityanalysis/dialog/SensitivityAnalysisDialog.java:534-543}
\fichakey{Comprobado}Leyendo el código y ejecutándolo
\fichakey{Esfuerzo}Pequeño (horas)
\fichakey{Decisión de equipo}\textcolor{decisionrojo}{\textbf{Sí.}} Si el diálogo debe proponer y ofrecer solo escenarios posibles con la evidencia, o dejarlos elegir y avisar al aceptar
\fichakey{Corregir antes}\hyperref[err:54]{error~54}: la comprobación necesita la probabilidad del escenario, y hoy sale 0 también con evidencia posible.
\fichakey{Relacionados}\hyperref[err:54]{error~54}, \hyperref[err:85]{error~85}, \hyperref[err:86]{error~86}
\end{ficha}
\paragraph{Qué pasa}
En el análisis de coste-efectividad con ámbito de decisión, el diálogo pide un escenario: un valor para cada predecesor informativo de la decisión elegida, es decir, las decisiones anteriores y las variables que se conocen al tomarla. Cada desplegable empieza en el primer estado; si la variable tiene un hallazgo previo, se fija en ese valor y se bloquea. Nada comprueba que el escenario sea posible.

Así, con un hallazgo sobre el resultado de una prueba, el escenario enseña el resultado bloqueado pero deja la decisión de hacer la prueba en su primera opción, «no», aunque el hallazgo diga que se hizo. El análisis recibe una evidencia imposible, todas las opciones salen con la partición de probabilidad cero y la ventana revienta con \texttt{Null\allowbreak{}Pointer\allowbreak{}Exception} (\hyperref[err:86]{error~86}).

Camino del usuario: abrir \texttt{0\_\allowbreak{}miscellaneous/\allowbreak{}networks/\allowbreak{}id/\allowbreak{}ID-CEA-test-2therapies-new-test.\allowbreak{}pgmx}; en modo edición, «Add finding» sobre «Test» → «negative»; menú Tools → «Cost-effectiveness analysis», aceptar las opciones; en el diálogo de ámbito, con ámbito de decisión, elegir «Dec: New test». El escenario enseña «Test = negative» bloqueado y «Dec: Test = no»; al aceptar sale el aviso de error inesperado.

El usuario también puede elegir a mano, sin ningún hallazgo, una combinación imposible. Según a qué variables afecte, la ventana revienta igual o el análisis termina con números que no corresponden a nada: en la misma red, para «Therapy» con «Dec: Test = yes» y «Test = not done», los costes salen 18, 20 018 y 70 018 y la efectividad 0 en las tres opciones; con el escenario posible «Dec: Test = no», la efectividad es 8,768, 9,074 y 8,908.

Por qué pasa: \texttt{Scope\allowbreak{}Selector\allowbreak{}Panel.\allowbreak{}get\allowbreak{}Decision\allowbreak{}Scenario\allowbreak{}Panel} crea un desplegable por predecesor con sus estados en orden, fija y bloquea los que tienen hallazgo, y \texttt{update\allowbreak{}Selected\allowbreak{}Scenario} convierte lo elegido en hallazgos sin mirar si son compatibles entre sí ni con la evidencia. \texttt{Cost\allowbreak{}Effectiveness\allowbreak{}Plugin.\allowbreak{}on\allowbreak{}Click} añade esos hallazgos a la evidencia previa y construye la ventana. El panel de ámbito del análisis de sensibilidad (\texttt{Scope\allowbreak{}Panel.\allowbreak{}get\allowbreak{}Decision\allowbreak{}Scenario\allowbreak{}Panel} y \texttt{Sensitivity\allowbreak{}Analysis\allowbreak{}Dialog}, líneas 534-543) hace lo mismo (por lectura).

Medido con el panel real (\texttt{Scope\allowbreak{}Selector\allowbreak{}Panel}, construido sin pantalla), la tarea \texttt{VECEAnalysis} y la evidencia que arma el menú:

\begin{Verbatim}[breaklines,breakanywhere,fontsize=\small]
ID-CEA-test-2therapies-new-test.pgmx, Test = negative, decisión Dec: New test
  escenario por omisión: Dec: Test = no           -> particiones [cero, cero]
  eligiendo Dec: Test = yes                       -> [normal, normal]
ID-mediastinet-ce.pgmx, TBNA = negative, decisión Dec:PET
  escenario por omisión: CT_scan = negative, Dec:TBNA = no   -> [cero, cero]
  eligiendo Dec:TBNA = yes                                   -> [normal, normal]
\end{Verbatim}

Recorriendo todas las decisiones de las redes de coste-efectividad de \texttt{0\_\allowbreak{}miscellaneous/\allowbreak{}networks/\allowbreak{}id/\allowbreak{}} sin hallazgos, el escenario por omisión es siempre posible: el defecto aparece con un hallazgo o eligiendo a mano.
\paragraph{El código}
gui/src/main/java/org/openmarkov/gui/dialog/inference/common/ScopeSelectorPanel.java:244-252

\begin{Shaded}
\begin{Highlighting}[]
\BuiltInTok{JComboBox}\OperatorTok{\textless{}}\BuiltInTok{String}\OperatorTok{\textgreater{}}\NormalTok{ stateSelector }\OperatorTok{=} \KeywordTok{new} \BuiltInTok{JComboBox}\OperatorTok{\textless{}\textgreater{}();}
\ControlFlowTok{for} \OperatorTok{(}\BuiltInTok{State}\NormalTok{ state }\OperatorTok{:}\NormalTok{ variable}\OperatorTok{.}\FunctionTok{getStates}\OperatorTok{())} \OperatorTok{\{}
\NormalTok{    stateSelector}\OperatorTok{.}\FunctionTok{addItem}\OperatorTok{(}\NormalTok{state}\OperatorTok{.}\FunctionTok{getName}\OperatorTok{());}
\OperatorTok{\}}
\ControlFlowTok{if} \OperatorTok{(}\NormalTok{preResolutionEvidence}\OperatorTok{.}\FunctionTok{contains}\OperatorTok{(}\NormalTok{variable}\OperatorTok{))} \OperatorTok{\{}
\NormalTok{    stateSelector}\OperatorTok{.}\FunctionTok{setSelectedItem}\OperatorTok{(}\NormalTok{preResolutionEvidence}\OperatorTok{.}\FunctionTok{getFinding}\OperatorTok{(}\NormalTok{variable}\OperatorTok{).}\FunctionTok{getState}\OperatorTok{());}
\NormalTok{    stateSelector}\OperatorTok{.}\FunctionTok{setEnabled}\OperatorTok{(}\KeywordTok{false}\OperatorTok{);}
\OperatorTok{\}}
\NormalTok{selectedScenario}\OperatorTok{.}\FunctionTok{put}\OperatorTok{(}\NormalTok{stateSelector}\OperatorTok{,}\NormalTok{ variable}\OperatorTok{);}
\end{Highlighting}
\end{Shaded}

costEffectiveness/src/main/java/org/openmarkov/costEffectiveness/CostEffectivenessPlugin.java:84-90

\begin{Shaded}
\begin{Highlighting}[]
\ControlFlowTok{case}\NormalTok{ DECISION }\OperatorTok{{-}\textgreater{}} \OperatorTok{\{}
\NormalTok{    EvidenceCase newPreResolutionEvidence }\OperatorTok{=} \KeywordTok{new} \FunctionTok{EvidenceCase}\OperatorTok{(}\NormalTok{preResolutionEvidence}\OperatorTok{);}
    \ControlFlowTok{for} \OperatorTok{(}\NormalTok{Finding finding }\OperatorTok{:}\NormalTok{ scopeSelectorPanel}\OperatorTok{.}\FunctionTok{getSelectedFindings}\OperatorTok{())} \OperatorTok{\{}
\NormalTok{        newPreResolutionEvidence}\OperatorTok{.}\FunctionTok{addFinding}\OperatorTok{(}\NormalTok{finding}\OperatorTok{);}
    \OperatorTok{\}}
    \ControlFlowTok{yield} \KeywordTok{new} \FunctionTok{CEDecisionResults}\OperatorTok{(}\NormalTok{parent}\OperatorTok{,}\NormalTok{ probNet}\OperatorTok{,}\NormalTok{ newPreResolutionEvidence}\OperatorTok{,}
\NormalTok{                                scopeSelectorPanel}\OperatorTok{.}\FunctionTok{getDecisionSelected}\OperatorTok{());}
\end{Highlighting}
\end{Shaded}
\paragraph{Cómo arreglarlo}
Hay que decidir el comportamiento:

\begin{itemize}
\tightlist
\item
  proponer por omisión un escenario compatible con la evidencia (para cada decisión del escenario, una opción con la que los hallazgos sean posibles) y ofrecer en cada desplegable solo los valores posibles con lo ya elegido; exige calcular probabilidades cada vez que cambia un desplegable;
\item
  dejar elegir cualquier combinación y comprobar al aceptar que la evidencia con el escenario tiene probabilidad mayor que cero, avisando si no.
\end{itemize}

Para esa comprobación no sirve hoy la probabilidad que devuelve \texttt{VECEAnalysis.\allowbreak{}get\allowbreak{}Probability}, que sale 0 también con evidencia posible (\hyperref[err:54]{error~54}). Aunque se arregle esto, la ventana debe resistir un resultado con todas las opciones a cero (\hyperref[err:86]{error~86}).

La misma regla debe cumplirse en el panel de escenario del análisis de sensibilidad (\texttt{Scope\allowbreak{}Panel}) y en \texttt{Sensitivity\allowbreak{}Analysis\allowbreak{}Dialog}, que arma allí la evidencia.

Prueba de regresión: con \texttt{0\_\allowbreak{}miscellaneous/\allowbreak{}networks/\allowbreak{}id/\allowbreak{}ID-CEA-test-2therapies-new-test.\allowbreak{}pgmx}, el hallazgo «Test = negative» y la decisión «Dec: New test», el escenario que propone \texttt{Scope\allowbreak{}Selector\allowbreak{}Panel} debe tener «Dec: Test = yes», o el diálogo debe rechazar «Dec: Test = no» con un aviso.

\setcounter{subsection}{87}
\subsection[Con valores todos negativos, el tornado dibuja barras hasta cero]{En el diagrama de tornado, un modelo cuyos valores son todos negativos dibuja barras que llegan hasta cero y ordena mal los parámetros}\label{err:88}
\begin{ficha}
\fichakey{Estado}Presente
\fichakey{Módulo}sensitivityAnalysis
\fichakey{Paquete}org.openmarkov.sensitivityanalysis.dialog
\fichakey{Ficheros}\path{sensitivityAnalysis/src/main/java/org/openmarkov/sensitivityanalysis/dialog/TornadoSpiderDialog.java:321-371} \newline \path{sensitivityAnalysis/src/main/java/org/openmarkov/sensitivityanalysis/dialog/TornadoSpiderDialog.java:532-533} \newline \path{sensitivityAnalysis/src/main/java/org/openmarkov/sensitivityanalysis/dialog/PlotDialog.java:213} \newline \path{sensitivityAnalysis/src/main/java/org/openmarkov/sensitivityanalysis/dialog/MapDialog.java:249} \newline \path{sensitivityAnalysis/src/main/java/org/openmarkov/sensitivityanalysis/model/TornadoBar.java:54-60}
\fichakey{Comprobado}Leyendo el código
\fichakey{Esfuerzo}Pequeño (horas)
\fichakey{Decisión de equipo}No
\fichakey{Corregir antes}\hyperref[err:83]{error~83}: a la vez, como mínimo, en \texttt{Map\allowbreak{}Dialog}: con este cambio solo, un mapa en el que gana la primera opción o hay empate en todas las casillas pasaría a fallar.
\fichakey{Tapado por}\hyperref[err:74]{error~74}: en el jar publicado; desde el IDE se ve.
\fichakey{Relacionados}\hyperref[err:74]{error~74}, \hyperref[err:83]{error~83}, \hyperref[err:89]{error~89}, \hyperref[err:90]{error~90}, \hyperref[err:91]{error~91}
\end{ficha}
\paragraph{Qué pasa}
El diagrama de tornado del análisis de sensibilidad dibuja, para cada parámetro incierto, una barra horizontal entre el menor y el mayor valor que toma el resultado del modelo (la utilidad esperada, o la diferencia entre dos opciones de una decisión) cuando ese parámetro se barre por su intervalo. Las barras se ordenan de más ancha a más estrecha (\texttt{Tornado\allowbreak{}Bar.\allowbreak{}compare\allowbreak{}To}), de modo que arriba quedan los parámetros que más influyen.

Camino de llegada: análisis de sensibilidad → tornado, sobre una red con parámetros inciertos cuyo resultado es negativo en todo el barrido (por ejemplo, un modelo de costes expresados como utilidades negativas).

Consecuencias, deducidas del código (el diálogo no se ha ejecutado porque es una ventana):

\begin{itemize}
\tightlist
\item
  Sin decisión: la barra va del valor mínimo real hasta 0 en lugar de hasta el máximo real. Por ejemplo, si la utilidad varía entre −1 010 y −990, la barra mide 1 010 en vez de 20.
\item
  Con decisión (comparar una opción con otra): \texttt{max\allowbreak{}Intervention\allowbreak{}Utility} y \texttt{max\allowbreak{}Reference\allowbreak{}Utility} también arrancan en \texttt{Double.\allowbreak{}MIN\_\allowbreak{}VALUE}, y con valores negativos las dos quedan en \textasciitilde0, así que los extremos de la barra salen mal.
\item
  Como el orden depende de la anchura, un parámetro poco influyente sobre un resultado muy negativo aparece arriba del todo, que es lo contrario de lo que el diagrama debe comunicar.
\end{itemize}

\texttt{Tornado\allowbreak{}Spider\allowbreak{}Dialog.\allowbreak{}get\allowbreak{}Tornado\allowbreak{}Bars} inicia cada mínimo con \texttt{Double.\allowbreak{}MAX\_\allowbreak{}VALUE} y cada máximo con \texttt{Double.\allowbreak{}MIN\_\allowbreak{}VALUE}. \texttt{Double.\allowbreak{}MIN\_\allowbreak{}VALUE} no es el número más negativo representable sino el positivo más pequeño (unos 4,9·10⁻³²⁴). Si todos los valores del barrido son negativos, \texttt{Math.\allowbreak{}max(\allowbreak{}Double.\allowbreak{}MIN\_\allowbreak{}VALUE,\allowbreak{}\ valor)} se queda en ese número positivo, prácticamente cero.

El mismo inicio aparece en el eje vertical del diagrama de araña de la misma clase (línea 533), en \texttt{Plot\allowbreak{}Dialog} (línea 213) y en \texttt{Map\allowbreak{}Dialog} (línea 249), con otra consecuencia que describe \hyperref[err:90]{error~90}.

Aparte, en la rama con decisión la barra no sale de la diferencia punto a punto entre las dos opciones; ese es otro defecto, \hyperref[err:89]{error~89}.
\paragraph{El código}
sensitivityAnalysis/src/main/java/org/openmarkov/sensitivityanalysis/dialog/TornadoSpiderDialog.java:321-365

\begin{Shaded}
\begin{Highlighting}[]
\KeywordTok{private} \BuiltInTok{List}\OperatorTok{\textless{}}\NormalTok{TornadoBar}\OperatorTok{\textgreater{}} \FunctionTok{getTornadoBars}\OperatorTok{()} \OperatorTok{\{}
    \KeywordTok{...}
        \DataTypeTok{double}\NormalTok{ minValue }\OperatorTok{=} \BuiltInTok{Double}\OperatorTok{.}\FunctionTok{MAX\_VALUE}\OperatorTok{;}
        \DataTypeTok{double}\NormalTok{ maxValue }\OperatorTok{=} \BuiltInTok{Double}\OperatorTok{.}\FunctionTok{MIN\_VALUE}\OperatorTok{;}
        \DataTypeTok{double}\NormalTok{ minInterventionUtility }\OperatorTok{=} \BuiltInTok{Double}\OperatorTok{.}\FunctionTok{MAX\_VALUE}\OperatorTok{;}
        \DataTypeTok{double}\NormalTok{ maxInterventionUtility }\OperatorTok{=} \BuiltInTok{Double}\OperatorTok{.}\FunctionTok{MIN\_VALUE}\OperatorTok{;}
        \DataTypeTok{double}\NormalTok{ minReferenceUtility }\OperatorTok{=} \BuiltInTok{Double}\OperatorTok{.}\FunctionTok{MAX\_VALUE}\OperatorTok{;}
        \DataTypeTok{double}\NormalTok{ maxReferenceUtility }\OperatorTok{=} \BuiltInTok{Double}\OperatorTok{.}\FunctionTok{MIN\_VALUE}\OperatorTok{;}
        \ControlFlowTok{if} \OperatorTok{(}\NormalTok{decisionVariable }\OperatorTok{!=} \KeywordTok{null}\OperatorTok{)} \OperatorTok{\{}
            \KeywordTok{...}
\NormalTok{            minValue }\OperatorTok{=} \BuiltInTok{Math}\OperatorTok{.}\FunctionTok{min}\OperatorTok{(}\NormalTok{minInterventionUtility }\OperatorTok{{-}}\NormalTok{ minReferenceUtility}\OperatorTok{,}\NormalTok{ maxInterventionUtility }\OperatorTok{{-}}\NormalTok{ maxReferenceUtility}\OperatorTok{);}
\NormalTok{            maxValue }\OperatorTok{=} \BuiltInTok{Math}\OperatorTok{.}\FunctionTok{max}\OperatorTok{(}\NormalTok{maxInterventionUtility }\OperatorTok{{-}}\NormalTok{ maxReferenceUtility}\OperatorTok{,}\NormalTok{ minInterventionUtility }\OperatorTok{{-}}\NormalTok{ minReferenceUtility}\OperatorTok{);}
        \OperatorTok{\}} \ControlFlowTok{else} \OperatorTok{\{}
            \ControlFlowTok{for} \OperatorTok{(}\DataTypeTok{double}\NormalTok{ value }\OperatorTok{:}\NormalTok{ potential}\OperatorTok{.}\FunctionTok{getValues}\OperatorTok{())} \OperatorTok{\{}
\NormalTok{                minValue }\OperatorTok{=} \BuiltInTok{Math}\OperatorTok{.}\FunctionTok{min}\OperatorTok{(}\NormalTok{minValue}\OperatorTok{,}\NormalTok{ value}\OperatorTok{);}
\NormalTok{                maxValue }\OperatorTok{=} \BuiltInTok{Math}\OperatorTok{.}\FunctionTok{max}\OperatorTok{(}\NormalTok{maxValue}\OperatorTok{,}\NormalTok{ value}\OperatorTok{);}
            \OperatorTok{\}}
        \OperatorTok{\}}
\end{Highlighting}
\end{Shaded}
\paragraph{Cómo arreglarlo}
Iniciar los máximos con \texttt{Double.\allowbreak{}NEGATIVE\_\allowbreak{}INFINITY} (o \texttt{-Double.\allowbreak{}MAX\_\allowbreak{}VALUE}, como ya hace \texttt{get\allowbreak{}Tornado\allowbreak{}Chart} en la línea 384), o arrancar mínimo y máximo con el primer valor del barrido.

Sitios donde debe cambiarse lo mismo: \texttt{Tornado\allowbreak{}Spider\allowbreak{}Dialog} líneas 328, 330, 332 y 533, \texttt{Plot\allowbreak{}Dialog} línea 213 y \texttt{Map\allowbreak{}Dialog} línea 249.

Conviene revisar a la vez el cálculo de la rama con decisión (\hyperref[err:89]{error~89}) y el rango del eje \texttt{min\allowbreak{}Range\allowbreak{}Utility\ *\ 0.\allowbreak{}999,\allowbreak{}\ max\allowbreak{}Range\allowbreak{}Utility\ *\ 1.\allowbreak{}001}, que con valores negativos recorta los extremos en vez de ampliarlos; los dos cambios de esas líneas tienen que ir juntos (\hyperref[err:90]{error~90}).

Prueba de regresión: extraer el cálculo de mínimos y máximos a un método sin interfaz y comprobar que, con valores todos negativos (−1 010 \ldots{} −990), la barra va de −1 010 a −990.

\setcounter{subsection}{88}
\subsection[El tornado por decisión no dibuja el rango real de la diferencia]{En el diagrama de tornado con ámbito de decisión, la barra de un parámetro no abarca lo que varía la diferencia entre las dos opciones, y puede quedar reducida a un punto}\label{err:89}
\begin{ficha}
\fichakey{Estado}Presente
\fichakey{Módulo}sensitivityAnalysis
\fichakey{Paquete}org.openmarkov.sensitivityanalysis.dialog
\fichakey{Ficheros}\path{sensitivityAnalysis/src/main/java/org/openmarkov/sensitivityanalysis/dialog/TornadoSpiderDialog.java:334-350} \newline \path{sensitivityAnalysis/src/main/java/org/openmarkov/sensitivityanalysis/dialog/TornadoSpiderDialog.java:561-571} \newline \path{sensitivityAnalysis/src/main/java/org/openmarkov/sensitivityanalysis/dialog/TornadoSpiderDialog.java:283-289}
\fichakey{Comprobado}Leyendo el código y ejecutándolo
\fichakey{Esfuerzo}Pequeño (horas)
\fichakey{Decisión de equipo}No
\fichakey{Tapado por}\hyperref[err:74]{error~74}: en el jar publicado la ventana no llega a abrirse; desde el IDE se ve.
\fichakey{Relacionados}\hyperref[err:74]{error~74}, \hyperref[err:88]{error~88}, \hyperref[err:90]{error~90}, \hyperref[err:91]{error~91}
\end{ficha}
\paragraph{Qué pasa}
En el análisis de sensibilidad «Tornado/Spider» con ámbito de decisión, el diagrama de tornado debería dibujar, para cada parámetro, una barra que vaya del menor al mayor valor que toma la diferencia de utilidad esperada entre la opción elegida como intervención y la de referencia al barrer el parámetro. La barra que dibuja no sale de esa diferencia, sino de restar por separado los extremos de cada opción, y puede ser más corta que el rango real, o no ser más que un punto, de modo que el parámetro baja en el orden del diagrama. El diagrama de araña de la misma ventana sí resta punto a punto, así que las dos pestañas no cuentan lo mismo.

Camino del usuario: menú Tools → «Sensitivity analysis», aceptar las opciones de inferencia, tipo «Tornado/Spider», ámbito de decisión, elegir la decisión; en la pestaña del tornado, «Intervention:» y «relative to:» son por omisión la segunda y la primera opción. Con el jar ejecutable publicado esta ventana hoy no llega a abrirse (\hyperref[err:74]{error~74}); desde el IDE sí.

Por qué pasa: \texttt{Tornado\allowbreak{}Spider\allowbreak{}Dialog.\allowbreak{}get\allowbreak{}Tornado\allowbreak{}Bars} recorre el barrido de la intervención y guarda su mínimo y su máximo; hace lo mismo con la referencia; y toma como extremos de la barra el menor y el mayor de \texttt{min\allowbreak{}Intervención\ −\ min\allowbreak{}Referencia} y \texttt{max\allowbreak{}Intervención\ −\ max\allowbreak{}Referencia}. Eso solo coincide con el rango de la diferencia cuando los extremos de las dos opciones caen en los mismos puntos del barrido y la diferencia no tiene su extremo en medio. Si una opción sube con el parámetro y la otra baja, los mínimos están en puntos opuestos y la resta mezcla valores de puntos distintos. En cambio, \texttt{get\allowbreak{}Spider\allowbreak{}Chart} calcula en cada punto \texttt{value\allowbreak{}Intervention\ -\ value\allowbreak{}Reference}.

Aparte de esto, los máximos arrancan en \texttt{Double.\allowbreak{}MIN\_\allowbreak{}VALUE}, que falla con utilidades negativas; eso es \hyperref[err:88]{error~88} y se arregla por separado.

Medido:

\begin{itemize}
\tightlist
\item
  Con \texttt{0\_\allowbreak{}miscellaneous/\allowbreak{}networks/\allowbreak{}id/\allowbreak{}ID-decide-test.\allowbreak{}pgmx}, la tarea real \texttt{VESens\allowbreak{}An\allowbreak{}Tornado\allowbreak{}Spider} (variación del 25 \%, 50 puntos) y \texttt{get\allowbreak{}Tornado\allowbreak{}Bars} con las opciones por omisión, en la decisión «Do test?» («yes» respecto a «no») la barra de «utility not treated» va de −0,2000 a 0,1109, y la diferencia punto a punto va de −0,2000 a 0,1991: la barra mide 0,31 en vez de 0,40. Las otras siete barras de esa decisión, y las ocho de «Therapy», coinciden.
\item
  Con barridos escritos a mano de cinco puntos: si «no» va de 0 a −1 y «yes» de 0 a 1, la barra sale {[}1, 1{]}, un punto, cuando la diferencia va de 0 a 2. Si la referencia es constante, o las dos opciones suben a la vez, la barra es correcta.
\end{itemize}
\paragraph{El código}
sensitivityAnalysis/src/main/java/org/openmarkov/sensitivityanalysis/dialog/TornadoSpiderDialog.java:338-350

\begin{Shaded}
\begin{Highlighting}[]
\ControlFlowTok{for} \OperatorTok{(}\DataTypeTok{int}\NormalTok{ i }\OperatorTok{=} \DecValTok{0}\OperatorTok{;}\NormalTok{ i }\OperatorTok{\textless{}=}\NormalTok{ iterations}\OperatorTok{;}\NormalTok{ i}\OperatorTok{++)} \OperatorTok{\{}
    \DataTypeTok{double}\NormalTok{ value }\OperatorTok{=}\NormalTok{ potential}\OperatorTok{.}\FunctionTok{getValues}\OperatorTok{()[}\NormalTok{i }\OperatorTok{+} \OperatorTok{(}\NormalTok{interventionDecisionIndex }\OperatorTok{*} \OperatorTok{(}\NormalTok{iterations }\OperatorTok{+} \DecValTok{1}\OperatorTok{))];}
\NormalTok{    minInterventionUtility }\OperatorTok{=} \BuiltInTok{Math}\OperatorTok{.}\FunctionTok{min}\OperatorTok{(}\NormalTok{minInterventionUtility}\OperatorTok{,}\NormalTok{ value}\OperatorTok{);}
\NormalTok{    maxInterventionUtility }\OperatorTok{=} \BuiltInTok{Math}\OperatorTok{.}\FunctionTok{max}\OperatorTok{(}\NormalTok{maxInterventionUtility}\OperatorTok{,}\NormalTok{ value}\OperatorTok{);}
\OperatorTok{\}}
\CommentTok{// For reference utility values}
\ControlFlowTok{for} \OperatorTok{(}\DataTypeTok{int}\NormalTok{ j }\OperatorTok{=} \DecValTok{0}\OperatorTok{;}\NormalTok{ j }\OperatorTok{\textless{}=}\NormalTok{ iterations}\OperatorTok{;}\NormalTok{ j}\OperatorTok{++)} \OperatorTok{\{}
    \DataTypeTok{double}\NormalTok{ value }\OperatorTok{=}\NormalTok{ potential}\OperatorTok{.}\FunctionTok{getValues}\OperatorTok{()[}\NormalTok{j }\OperatorTok{+} \OperatorTok{(}\NormalTok{relativeDecisionIndex }\OperatorTok{*} \OperatorTok{(}\NormalTok{iterations }\OperatorTok{+} \DecValTok{1}\OperatorTok{))];}
\NormalTok{    minReferenceUtility }\OperatorTok{=} \BuiltInTok{Math}\OperatorTok{.}\FunctionTok{min}\OperatorTok{(}\NormalTok{minReferenceUtility}\OperatorTok{,}\NormalTok{ value}\OperatorTok{);}
\NormalTok{    maxReferenceUtility }\OperatorTok{=} \BuiltInTok{Math}\OperatorTok{.}\FunctionTok{max}\OperatorTok{(}\NormalTok{maxReferenceUtility}\OperatorTok{,}\NormalTok{ value}\OperatorTok{);}
\OperatorTok{\}}
\NormalTok{minValue }\OperatorTok{=} \BuiltInTok{Math}\OperatorTok{.}\FunctionTok{min}\OperatorTok{(}\NormalTok{minInterventionUtility }\OperatorTok{{-}}\NormalTok{ minReferenceUtility}\OperatorTok{,}\NormalTok{ maxInterventionUtility }\OperatorTok{{-}}\NormalTok{ maxReferenceUtility}\OperatorTok{);}
\NormalTok{maxValue }\OperatorTok{=} \BuiltInTok{Math}\OperatorTok{.}\FunctionTok{max}\OperatorTok{(}\NormalTok{maxInterventionUtility }\OperatorTok{{-}}\NormalTok{ maxReferenceUtility}\OperatorTok{,}\NormalTok{ minInterventionUtility }\OperatorTok{{-}}\NormalTok{ minReferenceUtility}\OperatorTok{);}
\end{Highlighting}
\end{Shaded}

sensitivityAnalysis/src/main/java/org/openmarkov/sensitivityanalysis/dialog/TornadoSpiderDialog.java:565-567 (araña)

\begin{Shaded}
\begin{Highlighting}[]
\DataTypeTok{double}\NormalTok{ valueIntervention }\OperatorTok{=}\NormalTok{ potentials}\OperatorTok{[}\NormalTok{i }\OperatorTok{+} \OperatorTok{(}\NormalTok{interventionSpiderDecisionSelector}\OperatorTok{.}\FunctionTok{getSelectedIndex}\OperatorTok{()} \OperatorTok{*} \OperatorTok{(}\NormalTok{iterations }\OperatorTok{+} \DecValTok{1}\OperatorTok{))];}
\DataTypeTok{double}\NormalTok{ valueReference }\OperatorTok{=}\NormalTok{ potentials}\OperatorTok{[}\NormalTok{i }\OperatorTok{+} \OperatorTok{(}\NormalTok{relativeSpiderDecisionSelector}\OperatorTok{.}\FunctionTok{getSelectedIndex}\OperatorTok{()} \OperatorTok{*} \OperatorTok{(}\NormalTok{iterations }\OperatorTok{+} \DecValTok{1}\OperatorTok{))];}
\NormalTok{value }\OperatorTok{=}\NormalTok{ valueIntervention }\OperatorTok{{-}}\NormalTok{ valueReference}\OperatorTok{;}
\end{Highlighting}
\end{Shaded}
\paragraph{Cómo arreglarlo}
Calcular en cada punto del barrido la diferencia \texttt{intervención\ −\ referencia} y tomar el mínimo y el máximo de esas diferencias, como hace la araña. Conviene sacar ese cálculo (los valores de la curva de un parámetro para una pareja de opciones) a un método sin interfaz que usen \texttt{get\allowbreak{}Tornado\allowbreak{}Bars} y \texttt{get\allowbreak{}Spider\allowbreak{}Chart}, para que las dos pestañas no puedan volver a separarse.

Al tocarlo, arrancar el mínimo y el máximo con infinito o con el primer valor, que es el arreglo de \hyperref[err:88]{error~88} para las mismas líneas; el cálculo sin decisión de \texttt{get\allowbreak{}Tornado\allowbreak{}Bars} no tiene este defecto, solo el de \hyperref[err:88]{error~88}.

Prueba de regresión: con el método extraído, un barrido en que «no» baja de 0 a −1 y «yes» sube de 0 a 1 debe dar {[}0, 2{]}; y con \texttt{0\_\allowbreak{}miscellaneous/\allowbreak{}networks/\allowbreak{}id/\allowbreak{}ID-decide-test.\allowbreak{}pgmx}, decisión «Do test?», «yes» respecto a «no», la barra de «utility not treated» debe ir de −0,2000 a 0,1991.

\setcounter{subsection}{89}
\subsection[Con utilidades negativas, «Plot» y la araña aplastan y recortan las curvas]{Con utilidades todas negativas, las gráficas «Plot» y de araña del análisis de sensibilidad aplastan las curvas contra el borde inferior y recortan su punto más bajo, y una utilidad constante negativa no se ve}\label{err:90}
\begin{ficha}
\fichakey{Estado}Presente
\fichakey{Módulo}sensitivityAnalysis
\fichakey{Paquete}org.openmarkov.sensitivityanalysis.dialog
\fichakey{Ficheros}\path{sensitivityAnalysis/src/main/java/org/openmarkov/sensitivityanalysis/dialog/PlotDialog.java:212-213} \newline \path{sensitivityAnalysis/src/main/java/org/openmarkov/sensitivityanalysis/dialog/PlotDialog.java:279-281} \newline \path{sensitivityAnalysis/src/main/java/org/openmarkov/sensitivityanalysis/dialog/TornadoSpiderDialog.java:532-533} \newline \path{sensitivityAnalysis/src/main/java/org/openmarkov/sensitivityanalysis/dialog/TornadoSpiderDialog.java:615-617} \newline \path{sensitivityAnalysis/src/main/java/org/openmarkov/sensitivityanalysis/dialog/MapDialog.java:248-249} \newline \path{sensitivityAnalysis/src/main/java/org/openmarkov/sensitivityanalysis/dialog/MapDialog.java:349-350} \newline \path{sensitivityAnalysis/src/main/java/org/openmarkov/sensitivityanalysis/dialog/MapDialog.java:384}
\fichakey{Comprobado}Leyendo el código y ejecutándolo
\fichakey{Esfuerzo}Pequeño (horas)
\fichakey{Decisión de equipo}No
\fichakey{Corregir antes}\hyperref[err:83]{error~83}: a la vez, como mínimo, en \texttt{Map\allowbreak{}Dialog}: el caso del rango vacío es ese error.
\fichakey{Tapado por}\hyperref[err:74]{error~74}: en el jar publicado; desde el IDE se ve. \newline \hyperref[err:82]{error~82}: en el mapa de ámbito global la zona de datos sale de un solo color.
\fichakey{Relacionados}\hyperref[err:74]{error~74}, \hyperref[err:82]{error~82}, \hyperref[err:83]{error~83}, \hyperref[err:88]{error~88}, \hyperref[err:89]{error~89}
\end{ficha}
\paragraph{Qué pasa}
En el análisis de sensibilidad, cuando todos los valores que se dibujan son negativos, la gráfica de un parámetro («Plot») y el diagrama de araña («Tornado/Spider», pestaña de araña) ponen el eje vertical desde algo por encima del valor más bajo hasta 0. Las curvas quedan aplastadas contra el borde inferior, casi planas, y el punto más bajo de cada una queda fuera del eje y no se dibuja. Si la utilidad es constante y negativa, la curva entera queda fuera y la gráfica sale vacía. Los valores son todos negativos, por ejemplo, en un modelo cuyas utilidades son costes con signo negativo, o, con ámbito de decisión, cuando la opción elegida como intervención es peor que la de referencia en todo el barrido (la araña dibuja la diferencia).

En el mapa de dos parámetros con ámbito global, la escala de grises va del valor más bajo a 0, así que todas las casillas salen casi negras y la leyenda llega hasta 0; hoy esto queda tapado porque la zona de datos sale de un solo color (\hyperref[err:82]{error~82}).

Camino del usuario: menú Tools → «Sensitivity analysis», tipo «Plot», «Tornado/Spider» o «Map (two-way analysis)», sobre una red cuya utilidad esperada sea negativa en todo el intervalo de los parámetros. Desde el IDE; con el jar ejecutable publicado estas ventanas no llegan a abrirse (\hyperref[err:74]{error~74}).

Por qué pasa: \texttt{Plot\allowbreak{}Dialog.\allowbreak{}get\allowbreak{}Plot\allowbreak{}Chart} (línea 213), \texttt{Tornado\allowbreak{}Spider\allowbreak{}Dialog.\allowbreak{}get\allowbreak{}Spider\allowbreak{}Chart} (línea 533) y \texttt{Map\allowbreak{}Dialog.\allowbreak{}get\allowbreak{}Map\allowbreak{}Chart} (línea 249) arrancan el máximo con \texttt{Double.\allowbreak{}MIN\_\allowbreak{}VALUE}, que no es el número más negativo sino el positivo más pequeño (unos 4,9·10⁻³²⁴). Con todos los valores negativos, el máximo se queda en ese número, prácticamente 0. Es el mismo inicio que en las barras del tornado (\hyperref[err:88]{error~88}), pero aquí la consecuencia es el eje o la escala de color.

En «Plot» y en la araña se suma otra cuenta: el eje se fija en \texttt{min\allowbreak{}Range\allowbreak{}Utility\ *\ 0.\allowbreak{}999,\allowbreak{}\ max\allowbreak{}Range\allowbreak{}Utility\ *\ 1.\allowbreak{}001}, que ensancha el rango con valores positivos pero lo estrecha con negativos: el mínimo multiplicado por 0,999 queda por encima del valor más bajo, que se sale.

Medido construyendo las gráficas sin pantalla con JFreeChart 1.5.6, barrido de cinco puntos:

\begin{Verbatim}[breaklines,breakanywhere,fontsize=\small]
Plot y araña, utilidades de -1010 a -990: eje [-1009, 4.9e-324], 1 de 5 puntos fuera, la curva ocupa el 1,9 % de la altura
Plot y araña, utilidad constante -7:      eje [-6.993, 4.9e-324], 5 de 5 puntos fuera
Plot y araña, utilidades de 990 a 1010:   eje [989, 1011], ningún punto fuera, la curva ocupa el 90,9 %
Mapa global, utilidades de -1010 a -990:  escala [-1010, 4.9e-324], gris del valor más bajo 0 y del más alto 5 (de 255)
Mapa global, utilidades de 990 a 1010:    escala [990, 1010], gris 0 y 255
\end{Verbatim}
\paragraph{El código}
sensitivityAnalysis/src/main/java/org/openmarkov/sensitivityanalysis/dialog/PlotDialog.java:212-213, 279-281 (lo mismo en \texttt{Tornado\allowbreak{}Spider\allowbreak{}Dialog.\allowbreak{}java}, líneas 532-533 y 615-617)

\begin{Shaded}
\begin{Highlighting}[]
\DataTypeTok{double}\NormalTok{ minRangeUtility }\OperatorTok{=} \BuiltInTok{Double}\OperatorTok{.}\FunctionTok{MAX\_VALUE}\OperatorTok{;}
\DataTypeTok{double}\NormalTok{ maxRangeUtility }\OperatorTok{=} \BuiltInTok{Double}\OperatorTok{.}\FunctionTok{MIN\_VALUE}\OperatorTok{;}
\KeywordTok{...}
\CommentTok{// Set the range axis}
\NormalTok{NumberAxis rangeAxis }\OperatorTok{=} \OperatorTok{(}\NormalTok{NumberAxis}\OperatorTok{)}\NormalTok{ plot}\OperatorTok{.}\FunctionTok{getRangeAxis}\OperatorTok{();}
\NormalTok{rangeAxis}\OperatorTok{.}\FunctionTok{setRange}\OperatorTok{(}\NormalTok{minRangeUtility }\OperatorTok{*} \FloatTok{0.999}\OperatorTok{,}\NormalTok{ maxRangeUtility }\OperatorTok{*} \FloatTok{1.001}\OperatorTok{);}
\end{Highlighting}
\end{Shaded}
\paragraph{Cómo arreglarlo}
Arrancar el máximo con \texttt{Double.\allowbreak{}NEGATIVE\_\allowbreak{}INFINITY} (o mínimo y máximo con el primer valor) en los tres sitios, que es el mismo cambio que \hyperref[err:88]{error~88} pide para el tornado, y a la vez cambiar la cuenta del eje de «Plot» y de la araña por un margen que no dependa del signo, por ejemplo \texttt{min\ -\ 0.\allowbreak{}001\ *\ \textbar{}max\ -\ min\textbar{}} y \texttt{max\ +\ 0.\allowbreak{}001\ *\ \textbar{}max\ -\ min\textbar{}}, con un margen fijo cuando \texttt{max\ ==\ min}.

Las dos cosas tienen que ir juntas. Con solo el cambio del inicio, una utilidad constante negativa daría un límite inferior mayor que el superior y JFreeChart lanzaría \texttt{Illegal\allowbreak{}Argument\allowbreak{}Exception} (\texttt{Range(\allowbreak{}double,\allowbreak{}\ double):\ require\ lower\ \textless{}=\ upper}, comprobado leyendo \texttt{org.\allowbreak{}jfree.\allowbreak{}data.\allowbreak{}Range}), y una constante 0 daría un rango de longitud cero, que también rechaza. En \texttt{Map\allowbreak{}Dialog} pasa lo mismo con la escala y el eje de la leyenda; allí el caso de rango vacío es \hyperref[err:83]{error~83}, y su arreglo debe ir junto con este.

La araña con ámbito de decisión y la casilla de la línea del 0 marcada ya ensancha el eje hasta incluir el 0 (líneas 633-641); el arreglo no debe romper ese ajuste.

Prueba de regresión: construir sin pantalla la gráfica de \texttt{Plot\allowbreak{}Dialog} y la de la araña con un barrido de −1 010 a −990, otro constante en −7, otro constante en 0 y otro positivo, y comprobar que todos los puntos quedan dentro del eje y que ninguna lanza excepción; y, para el mapa global, que el valor más alto de un barrido negativo se pinta blanco.

\setcounter{subsection}{90}
\subsection[El tornado con dieciséis parámetros o más deja barras fuera del recuadro]{El diagrama de tornado con dieciséis parámetros o más dibuja barras fuera del recuadro, y con veinte algunas no se ven}\label{err:91}
\begin{ficha}
\fichakey{Estado}Presente
\fichakey{Módulo}sensitivityAnalysis
\fichakey{Paquete}org.openmarkov.sensitivityanalysis.dialog
\fichakey{Ficheros}\path{sensitivityAnalysis/src/main/java/org/openmarkov/sensitivityanalysis/dialog/TornadoSpiderDialog.java:373-376} \newline \path{sensitivityAnalysis/src/main/java/org/openmarkov/sensitivityanalysis/dialog/TornadoSpiderDialog.java:420-423}
\fichakey{Comprobado}Leyendo el código y ejecutándolo
\fichakey{Esfuerzo}Pequeño (horas)
\fichakey{Decisión de equipo}No
\fichakey{Tapado por}\hyperref[err:74]{error~74}: en el jar publicado; desde el IDE se ve.
\fichakey{Relacionados}\hyperref[err:74]{error~74}, \hyperref[err:88]{error~88}, \hyperref[err:89]{error~89}
\end{ficha}
\paragraph{Qué pasa}
El análisis de sensibilidad determinista de tipo «tornado/araña» (menú Tools → Sensitivity analysis, aceptar las opciones de inferencia, elegir el tipo de análisis «Tornado/Spider», marcar los parámetros inciertos y pulsar OK) abre \texttt{Tornado\allowbreak{}Spider\allowbreak{}Dialog}, que dibuja un diagrama de tornado: una barra horizontal por cada parámetro incierto marcado. No hay límite al número de parámetros que se pueden marcar: son las casillas de la lista de parámetros del diálogo.

Con 16 parámetros o más, las barras de los extremos se salen del área de dibujo, y con 20 algunas quedan enteras fuera y no se ven, sin que el diagrama avise de que faltan. Como las barras se ordenan de mayor a menor efecto (\texttt{Collections.\allowbreak{}sort(\allowbreak{}tornado\allowbreak{}Bars)}), las que desaparecen por arriba son precisamente las de los parámetros más influyentes.

El método privado \texttt{get\allowbreak{}Tornado\allowbreak{}Chart} reparte el espacio del eje de categorías (el eje vertical, donde van las barras) con una cuenta fija: el margen superior y el inferior valen \texttt{0.\allowbreak{}475\ -\ bars\ *\ 0.\allowbreak{}03}, donde \texttt{bars} es el número de parámetros. Con 15 parámetros el margen vale 0,025; con 16 vale −0,005 y a partir de ahí es negativo.

JFreeChart (la biblioteca de gráficas) no comprueba el signo del margen: con un margen negativo coloca la primera y la última categoría por encima y por debajo del área de dibujo, que recorta lo que se sale.

Medido con un pequeño programa de prueba sin ventana que construye el mismo \texttt{Category\allowbreak{}Plot} (mismos márgenes, \texttt{set\allowbreak{}Category\allowbreak{}Margin(\allowbreak{}0.\allowbreak{}01)}, \texttt{set\allowbreak{}Maximum\allowbreak{}Bar\allowbreak{}Width(\allowbreak{}0.\allowbreak{}05)}, orientación horizontal, título) en una imagen de 600×400 píxeles cuya área de datos va del píxel 30,5 al 378,4:

\begin{itemize}
\tightlist
\item
  5 parámetros: la primera barra empieza en 143,6 y la última acaba en 259,1; todas dentro.
\item
  15 parámetros: de 39,2 a 365,3; todas dentro.
\item
  16 parámetros: de 28,8 a 375,8; la primera barra ya asoma por encima del área.
\item
  20 parámetros: de −13,0 a 417,7; 4 de las 20 barras quedan enteras fuera del área.
\item
  25 parámetros: de −65,1 a 470,0; 8 barras enteras fuera.
\end{itemize}

Los píxeles exactos dependen del tamaño de la ventana y de los rótulos, pero el signo del margen no: con 16 o más parámetros el reparto deja barras fuera.
\paragraph{El código}
sensitivityAnalysis/src/main/java/org/openmarkov/sensitivityanalysis/dialog/TornadoSpiderDialog.java:373-376, 420-423

\begin{Shaded}
\begin{Highlighting}[]
\KeywordTok{private}\NormalTok{ JFreeChart }\FunctionTok{getTornadoChart}\OperatorTok{(}\BuiltInTok{List}\OperatorTok{\textless{}}\NormalTok{TornadoBar}\OperatorTok{\textgreater{}}\NormalTok{ tornadoBars}\OperatorTok{)} \OperatorTok{\{}
    \CommentTok{// JFreeChart attributes definition. Number of bars in the chart, one per parameter}
    \DataTypeTok{int}\NormalTok{ bars }\OperatorTok{=}\NormalTok{ uncertainParameters}\OperatorTok{.}\FunctionTok{size}\OperatorTok{();}
\KeywordTok{...}
    \DataTypeTok{double}\NormalTok{ space }\OperatorTok{=} \FloatTok{0.475} \OperatorTok{{-}} \OperatorTok{(}\NormalTok{bars }\OperatorTok{*} \FloatTok{0.03}\OperatorTok{);}
\NormalTok{    categoryAxis}\OperatorTok{.}\FunctionTok{setLowerMargin}\OperatorTok{(}\NormalTok{space}\OperatorTok{);}
\NormalTok{    categoryAxis}\OperatorTok{.}\FunctionTok{setUpperMargin}\OperatorTok{(}\NormalTok{space}\OperatorTok{);}
\NormalTok{    categoryAxis}\OperatorTok{.}\FunctionTok{setCategoryMargin}\OperatorTok{(}\FloatTok{0.01}\OperatorTok{);}
\end{Highlighting}
\end{Shaded}
\paragraph{Cómo arreglarlo}
Que el margen nunca sea negativo: acotarlo por abajo (por ejemplo \texttt{Math.\allowbreak{}max(\allowbreak{}0.\allowbreak{}0,\allowbreak{}\ .\allowbreak{}.\allowbreak{}.\allowbreak{})} o un mínimo pequeño positivo) o calcularlo de forma que el espacio total de márgenes más categorías no supere el área. Si con muchas barras no caben con una altura legible, la alternativa es hacer crecer la altura del panel con el número de barras (con barra de desplazamiento).

Es el único sitio del módulo con esta cuenta; \texttt{Map\allowbreak{}Dialog} pone márgenes a cero y no se ve afectado.

Prueba: construir el \texttt{JFree\allowbreak{}Chart} que devuelve \texttt{get\allowbreak{}Tornado\allowbreak{}Chart} con 16, 20 y 30 barras, dibujarlo con \texttt{create\allowbreak{}Buffered\allowbreak{}Image(\allowbreak{}600,\allowbreak{}\ 400,\allowbreak{}\ info)} y comprobar con \texttt{Chart\allowbreak{}Rendering\allowbreak{}Info} que todas las \texttt{Category\allowbreak{}Item\allowbreak{}Entity} quedan dentro de \texttt{get\allowbreak{}Plot\allowbreak{}Info(\allowbreak{}).\allowbreak{}get\allowbreak{}Data\allowbreak{}Area(\allowbreak{})}.

\setcounter{subsection}{91}
\subsection[El análisis de sensibilidad no se abre sobre una red sin decisiones]{El análisis de sensibilidad no se abre sobre una red sin decisiones: sale un error inesperado}\label{err:92}
\begin{ficha}
\fichakey{Estado}Presente
\fichakey{Módulo}sensitivityAnalysis
\fichakey{Paquete}org.openmarkov.sensitivityanalysis.dialog
\fichakey{Ficheros}\path{sensitivityAnalysis/src/main/java/org/openmarkov/sensitivityanalysis/dialog/ScopePanel.java:114-132} \newline \path{sensitivityAnalysis/src/main/java/org/openmarkov/sensitivityanalysis/dialog/ScopePanel.java:213-246} \newline \path{sensitivityAnalysis/src/main/java/org/openmarkov/sensitivityanalysis/dialog/SensitivityAnalysisDialog.java:154-181} \newline \path{sensitivityAnalysis/src/main/java/org/openmarkov/sensitivityanalysis/dialog/SensitivityAnalysisDialog.java:320-384} \newline \path{sensitivityAnalysis/src/main/java/org/openmarkov/sensitivityanalysis/SensitivityAnalysisPlugin.java:31-40} \newline \path{sensitivityAnalysis/src/main/java/org/openmarkov/sensitivityanalysis/dialog/SensitivityAnalysisFrameGenerator.java:33-47}
\fichakey{Comprobado}Leyendo el código y ejecutándolo
\fichakey{Esfuerzo}Pequeño (horas)
\fichakey{Decisión de equipo}No
\fichakey{Relacionados}\hyperref[err:65]{error~65}
\end{ficha}
\paragraph{Qué pasa}
Sobre una red sin nodos de decisión que tenga parámetros inciertos con nombre, el análisis de sensibilidad no llega a abrirse: sale un error inesperado. Si la red no tiene parámetros inciertos con nombre, el diálogo sí se abre (con el aviso de pocos parámetros), porque en ese caso no se construye el panel donde está el fallo.

Camino exacto desde el menú:

\begin{itemize}
\tightlist
\item
  Abrir una red sin nodos de decisión que tenga parámetros inciertos con nombre, por ejemplo \texttt{integration\allowbreak{}Tests/\allowbreak{}src/\allowbreak{}main/\allowbreak{}resources/\allowbreak{}networks/\allowbreak{}id/\allowbreak{}ID-one-chance-sa.\allowbreak{}pgmx} (diagrama de influencia con un nodo de azar, una utilidad, ninguna decisión y dos parámetros inciertos con nombre).
\item
  Menú Tools → «Sensitivity analysis». La entrada la añade \texttt{Sensitivity\allowbreak{}Analysis\allowbreak{}Plugin} y está activa en cuanto hay una red abierta.
\item
  Su acción llama a \texttt{Sensitivity\allowbreak{}Analysis\allowbreak{}Frame\allowbreak{}Generator.\allowbreak{}create}, que construye \texttt{Sensitivity\allowbreak{}Analysis\allowbreak{}Controller} (recoge los parámetros inciertos con nombre) y muestra el diálogo modal de opciones de inferencia (\texttt{Inference\allowbreak{}Options\allowbreak{}Dialog}).
\item
  Pulsar OK en ese diálogo. \texttt{create} construye \texttt{Sensitivity\allowbreak{}Analysis\allowbreak{}Dialog}; su constructor llama a \texttt{get\allowbreak{}Content\allowbreak{}Frame(\allowbreak{})}, que arma el selector de tipo de análisis (\texttt{get\allowbreak{}Analysis\allowbreak{}Type\allowbreak{}Panel}: selecciona el primer tipo, «Tornado/Spider» en un análisis de un criterio o «Spider CE» en coste-efectividad, llama a \texttt{update\allowbreak{}Flags(\allowbreak{})} y vuelve a llamar a \texttt{get\allowbreak{}Content\allowbreak{}Frame(\allowbreak{})}).
\item
  Como hay parámetros inciertos, \texttt{get\allowbreak{}Content\allowbreak{}Frame} añade \texttt{Parameters\allowbreak{}Panel}, \texttt{Axis\allowbreak{}Variation\allowbreak{}Panel} y \texttt{get\allowbreak{}Scope\allowbreak{}Panel(\allowbreak{})} → \texttt{new\ Scope\allowbreak{}Panel(\allowbreak{}controller)}.
\item
  Todos los tipos que se ofrecen ponen \texttt{set\allowbreak{}Can\allowbreak{}Be\allowbreak{}Decision(\allowbreak{}true)}, así que \texttt{Scope\allowbreak{}Panel.\allowbreak{}get\allowbreak{}Main\allowbreak{}Panel} llama a \texttt{get\allowbreak{}Decision\allowbreak{}Selector\allowbreak{}Panel(\allowbreak{})}, que llena el desplegable con las decisiones de la red (ninguna) y hace \texttt{decision\allowbreak{}Selector.\allowbreak{}set\allowbreak{}Selected\allowbreak{}Index(\allowbreak{}0)}: \texttt{Illegal\allowbreak{}Argument\allowbreak{}Exception}.
\end{itemize}

La comprobación de que no hay decisiones (líneas 122-127) llega después y solo desactiva controles. La excepción sube por el constructor del diálogo hasta la acción del menú; el manejador global la enseña como error inesperado y el diálogo del análisis de sensibilidad no llega a abrirse.

Medido con un pequeño programa de prueba que construye \texttt{Scope\allowbreak{}Panel} con la configuración que pone «Tornado/Spider» (\texttt{can\allowbreak{}Be\allowbreak{}Global} y \texttt{can\allowbreak{}Be\allowbreak{}Decision} verdaderos). Con \texttt{integration\allowbreak{}Tests/\allowbreak{}src/\allowbreak{}main/\allowbreak{}resources/\allowbreak{}networks/\allowbreak{}id/\allowbreak{}ID-one-chance-sa.\allowbreak{}pgmx} (0 decisiones, 2 parámetros con nombre):

\begin{Verbatim}[breaklines,breakanywhere,fontsize=\small]
java.lang.IllegalArgumentException: setSelectedIndex: 0 out of bounds
   at javax.swing.JComboBox.setSelectedIndex(JComboBox.java:652)
   at ScopePanel.getDecisionSelectorPanel(ScopePanel.java:233)
   at ScopePanel.getMainPanel(ScopePanel.java:120)
   at ScopePanel.<init>(ScopePanel.java:106)
\end{Verbatim}

Con \texttt{integration\allowbreak{}Tests/\allowbreak{}src/\allowbreak{}main/\allowbreak{}resources/\allowbreak{}networks/\allowbreak{}id/\allowbreak{}ID-decide-test-sv.\allowbreak{}pgmx} (2 decisiones) el panel se construye bien.
\paragraph{El código}
sensitivityAnalysis/src/main/java/org/openmarkov/sensitivityanalysis/dialog/ScopePanel.java:119-127

\begin{Shaded}
\begin{Highlighting}[]
\ControlFlowTok{if} \OperatorTok{(}\NormalTok{configuration}\OperatorTok{.}\FunctionTok{isCanBeDecision}\OperatorTok{())} \OperatorTok{\{}
\NormalTok{    mainPanel}\OperatorTok{.}\FunctionTok{add}\OperatorTok{(}\FunctionTok{getDecisionSelectorPanel}\OperatorTok{());}

    \BuiltInTok{List}\OperatorTok{\textless{}}\BuiltInTok{Node}\OperatorTok{\textgreater{}}\NormalTok{ decisionNodes }\OperatorTok{=}\NormalTok{ probNet}\OperatorTok{.}\FunctionTok{getNodes}\OperatorTok{(}\NormalTok{NodeType}\OperatorTok{.}\FunctionTok{DECISION}\OperatorTok{);}
    \ControlFlowTok{if} \OperatorTok{(}\NormalTok{decisionNodes }\OperatorTok{==} \KeywordTok{null} \OperatorTok{||}\NormalTok{ decisionNodes}\OperatorTok{.}\FunctionTok{isEmpty}\OperatorTok{())} \OperatorTok{\{}
        \ControlFlowTok{for} \OperatorTok{(}\BuiltInTok{Component}\NormalTok{ component }\OperatorTok{:}\NormalTok{ scopeTypePanel}\OperatorTok{.}\FunctionTok{getComponents}\OperatorTok{())} \OperatorTok{\{}
\NormalTok{            component}\OperatorTok{.}\FunctionTok{setEnabled}\OperatorTok{(}\KeywordTok{false}\OperatorTok{);}
        \OperatorTok{\}}
    \OperatorTok{\}}
\end{Highlighting}
\end{Shaded}

sensitivityAnalysis/src/main/java/org/openmarkov/sensitivityanalysis/dialog/ScopePanel.java:220-243

\begin{Shaded}
\begin{Highlighting}[]
\NormalTok{decisionSelector }\OperatorTok{=} \KeywordTok{new} \BuiltInTok{JComboBox}\OperatorTok{\textless{}\textgreater{}();}
\ControlFlowTok{for} \OperatorTok{(}\BuiltInTok{Node}\NormalTok{ node }\OperatorTok{:}\NormalTok{ probNet}\OperatorTok{.}\FunctionTok{getNodes}\OperatorTok{(}\NormalTok{NodeType}\OperatorTok{.}\FunctionTok{DECISION}\OperatorTok{))} \OperatorTok{\{}
    \KeywordTok{...}
\NormalTok{    decisionSelector}\OperatorTok{.}\FunctionTok{addItem}\OperatorTok{(}\NormalTok{node}\OperatorTok{.}\FunctionTok{getName}\OperatorTok{());}
\OperatorTok{\}}
\KeywordTok{...}
\NormalTok{decisionSelector}\OperatorTok{.}\FunctionTok{setSelectedIndex}\OperatorTok{(}\DecValTok{0}\OperatorTok{);}
\NormalTok{decisionSelectorPanel}\OperatorTok{.}\FunctionTok{add}\OperatorTok{(}\NormalTok{decisionSelector}\OperatorTok{);}

\ControlFlowTok{if} \OperatorTok{(}\NormalTok{scopeType }\OperatorTok{==}\NormalTok{ ScopeType}\OperatorTok{.}\FunctionTok{GLOBAL}\OperatorTok{)} \OperatorTok{\{}
    \KeywordTok{...}
\OperatorTok{\}} \ControlFlowTok{else} \OperatorTok{\{}
\NormalTok{    decisionSelected }\OperatorTok{=}\NormalTok{ probNet}\OperatorTok{.}\FunctionTok{getVariable}\OperatorTok{(}\NormalTok{decisionSelector}\OperatorTok{.}\FunctionTok{getSelectedItem}\OperatorTok{().}\FunctionTok{toString}\OperatorTok{());}
\OperatorTok{\}}
\end{Highlighting}
\end{Shaded}
\paragraph{Cómo arreglarlo}
Cuando la red no tiene decisiones, el ámbito debe quedarse en «global» y no ofrecer el de decisión. Sitios:

\begin{itemize}
\tightlist
\item
  \texttt{Scope\allowbreak{}Panel.\allowbreak{}get\allowbreak{}Decision\allowbreak{}Selector\allowbreak{}Panel}, línea 233: seleccionar solo si \texttt{decision\allowbreak{}Selector.\allowbreak{}get\allowbreak{}Item\allowbreak{}Count(\allowbreak{})\ \textgreater{}\ 0}, que es exactamente la guarda que ya tiene el panel equivalente de la interfaz general, \texttt{gui/\allowbreak{}src/\allowbreak{}main/\allowbreak{}java/\allowbreak{}org/\allowbreak{}openmarkov/\allowbreak{}gui/\allowbreak{}dialog/\allowbreak{}inference/\allowbreak{}common/\allowbreak{}Scope\allowbreak{}Selector\allowbreak{}Panel.\allowbreak{}java}, líneas 210-212.
\item
  Las lecturas de \texttt{decision\allowbreak{}Selector.\allowbreak{}get\allowbreak{}Selected\allowbreak{}Item(\allowbreak{}).\allowbreak{}to\allowbreak{}String(\allowbreak{})} de las líneas 166 y 242, que darían \texttt{Null\allowbreak{}Pointer\allowbreak{}Exception} con el desplegable vacío si el ámbito queda en «decisión».
\item
  Mejor aún, en \texttt{Sensitivity\allowbreak{}Analysis\allowbreak{}Dialog.\allowbreak{}update\allowbreak{}Flags} (líneas 320-384), poner \texttt{set\allowbreak{}Can\allowbreak{}Be\allowbreak{}Decision(\allowbreak{}false)} cuando la red no tenga decisiones, para que la comprobación se haga antes de construir el panel y no después.
\end{itemize}

Prueba: construir \texttt{Scope\allowbreak{}Panel} (o el diálogo, con una red abierta simulada) sobre \texttt{integration\allowbreak{}Tests/\allowbreak{}src/\allowbreak{}main/\allowbreak{}resources/\allowbreak{}networks/\allowbreak{}id/\allowbreak{}ID-one-chance-sa.\allowbreak{}pgmx} y comprobar que no lanza y que el ámbito queda en \texttt{Scope\allowbreak{}Type.\allowbreak{}GLOBAL}.

\setcounter{subsection}{92}
\subsection[Los parámetros inciertos salen en otro orden en cada apertura del análisis]{Los parámetros inciertos del análisis de sensibilidad salen en un orden distinto cada vez que se abre, y cambia el parámetro que el diálogo propone por omisión}\label{err:93}
\begin{ficha}
\fichakey{Estado}Presente
\fichakey{Módulo}core
\fichakey{Paquete}org.openmarkov.core.model.network.modelUncertainty
\fichakey{Ficheros}\path{core/src/main/java/org/openmarkov/core/model/network/modelUncertainty/SystematicSampling.java:46-96} \newline \path{sensitivityAnalysis/src/main/java/org/openmarkov/sensitivityanalysis/model/SensitivityAnalysisController.java:100-113} \newline \path{sensitivityAnalysis/src/main/java/org/openmarkov/sensitivityanalysis/dialog/ParametersPanel.java:96-166}
\fichakey{Comprobado}Leyendo el código y ejecutándolo
\fichakey{Esfuerzo}Pequeño (horas)
\fichakey{Decisión de equipo}No
\fichakey{Relacionados}\hyperref[err:107]{error~107}
\end{ficha}
\paragraph{Qué pasa}
La lista de parámetros inciertos del diálogo del análisis de sensibilidad aparece en un orden distinto cada vez que se abre, y con ella cambia el parámetro que el diálogo marca por omisión.

Camino del usuario: menú Tools → «Sensitivity analysis» con una red con parámetros inciertos con nombre. \texttt{Sensitivity\allowbreak{}Analysis\allowbreak{}Controller} pide \texttt{Systematic\allowbreak{}Sampling.\allowbreak{}get\allowbreak{}Uncertain\allowbreak{}Parameters(\allowbreak{}prob\allowbreak{}Net)} y guarda los nombres en \texttt{ordered\allowbreak{}Uncertain\allowbreak{}Parameters\allowbreak{}Keys} en el orden recibido; \texttt{Parameters\allowbreak{}Panel} crea una casilla o un botón de opción por parámetro en ese orden. En los análisis de un parámetro («Plot») marca el primero de la lista; en el mapa de dos parámetros («Map») marca el primero para el eje horizontal y el segundo para el vertical.

\texttt{get\allowbreak{}Uncertain\allowbreak{}Parameters(\allowbreak{}Prob\allowbreak{}Net)} recorre los potenciales de la red y, para cada uno, llama a la versión privada que mete los parámetros en un \texttt{Hash\allowbreak{}Set\textless{}Uncertain\allowbreak{}Parameter\textgreater{}} a partir de las claves de un \texttt{Hash\allowbreak{}Map\textless{}Uncertain\allowbreak{}Value,\allowbreak{}\ …\textgreater{}}. Ni \texttt{Uncertain\allowbreak{}Value} ni \texttt{Uncertain\allowbreak{}Parameter} redefinen \texttt{hash\allowbreak{}Code}, así que el orden dentro de cada potencial depende del código de identidad de los objetos, que cambia entre ejecuciones.

Medido con un pequeño programa de prueba sobre \texttt{integration\allowbreak{}Tests/\allowbreak{}src/\allowbreak{}main/\allowbreak{}resources/\allowbreak{}networks/\allowbreak{}id/\allowbreak{}ID-decide-test-sv.\allowbreak{}pgmx}, seis ejecuciones (los parámetros sin nombre, que el diálogo descarta, salen como \texttt{null}):

\begin{Verbatim}[breaklines,breakanywhere,fontsize=\small]
null | prevalence | null | null | sensitivity | specificity | null | null | utility non-treated disease | utility treated disease | utility not treated | cost of test | cost of therapy
null | prevalence | null | null | specificity | sensitivity | null | null | utility not treated | utility non-treated disease | utility treated disease | cost of test | cost of therapy
null | prevalence | null | sensitivity | specificity | null | null | null | utility treated disease | utility not treated | utility non-treated disease | cost of test | cost of therapy
null | prevalence | null | null | sensitivity | null | null | specificity | utility treated disease | utility non-treated disease | utility not treated | cost of test | cost of therapy
null | prevalence | null | null | sensitivity | null | specificity | null | utility not treated | utility non-treated disease | utility treated disease | cost of test | cost of therapy
null | prevalence | null | specificity | null | null | sensitivity | null | utility non-treated disease | utility not treated | utility treated disease | cost of test | cost of therapy
\end{Verbatim}

Los ocho parámetros con nombre son siempre los mismos; lo que cambia es su orden dentro de cada potencial. En esta red, el parámetro que el mapa de dos parámetros propone para el eje vertical es unas veces «sensitivity» y otras «specificity».
\paragraph{El código}
core/src/main/java/org/openmarkov/core/model/network/modelUncertainty/SystematicSampling.java:62-81

\begin{Shaded}
\begin{Highlighting}[]
\KeywordTok{private} \DataTypeTok{static} \BuiltInTok{Set}\OperatorTok{\textless{}}\NormalTok{UncertainParameter}\OperatorTok{\textgreater{}} \FunctionTok{getUncertainParameters}\OperatorTok{(}\NormalTok{Potential potential}\OperatorTok{)} \OperatorTok{\{}
    \BuiltInTok{Set}\OperatorTok{\textless{}}\NormalTok{UncertainParameter}\OperatorTok{\textgreater{}}\NormalTok{ uncertainParams }\OperatorTok{=} \KeywordTok{new} \BuiltInTok{HashSet}\OperatorTok{\textless{}\textgreater{}();}

    \BuiltInTok{Map}\OperatorTok{\textless{}}\NormalTok{UncertainValue}\OperatorTok{,}\NormalTok{ SubPotentialAndPositionInTablePotential}\OperatorTok{\textgreater{}}\NormalTok{ uncertainValues }\OperatorTok{=} \FunctionTok{getUncertainValues}\OperatorTok{(}\NormalTok{potential}\OperatorTok{);}

    \ControlFlowTok{for} \OperatorTok{(}\NormalTok{UncertainValue auxUncertainValue }\OperatorTok{:}\NormalTok{ uncertainValues}\OperatorTok{.}\FunctionTok{keySet}\OperatorTok{())} \OperatorTok{\{}
\NormalTok{        SubPotentialAndPositionInTablePotential subPotentialAndPosition }\OperatorTok{=}\NormalTok{ uncertainValues}\OperatorTok{.}\FunctionTok{get}\OperatorTok{(}\NormalTok{auxUncertainValue}\OperatorTok{);}
\NormalTok{        uncertainParams}
                \OperatorTok{.}\FunctionTok{add}\OperatorTok{(}\KeywordTok{new} \FunctionTok{UncertainParameter}\OperatorTok{(}\NormalTok{potential}\OperatorTok{,}\NormalTok{ auxUncertainValue}\OperatorTok{,}\NormalTok{ subPotentialAndPosition}\OperatorTok{.}\FunctionTok{getSubPotential}\OperatorTok{(),}
\NormalTok{                        subPotentialAndPosition}\OperatorTok{.}\FunctionTok{getPosition}\OperatorTok{()));}
    \OperatorTok{\}}
    \ControlFlowTok{return}\NormalTok{ uncertainParams}\OperatorTok{;}
\OperatorTok{\}}

\KeywordTok{private} \DataTypeTok{static} \BuiltInTok{Map}\OperatorTok{\textless{}}\NormalTok{UncertainValue}\OperatorTok{,}\NormalTok{ SubPotentialAndPositionInTablePotential}\OperatorTok{\textgreater{}} \FunctionTok{getUncertainValues}\OperatorTok{(}\NormalTok{Potential potential}\OperatorTok{)} \OperatorTok{\{}
    \BuiltInTok{Map}\OperatorTok{\textless{}}\NormalTok{UncertainValue}\OperatorTok{,}\NormalTok{ SubPotentialAndPositionInTablePotential}\OperatorTok{\textgreater{}}\NormalTok{ uncertainValuesHash }\OperatorTok{=} \KeywordTok{new} \BuiltInTok{HashMap}\OperatorTok{\textless{}\textgreater{}();}
\end{Highlighting}
\end{Shaded}
\paragraph{Cómo arreglarlo}
Usar \texttt{Linked\allowbreak{}Hash\allowbreak{}Map} en \texttt{get\allowbreak{}Uncertain\allowbreak{}Values} (línea 81) y \texttt{Linked\allowbreak{}Hash\allowbreak{}Set} en la versión privada de \texttt{get\allowbreak{}Uncertain\allowbreak{}Parameters} (línea 63), para conservar el orden en que se recorren las casillas (por índice) y las ramas de los árboles.

El orden entre potenciales lo da \texttt{Prob\allowbreak{}Net.\allowbreak{}get\allowbreak{}Potentials(\allowbreak{})}; conviene comprobar que también es estable (en la medida, el orden entre potenciales no cambió).

Es el mismo patrón (conjuntos y mapas sin orden sobre objetos sin \texttt{hash\allowbreak{}Code}) que causa \hyperref[err:4]{error~4} en las operaciones de potenciales, pero en otro sitio: arreglar aquél no arregla éste.

El otro llamador de \texttt{get\allowbreak{}Uncertain\allowbreak{}Parameters(\allowbreak{}Prob\allowbreak{}Net)} es \texttt{Systematic\allowbreak{}Sampling} línea 413 (búsqueda por nombre), al que el orden no le afecta.

Prueba: llamar varias veces a \texttt{get\allowbreak{}Uncertain\allowbreak{}Parameters} sobre \texttt{integration\allowbreak{}Tests/\allowbreak{}src/\allowbreak{}main/\allowbreak{}resources/\allowbreak{}networks/\allowbreak{}id/\allowbreak{}ID-decide-test-sv.\allowbreak{}pgmx} en procesos distintos (o con varias copias de la red) y comprobar que la lista de nombres es siempre la misma y sigue el orden de las casillas de cada tabla.

\setcounter{subsection}{93}
\subsection[Ocultar una opción en una pestaña deja marcada su casilla en la otra]{En el análisis probabilístico de coste-efectividad, ocultar una opción en una pestaña la oculta en las dos, pero la casilla de la otra pestaña sigue marcada}\label{err:94}
\begin{ficha}
\fichakey{Estado}Presente
\fichakey{Módulo}sensitivityAnalysis
\fichakey{Paquete}org.openmarkov.sensitivityanalysis.dialog
\fichakey{Ficheros}\path{sensitivityAnalysis/src/main/java/org/openmarkov/sensitivityanalysis/dialog/CEProbabilisticDialog.java:99,126,207-242,338-346,412-457}
\fichakey{Comprobado}Leyendo el código
\fichakey{Esfuerzo}Pequeño (horas)
\fichakey{Decisión de equipo}No
\fichakey{Relacionados}\hyperref[nota:132]{nota~132}
\end{ficha}
\paragraph{Qué pasa}
El análisis de sensibilidad probabilístico de coste-efectividad (menú de herramientas → «Sensitivity analysis» sobre una red con análisis de coste-efectividad, tipo «Cost-effectiveness plot») abre \texttt{CEProbabilistic\allowbreak{}Dialog}, con dos pestañas: el plano de coste-efectividad y la curva de aceptabilidad. Cada una tiene a la derecha un recuadro «Show/hide interventions» con una casilla por opción de la decisión.

Desmarcar una opción en la pestaña del plano la quita también de la curva, pero la casilla de esa opción en la pestaña de la curva sigue marcada: esa pestaña enseña una casilla marcada para una opción que no dibuja. Si el usuario vuelve a pulsar esa casilla, la desmarca (y su pulsación deja la opción en «oculto»), con lo que tiene que pulsar dos veces para volver a ver la opción.

\texttt{get\allowbreak{}Show\allowbreak{}Hide\allowbreak{}Panel(\allowbreak{})} fabrica casillas nuevas cada vez que se llama, y se llama dos veces: desde \texttt{get\allowbreak{}Abs\allowbreak{}Rel\allowbreak{}Show\allowbreak{}Hide\allowbreak{}Panel(\allowbreak{})} para la pestaña del plano y desde \texttt{get\allowbreak{}Acceptability\allowbreak{}Curve\allowbreak{}Panel(\allowbreak{})} para la de la curva. Cada casilla, al pulsarla, escribe su estado en el vector compartido \texttt{selected\allowbreak{}States} y llama a \texttt{refresh\allowbreak{}Chart\allowbreak{}Panels(\allowbreak{})}, que rehace las dos gráficas. Ninguno de los dos juegos actualiza las casillas del otro.

No es el mismo defecto que el ya corregido en el commit \texttt{dce0e1b} (el escenario del diálogo de configuración del análisis frente a la evidencia de la red): éste está en la ventana de resultados.

Comprobado leyendo; no se ejecutó porque construir la ventana exige ejecutar la simulación probabilística (\texttt{VECEPSA}) en su constructor.
\paragraph{El código}
sensitivityAnalysis/src/main/java/org/openmarkov/sensitivityanalysis/dialog/CEProbabilisticDialog.java:412-437

\begin{Shaded}
\begin{Highlighting}[]
\KeywordTok{public} \BuiltInTok{JScrollPane} \FunctionTok{getShowHidePanel}\OperatorTok{()} \OperatorTok{\{}
    \BuiltInTok{JPanel}\NormalTok{ showHidePanel }\OperatorTok{=} \KeywordTok{new} \BuiltInTok{JPanel}\OperatorTok{();}
    \CommentTok{// }\AlertTok{TODO}\CommentTok{ {-} LOCALIZE}
\NormalTok{    showHidePanel}\OperatorTok{.}\FunctionTok{setBorder}\OperatorTok{(}\KeywordTok{new} \BuiltInTok{TitledBorder}\OperatorTok{(}\StringTok{"Show/hide interventions"}\OperatorTok{));}
    \KeywordTok{...}
    \BuiltInTok{List}\OperatorTok{\textless{}}\BuiltInTok{JCheckBox}\OperatorTok{\textgreater{}}\NormalTok{ showHideCheckBoxes }\OperatorTok{=} \KeywordTok{new} \BuiltInTok{ArrayList}\OperatorTok{();}
    \ControlFlowTok{for} \OperatorTok{(}\DataTypeTok{int}\NormalTok{ stateIndex }\OperatorTok{=} \DecValTok{0}\OperatorTok{;}\NormalTok{ stateIndex }\OperatorTok{\textless{}}\NormalTok{ decisionVariable}\OperatorTok{.}\FunctionTok{getNumStates}\OperatorTok{();}\NormalTok{ stateIndex}\OperatorTok{++)} \OperatorTok{\{}
        \BuiltInTok{JCheckBox}\NormalTok{ stateCheckbox }\OperatorTok{=} \KeywordTok{new} \BuiltInTok{JCheckBox}\OperatorTok{(}\NormalTok{decisionVariable}\OperatorTok{.}\FunctionTok{getStateName}\OperatorTok{(}\NormalTok{stateIndex}\OperatorTok{));}
\NormalTok{        stateCheckbox}\OperatorTok{.}\FunctionTok{addActionListener}\OperatorTok{(}\KeywordTok{new} \BuiltInTok{ActionListener}\OperatorTok{()} \OperatorTok{\{}
            \AttributeTok{@Override} \KeywordTok{public} \DataTypeTok{void} \FunctionTok{actionPerformed}\OperatorTok{(}\BuiltInTok{ActionEvent}\NormalTok{ e}\OperatorTok{)} \OperatorTok{\{}
                \BuiltInTok{JCheckBox}\NormalTok{ checkBox }\OperatorTok{=} \OperatorTok{(}\BuiltInTok{JCheckBox}\OperatorTok{)}\NormalTok{ e}\OperatorTok{.}\FunctionTok{getSource}\OperatorTok{();}
                \ControlFlowTok{for} \OperatorTok{(}\DataTypeTok{int}\NormalTok{ stateIndex }\OperatorTok{=} \DecValTok{0}\OperatorTok{;}\NormalTok{ stateIndex }\OperatorTok{\textless{}}\NormalTok{ decisionVariable}\OperatorTok{.}\FunctionTok{getNumStates}\OperatorTok{();}\NormalTok{ stateIndex}\OperatorTok{++)} \OperatorTok{\{}
                    \ControlFlowTok{if} \OperatorTok{(}\NormalTok{decisionVariable}\OperatorTok{.}\FunctionTok{getStateName}\OperatorTok{(}\NormalTok{stateIndex}\OperatorTok{).}\FunctionTok{equals}\OperatorTok{(}\NormalTok{checkBox}\OperatorTok{.}\FunctionTok{getText}\OperatorTok{()))} \OperatorTok{\{}
\NormalTok{                        selectedStates}\OperatorTok{[}\NormalTok{stateIndex}\OperatorTok{]} \OperatorTok{=}\NormalTok{ checkBox}\OperatorTok{.}\FunctionTok{isSelected}\OperatorTok{();}
                    \OperatorTok{\}}
                \OperatorTok{\}}
                \FunctionTok{refreshChartPanels}\OperatorTok{();}
            \OperatorTok{\}}
        \OperatorTok{\});}
\NormalTok{        stateCheckbox}\OperatorTok{.}\FunctionTok{setSelected}\OperatorTok{(}\KeywordTok{true}\OperatorTok{);}
\NormalTok{        selectedStates}\OperatorTok{[}\NormalTok{stateIndex}\OperatorTok{]} \OperatorTok{=} \KeywordTok{true}\OperatorTok{;}
\end{Highlighting}
\end{Shaded}
\paragraph{Cómo arreglarlo}
Dos opciones:

\begin{itemize}
\tightlist
\item
  Un solo juego de casillas compartido por las dos pestañas. Swing no permite un mismo componente en dos contenedores, así que habría que moverlo a la pestaña visible al cambiar de pestaña.
\item
  Mantener dos juegos y, al pulsar una casilla, poner la casilla de la misma opción del otro juego en el mismo estado, guardando los dos juegos en campos, como hace \texttt{CEDecision\allowbreak{}Results} con \texttt{ce\allowbreak{}Plane\allowbreak{}Show\allowbreak{}Hide\allowbreak{}Check\allowbreak{}Boxes}.
\end{itemize}

Si se prefiere que cada pestaña filtre por su cuenta, entonces cada juego necesita su propio vector de estados y \texttt{refresh\allowbreak{}Chart\allowbreak{}Panels} debe rehacer solo su gráfica; es una decisión de interfaz.

La misma estructura aparece en \texttt{inference.\allowbreak{}des/\allowbreak{}src/\allowbreak{}main/\allowbreak{}java/\allowbreak{}org/\allowbreak{}openmarkov/\allowbreak{}inference/\allowbreak{}DES/\allowbreak{}CEDESDialog.\allowbreak{}java} (ventana del plano de la simulación de sucesos discretos), aunque allí la pestaña de la curva de aceptabilidad está comentada en \texttt{get\allowbreak{}Tabbed\allowbreak{}Pane}, \texttt{get\allowbreak{}Show\allowbreak{}Hide\allowbreak{}Panel} se usa una sola vez y no se produce el desajuste.

Prueba: construir el diálogo sin su constructor con resultados escritos a mano, desmarcar una casilla del plano y comprobar que la de la curva queda desmarcada.

% ══════════════════════════════════════════════════════════════════════
\section{Simulación de sucesos discretos}\label{sec:des}

Las redes de simulación de eventos discretos (DES), en las que el tiempo avanza de suceso en suceso, y la simulación de Monte Carlo y la propagación estocástica.

\setcounter{subsection}{94}
\subsection[Una utilidad acumulativa vale cero hasta el primer suceso que la alcanza]{En la simulación de sucesos, una utilidad acumulativa que depende de un nodo de azar vale cero hasta que llega el primer suceso, y el resultado sale más bajo}\label{err:95}
\begin{ficha}
\fichakey{Estado}Presente
\fichakey{Módulo}inference.des
\fichakey{Paquete}org.openmarkov.inference.DES
\fichakey{Ficheros}\path{inference.des/src/main/java/org/openmarkov/inference/DES/UtilityEvaluation.java:50-72} \newline \path{inference.des/src/main/java/org/openmarkov/inference/DES/UtilityEvaluation.java:80-103} \newline \path{inference.des/src/main/java/org/openmarkov/inference/DES/UtilityRecord.java:86-95} \newline \path{inference.des/src/main/java/org/openmarkov/inference/DES/ChanceEvaluation.java:51-89} \newline \path{inference.des/src/main/java/org/openmarkov/inference/DES/GenericEvaluation.java:83-86} \newline \path{inference.des/src/main/java/org/openmarkov/inference/DES/DESRecord.java:120-123} \newline \path{inference.des/src/main/java/org/openmarkov/inference/DES/DESRecord.java:195-197} \newline \path{inference.des/src/main/java/org/openmarkov/inference/DES/DESInference.java:315-333}
\fichakey{Comprobado}Leyendo el código y ejecutándolo
\fichakey{Esfuerzo}Medio (días)
\fichakey{Decisión de equipo}\textcolor{decisionrojo}{\textbf{Sí.}} Qué valor tiene en el instante cero un nodo de azar que tiene un suceso como antecesor (hoy, implícitamente, su primer estado o 0,0, sin muestrearlo), porque de ese valor depende la utilidad del primer tramo
\fichakey{Relacionados}\hyperref[err:96]{error~96}
\end{ficha}
\paragraph{Qué pasa}
Una utilidad acumulativa en una red DES (simulación de eventos discretos) es una tasa ---coste por año, calidad de vida--- que se acumula multiplicando su valor instantáneo por el tiempo que lo mantiene. En \texttt{Utility\allowbreak{}Record} es acumulativa toda utilidad sin sucesos padres.

Si una utilidad acumulativa depende de un nodo de azar que un suceso puede cambiar ---el caso normal: la calidad de vida depende de estar enfermo, y se enferma cuando ocurre el suceso---, desde el instante 0 hasta el primer suceso que la alcanza vale 0. La simulación no avisa: da una efectividad y un coste más bajos.

Camino de llegada: cualquier red DES con una utilidad acumulativa que cuelgue de un nodo de azar con un suceso antecesor, simulada desde el botón o menú de análisis de coste-efectividad (en una \texttt{DESNet}, \texttt{Main\allowbreak{}Panel\allowbreak{}Listener\allowbreak{}Assistant} desvía \texttt{COST\_\allowbreak{}EFFECTIVENESS\_\allowbreak{}DETERMINISTIC} a \texttt{monte\allowbreak{}Carlo\allowbreak{}Simulation}).

El valor instantáneo de una utilidad se fija en dos momentos:

\begin{itemize}
\tightlist
\item
  al empezar cada individuo, \texttt{Utility\allowbreak{}Evaluation.\allowbreak{}start\allowbreak{}Simulation} llama a \texttt{update\allowbreak{}Orphan\allowbreak{}Nodes}, que solo recorre \texttt{orphan\allowbreak{}Records}: las utilidades «huérfanas», las que no tienen ningún suceso como antecesor (\texttt{DESRecord.\allowbreak{}is\allowbreak{}Orphan\allowbreak{}Node});
\item
  cuando ocurre un suceso, \texttt{Utility\allowbreak{}Evaluation.\allowbreak{}update} recalcula las utilidades que ese suceso alcanza.
\end{itemize}

Una utilidad que depende de un nodo de azar con un suceso antecesor no es huérfana, así que nadie le da valor al empezar: conserva el 0 que le pone \texttt{clear(\allowbreak{})}. El tramo desde el instante 0 hasta el primer suceso que la alcanza se acumula multiplicado por 0.

Los nodos de azar tienen el mismo origen: \texttt{Chance\allowbreak{}Evaluation.\allowbreak{}start\allowbreak{}Simulation} solo evalúa los nodos de azar huérfanos (\texttt{add\allowbreak{}Values\allowbreak{}Orphan\allowbreak{}Nodes}). Un nodo de azar con un suceso antecesor empieza con \texttt{variable\allowbreak{}Value} = 0 de \texttt{clear(\allowbreak{})}, es decir, en su primer estado (o 0,0 si es numérico), sin muestrear su potencial. Para dar valor a la utilidad en el instante 0 hay que usar el valor de ese nodo en el instante 0, y ese valor hoy no lo define el algoritmo: de ahí la decisión de equipo.

El repositorio no tiene ninguna red de sucesos. La medida se hizo con una red de sucesos de prueba con:

\begin{itemize}
\tightlist
\item
  una decisión \texttt{Therapy};
\item
  un suceso \texttt{Onset} que hace pasar el nodo de azar \texttt{Ill} de «no» a «yes», y un suceso terminal \texttt{Death};
\item
  tres utilidades acumulativas: \texttt{QALY} = 1,0 si \texttt{Ill}=no y 0,6 si yes; \texttt{Care\allowbreak{}Cost} = 100 si no y 2000 si yes; \texttt{Monitoring} = 10, que solo depende de la decisión;
\item
  horizonte de 10 años, una simulación con la traza textual activada.
\end{itemize}

Para \texttt{Therapy}=none, \texttt{Onset} ocurre en el año 3 y \texttt{Death} en el 6. La traza dice:

\begin{Verbatim}[breaklines,breakanywhere,fontsize=\small]
Accrue cumulative utilities/payoffs from clock=0 to clock=3 (interval length 3)
    Monitoring   Instantaneous value in the interval: 10.0
    QALY         Instantaneous value in the interval: 0.0
    CareCost     Instantaneous value in the interval: 0.0
\end{Verbatim}

\texttt{Monitoring} es huérfana y vale lo que debe. \texttt{QALY} y \texttt{Care\allowbreak{}Cost} valen 0 durante tres años en los que el paciente está sano. El resumen da efectividad media 1,8 (solo 0,6 × 3 años enfermo) cuando el modelo dice 4,8 (3 × 1,0 + 3 × 0,6), y coste 6120 cuando serían 6420 (faltan 3 × 100).

En esa red, el valor inicial de \texttt{Ill} coincide con lo correcto porque «no» es su primer estado; con los estados en otro orden el paciente empezaría enfermo.

Para reproducirlo hay que construir en el editor una red de sucesos como la descrita.
\paragraph{El código}
inference.des/src/main/java/org/openmarkov/inference/DES/UtilityEvaluation.java:50-72

\begin{Shaded}
\begin{Highlighting}[]
    \AttributeTok{@Override}
    \KeywordTok{public} \DataTypeTok{void} \FunctionTok{startSimulation}\OperatorTok{(}\NormalTok{Finding decisionFinding}\OperatorTok{)} \KeywordTok{throws}\NormalTok{ OpenMarkovException }\OperatorTok{\{}
\NormalTok{        desRecordHashMap}\OperatorTok{.}\FunctionTok{values}\OperatorTok{().}\FunctionTok{forEach}\OperatorTok{(}\NormalTok{UtilityRecord}\OperatorTok{::}\NormalTok{clear}\OperatorTok{);}
        \KeywordTok{this}\OperatorTok{.}\FunctionTok{decisionFinding} \OperatorTok{=}\NormalTok{ decisionFinding}\OperatorTok{;}
        \FunctionTok{updateOrphanNodes}\OperatorTok{();}
\NormalTok{        previousEvaluationTime }\OperatorTok{=} \DecValTok{0}\OperatorTok{;}
    \OperatorTok{\}}

    \CommentTok{/**}
     \CommentTok{*}\NormalTok{ Update orphan UTILITY nodes}
     \CommentTok{*/}
    \DataTypeTok{void} \FunctionTok{updateOrphanNodes}\OperatorTok{()} \KeywordTok{throws}\NormalTok{ OpenMarkovException }\OperatorTok{\{}
        \ControlFlowTok{for} \OperatorTok{(}\DataTypeTok{var}\NormalTok{ utilityRecord }\OperatorTok{:}\NormalTok{ orphanRecords}\OperatorTok{)} \OperatorTok{\{}
                \BuiltInTok{Configuration}\NormalTok{ configuration }\OperatorTok{=} \KeywordTok{new} \BuiltInTok{Configuration}\OperatorTok{();}
                \ControlFlowTok{if} \OperatorTok{(}\NormalTok{utilityRecord}\OperatorTok{.}\FunctionTok{hasDecisionParent}\OperatorTok{())} \OperatorTok{\{}
\NormalTok{                    configuration}\OperatorTok{.}\FunctionTok{addFinding}\OperatorTok{(}\NormalTok{decisionFinding}\OperatorTok{);}
                \OperatorTok{\}}
                \ControlFlowTok{for} \OperatorTok{(}\BuiltInTok{Node}\NormalTok{ chanceParent }\OperatorTok{:}\NormalTok{ utilityRecord}\OperatorTok{.}\FunctionTok{getParentsByType}\OperatorTok{(}\NormalTok{NodeType}\OperatorTok{.}\FunctionTok{CHANCE}\OperatorTok{))} \OperatorTok{\{}
\NormalTok{                    Variable parentVariable }\OperatorTok{=}\NormalTok{ chanceParent}\OperatorTok{.}\FunctionTok{getVariable}\OperatorTok{();}
\NormalTok{                    configuration}\OperatorTok{.}\FunctionTok{addFinding}\OperatorTok{(}\NormalTok{parentVariable}\OperatorTok{,}\NormalTok{ desInference}\OperatorTok{.}\FunctionTok{getChanceEvaluation}\OperatorTok{().}\FunctionTok{getDESRecord}\OperatorTok{(}\NormalTok{chanceParent}\OperatorTok{).}\FunctionTok{getVariableValue}\OperatorTok{());}
                \OperatorTok{\}}
\NormalTok{                utilityRecord}\OperatorTok{.}\FunctionTok{setVariableValue}\OperatorTok{(}\NormalTok{configuration}\OperatorTok{);}
        \OperatorTok{\}}
    \OperatorTok{\}}
\end{Highlighting}
\end{Shaded}

inference.des/src/main/java/org/openmarkov/inference/DES/UtilityRecord.java:86-89

\begin{Shaded}
\begin{Highlighting}[]
    \KeywordTok{public} \DataTypeTok{void} \FunctionTok{accrueCumulativeUtility}\OperatorTok{(}\DataTypeTok{double}\NormalTok{ previousEvaluationTime}\OperatorTok{,} \DataTypeTok{double}\NormalTok{ clock}\OperatorTok{)} \OperatorTok{\{}
        \CommentTok{//Cumulative Utility}
        \DataTypeTok{double}\NormalTok{ utility }\OperatorTok{=}\NormalTok{ variableValue }\OperatorTok{*} \OperatorTok{(}\NormalTok{clock }\OperatorTok{{-}}\NormalTok{ previousEvaluationTime}\OperatorTok{);}
\NormalTok{        accruedUtility }\OperatorTok{+=}\NormalTok{ utility}\OperatorTok{;}
\end{Highlighting}
\end{Shaded}
\paragraph{Cómo arreglarlo}
Dar valor al empezar a todas las utilidades acumulativas, no solo a las huérfanas, con los valores de sus padres en el instante 0. Para eso hace falta antes que todos los nodos de azar tengan en el instante 0 un valor que el modelo respalde, que es la decisión de equipo: hoy los no huérfanos tienen el 0 implícito de \texttt{clear(\allowbreak{})}.

Sitios que tienen que cumplir la misma regla («en el instante 0 todo nodo con valor instantáneo tiene el valor que dicta el modelo»):

\begin{itemize}
\tightlist
\item
  \texttt{Chance\allowbreak{}Evaluation.\allowbreak{}start\allowbreak{}Simulation} / \texttt{add\allowbreak{}Values\allowbreak{}Orphan\allowbreak{}Nodes} (nodos de azar).
\item
  \texttt{Utility\allowbreak{}Evaluation.\allowbreak{}start\allowbreak{}Simulation} / \texttt{update\allowbreak{}Orphan\allowbreak{}Nodes} (utilidades).
\item
  \texttt{Event\allowbreak{}Evaluation.\allowbreak{}update\allowbreak{}Orphan\allowbreak{}Events}, que calcula los tiempos de los sucesos huérfanos con los valores de sus padres de azar en el instante 0.
\item
  \texttt{Generic\allowbreak{}Evaluation.\allowbreak{}set\allowbreak{}Orphan\allowbreak{}Nodes} y \texttt{DESRecord.\allowbreak{}is\allowbreak{}Orphan\allowbreak{}Node}, que definen hoy quién se evalúa al empezar.
\end{itemize}

\texttt{Utility\allowbreak{}Evaluation.\allowbreak{}accrue\allowbreak{}Immediate\allowbreak{}Utility} construye una configuración que no usa (líneas 116-124); no es este defecto, pero conviene no copiar ese patrón.

Pruebas de regresión: la red descrita (sano al empezar, suceso en el año 3, muerte en el 6, con números exactos) debe dar efectividad 4,8 y coste 6420 para \texttt{Therapy}=none; y otra con los estados de \texttt{Ill} en orden inverso debe dar el mismo resultado.

\setcounter{subsection}{95}
\subsection[La simulación se cae si un nodo espera a un padre de un suceso no ocurrido]{La simulación de sucesos se cae cuando un suceso actualiza un nodo de azar cuyo padre de azar depende de otro suceso que aún no ha ocurrido}\label{err:96}
\begin{ficha}
\fichakey{Estado}Presente
\fichakey{Módulo}inference.des
\fichakey{Paquete}org.openmarkov.inference.DES
\fichakey{Ficheros}\path{inference.des/src/main/java/org/openmarkov/inference/DES/ChanceEvaluation.java:63-89, 97-113} \newline \path{inference.des/src/main/java/org/openmarkov/inference/DES/ChanceRecord.java:90-123} \newline \path{gui/src/main/java/org/openmarkov/gui/window/InferenceHandler.java:245-297}
\fichakey{Comprobado}Leyendo el código y ejecutándolo
\fichakey{Esfuerzo}Pequeño (horas)
\fichakey{Decisión de equipo}No
\fichakey{Relacionados}\hyperref[err:95]{error~95}, \hyperref[err:97]{error~97}, \hyperref[err:98]{error~98}
\end{ficha}
\paragraph{Qué pasa}
En una red DES (simulación de eventos discretos), si un suceso E1 actualiza un nodo de azar A que tiene como padre de azar un nodo B que depende de otro suceso E2 aún no ocurrido, la simulación se cae con \texttt{Index\allowbreak{}Out\allowbreak{}Of\allowbreak{}Bounds\allowbreak{}Exception} en cuanto ocurre E1. El mensaje no dice qué nodos se esperan. No hace falta ningún ciclo en la red.

Camino de llegada: construir en el editor una red DES así (los enlaces y la tabla de transiciones con dos padres se crean desde la interfaz) y lanzar la simulación con el botón o menú de coste-efectividad. La excepción no la captura nadie en el hilo que lanza \texttt{Main\allowbreak{}Panel\allowbreak{}Listener\allowbreak{}Assistant.\allowbreak{}monte\allowbreak{}Carlo\allowbreak{}Simulation}, así que llega al manejador global \texttt{OMException\allowbreak{}Handler}, que abre el diálogo de error inesperado (esto último, por lectura).

Cuando ocurre un suceso, \texttt{Chance\allowbreak{}Evaluation.\allowbreak{}update} recalcula los nodos de azar descendientes de ese suceso. Para respetar el orden (un nodo después de sus padres de azar) repite: busca el primer pendiente cuyos padres de azar estén todos marcados como evaluados (\texttt{Chance\allowbreak{}Record.\allowbreak{}is\allowbreak{}Updatable}), lo calcula y lo quita de la lista. La búsqueda es \texttt{while\ (\allowbreak{}!remaining.\allowbreak{}get(\allowbreak{}i).\allowbreak{}is\allowbreak{}Updatable(\allowbreak{}))\ ++i;}, sin tope: si ningún pendiente está listo, \texttt{i} llega al tamaño de la lista y \texttt{get(\allowbreak{}i)} lanza la excepción.

\texttt{is\allowbreak{}Updatable} exige que \textbf{todos} los padres de azar estén marcados como evaluados, también los que no están entre los pendientes. Un nodo de azar con un suceso antecesor no se evalúa al empezar la simulación (ver \hyperref[err:95]{error~95}) y su marca \texttt{evaluated} queda a \texttt{false} desde \texttt{clear(\allowbreak{})} hasta que ocurre su propio suceso. Por eso A nunca está «listo» y la lista revienta.

El otro bucle igual, en \texttt{add\allowbreak{}Values\allowbreak{}Orphan\allowbreak{}Nodes} (línea 72), solo recorre nodos huérfanos, cuyos padres de azar son también huérfanos y están en la misma lista; allí solo se llegaría con un ciclo entre nodos de azar, y ese caso se cae antes en otro sitio (\hyperref[err:97]{error~97}).

El repositorio no tiene ninguna red de sucesos. La medida se hizo con una red de sucesos de prueba con dos sucesos: \texttt{Onset} (año 3), que hace pasar \texttt{Ill} a «yes», y \texttt{Accident} (año 5), que hace pasar \texttt{Injured} a «yes»; y un enlace \texttt{Injured\ →\ Ill}, con la tabla de \texttt{Ill} sobre \texttt{{[}Ill,\allowbreak{}\ Onset,\allowbreak{}\ Injured{]}}. Al ocurrir \texttt{Onset} la simulación termina con \texttt{Index\allowbreak{}Out\allowbreak{}Of\allowbreak{}Bounds\allowbreak{}Exception:\ Index\ 1\ out\ of\ bounds\ for\ length\ 1} en \texttt{Chance\allowbreak{}Evaluation.\allowbreak{}update} (línea 107).
\paragraph{El código}
inference.des/src/main/java/org/openmarkov/inference/DES/ChanceEvaluation.java:103-110

\begin{Shaded}
\begin{Highlighting}[]
\NormalTok{        remainingDescendants}\OperatorTok{.}\FunctionTok{forEach}\OperatorTok{(}\NormalTok{ChanceRecord}\OperatorTok{::}\NormalTok{resetEvaluationStatus}\OperatorTok{);}

        \ControlFlowTok{while} \OperatorTok{(!}\NormalTok{remainingDescendants}\OperatorTok{.}\FunctionTok{isEmpty}\OperatorTok{())} \OperatorTok{\{}
            \DataTypeTok{int}\NormalTok{ i }\OperatorTok{=} \DecValTok{0}\OperatorTok{;}
            \ControlFlowTok{while} \OperatorTok{(!}\NormalTok{remainingDescendants}\OperatorTok{.}\FunctionTok{get}\OperatorTok{(}\NormalTok{i}\OperatorTok{).}\FunctionTok{isUpdatable}\OperatorTok{())}
                \OperatorTok{++}\NormalTok{i}\OperatorTok{;}
\NormalTok{            ChanceRecord chanceRecord }\OperatorTok{=}\NormalTok{ remainingDescendants}\OperatorTok{.}\FunctionTok{remove}\OperatorTok{(}\NormalTok{i}\OperatorTok{);}
            \FunctionTok{computeChanceValue}\OperatorTok{(}\NormalTok{chanceRecord}\OperatorTok{,}\NormalTok{ eventHappened}\OperatorTok{);}
        \OperatorTok{\}}
\end{Highlighting}
\end{Shaded}

inference.des/src/main/java/org/openmarkov/inference/DES/ChanceRecord.java:120-123

\begin{Shaded}
\begin{Highlighting}[]
    \KeywordTok{public} \DataTypeTok{boolean} \FunctionTok{isUpdatable}\OperatorTok{()} \OperatorTok{\{}
        \CommentTok{//08/03/2023; Self{-}loop{-}{-}\textgreater{}first extracted parents different from node}
        \ControlFlowTok{return}\NormalTok{ chanceRecordParents}\OperatorTok{.}\FunctionTok{stream}\OperatorTok{().}\FunctionTok{filter}\OperatorTok{(}\NormalTok{chanceRecord }\OperatorTok{{-}\textgreater{}} \OperatorTok{!}\NormalTok{chanceRecord}\OperatorTok{.}\FunctionTok{equals}\OperatorTok{(}\KeywordTok{this}\OperatorTok{)).}\FunctionTok{allMatch}\OperatorTok{(}\NormalTok{ChanceRecord}\OperatorTok{::}\NormalTok{isEvaluated}\OperatorTok{);}
    \OperatorTok{\}}
\end{Highlighting}
\end{Shaded}
\paragraph{Cómo arreglarlo}
Dos cambios que van juntos:

\begin{itemize}
\tightlist
\item
  Que un padre de azar solo haga esperar si está entre los nodos que se van a recalcular en esta pasada; los demás ya tienen su valor vigente. Hoy la marca \texttt{evaluated} mezcla «calculado en esta pasada» con «calculado alguna vez».
\item
  Poner tope a la búsqueda y, si no queda ningún pendiente listo, lanzar una excepción que nombre los nodos que se esperan entre sí (solo posible con un ciclo).
\end{itemize}

Sitios con la misma regla: \texttt{Chance\allowbreak{}Evaluation.\allowbreak{}update} (líneas 105-110), \texttt{Chance\allowbreak{}Evaluation.\allowbreak{}add\allowbreak{}Values\allowbreak{}Orphan\allowbreak{}Nodes} (líneas 70-73), \texttt{Chance\allowbreak{}Record.\allowbreak{}is\allowbreak{}Updatable} / \texttt{reset\allowbreak{}Evaluation\allowbreak{}Status} / \texttt{clear}.

El valor que tiene un nodo de azar no huérfano antes de su primer suceso es la decisión de equipo de \hyperref[err:95]{error~95}. Este arreglo no la necesita (basta con no esperar al padre), pero la simulación usará ese valor.

Prueba de regresión: la red descrita (dos sucesos independientes, el nodo actualizado por el primero con un padre de azar que depende del segundo) debe simular sin excepción, y \texttt{Ill} debe pasar a «yes» en el año 3.

\setcounter{subsection}{96}
\subsection[Un ciclo entre nodos de azar de una red de sucesos impide simular]{Un ciclo entre nodos de azar en una red de sucesos impide simular: hoy desborda la pila y, detrás, hay un bucle infinito}\label{err:97}
\begin{ficha}
\fichakey{Estado}Presente
\fichakey{Módulo}inference.des
\fichakey{Paquete}org.openmarkov.inference.DES
\fichakey{Ficheros}\path{inference.des/src/main/java/org/openmarkov/inference/DES/DESRecord.java:153-178} \newline \path{inference.des/src/main/java/org/openmarkov/inference/DES/EventRecord.java:83-125} \newline \path{inference.des/src/main/java/org/openmarkov/inference/DES/GenericEvaluation.java:63-80} \newline \path{inference.des/src/main/java/org/openmarkov/inference/DES/ChanceEvaluation.java:70-73} \newline \path{core/src/main/java/org/openmarkov/core/model/network/type/DESNetworkType.java:30-43}
\fichakey{Comprobado}Leyendo el código y ejecutándolo
\fichakey{Esfuerzo}Medio (días)
\fichakey{Decisión de equipo}\textcolor{decisionrojo}{\textbf{Sí.}} Si una red DES puede tener ciclos entre nodos de azar (y entonces con qué orden se evalúan) o si debe prohibirlos una restricción del tipo de red
\fichakey{Relacionados}\hyperref[err:96]{error~96}, \hyperref[err:106]{error~106}
\end{ficha}
\paragraph{Qué pasa}
El tipo de red \texttt{DESNet} (red de simulación de eventos discretos, DES) desactiva a propósito las restricciones \texttt{No\allowbreak{}Cycle}, \texttt{No\allowbreak{}Loops} y \texttt{No\allowbreak{}Self\allowbreak{}Loop}, para que un suceso pueda volver a citarse a sí mismo. Por eso el editor deja crear dos nodos de azar que dependen el uno del otro: medido, en una red DES, \texttt{Add\allowbreak{}Link\allowbreak{}Edit} acepta \texttt{Ill\ →\ Mood} y después \texttt{Mood\ →\ Ill}.

Con esa red, la simulación no llega a empezar: sale un \texttt{Stack\allowbreak{}Overflow\allowbreak{}Error}, que en la aplicación llega al manejador global \texttt{OMException\allowbreak{}Handler} y se muestra como error inesperado (por lectura), sin decir que el problema es un ciclo.

Al construir los registros de los nodos de azar, el constructor de \texttt{DESRecord} llama a \texttt{set\allowbreak{}Event\allowbreak{}Ancestors}, y \texttt{add\allowbreak{}Event\allowbreak{}Ancestors} sube por los padres recursivamente sin llevar cuenta de los ya visitados; solo se salta el propio nodo (lazo sobre sí mismo). Con \texttt{Ill\ ↔\ Mood} va de uno a otro hasta desbordar la pila.

Medido: \texttt{Stack\allowbreak{}Overflow\allowbreak{}Error} en \texttt{DESRecord.\allowbreak{}add\allowbreak{}Event\allowbreak{}Ancestors} (líneas 166/174), que \texttt{Generic\allowbreak{}Evaluation} recibe envuelto en \texttt{Invocation\allowbreak{}Target\allowbreak{}Exception} y relanza como \texttt{Unreachable\allowbreak{}Exception}.

Detrás hay otros dos puntos con el mismo defecto, que hoy no se ejecutan porque la construcción ya ha fallado (comprobados leyendo):

\begin{itemize}
\tightlist
\item
  \texttt{Event\allowbreak{}Record.\allowbreak{}find\allowbreak{}Descendants} baja por los nodos de azar hijos con una lista de pendientes y solo comprueba si el hijo \textbf{ya está en la lista}, no si ya se visitó: con \texttt{A\ →\ B\ →\ A} saca A y mete B, saca B y vuelve a meter A, y así sin fin. La aplicación se quedaría colgada en el hilo de la simulación, sin excepción.
\item
  Los bucles sin tope de \texttt{Chance\allowbreak{}Evaluation} (\hyperref[err:96]{error~96}) lanzarían \texttt{Index\allowbreak{}Out\allowbreak{}Of\allowbreak{}Bounds\allowbreak{}Exception} porque ninguno de los dos nodos está listo antes que el otro.
\end{itemize}

Un ciclo que pase por un nodo de suceso (A → E → B → A) no provoca la recursión, porque \texttt{add\allowbreak{}Event\allowbreak{}Ancestors} se detiene en los sucesos.

El repositorio no tiene ninguna red de sucesos: para reproducirlo basta con construir una red de sucesos con dos nodos de azar enlazados en los dos sentidos (\texttt{Ill\ →\ Mood} y \texttt{Mood\ →\ Ill}) y simularla.
\paragraph{El código}
inference.des/src/main/java/org/openmarkov/inference/DES/DESRecord.java:165-178

\begin{Shaded}
\begin{Highlighting}[]
    \KeywordTok{protected} \DataTypeTok{void} \FunctionTok{addEventAncestors}\OperatorTok{(}\BuiltInTok{Node}\NormalTok{ node}\OperatorTok{)} \OperatorTok{\{}
        \BuiltInTok{List}\OperatorTok{\textless{}}\BuiltInTok{Node}\OperatorTok{\textgreater{}}\NormalTok{ parents }\OperatorTok{=}\NormalTok{ node}\OperatorTok{.}\FunctionTok{getParents}\OperatorTok{();}
        \ControlFlowTok{if} \OperatorTok{(}\NormalTok{parents}\OperatorTok{.}\FunctionTok{isEmpty}\OperatorTok{())} \OperatorTok{\{}
            \ControlFlowTok{return}\OperatorTok{;}
        \OperatorTok{\}}
        \ControlFlowTok{for} \OperatorTok{(}\BuiltInTok{Node}\NormalTok{ parent }\OperatorTok{:}\NormalTok{ parents}\OperatorTok{)} \OperatorTok{\{}
            \ControlFlowTok{if} \OperatorTok{(}\NormalTok{parent}\OperatorTok{.}\FunctionTok{getNodeType}\OperatorTok{()} \OperatorTok{==}\NormalTok{ NodeType}\OperatorTok{.}\FunctionTok{EVENT}\OperatorTok{)} \OperatorTok{\{}
\NormalTok{                eventAncestors}\OperatorTok{.}\FunctionTok{add}\OperatorTok{(}\NormalTok{parent}\OperatorTok{);}
            \OperatorTok{\}} \ControlFlowTok{else} \ControlFlowTok{if} \OperatorTok{(!}\NormalTok{parent}\OperatorTok{.}\FunctionTok{equals}\OperatorTok{(}\NormalTok{node}\OperatorTok{))} \OperatorTok{\{}
                \FunctionTok{addEventAncestors}\OperatorTok{(}\NormalTok{parent}\OperatorTok{);}
            \OperatorTok{\}}
        \OperatorTok{\}}
        
    \OperatorTok{\}}
\end{Highlighting}
\end{Shaded}

inference.des/src/main/java/org/openmarkov/inference/DES/EventRecord.java:106-117

\begin{Shaded}
\begin{Highlighting}[]
        \ControlFlowTok{while} \OperatorTok{(!}\NormalTok{pending}\OperatorTok{.}\FunctionTok{isEmpty}\OperatorTok{())} \OperatorTok{\{}
            \BuiltInTok{Node}\NormalTok{ chanceNode }\OperatorTok{=}\NormalTok{ pending}\OperatorTok{.}\FunctionTok{remove}\OperatorTok{(}\DecValTok{0}\OperatorTok{);}
\NormalTok{            chanceDescendants}\OperatorTok{.}\FunctionTok{add}\OperatorTok{((}\NormalTok{ChanceRecord}\OperatorTok{)}\NormalTok{ desInference}\OperatorTok{.}\FunctionTok{getChanceEvaluation}\OperatorTok{().}\FunctionTok{getDESRecord}\OperatorTok{(}\NormalTok{chanceNode}\OperatorTok{));}
            \ControlFlowTok{for} \OperatorTok{(}\BuiltInTok{Node}\NormalTok{ child }\OperatorTok{:}\NormalTok{ chanceNode}\OperatorTok{.}\FunctionTok{getChildren}\OperatorTok{())} \OperatorTok{\{}
                \CommentTok{//08/03/2023 {-} self{-}loop}
                \ControlFlowTok{if} \OperatorTok{(}\NormalTok{child}\OperatorTok{.}\FunctionTok{equals}\OperatorTok{(}\NormalTok{chanceNode}\OperatorTok{))} \ControlFlowTok{continue}\OperatorTok{;}
                \ControlFlowTok{switch} \OperatorTok{(}\NormalTok{child}\OperatorTok{.}\FunctionTok{getNodeType}\OperatorTok{())} \OperatorTok{\{}
                    \ControlFlowTok{case}\NormalTok{ CHANCE}\OperatorTok{:}
                        \ControlFlowTok{if} \OperatorTok{(!}\NormalTok{pending}\OperatorTok{.}\FunctionTok{contains}\OperatorTok{(}\NormalTok{child}\OperatorTok{))} \OperatorTok{\{}
\NormalTok{                            pending}\OperatorTok{.}\FunctionTok{add}\OperatorTok{(}\NormalTok{child}\OperatorTok{);}
                        \OperatorTok{\}}
                        \ControlFlowTok{break}\OperatorTok{;}
\end{Highlighting}
\end{Shaded}
\paragraph{Cómo arreglarlo}
Primero hay que tomar la decisión de equipo: si esos ciclos tienen sentido en un modelo DES.

Si no lo tienen: una restricción del tipo de red que prohíba ciclos que no pasen por un nodo de suceso. Como toda restricción, debe imponerse en los dos sitios: \texttt{PNConstraint.\allowbreak{}check\allowbreak{}Prob\allowbreak{}Net} (red completa, por ejemplo al abrir un fichero) y \texttt{PNEdit.\allowbreak{}check\allowbreak{}Constraints\allowbreak{}Will\allowbreak{}Be\allowbreak{}Met} (la edición que añadiría el enlace, \texttt{Add\allowbreak{}Link\allowbreak{}Edit}, y las que lo crean de otra forma, como invertir enlaces o pegar). Aun así conviene que la simulación rechace la red con un mensaje claro.

Si lo tienen: hay que definir el orden de evaluación y hacer que los tres recorridos terminen:

\begin{itemize}
\tightlist
\item
  \texttt{DESRecord.\allowbreak{}add\allowbreak{}Event\allowbreak{}Ancestors}: llevar un conjunto de nodos visitados.
\item
  \texttt{Event\allowbreak{}Record.\allowbreak{}find\allowbreak{}Descendants}: llevar un conjunto de visitados, no solo mirar la lista de pendientes.
\item
  \texttt{Chance\allowbreak{}Evaluation.\allowbreak{}add\allowbreak{}Values\allowbreak{}Orphan\allowbreak{}Nodes} y \texttt{Chance\allowbreak{}Evaluation.\allowbreak{}update}: tope y mensaje que nombre los nodos que se esperan (ver \hyperref[err:96]{error~96}).
\end{itemize}

Prueba de regresión: una red DES con \texttt{Ill\ ↔\ Mood} produce el mensaje decidido (rechazo del enlace o de la red), o simula y termina, según la decisión; en ningún caso \texttt{Stack\allowbreak{}Overflow\allowbreak{}Error} ni cuelgue (con tiempo límite en la prueba).

\setcounter{subsection}{97}
\subsection[Una tabla o valor exacto con un suceso padre hace caer la simulación]{En una red de sucesos, un nodo cuya tabla o valor exacto incluye un suceso padre hace caer la simulación en cuanto ocurre ese suceso}\label{err:98}
\begin{ficha}
\fichakey{Estado}Presente
\fichakey{Módulo}core
\fichakey{Paquete}org.openmarkov.core.model.network.potential
\fichakey{Ficheros}\path{core/src/main/java/org/openmarkov/core/model/network/Finding.java:82-91} \newline \path{core/src/main/java/org/openmarkov/core/model/network/potential/TablePotential.java:238-268} \newline \path{core/src/main/java/org/openmarkov/core/model/network/potential/TablePotential.java:721-738} \newline \path{core/src/main/java/org/openmarkov/core/model/network/potential/ExactDistrPotential.java:158-160} \newline \path{core/src/main/java/org/openmarkov/core/model/network/potential/ExactDistrPotential.java:219-223} \newline \path{core/src/main/java/org/openmarkov/core/action/base/linkEdits/AddLinkEdit.java:224-229} \newline \path{core/src/main/java/org/openmarkov/core/model/network/potential/plugin/PotentialUtils.java:149-161} \newline \path{gui/src/main/java/org/openmarkov/gui/action/ChangeNodeTypeEdit.java:48-63} \newline \path{inference.des/src/main/java/org/openmarkov/inference/DES/ChanceEvaluation.java:136-138} \newline \path{inference.des/src/main/java/org/openmarkov/inference/DES/UtilityEvaluation.java:90-92} \newline \path{inference.des/src/main/java/org/openmarkov/inference/DES/EventEvaluation.java:96-101}
\fichakey{Comprobado}Leyendo el código y ejecutándolo
\fichakey{Esfuerzo}Medio (días)
\fichakey{Decisión de equipo}\textcolor{decisionrojo}{\textbf{Sí.}} Cómo deben tratar las variables de suceso los potenciales que no las entienden (tabla, valor exacto): prohibirlas y convertir el potencial al enlazar, ignorarlas al leer la casilla, o dar al hallazgo de un suceso un estado válido
\fichakey{Relacionados}\hyperref[err:96]{error~96}, \hyperref[err:99]{error~99}
\end{ficha}
\paragraph{Qué pasa}
En una red DES (simulación de eventos discretos), un nodo cuyo potencial es una tabla (\texttt{Table\allowbreak{}Potential}) o un valor exacto (\texttt{Exact\allowbreak{}Distr\allowbreak{}Potential}) y que tiene un suceso entre sus variables hace caer la simulación con \texttt{Array\allowbreak{}Index\allowbreak{}Out\allowbreak{}Of\allowbreak{}Bounds\allowbreak{}Exception} en cuanto ocurre ese suceso. En los tres caminos medidos la excepción no la captura el hilo de la simulación y termina, por lectura, en el diálogo de error inesperado de \texttt{OMException\allowbreak{}Handler}.

Nada en la interfaz dice que un nodo con un suceso padre necesita un potencial de sucesos; se descubre al simular. La ventana de edición de potenciales ya evita crear estos potenciales: para un nodo cuya lista de variables incluye un suceso no ofrece ni «Table» ni «Exact» (medido con \texttt{Potential\allowbreak{}Utils.\allowbreak{}get\allowbreak{}Filtered\allowbreak{}Potential\allowbreak{}Classes}; \texttt{Exact\allowbreak{}Distr\allowbreak{}Potential.\allowbreak{}validate} rechaza explícitamente las variables de suceso en una \texttt{DESNet}). Pero no retira el potencial que ya tiene el nodo, y ni el enlace (\texttt{Add\allowbreak{}Link\allowbreak{}Edit}), ni el cambio de tipo (\texttt{Change\allowbreak{}Node\allowbreak{}Type\allowbreak{}Edit}), ni el lector de ficheros comprueban esa misma regla.

Cuando ocurre un suceso, los nodos hijos de ese suceso se recalculan con una configuración que incluye el suceso ocurrido: \texttt{Chance\allowbreak{}Evaluation.\allowbreak{}compute\allowbreak{}Chance\allowbreak{}Value}, \texttt{Utility\allowbreak{}Evaluation.\allowbreak{}update} y \texttt{Event\allowbreak{}Evaluation.\allowbreak{}update\allowbreak{}Child\allowbreak{}Event} llaman a \texttt{configuration.\allowbreak{}add\allowbreak{}Finding(\allowbreak{}variable\allowbreak{}Del\allowbreak{}Suceso,\allowbreak{}\ instante)}. Para una variable de suceso, \texttt{Finding(\allowbreak{}Variable,\allowbreak{}\ double)} guarda como índice de estado \texttt{Integer.\allowbreak{}MAX\_\allowbreak{}VALUE}.

Los potenciales que saben de sucesos (\texttt{Transition\allowbreak{}Table\allowbreak{}Potential}, \texttt{Distribution\allowbreak{}Table\allowbreak{}Potential}, \texttt{Univariate\allowbreak{}Distr\allowbreak{}Potential}, los árboles de sucesos) traducen ese hallazgo a su variable interna \texttt{Events}. Pero una \texttt{Table\allowbreak{}Potential} o un \texttt{Exact\allowbreak{}Distr\allowbreak{}Potential} cuya lista de variables incluya el suceso calculan la casilla como \texttt{desplazamiento\ ×\ índice}, y el entero desborda: la lectura cae fuera de la tabla.

El repositorio no tiene ninguna red de sucesos. Las medidas se hicieron con una red de sucesos de prueba con un suceso \texttt{Onset} que ocurre en el año 3, por tres caminos:

\begin{itemize}
\tightlist
\item
  \textbf{Nodo de azar enlazado desde un suceso sin cambiar su potencial.} Se crea un nodo de azar \texttt{Pain} (su potencial por omisión es una tabla) y se enlaza \texttt{Onset\ →\ Pain} con \texttt{Add\allowbreak{}Link\allowbreak{}Edit}, como hace el editor. \texttt{Table\allowbreak{}Potential.\allowbreak{}add\allowbreak{}Variable} mantiene la tabla, ahora sobre \texttt{{[}Pain,\allowbreak{}\ Onset{]}}. Al simular, cuando ocurre \texttt{Onset}: \texttt{Array\allowbreak{}Index\allowbreak{}Out\allowbreak{}Of\allowbreak{}Bounds\allowbreak{}Exception:\ Index\ -2\ out\ of\ bounds\ for\ length\ 2} en \texttt{Table\allowbreak{}Potential.\allowbreak{}sample\allowbreak{}Conditioned\allowbreak{}Variable}.
\item
  \textbf{Nodo cambiado a utilidad en la ventana del nodo.} Un nodo de azar enlazado desde \texttt{Onset} se cambia a utilidad con \texttt{Change\allowbreak{}Node\allowbreak{}Type\allowbreak{}Edit} (lo que ejecuta el panel de definición del nodo); como su tabla no vale para una utilidad, recibe el potencial por omisión de \texttt{Potential\allowbreak{}Utils.\allowbreak{}generate\allowbreak{}Default\allowbreak{}Potential}, un \texttt{Exact\allowbreak{}Distr\allowbreak{}Potential} sobre la variable y \textbf{todos} sus padres, suceso incluido. Al simular: \texttt{Array\allowbreak{}Index\allowbreak{}Out\allowbreak{}Of\allowbreak{}Bounds\allowbreak{}Exception:\ Index\ 2147483647\ out\ of\ bounds\ for\ length\ 1} en \texttt{Table\allowbreak{}Potential.\allowbreak{}get\allowbreak{}Value}, llamado desde \texttt{Utility\allowbreak{}Evaluation.\allowbreak{}update}.
\item
  \textbf{Fichero.} La misma red con el potencial exacto de la utilidad \texttt{Treatment\allowbreak{}Cost} escrito sobre \texttt{{[}Treatment\allowbreak{}Cost,\allowbreak{}\ Onset,\allowbreak{}\ Therapy{]}}: \texttt{Index\ 2147483647\ out\ of\ bounds\ for\ length\ 2}. La tabla sin \texttt{Onset} contesta lo que debe.
\end{itemize}

Para reproducir los dos primeros basta con construir en el editor una red de sucesos con un suceso y enlazarlo a un nodo de azar nuevo, o a uno que después se cambia a utilidad.
\paragraph{El código}
core/src/main/java/org/openmarkov/core/model/network/Finding.java:82-91

\begin{Shaded}
\begin{Highlighting}[]
    \KeywordTok{public} \FunctionTok{Finding}\OperatorTok{(}\NormalTok{Variable variable}\OperatorTok{,} \DataTypeTok{double}\NormalTok{ numericalValue}\OperatorTok{)} \OperatorTok{\{}
        \CommentTok{// }\AlertTok{TODO}\CommentTok{ Throw exception if numerical values is outside the domain of variable}
        \KeywordTok{this}\OperatorTok{.}\FunctionTok{variable} \OperatorTok{=}\NormalTok{ variable}\OperatorTok{;}
        \KeywordTok{this}\OperatorTok{.}\FunctionTok{numericalValue} \OperatorTok{=}\NormalTok{ numericalValue}\OperatorTok{;}
        \ControlFlowTok{if} \OperatorTok{(}\NormalTok{variable}\OperatorTok{.}\FunctionTok{getVariableType}\OperatorTok{()} \OperatorTok{==}\NormalTok{ VariableType}\OperatorTok{.}\FunctionTok{DISCRETIZED}\OperatorTok{)} \OperatorTok{\{}
\NormalTok{            stateIndex }\OperatorTok{=} \KeywordTok{this}\OperatorTok{.}\FunctionTok{variable}\OperatorTok{.}\FunctionTok{getStateIndex}\OperatorTok{(}\KeywordTok{this}\OperatorTok{.}\FunctionTok{numericalValue}\OperatorTok{);}
        \OperatorTok{\}} \ControlFlowTok{else} \OperatorTok{\{}
            \KeywordTok{this}\OperatorTok{.}\FunctionTok{stateIndex} \OperatorTok{=} \BuiltInTok{Integer}\OperatorTok{.}\FunctionTok{MAX\_VALUE}\OperatorTok{;}
        \OperatorTok{\}}
    \OperatorTok{\}}
\end{Highlighting}
\end{Shaded}

core/src/main/java/org/openmarkov/core/model/network/potential/TablePotential.java:721-729

\begin{Shaded}
\begin{Highlighting}[]
    \KeywordTok{public} \DataTypeTok{double} \FunctionTok{sampleConditionedVariable}\OperatorTok{(}\DataTypeTok{double}\OperatorTok{[]}\NormalTok{ randomNumbers}\OperatorTok{,}\NormalTok{ EvidenceCase parents}\OperatorTok{)} \OperatorTok{\{}
        \DataTypeTok{int}\NormalTok{ index }\OperatorTok{=} \DecValTok{0}\OperatorTok{;}
        \DataTypeTok{int}\NormalTok{ sampleIndex }\OperatorTok{=} \DecValTok{0}\OperatorTok{;}
        \CommentTok{// find index of first position for the given configuration}
        \ControlFlowTok{for} \OperatorTok{(}\DataTypeTok{int}\NormalTok{ i }\OperatorTok{=} \DecValTok{1}\OperatorTok{;}\NormalTok{ i }\OperatorTok{\textless{}}\NormalTok{ variables}\OperatorTok{.}\FunctionTok{size}\OperatorTok{();} \OperatorTok{++}\NormalTok{i}\OperatorTok{)} \OperatorTok{\{}
\NormalTok{            index }\OperatorTok{+=}\NormalTok{ parents}\OperatorTok{.}\FunctionTok{getState}\OperatorTok{(}\NormalTok{variables}\OperatorTok{.}\FunctionTok{get}\OperatorTok{(}\NormalTok{i}\OperatorTok{))} \OperatorTok{*}\NormalTok{ offsets}\OperatorTok{[}\NormalTok{i}\OperatorTok{];}
        \OperatorTok{\}}
        \DataTypeTok{double}\NormalTok{ accumulatedProbability }\OperatorTok{=}\NormalTok{ values}\OperatorTok{[}\NormalTok{index }\OperatorTok{+}\NormalTok{ sampleIndex}\OperatorTok{];}
\end{Highlighting}
\end{Shaded}

core/src/main/java/org/openmarkov/core/model/network/potential/ExactDistrPotential.java:219-223

\begin{Shaded}
\begin{Highlighting}[]
    \KeywordTok{public} \DataTypeTok{static} \DataTypeTok{boolean} \FunctionTok{validate}\OperatorTok{(}\BuiltInTok{Node}\NormalTok{ node}\OperatorTok{,} \BuiltInTok{List}\OperatorTok{\textless{}}\NormalTok{Variable}\OperatorTok{\textgreater{}}\NormalTok{ variables}\OperatorTok{,}\NormalTok{ PotentialRole role}\OperatorTok{)} \OperatorTok{\{}
        \CommentTok{//11/01/2023 }\AlertTok{FIXME}\CommentTok{ Provisional; validate for DESnets}
        \ControlFlowTok{if} \OperatorTok{((}\NormalTok{node}\OperatorTok{.}\FunctionTok{getProbNet}\OperatorTok{().}\FunctionTok{getNetworkType}\OperatorTok{()} \KeywordTok{instanceof}\NormalTok{ DESNetworkType}\OperatorTok{))\{}
            \ControlFlowTok{return} \OperatorTok{!}\NormalTok{variables}\OperatorTok{.}\FunctionTok{stream}\OperatorTok{().}\FunctionTok{anyMatch}\OperatorTok{(}\NormalTok{variable }\OperatorTok{{-}\textgreater{}}\NormalTok{ variable}\OperatorTok{.}\FunctionTok{getVariableType}\OperatorTok{()==}\NormalTok{VariableType}\OperatorTok{.}\FunctionTok{EVENT}\OperatorTok{)} \OperatorTok{\&\&}\NormalTok{ node}\OperatorTok{.}\FunctionTok{getVariable}\OperatorTok{().}\FunctionTok{getVariableType}\OperatorTok{()} \OperatorTok{==}\NormalTok{ VariableType}\OperatorTok{.}\FunctionTok{NUMERIC}\OperatorTok{;}
        \OperatorTok{\}}
\end{Highlighting}
\end{Shaded}
\paragraph{Cómo arreglarlo}
Hace falta la decisión de equipo sobre el tratamiento. Opciones, sin elegir:

\begin{itemize}
\tightlist
\item
  La regla que ya aplica \texttt{validate} ---un potencial que no entiende sucesos no puede tener variables de suceso--- se impone también al enlazar (sustituyendo el potencial por uno de sucesos o por uno válido), al cambiar el tipo de nodo (que \texttt{generate\allowbreak{}Default\allowbreak{}Potential} no incluya los sucesos o elija un potencial de sucesos en una \texttt{DESNet}) y al leer un fichero (error de lectura con el nombre del nodo).
\item
  Que \texttt{Table\allowbreak{}Potential} y \texttt{Exact\allowbreak{}Distr\allowbreak{}Potential} ignoren las variables de suceso al calcular la casilla.
\item
  Que el hallazgo de un suceso tenga un índice de estado válido.
\end{itemize}

Sitios donde tendría que cumplirse la regla elegida:

\begin{itemize}
\tightlist
\item
  \texttt{Add\allowbreak{}Link\allowbreak{}Edit.\allowbreak{}do\allowbreak{}Edit} (rama general, \texttt{old\allowbreak{}Potential.\allowbreak{}add\allowbreak{}Variable}), y las otras ediciones que añaden padres (invertir enlace, pegar).
\item
  \texttt{Change\allowbreak{}Node\allowbreak{}Type\allowbreak{}Edit.\allowbreak{}do\allowbreak{}Multi\allowbreak{}Step\allowbreak{}Edit} y \texttt{Potential\allowbreak{}Utils.\allowbreak{}generate\allowbreak{}Default\allowbreak{}Potential}.
\item
  \texttt{PGMXReader\_\allowbreak{}1\_\allowbreak{}0} al construir potenciales de una \texttt{DESNet}.
\item
  \texttt{Exact\allowbreak{}Distr\allowbreak{}Potential.\allowbreak{}validate} y la validación de \texttt{Table\allowbreak{}Potential} para \texttt{DESNet}.
\item
  En la simulación, \texttt{Chance\allowbreak{}Evaluation.\allowbreak{}compute\allowbreak{}Chance\allowbreak{}Value}, \texttt{Utility\allowbreak{}Evaluation.\allowbreak{}update} y \texttt{Event\allowbreak{}Evaluation.\allowbreak{}update\allowbreak{}Child\allowbreak{}Event}, que construyen el hallazgo del suceso; y, para no leer fuera de la tabla sin decirlo, \texttt{Table\allowbreak{}Potential.\allowbreak{}get\allowbreak{}Value} / \texttt{sample\allowbreak{}Conditioned\allowbreak{}Variable} con un índice de estado fuera de rango.
\end{itemize}

Pruebas de regresión: los tres casos medidos (nodo de azar con tabla por omisión enlazado desde un suceso, nodo cambiado a utilidad, fichero con valor exacto sobre el suceso) deben o bien ser rechazados con un mensaje en el momento de la edición o de la lectura, o bien simular sin excepción, según la decisión.

\setcounter{subsection}{98}
\subsection[Añadir o quitar un enlace deja una relación uniforme y la simulación falla]{En una red de sucesos, añadir o quitar un enlace cambia la relación del nodo hijo por una uniforme que la simulación no admite, y simular termina en un error inesperado que no dice qué nodo es}\label{err:99}
\begin{ficha}
\fichakey{Estado}Presente en parte. Con un nodo recién creado al que aún no se le ha abierto la relación, enlazarlo ya no le deja una relación uniforme sino ninguna, y antes de simular sale un mensaje normal que nombra el nodo sin potencial (commits \texttt{710de2c} y \texttt{83ff8ab}). Sigue sin corregir para todo nodo que ya tuviera relación: añadir o quitar un enlace, o un bucle, la cambia por \texttt{Uniform\allowbreak{}Potential} perdiendo sus números, y la simulación sigue terminando en la \texttt{Class\allowbreak{}Cast\allowbreak{}Exception} de \texttt{DESRecord} como error inesperado sin decir el nodo.
\fichakey{Módulo}core, inference.des
\fichakey{Paquete}org.openmarkov.core.action.base.linkEdits; org.openmarkov.core.model.network.potential; org.openmarkov.inference.DES
\fichakey{Ficheros}\path{core/src/main/java/org/openmarkov/core/action/base/linkEdits/AddLinkEdit.java:189-234} \newline \path{core/src/main/java/org/openmarkov/core/action/base/linkEdits/RemoveLinkEdit.java:92-116} \newline \path{core/src/main/java/org/openmarkov/core/model/network/potential/Potential.java:700-736} \newline \path{core/src/main/java/org/openmarkov/core/model/network/potential/plugin/PotentialUtils.java:149-156} \newline \path{inference.des/src/main/java/org/openmarkov/inference/DES/DESRecord.java:107} \newline \path{inference.des/src/main/java/org/openmarkov/inference/DES/GenericEvaluation.java:63-80} \newline \path{inference.des/src/main/java/org/openmarkov/inference/DES/DESInference.java:392-441} \newline \path{gui/src/main/java/org/openmarkov/gui/window/InferenceHandler.java:274-297} \newline \path{gui/src/main/java/org/openmarkov/gui/dialog/common/EmptyPotentialPanel.java:19} \newline \path{gui/src/main/java/org/openmarkov/gui/dialog/common/DistributionTablePotentialPanel.java:66-72} \newline \path{gui/src/main/java/org/openmarkov/gui/dialog/common/TransitionTablePotentialPanel.java:24-25}
\fichakey{Comprobado}Leyendo el código y ejecutándolo
\fichakey{Esfuerzo}Medio (días)
\fichakey{Decisión de equipo}\textcolor{decisionrojo}{\textbf{Sí.}} Qué relación debe quedar en un nodo de una red de sucesos cuando se le añade o se le quita un padre: una que la simulación admita y conserve los números que se puedan (como hace una tabla corriente), o la uniforme de hoy, obligando a editarla antes de simular
\fichakey{Relacionados}\hyperref[err:21]{error~21}, \hyperref[err:98]{error~98}, \hyperref[err:101]{error~101}, \hyperref[err:106]{error~106}
\end{ficha}
\paragraph{Qué pasa}
En una red DES (simulación de eventos discretos), el usuario enlaza en el editor un nodo con una utilidad o con un suceso. Por ejemplo, traza un enlace de un suceso a una utilidad nueva, o de un nodo de azar a la tabla de un suceso que ya tenía sus parámetros escritos. Si después lanza la simulación (botón o menú de coste-efectividad, que en una \texttt{DESNet} llama a \texttt{monte\allowbreak{}Carlo\allowbreak{}Simulation}) sin haber abierto y cambiado la relación de ese nodo, la simulación no empieza. Sale el diálogo de error inesperado con una \texttt{Class\allowbreak{}Cast\allowbreak{}Exception} que no dice qué nodo tiene la relación que no vale. Si el nodo tenía números (los parámetros de un suceso, los valores de una utilidad), se han perdido sin aviso.

Por qué pasa, en dos piezas:

\textbf{1. Las ediciones de enlaces dejan una relación uniforme.} \texttt{Add\allowbreak{}Link\allowbreak{}Edit.\allowbreak{}do\allowbreak{}Edit} (líneas 189-234) actualiza la relación del hijo con \texttt{old\allowbreak{}Potential.\allowbreak{}add\allowbreak{}Variable(\allowbreak{}padre)}. \texttt{Potential.\allowbreak{}add\allowbreak{}Variable} (líneas 700-717), que es la versión de la clase base, devuelve un \texttt{Uniform\allowbreak{}Potential} nuevo con todas las variables y ningún número. Solo algunas clases lo redefinen: \texttt{Table\allowbreak{}Potential} amplía su tabla, pero ninguna de las relaciones de sucesos lo hace. Es lo que les pasa a:

\begin{itemize}
\tightlist
\item
  las utilidades, cuya relación por omisión en una red de sucesos es \texttt{Exact\allowbreak{}Distr\allowbreak{}Potential} (\texttt{Potential\allowbreak{}Utils.\allowbreak{}generate\allowbreak{}Default\allowbreak{}Potential}, líneas 152-160);
\item
  los sucesos, que empiezan con \texttt{Delta\allowbreak{}Potential} y suelen tener \texttt{Distribution\allowbreak{}Table\allowbreak{}Potential}.
\end{itemize}

Lo mismo hacen otras dos ramas:

\begin{itemize}
\tightlist
\item
  el bucle sobre sí mismo de un suceso o de un nodo de azar, en \texttt{Add\allowbreak{}Link\allowbreak{}Edit}, crea a propósito un \texttt{Uniform\allowbreak{}Potential} (líneas 213-221);
\item
  \texttt{Remove\allowbreak{}Link\allowbreak{}Edit.\allowbreak{}do\allowbreak{}Edit}, al quitar un enlace hacia una utilidad (líneas 92-116).
\end{itemize}

Y por último \texttt{Potential.\allowbreak{}remove\allowbreak{}Variable} (líneas 726-736), que usa \texttt{Paste\allowbreak{}Edit} al pegar un nodo sin sus padres (\hyperref[err:21]{error~21}).

\textbf{2. La simulación no comprueba las relaciones.} \texttt{DESInference.\allowbreak{}initialize} (líneas 392-441) construye un registro por nodo, y el constructor de \texttt{DESRecord} (línea 107) hace \texttt{(\allowbreak{}DESSimulable\allowbreak{}Potential)\ record\allowbreak{}Node.\allowbreak{}get\allowbreak{}Potentials(\allowbreak{}).\allowbreak{}get\allowbreak{}First(\allowbreak{})}. \texttt{Uniform\allowbreak{}Potential} no implementa \texttt{DESSimulable\allowbreak{}Potential}, así que el molde falla. \texttt{Generic\allowbreak{}Evaluation} (líneas 63-80) envuelve la excepción en \texttt{Unreachable\allowbreak{}Exception} sin decir el nodo (guarda \texttt{last\allowbreak{}Node}, pero no lo usa). En la aplicación, el hilo que crea \texttt{monte\allowbreak{}Carlo\allowbreak{}Simulation} (líneas 435-459) solo captura \texttt{IOException} y \texttt{Open\allowbreak{}Markov\allowbreak{}Exception}, así que la excepción llega al manejador global y se muestra como error inesperado (esto último, por lectura).

La interfaz lo sabe a medias. Si el usuario abre el diálogo del nodo, ve la relación «Uniform» con un panel vacío (\texttt{Empty\allowbreak{}Potential\allowbreak{}Panel}) y puede elegir otra. Los comentarios de \texttt{Distribution\allowbreak{}Table\allowbreak{}Potential\allowbreak{}Panel} (líneas 66-72, con un \texttt{FIXME}) y de \texttt{Transition\allowbreak{}Table\allowbreak{}Potential\allowbreak{}Panel} (líneas 24-25) dicen que al añadir o quitar un enlace el nodo se queda con una relación uniforme. Pero nada avisa antes de simular.

El repositorio no tiene redes de sucesos. Medido con una red de sucesos de prueba que simula sin problemas (una decisión \texttt{Therapy}, un nodo de azar \texttt{Sex}, un suceso \texttt{Onset} con tabla de distribución sobre \texttt{Therapy} y \texttt{Sex}, un suceso terminal \texttt{Death} con tabla de distribución sobre \texttt{Onset}, y las utilidades \texttt{QALY} y \texttt{Monitoring}), aplicando las mismas ediciones que el editor y lanzando la simulación:

\begin{Verbatim}[breaklines,breakanywhere,fontsize=\small]
utilidad nueva U, sin enlaces:          ExactDistrPotential  -> simula
enlace Onset -> U (utilidad nueva):     UniformPotential [U, Onset]                     -> falla
enlace Sex -> Monitoring:               UniformPotential [Monitoring, Therapy, Sex]     -> falla
suceso nuevo Relapse, enlace Onset -> Relapse: DeltaPotential pasa a UniformPotential   -> falla
enlace Sex -> Death:                    UniformPotential [Death, Onset, Sex] (se pierde el 4,0) -> falla
quitar el enlace Therapy -> Monitoring: UniformPotential [Monitoring]                    -> falla
bucle Onset -> Onset:                   UniformPotential [Onset, Therapy, Sex, Onset]   -> falla
nodo de azar nuevo, enlace Onset -> nodo: TablePotential (ampliada), no uniforme
\end{Verbatim}

El último caso no queda uniforme, pero es el de \hyperref[err:98]{error~98}: esa tabla con un suceso hace caer la simulación cuando el suceso ocurre.

En todos los casos que fallan, el error es el mismo:

\begin{Verbatim}[breaklines,breakanywhere,fontsize=\small]
UnreachableException: InvocationTargetException
  causa: ClassCastException: class UniformPotential cannot be cast to class DESSimulablePotential
    at DESRecord.<init>(DESRecord.java:107)
    at UtilityRecord.<init>  /  EventRecord.<init>
    at GenericEvaluation.<init>(GenericEvaluation.java:72)
    at DESInference.initialize(DESInference.java:439 / 441)
\end{Verbatim}
\paragraph{El código}
core/src/main/java/org/openmarkov/core/action/base/linkEdits/AddLinkEdit.java:225-231

\begin{Shaded}
\begin{Highlighting}[]
            \OperatorTok{\}} \ControlFlowTok{else} \OperatorTok{\{}
                \ControlFlowTok{for} \OperatorTok{(}\NormalTok{Potential oldPotential }\OperatorTok{:}\NormalTok{ oldPotentials}\OperatorTok{)} \OperatorTok{\{}
                    \CommentTok{// Update potential}
\NormalTok{                    Potential newPotential }\OperatorTok{=}\NormalTok{ oldPotential}\OperatorTok{.}\FunctionTok{addVariable}\OperatorTok{(}\NormalTok{nodeFrom}\OperatorTok{.}\FunctionTok{getVariable}\OperatorTok{());}
\NormalTok{                    newPotentials}\OperatorTok{.}\FunctionTok{add}\OperatorTok{(}\NormalTok{newPotential}\OperatorTok{);}
                \OperatorTok{\}}
            \OperatorTok{\}}
\end{Highlighting}
\end{Shaded}

core/src/main/java/org/openmarkov/core/model/network/potential/Potential.java:700-706

\begin{Shaded}
\begin{Highlighting}[]
    \KeywordTok{public}\NormalTok{ Potential }\FunctionTok{addVariable}\OperatorTok{(}\NormalTok{Variable variable}\OperatorTok{)} \OperatorTok{\{}
\NormalTok{        Potential newPotential}\OperatorTok{;}
        \ControlFlowTok{if} \OperatorTok{(!}\NormalTok{variables}\OperatorTok{.}\FunctionTok{contains}\OperatorTok{(}\NormalTok{variable}\OperatorTok{))} \OperatorTok{\{}
            \BuiltInTok{List}\OperatorTok{\textless{}}\NormalTok{Variable}\OperatorTok{\textgreater{}}\NormalTok{ newVariables }\OperatorTok{=} \KeywordTok{new} \BuiltInTok{ArrayList}\OperatorTok{\textless{}\textgreater{}(}\NormalTok{variables}\OperatorTok{);}
\NormalTok{            newVariables}\OperatorTok{.}\FunctionTok{add}\OperatorTok{(}\NormalTok{variable}\OperatorTok{);}
\NormalTok{            newPotential }\OperatorTok{=} \KeywordTok{new} \FunctionTok{UniformPotential}\OperatorTok{(}\NormalTok{newVariables}\OperatorTok{,}\NormalTok{ role}\OperatorTok{);}
        \OperatorTok{\}} \ControlFlowTok{else} \OperatorTok{\{}
\end{Highlighting}
\end{Shaded}

inference.des/src/main/java/org/openmarkov/inference/DES/DESRecord.java:107

\begin{Shaded}
\begin{Highlighting}[]
\NormalTok{        recordPotential }\OperatorTok{=} \OperatorTok{(}\NormalTok{DESSimulablePotential}\OperatorTok{)}\NormalTok{ recordNode}\OperatorTok{.}\FunctionTok{getPotentials}\OperatorTok{().}\FunctionTok{getFirst}\OperatorTok{();}
\end{Highlighting}
\end{Shaded}
\paragraph{Cómo arreglarlo}
Dos partes.

\textbf{La simulación debe rechazar la red con un mensaje.} Esta parte no necesita decisión. Antes de construir los registros, \texttt{DESInference.\allowbreak{}initialize} debe comprobar que la relación de cada nodo de azar, suceso y utilidad es simulable (\texttt{DESSimulable\allowbreak{}Potential}) y válida para sus padres. Si no lo es, debe avisar nombrando los nodos, igual que ya avisa cuando faltan criterios o hay más de una decisión. \texttt{Generic\allowbreak{}Evaluation} no debería convertir en \texttt{Unreachable\allowbreak{}Exception} un fallo que depende de la red del usuario.

\textbf{Qué relación dejan las ediciones.} Esta parte necesita la decisión de equipo. Opciones:

\begin{itemize}
\tightlist
\item
  Que las relaciones de sucesos (\texttt{Exact\allowbreak{}Distr\allowbreak{}Potential}, \texttt{Distribution\allowbreak{}Table\allowbreak{}Potential}, \texttt{Delta\allowbreak{}Potential}\ldots) redefinan \texttt{add\allowbreak{}Variable} y \texttt{remove\allowbreak{}Variable} para ampliar o reducir su tabla conservando los números, como hace \texttt{Table\allowbreak{}Potential}.
\item
  Que en una \texttt{DESNet} la edición ponga la relación por omisión válida para los padres nuevos (la que ofrezca primero \texttt{Potential\allowbreak{}Utils.\allowbreak{}get\allowbreak{}Filtered\allowbreak{}Potential\allowbreak{}Classes}).
\item
  Dejar la uniforme y confiar en el aviso de la simulación.
\end{itemize}

Sitios que producen hoy la relación uniforme y deben seguir la regla que se elija:

\begin{itemize}
\tightlist
\item
  \texttt{Add\allowbreak{}Link\allowbreak{}Edit.\allowbreak{}do\allowbreak{}Edit}, en la rama general y en la del bucle sobre sí mismo (líneas 213-221).
\item
  \texttt{Remove\allowbreak{}Link\allowbreak{}Edit.\allowbreak{}do\allowbreak{}Edit}, en las ramas de utilidad (líneas 97-116) y en la general.
\item
  \texttt{Potential.\allowbreak{}add\allowbreak{}Variable} y \texttt{Potential.\allowbreak{}remove\allowbreak{}Variable}.
\item
  \texttt{Paste\allowbreak{}Edit}, al quitar los padres que no se copian (\hyperref[err:21]{error~21}).
\item
  Las ediciones que cambian los padres de otra forma (invertir un enlace).
\end{itemize}

Cuando \texttt{Sum\allowbreak{}Potential} sustituye a la relación de una utilidad cuyos padres son todos numéricos, en \texttt{Add\allowbreak{}Link\allowbreak{}Edit} (líneas 195-205) y en \texttt{Remove\allowbreak{}Link\allowbreak{}Edit} (líneas 99-106), tampoco es una relación simulable, por lectura.

Pruebas de regresión:

\begin{itemize}
\tightlist
\item
  Con la red de prueba descrita, cada una de las ediciones medidas debe dejar una relación que simula o, según lo decidido, hacer que la simulación se niegue con un mensaje que nombre el nodo; nunca \texttt{Class\allowbreak{}Cast\allowbreak{}Exception}.
\item
  Añadir un padre a un suceso con tabla de distribución conserva sus parámetros, si se elige esa opción.
\end{itemize}

\setcounter{subsection}{99}
\subsection[Una tabla de vida «con tasas» se simula como si fueran probabilidades]{Una tabla de vida marcada «con tasas» se simula como si fueran probabilidades}\label{err:100}
\begin{ficha}
\fichakey{Estado}Presente
\fichakey{Módulo}core
\fichakey{Paquete}org.openmarkov.core.model.network.potential
\fichakey{Ficheros}\path{core/src/main/java/org/openmarkov/core/model/network/potential/PiecewiseExponentialPotential.java:56-59,238-277,369-375} \newline \path{gui/src/main/java/org/openmarkov/gui/dialog/common/PiecewiseExponentialPanel.java:84-86, 114} \newline \path{io/src/main/java/org/openmarkov/io/probmodel/reader/PGMXReader_1_0.java:535-536} \newline \path{io/src/main/java/org/openmarkov/io/probmodel/writer/PGMXWriter_1_0.java:432-434}
\fichakey{Comprobado}Leyendo el código y ejecutándolo
\fichakey{Esfuerzo}Pequeño (horas)
\fichakey{Decisión de equipo}\textcolor{decisionrojo}{\textbf{Sí.}} Si se implementa la variante con tasas (y se enciende la casilla) o se retira el campo, que hoy promete algo que no hace
\fichakey{Relacionados}\hyperref[err:101]{error~101}
\end{ficha}
\paragraph{Qué pasa}
Un fichero con \texttt{\textless{}Rates\textgreater{}true\textless{}/\allowbreak{}Rates\textgreater{}} en una tabla de vida se abre, se enseña con la casilla «Use rates» marcada, se guarda igual y se simula tratando las tasas como probabilidades, sin aviso. Ninguna red del repositorio lleva \texttt{\textless{}Rates\textgreater{}}, y desde la interfaz el campo siempre queda en \texttt{false}, así que hoy solo se llega con un fichero escrito a mano o por otra herramienta.

\texttt{Piecewise\allowbreak{}Exponential\allowbreak{}Potential} (módulo \texttt{core}) es la relación «exponencial por tramos» o tabla de vida de la simulación de eventos discretos (DES): una tabla de pares (inicio del tramo, valor) con la que se sortea cuándo ocurre un suceso. Tiene un campo \texttt{use\allowbreak{}Rates} cuyo comentario (líneas 56-59) dice que es cierto «si la distribución está definida por tasas en vez de por probabilidades por unidad de tiempo (como una probabilidad anual o mensual)». Son dos cosas distintas: una tasa λ y una probabilidad por unidad de tiempo p se relacionan por p = 1 − e\^{}(−λ).

El sorteo (\texttt{sample\allowbreak{}Conditioned\allowbreak{}Variable}, líneas 238-277) y el precálculo de la tabla (\texttt{set\allowbreak{}Piecewise\allowbreak{}Table}) usan siempre la fórmula de probabilidades, \texttt{Math.\allowbreak{}log(\allowbreak{}1\ -\ p\allowbreak{}I)} y \texttt{Math.\allowbreak{}pow(\allowbreak{}1\ -\ p,\allowbreak{}\ …)}, y no leen \texttt{use\allowbreak{}Rates} en ningún sitio.

El campo solo se copia, se escribe en el fichero como \texttt{\textless{}Rates\textgreater{}} (\texttt{PGMXWriter\_\allowbreak{}1\_\allowbreak{}0}, líneas 432-434), se lee de vuelta (\texttt{PGMXReader\_\allowbreak{}1\_\allowbreak{}0}, líneas 535-536) y se enseña en el panel con una casilla «Use rates» que está apagada (\texttt{Piecewise\allowbreak{}Exponential\allowbreak{}Panel}, líneas 86-88), así que el usuario no puede cambiarlo.

Medido con un pequeño programa de prueba: con la misma tabla y los mismos números aleatorios, el sorteo da \texttt{7.\allowbreak{}094838467442498} con \texttt{use\allowbreak{}Rates=false} y con \texttt{use\allowbreak{}Rates=true}.
\paragraph{El código}
core/src/main/java/org/openmarkov/core/model/network/potential/PiecewiseExponentialPotential.java:275

\begin{Shaded}
\begin{Highlighting}[]
\DataTypeTok{double}\NormalTok{ sample }\OperatorTok{=}\NormalTok{ tI }\OperatorTok{+} \OperatorTok{(}\BuiltInTok{Math}\OperatorTok{.}\FunctionTok{log}\OperatorTok{(}\NormalTok{u}\OperatorTok{)} \OperatorTok{{-}} \BuiltInTok{Math}\OperatorTok{.}\FunctionTok{log}\OperatorTok{(}\NormalTok{prodAi}\OperatorTok{))/}\BuiltInTok{Math}\OperatorTok{.}\FunctionTok{log}\OperatorTok{(}\DecValTok{1} \OperatorTok{{-}}\NormalTok{ pI}\OperatorTok{);}
\end{Highlighting}
\end{Shaded}

gui/src/main/java/org/openmarkov/gui/dialog/common/PiecewiseExponentialPanel.java:86-88

\begin{Shaded}
\begin{Highlighting}[]
\NormalTok{ratesCheck }\OperatorTok{=} \KeywordTok{new} \BuiltInTok{JCheckBox}\OperatorTok{(}\StringTok{"Use rates"}\OperatorTok{);}
\NormalTok{ratesCheck}\OperatorTok{.}\FunctionTok{setSelected}\OperatorTok{(}\NormalTok{piecewiseExponentialPotential}\OperatorTok{.}\FunctionTok{isUseRates}\OperatorTok{());}
\NormalTok{ratesCheck}\OperatorTok{.}\FunctionTok{setEnabled}\OperatorTok{(}\KeywordTok{false}\OperatorTok{);}
\end{Highlighting}
\end{Shaded}
\paragraph{Cómo arreglarlo}
Exige decidir entre:

\begin{itemize}
\tightlist
\item
  Implementar las tasas: convertir cada valor con p = 1 − e\^{}(−λ) cuando \texttt{use\allowbreak{}Rates} es cierto, en \texttt{set\allowbreak{}Piecewise\allowbreak{}Table} (productos acumulados) y en \texttt{sample\allowbreak{}Conditioned\allowbreak{}Variable}, y encender la casilla. El campo debe tenerse en cuenta también en \texttt{equals}/copia si los hay, y la conversión debe hacerse en un solo sitio para que las dos cuentas no diverjan.
\item
  Retirar la casilla y dejar de aceptar \texttt{true} en el lector (o avisar al leerlo), para que el fichero no diga algo que el programa no hace.
\end{itemize}

Prueba: con una tabla de un tramo de tasa λ, comprobar que la media de muchos sorteos con \texttt{use\allowbreak{}Rates=true} se acerca a 1/λ, y con \texttt{false} y p se acerca a −1/ln(1 − p).

\setcounter{subsection}{100}
\subsection[La relación «External» se ofrece, elegirla falla y no trae valores de fuera]{La relación «External» se ofrece en el editor, elegirla lanza una excepción, y no trae ningún valor de fuera}\label{err:101}
\begin{ficha}
\fichakey{Estado}Presente
\fichakey{Módulo}core
\fichakey{Paquete}org.openmarkov.core.model.network.potential
\fichakey{Ficheros}\path{core/src/main/java/org/openmarkov/core/model/network/potential/ExternalPotential.java:17-24,67-70,137-144} \newline \path{inference.des/src/main/java/org/openmarkov/inference/DES/ChanceEvaluation.java:63-89,123-161} \newline \path{core/src/main/java/org/openmarkov/core/model/network/potential/plugin/PotentialUtils.java:117-134} \newline \path{gui/src/main/java/org/openmarkov/gui/dialog/node/PotentialEditPanel.java:228-253, 488-522, 584-587} \newline \path{gui/src/main/java/org/openmarkov/gui/dialog/common/PotentialPanelManager.java:95-117} \newline \path{io/src/main/java/org/openmarkov/io/probmodel/reader/PGMXReader_1_0.java:117,515-519} \newline \path{integrationTests/src/test/java/org/openmarkov/integrationTests/gui/AllPotentialsHaveAPanelTest.java:13-16}
\fichakey{Comprobado}Leyendo el código y ejecutándolo
\fichakey{Esfuerzo}Medio (días)
\fichakey{Decisión de equipo}\textcolor{decisionrojo}{\textbf{Sí.}} Si la relación «External» se completa (ligada al fichero de datos de Monte Carlo, con su panel) o se deja de ofrecer como tipo de relación
\fichakey{Relacionados}\hyperref[err:23]{error~23}, \hyperref[err:99]{error~99}, \hyperref[err:100]{error~100}, \hyperref[nota:71]{nota~71}, \hyperref[nota:90]{nota~90}, \hyperref[nota:96]{nota~96}, \hyperref[nota:98]{nota~98}, \hyperref[nota:162]{nota~162}
\end{ficha}
\paragraph{Qué pasa}
En una red de simulación de eventos discretos (DES), el selector de clase de relación del editor ofrece «External» para cualquier nodo de azar. Si el usuario la elige, sale un error no previsto (\texttt{Null\allowbreak{}Pointer\allowbreak{}Exception}) porque no tiene panel, y el nodo no puede editarse. Y un nodo con esta relación (por ejemplo, leído de un fichero) no trae ningún valor de fuera: cuando el fichero de datos de la simulación no le da valor, la simulación sigue con un valor sin sentido en vez de avisar.

\texttt{External\allowbreak{}Potential} (módulo \texttt{core}) se declara como tipo de relación con \texttt{@Potential\allowbreak{}Type(\allowbreak{}names\ =\ "External")} y su comentario dice que representa «un valor que se toma de una fuente externa, como un fichero». La clase no guarda ninguna fuente (ni fichero, ni columna, ni nombre) y su único método de simulación devuelve siempre \texttt{Double.\allowbreak{}MAX\_\allowbreak{}VALUE} (líneas 137-144), que el comentario de la interfaz \texttt{DESSimulable\allowbreak{}Potential} usa como «no hay muestra». Medido: \texttt{sample\allowbreak{}Conditioned\allowbreak{}Variable} devuelve \texttt{1.\allowbreak{}7976931348623157E308}. El lector de ProbModelXML sí la reconoce (\texttt{PGMXReader\_\allowbreak{}1\_\allowbreak{}0}, líneas 117 y 515-519), así que un fichero puede traerla.

La simulación DES sí tiene un fichero de datos de entrada, elegido en las opciones de Monte Carlo (\texttt{Monte\allowbreak{}Carlo\allowbreak{}Options.\allowbreak{}get\allowbreak{}Input\allowbreak{}File\allowbreak{}Path}, \texttt{DESInference}, línea 404). Al empezar cada individuo, \texttt{Chance\allowbreak{}Evaluation.\allowbreak{}add\allowbreak{}Values\allowbreak{}Orphan\allowbreak{}Nodes} (\texttt{inference.\allowbreak{}des}, líneas 63-89) toma del fichero el valor de cada nodo de azar \textbf{sin sucesos antecesores} cuyo nombre sea una columna del fichero, y solo si no lo encuentra muestrea su relación.

Esa lectura no depende de la clase de relación: vale igual para una tabla que para «External». Por tanto, «External» funciona de hecho como una etiqueta que no comprueba nada. Si no hay fichero, si el nombre del nodo no es una columna, o si el nodo tiene un suceso antecesor (entonces se recalcula con \texttt{compute\allowbreak{}Chance\allowbreak{}Value} sin mirar el fichero), la simulación usa \texttt{Double.\allowbreak{}MAX\_\allowbreak{}VALUE} como valor del nodo (para una variable de estados, \texttt{(\allowbreak{}int)\ Double.\allowbreak{}MAX\_\allowbreak{}VALUE} como índice de estado) sin avisar. Esto último es lectura del código; no se ha lanzado una simulación.

El selector de clase de relación del editor (\texttt{Potential\allowbreak{}Edit\allowbreak{}Panel.\allowbreak{}get\allowbreak{}Potential\allowbreak{}Type\allowbreak{}JCombobox}, líneas 227-241) ofrece todo tipo cuyo \texttt{validate} acepte el nodo (\texttt{Potential\allowbreak{}Utils.\allowbreak{}get\allowbreak{}Filtered\allowbreak{}Potential\allowbreak{}Classes}, líneas 121-138), sin mirar si tiene panel, y \texttt{External\allowbreak{}Potential.\allowbreak{}validate} (líneas 67-70) acepta cualquier nodo de azar de una red DES. Medido con un pequeño programa de prueba en una red \texttt{DESNetwork\allowbreak{}Type} con un nodo de azar: los tipos ofrecidos son \texttt{{[}External\allowbreak{}Potential,\allowbreak{}\ Delta\allowbreak{}Potential,\allowbreak{}\ Weibull\allowbreak{}Hazard\allowbreak{}Potential,\allowbreak{}\ Table\allowbreak{}Potential{]}}.

Si el usuario la elige, \texttt{potential\allowbreak{}Type\allowbreak{}Changed} (líneas 476-507) instala la relación nueva en el nodo y recarga el panel; \texttt{Potential\allowbreak{}Panel\allowbreak{}Manager.\allowbreak{}create\allowbreak{}Potential\allowbreak{}Panel} (líneas 95-114) busca el panel de esa clase, obtiene \texttt{null} y lanza \texttt{Null\allowbreak{}Pointer\allowbreak{}Exception}. Medido llamando a ese método con un nodo que tiene \texttt{External\allowbreak{}Potential}:

\begin{Verbatim}[breaklines,breakanywhere,fontsize=\small]
panel THREW java.lang.NullPointerException: Cannot invoke "java.lang.Class.getConstructor(java.lang.Class[])" because "potentialPanelClass" is null
\end{Verbatim}

La excepción no la caza nadie y llega al manejador global como error no previsto.

La prueba \texttt{All\allowbreak{}Potentials\allowbreak{}Have\allowbreak{}APanel\allowbreak{}Test} está desactivada precisamente porque este tipo, \texttt{Increment\allowbreak{}Potential} y \texttt{Table\allowbreak{}With\allowbreak{}Events} no tienen panel (líneas 13-16).
\paragraph{El código}
core/src/main/java/org/openmarkov/core/model/network/potential/ExternalPotential.java:67-70,141-144

\begin{Shaded}
\begin{Highlighting}[]
\KeywordTok{public} \DataTypeTok{static} \DataTypeTok{boolean} \FunctionTok{validate}\OperatorTok{(}\BuiltInTok{Node}\NormalTok{ node}\OperatorTok{,} \BuiltInTok{List}\OperatorTok{\textless{}}\NormalTok{Variable}\OperatorTok{\textgreater{}}\NormalTok{ variables}\OperatorTok{,}\NormalTok{ PotentialRole role}\OperatorTok{)} \OperatorTok{\{}
    \CommentTok{//28/08/2023 }\AlertTok{FIXME}\CommentTok{ Provisional; only for DESnets}
    \ControlFlowTok{return} \OperatorTok{(}\NormalTok{node}\OperatorTok{.}\FunctionTok{getProbNet}\OperatorTok{().}\FunctionTok{getNetworkType}\OperatorTok{()} \KeywordTok{instanceof}\NormalTok{ DESNetworkType}\OperatorTok{)} \OperatorTok{\&\&} \OperatorTok{(}\NormalTok{node}\OperatorTok{.}\FunctionTok{getNodeType}\OperatorTok{()} \OperatorTok{==}\NormalTok{ NodeType}\OperatorTok{.}\FunctionTok{CHANCE}\OperatorTok{);}
\OperatorTok{\}}
\KeywordTok{...}
\AttributeTok{@Override}
\KeywordTok{public} \DataTypeTok{double} \FunctionTok{sampleConditionedVariable}\OperatorTok{(}\DataTypeTok{double}\OperatorTok{[]}\NormalTok{ randomNumbers}\OperatorTok{,}\NormalTok{ EvidenceCase parents}\OperatorTok{)} \OperatorTok{\{}
    \ControlFlowTok{return} \BuiltInTok{Double}\OperatorTok{.}\FunctionTok{MAX\_VALUE}\OperatorTok{;}
\OperatorTok{\}}
\end{Highlighting}
\end{Shaded}
\paragraph{Cómo arreglarlo}
Necesita la decisión de equipo: completar la relación o retirarla como tipo ofrecido.

Completarla supone decidir su relación con el fichero de datos de Monte Carlo (por ejemplo, que la simulación exija una columna para todo nodo «External» y avise antes de empezar si falta, o si el nodo tiene sucesos antecesores), qué guarda el fichero \texttt{.\allowbreak{}pgmx} y qué panel la edita.

Mientras no se decida, lo mínimo para que el usuario no tropiece es que \texttt{validate} devuelva \texttt{false} (o que el selector no ofrezca tipos sin panel) y que \texttt{create\allowbreak{}Potential\allowbreak{}Panel} lance un error con mensaje en vez de \texttt{Null\allowbreak{}Pointer\allowbreak{}Exception}.

La regla «no ofrecer lo que no se puede editar» afecta a los tres tipos sin panel y debe cumplirse en \texttt{Potential\allowbreak{}Utils.\allowbreak{}get\allowbreak{}Filtered\allowbreak{}Potential\allowbreak{}Classes} o en \texttt{Potential\allowbreak{}Edit\allowbreak{}Panel.\allowbreak{}get\allowbreak{}Potential\allowbreak{}Type\allowbreak{}JCombobox}. El lector de ficheros seguirá pudiendo traerlos, así que el diálogo también debe abrirse sin error para un nodo que ya la tenga.

Reactivar \texttt{All\allowbreak{}Potentials\allowbreak{}Have\allowbreak{}APanel\allowbreak{}Test} es la prueba natural.

\setcounter{subsection}{101}
\subsection[Resultados de la simulación: la media descontada va al criterio equivocado]{La tabla de resultados de la simulación de sucesos pone la columna de media descontada al criterio que no tiene descuento, y la media descontada del coste no se ve}\label{err:102}
\begin{ficha}
\fichakey{Estado}Presente
\fichakey{Módulo}inference.des
\fichakey{Paquete}org.openmarkov.inference.DES
\fichakey{Ficheros}\path{inference.des/src/main/java/org/openmarkov/inference/DES/CriteriaValues.java:20,36-42,195-209} \newline \path{inference.des/src/main/java/org/openmarkov/inference/DES/SimulationSummaryResults.java:184-212} \newline \path{inference.des/src/main/java/org/openmarkov/inference/DES/DESInference.java:373-386} \newline \path{inference.des/src/main/java/org/openmarkov/inference/DES/DESResultsWindow.java:135-218}
\fichakey{Comprobado}Leyendo el código y ejecutándolo
\fichakey{Esfuerzo}Pequeño (horas)
\fichakey{Decisión de equipo}No
\fichakey{Relacionados}\hyperref[err:103]{error~103}
\end{ficha}
\paragraph{Qué pasa}
En la tabla de resultados de una simulación de sucesos discretos (DES ---simulación de eventos discretos---, redes de tipo «DESNet») con dos criterios, uno con descuento y otro sin él, las columnas pueden asignarse al criterio equivocado. Con criterios llamados «Cost» (con descuento) y «Effectiveness» (sin él), la efectividad sale con dos columnas y el coste con una: la media descontada del coste, que es el resultado que se busca al poner descuento, no aparece en la tabla, y bajo «Disc. Mean» se enseña un valor del criterio que no tiene descuento. Con los nombres por omisión de una red nueva («Cost (€)», «Effectiveness (QALY)») no se nota.

Camino del usuario: red de tipo DESNet con dos criterios de decisión, uno con descuento y otro sin él, cuyos nombres caigan en orden distinto en el \texttt{Hash\allowbreak{}Map} (por ejemplo «Cost» y «Effectiveness»); botón de coste-efectividad de la barra de herramientas (en una DESNet lanza la simulación, \texttt{Main\allowbreak{}Panel\allowbreak{}Listener\allowbreak{}Assistant}, líneas 308-311) → aceptar las opciones → tabla de resultados. En el repositorio no hay redes DESNet de ejemplo, así que no se ha ejecutado una simulación completa.

Al terminar la simulación, \texttt{DESInference.\allowbreak{}show\allowbreak{}Results\allowbreak{}Window} llama a \texttt{DESResults\allowbreak{}Window.\allowbreak{}present\allowbreak{}Results} pasándole dos cosas por separado: el texto \texttt{simulation\allowbreak{}Summary\allowbreak{}Results.\allowbreak{}to\allowbreak{}String(\allowbreak{})} y un vector \texttt{using\allowbreak{}Discount} que dice, criterio por criterio, si tiene descuento. La ventana no recibe objetos: trocea ese texto por puntos y comas (\texttt{generate\allowbreak{}Table}) y recorre a la vez los trozos y \texttt{using\allowbreak{}Discount}, poniendo dos columnas («Mean» y «Disc. Mean») a los criterios con descuento y una a los demás.

El problema es el orden. \texttt{using\allowbreak{}Discount} se construye recorriendo la lista de criterios de la red (\texttt{criteria.\allowbreak{}get(\allowbreak{}criterion\allowbreak{}Num)}). El texto lo escribe \texttt{Criteria\allowbreak{}Values.\allowbreak{}to\allowbreak{}String(\allowbreak{})} recorriendo \texttt{criteria\allowbreak{}Values\allowbreak{}Hash\allowbreak{}Map.\allowbreak{}key\allowbreak{}Set(\allowbreak{})}, un \texttt{Hash\allowbreak{}Map\textless{}Equal\allowbreak{}Criterion,\allowbreak{}\ Simulation\allowbreak{}Results\textgreater{}} cuya clave se dispersa por el nombre del criterio: el orden depende de los nombres. Cuando los dos órdenes no coinciden, cada columna se asigna al criterio equivocado.

Medido con un pequeño programa de prueba que arma un texto con el formato real (criterios en el orden que da el \texttt{Hash\allowbreak{}Map}) y llama por reflexión a \texttt{DESResults\allowbreak{}Window.\allowbreak{}generate\allowbreak{}Table}, con la red declarando «Cost» (con descuento: media 100, media descontada 90) y «Effectiveness» (sin descuento: media 2):

\begin{Verbatim}[breaklines,breakanywhere,fontsize=\small]
orden de la red [Cost (€), Effectiveness (QALY)] -> orden del texto [Cost (€), Effectiveness (QALY)]
orden de la red [cost, effectiveness]            -> orden del texto [effectiveness, cost]
orden de la red [Cost, Effectiveness]            -> orden del texto [Effectiveness, Cost]
con [Cost, Effectiveness]: cabeceras Effectiveness (dos columnas: Mean, Disc. Mean) | Cost (una columna)
fila de la serie 0, opción «none»: 2 | 2 | 100
\end{Verbatim}
\paragraph{El código}
inference.des/src/main/java/org/openmarkov/inference/DES/CriteriaValues.java:202-207

\begin{Shaded}
\begin{Highlighting}[]
\BuiltInTok{String}\NormalTok{ result }\OperatorTok{=} \StringTok{""}\OperatorTok{;}
\ControlFlowTok{for} \OperatorTok{(}\NormalTok{EqualCriterion key }\OperatorTok{:}\NormalTok{ criteriaValuesHashMap}\OperatorTok{.}\FunctionTok{keySet}\OperatorTok{())} \OperatorTok{\{}
\NormalTok{    result }\OperatorTok{+=}\NormalTok{ key}\OperatorTok{.}\FunctionTok{getCriterionName}\OperatorTok{()} \OperatorTok{+} \StringTok{";"} \OperatorTok{+}\NormalTok{ criteriaValuesHashMap}\OperatorTok{.}\FunctionTok{get}\OperatorTok{(}\NormalTok{key}\OperatorTok{).}\FunctionTok{toString}\OperatorTok{()} \OperatorTok{+} \StringTok{";"}\OperatorTok{;}
\OperatorTok{\}}
\ControlFlowTok{return}\NormalTok{ result}\OperatorTok{;}
\end{Highlighting}
\end{Shaded}

inference.des/src/main/java/org/openmarkov/inference/DES/DESInference.java:373-376

\begin{Shaded}
\begin{Highlighting}[]
\KeywordTok{private} \DataTypeTok{void} \FunctionTok{showResultsWindow}\OperatorTok{(}\DataTypeTok{double}\NormalTok{ simulationTime}\OperatorTok{,} \BuiltInTok{JPanel}\NormalTok{ psaPanel}\OperatorTok{)} \OperatorTok{\{}
    \DataTypeTok{boolean}\OperatorTok{[]}\NormalTok{ usingDiscount }\OperatorTok{=} \KeywordTok{new} \DataTypeTok{boolean}\OperatorTok{[}\NormalTok{criteria}\OperatorTok{.}\FunctionTok{size}\OperatorTok{()];}
    \ControlFlowTok{for} \OperatorTok{(}\DataTypeTok{int}\NormalTok{ criterionNum }\OperatorTok{=} \DecValTok{0}\OperatorTok{;}\NormalTok{ criterionNum }\OperatorTok{\textless{}}\NormalTok{ criteria}\OperatorTok{.}\FunctionTok{size}\OperatorTok{();}\NormalTok{ criterionNum}\OperatorTok{++)}
\NormalTok{        usingDiscount}\OperatorTok{[}\NormalTok{criterionNum}\OperatorTok{]} \OperatorTok{=}\NormalTok{ criteria}\OperatorTok{.}\FunctionTok{get}\OperatorTok{(}\NormalTok{criterionNum}\OperatorTok{).}\FunctionTok{getDiscount}\OperatorTok{()} \OperatorTok{\textgreater{}} \DecValTok{0}\OperatorTok{;}
\end{Highlighting}
\end{Shaded}
\paragraph{Cómo arreglarlo}
Que el texto recorra los criterios en el orden de la lista (\texttt{Criteria\allowbreak{}Values.\allowbreak{}criteria}, que ya existe) en vez de \texttt{key\allowbreak{}Set(\allowbreak{})} del mapa; o cambiar el mapa por un \texttt{Linked\allowbreak{}Hash\allowbreak{}Map} rellenado en el orden de la lista (constructor, líneas 36-42, y \texttt{set\allowbreak{}Criteria\allowbreak{}Values\allowbreak{}Hash\allowbreak{}Map}, que puede sustituirlo). La solución de fondo sería que la ventana reciba los objetos de resultados en vez de trocear un texto, pero es más trabajo.

Sitios afectados por el mismo orden: \texttt{generate\allowbreak{}Table} (tabla) y \texttt{save\allowbreak{}Csv} (exportación). \texttt{save\allowbreak{}Csv} toma los nombres de criterios de lo que \texttt{generate\allowbreak{}Table} dejó en el campo estático \texttt{criteria}, así que su cabecera y sus datos siguen el orden del texto y no se desalinean entre sí por esto.

\texttt{Simulation\allowbreak{}Summary\allowbreak{}Results.\allowbreak{}to\allowbreak{}String} recorre también un \texttt{Hash\allowbreak{}Map\textless{}String,\allowbreak{}\ Criteria\allowbreak{}Values\textgreater{}} por nombre de opción, pero la ventana toma los nombres de las opciones del propio texto, así que ese orden no causa errores.

Prueba de regresión: la situación medida arriba, con criterios «Cost» (con descuento) y «Effectiveness» (sin él), comprobando que la cabecera con dos columnas es la de «Cost» y que la media descontada 90 aparece.

\setcounter{subsection}{102}
\subsection[El CSV de la simulación de sucesos tiene los títulos corridos una columna]{El CSV exportado de la simulación de sucesos tiene la fila de títulos con una columna menos que las de datos, y una hoja de cálculo desalinea todas las columnas}\label{err:103}
\begin{ficha}
\fichakey{Estado}Presente
\fichakey{Módulo}inference.des
\fichakey{Paquete}org.openmarkov.inference.DES
\fichakey{Ficheros}\path{inference.des/src/main/java/org/openmarkov/inference/DES/DESResultsWindow.java:88-117}
\fichakey{Comprobado}Leyendo el código y ejecutándolo
\fichakey{Esfuerzo}Pequeño (horas)
\fichakey{Decisión de equipo}No
\fichakey{Corregir antes}\hyperref[err:102]{error~102}: si su arreglo reescribe la exportación a partir de los objetos de resultados, este error desaparece.
\fichakey{Relacionados}\hyperref[err:102]{error~102}
\end{ficha}
\paragraph{Qué pasa}
El botón «Export as CSV» (CSV: valores separados por comas) de la ventana de resultados de la simulación de sucesos discretos escribe un fichero cuya fila de títulos tiene una celda menos que las de datos. Al abrirlo en una hoja de cálculo, cada título queda encima de la columna de su vecina (el de «Mean cost for none» sobre la serie, etc.).

Camino del usuario: simular una red DESNet (botón de coste-efectividad de la barra de herramientas) → en la ventana de resultados, «Export as CSV» → abrir el fichero en una hoja de cálculo.

\texttt{DESResults\allowbreak{}Window.\allowbreak{}save\allowbreak{}Csv} escribe un fichero separado por comas. La fila de títulos empieza con \texttt{sb.\allowbreak{}append(\allowbreak{}"SERIES;")} y sigue con los títulos separados por comas. El punto y coma no separa nada en ese fichero, así que la primera celda es «SERIES;Mean cost for none» y la columna del número de serie queda sin título. Las filas de datos, en cambio, empiezan por el número de serie seguido de coma.

Comprobado con un pequeño programa de prueba que copia literalmente las líneas 89-117 (el método no se puede llamar sin pantalla porque abre un selector de ficheros) y las aplica a un texto con el formato que produce \texttt{Simulation\allowbreak{}Summary\allowbreak{}Results.\allowbreak{}to\allowbreak{}String(\allowbreak{})}:

\begin{Verbatim}[breaklines,breakanywhere,fontsize=\small]
SERIES;Mean cost for none,Disc. Mean cost for none,Mean effectiveness for none,...,Disc. Mean effectiveness for therapy,
0,100.75,90.675,200.75,180.675,300.5,270.45,400.5,360.45,
fila de títulos: 9 campos, el primero «SERIES;Mean cost for none»
fila de datos:   10 campos, el primero «0»
\end{Verbatim}
\paragraph{El código}
inference.des/src/main/java/org/openmarkov/inference/DES/DESResultsWindow.java:88-107

\begin{Shaded}
\begin{Highlighting}[]
\KeywordTok{private} \DataTypeTok{static} \DataTypeTok{void} \FunctionTok{saveCsv}\OperatorTok{(}\BuiltInTok{String}\NormalTok{ simulationResults}\OperatorTok{)} \OperatorTok{\{}
    \BuiltInTok{Scanner}\NormalTok{ scanner }\OperatorTok{=} \KeywordTok{new} \BuiltInTok{Scanner}\OperatorTok{(}\NormalTok{simulationResults}\OperatorTok{);}
    \BuiltInTok{StringBuilder}\NormalTok{ sb }\OperatorTok{=} \KeywordTok{new} \BuiltInTok{StringBuilder}\OperatorTok{();}
\NormalTok{    sb}\OperatorTok{.}\FunctionTok{append}\OperatorTok{(}\StringTok{"SERIES;"}\OperatorTok{);}
    \CommentTok{// }\AlertTok{TODO}\CommentTok{: skip discounted values when there is no discount.}
\NormalTok{    interventions}\OperatorTok{.}\FunctionTok{forEach}\OperatorTok{(}\NormalTok{intervention }\OperatorTok{{-}\textgreater{}} \OperatorTok{\{}
\NormalTok{        criteria}\OperatorTok{.}\FunctionTok{forEach}\OperatorTok{(}\NormalTok{criterion }\OperatorTok{{-}\textgreater{}} \OperatorTok{\{}
\NormalTok{            sb}\OperatorTok{.}\FunctionTok{append}\OperatorTok{(}\StringTok{"Mean "}\OperatorTok{)}
                    \OperatorTok{.}\FunctionTok{append}\OperatorTok{(}\NormalTok{criterion}\OperatorTok{)}
                    \OperatorTok{.}\FunctionTok{append}\OperatorTok{(}\StringTok{" for "}\OperatorTok{)}
                    \OperatorTok{.}\FunctionTok{append}\OperatorTok{(}\NormalTok{intervention}\OperatorTok{)}
                    \OperatorTok{.}\FunctionTok{append}\OperatorTok{(}\StringTok{","}\OperatorTok{);}
            \KeywordTok{...}
        \OperatorTok{\});}
    \OperatorTok{\});}
\NormalTok{    sb}\OperatorTok{.}\FunctionTok{append}\OperatorTok{(}\StringTok{"}\SpecialCharTok{\textbackslash{}n}\StringTok{"}\OperatorTok{);}
\end{Highlighting}
\end{Shaded}
\paragraph{Cómo arreglarlo}
Cambiar \texttt{"SERIES;"} por \texttt{"SERIES,\allowbreak{}"} en la línea 91. Las dos filas terminan en una coma final que crea un campo vacío; conviene quitarla en las dos a la vez para que sigan teniendo el mismo número de campos.

Hay otra fragilidad en la misma función: las filas de datos se obtienen borrando con \texttt{replace\allowbreak{}All(\allowbreak{}"{[}a-z\allowbreak{}A-Z{]}.\allowbreak{}*?;",\allowbreak{}\ "")} todo lo que empiece por una letra hasta el siguiente punto y coma; un nombre de opción o de criterio que empiece por un número o un símbolo (por ejemplo «10 mg») dejaría restos en la fila. Si se reescribe la exportación a partir de los objetos de resultados (ver \hyperref[err:102]{error~102}), desaparecen las dos cosas.

Prueba de regresión: extraer la construcción del texto del CSV a un método sin selector de ficheros y comprobar, con un texto de dos opciones y dos criterios, que la fila de títulos y las de datos tienen el mismo número de campos y que el primer título es «SERIES».

\setcounter{subsection}{103}
\subsection[Simular con traza una red de sucesos sin guardar falla antes de empezar]{Simular con la traza textual activada una red de sucesos que aún no se ha guardado falla antes de empezar}\label{err:104}
\begin{ficha}
\fichakey{Estado}Presente
\fichakey{Módulo}inference.des
\fichakey{Paquete}org.openmarkov.inference.DES
\fichakey{Ficheros}\path{inference.des/src/main/java/org/openmarkov/inference/DES/DESLogTextWriter.java:53-61} \newline \path{inference.des/src/main/java/org/openmarkov/inference/DES/DESLogTextWriter.java:221-238} \newline \path{inference.des/src/main/java/org/openmarkov/inference/DES/DESInference.java:433} \newline \path{gui/src/main/java/org/openmarkov/gui/window/NetworkFileHandler.java:79-101} \newline \path{gui/src/main/java/org/openmarkov/gui/dialog/inference/common/MonteCarloOptionsPanel.java:313-319} \newline \path{gui/src/main/java/org/openmarkov/gui/window/InferenceHandler.java:274-297}
\fichakey{Comprobado}Leyendo el código y ejecutándolo
\fichakey{Esfuerzo}Pequeño (horas)
\fichakey{Decisión de equipo}No
\fichakey{Relacionados}\hyperref[nota:111]{nota~111}
\end{ficha}
\paragraph{Qué pasa}
La simulación DES (simulación de eventos discretos) puede escribir una traza del algoritmo en un fichero de texto: casilla «Generate evaluation algorithm trace» del panel de opciones de Monte Carlo, que aparece al lanzar la simulación. Con esa casilla marcada, una red de sucesos que aún no se ha guardado no se simula: sale un error inesperado.

Camino de llegada: Archivo, Nuevo, tipo DESNet; construir el modelo; pulsar el botón o menú de coste-efectividad (en una \texttt{DESNet} lanza \texttt{monte\allowbreak{}Carlo\allowbreak{}Simulation}); marcar la traza textual y aceptar. La excepción sale del constructor de \texttt{DESInference} dentro del hilo que crea \texttt{monte\allowbreak{}Carlo\allowbreak{}Simulation}, que solo captura \texttt{IOException} y \texttt{Open\allowbreak{}Markov\allowbreak{}Exception}, así que (por lectura) llega al manejador global y se muestra como error inesperado; no se simula nada.

\texttt{DESInference.\allowbreak{}initialize} construye siempre un \texttt{DESLog\allowbreak{}Text\allowbreak{}Writer}; si la traza está activada, su constructor llama a \texttt{create\allowbreak{}Filename(\allowbreak{}prob\allowbreak{}Net.\allowbreak{}get\allowbreak{}Name(\allowbreak{}))}, que hace \texttt{prob\allowbreak{}Net\allowbreak{}Name.\allowbreak{}substring(\allowbreak{}0,\allowbreak{}\ prob\allowbreak{}Net\allowbreak{}Name.\allowbreak{}index\allowbreak{}Of(\allowbreak{}\textquotesingle{}.\allowbreak{}\textquotesingle{}))}. Si el nombre no tiene punto, \texttt{index\allowbreak{}Of} devuelve -1 y \texttt{substring(\allowbreak{}0,\allowbreak{}\ -1)} lanza \texttt{String\allowbreak{}Index\allowbreak{}Out\allowbreak{}Of\allowbreak{}Bounds\allowbreak{}Exception}.

Una red abierta desde fichero se llama como el fichero (\texttt{asia.\allowbreak{}pgmx}), pero una red creada con Archivo, Nuevo recibe el nombre del texto \texttt{Internal\allowbreak{}Frame.\allowbreak{}Title}, que es «Untitled» (\texttt{Network\allowbreak{}File\allowbreak{}Handler.\allowbreak{}create\allowbreak{}New\allowbreak{}Network}), y lo conserva hasta que se guarda con otro nombre.

Medido con una red de sucesos de prueba (no está en el repositorio): con el nombre cambiado a «Untitled» y la traza activada, \texttt{initialize} termina con \texttt{String\allowbreak{}Index\allowbreak{}Out\allowbreak{}Of\allowbreak{}Bounds\allowbreak{}Exception:\ Range\ {[}0,\allowbreak{}\ -1)\ out\ of\ bounds\ for\ length\ 8}. Con un nombre con extensión (\texttt{nombre.\allowbreak{}pgmx}), escribe la traza en \texttt{DESNet\allowbreak{}Files/\allowbreak{}Textual\allowbreak{}Log/\allowbreak{}nombre-AAAA-MM-DD-HH-mm.\allowbreak{}log}. Para reproducirlo basta cualquier red de sucesos nueva, sin guardar, siguiendo el camino de arriba.

Detalles menores del mismo método:

\begin{itemize}
\tightlist
\item
  Un nombre con varios puntos (\texttt{modelo.\allowbreak{}v2.\allowbreak{}pgmx}) se corta en el primero (\texttt{modelo}).
\item
  El nombre lleva la fecha con resolución de minutos, así que dos simulaciones de la misma red en el mismo minuto escriben en el mismo fichero y la segunda sobrescribe la primera.
\item
  Como el fichero se abre en \texttt{initialize} antes de construir los registros, una simulación que falla después (por ejemplo, la de \hyperref[err:97]{error~97}) deja un fichero de traza vacío.
\end{itemize}
\paragraph{El código}
inference.des/src/main/java/org/openmarkov/inference/DES/DESLogTextWriter.java:221-227

\begin{Shaded}
\begin{Highlighting}[]
    \KeywordTok{private} \BuiltInTok{File} \FunctionTok{createFilename}\OperatorTok{(}\BuiltInTok{String}\NormalTok{ probNetName}\OperatorTok{)} \OperatorTok{\{}

        \BuiltInTok{String}\NormalTok{ filename }\OperatorTok{=}\NormalTok{ probNetName}\OperatorTok{.}\FunctionTok{substring}\OperatorTok{(}\DecValTok{0}\OperatorTok{,}\NormalTok{ probNetName}\OperatorTok{.}\FunctionTok{indexOf}\OperatorTok{(}\CharTok{\textquotesingle{}.\textquotesingle{}}\OperatorTok{));}
\NormalTok{        filename }\OperatorTok{+=} \StringTok{"{-}"}\OperatorTok{;}
\NormalTok{        filename }\OperatorTok{+=} \KeywordTok{new} \BuiltInTok{SimpleDateFormat}\OperatorTok{(}\StringTok{"yyyy{-}MM{-}dd{-}HH{-}mm"}\OperatorTok{).}\FunctionTok{format}\OperatorTok{(}\KeywordTok{new} \BuiltInTok{Date}\OperatorTok{());}
\NormalTok{        filename }\OperatorTok{+=} \StringTok{".log"}\OperatorTok{;}
\end{Highlighting}
\end{Shaded}
\paragraph{Cómo arreglarlo}
Quitar la extensión solo si la hay (último punto, no el primero) y usar el nombre entero si no la hay.

Sitios con la misma suposición sobre el nombre de la red: \texttt{DESLog\allowbreak{}Text\allowbreak{}Writer.\allowbreak{}create\allowbreak{}Filename} es el único que corta en \texttt{index\allowbreak{}Of(\allowbreak{}\textquotesingle{}.\allowbreak{}\textquotesingle{})} dentro de \texttt{inference.\allowbreak{}des}; \texttt{DESResults\allowbreak{}Excel\allowbreak{}Writer.\allowbreak{}create\allowbreak{}Filename} (líneas 181-183) repite exactamente el mismo corte, aunque hoy esa clase no tiene ninguna referencia; si se vuelve a usar, necesita el mismo arreglo.

Prueba de regresión: construir \texttt{DESLog\allowbreak{}Text\allowbreak{}Writer} con la traza activada para redes llamadas «Untitled», «a.pgmx» y «a.b.pgmx» no debe lanzar y debe producir nombres que empiecen por «Untitled», «a» y «a.b».

\setcounter{subsection}{104}
\subsection[Un fallo de «Tree with Excluded Events» abre una ventana desde el núcleo]{Un fallo al muestrear un «Tree with Excluded Events» abre una ventana de aviso desde el núcleo, y sin pantalla la tapa con otra excepción}\label{err:105}
\begin{ficha}
\fichakey{Estado}Presente en parte. La relación hermana, \texttt{Tree\allowbreak{}With\allowbreak{}Events\allowbreak{}Potential}, ya no abre ninguna ventana: lanza \texttt{There\allowbreak{}Must\allowbreak{}Be\allowbreak{}Just\allowbreak{}One\allowbreak{}Event\allowbreak{}In\allowbreak{}Configuration} (commit \texttt{e3a4b1e}).
\fichakey{Módulo}core
\fichakey{Paquete}org.openmarkov.core.model.network.potential.treeadd
\fichakey{Ficheros}\path{core/src/main/java/org/openmarkov/core/model/network/potential/treeadd/TreeWithExcludedEventsPotential.java:28,125-143} \newline \path{core/src/main/java/module-info.java:20} \newline \path{core/src/main/java/org/openmarkov/core/action/core/PurposeEdit.java:18} \newline \path{core/src/main/java/org/openmarkov/core/action/core/MonteCarloOptionsEdit.java:8,32} \newline \path{inference.des/src/main/java/org/openmarkov/inference/DES/ChanceEvaluation.java:85,123-161} \newline \path{inference.des/src/main/java/org/openmarkov/inference/DES/DESRecord.java:228} \newline \path{inference.des/src/main/java/org/openmarkov/inference/DES/EventRecord.java:134}
\fichakey{Comprobado}Leyendo el código y ejecutándolo
\fichakey{Esfuerzo}Pequeño (horas)
\fichakey{Decisión de equipo}No
\fichakey{Relacionados}\hyperref[err:21]{error~21}, \hyperref[err:22]{error~22}, \hyperref[err:23]{error~23}, \hyperref[err:106]{error~106}, \hyperref[nota:68]{nota~68}, \hyperref[nota:74]{nota~74}, \hyperref[nota:170]{nota~170}
\end{ficha}
\paragraph{Qué pasa}
\texttt{Tree\allowbreak{}With\allowbreak{}Excluded\allowbreak{}Events\allowbreak{}Potential} es la relación «Tree with Excluded Events» de la simulación de eventos discretos (DES): un árbol sobre los padres del nodo del que se excluye el suceso. Si algo falla al sortear el valor del nodo, abre una ventana de aviso desde \texttt{core}, el módulo del modelo, que no debería mostrar ventanas:

\begin{itemize}
\tightlist
\item
  Con pantalla (la aplicación), la simulación se detiene en una ventana modal con el texto fijo «TreeWithExcludeEventsPotential:Getting sample exception: », con el nombre de la clase en lugar de una explicación, y queda parada hasta que el usuario pulsa Aceptar. Después llega la excepción original por el camino normal, así que el error se comunica dos veces y la primera sin sentido para el usuario.
\item
  Sin pantalla (una prueba, el módulo \texttt{resttemplate}, un servidor, cualquier ejecución \texttt{headless}), \texttt{JOption\allowbreak{}Pane} lanza \texttt{Headless\allowbreak{}Exception}, que \textbf{sustituye} a la excepción original y oculta la causa.
\end{itemize}

Su método \texttt{sample\allowbreak{}Conditioned\allowbreak{}Variable} (líneas 125-143), que usa la simulación para sortear el valor del nodo (\texttt{DESRecord.\allowbreak{}java:228}, \texttt{Event\allowbreak{}Record.\allowbreak{}java:134}), busca entre los padres de la configuración las variables de tipo suceso, quita el hallazgo del primero y pregunta a su árbol sin sucesos. Si algo falla, su \texttt{catch} (líneas 138-142, con un \texttt{/\allowbreak{}/\allowbreak{}FIXME\ when\ completed\ this\ method\ coding\ this\ catch\ will\ be\ removed}) abre un \texttt{JOption\allowbreak{}Pane.\allowbreak{}show\allowbreak{}Message\allowbreak{}Dialog} y después relanza la excepción.

El fallo que dispara la ventana no es artificial. El propio método tiene un comentario (líneas 131-133) que dice que esta relación se usa también con nodos de azar y que algunas configuraciones no llevan sucesos, y \texttt{Chance\allowbreak{}Evaluation} calcula los nodos de azar sin sucesos antecesores con una configuración sin suceso (\texttt{compute\allowbreak{}Chance\allowbreak{}Value(\allowbreak{}chance\allowbreak{}Record,\allowbreak{}\ null)}, línea 85). En ese caso \texttt{event\allowbreak{}Variables.\allowbreak{}get(\allowbreak{}0)} (línea 135) sobre una lista vacía lanza \texttt{Index\allowbreak{}Out\allowbreak{}Of\allowbreak{}Bounds\allowbreak{}Exception}, que es la que abre la ventana. Ese acceso a una lista vacía es un defecto aparte, que conviene arreglar a la vez.

Medido con un pequeño programa de prueba, sin pantalla, sobre una red de sucesos construida en código con un nodo Y cuyos padres son un suceso E y un nodo Age: muestrear Y con una configuración sin sucesos termina en \texttt{java.\allowbreak{}awt.\allowbreak{}Headless\allowbreak{}Exception}, y la causa real no aparece. Se llamó directamente al sorteo; no se ha ejecutado una simulación completa que llegue a ese \texttt{catch}, así que lo que ve el usuario en la aplicación se deduce de la lectura.

Además de esta llamada, en \texttt{core} quedan otros dos restos de escritorio: \texttt{Monte\allowbreak{}Carlo\allowbreak{}Options\allowbreak{}Edit}, que declara \texttt{javax.\allowbreak{}swing.\allowbreak{}undo.\allowbreak{}Cannot\allowbreak{}Undo\allowbreak{}Exception} en su \texttt{undo(\allowbreak{})} (líneas 8 y 32), y un \texttt{import\ javax.\allowbreak{}swing.\allowbreak{}JOption\allowbreak{}Pane} sin uso en \texttt{Purpose\allowbreak{}Edit} (línea 18). Por eso \texttt{core/\allowbreak{}src/\allowbreak{}main/\allowbreak{}java/\allowbreak{}module-info.\allowbreak{}java} sigue con \texttt{requires\ java.\allowbreak{}desktop} (línea 20).
\paragraph{El código}
core/src/main/java/org/openmarkov/core/model/network/potential/treeadd/TreeWithExcludedEventsPotential.java:125-143

\begin{Shaded}
\begin{Highlighting}[]
\AttributeTok{@Override}
\KeywordTok{public} \DataTypeTok{double} \FunctionTok{sampleConditionedVariable}\OperatorTok{(}\DataTypeTok{double}\OperatorTok{[]}\NormalTok{ randomNumbers}\OperatorTok{,}\NormalTok{ EvidenceCase parents}\OperatorTok{)} \KeywordTok{throws}\NormalTok{ OpenMarkovException }\OperatorTok{\{}
    \BuiltInTok{List}\OperatorTok{\textless{}}\NormalTok{Variable}\OperatorTok{\textgreater{}}\NormalTok{ eventVariables }\OperatorTok{=}\NormalTok{ parents}\OperatorTok{.}\FunctionTok{getVariables}\OperatorTok{()}
                                           \OperatorTok{.}\FunctionTok{stream}\OperatorTok{()}
                                           \OperatorTok{.}\FunctionTok{filter}\OperatorTok{(}\NormalTok{variable }\OperatorTok{{-}\textgreater{}}\NormalTok{ variable}\OperatorTok{.}\FunctionTok{getVariableType}\OperatorTok{()} \OperatorTok{==}\NormalTok{ VariableType}\OperatorTok{.}\FunctionTok{EVENT}\OperatorTok{)}
                                           \OperatorTok{.}\FunctionTok{toList}\OperatorTok{();}
    \KeywordTok{...}
    \ControlFlowTok{try} \OperatorTok{\{}
\NormalTok{        parents}\OperatorTok{.}\FunctionTok{removeFinding}\OperatorTok{(}\NormalTok{eventVariables}\OperatorTok{.}\FunctionTok{get}\OperatorTok{(}\DecValTok{0}\OperatorTok{));}
        \ControlFlowTok{return}\NormalTok{ noEventTree}\OperatorTok{.}\FunctionTok{sampleConditionedVariable}\OperatorTok{(}\NormalTok{randomNumbers}\OperatorTok{,}\NormalTok{ parents}\OperatorTok{);}
    \OperatorTok{\}} \ControlFlowTok{catch} \OperatorTok{(}\NormalTok{OpenMarkovException }\OperatorTok{|} \BuiltInTok{RuntimeException}\NormalTok{ e}\OperatorTok{)} \OperatorTok{\{}
        \CommentTok{//}\AlertTok{FIXME}\CommentTok{ when completed this method coding this catch will be removed}
        \BuiltInTok{JOptionPane}\OperatorTok{.}\FunctionTok{showMessageDialog}\OperatorTok{(}\KeywordTok{null}\OperatorTok{,} \StringTok{"TreeWithExcludeEventsPotential:Getting sample exception: "} \OperatorTok{+} \KeywordTok{this}\OperatorTok{.}\FunctionTok{getConditionedVariable}\OperatorTok{());}
        \ControlFlowTok{throw}\NormalTok{ e}\OperatorTok{;}
    \OperatorTok{\}}
\OperatorTok{\}}
\end{Highlighting}
\end{Shaded}
\paragraph{Cómo arreglarlo}
Quitar el \texttt{try/\allowbreak{}catch} con el diálogo y dejar salir la excepción, o envolverla en una excepción del dominio con un mensaje que diga qué nodo, qué configuración y qué falta; quien lanza la simulación en la interfaz decide cómo enseñarla. Es lo que ya hace \texttt{Tree\allowbreak{}With\allowbreak{}Events\allowbreak{}Potential.\allowbreak{}sample\allowbreak{}Conditioned\allowbreak{}Variable} (líneas 188-190), que comprueba antes cuántos sucesos hay en la configuración. Quitar también el import de la línea 28.

Quitar el hallazgo del suceso solo si hay alguno (\texttt{event\allowbreak{}Variables.\allowbreak{}is\allowbreak{}Empty(\allowbreak{})}), coherente con el comentario que admite configuraciones sin sucesos.

Ese mismo método modifica la configuración que recibe (\texttt{remove\allowbreak{}Finding}), cosa que describe \hyperref[nota:74]{nota~74}; allí no hay nada que corregir, pero el cambio cae en el mismo sitio.

Para que \texttt{core} deje de depender del escritorio: sustituir \texttt{Cannot\allowbreak{}Undo\allowbreak{}Exception} en \texttt{Monte\allowbreak{}Carlo\allowbreak{}Options\allowbreak{}Edit} por una excepción propia o ninguna, borrar el import sin uso de \texttt{Purpose\allowbreak{}Edit} y retirar \texttt{requires\ java.\allowbreak{}desktop} de \texttt{core/\allowbreak{}src/\allowbreak{}main/\allowbreak{}java/\allowbreak{}module-info.\allowbreak{}java}, de modo que cualquier diálogo nuevo en el núcleo sea un error de compilación. La propuesta de conjunto está en \hyperref[nota:170]{nota~170}, que es una nota: allí no hay ningún error que corregir.

Prueba: construir un \texttt{Tree\allowbreak{}With\allowbreak{}Excluded\allowbreak{}Events\allowbreak{}Potential}, pedirle un sorteo con una configuración sin suceso bajo \texttt{-Djava.\allowbreak{}awt.\allowbreak{}headless=true} y comprobar que no sale \texttt{Headless\allowbreak{}Exception}: debe salir la excepción del dominio o, si se decide admitir configuraciones sin sucesos, la muestra del árbol.

\setcounter{subsection}{105}
\subsection[Un nodo de azar discretizado impide simular con un error inesperado]{Un nodo de azar discretizado en una red de sucesos impide simular con un error inesperado, y el mensaje previsto diría que ese nodo «debe ser de azar»}\label{err:106}
\begin{ficha}
\fichakey{Estado}Presente
\fichakey{Módulo}inference.des, gui
\fichakey{Paquete}org.openmarkov.inference.DES; org.openmarkov.gui.dialog.node
\fichakey{Ficheros}\path{inference.des/src/main/java/org/openmarkov/inference/DES/ChanceRecord.java:76-84} \newline \path{inference.des/src/main/java/org/openmarkov/inference/DES/GenericEvaluation.java:63-80} \newline \path{inference.des/src/main/java/org/openmarkov/inference/DES/DESInference.java:105-113,437} \newline \path{inference.des/src/main/java/org/openmarkov/inference/DES/exception/NodeMustBeChance.java:6-13} \newline \path{inference.des/src/main/resources/inference_des/localize/inference_des_class_localizations_en.xml:5-6} \newline \path{core/src/main/java/org/openmarkov/core/exception/OpenMarkovException.java:27-35} \newline \path{core/src/main/java/org/openmarkov/core/model/network/VariableType.java:58-65} \newline \path{gui/src/main/java/org/openmarkov/gui/dialog/node/NodeDomainValuesTablePanel.java:856-868, 886-894} \newline \path{gui/src/main/java/org/openmarkov/gui/window/InferenceHandler.java:274-297}
\fichakey{Comprobado}Leyendo el código y ejecutándolo
\fichakey{Esfuerzo}Pequeño (horas)
\fichakey{Decisión de equipo}\textcolor{decisionrojo}{\textbf{Sí.}} Si una red de sucesos debe admitir nodos de azar discretizados (y entonces la simulación los trata como de estados) o si el editor no debe ofrecer ese tipo en una red de sucesos
\fichakey{Relacionados}\hyperref[err:97]{error~97}, \hyperref[err:99]{error~99}, \hyperref[err:105]{error~105}, \hyperref[nota:55]{nota~55}
\end{ficha}
\paragraph{Qué pasa}
En una red DES (simulación de eventos discretos), el usuario abre las propiedades de un nodo de azar y, en la pestaña del dominio, cambia el tipo de variable a «Discretized» (el selector ofrece a todo nodo de azar los tipos de \texttt{Variable\allowbreak{}Type.\allowbreak{}of(\allowbreak{}CHANCE)}: estados finitos, discretizada y numérica, en cualquier tipo de red). La edición se acepta. Al lanzar la simulación (botón o menú de coste-efectividad), no se simula nada y sale el diálogo de error inesperado.

La simulación no admite nodos de azar discretizados, y lo comprueba a propósito: el constructor de \texttt{Chance\allowbreak{}Record} (líneas 76-84) lanza \texttt{Node\allowbreak{}Must\allowbreak{}Be\allowbreak{}Chance} si la variable del nodo no es numérica ni de estados finitos. Pero ese aviso previsto no llega al usuario, por tres motivos que se encadenan:

\begin{itemize}
\tightlist
\item
  \texttt{Generic\allowbreak{}Evaluation} (líneas 63-80) construye los registros por reflexión y convierte cualquier \texttt{Invocation\allowbreak{}Target\allowbreak{}Exception} en \texttt{Unreachable\allowbreak{}Exception}, así que la excepción de la red del usuario se presenta como fallo del programa. En la aplicación, el hilo que crea \texttt{monte\allowbreak{}Carlo\allowbreak{}Simulation} no la captura y llega al manejador global, que abre el diálogo de error inesperado (por lectura).
\item
  Si llegara al \texttt{catch\ (\allowbreak{}Open\allowbreak{}Markov\allowbreak{}Exception\ e)} del constructor de \texttt{DESInference} (líneas 105-113), la ventana «Cannot simulate» enseñaría \texttt{e.\allowbreak{}get\allowbreak{}Message(\allowbreak{})}, y en las excepciones de OpenMarkov ese mensaje es nulo: el texto está en \texttt{to\allowbreak{}String(\allowbreak{})} (\texttt{Open\allowbreak{}Markov\allowbreak{}Exception}, líneas 27-35).
\item
  El texto de \texttt{Node\allowbreak{}Must\allowbreak{}Be\allowbreak{}Chance} es «\{node\} must be chance.», que para un nodo que sí es de azar no explica nada.
\end{itemize}

Medido con una red de sucesos de prueba que simula sin problemas, cambiando el nodo de azar \texttt{Sex} a discretizado con \texttt{Variable\allowbreak{}Type\allowbreak{}Edit} (la edición que ejecuta el selector del diálogo). Las restricciones de la red no lo impiden, y el nodo conserva su tabla. Al simular:

\begin{Verbatim}[breaklines,breakanywhere,fontsize=\small]
DESInference.initialize -> UnreachableException: java.lang.reflect.InvocationTargetException
  causa: chance node of Discretized variable Sex with states [state male, state female] must be chance.
    at ChanceRecord.<init>(ChanceRecord.java:81)
    at GenericEvaluation.<init>(GenericEvaluation.java:72)
    at ChanceEvaluation.<init>(ChanceEvaluation.java:40)
    at DESInference.initialize(DESInference.java:437)
NodeMustBeChance.getMessage() = null
\end{Verbatim}
\paragraph{El código}
inference.des/src/main/java/org/openmarkov/inference/DES/ChanceRecord.java:76-81

\begin{Shaded}
\begin{Highlighting}[]
    \KeywordTok{public} \FunctionTok{ChanceRecord}\OperatorTok{(}\BuiltInTok{Node}\NormalTok{ chanceNode}\OperatorTok{,} \DataTypeTok{double}\NormalTok{ variableValue}\OperatorTok{)} \KeywordTok{throws}\NormalTok{ OpenMarkovException }\OperatorTok{\{}
        \KeywordTok{super}\OperatorTok{(}\NormalTok{chanceNode}\OperatorTok{);}
        \ControlFlowTok{if} \OperatorTok{((}\NormalTok{chanceNode}\OperatorTok{.}\FunctionTok{getNodeType}\OperatorTok{()} \OperatorTok{!=}\NormalTok{ NodeType}\OperatorTok{.}\FunctionTok{CHANCE}\OperatorTok{)}
                \OperatorTok{||} \OperatorTok{((}\NormalTok{chanceNode}\OperatorTok{.}\FunctionTok{getVariable}\OperatorTok{().}\FunctionTok{getVariableType}\OperatorTok{()} \OperatorTok{!=}\NormalTok{ VariableType}\OperatorTok{.}\FunctionTok{NUMERIC}\OperatorTok{)}
                \OperatorTok{\&\&} \OperatorTok{(}\NormalTok{chanceNode}\OperatorTok{.}\FunctionTok{getVariable}\OperatorTok{().}\FunctionTok{getVariableType}\OperatorTok{()} \OperatorTok{!=}\NormalTok{ VariableType}\OperatorTok{.}\FunctionTok{FINITE\_STATES}\OperatorTok{)))}
            \ControlFlowTok{throw} \KeywordTok{new} \FunctionTok{NodeMustBeChance}\OperatorTok{(}\NormalTok{chanceNode}\OperatorTok{);}
\end{Highlighting}
\end{Shaded}

inference.des/src/main/java/org/openmarkov/inference/DES/GenericEvaluation.java:68-78

\begin{Shaded}
\begin{Highlighting}[]
        \ControlFlowTok{try} \OperatorTok{\{}
            \BuiltInTok{Constructor}\OperatorTok{\textless{}}\NormalTok{T}\OperatorTok{\textgreater{}}\NormalTok{ recordConstructor }\OperatorTok{=}\NormalTok{ recordClass}\OperatorTok{.}\FunctionTok{getDeclaredConstructor}\OperatorTok{(}\BuiltInTok{Node}\OperatorTok{.}\FunctionTok{class}\OperatorTok{);}
            \ControlFlowTok{for} \OperatorTok{(}\BuiltInTok{Node}\NormalTok{ node }\OperatorTok{:}\NormalTok{ nodeList}\OperatorTok{)} \OperatorTok{\{}
\NormalTok{                lastNode }\OperatorTok{=}\NormalTok{ node}\OperatorTok{;}
\NormalTok{                desRecordHashMap}\OperatorTok{.}\FunctionTok{put}\OperatorTok{(}\NormalTok{node}\OperatorTok{,}\NormalTok{ recordConstructor}\OperatorTok{.}\FunctionTok{newInstance}\OperatorTok{(}\NormalTok{node}\OperatorTok{));}
            \OperatorTok{\}}

        \OperatorTok{\}} \ControlFlowTok{catch} \OperatorTok{(}\BuiltInTok{NoSuchMethodException} \OperatorTok{|} \BuiltInTok{InvocationTargetException} \OperatorTok{|} \BuiltInTok{InstantiationException} \OperatorTok{|}
                 \BuiltInTok{IllegalAccessException}\NormalTok{ e}\OperatorTok{)} \OperatorTok{\{}
            \ControlFlowTok{throw} \KeywordTok{new} \FunctionTok{UnreachableException}\OperatorTok{(}\NormalTok{e}\OperatorTok{);}
        \OperatorTok{\}}
\end{Highlighting}
\end{Shaded}

inference.des/src/main/java/org/openmarkov/inference/DES/DESInference.java:107-113

\begin{Shaded}
\begin{Highlighting}[]
        \ControlFlowTok{try} \OperatorTok{\{}
            \ControlFlowTok{if} \OperatorTok{(!}\FunctionTok{initialize}\OperatorTok{(}\NormalTok{probNet}\OperatorTok{))}
                \ControlFlowTok{return}\OperatorTok{;}
        \OperatorTok{\}} \ControlFlowTok{catch} \OperatorTok{(}\NormalTok{OpenMarkovException e}\OperatorTok{)} \OperatorTok{\{}
            \BuiltInTok{JOptionPane}\OperatorTok{.}\FunctionTok{showMessageDialog}\OperatorTok{(}\KeywordTok{null}\OperatorTok{,}\NormalTok{ e}\OperatorTok{.}\FunctionTok{getMessage}\OperatorTok{(),} \StringTok{"Cannot simulate"}\OperatorTok{,} \BuiltInTok{JOptionPane}\OperatorTok{.}\FunctionTok{ERROR\_MESSAGE}\OperatorTok{);}
            \ControlFlowTok{return}\OperatorTok{;}
        \OperatorTok{\}}
\end{Highlighting}
\end{Shaded}
\paragraph{Cómo arreglarlo}
Lo que no necesita decisión:

\begin{itemize}
\tightlist
\item
  En \texttt{Generic\allowbreak{}Evaluation}, si la causa de la \texttt{Invocation\allowbreak{}Target\allowbreak{}Exception} es una \texttt{Open\allowbreak{}Markov\allowbreak{}Exception}, relanzar esa causa en lugar de envolverla en \texttt{Unreachable\allowbreak{}Exception} (el constructor ya declara \texttt{throws\ Open\allowbreak{}Markov\allowbreak{}Exception}). Solo los errores de reflexión propiamente dichos son inalcanzables.
\item
  En \texttt{DESInference}, enseñar la excepción con su texto localizado (\texttt{to\allowbreak{}String(\allowbreak{})}, o \texttt{Exception\allowbreak{}Dialog}) en vez de \texttt{get\allowbreak{}Message(\allowbreak{})}.
\item
  Dar a este caso un mensaje que diga lo que pasa: por ejemplo, que la variable del nodo de azar es discretizada y la simulación solo admite nodos de azar numéricos o de estados finitos. \texttt{Node\allowbreak{}Must\allowbreak{}Be\allowbreak{}Chance} se usa también para el caso de tipo de nodo, que hoy no se alcanza porque \texttt{Generic\allowbreak{}Evaluation} elige los nodos por su tipo.
\end{itemize}

Lo que necesita decisión: si la red de sucesos admite nodos de azar discretizados. Si no los admite, el selector de tipo de variable no debería ofrecer «Discretized» en una red de sucesos (\texttt{Node\allowbreak{}Domain\allowbreak{}Values\allowbreak{}Table\allowbreak{}Panel}, líneas 855-867, que hoy usa \texttt{Variable\allowbreak{}Type.\allowbreak{}of} sin mirar el tipo de red). La regla tendría que cumplirse también en \texttt{Variable\allowbreak{}Type\allowbreak{}Edit.\allowbreak{}check\allowbreak{}Constraints\allowbreak{}Will\allowbreak{}Be\allowbreak{}Met} y al abrir un fichero. Si los admite, \texttt{Chance\allowbreak{}Record} y el muestreo deben tratarlos como variables de estados.

Los mismos constructores por reflexión lanzan \texttt{Node\allowbreak{}Must\allowbreak{}Be\allowbreak{}Event}, \texttt{Node\allowbreak{}Must\allowbreak{}Be\allowbreak{}Utility} y \texttt{Event\allowbreak{}Is\allowbreak{}Not\allowbreak{}Parent\allowbreak{}Of}, y un fallo del potencial (como el de \hyperref[err:99]{error~99}) o un desbordamiento de pila (como el de \hyperref[err:97]{error~97}) acaba en el mismo \texttt{Unreachable\allowbreak{}Exception}. El cambio en \texttt{Generic\allowbreak{}Evaluation} los afecta a todos.

Prueba de regresión: con una red de sucesos cuyo nodo de azar es discretizado, \texttt{new\ DESInference(\allowbreak{}red,\allowbreak{}\ monitor)} (o \texttt{initialize}) no lanza \texttt{Unreachable\allowbreak{}Exception}: o bien rechaza con una \texttt{Open\allowbreak{}Markov\allowbreak{}Exception} cuyo texto nombra el nodo y el tipo de variable, o bien simula, según lo decidido.

\setcounter{subsection}{106}
\subsection[La hoja de la propagación estocástica cambia el orden de los hallazgos]{La hoja exportada de la propagación estocástica escribe los bloques de los hallazgos en un orden distinto en cada ejecución}\label{err:107}
\begin{ficha}
\fichakey{Estado}Presente
\fichakey{Módulo}stochasticPropagationOutput
\fichakey{Paquete}org.openmarkov.stochasticPropagationOutput
\fichakey{Ficheros}\path{stochasticPropagationOutput/src/main/java/org/openmarkov/stochasticPropagationOutput/XlsxWrite.java:123-125} \newline \path{stochasticPropagationOutput/src/main/java/org/openmarkov/stochasticPropagationOutput/XlsxWrite.java:362-364} \newline \path{core/src/main/java/org/openmarkov/core/model/network/EvidenceCase.java:37} \newline \path{core/src/main/java/org/openmarkov/core/model/network/EvidenceCase.java:196-198}
\fichakey{Comprobado}Leyendo el código y ejecutándolo
\fichakey{Esfuerzo}Pequeño (horas)
\fichakey{Decisión de equipo}No
\fichakey{Relacionados}\hyperref[err:93]{error~93}, \hyperref[err:108]{error~108}
\end{ficha}
\paragraph{Qué pasa}
Exportar dos veces la misma red con la misma semilla y los mismos hallazgos produce ficheros con los bloques de hallazgos en orden distinto, lo que impide compararlos. Con un solo hallazgo no se nota.

Camino del usuario: con una red abierta y hallazgos introducidos, menú Tools → «Export stochastic propagation data» (\texttt{Stochastic\allowbreak{}Propagation\allowbreak{}Output\allowbreak{}Plugin}), que abre \texttt{Stochastic\allowbreak{}Propagation\allowbreak{}Output\allowbreak{}Frame}; al exportar, un hilo con \texttt{Xlsx\allowbreak{}Write} escribe un fichero \texttt{.\allowbreak{}xlsx} con las muestras y un resumen por variable.

En el resumen, para cada variable con hallazgo que no se ha muestreado (en el muestreo ponderado por verosimilitud, las variables observadas no se muestrean), \texttt{Xlsx\allowbreak{}Write} escribe un bloque (nombre, estados, en cuál se observó). Recorre esas variables en el orden de \texttt{algorithm.\allowbreak{}get\allowbreak{}Fused\allowbreak{}Evidence(\allowbreak{}).\allowbreak{}get\allowbreak{}Variables(\allowbreak{})}. Las variables muestreadas (el grueso de la hoja) sí van en un orden estable (\texttt{get\allowbreak{}Variables\allowbreak{}To\allowbreak{}Sample(\allowbreak{})}).

\texttt{Evidence\allowbreak{}Case} guarda los hallazgos en un \texttt{Hash\allowbreak{}Map\textless{}Variable,\allowbreak{}\ Finding\textgreater{}} y \texttt{get\allowbreak{}Variables(\allowbreak{})} devuelve \texttt{new\ Array\allowbreak{}List\textless{}\textgreater{}(\allowbreak{}findings.\allowbreak{}key\allowbreak{}Set(\allowbreak{}))}: el orden es el de la tabla de dispersión. Como \texttt{Variable} no redefine \texttt{hash\allowbreak{}Code} ni \texttt{equals}, se usa el código de identidad del objeto, que la máquina virtual asigna de forma distinta en cada ejecución.

La causa está en \texttt{Evidence\allowbreak{}Case}, no en el exportador: \texttt{get\allowbreak{}Variables(\allowbreak{})} tiene muchos llamadores en el proyecto, y cualquiera que dependa del orden hereda el mismo comportamiento.

Medido con un pequeño programa de prueba que lee \texttt{inference/\allowbreak{}src/\allowbreak{}test/\allowbreak{}resources/\allowbreak{}asia.\allowbreak{}pgmx}, crea un \texttt{Evidence\allowbreak{}Case} con hallazgos en «Positive X-ray?», «Has bronchitis» y «Smoker?» (en ese orden) y escribe \texttt{get\allowbreak{}Variables(\allowbreak{})}, seis ejecuciones:

\begin{Verbatim}[breaklines,breakanywhere,fontsize=\small]
[Smoker?, Positive X-ray?, Has bronchitis]
[Smoker?, Has bronchitis, Positive X-ray?]
[Has bronchitis, Smoker?, Positive X-ray?]
[Has bronchitis, Smoker?, Positive X-ray?]
[Has bronchitis, Smoker?, Positive X-ray?]
[Positive X-ray?, Smoker?, Has bronchitis]
\end{Verbatim}
\paragraph{El código}
core/src/main/java/org/openmarkov/core/model/network/EvidenceCase.java:37, 196-198

\begin{Shaded}
\begin{Highlighting}[]
\KeywordTok{protected} \DataTypeTok{final} \BuiltInTok{HashMap}\OperatorTok{\textless{}}\NormalTok{Variable}\OperatorTok{,}\NormalTok{ Finding}\OperatorTok{\textgreater{}}\NormalTok{ findings}\OperatorTok{;}
\KeywordTok{...}
\KeywordTok{public} \BuiltInTok{List}\OperatorTok{\textless{}}\NormalTok{Variable}\OperatorTok{\textgreater{}} \FunctionTok{getVariables}\OperatorTok{()} \OperatorTok{\{}
    \ControlFlowTok{return} \KeywordTok{new} \BuiltInTok{ArrayList}\OperatorTok{\textless{}\textgreater{}(}\NormalTok{findings}\OperatorTok{.}\FunctionTok{keySet}\OperatorTok{());}
\OperatorTok{\}}
\end{Highlighting}
\end{Shaded}

stochasticPropagationOutput/src/main/java/org/openmarkov/stochasticPropagationOutput/XlsxWrite.java:123-125, 362-364

\begin{Shaded}
\begin{Highlighting}[]
\CommentTok{// Evidence variables}
\NormalTok{EvidenceCase evidence }\OperatorTok{=}\NormalTok{ algorithm}\OperatorTok{.}\FunctionTok{getFusedEvidence}\OperatorTok{();}
\BuiltInTok{List}\OperatorTok{\textless{}}\NormalTok{Variable}\OperatorTok{\textgreater{}}\NormalTok{ evidenceVariables }\OperatorTok{=}\NormalTok{ evidence}\OperatorTok{.}\FunctionTok{getVariables}\OperatorTok{();}
\KeywordTok{...}
\CommentTok{// Write about evidence variables that have not been sampled}
\ControlFlowTok{for} \OperatorTok{(}\NormalTok{Variable evidenceVariable }\OperatorTok{:}\NormalTok{ evidenceVariables}\OperatorTok{)} \OperatorTok{\{}
    \ControlFlowTok{if} \OperatorTok{(!}\NormalTok{sampledVariables}\OperatorTok{.}\FunctionTok{contains}\OperatorTok{(}\NormalTok{evidenceVariable}\OperatorTok{))} \OperatorTok{\{}
\end{Highlighting}
\end{Shaded}
\paragraph{Cómo arreglarlo}
Cambiar el mapa de \texttt{Evidence\allowbreak{}Case} a \texttt{Linked\allowbreak{}Hash\allowbreak{}Map} (orden de inserción), en la declaración del campo y en los constructores que lo crean (líneas 61 y 68; el constructor de la línea 46 recibe un \texttt{Hash\allowbreak{}Map} de fuera y habría que copiarlo a un \texttt{Linked\allowbreak{}Hash\allowbreak{}Map} o cambiar su tipo).

Hay que revisar que los constructores de copia de \texttt{Evidence\allowbreak{}Case} y \texttt{Stochastic\allowbreak{}Propagation} (línea 114, \texttt{new\ Evidence\allowbreak{}Case(\allowbreak{}post\allowbreak{}Resolution\allowbreak{}Evidence)}) conserven el orden.

Es el mismo patrón (mapas sin orden sobre objetos sin \texttt{hash\allowbreak{}Code}) que causa \hyperref[err:4]{error~4} en las operaciones de potenciales, pero en otro sitio: arreglar aquél no arregla éste.

Alternativa local, peor: ordenar en \texttt{Xlsx\allowbreak{}Write} por el orden de las variables de la red.

Prueba: crear un \texttt{Evidence\allowbreak{}Case} con hallazgos en tres variables en un orden dado y comprobar que \texttt{get\allowbreak{}Variables(\allowbreak{})} los devuelve en ese orden. Como el defecto depende de la dispersión por identidad, la prueba debe repetirse con varias instancias o en varios procesos para que tenga valor.

\setcounter{subsection}{107}
\subsection[La propagación estocástica no rechaza nodos numéricos: falla sin explicarlo]{«Export stochastic propagation data» sobre una red bayesiana con un nodo numérico de azar sale como error no previsto, o con un mensaje técnico, en vez de decir que la herramienta no admite nodos numéricos}\label{err:108}
\begin{ficha}
\fichakey{Estado}Presente
\fichakey{Módulo}stochasticPropagationOutput, inference
\fichakey{Paquete}org.openmarkov.stochasticPropagationOutput; org.openmarkov.inference.algorithm.likelihoodWeighting; org.openmarkov.inference.algorithm.huginPropagation
\fichakey{Ficheros}\path{stochasticPropagationOutput/src/main/java/org/openmarkov/stochasticPropagationOutput/StochasticPropagationOutputFrame.java:175-189, 194-244} \newline \path{stochasticPropagationOutput/src/main/java/org/openmarkov/stochasticPropagationOutput/StochasticPropagationOutputPlugin.java:36-49} \newline \path{inference/src/main/java/org/openmarkov/inference/algorithm/likelihoodWeighting/StochasticPropagation.java:77-86} \newline \path{inference/src/main/java/org/openmarkov/inference/algorithm/huginPropagation/HuginPropagation.java:56-65} \newline \path{core/src/main/java/org/openmarkov/core/model/network/potential/Potential.java:672-675} \newline \path{core/src/main/java/org/openmarkov/core/inference/InferenceAlgorithm.java:94-97,154-171}
\fichakey{Comprobado}Leyendo el código y ejecutándolo
\fichakey{Esfuerzo}Pequeño (horas)
\fichakey{Decisión de equipo}No
\fichakey{Relacionados}\hyperref[err:42]{error~42}, \hyperref[err:48]{error~48}, \hyperref[err:107]{error~107}, \hyperref[nota:109]{nota~109}, \hyperref[nota:168]{nota~168}
\end{ficha}
\paragraph{Qué pasa}
Una red bayesiana puede tener nodos numéricos de azar: por omisión, las propiedades de la red permiten variables discretas y continuas, y en la pestaña de dominio del diálogo de un nodo se puede cambiar su tipo a numérico. Ninguna red bayesiana del repositorio los tiene; se llega construyéndola en el editor.

Con una red así abierta, el menú Tools → «Export stochastic propagation data» abre la ventana «Stochastic propagation output». La ventana solo comprueba que haya una red, que sea bayesiana y que tenga nodos (\texttt{bound\allowbreak{}Error}, líneas 182-196). Al pulsar «Export to a .xlsx file», \texttt{on\allowbreak{}Action\allowbreak{}Request} construye el algoritmo de muestreo elegido y el algoritmo exacto de Hugin (propagación por agrupamiento), compila Hugin, muestrea y después calcula las probabilidades exactas. Ninguno de los dos algoritmos declara que no admite variables numéricas: solo exigen que la red sea bayesiana (\texttt{Stochastic\allowbreak{}Propagation.\allowbreak{}get\allowbreak{}Additional\allowbreak{}Constraints} devuelve una lista vacía y \texttt{Hugin\allowbreak{}Propagation.\allowbreak{}get\allowbreak{}Additional\allowbreak{}Constraints}, \texttt{null}). El fallo llega después, y según la relación del nodo:

\begin{itemize}
\tightlist
\item
  Si la relación no se puede convertir en tabla, falla la compilación de Hugin con \texttt{Non\allowbreak{}Projectable\allowbreak{}Potential\allowbreak{}Exception}. \texttt{action\allowbreak{}Performed} la envuelve y sale un error normal: «Potential P(N \textbar{} P) cannot be converted to a table». No dice qué hacer ni que el problema es el tipo del nodo.
\item
  Si se puede convertir en tabla, Hugin compila y falla el muestreo: \texttt{Potential.\allowbreak{}sample\allowbreak{}Conditioned\allowbreak{}Variable} lanza \texttt{Unsupported\allowbreak{}Operation\allowbreak{}Exception} para toda clase de potencial que no lo redefine. Es una excepción de tiempo de ejecución que \texttt{action\allowbreak{}Performed} no recoge, y sale la ventana de error no previsto que pide enviar el fallo a los desarrolladores.
\end{itemize}

Medido con un pequeño programa de prueba que construye la red P → N con las ediciones del editor (añadir dos nodos de azar, enlazarlos, \texttt{Variable\allowbreak{}Type\allowbreak{}Edit} a numérico), pone en N cada relación que el diálogo del nodo ofrece, instanciada como la instancia el diálogo, y repite las llamadas de \texttt{on\allowbreak{}Action\allowbreak{}Request} en su orden:

\begin{itemize}
\tightlist
\item
  sin tocar la relación (el editor deja un \texttt{Uniform\allowbreak{}Potential}), y con «Binomial», «DistributionTable», «Function», «Uniform» y «UnivariateDistr»: error normal, «Potential P(N \textbar{} P) cannot be converted to a table»;
\item
  con «Delta», «Exact» (también con valores 3 y 7), «Exponential», «Linear combination», «Product», «Sum» y «Tree/ADD»: error no previsto, \texttt{Unsupported\allowbreak{}Operation\allowbreak{}Exception:\ sample\allowbreak{}Conditioned\allowbreak{}Variable\ not\ implemented\ for\ …}.
\end{itemize}

Da lo mismo con «Likelihood weighting», el algoritmo marcado por omisión, que con «Logic sampling», y con la red P → N → C, en la que el nodo numérico es además padre de otro. Con la relación «Exact», la misma red tampoco se puede propagar en el modo inferencia de la ventana principal, por otra causa (\hyperref[err:42]{error~42}).
\paragraph{El código}
stochasticPropagationOutput/src/main/java/org/openmarkov/stochasticPropagationOutput/StochasticPropagationOutputFrame.java:182-196

\begin{Shaded}
\begin{Highlighting}[]
\KeywordTok{public} \BuiltInTok{Exception} \FunctionTok{boundError}\OperatorTok{()} \OperatorTok{\{}
    \CommentTok{// Stop if there is no net opened yet}
    \ControlFlowTok{if} \OperatorTok{(}\NormalTok{probNet }\OperatorTok{==} \KeywordTok{null}\OperatorTok{)} \OperatorTok{\{}
        \ControlFlowTok{return} \KeywordTok{new} \FunctionTok{NoNetOpenedException}\OperatorTok{();}
    \OperatorTok{\}}
    \CommentTok{// Stop if other than bayesian net}
    \ControlFlowTok{if} \OperatorTok{(!(}\NormalTok{probNet}\OperatorTok{.}\FunctionTok{getNetworkType}\OperatorTok{()} \KeywordTok{instanceof}\NormalTok{ BayesianNetworkType}\OperatorTok{))} \OperatorTok{\{}
        \ControlFlowTok{return} \KeywordTok{new}\NormalTok{ InvalidNetworkTypeException}\OperatorTok{.}\FunctionTok{NotAllowedType}\OperatorTok{(}\NormalTok{probNet}\OperatorTok{,} \BuiltInTok{Arrays}\OperatorTok{.}\FunctionTok{asList}\OperatorTok{(}\NormalTok{BayesianNetworkType}\OperatorTok{.}\FunctionTok{getUniqueInstance}\OperatorTok{()));}
    \OperatorTok{\}}
    \CommentTok{// Stop if no nodes in the net}
    \ControlFlowTok{if} \OperatorTok{(}\NormalTok{probNet}\OperatorTok{.}\FunctionTok{getNumNodes}\OperatorTok{()} \OperatorTok{==} \DecValTok{0}\OperatorTok{)} \OperatorTok{\{}
        \ControlFlowTok{return} \KeywordTok{new} \FunctionTok{NetworkHasNoNodesException}\OperatorTok{(}\NormalTok{probNet}\OperatorTok{);}
    \OperatorTok{\}}
    \ControlFlowTok{return} \KeywordTok{null}\OperatorTok{;}
\OperatorTok{\}}
\end{Highlighting}
\end{Shaded}

core/src/main/java/org/openmarkov/core/model/network/potential/Potential.java:672-675

\begin{Shaded}
\begin{Highlighting}[]
\KeywordTok{public} \DataTypeTok{int} \FunctionTok{sampleConditionedVariable}\OperatorTok{(}\BuiltInTok{Random}\NormalTok{ randomGenerator}\OperatorTok{,} \BuiltInTok{Map}\OperatorTok{\textless{}}\NormalTok{Variable}\OperatorTok{,} \BuiltInTok{Integer}\OperatorTok{\textgreater{}}\NormalTok{ sampledParents}\OperatorTok{)} \OperatorTok{\{}
    \ControlFlowTok{throw} \KeywordTok{new} \BuiltInTok{UnsupportedOperationException}\OperatorTok{(}
            \StringTok{"sampleConditionedVariable not implemented for "} \OperatorTok{+} \FunctionTok{getClass}\OperatorTok{().}\FunctionTok{getSimpleName}\OperatorTok{());}
\OperatorTok{\}}
\end{Highlighting}
\end{Shaded}
\paragraph{Cómo arreglarlo}
Rechazar la red antes de calcular, con un mensaje que diga que la herramienta solo admite variables discretas (de estados finitos o discretizadas). Hay dos sitios, y conviene usar los dos:

\begin{itemize}
\tightlist
\item
  En los algoritmos, para que cualquier llamador reciba la negativa: que \texttt{Stochastic\allowbreak{}Propagation.\allowbreak{}get\allowbreak{}Additional\allowbreak{}Constraints} y \texttt{Hugin\allowbreak{}Propagation.\allowbreak{}get\allowbreak{}Additional\allowbreak{}Constraints} exijan que todas las variables de azar sean discretas. Existe ya la restricción \texttt{Only\allowbreak{}Discrete\allowbreak{}Variables}, con su texto en inglés; \texttt{Inference\allowbreak{}Algorithm.\allowbreak{}check\allowbreak{}Constraints\allowbreak{}Applicability} la comprobaría al construir el algoritmo, y la ventana ya convierte \texttt{Constraint\allowbreak{}Violated\allowbreak{}Exception} en un error normal. El otro llamador de producción de Hugin, la maximización de la esperanza del aprendizaje (\texttt{EMAlgorithm}), trabaja con redes aprendidas de bases de casos, de estados finitos, y no le afectaría.
\item
  En la ventana, en \texttt{bound\allowbreak{}Error}, para avisar al abrir la herramienta en vez de al pulsar el botón, igual que ya se hace con el tipo de red.
\end{itemize}

Que el muestreo admita nodos numéricos sería otra funcionalidad, no el arreglo de este fallo.

Prueba: construir \texttt{Likelihood\allowbreak{}Weighting}, \texttt{Logic\allowbreak{}Sampling} y \texttt{Hugin\allowbreak{}Propagation} sobre una red bayesiana con un nodo numérico de azar con relación «Exact» debe lanzar la excepción declarada con el nombre del nodo, nunca \texttt{Unsupported\allowbreak{}Operation\allowbreak{}Exception}; y con \texttt{0\_\allowbreak{}miscellaneous/\allowbreak{}networks/\allowbreak{}bn/\allowbreak{}BN-asia.\allowbreak{}pgmx} deben seguir construyéndose.

\setcounter{subsection}{108}
\subsection[Las casillas de estadísticos del panel de Monte Carlo no hacen nada]{Las cuatro casillas de estadísticos del panel de Monte Carlo no reflejan el modelo ni cambian nada}\label{err:109}
\begin{ficha}
\fichakey{Estado}Presente
\fichakey{Módulo}gui
\fichakey{Paquete}org.openmarkov.gui.dialog.inference.common
\fichakey{Ficheros}\path{gui/src/main/java/org/openmarkov/gui/dialog/inference/common/MonteCarloOptionsPanel.java:189-201} \newline \path{gui/src/main/java/org/openmarkov/gui/dialog/inference/common/MonteCarloOptionsPanel.java:343-409} \newline \path{core/src/main/java/org/openmarkov/core/inference/MonteCarloOptions.java:58-72} \newline \path{gui/src/main/java/org/openmarkov/gui/dialog/inference/common/InferenceOptionsDialog.java:1069-1071}
\fichakey{Comprobado}Leyendo el código
\fichakey{Esfuerzo}Pequeño (horas)
\fichakey{Decisión de equipo}\textcolor{decisionrojo}{\textbf{Sí.}} Si los estadísticos de la simulación deben poder elegirse (y entonces conectarlos al cálculo y al fichero) o si las cuatro casillas sobran
\fichakey{Relacionados}\hyperref[err:164]{error~164}, \hyperref[nota:111]{nota~111}, \hyperref[nota:131]{nota~131}
\end{ficha}
\paragraph{Qué pasa}
Con una red de sucesos discretos (DES ---simulación de eventos discretos---) abierta, el diálogo de opciones de inferencia (\texttt{Inference\allowbreak{}Options\allowbreak{}Dialog}) incluye \texttt{Monte\allowbreak{}Carlo\allowbreak{}Options\allowbreak{}Panel} y, al aceptar, copia sus opciones a la red (\texttt{set\allowbreak{}Monte\allowbreak{}Carlo\allowbreak{}Options(\allowbreak{}monte\allowbreak{}Carlo\allowbreak{}Options\allowbreak{}Panel.\allowbreak{}get\allowbreak{}Monte\allowbreak{}Carlo\allowbreak{}Options(\allowbreak{}))}).

El recuadro «DES Inference Statistics» tiene cuatro casillas: Mean, Sum, Trimmed Mean y Median. Lo que el usuario ve es una casilla activa («Mean») que puede desmarcar sin efecto alguno y tres que nunca puede tocar; ninguna refleja lo que dice el modelo.

Las cuatro se construyen con el valor de \texttt{monte\allowbreak{}Carlo\allowbreak{}Options.\allowbreak{}is\allowbreak{}Mean(\allowbreak{})} (copiado y pegado; debería ser \texttt{is\allowbreak{}Sum}, \texttt{is\allowbreak{}Trimmed\allowbreak{}Mean}, \texttt{is\allowbreak{}Median}) y a continuación se les fuerza un estado fijo: Mean marcada y activa; las otras tres desmarcadas y desactivadas.

\texttt{get\allowbreak{}Monte\allowbreak{}Carlo\allowbreak{}Options(\allowbreak{})} escribe después los cuatro estados en el objeto de opciones, de modo que aceptar deja siempre mean=true, trimmedMean=false, median=false y sum=false (el valor por omisión de \texttt{trimmed\allowbreak{}Mean} en \texttt{Monte\allowbreak{}Carlo\allowbreak{}Options} es \texttt{true}).

Esos cuatro indicadores no los lee nadie. Buscando \texttt{is\allowbreak{}Mean(\allowbreak{})}, \texttt{is\allowbreak{}Sum(\allowbreak{})}, \texttt{is\allowbreak{}Trimmed\allowbreak{}Mean(\allowbreak{})} e \texttt{is\allowbreak{}Median(\allowbreak{})} en todo el código de producción solo aparecen en la propia clase \texttt{Monte\allowbreak{}Carlo\allowbreak{}Options} (constructor de copia) y en este panel; los escritores y lectores de ProbModelXML (\texttt{PGMXWriter\_\allowbreak{}1\_\allowbreak{}0}, \texttt{PGMXReader\_\allowbreak{}1\_\allowbreak{}0}) solo guardan el número de simulaciones, el de series y el fichero de entrada. Por tanto pisar los valores no cambia ningún resultado: la consecuencia es un control que no hace nada, no un resultado cambiado.
\paragraph{El código}
gui/src/main/java/org/openmarkov/gui/dialog/inference/common/MonteCarloOptionsPanel.java:359-381

\begin{Shaded}
\begin{Highlighting}[]
\KeywordTok{private} \BuiltInTok{JCheckBox} \FunctionTok{getJCheckBoxMean}\OperatorTok{()} \OperatorTok{\{}
    \ControlFlowTok{if} \OperatorTok{(}\NormalTok{meanCheckBox }\OperatorTok{==} \KeywordTok{null}\OperatorTok{)} \OperatorTok{\{}
        \CommentTok{//}\AlertTok{TODO}\CommentTok{ use stringDatabase}
\NormalTok{        meanCheckBox }\OperatorTok{=} \KeywordTok{new} \BuiltInTok{JCheckBox}\OperatorTok{(}\StringTok{"Mean"}\OperatorTok{,}\NormalTok{ monteCarloOptions}\OperatorTok{.}\FunctionTok{isMean}\OperatorTok{());}
\NormalTok{        meanCheckBox}\OperatorTok{.}\FunctionTok{setEnabled}\OperatorTok{(}\KeywordTok{true}\OperatorTok{);}
\NormalTok{        meanCheckBox}\OperatorTok{.}\FunctionTok{setSelected}\OperatorTok{(}\KeywordTok{true}\OperatorTok{);}
    \OperatorTok{\}}
    \ControlFlowTok{return}\NormalTok{ meanCheckBox}\OperatorTok{;}
\OperatorTok{\}}
\KeywordTok{...}
\KeywordTok{private} \BuiltInTok{JCheckBox} \FunctionTok{getJCheckBoxSum}\OperatorTok{()} \OperatorTok{\{}
    \ControlFlowTok{if} \OperatorTok{(}\NormalTok{sumCheckBox }\OperatorTok{==} \KeywordTok{null}\OperatorTok{)} \OperatorTok{\{}
        \CommentTok{//}\AlertTok{TODO}\CommentTok{ use stringDatabase}
\NormalTok{        sumCheckBox }\OperatorTok{=} \KeywordTok{new} \BuiltInTok{JCheckBox}\OperatorTok{(}\StringTok{"Sum"}\OperatorTok{,}\NormalTok{ monteCarloOptions}\OperatorTok{.}\FunctionTok{isMean}\OperatorTok{());}
\NormalTok{        sumCheckBox}\OperatorTok{.}\FunctionTok{setEnabled}\OperatorTok{(}\KeywordTok{false}\OperatorTok{);}
\NormalTok{        sumCheckBox}\OperatorTok{.}\FunctionTok{setSelected}\OperatorTok{(}\KeywordTok{false}\OperatorTok{);}
    \OperatorTok{\}}
    \ControlFlowTok{return}\NormalTok{ sumCheckBox}\OperatorTok{;}
\OperatorTok{\}}
\end{Highlighting}
\end{Shaded}

gui/src/main/java/org/openmarkov/gui/dialog/inference/common/MonteCarloOptionsPanel.java:195-198

\begin{Shaded}
\begin{Highlighting}[]
\KeywordTok{this}\OperatorTok{.}\FunctionTok{monteCarloOptions}\OperatorTok{.}\FunctionTok{setMean}\OperatorTok{(}\NormalTok{meanCheckBox}\OperatorTok{.}\FunctionTok{isSelected}\OperatorTok{());}
\KeywordTok{this}\OperatorTok{.}\FunctionTok{monteCarloOptions}\OperatorTok{.}\FunctionTok{setTrimmedMean}\OperatorTok{(}\NormalTok{trimmedMeanCheckBox}\OperatorTok{.}\FunctionTok{isSelected}\OperatorTok{());}
\KeywordTok{this}\OperatorTok{.}\FunctionTok{monteCarloOptions}\OperatorTok{.}\FunctionTok{setMedian}\OperatorTok{(}\NormalTok{medianCheckBox}\OperatorTok{.}\FunctionTok{isSelected}\OperatorTok{());}
\KeywordTok{this}\OperatorTok{.}\FunctionTok{monteCarloOptions}\OperatorTok{.}\FunctionTok{setSum}\OperatorTok{(}\NormalTok{sumCheckBox}\OperatorTok{.}\FunctionTok{isSelected}\OperatorTok{());}
\end{Highlighting}
\end{Shaded}
\paragraph{Cómo arreglarlo}
Depende de la decisión de equipo:

\begin{itemize}
\tightlist
\item
  Si los estadísticos deben poder elegirse: inicializar cada casilla con su propio indicador (\texttt{is\allowbreak{}Mean}, \texttt{is\allowbreak{}Sum}, \texttt{is\allowbreak{}Trimmed\allowbreak{}Mean}, \texttt{is\allowbreak{}Median}) sin forzar el estado, activarlas, hacer que la simulación (\texttt{inference.\allowbreak{}des}, \texttt{DESInference} y la ventana de resultados) use los indicadores, y guardarlos y leerlos en \texttt{PGMXWriter\_\allowbreak{}1\_\allowbreak{}0}/\texttt{PGMXReader\_\allowbreak{}1\_\allowbreak{}0}.
\item
  Si no: quitar el recuadro «DES Inference Statistics» y las cuatro escrituras de \texttt{get\allowbreak{}Monte\allowbreak{}Carlo\allowbreak{}Options(\allowbreak{})}, y los cuatro campos de \texttt{Monte\allowbreak{}Carlo\allowbreak{}Options} si nadie más los necesita (hay una prueba, \texttt{Options\allowbreak{}Are\allowbreak{}Copied\allowbreak{}Whole\allowbreak{}Test}, que los copia).
\end{itemize}

Prueba para la primera opción: con \texttt{Monte\allowbreak{}Carlo\allowbreak{}Options} con sum=true y mean=false, construir el panel y comprobar que las casillas lo reflejan y que \texttt{get\allowbreak{}Monte\allowbreak{}Carlo\allowbreak{}Options(\allowbreak{})} devuelve lo mismo sin tocar nada.

\setcounter{subsection}{109}
\subsection[«Sort probnet» falla en una red de sucesos con un nodo enlazado a sí mismo]{«Sort probnet» falla con una excepción en un modelo de sucesos discretos que tiene un enlace de un nodo a sí mismo}\label{err:110}
\begin{ficha}
\fichakey{Estado}Presente
\fichakey{Módulo}full
\fichakey{Paquete}org.openmarkov.full
\fichakey{Ficheros}\path{full/src/main/java/org/openmarkov/full/SortProbnetPlugin.java:52-63} \newline \path{full/src/main/java/org/openmarkov/full/SortProbnetPlugin.java:95-120} \newline \path{core/src/main/java/org/openmarkov/core/model/network/type/DESNetworkType.java:34-39}
\fichakey{Comprobado}Leyendo el código y ejecutándolo
\fichakey{Esfuerzo}Pequeño (horas)
\fichakey{Decisión de equipo}No
\end{ficha}
\paragraph{Qué pasa}
\texttt{Sort\allowbreak{}Probnet\allowbreak{}Plugin} (módulo \texttt{full}) añade al menú Herramientas la entrada «Sort probnet», que recoloca los nodos de la red abierta. En cualquier red con un enlace de un nodo a sí mismo, «Sort probnet» sale por una excepción no controlada desde la acción del menú y no mueve ningún nodo (la edición \texttt{Move\allowbreak{}Node\allowbreak{}Edit} va después).

En las redes de simulación de eventos discretos (DES, tipo \texttt{DESNetwork\allowbreak{}Type}) un enlace de un nodo a sí mismo está permitido: el tipo desactiva \texttt{No\allowbreak{}Self\allowbreak{}Loop} y activa \texttt{Only\allowbreak{}Self\allowbreak{}Loops\allowbreak{}With\allowbreak{}Event\allowbreak{}And\allowbreak{}Chance\allowbreak{}Nodes} (líneas 34-39), es decir, se admite en nodos de suceso y de azar. Se usa, por ejemplo, para un suceso que vuelve a programarse a sí mismo (una revisión que programa la siguiente), y el editor sabe dibujarlo (\texttt{Visual\allowbreak{}Link}, \texttt{Visual\allowbreak{}Arrow.\allowbreak{}is\allowbreak{}Self\allowbreak{}Loop}).

En las redes de ejemplo del proyecto no hay ahora ninguna red DES con un enlace de un nodo a sí mismo, pero el editor permite crearlas. Para reproducirlo en la aplicación hay que construir una red de sucesos con un enlace de un suceso a sí mismo y pedir «Sort probnet».

El complemento usa el algoritmo de disposición por fuerzas de Fruchterman-Reingold de la biblioteca JGraphT. En \texttt{get\allowbreak{}Sorted\allowbreak{}Coordinates} (líneas 95-120) mete los nodos en un \texttt{Simple\allowbreak{}Directed\allowbreak{}Graph} y, por cada enlace dibujado, llama a \texttt{graph.\allowbreak{}add\allowbreak{}Edge(\allowbreak{}origen,\allowbreak{}\ destino)}. \texttt{Simple\allowbreak{}Directed\allowbreak{}Graph} es un grafo «simple»: no admite un enlace de un vértice a sí mismo y lanza \texttt{Illegal\allowbreak{}Argument\allowbreak{}Exception:\ loops\ not\ allowed}.

Medido por partes con un pequeño programa de prueba: en una red DES construida en código con dos nodos de suceso, \texttt{new\ Add\allowbreak{}Link\allowbreak{}Edit(\allowbreak{}net,\allowbreak{}\ checkup,\allowbreak{}\ checkup,\allowbreak{}\ true).\allowbreak{}execute\allowbreak{}Edit(\allowbreak{})} se acepta; y \texttt{new\ Simple\allowbreak{}Directed\allowbreak{}Graph\textless{}\textgreater{}(\allowbreak{}Default\allowbreak{}Edge.\allowbreak{}class).\allowbreak{}add\allowbreak{}Edge(\allowbreak{}v,\allowbreak{}\ v)} lanza \texttt{java.\allowbreak{}lang.\allowbreak{}Illegal\allowbreak{}Argument\allowbreak{}Exception:\ loops\ not\ allowed} con el JGraphT del jar. No se ha ejecutado el complemento completo (necesita el panel del editor con la ventana principal), pero no hay nada entre el recorrido de enlaces y \texttt{add\allowbreak{}Edge} que filtre esos enlaces.
\paragraph{El código}
full/src/main/java/org/openmarkov/full/SortProbnetPlugin.java:97-113

\begin{Shaded}
\begin{Highlighting}[]
\NormalTok{Graph}\OperatorTok{\textless{}}\NormalTok{VisualNode}\OperatorTok{,}\NormalTok{ DefaultEdge}\OperatorTok{\textgreater{}}\NormalTok{ graph }\OperatorTok{=} \KeywordTok{new}\NormalTok{ SimpleDirectedGraph}\OperatorTok{\textless{}\textgreater{}(}\NormalTok{DefaultEdge}\OperatorTok{.}\FunctionTok{class}\OperatorTok{);}
\KeywordTok{...}
\NormalTok{visualNetwork}\OperatorTok{.}\FunctionTok{getAllNodes}\OperatorTok{().}\FunctionTok{forEach}\OperatorTok{(}\NormalTok{graph}\OperatorTok{::}\NormalTok{addVertex}\OperatorTok{);}
\ControlFlowTok{for} \OperatorTok{(}\DataTypeTok{var}\NormalTok{ link }\OperatorTok{:}\NormalTok{ visualNetwork}\OperatorTok{.}\FunctionTok{getVisualLinks}\OperatorTok{())} \OperatorTok{\{}
\NormalTok{    graph}\OperatorTok{.}\FunctionTok{addEdge}\OperatorTok{(}\NormalTok{link}\OperatorTok{.}\FunctionTok{getSourceNode}\OperatorTok{(),}\NormalTok{ link}\OperatorTok{.}\FunctionTok{getDestinationNode}\OperatorTok{());}
    \ControlFlowTok{if} \OperatorTok{(!}\NormalTok{link}\OperatorTok{.}\FunctionTok{isDirected}\OperatorTok{())} \OperatorTok{\{}
\NormalTok{        graph}\OperatorTok{.}\FunctionTok{addEdge}\OperatorTok{(}\NormalTok{link}\OperatorTok{.}\FunctionTok{getDestinationNode}\OperatorTok{(),}\NormalTok{ link}\OperatorTok{.}\FunctionTok{getSourceNode}\OperatorTok{());}
    \OperatorTok{\}}
\OperatorTok{\}}
\NormalTok{layoutAlgorithm}\OperatorTok{.}\FunctionTok{layout}\OperatorTok{(}\NormalTok{graph}\OperatorTok{,}\NormalTok{ model}\OperatorTok{);}
\end{Highlighting}
\end{Shaded}
\paragraph{Cómo arreglarlo}
Saltarse en el bucle los enlaces cuyo origen y destino son el mismo nodo (no aportan nada a una disposición por fuerzas: la fuerza de un nodo sobre sí mismo no está definida). Alternativa: usar un tipo de grafo de JGraphT que admita lazos (\texttt{Directed\allowbreak{}Pseudograph}), comprobando antes que \texttt{FRLayout\allowbreak{}Algorithm2D} los tolera.

El mismo método lo usan \texttt{graphically\allowbreak{}Sort\allowbreak{}Network} y \texttt{graphically\allowbreak{}Sort\allowbreak{}Network\allowbreak{}As\allowbreak{}Move\allowbreak{}Edit}, así que el arreglo en \texttt{get\allowbreak{}Sorted\allowbreak{}Coordinates} cubre las dos.

Prueba: construir el grafo con la lista de enlaces de una red DES con un enlace de un nodo a sí mismo (conviene extraer la construcción del grafo a un método que reciba nodos y enlaces, para probarlo sin ventana) y comprobar que la disposición se calcula y da una posición a cada nodo.

% ══════════════════════════════════════════════════════════════════════
\section{Abrir y guardar redes}\label{sec:ficheros}

La lectura y escritura de redes: el formato propio ProbModelXML (.pgmx) y la puerta de entrada de los demás formatos.

\setcounter{subsection}{110}
\subsection[Un .pgmx mal escrito se abre sin aviso y con todas las tablas uniformes]{Un \texttt{.\allowbreak{}pgmx} con un elemento mal escrito se abre sin ningún aviso y con todas sus tablas convertidas en uniformes}\label{err:111}
\begin{ficha}
\fichakey{Estado}Presente
\fichakey{Módulo}io
\fichakey{Paquete}org.openmarkov.core.io.format.annotation
\fichakey{Ficheros}\path{io/src/main/java/org/openmarkov/core/io/format/annotation/FormatManager.java:119-132,226-238,240-258} \newline \path{io/src/main/java/org/openmarkov/io/probmodel/reader/PGMXReader_0_2.java:134-150} \newline \path{gui/src/main/java/org/openmarkov/gui/dialog/io/NetsIO.java:168-180}
\fichakey{Comprobado}Leyendo el código y ejecutándolo
\fichakey{Esfuerzo}Pequeño (horas)
\fichakey{Decisión de equipo}\textcolor{decisionrojo}{\textbf{Sí.}} Rechazar el fichero que no cumple el esquema o abrirlo avisando (hoy hay dos redes del repositorio que no lo cumplen)
\fichakey{Relacionados}\hyperref[err:37]{error~37}, \hyperref[err:114]{error~114}, \hyperref[err:117]{error~117}, \hyperref[nota:50]{nota~50}, \hyperref[nota:181]{nota~181}
\end{ficha}
\paragraph{Qué pasa}
Un fichero \texttt{.\allowbreak{}pgmx} (el formato propio ProbModelXML) tiene un esquema XML (lenguaje de marcado extensible), \texttt{val\_\allowbreak{}v4.\allowbreak{}xsd}, que dice qué elementos lleva y en qué orden. Al abrir un fichero, el programa debería comprobarlo contra ese esquema y, si no lo cumple, decir en qué línea está el problema.

El usuario llega con Archivo → Abrir sobre un \texttt{.\allowbreak{}pgmx} editado a mano, generado por otro programa o dañado. La red aparece normal y el usuario trabaja con números que no son los del fichero, sin ninguna señal.

Los dos caminos de apertura llaman a \texttt{Format\allowbreak{}Manager.\allowbreak{}check\allowbreak{}Structure(\allowbreak{}URL)}: la ventana, a través de \texttt{Nets\allowbreak{}IO.\allowbreak{}open\allowbreak{}Network\allowbreak{}URL}/\texttt{open\allowbreak{}Network\allowbreak{}File} → \texttt{Format\allowbreak{}Manager.\allowbreak{}get\allowbreak{}Prob\allowbreak{}Net\allowbreak{}Reader} (línea 124), y el lector, en \texttt{PGMXReader\_\allowbreak{}0\_\allowbreak{}2.\allowbreak{}read} (línea 138). Ese método tiene todo su cuerpo dentro de un comentario (líneas 240-258), así que no comprueba nada.

Existe otra variante, \texttt{check\allowbreak{}Structure(\allowbreak{}String,\allowbreak{}\ Input\allowbreak{}Stream)} (líneas 226-238), que sí valida contra el esquema y lanza \texttt{Prob\allowbreak{}Net\allowbreak{}Parser\allowbreak{}Exception.\allowbreak{}Badly\allowbreak{}Structured\allowbreak{}File} con la línea y la columna del error, pero nadie la llama. \texttt{check\allowbreak{}Version}, que sí está activo, solo mira el atributo de versión.

Como el lector busca los elementos por nombre y, si no los encuentra, sigue con valores por defecto, un fichero con la estructura mal se lee «bien» con otro contenido.

Medido con un pequeño programa de prueba: se copió \texttt{0\_\allowbreak{}miscellaneous/\allowbreak{}networks/\allowbreak{}bn/\allowbreak{}BN-asia.\allowbreak{}pgmx} cambiando \texttt{\textless{}Values\textgreater{}} por \texttt{\textless{}Valores\textgreater{}} en las tablas. Por el mismo camino que usa la ventana (\texttt{get\allowbreak{}Prob\allowbreak{}Net\allowbreak{}Reader(\allowbreak{}url).\allowbreak{}read(\allowbreak{}url)}) se abre sin error, con sus 8 nodos, y \textbf{todas} las tablas salen uniformes (por ejemplo, X-ray \texttt{{[}0.\allowbreak{}5,\allowbreak{}\ 0.\allowbreak{}5,\allowbreak{}\ 0.\allowbreak{}5,\allowbreak{}\ 0.\allowbreak{}5{]}} en lugar de \texttt{{[}0.\allowbreak{}95,\allowbreak{}\ 0.\allowbreak{}05,\allowbreak{}\ 0.\allowbreak{}02,\allowbreak{}\ 0.\allowbreak{}98{]}}). La variante con flujo, sobre el mismo fichero, lo rechaza: «An error occurred at line 142, column 18: cvc-complex-type.2.4.a \ldots{} `Valores' \ldots».

Hace falta una decisión porque, validando con la variante que funciona las 150 redes \texttt{.\allowbreak{}pgmx} de \texttt{0\_\allowbreak{}miscellaneous/\allowbreak{}networks}, 148 pasan y 2 no:

\begin{itemize}
\tightlist
\item
  \texttt{0\_\allowbreak{}miscellaneous/\allowbreak{}networks/\allowbreak{}pomdp/\allowbreak{}Dec-POMDP-wireless-network.\allowbreak{}pgmx}: elemento \texttt{Agents} en un sitio que el esquema no admite, línea 349.
\item
  \texttt{0\_\allowbreak{}miscellaneous/\allowbreak{}networks/\allowbreak{}pomdp/\allowbreak{}POMDP-coffee-robot.\allowbreak{}pgmx}: un \texttt{States} vacío, línea 627.
\end{itemize}

Las dos son POMDP (procesos de decisión de Markov parcialmente observables) y hoy se abren. Si la validación rechaza, dejarán de abrirse; si avisa, se abrirán con un aviso. Es la misma pregunta que plantea \hyperref[err:37]{error~37} para las restricciones de la red.
\paragraph{El código}
io/src/main/java/org/openmarkov/core/io/format/annotation/FormatManager.java:240-258

\begin{Shaded}
\begin{Highlighting}[]
\KeywordTok{public} \DataTypeTok{void} \FunctionTok{checkStructure}\OperatorTok{(}\BuiltInTok{URL}\NormalTok{ url}\OperatorTok{)} \KeywordTok{throws}\NormalTok{ ProbNetParserException}\OperatorTok{.}\FunctionTok{BadlyStructuredFile} \OperatorTok{\{}
    \CommentTok{/*}
\CommentTok{    InputStream xsd = getClass().getClassLoader().getResourceAsStream("val\_v4.xsd");}
\CommentTok{    SchemaFactory factory = SchemaFactory.newInstance(XMLConstants.W3C\_XML\_SCHEMA\_NS\_URI);}
\CommentTok{    ...}
\CommentTok{    try \{}
\CommentTok{        schema.newValidator().validate(new StreamSource(url.openStream()));}
\CommentTok{    \} catch (SAXException e) \{}
\CommentTok{        throw new UnrecoverableException(new ProbNetParserException.CannotParseFile(e, url));}
\CommentTok{    \} ...}
\CommentTok{    */}
\OperatorTok{\}}
\end{Highlighting}
\end{Shaded}
\paragraph{Cómo arreglarlo}
Hacer que \texttt{check\allowbreak{}Structure(\allowbreak{}URL)} valide apoyándose en la variante con flujo: abrir \texttt{url.\allowbreak{}open\allowbreak{}Stream(\allowbreak{})} y delegar en \texttt{check\allowbreak{}Structure(\allowbreak{}String,\allowbreak{}\ Input\allowbreak{}Stream)} (líneas 226-238), que ya produce \texttt{Badly\allowbreak{}Structured\allowbreak{}File} con línea y columna. El bloque comentado convertía el error de validación en \texttt{Unrecoverable\allowbreak{}Exception}, que no es lo adecuado para un error del fichero del usuario.

Hay dos llamadores y los dos quedan cubiertos con ese único cambio: \texttt{Format\allowbreak{}Manager.\allowbreak{}get\allowbreak{}Prob\allowbreak{}Net\allowbreak{}Reader} (línea 124) y \texttt{PGMXReader\_\allowbreak{}0\_\allowbreak{}2.\allowbreak{}read} (línea 138). Conviene que solo valide uno de los dos para no leer el fichero dos veces.

Según lo que decida el equipo, \texttt{Nets\allowbreak{}IO} rechaza el fichero con el mensaje o lo abre y muestra el aviso. Si se elige rechazar, hay que corregir antes las dos redes POMDP citadas.

Prueba de regresión: el fichero de prueba descrito arriba (\texttt{BN-asia.\allowbreak{}pgmx} con las tablas en \texttt{\textless{}Valores\textgreater{}}) debe dar \texttt{Badly\allowbreak{}Structured\allowbreak{}File} (o el aviso) al abrirse por \texttt{get\allowbreak{}Prob\allowbreak{}Net\allowbreak{}Reader(\allowbreak{}url).\allowbreak{}read(\allowbreak{}url)}, y las 148 redes que hoy cumplen el esquema deben seguir abriéndose.

\setcounter{subsection}{111}
\subsection[Guardar en formato 1.0 deja la red sin poder evaluarse al reabrirla]{Guardar en formato ProbModelXML 1.0, el que la ventana propone para una red nueva, convierte las utilidades en distribuciones y la red ya no se puede evaluar al reabrirla}\label{err:112}
\begin{ficha}
\fichakey{Estado}Presente
\fichakey{Módulo}io, core, gui
\fichakey{Paquete}org.openmarkov.io.probmodel.writer; org.openmarkov.io.probmodel.reader; org.openmarkov.core.model.network.potential; org.openmarkov.gui.window
\fichakey{Ficheros}\path{io/src/main/java/org/openmarkov/io/probmodel/writer/PGMXWriter_1_0.java:191-199} \newline \path{io/src/main/java/org/openmarkov/io/probmodel/reader/PGMXReader_1_0.java:126-166} \newline \path{core/src/main/java/org/openmarkov/core/model/network/potential/UnivariateDistrPotential.java:356-364} \newline \path{core/src/main/java/org/openmarkov/core/model/network/potential/plugin/PotentialUtils.java:149-156} \newline \path{io/src/main/java/org/openmarkov/io/probmodel/reader/PGMXReader.java:32-42} \newline \path{gui/src/main/java/org/openmarkov/gui/window/NetworkFileHandler.java:79-101, 126-129, 204-205, 246-280, 480-510} \newline \path{gui/src/main/java/org/openmarkov/gui/dialog/io/NetworkOMFileChooser.java:73-82, 107-124} \newline \path{gui/src/main/java/org/openmarkov/gui/configuration/UserPreferences.java:113-115}
\fichakey{Comprobado}Leyendo el código y ejecutándolo
\fichakey{Esfuerzo}Medio (días)
\fichakey{Decisión de equipo}\textcolor{decisionrojo}{\textbf{Sí.}} Dónde se resuelve (al escribir, al leer o haciendo que una distribución exacta se pueda convertir en tabla) y qué pasa con los ficheros 1.0 ya guardados, entre ellos las copias de las carpetas \texttt{1-0} del repositorio
\fichakey{Relacionados}\hyperref[err:8]{error~8}, \hyperref[err:9]{error~9}, \hyperref[err:42]{error~42}, \hyperref[err:46]{error~46}, \hyperref[err:50]{error~50}, \hyperref[err:113]{error~113}, \hyperref[err:114]{error~114}, \hyperref[err:153]{error~153}
\end{ficha}
\paragraph{Qué pasa}
Un usuario crea un diagrama de influencia en el editor, lo evalúa sin problemas y lo guarda. Al reabrirlo no puede evaluarlo: entrar en modo inferencia, pedir la estrategia óptima o lanzar el análisis de coste-efectividad terminan en una ventana de error con «Potential P(U \textbar{} \ldots) cannot be converted to a table»; si la red tiene nodos de utilidad que suman o multiplican otros, la estrategia óptima sale en la ventana de error no previsto. Guardar otra vez no lo arregla.

Qué formato usa cada gesto de guardar:

\begin{itemize}
\tightlist
\item
  Al abrir un \texttt{.\allowbreak{}pgmx}, \texttt{PGMXReader.\allowbreak{}read} (líneas 37-41) elige el escritor según la versión del fichero, y la ventana lo apunta en la pestaña de la red (\texttt{Network\allowbreak{}File\allowbreak{}Handler}, líneas 208-209). «Save» usa ese escritor (líneas 255-257): una red abierta desde un fichero 0.2 se guarda en 0.2, y una abierta desde un 1.0, en 1.0.
\item
  Una red que no tiene fichero no tiene escritor: una red nueva (línea 107), una red expandida (\texttt{Inference\allowbreak{}Handler}, línea 181), una aprendida (\texttt{Learning\allowbreak{}Dialog}, línea 904) o una abierta desde el árbol de decisión. Para ella «Save» hace «Save network as\ldots».
\item
  «Save network as\ldots» abre la ventana de ficheros con los formatos «Elvira», «OpenMarkov.0.2» y «OpenMarkov.1.0». Marca de entrada el formato del escritor de la red, si lo tiene (líneas 491-500); si no, el último que el usuario eligió al guardar, y la primera vez \texttt{PGMXWriter\_\allowbreak{}1\_\allowbreak{}0} (\texttt{User\allowbreak{}Preferences}, líneas 107-109; \texttt{Network\allowbreak{}OMFile\allowbreak{}Chooser}, líneas 76-85).
\end{itemize}

Así, llegan al formato 1.0 una red nueva (o expandida, o aprendida) guardada tal como propone la ventana, cualquier red guardada con «Save network as\ldots» eligiendo «OpenMarkov.1.0», y cualquier red abierta desde un fichero 1.0.

Por qué pasa. El editor da a todo nodo de utilidad nuevo una relación «Exact», un \texttt{Exact\allowbreak{}Distr\allowbreak{}Potential} (\texttt{Potential\allowbreak{}Utils.\allowbreak{}generate\allowbreak{}Default\allowbreak{}Potential}, línea 157), y es la relación que el diálogo del nodo llama «Exact»: una tabla con el valor del nodo para cada combinación de sus padres. Los ficheros 0.2 la guardan como tabla y el lector 0.2 la reconstruye como \texttt{Exact\allowbreak{}Distr\allowbreak{}Potential}. \texttt{PGMXWriter\_\allowbreak{}1\_\allowbreak{}0} la escribe en cambio como \texttt{\textless{}Potential\ type="Univariate\allowbreak{}Distr"\ distribution="Exact"\textgreater{}} (líneas 191-199), y \texttt{PGMXReader\_\allowbreak{}1\_\allowbreak{}0.\allowbreak{}get\allowbreak{}Univariate\allowbreak{}Distr\allowbreak{}Potential} la lee como \texttt{Univariate\allowbreak{}Distr\allowbreak{}Potential} (líneas 142-165), salvo en las redes de simulación de sucesos, donde ya la devuelve como \texttt{Exact\allowbreak{}Distr\allowbreak{}Potential} (líneas 136-141). \texttt{Univariate\allowbreak{}Distr\allowbreak{}Potential.\allowbreak{}table\allowbreak{}Project} lanza siempre \texttt{Potential\allowbreak{}Cannot\allowbreak{}Be\allowbreak{}Converted\allowbreak{}To\allowbreak{}ATable} (líneas 356-364), y todos los cálculos exactos necesitan convertir en tabla cada potencial: la evaluación, la propagación, el coste-efectividad, el cálculo de rangos del modo inferencia (\hyperref[err:46]{error~46}), «Absorb parents» (\hyperref[err:9]{error~9}) y el árbol de decisión (\hyperref[err:50]{error~50}).

Medido con un pequeño programa de prueba:

\begin{itemize}
\tightlist
\item
  Un diagrama de influencia construido con las mismas ediciones que el editor (Disease → U ← Therapy, U con relación «Exact» y valores 10, 3, 8 y 9) da una utilidad esperada de 9,02. Guardado con \texttt{PGMXWriter\_\allowbreak{}0\_\allowbreak{}2} y reabierto, sigue dando 9,02. Guardado con \texttt{PGMXWriter\_\allowbreak{}1\_\allowbreak{}0} y reabierto, U es un \texttt{Univariate\allowbreak{}Distr\allowbreak{}Potential} y la evaluación y la propagación del modo inferencia fallan con «Potential P(U \textbar{} Disease, Therapy) cannot be converted to a table». Guardado otra vez con «Save» (1.0) y reabierto, igual; el fichero lleva además \texttt{\textless{}Functions\textgreater{}"1"\ "1"\ "1"\ "1"\textless{}/\allowbreak{}Functions\textgreater{}} junto a los números, de modo que el editor de la relación enseña «1» en cada casilla (\hyperref[err:153]{error~153}).
\item
  Guardando y reabriendo cada red en formato 0.2 de \texttt{0\_\allowbreak{}miscellaneous/\allowbreak{}networks/\allowbreak{}id}, \texttt{dan}, \texttt{mid}, \texttt{limids} y \texttt{pomdp}: de las 53 cuya evaluación (\texttt{VEEvaluation}) termina, tras el formato 0.2 se siguen evaluando 44 (las otras nueve fallan por defectos distintos de este: siete, las que tienen nodos de producto, por \hyperref[err:113]{error~113}, y dos no llegan a guardarse por \hyperref[err:114]{error~114}) y tras el formato 1.0 ninguna. De las 41 en las que el cálculo del modo inferencia termina, tras el formato 1.0 ninguna. En \texttt{0\_\allowbreak{}miscellaneous/\allowbreak{}networks/\allowbreak{}id} son 14 redes evaluables que dejan de serlo; el análisis de coste-efectividad (\texttt{VECEAnalysis}) de \texttt{ID-CEA-minimal.\allowbreak{}pgmx} y de \texttt{ID-arthronet-ce.\allowbreak{}pgmx} también deja de funcionar.
\item
  Las 15 redes de \texttt{0\_\allowbreak{}miscellaneous/\allowbreak{}networks/\allowbreak{}id/\allowbreak{}1-0/\allowbreak{}} fallan al evaluarse; sus versiones 0.2 de \texttt{0\_\allowbreak{}miscellaneous/\allowbreak{}networks/\allowbreak{}id/\allowbreak{}} se evalúan, salvo \texttt{ID-delayed-result-of-test.\allowbreak{}pgmx}, que falla en los dos formatos por tener utilidades uniformes.
\item
  No hay vuelta atrás desde la ventana: guardar con «Save network as\ldots» en «OpenMarkov.0.2» una red reabierta desde un fichero 1.0 escribe el \texttt{Univariate\allowbreak{}Distr\allowbreak{}Potential} tal cual, y al abrir ese fichero el lector 0.2 lo rechaza (\texttt{PGMXParser\allowbreak{}Exception\$Potential\allowbreak{}Type\allowbreak{}Not\allowbreak{}Supported}). Queda volver a elegir «Exact» en cada nodo y teclear otra vez sus números (por lectura: el diálogo crea la relación nueva vacía).
\end{itemize}
\paragraph{El código}
io/src/main/java/org/openmarkov/io/probmodel/writer/PGMXWriter\_1\_0.java:191-199

\begin{Shaded}
\begin{Highlighting}[]
\BuiltInTok{String}\NormalTok{ potentialType }\OperatorTok{=}\NormalTok{ PotentialUtils}\OperatorTok{.}\FunctionTok{getPotentialName}\OperatorTok{(}\NormalTok{potential}\OperatorTok{.}\FunctionTok{getClass}\OperatorTok{());}
\ControlFlowTok{if} \OperatorTok{(}\NormalTok{potential}\OperatorTok{.}\FunctionTok{getClass}\OperatorTok{()} \OperatorTok{==}\NormalTok{ ExactDistrPotential}\OperatorTok{.}\FunctionTok{class}\OperatorTok{)} \OperatorTok{\{}
\NormalTok{    potentialType }\OperatorTok{=}\NormalTok{ PotentialUtils}\OperatorTok{.}\FunctionTok{getPotentialName}\OperatorTok{(}\NormalTok{UnivariateDistrPotential}\OperatorTok{.}\FunctionTok{class}\OperatorTok{);}
    \BuiltInTok{String}\NormalTok{ distribution }\OperatorTok{=} \StringTok{"Exact"}\OperatorTok{;}
\NormalTok{    potentialElement}\OperatorTok{.}\FunctionTok{setAttribute}\OperatorTok{(}\NormalTok{XMLAttributes}\OperatorTok{.}\FunctionTok{TYPE}\OperatorTok{.}\FunctionTok{toString}\OperatorTok{(),}\NormalTok{ potentialType}\OperatorTok{);}
\NormalTok{    potentialElement}\OperatorTok{.}\FunctionTok{setAttribute}\OperatorTok{(}\NormalTok{XMLAttributes}\OperatorTok{.}\FunctionTok{DISTRIBUTION}\OperatorTok{.}\FunctionTok{toString}\OperatorTok{(),}\NormalTok{ distribution}\OperatorTok{);}
\OperatorTok{\}} \ControlFlowTok{else} \OperatorTok{\{}
\NormalTok{    potentialElement}\OperatorTok{.}\FunctionTok{setAttribute}\OperatorTok{(}\NormalTok{XMLAttributes}\OperatorTok{.}\FunctionTok{TYPE}\OperatorTok{.}\FunctionTok{toString}\OperatorTok{(),}\NormalTok{ potentialType}\OperatorTok{);}
\OperatorTok{\}}
\end{Highlighting}
\end{Shaded}

io/src/main/java/org/openmarkov/io/probmodel/reader/PGMXReader\_1\_0.java:136-144

\begin{Shaded}
\begin{Highlighting}[]
\ControlFlowTok{if} \OperatorTok{(}\NormalTok{probNet}\OperatorTok{.}\FunctionTok{getNetworkType}\OperatorTok{()} \KeywordTok{instanceof}\NormalTok{ DESNetworkType}\OperatorTok{)\{}
    \ControlFlowTok{if} \OperatorTok{(}\NormalTok{Objects}\OperatorTok{.}\FunctionTok{requireNonNull}\OperatorTok{(}\NormalTok{probNet}\OperatorTok{.}\FunctionTok{getNode}\OperatorTok{(}\NormalTok{variables}\OperatorTok{.}\FunctionTok{getFirst}\OperatorTok{())).}\FunctionTok{getNodeType}\OperatorTok{()} \OperatorTok{==}\NormalTok{ NodeType}\OperatorTok{.}\FunctionTok{CHANCE}\OperatorTok{)\{}
        \ControlFlowTok{return}\NormalTok{ PGMXPotentialParsers}\OperatorTok{.}\FunctionTok{getTablePotential}\OperatorTok{(}\NormalTok{xmlPotential}\OperatorTok{,}\NormalTok{probNet}\OperatorTok{,}\NormalTok{xmlRole}\OperatorTok{,}\NormalTok{variables}\OperatorTok{);}
    \OperatorTok{\}}
    \ControlFlowTok{return}\NormalTok{ PGMXPotentialParsers}\OperatorTok{.}\FunctionTok{getExactDistrPotential}\OperatorTok{(}\NormalTok{xmlPotential}\OperatorTok{,}\NormalTok{probNet}\OperatorTok{,}\NormalTok{xmlRole}\OperatorTok{,}\NormalTok{variables}\OperatorTok{);}
\OperatorTok{\}}
\NormalTok{UnivariateDistrPotential potential}\OperatorTok{;}
\ControlFlowTok{if} \OperatorTok{(}\NormalTok{parametrization }\OperatorTok{!=} \KeywordTok{null}\OperatorTok{)} \OperatorTok{\{}
\NormalTok{    potential }\OperatorTok{=} \KeywordTok{new} \FunctionTok{UnivariateDistrPotential}\OperatorTok{(}\NormalTok{variables}\OperatorTok{,}\NormalTok{ univariateName}\OperatorTok{,}\NormalTok{ parametrization}\OperatorTok{,}\NormalTok{ xmlRole}\OperatorTok{);}
\end{Highlighting}
\end{Shaded}

core/src/main/java/org/openmarkov/core/model/network/potential/UnivariateDistrPotential.java:356-359

\begin{Shaded}
\begin{Highlighting}[]
\AttributeTok{@Override}
\KeywordTok{public} \AttributeTok{@NotNull}\NormalTok{ TablePotential }\FunctionTok{tableProject}\OperatorTok{(}\NormalTok{EvidenceCase evidenceCase}\OperatorTok{,}\NormalTok{ InferenceOptions inferenceOptions}\OperatorTok{)} \KeywordTok{throws}\NormalTok{ NonProjectablePotentialException}\OperatorTok{.}\FunctionTok{PotentialCannotBeConvertedToATable} \OperatorTok{\{}
    \ControlFlowTok{throw} \KeywordTok{new}\NormalTok{ NonProjectablePotentialException}\OperatorTok{.}\FunctionTok{PotentialCannotBeConvertedToATable}\OperatorTok{(}\KeywordTok{this}\OperatorTok{);}
\OperatorTok{\}}
\end{Highlighting}
\end{Shaded}

gui/src/main/java/org/openmarkov/gui/configuration/UserPreferences.java:107-109

\begin{Shaded}
\begin{Highlighting}[]
\KeywordTok{public} \DataTypeTok{static} \DataTypeTok{final}\NormalTok{ UserPreference}\OperatorTok{\textless{}}\BuiltInTok{Class}\OperatorTok{\textless{}?} \KeywordTok{extends}\NormalTok{ ProbNetWriter}\OperatorTok{\textgreater{}\textgreater{}}\NormalTok{ LATEST\_SAVED\_NETWORK\_WRITER\_CLASS }\OperatorTok{=}\NormalTok{ UserPreference}
        \OperatorTok{.}\FunctionTok{of}\OperatorTok{(}\StringTok{"formats/latest\_saved\_network\_writer"}\OperatorTok{,} \OperatorTok{()} \OperatorTok{{-}\textgreater{}}\NormalTok{ PGMXWriter\_1\_0}\OperatorTok{.}\FunctionTok{class}\OperatorTok{,} \KeywordTok{null}\OperatorTok{,} \KeywordTok{new}\NormalTok{ TypeToken}\OperatorTok{\textless{}\textgreater{}()} \OperatorTok{\{}
        \OperatorTok{\});}
\end{Highlighting}
\end{Shaded}
\paragraph{Cómo arreglarlo}
Una relación «Exact» tiene que volver de un fichero 1.0 como algo que los cálculos exactos puedan convertir en tabla. Hay tres sitios posibles, y elegir entre ellos es una decisión:

\begin{itemize}
\tightlist
\item
  Al escribir: que \texttt{PGMXWriter\_\allowbreak{}1\_\allowbreak{}0} guarde un \texttt{Exact\allowbreak{}Distr\allowbreak{}Potential} de forma que el lector 1.0 lo reconstruya como \texttt{Exact\allowbreak{}Distr\allowbreak{}Potential}, como ya ocurre entre el escritor y el lector 0.2. No arregla los ficheros 1.0 ya guardados, entre ellos las copias de \texttt{0\_\allowbreak{}miscellaneous/\allowbreak{}networks/\allowbreak{}*/\allowbreak{}1-0/\allowbreak{}}.
\item
  Al leer: que \texttt{PGMXReader\_\allowbreak{}1\_\allowbreak{}0.\allowbreak{}get\allowbreak{}Univariate\allowbreak{}Distr\allowbreak{}Potential} devuelva un \texttt{Exact\allowbreak{}Distr\allowbreak{}Potential} cuando la distribución es «Exact» y sus parámetros son números (sin fórmulas en \texttt{\textless{}Functions\textgreater{}} que dependan de otras variables), como ya hace para las redes de sucesos (líneas 136-141). Arregla también los ficheros existentes; el editor enseñaría «Exact» en lugar de «UnivariateDistr».
\item
  En el potencial: que \texttt{Univariate\allowbreak{}Distr\allowbreak{}Potential.\allowbreak{}table\allowbreak{}Project} (las dos versiones, líneas 356-364) construya la tabla cuando la distribución es «Exact» con parámetros numéricos, igual que \texttt{Exact\allowbreak{}Distr\allowbreak{}Potential}. Arregla los ficheros existentes y todos los llamadores a la vez, sin cambiar la relación que ve el usuario.
\end{itemize}

Con cualquiera de las tres, lo que hoy falla por esta causa en otras fichas deja de fallar por ella: los rangos del modo inferencia de las copias 1.0 (\hyperref[err:46]{error~46}), «Absorb parents» en sus nodos de suma y producto (\hyperref[err:9]{error~9}; las versiones 0.2 de esas redes, con \texttt{Exact\allowbreak{}Distr\allowbreak{}Potential}, ya se absorben) y la proyección de \texttt{DANOperations.\allowbreak{}project\allowbreak{}Potentials} del árbol de decisión (\hyperref[err:50]{error~50}, por lectura: pasaría por su rama de \texttt{table\allowbreak{}Project}). Conviene comprobar esas tres fichas después del arreglo.

Aparte de dónde se arregle: \texttt{PGMXWriter\_\allowbreak{}0\_\allowbreak{}2} no debería escribir un tipo de potencial que su propio lector rechaza.

Prueba de regresión: para cada red de \texttt{0\_\allowbreak{}miscellaneous/\allowbreak{}networks/\allowbreak{}id/\allowbreak{}*.\allowbreak{}pgmx} en formato 0.2, guardar con \texttt{PGMXWriter\_\allowbreak{}1\_\allowbreak{}0}, reabrir con \texttt{PGMXReader} y comprobar que \texttt{VEEvaluation.\allowbreak{}get\allowbreak{}Utility(\allowbreak{})} da lo mismo que antes de guardar; lo mismo con un diagrama de influencia creado con \texttt{Add\allowbreak{}Node\allowbreak{}Edit} y \texttt{Add\allowbreak{}Link\allowbreak{}Edit} y un nodo de utilidad «Exact». Si se arregla al leer o en el potencial, además, \texttt{0\_\allowbreak{}miscellaneous/\allowbreak{}networks/\allowbreak{}id/\allowbreak{}1-0/\allowbreak{}ID-decide-test-1-0.\allowbreak{}pgmx} debe dar la misma utilidad esperada que \texttt{0\_\allowbreak{}miscellaneous/\allowbreak{}networks/\allowbreak{}id/\allowbreak{}ID-decide-test.\allowbreak{}pgmx} (9,3289).

\setcounter{subsection}{112}
\subsection[Guardar en formato 0.2 convierte los nodos de producto en uniformes]{Guardar en formato ProbModelXML 0.2, el que usa «Save» con una red abierta desde un fichero 0.2, convierte en uniforme todo nodo de utilidad de producto, y al reabrirla la red ya no se evalúa}\label{err:113}
\begin{ficha}
\fichakey{Estado}Presente
\fichakey{Módulo}io
\fichakey{Paquete}org.openmarkov.io.probmodel.writer; org.openmarkov.io.probmodel.reader
\fichakey{Ficheros}\path{io/src/main/java/org/openmarkov/io/probmodel/writer/PGMXWriter_0_2.java:770-798,833-851} \newline \path{core/src/main/java/org/openmarkov/core/model/network/potential/Potential.java:726-736} \newline \path{core/src/main/java/org/openmarkov/core/model/network/potential/SumPotential.java:141-144} \newline \path{core/src/main/java/org/openmarkov/core/model/network/potential/ProductPotential.java:68-77} \newline \path{io/src/main/java/org/openmarkov/io/probmodel/reader/PGMXReader_0_2.java:1397-1409} \newline \path{io/src/main/java/org/openmarkov/io/probmodel/reader/PGMXReader.java:32-41} \newline \path{gui/src/main/java/org/openmarkov/gui/window/NetworkFileHandler.java:246-254} \newline \path{core/src/main/java/org/openmarkov/core/inference/BasicOperations.java:42-60}
\fichakey{Comprobado}Leyendo el código y ejecutándolo
\fichakey{Esfuerzo}Pequeño (horas)
\fichakey{Decisión de equipo}No
\fichakey{Relacionados}\hyperref[err:8]{error~8}, \hyperref[err:42]{error~42}, \hyperref[err:46]{error~46}, \hyperref[err:63]{error~63}, \hyperref[err:112]{error~112}, \hyperref[err:114]{error~114}
\end{ficha}
\paragraph{Qué pasa}
Un nodo de utilidad de producto multiplica las utilidades de sus padres: es la relación «Product» del diálogo del nodo, que se ofrece a los nodos de utilidad cuyos padres son otros nodos de utilidad (\texttt{Product\allowbreak{}Potential.\allowbreak{}validate}, líneas 68-77). Siete redes de \texttt{0\_\allowbreak{}miscellaneous/\allowbreak{}networks} en formato 0.2 los tienen: \texttt{id/\allowbreak{}ID-arthronet.\allowbreak{}pgmx}, \texttt{id/\allowbreak{}ID-mediastinet.\allowbreak{}pgmx}, \texttt{dan/\allowbreak{}DAN-arthronet.\allowbreak{}pgmx}, \texttt{dan/\allowbreak{}DAN-economic-mediastinet.\allowbreak{}pgmx}, \texttt{dan/\allowbreak{}DAN-mediastinet.\allowbreak{}pgmx}, \texttt{dan/\allowbreak{}DAN-qale-mediastinet.\allowbreak{}pgmx} y \texttt{mid/\allowbreak{}MID-HPV.\allowbreak{}pgmx}.

Camino del usuario: abrir \texttt{0\_\allowbreak{}miscellaneous/\allowbreak{}networks/\allowbreak{}id/\allowbreak{}ID-arthronet.\allowbreak{}pgmx}; en modo edición, «Show optimal strategy» enseña la estrategia óptima. Hacer cualquier cambio («Save» solo se enciende con la red modificada), pulsar «Save», cerrar la red y volver a abrir el fichero. Ahora «Show optimal strategy» termina en la ventana de error no previsto con \texttt{Class\allowbreak{}Cast\allowbreak{}Exception:\ Uniform\allowbreak{}Potential\ cannot\ be\ cast\ to\ Function\allowbreak{}Potential}, y la relación del nodo «Coste ajustado» aparece como «Uniform». Guardar otra vez no lo arregla: el fichero ya no dice que el nodo era un producto.

«Save» usa el formato del fichero abierto: \texttt{PGMXReader.\allowbreak{}read} (líneas 37-41) devuelve el escritor de su versión y \texttt{Network\allowbreak{}File\allowbreak{}Handler.\allowbreak{}save\allowbreak{}Network} (líneas 250-258) lo usa, así que toda red abierta desde un fichero 0.2 se guarda con \texttt{PGMXWriter\_\allowbreak{}0\_\allowbreak{}2}. Llega también «Save network as\ldots» eligiendo el formato «OpenMarkov.0.2», con cualquier red. El formato 1.0 escribe bien el producto; su problema con las utilidades es otro (\hyperref[err:112]{error~112}).

Por qué pasa. \texttt{PGMXWriter\_\allowbreak{}0\_\allowbreak{}2.\allowbreak{}get\allowbreak{}Potential\allowbreak{}Attributes\allowbreak{}And\allowbreak{}Variables} (módulo \texttt{io}, líneas 770-798) trata juntos los potenciales de suma y de producto de un nodo de utilidad: escribe la variable del nodo como \texttt{\textless{}Utility\allowbreak{}Variable\textgreater{}} y, para quitarla de la lista, llama a \texttt{potential.\allowbreak{}remove\allowbreak{}Variable(\allowbreak{}utility\allowbreak{}Variable)} y sigue con lo que devuelve. \texttt{Sum\allowbreak{}Potential} redefine \texttt{remove\allowbreak{}Variable} y devuelve el mismo potencial; \texttt{Product\allowbreak{}Potential} no, y la versión de \texttt{Potential} (líneas 726-736) crea un \texttt{Uniform\allowbreak{}Potential} con las variables que quedan. El tipo que se escribe sale de la clase de ese objeto: \texttt{type="Uniform"}.

En el mismo método hay un segundo fallo, que afecta a sumas y productos. Después, \texttt{get\allowbreak{}Potential\allowbreak{}Variables} (líneas 833-851) mira si la primera variable que queda es de un nodo de utilidad y, si lo es, la escribe como otra \texttt{\textless{}Utility\allowbreak{}Variable\textgreater{}} en lugar de dentro de \texttt{\textless{}Variables\textgreater{}}. En un nodo de suma o de producto esa variable es su primer padre, que es otro nodo de utilidad. El lector (\texttt{PGMXReader\_\allowbreak{}0\_\allowbreak{}2.\allowbreak{}get\allowbreak{}Potential}, líneas 1403-1409) solo lee la primera \texttt{\textless{}Utility\allowbreak{}Variable\textgreater{}}, así que ese padre desaparece de las variables del potencial al reabrir.

Al calcular, el preproceso de los algoritmos exactos absorbe el nodo uniforme como si fuera una fórmula y falla con el \texttt{Class\allowbreak{}Cast\allowbreak{}Exception} que describe \hyperref[err:8]{error~8} (\texttt{Basic\allowbreak{}Operations.\allowbreak{}absorb\allowbreak{}Parent\allowbreak{}Potentials}, línea 54). Arreglado ese error, el nodo seguiría sin ser un producto.

Medido con un pequeño programa de prueba que guarda cada red con el escritor que usaría «Save», la reabre y repite los cálculos de «Show optimal strategy» (\texttt{VEOptimal\allowbreak{}Intervention} en diagramas de influencia y en MID ---diagramas de influencia con memoria---, \texttt{DANDecomposition\allowbreak{}Into\allowbreak{}Symmetric\allowbreak{}DANs\allowbreak{}Evaluation} en las DAN ---redes de análisis de decisiones---) y, en \texttt{MID-HPV.\allowbreak{}pgmx}, el análisis de coste-efectividad global:

\begin{itemize}
\tightlist
\item
  \texttt{ID-arthronet.\allowbreak{}pgmx}: utilidad esperada 0,4961 antes de guardar; tras reabrir, \texttt{Class\allowbreak{}Cast\allowbreak{}Exception}. «Coste ajustado» pasa de producto de «Coste total» y «C2E» a uniforme sobre «C2E».
\item
  \texttt{ID-mediastinet.\allowbreak{}pgmx}: 3,9844 → \texttt{Class\allowbreak{}Cast\allowbreak{}Exception}; se estropean «Net\_QALE» y «Weighted\_Economic\_Cost».
\item
  \texttt{DAN-economic-mediastinet.\allowbreak{}pgmx} (−0,1223), \texttt{DAN-mediastinet.\allowbreak{}pgmx} (1,4710) y \texttt{DAN-qale-mediastinet.\allowbreak{}pgmx} (2,1154): → \texttt{Class\allowbreak{}Cast\allowbreak{}Exception}.
\item
  \texttt{MID-HPV.\allowbreak{}pgmx}: utilidad 3480,73 y coste-efectividad global con tres tramos → \texttt{Class\allowbreak{}Cast\allowbreak{}Exception} en los dos; se estropea «QoL {[}0{]}».
\item
  \texttt{DAN-arthronet.\allowbreak{}pgmx}: su estrategia ya falla antes de guardar (\hyperref[err:63]{error~63}); tras guardar, «Coste ajustado» queda también uniforme.
\end{itemize}

Para ver qué exige el arreglo se corrigieron a mano los ficheros guardados. Cambiando solo \texttt{type="Uniform"} por \texttt{type="Product"}, los dos diagramas de influencia y las tres DAN vuelven a dar sus números, pero \texttt{MID-HPV.\allowbreak{}pgmx} da otros sin avisar: utilidad 3563,98 y efectividad 143,5 en cada tramo, en vez de 3480,73 y 60,16, porque a «QoL {[}0{]}» le falta el factor «HSR QoL {[}0{]}». Pasando además el primer padre a \texttt{\textless{}Variables\textgreater{}}, las seis dan lo mismo que antes de guardar. En las sumas, la segunda \texttt{\textless{}Utility\allowbreak{}Variable\textgreater{}} no ha cambiado ningún resultado medido (la absorción toma los sumandos de los enlaces del nodo), pero el potencial reabierto tampoco tiene todas sus variables.
\paragraph{El código}
io/src/main/java/org/openmarkov/io/probmodel/writer/PGMXWriter\_0\_2.java:777-789

\begin{Shaded}
\begin{Highlighting}[]
\OperatorTok{\}} \ControlFlowTok{else} \ControlFlowTok{if} \OperatorTok{(}\NormalTok{potential }\KeywordTok{instanceof}\NormalTok{ SumPotential }\OperatorTok{||}\NormalTok{ potential }\KeywordTok{instanceof}\NormalTok{ ProductPotential}\OperatorTok{)} \OperatorTok{\{}
    \ControlFlowTok{if} \OperatorTok{(}\NormalTok{potential}\OperatorTok{.}\FunctionTok{getPotentialRole}\OperatorTok{()} \OperatorTok{!=}\NormalTok{ PotentialRole}\OperatorTok{.}\FunctionTok{CONDITIONAL\_PROBABILITY}\OperatorTok{)} \OperatorTok{\{}
\NormalTok{        Variable utilityVariable }\OperatorTok{=}\NormalTok{ potential}\OperatorTok{.}\FunctionTok{getVariable}\OperatorTok{(}\DecValTok{0}\OperatorTok{);} \CommentTok{// it could be null in Branches potentials}
\NormalTok{        potential }\OperatorTok{=}\NormalTok{ potential}\OperatorTok{.}\FunctionTok{removeVariable}\OperatorTok{(}\NormalTok{utilityVariable}\OperatorTok{);}
        \FunctionTok{getUtilityElement}\OperatorTok{(}\NormalTok{potentialElement}\OperatorTok{,}\NormalTok{ utilityVariable}\OperatorTok{);}
    \OperatorTok{\}}
\OperatorTok{\}}

\BuiltInTok{String}\NormalTok{ potentialType }\OperatorTok{=}\NormalTok{ PotentialUtils}\OperatorTok{.}\FunctionTok{getPotentialName}\OperatorTok{(}\NormalTok{potential}\OperatorTok{.}\FunctionTok{getClass}\OperatorTok{());}
\KeywordTok{...}
\NormalTok{potentialElement}\OperatorTok{.}\FunctionTok{setAttribute}\OperatorTok{(}\NormalTok{XMLAttributes}\OperatorTok{.}\FunctionTok{TYPE}\OperatorTok{.}\FunctionTok{toString}\OperatorTok{(),}\NormalTok{ potentialType}\OperatorTok{);}
\end{Highlighting}
\end{Shaded}

io/src/main/java/org/openmarkov/io/probmodel/writer/PGMXWriter\_0\_2.java:833-845

\begin{Shaded}
\begin{Highlighting}[]
\KeywordTok{protected} \DataTypeTok{static} \DataTypeTok{void} \FunctionTok{getPotentialVariables}\OperatorTok{(}\NormalTok{Potential potential}\OperatorTok{,} \BuiltInTok{Element}\NormalTok{ potentialElement}\OperatorTok{,}\NormalTok{ ProbNet probNet}\OperatorTok{)} \OperatorTok{\{}
    \BuiltInTok{List}\OperatorTok{\textless{}}\NormalTok{Variable}\OperatorTok{\textgreater{}}\NormalTok{ potentialVariables }\OperatorTok{=}\NormalTok{ potential}\OperatorTok{.}\FunctionTok{getVariables}\OperatorTok{();}
    \ControlFlowTok{if} \OperatorTok{(!}\NormalTok{potentialVariables}\OperatorTok{.}\FunctionTok{isEmpty}\OperatorTok{()}
            \OperatorTok{\&\&}\NormalTok{ probNet}\OperatorTok{.}\FunctionTok{getNode}\OperatorTok{(}\NormalTok{potentialVariables}\OperatorTok{.}\FunctionTok{get}\OperatorTok{(}\DecValTok{0}\OperatorTok{)).}\FunctionTok{getNodeType}\OperatorTok{()} \OperatorTok{==}\NormalTok{ NodeType}\OperatorTok{.}\FunctionTok{UTILITY}\OperatorTok{)} \OperatorTok{\{}
\NormalTok{        Variable utilityVariable }\OperatorTok{=}\NormalTok{ potential}\OperatorTok{.}\FunctionTok{getVariable}\OperatorTok{(}\DecValTok{0}\OperatorTok{);}
        \KeywordTok{...}
            \BuiltInTok{Element}\NormalTok{ utilityElement }\OperatorTok{=} \KeywordTok{new} \BuiltInTok{Element}\OperatorTok{(}\NormalTok{XMLTags}\OperatorTok{.}\FunctionTok{UTILITY\_VARIABLE}\OperatorTok{.}\FunctionTok{toString}\OperatorTok{());}
            \FunctionTok{writeVariableName}\OperatorTok{(}\NormalTok{utilityVariable}\OperatorTok{,}\NormalTok{ utilityElement}\OperatorTok{);}
\NormalTok{            potentialElement}\OperatorTok{.}\FunctionTok{addContent}\OperatorTok{(}\NormalTok{utilityElement}\OperatorTok{);}
        \KeywordTok{...}
\NormalTok{        potentialVariables}\OperatorTok{.}\FunctionTok{remove}\OperatorTok{(}\NormalTok{utilityVariable}\OperatorTok{);}
\end{Highlighting}
\end{Shaded}

core/src/main/java/org/openmarkov/core/model/network/potential/Potential.java:726-736

\begin{Shaded}
\begin{Highlighting}[]
\KeywordTok{public}\NormalTok{ Potential }\FunctionTok{removeVariable}\OperatorTok{(}\NormalTok{Variable variable}\OperatorTok{)} \OperatorTok{\{}
\NormalTok{    Potential newPotential}\OperatorTok{;}
    \ControlFlowTok{if} \OperatorTok{(}\NormalTok{variables}\OperatorTok{.}\FunctionTok{contains}\OperatorTok{(}\NormalTok{variable}\OperatorTok{))} \OperatorTok{\{}
        \BuiltInTok{List}\OperatorTok{\textless{}}\NormalTok{Variable}\OperatorTok{\textgreater{}}\NormalTok{ newVariables }\OperatorTok{=} \KeywordTok{new} \BuiltInTok{ArrayList}\OperatorTok{\textless{}\textgreater{}(}\NormalTok{variables}\OperatorTok{);}
\NormalTok{        newVariables}\OperatorTok{.}\FunctionTok{remove}\OperatorTok{(}\NormalTok{variable}\OperatorTok{);}
\NormalTok{        newPotential }\OperatorTok{=} \KeywordTok{new} \FunctionTok{UniformPotential}\OperatorTok{(}\NormalTok{newVariables}\OperatorTok{,}\NormalTok{ role}\OperatorTok{);}
    \OperatorTok{\}} \ControlFlowTok{else} \OperatorTok{\{}
\NormalTok{        newPotential }\OperatorTok{=} \KeywordTok{this}\OperatorTok{;}
    \OperatorTok{\}}
    \ControlFlowTok{return}\NormalTok{ newPotential}\OperatorTok{;}
\OperatorTok{\}}
\end{Highlighting}
\end{Shaded}

0\_miscellaneous/networks/id/ID-arthronet.pgmx:734-740 («Coste ajustado» en el fichero original)

\begin{Shaded}
\begin{Highlighting}[]
\NormalTok{\textless{}}\KeywordTok{Potential}\OtherTok{ type=}\StringTok{"Product"}\OtherTok{ role=}\StringTok{"utility"}\NormalTok{\textgreater{}}
\NormalTok{  \textless{}}\KeywordTok{UtilityVariable}\OtherTok{ name=}\StringTok{"Coste ajustado"}\NormalTok{ /\textgreater{}}
\NormalTok{  \textless{}}\KeywordTok{Variables}\NormalTok{\textgreater{}}
\NormalTok{    \textless{}}\KeywordTok{Variable}\OtherTok{ name=}\StringTok{"Coste total"}\NormalTok{ /\textgreater{}}
\NormalTok{    \textless{}}\KeywordTok{Variable}\OtherTok{ name=}\StringTok{"C2E"}\NormalTok{ /\textgreater{}}
\NormalTok{  \textless{}/}\KeywordTok{Variables}\NormalTok{\textgreater{}}
\NormalTok{\textless{}/}\KeywordTok{Potential}\NormalTok{\textgreater{}}
\end{Highlighting}
\end{Shaded}

El mismo nodo en el fichero que escribe «Save»:

\begin{Shaded}
\begin{Highlighting}[]
\NormalTok{\textless{}}\KeywordTok{Potential}\OtherTok{ type=}\StringTok{"Uniform"}\OtherTok{ role=}\StringTok{"utility"}\NormalTok{\textgreater{}}
\NormalTok{  \textless{}}\KeywordTok{UtilityVariable}\OtherTok{ name=}\StringTok{"Coste ajustado"}\NormalTok{ /\textgreater{}}
\NormalTok{  \textless{}}\KeywordTok{UtilityVariable}\OtherTok{ name=}\StringTok{"Coste total"}\NormalTok{ /\textgreater{}}
\NormalTok{  \textless{}}\KeywordTok{Variables}\NormalTok{\textgreater{}}
\NormalTok{    \textless{}}\KeywordTok{Variable}\OtherTok{ name=}\StringTok{"C2E"}\NormalTok{ /\textgreater{}}
\NormalTok{  \textless{}/}\KeywordTok{Variables}\NormalTok{\textgreater{}}
\NormalTok{\textless{}/}\KeywordTok{Potential}\NormalTok{\textgreater{}}
\end{Highlighting}
\end{Shaded}
\paragraph{Cómo arreglarlo}
En \texttt{PGMXWriter\_\allowbreak{}0\_\allowbreak{}2.\allowbreak{}get\allowbreak{}Potential\allowbreak{}Attributes\allowbreak{}And\allowbreak{}Variables}, para la suma o el producto de un nodo de utilidad: escribir la variable del nodo una sola vez como \texttt{\textless{}Utility\allowbreak{}Variable\textgreater{}}, el tipo de la clase original (\texttt{Sum} o \texttt{Product}) y todos los padres dentro de \texttt{\textless{}Variables\textgreater{}}. Sin llamar a \texttt{remove\allowbreak{}Variable}, que en \texttt{Potential} devuelve un \texttt{Uniform\allowbreak{}Potential} y en \texttt{Sum\allowbreak{}Potential} quita la variable del potencial de la propia red (por lectura), y sin que \texttt{get\allowbreak{}Potential\allowbreak{}Variables} vuelva a sacar el primer padre; por ejemplo, trabajando con una copia de la lista de variables sin la del nodo.

Hacen falta las dos partes: con solo el tipo, \texttt{MID-HPV.\allowbreak{}pgmx} da otros números sin avisar (medido arriba).

Otros sitios: las ramas de un árbol se escriben con el mismo \texttt{get\allowbreak{}Potential} (\texttt{get\allowbreak{}Tree\allowbreak{}ADDBranch}), así que el cambio les vale también; \texttt{PGMXWriter\_\allowbreak{}1\_\allowbreak{}0} redefine este método y no tiene el problema. Los ficheros ya guardados no conservan que el nodo era un producto, y habrá que volver a elegir «Product» en su diálogo. Si se quiere que el lector recupere al menos el padre perdido de esos ficheros, \texttt{PGMXReader\_\allowbreak{}0\_\allowbreak{}2.\allowbreak{}get\allowbreak{}Potential} puede tomar una segunda \texttt{\textless{}Utility\allowbreak{}Variable\textgreater{}} como una variable más del potencial.

Prueba de regresión: para las siete redes citadas, guardar con \texttt{PGMXWriter\_\allowbreak{}0\_\allowbreak{}2}, reabrir y comprobar que cada nodo de utilidad conserva su clase y sus variables, y que la utilidad esperada (y, en \texttt{0\_\allowbreak{}miscellaneous/\allowbreak{}networks/\allowbreak{}mid/\allowbreak{}MID-HPV.\allowbreak{}pgmx}, el análisis de coste-efectividad global) es la misma que antes de guardar.

\setcounter{subsection}{113}
\subsection[«Save» falla en tres redes con un papel de potencial mal escrito]{«Save» falla con NullPointerException en ID-CEA-test-2therapies-new-test, MID-CHAP-Ryan-Griffin y MID-Cochlear: sus ficheros traen un papel de potencial mal escrito que el lector acepta sin avisar}\label{err:114}
\begin{ficha}
\fichakey{Estado}Presente
\fichakey{Módulo}io
\fichakey{Paquete}org.openmarkov.io.probmodel.reader; org.openmarkov.io.probmodel.writer
\fichakey{Ficheros}\path{io/src/main/java/org/openmarkov/io/probmodel/reader/PGMXReader_0_2.java:1375-1387,1563-1571} \newline \path{io/src/main/java/org/openmarkov/io/probmodel/writer/PGMXWriter_0_2.java:70-97,733-753} \newline \path{core/src/main/java/org/openmarkov/core/model/network/potential/PotentialRole.java:31-42} \newline \path{gui/src/main/java/org/openmarkov/gui/window/NetworkFileHandler.java:246-254, 308-324, 368-388} \newline \path{0_miscellaneous/networks/id/ID-CEA-test-2therapies-new-test.pgmx:148} \newline \path{0_miscellaneous/networks/mid/MID-CHAP-Ryan-Griffin.pgmx:808} \newline \path{0_miscellaneous/networks/mid/MID-Cochlear.pgmx:4854} \newline \path{integrationTests/src/main/resources/networks/mid/MID-Cochlear.pgmx}
\fichakey{Comprobado}Leyendo el código y ejecutándolo
\fichakey{Esfuerzo}Pequeño (horas)
\fichakey{Decisión de equipo}\textcolor{decisionrojo}{\textbf{Sí.}} Si el lector rechaza un papel de potencial desconocido (corrigiendo antes los ficheros que lo traen) o lo acepta como sinónimo avisando
\fichakey{Relacionados}\hyperref[err:71]{error~71}, \hyperref[err:111]{error~111}, \hyperref[err:112]{error~112}, \hyperref[err:113]{error~113}
\end{ficha}
\paragraph{Qué pasa}
Abrir \texttt{0\_\allowbreak{}miscellaneous/\allowbreak{}networks/\allowbreak{}id/\allowbreak{}ID-CEA-test-2therapies-new-test.\allowbreak{}pgmx}, \texttt{0\_\allowbreak{}miscellaneous/\allowbreak{}networks/\allowbreak{}mid/\allowbreak{}MID-CHAP-Ryan-Griffin.\allowbreak{}pgmx} o \texttt{0\_\allowbreak{}miscellaneous/\allowbreak{}networks/\allowbreak{}mid/\allowbreak{}MID-Cochlear.\allowbreak{}pgmx}, cambiar algo y pulsar «Save» termina en la ventana de error no previsto con \texttt{Null\allowbreak{}Pointer\allowbreak{}Exception:\ Cannot\ invoke\ "Potential\allowbreak{}Role.\allowbreak{}to\allowbreak{}String(\allowbreak{})"\ because\ "potential\allowbreak{}Role"\ is\ null}. La red no se guarda y sigue marcada como modificada. Al cerrarla contestando que sí a guardar, falla igual y la pestaña no se cierra (por lectura: \texttt{Network\allowbreak{}File\allowbreak{}Handler.\allowbreak{}network\allowbreak{}Can\allowbreak{}Be\allowbreak{}Closed}, líneas 372-392). «Save network as\ldots» con el formato «OpenMarkov.0.2» falla también; con «OpenMarkov.1.0» guarda, pero las utilidades quedan como describe \hyperref[err:112]{error~112}.

El fichero no se estropea: \texttt{PGMXWriter\_\allowbreak{}0\_\allowbreak{}2.\allowbreak{}write} (líneas 70-97) construye el documento entero antes de abrir el fichero, y la excepción salta al construirlo. Antes de escribir, «Save» ya ha copiado el fichero a uno \texttt{.\allowbreak{}bak} junto a él (\texttt{create\allowbreak{}Back\allowbreak{}Up\allowbreak{}Network\allowbreak{}File}, líneas 312-328).

Por qué pasa. Los tres ficheros, en formato 0.2, traen un potencial con \texttt{role="join\allowbreak{}Probability"}, y el nombre que usa el programa es \texttt{joint\allowbreak{}Probability} (\texttt{Potential\allowbreak{}Role.\allowbreak{}JOINT\_\allowbreak{}PROBABILITY} escrito con su \texttt{to\allowbreak{}String}). Es la tabla de probabilidad condicionada de un nodo de azar: «New test» en la primera red (línea 148 del fichero), «AIDS {[}0{]}» en la segunda (línea 808) e «Intervention applied» en la tercera (línea 4854). La copia \texttt{integration\allowbreak{}Tests/\allowbreak{}src/\allowbreak{}main/\allowbreak{}resources/\allowbreak{}networks/\allowbreak{}mid/\allowbreak{}MID-Cochlear.\allowbreak{}pgmx} tiene la misma errata.

\texttt{PGMXReader\_\allowbreak{}0\_\allowbreak{}2.\allowbreak{}get\allowbreak{}Potential\allowbreak{}Role} (módulo \texttt{io}, líneas 1375-1387) rechaza un potencial sin papel con \texttt{Potential\allowbreak{}Without\allowbreak{}Role}, pero un papel que no reconoce lo pasa a \texttt{get\allowbreak{}Potential\allowbreak{}Rol\allowbreak{}By\allowbreak{}Label} (líneas 1563-1571), que devuelve \texttt{null} sin avisar. Los cálculos no usan ese papel, así que las redes se abren y se evalúan con normalidad. Al guardar, \texttt{PGMXWriter\_\allowbreak{}0\_\allowbreak{}2.\allowbreak{}set\allowbreak{}Potential\allowbreak{}Role} (líneas 746-753) llama a \texttt{potential\allowbreak{}Role.\allowbreak{}to\allowbreak{}String(\allowbreak{})} y falla. \texttt{PGMXWriter\_\allowbreak{}1\_\allowbreak{}0} no escribe el papel dentro del potencial, y por eso ese formato sí guarda. Validar el fichero contra su esquema (\hyperref[err:111]{error~111}) tampoco lo detectaría: el esquema no comprueba los atributos del potencial.

Medido con un pequeño programa de prueba:

\begin{itemize}
\tightlist
\item
  Recorriendo todas las redes de \texttt{0\_\allowbreak{}miscellaneous/\allowbreak{}networks}, estas tres son las únicas con un potencial de papel nulo y las únicas que \texttt{PGMXWriter\_\allowbreak{}0\_\allowbreak{}2} no puede guardar.
\item
  «Save» sobre una copia de cada una, como lo hace la ventana (el escritor que devuelve el lector, sobre el mismo fichero): \texttt{Null\allowbreak{}Pointer\allowbreak{}Exception} en \texttt{set\allowbreak{}Potential\allowbreak{}Role}, y el fichero con el mismo tamaño que antes.
\item
  Cambiando en otra copia \texttt{join\allowbreak{}Probability} por \texttt{conditional\allowbreak{}Probability}, o por \texttt{joint\allowbreak{}Probability}, la red se guarda, se reabre y la evaluación sin hallazgos da lo mismo que la original: 271 387,976 en \texttt{ID-CEA-test-2therapies-new-test.\allowbreak{}pgmx} y 156 890,28 en \texttt{MID-Cochlear.\allowbreak{}pgmx}. \texttt{MID-CHAP-Ryan-Griffin.\allowbreak{}pgmx} no se evalúa ni antes ni después, por \hyperref[err:71]{error~71}.
\end{itemize}
\paragraph{El código}
io/src/main/java/org/openmarkov/io/probmodel/reader/PGMXReader\_0\_2.java:1375-1386

\begin{Shaded}
\begin{Highlighting}[]
\KeywordTok{protected}\NormalTok{ PotentialRole }\FunctionTok{getPotentialRole}\OperatorTok{(}\BuiltInTok{Element}\NormalTok{ xmlPotential}\OperatorTok{)} \KeywordTok{throws}\NormalTok{ PGMXParserException }\OperatorTok{\{}
    \BuiltInTok{String}\NormalTok{ xmlPotentialRole }\OperatorTok{=}\NormalTok{ xmlPotential}\OperatorTok{.}\FunctionTok{getAttributeValue}\OperatorTok{(}\NormalTok{XMLAttributes}\OperatorTok{.}\FunctionTok{ROLE}\OperatorTok{.}\FunctionTok{toString}\OperatorTok{());}
    \ControlFlowTok{if} \OperatorTok{(}\NormalTok{xmlPotentialRole }\OperatorTok{==} \KeywordTok{null}\OperatorTok{)} \OperatorTok{\{}
        \ControlFlowTok{throw} \KeywordTok{new}\NormalTok{ PGMXParserException}\OperatorTok{.}\FunctionTok{PotentialWithoutRole}\OperatorTok{(}\NormalTok{xmlPotential}\OperatorTok{);}
    \OperatorTok{\}}
\NormalTok{    PotentialRole xmlRole}\OperatorTok{;}
    \ControlFlowTok{if} \OperatorTok{(}\NormalTok{xmlPotentialRole}\OperatorTok{.}\FunctionTok{equalsIgnoreCase}\OperatorTok{(}\StringTok{"utility"}\OperatorTok{))} \OperatorTok{\{}
\NormalTok{        xmlRole }\OperatorTok{=}\NormalTok{ PotentialRole}\OperatorTok{.}\FunctionTok{UNSPECIFIED}\OperatorTok{;}
    \OperatorTok{\}} \ControlFlowTok{else} \OperatorTok{\{}
\NormalTok{        xmlRole }\OperatorTok{=}\NormalTok{ PGMXReader\_0\_2}\OperatorTok{.}\FunctionTok{getPotentialRolByLabel}\OperatorTok{(}\NormalTok{xmlPotentialRole}\OperatorTok{);}
    \OperatorTok{\}}
    \ControlFlowTok{return}\NormalTok{ xmlRole}\OperatorTok{;}
\end{Highlighting}
\end{Shaded}

io/src/main/java/org/openmarkov/io/probmodel/reader/PGMXReader\_0\_2.java:1563-1571

\begin{Shaded}
\begin{Highlighting}[]
\KeywordTok{protected} \DataTypeTok{static}\NormalTok{ PotentialRole }\FunctionTok{getPotentialRolByLabel}\OperatorTok{(}\BuiltInTok{String}\NormalTok{ auxLabel}\OperatorTok{)} \OperatorTok{\{}
    \ControlFlowTok{for} \OperatorTok{(}\NormalTok{PotentialRole role }\OperatorTok{:}\NormalTok{ PotentialRole}\OperatorTok{.}\FunctionTok{values}\OperatorTok{())} \OperatorTok{\{}
        \BuiltInTok{String}\NormalTok{ u }\OperatorTok{=}\NormalTok{ role}\OperatorTok{.}\FunctionTok{toString}\OperatorTok{();}
        \ControlFlowTok{if} \OperatorTok{(}\NormalTok{u}\OperatorTok{.}\FunctionTok{equals}\OperatorTok{(}\NormalTok{auxLabel}\OperatorTok{))} \OperatorTok{\{}
            \ControlFlowTok{return}\NormalTok{ role}\OperatorTok{;}
        \OperatorTok{\}}
    \OperatorTok{\}}
    \ControlFlowTok{return} \KeywordTok{null}\OperatorTok{;}
\OperatorTok{\}}
\end{Highlighting}
\end{Shaded}

io/src/main/java/org/openmarkov/io/probmodel/writer/PGMXWriter\_0\_2.java:746-748

\begin{Shaded}
\begin{Highlighting}[]
\KeywordTok{protected} \DataTypeTok{static} \DataTypeTok{void} \FunctionTok{setPotentialRole}\OperatorTok{(}\BuiltInTok{Element}\NormalTok{ potentialElement}\OperatorTok{,}\NormalTok{ Potential potential}\OperatorTok{)} \OperatorTok{\{}
\NormalTok{    PotentialRole potentialRole }\OperatorTok{=}\NormalTok{ potential}\OperatorTok{.}\FunctionTok{getPotentialRole}\OperatorTok{();}
    \BuiltInTok{String}\NormalTok{ oldPotentialRoleString }\OperatorTok{=}\NormalTok{ potentialRole}\OperatorTok{.}\FunctionTok{toString}\OperatorTok{();}
\end{Highlighting}
\end{Shaded}
\paragraph{Cómo arreglarlo}
Hay tres cosas que hacer, y la primera exige decidir:

\begin{itemize}
\tightlist
\item
  Qué hace el lector con un papel desconocido. Una opción es rechazar el fichero con una excepción que nombre el papel y el nodo, como ya hace con un potencial sin papel; entonces hay que corregir antes los cuatro ficheros citados, o dejarán de abrirse. La otra es aceptar \texttt{join\allowbreak{}Probability} como sinónimo y avisar al abrir. En los dos casos, lo que no puede hacer es dejar el papel en \texttt{null} en silencio. Para estas tablas de nodos de azar, el papel que corresponde es el de probabilidad condicionada.
\item
  El escritor no debería fallar con un papel nulo: o deduce el papel del tipo de nodo, o lanza \texttt{Writer\allowbreak{}Exception} con un mensaje que nombre el nodo, que la ventana enseñaría como error normal.
\item
  Corregir la errata en los cuatro ficheros.
\end{itemize}

Prueba: leer las tres redes, guardarlas con \texttt{PGMXWriter\_\allowbreak{}0\_\allowbreak{}2} en un fichero temporal, reabrirlas y comprobar que \texttt{0\_\allowbreak{}miscellaneous/\allowbreak{}networks/\allowbreak{}id/\allowbreak{}ID-CEA-test-2therapies-new-test.\allowbreak{}pgmx} sigue dando una utilidad esperada de 271 387,976; y una prueba del lector con un papel inventado que dé el comportamiento que se decida.

\setcounter{subsection}{114}
\subsection[Una distribución desconocida en una red de sucesos falla lejos de la causa]{Abrir una red de sucesos con un nombre de distribución o parametrización que no existe falla con un NullPointerException lejos de la causa}\label{err:115}
\begin{ficha}
\fichakey{Estado}Presente
\fichakey{Módulo}core
\fichakey{Paquete}org.openmarkov.core.model.network.potential
\fichakey{Ficheros}\path{core/src/main/java/org/openmarkov/core/model/network/potential/DistributionTablePotential.java:93-96} \newline \path{core/src/main/java/org/openmarkov/core/model/network/potential/DistributionTablePotential.java:141-165} \newline \path{io/src/main/java/org/openmarkov/io/probmodel/reader/PGMXReader_1_0.java:310-324} \newline \path{gui/src/main/java/org/openmarkov/gui/dialog/common/DistributionTablePotentialPanel.java:217-236}
\fichakey{Comprobado}Leyendo el código y ejecutándolo
\fichakey{Esfuerzo}Pequeño (horas)
\fichakey{Decisión de equipo}\textcolor{decisionrojo}{\textbf{Sí.}} Si un fichero cuyo potencial de suceso no trae el atributo \texttt{parametrization} debe leerse con una parametrización por omisión de esa distribución o rechazarse con un mensaje
\fichakey{Relacionados}\hyperref[err:116]{error~116}
\end{ficha}
\paragraph{Qué pasa}
El usuario abre un fichero \texttt{.\allowbreak{}pgmx} de una red DES (simulación de eventos discretos) en la que el potencial de un suceso nombra una distribución o una parametrización que el programa no conoce, o no trae el atributo \texttt{parametrization}. La lectura termina con un \texttt{Null\allowbreak{}Pointer\allowbreak{}Exception} cuyo mensaje no dice qué nodo ni qué nombre no se reconoció.

Lo que ve el usuario es el aviso de fichero estropeado: \texttt{Nets\allowbreak{}IO.\allowbreak{}open\allowbreak{}Network\allowbreak{}URL} (módulo \texttt{gui}, líneas 173-178) convierte cualquier excepción de tiempo de ejecución de la lectura en \texttt{Corrupt\allowbreak{}Network\allowbreak{}File}, cuyo mensaje es «Network from \ldots{} is corrupted and could not be loaded.». El aviso no dice qué nodo ni qué nombre no se reconoció, y la causa solo queda en la traza de la consola.

\texttt{Distribution\allowbreak{}Table\allowbreak{}Potential} es el potencial de un nodo de suceso: una distribución de probabilidad cuyos parámetros se leen de una tabla. Construye su mitad interna en \texttt{change\allowbreak{}Distribution}: pide al gestor \texttt{Parametrized\allowbreak{}Function\allowbreak{}Manager} la lista de parámetros de la pareja (distribución, parametrización), fabrica la variable de parámetros y la tabla \texttt{Table\allowbreak{}With\allowbreak{}Events}. Todo va dentro de un \texttt{try} cuyo \texttt{catch\ (\allowbreak{}Exception\ e)} solo imprime la traza.

Si el gestor no conoce la pareja, \texttt{get\allowbreak{}Parameters} devuelve \texttt{null}, la línea siguiente lanza \texttt{Null\allowbreak{}Pointer\allowbreak{}Exception}, el \texttt{catch} lo imprime por la consola y el constructor termina con normalidad, dejando \texttt{distribution\allowbreak{}Parameters}, \texttt{distribution\allowbreak{}Variable} y \texttt{table\allowbreak{}With\allowbreak{}Events} a \texttt{null}. El siguiente que use el potencial falla lejos de la causa.

Al abrir, \texttt{PGMXReader\_\allowbreak{}1\_\allowbreak{}0.\allowbreak{}get\allowbreak{}Distribution\allowbreak{}Table\allowbreak{}Potential} toma los nombres de los atributos \texttt{distribution} y \texttt{parametrization} del elemento \texttt{\textless{}Value\textgreater{}} sin validarlos, construye el potencial y en la línea siguiente llama a \texttt{potential.\allowbreak{}get\allowbreak{}Table\allowbreak{}With\allowbreak{}Events(\allowbreak{}).\allowbreak{}set\allowbreak{}Values(\allowbreak{}table)}.

Medido:

\begin{itemize}
\tightlist
\item
  \texttt{new\ Distribution\allowbreak{}Table\allowbreak{}Potential(\allowbreak{}vars,\allowbreak{}\ rol,\allowbreak{}\ "Beta",\allowbreak{}\ "Alpha\ /\allowbreak{}\ Beta")} imprime la traza y devuelve un objeto con \texttt{get\allowbreak{}Table\allowbreak{}With\allowbreak{}Events(\allowbreak{})} y \texttt{get\allowbreak{}Distribution\allowbreak{}Variable(\allowbreak{})} a \texttt{null}.
\item
  Con dos variantes de una red DES de prueba, (a) con \texttt{parametrization="Alpha\ /\allowbreak{}\ Beta"} en el potencial del suceso \texttt{Death} y (b) sin atributo \texttt{parametrization}, la consola muestra la primera traza y la lectura termina con \texttt{Null\allowbreak{}Pointer\allowbreak{}Exception:\ Cannot\ invoke\ "Table\allowbreak{}With\allowbreak{}Events.\allowbreak{}set\allowbreak{}Values(\allowbreak{}double{[}{]})"\ because\ the\ return\ value\ of\ "Distribution\allowbreak{}Table\allowbreak{}Potential.\allowbreak{}get\allowbreak{}Table\allowbreak{}With\allowbreak{}Events(\allowbreak{})"\ is\ null}.
\end{itemize}

Desde la ventana del potencial (\texttt{Distribution\allowbreak{}Table\allowbreak{}Potential\allowbreak{}Panel}) los desplegables solo ofrecen distribuciones y parametrizaciones que el gestor conoce, así que por ahí no se ha encontrado manera de llegar (comprobado leyendo, no ejecutando la ventana).
\paragraph{El código}
core/src/main/java/org/openmarkov/core/model/network/potential/DistributionTablePotential.java:141-165

\begin{Shaded}
\begin{Highlighting}[]
    \KeywordTok{public} \DataTypeTok{void} \FunctionTok{changeDistribution}\OperatorTok{(}\BuiltInTok{String}\NormalTok{ distributionName}\OperatorTok{,} \BuiltInTok{String}\NormalTok{ distributionParametrization}\OperatorTok{)} \OperatorTok{\{}
        \ControlFlowTok{try} \OperatorTok{\{}
            \KeywordTok{this}\OperatorTok{.}\FunctionTok{distributionName} \OperatorTok{=}\NormalTok{ distributionName}\OperatorTok{;}
            \KeywordTok{this}\OperatorTok{.}\FunctionTok{parametrizationName} \OperatorTok{=}\NormalTok{ distributionParametrization}\OperatorTok{;}
            \BuiltInTok{List}\OperatorTok{\textless{}}\BuiltInTok{String}\OperatorTok{\textgreater{}}\NormalTok{ distributionParametersList }\OperatorTok{=}\NormalTok{ ParametrizedFunctionManager}\OperatorTok{.}\FunctionTok{getUniqueInstance}\OperatorTok{().}\FunctionTok{getParameters}\OperatorTok{(}\NormalTok{distributionName}\OperatorTok{,}\NormalTok{ distributionParametrization}\OperatorTok{);}
\NormalTok{            distributionParameters }\OperatorTok{=}\NormalTok{ distributionParametersList}\OperatorTok{.}\FunctionTok{toArray}\OperatorTok{(}\KeywordTok{new} \BuiltInTok{String}\OperatorTok{[}\NormalTok{distributionParametersList}\OperatorTok{.}\FunctionTok{size}\OperatorTok{()]);}
            \KeywordTok{...}
            \KeywordTok{this}\OperatorTok{.}\FunctionTok{tableWithEvents} \OperatorTok{=} \KeywordTok{new} \FunctionTok{TableWithEvents}\OperatorTok{(}\NormalTok{tableParents}\OperatorTok{,}\NormalTok{ role}\OperatorTok{,} \OperatorTok{(}\NormalTok{numericVariables}\OperatorTok{.}\FunctionTok{size}\OperatorTok{()} \OperatorTok{\textgreater{}} \DecValTok{0}\OperatorTok{));}
        \OperatorTok{\}} \ControlFlowTok{catch} \OperatorTok{(}\BuiltInTok{Exception}\NormalTok{ e}\OperatorTok{)} \OperatorTok{\{}
\NormalTok{            e}\OperatorTok{.}\FunctionTok{printStackTrace}\OperatorTok{();}
        \OperatorTok{\}}
    \OperatorTok{\}}
\end{Highlighting}
\end{Shaded}

io/src/main/java/org/openmarkov/io/probmodel/reader/PGMXReader\_1\_0.java:313-324

\begin{Shaded}
\begin{Highlighting}[]
        \BuiltInTok{String}\NormalTok{ distributionName }\OperatorTok{=}\NormalTok{ distribution}\OperatorTok{.}\FunctionTok{getAttributeValue}\OperatorTok{(}\NormalTok{XMLAttributes}\OperatorTok{.}\FunctionTok{DISTRIBUTION}\OperatorTok{.}\FunctionTok{toString}\OperatorTok{());}
        \CommentTok{//08/10/2020 {-}parametrization added}
        \BuiltInTok{String}\NormalTok{ distributionParametrization }\OperatorTok{=}\NormalTok{ distribution}\OperatorTok{.}\FunctionTok{getAttributeValue}\OperatorTok{(}\NormalTok{XMLAttributes}\OperatorTok{.}\FunctionTok{PARAMETRIZATION}\OperatorTok{.}\FunctionTok{toString}\OperatorTok{());}
        \KeywordTok{...}
\NormalTok{        potential }\OperatorTok{=} \KeywordTok{new} \FunctionTok{DistributionTablePotential}\OperatorTok{(}\NormalTok{variables}\OperatorTok{,}\NormalTok{ xmlRole}\OperatorTok{,}\NormalTok{ distributionName}\OperatorTok{,}\NormalTok{ distributionParametrization}\OperatorTok{);}
\NormalTok{        potential}\OperatorTok{.}\FunctionTok{getTableWithEvents}\OperatorTok{().}\FunctionTok{setValues}\OperatorTok{(}\NormalTok{table}\OperatorTok{);}
\end{Highlighting}
\end{Shaded}
\paragraph{Cómo arreglarlo}
Quitar el \texttt{catch} genérico de \texttt{change\allowbreak{}Distribution} y rechazar en el sitio una pareja desconocida con una excepción declarada que nombre la distribución y la parametrización.

Sitios donde la regla «la pareja distribución/parametrización debe existir» tiene que cumplirse:

\begin{itemize}
\tightlist
\item
  \texttt{Distribution\allowbreak{}Table\allowbreak{}Potential.\allowbreak{}change\allowbreak{}Distribution}, que llaman el constructor de cuatro argumentos, el de dos, el de copia y la ventana en \texttt{Distribution\allowbreak{}Table\allowbreak{}Potential\allowbreak{}Panel.\allowbreak{}item\allowbreak{}State\allowbreak{}Changed} (línea 229).
\item
  \texttt{PGMXReader\_\allowbreak{}1\_\allowbreak{}0.\allowbreak{}get\allowbreak{}Distribution\allowbreak{}Table\allowbreak{}Potential}: traducir esa excepción a la excepción de lectura del formato (\texttt{Prob\allowbreak{}Net\allowbreak{}Parser\allowbreak{}Exception} o la que use el lector) con el nombre del nodo, para que el usuario vea un mensaje de fichero inválido.
\item
  \texttt{Distribution\allowbreak{}Table\allowbreak{}Potential.\allowbreak{}sample\allowbreak{}Conditioned\allowbreak{}Variable} (línea 291) vuelve a pedir la clase al gestor con los mismos nombres; con la validación en la construcción ya no puede recibir una pareja desconocida.
\end{itemize}

Queda por decidir qué hacer con los ficheros sin atributo \texttt{parametrization}. Opciones: rechazarlos con un mensaje, o elegir una parametrización por omisión para la distribución nombrada.

Prueba de regresión: leer un \texttt{.\allowbreak{}pgmx} con una parametrización inexistente produce la excepción de lectura con el nombre del nodo y de la pareja, sin trazas en la consola; construir el potencial con esa pareja lanza la excepción declarada.

\setcounter{subsection}{115}
\subsection[Una tabla de sucesos con incertidumbre en el fichero impide abrir la red]{Una red de sucesos cuyo fichero trae incertidumbre en la tabla de parámetros de un suceso no se puede abrir, y el aviso solo dice que el fichero está dañado}\label{err:116}
\begin{ficha}
\fichakey{Estado}Presente
\fichakey{Módulo}io, core
\fichakey{Paquete}org.openmarkov.io.probmodel.reader; org.openmarkov.core.model.network.potential
\fichakey{Ficheros}\path{io/src/main/java/org/openmarkov/io/probmodel/reader/PGMXReader_1_0.java:311-341} \newline \path{core/src/main/java/org/openmarkov/core/model/network/potential/TablePotential.java:297-316} \newline \path{core/src/main/java/org/openmarkov/core/model/network/potential/TableWithEvents.java:71,114,233-239} \newline \path{io/src/main/java/org/openmarkov/io/probmodel/reader/PGMXPotentialParsers.java:97-120} \newline \path{io/src/main/java/org/openmarkov/io/probmodel/writer/PGMXWriter_1_0.java:317-324} \newline \path{io/src/main/java/org/openmarkov/io/probmodel/writer/PGMXWriter_0_2.java:893-902} \newline \path{core/src/main/java/org/openmarkov/core/action/core/UncertainValuesEdit.java:151-174} \newline \path{gui/src/main/java/org/openmarkov/gui/dialog/io/NetsIO.java:168-180}
\fichakey{Comprobado}Leyendo el código y ejecutándolo
\fichakey{Esfuerzo}Pequeño (horas)
\fichakey{Decisión de equipo}\textcolor{decisionrojo}{\textbf{Sí.}} Si las tablas de parámetros de los sucesos pueden llevar incertidumbre (y entonces se conserva y se vuelve a escribir, aunque hoy la simulación no la use) o si el lector debe rechazarla con un mensaje que nombre el nodo, o ignorarla avisando
\fichakey{Relacionados}\hyperref[err:76]{error~76}, \hyperref[err:115]{error~115}, \hyperref[nota:58]{nota~58}
\end{ficha}
\paragraph{Qué pasa}
En una red DES (simulación de eventos discretos), el tiempo hasta que ocurre un suceso lo da un \texttt{Distribution\allowbreak{}Table\allowbreak{}Potential}: una distribución (exacta, exponencial\ldots) cuyos parámetros se leen de una tabla interna, un \texttt{Table\allowbreak{}With\allowbreak{}Events} que a su vez guarda un \texttt{Table\allowbreak{}Potential}.

El usuario abre con Archivo → Abrir un \texttt{.\allowbreak{}pgmx} de una red de sucesos en el que el potencial de un suceso trae, además de \texttt{\textless{}Values\textgreater{}}, un elemento \texttt{\textless{}Uncertain\allowbreak{}Values\textgreater{}} con distribuciones para sus parámetros. La red no se abre. La ventana dice «Network from file:\ldots{} is corrupted and could not be loaded.», sin decir qué nodo ni qué elemento del fichero es el problema.

El formato lo admite: el esquema \texttt{val\_\allowbreak{}v4.\allowbreak{}xsd} permite \texttt{\textless{}Uncertain\allowbreak{}Values\textgreater{}} en cualquier \texttt{\textless{}Potential\textgreater{}}, y el lector de la versión 1.0 lo busca a propósito en estos potenciales. \texttt{PGMXReader\_\allowbreak{}1\_\allowbreak{}0.\allowbreak{}get\allowbreak{}Distribution\allowbreak{}Table\allowbreak{}Potential} (líneas 326-331) lo lee y lo pasa a \texttt{potential.\allowbreak{}get\allowbreak{}Table\allowbreak{}With\allowbreak{}Events(\allowbreak{}).\allowbreak{}set\allowbreak{}Uncertain\allowbreak{}Values(\allowbreak{}.\allowbreak{}.\allowbreak{}.\allowbreak{})}.

Por qué falla: la tabla interna es siempre un \texttt{Table\allowbreak{}Potential} a secas (los constructores de \texttt{Table\allowbreak{}With\allowbreak{}Events} la crean con \texttt{new\ Table\allowbreak{}Potential(\allowbreak{}.\allowbreak{}.\allowbreak{}.\allowbreak{})}, líneas 71 y 114), y \texttt{Table\allowbreak{}Potential.\allowbreak{}set\allowbreak{}Uncertain\allowbreak{}Values} (líneas 310-316) lanza \texttt{Unsupported\allowbreak{}Operation\allowbreak{}Exception} con cualquier lista que no sea nula: solo su subclase \texttt{Uncertain\allowbreak{}Table\allowbreak{}Potential} puede guardar incertidumbre. Para las tablas corrientes el lector ya elige \texttt{Uncertain\allowbreak{}Table\allowbreak{}Potential} cuando el fichero trae incertidumbre (\texttt{PGMXPotential\allowbreak{}Parsers.\allowbreak{}get\allowbreak{}Table\allowbreak{}Potential}, líneas 97-120); para las tablas de sucesos nadie lo hace.

\texttt{Nets\allowbreak{}IO.\allowbreak{}open\allowbreak{}Network\allowbreak{}URL} convierte la excepción en \texttt{Corrupt\allowbreak{}Network\allowbreak{}File} sin guardar la causa, y la ventana la presenta con su texto fijo (esto último, por lectura).

Medido con una red de sucesos de prueba que se abre sin problemas (una decisión \texttt{Therapy}, un suceso \texttt{Onset}, un nodo de azar \texttt{Ill} y un suceso terminal \texttt{Death} con un \texttt{Distribution\allowbreak{}Table} exacto sobre \texttt{{[}Death,\allowbreak{}\ Onset,\allowbreak{}\ Ill{]}}), añadiendo al potencial de \texttt{Death} un \texttt{\textless{}Uncertain\allowbreak{}Values\textgreater{}} con dos distribuciones normales:

\begin{Verbatim}[breaklines,breakanywhere,fontsize=\small]
FormatManager.getProbNetReader(url).read(url) -> UnsupportedOperationException: Cannot set uncertain values on a plain TablePotential. Use UncertainTablePotential instead.
     at TablePotential.setUncertainValues(TablePotential.java:312)
     at TableWithEvents.setUncertainValues(TableWithEvents.java:238)
     at PGMXReader_1_0.getDistributionTablePotential(PGMXReader_1_0.java:330)
NetsIO.openNetworkURL(url) -> CorruptNetworkFile: Network from file:…/red.pgmx is corrupted and could not be loaded. (causa: null)
\end{Verbatim}

Como comparación, el mismo \texttt{\textless{}Uncertain\allowbreak{}Values\textgreater{}} (una Beta y un complemento) puesto en la tabla corriente del nodo de azar \texttt{Sex} de esa red se abre bien, y el potencial queda como \texttt{Uncertain\allowbreak{}Table\allowbreak{}Potential}.

Camino de llegada: solo un fichero escrito a mano o por otra herramienta. OpenMarkov no escribe nunca esa incertidumbre: \texttt{PGMXWriter\_\allowbreak{}1\_\allowbreak{}0.\allowbreak{}get\allowbreak{}Distribution\allowbreak{}Table\allowbreak{}Potential} escribe la tabla interna con \texttt{get\allowbreak{}Values\allowbreak{}Table\allowbreak{}Potential} (\texttt{PGMXWriter\_\allowbreak{}0\_\allowbreak{}2}, líneas 893-902), que solo añade \texttt{\textless{}Uncertain\allowbreak{}Values\textgreater{}} si \texttt{get\allowbreak{}Uncertain\allowbreak{}Values(\allowbreak{})} no es nulo, y en esta tabla lo es siempre. Tampoco se puede poner desde la interfaz: el menú contextual de la tabla de un suceso (\texttt{Table\allowbreak{}With\allowbreak{}Events\allowbreak{}Panel}) solo ofrece «Add function», y \texttt{Uncertain\allowbreak{}Values\allowbreak{}Edit.\allowbreak{}prepare\allowbreak{}Uncertain\allowbreak{}Potential} (líneas 151-174) rechaza expresamente una tabla que va dentro de otro potencial. El repositorio no tiene redes de sucesos.

Aunque se leyera, esa incertidumbre no tendría hoy ningún efecto: la rama de análisis probabilístico de sensibilidad de la simulación no se puede activar, como se explica en \hyperref[nota:58]{nota~58}, donde no hay nada que corregir.
\paragraph{El código}
io/src/main/java/org/openmarkov/io/probmodel/reader/PGMXReader\_1\_0.java:320-331

\begin{Shaded}
\begin{Highlighting}[]
        \BuiltInTok{Element}\NormalTok{ xmlUncertain }\OperatorTok{=}\NormalTok{ xmlPotential}\OperatorTok{.}\FunctionTok{getChild}\OperatorTok{(}\NormalTok{XMLTags}\OperatorTok{.}\FunctionTok{UNCERTAIN\_VALUES}\OperatorTok{.}\FunctionTok{toString}\OperatorTok{());}
        \KeywordTok{...}
\NormalTok{        potential }\OperatorTok{=} \KeywordTok{new} \FunctionTok{DistributionTablePotential}\OperatorTok{(}\NormalTok{variables}\OperatorTok{,}\NormalTok{ xmlRole}\OperatorTok{,}\NormalTok{ distributionName}\OperatorTok{,}\NormalTok{ distributionParametrization}\OperatorTok{);}
\NormalTok{        potential}\OperatorTok{.}\FunctionTok{getTableWithEvents}\OperatorTok{().}\FunctionTok{setValues}\OperatorTok{(}\NormalTok{table}\OperatorTok{);}
        \CommentTok{// 20/05/2024 Begin PSA}
\NormalTok{        UncertainValue}\OperatorTok{[]}\NormalTok{ uncertainValues}\OperatorTok{;}
        \ControlFlowTok{if} \OperatorTok{(}\NormalTok{xmlUncertain }\OperatorTok{!=} \KeywordTok{null}\OperatorTok{)} \OperatorTok{\{}
\NormalTok{            uncertainValues }\OperatorTok{=} \FunctionTok{getUncertainValues}\OperatorTok{(}\NormalTok{xmlUncertain}\OperatorTok{);}
\NormalTok{            potential}\OperatorTok{.}\FunctionTok{getTableWithEvents}\OperatorTok{().}\FunctionTok{setUncertainValues}\OperatorTok{(}\NormalTok{uncertainValues}\OperatorTok{);}
        \OperatorTok{\}}
\end{Highlighting}
\end{Shaded}

core/src/main/java/org/openmarkov/core/model/network/potential/TablePotential.java:310-316

\begin{Shaded}
\begin{Highlighting}[]
    \KeywordTok{public} \DataTypeTok{void} \FunctionTok{setUncertainValues}\OperatorTok{(}\NormalTok{UncertainValue}\OperatorTok{[]}\NormalTok{ uv}\OperatorTok{)} \OperatorTok{\{}
        \ControlFlowTok{if} \OperatorTok{(}\NormalTok{uv }\OperatorTok{!=} \KeywordTok{null}\OperatorTok{)} \OperatorTok{\{}
            \ControlFlowTok{throw} \KeywordTok{new} \BuiltInTok{UnsupportedOperationException}\OperatorTok{(}
                    \StringTok{"Cannot set uncertain values on a plain TablePotential. Use UncertainTablePotential instead."}\OperatorTok{);}
        \OperatorTok{\}}
        \CommentTok{// null is a no{-}op: removing uncertainty from a non{-}uncertain potential is harmless}
    \OperatorTok{\}}
\end{Highlighting}
\end{Shaded}
\paragraph{Cómo arreglarlo}
Hace falta decidir si las tablas de parámetros de los sucesos admiten incertidumbre. Opciones, sin elegir:

\begin{itemize}
\tightlist
\item
  \textbf{Admitirla.} Que la tabla interna de \texttt{Table\allowbreak{}With\allowbreak{}Events} sea un \texttt{Uncertain\allowbreak{}Table\allowbreak{}Potential} cuando tenga que llevar incertidumbre, o que \texttt{Table\allowbreak{}With\allowbreak{}Events.\allowbreak{}set\allowbreak{}Uncertain\allowbreak{}Values} la sustituya por uno con las mismas variables y valores, como hace \texttt{Uncertain\allowbreak{}Values\allowbreak{}Edit.\allowbreak{}upgrade\allowbreak{}To\allowbreak{}Uncertain} con la tabla de un nodo. Entonces el constructor de copia de \texttt{Table\allowbreak{}With\allowbreak{}Events} (línea 114, \texttt{new\ Table\allowbreak{}Potential(\allowbreak{}.\allowbreak{}.\allowbreak{}.\allowbreak{})}) tiene que conservar la clase y la incertidumbre, y el escritor la volverá a escribir sin más cambios (\texttt{get\allowbreak{}Values\allowbreak{}Table\allowbreak{}Potential}). La simulación seguiría sin usarla.
\item
  \textbf{Rechazarla.} En \texttt{PGMXReader\_\allowbreak{}1\_\allowbreak{}0.\allowbreak{}get\allowbreak{}Distribution\allowbreak{}Table\allowbreak{}Potential}, lanzar la excepción de lectura del formato (\texttt{PGMXParser\allowbreak{}Exception} o la que use el lector) con el nombre del nodo, diciendo que las tablas de sucesos no admiten incertidumbre.
\item
  \textbf{Ignorarla avisando.} Leer el resto del potencial y avisar de que se ha descartado \texttt{\textless{}Uncertain\allowbreak{}Values\textgreater{}}.
\end{itemize}

Sitios donde tiene que cumplirse la regla que se elija:

\begin{itemize}
\tightlist
\item
  \texttt{PGMXReader\_\allowbreak{}1\_\allowbreak{}0.\allowbreak{}get\allowbreak{}Distribution\allowbreak{}Table\allowbreak{}Potential}, el único que hoy lee incertidumbre para una tabla de sucesos.
\item
  \texttt{PGMXReader\_\allowbreak{}1\_\allowbreak{}0.\allowbreak{}get\allowbreak{}Transition\allowbreak{}Table\allowbreak{}Potential} (líneas 288-298), que no mira \texttt{\textless{}Uncertain\allowbreak{}Values\textgreater{}} y, por lectura, la pierde sin avisar; debería tratarla igual.
\item
  \texttt{Table\allowbreak{}With\allowbreak{}Events.\allowbreak{}set\allowbreak{}Uncertain\allowbreak{}Values} y el constructor de copia de \texttt{Table\allowbreak{}With\allowbreak{}Events}.
\item
  \texttt{Uncertain\allowbreak{}Values\allowbreak{}Edit.\allowbreak{}prepare\allowbreak{}Uncertain\allowbreak{}Potential}, que ya rechaza la incertidumbre en tablas internas desde la interfaz; si se decide admitirla, es el otro sitio que cambia.
\end{itemize}

Prueba de regresión: en \texttt{io/\allowbreak{}src/\allowbreak{}test}, leer un \texttt{.\allowbreak{}pgmx} de sucesos con \texttt{\textless{}Uncertain\allowbreak{}Values\textgreater{}} en un \texttt{Distribution\allowbreak{}Table}. Según lo decidido, se abre y conserva la incertidumbre (y al guardar y releer sigue ahí), o lanza una \texttt{Prob\allowbreak{}Net\allowbreak{}Parser\allowbreak{}Exception} con el nombre del nodo; en ningún caso \texttt{Unsupported\allowbreak{}Operation\allowbreak{}Exception}.

\setcounter{subsection}{116}
\subsection[Abrir un .xml que no es ProbModelXML falla con una excepción de índice]{Abrir un fichero \texttt{.\allowbreak{}xml} que no es ProbModelXML (por ejemplo un XMLBIF) desde «Abrir» falla con un \texttt{String\allowbreak{}Index\allowbreak{}Out\allowbreak{}Of\allowbreak{}Bounds\allowbreak{}Exception} sin decir que el formato no se reconoce}\label{err:117}
\begin{ficha}
\fichakey{Estado}Presente
\fichakey{Módulo}io, gui
\fichakey{Paquete}org.openmarkov.core.io.format.annotation; org.openmarkov.io.probmodel.reader; org.openmarkov.gui.dialog.io
\fichakey{Ficheros}\path{io/src/main/java/org/openmarkov/core/io/format/annotation/FormatManager.java:119-160} \newline \path{io/src/main/java/org/openmarkov/io/probmodel/reader/PGMXReader.java:29-41} \newline \path{io/src/main/java/org/openmarkov/io/probmodel/reader/ReaderFactory.java:24-35} \newline \path{gui/src/main/java/org/openmarkov/gui/dialog/io/NetsIO.java:168-180} \newline \path{gui/src/main/java/org/openmarkov/gui/dialog/io/NetworkOMFileChooser.java:47-59}
\fichakey{Comprobado}Leyendo el código y ejecutándolo
\fichakey{Esfuerzo}Pequeño (horas)
\fichakey{Decisión de equipo}No
\fichakey{Relacionados}\hyperref[err:111]{error~111}, \hyperref[nota:173]{nota~173}
\end{ficha}
\paragraph{Qué pasa}
El usuario elige en el diálogo de abrir red un fichero \texttt{.\allowbreak{}xml} que no es ProbModelXML, por ejemplo uno en XMLBIF (un formato de intercambio de redes bayesianas en XML, lenguaje de marcado extensible), que usa precisamente esa extensión. Recibe una excepción de índice sin relación con el problema, en vez de un mensaje del tipo «este fichero no es ProbModelXML» (o «XMLBIF no está soportado»). No se ha comprobado cómo presenta la ventana esta excepción no controlada. Lo mismo ocurre con cualquier \texttt{.\allowbreak{}xml} sin atributo \texttt{format\allowbreak{}Version}, o con un valor sin punto.

El diálogo (\texttt{Network\allowbreak{}OMFile\allowbreak{}Chooser}) construye sus filtros a partir de los lectores dados de alta. El de ProbModelXML (\texttt{PGMXReader}) declara las extensiones \texttt{pgmx} y \texttt{xml}, así que el usuario puede elegir cualquier fichero \texttt{.\allowbreak{}xml}.

Al abrir, \texttt{Nets\allowbreak{}IO.\allowbreak{}open\allowbreak{}Network\allowbreak{}URL} llama a \texttt{Format\allowbreak{}Manager.\allowbreak{}get\allowbreak{}Prob\allowbreak{}Net\allowbreak{}Reader(\allowbreak{}url)}, que antes de elegir lector calcula la versión del fichero en \texttt{get\allowbreak{}File\allowbreak{}Version}: lee el atributo \texttt{format\allowbreak{}Version} del elemento raíz y le quita la última parte con \texttt{substring(\allowbreak{}0,\allowbreak{}\ last\allowbreak{}Index\allowbreak{}Of(\allowbreak{}\textquotesingle{}.\allowbreak{}\textquotesingle{}))}. En un fichero sin ese atributo, \texttt{get\allowbreak{}Attribute} devuelve la cadena vacía, \texttt{last\allowbreak{}Index\allowbreak{}Of(\allowbreak{}\textquotesingle{}.\allowbreak{}\textquotesingle{})} es −1 y \texttt{substring(\allowbreak{}0,\allowbreak{}\ -1)} lanza. Esa llamada está fuera del \texttt{try} de \texttt{open\allowbreak{}Network\allowbreak{}URL}, así que la excepción no se convierte en \texttt{Corrupt\allowbreak{}Network\allowbreak{}File} y sube tal cual.

Ejecutado con \texttt{io/\allowbreak{}src/\allowbreak{}test/\allowbreak{}resources/\allowbreak{}net\allowbreak{}AB.\allowbreak{}xml} (un XMLBIF correcto):

\begin{Verbatim}[breaklines,breakanywhere,fontsize=\small]
FormatManager.getProbNetReader -> java.lang.StringIndexOutOfBoundsException: Range [0, -1) out of bounds for length 0
PGMXReader.read               -> java.lang.NullPointerException: Cannot invoke "String.startsWith(String)" because "strVersion" is null
                                 at ReaderFactory.lambda$getReader$0(ReaderFactory.java:27)
\end{Verbatim}

La segunda línea es lo que pasaría si se llegara al lector: \texttt{PGMXReader.\allowbreak{}read} pasa a \texttt{Reader\allowbreak{}Factory.\allowbreak{}get\allowbreak{}Reader} el atributo \texttt{format\allowbreak{}Version} leído con la biblioteca JDOM, que es \texttt{null}, y \texttt{str\allowbreak{}Version.\allowbreak{}starts\allowbreak{}With} revienta antes de llegar al \texttt{case\ null} que lanzaría la excepción prevista \texttt{Prob\allowbreak{}Net\allowbreak{}Parser\allowbreak{}Exception.\allowbreak{}Wrong\allowbreak{}Version}.

El estado del propio lector de XMLBIF se describe aparte, en \hyperref[nota:50]{nota~50}, donde no hay nada que corregir.
\paragraph{El código}
io/src/main/java/org/openmarkov/core/io/format/annotation/FormatManager.java:134-159

\begin{Shaded}
\begin{Highlighting}[]
\KeywordTok{private} \DataTypeTok{static} \AttributeTok{@NotNull} \BuiltInTok{String} \FunctionTok{getFileVersion}\OperatorTok{(}\BuiltInTok{URL}\NormalTok{ url}\OperatorTok{,} \BuiltInTok{String}\NormalTok{ fileExtension}\OperatorTok{)} \KeywordTok{throws}\NormalTok{ ProbNetParserException}\OperatorTok{.}\FunctionTok{BadlyStructuredFile} \OperatorTok{\{}
    \BuiltInTok{String}\NormalTok{ fileVersion}\OperatorTok{;}
    \ControlFlowTok{if} \OperatorTok{(}\NormalTok{fileExtension}\OperatorTok{.}\FunctionTok{equals}\OperatorTok{(}\StringTok{"elv"}\OperatorTok{))} \OperatorTok{\{}
\NormalTok{        fileVersion }\OperatorTok{=} \StringTok{""}\OperatorTok{;}
    \OperatorTok{\}} \ControlFlowTok{else} \OperatorTok{\{}
        \KeywordTok{...}
\NormalTok{        fileVersion }\OperatorTok{=}\NormalTok{ doc}\OperatorTok{.}\FunctionTok{getDocumentElement}\OperatorTok{().}\FunctionTok{getAttribute}\OperatorTok{(}\StringTok{"formatVersion"}\OperatorTok{);}
        \CommentTok{//Removing the last index of the version}
\NormalTok{        fileVersion }\OperatorTok{=}\NormalTok{ fileVersion}\OperatorTok{.}\FunctionTok{substring}\OperatorTok{(}\DecValTok{0}\OperatorTok{,}\NormalTok{ fileVersion}\OperatorTok{.}\FunctionTok{lastIndexOf}\OperatorTok{(}\CharTok{\textquotesingle{}.\textquotesingle{}}\OperatorTok{));}
    \OperatorTok{\}}
    \ControlFlowTok{return}\NormalTok{ fileVersion}\OperatorTok{;}
\OperatorTok{\}}
\end{Highlighting}
\end{Shaded}

io/src/main/java/org/openmarkov/io/probmodel/reader/ReaderFactory.java:24-34

\begin{Shaded}
\begin{Highlighting}[]
\KeywordTok{public} \DataTypeTok{static}\NormalTok{ PGMXReader\_0\_2 }\FunctionTok{getReader}\OperatorTok{(}\BuiltInTok{String}\NormalTok{ strVersion}\OperatorTok{)} \KeywordTok{throws}\NormalTok{ ProbNetParserException}\OperatorTok{.}\FunctionTok{WrongVersion} \OperatorTok{\{}
\NormalTok{    Version version }\OperatorTok{=} \BuiltInTok{Arrays}\OperatorTok{.}\FunctionTok{stream}\OperatorTok{(}\NormalTok{Version}\OperatorTok{.}\FunctionTok{values}\OperatorTok{())}
                            \OperatorTok{.}\FunctionTok{filter}\OperatorTok{(}
\NormalTok{                                    iteratorVersion }\OperatorTok{{-}\textgreater{}}\NormalTok{ strVersion}\OperatorTok{.}\FunctionTok{startsWith}\OperatorTok{(}\NormalTok{iteratorVersion}\OperatorTok{.}\FunctionTok{toString}\OperatorTok{()))}
                            \OperatorTok{.}\FunctionTok{findFirst}\OperatorTok{()}
                            \OperatorTok{.}\FunctionTok{orElse}\OperatorTok{(}\KeywordTok{null}\OperatorTok{);}
    \ControlFlowTok{return} \ControlFlowTok{switch} \OperatorTok{(}\NormalTok{version}\OperatorTok{)} \OperatorTok{\{}
        \ControlFlowTok{case}\NormalTok{ V02 }\OperatorTok{{-}\textgreater{}} \KeywordTok{new} \FunctionTok{PGMXReader\_0\_2}\OperatorTok{();}
        \ControlFlowTok{case}\NormalTok{ V10 }\OperatorTok{{-}\textgreater{}} \KeywordTok{new} \FunctionTok{PGMXReader\_1\_0}\OperatorTok{();}
        \ControlFlowTok{case} \KeywordTok{null} \OperatorTok{{-}\textgreater{}} \ControlFlowTok{throw} \KeywordTok{new}\NormalTok{ ProbNetParserException}\OperatorTok{.}\FunctionTok{WrongVersion}\OperatorTok{(}\NormalTok{strVersion}\OperatorTok{);}
    \OperatorTok{\};}
\OperatorTok{\}}
\end{Highlighting}
\end{Shaded}
\paragraph{Cómo arreglarlo}
Hacer que un fichero sin versión (o que no es ProbModelXML) produzca una excepción de lectura con mensaje, en los dos sitios donde se interpreta la versión:

\begin{itemize}
\tightlist
\item
  \texttt{Format\allowbreak{}Manager.\allowbreak{}get\allowbreak{}File\allowbreak{}Version}: si el atributo está vacío o no tiene punto, lanzar \texttt{Prob\allowbreak{}Net\allowbreak{}Parser\allowbreak{}Exception.\allowbreak{}Wrong\allowbreak{}Version} (o \texttt{Badly\allowbreak{}Structured\allowbreak{}File}) en vez de cortar la cadena; mejor aún, comprobar primero que el elemento raíz se llama \texttt{Prob\allowbreak{}Model\allowbreak{}XML}. La versión que calcula ni siquiera se usa para elegir lector: solo va en el mensaje de \texttt{No\allowbreak{}Reader\allowbreak{}For\allowbreak{}File\allowbreak{}Exception}.
\item
  \texttt{Reader\allowbreak{}Factory.\allowbreak{}get\allowbreak{}Reader}: tratar \texttt{str\allowbreak{}Version\ ==\ null} antes del \texttt{stream}, para que llegue al \texttt{Wrong\allowbreak{}Version} previsto.
\end{itemize}

La solución de fondo, reconocer el formato por el contenido y no solo por la extensión, es la que describe \hyperref[nota:173]{nota~173}; este arreglo local es independiente y no la bloquea.

Prueba de regresión: en \texttt{io/\allowbreak{}src/\allowbreak{}test}, abrir \texttt{net\allowbreak{}AB.\allowbreak{}xml} con \texttt{Format\allowbreak{}Manager.\allowbreak{}get\allowbreak{}Instance(\allowbreak{}).\allowbreak{}get\allowbreak{}Prob\allowbreak{}Net\allowbreak{}Reader(\allowbreak{}url)} y con \texttt{new\ PGMXReader(\allowbreak{}).\allowbreak{}read(\allowbreak{}url)} y comprobar que ambos lanzan una \texttt{Prob\allowbreak{}Net\allowbreak{}Parser\allowbreak{}Exception} (no una excepción de índice ni de nulo).

% ══════════════════════════════════════════════════════════════════════
\section{Bases de casos}\label{sec:casos}

Los ficheros de casos (CSV y Weka ARFF) con los que se aprende y se evalúa una red.

\setcounter{subsection}{117}
\subsection[Un «?» en una base de Weka descuadra las filas e inventa casos sin avisar]{Una base de casos de Weka con un valor desconocido («?») se lee con filas descuadradas y casos inventados, sin ningún aviso}\label{err:118}
\begin{ficha}
\fichakey{Estado}Presente
\fichakey{Módulo}io
\fichakey{Paquete}org.openmarkov.io.database.weka
\fichakey{Ficheros}\path{io/src/main/java/org/openmarkov/io/database/weka/ArffLexer.java:181-205} \newline \path{io/src/main/java/org/openmarkov/io/database/weka/ArffParser.java:517-550} \newline \path{io/src/main/java/org/openmarkov/io/database/weka/ArffParser.java:586-596} \newline \path{io/src/main/java/org/openmarkov/io/database/weka/ArffDataBaseIO.java:53-70} \newline \path{io/src/main/java/org/openmarkov/io/database/weka/arff.g:15,242-263}
\fichakey{Comprobado}Leyendo el código y ejecutándolo
\fichakey{Esfuerzo}Medio (días)
\fichakey{Decisión de equipo}No
\fichakey{Relacionados}\hyperref[err:119]{error~119}, \hyperref[err:120]{error~120}, \hyperref[err:121]{error~121}, \hyperref[err:122]{error~122}, \hyperref[err:123]{error~123}
\end{ficha}
\paragraph{Qué pasa}
El formato ARFF (Attribute-Relation File Format, el de la herramienta de aprendizaje automático Weka) marca un valor que no se conoce con una interrogación: \texttt{?}. Una base de casos con algún \texttt{?} se carga sin ningún error: \texttt{Arff\allowbreak{}Data\allowbreak{}Base\allowbreak{}IO.\allowbreak{}load} devuelve una base con filas descuadradas y casos que no están en el fichero, la ventana la acepta y se aprende o se evalúa con otros datos. El único aviso sale por la salida de errores.

OpenMarkov lee estos ficheros con \texttt{Arff\allowbreak{}Data\allowbreak{}Base\allowbreak{}IO} (módulo \texttt{io}), registrado como formato de base de casos «WekaDB» con la extensión \texttt{arff}. El usuario llega a él desde la ventana de aprender una red (botón de abrir la base de casos, que ofrece el filtro «WekaDB») y desde las herramientas de evaluación de redes del módulo \texttt{bn\allowbreak{}Evaluation} (panel \texttt{DBOpener\allowbreak{}Panel}, que usan la validación cruzada, la evaluación de una red y partir un banco de casos).

El lector está generado con ANTLR 2 (un generador de analizadores sintácticos) a partir de \texttt{arff.\allowbreak{}g}. El analizador léxico (\texttt{Arff\allowbreak{}Lexer}) está declarado como filtro (\texttt{options\ \{filter=true;\}}): todo carácter que no empieza ninguna de sus reglas se descarta en silencio (rama \texttt{default} de \texttt{next\allowbreak{}Token}, que hace \texttt{consume(\allowbreak{});\ continue\ try\allowbreak{}Again;}). La interrogación no está en ninguna regla (\texttt{IDENT} solo admite letras, dígitos y unos pocos signos), así que desaparece.

Además, los finales de línea se descartan también (regla \texttt{SEPARATOR} con \texttt{Token.\allowbreak{}SKIP}), de modo que el analizador sintáctico no sabe dónde acaba una fila: decide que un caso termina cuando detrás de un valor no viene una coma (\texttt{example(\allowbreak{})}, línea 527).

Con \texttt{a,\allowbreak{}?,\allowbreak{}e} el analizador recibe \texttt{a} \texttt{,\allowbreak{}} \texttt{,\allowbreak{}} \texttt{e}: la segunda coma no encaja, \texttt{report\allowbreak{}Error} escribe «line 7:4: unexpected token: ,» en la salida de errores, \texttt{recover} se salta elementos y la lectura continúa.

Medido con un pequeño programa de prueba, sobre tres columnas X\{a,b\}, Y\{c,d\}, Z\{e,f\}:

\begin{itemize}
\tightlist
\item
  Fichero: \texttt{a,\allowbreak{}c,\allowbreak{}e} / \texttt{a,\allowbreak{}?,\allowbreak{}e} / \texttt{b,\allowbreak{}d,\allowbreak{}f} / \texttt{b,\allowbreak{}d,\allowbreak{}e} (cuatro filas).
\item
  Base leída: \textbf{cinco casos}, \texttt{{[}0,\allowbreak{}0,\allowbreak{}0{]}\ {[}0,\allowbreak{}0,\allowbreak{}0{]}\ {[}0,\allowbreak{}0,\allowbreak{}0{]}\ {[}1,\allowbreak{}1,\allowbreak{}1{]}\ {[}1,\allowbreak{}1,\allowbreak{}0{]}}. La segunda fila del fichero (\texttt{a,\allowbreak{}d,\allowbreak{}e} sin el «?» sería \texttt{{[}0,\allowbreak{}1,\allowbreak{}0{]}}) no aparece y hay un caso inventado.
\item
  Con el «?» en la última columna (\texttt{a,\allowbreak{}c,\allowbreak{}e} / \texttt{a,\allowbreak{}c,\allowbreak{}?} / \texttt{b,\allowbreak{}d,\allowbreak{}f}): \textbf{dos casos}, \texttt{{[}0,\allowbreak{}0,\allowbreak{}0{]}\ {[}0,\allowbreak{}0,\allowbreak{}0{]}}; la tercera fila se pierde.
\end{itemize}

Del mismo mecanismo sale un segundo efecto, también silencioso: en \texttt{case\allowbreak{}Data(\allowbreak{})} (líneas 586-596), si un valor de la sección de datos no es ninguno de los estados declarados, \texttt{get\allowbreak{}State\allowbreak{}Index} devuelve -1 y el caso se queda con un 0, es decir, con el primer estado de esa columna.

El «?» es además la marca de valor ausente que usa el resto del proyecto: el preprocesado de aprendizaje (\texttt{learning.\allowbreak{}core}, \texttt{Discretization}) busca un estado llamado \texttt{?} para tratar los valores ausentes, y el lector de CSV (valores separados por comas) trata \texttt{?} como cualquier otro valor y lo convierte en un estado de la columna. El lector de Weka es el único que lo pierde.

El commit \texttt{c5f5014} añadió a \texttt{load} la comprobación de columnas repetidas (\texttt{Case\allowbreak{}Database\allowbreak{}Reader.\allowbreak{}with\allowbreak{}Distinct\allowbreak{}Variable\allowbreak{}Names}); no toca el análisis de los datos y no cambia este fallo.
\paragraph{El código}
io/src/main/java/org/openmarkov/io/database/weka/ArffLexer.java:181-205

\begin{Shaded}
\begin{Highlighting}[]
\KeywordTok{default}\OperatorTok{:}
    \ControlFlowTok{if} \OperatorTok{((}\FunctionTok{LA}\OperatorTok{(}\DecValTok{1}\OperatorTok{)} \OperatorTok{==} \CharTok{\textquotesingle{}"\textquotesingle{}} \OperatorTok{||} \FunctionTok{LA}\OperatorTok{(}\DecValTok{1}\OperatorTok{)} \OperatorTok{==} \CharTok{\textquotesingle{}\textbackslash{}\textquotesingle{}\textquotesingle{}}\OperatorTok{)} \OperatorTok{\&\&} \OperatorTok{(}\NormalTok{\_tokenSet\_0}\OperatorTok{.}\FunctionTok{member}\OperatorTok{(}\FunctionTok{LA}\OperatorTok{(}\DecValTok{2}\OperatorTok{))))} \OperatorTok{\{}
        \FunctionTok{mSTRING}\OperatorTok{(}\KeywordTok{true}\OperatorTok{);}
        \KeywordTok{...}
    \OperatorTok{\}} \ControlFlowTok{else} \OperatorTok{\{}
        \ControlFlowTok{if} \OperatorTok{(}\FunctionTok{LA}\OperatorTok{(}\DecValTok{1}\OperatorTok{)} \OperatorTok{==}\NormalTok{ EOF\_CHAR}\OperatorTok{)} \OperatorTok{\{}
            \FunctionTok{uponEOF}\OperatorTok{();}
\NormalTok{            \_returnToken }\OperatorTok{=} \FunctionTok{makeToken}\OperatorTok{(}\NormalTok{Token}\OperatorTok{.}\FunctionTok{EOF\_TYPE}\OperatorTok{);}
        \OperatorTok{\}} \ControlFlowTok{else} \OperatorTok{\{}
            \FunctionTok{consume}\OperatorTok{();}
            \ControlFlowTok{continue}\NormalTok{ tryAgain}\OperatorTok{;}
        \OperatorTok{\}}
    \OperatorTok{\}}
\end{Highlighting}
\end{Shaded}

io/src/main/java/org/openmarkov/io/database/weka/ArffParser.java:586-596

\begin{Shaded}
\begin{Highlighting}[]
\ControlFlowTok{try} \OperatorTok{\{}
\NormalTok{    Token lcommon }\OperatorTok{=}\NormalTok{ l1 }\OperatorTok{!=} \KeywordTok{null} \OperatorTok{?}\NormalTok{ l1 }\OperatorTok{:}\NormalTok{ l2}\OperatorTok{;}
    \DataTypeTok{int}\NormalTok{ stateIndex }\OperatorTok{=}\NormalTok{ fsVariables}\OperatorTok{.}\FunctionTok{get}\OperatorTok{(}\NormalTok{index}\OperatorTok{).}\FunctionTok{getStateIndex}\OperatorTok{(}\NormalTok{lcommon}\OperatorTok{.}\FunctionTok{getText}\OperatorTok{());}
    \ControlFlowTok{if} \OperatorTok{(}\NormalTok{stateIndex }\OperatorTok{!=} \OperatorTok{{-}}\DecValTok{1}\OperatorTok{)} \OperatorTok{\{}
\NormalTok{        example}\OperatorTok{[}\NormalTok{index}\OperatorTok{]} \OperatorTok{=}\NormalTok{ stateIndex}\OperatorTok{;}
    \OperatorTok{\}}
\NormalTok{    index}\OperatorTok{++;}
\OperatorTok{\}} \ControlFlowTok{catch} \OperatorTok{(}\BuiltInTok{IndexOutOfBoundsException}\NormalTok{ ignored}\OperatorTok{)} \OperatorTok{\{}
\OperatorTok{\}}
\end{Highlighting}
\end{Shaded}
\paragraph{Cómo arreglarlo}
El fichero \texttt{arff.\allowbreak{}g} no corresponde ya al Java generado: declara \texttt{package\ org.\allowbreak{}openmarkov.\allowbreak{}learning.\allowbreak{}io}, llama a \texttt{add\allowbreak{}Prob\allowbreak{}Node} y usa \texttt{prob\allowbreak{}Node.\allowbreak{}properties}, mientras que \texttt{Arff\allowbreak{}Parser.\allowbreak{}java} tiene \texttt{add\allowbreak{}Node} y \texttt{set\allowbreak{}Additional\allowbreak{}Properties}. Regenerar desde la gramática no es directo.

Hay dos caminos: corregir a mano el código generado, o sustituir el analizador de ANTLR por un lector escrito a mano. El formato es de líneas (cabecera \texttt{@relation}, \texttt{@attribute}, \texttt{@data}, y una fila por línea separada por comas), y un lector a mano resuelve a la vez este error, \hyperref[err:119]{error~119}, \hyperref[err:120]{error~120} y \hyperref[err:121]{error~121}.

Lo que el arreglo debe cumplir:

\begin{itemize}
\tightlist
\item
  Una fila del fichero es un caso: los finales de línea cuentan.
\item
  \texttt{?} se reconoce como valor ausente y se trata como en el resto del proyecto: un estado \texttt{?} en esa columna, como hace \texttt{CSVData\allowbreak{}Base\allowbreak{}IO}.
\item
  Un valor no declarado en la cabecera no se convierte en el estado 0 en silencio: o se añade como estado, o se rechaza el fichero con un mensaje que lo nombre.
\item
  Un error de análisis no se escribe en la salida de errores para seguir leyendo, sino que llega al usuario como excepción de análisis (\texttt{Parsing\allowbreak{}Source\allowbreak{}Exception.\allowbreak{}Could\allowbreak{}Not\allowbreak{}Parse\allowbreak{}Source\allowbreak{}Exception}, que \texttt{load} ya declara).
\end{itemize}

Prueba de regresión en \texttt{io/\allowbreak{}src/\allowbreak{}test/\allowbreak{}java/\allowbreak{}org/\allowbreak{}openmarkov/\allowbreak{}io/\allowbreak{}database/\allowbreak{}} (junto a \texttt{Repeated\allowbreak{}Column\allowbreak{}Names\allowbreak{}Are\allowbreak{}Refused\allowbreak{}Test}, que ya usa ficheros ARFF temporales): las cuatro filas del ejemplo dan cuatro casos, con el «?» en la columna que toca; un «?» en la última columna no pierde filas; un valor no declarado no da el estado 0.

\setcounter{subsection}{118}
\subsection[Un fichero de Weka en UTF-8 con tildes o «ñ» pierde estados y lee ceros]{Un fichero de Weka en UTF-8 con acentos o con «ñ» pierde estados y lee todos sus casos con ceros}\label{err:119}
\begin{ficha}
\fichakey{Estado}Presente
\fichakey{Módulo}io
\fichakey{Paquete}org.openmarkov.io.database.weka
\fichakey{Ficheros}\path{io/src/main/java/org/openmarkov/io/database/weka/ArffLexer.java:36} \newline \path{io/src/main/java/org/openmarkov/io/database/weka/ArffLexer.java:432-685} \newline \path{io/src/main/java/org/openmarkov/io/database/weka/ArffLexer.java:687-1000} \newline \path{io/src/main/java/org/openmarkov/io/database/weka/arff.g:36-45}
\fichakey{Comprobado}Leyendo el código y ejecutándolo
\fichakey{Esfuerzo}Medio (días)
\fichakey{Decisión de equipo}No
\fichakey{Relacionados}\hyperref[err:118]{error~118}, \hyperref[err:120]{error~120}, \hyperref[err:121]{error~121}, \hyperref[err:122]{error~122}, \hyperref[err:123]{error~123}
\end{ficha}
\paragraph{Qué pasa}
Un fichero de Weka (formato ARFF, Attribute-Relation File Format, extensión \texttt{arff}) guardado en UTF-8, la codificación de texto habitual hoy, cuyos estados llevan tildes o «ñ» se carga sin ningún error, pero con una columna que no tiene sus estados y con todos los casos en el estado 0. Se aprende o se evalúa con esos datos. Se llega desde la ventana de aprender (abrir base de casos, filtro «WekaDB») y desde las herramientas de \texttt{bn\allowbreak{}Evaluation} que abren un banco de casos.

El lector de bases de casos de Weka (\texttt{Arff\allowbreak{}Data\allowbreak{}Base\allowbreak{}IO}) lee el fichero \textbf{byte a byte}: el analizador léxico se construye con \texttt{new\ Byte\allowbreak{}Buffer(\allowbreak{}in)} (línea 36), que convierte cada byte en un carácter.

Sus reglas de palabras (\texttt{IDENT}) y de cadenas entre comillas (\texttt{STRING}) enumeran los caracteres aceptados, y entre ellos solo hay estas letras que no son ASCII (el juego básico de 128 caracteres del inglés): \texttt{á\ é\ í\ ó\ ú\ ü} y sus mayúsculas, escritas en el código como \texttt{á}, \texttt{é}, \texttt{í}, \texttt{ó}, \texttt{ú}, \texttt{ü}, \texttt{Á}\ldots{} \texttt{Ü}.

Cada una de esas letras ocupa un byte en Latin-1 (ISO-8859-1) y dos en UTF-8; en UTF-8 ninguno de los dos bytes está en la lista. Como el analizador léxico es de filtro, los bytes que no reconoce se descartan en silencio (ver \hyperref[err:118]{error~118}) y la palabra se corta.

Medido con pequeños programas de prueba:

\begin{itemize}
\tightlist
\item
  UTF-8, \texttt{@ATTRIBUTE\ W\ \{camión,\allowbreak{}seco\}} y datos \texttt{camión,\allowbreak{}a} / \texttt{seco,\allowbreak{}b} / \texttt{camión,\allowbreak{}b}: aviso «expecting RIGHTB, found `n'» por la salida de errores; la variable W queda con \textbf{un solo estado, \texttt{cami}}; salen \textbf{cinco casos} (en vez de tres), todos con W en el estado 0.
\item
  Latin-1, \texttt{camión}: se lee bien.
\item
  Latin-1 y UTF-8, \texttt{@ATTRIBUTE\ W\ \{niño,\allowbreak{}adulto\}}: W queda con el estado \texttt{ni}, tres casos en vez de dos, todos a 0. La «ñ» no está en la lista, de modo que falla en cualquier codificación; lo mismo cualquier otra letra fuera de la lista (\texttt{ç}, \texttt{à}\ldots).
\end{itemize}
\paragraph{El código}
io/src/main/java/org/openmarkov/io/database/weka/ArffLexer.java:35-37

\begin{Shaded}
\begin{Highlighting}[]
\KeywordTok{public} \FunctionTok{ArffLexer}\OperatorTok{(}\BuiltInTok{InputStream}\NormalTok{ in}\OperatorTok{)} \OperatorTok{\{}
    \KeywordTok{this}\OperatorTok{(}\KeywordTok{new} \BuiltInTok{ByteBuffer}\OperatorTok{(}\NormalTok{in}\OperatorTok{));}
\OperatorTok{\}}
\end{Highlighting}
\end{Shaded}

io/src/main/java/org/openmarkov/io/database/weka/arff.g:36-39 (la misma lista está expandida en \texttt{m\allowbreak{}IDENT} y \texttt{m\allowbreak{}STRING} de \texttt{Arff\allowbreak{}Lexer.\allowbreak{}java})

\begin{Verbatim}[breaklines,breakanywhere,fontsize=\small]
IDENT
    options {testLiterals=true;}
    :   ('a'..'z'|'A'..'Z'|'_'|'$'|'0'..'9')
 ('a'..'z'|'á'|'é'|'í'|'ó'|'ú'|'ü'|'A'..'Z'|'Á'|'É'|'Í'|'Ó'|'Ú'|'Ü'|'_'|'$'|'&'|'-'|'0'..'9'|'%'|'.')*;
\end{Verbatim}
\paragraph{Cómo arreglarlo}
Leer el fichero como texto en UTF-8 (un \texttt{Reader} con \texttt{Standard\allowbreak{}Charsets.\allowbreak{}UTF\_\allowbreak{}8}, no bytes sueltos) y no enumerar letras: un nombre es todo lo que no sea uno de los separadores del formato, que son todos ASCII (coma, llaves, comillas, espacio, tabulador, fin de línea, \texttt{\%} de comentario). Si se reescribe el lector a mano (ver \hyperref[err:118]{error~118}), esto va incluido.

El escritor \texttt{Arff\allowbreak{}Data\allowbreak{}Base\allowbreak{}IO.\allowbreak{}save} usa \texttt{new\ Output\allowbreak{}Stream\allowbreak{}Writer(\allowbreak{}stream)} sin codificación explícita; conviene fijarla también a UTF-8 para que lectura y escritura casen en cualquier máquina.

Prueba de regresión en \texttt{io/\allowbreak{}src/\allowbreak{}test/\allowbreak{}java/\allowbreak{}org/\allowbreak{}openmarkov/\allowbreak{}io/\allowbreak{}database/\allowbreak{}}: un fichero UTF-8 con \texttt{camión} y \texttt{niño} como estados conserva los nombres completos y asigna a cada caso su estado.

\setcounter{subsection}{119}
\subsection[Una columna numérica, de texto o de fecha en Weka da NullPointerException]{Un fichero de Weka con una columna numérica, de texto o de fecha no se puede abrir: falla con NullPointerException}\label{err:120}
\begin{ficha}
\fichakey{Estado}Presente
\fichakey{Módulo}io
\fichakey{Paquete}org.openmarkov.io.database.weka
\fichakey{Ficheros}\path{io/src/main/java/org/openmarkov/io/database/weka/ArffParser.java:148-159} \newline \path{io/src/main/java/org/openmarkov/io/database/weka/ArffParser.java:340-470} \newline \path{io/src/main/java/org/openmarkov/io/database/weka/arff.g:192,233}
\fichakey{Comprobado}Leyendo el código y ejecutándolo
\fichakey{Esfuerzo}Medio (días)
\fichakey{Decisión de equipo}\textcolor{decisionrojo}{\textbf{Sí.}} Qué hacer con una columna numérica (tomar como estados los valores vistos, discretizarla, o rechazar el fichero con un mensaje claro)
\fichakey{Relacionados}\hyperref[err:118]{error~118}, \hyperref[err:119]{error~119}, \hyperref[err:121]{error~121}, \hyperref[err:122]{error~122}, \hyperref[err:123]{error~123}
\end{ficha}
\paragraph{Qué pasa}
El formato ARFF (Attribute-Relation File Format) de Weka admite cuatro clases de columna: \texttt{numeric} (también \texttt{real} e \texttt{integer}), \texttt{string}, \texttt{date} y la lista de valores entre llaves (\texttt{\{a,\allowbreak{}b,\allowbreak{}c\}}). Un fichero con una columna de cualquiera de las tres primeras clases no se abre: en la ventana de aprender sale una excepción no controlada y la base no se carga.

Es el caso de muchos ficheros ARFF reales: los conjuntos de datos de ejemplo de Weka con variables continuas (por ejemplo, edad o presión) declaran \texttt{numeric} o \texttt{real}.

La gramática del lector de OpenMarkov (\texttt{Arff\allowbreak{}Data\allowbreak{}Base\allowbreak{}IO}, analizador generado \texttt{Arff\allowbreak{}Parser}) solo acepta la lista entre llaves: la regla \texttt{attribute} exige detrás del nombre la regla \texttt{states}, que empieza por \texttt{match(\allowbreak{}LEFTB)} (la llave de apertura).

Con cualquiera de las otras tres clases, \texttt{states(\allowbreak{})} lanza un error de reconocimiento que se imprime por la salida de errores («line 2:14: expecting LEFTB, found `numeric'») y se «recupera»; la lista de estados queda vacía. \texttt{add\allowbreak{}Node} (línea 148) solo crea la variable si hay estados (\texttt{if\ (\allowbreak{}states.\allowbreak{}length\ !=\ 0)}), así que llama a \texttt{prob\allowbreak{}Net.\allowbreak{}add\allowbreak{}Node(\allowbreak{}null,\allowbreak{}\ …)} y el constructor de \texttt{Node} falla con \texttt{Null\allowbreak{}Pointer\allowbreak{}Exception:\ Cannot\ invoke\ "Object.\allowbreak{}hash\allowbreak{}Code(\allowbreak{})"\ because\ "variable"\ is\ null}. La excepción es de ejecución y no es de las que \texttt{load} declara.

Medido con \texttt{@ATTRIBUTE\ N\ numeric} y \texttt{@ATTRIBUTE\ S\ string}: el aviso «expecting LEFTB, found `numeric'» y la \texttt{Null\allowbreak{}Pointer\allowbreak{}Exception}, con la pila \texttt{Node.\allowbreak{}\textless{}init\textgreater{}} ← \texttt{Prob\allowbreak{}Net.\allowbreak{}add\allowbreak{}Node} ← \texttt{Arff\allowbreak{}Parser.\allowbreak{}add\allowbreak{}Node} ← \texttt{Arff\allowbreak{}Parser.\allowbreak{}attribute}.

Las columnas numéricas son además las que la ventana de aprender sabe tratar: su panel de variables ofrece discretizar las variables cuyos estados son números. Con el lector de CSV (valores separados por comas), donde ninguna columna declara nada, los valores vistos se convierten en estados y esa discretización puede aplicarse; con el de Weka el fichero no llega a abrirse.
\paragraph{El código}
io/src/main/java/org/openmarkov/io/database/weka/ArffParser.java:148-159

\begin{Shaded}
\begin{Highlighting}[]
\KeywordTok{public} \BuiltInTok{Node} \FunctionTok{addNode}\OperatorTok{(}\BuiltInTok{LinkedHashMap}\OperatorTok{\textless{}}\BuiltInTok{String}\OperatorTok{,} \BuiltInTok{String}\OperatorTok{\textgreater{}}\NormalTok{ infoNode}\OperatorTok{,} \BuiltInTok{State}\OperatorTok{[]}\NormalTok{ states}\OperatorTok{,}\NormalTok{ NodeType nodeType}\OperatorTok{,} \BuiltInTok{String}\NormalTok{ nodeName}\OperatorTok{)} \OperatorTok{\{}
\NormalTok{    Variable fsVariable }\OperatorTok{=} \KeywordTok{null}\OperatorTok{;}
    \KeywordTok{...}
    \ControlFlowTok{if} \OperatorTok{(}\NormalTok{states}\OperatorTok{.}\FunctionTok{length} \OperatorTok{!=} \DecValTok{0}\OperatorTok{)} \OperatorTok{\{}
\NormalTok{        fsVariable }\OperatorTok{=} \KeywordTok{new} \FunctionTok{Variable}\OperatorTok{(}\NormalTok{nodeName}\OperatorTok{,}\NormalTok{ states}\OperatorTok{);}
    \OperatorTok{\}}
    \BuiltInTok{Node}\NormalTok{ node }\OperatorTok{=}\NormalTok{ probNet}\OperatorTok{.}\FunctionTok{addNode}\OperatorTok{(}\NormalTok{fsVariable}\OperatorTok{,}\NormalTok{ nodeType}\OperatorTok{);}
\NormalTok{    node}\OperatorTok{.}\FunctionTok{setAdditionalProperties}\OperatorTok{(}\NormalTok{infoNode}\OperatorTok{);}
    \ControlFlowTok{return}\NormalTok{ node}\OperatorTok{;}
\OperatorTok{\}}
\end{Highlighting}
\end{Shaded}

io/src/main/java/org/openmarkov/io/database/weka/arff.g:192,233

\begin{Verbatim}[breaklines,breakanywhere,fontsize=\small]
n:ATTRIBUTE^ (id1:IDENT! | id2:STRING!) states {
...
states : lb:LEFTB! state! (c:COMMA! state!)* rb:RIGHTB! {};
\end{Verbatim}
\paragraph{Cómo arreglarlo}
Hay que decidir qué significa en OpenMarkov una columna \texttt{numeric}, \texttt{string} o \texttt{date}. Opciones:

\begin{itemize}
\tightlist
\item
  \begin{enumerate}
  \def\labelenumi{(\alph{enumi})}
  \tightlist
  \item
    Tomar como estados los valores que aparecen en la sección de datos, que es lo que el lector de CSV hace con todas las columnas (\texttt{CSVData\allowbreak{}Base\allowbreak{}IO.\allowbreak{}load}, que además los ordena alfabéticamente) y deja la columna lista para la discretización de la ventana de aprender.
  \end{enumerate}
\item
  \begin{enumerate}
  \def\labelenumi{(\alph{enumi})}
  \setcounter{enumi}{1}
  \tightlist
  \item
    Rechazar el fichero con una \texttt{Parsing\allowbreak{}Source\allowbreak{}Exception} que nombre la columna y su clase.
  \end{enumerate}
\end{itemize}

En ningún caso debe salir una \texttt{Null\allowbreak{}Pointer\allowbreak{}Exception}. En el mismo cambio, \texttt{add\allowbreak{}Node} no debe llamar a \texttt{prob\allowbreak{}Net.\allowbreak{}add\allowbreak{}Node} con una variable nula.

Si se reescribe el lector a mano (ver \hyperref[err:118]{error~118}), el tratamiento de las cuatro clases entra en la lectura de la cabecera. Conviene aceptar también \texttt{real} e \texttt{integer}, sinónimos de \texttt{numeric} en el formato.

Prueba de regresión: un fichero con una columna \texttt{numeric} y otra \texttt{string} se abre (o se rechaza con la excepción declarada, según lo que se decida) y nunca lanza \texttt{Null\allowbreak{}Pointer\allowbreak{}Exception}.

\setcounter{subsection}{120}
\subsection[Weka con @Relation, @Attribute o @Data en mayúsculas mezcladas no se abre]{Un fichero de Weka que escribe \texttt{@Relation}, \texttt{@Attribute} o \texttt{@Data} con mayúsculas y minúsculas mezcladas no se abre}\label{err:121}
\begin{ficha}
\fichakey{Estado}Presente
\fichakey{Módulo}io
\fichakey{Paquete}org.openmarkov.io.database.weka
\fichakey{Ficheros}\path{io/src/main/java/org/openmarkov/io/database/weka/ArffLexer.java:1020-1097} \newline \path{io/src/main/java/org/openmarkov/io/database/weka/ArffLexer.java:186-194} \newline \path{io/src/main/java/org/openmarkov/io/database/weka/ArffDataBaseIO.java:56-59} \newline \path{io/src/main/java/org/openmarkov/io/database/weka/arff.g:47-51}
\fichakey{Comprobado}Leyendo el código y ejecutándolo
\fichakey{Esfuerzo}Pequeño (horas)
\fichakey{Decisión de equipo}No
\fichakey{Relacionados}\hyperref[err:118]{error~118}, \hyperref[err:119]{error~119}, \hyperref[err:120]{error~120}, \hyperref[err:122]{error~122}, \hyperref[err:123]{error~123}
\end{ficha}
\paragraph{Qué pasa}
En el formato ARFF (Attribute-Relation File Format) de Weka las palabras de cabecera \texttt{@relation}, \texttt{@attribute} y \texttt{@data} no distinguen mayúsculas de minúsculas; en la práctica se ven escritas \texttt{@relation}, \texttt{@RELATION}, \texttt{@Relation}\ldots{} Un fichero que escriba \texttt{@Relation}, \texttt{@Attribute} o \texttt{@Data} no se abre: en la ventana de aprender o en las herramientas de evaluación de \texttt{bn\allowbreak{}Evaluation} sale una excepción no controlada y la base no se carga.

El analizador léxico del lector de OpenMarkov (\texttt{Arff\allowbreak{}Lexer}, generado desde \texttt{arff.\allowbreak{}g}) solo acepta dos formas de cada palabra, todo en minúsculas o todo en mayúsculas: \texttt{m\allowbreak{}RELATION} mira la segunda letra y, si es \texttt{r}, exige \texttt{"@relation"}; si es \texttt{R}, exige \texttt{"@RELATION"}.

Con \texttt{@Relation}, \texttt{match(\allowbreak{}"@RELATION")} falla en la tercera letra y el error se imprime por la salida de errores. El analizador sintáctico abandona la regla \texttt{relation(\allowbreak{})} sin rellenar su resultado, \texttt{parser.\allowbreak{}relation(\allowbreak{})} devuelve \texttt{null}, y \texttt{Arff\allowbreak{}Data\allowbreak{}Base\allowbreak{}IO.\allowbreak{}load} falla en la línea 59 al hacer \texttt{io\allowbreak{}Net.\allowbreak{}get(\allowbreak{}"Prob\allowbreak{}Net")}. La excepción es de ejecución y \texttt{load} no la declara.

Medido con \texttt{@Relation\ r}: «line 1:4: expecting RELATION, found `lation'» y \texttt{Null\allowbreak{}Pointer\allowbreak{}Exception:\ Cannot\ invoke\ "java.\allowbreak{}util.\allowbreak{}Hash\allowbreak{}Map.\allowbreak{}get(\allowbreak{}Object)"\ because\ "this.\allowbreak{}io\allowbreak{}Net"\ is\ null}. Igual con \texttt{@Attribute} o \texttt{@Data}.
\paragraph{El código}
io/src/main/java/org/openmarkov/io/database/weka/ArffLexer.java:1028-1037

\begin{Shaded}
\begin{Highlighting}[]
\ControlFlowTok{if} \OperatorTok{((}\FunctionTok{LA}\OperatorTok{(}\DecValTok{1}\OperatorTok{)} \OperatorTok{==} \CharTok{\textquotesingle{}@\textquotesingle{}}\OperatorTok{)} \OperatorTok{\&\&} \OperatorTok{(}\FunctionTok{LA}\OperatorTok{(}\DecValTok{2}\OperatorTok{)} \OperatorTok{==} \CharTok{\textquotesingle{}r\textquotesingle{}}\OperatorTok{))} \OperatorTok{\{}
    \FunctionTok{match}\OperatorTok{(}\StringTok{"@relation"}\OperatorTok{);}
\OperatorTok{\}} \ControlFlowTok{else} \ControlFlowTok{if} \OperatorTok{((}\FunctionTok{LA}\OperatorTok{(}\DecValTok{1}\OperatorTok{)} \OperatorTok{==} \CharTok{\textquotesingle{}@\textquotesingle{}}\OperatorTok{)} \OperatorTok{\&\&} \OperatorTok{(}\FunctionTok{LA}\OperatorTok{(}\DecValTok{2}\OperatorTok{)} \OperatorTok{==} \CharTok{\textquotesingle{}R\textquotesingle{}}\OperatorTok{))} \OperatorTok{\{}
    \FunctionTok{match}\OperatorTok{(}\StringTok{"@RELATION"}\OperatorTok{);}
\OperatorTok{\}} \ControlFlowTok{else} \OperatorTok{\{}
    \ControlFlowTok{throw} \KeywordTok{new} \FunctionTok{NoViableAltForCharException}\OperatorTok{((}\DataTypeTok{char}\OperatorTok{)} \FunctionTok{LA}\OperatorTok{(}\DecValTok{1}\OperatorTok{),} \FunctionTok{getFilename}\OperatorTok{(),} \FunctionTok{getLine}\OperatorTok{(),} \FunctionTok{getColumn}\OperatorTok{());}
\OperatorTok{\}}
\end{Highlighting}
\end{Shaded}

io/src/main/java/org/openmarkov/io/database/weka/ArffDataBaseIO.java:56-59

\begin{Shaded}
\begin{Highlighting}[]
\ControlFlowTok{try} \OperatorTok{(}\BuiltInTok{FileInputStream}\NormalTok{ fileStream }\OperatorTok{=} \KeywordTok{new} \BuiltInTok{FileInputStream}\OperatorTok{(}\NormalTok{file}\OperatorTok{))} \OperatorTok{\{}
\NormalTok{    ArffParser parser }\OperatorTok{=} \KeywordTok{new} \FunctionTok{ArffParser}\OperatorTok{(}\KeywordTok{new} \FunctionTok{ArffLexer}\OperatorTok{(}\NormalTok{fileStream}\OperatorTok{));}
\NormalTok{    ioNet }\OperatorTok{=}\NormalTok{ parser}\OperatorTok{.}\FunctionTok{relation}\OperatorTok{();}
\NormalTok{    ProbNet probNet }\OperatorTok{=} \OperatorTok{(}\NormalTok{ProbNet}\OperatorTok{)}\NormalTok{ ioNet}\OperatorTok{.}\FunctionTok{get}\OperatorTok{(}\StringTok{"ProbNet"}\OperatorTok{);}
\end{Highlighting}
\end{Shaded}
\paragraph{Cómo arreglarlo}
Reconocer las tres palabras sin distinguir mayúsculas: en ANTLR 2 (el generador con que está hecho el analizador), con \texttt{case\allowbreak{}Sensitive=false} limitado a esas reglas; si se reescribe el lector a mano (ver \hyperref[err:118]{error~118}), comparando en minúsculas.

Además, \texttt{load} debe comprobar si \texttt{parser.\allowbreak{}relation(\allowbreak{})} devuelve \texttt{null} y lanzar \texttt{Parsing\allowbreak{}Source\allowbreak{}Exception.\allowbreak{}Could\allowbreak{}Not\allowbreak{}Parse\allowbreak{}Source\allowbreak{}Exception} en vez de dejar que salga una \texttt{Null\allowbreak{}Pointer\allowbreak{}Exception}; eso protege también de cualquier otro error de cabecera.

Prueba de regresión: \texttt{@Relation}, \texttt{@Attribute} y \texttt{@Data} en un mismo fichero se leen igual que en minúsculas; una cabecera ilegible produce la excepción declarada.

\setcounter{subsection}{121}
\subsection[Una base con una columna de números guardada en Weka no se puede reabrir]{Una base de casos con una columna de números, guardada en formato Weka, ya no la puede abrir el propio OpenMarkov}\label{err:122}
\begin{ficha}
\fichakey{Estado}Presente en parte. No hay nada corregido por un commit; la otra parte revisada no es un fallo: la cuenta de comas con que el escritor evita declarar \texttt{?} como estado (líneas 114-129) da \texttt{\{a,\allowbreak{}b\}} con el \texttt{?} al principio, en medio o al final (medido), y solo fallaría con dos estados \texttt{?} en una misma variable, que no pueden existir. El nombre de relación con la ruta completa se trata en \hyperref[err:123]{error~123}.
\fichakey{Módulo}io
\fichakey{Paquete}org.openmarkov.io.database.weka
\fichakey{Ficheros}\path{io/src/main/java/org/openmarkov/io/database/weka/ArffDataBaseIO.java:80-132} \newline \path{io/src/main/java/org/openmarkov/io/database/weka/ArffDataBaseIO.java:137-156}
\fichakey{Comprobado}Leyendo el código y ejecutándolo
\fichakey{Esfuerzo}Pequeño (horas)
\fichakey{Decisión de equipo}No
\fichakey{Relacionados}\hyperref[err:118]{error~118}, \hyperref[err:119]{error~119}, \hyperref[err:120]{error~120}, \hyperref[err:121]{error~121}
\end{ficha}
\paragraph{Qué pasa}
Una base de casos que tenga una variable cuyos estados se llaman con números (\texttt{0}, \texttt{1}, \texttt{2}\ldots, cosa habitual), guardada en formato «WekaDB», no la puede volver a abrir el propio OpenMarkov.

Se llega desde las ventanas que guardan bases de casos y ofrecen el formato «WekaDB» en su selector (\texttt{DBWriter\allowbreak{}OMFile\allowbreak{}Chooser}): el generador de bases de casos (\texttt{DBGenerator\allowbreak{}GUI}, módulo \texttt{db\allowbreak{}Generator}), el preprocesado de datos (\texttt{Data\allowbreak{}Preprocessing\allowbreak{}Dialog}) y partir un banco de casos (\texttt{Split\allowbreak{}Dataset\allowbreak{}Dialog}), estos dos últimos en \texttt{bn\allowbreak{}Evaluation}.

\texttt{Arff\allowbreak{}Data\allowbreak{}Base\allowbreak{}IO.\allowbreak{}save} (módulo \texttt{io}) escribe la base en formato ARFF (Attribute-Relation File Format, el de la herramienta Weka). Antes de escribir los estados de cada variable mira si todos sus nombres (salvo \texttt{?}) son números enteros; si lo son, decide que la columna es numérica y escribe \textbf{a la vez} la palabra \texttt{numeric} y la lista de valores: \texttt{@ATTRIBUTE\ N\ numeric\ \{1,\allowbreak{}2\}}. En el formato ARFF una columna es \texttt{numeric} o es una lista de valores entre llaves, no las dos cosas.

El lector del propio OpenMarkov rechaza esa línea, porque su gramática exige una llave detrás del nombre (ver \hyperref[err:120]{error~120}): «expecting LEFTB, found `numeric'» y \texttt{Null\allowbreak{}Pointer\allowbreak{}Exception} al crear el nodo.

Medido: una base con N\{1,2\} y X\{a,b\}, guardada con \texttt{save} y vuelta a abrir con \texttt{load}, no se abre.

Un segundo efecto del mismo escritor: si un caso tiene en alguna columna el estado \texttt{?} (la marca de valor ausente del proyecto), la sección de datos escribe \texttt{?} y el lector descuadra las filas (ver \hyperref[err:118]{error~118}). Medido: tres casos escritos, dos leídos.
\paragraph{El código}
io/src/main/java/org/openmarkov/io/database/weka/ArffDataBaseIO.java:94-112

\begin{Shaded}
\begin{Highlighting}[]
\DataTypeTok{boolean}\NormalTok{ numeric }\OperatorTok{=} \KeywordTok{true}\OperatorTok{;}
\ControlFlowTok{for} \OperatorTok{(}\DataTypeTok{int}\NormalTok{ i }\OperatorTok{=} \DecValTok{0}\OperatorTok{;}\NormalTok{ i }\OperatorTok{\textless{}}\NormalTok{ states}\OperatorTok{.}\FunctionTok{length}\OperatorTok{;}\NormalTok{ i}\OperatorTok{++)} \OperatorTok{\{}
    \ControlFlowTok{try} \OperatorTok{\{}
        \ControlFlowTok{if} \OperatorTok{(!}\NormalTok{states}\OperatorTok{[}\NormalTok{i}\OperatorTok{].}\FunctionTok{getName}\OperatorTok{().}\FunctionTok{equals}\OperatorTok{(}\StringTok{"?"}\OperatorTok{))} \OperatorTok{\{}
            \BuiltInTok{Integer}\OperatorTok{.}\FunctionTok{parseInt}\OperatorTok{(}\NormalTok{states}\OperatorTok{[}\NormalTok{i}\OperatorTok{].}\FunctionTok{getName}\OperatorTok{());}
        \OperatorTok{\}}
    \OperatorTok{\}} \ControlFlowTok{catch} \OperatorTok{(}\BuiltInTok{NumberFormatException}\NormalTok{ e}\OperatorTok{)} \OperatorTok{\{}
\NormalTok{        numeric }\OperatorTok{=} \KeywordTok{false}\OperatorTok{;}
    \OperatorTok{\}}
\OperatorTok{\}}
\ControlFlowTok{if} \OperatorTok{(}\NormalTok{numeric}\OperatorTok{)} \OperatorTok{\{}
\NormalTok{    output}\OperatorTok{.}\FunctionTok{write}\OperatorTok{(}\StringTok{"numeric \{"}\OperatorTok{);}
\OperatorTok{\}} \ControlFlowTok{else} \OperatorTok{\{}
\NormalTok{    output}\OperatorTok{.}\FunctionTok{write}\OperatorTok{(}\StringTok{"\{"}\OperatorTok{);}
\OperatorTok{\}}
\end{Highlighting}
\end{Shaded}
\paragraph{Cómo arreglarlo}
Escribir siempre la lista de valores entre llaves, sin \texttt{numeric}: los estados de una variable de OpenMarkov son una lista finita, y eso es lo que representa \texttt{\{…\}}. Aunque el lector llegue a aceptar columnas \texttt{numeric} (ver \hyperref[err:120]{error~120}), la línea \texttt{numeric\ \{…\}} seguiría siendo inválida para cualquier otro programa que lea ARFF, incluido Weka.

El valor ausente debe escribirse y leerse con el mismo convenio en los dos lados: si el escritor no declara \texttt{?} como estado y escribe \texttt{?} en los datos, el lector tiene que reconocer \texttt{?} como ausente (ver \hyperref[err:118]{error~118}).

Prueba de regresión de ida y vuelta en \texttt{io/\allowbreak{}src/\allowbreak{}test/\allowbreak{}java/\allowbreak{}org/\allowbreak{}openmarkov/\allowbreak{}io/\allowbreak{}database/\allowbreak{}}: una base con una variable de estados \texttt{1},\texttt{2} y otra con un caso en \texttt{?}, guardada y vuelta a leer, da las mismas variables, los mismos estados y los mismos casos.

\setcounter{subsection}{122}
\subsection[Un fichero de Weka en una carpeta con tilde, «ñ» o «+» no se puede reabrir]{Un fichero de Weka guardado en una carpeta cuyo nombre lleva acentos, «ñ» o «+» no lo puede volver a abrir OpenMarkov}\label{err:123}
\begin{ficha}
\fichakey{Estado}Presente
\fichakey{Módulo}io
\fichakey{Paquete}org.openmarkov.io.database.weka
\fichakey{Ficheros}\path{io/src/main/java/org/openmarkov/io/database/weka/ArffDataBaseIO.java:83-85} \newline \path{io/src/main/java/org/openmarkov/io/database/weka/arff.g:41-45}
\fichakey{Comprobado}Leyendo el código y ejecutándolo
\fichakey{Esfuerzo}Pequeño (horas)
\fichakey{Decisión de equipo}No
\fichakey{Relacionados}\hyperref[err:118]{error~118}, \hyperref[err:119]{error~119}, \hyperref[err:120]{error~120}, \hyperref[err:121]{error~121}
\end{ficha}
\paragraph{Qué pasa}
Una base de casos guardada en formato «WekaDB» en una carpeta cuyo nombre lleva tilde, «ñ», \texttt{+} u otros signos no se puede volver a abrir: la lectura termina con \texttt{Null\allowbreak{}Pointer\allowbreak{}Exception}. Se llega guardando la base desde el generador de bases de casos, el preprocesado de datos o partir un banco de casos, en cualquier carpeta con uno de esos caracteres (en castellano, cualquier carpeta con tilde o eñe).

Es otro defecto del mismo escritor que \hyperref[err:122]{error~122}, esta vez en el nombre de la relación que escribe \texttt{Arff\allowbreak{}Data\allowbreak{}Base\allowbreak{}IO.\allowbreak{}save}.

La primera línea de un fichero ARFF (Attribute-Relation File Format, el formato de Weka), \texttt{@RELATION\ nombre}, da nombre al conjunto de datos. El escritor usa la \textbf{ruta absoluta} del fichero: \texttt{@RELATION\ "/\allowbreak{}home/\allowbreak{}usuario/\allowbreak{}Investigación/\allowbreak{}casos.\allowbreak{}arff"}. Además de dejar escrita en el fichero la estructura de carpetas del usuario, eso hace que el nombre dependa de caracteres que el lector no admite.

La regla de cadenas entre comillas del lector (\texttt{STRING}) solo acepta letras ASCII (el juego básico de 128 caracteres del inglés), \texttt{á\ é\ í\ ó\ ú\ ü} en un byte (codificación Latin-1) y una lista corta de signos, entre los que no están \texttt{+}, \texttt{\textasciitilde{}}, \texttt{\#}, \texttt{ñ} ni ninguna letra acentuada escrita en UTF-8 (ver \hyperref[err:119]{error~119}).

Medido guardando la misma base de dos variables en tres carpetas:

\begin{itemize}
\tightlist
\item
  \texttt{…/\allowbreak{}plain/\allowbreak{}casos.\allowbreak{}arff}: se vuelve a leer (aparte del defecto del \texttt{?} de \hyperref[err:118]{error~118}).
\item
  \texttt{…/\allowbreak{}Investigación/\allowbreak{}casos.\allowbreak{}arff}: «unexpected token: casos.arff» y \texttt{Null\allowbreak{}Pointer\allowbreak{}Exception:\ …\ because\ "this.\allowbreak{}cases\allowbreak{}Aux"\ is\ null}.
\item
  \texttt{…/\allowbreak{}a+b/\allowbreak{}casos.\allowbreak{}arff}: «unexpected token: ``» y la misma \texttt{Null\allowbreak{}Pointer\allowbreak{}Exception}.
\end{itemize}
\paragraph{El código}
io/src/main/java/org/openmarkov/io/database/weka/ArffDataBaseIO.java:80-85

\begin{Shaded}
\begin{Highlighting}[]
\AttributeTok{@Override} \KeywordTok{public} \DataTypeTok{void} \FunctionTok{save}\OperatorTok{(}\BuiltInTok{File}\NormalTok{ file}\OperatorTok{,}\NormalTok{ CaseDatabase database}\OperatorTok{)} \KeywordTok{throws} \BuiltInTok{IOException} \OperatorTok{\{}
    \BuiltInTok{FileOutputStream}\NormalTok{ stream }\OperatorTok{=} \KeywordTok{new} \BuiltInTok{FileOutputStream}\OperatorTok{(}\NormalTok{file}\OperatorTok{);}
    \BuiltInTok{OutputStreamWriter}\NormalTok{ output }\OperatorTok{=} \KeywordTok{new} \BuiltInTok{OutputStreamWriter}\OperatorTok{(}\NormalTok{stream}\OperatorTok{);}
    \CommentTok{/* Relation name */}
    \BuiltInTok{String}\NormalTok{ relationName }\OperatorTok{=}\NormalTok{ file}\OperatorTok{.}\FunctionTok{getAbsolutePath}\OperatorTok{();}
\NormalTok{    output}\OperatorTok{.}\FunctionTok{write}\OperatorTok{(}\StringTok{"}\SpecialCharTok{\textbackslash{}n}\StringTok{@RELATION }\SpecialCharTok{\textbackslash{}"}\StringTok{"} \OperatorTok{+}\NormalTok{ relationName }\OperatorTok{+} \StringTok{"}\SpecialCharTok{\textbackslash{}"\textbackslash{}n}\StringTok{"}\OperatorTok{);}
\end{Highlighting}
\end{Shaded}
\paragraph{Cómo arreglarlo}
Usar como nombre de relación el nombre del fichero sin extensión (o un nombre neutro), escapando comillas si las lleva, y no la ruta absoluta.

Del lado del lector, aceptar en una cadena entre comillas cualquier carácter salvo la comilla de cierre (ver \hyperref[err:119]{error~119}). Cualquiera de las dos cosas por separado ya evita el fallo; juntas hacen que la ida y vuelta no dependa de la carpeta.

Prueba de regresión: guardar una base en una carpeta temporal cuyo nombre lleve tilde, eñe y \texttt{+}, y volver a leerla con el mismo número de variables y casos; y comprobar que el fichero no contiene la ruta de la carpeta.

\setcounter{subsection}{123}
\subsection[Una línea de CSV con campos de menos, de más o en blanco inventa casos]{Una línea de un CSV con campos de menos, de más o en blanco se lee como un caso inventado o tumba la lectura}\label{err:124}
\begin{ficha}
\fichakey{Estado}Presente
\fichakey{Módulo}io
\fichakey{Paquete}org.openmarkov.io.database.excel
\fichakey{Ficheros}\path{io/src/main/java/org/openmarkov/io/database/excel/CSVDataBaseIO.java:106-165} \newline \path{io/src/main/java/org/openmarkov/io/database/excel/CSVDataBaseIO.java:233-261}
\fichakey{Comprobado}Leyendo el código y ejecutándolo
\fichakey{Esfuerzo}Pequeño (horas)
\fichakey{Decisión de equipo}No
\end{ficha}
\paragraph{Qué pasa}
Una base de casos en formato CSV (valores separados por comas) con alguna línea que no trae tantos campos como columnas tiene la cabecera, o con una línea en blanco, se lee mal. Si a la línea le faltan campos o está en blanco, el usuario aprende una red (o carga evidencia) con casos que el fichero no dice, sin ningún aviso. Si le sobran campos, o si una columna no la rellena ninguna línea, recibe una excepción genérica que no señala la línea culpable.

El lector, \texttt{CSVData\allowbreak{}Base\allowbreak{}IO}, se da de alta con \texttt{@Case\allowbreak{}Database\allowbreak{}Format(\allowbreak{}extension\ =\ "csv")} y lo usan todos los sitios que abren una base de casos a través de \texttt{Case\allowbreak{}Database\allowbreak{}Manager}:

\begin{itemize}
\tightlist
\item
  el diálogo de aprendizaje (\texttt{Learning\allowbreak{}Dialog.\allowbreak{}load\allowbreak{}Case\allowbreak{}File}, que se abre desde el menú de herramientas de aprendizaje);
\item
  «cargar evidencia» del editor (\texttt{Network\allowbreak{}File\allowbreak{}Handler.\allowbreak{}load\allowbreak{}Evidence});
\item
  el generador de bases de casos (\texttt{DBGenerator\allowbreak{}GUI});
\item
  los diálogos de evaluación de redes (\texttt{DBOpener\allowbreak{}Panel}, \texttt{Split\allowbreak{}Dataset\allowbreak{}Dialog}, \texttt{Data\allowbreak{}Preprocessing\allowbreak{}Dialog}).
\end{itemize}

El lector toma la primera línea como la lista de nombres de las columnas y cada línea siguiente como un caso. Para cada línea, \texttt{get\allowbreak{}Data\allowbreak{}Line} reserva una fila de enteros del ancho de la cabecera (un entero por columna, que es el número de estado) y recorre \textbf{los campos que traiga la línea}, sin compararlos con el número de columnas. Luego \texttt{load} ordena alfabéticamente los estados de cada columna y renumera los casos. Nada avisa de que una línea no tiene la forma esperada.

Las cuatro consecuencias, ejecutadas con un pequeño programa de prueba:

\begin{enumerate}
\def\labelenumi{\arabic{enumi}.}
\item
  \textbf{Línea con campos de menos.} Las columnas a las que la línea no llega se quedan con el valor inicial del array, 0. Ese 0 no significa «falta el dato»: es el número del primer estado que se vio en esa columna. Con el fichero \texttt{A,\allowbreak{}B,\allowbreak{}C\ /\allowbreak{}\ yes,\allowbreak{}no,\allowbreak{}yes\ /\allowbreak{}\ no,\allowbreak{}no\ /\allowbreak{}\ yes,\allowbreak{}yes,\allowbreak{}no}, el segundo caso sale \texttt{no\ no\ yes}: la C que la línea no traía se lee como «yes», porque «yes» fue lo primero que apareció en la columna C. Qué valor se inventa depende del orden de las líneas del fichero.
\item
  \textbf{Línea con campos de más.} \texttt{variables\allowbreak{}States.\allowbreak{}get(\allowbreak{}num\allowbreak{}Variable)} se sale de la lista y la lectura muere con \texttt{Index\allowbreak{}Out\allowbreak{}Of\allowbreak{}Bounds\allowbreak{}Exception:\ Index\ 3\ out\ of\ bounds\ for\ length\ 3}, sin decir qué línea del fichero es ni cuántos campos traía.
\item
  \textbf{Línea en blanco.} No se salta: cuenta como un caso. Su único campo, la cadena vacía, se da de alta como un estado más de la primera columna, y las demás columnas se rellenan con el 0 del punto 1. Con \texttt{A,\allowbreak{}B\ /\allowbreak{}\ yes,\allowbreak{}no\ /\allowbreak{}\ (\allowbreak{}línea\ vacía)\ /\allowbreak{}\ no,\allowbreak{}yes} salen tres casos, la variable A queda con tres estados (\texttt{""}, \texttt{no}, \texttt{yes}) y el caso inventado es \texttt{""\ no}. Basta con que el fichero termine con una línea vacía de más para que ocurra.
\item
  \textbf{Columna que ninguna línea llega a rellenar} (caso extremo del punto 1: todas las líneas traen menos campos que la cabecera). Esa columna se queda sin ningún estado y la renumeración de \texttt{load} indexa un array de longitud cero: \texttt{Array\allowbreak{}Index\allowbreak{}Out\allowbreak{}Of\allowbreak{}Bounds\allowbreak{}Exception:\ Index\ 0\ out\ of\ bounds\ for\ length\ 0} en \texttt{CSVData\allowbreak{}Base\allowbreak{}IO.\allowbreak{}java:149}.
\end{enumerate}
\paragraph{El código}
io/src/main/java/org/openmarkov/io/database/excel/CSVDataBaseIO.java:128-131

\begin{Shaded}
\begin{Highlighting}[]
\ControlFlowTok{while} \OperatorTok{(}\NormalTok{scanner}\OperatorTok{.}\FunctionTok{hasNextLine}\OperatorTok{())} \OperatorTok{\{}
\NormalTok{    data}\OperatorTok{.}\FunctionTok{add}\OperatorTok{(}\FunctionTok{getDataLine}\OperatorTok{(}\NormalTok{scanner}\OperatorTok{.}\FunctionTok{nextLine}\OperatorTok{(),}\NormalTok{ variablesStates}\OperatorTok{,}\NormalTok{ variablesStatesNames}\OperatorTok{));}
\NormalTok{    numRows}\OperatorTok{++;}
\OperatorTok{\}}
\end{Highlighting}
\end{Shaded}

io/src/main/java/org/openmarkov/io/database/excel/CSVDataBaseIO.java:233-259

\begin{Shaded}
\begin{Highlighting}[]
\KeywordTok{private} \DataTypeTok{static} \DataTypeTok{int}\OperatorTok{[]} \FunctionTok{getDataLine}\OperatorTok{(}\BuiltInTok{String}\NormalTok{ line}\OperatorTok{,} \BuiltInTok{List}\OperatorTok{\textless{}}\BuiltInTok{Map}\OperatorTok{\textless{}}\BuiltInTok{String}\OperatorTok{,} \BuiltInTok{Integer}\OperatorTok{\textgreater{}\textgreater{}}\NormalTok{ variablesStates}\OperatorTok{,}
                                 \BuiltInTok{List}\OperatorTok{\textless{}}\BuiltInTok{List}\OperatorTok{\textless{}}\BuiltInTok{String}\OperatorTok{\textgreater{}\textgreater{}}\NormalTok{ variablesStatesNames}\OperatorTok{)} \OperatorTok{\{}
    \DataTypeTok{int}\OperatorTok{[]}\NormalTok{ statesLines }\OperatorTok{=} \KeywordTok{new} \DataTypeTok{int}\OperatorTok{[}\NormalTok{variablesStates}\OperatorTok{.}\FunctionTok{size}\OperatorTok{()];}
    \KeywordTok{...}
    \BuiltInTok{String}\OperatorTok{[]}\NormalTok{ fields }\OperatorTok{=}\NormalTok{ line}\OperatorTok{.}\FunctionTok{split}\OperatorTok{(}\StringTok{"[,;]"}\OperatorTok{,} \OperatorTok{{-}}\DecValTok{1}\OperatorTok{);}
    \ControlFlowTok{for} \OperatorTok{(}\DataTypeTok{int}\NormalTok{ numVariable }\OperatorTok{=} \DecValTok{0}\OperatorTok{;}\NormalTok{ numVariable }\OperatorTok{\textless{}}\NormalTok{ fields}\OperatorTok{.}\FunctionTok{length}\OperatorTok{;}\NormalTok{ numVariable}\OperatorTok{++)} \OperatorTok{\{}
        \BuiltInTok{String}\NormalTok{ stateVariable }\OperatorTok{=}\NormalTok{ fields}\OperatorTok{[}\NormalTok{numVariable}\OperatorTok{];}
        \KeywordTok{...}
        \BuiltInTok{Map}\OperatorTok{\textless{}}\BuiltInTok{String}\OperatorTok{,} \BuiltInTok{Integer}\OperatorTok{\textgreater{}}\NormalTok{ variableStates }\OperatorTok{=}\NormalTok{ variablesStates}\OperatorTok{.}\FunctionTok{get}\OperatorTok{(}\NormalTok{numVariable}\OperatorTok{);}
        \BuiltInTok{Integer}\NormalTok{ stateNumber }\OperatorTok{=}\NormalTok{ variableStates}\OperatorTok{.}\FunctionTok{get}\OperatorTok{(}\NormalTok{stateVariable}\OperatorTok{);}
        \ControlFlowTok{if} \OperatorTok{(}\NormalTok{stateNumber }\OperatorTok{==} \KeywordTok{null}\OperatorTok{)} \OperatorTok{\{}
\NormalTok{            stateNumber }\OperatorTok{=}\NormalTok{ variableStates}\OperatorTok{.}\FunctionTok{size}\OperatorTok{();}
\NormalTok{            variableStates}\OperatorTok{.}\FunctionTok{put}\OperatorTok{(}\NormalTok{stateVariable}\OperatorTok{,}\NormalTok{ stateNumber}\OperatorTok{);}
\NormalTok{            variablesStatesNames}\OperatorTok{.}\FunctionTok{get}\OperatorTok{(}\NormalTok{numVariable}\OperatorTok{).}\FunctionTok{add}\OperatorTok{(}\NormalTok{stateVariable}\OperatorTok{);}
        \OperatorTok{\}}
\NormalTok{        statesLines}\OperatorTok{[}\NormalTok{numVariable}\OperatorTok{]} \OperatorTok{=}\NormalTok{ stateNumber}\OperatorTok{;}
    \OperatorTok{\}}
    \ControlFlowTok{return}\NormalTok{ statesLines}\OperatorTok{;}
\OperatorTok{\}}
\end{Highlighting}
\end{Shaded}
\paragraph{Cómo arreglarlo}
En \texttt{load}, antes de llamar a \texttt{get\allowbreak{}Data\allowbreak{}Line}, decidir qué hacer con cada línea según su número de campos, llevando la cuenta del número de línea del fichero para los mensajes:

\begin{itemize}
\tightlist
\item
  línea vacía (o solo espacios): saltarla sin contarla como caso;
\item
  número de campos distinto del de la cabecera: rechazar el fichero con una \texttt{Parsing\allowbreak{}Source\allowbreak{}Exception} (o la excepción de lectura que ya declara el método) cuyo mensaje diga el número de línea, cuántos campos traía y cuántos esperaba.
\end{itemize}

Rellenar los campos que faltan con un estado de «valor desconocido» sería otra opción, pero cambia lo que el aprendizaje recibe y conviene no inventarlo sin avisar.

Es el único sitio: los lectores de Excel del mismo paquete (\texttt{Excel\allowbreak{}Data\allowbreak{}Base\allowbreak{}IO}, \texttt{Old\allowbreak{}Excel\allowbreak{}Data\allowbreak{}Base\allowbreak{}IO}) recorren las columnas de la cabecera y no los campos de la fila, así que no tienen esta forma del defecto.

Prueba de regresión: ampliar \texttt{CSVData\allowbreak{}Base\allowbreak{}IOEmpty\allowbreak{}Fields\allowbreak{}Test} (\texttt{io/\allowbreak{}src/\allowbreak{}test/\allowbreak{}java/\allowbreak{}org/\allowbreak{}openmarkov/\allowbreak{}io/\allowbreak{}database/\allowbreak{}excel/\allowbreak{}}) con cuatro ficheros ---línea corta, línea larga, línea en blanco intermedia y línea en blanco final--- y comprobar que las líneas en blanco no crean casos ni estados y que las otras dos producen una excepción cuyo mensaje contiene el número de línea.

% ══════════════════════════════════════════════════════════════════════
\section{Aprendizaje y evaluación de redes}\label{sec:aprendizaje}

El aprendizaje de estructura y parámetros (PC, maximización de la esperanza, clasificadores) y las herramientas que dividen, discretizan y evalúan bases de casos.

\setcounter{subsection}{124}
\subsection[La exactitud del Bayes ingenuo selectivo usa el caso evaluado al entrenar]{La exactitud del Bayes ingenuo selectivo se calcula con el caso evaluado dentro de los datos de entrenamiento, y con muchas variables predice siempre la primera clase}\label{err:125}
\begin{ficha}
\fichakey{Estado}Presente
\fichakey{Módulo}learning.metric
\fichakey{Paquete}org.openmarkov.learning.metric.cmi.accuracy
\fichakey{Ficheros}\path{learning.metric/src/main/java/org/openmarkov/learning/metric/cmi/accuracy/Accuracy.java:130-142} \newline \path{learning.metric/src/main/java/org/openmarkov/learning/metric/cmi/accuracy/Accuracy.java:172-185} \newline \path{learning.metric/src/main/java/org/openmarkov/learning/metric/cmi/util/MetricUtils.java:95-107}
\fichakey{Comprobado}Leyendo el código y ejecutándolo
\fichakey{Esfuerzo}Pequeño (horas)
\fichakey{Decisión de equipo}No
\fichakey{Relacionados}\hyperref[err:126]{error~126}, \hyperref[err:128]{error~128}, \hyperref[err:131]{error~131}, \hyperref[err:135]{error~135}
\end{ficha}
\paragraph{Qué pasa}
El usuario llega por el menú de aprendizaje → tipo discriminativo → «Selective naive bayes» (hacia delante o hacia atrás). La exactitud (proporción de casos cuya clase se acierta) es la que decide qué característica se añade o se quita, y se calcula mal: el clasificador se evalúa con casos que ha visto al entrenar, y con muchas características predice siempre la primera clase.

En modo Bayes ingenuo simple, el que usa «Selective naive bayes», \texttt{Accuracy.\allowbreak{}compute\allowbreak{}NBNet\allowbreak{}Accuracy} recorre las diez ventanas de prueba y para cada caso llama a \texttt{predict\allowbreak{}Class\allowbreak{}Value(\allowbreak{}row,\allowbreak{}\ variables)}. Ese método no recibe la iteración a la que pertenece el caso: para cada característica multiplica un factor por \textbf{cada uno de los diez conjuntos de entrenamiento} (\texttt{cross\allowbreak{}Tabs.\allowbreak{}get(\allowbreak{}variable){[}it{]}} con \texttt{it} de 0 a 9).

El caso evaluado solo está excluido del entrenamiento de su propia ventana, así que está dentro de los demás. Es justo lo que la división en pliegues pretendía evitar.

Además, el producto no es una probabilidad. Cada factor es una cuenta de un conjunto de entrenamiento (unos nueve décimos de los casos) dividida por la frecuencia de la clase en \textbf{toda} la base, y se multiplican diez por característica, con lo que las características pesan como si estuvieran elevadas a la décima potencia frente a la probabilidad a priori de la clase.

Y no se trabaja con logaritmos: con muchas características el producto baja por debajo del menor \texttt{double} positivo para todas las clases, y \texttt{Metric\allowbreak{}Utils.\allowbreak{}get\allowbreak{}Index\allowbreak{}Of\allowbreak{}Max\allowbreak{}Value}, que empieza con máximo 0,0 y exige «mayor que», devuelve siempre el índice 0.

Medido:

\begin{itemize}
\tightlist
\item
  Sobre \texttt{learning.\allowbreak{}algorithm/\allowbreak{}src/\allowbreak{}test/\allowbreak{}resources/\allowbreak{}network/\allowbreak{}BN-asia10k.\allowbreak{}csv}: de las 10 000 evaluaciones, 2 073 casos están en 7 de los 10 conjuntos de entrenamiento, 4 618 en 8 y 3 309 en 9 (por el solapamiento de las ventanas que describe \hyperref[err:131]{error~131}).
\item
  Con una base sintética de 2 000 casos, clase binaria equilibrada y características binarias que dependen de la clase (probabilidad 0,7 frente a 0,3): con 20 características la exactitud calculada es 0,9625; con 60, 0,9995; con 100 baja a 0,902; con 150 predice la clase 0 en los 2 000 casos y la exactitud es 0,4825 (la proporción de la clase 0).
\item
  Sobre \texttt{BN-asia10k.\allowbreak{}csv}, con sus 7 características, el efecto es pequeño: la exactitud calculada es 0,843 frente a 0,8445 de una evaluación correcta con las mismas ventanas, y difieren 49 predicciones de 10 000.
\end{itemize}

La otra rama de la clase, \texttt{predict\allowbreak{}Class\allowbreak{}Value\allowbreak{}Augmented\allowbreak{}NB} (la del superpadre), usa solo el conjunto de entrenamiento de la iteración que corresponde (\texttt{it}), así que no tiene este fallo (tiene otros: ver \hyperref[err:127]{error~127} y \hyperref[err:128]{error~128}).
\paragraph{El código}
learning.metric/src/main/java/org/openmarkov/learning/metric/cmi/accuracy/Accuracy.java:172-185

\begin{Shaded}
\begin{Highlighting}[]
\KeywordTok{public} \DataTypeTok{int} \FunctionTok{predictClassValue}\OperatorTok{(}\DataTypeTok{int}\OperatorTok{[]}\NormalTok{ row}\OperatorTok{,} \BuiltInTok{Collection}\OperatorTok{\textless{}}\NormalTok{Variable}\OperatorTok{\textgreater{}}\NormalTok{ variables}\OperatorTok{)\{}
    \DataTypeTok{double}\OperatorTok{[]}\NormalTok{ probs }\OperatorTok{=} \FunctionTok{getProbRootNode}\OperatorTok{();}
    \ControlFlowTok{for} \OperatorTok{(}\BuiltInTok{String}\NormalTok{ variable }\OperatorTok{:}\NormalTok{ variables}\OperatorTok{.}\FunctionTok{stream}\OperatorTok{().}\FunctionTok{map}\OperatorTok{(}\NormalTok{Variable}\OperatorTok{::}\NormalTok{getName}\OperatorTok{).}\FunctionTok{toList}\OperatorTok{())} \OperatorTok{\{}
        \DataTypeTok{int}\NormalTok{ index }\OperatorTok{=} \FunctionTok{getIndexVariable}\OperatorTok{(}\NormalTok{variable}\OperatorTok{);}
        \ControlFlowTok{for} \OperatorTok{(}\DataTypeTok{int}\NormalTok{ it}\OperatorTok{=}\DecValTok{0}\OperatorTok{;}\NormalTok{ it}\OperatorTok{\textless{}}\NormalTok{KFOLD}\OperatorTok{;}\NormalTok{ it}\OperatorTok{++)\{}
            \ControlFlowTok{for}\OperatorTok{(}\DataTypeTok{int}\NormalTok{ j}\OperatorTok{=}\DecValTok{0}\OperatorTok{;}\NormalTok{ j}\OperatorTok{\textless{}}\NormalTok{freqRootNode}\OperatorTok{.}\FunctionTok{length}\OperatorTok{;}\NormalTok{j}\OperatorTok{++)\{}
\NormalTok{                probs}\OperatorTok{[}\NormalTok{j}\OperatorTok{]*=} \OperatorTok{((}\NormalTok{crossTabs}\OperatorTok{.}\FunctionTok{get}\OperatorTok{(}\NormalTok{variable}\OperatorTok{)[}\NormalTok{it}\OperatorTok{][}\NormalTok{row}\OperatorTok{[}\NormalTok{index}\OperatorTok{]][}\NormalTok{j}\OperatorTok{]} \OperatorTok{+}\NormalTok{ alpha}\OperatorTok{)/(}\NormalTok{freqRootNode}\OperatorTok{[}\NormalTok{j}\OperatorTok{]} \OperatorTok{+}\NormalTok{ alpha}\OperatorTok{*}\NormalTok{ variables}\OperatorTok{.}\FunctionTok{size}\OperatorTok{()));}
            \OperatorTok{\}}
        \OperatorTok{\}}
    \OperatorTok{\}}
    \ControlFlowTok{return} \FunctionTok{getIndexOfMaxValue}\OperatorTok{(}\NormalTok{probs}\OperatorTok{);}
\OperatorTok{\}}
\end{Highlighting}
\end{Shaded}
\paragraph{Cómo arreglarlo}
Pasar a \texttt{predict\allowbreak{}Class\allowbreak{}Value} la iteración del caso (como ya hace \texttt{predict\allowbreak{}Class\allowbreak{}Value\allowbreak{}Augmented\allowbreak{}NB}) y usar solo \texttt{cross\allowbreak{}Tabs.\allowbreak{}get(\allowbreak{}variable){[}it{]}}. Dividir por la frecuencia de la clase en ese mismo conjunto de entrenamiento, y tomar la probabilidad a priori también de ese conjunto (hoy \texttt{get\allowbreak{}Prob\allowbreak{}Root\allowbreak{}Node} usa toda la base).

Sumar logaritmos en vez de multiplicar probabilidades, o normalizar tras cada factor, para que el producto no se anule.

El suavizado usa \texttt{alpha\ *\ variables.\allowbreak{}size(\allowbreak{})} en el denominador; lo habitual es \texttt{alpha} por el número de estados de la característica, pero eso es otra cuestión (y hoy \texttt{alpha} vale 0 en esta rama: ver \hyperref[err:126]{error~126}).

Prueba de regresión: con la base sintética de 150 características descrita arriba, la exactitud debe ser próxima a la de un Bayes ingenuo evaluado por separado, no la proporción de la clase 0; y con una característica que identifica cada caso (un estado distinto por caso), la exactitud no debe superar la de adivinar.

\setcounter{subsection}{125}
\subsection[El Bayes ingenuo selectivo elige características sin el alfa del diálogo]{El Bayes ingenuo selectivo elige las características sin el suavizado que el usuario pone en el diálogo: la métrica trabaja siempre con alfa 0}\label{err:126}
\begin{ficha}
\fichakey{Estado}Presente
\fichakey{Módulo}learning.algorithm
\fichakey{Paquete}org.openmarkov.learning.algorithm.nbderived.snb
\fichakey{Ficheros}\path{learning.algorithm/src/main/java/org/openmarkov/learning/algorithm/nbderived/snb/SelectiveNBAlgorithm.java:40-50} \newline \path{learning.metric/src/main/java/org/openmarkov/learning/metric/cmi/accuracy/Accuracy.java:42} \newline \path{learning.metric/src/main/java/org/openmarkov/learning/metric/cmi/accuracy/Accuracy.java:172-185} \newline \path{learning.algorithm/src/main/java/org/openmarkov/learning/algorithm/nbderived/spnb/SuperParentNBAlgorithm.java:59-65} \newline \path{learning.gui/src/main/java/org/openmarkov/learning/algorithm/nbderived/snb/gui/SelectiveNBParametersDialog.java:45-51}
\fichakey{Comprobado}Leyendo el código y ejecutándolo
\fichakey{Esfuerzo}Pequeño (horas)
\fichakey{Decisión de equipo}No
\fichakey{Relacionados}\hyperref[err:125]{error~125}, \hyperref[err:128]{error~128}, \hyperref[err:135]{error~135}
\end{ficha}
\paragraph{Qué pasa}
El usuario llega por el menú de aprendizaje → tipo discriminativo → «Selective naive bayes» → opciones, donde el diálogo ofrece «Alpha parameter (Laplace-like correction)» con 0,5 por omisión, y también por la validación cruzada de \texttt{bn\allowbreak{}Evaluation}. Ese valor no interviene en qué características conserva o quita el algoritmo: se eligen con alfa 0. El alfa solo se aplica después, al calcular las tablas de la red resultante.

El alfa es el suavizado de Laplace: se suma a cada cuenta para que una combinación de valores que no aparece en los datos no dé probabilidad cero.

«Selective naive bayes» elige las características con la métrica \texttt{Accuracy} (exactitud: proporción de casos cuya clase se acierta), que tiene su propio campo \texttt{alpha}, inicializado a 0 y que solo cambia con \texttt{set\allowbreak{}Alpha}. \texttt{Selective\allowbreak{}NBParameters\allowbreak{}Dialog.\allowbreak{}get\allowbreak{}Instance} pasa el alfa al constructor del algoritmo, que lo guarda en \texttt{Learning\allowbreak{}Algorithm.\allowbreak{}alpha} para el aprendizaje de parámetros. \texttt{Selective\allowbreak{}NBAlgorithm.\allowbreak{}init} le pasa a la métrica la variable clase, pero no el alfa. \texttt{Super\allowbreak{}Parent\allowbreak{}NBAlgorithm.\allowbreak{}init}, que usa la misma métrica, sí llama a \texttt{set\allowbreak{}Alpha(\allowbreak{}alpha)}.

Con alfa 0, en \texttt{Accuracy.\allowbreak{}predict\allowbreak{}Class\allowbreak{}Value} cada factor es cuenta / frecuencia de la clase: un valor de una característica que no aparece con una clase en uno de los diez conjuntos de entrenamiento anula esa clase para el caso. Eso cambia qué características parecen útiles.

Medido ejecutando «Selective naive bayes» hacia atrás (quitando características desde el Bayes ingenuo completo; hacia delante el resultado no es repetible, ver \hyperref[err:135]{error~135}), con alfa 0,5 en el algoritmo, primero como hoy y después llamando también a \texttt{set\allowbreak{}Alpha(\allowbreak{}0,\allowbreak{}5)} en la métrica:

\begin{itemize}
\tightlist
\item
  \texttt{core/\allowbreak{}src/\allowbreak{}test/\allowbreak{}resources/\allowbreak{}nets/\allowbreak{}Cataract-RRDdataset\allowbreak{}Pre\allowbreak{}Proc3.\allowbreak{}csv} (537 casos), clase \texttt{success}: hoy quita las seis características y deja la clase sola; con el alfa en la métrica quita solo \texttt{loc57} y conserva \texttt{break\allowbreak{}Type}, \texttt{loc48}, \texttt{loc48rrd}, \texttt{loc57rrd} y \texttt{op\allowbreak{}Type}.
\item
  \texttt{integration\allowbreak{}Tests/\allowbreak{}src/\allowbreak{}main/\allowbreak{}resources/\allowbreak{}networks/\allowbreak{}learning/\allowbreak{}alarm500.\allowbreak{}csv} (500 casos, 36 características): con clase \texttt{HYPOVOLEMIA} hoy conserva 8 y con el alfa en la métrica 28; con clase \texttt{INTUBATION}, 7 frente a 12; con clase \texttt{BP}, 21 frente a 29.
\item
  \texttt{learning.\allowbreak{}algorithm/\allowbreak{}src/\allowbreak{}test/\allowbreak{}resources/\allowbreak{}network/\allowbreak{}BN-asia10k.\allowbreak{}csv} (clase \texttt{Bronchitis}), \texttt{integration\allowbreak{}Tests/\allowbreak{}src/\allowbreak{}main/\allowbreak{}resources/\allowbreak{}networks/\allowbreak{}learning/\allowbreak{}learn\allowbreak{}Test\allowbreak{}Data\allowbreak{}Base.\allowbreak{}csv} (clase \texttt{A}) y \texttt{learning.\allowbreak{}algorithm/\allowbreak{}src/\allowbreak{}test/\allowbreak{}resources/\allowbreak{}network/\allowbreak{}Head\allowbreak{}To\allowbreak{}Head2.\allowbreak{}csv} (clase \texttt{C}): el mismo resultado en los dos casos.
\end{itemize}
\paragraph{El código}
learning.algorithm/src/main/java/org/openmarkov/learning/algorithm/nbderived/snb/SelectiveNBAlgorithm.java:40-50

\begin{Shaded}
\begin{Highlighting}[]
\AttributeTok{@Override} \KeywordTok{public} \DataTypeTok{void} \FunctionTok{init}\OperatorTok{(}\NormalTok{ModelNetUse modelNetUse}\OperatorTok{)} \OperatorTok{\{}
    \ControlFlowTok{if} \OperatorTok{(}\NormalTok{metric }\KeywordTok{instanceof}\NormalTok{ Accuracy}\OperatorTok{)} \OperatorTok{\{}
        \OperatorTok{((}\NormalTok{Accuracy}\OperatorTok{)}\NormalTok{ metric}\OperatorTok{).}\FunctionTok{setClassVariable}\OperatorTok{(}\KeywordTok{this}\OperatorTok{.}\FunctionTok{classVariableName}\OperatorTok{);}
    \OperatorTok{\}}
    \ControlFlowTok{if} \OperatorTok{(!}\NormalTok{forward}\OperatorTok{)} \OperatorTok{\{}
        \FunctionTok{setRelationsForRootVariable}\OperatorTok{();}
    \OperatorTok{\}}
    \KeywordTok{...}
\OperatorTok{\}}
\end{Highlighting}
\end{Shaded}

learning.algorithm/src/main/java/org/openmarkov/learning/algorithm/nbderived/spnb/SuperParentNBAlgorithm.java:59-65

\begin{Shaded}
\begin{Highlighting}[]
\AttributeTok{@Override} \KeywordTok{public} \DataTypeTok{void} \FunctionTok{init}\OperatorTok{(}\NormalTok{ModelNetUse modelNetUse}\OperatorTok{)} \OperatorTok{\{}
    \ControlFlowTok{if} \OperatorTok{(}\NormalTok{metric }\KeywordTok{instanceof}\NormalTok{ Accuracy}\OperatorTok{)} \OperatorTok{\{}
        \OperatorTok{((}\NormalTok{Accuracy}\OperatorTok{)}\NormalTok{ metric}\OperatorTok{).}\FunctionTok{setClassVariable}\OperatorTok{(}\KeywordTok{this}\OperatorTok{.}\FunctionTok{classVariableName}\OperatorTok{);}
        \OperatorTok{((}\NormalTok{Accuracy}\OperatorTok{)}\NormalTok{ metric}\OperatorTok{).}\FunctionTok{setAugmentedNet}\OperatorTok{(}\KeywordTok{true}\OperatorTok{);}
        \OperatorTok{((}\NormalTok{Accuracy}\OperatorTok{)}\NormalTok{ metric}\OperatorTok{).}\FunctionTok{setAlpha}\OperatorTok{(}\NormalTok{alpha}\OperatorTok{);}
    \OperatorTok{\}}
\end{Highlighting}
\end{Shaded}
\paragraph{Cómo arreglarlo}
Llamar a \texttt{(\allowbreak{}(\allowbreak{}Accuracy)\ metric).\allowbreak{}set\allowbreak{}Alpha(\allowbreak{}alpha)} en \texttt{Selective\allowbreak{}NBAlgorithm.\allowbreak{}init}, como ya hace \texttt{Super\allowbreak{}Parent\allowbreak{}NBAlgorithm.\allowbreak{}init}.

Al activarse el suavizado empieza a contar la forma del denominador de \texttt{predict\allowbreak{}Class\allowbreak{}Value}, \texttt{freq\allowbreak{}Root\allowbreak{}Node{[}j{]}\ +\ alpha\ *\ variables.\allowbreak{}size(\allowbreak{})}, que multiplica el alfa por el número de características en vez de por el número de estados de la característica (ver \hyperref[err:125]{error~125}); conviene revisarlo a la vez.

Los únicos algoritmos que usan \texttt{Accuracy} son estos dos, y en \texttt{Super\allowbreak{}Parent\allowbreak{}NBAlgorithm} la regla ya se cumple.

Prueba de regresión: tras \texttt{init}, el alfa de la métrica debe ser el del algoritmo; y con \texttt{Cataract-RRDdataset\allowbreak{}Pre\allowbreak{}Proc3.\allowbreak{}csv}, clase \texttt{success}, hacia atrás, el modelo aprendido con alfa 0 y con alfa 0,5 no debe ser el mismo (hoy, con los dos, la clase sola).

\setcounter{subsection}{126}
\subsection[El superpadre pesa sus enlaces con cuentas de casillas equivocadas]{El clasificador de superpadre pesa sus enlaces con cuentas leídas en casillas equivocadas de la tabla de tres variables}\label{err:127}
\begin{ficha}
\fichakey{Estado}Presente
\fichakey{Módulo}learning.metric
\fichakey{Paquete}org.openmarkov.learning.metric.cmi.accuracy
\fichakey{Ficheros}\path{learning.metric/src/main/java/org/openmarkov/learning/metric/cmi/accuracy/Accuracy.java:292-310} \newline \path{learning.metric/src/main/java/org/openmarkov/learning/metric/cmi/accuracy/Accuracy.java:193-222} \newline \path{learning.metric/src/main/java/org/openmarkov/learning/metric/cmi/util/MetricUtils.java:173-211}
\fichakey{Comprobado}Leyendo el código y ejecutándolo
\fichakey{Esfuerzo}Pequeño (horas)
\fichakey{Decisión de equipo}No
\fichakey{Relacionados}\hyperref[err:128]{error~128}, \hyperref[err:129]{error~129}, \hyperref[err:130]{error~130}, \hyperref[err:131]{error~131}, \hyperref[err:132]{error~132}, \hyperref[err:133]{error~133}, \hyperref[nota:8]{nota~8}
\end{ficha}
\paragraph{Qué pasa}
El usuario llega por el menú de aprendizaje → tipo discriminativo → «Superparent naive bayes» (con la casilla del mismo superpadre marcada o no), y también desde la validación cruzada de \texttt{bn\allowbreak{}Evaluation}. El algoritmo elige el superpadre y los enlaces que cuelgan de él con probabilidades que no corresponden a los datos.

En modo «red aumentada» (el que activa \texttt{Super\allowbreak{}Parent\allowbreak{}NBAlgorithm.\allowbreak{}init} con \texttt{set\allowbreak{}Augmented\allowbreak{}Net(\allowbreak{}true)}), la métrica \texttt{Accuracy} necesita, para cada pareja de características (n1, n2), cuántos casos de entrenamiento tienen cada combinación (valor de n1, valor de n2, valor de la clase).

\texttt{build2nd\allowbreak{}Level\allowbreak{}Prob\allowbreak{}Distribution} pide esas cuentas a \texttt{Metric\allowbreak{}Utils.\allowbreak{}build\allowbreak{}Absolute\allowbreak{}Freq\allowbreak{}Cross\allowbreak{}Tab}, que devuelve un \texttt{Table\allowbreak{}Potential} con variables (n1, clase, n2), rellenado según la convención de OpenMarkov: la primera variable cambia en el paso más corto, es decir, casilla = x1 + s1·(c + sc·x2) (s = número de estados).

Después lo reparte en \texttt{cross\allowbreak{}Tab{[}it{]}{[}i{]}{[}j{]}{[}k{]}} cortando el array como si la casilla fuera k + sc·(j + s2·i): n1 en el paso más largo, n2 en medio y la clase en el más corto. Las dimensiones cuadran, así que no hay error de índice, pero cada \texttt{cross\allowbreak{}Tab{[}it{]}{[}i{]}{[}j{]}{[}k{]}} contiene la cuenta de otra combinación.

Esas cuentas son las que \texttt{get\allowbreak{}Prob(\allowbreak{}tail,\allowbreak{}\ head,\allowbreak{}\ row,\allowbreak{}\ it)} usa para estimar P(tail \textbar{} head, clase), y con ellas \texttt{predict\allowbreak{}Class\allowbreak{}Value\allowbreak{}Augmented\allowbreak{}NB} calcula la exactitud con la que \texttt{Super\allowbreak{}Parent\allowbreak{}NBAlgorithm} elige el superpadre y sus enlaces.

Medido sobre \texttt{learning.\allowbreak{}algorithm/\allowbreak{}src/\allowbreak{}test/\allowbreak{}resources/\allowbreak{}network/\allowbreak{}BN-asia10k.\allowbreak{}csv}, primer conjunto de entrenamiento, n1 = \texttt{Dyspnea}, n2 = \texttt{Smoker}, clase = \texttt{Bronchitis}, comparando lo leído con las cuentas hechas directamente sobre los casos: 6 de las 8 casillas no coinciden. Por ejemplo, (0,0,1) lee 340 y son 248; (0,1,0) lee 248 y son 1 489; (1,0,0) lee 1 489 y son 340. Solo coinciden (0,0,0) y (1,1,1).

\texttt{predict\allowbreak{}Class\allowbreak{}Value\allowbreak{}Augmented\allowbreak{}NB}, el método que consume estas cuentas, tiene además otro defecto, descrito en \hyperref[err:128]{error~128}.
\paragraph{El código}
learning.metric/src/main/java/org/openmarkov/learning/metric/cmi/accuracy/Accuracy.java:292-310

\begin{Shaded}
\begin{Highlighting}[]
\KeywordTok{protected} \DataTypeTok{double}\OperatorTok{[][][][]} \FunctionTok{build2ndLevelProbDistribution}\OperatorTok{(}\BuiltInTok{Node}\NormalTok{ n1}\OperatorTok{,} \BuiltInTok{Node}\NormalTok{ n2}\OperatorTok{)\{}
    \DataTypeTok{int}\NormalTok{ rootStates }\OperatorTok{=} \FunctionTok{getRootNode}\OperatorTok{().}\FunctionTok{getVariable}\OperatorTok{().}\FunctionTok{getNumStates}\OperatorTok{();}
    \DataTypeTok{int}\NormalTok{ fstStates }\OperatorTok{=}\NormalTok{ n1}\OperatorTok{.}\FunctionTok{getVariable}\OperatorTok{().}\FunctionTok{getNumStates}\OperatorTok{();}
    \DataTypeTok{int}\NormalTok{ sndStates }\OperatorTok{=}\NormalTok{ n2}\OperatorTok{.}\FunctionTok{getVariable}\OperatorTok{().}\FunctionTok{getNumStates}\OperatorTok{();}
    \DataTypeTok{double}\OperatorTok{[][][][]}\NormalTok{ crossTab }\OperatorTok{=} \KeywordTok{new} \DataTypeTok{double}\OperatorTok{[}\NormalTok{KFOLD}\OperatorTok{][}\NormalTok{fstStates}\OperatorTok{][}\NormalTok{sndStates}\OperatorTok{][}\NormalTok{rootStates}\OperatorTok{];}
    \ControlFlowTok{for}\OperatorTok{(}\DataTypeTok{int}\NormalTok{ it}\OperatorTok{=}\DecValTok{0}\OperatorTok{;}\NormalTok{ it}\OperatorTok{\textless{}}\NormalTok{KFOLD}\OperatorTok{;}\NormalTok{ it}\OperatorTok{++)\{}
\NormalTok{        TablePotential potential }\OperatorTok{=} \FunctionTok{buildAbsoluteFreqCrossTab}\OperatorTok{(}\NormalTok{probNet}\OperatorTok{,}\NormalTok{ caseDatabase}\OperatorTok{,}\NormalTok{ n1}\OperatorTok{,}\NormalTok{ n2}\OperatorTok{,} \FunctionTok{getRootNode}\OperatorTok{(),}\NormalTok{  dataset}\OperatorTok{.}\FunctionTok{getTraining}\OperatorTok{()[}\NormalTok{it}\OperatorTok{]);}
        \ControlFlowTok{for}\OperatorTok{(}\DataTypeTok{int}\NormalTok{ i}\OperatorTok{=}\DecValTok{0}\OperatorTok{;}\NormalTok{ i}\OperatorTok{\textless{}}\NormalTok{fstStates}\OperatorTok{;}\NormalTok{i}\OperatorTok{++)\{}
            \DataTypeTok{double}\OperatorTok{[][]}\NormalTok{ tmp }\OperatorTok{=} \KeywordTok{new} \DataTypeTok{double}\OperatorTok{[}\NormalTok{sndStates}\OperatorTok{][}\NormalTok{rootStates}\OperatorTok{];}
            \ControlFlowTok{for}\OperatorTok{(}\DataTypeTok{int}\NormalTok{ j}\OperatorTok{=}\DecValTok{0}\OperatorTok{;}\NormalTok{ j}\OperatorTok{\textless{}}\NormalTok{sndStates}\OperatorTok{;}\NormalTok{ j}\OperatorTok{++)\{}
\NormalTok{                tmp}\OperatorTok{[}\NormalTok{j}\OperatorTok{]} \OperatorTok{=} \BuiltInTok{Arrays}\OperatorTok{.}\FunctionTok{copyOfRange}\OperatorTok{(}\NormalTok{potential}\OperatorTok{.}\FunctionTok{getValues}\OperatorTok{(),} \OperatorTok{(}\NormalTok{i }\OperatorTok{*}\NormalTok{ sndStates }\OperatorTok{+}\NormalTok{ j}\OperatorTok{)} \OperatorTok{*}\NormalTok{ rootStates}\OperatorTok{,} \OperatorTok{(}\NormalTok{i }\OperatorTok{*}\NormalTok{ sndStates }\OperatorTok{+}\NormalTok{ j }\OperatorTok{+} \DecValTok{1}\OperatorTok{)} \OperatorTok{*}\NormalTok{ rootStates}\OperatorTok{);}
            \OperatorTok{\}}
\NormalTok{           crossTab}\OperatorTok{[}\NormalTok{it}\OperatorTok{][}\NormalTok{i}\OperatorTok{]} \OperatorTok{=}\NormalTok{ tmp}\OperatorTok{;}
        \OperatorTok{\}}
    \OperatorTok{\}}
    \KeywordTok{...}
\end{Highlighting}
\end{Shaded}
\paragraph{Cómo arreglarlo}
Leer cada casilla con el orden real de la tabla: \texttt{cross\allowbreak{}Tab{[}it{]}{[}i{]}{[}j{]}{[}k{]}\ =\ values{[}i\ +\ fst\allowbreak{}States\ *\ (\allowbreak{}k\ +\ root\allowbreak{}States\ *\ j){]}}, o pedir la tabla con las variables en el orden que el corte espera.

La otra tabla de la clase, \texttt{build\allowbreak{}Prob\allowbreak{}Distribution\allowbreak{}For\allowbreak{}Node}, sí respeta el orden de \texttt{get\allowbreak{}Absolute\allowbreak{}Frequencies} y no necesita cambio.

Conviene hacerlo junto con \hyperref[err:128]{error~128}, que corrige cómo \texttt{predict\allowbreak{}Class\allowbreak{}Value\allowbreak{}Augmented\allowbreak{}NB} usa estas cuentas.

Prueba de regresión: con variables de números de estados distintos (por ejemplo 3, 2 y 4), comprobar que \texttt{cross\allowbreak{}Tab{[}0{]}{[}i{]}{[}j{]}{[}k{]}} coincide con las cuentas directas sobre \texttt{dataset.\allowbreak{}get\allowbreak{}Training(\allowbreak{}){[}0{]}} en todas las casillas.

\setcounter{subsection}{127}
\subsection[El superpadre penaliza la clase verdadera: ningún enlace sube la exactitud]{El clasificador de superpadre mide cada enlace penalizando solo la clase verdadera del caso, así que ningún enlace puede mejorar la exactitud y elige enlaces que la empeoran}\label{err:128}
\begin{ficha}
\fichakey{Estado}Presente
\fichakey{Módulo}learning.metric
\fichakey{Paquete}org.openmarkov.learning.metric.cmi.accuracy
\fichakey{Ficheros}\path{learning.metric/src/main/java/org/openmarkov/learning/metric/cmi/accuracy/Accuracy.java:193-222} \newline \path{learning.metric/src/main/java/org/openmarkov/learning/metric/cmi/accuracy/Accuracy.java:150-163} \newline \path{learning.algorithm/src/main/java/org/openmarkov/learning/algorithm/nbderived/spnb/SuperParentNBAlgorithm.java:84-142}
\fichakey{Comprobado}Leyendo el código y ejecutándolo
\fichakey{Esfuerzo}Pequeño (horas)
\fichakey{Decisión de equipo}No
\fichakey{Tapado por}\hyperref[err:132]{error~132}: con «Unique superparent» marcada, que es la opción de fábrica, no se aprende ningún enlace.
\fichakey{Relacionados}\hyperref[err:125]{error~125}, \hyperref[err:126]{error~126}, \hyperref[err:127]{error~127}, \hyperref[err:129]{error~129}, \hyperref[err:130]{error~130}, \hyperref[err:131]{error~131}, \hyperref[err:132]{error~132}, \hyperref[err:133]{error~133}
\end{ficha}
\paragraph{Qué pasa}
El usuario llega por el menú de aprendizaje → tipo discriminativo → «Superparent naive bayes», y también por la validación cruzada de \texttt{bn\allowbreak{}Evaluation}. Con la opción «Unique superparent» marcada, que es la de fábrica, hoy no se aprende ningún enlace por otro motivo (\hyperref[err:132]{error~132}); desmarcándola, el algoritmo añade enlaces que, según su propia medida, empeoran el clasificador, y no puede encontrar ninguno que lo mejore.

El clasificador de superpadre parte del Bayes ingenuo y cuelga enlaces entre características de una característica elegida («superpadre»). Debe quedarse con un enlace solo si mejora la exactitud (la proporción de casos cuya clase se acierta). Esa exactitud la calcula la métrica \texttt{Accuracy} en modo «red aumentada» con \texttt{compute\allowbreak{}Augmented\allowbreak{}Net\allowbreak{}Accuracy}, que para cada caso de prueba llama a \texttt{predict\allowbreak{}Class\allowbreak{}Value\allowbreak{}Augmented\allowbreak{}NB}.

\texttt{predict\allowbreak{}Class\allowbreak{}Value\allowbreak{}Augmented\allowbreak{}NB} calcula para cada clase c la probabilidad a priori por el producto de P(característica \textbar{} c), y después, para cada enlace, multiplica \textbf{solo la casilla de la clase verdadera del caso} (\texttt{probs{[}row{[}get\allowbreak{}Class\allowbreak{}Variable\allowbreak{}Index(\allowbreak{}){]}{]}}) por \texttt{get\allowbreak{}Prob}, que además lee la cuenta de esa misma clase verdadera. Lo correcto es multiplicar cada clase c por la probabilidad del enlace calculada con c.~El programa usa la respuesta que tiene que adivinar.

El efecto es sistemático, no aleatorio. \texttt{get\allowbreak{}Prob} devuelve (cuenta + α)/(total + α·n), que nunca pasa de 1, así que multiplicar solo la clase verdadera solo puede bajarla frente a las demás: un caso bien clasificado puede pasar a mal clasificado y nunca al revés. La exactitud con cualquier conjunto de enlaces es, por tanto, menor o igual que sin enlaces. Con α = 0 y una combinación no vista, \texttt{get\allowbreak{}Prob} da 0/0 y la clase verdadera no puede ganar nunca.

En \texttt{Super\allowbreak{}Parent\allowbreak{}NBAlgorithm.\allowbreak{}get\allowbreak{}Optimal\allowbreak{}Edit} eso hace que el superpadre y cada enlace se elijan por cuánto \textbf{menos} empeoran la exactitud, y que ningún enlace pueda mejorar de verdad al Bayes ingenuo (el algoritmo tampoco lo comprueba: ver \hyperref[err:133]{error~133}).

Medido sobre \texttt{learning.\allowbreak{}algorithm/\allowbreak{}src/\allowbreak{}test/\allowbreak{}resources/\allowbreak{}network/\allowbreak{}BN-asia10k.\allowbreak{}csv}, clase \texttt{Bronchitis}, α = 0,5, con las mismas ventanas de prueba que usa la métrica:

\begin{itemize}
\tightlist
\item
  Con el código de hoy, la exactitud sin enlaces es 0,8444, y cada uno de los 42 enlaces posibles (parejas ordenadas de características) la baja, a valores entre 0,4264 y 0,8376.
\item
  Repitiendo la cuenta con cuentas hechas directamente sobre los casos (para apartar el defecto de \hyperref[err:127]{error~127}): multiplicando solo la clase verdadera, ningún enlace pasa de 0,8444 (el mejor, 0,8402); multiplicando cada clase con su propia cuenta, 7 de los 42 enlaces la mejoran (el mejor, el del factor P(Dyspnea \textbar{} X-ray, clase), con 0,8458). Esta cuenta conserva la doble cuenta de la característica hija (\hyperref[err:130]{error~130}); sin ella, los enlaces que mejoran son 14.
\item
  El algoritmo de superpadre con «Unique superparent» desmarcada añade el enlace X-ray → Smoker con una exactitud medida de 0,8253, por debajo de la del Bayes ingenuo sin enlaces, y se detiene; que lo acepte aunque quede por debajo es otro defecto, descrito en \hyperref[err:133]{error~133}.
\end{itemize}
\paragraph{El código}
learning.metric/src/main/java/org/openmarkov/learning/metric/cmi/accuracy/Accuracy.java:193-222

\begin{Shaded}
\begin{Highlighting}[]
\KeywordTok{public} \DataTypeTok{int} \FunctionTok{predictClassValueAugmentedNB}\OperatorTok{(}\DataTypeTok{int}\OperatorTok{[]}\NormalTok{ row}\OperatorTok{,} \BuiltInTok{List}\OperatorTok{\textless{}}\NormalTok{Variable}\OperatorTok{[]\textgreater{}}\NormalTok{ van}\OperatorTok{,} \DataTypeTok{int}\NormalTok{ it}\OperatorTok{)\{}
    \DataTypeTok{double}\OperatorTok{[]}\NormalTok{ probs }\OperatorTok{=} \FunctionTok{getProbRootNode}\OperatorTok{();}
    \ControlFlowTok{for} \OperatorTok{(}\BuiltInTok{String}\NormalTok{ variable }\OperatorTok{:} \FunctionTok{getNonRootNodes}\OperatorTok{().}\FunctionTok{stream}\OperatorTok{().}\FunctionTok{map}\OperatorTok{(}\BuiltInTok{Node}\OperatorTok{::}\NormalTok{getVariable}\OperatorTok{).}\FunctionTok{map}\OperatorTok{(}\NormalTok{Variable}\OperatorTok{::}\NormalTok{getName}\OperatorTok{).}\FunctionTok{toList}\OperatorTok{())} \OperatorTok{\{}
        \DataTypeTok{int}\NormalTok{ index }\OperatorTok{=} \FunctionTok{getIndexVariable}\OperatorTok{(}\NormalTok{variable}\OperatorTok{);}
        \ControlFlowTok{for}\OperatorTok{(}\DataTypeTok{int}\NormalTok{ i}\OperatorTok{=}\DecValTok{0}\OperatorTok{;}\NormalTok{ i}\OperatorTok{\textless{}}\NormalTok{freqRootNode}\OperatorTok{.}\FunctionTok{length}\OperatorTok{;}\NormalTok{i}\OperatorTok{++)\{}
\NormalTok{            probs}\OperatorTok{[}\NormalTok{i}\OperatorTok{]*=} \OperatorTok{((}\NormalTok{crossTabs}\OperatorTok{.}\FunctionTok{get}\OperatorTok{(}\NormalTok{variable}\OperatorTok{)[}\NormalTok{it}\OperatorTok{][}\NormalTok{row}\OperatorTok{[}\NormalTok{index}\OperatorTok{]][}\NormalTok{i}\OperatorTok{]+}\NormalTok{ alpha}\OperatorTok{)/(}\NormalTok{freqRootNode}\OperatorTok{[}\NormalTok{i}\OperatorTok{]+}\NormalTok{ alpha}\OperatorTok{*}\FunctionTok{getNonRootNodes}\OperatorTok{().}\FunctionTok{size}\OperatorTok{()))}
            \OperatorTok{;}
        \OperatorTok{\}}
    \OperatorTok{\}}
\NormalTok{    van}\OperatorTok{.}\FunctionTok{forEach}\OperatorTok{(}\NormalTok{arc}\OperatorTok{{-}\textgreater{}\{}
\NormalTok{        probs}\OperatorTok{[}\NormalTok{row}\OperatorTok{[}\FunctionTok{getClassVariableIndex}\OperatorTok{()]]} \OperatorTok{*=} \FunctionTok{getProb}\OperatorTok{(}\NormalTok{arc}\OperatorTok{[}\DecValTok{0}\OperatorTok{],}\NormalTok{ arc}\OperatorTok{[}\DecValTok{1}\OperatorTok{],}\NormalTok{ row}\OperatorTok{,}\NormalTok{ it}\OperatorTok{);}
    \OperatorTok{\});}
    \ControlFlowTok{return} \FunctionTok{getIndexOfMaxValue}\OperatorTok{(}\NormalTok{probs}\OperatorTok{);}
\OperatorTok{\}}

\KeywordTok{protected} \DataTypeTok{double} \FunctionTok{getProb}\OperatorTok{(}\NormalTok{Variable tail}\OperatorTok{,}\NormalTok{ Variable head}\OperatorTok{,} \DataTypeTok{int}\OperatorTok{[]}\NormalTok{ row}\OperatorTok{,} \DataTypeTok{int}\NormalTok{ it}\OperatorTok{)\{}
    \DataTypeTok{double}\OperatorTok{[][][][]}\NormalTok{ ct }\OperatorTok{=}\NormalTok{ \_2ndLevelCrosstab}\OperatorTok{.}\FunctionTok{get}\OperatorTok{(}\NormalTok{tail}\OperatorTok{.}\FunctionTok{getName}\OperatorTok{()+}\StringTok{"{-}"}\OperatorTok{+}\NormalTok{head}\OperatorTok{.}\FunctionTok{getName}\OperatorTok{());}
    \DataTypeTok{double}\NormalTok{ count }\OperatorTok{=} \DecValTok{0}\OperatorTok{;}
    \DataTypeTok{double}\NormalTok{ value }\OperatorTok{=}\NormalTok{ ct}\OperatorTok{[}\NormalTok{it}\OperatorTok{][}\NormalTok{row}\OperatorTok{[}\FunctionTok{getIndexVariable}\OperatorTok{(}\NormalTok{tail}\OperatorTok{)]][}\NormalTok{row}\OperatorTok{[}\FunctionTok{getIndexVariable}\OperatorTok{(}\NormalTok{head}\OperatorTok{)]][}\NormalTok{row}\OperatorTok{[}\FunctionTok{getClassVariableIndex}\OperatorTok{()]];}
    \ControlFlowTok{for}\OperatorTok{(}\DataTypeTok{int}\NormalTok{ i}\OperatorTok{=}\DecValTok{0}\OperatorTok{;}\NormalTok{ i }\OperatorTok{\textless{}}\NormalTok{ tail}\OperatorTok{.}\FunctionTok{getNumStates}\OperatorTok{();}\NormalTok{ i}\OperatorTok{++)\{}
\NormalTok{        count}\OperatorTok{+=}\NormalTok{ct}\OperatorTok{[}\NormalTok{it}\OperatorTok{][}\NormalTok{i}\OperatorTok{][}\NormalTok{row}\OperatorTok{[}\FunctionTok{getIndexVariable}\OperatorTok{(}\NormalTok{head}\OperatorTok{)]][}\NormalTok{row}\OperatorTok{[}\FunctionTok{getClassVariableIndex}\OperatorTok{()]];}
    \OperatorTok{\}}
    \ControlFlowTok{return} \OperatorTok{(}\NormalTok{value }\OperatorTok{+}\NormalTok{ alpha}\OperatorTok{)} \OperatorTok{/} \OperatorTok{(}\NormalTok{count }\OperatorTok{+}\NormalTok{ alpha }\OperatorTok{*} \FunctionTok{getNonRootNodes}\OperatorTok{().}\FunctionTok{size}\OperatorTok{());}
\OperatorTok{\}}
\end{Highlighting}
\end{Shaded}
\paragraph{Cómo arreglarlo}
Pasar a \texttt{get\allowbreak{}Prob} la clase como parámetro en vez de leerla del caso, y en \texttt{predict\allowbreak{}Class\allowbreak{}Value\allowbreak{}Augmented\allowbreak{}NB} multiplicar cada \texttt{probs{[}c{]}} por \texttt{get\allowbreak{}Prob(\allowbreak{}arc{[}0{]},\allowbreak{}\ arc{[}1{]},\allowbreak{}\ row,\allowbreak{}\ it,\allowbreak{}\ c)} para todas las clases c.

\texttt{predict\allowbreak{}Class\allowbreak{}Value\allowbreak{}Augmented\allowbreak{}NB} es el único sitio que usa \texttt{get\allowbreak{}Prob}; \texttt{compute\allowbreak{}Augmented\allowbreak{}Net\allowbreak{}Accuracy} es su único llamador, y a él llegan \texttt{score\allowbreak{}Augmented\allowbreak{}Net} y \texttt{get\allowbreak{}Score} en modo aumentado.

Conviene hacerlo junto con \hyperref[err:127]{error~127}, que da a \texttt{get\allowbreak{}Prob} cuentas de casillas equivocadas, y con \hyperref[err:131]{error~131} y \hyperref[err:132]{error~132}, porque los cuatro cambian los enlaces que aprende el clasificador de superpadre.

Prueba de regresión: sobre \texttt{BN-asia10k.\allowbreak{}csv} con clase \texttt{Bronchitis}, la exactitud de \texttt{compute\allowbreak{}Augmented\allowbreak{}Net\allowbreak{}Accuracy} con un enlace debe poder superar a la de sin enlaces (con cuentas correctas lo hace el enlace del factor P(LungCancer \textbar{} Smoker, clase), con y sin la doble cuenta de \hyperref[err:130]{error~130}); y, en general, multiplicar el factor de un enlace por una misma constante en todas las clases no debe cambiar la predicción.

\setcounter{subsection}{128}
\subsection[El superpadre puntúa los enlaces candidatos con la probabilidad al revés]{El clasificador de superpadre puntúa el superpadre y cada enlace candidato con la probabilidad condicionada al revés, P(padre \textbar{} hijo, clase), y a los enlaces ya añadidos les aplica la del sentido correcto, así que decide con una exactitud que no es la de la red que resulta}\label{err:129}
\begin{ficha}
\fichakey{Estado}Presente
\fichakey{Módulo}learning.metric
\fichakey{Paquete}org.openmarkov.learning.metric.cmi.accuracy
\fichakey{Ficheros}\path{learning.metric/src/main/java/org/openmarkov/learning/metric/cmi/accuracy/Accuracy.java:64-97} \newline \path{learning.metric/src/main/java/org/openmarkov/learning/metric/cmi/accuracy/Accuracy.java:193-222} \newline \path{learning.algorithm/src/main/java/org/openmarkov/learning/algorithm/nbderived/spnb/SuperParentNBAlgorithm.java:84-139}
\fichakey{Comprobado}Leyendo el código y ejecutándolo
\fichakey{Esfuerzo}Pequeño (horas)
\fichakey{Decisión de equipo}No
\fichakey{Tapado por}\hyperref[err:132]{error~132}: con «Unique superparent» marcada, que es la opción de fábrica, no se aprende ningún enlace.
\fichakey{Relacionados}\hyperref[err:127]{error~127}, \hyperref[err:128]{error~128}, \hyperref[err:130]{error~130}, \hyperref[err:131]{error~131}, \hyperref[err:132]{error~132}, \hyperref[err:133]{error~133}
\end{ficha}
\paragraph{Qué pasa}
El usuario llega por Herramientas («Tools») → «Learning» → tipo «Discriminative» → «Superparent naive bayes», y también por la validación cruzada de \texttt{bn\allowbreak{}Evaluation}, que ofrece el mismo algoritmo. El algoritmo elige el superpadre y cada enlace por la exactitud (proporción de casos cuya clase se acierta) que la métrica \texttt{Accuracy} da a la red con el cambio propuesto; esa exactitud se calcula con un factor distinto del que tendrá la red una vez hecho el cambio, así que acepta o rechaza enlaces por un número que no es el de la red resultante.

El clasificador de superpadre parte del Bayes ingenuo, elige una característica como «superpadre» y le cuelga como hijas otras características. En modo «red aumentada», \texttt{Accuracy.\allowbreak{}score\allowbreak{}Augmented\allowbreak{}Net} reúne una lista de parejas de variables y \texttt{predict\allowbreak{}Class\allowbreak{}Value\allowbreak{}Augmented\allowbreak{}NB} multiplica, por cada pareja, \texttt{get\allowbreak{}Prob(\allowbreak{}pareja{[}0{]},\allowbreak{}\ pareja{[}1{]},\allowbreak{}\ …)}. \texttt{get\allowbreak{}Prob(\allowbreak{}tail,\allowbreak{}\ head,\allowbreak{}\ …)} divide la cuenta de (tail, head, clase) entre la de (head, clase): es P(pareja{[}0{]} \textbar{} pareja{[}1{]}, clase). Para que sea la probabilidad de la hija dado su padre, la pareja tiene que ir como \{hija, padre\}.

\texttt{score\allowbreak{}Augmented\allowbreak{}Net} no ordena igual todas las parejas:

\begin{itemize}
\tightlist
\item
  Enlaces ya añadidos (líneas 83-86): para cada característica con un segundo padre pone \{característica, padre\}, y el factor es P(hija \textbar{} padre, clase), el correcto.
\item
  Enlace candidato entre dos características, X→Y (línea 89): pone la pareja tal como viene en el cambio, \{X, Y\}, y el factor es P(X \textbar{} Y, clase): el del padre dado la hija.
\item
  Elección del superpadre S, que se puntúa con el cambio clase→S (líneas 91-93): pone \{S, característica\} para cada característica sin segundo padre, y el factor es P(S \textbar{} característica, clase), también al revés.
\end{itemize}

\texttt{Super\allowbreak{}Parent\allowbreak{}NBAlgorithm.\allowbreak{}get\allowbreak{}Optimal\allowbreak{}Edit} elige el superpadre con la mejor de esas puntuaciones y acepta cada enlace S→hija si su puntuación supera a la anterior. En el paso siguiente, el enlace ya añadido se puntúa en el otro sentido, con otro número.

Medido sobre \texttt{learning.\allowbreak{}algorithm/\allowbreak{}src/\allowbreak{}test/\allowbreak{}resources/\allowbreak{}network/\allowbreak{}BN-asia10k.\allowbreak{}csv}, clase \texttt{Bronchitis}, alfa 0,5, con «Unique superparent» desmarcada (con la opción marcada, que es la de fábrica, hoy no se aprende nada por \hyperref[err:132]{error~132}):

\begin{itemize}
\tightlist
\item
  Con el código de hoy, el algoritmo añade X-ray→Smoker con una exactitud de 0,8253. La misma métrica, sobre la red con ese enlace tal como lo puntúa a partir del paso siguiente (la pareja \{Smoker, X-ray\}), da 0,6735.
\item
  Repitiendo la cuenta con cuentas hechas directamente sobre los casos y multiplicando cada clase por su factor (para apartar \hyperref[err:127]{error~127} y \hyperref[err:128]{error~128}), y con el resto del producto tal como lo escribe hoy el método, que cuenta dos veces la característica hija (\hyperref[err:130]{error~130}): con el factor al revés el superpadre elegido sería Dyspnea (0,8458), y con el correcto, TuberculosisOrCancer, empatado con Smoker a 0,8456 (en un empate el algoritmo se queda con el último). Los seis enlaces que salen de Dyspnea, que con el factor al revés dan 0,8458, por encima del Bayes ingenuo (0,8444), con el factor correcto quedan todos por debajo (entre 0,8426 y 0,8443). De los 42 enlaces posibles entre dos características, 40 dan distinta exactitud según el sentido del factor.
\end{itemize}
\paragraph{El código}
learning.metric/src/main/java/org/openmarkov/learning/metric/cmi/accuracy/Accuracy.java:80-97

\begin{Shaded}
\begin{Highlighting}[]
\KeywordTok{private} \DataTypeTok{double} \FunctionTok{scoreAugmentedNet}\OperatorTok{(}\NormalTok{Variable}\OperatorTok{[]}\NormalTok{ v}\OperatorTok{)\{}
    \BuiltInTok{LinkedList}\OperatorTok{\textless{}}\NormalTok{Variable}\OperatorTok{[]\textgreater{}}\NormalTok{ augmentedNet }\OperatorTok{=} \KeywordTok{new} \BuiltInTok{LinkedList}\OperatorTok{\textless{}\textgreater{}();}

    \FunctionTok{getNonRootNodes}\OperatorTok{().}\FunctionTok{stream}\OperatorTok{().}\FunctionTok{filter}\OperatorTok{(}\NormalTok{n }\OperatorTok{{-}\textgreater{}}\NormalTok{ n}\OperatorTok{.}\FunctionTok{getNumParents}\OperatorTok{()} \OperatorTok{\textgreater{}} \DecValTok{1}\OperatorTok{).}\FunctionTok{toList}\OperatorTok{().}\FunctionTok{forEach}\OperatorTok{(}\NormalTok{node }\OperatorTok{{-}\textgreater{}} \OperatorTok{\{}
\NormalTok{        augmentedNet}\OperatorTok{.}\FunctionTok{add}\OperatorTok{(}\KeywordTok{new}\NormalTok{ Variable}\OperatorTok{[]\{}\NormalTok{node}\OperatorTok{.}\FunctionTok{getVariable}\OperatorTok{(),}
\NormalTok{                node}\OperatorTok{.}\FunctionTok{getParents}\OperatorTok{().}\FunctionTok{stream}\OperatorTok{().}\FunctionTok{filter}\OperatorTok{(}\NormalTok{n }\OperatorTok{{-}\textgreater{}}\NormalTok{ n }\OperatorTok{!=} \FunctionTok{getRootNode}\OperatorTok{()).}\FunctionTok{toList}\OperatorTok{().}\FunctionTok{get}\OperatorTok{(}\DecValTok{0}\OperatorTok{).}\FunctionTok{getVariable}\OperatorTok{()\});}
    \OperatorTok{\});}

    \ControlFlowTok{if}\OperatorTok{(!}\NormalTok{v}\OperatorTok{[}\DecValTok{0}\OperatorTok{].}\FunctionTok{getName}\OperatorTok{().}\FunctionTok{equals}\OperatorTok{(}\FunctionTok{getRootNode}\OperatorTok{().}\FunctionTok{getName}\OperatorTok{())} \OperatorTok{\&\&} \OperatorTok{!}\NormalTok{v}\OperatorTok{[}\DecValTok{1}\OperatorTok{].}\FunctionTok{getName}\OperatorTok{().}\FunctionTok{equals}\OperatorTok{(}\FunctionTok{getRootNode}\OperatorTok{().}\FunctionTok{getName}\OperatorTok{()))\{}
\NormalTok{        augmentedNet}\OperatorTok{.}\FunctionTok{add}\OperatorTok{(}\NormalTok{v}\OperatorTok{);}
    \OperatorTok{\}}\ControlFlowTok{else}\OperatorTok{\{}
        \FunctionTok{getNonRootNodes}\OperatorTok{().}\FunctionTok{stream}\OperatorTok{().}\FunctionTok{filter}\OperatorTok{(}\NormalTok{n}\OperatorTok{{-}\textgreater{}}\NormalTok{n}\OperatorTok{.}\FunctionTok{getNumParents}\OperatorTok{()\textless{}}\DecValTok{2} \OperatorTok{\&\&} \OperatorTok{!}\NormalTok{n}\OperatorTok{.}\FunctionTok{getName}\OperatorTok{().}\FunctionTok{equals}\OperatorTok{(}\NormalTok{v}\OperatorTok{[}\DecValTok{1}\OperatorTok{].}\FunctionTok{getName}\OperatorTok{())).}\FunctionTok{forEach}\OperatorTok{(}\NormalTok{tail }\OperatorTok{{-}\textgreater{}\{}
\NormalTok{            augmentedNet}\OperatorTok{.}\FunctionTok{add}\OperatorTok{(}\KeywordTok{new}\NormalTok{ Variable}\OperatorTok{[]\{}\NormalTok{v}\OperatorTok{[}\DecValTok{1}\OperatorTok{],}\NormalTok{ tail}\OperatorTok{.}\FunctionTok{getVariable}\OperatorTok{()\});}
        \OperatorTok{\});}
    \OperatorTok{\}}

    \ControlFlowTok{return} \FunctionTok{computeAugmentedNetAccuracy}\OperatorTok{(}\NormalTok{augmentedNet}\OperatorTok{);}
\OperatorTok{\}}
\end{Highlighting}
\end{Shaded}

learning.metric/src/main/java/org/openmarkov/learning/metric/cmi/accuracy/Accuracy.java:212-222

\begin{Shaded}
\begin{Highlighting}[]
\KeywordTok{protected} \DataTypeTok{double} \FunctionTok{getProb}\OperatorTok{(}\NormalTok{Variable tail}\OperatorTok{,}\NormalTok{ Variable head}\OperatorTok{,} \DataTypeTok{int}\OperatorTok{[]}\NormalTok{ row}\OperatorTok{,} \DataTypeTok{int}\NormalTok{ it}\OperatorTok{)\{}
    \DataTypeTok{double}\OperatorTok{[][][][]}\NormalTok{ ct }\OperatorTok{=}\NormalTok{ \_2ndLevelCrosstab}\OperatorTok{.}\FunctionTok{get}\OperatorTok{(}\NormalTok{tail}\OperatorTok{.}\FunctionTok{getName}\OperatorTok{()+}\StringTok{"{-}"}\OperatorTok{+}\NormalTok{head}\OperatorTok{.}\FunctionTok{getName}\OperatorTok{());}
    \DataTypeTok{double}\NormalTok{ count }\OperatorTok{=} \DecValTok{0}\OperatorTok{;}

    \DataTypeTok{double}\NormalTok{ value }\OperatorTok{=}\NormalTok{ ct}\OperatorTok{[}\NormalTok{it}\OperatorTok{][}\NormalTok{row}\OperatorTok{[}\FunctionTok{getIndexVariable}\OperatorTok{(}\NormalTok{tail}\OperatorTok{)]][}\NormalTok{row}\OperatorTok{[}\FunctionTok{getIndexVariable}\OperatorTok{(}\NormalTok{head}\OperatorTok{)]][}\NormalTok{row}\OperatorTok{[}\FunctionTok{getClassVariableIndex}\OperatorTok{()]];}
    \ControlFlowTok{for}\OperatorTok{(}\DataTypeTok{int}\NormalTok{ i}\OperatorTok{=}\DecValTok{0}\OperatorTok{;}\NormalTok{ i }\OperatorTok{\textless{}}\NormalTok{ tail}\OperatorTok{.}\FunctionTok{getNumStates}\OperatorTok{();}\NormalTok{ i}\OperatorTok{++)\{}
\NormalTok{        count}\OperatorTok{+=}\NormalTok{ct}\OperatorTok{[}\NormalTok{it}\OperatorTok{][}\NormalTok{i}\OperatorTok{][}\NormalTok{row}\OperatorTok{[}\FunctionTok{getIndexVariable}\OperatorTok{(}\NormalTok{head}\OperatorTok{)]][}\NormalTok{row}\OperatorTok{[}\FunctionTok{getClassVariableIndex}\OperatorTok{()]];}
    \OperatorTok{\}}

    \ControlFlowTok{return} \OperatorTok{(}\NormalTok{value }\OperatorTok{+}\NormalTok{ alpha}\OperatorTok{)} \OperatorTok{/} \OperatorTok{(}\NormalTok{count }\OperatorTok{+}\NormalTok{ alpha }\OperatorTok{*} \FunctionTok{getNonRootNodes}\OperatorTok{().}\FunctionTok{size}\OperatorTok{());}
\OperatorTok{\}}
\end{Highlighting}
\end{Shaded}
\paragraph{Cómo arreglarlo}
Usar un solo convenio para todas las parejas. Lo más corto es construir las del candidato con la hija primero, como ya se hace con los enlaces añadidos: \texttt{new\ Variable{[}{]}\{v{[}1{]},\allowbreak{}\ v{[}0{]}\}} para un enlace entre características y \texttt{new\ Variable{[}{]}\{tail.\allowbreak{}get\allowbreak{}Variable(\allowbreak{}),\allowbreak{}\ v{[}1{]}\}} para el superpadre. La otra opción es que \texttt{get\allowbreak{}Prob} reciba (padre, hija) y dar la vuelta a las líneas 83-86. En los dos casos conviene renombrar los parámetros de \texttt{get\allowbreak{}Prob}: \texttt{tail} y \texttt{head} suelen nombrar el origen y el destino de un enlace, y aquí se usan al contrario.

\texttt{score\allowbreak{}Augmented\allowbreak{}Net} es el único sitio que construye las parejas; las consumen \texttt{compute\allowbreak{}Augmented\allowbreak{}Net\allowbreak{}Accuracy}, \texttt{predict\allowbreak{}Class\allowbreak{}Value\allowbreak{}Augmented\allowbreak{}NB} y \texttt{get\allowbreak{}Prob}. \texttt{Super\allowbreak{}Parent\allowbreak{}NBAlgorithm} es el único algoritmo que activa el modo aumentado.

Conviene hacerlo a la vez que \hyperref[err:127]{error~127} (cuentas leídas de casillas equivocadas en la misma tabla), \hyperref[err:128]{error~128} (el factor solo se aplica a la clase verdadera) y \hyperref[err:133]{error~133} (el algoritmo no compara con el Bayes ingenuo de partida): los cuatro cambian qué enlaces aprende el clasificador de superpadre.

Prueba de regresión: sobre \texttt{BN-asia10k.\allowbreak{}csv} con clase \texttt{Bronchitis}, la puntuación que \texttt{get\allowbreak{}Score} da a \texttt{Add\allowbreak{}Link\allowbreak{}Edit(\allowbreak{}X-ray\ →\ Smoker)} sin enlaces entre características debe ser igual a la que la métrica da a la red una vez añadido ese enlace; y, en general, un enlace debe puntuarse igual como candidato que como enlace ya añadido.

\setcounter{subsection}{129}
\subsection[El superpadre cuenta dos veces la hija de cada enlace al medir la exactitud]{El clasificador de superpadre cuenta dos veces cada característica que recibe un segundo padre: al medir la exactitud multiplica P(hija \textbar{} clase) y además P(hija \textbar{} padre, clase), así que elige enlaces con un clasificador que no es el que construye}\label{err:130}
\begin{ficha}
\fichakey{Estado}Presente
\fichakey{Módulo}learning.metric
\fichakey{Paquete}org.openmarkov.learning.metric.cmi.accuracy
\fichakey{Ficheros}\path{learning.metric/src/main/java/org/openmarkov/learning/metric/cmi/accuracy/Accuracy.java:193-209} \newline \path{learning.metric/src/main/java/org/openmarkov/learning/metric/cmi/accuracy/Accuracy.java:80-97} \newline \path{learning.metric/src/main/java/org/openmarkov/learning/metric/cmi/accuracy/Accuracy.java:212-222} \newline \path{learning.algorithm/src/main/java/org/openmarkov/learning/algorithm/nbderived/spnb/SuperParentNBAlgorithm.java:84-142} \newline \path{learning.core/src/main/java/org/openmarkov/learning/core/algorithm/LearningAlgorithm.java:172-183}
\fichakey{Comprobado}Leyendo el código y ejecutándolo
\fichakey{Esfuerzo}Pequeño (horas)
\fichakey{Decisión de equipo}No
\fichakey{Tapado por}\hyperref[err:132]{error~132}: con «Unique superparent» marcada, que es la opción de fábrica, no se aprende ningún enlace.
\fichakey{Relacionados}\hyperref[err:127]{error~127}, \hyperref[err:128]{error~128}, \hyperref[err:129]{error~129}, \hyperref[err:132]{error~132}, \hyperref[err:133]{error~133}
\end{ficha}
\paragraph{Qué pasa}
El usuario llega por Herramientas («Tools») → «Learning» → tipo «Discriminative» → «Superparent naive bayes», y también por la validación cruzada de \texttt{bn\allowbreak{}Evaluation}, que ofrece el mismo algoritmo. Para elegir el superpadre y los enlaces, el algoritmo mide la exactitud (la proporción de casos cuya clase se acierta) de un clasificador en el que cada característica que recibe un segundo padre cuenta dos veces. La red que devuelve no es así, de modo que los enlaces se eligen por un número que no corresponde a la red resultante. Con «Unique superparent» marcada, que es la opción de fábrica, hoy no se aprende ningún enlace por otro motivo (\hyperref[err:132]{error~132}); desmarcándola, quitar solo esta doble cuenta cambia el enlace aprendido en tres de las cuatro bases probadas.

En un Bayes ingenuo aumentado, la probabilidad de cada clase c es P(c) por el producto, para cada característica X, de su probabilidad dada su familia: P(X \textbar{} c) si solo tiene la clase como padre, y P(X \textbar{} padre, c) si tiene además otro padre. La red que resulta del aprendizaje es así: el aprendizaje de parámetros (\texttt{Learning\allowbreak{}Algorithm.\allowbreak{}parametric\allowbreak{}Learning}) da a cada característica una sola tabla, condicionada a todos sus padres.

\texttt{Accuracy.\allowbreak{}predict\allowbreak{}Class\allowbreak{}Value\allowbreak{}Augmented\allowbreak{}NB} multiplica primero P(X \textbar{} c) para \textbf{todas} las características (\texttt{get\allowbreak{}Non\allowbreak{}Root\allowbreak{}Nodes}) y después, por cada pareja de la lista de enlaces, el factor de \texttt{get\allowbreak{}Prob}, que es la probabilidad de un miembro de la pareja dado el otro y la clase. No quita el factor P(X \textbar{} c) de la característica que en ese enlace hace de hija, así que esa característica aporta dos factores, P(hija \textbar{} c) · P(hija \textbar{} padre, c). Pasa con todas las parejas que construye \texttt{score\allowbreak{}Augmented\allowbreak{}Net}: las de los enlaces ya añadidos, la del enlace candidato y las del superpadre que se está puntuando con todas sus hijas.

Hoy el cálculo tiene además otros tres defectos: cuentas leídas en casillas equivocadas (\hyperref[err:127]{error~127}), el factor de cada enlace aplicado solo a la clase verdadera del caso (\hyperref[err:128]{error~128}) y la probabilidad de las parejas candidatas tomada al revés (\hyperref[err:129]{error~129}). Las cuentas con que \hyperref[err:128]{error~128} y \hyperref[err:129]{error~129} apartan esos defectos conservan esta doble cuenta.

Medido sobre cuatro bases de casos, con alfa 0,5 y «Unique superparent» desmarcada, ejecutando el algoritmo con la métrica de hoy y con una copia de la métrica que solo deja fuera P(hija \textbar{} c) de la hija de cada pareja (la hija de un enlace ya añadido y la del enlace o del superpadre que se puntúa), sin tocar nada más:

\begin{itemize}
\tightlist
\item
  \texttt{learning.\allowbreak{}algorithm/\allowbreak{}src/\allowbreak{}test/\allowbreak{}resources/\allowbreak{}network/\allowbreak{}BN-asia10k.\allowbreak{}csv}, clase \texttt{Bronchitis}: hoy aprende X-ray→Smoker (exactitud medida 0,8253); sin la doble cuenta, Dyspnea→Smoker (0,8394).
\item
  \texttt{integration\allowbreak{}Tests/\allowbreak{}src/\allowbreak{}main/\allowbreak{}resources/\allowbreak{}networks/\allowbreak{}learning/\allowbreak{}learn\allowbreak{}Test\allowbreak{}Data\allowbreak{}Base.\allowbreak{}csv}, clase \texttt{A}: B→E (0,7220) frente a D→B (0,7290).
\item
  \texttt{learning.\allowbreak{}algorithm/\allowbreak{}src/\allowbreak{}test/\allowbreak{}resources/\allowbreak{}network/\allowbreak{}Head\allowbreak{}To\allowbreak{}Head2.\allowbreak{}csv}, clase \texttt{C}: B→F (0,4800) frente a B→E (0,4116).
\item
  \texttt{core/\allowbreak{}src/\allowbreak{}test/\allowbreak{}resources/\allowbreak{}nets/\allowbreak{}Cataract-RRDdataset\allowbreak{}Pre\allowbreak{}Proc3.\allowbreak{}csv}, clase \texttt{success}: opType→breakType en los dos casos (0,8302 frente a 0,8453).
\end{itemize}

Medido también con los otros tres defectos apartados, con un cálculo aparte sobre las mismas ventanas de prueba y el mismo suavizado que la métrica (cuentas hechas directamente sobre los casos, el factor de cada enlace en cada clase y la probabilidad de la hija dado el padre), con y sin la doble cuenta:

\begin{itemize}
\tightlist
\item
  \texttt{BN-asia10k.\allowbreak{}csv}: el Bayes ingenuo acierta 0,8444. Con la doble cuenta, 7 de los 42 enlaces posibles entre dos características lo mejoran, entre ellos los seis que llegan a Dyspnea; sin ella, 14, y ninguno de esos seis. De los 42 enlaces, 36 dan distinta exactitud con y sin la doble cuenta. Siguiendo el algoritmo publicado desde el Bayes ingenuo, con la doble cuenta se añade Smoker→Dyspnea (0,8458), que sin ella da 0,8438, por debajo del Bayes ingenuo; sin ella se añaden TuberculosisOrCancer→Smoker y TuberculosisOrCancer→X-ray (0,8458).
\item
  \texttt{learn\allowbreak{}Test\allowbreak{}Data\allowbreak{}Base.\allowbreak{}csv}: con la doble cuenta, E→C, E→B y B→D (0,8400, y 0,8070 sin la doble cuenta); sin ella, D→B, D→C y D→E (0,8360).
\item
  \texttt{Head\allowbreak{}To\allowbreak{}Head2.\allowbreak{}csv}: con la doble cuenta, F→E y B→A (0,7452); sin ella, F→E y A→B (0,7480).
\item
  \texttt{Cataract-RRDdataset\allowbreak{}Pre\allowbreak{}Proc3.\allowbreak{}csv}: en los dos casos, ningún enlace mejora el Bayes ingenuo (0,9642).
\end{itemize}
\paragraph{El código}
learning.metric/src/main/java/org/openmarkov/learning/metric/cmi/accuracy/Accuracy.java:193-209

\begin{Shaded}
\begin{Highlighting}[]
\KeywordTok{public} \DataTypeTok{int} \FunctionTok{predictClassValueAugmentedNB}\OperatorTok{(}\DataTypeTok{int}\OperatorTok{[]}\NormalTok{ row}\OperatorTok{,} \BuiltInTok{List}\OperatorTok{\textless{}}\NormalTok{Variable}\OperatorTok{[]\textgreater{}}\NormalTok{ van}\OperatorTok{,} \DataTypeTok{int}\NormalTok{ it}\OperatorTok{)\{}
    \DataTypeTok{double}\OperatorTok{[]}\NormalTok{ probs }\OperatorTok{=} \FunctionTok{getProbRootNode}\OperatorTok{();}

    \ControlFlowTok{for} \OperatorTok{(}\BuiltInTok{String}\NormalTok{ variable }\OperatorTok{:} \FunctionTok{getNonRootNodes}\OperatorTok{().}\FunctionTok{stream}\OperatorTok{().}\FunctionTok{map}\OperatorTok{(}\BuiltInTok{Node}\OperatorTok{::}\NormalTok{getVariable}\OperatorTok{).}\FunctionTok{map}\OperatorTok{(}\NormalTok{Variable}\OperatorTok{::}\NormalTok{getName}\OperatorTok{).}\FunctionTok{toList}\OperatorTok{())} \OperatorTok{\{}
        \DataTypeTok{int}\NormalTok{ index }\OperatorTok{=} \FunctionTok{getIndexVariable}\OperatorTok{(}\NormalTok{variable}\OperatorTok{);}
        \ControlFlowTok{for}\OperatorTok{(}\DataTypeTok{int}\NormalTok{ i}\OperatorTok{=}\DecValTok{0}\OperatorTok{;}\NormalTok{ i}\OperatorTok{\textless{}}\NormalTok{freqRootNode}\OperatorTok{.}\FunctionTok{length}\OperatorTok{;}\NormalTok{i}\OperatorTok{++)\{}
\NormalTok{            probs}\OperatorTok{[}\NormalTok{i}\OperatorTok{]*=} \OperatorTok{((}\NormalTok{crossTabs}\OperatorTok{.}\FunctionTok{get}\OperatorTok{(}\NormalTok{variable}\OperatorTok{)[}\NormalTok{it}\OperatorTok{][}\NormalTok{row}\OperatorTok{[}\NormalTok{index}\OperatorTok{]][}\NormalTok{i}\OperatorTok{]+}\NormalTok{ alpha}\OperatorTok{)/(}\NormalTok{freqRootNode}\OperatorTok{[}\NormalTok{i}\OperatorTok{]+}\NormalTok{ alpha}\OperatorTok{*}\FunctionTok{getNonRootNodes}\OperatorTok{().}\FunctionTok{size}\OperatorTok{()))}
            \OperatorTok{;}
        \OperatorTok{\}}
    \OperatorTok{\}}

\NormalTok{    van}\OperatorTok{.}\FunctionTok{forEach}\OperatorTok{(}\NormalTok{arc}\OperatorTok{{-}\textgreater{}\{}
\NormalTok{        probs}\OperatorTok{[}\NormalTok{row}\OperatorTok{[}\FunctionTok{getClassVariableIndex}\OperatorTok{()]]} \OperatorTok{*=} \FunctionTok{getProb}\OperatorTok{(}\NormalTok{arc}\OperatorTok{[}\DecValTok{0}\OperatorTok{],}\NormalTok{ arc}\OperatorTok{[}\DecValTok{1}\OperatorTok{],}\NormalTok{ row}\OperatorTok{,}\NormalTok{ it}\OperatorTok{);}
    \OperatorTok{\});}

    \ControlFlowTok{return} \FunctionTok{getIndexOfMaxValue}\OperatorTok{(}\NormalTok{probs}\OperatorTok{);}
\OperatorTok{\}}
\end{Highlighting}
\end{Shaded}
\paragraph{Cómo arreglarlo}
En \texttt{predict\allowbreak{}Class\allowbreak{}Value\allowbreak{}Augmented\allowbreak{}NB}, reunir antes del primer bucle las hijas de las parejas de \texttt{van} y no multiplicar P(X \textbar{} c) para ellas; su única contribución debe ser el factor del enlace.

Saber cuál de los dos miembros de la pareja es la hija depende de un convenio que hoy no es uniforme: \texttt{score\allowbreak{}Augmented\allowbreak{}Net} pone la hija primero en los enlaces ya añadidos y la segunda en los candidatos (\hyperref[err:129]{error~129}). Con el arreglo de esa ficha (todas las parejas como \{hija, padre\}), la hija es siempre \texttt{arc{[}0{]}}; por eso conviene hacer los dos cambios a la vez.

\texttt{predict\allowbreak{}Class\allowbreak{}Value\allowbreak{}Augmented\allowbreak{}NB} es el único sitio que combina los dos tipos de factor; \texttt{predict\allowbreak{}Class\allowbreak{}Value}, que usa el Bayes ingenuo selectivo, no tiene enlaces entre características y no está afectado. Conviene arreglarlo junto con \hyperref[err:127]{error~127}, \hyperref[err:128]{error~128}, \hyperref[err:129]{error~129} y \hyperref[err:133]{error~133}, porque todos cambian qué enlaces aprende el clasificador de superpadre.

Prueba de regresión: con un enlace cuyo factor P(hija \textbar{} padre, c) valga lo mismo para todas las clases, la predicción de \texttt{predict\allowbreak{}Class\allowbreak{}Value\allowbreak{}Augmented\allowbreak{}NB} debe ser la del Bayes ingenuo sin esa característica; y sobre \texttt{BN-asia10k.\allowbreak{}csv} con clase \texttt{Bronchitis}, la exactitud con el enlace Smoker→Dyspnea debe ser la del clasificador en que Dyspnea solo aporta P(Dyspnea \textbar{} Smoker, c).

\setcounter{subsection}{130}
\subsection[La validación cruzada de dos clasificadores no prueba un tercio de casos]{La «validación cruzada de diez pliegues» con que los clasificadores selectivo y de superpadre eligen enlaces aparta diez ventanas al azar que se solapan y deja sin evaluar más de un tercio de los casos}\label{err:131}
\begin{ficha}
\fichakey{Estado}Presente
\fichakey{Módulo}learning.metric
\fichakey{Paquete}org.openmarkov.learning.metric.cmi.accuracy
\fichakey{Ficheros}\path{learning.metric/src/main/java/org/openmarkov/learning/metric/cmi/accuracy/Dataset.java:42-68} \newline \path{learning.metric/src/main/java/org/openmarkov/learning/metric/cmi/accuracy/Accuracy.java:130-162}
\fichakey{Comprobado}Leyendo el código y ejecutándolo
\fichakey{Esfuerzo}Pequeño (horas)
\fichakey{Decisión de equipo}No
\fichakey{Relacionados}\hyperref[err:125]{error~125}, \hyperref[err:127]{error~127}, \hyperref[err:128]{error~128}, \hyperref[err:129]{error~129}, \hyperref[err:132]{error~132}, \hyperref[err:135]{error~135}
\end{ficha}
\paragraph{Qué pasa}
El usuario llega por el menú de aprendizaje → tipo discriminativo → «Selective naive bayes» o «Superparent naive bayes», y también por la validación cruzada de \texttt{bn\allowbreak{}Evaluation} (\texttt{Cross\allowbreak{}Validation\allowbreak{}Dialog}), que ofrece los mismos algoritmos. Los enlaces que se añaden o se quitan se eligen con una exactitud que no es la de una validación cruzada: más de un tercio de los casos no se prueba nunca y otros se prueban varias veces.

La métrica \texttt{Accuracy} (exactitud: proporción de casos cuya clase se acierta) es la que usan esos dos clasificadores para decidir qué enlaces añadir o quitar. Para medir la exactitud sin examinar al clasificador con los mismos casos con que se entrena, construye un \texttt{Dataset} que, según su javadoc, divide la base de casos en k pliegues (k = \texttt{Accuracy.\allowbreak{}KFOLD} = 10). En una validación cruzada de k pliegues cada caso cae en exactamente una de las k partes de prueba.

\texttt{Dataset.\allowbreak{}initialize\allowbreak{}Datasets} no hace eso. Para cada una de las diez iteraciones elige un inicio aleatorio \texttt{max\allowbreak{}Case} (generador con semilla fija 42) y toma como prueba la ventana contigua \texttt{{[}max\allowbreak{}Case,\allowbreak{}\ max\allowbreak{}Case\ +\ n/\allowbreak{}10)}; el entrenamiento es el resto.

Cada ventana se sortea independientemente de las demás, así que se solapan: hay casos que caen en varias y casos que no caen en ninguna. Además, \texttt{next\allowbreak{}Int(\allowbreak{}n\ -\ sample\allowbreak{}Size\ -\ 1)} impide que la ventana llegue a los dos últimos casos, y al ser ventanas contiguas, si la base viene ordenada (por ejemplo, por clase) cada ventana es una muestra sesgada.

\texttt{compute\allowbreak{}NBNet\allowbreak{}Accuracy} y \texttt{compute\allowbreak{}Augmented\allowbreak{}Net\allowbreak{}Accuracy} dividen los aciertos entre \texttt{test{[}0{]}.\allowbreak{}length\ *\ KFOLD}. Como todas las ventanas tienen el mismo tamaño, el número resultante es la media de diez exactitudes medidas en ventanas aleatorias (un caso probado dos veces cuenta dos veces); no es la exactitud de una validación cruzada. Con menos de diez casos la ventana tiene tamaño cero y la división da \texttt{Na\allowbreak{}N}.

Medido sobre \texttt{learning.\allowbreak{}algorithm/\allowbreak{}src/\allowbreak{}test/\allowbreak{}resources/\allowbreak{}network/\allowbreak{}BN-asia10k.\allowbreak{}csv} (10 000 casos): las ventanas empiezan en 6752, 808, 5289, 324, 4641, 4337, 7625, 803, 7342 y 2886; 3 691 casos no se prueban nunca, 3 309 una vez y 3 000 dos o más veces.
\paragraph{El código}
learning.metric/src/main/java/org/openmarkov/learning/metric/cmi/accuracy/Dataset.java:53-68

\begin{Shaded}
\begin{Highlighting}[]
\KeywordTok{private} \DataTypeTok{void} \FunctionTok{initializeDatasets}\OperatorTok{()\{}
\NormalTok{    IntStream}\OperatorTok{.}\FunctionTok{range}\OperatorTok{(}\DecValTok{0}\OperatorTok{,}\NormalTok{ numOfSamples}\OperatorTok{).}\FunctionTok{forEach}\OperatorTok{(}\NormalTok{it }\OperatorTok{{-}\textgreater{}\{}
        \DataTypeTok{int}\NormalTok{ maxCase }\OperatorTok{=}\NormalTok{ random}\OperatorTok{.}\FunctionTok{nextInt}\OperatorTok{(}\NormalTok{cases}\OperatorTok{.}\FunctionTok{length}\OperatorTok{{-}}\NormalTok{sampleSize}\OperatorTok{{-}}\DecValTok{1}\OperatorTok{);}
\NormalTok{        test}\OperatorTok{[}\NormalTok{it}\OperatorTok{]=} \BuiltInTok{Arrays}\OperatorTok{.}\FunctionTok{copyOfRange}\OperatorTok{(}\NormalTok{cases}\OperatorTok{,}\NormalTok{ maxCase}\OperatorTok{,}\NormalTok{ maxCase}\OperatorTok{+}\NormalTok{sampleSize}\OperatorTok{);}
\NormalTok{        training}\OperatorTok{[}\NormalTok{it}\OperatorTok{]} \OperatorTok{=} \KeywordTok{new} \DataTypeTok{int}\OperatorTok{[}\NormalTok{cases}\OperatorTok{.}\FunctionTok{length}\OperatorTok{{-}}\NormalTok{sampleSize}\OperatorTok{][];}
        \ControlFlowTok{for}\OperatorTok{(}\DataTypeTok{int}\NormalTok{ i }\OperatorTok{=} \DecValTok{0}\OperatorTok{;}\NormalTok{ i}\OperatorTok{\textless{}}\NormalTok{maxCase}\OperatorTok{;}\NormalTok{i}\OperatorTok{++)\{}
\NormalTok{            training}\OperatorTok{[}\NormalTok{it}\OperatorTok{][}\NormalTok{i}\OperatorTok{]=}\NormalTok{cases}\OperatorTok{[}\NormalTok{i}\OperatorTok{];}
        \OperatorTok{\}}
        \ControlFlowTok{for}\OperatorTok{(}\DataTypeTok{int}\NormalTok{ i}\OperatorTok{=}\NormalTok{maxCase}\OperatorTok{+}\NormalTok{sampleSize}\OperatorTok{;}\NormalTok{ i}\OperatorTok{\textless{}}\NormalTok{cases}\OperatorTok{.}\FunctionTok{length}\OperatorTok{;}\NormalTok{i}\OperatorTok{++)\{}
\NormalTok{            training}\OperatorTok{[}\NormalTok{it}\OperatorTok{][}\NormalTok{i}\OperatorTok{{-}}\NormalTok{sampleSize}\OperatorTok{]=}\NormalTok{cases}\OperatorTok{[}\NormalTok{i}\OperatorTok{];}
        \OperatorTok{\}}
    \OperatorTok{\});}
\OperatorTok{\}}
\end{Highlighting}
\end{Shaded}
\paragraph{Cómo arreglarlo}
Construir una partición de verdad: barajar los índices de los casos una vez (con semilla fija, para que sea repetible) y repartirlos en diez partes disjuntas que cubran todos los casos, repartiendo el resto de la división entre las primeras partes. El entrenamiento de cada iteración es la unión de las otras nueve.

Entonces la exactitud total es aciertos / número de casos, y no \texttt{test{[}0{]}.\allowbreak{}length\ *\ KFOLD}, que ya no sería igual al total si las partes no tienen el mismo tamaño. Hay que cambiarlo en \texttt{compute\allowbreak{}NBNet\allowbreak{}Accuracy} y en \texttt{compute\allowbreak{}Augmented\allowbreak{}Net\allowbreak{}Accuracy}.

Hay que decidir también qué hacer con bases de menos de diez casos (menos pliegues o un aviso).

El cambio alterará los modelos que aprenden los dos clasificadores; conviene hacerlo junto con \hyperref[err:125]{error~125} y \hyperref[err:127]{error~127}, que usan las mismas partes.

Prueba: con una base de 1 003 casos, cada índice aparece en exactamente una parte de prueba y en nueve de entrenamiento.

\setcounter{subsection}{131}
\subsection[El superpadre por defecto no aprende nada: propone un enlace a sí mismo]{El clasificador de superpadre, con su opción por defecto, no aprende nada: propone un enlace del superpadre a sí mismo y el cambio entero se rechaza}\label{err:132}
\begin{ficha}
\fichakey{Estado}Presente
\fichakey{Módulo}learning.algorithm
\fichakey{Paquete}org.openmarkov.learning.algorithm.nbderived.spnb
\fichakey{Ficheros}\path{learning.algorithm/src/main/java/org/openmarkov/learning/algorithm/nbderived/spnb/SuperParentNBAlgorithm.java:59-142} \newline \path{learning.core/src/main/java/org/openmarkov/learning/core/algorithm/LearningAlgorithm.java} \newline \path{learning.gui/src/main/java/org/openmarkov/learning/algorithm/nbderived/spnb/gui/SuperParentNaiveBayesParametersDialog.java:26-52}
\fichakey{Comprobado}Leyendo el código y ejecutándolo
\fichakey{Esfuerzo}Pequeño (horas)
\fichakey{Decisión de equipo}No
\fichakey{Corregir antes}\hyperref[err:133]{error~133}: a la vez, como mínimo: arreglado solo este, con la opción de fábrica el superpadre se aplicaría sin compararlo con el Bayes ingenuo.
\fichakey{Relacionados}\hyperref[err:127]{error~127}, \hyperref[err:128]{error~128}, \hyperref[err:129]{error~129}, \hyperref[err:130]{error~130}, \hyperref[err:131]{error~131}, \hyperref[err:133]{error~133}, \hyperref[err:134]{error~134}
\end{ficha}
\paragraph{Qué pasa}
El usuario llega por el menú de aprendizaje → tipo discriminativo → «Superparent naive bayes» con los parámetros por defecto, y también desde la validación cruzada de \texttt{bn\allowbreak{}Evaluation}. En el aprendizaje automático recibe el Bayes ingenuo de partida, sin ningún enlace nuevo y sin aviso.

«Superparent naive bayes» parte del Bayes ingenuo, elige una característica como «superpadre» (la que más exactitud da) y cuelga de ella otras características que todavía no tienen más padre que la clase (las «huérfanas»). En \texttt{Super\allowbreak{}Parent\allowbreak{}NBAlgorithm.\allowbreak{}init} la lista \texttt{orphans} se llena con \textbf{todas} las características, y en \texttt{get\allowbreak{}Optimal\allowbreak{}Edit}, al elegir el superpadre, se añade a \texttt{super\allowbreak{}Parents} pero no se quita de \texttt{orphans}.

Con la opción «Unique superparent» (\texttt{same\allowbreak{}SP}), que es la que el diálogo de parámetros trae marcada por defecto, se construye un único cambio compuesto (\texttt{List\allowbreak{}PNEdit}) con un \texttt{Add\allowbreak{}Link\allowbreak{}Edit} del superpadre a cada huérfana, incluida ella misma. Una red no admite un enlace de un nodo consigo mismo, así que el cambio entero se rechaza. Como \texttt{orphans} ya se ha vaciado, en la siguiente llamada no hay nada que proponer y el aprendizaje termina con el Bayes ingenuo inicial.

En el aprendizaje automático (\texttt{Learning\allowbreak{}Algorithm.\allowbreak{}run}) el rechazo se captura y solo se escribe en el registro a nivel de depuración. En el aprendizaje interactivo, aplicar esa propuesta hace que \texttt{Interactive\allowbreak{}Learning\allowbreak{}Dialog} convierta el \texttt{Do\allowbreak{}Edit\allowbreak{}Exception} en \texttt{Unrecoverable\allowbreak{}Exception} (esto último comprobado solo leyendo).

Medido sobre \texttt{learning.\allowbreak{}algorithm/\allowbreak{}src/\allowbreak{}test/\allowbreak{}resources/\allowbreak{}network/\allowbreak{}BN-asia10k.\allowbreak{}csv} (clase \texttt{Bronchitis}, alfa 0,5, \texttt{same\allowbreak{}SP\ =\ true}):

\begin{Verbatim}[breaklines,breakanywhere,fontsize=\small]
step 0 [Add link ... X-ray ... to ... X-ray ..., Add link ... X-ray ... to ... Dyspnea ..., ...] {0.432300} REFUSED CannotDoEditException
final links(7) [Bronchitis->Dyspnea, Bronchitis->LungCancer, Bronchitis->Smoker, Bronchitis->Tuberculosis, Bronchitis->TuberculosisOrCancer, Bronchitis->VisitToAsia, Bronchitis->X-ray]
\end{Verbatim}

Con \texttt{same\allowbreak{}SP\ =\ false} (de una en una) el enlace del superpadre consigo mismo también se puntúa como candidato. En el aprendizaje automático lo descarta la comprobación de restricciones (\texttt{is\allowbreak{}Allowed}), y en la misma base el algoritmo aprendió \texttt{X-ray\ -\textgreater{}\ Smoker} sin problema. Podría aparecer como propuesta en el interactivo si el usuario desmarca «solo cambios permitidos».
\paragraph{El código}
learning.algorithm/src/main/java/org/openmarkov/learning/algorithm/nbderived/spnb/SuperParentNBAlgorithm.java:104-117

\begin{Shaded}
\begin{Highlighting}[]
\ControlFlowTok{if} \OperatorTok{(}\NormalTok{bestParent }\OperatorTok{!=} \KeywordTok{null} \OperatorTok{\&\&} \OperatorTok{!}\NormalTok{orphans}\OperatorTok{.}\FunctionTok{isEmpty}\OperatorTok{())} \OperatorTok{\{}
\NormalTok{    superParents}\OperatorTok{.}\FunctionTok{add}\OperatorTok{(}\NormalTok{bestParent}\OperatorTok{);}
    \DataTypeTok{final} \BuiltInTok{Node}\NormalTok{ selectedParent }\OperatorTok{=}\NormalTok{ bestParent}\OperatorTok{;}
    \ControlFlowTok{if} \OperatorTok{(}\KeywordTok{this}\OperatorTok{.}\FunctionTok{sameSP}\OperatorTok{)} \OperatorTok{\{}
        \CommentTok{// sameSP mode: connect bestParent to ALL orphans via a compound edit}
        \BuiltInTok{List}\OperatorTok{\textless{}}\NormalTok{PNEdit}\OperatorTok{\textgreater{}}\NormalTok{ linkEdits }\OperatorTok{=}\NormalTok{ orphans}\OperatorTok{.}\FunctionTok{stream}\OperatorTok{()}
                \OperatorTok{.}\FunctionTok{map}\OperatorTok{(}\NormalTok{orphan }\OperatorTok{{-}\textgreater{}} \OperatorTok{(}\NormalTok{PNEdit}\OperatorTok{)} \KeywordTok{new} \FunctionTok{AddLinkEdit}\OperatorTok{(}\NormalTok{probNet}\OperatorTok{,}\NormalTok{ selectedParent}\OperatorTok{.}\FunctionTok{getVariable}\OperatorTok{(),}
\NormalTok{                        orphan}\OperatorTok{.}\FunctionTok{getVariable}\OperatorTok{(),} \KeywordTok{true}\OperatorTok{))}
                \OperatorTok{.}\FunctionTok{toList}\OperatorTok{();}
\NormalTok{        ListPNEdit compoundEdit }\OperatorTok{=} \KeywordTok{new} \FunctionTok{ListPNEdit}\OperatorTok{(}\NormalTok{probNet}\OperatorTok{,}\NormalTok{ linkEdits}\OperatorTok{);}
\NormalTok{        bestEditProposal }\OperatorTok{=} \KeywordTok{new} \FunctionTok{LearningEditProposal}\OperatorTok{(}\NormalTok{compoundEdit}\OperatorTok{,}
                \KeywordTok{new} \FunctionTok{ScoreEditMotivation}\OperatorTok{(}\NormalTok{bestPartialScore}\OperatorTok{));}
\NormalTok{        orphans}\OperatorTok{.}\FunctionTok{clear}\OperatorTok{();}
\NormalTok{        currentAccuracy }\OperatorTok{=}\NormalTok{ bestPartialScore}\OperatorTok{;}
\end{Highlighting}
\end{Shaded}
\paragraph{Cómo arreglarlo}
Quitar el superpadre de \texttt{orphans} al elegirlo (antes de construir el cambio compuesto y antes del bucle de una en una), o filtrar \texttt{orphan\ !=\ selected\allowbreak{}Parent} en los dos modos.

Revisar a la vez qué pasa si el cambio compuesto se rechaza por otra razón: hoy \texttt{orphans.\allowbreak{}clear(\allowbreak{})} se hace antes de saber si se aplicó.

Prueba de regresión: ejecutar el algoritmo con \texttt{same\allowbreak{}SP\ =\ true} sobre \texttt{BN-asia10k.\allowbreak{}csv} y comprobar que la red final tiene enlaces entre características y ninguno de un nodo consigo mismo; y con \texttt{same\allowbreak{}SP\ =\ false}, que la lista de candidatas nunca contiene el enlace superpadre → superpadre.

\setcounter{subsection}{132}
\subsection[El superpadre acepta el primer enlace aunque empeore al Bayes ingenuo]{El clasificador de superpadre no mide nunca el Bayes ingenuo de partida: acepta el primer superpadre y el primer enlace aunque den menos exactitud que no añadir nada}\label{err:133}
\begin{ficha}
\fichakey{Estado}Presente
\fichakey{Módulo}learning.algorithm
\fichakey{Paquete}org.openmarkov.learning.algorithm.nbderived.spnb
\fichakey{Ficheros}\path{learning.algorithm/src/main/java/org/openmarkov/learning/algorithm/nbderived/spnb/SuperParentNBAlgorithm.java:31-47} \newline \path{learning.algorithm/src/main/java/org/openmarkov/learning/algorithm/nbderived/spnb/SuperParentNBAlgorithm.java:59-74} \newline \path{learning.algorithm/src/main/java/org/openmarkov/learning/algorithm/nbderived/spnb/SuperParentNBAlgorithm.java:84-142}
\fichakey{Comprobado}Leyendo el código y ejecutándolo
\fichakey{Esfuerzo}Pequeño (horas)
\fichakey{Decisión de equipo}No
\fichakey{Corregir antes}\hyperref[err:128]{error~128}: con él, todos los enlaces bajan la exactitud medida, así que comparando con el Bayes ingenuo no se añadiría ninguno.
\fichakey{Tapado por}\hyperref[err:132]{error~132}: con «Unique superparent» marcada, que es la opción de fábrica, no se aprende ningún enlace.
\fichakey{Relacionados}\hyperref[err:127]{error~127}, \hyperref[err:128]{error~128}, \hyperref[err:129]{error~129}, \hyperref[err:130]{error~130}, \hyperref[err:132]{error~132}, \hyperref[err:134]{error~134}, \hyperref[err:141]{error~141}
\end{ficha}
\paragraph{Qué pasa}
El usuario llega por Herramientas («Tools») → «Learning» → tipo «Discriminative» → «Superparent naive bayes», y también por la validación cruzada de \texttt{bn\allowbreak{}Evaluation}. Con «Unique superparent» desmarcada, el algoritmo añade un enlace aunque, según su propia medida, el clasificador con ese enlace acierte menos casos que el Bayes ingenuo del que partió. Con la opción marcada, que es la de fábrica, hoy no se aprende nada por otro motivo (\hyperref[err:132]{error~132}); arreglado eso, el superpadre con todas sus hijas se aplicaría también sin compararlo con el Bayes ingenuo.

El javadoc de la clase describe el algoritmo publicado: partir del Bayes ingenuo, \textbf{evaluarlo}, elegir el superpadre que más aumente la exactitud (la proporción de casos cuya clase se acierta), probar un enlace de él a cada característica huérfana y quedarse con el mejor solo si mejora la exactitud; si no, devolver el clasificador. El paso de evaluar el clasificador de partida no está en el código: \texttt{current\allowbreak{}Accuracy} empieza en 0,0 (línea 47) y nada le asigna la exactitud del Bayes ingenuo; \texttt{init} tampoco la pone a cero, así que una segunda llamada a \texttt{init} (por ejemplo, «Finish» en el aprendizaje interactivo, \hyperref[err:141]{error~141}) conserva la de antes, que en el aprendizaje interactivo ya ha cambiado solo con llenar la tabla (\hyperref[err:134]{error~134}).

En \texttt{get\allowbreak{}Optimal\allowbreak{}Edit}, \texttt{best\allowbreak{}Partial\allowbreak{}Score} empieza valiendo \texttt{current\allowbreak{}Accuracy}:

\begin{itemize}
\tightlist
\item
  En la primera llamada cualquier superpadre cumple \texttt{add\allowbreak{}Score\ \textgreater{}=\ best\allowbreak{}Partial\allowbreak{}Score}, porque la exactitud no es negativa, y se elige el mejor aunque quede por debajo del Bayes ingenuo. \texttt{best\allowbreak{}Partial\allowbreak{}Score} pasa a ser la exactitud de la red con ese superpadre como padre de todas las características que aún no tienen segundo padre.
\item
  Con «Unique superparent» se propone el cambio con todos esos enlaces y \texttt{current\allowbreak{}Accuracy} toma esa exactitud.
\item
  Sin ella, cada enlace S→hija se acepta si su exactitud supera a \texttt{best\allowbreak{}Partial\allowbreak{}Score}, es decir, a la de la red con el superpadre como padre de todas ellas, no a la del Bayes ingenuo.
\end{itemize}

En los pasos siguientes la comparación es con la exactitud del paso anterior, así que lo aceptado de más en el primero ya no se corrige.

Medido con el código de hoy, alfa 0,5 y «Unique superparent» desmarcada, calculando con la misma métrica la exactitud del Bayes ingenuo sin enlaces (\texttt{Accuracy.\allowbreak{}compute\allowbreak{}Augmented\allowbreak{}Net\allowbreak{}Accuracy} con la lista de enlaces vacía):

\begin{itemize}
\tightlist
\item
  \texttt{learning.\allowbreak{}algorithm/\allowbreak{}src/\allowbreak{}test/\allowbreak{}resources/\allowbreak{}network/\allowbreak{}BN-asia10k.\allowbreak{}csv}, clase \texttt{Bronchitis}: Bayes ingenuo 0,8444; el mejor superpadre, X-ray, 0,4323; añade X-ray→Smoker con 0,8253.
\item
  \texttt{integration\allowbreak{}Tests/\allowbreak{}src/\allowbreak{}main/\allowbreak{}resources/\allowbreak{}networks/\allowbreak{}learning/\allowbreak{}learn\allowbreak{}Test\allowbreak{}Data\allowbreak{}Base.\allowbreak{}csv}, clase \texttt{A}: Bayes ingenuo 0,7860; superpadre B, 0,3670; añade B→E con 0,7220.
\item
  \texttt{learning.\allowbreak{}algorithm/\allowbreak{}src/\allowbreak{}test/\allowbreak{}resources/\allowbreak{}network/\allowbreak{}Head\allowbreak{}To\allowbreak{}Head2.\allowbreak{}csv}, clase \texttt{C}: Bayes ingenuo 0,6756; superpadre B, 0,0290; añade B→F con 0,4800.
\item
  \texttt{core/\allowbreak{}src/\allowbreak{}test/\allowbreak{}resources/\allowbreak{}nets/\allowbreak{}Cataract-RRDdataset\allowbreak{}Pre\allowbreak{}Proc3.\allowbreak{}csv}, clase \texttt{success}: Bayes ingenuo 0,9642; superpadre opType, 0,0321; añade opType→breakType con 0,8302.
\end{itemize}

En las cuatro el algoritmo se detiene después de ese enlace. Hoy todos los enlaces bajan la exactitud medida por \hyperref[err:128]{error~128}, así que, partiendo de la exactitud del Bayes ingenuo, en estas bases no se añadiría ninguno. El defecto no depende de eso: con la métrica corregida, un primer superpadre o un primer enlace que no mejore el Bayes ingenuo se seguiría aceptando (esto último está leído en las comparaciones de \texttt{get\allowbreak{}Optimal\allowbreak{}Edit}).
\paragraph{El código}
learning.algorithm/src/main/java/org/openmarkov/learning/algorithm/nbderived/spnb/SuperParentNBAlgorithm.java:31-47

\begin{Shaded}
\begin{Highlighting}[]
\CommentTok{/**}
 \CommentTok{*}\NormalTok{         SUPERPARENT ALGORITHM}
 \CommentTok{*}             \CommentTok{0. }\NormalTok{Initialize network to naive Bayes}\CommentTok{.}
 \CommentTok{*}             \CommentTok{1.}\NormalTok{ Evaluate the current classifier}\CommentTok{.}
 \CommentTok{*}             \CommentTok{2.}\NormalTok{ Consider making each node a SuperParent}\CommentTok{.}\NormalTok{  Let ASP be the SuperParent which increases accuracy the most}\CommentTok{.}
 \CommentTok{*}             \CommentTok{3.}\NormalTok{ Consider an arc from ASP to each orphan}\CommentTok{.}\NormalTok{ If the best such arc improves accuracy}\CommentTok{,}\NormalTok{ keep it and go to }\CommentTok{2.}
 \CommentTok{*}\NormalTok{             Else}\CommentTok{:}\NormalTok{ Return the current classifier}\CommentTok{.}
 \CommentTok{*/}
\KeywordTok{...}
\KeywordTok{private} \DataTypeTok{double}\NormalTok{ currentAccuracy }\OperatorTok{=} \FloatTok{0.0}\OperatorTok{;}
\end{Highlighting}
\end{Shaded}

learning.algorithm/src/main/java/org/openmarkov/learning/algorithm/nbderived/spnb/SuperParentNBAlgorithm.java:86-129

\begin{Shaded}
\begin{Highlighting}[]
\DataTypeTok{double}\NormalTok{ bestPartialScore }\OperatorTok{=}\NormalTok{ currentAccuracy}\OperatorTok{;}
\KeywordTok{...}
\ControlFlowTok{for} \OperatorTok{(}\BuiltInTok{Node}\NormalTok{ nodeParent }\OperatorTok{:} \FunctionTok{subtractListFromNonRootNode}\OperatorTok{(}\NormalTok{superParents}\OperatorTok{))} \OperatorTok{\{}
\NormalTok{    AddLinkEdit addLink }\OperatorTok{=} \KeywordTok{new} \FunctionTok{AddLinkEdit}\OperatorTok{(}\NormalTok{probNet}\OperatorTok{,} \FunctionTok{getRootNode}\OperatorTok{().}\FunctionTok{getVariable}\OperatorTok{(),}\NormalTok{ nodeParent}\OperatorTok{.}\FunctionTok{getVariable}\OperatorTok{(),} \KeywordTok{true}\OperatorTok{);}
    \DataTypeTok{double}\NormalTok{ addScore }\OperatorTok{=}\NormalTok{ metric}\OperatorTok{.}\FunctionTok{getScore}\OperatorTok{(}\NormalTok{addLink}\OperatorTok{);}
    \ControlFlowTok{if} \OperatorTok{(}\KeywordTok{...} \OperatorTok{\&\&} \OperatorTok{(}\NormalTok{addScore }\OperatorTok{\textgreater{}=}\NormalTok{ bestPartialScore }\OperatorTok{||} \OperatorTok{!}\NormalTok{onlyPositiveEdits}\OperatorTok{))} \OperatorTok{\{}
\NormalTok{        bestPartialScore }\OperatorTok{=}\NormalTok{ addScore}\OperatorTok{;}
\NormalTok{        bestParent }\OperatorTok{=}\NormalTok{ nodeParent}\OperatorTok{;}
    \OperatorTok{\}}
\OperatorTok{\}}
\KeywordTok{...}
    \ControlFlowTok{for} \OperatorTok{(}\BuiltInTok{Node}\NormalTok{ nodeChild }\OperatorTok{:}\NormalTok{ orphans}\OperatorTok{)} \OperatorTok{\{}
\NormalTok{        AddLinkEdit addLink }\OperatorTok{=} \KeywordTok{new} \FunctionTok{AddLinkEdit}\OperatorTok{(}\NormalTok{probNet}\OperatorTok{,}\NormalTok{ bestParent}\OperatorTok{.}\FunctionTok{getVariable}\OperatorTok{(),}\NormalTok{ nodeChild}\OperatorTok{.}\FunctionTok{getVariable}\OperatorTok{(),} \KeywordTok{true}\OperatorTok{);}
        \DataTypeTok{double}\NormalTok{ addScore }\OperatorTok{=}\NormalTok{ metric}\OperatorTok{.}\FunctionTok{getScore}\OperatorTok{(}\NormalTok{addLink}\OperatorTok{);}
        \ControlFlowTok{if} \OperatorTok{(}\KeywordTok{...} \OperatorTok{\&\&} \OperatorTok{(}\NormalTok{addScore }\OperatorTok{\textgreater{}}\NormalTok{ bestPartialScore }\OperatorTok{||} \OperatorTok{!}\NormalTok{onlyPositiveEdits}\OperatorTok{))} \OperatorTok{\{}
\NormalTok{            bestEdit }\OperatorTok{=}\NormalTok{ addLink}\OperatorTok{;}
\NormalTok{            bestPartialScore }\OperatorTok{=}\NormalTok{ addScore}\OperatorTok{;}
        \OperatorTok{\}}
    \OperatorTok{\}}
\end{Highlighting}
\end{Shaded}
\paragraph{Cómo arreglarlo}
Medir el Bayes ingenuo al empezar y guardarlo en \texttt{current\allowbreak{}Accuracy}, en \texttt{init} (después de \texttt{set\allowbreak{}Relations\allowbreak{}For\allowbreak{}Root\allowbreak{}Variable} y de configurar la métrica), por ejemplo con \texttt{Accuracy.\allowbreak{}compute\allowbreak{}Augmented\allowbreak{}Net\allowbreak{}Accuracy} sobre la lista vacía de enlaces; así también se repone en cada \texttt{init}.

Separar la puntuación del superpadre de la del clasificador actual: el superpadre se elige por su exactitud, pero un enlace S→hija solo debe aceptarse si su exactitud supera a \texttt{current\allowbreak{}Accuracy}, y no a la del superpadre con todas sus hijas. En el modo «Unique superparent», que añade todos los enlaces de una vez, el cambio compuesto solo debe proponerse si su exactitud supera a \texttt{current\allowbreak{}Accuracy}. Hay que mantener la condición \texttt{\textbar{}\textbar{}\ !only\allowbreak{}Positive\allowbreak{}Edits}, que en el aprendizaje interactivo deja ver también los cambios que no mejoran.

Conviene hacerlo con \hyperref[err:128]{error~128}, \hyperref[err:129]{error~129} y \hyperref[err:127]{error~127}, que cambian las exactitudes con las que se compara, y con \hyperref[err:132]{error~132}, que hoy impide aprender nada con la opción de fábrica.

Prueba de regresión: con una métrica de prueba (una subclase de \texttt{Accuracy}) que dé al Bayes ingenuo una exactitud mayor que la de cualquier superpadre y cualquier enlace, el algoritmo, en los dos modos, debe devolver el Bayes ingenuo sin enlaces entre características.

\setcounter{subsection}{133}
\subsection[Superpadre interactivo: llenar la tabla avanza el algoritmo y la vacía]{En el aprendizaje interactivo con el clasificador de superpadre, llenar la tabla hace avanzar el algoritmo como si se hubieran aplicado las propuestas: «Block edit» o volver a marcar «Only positive edits» dejan la tabla vacía y dan el aprendizaje por terminado, y «Finish» nada más abrir no aprende ningún enlace}\label{err:134}
\begin{ficha}
\fichakey{Estado}Presente
\fichakey{Módulo}learning.algorithm
\fichakey{Paquete}org.openmarkov.learning.algorithm.nbderived.spnb
\fichakey{Ficheros}\path{learning.algorithm/src/main/java/org/openmarkov/learning/algorithm/nbderived/spnb/SuperParentNBAlgorithm.java:41-47} \newline \path{learning.algorithm/src/main/java/org/openmarkov/learning/algorithm/nbderived/spnb/SuperParentNBAlgorithm.java:59-74} \newline \path{learning.algorithm/src/main/java/org/openmarkov/learning/algorithm/nbderived/spnb/SuperParentNBAlgorithm.java:84-142} \newline \path{learning.algorithm/src/main/java/org/openmarkov/learning/algorithm/nbderived/common/DiscriminativeAlgorithm.java:81-90} \newline \path{learning.gui/src/main/java/org/openmarkov/learning/gui/interactive/InteractiveLearningDialog.java:298-352} \newline \path{learning.gui/src/main/java/org/openmarkov/learning/gui/interactive/InteractiveLearningDialog.java:403-427} \newline \path{learning.core/src/main/java/org/openmarkov/learning/core/algorithm/LearningAlgorithm.java:88-103}
\fichakey{Comprobado}Leyendo el código y ejecutándolo
\fichakey{Esfuerzo}Medio (días)
\fichakey{Decisión de equipo}\textcolor{decisionrojo}{\textbf{Sí.}} Si la tabla del aprendizaje interactivo con el superpadre debe enseñar solo alternativas para el paso en curso o también los pasos siguientes
\fichakey{Tapado por}\hyperref[err:132]{error~132}: con «Unique superparent» marcada, que es la opción de fábrica, no se aprende ningún enlace.
\fichakey{Relacionados}\hyperref[err:132]{error~132}, \hyperref[err:133]{error~133}, \hyperref[err:140]{error~140}, \hyperref[err:141]{error~141}
\end{ficha}
\paragraph{Qué pasa}
El usuario llega por Herramientas («Tools») → «Learning» → tipo «Discriminative» → «Superparent naive bayes» → «Options», desmarca «Unique superparent» (con la opción marcada, que es la de fábrica, hoy no se aprende nada por \hyperref[err:132]{error~132}), deja «Interactive learning», que es la opción marcada por omisión, y pulsa «Learn». La ventana «Interactive learning» se abre con una propuesta: el primer enlace que añadiría el aprendizaje automático. A partir de ahí:

\begin{itemize}
\tightlist
\item
  Si pulsa «Block edit» sobre esa propuesta, la tabla se queda vacía y la ventana da el aprendizaje por terminado: desactiva «Apply edit» y «Complete phase», y «Finish» ya no aprende. Hay otros enlaces que el algoritmo ofrecería.
\item
  Si desmarca «Only positive edits», la tabla enseña varios enlaces, cada uno desde un superpadre distinto y ninguno desde el de la propuesta inicial; si la vuelve a marcar, la tabla se queda vacía y la ventana da el aprendizaje por terminado.
\item
  Si pulsa «Finish» nada más abrir, la red no recibe ningún enlace entre características, cuando el aprendizaje automático sí añade uno.
\end{itemize}

El clasificador de superpadre guarda el estado de su búsqueda en tres campos: \texttt{super\allowbreak{}Parents} (las características ya usadas como superpadre), \texttt{orphans} (las que aún no tienen un segundo padre) y \texttt{current\allowbreak{}Accuracy} (la exactitud que hay que superar). \texttt{get\allowbreak{}Optimal\allowbreak{}Edit} los cambia en cuanto encuentra una propuesta, antes de que se aplique nada: añade el superpadre elegido a \texttt{super\allowbreak{}Parents} (aunque después ningún enlace desde él pase la comparación), quita la hija propuesta de \texttt{orphans} (o vacía la lista en el modo «Unique superparent») y pone en \texttt{current\allowbreak{}Accuracy} la exactitud de la propuesta.

El aprendizaje automático (\texttt{Learning\allowbreak{}Algorithm.\allowbreak{}run}) aplica cada propuesta nada más pedirla, y no lo nota. La ventana, en cambio, cada vez que rehace la tabla (\texttt{Interactive\allowbreak{}Learning\allowbreak{}Dialog.\allowbreak{}update\allowbreak{}Editions\allowbreak{}Table}: al abrirse, tras «Apply edit», «Block edit» o «Show blocked» y al cambiar cada casilla) llama a \texttt{get\allowbreak{}Best\allowbreak{}Edit} y después a \texttt{get\allowbreak{}Next\allowbreak{}Edit} hasta tener 50 propuestas. \texttt{Discriminative\allowbreak{}Algorithm.\allowbreak{}get\allowbreak{}Best\allowbreak{}Edit} solo vacía el historial de propuestas ya consideradas, y \texttt{get\allowbreak{}Next\allowbreak{}Edit} vuelve a llamar a \texttt{get\allowbreak{}Optimal\allowbreak{}Edit}, así que cada llamada avanza el estado como si la propuesta anterior se hubiera aplicado. El refresco siguiente parte de ese estado: deja fuera los superpadres ya propuestos, no ofrece las hijas ya propuestas y exige superar la exactitud de la última propuesta.

«Finish» pasa por \texttt{init} (\hyperref[err:141]{error~141}), que rehace \texttt{super\allowbreak{}Parents} y \texttt{orphans} pero no \texttt{current\allowbreak{}Accuracy} (\hyperref[err:133]{error~133}): se queda la exactitud que dejó el llenado de la tabla, y ningún superpadre la alcanza.

Es el mismo patrón que el del aprendizaje interactivo con PC (\hyperref[err:140]{error~140}), con otro estado y en otra clase.

Medido con un pequeño programa de prueba que construye el algoritmo como la ventana de aprendizaje (alfa 0,5, «Unique superparent» desmarcada) y repite las llamadas de la ventana (llenar la tabla, bloquear con \texttt{block\allowbreak{}Edit}, volver a llenarla con la casilla desmarcada y marcada, \texttt{learn} como «Finish»), sobre cuatro bases:

\begin{itemize}
\tightlist
\item
  \texttt{learning.\allowbreak{}algorithm/\allowbreak{}src/\allowbreak{}test/\allowbreak{}resources/\allowbreak{}network/\allowbreak{}BN-asia10k.\allowbreak{}csv}, clase \texttt{Bronchitis}: la tabla se abre con X-ray→Smoker (0,8253), igual que el aprendizaje automático, que solo añade ese enlace. Bloqueándolo, la tabla queda vacía; el mismo algoritmo, con esa propuesta bloqueada antes de llenar la tabla, ofrece X-ray→LungCancer (0,8124). Desmarcando «Only positive edits», la tabla enseña seis enlaces desde seis superpadres distintos (TuberculosisOrCancer→Tuberculosis, Tuberculosis→TuberculosisOrCancer, LungCancer→VisitToAsia, Smoker→LungCancer, VisitToAsia→Dyspnea y Dyspnea→X-ray) y ninguno desde X-ray; volviéndola a marcar, la tabla queda vacía. «Finish» nada más abrir: ningún enlace.
\item
  \texttt{integration\allowbreak{}Tests/\allowbreak{}src/\allowbreak{}main/\allowbreak{}resources/\allowbreak{}networks/\allowbreak{}learning/\allowbreak{}learn\allowbreak{}Test\allowbreak{}Data\allowbreak{}Base.\allowbreak{}csv}, clase \texttt{A}: se abre con B→E, lo mismo que aprende el automático; bloqueándolo, tabla vacía (sin el llenado previo se ofrecería B→C); con la casilla desmarcada, tres enlaces desde otros superpadres; al volver a marcarla, vacía; «Finish» nada más abrir, ningún enlace.
\item
  \texttt{learning.\allowbreak{}algorithm/\allowbreak{}src/\allowbreak{}test/\allowbreak{}resources/\allowbreak{}network/\allowbreak{}Head\allowbreak{}To\allowbreak{}Head2.\allowbreak{}csv}, clase \texttt{C}: se abre con B→F, lo mismo que el automático; bloqueándolo, vacía (se ofrecería B→E); desmarcada, dos enlaces; marcada otra vez, vacía; «Finish», ningún enlace.
\item
  \texttt{core/\allowbreak{}src/\allowbreak{}test/\allowbreak{}resources/\allowbreak{}nets/\allowbreak{}Cataract-RRDdataset\allowbreak{}Pre\allowbreak{}Proc3.\allowbreak{}csv}, clase \texttt{success}: se abre con opType→breakType, lo mismo que el automático; bloqueándolo, vacía (se ofrecería opType→loc57); desmarcada, cinco enlaces; marcada otra vez, vacía; «Finish», ningún enlace.
\end{itemize}

En las cuatro bases, aplicar la propuesta de la tabla con «Apply edit» da lo mismo que el aprendizaje automático, que se detiene tras el primer enlace: el adelanto se ve al bloquear, al cambiar las casillas y al pulsar «Finish».

Con «Unique superparent» marcada, en \texttt{BN-asia10k.\allowbreak{}csv} la tabla se abre con la propuesta compuesta que describe \hyperref[err:132]{error~132} y cualquier refresco la deja vacía, porque la primera propuesta ya vació \texttt{orphans}.

La ventana es un \texttt{JDialog} y no se ha abierto: lo que hace con la tabla vacía (botones desactivados, «Finish» sin efecto) está leído en \texttt{update\allowbreak{}Editions\allowbreak{}Table} y \texttt{btn\allowbreak{}Run\allowbreak{}Action\allowbreak{}Performed}.
\paragraph{El código}
learning.algorithm/src/main/java/org/openmarkov/learning/algorithm/nbderived/spnb/SuperParentNBAlgorithm.java:86-139

\begin{Shaded}
\begin{Highlighting}[]
\DataTypeTok{double}\NormalTok{ bestPartialScore }\OperatorTok{=}\NormalTok{ currentAccuracy}\OperatorTok{;}
\KeywordTok{...}
\ControlFlowTok{if} \OperatorTok{(}\NormalTok{bestParent }\OperatorTok{!=} \KeywordTok{null} \OperatorTok{\&\&} \OperatorTok{!}\NormalTok{orphans}\OperatorTok{.}\FunctionTok{isEmpty}\OperatorTok{())} \OperatorTok{\{}
\NormalTok{    superParents}\OperatorTok{.}\FunctionTok{add}\OperatorTok{(}\NormalTok{bestParent}\OperatorTok{);}
    \DataTypeTok{final} \BuiltInTok{Node}\NormalTok{ selectedParent }\OperatorTok{=}\NormalTok{ bestParent}\OperatorTok{;}
    \ControlFlowTok{if} \OperatorTok{(}\KeywordTok{this}\OperatorTok{.}\FunctionTok{sameSP}\OperatorTok{)} \OperatorTok{\{}
        \KeywordTok{...}
\NormalTok{        orphans}\OperatorTok{.}\FunctionTok{clear}\OperatorTok{();}
\NormalTok{        currentAccuracy }\OperatorTok{=}\NormalTok{ bestPartialScore}\OperatorTok{;}
    \OperatorTok{\}} \ControlFlowTok{else} \OperatorTok{\{}
        \ControlFlowTok{for} \OperatorTok{(}\BuiltInTok{Node}\NormalTok{ nodeChild }\OperatorTok{:}\NormalTok{ orphans}\OperatorTok{)} \OperatorTok{\{}
            \KeywordTok{...}
        \OperatorTok{\}}
    \OperatorTok{\}}
\OperatorTok{\}}

\ControlFlowTok{if} \OperatorTok{(}\NormalTok{bestEdit }\OperatorTok{!=} \KeywordTok{null}\OperatorTok{)} \OperatorTok{\{}
\NormalTok{    bestEditProposal }\OperatorTok{=}\NormalTok{ LearningEditProposal}\OperatorTok{.}\FunctionTok{scored}\OperatorTok{(}\NormalTok{bestEdit}\OperatorTok{,}\NormalTok{ bestPartialScore}\OperatorTok{);}
    \FunctionTok{markEditAsConsidered}\OperatorTok{(}\NormalTok{bestEdit}\OperatorTok{);}
\NormalTok{    orphans}\OperatorTok{.}\FunctionTok{remove}\OperatorTok{(}\NormalTok{probNet}\OperatorTok{.}\FunctionTok{getNode}\OperatorTok{(}\NormalTok{bestEdit}\OperatorTok{.}\FunctionTok{getVariableTo}\OperatorTok{()));}
\NormalTok{    currentAccuracy }\OperatorTok{=}\NormalTok{ bestPartialScore}\OperatorTok{;}
\OperatorTok{\}}
\end{Highlighting}
\end{Shaded}

learning.algorithm/src/main/java/org/openmarkov/learning/algorithm/nbderived/common/DiscriminativeAlgorithm.java:81-90

\begin{Shaded}
\begin{Highlighting}[]
\AttributeTok{@Override}
\KeywordTok{public}\NormalTok{ LearningEditProposal }\FunctionTok{getBestEdit}\OperatorTok{(}\DataTypeTok{boolean}\NormalTok{ onlyAllowedEdits}\OperatorTok{,} \DataTypeTok{boolean}\NormalTok{ onlyPositiveEdits}\OperatorTok{)} \OperatorTok{\{}
    \FunctionTok{resetHistory}\OperatorTok{();}
    \ControlFlowTok{return} \FunctionTok{getNextEdit}\OperatorTok{(}\NormalTok{onlyAllowedEdits}\OperatorTok{,}\NormalTok{ onlyPositiveEdits}\OperatorTok{);}
\OperatorTok{\}}

\AttributeTok{@Override}
\KeywordTok{public}\NormalTok{ LearningEditProposal }\FunctionTok{getNextEdit}\OperatorTok{(}\DataTypeTok{boolean}\NormalTok{ onlyAllowedEdits}\OperatorTok{,} \DataTypeTok{boolean}\NormalTok{ onlyPositiveEdits}\OperatorTok{)} \OperatorTok{\{}
    \ControlFlowTok{return} \FunctionTok{getOptimalEdit}\OperatorTok{(}\NormalTok{onlyAllowedEdits}\OperatorTok{,}\NormalTok{ onlyPositiveEdits}\OperatorTok{);}
\OperatorTok{\}}
\end{Highlighting}
\end{Shaded}
\paragraph{Cómo arreglarlo}
Buscar propuestas no debe cambiar el estado del algoritmo; \texttt{super\allowbreak{}Parents}, \texttt{orphans} y \texttt{current\allowbreak{}Accuracy} deben cambiar cuando se aplica (o se deshace) una edición. Dos caminos:

\begin{itemize}
\tightlist
\item
  Actualizarlos al ejecutarse o deshacerse una edición, escuchando las ediciones de la red, y dejar \texttt{get\allowbreak{}Optimal\allowbreak{}Edit} como una búsqueda que solo lee.
\item
  Deducirlos de la red en cada búsqueda: huérfanas son las características cuyo único padre es la clase, superpadres las que ya tienen hijas entre las características, y la exactitud actual la de la red tal como está.
\end{itemize}

Qué debe enseñar la tabla depende de lo que decida el equipo: solo alternativas para el paso en curso (otros superpadres y otras hijas, con el historial de propuestas ya consideradas para no repetirlas), o también los pasos siguientes, en cuyo caso el llenado debe hacerse sobre una copia del estado y separar en la tabla lo que depende de aplicar antes otra fila.

Hay que mantener el resultado del aprendizaje automático y de la validación cruzada de \texttt{bn\allowbreak{}Evaluation}, que usan \texttt{run}. Si una edición se rechaza, el estado no debe cambiar (hoy \texttt{orphans.\allowbreak{}clear(\allowbreak{})} se hace antes de saber si se aplicó, como señala \hyperref[err:132]{error~132}). \texttt{init} debe reponer también \texttt{current\allowbreak{}Accuracy}, que es lo que propone \hyperref[err:133]{error~133}.

Pruebas de regresión, repitiendo las llamadas de la ventana sobre \texttt{BN-asia10k.\allowbreak{}csv} con clase \texttt{Bronchitis} y «Unique superparent» desmarcada: tras llenar la tabla, bloquear la primera propuesta y volver a llenarla no debe dejarla vacía (hoy, sin el llenado previo, se ofrece X-ray→LungCancer); desmarcar y volver a marcar «Only positive edits» debe dar la misma tabla que al abrir; y «Finish» nada más abrir debe dar la misma red que el aprendizaje automático.

\setcounter{subsection}{134}
\subsection[El Bayes ingenuo selectivo hacia delante no da dos veces el mismo modelo]{El Bayes ingenuo selectivo hacia delante aprende modelos distintos en ejecuciones distintas sobre la misma base}\label{err:135}
\begin{ficha}
\fichakey{Estado}Presente
\fichakey{Módulo}learning.algorithm
\fichakey{Paquete}org.openmarkov.learning.algorithm.nbderived.snb
\fichakey{Ficheros}\path{learning.algorithm/src/main/java/org/openmarkov/learning/algorithm/nbderived/snb/SelectiveNBAlgorithm.java:59-96}
\fichakey{Comprobado}Leyendo el código y ejecutándolo
\fichakey{Esfuerzo}Pequeño (horas)
\fichakey{Decisión de equipo}No
\fichakey{Relacionados}\hyperref[err:125]{error~125}, \hyperref[err:126]{error~126}, \hyperref[err:131]{error~131}, \hyperref[err:136]{error~136}
\end{ficha}
\paragraph{Qué pasa}
El usuario llega por el menú de aprendizaje → tipo discriminativo → «Selective naive bayes» con la casilla «Forward» marcada, que es el valor por defecto del diálogo de parámetros; también desde la validación cruzada de \texttt{bn\allowbreak{}Evaluation}. Sobre la misma base de casos, cada ejecución puede dar un modelo final distinto.

«Selective naive bayes» elige qué características conserva: hacia delante parte de la clase sola y en cada paso añade la característica que más exactitud da; hacia atrás parte del Bayes ingenuo completo y quita.

En \texttt{Selective\allowbreak{}NBAlgorithm.\allowbreak{}get\allowbreak{}Optimal\allowbreak{}Edit}, hacia delante, las candidatas se reúnen con \texttt{Collectors.\allowbreak{}to\allowbreak{}Set(\allowbreak{})} (un \texttt{Hash\allowbreak{}Set}), y el método se queda con la última que cumpla \texttt{add\allowbreak{}Score\ \textgreater{}=\ best\allowbreak{}Partial\allowbreak{}Score}. El orden de un \texttt{Hash\allowbreak{}Set} de nodos depende de \texttt{Node.\allowbreak{}hash\allowbreak{}Code}, que depende de \texttt{Variable}, que no redefine \texttt{hash\allowbreak{}Code} (usa la identidad del objeto, es decir, las señas de memoria).

Como la exactitud es un número de aciertos sobre un número fijo de casos, los empates son frecuentes, y cuál gana depende de ese orden.

Medido ejecutando el algoritmo hacia delante sobre \texttt{learning.\allowbreak{}algorithm/\allowbreak{}src/\allowbreak{}test/\allowbreak{}resources/\allowbreak{}network/\allowbreak{}BN-asia10k.\allowbreak{}csv} (clase \texttt{Bronchitis}) en 11 procesos Java distintos. El orden de los pasos cambia de una ejecución a otra (todos con exactitud 0,845799), por ejemplo:

\begin{Verbatim}[breaklines,breakanywhere,fontsize=\small]
ejecución 1: Dyspnea, Tuberculosis, TuberculosisOrCancer, LungCancer, VisitToAsia, X-ray
ejecución 2: Dyspnea, VisitToAsia, TuberculosisOrCancer, X-ray, Tuberculosis, LungCancer
\end{Verbatim}

Y el modelo final también cambia. En 8 ejecuciones seguidas salieron cuatro modelos:

\begin{Verbatim}[breaklines,breakanywhere,fontsize=\small]
5 veces: Dyspnea, LungCancer, Tuberculosis, TuberculosisOrCancer, VisitToAsia, X-ray   (sin Smoker)
1 vez:   Dyspnea, Smoker, TuberculosisOrCancer, VisitToAsia, X-ray
1 vez:   Dyspnea, LungCancer, Smoker, Tuberculosis, VisitToAsia
1 vez:   Dyspnea, Smoker, Tuberculosis, VisitToAsia, X-ray
\end{Verbatim}

Hacia atrás las candidatas son \texttt{get\allowbreak{}Root\allowbreak{}Node(\allowbreak{}).\allowbreak{}get\allowbreak{}Children(\allowbreak{})}, una lista con el orden de inserción, y en tres ejecuciones salió siempre lo mismo: quita Smoker, TuberculosisOrCancer, Tuberculosis, LungCancer, VisitToAsia y X-ray, y deja solo Dyspnea.
\paragraph{El código}
learning.algorithm/src/main/java/org/openmarkov/learning/algorithm/nbderived/snb/SelectiveNBAlgorithm.java:67-89

\begin{Shaded}
\begin{Highlighting}[]
\BuiltInTok{Collection}\OperatorTok{\textless{}}\BuiltInTok{Node}\OperatorTok{\textgreater{}}\NormalTok{ candidates }\OperatorTok{=} \OperatorTok{(!}\NormalTok{forward}\OperatorTok{)} \OperatorTok{?} \FunctionTok{getRootNode}\OperatorTok{().}\FunctionTok{getChildren}\OperatorTok{()} \OperatorTok{:}
        \FunctionTok{getNonRootNodes}\OperatorTok{().}\FunctionTok{stream}\OperatorTok{()}
                         \OperatorTok{.}\FunctionTok{filter}\OperatorTok{(}\NormalTok{n }\OperatorTok{{-}\textgreater{}} \OperatorTok{!}\FunctionTok{getRootNode}\OperatorTok{().}\FunctionTok{getChildren}\OperatorTok{().}\FunctionTok{contains}\OperatorTok{(}\NormalTok{n}\OperatorTok{))}
                         \OperatorTok{.}\FunctionTok{collect}\OperatorTok{(}\NormalTok{Collectors}\OperatorTok{.}\FunctionTok{toSet}\OperatorTok{());}

\ControlFlowTok{for} \OperatorTok{(}\BuiltInTok{Node}\NormalTok{ n1 }\OperatorTok{:}\NormalTok{ candidates}\OperatorTok{)} \OperatorTok{\{}
    \KeywordTok{...}
    \DataTypeTok{double}\NormalTok{ addScore }\OperatorTok{=}\NormalTok{ metric}\OperatorTok{.}\FunctionTok{getScore}\OperatorTok{(}\NormalTok{edit}\OperatorTok{);}
    \ControlFlowTok{if} \OperatorTok{(!}\FunctionTok{isEditAlreadyConsidered}\OperatorTok{(}\NormalTok{edit}\OperatorTok{)} \OperatorTok{\&\&} \OperatorTok{!}\FunctionTok{isBlocked}\OperatorTok{(}\NormalTok{edit}\OperatorTok{)}
            \OperatorTok{\&\&} \OperatorTok{(!}\NormalTok{onlyAllowedEdits }\OperatorTok{||} \FunctionTok{isAllowed}\OperatorTok{(}\NormalTok{edit}\OperatorTok{))}
            \OperatorTok{\&\&} \OperatorTok{(}\NormalTok{addScore }\OperatorTok{\textgreater{}=}\NormalTok{ bestPartialScore }\OperatorTok{||} \OperatorTok{!}\NormalTok{onlyPositiveEdits}\OperatorTok{)}
    \OperatorTok{)} \OperatorTok{\{}
\NormalTok{        bestEdit }\OperatorTok{=}\NormalTok{ edit}\OperatorTok{;}
\NormalTok{        bestPartialScore }\OperatorTok{=}\NormalTok{ addScore}\OperatorTok{;}
    \OperatorTok{\}}
\OperatorTok{\}}
\end{Highlighting}
\end{Shaded}
\paragraph{Cómo arreglarlo}
Recorrer las candidatas en un orden determinista (una lista en el orden de la base de casos, o ordenada por nombre) en las dos direcciones, y fijar una regla de desempate explícita: por ejemplo, la primera en ese orden, con \texttt{\textgreater{}} en vez de \texttt{\textgreater{}=}. Hay que tener en cuenta que \texttt{\textgreater{}=} también sirve hoy para aceptar pasos que no empeoran la exactitud; si se cambia, hay que conservar ese criterio con otra condición.

El mismo tipo de dependencia del orden de un \texttt{Hash\allowbreak{}Map} está en \texttt{Maximum\allowbreak{}Weight\allowbreak{}Spanning\allowbreak{}Tree.\allowbreak{}build} (ver \hyperref[err:136]{error~136}).

Prueba de regresión: ejecutar el algoritmo dos veces sobre la misma base y comprobar que los pasos y el modelo final coinciden. Como el orden de un \texttt{Hash\allowbreak{}Set} puede coincidir dentro de un mismo proceso, conviene crear las variables en órdenes distintos en las dos ejecuciones.

\setcounter{subsection}{135}
\subsection[El árbol aumentado (TAN) orienta sus enlaces distinto en cada aprendizaje]{El clasificador de árbol aumentado devuelve un modelo con los enlaces dirigidos de otra manera cada vez que se aprende}\label{err:136}
\begin{ficha}
\fichakey{Estado}Presente
\fichakey{Módulo}learning.algorithm
\fichakey{Paquete}org.openmarkov.learning.algorithm.nbderived.treeaugmentednb
\fichakey{Ficheros}\path{learning.algorithm/src/main/java/org/openmarkov/learning/algorithm/nbderived/treeaugmentednb/TreeAugmentedNBAlgorithm.java:34-56} \newline \path{learning.algorithm/src/main/java/org/openmarkov/learning/algorithm/nbderived/common/MaximumWeightSpanningTree.java:38-100}
\fichakey{Comprobado}Leyendo el código y ejecutándolo
\fichakey{Esfuerzo}Pequeño (horas)
\fichakey{Decisión de equipo}\textcolor{decisionrojo}{\textbf{Sí.}} Qué regla fija la raíz del árbol (la primera característica de la base de casos, la primera por nombre, una que elija el usuario, o un sorteo con semilla fija)
\fichakey{Relacionados}\hyperref[err:4]{error~4}, \hyperref[err:77]{error~77}, \hyperref[err:135]{error~135}
\end{ficha}
\paragraph{Qué pasa}
El usuario llega por el menú de aprendizaje → pestaña general → tipo de algoritmo «discriminativo» → algoritmo «Tree augmented naive bayes», con aprendizaje automático o interactivo, y también desde la validación cruzada de \texttt{bn\allowbreak{}Evaluation} (\texttt{Cross\allowbreak{}Validation\allowbreak{}Dialog} ofrece el mismo algoritmo). Cada vez que aprende sobre la misma base, la red tiene las mismas parejas unidas pero los enlaces dirigidos de otra manera, y nada lo advierte.

Lo que el clasificador responde sobre la clase apenas cambia: todas las orientaciones de un mismo árbol son equivalentes en lo que dicen sobre la clase, y con el suavizado de los parámetros puede haber diferencias mínimas en los números. Pero la red que el usuario ve, guarda y compara es distinta en cada aprendizaje.

El clasificador TAN (Tree Augmented Naive Bayes, Bayes ingenuo aumentado con un árbol) parte del Bayes ingenuo ---la variable clase es padre de todas las demás, que se llaman características--- y añade entre las características un árbol de peso máximo, donde el peso de cada pareja es su información mutua condicionada a la clase. Ese árbol se construye sin dirección y después hay que orientarlo: se elige una característica como raíz y todos los enlaces se dirigen alejándose de ella.

La raíz se elige al azar con un generador sembrado con el reloj, en \texttt{Tree\allowbreak{}Augmented\allowbreak{}NBAlgorithm.\allowbreak{}get\allowbreak{}Random\allowbreak{}Variable(\allowbreak{})}, llamado desde \texttt{init}. El algoritmo publicado (Friedman, Geiger y Goldszmidt, 1997) permite cualquier raíz; el problema no es la regla sino que no sea repetible.

Medido aprendiendo el TAN tres veces seguidas sobre \texttt{learning.\allowbreak{}algorithm/\allowbreak{}src/\allowbreak{}test/\allowbreak{}resources/\allowbreak{}network/\allowbreak{}BN-asia10k.\allowbreak{}csv} con \texttt{Bronchitis} como clase (enlaces entre características, ordenados):

\begin{Verbatim}[breaklines,breakanywhere,fontsize=\small]
run 0 [Dyspnea->TuberculosisOrCancer, Dyspnea->VisitToAsia, LungCancer->Smoker, TuberculosisOrCancer->LungCancer, TuberculosisOrCancer->Tuberculosis, TuberculosisOrCancer->X-ray]
run 1 [Dyspnea->VisitToAsia, LungCancer->Smoker, LungCancer->TuberculosisOrCancer, TuberculosisOrCancer->Dyspnea, TuberculosisOrCancer->Tuberculosis, TuberculosisOrCancer->X-ray]
run 2 [Dyspnea->VisitToAsia, LungCancer->Smoker, TuberculosisOrCancer->Dyspnea, TuberculosisOrCancer->LungCancer, TuberculosisOrCancer->Tuberculosis, TuberculosisOrCancer->X-ray]
\end{Verbatim}

Hay una segunda fuente posible de variación, no medida: \texttt{Maximum\allowbreak{}Weight\allowbreak{}Spanning\allowbreak{}Tree.\allowbreak{}build} guarda los pesos en un \texttt{Hash\allowbreak{}Map} cuyas claves son conjuntos de nodos, y \texttt{Node.\allowbreak{}hash\allowbreak{}Code} depende de \texttt{Variable}, que no redefine \texttt{hash\allowbreak{}Code} (usa la identidad del objeto). Si dos parejas empatan en peso, el orden en que el algoritmo de Kruskal las considera depende de las señas de memoria, y con él qué árbol sale. Con \texttt{BN-asia10k.\allowbreak{}csv} el árbol sin dirigir salió igual en las tres ejecuciones.
\paragraph{El código}
learning.algorithm/src/main/java/org/openmarkov/learning/algorithm/nbderived/treeaugmentednb/TreeAugmentedNBAlgorithm.java:34-56

\begin{Shaded}
\begin{Highlighting}[]
\AttributeTok{@Override} \KeywordTok{public} \DataTypeTok{void} \FunctionTok{init}\OperatorTok{(}\NormalTok{ModelNetUse modelNetUse}\OperatorTok{)} \OperatorTok{\{}
    \KeywordTok{...}
    \ControlFlowTok{if} \OperatorTok{(!}\NormalTok{mwst}\OperatorTok{.}\FunctionTok{isBuilt}\OperatorTok{())} \OperatorTok{\{}
\NormalTok{        mwst}\OperatorTok{.}\FunctionTok{build}\OperatorTok{(}\NormalTok{probNet}\OperatorTok{,}\NormalTok{ metric}\OperatorTok{,} \FunctionTok{getNonRootNodes}\OperatorTok{());}
\NormalTok{        Variable randomRoot }\OperatorTok{=} \FunctionTok{getRandomVariable}\OperatorTok{();}
\NormalTok{        mwst}\OperatorTok{.}\FunctionTok{redirect}\OperatorTok{(}\NormalTok{randomRoot}\OperatorTok{,}\NormalTok{ probNet}\OperatorTok{);}
    \OperatorTok{\}}
    \KeywordTok{...}
\OperatorTok{\}}

\KeywordTok{protected}\NormalTok{ Variable }\FunctionTok{getRandomVariable}\OperatorTok{()} \OperatorTok{\{}
    \ControlFlowTok{return} \FunctionTok{getNonRootVariables}\OperatorTok{().}\FunctionTok{get}\OperatorTok{(}\KeywordTok{new} \BuiltInTok{Random}\OperatorTok{().}\FunctionTok{nextInt}\OperatorTok{(}\FunctionTok{getNonRootVariables}\OperatorTok{().}\FunctionTok{size}\OperatorTok{()));}
\OperatorTok{\}}
\end{Highlighting}
\end{Shaded}
\paragraph{Cómo arreglarlo}
Sustituir el sorteo por una regla determinista (la que decida el equipo) en \texttt{get\allowbreak{}Random\allowbreak{}Variable}, y renombrar el método si deja de sortear. Las pruebas \texttt{Tree\allowbreak{}Augmented\allowbreak{}NBAlgorithm\allowbreak{}Test} y \texttt{Get\allowbreak{}Random\allowbreak{}Variable\allowbreak{}Regression\allowbreak{}Test} redefinen o llaman a ese método y habrá que ajustarlas.

Para que el resultado sea repetible del todo, \texttt{Maximum\allowbreak{}Weight\allowbreak{}Spanning\allowbreak{}Tree.\allowbreak{}build} también debería desempatar de forma determinista: por ejemplo, ordenando por peso y, a igualdad, por los nombres de la pareja, y usando un \texttt{Linked\allowbreak{}Hash\allowbreak{}Map} o una lista en vez de un \texttt{Hash\allowbreak{}Map} con claves \texttt{Set\textless{}Node\textgreater{}}.

La misma familia de problema aparece en \texttt{Selective\allowbreak{}NBAlgorithm} (ver \hyperref[err:135]{error~135}).

Prueba de regresión: aprender dos veces el TAN sobre \texttt{BN-asia10k.\allowbreak{}csv} con la misma clase y comprobar que los enlaces dirigidos coinciden.

\setcounter{subsection}{136}
\subsection[El EM deja sin aprender los nodos con potencial uniforme]{La maximización de la esperanza no devuelve lo aprendido para los nodos cuyo potencial es uniforme: siguen uniformes en la red resultante}\label{err:137}
\begin{ficha}
\fichakey{Estado}Presente
\fichakey{Módulo}learning.algorithm
\fichakey{Paquete}org.openmarkov.learning.algorithm.em
\fichakey{Ficheros}\path{learning.algorithm/src/main/java/org/openmarkov/learning/algorithm/em/EMAlgorithm.java:119-235} \newline \path{learning.algorithm/src/main/java/org/openmarkov/learning/algorithm/em/EMAlgorithm.java:246-290} \newline \path{core/src/main/java/org/openmarkov/core/model/network/ProbNet.java:396-398}
\fichakey{Comprobado}Leyendo el código y ejecutándolo
\fichakey{Esfuerzo}Pequeño (horas)
\fichakey{Decisión de equipo}No
\end{ficha}
\paragraph{Qué pasa}
El usuario llega por el menú de aprendizaje → algoritmo EM → pestaña de red modelo → partir de la red modelo, con una red modelo en la que algún nodo tiene potencial uniforme. Recibe la red aprendida con esos nodos todavía uniformes, sin aviso.

El algoritmo EM (Expectation Maximization, maximización de la esperanza; «Expectation maximization (EM)» en el menú de aprendizaje) aprende los parámetros de una red con estructura fija, aunque falten datos.

\texttt{EMAlgorithm.\allowbreak{}parametric\allowbreak{}Learning} llama a \texttt{adapt\allowbreak{}Network}, que hace \texttt{prob\allowbreak{}Net.\allowbreak{}copy(\allowbreak{})} ---una copia superficial: nodos nuevos que comparten los mismos objetos potencial--- y recoge en la lista \texttt{potentials} las tablas a aprender. En cada iteración escribe los números nuevos en \texttt{potential.\allowbreak{}get\allowbreak{}Values(\allowbreak{})}. Como para un \texttt{Table\allowbreak{}Potential} ese array es el mismo objeto que tiene la red original, lo aprendido llega a la red que se devuelve (\texttt{prob\allowbreak{}Net}).

Un \texttt{Uniform\allowbreak{}Potential} es el potencial «uniforme»: no es una tabla sino una clase hermana, y se usa cuando nadie ha escrito los números; lo crean, por ejemplo, \texttt{Potential.\allowbreak{}add\allowbreak{}Variable} al añadir un enlace y los ficheros \texttt{.\allowbreak{}pgmx} con \texttt{type="Uniform"}. Para un nodo con ese potencial, \texttt{adapt\allowbreak{}Network} construye un \texttt{Table\allowbreak{}Potential} nuevo, lo añade a \texttt{potentials} y lo pone en el nodo de \textbf{la copia} (\texttt{expanded\allowbreak{}Net}). El EM lo aprende, pero el nodo de \texttt{prob\allowbreak{}Net} sigue con su \texttt{Uniform\allowbreak{}Potential}, y \texttt{parametric\allowbreak{}Learning} devuelve \texttt{prob\allowbreak{}Net}.

Medido con un pequeño programa de prueba: red A → B, A con una tabla y B con \texttt{Uniform\allowbreak{}Potential}, 100 casos en los que B depende claramente de A, alfa 0. Tras \texttt{parametric\allowbreak{}Learning}, A tiene los valores aprendidos \texttt{{[}0.\allowbreak{}8,\allowbreak{}\ 0.\allowbreak{}2{]}} (el mismo objeto, actualizado) y B sigue siendo un \texttt{Uniform\allowbreak{}Potential}.

La prueba existente \texttt{EMAlgorithm\allowbreak{}Tests.\allowbreak{}parametric\allowbreak{}Learning\allowbreak{}With\allowbreak{}Uniform\allowbreak{}Potential\allowbreak{}Cpt\allowbreak{}Does\allowbreak{}Not\allowbreak{}Throw} solo comprueba que no se lanza excepción y dice expresamente que no mira los parámetros aprendidos.
\paragraph{El código}
learning.algorithm/src/main/java/org/openmarkov/learning/algorithm/em/EMAlgorithm.java:246-261

\begin{Shaded}
\begin{Highlighting}[]
\KeywordTok{private} \DataTypeTok{static}\NormalTok{ ProbNet }\FunctionTok{adaptNetwork}\OperatorTok{(}\NormalTok{ProbNet probNet}\OperatorTok{,} \BuiltInTok{List}\OperatorTok{\textless{}}\NormalTok{TablePotential}\OperatorTok{\textgreater{}}\NormalTok{ potentials}\OperatorTok{,}
                                    \BuiltInTok{Map}\OperatorTok{\textless{}}\NormalTok{ICIPotential}\OperatorTok{,} \BuiltInTok{List}\OperatorTok{\textless{}}\NormalTok{TablePotential}\OperatorTok{\textgreater{}\textgreater{}}\NormalTok{ iciSubpotentials}\OperatorTok{)} \OperatorTok{\{}
\NormalTok{    ProbNet expandedNet }\OperatorTok{=}\NormalTok{ probNet}\OperatorTok{.}\FunctionTok{copy}\OperatorTok{();}
    \ControlFlowTok{for} \OperatorTok{(}\NormalTok{Potential potential }\OperatorTok{:}\NormalTok{ expandedNet}\OperatorTok{.}\FunctionTok{getPotentials}\OperatorTok{())} \OperatorTok{\{}
        \ControlFlowTok{if} \OperatorTok{(}\NormalTok{potential}\OperatorTok{.}\FunctionTok{getPotentialRole}\OperatorTok{()} \OperatorTok{==}\NormalTok{ PotentialRole}\OperatorTok{.}\FunctionTok{CONDITIONAL\_PROBABILITY}\OperatorTok{)} \OperatorTok{\{}
            \ControlFlowTok{switch} \OperatorTok{(}\NormalTok{potential}\OperatorTok{)} \OperatorTok{\{}
                \ControlFlowTok{case}\NormalTok{ UniformPotential uniformPotential }\OperatorTok{{-}\textgreater{}} \OperatorTok{\{}
                    \KeywordTok{...}
\NormalTok{                    TablePotential newPotential }\OperatorTok{=} \KeywordTok{new} \FunctionTok{TablePotential}\OperatorTok{(}\NormalTok{uniformPotential}\OperatorTok{.}\FunctionTok{getVariables}\OperatorTok{(),}
\NormalTok{                                                                     uniformPotential}\OperatorTok{.}\FunctionTok{getPotentialRole}\OperatorTok{());}
\NormalTok{                    potentials}\OperatorTok{.}\FunctionTok{add}\OperatorTok{(}\NormalTok{newPotential}\OperatorTok{);}
\NormalTok{                    expandedNet}\OperatorTok{.}\FunctionTok{getNode}\OperatorTok{(}\NormalTok{potential}\OperatorTok{.}\FunctionTok{getVariable}\OperatorTok{(}\DecValTok{0}\OperatorTok{)).}\FunctionTok{setPotential}\OperatorTok{(}\NormalTok{newPotential}\OperatorTok{);}
                \OperatorTok{\}}
\end{Highlighting}
\end{Shaded}
\paragraph{Cómo arreglarlo}
Al terminar \texttt{parametric\allowbreak{}Learning}, poner en el nodo correspondiente de \texttt{prob\allowbreak{}Net} cada tabla creada para sustituir a un \texttt{Uniform\allowbreak{}Potential}. Para ello, guardar en \texttt{adapt\allowbreak{}Network} la correspondencia variable → tabla nueva, igual que ya se guarda \texttt{ici\allowbreak{}Subpotentials} para los modelos canónicos y se reinstala al final.

Si la instalación en la red devuelta debe poder deshacerse, hacerla con una edición.

Prueba de regresión: la misma red A → B con B uniforme; tras aprender, B debe tener un \texttt{Table\allowbreak{}Potential} con P(B \textbar{} A) próximo a las frecuencias de los casos.

\setcounter{subsection}{137}
\subsection[El esqueleto del algoritmo PC cambia al renombrar una variable]{El esqueleto que aprende el algoritmo PC depende de los nombres de las variables: cambiar el nombre de una columna de la base de casos añade o quita enlaces}\label{err:138}
\begin{ficha}
\fichakey{Estado}Presente
\fichakey{Módulo}learning.algorithm
\fichakey{Paquete}org.openmarkov.learning.algorithm.pc
\fichakey{Ficheros}\path{learning.algorithm/src/main/java/org/openmarkov/learning/algorithm/pc/SkeletonDiscovery.java:18-27} \newline \path{learning.algorithm/src/main/java/org/openmarkov/learning/algorithm/pc/SkeletonDiscovery.java:173-216} \newline \path{learning.algorithm/src/main/java/org/openmarkov/learning/algorithm/pc/util/NodePair.java:38-51} \newline \path{learning.algorithm/src/main/java/org/openmarkov/learning/algorithm/pc/PCAlgorithm.java:412-425}
\fichakey{Comprobado}Leyendo el código y ejecutándolo
\fichakey{Esfuerzo}Pequeño (horas)
\fichakey{Decisión de equipo}No
\fichakey{Relacionados}\hyperref[err:139]{error~139}, \hyperref[err:140]{error~140}, \hyperref[nota:9]{nota~9}
\end{ficha}
\paragraph{Qué pasa}
El usuario llega por el menú de aprendizaje → algoritmo PC (el que aprende la estructura de la red con pruebas de independencia condicional), en aprendizaje automático o interactivo, y también por la validación cruzada de \texttt{bn\allowbreak{}Evaluation} (\texttt{Cross\allowbreak{}Validation\allowbreak{}Dialog}), que ofrece el mismo algoritmo. Con la misma base de casos, basta cambiar el nombre de una columna para que la red aprendida tenga enlaces de más o de menos. Nada lo advierte.

La primera fase del PC construye el esqueleto (los enlaces sin dirección). Parte del grafo completo y quita el enlace X--Y cuando encuentra un conjunto de variables que separa X de Y (una prueba de independencia condicional con p mayor que el nivel de significación). En el algoritmo publicado, en la profundidad d (el tamaño de los conjuntos que se prueban) se prueban los subconjuntos de d vecinos de X \textbf{y} los de d vecinos de Y. El javadoc de \texttt{Skeleton\allowbreak{}Discovery} dice que sigue la variante PC-Stable y que el esqueleto no depende del orden.

\texttt{Skeleton\allowbreak{}Discovery.\allowbreak{}separation\allowbreak{}Sets\allowbreak{}Logic} recorre los nodos por orden alfabético (\texttt{PCAlgorithm.\allowbreak{}sorted\allowbreak{}Nodes}) y, para cada nodo X, sus vecinos Y, así que visita cada pareja dos veces: una desde cada extremo. La primera visita, desde el nodo cuyo nombre va antes, prueba los subconjuntos de vecinos de ese nodo y guarda el mejor resultado en la caché con la llave \texttt{Node\allowbreak{}Pair}, que ordena los dos nodos por nombre y es la misma para (X, Y) y (Y, X).

En la segunda visita, desde el otro extremo, la condición \texttt{motivation.\allowbreak{}get\allowbreak{}Separation\allowbreak{}Set(\allowbreak{}).\allowbreak{}size(\allowbreak{})\ \textless{}\ adjacency\allowbreak{}Size} ya es falsa, porque lo guardado tiene tamaño d, y no se prueba nada. Los vecinos del segundo nodo solo se prueban en una profundidad en la que el primero no tiene bastantes vecinos. Si el conjunto que separa la pareja solo está entre los vecinos del segundo, el enlace falso se queda.

Medido llamando a las mismas clases que la ventana de aprendizaje con sus opciones por defecto (prueba de entropía cruzada, nivel de significación 0,05, alfa 0,5, prueba de dirección causal activada, sin red modelo, valores ausentes conservados y sin discretizar), y comparando los esqueletos por los nombres originales:

\begin{itemize}
\tightlist
\item
  Con una red sintética F→A, A→C, E→C, C→B, E→B y 20 000 casos, A y B quedan separados por \{C, E\} (p = 0,187), que son vecinos de B. Con el nodo llamado «A» ese conjunto no se prueba nunca (desde A solo se prueba \{F, C\}, con p = 2,6·10⁻³⁸) y sobra el enlace A--B. Con el mismo nodo llamado «G», el esqueleto sale bien.
\item
  \texttt{integration\allowbreak{}Tests/\allowbreak{}src/\allowbreak{}main/\allowbreak{}resources/\allowbreak{}networks/\allowbreak{}learning/\allowbreak{}alarm10k.\allowbreak{}csv} (37 variables, 10 000 casos): con los nombres del fichero el esqueleto tiene 44 enlaces; invirtiendo el orden alfabético de todos los nombres salen 48, con cinco de más (BP--STROKEVOLUME, DISCONNECT--KINKEDTUBE, HISTORY--PAP, MINVOL--TPR, MINVOL--VENTTUBE) y uno de menos (TPR--VENTMACH). De 74 cambios de nombre de una sola variable (anteponer «zz» o «AA»), 11 cambian el esqueleto.
\item
  En esa misma base, con los nombres del fichero se queda TPR→VENTMACH: el mejor conjunto entre los vecinos de TPR, \{ANAPHYLAXIS, BP, CATECHOL\}, da p = 3,5·10⁻⁶, pero VENTTUBE, vecino de VENTMACH, los separa con p = 0,146 y no se prueba porque «TPR» va antes que «VENTMACH». Renombrando TPR como «zzTPR», el enlace desaparece.
\item
  \texttt{integration\allowbreak{}Tests/\allowbreak{}src/\allowbreak{}main/\allowbreak{}resources/\allowbreak{}networks/\allowbreak{}learning/\allowbreak{}alarm500.\allowbreak{}csv} (500 casos): 35 enlaces con los nombres del fichero; invirtiendo el orden, 36 (sobra VENTLUNG--VENTTUBE); 6 de los 74 cambios de nombre de una variable cambian el esqueleto. Con VENTTUBE renombrada «AAVENTTUBE» en la cabecera del fichero, VENTLUNG--VENTTUBE se queda: desde los vecinos de AAVENTTUBE el mejor p es 2,6·10⁻¹¹, y \{EXPCO2, VENTALV\}, vecinos de VENTLUNG, da p = 0,30 pero no se prueba.
\item
  En las otras 13 bases probadas (\texttt{learning.\allowbreak{}algorithm/\allowbreak{}src/\allowbreak{}test/\allowbreak{}resources/\allowbreak{}network/\allowbreak{}BN-asia10k.\allowbreak{}csv}, \texttt{Head\allowbreak{}To\allowbreak{}Head1.\allowbreak{}csv}, \texttt{Head\allowbreak{}To\allowbreak{}Head2.\allowbreak{}csv}, \texttt{Head\allowbreak{}To\allowbreak{}Head2-10k.\allowbreak{}csv}, \texttt{Three\allowbreak{}Nodes\allowbreak{}Ato\allowbreak{}Bto\allowbreak{}C.\allowbreak{}csv}, \texttt{Two\allowbreak{}Nodes\allowbreak{}Ato\allowbreak{}B-1000.\allowbreak{}csv}; \texttt{integration\allowbreak{}Tests/\allowbreak{}src/\allowbreak{}main/\allowbreak{}resources/\allowbreak{}networks/\allowbreak{}learning/\allowbreak{}asia10K.\allowbreak{}csv}, \texttt{BN-A-B-C-E.\allowbreak{}csv}, \texttt{learn\allowbreak{}Test\allowbreak{}Data\allowbreak{}Base.\allowbreak{}csv}, \texttt{line\allowbreak{}And\allowbreak{}Star.\allowbreak{}xlsx}; \texttt{core/\allowbreak{}src/\allowbreak{}test/\allowbreak{}resources/\allowbreak{}nets/\allowbreak{}Cataract-RRDdataset\allowbreak{}Pre\allowbreak{}Proc3.\allowbreak{}csv}, \texttt{Cataract-RRDdataset\allowbreak{}Pre\allowbreak{}Proc4.\allowbreak{}csv}; \texttt{0\_\allowbreak{}miscellaneous/\allowbreak{}networks/\allowbreak{}bn/\allowbreak{}BN-asia-5k.\allowbreak{}xls}) ningún cambio de nombre cambió el esqueleto. No se probaron las bases de \texttt{Cataract} con columnas numéricas sin discretizar ni la de Elvira (\texttt{.\allowbreak{}dbc}), que la ventana no abre.
\end{itemize}

El resultado es repetible: dos ejecuciones con los mismos nombres dan la misma red. Cambiar el nombre en la cabecera del fichero da lo mismo que cambiarlo después de leer la base.
\paragraph{El código}
learning.algorithm/src/main/java/org/openmarkov/learning/algorithm/pc/SkeletonDiscovery.java:173-192

\begin{Shaded}
\begin{Highlighting}[]
\KeywordTok{private} \DataTypeTok{void} \FunctionTok{separationSetsLogic}\OperatorTok{(}\DataTypeTok{int}\NormalTok{ adjacencySize}\OperatorTok{)} \OperatorTok{\{}
    \ControlFlowTok{for} \OperatorTok{(}\BuiltInTok{Node}\NormalTok{ nodeX }\OperatorTok{:}\NormalTok{ pc}\OperatorTok{.}\FunctionTok{sortedNodes}\OperatorTok{())} \OperatorTok{\{}
        \ControlFlowTok{for} \OperatorTok{(}\BuiltInTok{Node}\NormalTok{ nodeY }\OperatorTok{:}\NormalTok{ PCAlgorithm}\OperatorTok{.}\FunctionTok{sorted}\OperatorTok{(}\NormalTok{nodeX}\OperatorTok{.}\FunctionTok{getSiblings}\OperatorTok{()))} \OperatorTok{\{}
            \BuiltInTok{List}\OperatorTok{\textless{}}\BuiltInTok{Node}\OperatorTok{\textgreater{}}\NormalTok{ snapshotNeighbors }\OperatorTok{=}\NormalTok{ stableAdjSnapshot}\OperatorTok{.}\FunctionTok{getOrDefault}\OperatorTok{(}\NormalTok{nodeX}\OperatorTok{,} \BuiltInTok{Collections}\OperatorTok{.}\FunctionTok{emptyList}\OperatorTok{());}
            \BuiltInTok{List}\OperatorTok{\textless{}}\BuiltInTok{Node}\OperatorTok{\textgreater{}}\NormalTok{ adjacencySubset }\OperatorTok{=} \KeywordTok{new} \BuiltInTok{ArrayList}\OperatorTok{\textless{}\textgreater{}(}\NormalTok{snapshotNeighbors}\OperatorTok{);}
\NormalTok{            adjacencySubset}\OperatorTok{.}\FunctionTok{remove}\OperatorTok{(}\NormalTok{nodeY}\OperatorTok{);}
            \KeywordTok{...}
            \ControlFlowTok{if} \OperatorTok{(!}\NormalTok{pc}\OperatorTok{.}\FunctionTok{alreadyConsidered}\OperatorTok{(}\NormalTok{removeLinkEdit}\OperatorTok{,}\NormalTok{ lastRemovedEdits}\OperatorTok{))} \OperatorTok{\{}
\NormalTok{                PCEditMotivation motivation }\OperatorTok{=}\NormalTok{ pc}\OperatorTok{.}\FunctionTok{cache}\OperatorTok{.}\FunctionTok{get}\OperatorTok{(}\KeywordTok{new} \FunctionTok{NodePair}\OperatorTok{(}\NormalTok{nodeX}\OperatorTok{,}\NormalTok{ nodeY}\OperatorTok{));}
                \ControlFlowTok{if} \OperatorTok{(}\NormalTok{motivation }\OperatorTok{==} \KeywordTok{null} \OperatorTok{||} \OperatorTok{(}\NormalTok{motivation}\OperatorTok{.}\FunctionTok{getScore}\OperatorTok{()} \OperatorTok{!=}\NormalTok{ ALREADY\_DONE}
                        \OperatorTok{\&\&}\NormalTok{ motivation}\OperatorTok{.}\FunctionTok{getSeparationSet}\OperatorTok{().}\FunctionTok{size}\OperatorTok{()} \OperatorTok{\textless{}}\NormalTok{ adjacencySize}\OperatorTok{))} \OperatorTok{\{}
                    \FunctionTok{evaluateSeparationSets}\OperatorTok{(}\NormalTok{nodeX}\OperatorTok{,}\NormalTok{ nodeY}\OperatorTok{,}\NormalTok{ adjacencySubset}\OperatorTok{,}\NormalTok{ adjacencySize}\OperatorTok{);}
                \OperatorTok{\}}
            \OperatorTok{\}}
        \OperatorTok{\}}
    \OperatorTok{\}}
\OperatorTok{\}}
\end{Highlighting}
\end{Shaded}

learning.algorithm/src/main/java/org/openmarkov/learning/algorithm/pc/SkeletonDiscovery.java:199-216

\begin{Shaded}
\begin{Highlighting}[]
\KeywordTok{private} \DataTypeTok{void} \FunctionTok{evaluateSeparationSets}\OperatorTok{(}\BuiltInTok{Node}\NormalTok{ nodeX}\OperatorTok{,} \BuiltInTok{Node}\NormalTok{ nodeY}\OperatorTok{,}
                                    \BuiltInTok{List}\OperatorTok{\textless{}}\BuiltInTok{Node}\OperatorTok{\textgreater{}}\NormalTok{ adjacencySubset}\OperatorTok{,} \DataTypeTok{int}\NormalTok{ adjacencySize}\OperatorTok{)} \OperatorTok{\{}
    \DataTypeTok{double}\NormalTok{ bestScore }\OperatorTok{=} \FloatTok{0.0}\OperatorTok{;}
    \BuiltInTok{List}\OperatorTok{\textless{}}\BuiltInTok{Node}\OperatorTok{\textgreater{}}\NormalTok{ bestScoreSeparationSet }\OperatorTok{=} \KeywordTok{null}\OperatorTok{;}
    \ControlFlowTok{for} \OperatorTok{(}\BuiltInTok{List}\OperatorTok{\textless{}}\BuiltInTok{Node}\OperatorTok{\textgreater{}}\NormalTok{ separationSet }\OperatorTok{:}\NormalTok{ PCAlgorithm}\OperatorTok{.}\FunctionTok{subSetsOfSize}\OperatorTok{(}\NormalTok{adjacencySubset}\OperatorTok{,}\NormalTok{ adjacencySize}\OperatorTok{))} \OperatorTok{\{}
        \DataTypeTok{double}\NormalTok{ linkScore }\OperatorTok{=}\NormalTok{ pc}\OperatorTok{.}\FunctionTok{independenceTester}\OperatorTok{.}\FunctionTok{test}\OperatorTok{(}\NormalTok{pc}\OperatorTok{.}\FunctionTok{database}\OperatorTok{(),}\NormalTok{ nodeX}\OperatorTok{,}\NormalTok{ nodeY}\OperatorTok{,}\NormalTok{ separationSet}\OperatorTok{);}
        \KeywordTok{...}
    \OperatorTok{\}}
    \ControlFlowTok{if} \OperatorTok{(}\NormalTok{bestScoreSeparationSet }\OperatorTok{!=} \KeywordTok{null}\OperatorTok{)} \OperatorTok{\{}
\NormalTok{        pc}\OperatorTok{.}\FunctionTok{cache}\OperatorTok{.}\FunctionTok{put}\OperatorTok{(}\KeywordTok{new} \FunctionTok{NodePair}\OperatorTok{(}\NormalTok{nodeX}\OperatorTok{,}\NormalTok{ nodeY}\OperatorTok{),}
                \KeywordTok{new} \FunctionTok{PCEditMotivation}\OperatorTok{(}\NormalTok{bestScore}\OperatorTok{,}\NormalTok{ bestScoreSeparationSet}\OperatorTok{));}
    \OperatorTok{\}}
\OperatorTok{\}}
\end{Highlighting}
\end{Shaded}
\paragraph{Cómo arreglarlo}
En cada profundidad d, evaluar cada pareja una sola vez con los subconjuntos de tamaño d de los vecinos de X (sin Y) \textbf{y} de los vecinos de Y (sin X), tomados de la misma instantánea de vecinos, y guardar en la caché el mejor de los dos lados. Por ejemplo, recorrer solo las parejas con X antes que Y y pasar a \texttt{evaluate\allowbreak{}Separation\allowbreak{}Sets} las dos listas.

La caché tiene que saber que en la profundidad d se probaron los dos lados, para que la condición de \texttt{separation\allowbreak{}Sets\allowbreak{}Logic} no deje fuera a ninguno; hoy el tamaño del conjunto guardado hace ese papel y no basta.

El conjunto separador guardado es el que después leen \texttt{Collider\allowbreak{}Orientation.\allowbreak{}find\allowbreak{}Edit} y \texttt{creates\allowbreak{}Contradictory\allowbreak{}Collider} para orientar las figuras en cabeza, \texttt{on\allowbreak{}Edit\allowbreak{}Undone} para repetir la prueba al deshacer e \texttt{invalidate\allowbreak{}Neighbor\allowbreak{}Cache} para descartar pruebas tras un borrado; todos deben seguir recibiendo el conjunto que separó la pareja, sea del lado que sea.

Probar los dos lados puede duplicar el número de pruebas de independencia en redes con muchos vecinos.

Pruebas de regresión: con la red sintética de cinco variables descrita arriba, el esqueleto debe ser el mismo con el nodo llamado «A» o «G»; con \texttt{alarm10k.\allowbreak{}csv}, renombrar TPR como «zzTPR» no debe cambiar el esqueleto; y, en general, aplicar a los nombres una permutación (por ejemplo, invertir su orden alfabético) no debe cambiar el esqueleto.

\setcounter{subsection}{138}
\subsection[El PC interactivo ofrece las dos orientaciones de un enlace, contra Meek]{El aprendizaje interactivo con PC ofrece las dos orientaciones de un mismo enlace, incluida la contraria a la que obliga una regla de Meek}\label{err:139}
\begin{ficha}
\fichakey{Estado}Presente
\fichakey{Módulo}learning.algorithm
\fichakey{Paquete}org.openmarkov.learning.algorithm.pc
\fichakey{Ficheros}\path{learning.algorithm/src/main/java/org/openmarkov/learning/algorithm/pc/PCAlgorithm.java:544-548} \newline \path{learning.algorithm/src/main/java/org/openmarkov/learning/algorithm/pc/MeekOrientation.java:177-245} \newline \path{core/src/main/java/org/openmarkov/core/action/base/linkEdits/BaseLinkEdit.java:71-85} \newline \path{learning.gui/src/main/java/org/openmarkov/learning/gui/interactive/InteractiveLearningDialog.java:403-410}
\fichakey{Comprobado}Leyendo el código y ejecutándolo
\fichakey{Esfuerzo}Pequeño (horas)
\fichakey{Decisión de equipo}\textcolor{decisionrojo}{\textbf{Sí.}} Si la lista interactiva debe ofrecer las dos direcciones de un enlace que ninguna regla orienta (hoy lo hace) o solo una
\fichakey{Relacionados}\hyperref[err:138]{error~138}, \hyperref[err:140]{error~140}, \hyperref[err:141]{error~141}, \hyperref[err:142]{error~142}
\end{ficha}
\paragraph{Qué pasa}
El usuario llega por el menú de aprendizaje → algoritmo PC (el que aprende la estructura de la red con pruebas de independencia condicional) → aprendizaje interactivo → pasar a la fase de orientación. La tabla de propuestas ofrece las dos direcciones de un mismo enlace.

Entre ellas aparece a veces, junto a una orientación obligada por una regla de Meek (las reglas que orientan los enlaces que quedan sin dirección), su contraria con el motivo inocuo «Do not create cycles». Si el usuario elige esa fila, la red aprendida contradice la regla sin aviso.

\texttt{PCAlgorithm.\allowbreak{}already\allowbreak{}Considered(\allowbreak{}edit,\allowbreak{}\ considered\allowbreak{}Edits)} sirve para no proponer dos veces el mismo cambio sobre un enlace: pregunta si el cambio ya está en el conjunto de los considerados y si está «el mismo enlace al revés». Pero el cambio al revés lo construye siempre como \texttt{Remove\allowbreak{}Link\allowbreak{}Edit}, sea cual sea el tipo del cambio recibido, y \texttt{Base\allowbreak{}Link\allowbreak{}Edit.\allowbreak{}equals} exige que las dos clases sean iguales.

En la fase del esqueleto (\texttt{Skeleton\allowbreak{}Discovery}), donde todo son borrados, funciona. En la fase de Meek (\texttt{Meek\allowbreak{}Orientation}), los cambios son \texttt{Orient\allowbreak{}Link\allowbreak{}Edit} y el conjunto está lleno de orientaciones, así que la segunda mitad de la pregunta no puede dar «sí» nunca.

En aprendizaje automático no se nota, porque \texttt{run} pide cada vez solo la mejor propuesta y el historial se vacía. Sí se nota en el aprendizaje interactivo: \texttt{Interactive\allowbreak{}Learning\allowbreak{}Dialog.\allowbreak{}update\allowbreak{}Editions\allowbreak{}Table} llena la tabla llamando a \texttt{get\allowbreak{}Best\allowbreak{}Edit} y luego a \texttt{get\allowbreak{}Next\allowbreak{}Edit} hasta 50 veces, y cada llamada añade la propuesta al historial.

Para un enlace X--Z sin dirección, la regla de reserva (\texttt{try\allowbreak{}Orient\allowbreak{}Unoriented\allowbreak{}Without\allowbreak{}Path}) propone primero X→Z; en la llamada siguiente X→Z ya está considerada, pero Z→X no se reconoce como «el mismo enlace al revés» y se propone también. Y lo mismo ocurre después de una orientación obligada por una regla de Meek: se ofrece la contraria como si fuera neutra.

La regla R1 orienta B→C cuando hay A→B y A no es vecino de C, porque la dirección contraria crearía una figura en cabeza que los datos no apoyan.

Medido con un pequeño programa de prueba que reproduce el llenado de la tabla (\texttt{get\allowbreak{}Best\allowbreak{}Edit(\allowbreak{}true,\allowbreak{}\ true)} y \texttt{get\allowbreak{}Next\allowbreak{}Edit(\allowbreak{}true,\allowbreak{}\ true)} hasta 50) sobre \texttt{learning.\allowbreak{}algorithm/\allowbreak{}src/\allowbreak{}test/\allowbreak{}resources/\allowbreak{}network/\allowbreak{}BN-asia10k.\allowbreak{}csv} con el PC, prueba de entropía cruzada y nivel de significación 0,05, ejecutando cada vez la primera propuesta. Casi todas las listas de la fase de orientación contienen las dos direcciones de cada enlace; en dos pasos aparecen juntas una orientación obligada y su contraria:

\begin{Verbatim}[breaklines,breakanywhere,fontsize=\small]
Orient link: TuberculosisOrCancer --> X-ray  {Meek R1}
...
Orient link: X-ray --> TuberculosisOrCancer  {Do not create cycles}

Orient link: Smoker --> LungCancer  {Meek R1}
Orient link: LungCancer --> Smoker  {Do not create cycles}
\end{Verbatim}
\paragraph{El código}
learning.algorithm/src/main/java/org/openmarkov/learning/algorithm/pc/PCAlgorithm.java:544-548

\begin{Shaded}
\begin{Highlighting}[]
\KeywordTok{public} \DataTypeTok{boolean} \FunctionTok{alreadyConsidered}\OperatorTok{(}\NormalTok{BaseLinkEdit edit}\OperatorTok{,} \BuiltInTok{Set}\OperatorTok{\textless{}}\NormalTok{PNEdit}\OperatorTok{\textgreater{}}\NormalTok{ consideredEdits}\OperatorTok{)} \OperatorTok{\{}
\NormalTok{    BaseLinkEdit inverseEdit }\OperatorTok{=} \KeywordTok{new} \FunctionTok{RemoveLinkEdit}\OperatorTok{(}\NormalTok{probNet}\OperatorTok{,}\NormalTok{ edit}\OperatorTok{.}\FunctionTok{getVariableTo}\OperatorTok{(),}\NormalTok{ edit}\OperatorTok{.}\FunctionTok{getVariableFrom}\OperatorTok{(),}
\NormalTok{                                                  edit}\OperatorTok{.}\FunctionTok{isDirected}\OperatorTok{());}
    \ControlFlowTok{return}\NormalTok{ consideredEdits}\OperatorTok{.}\FunctionTok{contains}\OperatorTok{(}\NormalTok{edit}\OperatorTok{)} \OperatorTok{||}\NormalTok{ consideredEdits}\OperatorTok{.}\FunctionTok{contains}\OperatorTok{(}\NormalTok{inverseEdit}\OperatorTok{);}
\OperatorTok{\}}
\end{Highlighting}
\end{Shaded}

core/src/main/java/org/openmarkov/core/action/base/linkEdits/BaseLinkEdit.java:75-79

\begin{Shaded}
\begin{Highlighting}[]
\AttributeTok{@Override} \KeywordTok{public} \DataTypeTok{boolean} \FunctionTok{equals}\OperatorTok{(}\BuiltInTok{Object}\NormalTok{ obj}\OperatorTok{)} \OperatorTok{\{}
    \ControlFlowTok{if} \OperatorTok{(}\KeywordTok{this} \OperatorTok{==}\NormalTok{ obj}\OperatorTok{)}
        \ControlFlowTok{return} \KeywordTok{true}\OperatorTok{;}
    \ControlFlowTok{if} \OperatorTok{((}\NormalTok{obj }\OperatorTok{==} \KeywordTok{null}\OperatorTok{)} \OperatorTok{||} \OperatorTok{(}\NormalTok{obj}\OperatorTok{.}\FunctionTok{getClass}\OperatorTok{()} \OperatorTok{!=} \KeywordTok{this}\OperatorTok{.}\FunctionTok{getClass}\OperatorTok{()))}
        \ControlFlowTok{return} \KeywordTok{false}\OperatorTok{;}
    \KeywordTok{...}
\end{Highlighting}
\end{Shaded}
\paragraph{Cómo arreglarlo}
Si el equipo decide que no deben ofrecerse las dos direcciones del mismo enlace, construir el cambio al revés con la misma clase que el recibido (un \texttt{Orient\allowbreak{}Link\allowbreak{}Edit} para una orientación, un \texttt{Remove\allowbreak{}Link\allowbreak{}Edit} para un borrado); por ejemplo, con un método por tipo en vez de uno genérico.

Los sitios que usan la comprobación son \texttt{Skeleton\allowbreak{}Discovery.\allowbreak{}separation\allowbreak{}Sets\allowbreak{}Logic} e \texttt{is\allowbreak{}Valid\allowbreak{}Edit} (borrados) y \texttt{Meek\allowbreak{}Orientation.\allowbreak{}build\allowbreak{}Orient\allowbreak{}Proposal} y \texttt{try\allowbreak{}Orient\allowbreak{}Direction} (orientaciones).

Hay que revisar a la vez la reserva de \texttt{try\allowbreak{}Orient\allowbreak{}Unoriented\allowbreak{}Without\allowbreak{}Path}, que intenta la dirección contraria cuando la preferida ya se propuso: con la comprobación arreglada, esa segunda llamada solo serviría cuando la preferida crea un ciclo o está bloqueada, que es lo que su motivo dice.

Si el equipo decide que las dos direcciones deben ofrecerse cuando ninguna regla obliga, al menos no debe ofrecerse la contraria de una orientación de Meek.

Prueba de regresión: llenar la lista como hace el diálogo sobre una red con A→B--C (A no vecino de C) y comprobar que no aparece C→B.

\setcounter{subsection}{139}
\subsection[El PC interactivo adelanta la fase al llenar la tabla y descoloca botones]{En el aprendizaje interactivo con PC, llenar la tabla adelanta la fase del algoritmo: la tabla mezcla orientaciones con enlaces aún por quitar, «Complete phase» orienta sin quitar ningún enlace y marcar una casilla deja la tabla sin borrados o vacía}\label{err:140}
\begin{ficha}
\fichakey{Estado}Presente
\fichakey{Módulo}learning.algorithm
\fichakey{Paquete}org.openmarkov.learning.algorithm.pc
\fichakey{Ficheros}\path{learning.algorithm/src/main/java/org/openmarkov/learning/algorithm/pc/PCAlgorithm.java:151-239} \newline \path{learning.algorithm/src/main/java/org/openmarkov/learning/algorithm/pc/PCAlgorithm.java:288-387} \newline \path{learning.gui/src/main/java/org/openmarkov/learning/gui/interactive/InteractiveLearningDialog.java:92} \newline \path{learning.gui/src/main/java/org/openmarkov/learning/gui/interactive/InteractiveLearningDialog.java:298-312} \newline \path{learning.gui/src/main/java/org/openmarkov/learning/gui/interactive/InteractiveLearningDialog.java:354-361} \newline \path{learning.gui/src/main/java/org/openmarkov/learning/gui/interactive/InteractiveLearningDialog.java:403-427} \newline \path{learning.core/src/main/java/org/openmarkov/learning/core/algorithm/LearningAlgorithm.java:108-122}
\fichakey{Comprobado}Leyendo el código y ejecutándolo
\fichakey{Esfuerzo}Medio (días)
\fichakey{Decisión de equipo}\textcolor{decisionrojo}{\textbf{Sí.}} Si la tabla del aprendizaje interactivo debe enseñar solo las propuestas de la fase en curso o también las de las fases siguientes
\fichakey{Relacionados}\hyperref[err:134]{error~134}, \hyperref[err:138]{error~138}, \hyperref[err:139]{error~139}, \hyperref[err:141]{error~141}, \hyperref[err:142]{error~142}
\end{ficha}
\paragraph{Qué pasa}
El usuario llega por el menú de aprendizaje → algoritmo PC (el que aprende la estructura de la red con pruebas de independencia condicional) → «Interactive learning» → «Learn». Se abre la ventana «Interactive learning» con la tabla de propuestas. En cuanto quedan menos de 50 enlaces por quitar (desde el primer momento en una base de pocas variables, y hacia el final del esqueleto en las demás):

\begin{itemize}
\tightlist
\item
  Debajo de los enlaces por quitar, la tabla ofrece orientaciones de enlaces que todavía no se sabe si se quedan, incluidos los que la misma tabla propone quitar más arriba.
\item
  Si el usuario aplica una de esas orientaciones, ya no se le ofrece ningún borrado más, aunque quedaban.
\item
  Si desmarca o vuelve a marcar «Only positive edits» u «Only allowed edits», o pulsa «Block edit» o «Show blocked», la tabla pierde todos los borrados; a veces se queda vacía, y entonces la ventana da el aprendizaje por terminado y desactiva «Apply edit» y «Complete phase».
\item
  Si pulsa «Complete phase», el algoritmo orienta todos los enlaces que haya sin quitar ninguno, o no hace nada y da el aprendizaje por terminado.
\end{itemize}

El PC trabaja en fases: primero quita enlaces (esqueleto), después orienta figuras en cabeza (colisionadores) y por último orienta el resto con las reglas de Meek. \texttt{PCAlgorithm} guarda la fase en curso en el campo \texttt{phase}.

\texttt{Interactive\allowbreak{}Learning\allowbreak{}Dialog.\allowbreak{}update\allowbreak{}Editions\allowbreak{}Table}, al abrirse la ventana y cada vez que se actualiza, llama a \texttt{get\allowbreak{}Best\allowbreak{}Edit} y después a \texttt{get\allowbreak{}Next\allowbreak{}Edit} hasta tener 50 propuestas. Cada llamada va a \texttt{get\allowbreak{}Optimal\allowbreak{}Edit}, y cuando la fase en curso no tiene más propuestas, \texttt{transition\allowbreak{}To\allowbreak{}Next\allowbreak{}Phase} \textbf{cambia \texttt{phase}} y sigue buscando en la siguiente, sin haber aplicado nada a la red. Al terminar de llenar la tabla, el algoritmo se ha quedado en la fase de orientación o en \texttt{ORIENTATION\_\allowbreak{}FINISHED}.

La fase solo vuelve al principio al aplicar o deshacer un borrado, al añadir un enlace o al deshacer la orientación de un colisionador (\texttt{after\allowbreak{}Edit\allowbreak{}Executes} y \texttt{after\allowbreak{}Undoing\allowbreak{}Edit}); \texttt{get\allowbreak{}Best\allowbreak{}Edit} solo vacía el historial de propuestas, no la fase. Por eso todo lo que vuelve a llenar la tabla sin pasar por esas ediciones (las dos casillas, «Block edit», «Show blocked») empieza desde la fase en que se quedó el llenado anterior. «Complete phase» llama a \texttt{Learning\allowbreak{}Algorithm.\allowbreak{}run\allowbreak{}Till\allowbreak{}Next\allowbreak{}Phase}, que toma como fase de partida la adelantada y aplica propuestas mientras no cambie.

Aplicar una orientación de las reglas de Meek pone la fase en \texttt{REMAINING\_\allowbreak{}LINKS\_\allowbreak{}ORIENTATION}, y deshacerla la deja en \texttt{HEAD\_\allowbreak{}TO\_\allowbreak{}HEAD\_\allowbreak{}ORIENTATION}, así que los borrados pendientes no vuelven a ofrecerse.

«Finish» no arrastra la fase adelantada, porque \texttt{learn} pasa por \texttt{init}, que la devuelve al principio (y que además descarta lo que el usuario había aplicado: ver \hyperref[err:141]{error~141}); pero con la tabla vacía la ventana ya no llama a \texttt{learn}.

Medido con un pequeño programa de prueba que construye el algoritmo como la ventana de aprendizaje (prueba de entropía cruzada, nivel de significación 0,05, alfa 0,5, prueba de dirección causal activada) y repite las llamadas de la ventana:

\begin{itemize}
\tightlist
\item
  \texttt{learning.\allowbreak{}algorithm/\allowbreak{}src/\allowbreak{}test/\allowbreak{}resources/\allowbreak{}network/\allowbreak{}BN-asia10k.\allowbreak{}csv} (8 variables): la tabla inicial tiene 50 filas: 22 borrados y, desde la fila 23, 28 orientaciones de enlaces del grafo completo; al terminar de llenarla la fase es la de las reglas de Meek. Desmarcar «Only positive edits» deja 50 orientaciones y ningún borrado, y volver a marcarla, lo mismo. Aplicar la fila 23 («Orient link: Dyspnea --\textgreater{} Bronchitis», motivo «Causal direction test (ANM)») deja la tabla siguiente con 50 orientaciones y ningún borrado. «Complete phase» nada más abrir deja la red con sus 28 enlaces, todos orientados, y la tabla vacía.
\item
  \texttt{learning.\allowbreak{}algorithm/\allowbreak{}src/\allowbreak{}test/\allowbreak{}resources/\allowbreak{}network/\allowbreak{}Head\allowbreak{}To\allowbreak{}Head2.\allowbreak{}csv} (5 variables): la tabla inicial tiene 25 filas (5 borrados y 20 orientaciones) y el llenado termina en \texttt{ORIENTATION\_\allowbreak{}FINISHED}. Desmarcar «Only positive edits» deja la tabla vacía. «Complete phase» nada más abrir no hace nada y la tabla queda vacía, con la red en el grafo completo sin dirección (10 enlaces). Con la tabla vacía, la ventana marca el aprendizaje como terminado y «Finish» ya no aprende.
\item
  \texttt{integration\allowbreak{}Tests/\allowbreak{}src/\allowbreak{}main/\allowbreak{}resources/\allowbreak{}networks/\allowbreak{}learning/\allowbreak{}alarm500.\allowbreak{}csv} (37 variables): la tabla inicial son 50 borrados y la fase no se mueve; ahí el adelanto aparece cuando quedan menos de 50 enlaces por quitar.
\end{itemize}

Solo se midió la casilla «Only positive edits»; «Only allowed edits», «Block edit» y «Show blocked» llaman al mismo \texttt{update\allowbreak{}UI}. La ventana es un \texttt{JDialog} y no se ha abierto; lo que hace con la tabla vacía (botones desactivados, «Finish» sin efecto) está leído en \texttt{update\allowbreak{}Editions\allowbreak{}Table} y \texttt{btn\allowbreak{}Run\allowbreak{}Action\allowbreak{}Performed}.
\paragraph{El código}
learning.algorithm/src/main/java/org/openmarkov/learning/algorithm/pc/PCAlgorithm.java:193-200

\begin{Shaded}
\begin{Highlighting}[]
\KeywordTok{public}\NormalTok{ LearningEditProposal }\FunctionTok{getOptimalEdit}\OperatorTok{(}\DataTypeTok{boolean}\NormalTok{ onlyAllowedEdits}\OperatorTok{,} \DataTypeTok{boolean}\NormalTok{ onlyPositiveEdits}\OperatorTok{)} \OperatorTok{\{}
    \ControlFlowTok{if} \OperatorTok{(}\NormalTok{phase }\OperatorTok{==}\NormalTok{ Phase}\OperatorTok{.}\FunctionTok{INITIAL\_PHASE}\OperatorTok{)} \OperatorTok{\{}
\NormalTok{        LearningEditProposal proposal }\OperatorTok{=}\NormalTok{ skeleton}\OperatorTok{.}\FunctionTok{findEdit}\OperatorTok{(}\NormalTok{onlyAllowedEdits}\OperatorTok{,}\NormalTok{ onlyPositiveEdits}\OperatorTok{);}
        \ControlFlowTok{if} \OperatorTok{(}\NormalTok{proposal }\OperatorTok{!=} \KeywordTok{null}\OperatorTok{)} \OperatorTok{\{}
            \ControlFlowTok{return}\NormalTok{ proposal}\OperatorTok{;}
        \OperatorTok{\}}
        \ControlFlowTok{return} \FunctionTok{transitionToNextPhase}\OperatorTok{(}\NormalTok{onlyAllowedEdits}\OperatorTok{);}
    \OperatorTok{\}}
\end{Highlighting}
\end{Shaded}

learning.algorithm/src/main/java/org/openmarkov/learning/algorithm/pc/PCAlgorithm.java:367-372

\begin{Shaded}
\begin{Highlighting}[]
\KeywordTok{private}\NormalTok{ LearningEditProposal }\FunctionTok{transitionToNextPhase}\OperatorTok{(}\DataTypeTok{boolean}\NormalTok{ onlyAllowedEdits}\OperatorTok{)} \OperatorTok{\{}
    \ControlFlowTok{return} \ControlFlowTok{switch} \OperatorTok{(}\NormalTok{phase}\OperatorTok{)} \OperatorTok{\{}
        \ControlFlowTok{case}\NormalTok{ INITIAL\_PHASE }\OperatorTok{{-}\textgreater{}} \OperatorTok{\{}
\NormalTok{            phase }\OperatorTok{=}\NormalTok{ Phase}\OperatorTok{.}\FunctionTok{HEAD\_TO\_HEAD\_ORIENTATION}\OperatorTok{;}
            \ControlFlowTok{yield} \FunctionTok{getOrientationEdit}\OperatorTok{(}\NormalTok{onlyAllowedEdits}\OperatorTok{);}
        \OperatorTok{\}}
\end{Highlighting}
\end{Shaded}

learning.gui/src/main/java/org/openmarkov/learning/gui/interactive/InteractiveLearningDialog.java:403-419

\begin{Shaded}
\begin{Highlighting}[]
\KeywordTok{private} \DataTypeTok{void} \FunctionTok{updateEditionsTable}\OperatorTok{(}\DataTypeTok{boolean}\NormalTok{ onlyAllowed}\OperatorTok{,} \DataTypeTok{boolean}\NormalTok{ onlyPositive}\OperatorTok{,}
                                 \BuiltInTok{List}\OperatorTok{\textless{}}\NormalTok{LearningEditProposal}\OperatorTok{\textgreater{}}\NormalTok{ blockedEdits}\OperatorTok{)} \OperatorTok{\{}
    \BuiltInTok{ArrayList}\OperatorTok{\textless{}}\NormalTok{LearningEditProposal}\OperatorTok{\textgreater{}}\NormalTok{ showableEditProposals }\OperatorTok{=} \KeywordTok{new} \BuiltInTok{ArrayList}\OperatorTok{\textless{}\textgreater{}();}
\NormalTok{    LearningEditProposal bestEdition }\OperatorTok{=}\NormalTok{ learningManager}\OperatorTok{.}\FunctionTok{getLearningAlgorithm}\OperatorTok{().}\FunctionTok{getBestEdit}\OperatorTok{(}\NormalTok{onlyAllowed}\OperatorTok{,}\NormalTok{ onlyPositive}\OperatorTok{);}
    \ControlFlowTok{while} \OperatorTok{(}\NormalTok{showableEditProposals}\OperatorTok{.}\FunctionTok{size}\OperatorTok{()} \OperatorTok{\textless{}}\NormalTok{ SHOWING\_EDIT\_NUM }\OperatorTok{\&\&}\NormalTok{ bestEdition }\OperatorTok{!=} \KeywordTok{null}\OperatorTok{)} \OperatorTok{\{}
\NormalTok{        showableEditProposals}\OperatorTok{.}\FunctionTok{add}\OperatorTok{(}\NormalTok{bestEdition}\OperatorTok{);}
\NormalTok{        bestEdition }\OperatorTok{=}\NormalTok{ learningManager}\OperatorTok{.}\FunctionTok{getLearningAlgorithm}\OperatorTok{().}\FunctionTok{getNextEdit}\OperatorTok{(}\NormalTok{onlyAllowed}\OperatorTok{,}\NormalTok{ onlyPositive}\OperatorTok{);}
    \OperatorTok{\}}
\NormalTok{    editionsTable}\OperatorTok{.}\FunctionTok{fill}\OperatorTok{(}\NormalTok{showableEditProposals}\OperatorTok{);}
    \ControlFlowTok{if} \OperatorTok{(}\NormalTok{showableEditProposals}\OperatorTok{.}\FunctionTok{isEmpty}\OperatorTok{())} \OperatorTok{\{}
        \KeywordTok{...}
\NormalTok{        learningFinished }\OperatorTok{=} \KeywordTok{true}\OperatorTok{;}
    \OperatorTok{\}}
\end{Highlighting}
\end{Shaded}
\paragraph{Cómo arreglarlo}
Hay que separar la fase en que está la red (la que marcan las ediciones aplicadas) de la fase hasta la que se ha asomado el llenado de la tabla. Dos caminos, según lo que decida el equipo sobre qué debe enseñar la tabla:

\begin{itemize}
\tightlist
\item
  Si la tabla solo debe enseñar la fase en curso: que \texttt{get\allowbreak{}Next\allowbreak{}Edit} no pase de fase (o que \texttt{update\allowbreak{}Editions\allowbreak{}Table} deje de pedir cuando \texttt{get\allowbreak{}Phase(\allowbreak{})} cambie), y que la fase avance solo cuando la fase en curso se quede sin propuestas al pedir la mejor.
\item
  Si debe enseñar también las fases siguientes: que el llenado no deje rastro, guardando \texttt{phase} (y el estado de \texttt{Skeleton\allowbreak{}Discovery}: profundidad e instantánea de vecinos) antes de pedir la primera propuesta y reponiéndolo al terminar; por ejemplo, con un método del algoritmo que devuelva de una vez la lista de propuestas. Las orientaciones que se ofrezcan con borrados pendientes deben quedar claramente separadas, porque se calculan sobre un esqueleto que aún no es el final.
\end{itemize}

Los sitios que asignan \texttt{phase} son \texttt{get\allowbreak{}Optimal\allowbreak{}Edit}, \texttt{get\allowbreak{}Orientation\allowbreak{}Edit}, \texttt{find\allowbreak{}Best\allowbreak{}Edit\allowbreak{}In\allowbreak{}Current\allowbreak{}Phase}, \texttt{transition\allowbreak{}To\allowbreak{}Next\allowbreak{}Phase}, \texttt{init}, \texttt{after\allowbreak{}Edit\allowbreak{}Executes} y \texttt{after\allowbreak{}Undoing\allowbreak{}Edit}. Los comentarios de \texttt{init} y de \texttt{after\allowbreak{}Edit\allowbreak{}Executes} ya describen el adelanto y lo compensan solo tras una edición; con el arreglo esas compensaciones deben revisarse. \texttt{Learning\allowbreak{}Algorithm.\allowbreak{}run\allowbreak{}Till\allowbreak{}Next\allowbreak{}Phase} («Complete phase») debe partir de la fase real.

Pruebas de regresión, repitiendo las llamadas de la ventana sobre \texttt{BN-asia10k.\allowbreak{}csv}: tras llenar la tabla, un segundo llenado con otras casillas sigue ofreciendo los 22 borrados; «Complete phase» nada más abrir quita enlaces y no orienta ninguno; y con \texttt{Head\allowbreak{}To\allowbreak{}Head2.\allowbreak{}csv} un segundo llenado no sale vacío.

\setcounter{subsection}{140}
\subsection[«Finish» del PC interactivo empieza de cero y descarta lo aplicado a mano]{En el aprendizaje interactivo con PC, «Finish» vuelve a aprender desde el grafo completo y descarta los borrados y orientaciones que el usuario ya había aplicado, aunque su ayuda dice que sigue desde el estado actual}\label{err:141}
\begin{ficha}
\fichakey{Estado}Presente
\fichakey{Módulo}learning.algorithm
\fichakey{Paquete}org.openmarkov.learning.algorithm.pc
\fichakey{Ficheros}\path{learning.gui/src/main/java/org/openmarkov/learning/gui/interactive/InteractiveLearningDialog.java:208-219} \newline \path{learning.gui/src/main/java/org/openmarkov/learning/gui/interactive/InteractiveLearningDialog.java:343-352} \newline \path{learning.gui/src/main/resources/learning/gui/localize/Learning_en.xml:67-69} \newline \path{learning.core/src/main/java/org/openmarkov/learning/core/LearningManager.java:107-109} \newline \path{learning.core/src/main/java/org/openmarkov/learning/core/algorithm/LearningAlgorithm.java:88-121} \newline \path{learning.algorithm/src/main/java/org/openmarkov/learning/algorithm/pc/PCAlgorithm.java:127-139} \newline \path{learning.algorithm/src/main/java/org/openmarkov/learning/algorithm/pc/IndependenceRelationsAlgorithm.java:51-77} \newline \path{learning.gui/src/main/java/org/openmarkov/learning/gui/LearningController.java:83} \newline \path{learning.gui/src/main/java/org/openmarkov/learning/gui/LearningDialog.java:919-926}
\fichakey{Comprobado}Leyendo el código y ejecutándolo
\fichakey{Esfuerzo}Medio (días)
\fichakey{Decisión de equipo}\textcolor{decisionrojo}{\textbf{Sí.}} Si «Finish» debe seguir desde la red tal como la ha dejado el usuario, como dice su ayuda emergente, o volver a aprender desde el principio
\fichakey{Corregir antes}\hyperref[err:140]{error~140}: a la vez, como mínimo: «Finish» tiene que partir de la fase real del algoritmo, que hoy corrige \texttt{init}.
\fichakey{Relacionados}\hyperref[err:133]{error~133}, \hyperref[err:134]{error~134}, \hyperref[err:139]{error~139}, \hyperref[err:140]{error~140}, \hyperref[err:142]{error~142}
\end{ficha}
\paragraph{Qué pasa}
El usuario llega por Herramientas («Tools») → «Learning» → algoritmo PC (el que aprende la estructura de la red con pruebas de independencia condicional) → «Interactive learning», que es la opción marcada por omisión → «Learn». En la ventana «Interactive learning» aplica propuestas de la tabla con «Apply edit»: quita enlaces o elige la orientación de otros. La ventana no es modal y la red aprendida está abierta en el editor, así que también puede borrar un enlace en el lienzo. Al final pulsa «Finish», cuya ayuda emergente dice «Launches the algorithm from the current state». La ventana se cierra y la red que queda es la del aprendizaje automático hecho desde el principio: lo que el usuario decidió a mano se pierde, sin aviso.

\texttt{Interactive\allowbreak{}Learning\allowbreak{}Dialog.\allowbreak{}btn\allowbreak{}Run\allowbreak{}Action\allowbreak{}Performed} llama a \texttt{Learning\allowbreak{}Manager.\allowbreak{}learn}, que llama a \texttt{Learning\allowbreak{}Algorithm.\allowbreak{}run}. \texttt{run} empieza por \texttt{init(\allowbreak{}model\allowbreak{}Net\allowbreak{}Use)}, después aplica la mejor propuesta mientras haya y termina con el aprendizaje de parámetros.

Para el PC, \texttt{init} es \texttt{PCAlgorithm.\allowbreak{}init}: devuelve la fase al principio, vacía la caché de pruebas de independencia, el estado del esqueleto y el historial de propuestas, y antes llama a \texttt{Independence\allowbreak{}Relations\allowbreak{}Algorithm.\allowbreak{}init}. Este quita todos los enlaces de la red y une todos los nodos entre sí (\texttt{prob\allowbreak{}Net.\allowbreak{}marry}), es decir, vuelve al grafo completo sin dirección. Lo hace sin red modelo, y con red modelo cuando están marcadas «Use the information of the nodes» (queda marcada al cargar la red modelo) o «Allow link addition». Si con la red modelo solo se permite borrar o invertir enlaces, conserva los enlaces que haya pero les quita la dirección, con lo que se pierden las orientaciones.

Ese mismo \texttt{init} ya se ejecutó antes de abrir la ventana (\texttt{Learning\allowbreak{}Controller.\allowbreak{}init\allowbreak{}Learning} llama a \texttt{Learning\allowbreak{}Manager.\allowbreak{}init}), así que la segunda llamada no prepara nada nuevo: solo deshace lo hecho. Lo único que no toca es la lista de propuestas bloqueadas. «Complete phase» (\texttt{Learning\allowbreak{}Algorithm.\allowbreak{}run\allowbreak{}Till\allowbreak{}Next\allowbreak{}Phase}) no pasa por \texttt{init}.

El comentario de \texttt{PCAlgorithm.\allowbreak{}init} presenta «Finish» como una ejecución nueva que no debe arrastrar la fase adelantada por el llenado de la tabla (\hyperref[err:140]{error~140}); la ayuda del botón promete lo contrario.

Medido con un pequeño programa de prueba que construye el PC como la ventana de aprendizaje con sus opciones por defecto (prueba de entropía cruzada, nivel de significación 0,05, alfa 0,5, prueba de dirección causal activada), repite las llamadas de la ventana (llenar la tabla tras cada cambio, aplicar con \texttt{Learning\allowbreak{}Manager.\allowbreak{}apply\allowbreak{}Edit}) y termina llamando a \texttt{learn} como «Finish», sobre \texttt{learning.\allowbreak{}algorithm/\allowbreak{}src/\allowbreak{}test/\allowbreak{}resources/\allowbreak{}network/\allowbreak{}BN-asia10k.\allowbreak{}csv}:

\begin{itemize}
\tightlist
\item
  El aprendizaje automático da 7 enlaces: Bronchitis→Dyspnea, Bronchitis→Smoker, LungCancer→TuberculosisOrCancer, Smoker→LungCancer, Tuberculosis→TuberculosisOrCancer, TuberculosisOrCancer→Dyspnea y TuberculosisOrCancer→X-ray.
\item
  Aplicando a mano los 21 borrados que ofrece la tabla (siempre la primera fila) y después la fila 3, «Orient link: Dyspnea --\textgreater{} Bronchitis» (motivo «Causal direction test (ANM)»), la red queda con el esqueleto de 7 enlaces y Dyspnea→Bronchitis. Tras «Finish» la red es exactamente la del aprendizaje automático, con Bronchitis→Dyspnea.
\item
  Borrando nada más abrir el enlace Smoker--LungCancer, con la misma edición \texttt{Remove\allowbreak{}Link\allowbreak{}Edit} que construye el editor al borrar un enlace, tras «Finish» el enlace vuelve (Smoker→LungCancer) y la red es otra vez la del aprendizaje automático.
\end{itemize}

La ventana es un \texttt{JDialog} y no se ha abierto: que el botón llama a \texttt{learn} y después se cierra está leído en \texttt{btn\allowbreak{}Run\allowbreak{}Action\allowbreak{}Performed}.
\paragraph{El código}
learning.gui/src/main/java/org/openmarkov/learning/gui/interactive/InteractiveLearningDialog.java:343-352

\begin{Shaded}
\begin{Highlighting}[]
\KeywordTok{private} \DataTypeTok{void} \FunctionTok{btnRunActionPerformed}\OperatorTok{(}\BuiltInTok{ActionEvent}\NormalTok{ evt}\OperatorTok{)} \KeywordTok{throws} \KeywordTok{...} \OperatorTok{\{}
    \ControlFlowTok{try} \OperatorTok{\{}
        \ControlFlowTok{if} \OperatorTok{(!}\NormalTok{learningFinished}\OperatorTok{)} \OperatorTok{\{}
\NormalTok{            learningManager}\OperatorTok{.}\FunctionTok{learn}\OperatorTok{();}
            \KeywordTok{this}\OperatorTok{.}\FunctionTok{setVisible}\OperatorTok{(}\KeywordTok{false}\OperatorTok{);}
        \OperatorTok{\}}
    \OperatorTok{\}} \ControlFlowTok{finally} \OperatorTok{\{}
        \KeywordTok{this}\OperatorTok{.}\FunctionTok{setVisible}\OperatorTok{(}\KeywordTok{false}\OperatorTok{);}
    \OperatorTok{\}}
\OperatorTok{\}}
\end{Highlighting}
\end{Shaded}

learning.core/src/main/java/org/openmarkov/learning/core/algorithm/LearningAlgorithm.java:88-92

\begin{Shaded}
\begin{Highlighting}[]
\KeywordTok{public} \DataTypeTok{void} \FunctionTok{run}\OperatorTok{(}\NormalTok{ModelNetUse modelNetUse}\OperatorTok{)} \KeywordTok{throws} \KeywordTok{...} \OperatorTok{\{}
    \FunctionTok{init}\OperatorTok{(}\NormalTok{modelNetUse}\OperatorTok{);}
    \CommentTok{/* Main loop */}
\NormalTok{    LearningEditProposal bestEdition }\OperatorTok{=} \FunctionTok{getBestEdit}\OperatorTok{(}\KeywordTok{true}\OperatorTok{,} \KeywordTok{true}\OperatorTok{);}
    \ControlFlowTok{while} \OperatorTok{(}\NormalTok{bestEdition }\OperatorTok{!=} \KeywordTok{null}\OperatorTok{)} \OperatorTok{\{}
\end{Highlighting}
\end{Shaded}

learning.algorithm/src/main/java/org/openmarkov/learning/algorithm/pc/IndependenceRelationsAlgorithm.java:51-75

\begin{Shaded}
\begin{Highlighting}[]
\AttributeTok{@Override} \KeywordTok{public} \DataTypeTok{void} \FunctionTok{init}\OperatorTok{(}\NormalTok{ModelNetUse modelNetUse}\OperatorTok{)} \OperatorTok{\{}
    \ControlFlowTok{if} \OperatorTok{(}\NormalTok{modelNetUse }\OperatorTok{!=} \KeywordTok{null} \OperatorTok{\&\&} \OperatorTok{(}\NormalTok{modelNetUse}\OperatorTok{.}\FunctionTok{isLinkAdditionAllowed}\OperatorTok{()} \OperatorTok{||}\NormalTok{ modelNetUse}\OperatorTok{.}\FunctionTok{isUseNodePositions}\OperatorTok{()))} \OperatorTok{\{}
        \ControlFlowTok{for} \OperatorTok{(}\NormalTok{Link}\OperatorTok{\textless{}}\BuiltInTok{Node}\OperatorTok{\textgreater{}}\NormalTok{ link }\OperatorTok{:}\NormalTok{ probNet}\OperatorTok{.}\FunctionTok{getLinks}\OperatorTok{())} \OperatorTok{\{}
\NormalTok{            probNet}\OperatorTok{.}\FunctionTok{removeLink}\OperatorTok{(}\NormalTok{link}\OperatorTok{);}
        \OperatorTok{\}}
\NormalTok{        probNet}\OperatorTok{.}\FunctionTok{marry}\OperatorTok{(}\NormalTok{probNet}\OperatorTok{.}\FunctionTok{getNodes}\OperatorTok{());}
    \OperatorTok{\}}
    \ControlFlowTok{else} \ControlFlowTok{if} \OperatorTok{(}\NormalTok{modelNetUse }\OperatorTok{!=} \KeywordTok{null} \OperatorTok{\&\&} \OperatorTok{(}\NormalTok{modelNetUse}\OperatorTok{.}\FunctionTok{isLinkRemovalAllowed}\OperatorTok{()} \OperatorTok{||}\NormalTok{ modelNetUse}\OperatorTok{.}\FunctionTok{isLinkInversionAllowed}\OperatorTok{()))} \OperatorTok{\{}
        \KeywordTok{...}
    \OperatorTok{\}} \ControlFlowTok{else} \OperatorTok{\{}
        \ControlFlowTok{for} \OperatorTok{(}\NormalTok{Link}\OperatorTok{\textless{}}\BuiltInTok{Node}\OperatorTok{\textgreater{}}\NormalTok{ link }\OperatorTok{:}\NormalTok{ probNet}\OperatorTok{.}\FunctionTok{getLinks}\OperatorTok{())} \OperatorTok{\{}
\NormalTok{            probNet}\OperatorTok{.}\FunctionTok{removeLink}\OperatorTok{(}\NormalTok{link}\OperatorTok{);}
        \OperatorTok{\}}
\NormalTok{        probNet}\OperatorTok{.}\FunctionTok{marry}\OperatorTok{(}\NormalTok{probNet}\OperatorTok{.}\FunctionTok{getNodes}\OperatorTok{());}
    \OperatorTok{\}}
\end{Highlighting}
\end{Shaded}

learning.gui/src/main/resources/learning/gui/localize/Learning\_en.xml:67-69

\begin{Shaded}
\begin{Highlighting}[]
\NormalTok{\textless{}}\KeywordTok{Run}\OtherTok{ value=}\StringTok{"Finish"}\NormalTok{\textgreater{}}
\NormalTok{    \textless{}}\KeywordTok{ToolTip}\OtherTok{ value=}\StringTok{"Launches the algorithm from the current state"}\NormalTok{/\textgreater{}}
\NormalTok{\textless{}/}\KeywordTok{Run}\NormalTok{\textgreater{}}
\end{Highlighting}
\end{Shaded}
\paragraph{Cómo arreglarlo}
Depende de lo que decida el equipo sobre «Finish»:

\begin{itemize}
\tightlist
\item
  Si debe seguir desde la red tal como está, como dice su ayuda: que el botón no pase por \texttt{init}. Por ejemplo, un método de \texttt{Learning\allowbreak{}Algorithm} que haga el bucle de \texttt{run} y el aprendizaje de parámetros sin llamar a \texttt{init}, y que \texttt{Learning\allowbreak{}Manager} lo ofrezca para el aprendizaje interactivo. Para el PC hay que partir entonces de la fase real y no de la adelantada por el llenado de la tabla: hoy es \texttt{init} quien la corrige, así que este arreglo va junto con el de \hyperref[err:140]{error~140}. La caché de pruebas sigue valiendo, porque los borrados aplicados quedan marcados como hechos.
\item
  Si debe volver a empezar: que el rótulo y la ayuda lo digan, y que la ventana avise antes de descartar lo aplicado.
\end{itemize}

\texttt{Learning\allowbreak{}Manager.\allowbreak{}learn} lo usan también el aprendizaje automático de \texttt{Learning\allowbreak{}Dialog} y la validación cruzada de \texttt{bn\allowbreak{}Evaluation} (\texttt{Learning\allowbreak{}Evaluator}), donde la red está recién creada y \texttt{init} no descarta nada; el cambio no debe alterarlos.

Los demás algoritmos que ofrecen aprendizaje interactivo también vuelven a ejecutar su \texttt{init} al pulsar «Finish» (por ejemplo, \texttt{Super\allowbreak{}Parent\allowbreak{}NBAlgorithm.\allowbreak{}init} vuelve a llenar su lista de características huérfanas y \texttt{Hill\allowbreak{}Climbing\allowbreak{}Algorithm.\allowbreak{}init} vuelve a añadir la restricción de número máximo de padres, si se fijó). No se ha medido qué efecto tiene en cada uno, pero la regla que se elija debe cumplirse para todos.

Prueba de regresión, repitiendo las llamadas de la ventana sobre \texttt{BN-asia10k.\allowbreak{}csv}: aplicar los 21 borrados y la orientación Dyspnea→Bronchitis y llamar al método que use «Finish». Si se decide seguir desde el estado actual, la red final debe conservar Dyspnea→Bronchitis; y, borrando a mano Smoker--LungCancer, el enlace no debe volver.

\setcounter{subsection}{141}
\subsection[PC interactivo: con un borrado bloqueado, «Finish» da NullPointerException]{En el aprendizaje interactivo con PC, después de bloquear una propuesta de borrado, «Finish» termina con un NullPointerException y deja la red a medio aprender; en bases con muchas variables fallan también «Complete phase» y «Apply edit»}\label{err:142}
\begin{ficha}
\fichakey{Estado}Presente
\fichakey{Módulo}learning.algorithm
\fichakey{Paquete}org.openmarkov.learning.algorithm.pc
\fichakey{Ficheros}\path{learning.algorithm/src/main/java/org/openmarkov/learning/algorithm/pc/SkeletonDiscovery.java:221-269} \newline \path{learning.core/src/main/java/org/openmarkov/learning/core/util/LearningEditProposal.java:70-83} \newline \path{learning.core/src/main/java/org/openmarkov/learning/core/algorithm/LearningAlgorithm.java:88-122} \newline \path{learning.core/src/main/java/org/openmarkov/learning/core/algorithm/LearningAlgorithm.java:190-233} \newline \path{learning.algorithm/src/main/java/org/openmarkov/learning/algorithm/pc/PCAlgorithm.java:170-180} \newline \path{learning.algorithm/src/main/java/org/openmarkov/learning/algorithm/pc/PCAlgorithm.java:405} \newline \path{learning.gui/src/main/java/org/openmarkov/learning/gui/interactive/InteractiveLearningDialog.java:208-218} \newline \path{learning.gui/src/main/java/org/openmarkov/learning/gui/interactive/InteractiveLearningDialog.java:327-361}
\fichakey{Comprobado}Leyendo el código y ejecutándolo
\fichakey{Esfuerzo}Pequeño (horas)
\fichakey{Decisión de equipo}No
\fichakey{Relacionados}\hyperref[err:139]{error~139}, \hyperref[err:140]{error~140}, \hyperref[err:141]{error~141}
\end{ficha}
\paragraph{Qué pasa}
El usuario llega por Herramientas («Tools») → «Learning» → algoritmo PC (el que aprende la estructura de la red con pruebas de independencia condicional) → «Interactive learning», que es la opción marcada por omisión → «Learn». En la ventana «Interactive learning» selecciona una propuesta de borrado («Remove link between \ldots») y pulsa «Block edit», para que el algoritmo no quite ese enlace. Si después pulsa «Finish», sale el diálogo de «error que un desarrollador no previó» con un \texttt{Null\allowbreak{}Pointer\allowbreak{}Exception}, la ventana interactiva se cierra y la red queda en el editor a medio aprender: con parte de los enlaces del grafo completo, ninguno orientado y sin el aprendizaje de parámetros del final.

En una base con muchas variables falla igual «Complete phase», y también seguir aplicando propuestas con «Apply edit», en cuanto se han quitado los enlaces que el algoritmo mira antes que el bloqueado. Bloquear una propuesta de orientación no falla.

Por qué pasa: para buscar el siguiente borrado, \texttt{Skeleton\allowbreak{}Discovery.\allowbreak{}get\allowbreak{}Optimal\allowbreak{}Edit\allowbreak{}From\allowbreak{}Cache} recorre las parejas de nodos por orden alfabético y, para cada candidata, \texttt{is\allowbreak{}Valid\allowbreak{}Edit} pregunta si está bloqueada construyendo \texttt{new\ Learning\allowbreak{}Edit\allowbreak{}Proposal(\allowbreak{}remove\allowbreak{}Link\allowbreak{}Edit,\allowbreak{}\ best\allowbreak{}Motivation)}. El motivo que le pasa no es el de esa pareja, sino el de la mejor candidata encontrada hasta ese momento en el recorrido, que vale \texttt{null} para la primera. \texttt{Learning\allowbreak{}Algorithm.\allowbreak{}is\allowbreak{}Blocked(\allowbreak{}Learning\allowbreak{}Edit\allowbreak{}Proposal)} la busca con \texttt{blocked\allowbreak{}Edits.\allowbreak{}contains}, que llama al \texttt{equals} de la propuesta nueva contra cada bloqueada; \texttt{Learning\allowbreak{}Edit\allowbreak{}Proposal.\allowbreak{}equals} compara primero las ediciones y, si son la misma, llama a \texttt{this.\allowbreak{}motivation.\allowbreak{}equals(\allowbreak{}.\allowbreak{}.\allowbreak{}.\allowbreak{})} con el motivo \texttt{null}.

Falla, por tanto, cuando el borrado bloqueado es la primera candidata de un recorrido, lo que ocurre en cuanto las parejas que van antes por orden alfabético ya se han quitado o no cumplen las condiciones. Por eso no falla al bloquear, porque llenar la tabla no quita enlaces. «Finish» vuelve a aprender desde el grafo completo (\hyperref[err:141]{error~141}) y va quitando enlaces hasta llegar al bloqueado; «Complete phase» (\texttt{run\allowbreak{}Till\allowbreak{}Next\allowbreak{}Phase}) y «Apply edit» también llegan, si el algoritmo sigue en la fase de quitar enlaces. En las bases pequeñas llenar la tabla ya ha llevado el algoritmo a la fase de orientación (\hyperref[err:140]{error~140}), así que «Complete phase» no pasa por los borrados; y con \texttt{Head\allowbreak{}To\allowbreak{}Head2.\allowbreak{}csv} la tabla se queda vacía al bloquear y la ventana ya no llama a \texttt{learn}.

Cuando \texttt{best\allowbreak{}Motivation} no es \texttt{null}, la misma pregunta compara el borrado bloqueado con el valor p de otra pareja y no lo reconoce; lo que evita que se proponga es la segunda comprobación, en \texttt{PCAlgorithm.\allowbreak{}get\allowbreak{}Next\allowbreak{}Edit}, que sí usa el motivo de la propuesta.

Tras el fallo, \texttt{btn\allowbreak{}Run\allowbreak{}Action\allowbreak{}Performed} («Finish») esconde la ventana en su \texttt{finally}, y el manejador global (\texttt{OMException\allowbreak{}Handler}) presenta cualquier \texttt{Runtime\allowbreak{}Exception} como error no previsto. En «Complete phase», \texttt{btn\allowbreak{}Next\allowbreak{}Phase\allowbreak{}Action\allowbreak{}Performed} quita la ventana de la lista de oyentes de las ediciones de la red antes de llamar a \texttt{run\allowbreak{}Till\allowbreak{}Next\allowbreak{}Phase} y la vuelve a añadir después, sin \texttt{finally}: por lectura, tras el fallo la ventana sigue abierta con la tabla vieja y ya no se actualiza.

Medido con un pequeño programa de prueba que construye el PC como la ventana de aprendizaje con sus opciones por defecto (prueba de entropía cruzada, nivel de significación 0,05, alfa 0,5, prueba de dirección causal activada), llena la tabla como la ventana (\texttt{get\allowbreak{}Best\allowbreak{}Edit} y después \texttt{get\allowbreak{}Next\allowbreak{}Edit} hasta 50 filas), bloquea una propuesta con \texttt{block\allowbreak{}Edit}, vuelve a llenar la tabla como hace la ventana tras «Block edit», y después llama a lo mismo que «Finish» (\texttt{Learning\allowbreak{}Manager.\allowbreak{}learn}) o «Complete phase» (\texttt{run\allowbreak{}Till\allowbreak{}Next\allowbreak{}Phase}), o aplica cada vez la primera propuesta con \texttt{Learning\allowbreak{}Manager.\allowbreak{}apply\allowbreak{}Edit}:

\begin{itemize}
\tightlist
\item
  \texttt{learning.\allowbreak{}algorithm/\allowbreak{}src/\allowbreak{}test/\allowbreak{}resources/\allowbreak{}network/\allowbreak{}BN-asia10k.\allowbreak{}csv} (8 variables; la tabla inicial tiene 22 borrados y 28 orientaciones): bloqueando cualquiera de 21 de los 22 borrados, «Finish» lanza \texttt{Null\allowbreak{}Pointer\allowbreak{}Exception:\ Cannot\ invoke\ "Object.\allowbreak{}equals(\allowbreak{}Object)"\ because\ "this.\allowbreak{}motivation"\ is\ null} en \texttt{Learning\allowbreak{}Edit\allowbreak{}Proposal.\allowbreak{}equals}, llamado desde \texttt{Skeleton\allowbreak{}Discovery.\allowbreak{}is\allowbreak{}Valid\allowbreak{}Edit}, y la red se queda con entre 8 y 28 enlaces sin dirección (el grafo completo tiene 28; el aprendizaje automático da 7, todos orientados). El otro borrado, LungCancer--Smoker, es un enlace que el aprendizaje automático no quita, y con él «Finish» termina bien. Bloqueando cualquiera de las 28 orientaciones, «Finish» termina bien.
\item
  \texttt{integration\allowbreak{}Tests/\allowbreak{}src/\allowbreak{}main/\allowbreak{}resources/\allowbreak{}networks/\allowbreak{}learning/\allowbreak{}alarm500.\allowbreak{}csv} (37 variables; la tabla inicial son 50 borrados): bloqueando cualquiera de los 50, «Finish» y «Complete phase» lanzan la misma excepción; tras «Finish» la red se queda con 220 o 235 enlaces sin dirección de los 666 del grafo completo. Bloqueando el primero y aplicando cada vez la primera propuesta, la excepción salta al volver a llenar la tabla tras 396 borrados aplicados.
\item
  \texttt{learning.\allowbreak{}algorithm/\allowbreak{}src/\allowbreak{}test/\allowbreak{}resources/\allowbreak{}network/\allowbreak{}Head\allowbreak{}To\allowbreak{}Head2.\allowbreak{}csv}: tras bloquear, la tabla queda vacía y la ventana ya no llama a \texttt{learn}, así que no falla.
\end{itemize}

La ventana es un \texttt{JDialog} y no se ha abierto: lo que enseña tras la excepción está leído en su código y en \texttt{OMException\allowbreak{}Handler}.
\paragraph{El código}
learning.algorithm/src/main/java/org/openmarkov/learning/algorithm/pc/SkeletonDiscovery.java:232-238

\begin{Shaded}
\begin{Highlighting}[]
\NormalTok{RemoveLinkEdit removeLinkEdit }\OperatorTok{=}
        \KeywordTok{new} \FunctionTok{RemoveLinkEdit}\OperatorTok{(}\NormalTok{pc}\OperatorTok{.}\FunctionTok{net}\OperatorTok{(),}\NormalTok{ nodeX}\OperatorTok{.}\FunctionTok{getVariable}\OperatorTok{(),}\NormalTok{ nodeY}\OperatorTok{.}\FunctionTok{getVariable}\OperatorTok{(),} \KeywordTok{false}\OperatorTok{);}

\ControlFlowTok{if} \OperatorTok{(}\FunctionTok{isValidEdit}\OperatorTok{(}\NormalTok{removeLinkEdit}\OperatorTok{,}\NormalTok{ bestMotivation}\OperatorTok{,}\NormalTok{ onlyAllowedEdits}\OperatorTok{))} \OperatorTok{\{}
\NormalTok{    bestMotivation }\OperatorTok{=}\NormalTok{ motivation}\OperatorTok{;}
\NormalTok{    bestEditProposal }\OperatorTok{=} \KeywordTok{new} \FunctionTok{LearningEditProposal}\OperatorTok{(}\NormalTok{removeLinkEdit}\OperatorTok{,}\NormalTok{ motivation}\OperatorTok{);}
\OperatorTok{\}}
\end{Highlighting}
\end{Shaded}

learning.algorithm/src/main/java/org/openmarkov/learning/algorithm/pc/SkeletonDiscovery.java:264-269

\begin{Shaded}
\begin{Highlighting}[]
\KeywordTok{private} \DataTypeTok{boolean} \FunctionTok{isValidEdit}\OperatorTok{(}\NormalTok{RemoveLinkEdit removeLinkEdit}\OperatorTok{,}\NormalTok{ PCEditMotivation bestMotivation}\OperatorTok{,}
                            \DataTypeTok{boolean}\NormalTok{ onlyAllowedEdits}\OperatorTok{)} \OperatorTok{\{}
    \ControlFlowTok{return} \OperatorTok{!}\NormalTok{pc}\OperatorTok{.}\FunctionTok{checkBlocked}\OperatorTok{(}\KeywordTok{new} \FunctionTok{LearningEditProposal}\OperatorTok{(}\NormalTok{removeLinkEdit}\OperatorTok{,}\NormalTok{ bestMotivation}\OperatorTok{))}
            \OperatorTok{\&\&} \OperatorTok{!}\NormalTok{pc}\OperatorTok{.}\FunctionTok{alreadyConsidered}\OperatorTok{(}\NormalTok{removeLinkEdit}\OperatorTok{,}\NormalTok{ lastRemovedEdits}\OperatorTok{)}
            \OperatorTok{\&\&} \OperatorTok{(!}\NormalTok{onlyAllowedEdits }\OperatorTok{||}\NormalTok{ pc}\OperatorTok{.}\FunctionTok{checkAllowed}\OperatorTok{(}\NormalTok{removeLinkEdit}\OperatorTok{));}
\OperatorTok{\}}
\end{Highlighting}
\end{Shaded}

learning.core/src/main/java/org/openmarkov/learning/core/util/LearningEditProposal.java:70-78

\begin{Shaded}
\begin{Highlighting}[]
\AttributeTok{@Override}
\KeywordTok{public} \DataTypeTok{boolean} \FunctionTok{equals}\OperatorTok{(}\BuiltInTok{Object}\NormalTok{ obj}\OperatorTok{)} \OperatorTok{\{}
    \ControlFlowTok{if} \OperatorTok{(}\KeywordTok{this} \OperatorTok{==}\NormalTok{ obj}\OperatorTok{)}
        \ControlFlowTok{return} \KeywordTok{true}\OperatorTok{;}
    \ControlFlowTok{if} \OperatorTok{((}\NormalTok{obj }\OperatorTok{==} \KeywordTok{null}\OperatorTok{)} \OperatorTok{||} \OperatorTok{(}\NormalTok{obj}\OperatorTok{.}\FunctionTok{getClass}\OperatorTok{()} \OperatorTok{!=} \KeywordTok{this}\OperatorTok{.}\FunctionTok{getClass}\OperatorTok{()))}
        \ControlFlowTok{return} \KeywordTok{false}\OperatorTok{;}
\NormalTok{    LearningEditProposal other }\OperatorTok{=} \OperatorTok{(}\NormalTok{LearningEditProposal}\OperatorTok{)}\NormalTok{ obj}\OperatorTok{;}
    \ControlFlowTok{return} \KeywordTok{this}\OperatorTok{.}\FunctionTok{edit}\OperatorTok{.}\FunctionTok{equals}\OperatorTok{(}\NormalTok{other}\OperatorTok{.}\FunctionTok{edit}\OperatorTok{)} \OperatorTok{\&\&} \KeywordTok{this}\OperatorTok{.}\FunctionTok{motivation}\OperatorTok{.}\FunctionTok{equals}\OperatorTok{(}\NormalTok{other}\OperatorTok{.}\FunctionTok{motivation}\OperatorTok{);}
\OperatorTok{\}}
\end{Highlighting}
\end{Shaded}
\paragraph{Cómo arreglarlo}
En \texttt{Skeleton\allowbreak{}Discovery.\allowbreak{}is\allowbreak{}Valid\allowbreak{}Edit}, preguntar si el borrado está bloqueado por la edición sola, con \texttt{pc.\allowbreak{}is\allowbreak{}Blocked(\allowbreak{}remove\allowbreak{}Link\allowbreak{}Edit)}: la variante \texttt{Learning\allowbreak{}Algorithm.\allowbreak{}is\allowbreak{}Blocked(\allowbreak{}PNEdit)} (líneas 226-233) compara solo las ediciones y es la que ya usan el ascenso de colinas y los clasificadores. Así no se construye una propuesta con el motivo de otra pareja.

Hacer que \texttt{Learning\allowbreak{}Edit\allowbreak{}Proposal.\allowbreak{}equals} admita un motivo nulo (con \texttt{Objects.\allowbreak{}equals}, como ya hace \texttt{hash\allowbreak{}Code}) quitaría la excepción, pero dejaría la pregunta comparando con el motivo de otra pareja.

Los sitios del PC que preguntan si una propuesta está bloqueada son \texttt{Skeleton\allowbreak{}Discovery.\allowbreak{}is\allowbreak{}Valid\allowbreak{}Edit}, \texttt{PCAlgorithm.\allowbreak{}get\allowbreak{}Next\allowbreak{}Edit} (línea 178), \texttt{Collider\allowbreak{}Orientation} (línea 141) y \texttt{Meek\allowbreak{}Orientation} (líneas 220 y 239); los cuatro últimos comparan edición y motivo a la vez. Conviene que todos usen el mismo criterio.

En \texttt{Interactive\allowbreak{}Learning\allowbreak{}Dialog.\allowbreak{}btn\allowbreak{}Next\allowbreak{}Phase\allowbreak{}Action\allowbreak{}Performed}, volver a añadir la ventana como oyente en un \texttt{finally}, para que un fallo no la deje sin actualizar.

Pruebas de regresión, repitiendo las llamadas de la ventana: con \texttt{BN-asia10k.\allowbreak{}csv}, llenar la tabla, bloquear el borrado Tuberculosis--VisitToAsia y llamar a \texttt{learn}, que no debe lanzar y debe dejar ese enlace en la red; con \texttt{alarm500.\allowbreak{}csv}, bloquear el primer borrado de la tabla y llamar a \texttt{run\allowbreak{}Till\allowbreak{}Next\allowbreak{}Phase}, que tampoco debe lanzar ni quitar el enlace bloqueado.

\setcounter{subsection}{142}
\subsection[El bosque aumentado (FAN) falla sin características ligadas a la clase]{El clasificador de bosque aumentado lanza un puntero nulo cuando ninguna característica tiene información mutua positiva con la clase (por ejemplo, con una base de una sola columna)}\label{err:143}
\begin{ficha}
\fichakey{Estado}Presente
\fichakey{Módulo}learning.algorithm
\fichakey{Paquete}org.openmarkov.learning.algorithm.nbderived.fanb
\fichakey{Ficheros}\path{learning.algorithm/src/main/java/org/openmarkov/learning/algorithm/nbderived/fanb/ForestAugmentedNBAlgorithm.java:85-157}
\fichakey{Comprobado}Leyendo el código y ejecutándolo
\fichakey{Esfuerzo}Pequeño (horas)
\fichakey{Decisión de equipo}No
\end{ficha}
\paragraph{Qué pasa}
El usuario llega por el menú de aprendizaje → tipo discriminativo → «Forest augmented naive bayes», con una base cuyas características no tienen relación con la clase, o marcando como variables solo la clase. Al pulsar el botón de aprender, después de haberse abierto ya la ventana de la red nueva, sale una \texttt{Null\allowbreak{}Pointer\allowbreak{}Exception} que recoge el manejador general de excepciones de la aplicación.

El clasificador FAN (Forest Augmented Naive Bayes, Bayes ingenuo aumentado con un bosque) empieza eligiendo la característica de la que colgar el primer árbol. \texttt{get\allowbreak{}Best\allowbreak{}Root\allowbreak{}For\allowbreak{}Subtree} busca la de mayor información mutua con la clase, pero parte de \texttt{best\allowbreak{}Score\ =\ 0.\allowbreak{}0} y exige \texttt{add\allowbreak{}Score\ \textgreater{}\ best\allowbreak{}Score}, así que devuelve \texttt{null} si no hay características o si todas tienen información mutua exactamente cero con la clase. \texttt{get\allowbreak{}Optimal\allowbreak{}Edit} hace a continuación \texttt{best\allowbreak{}Edit.\allowbreak{}get\allowbreak{}Variable\allowbreak{}To(\allowbreak{})} sin comprobarlo.

\texttt{KDBAlgorithm.\allowbreak{}get\allowbreak{}Optimal\allowbreak{}Edit}, con una estructura parecida, sí comprueba \texttt{best\allowbreak{}Edit\ !=\ null}.

Medido con un pequeño programa de prueba:

\begin{Verbatim}[breaklines,breakanywhere,fontsize=\small]
OneNode-1000.csv (una sola columna, que es la clase)
  -> java.lang.NullPointerException: Cannot invoke "...BaseLinkEdit.getVariableTo()" because "bestEdit" is null
     at ...ForestAugmentedNBAlgorithm.getOptimalEdit(ForestAugmentedNBAlgorithm.java:95)
base sintética de 100 casos: clase binaria + una característica con un único estado
  -> la misma excepción en la misma línea
\end{Verbatim}

El fichero \texttt{One\allowbreak{}Node-1000.\allowbreak{}csv} es \texttt{learning.\allowbreak{}algorithm/\allowbreak{}src/\allowbreak{}test/\allowbreak{}resources/\allowbreak{}network/\allowbreak{}One\allowbreak{}Node-1000.\allowbreak{}csv}. El segundo caso muestra que no hace falta una base de una columna: basta con que ninguna característica esté relacionada con la clase en la muestra (por ejemplo, columnas constantes).

Además, con menos de dos características, \texttt{compute\allowbreak{}Averaged\allowbreak{}Conditional\allowbreak{}Mutual\allowbreak{}Information} divide por \texttt{n·(\allowbreak{}n−1)\ =\ 0} y el umbral sale \texttt{Na\allowbreak{}N}, aunque eso por sí solo no produce error (ninguna comparación con \texttt{Na\allowbreak{}N} es cierta).
\paragraph{El código}
learning.algorithm/src/main/java/org/openmarkov/learning/algorithm/nbderived/fanb/ForestAugmentedNBAlgorithm.java:93-97

\begin{Shaded}
\begin{Highlighting}[]
\ControlFlowTok{if} \OperatorTok{(}\NormalTok{subtreeRoot }\OperatorTok{==} \KeywordTok{null}\OperatorTok{)} \OperatorTok{\{}
\NormalTok{    bestEdit }\OperatorTok{=} \FunctionTok{getBestRootForSubtree}\OperatorTok{();}
\NormalTok{    subtreeRoot }\OperatorTok{=}\NormalTok{ probNet}\OperatorTok{.}\FunctionTok{getNode}\OperatorTok{(}\NormalTok{bestEdit}\OperatorTok{.}\FunctionTok{getVariableTo}\OperatorTok{());}
\NormalTok{    mwst}\OperatorTok{.}\FunctionTok{redirect}\OperatorTok{(}\NormalTok{subtreeRoot}\OperatorTok{.}\FunctionTok{getVariable}\OperatorTok{(),}\NormalTok{ probNet}\OperatorTok{);}
\NormalTok{    bestPartialScore }\OperatorTok{=}\NormalTok{ unconditionedMetric}\OperatorTok{.}\FunctionTok{getScore}\OperatorTok{(}\NormalTok{bestEdit}\OperatorTok{);}
\end{Highlighting}
\end{Shaded}
\paragraph{Cómo arreglarlo}
Si \texttt{get\allowbreak{}Best\allowbreak{}Root\allowbreak{}For\allowbreak{}Subtree} devuelve \texttt{null}, no proponer nada (devolver \texttt{null}, como hacen los demás clasificadores de la familia cuando no tienen propuesta, por ejemplo \texttt{KDBAlgorithm}), sin asignar \texttt{subtree\allowbreak{}Root}.

Proteger también la división de \texttt{compute\allowbreak{}Averaged\allowbreak{}Conditional\allowbreak{}Mutual\allowbreak{}Information} para menos de dos características.

Prueba de regresión: aprender FAN con \texttt{One\allowbreak{}Node-1000.\allowbreak{}csv} y con una base cuya única característica es constante; el resultado debe ser el Bayes ingenuo (o la clase sola) sin excepción.

\setcounter{subsection}{143}
\subsection[Discretizar por frecuencias iguales con decimales falla o cuenta mal]{Discretizar por frecuencias iguales una columna numérica escrita con decimales o notación científica lanza una excepción o cuenta mal los casos}\label{err:144}
\begin{ficha}
\fichakey{Estado}Presente
\fichakey{Módulo}learning.core
\fichakey{Paquete}org.openmarkov.learning.core.preprocess
\fichakey{Ficheros}\path{learning.core/src/main/java/org/openmarkov/learning/core/preprocess/Discretization.java:271-358} \newline \path{learning.gui/src/main/java/org/openmarkov/learning/gui/LearningDialog.java:1213-1214} \newline \path{bnEvaluation/src/main/java/org/openmarkov/bnEvaluation/DataPreprocessor.java:78-83}
\fichakey{Comprobado}Leyendo el código y ejecutándolo
\fichakey{Esfuerzo}Pequeño (horas)
\fichakey{Decisión de equipo}No
\fichakey{Relacionados}\hyperref[nota:54]{nota~54}
\end{ficha}
\paragraph{Qué pasa}
El usuario llega por el preprocesado de datos de \texttt{bn\allowbreak{}Evaluation} (\texttt{Data\allowbreak{}Preprocessing\allowbreak{}Dialog} → \texttt{Data\allowbreak{}Preprocessor.\allowbreak{}process}), eligiendo la discretización «Equal frequency intervals» para una variable numérica de un CSV (valores separados por comas) con valores como los de más abajo; el diálogo solo habilita la discretización para variables cuyos estados son todos números, y estas lo son. La ventana de aprender una red tiene los mismos controles, pero su pestaña de preprocesado no se añade a la ventana (\texttt{Learning\allowbreak{}Dialog.\allowbreak{}java}, línea 544, comentada), así que por ahí no se llega. Sale una \texttt{Array\allowbreak{}Index\allowbreak{}Out\allowbreak{}Of\allowbreak{}Bounds\allowbreak{}Exception} al guardar el preprocesado, o, sin ningún aviso, unos intervalos distintos de los que corresponden a los datos.

La discretización por «frecuencias iguales» convierte una variable numérica (sus estados son los valores leídos del fichero, por ejemplo «1000000000.0») en otra de intervalos que contienen aproximadamente el mismo número de casos.

\texttt{Discretization.\allowbreak{}discretize\allowbreak{}Equal\allowbreak{}Freq} convierte los nombres de los estados en números y los ordena, pero al ordenarlos pierde a qué estado pertenece cada número. Para saber cuántos casos tiene cada valor vuelve a buscar el estado por su nombre escribiendo el número: con \texttt{Double.\allowbreak{}to\allowbreak{}String} o, si contiene «E», con un \texttt{Decimal\allowbreak{}Format(\allowbreak{}"\#\#\#.\allowbreak{}\#\#\#\#\#\#\#\#")}, y si no lo encuentra, con la parte entera.

Eso solo funciona si alguna de esas escrituras reproduce exactamente el texto del fichero. Si ninguna lo hace, \texttt{state\allowbreak{}Index} vale −1 y \texttt{histogram{[}state\allowbreak{}Index{]}} lanza \texttt{Array\allowbreak{}Index\allowbreak{}Out\allowbreak{}Of\allowbreak{}Bounds\allowbreak{}Exception}.

Medido aplicando la discretización por frecuencias iguales (3 intervalos) a cada variable numérica de las bases CSV de los recursos de prueba, en \texttt{core/\allowbreak{}src/\allowbreak{}test/\allowbreak{}resources/\allowbreak{}nets/\allowbreak{}}:

\begin{Verbatim}[breaklines,breakanywhere,fontsize=\small]
Cataract-Data2.csv               var=id  -> ArrayIndexOutOfBoundsException: Index -1 out of bounds for length 537  (estado «1000000000.0», toString «1.0E9»)
Cataract-RRDdatasetPreProc1.csv  var=ID  -> la misma excepción (estado «1.000291517E9»)
Cataract-RRDdatasetPreProc2.csv  var=ID  -> la misma excepción (estado «1.00E+09»)
\end{Verbatim}

Hay además un caso silencioso: si la escritura no coincide pero la parte entera sí coincide con \textbf{otro} estado, se usa la frecuencia de ese otro estado. Medido con una variable de estados «1», «1.50» y «3» con 10, 80 y 10 casos, pidiendo 2 intervalos: al buscar 1,5 se encuentra «1» y se suma 10 en vez de 80, no se alcanza la frecuencia objetivo y sale un único intervalo \texttt{(\allowbreak{}-Infinity\ ,\allowbreak{}\ Infinity)} con los 100 casos, cuando con las cuentas correctas habría un corte en 1,5.
\paragraph{El código}
learning.core/src/main/java/org/openmarkov/learning/core/preprocess/Discretization.java:280-321

\begin{Shaded}
\begin{Highlighting}[]
\BuiltInTok{List}\OperatorTok{\textless{}}\BuiltInTok{Double}\OperatorTok{\textgreater{}}\NormalTok{ orderedStates }\OperatorTok{=} \KeywordTok{new} \BuiltInTok{ArrayList}\OperatorTok{\textless{}}\BuiltInTok{Double}\OperatorTok{\textgreater{}();}
\ControlFlowTok{for} \OperatorTok{(}\DataTypeTok{int}\NormalTok{ i }\OperatorTok{=} \DecValTok{0}\OperatorTok{;}\NormalTok{ i }\OperatorTok{\textless{}}\NormalTok{ states}\OperatorTok{.}\FunctionTok{length}\OperatorTok{;}\NormalTok{ i}\OperatorTok{++)} \OperatorTok{\{}
    \ControlFlowTok{if} \OperatorTok{(!}\NormalTok{states}\OperatorTok{[}\NormalTok{i}\OperatorTok{].}\FunctionTok{getName}\OperatorTok{().}\FunctionTok{equals}\OperatorTok{(}\StringTok{"?"}\OperatorTok{))}
\NormalTok{        orderedStates}\OperatorTok{.}\FunctionTok{add}\OperatorTok{(}\BuiltInTok{Double}\OperatorTok{.}\FunctionTok{parseDouble}\OperatorTok{(}\NormalTok{states}\OperatorTok{[}\NormalTok{i}\OperatorTok{].}\FunctionTok{getName}\OperatorTok{()));}
\OperatorTok{\}}
\BuiltInTok{Collections}\OperatorTok{.}\FunctionTok{sort}\OperatorTok{(}\NormalTok{orderedStates}\OperatorTok{);}
\KeywordTok{...}
\ControlFlowTok{for} \OperatorTok{(}\BuiltInTok{Double}\NormalTok{ state }\OperatorTok{:}\NormalTok{ orderedStates}\OperatorTok{)} \OperatorTok{\{}
\NormalTok{    stateName }\OperatorTok{=}\NormalTok{ state}\OperatorTok{.}\FunctionTok{toString}\OperatorTok{();}
    \BuiltInTok{String}\NormalTok{ stateToSearch }\OperatorTok{=}\NormalTok{ stateName}\OperatorTok{.}\FunctionTok{contains}\OperatorTok{(}\StringTok{"E"}\OperatorTok{)} \OperatorTok{?}
\NormalTok{            decimalFormat}\OperatorTok{.}\FunctionTok{format}\OperatorTok{(}\NormalTok{state}\OperatorTok{.}\FunctionTok{doubleValue}\OperatorTok{())} \OperatorTok{:}
\NormalTok{            state}\OperatorTok{.}\FunctionTok{toString}\OperatorTok{();}
\NormalTok{    stateIndex }\OperatorTok{=}\NormalTok{ variable}\OperatorTok{.}\FunctionTok{getStateIndex}\OperatorTok{(}\NormalTok{stateToSearch}\OperatorTok{);}
    \ControlFlowTok{if} \OperatorTok{(}\NormalTok{stateIndex }\OperatorTok{==} \OperatorTok{{-}}\DecValTok{1}\OperatorTok{)} \OperatorTok{\{}
\NormalTok{        stateIndex }\OperatorTok{=}\NormalTok{ variable}\OperatorTok{.}\FunctionTok{getStateIndex}\OperatorTok{(}\StringTok{""} \OperatorTok{+}\NormalTok{ state}\OperatorTok{.}\FunctionTok{intValue}\OperatorTok{());}
    \OperatorTok{\}}
\NormalTok{    stateFreq }\OperatorTok{=}\NormalTok{ histogram}\OperatorTok{[}\NormalTok{stateIndex}\OperatorTok{];}
\end{Highlighting}
\end{Shaded}
\paragraph{Cómo arreglarlo}
No volver a buscar el estado: construir pares (valor numérico, índice del estado) al leer los nombres, ordenar los pares por valor y tomar la frecuencia con el índice guardado. Si dos nombres distintos representan el mismo número («2» y «2.0»), sus frecuencias deben sumarse en el mismo valor.

Las demás discretizaciones (\texttt{discretize\allowbreak{}MDLP}, \texttt{discretize\allowbreak{}Chi\allowbreak{}Merge}, \texttt{discretize\allowbreak{}KMeans}) recorren los casos y leen el nombre del estado de cada caso, así que no tienen este problema.

Pruebas de regresión: discretizar por frecuencias iguales la columna \texttt{id} de \texttt{Cataract-Data2.\allowbreak{}csv} sin excepción; y la variable sintética de estados «1», «1.50», «3» (10, 80, 10 casos, 2 intervalos) debe dar un corte en 1,5.

\setcounter{subsection}{144}
\subsection[«Split Dataset» lanza siempre una excepción, elija la opción que elija]{Dividir un banco de casos en entrenamiento y prueba desde «Split Dataset» lanza siempre una excepción, elija el usuario la opción que elija}\label{err:145}
\begin{ficha}
\fichakey{Estado}Presente
\fichakey{Módulo}bnEvaluation
\fichakey{Paquete}org.openmarkov.bnEvaluation
\fichakey{Ficheros}\path{bnEvaluation/src/main/java/org/openmarkov/bnEvaluation/SplitSetManager.java:53-55} \newline \path{bnEvaluation/src/main/java/org/openmarkov/bnEvaluation/SplitSetManager.java:164-165} \newline \path{bnEvaluation/src/main/java/org/openmarkov/bnEvaluation/SplitSetManager.java:196-197} \newline \path{core/src/main/java/org/openmarkov/core/model/database/CaseDatabase.java:103-105} \newline \path{bnEvaluation/src/main/java/org/openmarkov/bnEvaluation/dialog/SplitDatasetDialog.java:265-314} \newline \path{bnEvaluation/src/main/java/org/openmarkov/bnEvaluation/dialog/ResultsDialog.java:180,276} \newline \path{bnEvaluation/src/main/java/org/openmarkov/bnEvaluation/SplitSet.java:49}
\fichakey{Comprobado}Leyendo el código y ejecutándolo
\fichakey{Esfuerzo}Pequeño (horas)
\fichakey{Decisión de equipo}No
\end{ficha}
\paragraph{Qué pasa}
El usuario abre el menú Herramientas («Tools») → «Split Dataset» (\texttt{Split\allowbreak{}Dataset\allowbreak{}Plugin}), elige una base de casos y cualquiera de las tres opciones, y pulsa Aceptar. La ventana no puede terminar con ninguna opción: la excepción sale como error no controlado, no se piden los ficheros de destino ni se guarda nada.

La ventana es \texttt{Split\allowbreak{}Dataset\allowbreak{}Dialog} (módulo \texttt{bn\allowbreak{}Evaluation}), que parte un banco de casos en uno de prueba y otro de entrenamiento. Ofrece tres maneras de elegir los casos de prueba: al azar, los primeros o los últimos. Al pulsar Aceptar, \texttt{do\allowbreak{}Ok\allowbreak{}Click\allowbreak{}Before\allowbreak{}Hide} (líneas 265-284) llama, según la opción, a \texttt{Split\allowbreak{}Set\allowbreak{}Manager.\allowbreak{}generate\allowbreak{}Random\allowbreak{}Test\allowbreak{}Set}, \texttt{generate\allowbreak{}First\allowbreak{}Test\allowbreak{}Set} o \texttt{generate\allowbreak{}Last\allowbreak{}Test\allowbreak{}Set}.

Los tres métodos empiezan convirtiendo la lista de variables a \texttt{Array\allowbreak{}List}: \texttt{Array\allowbreak{}List\textless{}Variable\textgreater{}\ lista\allowbreak{}Variables\ =\ (\allowbreak{}Array\allowbreak{}List\textless{}Variable\textgreater{})\ case\allowbreak{}Database.\allowbreak{}get\allowbreak{}Variables(\allowbreak{});}. \texttt{Case\allowbreak{}Database.\allowbreak{}get\allowbreak{}Variables(\allowbreak{})} devuelve siempre \texttt{Collections.\allowbreak{}unmodifiable\allowbreak{}List(\allowbreak{}variables)}, que no es un \texttt{Array\allowbreak{}List}, así que la conversión falla con \texttt{Class\allowbreak{}Cast\allowbreak{}Exception} antes de hacer nada.

Los otros dos métodos de \texttt{Split\allowbreak{}Set\allowbreak{}Manager}, \texttt{cross\allowbreak{}Validation} y \texttt{multiple\allowbreak{}Samples}, declaran \texttt{List\textless{}Variable\textgreater{}} y funcionan. Son los que usa la validación cruzada de \texttt{bn\allowbreak{}Evaluation} (\texttt{Cross\allowbreak{}Validation\allowbreak{}Controller}), que no se ve afectada.

Tampoco se llega nunca a la ventana de resultados de una partición (\texttt{Results\allowbreak{}Dialog} con un \texttt{Split\allowbreak{}Set}) ni a \texttt{Split\allowbreak{}Set.\allowbreak{}to\allowbreak{}Table}, que solo se usan desde esta ventana.

Medido:

\begin{itemize}
\tightlist
\item
  Sobre \texttt{learning.\allowbreak{}algorithm/\allowbreak{}src/\allowbreak{}test/\allowbreak{}resources/\allowbreak{}network/\allowbreak{}BN-asia10k.\allowbreak{}csv}: \texttt{generate\allowbreak{}Random\allowbreak{}Test\allowbreak{}Set(\allowbreak{}100)}, \texttt{generate\allowbreak{}First\allowbreak{}Test\allowbreak{}Set(\allowbreak{}100)} y \texttt{generate\allowbreak{}Last\allowbreak{}Test\allowbreak{}Set(\allowbreak{}100)} lanzan \texttt{java.\allowbreak{}lang.\allowbreak{}Class\allowbreak{}Cast\allowbreak{}Exception}; \texttt{cross\allowbreak{}Validation(\allowbreak{}10)} termina bien.
\item
  Sobre una base de 10 casos y 2 variables: los tres métodos lanzan \texttt{java.\allowbreak{}lang.\allowbreak{}Class\allowbreak{}Cast\allowbreak{}Exception:\ class\ java.\allowbreak{}util.\allowbreak{}Collections\$Unmodifiable\allowbreak{}Random\allowbreak{}Access\allowbreak{}List\ cannot\ be\ cast\ to\ class\ java.\allowbreak{}util.\allowbreak{}Array\allowbreak{}List}, en las líneas 55, 165 y 197 de \texttt{Split\allowbreak{}Set\allowbreak{}Manager}.
\end{itemize}
\paragraph{El código}
bnEvaluation/src/main/java/org/openmarkov/bnEvaluation/SplitSetManager.java:53-56

\begin{Shaded}
\begin{Highlighting}[]
\KeywordTok{public}\NormalTok{ SplitSet }\FunctionTok{generateRandomTestSet}\OperatorTok{(}\DataTypeTok{int}\NormalTok{ numTest}\OperatorTok{)} \OperatorTok{\{}
    \CommentTok{//database information}
    \BuiltInTok{ArrayList}\OperatorTok{\textless{}}\NormalTok{Variable}\OperatorTok{\textgreater{}}\NormalTok{ listaVariables }\OperatorTok{=} \OperatorTok{(}\BuiltInTok{ArrayList}\OperatorTok{\textless{}}\NormalTok{Variable}\OperatorTok{\textgreater{})}\NormalTok{ caseDatabase}\OperatorTok{.}\FunctionTok{getVariables}\OperatorTok{();}
    \DataTypeTok{int}\NormalTok{ numVariables }\OperatorTok{=}\NormalTok{ listaVariables}\OperatorTok{.}\FunctionTok{size}\OperatorTok{();}
\end{Highlighting}
\end{Shaded}

core/src/main/java/org/openmarkov/core/model/database/CaseDatabase.java:103-105

\begin{Shaded}
\begin{Highlighting}[]
\KeywordTok{public} \BuiltInTok{List}\OperatorTok{\textless{}}\NormalTok{Variable}\OperatorTok{\textgreater{}} \FunctionTok{getVariables}\OperatorTok{()} \OperatorTok{\{}
    \ControlFlowTok{return} \BuiltInTok{Collections}\OperatorTok{.}\FunctionTok{unmodifiableList}\OperatorTok{(}\NormalTok{variables}\OperatorTok{);}
\OperatorTok{\}}
\end{Highlighting}
\end{Shaded}

bnEvaluation/src/main/java/org/openmarkov/bnEvaluation/dialog/SplitDatasetDialog.java:274-284

\begin{Shaded}
\begin{Highlighting}[]
\ControlFlowTok{if} \OperatorTok{(}\NormalTok{randomCasesButton}\OperatorTok{.}\FunctionTok{isSelected}\OperatorTok{())} \OperatorTok{\{}
\NormalTok{    sets }\OperatorTok{=}\NormalTok{ splitSetManager}\OperatorTok{.}\FunctionTok{generateRandomTestSet}\OperatorTok{(}\NormalTok{numTest}\OperatorTok{);}
\NormalTok{    sets}\OperatorTok{.}\FunctionTok{setTitle}\OperatorTok{(}\NormalTok{title }\OperatorTok{+} \StringTok{"Random selected.}\SpecialCharTok{\textbackslash{}n}\StringTok{"}\OperatorTok{);}
\OperatorTok{\}} \ControlFlowTok{else} \OperatorTok{\{}
    \ControlFlowTok{if} \OperatorTok{(}\NormalTok{firstCasesButton}\OperatorTok{.}\FunctionTok{isSelected}\OperatorTok{())} \OperatorTok{\{}
\NormalTok{        sets }\OperatorTok{=}\NormalTok{ splitSetManager}\OperatorTok{.}\FunctionTok{generateFirstTestSet}\OperatorTok{(}\NormalTok{numTest}\OperatorTok{);}
        \KeywordTok{...}
    \OperatorTok{\}} \ControlFlowTok{else} \OperatorTok{\{}
\NormalTok{        sets }\OperatorTok{=}\NormalTok{ splitSetManager}\OperatorTok{.}\FunctionTok{generateLastTestSet}\OperatorTok{(}\NormalTok{numTest}\OperatorTok{);}
\end{Highlighting}
\end{Shaded}
\paragraph{Cómo arreglarlo}
En los tres métodos, declarar \texttt{List\textless{}Variable\textgreater{}\ lista\allowbreak{}Variables\ =\ case\allowbreak{}Database.\allowbreak{}get\allowbreak{}Variables(\allowbreak{});} sin conversión, como ya hacen \texttt{cross\allowbreak{}Validation} y \texttt{multiple\allowbreak{}Samples}, y comprobar que después no se modifica la lista; si hace falta una lista propia, copiarla con \texttt{new\ Array\allowbreak{}List\textless{}\textgreater{}(\allowbreak{}.\allowbreak{}.\allowbreak{}.\allowbreak{})}. Los tres sitios son las líneas 55, 165 y 197 de \texttt{Split\allowbreak{}Set\allowbreak{}Manager}.

Revisar después el resto de la ventana, que hoy no se ha podido ejecutar nunca: el guardado de los dos ficheros (\texttt{Case\allowbreak{}Database\allowbreak{}Manager.\allowbreak{}get\allowbreak{}Writer(\allowbreak{}.\allowbreak{}.\allowbreak{}.\allowbreak{}).\allowbreak{}save}) y la ventana de resultados con \texttt{Split\allowbreak{}Set.\allowbreak{}to\allowbreak{}Table}.

Prueba de regresión en \texttt{bn\allowbreak{}Evaluation/\allowbreak{}src/\allowbreak{}test/\allowbreak{}java/\allowbreak{}org/\allowbreak{}openmarkov/\allowbreak{}bn\allowbreak{}Evaluation/\allowbreak{}}: para una base de N casos, los tres métodos devuelven un \texttt{Split\allowbreak{}Set} cuyo banco de prueba tiene \texttt{num\allowbreak{}Test} casos y el de entrenamiento N − \texttt{num\allowbreak{}Test}, de modo que entre los dos suman todos los casos; los primeros y los últimos casos coinciden con las filas esperadas; y \texttt{Split\allowbreak{}Set.\allowbreak{}to\allowbreak{}Table(\allowbreak{})} se construye sin error.

\setcounter{subsection}{145}
\subsection[Validación cruzada: casos por pliegue truncados y pérdida por caso inflada]{En la validación cruzada, si el número de casos no es múltiplo del número de pliegues, la ventana de resultados dice que las puntuaciones se calcularon con menos casos y la pérdida por caso de la log-verosimilitud sale más alta}\label{err:146}
\begin{ficha}
\fichakey{Estado}Presente
\fichakey{Módulo}bnEvaluation
\fichakey{Paquete}org.openmarkov.bnEvaluation.measures
\fichakey{Ficheros}\path{bnEvaluation/src/main/java/org/openmarkov/bnEvaluation/measures/MeasureValue.java:48-58} \newline \path{bnEvaluation/src/main/java/org/openmarkov/bnEvaluation/measures/MeasuresSet.java:115-161} \newline \path{bnEvaluation/src/main/java/org/openmarkov/bnEvaluation/measures/MeasuresSet.java:168-187} \newline \path{bnEvaluation/src/main/java/org/openmarkov/bnEvaluation/LearningEvaluator.java:114-136} \newline \path{bnEvaluation/src/main/java/org/openmarkov/bnEvaluation/SplitSetManager.java:229-286} \newline \path{bnEvaluation/src/main/java/org/openmarkov/bnEvaluation/export/ExcelExporter.java:221-251} \newline \path{bnEvaluation/src/main/java/org/openmarkov/bnEvaluation/dialog/ResultsDialog.java:325-332}
\fichakey{Comprobado}Leyendo el código y ejecutándolo
\fichakey{Esfuerzo}Pequeño (horas)
\fichakey{Decisión de equipo}No
\fichakey{Relacionados}\hyperref[nota:20]{nota~20}
\end{ficha}
\paragraph{Qué pasa}
El usuario llega por el menú Herramientas («Tools») → «Cross-validation» (\texttt{Cross\allowbreak{}Validation\allowbreak{}Dialog}): elige una base de casos, un algoritmo de aprendizaje, el número de pliegues («Folds», 10 por omisión) y la medida «Goodness of fit Log-Likelihood measures». La ventana de resultados (\texttt{Results\allowbreak{}Dialog}) dice «Cross validation with k=\ldots{} folders. Scores are calculated with N cases.», y en la pestaña «Scores» enseña «LOGLIKELIHOOD score» y «LOGLIKELIHOOD Loss», cuya ayuda dice «average negative log-likelihood per case». La exportación a Excel escribe lo mismo.

En la validación cruzada de k pliegues, \texttt{Learning\allowbreak{}Evaluator.\allowbreak{}run\allowbreak{}Evaluator} aprende una red con cada conjunto de entrenamiento y la puntúa con su pliegue de prueba; la puntuación de un pliegue es una suma sobre sus casos. \texttt{Measures\allowbreak{}Set.\allowbreak{}accumulate\allowbreak{}Measure\allowbreak{}Set} suma las puntuaciones y los números de casos de los k pliegues, y \texttt{set\allowbreak{}Averaged} los divide por k. Lo que se muestra es, por tanto, la puntuación media de un pliegue, y N el número medio de casos por pliegue; la hoja de Excel lo dice («Measures calculated as an average of k iterations»). La pérdida es −puntuación / N.

Cuando el número de casos n es múltiplo de k, todo cuadra: con 12 casos y 4 pliegues sale «3 cases» y la pérdida es la media por caso.

\texttt{Measure\allowbreak{}Value.\allowbreak{}average\allowbreak{}Value} divide el número de casos con división entera (\texttt{get\allowbreak{}Num\allowbreak{}Cases(\allowbreak{})\ /\allowbreak{}\ num\allowbreak{}Iterations}). Si n no es múltiplo de k, \texttt{Split\allowbreak{}Set\allowbreak{}Manager.\allowbreak{}cross\allowbreak{}Validation} hace pliegues de ⌊n/k⌋ + 1 y ⌊n/k⌋ casos, la media verdadera tiene decimales y se trunca. La ventana dice menos casos de los que tiene un pliegue medio, y la pérdida se divide por ese número truncado: sale multiplicada por n / (k·⌊n/k⌋).

Medido con un algoritmo que solo aprende parámetros (el factor no depende del algoritmo), comparando la pérdida mostrada con la suma de las log-verosimilitudes de los pliegues dividida por n:

\begin{itemize}
\tightlist
\item
  12 casos, 4 pliegues (3, 3, 3, 3): «3 cases», pérdida 1,7019, igual a la verdadera.
\item
  10 casos, 4 pliegues (3, 3, 2, 2): «2 cases», pérdida 2,0864 frente a 1,6692 (un 25 \% más).
\item
  13 casos, 4 pliegues (4, 3, 3, 3): «3 cases», pérdida 1,6879 frente a 1,5580 (un 8,3 \% más).
\item
  \texttt{core/\allowbreak{}src/\allowbreak{}test/\allowbreak{}resources/\allowbreak{}nets/\allowbreak{}Cataract-RRDdataset\allowbreak{}Pre\allowbreak{}Proc3.\allowbreak{}csv} (537 casos), 10 pliegues (siete de 54 y tres de 53): «53 cases», pérdida 2,4318 frente a 2,4001 (un 1,3 \% más).
\end{itemize}

Las bases en CSV (valores separados por comas) de \texttt{learning.\allowbreak{}algorithm/\allowbreak{}src/\allowbreak{}test/\allowbreak{}resources/\allowbreak{}network/\allowbreak{}} e \texttt{integration\allowbreak{}Tests/\allowbreak{}src/\allowbreak{}main/\allowbreak{}resources/\allowbreak{}networks/\allowbreak{}learning/\allowbreak{}} tienen 500, 1 000, 5 000 o 10 000 casos, así que con los 10 pliegues por omisión no se nota; las de \texttt{Cataract} en \texttt{core/\allowbreak{}src/\allowbreak{}test/\allowbreak{}resources/\allowbreak{}nets/\allowbreak{}} (537 casos) sí. La evaluación de una red con una sola base de prueba y el muestreo repetido («multiple samples», pruebas del mismo tamaño) no se ven afectados.
\paragraph{El código}
bnEvaluation/src/main/java/org/openmarkov/bnEvaluation/measures/MeasureValue.java:48-58

\begin{Shaded}
\begin{Highlighting}[]
\KeywordTok{public} \DataTypeTok{void} \FunctionTok{averageValue}\OperatorTok{(}\DataTypeTok{int}\NormalTok{ numIterations}\OperatorTok{)} \OperatorTok{\{}
\NormalTok{    value }\OperatorTok{=}\NormalTok{ value }\OperatorTok{/} \OperatorTok{(}\DataTypeTok{double}\OperatorTok{)}\NormalTok{ numIterations}\OperatorTok{;}
    \KeywordTok{super}\OperatorTok{.}\FunctionTok{setNumCases}\OperatorTok{(}\KeywordTok{super}\OperatorTok{.}\FunctionTok{getNumCases}\OperatorTok{()} \OperatorTok{/}\NormalTok{ numIterations}\OperatorTok{);}
\OperatorTok{\}}

\AttributeTok{@Override}
\KeywordTok{public} \DataTypeTok{void} \FunctionTok{accumulate}\OperatorTok{(}\NormalTok{Measure measure}\OperatorTok{)} \OperatorTok{\{}
    \DataTypeTok{double}\NormalTok{ valueToAdd }\OperatorTok{=} \OperatorTok{((}\NormalTok{MeasureValue}\OperatorTok{)}\NormalTok{ measure}\OperatorTok{).}\FunctionTok{getValue}\OperatorTok{();}
\NormalTok{    value }\OperatorTok{=}\NormalTok{ value }\OperatorTok{+}\NormalTok{ valueToAdd}\OperatorTok{;}
    \KeywordTok{super}\OperatorTok{.}\FunctionTok{setNumCases}\OperatorTok{(}\KeywordTok{super}\OperatorTok{.}\FunctionTok{getNumCases}\OperatorTok{()} \OperatorTok{+}\NormalTok{ measure}\OperatorTok{.}\FunctionTok{getNumCases}\OperatorTok{());}
\OperatorTok{\}}
\end{Highlighting}
\end{Shaded}

bnEvaluation/src/main/java/org/openmarkov/bnEvaluation/measures/MeasuresSet.java:176-180

\begin{Shaded}
\begin{Highlighting}[]
\ControlFlowTok{for} \OperatorTok{(}\NormalTok{MeasureValue measure }\OperatorTok{:}\NormalTok{ measures}\OperatorTok{)} \OperatorTok{\{}
\NormalTok{    rows}\OperatorTok{.}\FunctionTok{add}\OperatorTok{(}\KeywordTok{new}\NormalTok{ ScoresRow}\OperatorTok{.}\FunctionTok{Data}\OperatorTok{(}\NormalTok{measure}\OperatorTok{.}\FunctionTok{getMeasureType}\OperatorTok{()} \OperatorTok{+} \StringTok{" score"}\OperatorTok{,}\NormalTok{ measure}\OperatorTok{.}\FunctionTok{getValue}\OperatorTok{()));}
    \ControlFlowTok{if} \OperatorTok{(}\NormalTok{measure}\OperatorTok{.}\FunctionTok{getMeasureType}\OperatorTok{()} \OperatorTok{==}\NormalTok{ LOGLIKELIHOOD}\OperatorTok{)} \OperatorTok{\{}
        \DataTypeTok{double}\NormalTok{ loss }\OperatorTok{=} \OperatorTok{{-}}\NormalTok{measure}\OperatorTok{.}\FunctionTok{getValue}\OperatorTok{()} \OperatorTok{/}\NormalTok{ measure}\OperatorTok{.}\FunctionTok{getNumCases}\OperatorTok{();}
\NormalTok{        rows}\OperatorTok{.}\FunctionTok{add}\OperatorTok{(}\KeywordTok{new}\NormalTok{ ScoresRow}\OperatorTok{.}\FunctionTok{Data}\OperatorTok{(}\NormalTok{measure}\OperatorTok{.}\FunctionTok{getMeasureType}\OperatorTok{()} \OperatorTok{+} \StringTok{" Loss"}\OperatorTok{,}\NormalTok{ loss}\OperatorTok{));}
\end{Highlighting}
\end{Shaded}
\paragraph{Cómo arreglarlo}
No truncar: guardar el total de casos de prueba de los k pliegues y calcular la pérdida por caso como −(suma de las log-verosimilitudes de los pliegues) / (total de casos), o, lo que es lo mismo, −puntuación media · k / total. El número medio de casos, si se quiere enseñar, debe ser un número con decimales.

Los sitios que usan ese número son \texttt{Measures\allowbreak{}Set.\allowbreak{}get\allowbreak{}Measure\allowbreak{}Information} (el texto de la ventana), \texttt{Measures\allowbreak{}Set.\allowbreak{}build\allowbreak{}Scores\allowbreak{}Rows} (la fila «Loss», que también usa la exportación) y \texttt{Excel\allowbreak{}Exporter.\allowbreak{}write\allowbreak{}Scores} (la nota «Scores are calculated with N cases»). Conviene que el texto diga a la vez que las puntuaciones son medias de k pliegues y cuántos casos hay en total, porque en la misma ventana la matriz de confusión se da con el total de casos (\texttt{Measure\allowbreak{}Matrix.\allowbreak{}accumulate} no divide) y las puntuaciones con la media por pliegue.

Prueba de regresión: validación cruzada con 10 casos y 4 pliegues; la fila «LOGLIKELIHOOD Loss» debe ser igual a −(suma de las log-verosimilitudes de los cuatro pliegues) / 10.

\setcounter{subsection}{146}
\subsection[Quedarse sin memoria al aprender sale como error no previsto]{Quedarse sin memoria al aprender una red se presenta como un error que el programa no había previsto}\label{err:147}
\begin{ficha}
\fichakey{Estado}Presente en parte. La normalización imposible ya se muestra bien desde el commit \texttt{87129d6}: el botón de aprender enseña un diálogo con el mensaje y el consejo de usar alfa mayor que cero (\texttt{Learning.\allowbreak{}Alpha.\allowbreak{}Smoothing\allowbreak{}Advice}), y el aprendizaje interactivo la envuelve en \texttt{Unrecoverable\allowbreak{}Exception}, que el manejador global muestra como error normal (sin el consejo).
\fichakey{Módulo}learning.gui
\fichakey{Paquete}org.openmarkov.learning.gui
\fichakey{Ficheros}\path{learning.gui/src/main/java/org/openmarkov/learning/gui/LearningDialog.java:936-945} \newline \path{learning.gui/src/main/java/org/openmarkov/learning/gui/interactive/InteractiveLearningDialog.java:210-218} \newline \path{gui/src/main/java/org/openmarkov/gui/dialog/OMExceptionHandler.java:17-58}
\fichakey{Comprobado}Leyendo el código
\fichakey{Esfuerzo}Pequeño (horas)
\fichakey{Decisión de equipo}No
\fichakey{Relacionados}\hyperref[err:48]{error~48}, \hyperref[err:148]{error~148}, \hyperref[err:149]{error~149}, \hyperref[nota:174]{nota~174}
\end{ficha}
\paragraph{Qué pasa}
Quedarse sin memoria al aprender una red (bases de casos grandes, muchas variables o muchos padres) es algo a lo que el usuario puede llegar desde el botón de aprender de \texttt{Learning\allowbreak{}Dialog}. En ese caso se abre el diálogo de «error que un desarrollador no previó», en vez de un mensaje de error normal. No se ha provocado el agotamiento de memoria.

En \texttt{Learning\allowbreak{}Dialog.\allowbreak{}learn\allowbreak{}Button\allowbreak{}Action\allowbreak{}Performed}, el mismo \texttt{try} que ya trata bien la normalización imposible tiene \texttt{catch\ (\allowbreak{}Out\allowbreak{}Of\allowbreak{}Memory\allowbreak{}Error\ e1)\ \{\ throw\ new\ Unreachable\allowbreak{}Exception(\allowbreak{}new\ Not\allowbreak{}Enough\allowbreak{}Memory\allowbreak{}Exception(\allowbreak{}e1));\ \}}.

\texttt{Not\allowbreak{}Enough\allowbreak{}Memory\allowbreak{}Exception} es la excepción del dominio para este caso, y en la propagación de evidencia se muestra como error normal. Aquí, al ir envuelta en \texttt{Unreachable\allowbreak{}Exception}, el manejador global (\texttt{OMException\allowbreak{}Handler}) la presenta como un error no previsto.

En la misma línea del bloque se declaran inalcanzables otras excepciones; la de las variables no observadas sí se alcanza y se describe en \hyperref[err:148]{error~148}.
\paragraph{El código}
learning.gui/src/main/java/org/openmarkov/learning/gui/LearningDialog.java:936-945

\begin{Shaded}
\begin{Highlighting}[]
\OperatorTok{\}} \ControlFlowTok{catch} \OperatorTok{(}\BuiltInTok{OutOfMemoryError}\NormalTok{ e1}\OperatorTok{)} \OperatorTok{\{}
    \ControlFlowTok{throw} \KeywordTok{new} \FunctionTok{UnreachableException}\OperatorTok{(}\KeywordTok{new} \FunctionTok{NotEnoughMemoryException}\OperatorTok{(}\NormalTok{e1}\OperatorTok{));}
\OperatorTok{\}} \ControlFlowTok{catch} \OperatorTok{(}\NormalTok{CannotNormalizePotentialException e}\OperatorTok{)} \OperatorTok{\{}
    \BuiltInTok{JOptionPane}\OperatorTok{.}\FunctionTok{showMessageDialog}\OperatorTok{(}\KeywordTok{this}\OperatorTok{,}
\NormalTok{            e}\OperatorTok{.}\FunctionTok{getExceptionMessage}\OperatorTok{()} \OperatorTok{+} \StringTok{"}\SpecialCharTok{\textbackslash{}n}\StringTok{"} \OperatorTok{+}\NormalTok{ stringDatabase}\OperatorTok{.}\FunctionTok{getString}\OperatorTok{(}\StringTok{"Learning.Alpha.SmoothingAdvice"}\OperatorTok{),}
\NormalTok{            e}\OperatorTok{.}\FunctionTok{getExceptionTitle}\OperatorTok{(),} \BuiltInTok{JOptionPane}\OperatorTok{.}\FunctionTok{ERROR\_MESSAGE}\OperatorTok{);}
\OperatorTok{\}} \ControlFlowTok{catch} \OperatorTok{(}\NormalTok{UnobservedVariablesException }\OperatorTok{|}\NormalTok{ EmptyModelNetException }\OperatorTok{|}\NormalTok{ DoEditException e}\OperatorTok{)} \OperatorTok{\{}
    \ControlFlowTok{throw} \KeywordTok{new} \FunctionTok{UnreachableException}\OperatorTok{(}\NormalTok{e}\OperatorTok{);}
\OperatorTok{\}}
\end{Highlighting}
\end{Shaded}
\paragraph{Cómo arreglarlo}
Envolver la \texttt{Not\allowbreak{}Enough\allowbreak{}Memory\allowbreak{}Exception} en \texttt{Unrecoverable\allowbreak{}Exception} (error normal), como hacen los demás sitios que la capturan, o mostrarla con un diálogo propio.

El último \texttt{catch} tiene el mismo problema con la excepción de variables no observadas (\hyperref[err:148]{error~148}); conviene arreglarlos a la vez.

Prueba: difícil de automatizar para la memoria.

\setcounter{subsection}{147}
\subsection[Una variable de la red modelo que falta en los datos da error no previsto]{Aprender con una red modelo que tiene una variable que no está en la base de casos sale como un error que el programa no había previsto, aunque la excepción trae su propio mensaje}\label{err:148}
\begin{ficha}
\fichakey{Estado}Presente
\fichakey{Módulo}learning.gui
\fichakey{Paquete}org.openmarkov.learning.gui
\fichakey{Ficheros}\path{learning.gui/src/main/java/org/openmarkov/learning/gui/LearningDialog.java:865-895} \newline \path{learning.gui/src/main/java/org/openmarkov/learning/gui/LearningDialog.java:936-944} \newline \path{learning.gui/src/main/java/org/openmarkov/learning/gui/LearningDialog.java:544} \newline \path{learning.gui/src/main/java/org/openmarkov/learning/gui/LearningDialog.java:1031-1041} \newline \path{bnEvaluation/src/main/java/org/openmarkov/bnEvaluation/LearningEvaluator.java:139-144} \newline \path{learning.core/src/main/java/org/openmarkov/learning/core/LearningManager.java:75-93} \newline \path{learning.core/src/main/java/org/openmarkov/learning/core/LearningManager.java:165-176} \newline \path{learning.gui/src/main/java/org/openmarkov/learning/gui/LearningController.java:66-85} \newline \path{core/src/main/resources/core/localize/core_class_localizations_en.xml:446-447} \newline \path{learning.gui/src/main/resources/learning/gui/localize/Learning_en.xml:122-127}
\fichakey{Comprobado}Leyendo el código y ejecutándolo
\fichakey{Esfuerzo}Pequeño (horas)
\fichakey{Decisión de equipo}No
\fichakey{Relacionados}\hyperref[err:147]{error~147}, \hyperref[err:149]{error~149}, \hyperref[nota:10]{nota~10}
\end{ficha}
\paragraph{Qué pasa}
En la ventana de aprender una red (\texttt{Learning\allowbreak{}Dialog}), el usuario abre una base de casos y, en la pestaña de la red modelo, carga una red desde fichero; al cargarla queda marcada «Use the information of the nodes». Si la red modelo tiene alguna variable que no es una columna de la base de casos, al pulsar «Learn» con cualquier algoritmo salvo «Expectation maximization (EM)» sale el diálogo de «error que un desarrollador no previó», en vez de un mensaje que diga qué variable falta. Pasa en aprendizaje automático y en interactivo.

La ventana tiene código para elegir qué variables usar, pero ese panel no se muestra (la línea que lo añade a las pestañas está comentada), así que el usuario no puede desmarcar una variable; se llega con una base a la que le falta una columna.

\texttt{Learning\allowbreak{}Controller.\allowbreak{}init\allowbreak{}Learning} construye el \texttt{Learning\allowbreak{}Manager}, y su constructor, cuando se usa la red modelo, llama a \texttt{apply\allowbreak{}Model\allowbreak{}Net}. Si el algoritmo no admite variables no observadas (\texttt{supports\allowbreak{}Unobserved\allowbreak{}Variables\ =\ false} en \texttt{@Learning\allowbreak{}Algorithm\allowbreak{}Type}, que es el caso de todos salvo EM) y a la base le falta alguna variable de la red modelo, lanza \texttt{Unobserved\allowbreak{}Variables\allowbreak{}Exception}.

\texttt{Learning\allowbreak{}Dialog.\allowbreak{}learn\allowbreak{}Button\allowbreak{}Action\allowbreak{}Performed} la captura en \texttt{catch\ (\allowbreak{}Unobserved\allowbreak{}Variables\allowbreak{}Exception\ \textbar{}\ Empty\allowbreak{}Model\allowbreak{}Net\allowbreak{}Exception\ \textbar{}\ Do\allowbreak{}Edit\allowbreak{}Exception\ e)\ \{\ throw\ new\ Unreachable\allowbreak{}Exception(\allowbreak{}e);\ \}}, y el manejador global (\texttt{OMException\allowbreak{}Handler}) presenta cualquier \texttt{Unreachable\allowbreak{}Exception} como error no previsto.

La excepción ya tiene un texto para el usuario en \texttt{core\_\allowbreak{}class\_\allowbreak{}localizations\_\allowbreak{}en.\allowbreak{}xml} («Algorithm \{algorithmType.name\} does not support unobserved variables, which are: \{unobservedVariables\}»), que es el que se vería si llegara como error normal; y \texttt{Learning\_\allowbreak{}en.\allowbreak{}xml} tiene otro, \texttt{Learning.\allowbreak{}Error.\allowbreak{}Unobserved\allowbreak{}Variables}, que ninguna clase usa.

Medido llamando a las mismas clases que la ventana con \texttt{integration\allowbreak{}Tests/\allowbreak{}src/\allowbreak{}main/\allowbreak{}resources/\allowbreak{}networks/\allowbreak{}learning/\allowbreak{}asia10K.\allowbreak{}csv} sin la variable \texttt{Smoker} y la red modelo \texttt{integration\allowbreak{}Tests/\allowbreak{}src/\allowbreak{}main/\allowbreak{}resources/\allowbreak{}networks/\allowbreak{}learning/\allowbreak{}BN-asia.\allowbreak{}pgmx}, con «Use the information of the nodes»: el \texttt{Learning\allowbreak{}Manager} lanza, con el PC y con «Hill climbing», \texttt{Unobserved\allowbreak{}Variables\allowbreak{}Exception} con el texto «Algorithm PC does not support unobserved variables, which are: {[}Finite states variable Smoker with states {[}state no, state yes{]}{]}». La ventana no se ha abierto: que acaba en el diálogo de error no previsto está leído en el \texttt{catch} y en \texttt{OMException\allowbreak{}Handler}.

Las otras dos excepciones del mismo \texttt{catch} sí son inalcanzables desde esta ventana. \texttt{Empty\allowbreak{}Model\allowbreak{}Net\allowbreak{}Exception} solo se lanza si se pide usar la red modelo sin haberla, y la ventana solo construye \texttt{Model\allowbreak{}Net\allowbreak{}Use} cuando hay red modelo. \texttt{Do\allowbreak{}Edit\allowbreak{}Exception} solo podría venir de \texttt{Auto\allowbreak{}Arrange\allowbreak{}Edit}, que no comprueba ninguna restricción.
\paragraph{El código}
learning.gui/src/main/java/org/openmarkov/learning/gui/LearningDialog.java:936-944

\begin{Shaded}
\begin{Highlighting}[]
\OperatorTok{\}} \ControlFlowTok{catch} \OperatorTok{(}\BuiltInTok{OutOfMemoryError}\NormalTok{ e1}\OperatorTok{)} \OperatorTok{\{}
    \ControlFlowTok{throw} \KeywordTok{new} \FunctionTok{UnreachableException}\OperatorTok{(}\KeywordTok{new} \FunctionTok{NotEnoughMemoryException}\OperatorTok{(}\NormalTok{e1}\OperatorTok{));}
\OperatorTok{\}} \ControlFlowTok{catch} \OperatorTok{(}\NormalTok{CannotNormalizePotentialException e}\OperatorTok{)} \OperatorTok{\{}
    \BuiltInTok{JOptionPane}\OperatorTok{.}\FunctionTok{showMessageDialog}\OperatorTok{(}\KeywordTok{this}\OperatorTok{,}
\NormalTok{            e}\OperatorTok{.}\FunctionTok{getExceptionMessage}\OperatorTok{()} \OperatorTok{+} \StringTok{"}\SpecialCharTok{\textbackslash{}n}\StringTok{"} \OperatorTok{+}\NormalTok{ stringDatabase}\OperatorTok{.}\FunctionTok{getString}\OperatorTok{(}\StringTok{"Learning.Alpha.SmoothingAdvice"}\OperatorTok{),}
\NormalTok{            e}\OperatorTok{.}\FunctionTok{getExceptionTitle}\OperatorTok{(),} \BuiltInTok{JOptionPane}\OperatorTok{.}\FunctionTok{ERROR\_MESSAGE}\OperatorTok{);}
\OperatorTok{\}} \ControlFlowTok{catch} \OperatorTok{(}\NormalTok{UnobservedVariablesException }\OperatorTok{|}\NormalTok{ EmptyModelNetException }\OperatorTok{|}\NormalTok{ DoEditException e}\OperatorTok{)} \OperatorTok{\{}
    \ControlFlowTok{throw} \KeywordTok{new} \FunctionTok{UnreachableException}\OperatorTok{(}\NormalTok{e}\OperatorTok{);}
\OperatorTok{\}}
\end{Highlighting}
\end{Shaded}

learning.core/src/main/java/org/openmarkov/learning/core/LearningManager.java:169-176

\begin{Shaded}
\begin{Highlighting}[]
\BuiltInTok{List}\OperatorTok{\textless{}}\NormalTok{Variable}\OperatorTok{\textgreater{}}\NormalTok{ missingVariables }\OperatorTok{=} \FunctionTok{getMissingVariables}\OperatorTok{(}\NormalTok{database}\OperatorTok{.}\FunctionTok{getVariables}\OperatorTok{(),}\NormalTok{ modelNet}\OperatorTok{.}\FunctionTok{getVariables}\OperatorTok{());}
\ControlFlowTok{if} \OperatorTok{(!}\NormalTok{algorithmClass}\OperatorTok{.}\FunctionTok{getAnnotation}\OperatorTok{(}\NormalTok{LearningAlgorithmType}\OperatorTok{.}\FunctionTok{class}\OperatorTok{)}
                   \OperatorTok{.}\FunctionTok{supportsUnobservedVariables}\OperatorTok{()} \OperatorTok{\&\&} \OperatorTok{!}\NormalTok{missingVariables}\OperatorTok{.}\FunctionTok{isEmpty}\OperatorTok{())} \OperatorTok{\{}
    \KeywordTok{...}
    \ControlFlowTok{throw} \KeywordTok{new} \FunctionTok{UnobservedVariablesException}\OperatorTok{(}\NormalTok{algorithmClass}\OperatorTok{,}\NormalTok{ missingVariables}\OperatorTok{);}
\OperatorTok{\}}
\end{Highlighting}
\end{Shaded}
\paragraph{Cómo arreglarlo}
Separar \texttt{Unobserved\allowbreak{}Variables\allowbreak{}Exception} del \texttt{catch} que la declara inalcanzable y mostrarla como error de entrada del usuario: envolverla en \texttt{Unrecoverable\allowbreak{}Exception} (el manejador global la enseña entonces con su texto) o con un diálogo propio, como ya se hace con la normalización imposible en el mismo \texttt{try}. Si se quiere, la ventana puede además avisar antes, al cargar la base de casos o la red modelo.

El texto de la excepción enseña cada variable con \texttt{Variable.\allowbreak{}to\allowbreak{}String} («Finite states variable Smoker with states {[}\ldots{]}»); conviene que liste solo los nombres. Hay que decidir también qué hacer con la clave \texttt{Learning.\allowbreak{}Error.\allowbreak{}Unobserved\allowbreak{}Variables}, que hoy no se usa.

La validación cruzada de \texttt{bn\allowbreak{}Evaluation} (\texttt{Learning\allowbreak{}Evaluator.\allowbreak{}learn\allowbreak{}Train\allowbreak{}Net}) también declara inalcanzable esta excepción, y ahí sí lo es, porque no usa red modelo.

Es el mismo \texttt{catch} que el de \hyperref[err:147]{error~147}, y conviene arreglarlos a la vez.

Prueba de regresión: construir el \texttt{Learning\allowbreak{}Manager} con una base sin la variable \texttt{Smoker}, la red modelo \texttt{BN-asia.\allowbreak{}pgmx} y el PC, comprobar que lanza \texttt{Unobserved\allowbreak{}Variables\allowbreak{}Exception}, y que la ventana la trata como error de entrada (no como \texttt{Unreachable\allowbreak{}Exception}).

\setcounter{subsection}{148}
\subsection[Red modelo sin alguna columna de la base: error no previsto al cargarla]{Usar en la ventana de aprendizaje una red modelo que no tiene todas las columnas de la base de casos saca un error no previsto al cargarla y deja la ventana a medio preparar}\label{err:149}
\begin{ficha}
\fichakey{Estado}Presente
\fichakey{Módulo}learning.gui
\fichakey{Paquete}org.openmarkov.learning.gui
\fichakey{Ficheros}\path{learning.gui/src/main/java/org/openmarkov/learning/gui/LearningDialog.java:1154-1241} \newline \path{learning.gui/src/main/java/org/openmarkov/learning/gui/LearningDialog.java:736-761} \newline \path{learning.gui/src/main/java/org/openmarkov/learning/gui/LearningDialog.java:773-790} \newline \path{learning.gui/src/main/java/org/openmarkov/learning/gui/LearningDialog.java:1019-1041} \newline \path{learning.gui/src/main/java/org/openmarkov/learning/gui/LearningDialog.java:544} \newline \path{learning.gui/src/main/java/org/openmarkov/learning/gui/LearningDialog.java:1243-1248} \newline \path{core/src/main/java/org/openmarkov/core/model/network/ProbNet.java:807-809}
\fichakey{Comprobado}Leyendo el código y ejecutándolo
\fichakey{Esfuerzo}Pequeño (horas)
\fichakey{Decisión de equipo}No
\fichakey{Relacionados}\hyperref[err:147]{error~147}, \hyperref[err:148]{error~148}
\end{ficha}
\paragraph{Qué pasa}
El usuario abre Herramientas («Tools») → «Learning» (\texttt{Learning\allowbreak{}Dialog}), carga una base de casos y, en la pestaña «Model network», carga una red desde fichero («Load model network from file» → «Open») o toma la red abierta («Use the already open network»). Si la base tiene alguna columna que no es una variable de la red modelo, sale el diálogo de «error que un desarrollador no previó» con un \texttt{Null\allowbreak{}Pointer\allowbreak{}Exception}. Si carga primero la red modelo y después la base, el mismo error sale al cargar la base.

Aprender con una red modelo que solo tiene parte de las columnas es un uso previsto: al cargarla, la ventana elige por código «Use only the variables in the model network» y deja sin marcar las columnas que la red no tiene; y \texttt{Learning\allowbreak{}Manager.\allowbreak{}apply\allowbreak{}Model\allowbreak{}Net}, cuando se parte de la red modelo, le añade las variables de la base que le faltan.

\texttt{update\allowbreak{}Variable\allowbreak{}Selection\allowbreak{}Panel} rehace el panel de preprocesado: para cada columna de la base crea su casilla y sus listas de opciones. Con red modelo, las líneas 1216-1223 buscan la columna en la red modelo para proponer «Use model network» si allí está discretizada, con \texttt{model\allowbreak{}Net.\allowbreak{}get\allowbreak{}Node(\allowbreak{}variable.\allowbreak{}get\allowbreak{}Name(\allowbreak{})).\allowbreak{}get\allowbreak{}Variable(\allowbreak{})}. \texttt{Prob\allowbreak{}Net.\allowbreak{}get\allowbreak{}Node(\allowbreak{}String)} devuelve \texttt{null} cuando la red no tiene esa variable, y la llamada lanza \texttt{Null\allowbreak{}Pointer\allowbreak{}Exception}. El comentario de encima dice «if the variable is in the model net», y más arriba (línea 1186) la casilla ya pregunta \texttt{model\allowbreak{}Net.\allowbreak{}contains\allowbreak{}Variable(\allowbreak{}.\allowbreak{}.\allowbreak{}.\allowbreak{})} para la misma columna; aquí no se comprueba.

La pestaña de preprocesado no se añade a la ventana (la línea 544 está comentada), pero el método se ejecuta igual: lo llaman \texttt{load\allowbreak{}Case\allowbreak{}File}, al terminar de leer la base; \texttt{model\allowbreak{}Net\allowbreak{}Selected}, al cargar la red modelo o tomar la abierta; y el botón de variables de la red modelo, que está en ese mismo panel. El usuario ve el error sin ver el panel.

Lo que queda después depende del orden:

\begin{itemize}
\tightlist
\item
  Base de casos y después red modelo: la excepción sale de \texttt{model\allowbreak{}Net\allowbreak{}Selected}, que se queda a medias. La red modelo queda cargada y el panel invisible solo tiene las columnas anteriores a la primera que falta en la red modelo. \texttt{get\allowbreak{}Selected\allowbreak{}Variables} lee ese panel, así que «Learn» aprende solo con esas columnas. Si las que faltan van todas detrás, el resultado es el que se buscaba; si alguna va antes, las columnas siguientes se pierden: con PC o ascenso de colinas («Hill climbing») sale \texttt{Unobserved\allowbreak{}Variables\allowbreak{}Exception}, que también se enseña como error no previsto (\hyperref[err:148]{error~148}), y con EM (maximización de la esperanza) se aprende una red con solo las primeras columnas, sin más aviso.
\item
  Red modelo y después base de casos: \texttt{load\allowbreak{}Case\allowbreak{}File} captura la excepción, borra el nombre del fichero, desactiva «Learn» (\texttt{reset\allowbreak{}Case\allowbreak{}Database\allowbreak{}File}) y la vuelve a lanzar. Esa base no se puede usar con esa red modelo.
\end{itemize}

Medido con un pequeño programa de prueba que ejecuta, en el mismo orden, la parte del bucle de \texttt{update\allowbreak{}Variable\allowbreak{}Selection\allowbreak{}Panel} que usa la red modelo, y después construye el \texttt{Learning\allowbreak{}Manager} como lo hace «Learn» con las columnas que quedan (red modelo con «Use the information of the nodes» marcada y lo demás sin marcar):

\begin{itemize}
\tightlist
\item
  \texttt{learning.\allowbreak{}algorithm/\allowbreak{}src/\allowbreak{}test/\allowbreak{}resources/\allowbreak{}network/\allowbreak{}Head\allowbreak{}To\allowbreak{}Head2.\allowbreak{}csv} (columnas A, B, C, E, F) con la red modelo \texttt{learning.\allowbreak{}algorithm/\allowbreak{}src/\allowbreak{}test/\allowbreak{}resources/\allowbreak{}network/\allowbreak{}Three\allowbreak{}Nodes\allowbreak{}Ato\allowbreak{}Bto\allowbreak{}C.\allowbreak{}pgmx} (A, B, C): \texttt{Null\allowbreak{}Pointer\allowbreak{}Exception} en la columna E («Cannot invoke ``\ldots Node.getVariable()'' because the return value of ``\ldots ProbNet.getNode(String)'' is null»). Quedan A, B y C, y con PC, ascenso de colinas y EM el \texttt{Learning\allowbreak{}Manager} se construye con esas tres variables.
\item
  \texttt{integration\allowbreak{}Tests/\allowbreak{}src/\allowbreak{}main/\allowbreak{}resources/\allowbreak{}networks/\allowbreak{}learning/\allowbreak{}asia10K.\allowbreak{}csv} (primera columna Smoker, segunda LungCancer) con \texttt{integration\allowbreak{}Tests/\allowbreak{}src/\allowbreak{}main/\allowbreak{}resources/\allowbreak{}networks/\allowbreak{}learning/\allowbreak{}BN-asia.\allowbreak{}pgmx} sin el nodo LungCancer: la excepción salta en LungCancer y solo queda Smoker. Con PC y ascenso de colinas, \texttt{Unobserved\allowbreak{}Variables\allowbreak{}Exception}; con EM, el \texttt{Learning\allowbreak{}Manager} se construye con una red de una sola variable, Smoker.
\end{itemize}

La ventana es un \texttt{JDialog} y no se ha abierto: qué diálogo enseña el manejador global y en qué estado quedan el nombre del fichero y «Learn» está leído.
\paragraph{El código}
learning.gui/src/main/java/org/openmarkov/learning/gui/LearningDialog.java:1184-1187

\begin{Shaded}
\begin{Highlighting}[]
\DataTypeTok{boolean}\NormalTok{ select }\OperatorTok{=} \KeywordTok{true}\OperatorTok{;}
\ControlFlowTok{if} \OperatorTok{(}\NormalTok{modelNetVariablesRadioButton}\OperatorTok{.}\FunctionTok{isSelected}\OperatorTok{())} \OperatorTok{\{}
\NormalTok{    select }\OperatorTok{=}\NormalTok{ modelNet}\OperatorTok{.}\FunctionTok{containsVariable}\OperatorTok{(}\NormalTok{variable}\OperatorTok{.}\FunctionTok{getName}\OperatorTok{());}
\OperatorTok{\}}
\end{Highlighting}
\end{Shaded}

learning.gui/src/main/java/org/openmarkov/learning/gui/LearningDialog.java:1214-1223

\begin{Shaded}
\begin{Highlighting}[]
\CommentTok{/* If there is a model net and the variable is in the model net}
\CommentTok{ * and is discretized, then the default option should be \textquotesingle{}As in model\textquotesingle{}*/}
\ControlFlowTok{if} \OperatorTok{(}\NormalTok{modelNet }\OperatorTok{!=} \KeywordTok{null}\OperatorTok{)} \OperatorTok{\{}
\NormalTok{    VariableType variableTypeInModel }\OperatorTok{=}\NormalTok{ modelNet}\OperatorTok{.}\FunctionTok{getNode}\OperatorTok{(}\NormalTok{variable}\OperatorTok{.}\FunctionTok{getName}\OperatorTok{()).}\FunctionTok{getVariable}\OperatorTok{()}
                                               \OperatorTok{.}\FunctionTok{getVariableType}\OperatorTok{();}
    \ControlFlowTok{if} \OperatorTok{(}\NormalTok{variableTypeInModel }\OperatorTok{==}\NormalTok{ VariableType}\OperatorTok{.}\FunctionTok{DISCRETIZED}\OperatorTok{)} \OperatorTok{\{}
\NormalTok{        discretizeOptions}\OperatorTok{.}\FunctionTok{setSelectedItem}\OperatorTok{(}\NormalTok{stringDatabase}\OperatorTok{.}\FunctionTok{getString}\OperatorTok{(}\StringTok{"Learning.Discretize.ModelNet"}\OperatorTok{));}
    \OperatorTok{\}}
\OperatorTok{\}}
\end{Highlighting}
\end{Shaded}

learning.gui/src/main/java/org/openmarkov/learning/gui/LearningDialog.java:1031-1037

\begin{Shaded}
\begin{Highlighting}[]
\KeywordTok{private} \DataTypeTok{void} \FunctionTok{modelNetSelected}\OperatorTok{()} \OperatorTok{\{}
\NormalTok{    modelNetVariablesRadioButton}\OperatorTok{.}\FunctionTok{setSelected}\OperatorTok{(}\KeywordTok{true}\OperatorTok{);}
\NormalTok{    modelNetTabPanel}\OperatorTok{.}\FunctionTok{useNodePositionsCheckBox}\OperatorTok{.}\FunctionTok{setEnabled}\OperatorTok{(}\KeywordTok{true}\OperatorTok{);}
\NormalTok{    modelNetTabPanel}\OperatorTok{.}\FunctionTok{useNodePositionsCheckBox}\OperatorTok{.}\FunctionTok{setSelected}\OperatorTok{(}\KeywordTok{true}\OperatorTok{);}
\NormalTok{    modelNetTabPanel}\OperatorTok{.}\FunctionTok{startFromModelNetCheckBox}\OperatorTok{.}\FunctionTok{setEnabled}\OperatorTok{(}\KeywordTok{true}\OperatorTok{);}

    \FunctionTok{updateVariableSelectionPanel}\OperatorTok{();}
\end{Highlighting}
\end{Shaded}
\paragraph{Cómo arreglarlo}
Comprobar que la red modelo tiene la columna antes de pedir su tipo, como dice el comentario: tomar \texttt{Node\ model\allowbreak{}Node\ =\ model\allowbreak{}Net.\allowbreak{}get\allowbreak{}Node(\allowbreak{}variable.\allowbreak{}get\allowbreak{}Name(\allowbreak{}))} y proponer «Use model network» solo si \texttt{model\allowbreak{}Node\ !=\ null} y su variable está discretizada.

Es el único sitio de la ventana que pide a la red modelo el nodo de una columna sin comprobarlo. Lo que viene después no necesita cambio: con «Use only the variables in the model network» las columnas que la red no tiene quedan sin marcar, \texttt{Filter\allowbreak{}Database} las quita antes de discretizar, y \texttt{Learning\allowbreak{}Manager.\allowbreak{}adapt\allowbreak{}Database\allowbreak{}To\allowbreak{}Model\allowbreak{}Net} ya comprueba que la variable exista en la base.

Prueba: una prueba de interfaz con pantalla (AssertJ-Swing) que cargue \texttt{Head\allowbreak{}To\allowbreak{}Head2.\allowbreak{}csv} y después \texttt{Three\allowbreak{}Nodes\allowbreak{}Ato\allowbreak{}Bto\allowbreak{}C.\allowbreak{}pgmx} sin excepción y compruebe que \texttt{get\allowbreak{}Selected\allowbreak{}Variables} da A, B y C; y en el orden contrario, que la base queda cargada y «Learn» activo. Si se prefiere una prueba sin pantalla, sacar el cuerpo del bucle a un método que no dependa de Swing.

% ══════════════════════════════════════════════════════════════════════
\section{Interfaz: editor de red y paneles de las relaciones}\label{sec:editor}

El lienzo donde se dibuja la red, el diálogo de propiedades de un nodo y los paneles con que se edita cada tipo de relación.

\setcounter{subsection}{149}
\subsection[Volver a elegir la clase de relación que ya tiene un nodo borra sus números]{Volver a elegir en el diálogo del nodo la clase de relación que ya tiene le borra todos sus números}\label{err:150}
\begin{ficha}
\fichakey{Estado}Presente
\fichakey{Módulo}gui
\fichakey{Paquete}org.openmarkov.gui.dialog.node
\fichakey{Ficheros}\path{gui/src/main/java/org/openmarkov/gui/dialog/node/PotentialEditPanel.java:228-256, 488-522, 658-669, 675-683} \newline \path{gui/src/main/java/org/openmarkov/gui/dialog/node/NodePropertiesDialog.java:236-262, 307-332, 399-430} \newline \path{core/src/main/java/org/openmarkov/core/model/network/potential/TablePotential.java:55-59}
\fichakey{Comprobado}Leyendo el código y ejecutándolo
\fichakey{Esfuerzo}Pequeño (horas)
\fichakey{Decisión de equipo}No
\end{ficha}
\paragraph{Qué pasa}
En modo edición, al abrir las propiedades de un nodo (doble clic sobre él), la pestaña «Probability» (o «Utility») tiene el selector «Relation Type:», que ofrece las clases de relación aplicables al nodo («Table», «Tree/ADD», «Hazard (Weibull)»\ldots). Si el usuario despliega el selector y pulsa la opción que ya aparece seleccionada, la tabla del nodo pasa a ser uniforme; en las demás clases, los parámetros vuelven a sus valores por omisión. Es un gesto fácil de hacer sin querer: abrir la lista para mirar qué opciones hay y cerrarla pinchando en la actual.

Qué pasa luego depende de lo que haga el usuario en el diálogo (\texttt{Node\allowbreak{}Properties\allowbreak{}Dialog}):

\begin{itemize}
\tightlist
\item
  Con «OK», y también con solo \textbf{cambiar a otra pestaña} (\texttt{handle\allowbreak{}Change\allowbreak{}Tab} → \texttt{accept\allowbreak{}Potential\allowbreak{}Edit\allowbreak{}Changes}), \texttt{commit\allowbreak{}Changes} compara con el potencial original, ve que ha cambiado y ejecuta un \texttt{Set\allowbreak{}Potential\allowbreak{}Edit}. Los números quedan perdidos en la red; la edición se puede deshacer desde el menú de edición si el usuario se da cuenta.
\item
  Para no perderlos, tiene que pulsar el botón «Reset Probability» de la pestaña (que llama a \texttt{uncommit\allowbreak{}Changes} y repone el potencial original) o «Cancel» en el diálogo (que, según el código, deshace las ediciones hechas dentro de él con \texttt{cancel\allowbreak{}Last\allowbreak{}Sub\allowbreak{}Edit\allowbreak{}History}; esto no se ha ejecutado).
\end{itemize}

La subclase \texttt{Impose\allowbreak{}Policy\allowbreak{}Panel} (imponer una política a una decisión) hereda el mismo selector y el mismo método.

La pestaña contiene un \texttt{Potential\allowbreak{}Edit\allowbreak{}Panel}, cuyo selector (\texttt{potential\allowbreak{}Type\allowbreak{}Combo\allowbreak{}Box}) se suscribe con \texttt{add\allowbreak{}Action\allowbreak{}Listener(\allowbreak{}evt\ -\textgreater{}\ this.\allowbreak{}potential\allowbreak{}Type\allowbreak{}Changed(\allowbreak{}))}. Un \texttt{JCombo\allowbreak{}Box} de Swing lanza el evento de acción cada vez que se elige un elemento, \textbf{también cuando es el que ya estaba elegido} (medido: dos eventos al volver a elegir dos veces el mismo elemento).

\texttt{potential\allowbreak{}Type\allowbreak{}Changed} no mira si la clase elegida es la que el nodo ya tiene. Si es «Table» y el potencial actual es una tabla (o uniforme), construye con \texttt{instanciate\allowbreak{}Potential} un \texttt{Table\allowbreak{}Potential} nuevo, cuyo constructor lo deja uniforme (\texttt{set\allowbreak{}Uniform(\allowbreak{})}). Para cualquier otra clase, la última rama hace lo mismo con el constructor de esa clase, que deja los parámetros por omisión. Después pone ese potencial en el nodo (\texttt{set\allowbreak{}Potential\allowbreak{}In\allowbreak{}Node}) y redibuja el panel: los números que tenía el nodo desaparecen de la vista y del nodo.

Medido con un pequeño programa de prueba sobre la clase real del panel. El programa fabrica el \texttt{Potential\allowbreak{}Edit\allowbreak{}Panel} sin ejecutar su constructor (el constructor crea un editor de comentarios que es un \texttt{JDialog}, y eso no se puede hacer sin pantalla), le da el nodo y elige en su selector real la clase que el nodo ya tiene:

\begin{Verbatim}[breaklines,breakanywhere,fontsize=\small]
0_miscellaneous/networks/bn/BN-hepar.pgmx, nodo ggtp: TablePotential, 384 valores, 341 distintos, primeros [0.68, 0.144, 0.096, 0.08, 0.40357143, 0.13214286]
el selector muestra TablePotential; se vuelve a elegir TablePotential
después: TablePotential, 1 valor distinto, primeros [0.25, 0.25, 0.25, 0.25, 0.25, 0.25]
potentialHasChanged (lo que decide si se guarda al aceptar): true

0_miscellaneous/networks/mid/MID-hip-Briggs.pgmx, nodo Failure [1]: coeficientes [0.3740968, -5.490935, -0.0367, 0.7685, -1.344474], con matriz de covarianzas
se vuelve a elegir WeibullHazardPotential
después: coeficientes [0, 0, 0, 0, 0], sin matriz de covarianzas
\end{Verbatim}

El redibujado posterior falla en el programa de prueba por no tener pantalla; el cambio del nodo ya se ha hecho antes.
\paragraph{El código}
gui/src/main/java/org/openmarkov/gui/dialog/node/PotentialEditPanel.java:476-506

\begin{Shaded}
\begin{Highlighting}[]
\KeywordTok{private} \DataTypeTok{void} \FunctionTok{potentialTypeChanged}\OperatorTok{()} \OperatorTok{\{}
    \BuiltInTok{Class}\OperatorTok{\textless{}?} \KeywordTok{extends}\NormalTok{ Potential}\OperatorTok{\textgreater{}}\NormalTok{ potentialType }\OperatorTok{=} \OperatorTok{(}\BuiltInTok{Class}\OperatorTok{\textless{}?} \KeywordTok{extends}\NormalTok{ Potential}\OperatorTok{\textgreater{})} \KeywordTok{this}\OperatorTok{.}\FunctionTok{potentialTypeComboBox}\OperatorTok{.}\FunctionTok{getSelectedItem}\OperatorTok{();}
\NormalTok{    Potential newPotential }\OperatorTok{=} \KeywordTok{null}\OperatorTok{;}
    \DataTypeTok{var}\NormalTok{ currentPotential }\OperatorTok{=} \KeywordTok{this}\OperatorTok{.}\FunctionTok{node}\OperatorTok{.}\FunctionTok{getPotential}\OperatorTok{();}
    \ControlFlowTok{if} \OperatorTok{(}\NormalTok{potentialType }\OperatorTok{==}\NormalTok{ TablePotential}\OperatorTok{.}\FunctionTok{class} \OperatorTok{\&\&} \OperatorTok{(}\NormalTok{currentPotential }\KeywordTok{instanceof}\NormalTok{ TablePotential }\OperatorTok{||}\NormalTok{ currentPotential }\KeywordTok{instanceof}\NormalTok{ UniformPotential}\OperatorTok{))} \OperatorTok{\{}
        \DataTypeTok{var}\NormalTok{ variables }\OperatorTok{=} \KeywordTok{this}\OperatorTok{.}\FunctionTok{node}\OperatorTok{.}\FunctionTok{getParents}\OperatorTok{()}
                                 \OperatorTok{.}\FunctionTok{stream}\OperatorTok{()}
                                 \OperatorTok{.}\FunctionTok{map}\OperatorTok{(}\BuiltInTok{Node}\OperatorTok{::}\NormalTok{getVariable}\OperatorTok{)}
                                 \OperatorTok{.}\FunctionTok{collect}\OperatorTok{(}\NormalTok{Collectors}\OperatorTok{.}\FunctionTok{toCollection}\OperatorTok{(}\BuiltInTok{ArrayList}\OperatorTok{::}\KeywordTok{new}\OperatorTok{));}
\NormalTok{        variables}\OperatorTok{.}\FunctionTok{addFirst}\OperatorTok{(}\KeywordTok{this}\OperatorTok{.}\FunctionTok{node}\OperatorTok{.}\FunctionTok{getVariable}\OperatorTok{());}
\NormalTok{        newPotential }\OperatorTok{=} \KeywordTok{this}\OperatorTok{.}\FunctionTok{instanciatePotential}\OperatorTok{(}\NormalTok{potentialType}\OperatorTok{,}\NormalTok{ variables}\OperatorTok{);}
    \OperatorTok{\}}
    \KeywordTok{...}
    \ControlFlowTok{if} \OperatorTok{(}\NormalTok{newPotential }\OperatorTok{==} \KeywordTok{null}\OperatorTok{)} \OperatorTok{\{}
\NormalTok{        newPotential }\OperatorTok{=} \KeywordTok{this}\OperatorTok{.}\FunctionTok{instanciatePotential}\OperatorTok{(}\NormalTok{potentialType}\OperatorTok{,}\NormalTok{ currentPotential}\OperatorTok{.}\FunctionTok{getVariables}\OperatorTok{());}
    \OperatorTok{\}}
    \KeywordTok{this}\OperatorTok{.}\FunctionTok{setPotentialInNode}\OperatorTok{(}\NormalTok{newPotential}\OperatorTok{);}
    \FunctionTok{reloadPotentialPanel}\OperatorTok{();}
\end{Highlighting}
\end{Shaded}
\paragraph{Cómo arreglarlo}
Al principio de \texttt{potential\allowbreak{}Type\allowbreak{}Changed}, salir sin hacer nada si la clase elegida es exactamente la del potencial actual del nodo (\texttt{potential\allowbreak{}Type\ ==\ current\allowbreak{}Potential.\allowbreak{}get\allowbreak{}Class(\allowbreak{})}). El caso «Table» con un \texttt{Uniform\allowbreak{}Potential} actual sí es un cambio de clase (esa rama convierte el uniforme en tabla), así que debe seguir haciéndose.

La comprobación debe estar en \texttt{potential\allowbreak{}Type\allowbreak{}Changed} y no en el oyente, para que cubra también a \texttt{Impose\allowbreak{}Policy\allowbreak{}Panel}, que hereda el método, y a cualquier llamada programática a \texttt{set\allowbreak{}Selected\allowbreak{}Item} sobre el selector (hoy \texttt{set\allowbreak{}Enabled\allowbreak{}Decision\allowbreak{}Options}, línea 616, lo hace para políticas probabilistas). Como alternativa, escuchar con un \texttt{Item\allowbreak{}Listener} que solo actúe en \texttt{Item\allowbreak{}Event.\allowbreak{}SELECTED}, que Swing no lanza si el elemento no cambia; pero la comprobación explícita protege también de llamadas directas.

Prueba de regresión: construir el panel sobre el nodo \texttt{ggtp} de \texttt{0\_\allowbreak{}miscellaneous/\allowbreak{}networks/\allowbreak{}bn/\allowbreak{}BN-hepar.\allowbreak{}pgmx}, elegir otra vez «Table» y comprobar que los 384 valores no cambian y que \texttt{potential\allowbreak{}Has\allowbreak{}Changed(\allowbreak{})} devuelve \texttt{false}; y lo mismo con un nodo «Hazard (Weibull)», como el de la medida.

\setcounter{subsection}{150}
\subsection[En una regresión, lo escrito como Cholesky se pierde al guardar la red]{Cambiar en el panel de una regresión de la matriz de covarianzas a su descomposición de Cholesky no borra la matriz vieja, y al guardar la red se pierde lo escrito}\label{err:151}
\begin{ficha}
\fichakey{Estado}Presente
\fichakey{Módulo}gui
\fichakey{Paquete}org.openmarkov.gui.dialog.common
\fichakey{Ficheros}\path{gui/src/main/java/org/openmarkov/gui/dialog/common/GLMPotentialPanel.java:123-127,141-173} \newline \path{core/src/main/java/org/openmarkov/core/model/network/potential/GLMPotential.java:180-191} \newline \path{io/src/main/java/org/openmarkov/io/probmodel/writer/PGMXWriter_0_2.java:1107-1115}
\fichakey{Comprobado}Leyendo el código y ejecutándolo
\fichakey{Esfuerzo}Pequeño (horas)
\fichakey{Decisión de equipo}No
\end{ficha}
\paragraph{Qué pasa}
Camino del usuario: abrir las propiedades de un nodo con potencial «Linear combination» o «Exponential» que tenga matriz de covarianzas → marcar «Uncertainty» si no lo está → cambiar el selector a «Cholesky decomposition» → escribir los números → aceptar (y guardar la red). Consecuencias:

\begin{enumerate}
\def\labelenumi{\arabic{enumi}.}
\tightlist
\item
  Al volver a abrir el panel, el selector vuelve a «Covariance matrix» con los números viejos. Lo que se ve y lo que el modelo usa para muestrear dejan de coincidir.
\item
  Al guardar la red, el fichero lleva la matriz vieja; al volver a abrirlo, la descomposición se recalcula a partir de ella. \textbf{Lo que escribió el usuario se pierde.}
\item
  Con el selector en «Cholesky decomposition», desmarcar «Uncertainty» no quita la incertidumbre: la matriz de covarianzas vieja sobrevive.
\end{enumerate}

Un potencial de regresión (\texttt{GLMPotential}) puede llevar incertidumbre sobre sus coeficientes dada de dos maneras: la matriz de covarianzas (\texttt{covariance\allowbreak{}Matrix}) o su descomposición de Cholesky (\texttt{cholesky\allowbreak{}Decomposition}, una matriz triangular L tal que L·Lᵀ es la de covarianzas). El muestreo usa siempre la de Cholesky. \texttt{set\allowbreak{}Covariance\allowbreak{}Matrix} guarda la matriz y recalcula la descomposición; \texttt{set\allowbreak{}Cholesky\allowbreak{}Decomposition} guarda la descomposición pero \textbf{no toca} la matriz de covarianzas.

\texttt{GLMPotential\allowbreak{}Panel} (el panel de «Linear combination»/«Linear regression» y «Exponential») ofrece el selector «Covariance matrix» / «Cholesky decomposition». Al aceptar, \texttt{save\allowbreak{}Changes} hace una copia del potencial actual ---que conserva las dos matrices--- y escribe solo la del tipo elegido. Si el nodo tenía matriz de covarianzas y el usuario cambia a «Cholesky decomposition» y escribe la descomposición, la copia se queda con la matriz de covarianzas vieja y la descomposición nueva. De ahí salen las tres consecuencias:

\begin{enumerate}
\def\labelenumi{\arabic{enumi}.}
\tightlist
\item
  Al reabrir el panel, \texttt{set\allowbreak{}Data} mira primero la matriz de covarianzas, y por eso el selector vuelve a «Covariance matrix».
\item
  El escritor ProbModelXML escribe la matriz de covarianzas si existe, y solo si no existe, la descomposición.
\item
  Desmarcar «Uncertainty» con el selector en «Cholesky decomposition» escribe \texttt{set\allowbreak{}Cholesky\allowbreak{}Decomposition(\allowbreak{}null)}, que deja intacta la matriz de covarianzas.
\end{enumerate}

Medido de extremo a extremo con un pequeño programa de prueba sobre el nodo «Age at state entry {[}0{]}» de \texttt{0\_\allowbreak{}miscellaneous/\allowbreak{}networks/\allowbreak{}mid/\allowbreak{}MID-CHAP-Ryan-Griffin.\allowbreak{}pgmx} (potencial «Linear combination», tres covariables), al que se le pone antes una matriz de covarianzas identidad; se construye el panel sin pantalla, se cambia el selector, se escribe la descomposición y se acepta:

\begin{Verbatim}[breaklines,breakanywhere,fontsize=\small]
abierto: selector=Covariance matrix  celda(1,1)=1.0
tras aceptar: covarianzas=[1,0,1,0,0,1]  Cholesky=[2,0.5,3,0.25,0.125,4]
reabierto: selector=Covariance matrix  triángulo=1 0 1 0 0 1
guardado con PGMXWriter_1_0 y leído: covarianzas=[1,0,1,0,0,1]  Cholesky=[1,0,1,0,0,1]
\end{Verbatim}

\texttt{Weibull\allowbreak{}Potential\allowbreak{}Panel} (paneles de «Hazard (Weibull)» y «Hazard (Exponential)») no tiene este defecto: construye un potencial nuevo con el constructor que recibe el tipo de matriz, y así solo existe la que se eligió.
\paragraph{El código}
gui/src/main/java/org/openmarkov/gui/dialog/common/GLMPotentialPanel.java:141-169

\begin{Shaded}
\begin{Highlighting}[]
\KeywordTok{public} \DataTypeTok{boolean} \FunctionTok{saveChanges}\OperatorTok{()} \KeywordTok{throws}\NormalTok{ DoEditException }\OperatorTok{\{}
\NormalTok{    GLMPotential newPotential }\OperatorTok{=} \OperatorTok{(}\NormalTok{GLMPotential}\OperatorTok{)} \KeywordTok{this}\OperatorTok{.}\FunctionTok{potential}\OperatorTok{.}\FunctionTok{copy}\OperatorTok{();}
    \KeywordTok{...}
    \BuiltInTok{String}\NormalTok{ matrixTypeName }\OperatorTok{=}\NormalTok{ matrixTypeComboBox}\OperatorTok{.}\FunctionTok{getSelectedItem}\OperatorTok{().}\FunctionTok{toString}\OperatorTok{();}
\NormalTok{    newPotential}\OperatorTok{.}\FunctionTok{setCovariates}\OperatorTok{(}\NormalTok{covariates}\OperatorTok{);}
\NormalTok{    newPotential}\OperatorTok{.}\FunctionTok{setCoefficients}\OperatorTok{(}\NormalTok{coefficients}\OperatorTok{);}
    \ControlFlowTok{if} \OperatorTok{(}\NormalTok{matrixTypeName}\OperatorTok{.}\FunctionTok{equals}\OperatorTok{(}\NormalTok{MATRIX\_TYPE\_COVARIANCE}\OperatorTok{))} \OperatorTok{\{}
\NormalTok{        newPotential}\OperatorTok{.}\FunctionTok{setCovarianceMatrix}\OperatorTok{(}\NormalTok{uncertaintyMatrix}\OperatorTok{);}
    \OperatorTok{\}} \ControlFlowTok{else} \OperatorTok{\{}
\NormalTok{        newPotential}\OperatorTok{.}\FunctionTok{setCholeskyDecomposition}\OperatorTok{(}\NormalTok{uncertaintyMatrix}\OperatorTok{);}
    \OperatorTok{\}}
\end{Highlighting}
\end{Shaded}

core/src/main/java/org/openmarkov/core/model/network/potential/GLMPotential.java:180-191

\begin{Shaded}
\begin{Highlighting}[]
\KeywordTok{public} \DataTypeTok{void} \FunctionTok{setCovarianceMatrix}\OperatorTok{(}\DataTypeTok{double}\OperatorTok{[]}\NormalTok{ covarianceMatrix}\OperatorTok{)} \OperatorTok{\{}
    \KeywordTok{this}\OperatorTok{.}\FunctionTok{covarianceMatrix} \OperatorTok{=}\NormalTok{ covarianceMatrix}\OperatorTok{;}
    \KeywordTok{this}\OperatorTok{.}\FunctionTok{choleskyDecomposition} \OperatorTok{=} \FunctionTok{calculateCholesky}\OperatorTok{(}\NormalTok{covarianceMatrix}\OperatorTok{);}
\OperatorTok{\}}
\KeywordTok{...}
\KeywordTok{public} \DataTypeTok{void} \FunctionTok{setCholeskyDecomposition}\OperatorTok{(}\DataTypeTok{double}\OperatorTok{[]}\NormalTok{ choleskyDecomposition}\OperatorTok{)} \OperatorTok{\{}
    \KeywordTok{this}\OperatorTok{.}\FunctionTok{choleskyDecomposition} \OperatorTok{=}\NormalTok{ choleskyDecomposition}\OperatorTok{;}
\OperatorTok{\}}
\end{Highlighting}
\end{Shaded}
\paragraph{Cómo arreglarlo}
Que dar la incertidumbre de una manera borre la otra. Hay dos lugares posibles, con consecuencias distintas:

\begin{itemize}
\tightlist
\item
  \begin{enumerate}
  \def\labelenumi{(\alph{enumi})}
  \tightlist
  \item
    En el panel: poner a \texttt{null} la matriz de covarianzas antes de \texttt{set\allowbreak{}Cholesky\allowbreak{}Decomposition}, o construir el potencial nuevo con el constructor que recibe \texttt{Matrix\allowbreak{}Type}, como ya hace \texttt{Weibull\allowbreak{}Potential\allowbreak{}Panel}.
  \end{enumerate}
\item
  \begin{enumerate}
  \def\labelenumi{(\alph{enumi})}
  \setcounter{enumi}{1}
  \tightlist
  \item
    En el modelo: que \texttt{set\allowbreak{}Cholesky\allowbreak{}Decomposition} ponga \texttt{covariance\allowbreak{}Matrix\ =\ null}. Esta opción cubre a cualquier llamador, pero el lector ProbModelXML (\texttt{PGMXPotential\allowbreak{}Parsers.\allowbreak{}java:219-224}) y el constructor de copia de \texttt{GLMPotential} (líneas 82-90) llaman a esos métodos y hay que confirmar que siguen haciendo lo mismo (hoy solo llaman a uno de los dos, así que no cambian).
  \end{enumerate}
\end{itemize}

Prueba de regresión: la misma secuencia de la medida ---potencial con covarianzas, cambiar a Cholesky, aceptar--- y comprobar que \texttt{get\allowbreak{}Covariance\allowbreak{}Matrix(\allowbreak{})} es \texttt{null}, que al reabrir el selector dice «Cholesky decomposition» con los números escritos, y que tras escribir y leer el fichero la descomposición es la escrita.

\setcounter{subsection}{151}
\subsection[Añadir una fila a la tabla de vida y aceptar borra el primer tramo]{Añadir una fila a la tabla de vida y aceptar sin tocarla borra la probabilidad del primer tramo}\label{err:152}
\begin{ficha}
\fichakey{Estado}Presente
\fichakey{Módulo}gui
\fichakey{Paquete}org.openmarkov.gui.dialog.common
\fichakey{Ficheros}\path{gui/src/main/java/org/openmarkov/gui/dialog/common/PiecewiseExponentialTablePanel.java:32,70-90,145-147} \newline \path{gui/src/main/java/org/openmarkov/gui/dialog/common/PiecewiseExponentialPanel.java:103-121} \newline \path{gui/src/main/java/org/openmarkov/gui/dialog/common/KeyTablePanel.java:400-413}
\fichakey{Comprobado}Leyendo el código y ejecutándolo
\fichakey{Esfuerzo}Pequeño (horas)
\fichakey{Decisión de equipo}No
\fichakey{Relacionados}\hyperref[err:164]{error~164}, \hyperref[err:184]{error~184}, \hyperref[nota:13]{nota~13}
\end{ficha}
\paragraph{Qué pasa}
El panel de la tabla de vida (\texttt{Piecewise\allowbreak{}Exponential\allowbreak{}Panel} con su tabla \texttt{Piecewise\allowbreak{}Exponential\allowbreak{}Table\allowbreak{}Panel}, módulo \texttt{gui}) enseña una fila por tramo con dos columnas visibles, «Time» y «Probability», y botones para añadir, quitar, subir y bajar filas.

\begin{itemize}
\tightlist
\item
  \textbf{Añadir} una fila y aceptar sin editarla deja a 0 la probabilidad del primer tramo, sin aviso. Lo mismo ocurre si se añaden dos filas y se edita solo una.
\item
  \textbf{Subir y bajar} cambian el orden de las filas en la tabla, pero ese orden no se conserva: al reabrir el diálogo las filas vuelven a salir ordenadas por tiempo. Son dos botones que no hacen nada que se conserve.
\end{itemize}

Al aceptar, \texttt{save\allowbreak{}Changes} (líneas 103-121) recorre las filas y las mete en un \texttt{Tree\allowbreak{}Map\textless{}Double,\allowbreak{}\ Double\textgreater{}} con el tiempo como clave.

\texttt{action\allowbreak{}Performed\allowbreak{}Add\allowbreak{}Value} (líneas 145-147) añade la fila \texttt{\{0,\allowbreak{}\ 0,\allowbreak{}\ 0\}}, es decir, tiempo 0 y probabilidad 0. Al volcar las filas en el mapa, la fila nueva (que va la última) sustituye a la que ya empezaba en el tiempo 0, que es lo habitual en una tabla de vida.

Los botones de subir y bajar se heredan de \texttt{Key\allowbreak{}Table\allowbreak{}Panel} (líneas 400-413); el orden que dan se pierde al volcar las filas en el mapa ordenado por tiempo.

Medido con un pequeño programa de prueba sobre el panel real, con un nodo de suceso de una red de simulación de eventos discretos (DES) cuya tabla es \texttt{\{0.\allowbreak{}0=0.\allowbreak{}1,\allowbreak{}\ 5.\allowbreak{}0=0.\allowbreak{}2\}}: tras añadir una fila y guardar, la tabla queda \texttt{\{0.\allowbreak{}0=0.\allowbreak{}0,\allowbreak{}\ 5.\allowbreak{}0=0.\allowbreak{}2\}}.

Los textos del panel («Use rates», «Use first interval», «Life Table», «Time», «Probability», «Load data file», «Choose a CSV data file») están escritos en el código, pero eso no forma parte de este error.
\paragraph{El código}
gui/src/main/java/org/openmarkov/gui/dialog/common/PiecewiseExponentialTablePanel.java:145-147

\begin{Shaded}
\begin{Highlighting}[]
\KeywordTok{protected} \DataTypeTok{void} \FunctionTok{actionPerformedAddValue}\OperatorTok{(}\BuiltInTok{ActionEvent}\NormalTok{ e}\OperatorTok{)} \OperatorTok{\{}
\NormalTok{    tableModel}\OperatorTok{.}\FunctionTok{addRow}\OperatorTok{(}\KeywordTok{new} \BuiltInTok{Object}\OperatorTok{[]\{}\DecValTok{0}\OperatorTok{,} \DecValTok{0}\OperatorTok{,} \DecValTok{0}\OperatorTok{\});}
\OperatorTok{\}}
\end{Highlighting}
\end{Shaded}

gui/src/main/java/org/openmarkov/gui/dialog/common/PiecewiseExponentialPanel.java:109-115

\begin{Shaded}
\begin{Highlighting}[]
\BuiltInTok{TreeMap}\OperatorTok{\textless{}}\BuiltInTok{Double}\OperatorTok{,} \BuiltInTok{Double}\OperatorTok{\textgreater{}}\NormalTok{ treeMapTable }\OperatorTok{=} \KeywordTok{new} \BuiltInTok{TreeMap}\OperatorTok{\textless{}\textgreater{}();}
\ControlFlowTok{for} \OperatorTok{(}\DataTypeTok{int}\NormalTok{ i }\OperatorTok{=} \DecValTok{0}\OperatorTok{;}\NormalTok{ i }\OperatorTok{\textless{}}\NormalTok{ piecewiseExponentialTablePanel}\OperatorTok{.}\FunctionTok{valuesTable}\OperatorTok{.}\FunctionTok{getRowCount}\OperatorTok{();}\NormalTok{ i}\OperatorTok{++)} \OperatorTok{\{}
\NormalTok{    treeMapTable}\OperatorTok{.}\FunctionTok{put}\OperatorTok{(}\BuiltInTok{Double}\OperatorTok{.}\FunctionTok{parseDouble}\OperatorTok{(}\KeywordTok{...} \FunctionTok{getValueAt}\OperatorTok{(}\NormalTok{i}\OperatorTok{,} \DecValTok{1}\OperatorTok{,} \KeywordTok{null}\OperatorTok{).}\FunctionTok{toString}\OperatorTok{()),}
                     \BuiltInTok{Double}\OperatorTok{.}\FunctionTok{parseDouble}\OperatorTok{(}\KeywordTok{...} \FunctionTok{getValueAt}\OperatorTok{(}\NormalTok{i}\OperatorTok{,} \DecValTok{2}\OperatorTok{,} \KeywordTok{null}\OperatorTok{).}\FunctionTok{toString}\OperatorTok{()));}
\OperatorTok{\}}
\end{Highlighting}
\end{Shaded}
\paragraph{Cómo arreglarlo}
Que la fila nueva proponga un tiempo que no exista (por ejemplo, el último tiempo más uno) y que \texttt{save\allowbreak{}Changes} rechace con un mensaje dos filas con el mismo tiempo en vez de quedarse con la última.

Los botones de subir y bajar pueden quitarse del panel (pasando \texttt{reorderable=false} al constructor de \texttt{Key\allowbreak{}Table\allowbreak{}Panel}, línea 32), o la tabla puede ordenarse sola por tiempo.

Como en \hyperref[err:164]{error~164}, \texttt{save\allowbreak{}Changes} usa \texttt{Double.\allowbreak{}parse\allowbreak{}Double} sin comprobar el texto.

Prueba: construir el panel sobre una tabla que empiece en 0, añadir una fila, guardar y comprobar que la probabilidad del tramo 0 no cambia (o que el guardado se rechaza por tiempo repetido).

\setcounter{subsection}{152}
\subsection[El editor de una distribución univariante enseña «1» en cada parámetro]{El editor de una distribución univariante enseña «1» en cada parámetro cuando el fichero solo trae los números, y editarlo separa las dos copias de los parámetros}\label{err:153}
\begin{ficha}
\fichakey{Estado}Presente
\fichakey{Módulo}core, gui, io
\fichakey{Paquete}org.openmarkov.core.model.network.potential; org.openmarkov.gui.dialog.common; org.openmarkov.gui.action; org.openmarkov.io.probmodel.reader
\fichakey{Ficheros}\path{core/src/main/java/org/openmarkov/core/model/network/potential/UnivariateDistrPotential.java:48,330-343} \newline \path{io/src/main/java/org/openmarkov/io/probmodel/reader/PGMXReader_1_0.java:126-166} \newline \path{gui/src/main/java/org/openmarkov/gui/dialog/common/UnivariateDistrPotentialPanel.java:231-248} \newline \path{gui/src/main/java/org/openmarkov/gui/action/AugmentedPotentialValueEdit.java:97-118,236-238} \newline \path{io/src/main/java/org/openmarkov/io/probmodel/writer/PGMXWriter_1_0.java:266-289}
\fichakey{Comprobado}Leyendo el código y ejecutándolo
\fichakey{Esfuerzo}Medio (días)
\fichakey{Decisión de equipo}No
\fichakey{Relacionados}\hyperref[err:9]{error~9}, \hyperref[err:112]{error~112}
\end{ficha}
\paragraph{Qué pasa}
Al abrir el editor de la relación de un nodo cuya distribución univariante se leyó de un fichero que no trae \texttt{\textless{}Functions\textgreater{}}, el usuario ve «1» en cada casilla de parámetro en lugar de los números del fichero. Por ejemplo, en \texttt{0\_\allowbreak{}miscellaneous/\allowbreak{}networks/\allowbreak{}id/\allowbreak{}1-0/\allowbreak{}ID-CEA-minimal-1-0.\allowbreak{}pgmx}, la relación de «Effectiveness» enseña «1» en las dos casillas en lugar de 0.8 y 1.3. No se ha abierto el diálogo: la conclusión sale de la lectura del panel y de la medida sobre el potencial leído.

Además, si el usuario escribe en una casilla, cambia solo una de las dos copias de los parámetros, y el fichero guardado puede acabar diciendo cosas distintas en cada una.

Un \texttt{Univariate\allowbreak{}Distr\allowbreak{}Potential} (módulo \texttt{core}) es una relación cuyo valor sigue una distribución de probabilidad sobre un número (por ejemplo «Exact», un valor fijo), con los parámetros por cada combinación de estados de los padres. Los parámetros se guardan dos veces en su tabla aumentada (\texttt{Augmented\allowbreak{}Prob\allowbreak{}Table}): como números (\texttt{get\allowbreak{}Values(\allowbreak{})}) y como expresiones (\texttt{get\allowbreak{}Function\allowbreak{}Values(\allowbreak{})}), que permiten escribir fórmulas que nombran otras variables. Al crear la tabla, \texttt{initialize\allowbreak{}Augmented\allowbreak{}Prob\allowbreak{}Table} pone todas las expresiones a «1» (\texttt{INITIALIZATION\_\allowbreak{}VALUE}, línea 48).

El lector (\texttt{PGMXReader\_\allowbreak{}1\_\allowbreak{}0.\allowbreak{}get\allowbreak{}Univariate\allowbreak{}Distr\allowbreak{}Potential}, líneas 126-166) solo rellena las expresiones si el fichero trae \texttt{\textless{}Functions\textgreater{}}, y siempre pone en la tabla los números leídos de \texttt{\textless{}Parameters\textgreater{}} o \texttt{\textless{}Values\textgreater{}} (el orden de estas dos cosas se corrigió en el commit \texttt{51a0cd6}). Si el fichero no trae \texttt{\textless{}Functions\textgreater{}}, las expresiones siguen valiendo «1».

\texttt{Univariate\allowbreak{}Distr\allowbreak{}Potential\allowbreak{}Panel} (el panel registrado para esta clase) llena la tabla con \texttt{get\allowbreak{}Table\allowbreak{}Potential(\allowbreak{}).\allowbreak{}get\allowbreak{}Function\allowbreak{}Values(\allowbreak{})} (línea 237), es decir, con las expresiones.

Al escribir en una casilla, \texttt{Augmented\allowbreak{}Potential\allowbreak{}Value\allowbreak{}Edit} modifica solo las expresiones (\texttt{new\allowbreak{}Augmented\allowbreak{}Values\ =\ new\allowbreak{}Augmented\allowbreak{}Prob\allowbreak{}Table.\allowbreak{}get\allowbreak{}Function\allowbreak{}Values(\allowbreak{})}, líneas 108-116 y 238) y deja los números como estaban. El escritor (\texttt{PGMXWriter\_\allowbreak{}1\_\allowbreak{}0.\allowbreak{}get\allowbreak{}Univariate\allowbreak{}Distr\allowbreak{}Potential}, líneas 266-274) guarda las dos copias: \texttt{\textless{}Functions\textgreater{}} con las expresiones y \texttt{\textless{}Parameters\textgreater{}} con los números. Qué copia usa cada consumidor (inferencia, simulación, análisis de sensibilidad) no se ha revisado.

Medido con un pequeño programa de prueba sobre \texttt{0\_\allowbreak{}miscellaneous/\allowbreak{}networks/\allowbreak{}id/\allowbreak{}1-0/\allowbreak{}ID-CEA-minimal-1-0.\allowbreak{}pgmx}, cuyo potencial de «Effectiveness» es \texttt{\textless{}Potential\ type="Univariate\allowbreak{}Distr"\ distribution="Exact"\textgreater{}} con \texttt{\textless{}Values\textgreater{}0.\allowbreak{}8\ 1.\allowbreak{}3\textless{}/\allowbreak{}Values\textgreater{}}: tras leer, números \texttt{{[}0.\allowbreak{}8,\allowbreak{}\ 1.\allowbreak{}3{]}} y expresiones \texttt{{[}1,\allowbreak{}\ 1{]}}; en «Cost», números \texttt{{[}0.\allowbreak{}0,\allowbreak{}\ 14000.\allowbreak{}0{]}} y expresiones \texttt{{[}1,\allowbreak{}\ 1{]}}.
\paragraph{El código}
gui/src/main/java/org/openmarkov/gui/dialog/common/UnivariateDistrPotentialPanel.java:231-246

\begin{Shaded}
\begin{Highlighting}[]
    \AttributeTok{@Override} \KeywordTok{protected} \BuiltInTok{Object}\OperatorTok{[][]} \FunctionTok{setPotentialDataInCentreArea}\OperatorTok{(}\BuiltInTok{Object}\OperatorTok{[][]}\NormalTok{ oldValues}\OperatorTok{)} \OperatorTok{\{}
        \BuiltInTok{Object}\OperatorTok{[][]}\NormalTok{ values }\OperatorTok{=}\NormalTok{ oldValues}\OperatorTok{;}
        \DataTypeTok{int}\NormalTok{ numColumns }\OperatorTok{=}\NormalTok{ values}\OperatorTok{[}\DecValTok{0}\OperatorTok{].}\FunctionTok{length}\OperatorTok{;}
        \CommentTok{// rounding initial values}
\NormalTok{        VariableExpression}\OperatorTok{[]}\NormalTok{ initialValues }\OperatorTok{=} \FunctionTok{getTablePotential}\OperatorTok{().}\FunctionTok{getFunctionValues}\OperatorTok{();}
        \ControlFlowTok{for} \OperatorTok{(}\DataTypeTok{int}\NormalTok{ j }\OperatorTok{=} \DecValTok{1}\OperatorTok{;}\NormalTok{ j }\OperatorTok{\textless{}=}\NormalTok{ numColumns }\OperatorTok{{-}} \DecValTok{1}\OperatorTok{;}\NormalTok{ j}\OperatorTok{++)} \OperatorTok{\{}
            \ControlFlowTok{for} \OperatorTok{(}\DataTypeTok{int}\NormalTok{ i }\OperatorTok{=} \FunctionTok{getLastEditableRow}\OperatorTok{();}\NormalTok{ i }\OperatorTok{\textgreater{}=} \FunctionTok{getFirstEditableRow}\OperatorTok{();}\NormalTok{ i}\OperatorTok{{-}{-})} \OperatorTok{\{}
                \DataTypeTok{int}\NormalTok{ potentialIndex }\OperatorTok{=}\NormalTok{ PotentialsTablePanelOperations}\OperatorTok{.}\FunctionTok{getPotentialIndex}\OperatorTok{(}\NormalTok{i}\OperatorTok{,}\NormalTok{ j}\OperatorTok{,} \FunctionTok{getTablePotential}\OperatorTok{());}
\NormalTok{                VariableExpression value }\OperatorTok{=}\NormalTok{ initialValues}\OperatorTok{[}\NormalTok{potentialIndex}\OperatorTok{];}
\NormalTok{                values}\OperatorTok{[}\NormalTok{i}\OperatorTok{][}\NormalTok{j}\OperatorTok{]} \OperatorTok{=}\NormalTok{ value}\OperatorTok{;}
            \OperatorTok{\}}
        \OperatorTok{\}}
        \ControlFlowTok{return}\NormalTok{ values}\OperatorTok{;}
    \OperatorTok{\}}
\end{Highlighting}
\end{Shaded}
\paragraph{Cómo arreglarlo}
Hay que fijar una sola regla para las dos copias y cumplirla en los tres sitios:

\begin{itemize}
\tightlist
\item
  Al leer (\texttt{PGMXReader\_\allowbreak{}1\_\allowbreak{}0.\allowbreak{}get\allowbreak{}Univariate\allowbreak{}Distr\allowbreak{}Potential}): sin \texttt{\textless{}Functions\textgreater{}}, construir las expresiones a partir de los números leídos, en lugar de dejar «1».
\item
  Al enseñar (\texttt{Univariate\allowbreak{}Distr\allowbreak{}Potential\allowbreak{}Panel.\allowbreak{}set\allowbreak{}Potential\allowbreak{}Data\allowbreak{}In\allowbreak{}Centre\allowbreak{}Area}).
\item
  Al editar (\texttt{Augmented\allowbreak{}Potential\allowbreak{}Value\allowbreak{}Edit}): cuando la expresión escrita es un número, actualizar también el número de esa casilla.
\end{itemize}

Conviene revisar \texttt{Augmented\allowbreak{}Prob\allowbreak{}Table\allowbreak{}Potential\allowbreak{}Panel} (línea 279), que lee las expresiones igual.

Prueba: leer \texttt{0\_\allowbreak{}miscellaneous/\allowbreak{}networks/\allowbreak{}id/\allowbreak{}1-0/\allowbreak{}ID-CEA-minimal-1-0.\allowbreak{}pgmx} y comprobar que expresiones y números de «Effectiveness» coinciden (0.8 y 1.3); editar una casilla con «2» y comprobar que las dos copias valen 2 y que el fichero guardado las escribe iguales.

\setcounter{subsection}{153}
\subsection[En la tabla de utilidad, una restricción de enlace bloquea otras casillas]{En la tabla de un nodo de utilidad, una restricción de enlace bloquea las casillas equivocadas y deja editables las prohibidas}\label{err:154}
\begin{ficha}
\fichakey{Estado}Presente
\fichakey{Módulo}core, gui
\fichakey{Paquete}org.openmarkov.core.model.network.potential.operation; org.openmarkov.gui.dialog.common
\fichakey{Ficheros}\path{core/src/main/java/org/openmarkov/core/model/network/potential/operation/LinkRestrictionPotentialOperations.java:70-113, 403-416} \newline \path{gui/src/main/java/org/openmarkov/gui/dialog/common/TablePotentialPanel.java:579-602, 616-640} \newline \path{gui/src/main/java/org/openmarkov/gui/validator/LinkRestrictionValidator.java:35-58}
\fichakey{Comprobado}Leyendo el código y ejecutándolo
\fichakey{Esfuerzo}Pequeño (horas)
\fichakey{Decisión de equipo}No
\fichakey{Relacionados}\hyperref[err:25]{error~25}, \hyperref[err:155]{error~155}
\end{ficha}
\paragraph{Qué pasa}
Una restricción de enlace declara qué combinaciones de estados de un padre y del hijo son incompatibles. En la tabla del hijo, las casillas de esas combinaciones salen en rojo y bloqueadas. Cuando el hijo es un nodo de utilidad, el usuario ve en rojo y bloqueadas casillas que no tienen nada que ver, y puede escribir en las que la restricción prohíbe.

Camino de llegada: \texttt{Link\allowbreak{}Restriction\allowbreak{}Validator.\allowbreak{}validate} permite la restricción en un enlace de un padre de estados finitos a un nodo de utilidad. La restricción \texttt{No\allowbreak{}Link\allowbreak{}Restriction} está activa por defecto y solo la desactivan las redes de análisis de decisiones (DAN), las DAN de Markov y las redes de ajuste. En esas redes: botón derecho sobre el enlace → editar la restricción, marcar una combinación incompatible, y abrir la tabla del nodo de utilidad. No se ha buscado si alguna red del repositorio trae ya una restricción así.

El panel de la tabla (\texttt{Table\allowbreak{}Potential\allowbreak{}Panel.\allowbreak{}get\allowbreak{}Not\allowbreak{}Editable\allowbreak{}Positions}, líneas 597-611) pide a \texttt{Link\allowbreak{}Restriction\allowbreak{}Potential\allowbreak{}Operations.\allowbreak{}get\allowbreak{}State\allowbreak{}Combinations\allowbreak{}With\allowbreak{}Link\allowbreak{}Restriction} (módulo \texttt{core}) las combinaciones prohibidas, las traduce a fila y columna con \texttt{get\allowbreak{}Row\allowbreak{}And\allowbreak{}Column\allowbreak{}For\allowbreak{}State\allowbreak{}Combination} (líneas 560-583), las bloquea y el pintor las pone en rojo.

En un nodo de utilidad la tabla es un \texttt{Exact\allowbreak{}Distr\allowbreak{}Potential}, y la función toma su tabla interior, cuyas variables son solo los padres (la variable de utilidad no está). Así lo tienen todos los nodos de utilidad de tabla de las redes del repositorio que se miraron.

\texttt{get\allowbreak{}State\allowbreak{}Combinations\allowbreak{}With\allowbreak{}Link\allowbreak{}Restriction} busca los dos extremos del enlace en esa lista con índices que empiezan en 0 y solo se asignan si la variable aparece. La variable de utilidad (\texttt{var2}) no aparece, así que \texttt{var2Index} se queda en 0 (el primer padre), y \texttt{get\allowbreak{}State\allowbreak{}Combinations} hace \texttt{combination{[}var1Index{]}\ =\ var1Value;\ combination{[}var2Index{]}\ =\ var2Value;}, que escribe el estado de la utilidad (siempre 0) encima del primer padre, en un vector con una posición de más.

Además, \texttt{get\allowbreak{}Row\allowbreak{}And\allowbreak{}Column\allowbreak{}For\allowbreak{}State\allowbreak{}Combination} calcula la columna recorriendo las variables desde la segunda, porque supone que la primera es la del nodo; en la tabla interior eso salta el primer padre.

Medido con un pequeño programa de prueba sobre \texttt{0\_\allowbreak{}miscellaneous/\allowbreak{}networks/\allowbreak{}dan/\allowbreak{}DAN-3-test-problem.\allowbreak{}pgmx}, nodo de utilidad «Quality of life», padres Disease (absent, present) y Therapy (no, yes); las columnas de la tabla son (absent,no), (absent,yes), (present,no), (present,yes):

\begin{itemize}
\tightlist
\item
  restricción «Therapy = yes incompatible»: se bloquea solo (present,no); deberían bloquearse (absent,yes) y (present,yes);
\item
  restricción «Disease = present incompatible»: se bloquean (absent,no) y (present,no); deberían bloquearse (present,no) y (present,yes);
\item
  restricción «Disease = absent incompatible»: sale lo mismo que con «present»; el estado del padre se pierde.
\end{itemize}
\paragraph{El código}
core/src/main/java/org/openmarkov/core/model/network/potential/operation/LinkRestrictionPotentialOperations.java:89-107

\begin{Shaded}
\begin{Highlighting}[]
            \BuiltInTok{Map}\OperatorTok{\textless{}}\BuiltInTok{Integer}\OperatorTok{,} \BuiltInTok{Integer}\OperatorTok{\textgreater{}}\NormalTok{ independentVariables }\OperatorTok{=} \KeywordTok{new} \BuiltInTok{HashMap}\OperatorTok{\textless{}\textgreater{}();}
            \DataTypeTok{int}\NormalTok{ var1Index }\OperatorTok{=} \DecValTok{0}\OperatorTok{,}\NormalTok{ var2Index }\OperatorTok{=} \DecValTok{0}\OperatorTok{;}
            \ControlFlowTok{for} \OperatorTok{(}\DataTypeTok{int}\NormalTok{ i }\OperatorTok{=} \DecValTok{0}\OperatorTok{;}\NormalTok{ i }\OperatorTok{\textless{}}\NormalTok{ nodeVariables}\OperatorTok{.}\FunctionTok{size}\OperatorTok{();}\NormalTok{ i}\OperatorTok{++)} \OperatorTok{\{}
\NormalTok{                Variable variable }\OperatorTok{=}\NormalTok{ nodeVariables}\OperatorTok{.}\FunctionTok{get}\OperatorTok{(}\NormalTok{i}\OperatorTok{);}
                \ControlFlowTok{if} \OperatorTok{(}\NormalTok{variable}\OperatorTok{.}\FunctionTok{equals}\OperatorTok{(}\NormalTok{var1}\OperatorTok{))} \OperatorTok{\{}
\NormalTok{                    var1Index }\OperatorTok{=}\NormalTok{ i}\OperatorTok{;}
                \OperatorTok{\}} \ControlFlowTok{else} \OperatorTok{\{}
                    \ControlFlowTok{if} \OperatorTok{(}\NormalTok{variable}\OperatorTok{.}\FunctionTok{equals}\OperatorTok{(}\NormalTok{var2}\OperatorTok{))} \OperatorTok{\{}
\NormalTok{                        var2Index }\OperatorTok{=}\NormalTok{ i}\OperatorTok{;}
                    \OperatorTok{\}} \ControlFlowTok{else} \OperatorTok{\{}
\NormalTok{                        independentVariables}\OperatorTok{.}\FunctionTok{put}\OperatorTok{(}\NormalTok{i}\OperatorTok{,}\NormalTok{ variable}\OperatorTok{.}\FunctionTok{getNumStates}\OperatorTok{());}
                    \OperatorTok{\}}
                \OperatorTok{\}}
            \OperatorTok{\}}
            \KeywordTok{...}
\NormalTok{                        stateList}\OperatorTok{.}\FunctionTok{addAll}\OperatorTok{(}\NormalTok{LinkRestrictionPotentialOperations}
                                \OperatorTok{.}\FunctionTok{getStateCombinations}\OperatorTok{(}\NormalTok{independentVariables}\OperatorTok{,}\NormalTok{ nodeVariables}\OperatorTok{,}\NormalTok{ i}\OperatorTok{,}\NormalTok{ var1Index}\OperatorTok{,}\NormalTok{ j}\OperatorTok{,}\NormalTok{ var2Index}\OperatorTok{));}
\end{Highlighting}
\end{Shaded}
\paragraph{Cómo arreglarlo}
Hay que corregir los dos sitios a la vez:

\begin{itemize}
\tightlist
\item
  En \texttt{get\allowbreak{}State\allowbreak{}Combinations\allowbreak{}With\allowbreak{}Link\allowbreak{}Restriction}/\texttt{get\allowbreak{}State\allowbreak{}Combinations}, construir las combinaciones con tantas coordenadas como variables tenga la tabla que se mira, y fijar solo las de las variables que están en ella (la de utilidad no aporta coordenada, así que cada estado prohibido del padre bloquea la columna entera).
\item
  En \texttt{Table\allowbreak{}Potential\allowbreak{}Panel.\allowbreak{}get\allowbreak{}Row\allowbreak{}And\allowbreak{}Column\allowbreak{}For\allowbreak{}State\allowbreak{}Combination}, calcular la columna a partir de las variables de los padres de la tabla que se mira, sin suponer que la primera es la del nodo cuando la tabla es la interior de un \texttt{Exact\allowbreak{}Distr\allowbreak{}Potential} (la fila, en utilidad, es siempre la única).
\end{itemize}

Prueba: el caso medido arriba, comprobando que «Therapy = yes» bloquea exactamente las columnas (absent,yes) y (present,yes), y una prueba equivalente con un nodo de azar para no romper el caso que hoy funciona.

\setcounter{subsection}{154}
\subsection[La ventana de restricción de un enlace edita la tabla del padre]{Pulsar una casilla de la ventana de restricción de un enlace edita la tabla del padre: cambia en silencio una de sus probabilidades o lanza una excepción, y si el padre es una decisión la ventana ni siquiera se abre}\label{err:155}
\begin{ficha}
\fichakey{Estado}Presente
\fichakey{Módulo}gui
\fichakey{Paquete}org.openmarkov.gui.component; org.openmarkov.gui.action; org.openmarkov.gui.dialog.link
\fichakey{Ficheros}\path{gui/src/main/java/org/openmarkov/gui/component/LinkRestrictionValuesTable.java:48-87} \newline \path{gui/src/main/java/org/openmarkov/gui/component/ValuesTable.java:142-149,170-207} \newline \path{gui/src/main/java/org/openmarkov/gui/action/TablePotentialValueEdit.java:144-157} \newline \path{gui/src/main/java/org/openmarkov/gui/component/PotentialsTablePanelOperations.java:48-75,204-207} \newline \path{gui/src/main/java/org/openmarkov/gui/dialog/link/LinkRestrictionPanel.java:85-91,140-154} \newline \path{gui/src/main/java/org/openmarkov/gui/dialog/link/LinkRestrictionEditDialog.java:39-48,99-108} \newline \path{gui/src/main/java/org/openmarkov/gui/window/MainPanelListenerAssistant.java:251-263} \newline \path{gui/src/main/java/org/openmarkov/gui/validator/LinkRestrictionValidator.java:35-58}
\fichakey{Comprobado}Leyendo el código y ejecutándolo
\fichakey{Esfuerzo}Pequeño (horas)
\fichakey{Decisión de equipo}No
\fichakey{Relacionados}\hyperref[err:25]{error~25}, \hyperref[err:27]{error~27}, \hyperref[err:154]{error~154}, \hyperref[nota:23]{nota~23}
\end{ficha}
\paragraph{Qué pasa}
En el menú contextual de un enlace, «Add restrictions» (o «Edit restrictions», si ya tiene) abre la ventana «Link Restriction: Link between \ldots{} and \ldots»: una casilla por cada pareja de un estado del padre (columnas) y un estado del hijo (filas), que se marca como incompatible (0) o compatible (1) al pulsarla. La ventana se ofrece para enlaces entre nodos de azar o de decisión de estados finitos y para enlaces hacia un nodo de utilidad (\texttt{Link\allowbreak{}Restriction\allowbreak{}Validator.\allowbreak{}validate}). Pulsar una casilla, además de cambiar la restricción, edita la tabla del \textbf{padre}, en una casilla que depende de dónde se pulsó. Lo que ve el usuario:

\begin{itemize}
\tightlist
\item
  Si el padre es una decisión sin política impuesta, la ventana no se abre: sale la ventana de error no previsto con \texttt{No\allowbreak{}Such\allowbreak{}Element\allowbreak{}Exception}. Los nodos de decisión no tienen potencial, ni los que se crean en el editor ni los que se leen de un fichero. Pasa, por ejemplo, con «Do test?» → «Result of test» en \texttt{0\_\allowbreak{}miscellaneous/\allowbreak{}networks/\allowbreak{}id/\allowbreak{}ID-decide-test.\allowbreak{}pgmx} y, con «Edit restrictions», en el mismo enlace de \texttt{0\_\allowbreak{}miscellaneous/\allowbreak{}networks/\allowbreak{}dan/\allowbreak{}DAN-decide-test.\allowbreak{}pgmx}. En las redes en formato 0.2 de \texttt{0\_\allowbreak{}miscellaneous/\allowbreak{}networks/\allowbreak{}dan} hay 107 enlaces con restricción que salen de una decisión, repartidos en 31 redes, y no se puede abrir la ventana de ninguno.
\item
  Si el padre tiene tabla, según la casilla pulsada, una probabilidad del padre pasa a 0 o a 1 sin ningún aviso, o sale la ventana de error no previsto con \texttt{Array\allowbreak{}Index\allowbreak{}Out\allowbreak{}Of\allowbreak{}Bounds\allowbreak{}Exception}. En este segundo caso la restricción ya ha cambiado, pero la tabla del hijo no llega a tocarse (lo que sí pasa en las demás casillas es \hyperref[err:25]{error~25}).
\end{itemize}

El cambio de la tabla del padre es una edición normal dentro de la sub-historia de ediciones que abre la ventana: «Cancel», o deshacer después de «OK», lo devuelven. Con «OK» se queda, la red calcula con él y se guarda así.

Medido con un pequeño programa de prueba que construye el panel y la tabla reales de la ventana (\texttt{Link\allowbreak{}Restriction\allowbreak{}Panel}, sin construir la ventana) y pulsa cada casilla como lo hace el panel:

\begin{itemize}
\tightlist
\item
  \texttt{0\_\allowbreak{}miscellaneous/\allowbreak{}networks/\allowbreak{}bn/\allowbreak{}BN-asia.\allowbreak{}pgmx}, VisitToAsia → Tuberculosis. En la fila de «Tuberculosis = yes», con cualquiera de las dos columnas, la tabla de VisitToAsia pasa de \texttt{{[}0.\allowbreak{}99,\allowbreak{}\ 0.\allowbreak{}01{]}} a \texttt{{[}0.\allowbreak{}0,\allowbreak{}\ 1.\allowbreak{}0{]}}; tras «OK», la propagación da P(VisitToAsia = yes) = 1; tras «Cancel», o tras «OK» y deshacer, vuelve a \texttt{{[}0.\allowbreak{}99,\allowbreak{}\ 0.\allowbreak{}01{]}}. En la fila de «Tuberculosis = no», las dos casillas lanzan \texttt{Array\allowbreak{}Index\allowbreak{}Out\allowbreak{}Of\allowbreak{}Bounds\allowbreak{}Exception:\ Index\ -1\ out\ of\ bounds\ for\ length\ 2}.
\item
  La misma red, Tuberculosis → TuberculosisOrCancer. Cada casilla pone a 0 o a 1 una columna de la tabla de Tuberculosis, \texttt{{[}0.\allowbreak{}99,\allowbreak{}\ 0.\allowbreak{}01,\allowbreak{}\ 0.\allowbreak{}95,\allowbreak{}\ 0.\allowbreak{}05{]}} (P(Tuberculosis \textbar{} VisitToAsia)). Por ejemplo, pulsar «Tuberculosis = no, TuberculosisOrCancer = no» la deja en \texttt{{[}0.\allowbreak{}0,\allowbreak{}\ 1.\allowbreak{}0,\allowbreak{}\ 0.\allowbreak{}95,\allowbreak{}\ 0.\allowbreak{}05{]}}: quien no ha estado en Asia tiene tuberculosis con probabilidad 1.
\item
  \texttt{0\_\allowbreak{}miscellaneous/\allowbreak{}networks/\allowbreak{}id/\allowbreak{}ID-decide-test.\allowbreak{}pgmx}, Disease → Health state (un nodo de utilidad): las dos casillas dejan Disease en \texttt{{[}0.\allowbreak{}0,\allowbreak{}\ 1.\allowbreak{}0{]}} en lugar de \texttt{{[}0.\allowbreak{}86,\allowbreak{}\ 0.\allowbreak{}14{]}}.
\item
  La misma red, Result of test → Therapy (el hijo es una decisión): cuatro de las seis casillas no cambian nada; «positive, yes» deja a cero una columna entera de la tabla de Result of test (de \texttt{{[}1.\allowbreak{}0,\allowbreak{}\ 0.\allowbreak{}0,\allowbreak{}\ 0.\allowbreak{}0{]}} a \texttt{{[}0.\allowbreak{}0,\allowbreak{}\ 0.\allowbreak{}0,\allowbreak{}\ 0.\allowbreak{}0{]}}) y «positive, no» cambia otra de \texttt{{[}0.\allowbreak{}0,\allowbreak{}\ 0.\allowbreak{}97,\allowbreak{}\ 0.\allowbreak{}03{]}} a \texttt{{[}0.\allowbreak{}03,\allowbreak{}\ 0.\allowbreak{}97,\allowbreak{}\ 0.\allowbreak{}0{]}}.
\item
  «Do test?» → «Result of test» en \texttt{ID-decide-test.\allowbreak{}pgmx} (sin restricción) y en \texttt{DAN-decide-test.\allowbreak{}pgmx} (con restricción): construir el panel lanza \texttt{No\allowbreak{}Such\allowbreak{}Element\allowbreak{}Exception} en \texttt{Values\allowbreak{}Table}, línea 149.
\end{itemize}

Por qué pasa. \texttt{Link\allowbreak{}Restriction\allowbreak{}Values\allowbreak{}Table} (módulo \texttt{gui}) es una subclase de \texttt{Values\allowbreak{}Table}, la tabla con que se editan los valores de un nodo, y le pasa como nodo el padre del enlace (\texttt{super(\allowbreak{}link.\allowbreak{}get\allowbreak{}From(\allowbreak{}),\allowbreak{}\ table\allowbreak{}Model,\allowbreak{}\ modifiable)}, línea 49):

\begin{itemize}
\tightlist
\item
  El constructor de \texttt{Values\allowbreak{}Table} toma el primer potencial de ese nodo (\texttt{node.\allowbreak{}get\allowbreak{}Potentials(\allowbreak{}).\allowbreak{}get\allowbreak{}First(\allowbreak{})}, línea 149); con un padre sin potencial lanza \texttt{No\allowbreak{}Such\allowbreak{}Element\allowbreak{}Exception}. \texttt{Link\allowbreak{}Restriction\allowbreak{}Panel.\allowbreak{}get\allowbreak{}Link\allowbreak{}Restriction\allowbreak{}Values\allowbreak{}Table} (línea 87) construye la tabla al montar la ventana, así que la ventana no llega a verse.
\item
  \texttt{Link\allowbreak{}Restriction\allowbreak{}Values\allowbreak{}Table.\allowbreak{}set\allowbreak{}Value\allowbreak{}At} (líneas 54-81) ejecuta la edición de la restricción y luego llama a \texttt{super.\allowbreak{}set\allowbreak{}Value\allowbreak{}At(\allowbreak{}new\allowbreak{}Numeric\allowbreak{}Value,\allowbreak{}\ row,\allowbreak{}\ column,\allowbreak{}\ source)} (línea 64). Es la versión de \texttt{Values\allowbreak{}Table} (líneas 170-207), que crea un \texttt{Table\allowbreak{}Potential\allowbreak{}Value\allowbreak{}Edit} sobre \texttt{this.\allowbreak{}node}, el padre, con la fila y la columna de la ventana de restricción. \texttt{Table\allowbreak{}Potential\allowbreak{}Value\allowbreak{}Edit} (líneas 144-157) las traduce a una posición de la tabla del padre como si vinieran de la ventana de la tabla del propio padre (\texttt{Potentials\allowbreak{}Table\allowbreak{}Panel\allowbreak{}Operations.\allowbreak{}get\allowbreak{}Potential\allowbreak{}Index}, líneas 67-75: el comienzo de la columna más la última fila editable del padre menos la fila pulsada). El resultado es una casilla cualquiera del padre, o −1, y \texttt{redistribute\allowbreak{}Probabilities} pone ahí el valor pulsado y reparte la columna, o falla en su línea 207.
\end{itemize}

La casilla de la ventana no necesita esa llamada para actualizarse: la tabla escucha las ediciones de la red (\texttt{Values\allowbreak{}Table}, línea 146) y \texttt{after\allowbreak{}Edit\allowbreak{}Executes} (líneas 83-87) ya pone el valor nuevo en el modelo (por lectura).
\paragraph{El código}
gui/src/main/java/org/openmarkov/gui/component/LinkRestrictionValuesTable.java:48-64

\begin{Shaded}
\begin{Highlighting}[]
\KeywordTok{public} \FunctionTok{LinkRestrictionValuesTable}\OperatorTok{(}\NormalTok{Link}\OperatorTok{\textless{}}\BuiltInTok{Node}\OperatorTok{\textgreater{}}\NormalTok{ link}\OperatorTok{,}\NormalTok{ ValuesTableModel tableModel}\OperatorTok{,} \DataTypeTok{final} \DataTypeTok{boolean}\NormalTok{ modifiable}\OperatorTok{)} \OperatorTok{\{}
    \KeywordTok{super}\OperatorTok{(}\NormalTok{link}\OperatorTok{.}\FunctionTok{getFrom}\OperatorTok{(),}\NormalTok{ tableModel}\OperatorTok{,}\NormalTok{ modifiable}\OperatorTok{);}
    \KeywordTok{this}\OperatorTok{.}\FunctionTok{link} \OperatorTok{=}\NormalTok{ link}\OperatorTok{;}
\NormalTok{    node2 }\OperatorTok{=}\NormalTok{ link}\OperatorTok{.}\FunctionTok{getTo}\OperatorTok{();}
\OperatorTok{\}}

\AttributeTok{@Override} \KeywordTok{public} \DataTypeTok{void} \FunctionTok{setValueAt}\OperatorTok{(}\BuiltInTok{Object}\NormalTok{ newValue}\OperatorTok{,} \DataTypeTok{int}\NormalTok{ row}\OperatorTok{,} \DataTypeTok{int}\NormalTok{ column}\OperatorTok{,} \BuiltInTok{Object}\NormalTok{ source}\OperatorTok{)} \OperatorTok{\{}
    \KeywordTok{...}
\NormalTok{    LinkRestrictionPotentialValueEdit linkPotentialEdit }\OperatorTok{=}
            \KeywordTok{new} \FunctionTok{LinkRestrictionPotentialValueEdit}\OperatorTok{(}\NormalTok{link}\OperatorTok{,}\NormalTok{ newNumericValue}\OperatorTok{,}\NormalTok{ row}\OperatorTok{,}\NormalTok{ column}\OperatorTok{);}
    \ControlFlowTok{try} \OperatorTok{\{}
\NormalTok{        linkPotentialEdit}\OperatorTok{.}\FunctionTok{executeEdit}\OperatorTok{();}
        \KeywordTok{super}\OperatorTok{.}\FunctionTok{setValueAt}\OperatorTok{(}\NormalTok{newNumericValue}\OperatorTok{,}\NormalTok{ row}\OperatorTok{,}\NormalTok{ column}\OperatorTok{,}\NormalTok{ source}\OperatorTok{);}
\end{Highlighting}
\end{Shaded}

gui/src/main/java/org/openmarkov/gui/component/ValuesTable.java:142-149, 192-195

\begin{Shaded}
\begin{Highlighting}[]
\KeywordTok{public} \FunctionTok{ValuesTable}\OperatorTok{(}\BuiltInTok{Node}\NormalTok{ node}\OperatorTok{,}\NormalTok{ ValuesTableModel tableModel}\OperatorTok{,} \DataTypeTok{final} \DataTypeTok{boolean}\NormalTok{ modifiable}\OperatorTok{)} \OperatorTok{\{}
    \KeywordTok{super}\OperatorTok{(}\NormalTok{tableModel}\OperatorTok{,}\NormalTok{ modifiable}\OperatorTok{,} \KeywordTok{false}\OperatorTok{,} \KeywordTok{true}\OperatorTok{);}
    \KeywordTok{this}\OperatorTok{.}\FunctionTok{tableModel} \OperatorTok{=}\NormalTok{ tableModel}\OperatorTok{;}
    \ControlFlowTok{if} \OperatorTok{(}\NormalTok{node }\OperatorTok{!=} \KeywordTok{null}\OperatorTok{)} \OperatorTok{\{}
\NormalTok{        node}\OperatorTok{.}\FunctionTok{getProbNet}\OperatorTok{().}\FunctionTok{getPNESupport}\OperatorTok{().}\FunctionTok{addListener}\OperatorTok{(}\KeywordTok{this}\OperatorTok{);}
        \KeywordTok{this}\OperatorTok{.}\FunctionTok{node} \OperatorTok{=}\NormalTok{ node}\OperatorTok{;}
        \KeywordTok{this}\OperatorTok{.}\FunctionTok{probNet} \OperatorTok{=}\NormalTok{ node}\OperatorTok{.}\FunctionTok{getProbNet}\OperatorTok{();}
        \KeywordTok{this}\OperatorTok{.}\FunctionTok{potential} \OperatorTok{=}\NormalTok{ node}\OperatorTok{.}\FunctionTok{getPotentials}\OperatorTok{().}\FunctionTok{getFirst}\OperatorTok{();}
    \KeywordTok{...}
        \KeywordTok{new} \FunctionTok{TablePotentialValueEdit}\OperatorTok{(}\KeywordTok{this}\OperatorTok{.}\FunctionTok{node}\OperatorTok{,} \OperatorTok{(}\BuiltInTok{Double}\OperatorTok{)}\NormalTok{ newValue}\OperatorTok{,}\NormalTok{ row}\OperatorTok{,}\NormalTok{ column}\OperatorTok{,} \KeywordTok{this}\OperatorTok{.}\FunctionTok{priorityList}\OperatorTok{,}
                                    \KeywordTok{this}\OperatorTok{.}\FunctionTok{getTableModel}\OperatorTok{().}\FunctionTok{getNotEditablePositions}\OperatorTok{())}
                \OperatorTok{.}\FunctionTok{executeEdit}\OperatorTok{();}
\end{Highlighting}
\end{Shaded}
\paragraph{Cómo arreglarlo}
La tabla de la ventana de restricción no debe estar ligada al padre como si enseñara sus valores:

\begin{itemize}
\tightlist
\item
  No pasar el padre al constructor de \texttt{Values\allowbreak{}Table}. Con \texttt{null}, el constructor ya no toma potenciales, pero tampoco registra la tabla como oyente de las ediciones, así que habría que registrarla en la red del enlace para que siga recibiendo \texttt{after\allowbreak{}Edit\allowbreak{}Executes} y \texttt{after\allowbreak{}Undoing\allowbreak{}Edit}. La otra forma es que \texttt{Link\allowbreak{}Restriction\allowbreak{}Values\allowbreak{}Table} herede de una tabla que no sea la de los valores de un nodo.
\item
  Sustituir \texttt{super.\allowbreak{}set\allowbreak{}Value\allowbreak{}At(\allowbreak{}.\allowbreak{}.\allowbreak{}.\allowbreak{})} (línea 64) por poner el valor en el modelo de la tabla, como ya hace \texttt{after\allowbreak{}Edit\allowbreak{}Executes}.
\end{itemize}

\texttt{Link\allowbreak{}Restriction\allowbreak{}Values\allowbreak{}Table} es la única subclase de \texttt{Values\allowbreak{}Table} que recibe un nodo distinto del que enseña. En el mismo método está \hyperref[err:25]{error~25} (la tabla del hijo, cambiada fuera del historial): conviene arreglar los dos a la vez.

Prueba de regresión: construir \texttt{Link\allowbreak{}Restriction\allowbreak{}Panel} sobre «Do test?» → «Result of test» de \texttt{0\_\allowbreak{}miscellaneous/\allowbreak{}networks/\allowbreak{}dan/\allowbreak{}DAN-decide-test.\allowbreak{}pgmx} sin excepción; sobre VisitToAsia → Tuberculosis de \texttt{0\_\allowbreak{}miscellaneous/\allowbreak{}networks/\allowbreak{}bn/\allowbreak{}BN-asia.\allowbreak{}pgmx}, pulsar las cuatro casillas: ninguna debe lanzar excepción y la tabla de VisitToAsia debe seguir en \texttt{{[}0.\allowbreak{}99,\allowbreak{}\ 0.\allowbreak{}01{]}}; y sobre Disease → Health state de \texttt{0\_\allowbreak{}miscellaneous/\allowbreak{}networks/\allowbreak{}id/\allowbreak{}ID-decide-test.\allowbreak{}pgmx}, Disease debe seguir en \texttt{{[}0.\allowbreak{}86,\allowbreak{}\ 0.\allowbreak{}14{]}}.

\setcounter{subsection}{155}
\subsection[La tabla de utilidad enseña 922 337 203,685478 en vez de valores mayores]{La tabla de un nodo de utilidad enseña 922 337 203,685478 en lugar de cualquier valor mayor, en valor absoluto, que unos 922 millones}\label{err:156}
\begin{ficha}
\fichakey{Estado}Presente
\fichakey{Módulo}core, gui
\fichakey{Paquete}org.openmarkov.core.model.network; org.openmarkov.gui.dialog.common
\fichakey{Ficheros}\path{core/src/main/java/org/openmarkov/core/model/network/Util.java:288-303, 318-327} \newline \path{gui/src/main/java/org/openmarkov/gui/dialog/common/TablePotentialPanel.java:521-533} \newline \path{gui/src/main/java/org/openmarkov/gui/dialog/common/TableWithEventsPanel.java:495}
\fichakey{Comprobado}Leyendo el código y ejecutándolo
\fichakey{Esfuerzo}Pequeño (horas)
\fichakey{Decisión de equipo}No
\end{ficha}
\paragraph{Qué pasa}
Camino de llegada: propiedades de un nodo de utilidad → pestaña de la tabla, en una red cuyo valor de utilidad supere unos 922 millones en valor absoluto (un coste en euros de un programa, por ejemplo). Todas esas casillas enseñan 922337203.685478 (o −922337203.685478 si el valor es negativo).

El modelo no cambia: al teclear 2000000000 en una casilla, el potencial guarda 2,0·10⁹ y la casilla lo enseña bien hasta que la tabla se vuelve a construir (por ejemplo, al reabrir el diálogo del nodo), cuando vuelve a salir 922337203.685478.

\texttt{Util.\allowbreak{}round\allowbreak{}With\allowbreak{}Precision(\allowbreak{}double\ x,\allowbreak{}\ int\ num\allowbreak{}Decimals)} (módulo \texttt{core}) redondea calculando \texttt{Math.\allowbreak{}round(\allowbreak{}x\ *\ 10\^{}num\allowbreak{}Decimals)\ /\allowbreak{}\ 10\^{}num\allowbreak{}Decimals}. \texttt{Math.\allowbreak{}round(\allowbreak{}double)} devuelve un \texttt{long} y, cuando el producto pasa de 9 223 372 036 854 775 807, devuelve ese máximo en lugar de desbordar. \texttt{Util.\allowbreak{}round\allowbreak{}And\allowbreak{}Reduce} empieza probando con 10 decimales, así que todo valor cuyo valor absoluto supere unos 922 337 203 (9,22·10¹⁸ / 10¹⁰) sale convertido en ±922 337 203,6854776.

\texttt{Table\allowbreak{}Potential\allowbreak{}Panel.\allowbreak{}set\allowbreak{}Potential\allowbreak{}Data\allowbreak{}In\allowbreak{}Centre\allowbreak{}Area} (líneas 502-514) pasa por \texttt{round\allowbreak{}And\allowbreak{}Reduce(\allowbreak{}valor,\allowbreak{}\ 10⁻¹²,\allowbreak{}\ 10)} todos los valores antes de ponerlos en la tabla. Esto afecta a la tabla de cualquier nodo que admita números grandes; en la práctica, a los nodos de utilidad, cuya tabla es un \texttt{Exact\allowbreak{}Distr\allowbreak{}Potential} que este mismo panel enseña.

\texttt{Table\allowbreak{}With\allowbreak{}Events\allowbreak{}Panel} (línea 495) redondea igual los valores de las tablas con sucesos. Los demás usos de \texttt{round\allowbreak{}And\allowbreak{}Reduce} (\texttt{Potentials\allowbreak{}Table\allowbreak{}Panel\allowbreak{}Operations.\allowbreak{}redistribute\allowbreak{}Probabilities}, \texttt{Event\allowbreak{}Table\allowbreak{}Potential\allowbreak{}Value\allowbreak{}Edit}, \texttt{Augmented\allowbreak{}Potential\allowbreak{}Value\allowbreak{}Edit}) trabajan con probabilidades, que no llegan a esos tamaños.

Medido con un pequeño programa de prueba sobre \texttt{0\_\allowbreak{}miscellaneous/\allowbreak{}networks/\allowbreak{}dan/\allowbreak{}DAN-3-test-problem.\allowbreak{}pgmx}, nodo de utilidad «Quality of life», poniendo en su tabla 1 000 000 000, 1 234 567 890, 10¹⁵ y −1 000 000 000: el panel enseña 922337203.685478 en las tres primeras casillas y −922337203.685478 en la cuarta.
\paragraph{El código}
core/src/main/java/org/openmarkov/core/model/network/Util.java:318-327

\begin{Shaded}
\begin{Highlighting}[]
    \KeywordTok{public} \DataTypeTok{static} \DataTypeTok{double} \FunctionTok{roundWithPrecision}\OperatorTok{(}\DataTypeTok{double}\NormalTok{ x}\OperatorTok{,} \DataTypeTok{int}\NormalTok{ numDecimals}\OperatorTok{)} \OperatorTok{\{}
        \DataTypeTok{double}\NormalTok{ xRounded}\OperatorTok{;}
        \DataTypeTok{double}\NormalTok{ scale }\OperatorTok{=} \BuiltInTok{Math}\OperatorTok{.}\FunctionTok{pow}\OperatorTok{(}\DecValTok{10}\OperatorTok{,}\NormalTok{ numDecimals}\OperatorTok{);}
\NormalTok{        xRounded }\OperatorTok{=} \BuiltInTok{Math}\OperatorTok{.}\FunctionTok{round}\OperatorTok{(}\NormalTok{x }\OperatorTok{*}\NormalTok{ scale}\OperatorTok{);}
\NormalTok{        xRounded }\OperatorTok{=}\NormalTok{ xRounded }\OperatorTok{/}\NormalTok{ scale}\OperatorTok{;}
        \ControlFlowTok{return}\NormalTok{ xRounded}\OperatorTok{;}
    \OperatorTok{\}}
\end{Highlighting}
\end{Shaded}
\paragraph{Cómo arreglarlo}
Arreglarlo en \texttt{Util.\allowbreak{}round\allowbreak{}With\allowbreak{}Precision}, que es donde está el defecto y lo heredan todos sus usos. Dos formas:

\begin{itemize}
\tightlist
\item
  Si \texttt{\textbar{}x\ *\ scale\textbar{}} no cabe en un \texttt{long} (o el valor no tiene parte decimal representable a esa escala), devolver \texttt{x} sin tocar.
\item
  Redondear con \texttt{Big\allowbreak{}Decimal.\allowbreak{}value\allowbreak{}Of(\allowbreak{}x).\allowbreak{}set\allowbreak{}Scale(\allowbreak{}num\allowbreak{}Decimals,\allowbreak{}\ Rounding\allowbreak{}Mode.\allowbreak{}HALF\_\allowbreak{}UP)}, como ya hace \texttt{round\allowbreak{}With\allowbreak{}Significant\allowbreak{}Figures}.
\end{itemize}

Sitios que dependen de ello: \texttt{Table\allowbreak{}Potential\allowbreak{}Panel.\allowbreak{}set\allowbreak{}Potential\allowbreak{}Data\allowbreak{}In\allowbreak{}Centre\allowbreak{}Area} y \texttt{Table\allowbreak{}With\allowbreak{}Events\allowbreak{}Panel} (línea 495).

Prueba: \texttt{round\allowbreak{}And\allowbreak{}Reduce(\allowbreak{}1e9,\allowbreak{}\ 1e-12,\allowbreak{}\ 10)\ ==\ 1e9}, \texttt{round\allowbreak{}And\allowbreak{}Reduce(\allowbreak{}-1e15,\allowbreak{}\ 1e-12,\allowbreak{}\ 10)\ ==\ -1e15}, y que los redondeos pequeños de siempre no cambian (por ejemplo \texttt{round\allowbreak{}And\allowbreak{}Reduce(\allowbreak{}0.\allowbreak{}30000000000000004,\allowbreak{}\ 1e-12,\allowbreak{}\ 10)\ ==\ 0.\allowbreak{}3}).

\setcounter{subsection}{156}
\subsection[El panel de Weibull deja borrar «Gamma» y la relación revienta al calcular]{\textcolor{corregidoverde}{[Corregido en parte]} El panel de un riesgo de Weibull deja borrar la covariable «Gamma», y la relación se queda sin forma y revienta al calcular}\label{err:157}
\begin{ficha}
\fichakey{Estado}Presente en parte. El commit \texttt{522d926} (23-09-2026) hizo que el panel pregunte al potencial que edita por sus covariables obligatorias, así que ya no deja quitar ni reescribir \texttt{Gamma} en un riesgo de Weibull; sigue apagándose el botón de añadir al elegir una fila obligatoria.
\fichakey{Módulo}gui
\fichakey{Paquete}org.openmarkov.gui.dialog.common
\fichakey{Ficheros}\path{gui/src/main/java/org/openmarkov/gui/dialog/common/GLMPanel.java:130-143,158-178} \newline \path{core/src/main/java/org/openmarkov/core/model/network/potential/GLMPotential.java:128-130} \newline \path{core/src/main/java/org/openmarkov/core/model/network/potential/WeibullHazardPotential.java:40-43,122-124,137,176,278-281}
\fichakey{Comprobado}Leyendo el código y ejecutándolo
\fichakey{Esfuerzo}Pequeño (horas)
\fichakey{Decisión de equipo}No
\fichakey{Relacionados}\hyperref[err:14]{error~14}, \hyperref[err:158]{error~158}
\end{ficha}
\paragraph{Qué pasa}
Camino del usuario: propiedades del nodo «Failure» de \texttt{0\_\allowbreak{}miscellaneous/\allowbreak{}networks/\allowbreak{}mid/\allowbreak{}MID-hip-Briggs.\allowbreak{}pgmx} (o de cualquier nodo con «Hazard (Weibull)») → elegir la fila «Gamma» → quitar → aceptar → evaluar la red. Pasa lo siguiente:

\begin{enumerate}
\def\labelenumi{\arabic{enumi}.}
\tightlist
\item
  Con la fila «Gamma» elegida, el botón de quitar está encendido. Si el usuario la quita y acepta, \texttt{Weibull\allowbreak{}Potential\allowbreak{}Panel.\allowbreak{}save\allowbreak{}Changes} crea un Weibull sin «Gamma». Después, \texttt{get\allowbreak{}Gamma\allowbreak{}Index} devuelve −1: \texttt{get\allowbreak{}Gamma(\allowbreak{})} y \texttt{table\allowbreak{}Project} (línea 176, \texttt{Math.\allowbreak{}exp(\allowbreak{}coefficients{[}gamma\allowbreak{}Index{]})}) lanzan \texttt{Array\allowbreak{}Index\allowbreak{}Out\allowbreak{}Of\allowbreak{}Bounds\allowbreak{}Exception} al evaluar la red, y el muestreo lanza \texttt{Invalid\allowbreak{}Argument\allowbreak{}Exception} («no covariate declares the shape parameter gamma»).
\item
  Con una covariable obligatoria elegida («Constant»), además del botón de quitar se apaga el de \textbf{añadir}, que no tiene relación con la fila elegida: añadir una covariable no le hace nada a la que esté seleccionada.
\end{enumerate}

Cada clase de regresión tiene covariables que no pueden faltar. Para \texttt{GLMPotential} es «Constant» (método estático \texttt{get\allowbreak{}Mandatory\allowbreak{}Covariates}). Para \texttt{Weibull\allowbreak{}Hazard\allowbreak{}Potential} («Hazard (Weibull)») son «Gamma» y «Constant» (campo estático \texttt{MANDATORY\_\allowbreak{}COVARIATES}): el coeficiente de «Gamma» es el logaritmo del parámetro de forma de la distribución de Weibull.

\texttt{GLMPanel.\allowbreak{}value\allowbreak{}Changed} decide, al elegir una fila, si se pueden usar los botones, y pregunta siempre a \texttt{GLMPotential.\allowbreak{}get\allowbreak{}Mandatory\allowbreak{}Covariates(\allowbreak{})}, aunque el potencial que edita sea un Weibull.

El oyente de doble clic sobre la columna de covariables (\texttt{Covariates\allowbreak{}Table\allowbreak{}Mouse\allowbreak{}Listener}, que permite reescribir una covariable) usa la misma lista, así que también dejaría reescribir «Gamma»; hoy esa columna no se ve (\hyperref[err:158]{error~158}).

Medido con un pequeño programa de prueba que construye sin pantalla el \texttt{Weibull\allowbreak{}Potential\allowbreak{}Panel} del nodo «Failure {[}1{]}» de \texttt{0\_\allowbreak{}miscellaneous/\allowbreak{}networks/\allowbreak{}mid/\allowbreak{}MID-hip-Briggs.\allowbreak{}pgmx} y elige cada fila:

\begin{Verbatim}[breaklines,breakanywhere,fontsize=\small]
fila 0 Gamma            -> añadir=on  quitar=on  subir=off bajar=on
fila 1 Constant         -> añadir=off quitar=off subir=on  bajar=on
fila 2 {Age at entry}   -> añadir=on  quitar=on  subir=on  bajar=on
getGamma de un Weibull sin Gamma: ArrayIndexOutOfBoundsException: Index -1 out of bounds for length 4
\end{Verbatim}
\paragraph{El código}
gui/src/main/java/org/openmarkov/gui/dialog/common/GLMPanel.java:130-143

\begin{Shaded}
\begin{Highlighting}[]
\AttributeTok{@Override} \KeywordTok{public} \DataTypeTok{void} \FunctionTok{valueChanged}\OperatorTok{(}\BuiltInTok{ListSelectionEvent}\NormalTok{ e}\OperatorTok{)} \OperatorTok{\{}
    \KeywordTok{super}\OperatorTok{.}\FunctionTok{valueChanged}\OperatorTok{(}\NormalTok{e}\OperatorTok{);}
    \DataTypeTok{int}\NormalTok{ row }\OperatorTok{=}\NormalTok{ valuesTable}\OperatorTok{.}\FunctionTok{getSelectedRow}\OperatorTok{();}
    \ControlFlowTok{if} \OperatorTok{(}\NormalTok{row }\OperatorTok{\textgreater{}=} \DecValTok{0} \OperatorTok{\&\&}\NormalTok{ row }\OperatorTok{\textless{}}\NormalTok{ tableModel}\OperatorTok{.}\FunctionTok{getRowCount}\OperatorTok{())} \OperatorTok{\{}
\NormalTok{        VariableExpression covariate }\OperatorTok{=} \OperatorTok{(}\NormalTok{VariableExpression}\OperatorTok{)}\NormalTok{ tableModel}\OperatorTok{.}\FunctionTok{getValueAt}\OperatorTok{(}\NormalTok{row}\OperatorTok{,} \DecValTok{0}\OperatorTok{);}
        \DataTypeTok{boolean}\NormalTok{ isMandatory }\OperatorTok{=} \KeywordTok{false}\OperatorTok{;}
\NormalTok{        VariableExpression}\OperatorTok{[]}\NormalTok{ mandatoryCovariates }\OperatorTok{=}\NormalTok{ GLMPotential}\OperatorTok{.}\FunctionTok{getMandatoryCovariates}\OperatorTok{();}
        \ControlFlowTok{for} \OperatorTok{(}\NormalTok{VariableExpression mandatoryCovariate }\OperatorTok{:}\NormalTok{ mandatoryCovariates}\OperatorTok{)} \OperatorTok{\{}
\NormalTok{            isMandatory }\OperatorTok{|=}\NormalTok{ mandatoryCovariate}\OperatorTok{.}\FunctionTok{asStringExpression}\OperatorTok{().}\FunctionTok{equals}\OperatorTok{(}\NormalTok{covariate}\OperatorTok{.}\FunctionTok{asStringExpression}\OperatorTok{());}
        \OperatorTok{\}}
        \FunctionTok{setEnabledRemoveValue}\OperatorTok{(!}\NormalTok{isMandatory}\OperatorTok{);}
        \FunctionTok{setEnabledAddValue}\OperatorTok{(!}\NormalTok{isMandatory}\OperatorTok{);}
    \OperatorTok{\}}
\OperatorTok{\}}
\end{Highlighting}
\end{Shaded}

core/src/main/java/org/openmarkov/core/model/network/potential/WeibullHazardPotential.java:40-43

\begin{Shaded}
\begin{Highlighting}[]
\KeywordTok{protected} \DataTypeTok{static} \DataTypeTok{final}\NormalTok{ VariableExpression}\OperatorTok{[]}\NormalTok{ MANDATORY\_COVARIATES }\OperatorTok{=} \KeywordTok{new}\NormalTok{ VariableExpression}\OperatorTok{[]\{}
\NormalTok{        VariableExpression}\OperatorTok{.}\FunctionTok{Common}\OperatorTok{.}\FunctionTok{GAMMA}\OperatorTok{,}
\NormalTok{        VariableExpression}\OperatorTok{.}\FunctionTok{Common}\OperatorTok{.}\FunctionTok{CONSTANT}
\OperatorTok{\};}
\end{Highlighting}
\end{Shaded}
\paragraph{Cómo arreglarlo}
Que la lista de covariables obligatorias dependa del potencial que se edita (el campo \texttt{potential} de \texttt{GLMPanel} ya lo tiene): por ejemplo, un método de instancia en \texttt{GLMPotential}, redefinido en \texttt{Weibull\allowbreak{}Hazard\allowbreak{}Potential} y en \texttt{Exponential\allowbreak{}Hazard\allowbreak{}Potential} (esta última hereda de Weibull, pero su única obligatoria es «Constant», línea 32 de su fichero). Un método estático no sirve, porque no se puede redefinir.

Sitios donde la regla tiene que cumplirse: \texttt{GLMPanel.\allowbreak{}value\allowbreak{}Changed} (botón de quitar) y \texttt{GLMPanel.\allowbreak{}Covariates\allowbreak{}Table\allowbreak{}Mouse\allowbreak{}Listener.\allowbreak{}mouse\allowbreak{}Clicked} (reescribir la covariable). Además, quitar la línea \texttt{set\allowbreak{}Enabled\allowbreak{}Add\allowbreak{}Value(\allowbreak{}!is\allowbreak{}Mandatory)}.

Si se quiere que el modelo también se defienda, el constructor de \texttt{Weibull\allowbreak{}Hazard\allowbreak{}Potential} que recibe covariables podría rechazar una lista sin «Gamma»; es una decisión aparte.

Prueba de regresión: construir \texttt{Weibull\allowbreak{}Potential\allowbreak{}Panel} sin pantalla sobre el nodo «Failure {[}1{]}» de \texttt{0\_\allowbreak{}miscellaneous/\allowbreak{}networks/\allowbreak{}mid/\allowbreak{}MID-hip-Briggs.\allowbreak{}pgmx}, elegir la fila «Gamma» y comprobar que quitar está apagado y añadir encendido; con un «Hazard (Exponential)», que la fila «Constant» apaga quitar y ninguna otra lo hace.

\setcounter{subsection}{157}
\subsection[En el panel de una regresión no se ve la columna de las covariables]{En el panel de una regresión la columna de las covariables no se ve: solo aparecen los coeficientes, sin decir de qué son}\label{err:158}
\begin{ficha}
\fichakey{Estado}Presente en parte. El diálogo «Reorder variables» (\texttt{Reorder\allowbreak{}Variables\allowbreak{}Panel}) ya no oculta su primera columna desde el commit \texttt{8055530}; queda el panel de las regresiones.
\fichakey{Módulo}gui
\fichakey{Paquete}org.openmarkov.gui.dialog.common
\fichakey{Ficheros}\path{gui/src/main/java/org/openmarkov/gui/dialog/common/GLMPanel.java:36-42} \newline \path{gui/src/main/java/org/openmarkov/gui/dialog/common/KeyTablePanel.java:132-140,190-197} \newline \path{gui/src/main/java/org/openmarkov/gui/dialog/common/KeyTable.java:176-203} \newline \path{gui/src/main/java/org/openmarkov/gui/component/StickyColumnsTablePane.java:111-114,124-146} \newline \path{gui/src/main/java/org/openmarkov/gui/component/OMTableModel.java:21-43}
\fichakey{Comprobado}Leyendo el código y ejecutándolo
\fichakey{Esfuerzo}Pequeño (horas)
\fichakey{Decisión de equipo}No
\fichakey{Relacionados}\hyperref[err:157]{error~157}, \hyperref[err:165]{error~165}
\end{ficha}
\paragraph{Qué pasa}
Camino del usuario: abrir las propiedades de un nodo cuyo potencial sea de regresión (por ejemplo «Failure» en \texttt{0\_\allowbreak{}miscellaneous/\allowbreak{}networks/\allowbreak{}mid/\allowbreak{}MID-hip-Briggs.\allowbreak{}pgmx}, o «Age at state entry» en \texttt{0\_\allowbreak{}miscellaneous/\allowbreak{}networks/\allowbreak{}mid/\allowbreak{}MID-CHAP-Ryan-Griffin.\allowbreak{}pgmx}) y mirar la pestaña de la relación.

El usuario ve una columna de números sueltos sin saber a qué covariable corresponde cada uno. Además, no puede usar la edición de una covariable con doble clic en su celda (\texttt{Covariates\allowbreak{}Table\allowbreak{}Mouse\allowbreak{}Listener}, que abre la ventanilla de expresiones), porque la columna no está en pantalla.

Los potenciales de regresión ---«Linear combination»/«Linear regression», «Exponential», «Hazard (Weibull)» y «Hazard (Exponential)»--- se editan en el diálogo del nodo con dos paneles, \texttt{GLMPotential\allowbreak{}Panel} y \texttt{Weibull\allowbreak{}Potential\allowbreak{}Panel}. Los dos meten un \texttt{GLMPanel}: una tabla de dos columnas, «Covariate» (la covariable: una variable o una expresión) y «Coefficient» (su coeficiente), con los botones de añadir, quitar, subir y bajar.

\texttt{GLMPanel} hereda de \texttt{Key\allowbreak{}Table\allowbreak{}Panel}, pensado para listas cuya primera columna es una clave interna que no se enseña. Su último argumento, \texttt{first\allowbreak{}Column\allowbreak{}Hidden}, hace que \texttt{Key\allowbreak{}Table.\allowbreak{}default\allowbreak{}Configuration(\allowbreak{})} ponga la primera columna a ancho cero y además quite de la vista la tabla fija de la izquierda (\texttt{remove\allowbreak{}Sticky\allowbreak{}Table\allowbreak{}From\allowbreak{}View}), que es la que contiene esa columna. \texttt{GLMPanel} pasa \texttt{true}, así que la columna que desaparece es precisamente «Covariate». Otros paneles que usan esta clase (\texttt{Piecewise\allowbreak{}Exponential\allowbreak{}Table\allowbreak{}Panel}) declaran una columna extra «Hidden» para que se oculte ésa; \texttt{GLMPanel} no.

Medido con un pequeño programa de prueba que abre \texttt{0\_\allowbreak{}miscellaneous/\allowbreak{}networks/\allowbreak{}mid/\allowbreak{}MID-hip-Briggs.\allowbreak{}pgmx}, construye sin pantalla el \texttt{Weibull\allowbreak{}Potential\allowbreak{}Panel} del nodo «Failure {[}1{]}» y mira su tabla:

\begin{Verbatim}[breaklines,breakanywhere,fontsize=\small]
columnas del modelo: 2: Covariate | Coefficient
tabla fija (columna 0) en la vista: false
tabla visible: 1 columna: Coefficient
filas: Gamma 0.3740968 / Constant -5.490935 / {Age at entry} -0.0367 / {Sex} 0.7685 / {Prosthesis type} -1.344474
\end{Verbatim}
\paragraph{El código}
gui/src/main/java/org/openmarkov/gui/dialog/common/GLMPanel.java:36-37

\begin{Shaded}
\begin{Highlighting}[]
\KeywordTok{public} \FunctionTok{GLMPanel}\OperatorTok{()} \OperatorTok{\{}
    \KeywordTok{super}\OperatorTok{(}\KeywordTok{new} \BuiltInTok{String}\OperatorTok{[]\{}\StringTok{"Covariate"}\OperatorTok{,} \StringTok{"Coefficient"}\OperatorTok{\},} \KeywordTok{new} \BuiltInTok{Object}\OperatorTok{[}\DecValTok{0}\OperatorTok{][}\DecValTok{2}\OperatorTok{],} \KeywordTok{true}\OperatorTok{,} \KeywordTok{true}\OperatorTok{,} \KeywordTok{true}\OperatorTok{,} \KeywordTok{true}\OperatorTok{);}
\end{Highlighting}
\end{Shaded}

gui/src/main/java/org/openmarkov/gui/dialog/common/KeyTable.java:176-199

\begin{Shaded}
\begin{Highlighting}[]
\ControlFlowTok{if} \OperatorTok{(}\KeywordTok{this}\OperatorTok{.}\FunctionTok{firstColumnHidden} \OperatorTok{\&\&} \KeywordTok{this}\OperatorTok{.}\FunctionTok{getColumnCount}\OperatorTok{()} \OperatorTok{\textgreater{}} \DecValTok{0}\OperatorTok{)} \OperatorTok{\{}
    \BuiltInTok{TableColumn}\NormalTok{ column }\OperatorTok{=} \KeywordTok{this}\OperatorTok{.}\FunctionTok{getColumn}\OperatorTok{(}\DecValTok{0}\OperatorTok{);}
\NormalTok{    column}\OperatorTok{.}\FunctionTok{setMinWidth}\OperatorTok{(}\DecValTok{0}\OperatorTok{);}
\NormalTok{    column}\OperatorTok{.}\FunctionTok{setMaxWidth}\OperatorTok{(}\DecValTok{0}\OperatorTok{);}
\NormalTok{    column}\OperatorTok{.}\FunctionTok{setWidth}\OperatorTok{(}\DecValTok{0}\OperatorTok{);}
\NormalTok{    column}\OperatorTok{.}\FunctionTok{setPreferredWidth}\OperatorTok{(}\DecValTok{0}\OperatorTok{);}
\NormalTok{    column}\OperatorTok{.}\FunctionTok{setResizable}\OperatorTok{(}\KeywordTok{false}\OperatorTok{);}
    \KeywordTok{...}
    \KeywordTok{this}\OperatorTok{.}\FunctionTok{removeStickyTableFromView}\OperatorTok{();}
    \KeywordTok{this}\OperatorTok{.}\FunctionTok{getHeaderColumn}\OperatorTok{(}\DecValTok{0}\OperatorTok{).}\FunctionTok{setMinWidth}\OperatorTok{(}\DecValTok{0}\OperatorTok{);}
\end{Highlighting}
\end{Shaded}
\paragraph{Cómo arreglarlo}
Pasar \texttt{false} como \texttt{first\allowbreak{}Column\allowbreak{}Hidden} en el constructor de \texttt{GLMPanel} \textbf{no basta}. \texttt{Key\allowbreak{}Table\allowbreak{}Panel.\allowbreak{}get\allowbreak{}Table\allowbreak{}Model} (línea 206) pasa el mismo indicador a \texttt{OMTable\allowbreak{}Model} como \texttt{first\allowbreak{}Column\allowbreak{}Is\allowbreak{}Header}, y \texttt{OMTable\allowbreak{}Model.\allowbreak{}split} (líneas 26-43), que es lo que usa \texttt{Sticky\allowbreak{}Columns\allowbreak{}Table\allowbreak{}Pane.\allowbreak{}set\allowbreak{}Model} para repartir las columnas entre la tabla fija y la desplazable, solo conserva los títulos de las columnas cuando ese indicador es \texttt{true}. Medido sobre un \texttt{Key\allowbreak{}Table\allowbreak{}Panel} de dos columnas «Covariate»/«Coefficient»:

\begin{Verbatim}[breaklines,breakanywhere,fontsize=\small]
firstColumnHidden=true:  tabla fija en la vista=false  títulos: Covariate | Coefficient
firstColumnHidden=false: tabla fija en la vista=true   títulos: A | A   (se pierden)
\end{Verbatim}

Opciones:

\begin{itemize}
\tightlist
\item
  \begin{enumerate}
  \def\labelenumi{(\alph{enumi})}
  \tightlist
  \item
    Separar en \texttt{Key\allowbreak{}Table\allowbreak{}Panel} los dos significados ---ocultar la primera columna y conservar los títulos---, por ejemplo pasando siempre \texttt{true} al modelo. Hay que revisar entonces a los otros llamadores que pasan \texttt{false} (hoy solo \texttt{Reorder\allowbreak{}Variables\allowbreak{}Panel}, que no enseña cabecera).
  \end{enumerate}
\item
  \begin{enumerate}
  \def\labelenumi{(\alph{enumi})}
  \setcounter{enumi}{1}
  \tightlist
  \item
    Hacer como \texttt{Piecewise\allowbreak{}Exponential\allowbreak{}Table\allowbreak{}Panel} y declarar una columna oculta extra. Esto obliga a desplazar todos los índices de columna que leen \texttt{GLMPanel} (\texttt{get\allowbreak{}Coefficients}, \texttt{get\allowbreak{}Covariates}, \texttt{value\allowbreak{}Changed}, el oyente de doble clic), \texttt{GLMPotential\allowbreak{}Panel.\allowbreak{}action\allowbreak{}Performed} y \texttt{Weibull\allowbreak{}Potential\allowbreak{}Panel.\allowbreak{}item\allowbreak{}State\allowbreak{}Changed}, que leen el modelo por posición.
  \end{enumerate}
\end{itemize}

Al hacerse visible la columna, el doble clic sobre una covariable vuelve a estar disponible, y con él la posibilidad de reescribir una covariable obligatoria: la comprobación de ese oyente usa la misma lista equivocada que describe \hyperref[err:157]{error~157}, así que conviene arreglar los dos juntos.

Prueba de regresión: construir \texttt{GLMPanel} (o \texttt{Weibull\allowbreak{}Potential\allowbreak{}Panel} sobre el nodo «Failure {[}1{]}» de \texttt{0\_\allowbreak{}miscellaneous/\allowbreak{}networks/\allowbreak{}mid/\allowbreak{}MID-hip-Briggs.\allowbreak{}pgmx}) sin pantalla y comprobar que la tabla visible enseña la columna «Covariate» con los nombres de las covariables.

\setcounter{subsection}{158}
\subsection[«Add variables to potential» en una hoja del árbol no añade ningún padre]{«Add variables to potential» en una hoja del árbol no añade nada salvo que se elija la primera variable del árbol, que es la condicionada}\label{err:159}
\begin{ficha}
\fichakey{Estado}Presente
\fichakey{Módulo}gui
\fichakey{Paquete}org.openmarkov.gui.dialog.treeadd
\fichakey{Ficheros}\path{gui/src/main/java/org/openmarkov/gui/dialog/treeadd/TreeADDEditorPanel.java:769-797} \newline \path{gui/src/main/java/org/openmarkov/gui/dialog/treeadd/AddVariablesCheckBoxPanel.java:36-43}
\fichakey{Comprobado}Leyendo el código y ejecutándolo
\fichakey{Esfuerzo}Pequeño (horas)
\fichakey{Decisión de equipo}No
\fichakey{Corregir antes}\hyperref[err:18]{error~18}: a la vez, como mínimo: en cuanto se puedan añadir variables, la del propio nodo seguiría ofrecida y se duplicaría.
\fichakey{Relacionados}\hyperref[err:18]{error~18}, \hyperref[nota:30]{nota~30}
\end{ficha}
\paragraph{Qué pasa}
En el editor de árboles (\texttt{Tree\allowbreak{}ADDEditor\allowbreak{}Panel}, módulo \texttt{gui}), el botón derecho sobre una hoja (el potencial que cuelga de una rama) ofrece «Add variables to potential». Abre \texttt{Add\allowbreak{}Variables\allowbreak{}Dialog}, con una casilla por cada variable que se puede añadir (las devuelve \texttt{Tree\allowbreak{}ADDBranch.\allowbreak{}get\allowbreak{}Addable\allowbreak{}Variables}).

\begin{itemize}
\tightlist
\item
  Si el usuario marca cualquier variable que no sea la condicionada (un padre), no se añade nada. El diálogo se cierra, no hay aviso y la hoja queda igual.
\item
  Si marca la variable condicionada (que el diálogo sí ofrece, ver \hyperref[err:18]{error~18}), es la única que se encuentra y se añade, y la relación de la hoja queda con esa variable repetida.
\end{itemize}

Es decir: la operación nunca sirve para lo que dice (añadir un padre a la relación de la hoja), y la única elección que hace algo produce una relación mal formada.

Al aceptar, \texttt{add\allowbreak{}Variables\allowbreak{}To\allowbreak{}Potential} busca, por cada casilla marcada, la variable con ese nombre entre las del árbol. El \texttt{break} está fuera del \texttt{if}, al final del cuerpo del bucle, así que el bucle compara \textbf{solo con la primera variable} de \texttt{parent\allowbreak{}Tree\allowbreak{}ADD.\allowbreak{}get\allowbreak{}Variables(\allowbreak{})} y sale. En un árbol de probabilidad condicionada la primera variable es siempre la condicionada (la del propio nodo).

Medido con un pequeño programa de prueba sobre \texttt{0\_\allowbreak{}miscellaneous/\allowbreak{}networks/\allowbreak{}id/\allowbreak{}ID-Monty-Hall-spanish.\allowbreak{}pgmx}, nodo «Puerta escogida 2», hoja de la rama «no» (potencial sobre «Puerta escogida 2» y «Puerta escogida 1»). El menú contextual de la hoja ofrece «Add variables to potential»; el diálogo ofrece «Puerta escogida por el presentador» y «Puerta escogida 2». Reproduciendo la búsqueda del método:

\begin{itemize}
\tightlist
\item
  elegir «Puerta escogida por el presentador» añade \texttt{{[}{]}} y la hoja sigue igual;
\item
  elegir «Puerta escogida 2» la añade y la hoja queda sobre \texttt{{[}Puerta\ escogida\ 2,\allowbreak{}\ Puerta\ escogida\ 1,\allowbreak{}\ Puerta\ escogida\ 2{]}}.
\end{itemize}

La búsqueda se reprodujo copiando sus líneas en el programa de prueba, porque el diálogo no se puede abrir sin pantalla; la lista de casillas y el menú sí se obtuvieron de las clases reales.
\paragraph{El código}
gui/src/main/java/org/openmarkov/gui/dialog/treeadd/TreeADDEditorPanel.java:771-787

\begin{Shaded}
\begin{Highlighting}[]
\BuiltInTok{List}\OperatorTok{\textless{}}\NormalTok{Variable}\OperatorTok{\textgreater{}}\NormalTok{ newVariables }\OperatorTok{=} \KeywordTok{new} \BuiltInTok{ArrayList}\OperatorTok{\textless{}}\NormalTok{Variable}\OperatorTok{\textgreater{}();}
\ControlFlowTok{for} \OperatorTok{(}\BuiltInTok{JCheckBox}\NormalTok{ checkBox }\OperatorTok{:}\NormalTok{ checkBoxes}\OperatorTok{)} \OperatorTok{\{}
    \ControlFlowTok{if} \OperatorTok{(}\NormalTok{checkBox}\OperatorTok{.}\FunctionTok{isSelected}\OperatorTok{())} \OperatorTok{\{}
        \BuiltInTok{String}\NormalTok{ variableName }\OperatorTok{=}\NormalTok{ checkBox}\OperatorTok{.}\FunctionTok{getText}\OperatorTok{();}
        \ControlFlowTok{for} \OperatorTok{(}\NormalTok{Variable variable }\OperatorTok{:}\NormalTok{ parentTreeADD}\OperatorTok{.}\FunctionTok{getVariables}\OperatorTok{())} \OperatorTok{\{}
            \ControlFlowTok{if} \OperatorTok{(}\NormalTok{variable}\OperatorTok{.}\FunctionTok{getName}\OperatorTok{().}\FunctionTok{equals}\OperatorTok{(}\NormalTok{variableName}\OperatorTok{))} \OperatorTok{\{}
\NormalTok{                newVariables}\OperatorTok{.}\FunctionTok{add}\OperatorTok{(}\NormalTok{variable}\OperatorTok{);}
            \OperatorTok{\}}
            \ControlFlowTok{break}\OperatorTok{;}
        \OperatorTok{\}}
    \OperatorTok{\}}
\OperatorTok{\}}
\NormalTok{Potential potential }\OperatorTok{=}\NormalTok{ branch}\OperatorTok{.}\FunctionTok{getPotential}\OperatorTok{();}
\ControlFlowTok{for} \OperatorTok{(}\NormalTok{Variable newVariable }\OperatorTok{:}\NormalTok{ newVariables}\OperatorTok{)} \OperatorTok{\{}
\NormalTok{    potential }\OperatorTok{=}\NormalTok{ potential}\OperatorTok{.}\FunctionTok{addVariable}\OperatorTok{(}\NormalTok{newVariable}\OperatorTok{);}
\OperatorTok{\}}
\end{Highlighting}
\end{Shaded}
\paragraph{Cómo arreglarlo}
Meter el \texttt{break} dentro del \texttt{if} (o quitarlo y buscar con un \texttt{stream(\allowbreak{}).\allowbreak{}filter(\allowbreak{}.\allowbreak{}.\allowbreak{}.\allowbreak{}).\allowbreak{}find\allowbreak{}First(\allowbreak{})}). Mejor aún, asociar a cada casilla su objeto \texttt{Variable} y no buscar por nombre; las búsquedas por nombre de este editor se describen en \hyperref[nota:30]{nota~30}, donde no hay nada que corregir.

Este arreglo debe ir junto con el de \hyperref[err:18]{error~18}: en cuanto las demás variables se puedan añadir, la variable condicionada seguirá ofrecida y seguirá duplicándose si no se corrige \texttt{get\allowbreak{}Addable\allowbreak{}Variables}.

Prueba: sobre la hoja «no» de \texttt{0\_\allowbreak{}miscellaneous/\allowbreak{}networks/\allowbreak{}id/\allowbreak{}ID-Monty-Hall-spanish.\allowbreak{}pgmx}, marcar «Puerta escogida por el presentador» y comprobar que la hoja pasa a tener tres variables distintas, con la tabla del tamaño correcto (27 casillas, 3×3×3) y los valores originales replicados.

\setcounter{subsection}{159}
\subsection[El árbol ofrece repetir una variable discretizada y elegirla da un error]{Bajo una rama de una variable discretizada el editor vuelve a ofrecer esa variable para un subárbol, y elegirla lanza NullPointerException}\label{err:160}
\begin{ficha}
\fichakey{Estado}Presente
\fichakey{Módulo}gui
\fichakey{Paquete}org.openmarkov.gui.dialog.treeadd
\fichakey{Ficheros}\path{gui/src/main/java/org/openmarkov/gui/dialog/treeadd/TreeADDEditorPanel.java:301-306, 352-357, 1045-1071}
\fichakey{Comprobado}Leyendo el código y ejecutándolo
\fichakey{Esfuerzo}Pequeño (horas)
\fichakey{Decisión de equipo}No
\fichakey{Relacionados}\hyperref[nota:17]{nota~17}
\end{ficha}
\paragraph{Qué pasa}
Se llega desde el diálogo de propiedades de un nodo cuyo potencial es un árbol que ramifica por una variable discretizada (variable numérica troceada en intervalos con nombre): botón derecho sobre una rama de un subárbol colgado de una rama de esa variable que se ha quedado con un solo estado. El menú ofrece «Add subtree» con esa misma variable, y al elegirlo sale el diálogo de «error inesperado» que pide enviar una captura: la excepción no controlada (\texttt{Null\allowbreak{}Pointer\allowbreak{}Exception}) llega al manejador global (\texttt{OMException\allowbreak{}Handler}).

En el menú de una hoja, la misma condición hace que se ofrezca añadir esa variable a la relación de la hoja, lo que por \hyperref[err:159]{error~159} no tiene efecto salvo en el caso de la variable condicionada.

En un árbol, una rama que ya ha preguntado por una variable y se ha quedado con \textbf{un solo estado} suyo no gana nada volviendo a preguntar por ella más abajo. El editor de árboles (\texttt{Tree\allowbreak{}ADDEditor\allowbreak{}Panel}, módulo \texttt{gui}) lo tiene en cuenta: al montar el menú contextual sube por el camino de ramas y quita de las opciones la variable de cada rama de un solo estado. Lo hace en dos sitios:

\begin{itemize}
\tightlist
\item
  \texttt{possible\allowbreak{}Root\allowbreak{}Variables} (líneas 351-356), usado por el menú de una rama, opción «Add subtree»;
\item
  \texttt{set\allowbreak{}Contextual\allowbreak{}Menu\allowbreak{}Potential} (líneas 300-305), menú de una hoja, opción «Add variables to potential».
\end{itemize}

En los dos, la condición compara el tipo de variable con \texttt{FINITE\_\allowbreak{}STATES} dos veces; la segunda mitad debía ser \texttt{DISCRETIZED}, que también ramifica por estados, como hacen \texttt{set\allowbreak{}Contextual\allowbreak{}Menu\allowbreak{}Branch} en la línea 240 y \texttt{add\allowbreak{}Subtree} en la línea 1043.

Al elegir «Add subtree» con la variable ya usada, \texttt{add\allowbreak{}Subtree} lo detecta (\texttt{is\allowbreak{}Root\allowbreak{}Variable\allowbreak{}Used\allowbreak{}Before}), busca hacia arriba una rama de esa variable con más de un estado para restringir los estados del subárbol nuevo, no la encuentra (todas tienen uno) y deja \texttt{grouped\allowbreak{}States} a \texttt{null}; la línea 1063 (\texttt{grouped\allowbreak{}States.\allowbreak{}size(\allowbreak{})}) lanza \texttt{Null\allowbreak{}Pointer\allowbreak{}Exception}.

Medido con un pequeño programa de prueba sobre dos redes gemelas construidas a mano: nodo X con padres A (de estados) y D, árbol de X que pregunta primero por D, subárbol de A colgado de la rama «D = high» (un solo estado).

\begin{itemize}
\tightlist
\item
  Con D \textbf{de estados finitos}: bajo la rama «A = \ldots» el menú no ofrece «Add subtree» con D. Correcto.
\item
  Con D \textbf{discretizada}: el menú de esa rama ofrece «Add subtree: D», y al elegirlo \texttt{add\allowbreak{}Subtree} muere con \texttt{Null\allowbreak{}Pointer\allowbreak{}Exception}.
\end{itemize}
\paragraph{El código}
gui/src/main/java/org/openmarkov/gui/dialog/treeadd/TreeADDEditorPanel.java:347-359

\begin{Shaded}
\begin{Highlighting}[]
\BuiltInTok{TreePath}\NormalTok{ parentPath }\OperatorTok{=}\NormalTok{ branchPath}\OperatorTok{.}\FunctionTok{getParentPath}\OperatorTok{();} \CommentTok{// treeADD}
\ControlFlowTok{while} \OperatorTok{(}\NormalTok{parentPath}\OperatorTok{.}\FunctionTok{getLastPathComponent}\OperatorTok{()} \OperatorTok{!=}\NormalTok{ rootTreeADDPotential}\OperatorTok{)} \OperatorTok{\{}
    \BuiltInTok{TreePath}\NormalTok{ grandParentPath }\OperatorTok{=}\NormalTok{ parentPath}\OperatorTok{.}\FunctionTok{getParentPath}\OperatorTok{();}\CommentTok{// branch}
    \ControlFlowTok{if} \OperatorTok{(}\NormalTok{grandParentPath}\OperatorTok{.}\FunctionTok{getLastPathComponent}\OperatorTok{()} \KeywordTok{instanceof}\NormalTok{ TreeADDBranch treeBranch}\OperatorTok{)} \OperatorTok{\{}
        \ControlFlowTok{if} \OperatorTok{((}
\NormalTok{                treeBranch}\OperatorTok{.}\FunctionTok{getRootVariable}\OperatorTok{().}\FunctionTok{getVariableType}\OperatorTok{()} \OperatorTok{==}\NormalTok{ VariableType}\OperatorTok{.}\FunctionTok{FINITE\_STATES}
                        \OperatorTok{||}\NormalTok{ treeBranch}\OperatorTok{.}\FunctionTok{getRootVariable}\OperatorTok{().}\FunctionTok{getVariableType}\OperatorTok{()} \OperatorTok{==}\NormalTok{ VariableType}\OperatorTok{.}\FunctionTok{FINITE\_STATES}
        \OperatorTok{)} \OperatorTok{\&\&}\NormalTok{ treeBranch}\OperatorTok{.}\FunctionTok{getBranchStates}\OperatorTok{().}\FunctionTok{size}\OperatorTok{()} \OperatorTok{==} \DecValTok{1}\OperatorTok{)} \OperatorTok{\{}
\NormalTok{            possibleRootVariables}\OperatorTok{.}\FunctionTok{remove}\OperatorTok{(}\NormalTok{treeBranch}\OperatorTok{.}\FunctionTok{getRootVariable}\OperatorTok{());}
        \OperatorTok{\}}
    \OperatorTok{\}}
\NormalTok{    parentPath }\OperatorTok{=}\NormalTok{ grandParentPath}\OperatorTok{;}
\OperatorTok{\}}
\end{Highlighting}
\end{Shaded}
\paragraph{Cómo arreglarlo}
Cambiar la segunda mitad por \texttt{Variable\allowbreak{}Type.\allowbreak{}DISCRETIZED} en los dos sitios (líneas 302 y 353); mejor, sacar el recorrido a un método único que usen los dos menús, porque es la misma regla escrita dos veces.

Además, \texttt{add\allowbreak{}Subtree} (líneas 1045-1066) no debería depender de que el menú haya filtrado bien: si la variable ya se usó y no hay rama agrupada, no debe usar \texttt{grouped\allowbreak{}States} nulo (rechazar con un mensaje o crear el subárbol completo, según lo que se quiera).

Prueba: los dos árboles gemelos descritos; comprobar que con D discretizada el menú de la rama interior no ofrece D, y que \texttt{add\allowbreak{}Subtree} invocado con D en esa situación no lanza \texttt{Null\allowbreak{}Pointer\allowbreak{}Exception}.

\setcounter{subsection}{160}
\subsection[«Split interval» con el campo vacío o sin un número da un error inesperado]{En el editor de árboles, aceptar «Split interval» con el campo vacío (tal como se abre) o con un texto que no es un número sale como error inesperado, y lo mismo en «Change interval»}\label{err:161}
\begin{ficha}
\fichakey{Estado}Presente
\fichakey{Módulo}gui
\fichakey{Paquete}org.openmarkov.gui.dialog.treeadd
\fichakey{Ficheros}\path{gui/src/main/java/org/openmarkov/gui/dialog/treeadd/TreeADDEditorPanel.java:411-425, 450-455, 636-648} \newline \path{gui/src/main/java/org/openmarkov/gui/dialog/treeadd/SplitIntervalDialog.java:79-81} \newline \path{gui/src/main/java/org/openmarkov/gui/dialog/treeadd/SplitIntervalPanel.java:31-32} \newline \path{gui/src/main/java/org/openmarkov/gui/dialog/treeadd/ChangeIntervalDialog.java:87-89} \newline \path{gui/src/main/java/org/openmarkov/gui/dialog/treeadd/ChangeIntervalPanel.java:41-45}
\fichakey{Comprobado}Leyendo el código y ejecutándolo
\fichakey{Esfuerzo}Pequeño (horas)
\fichakey{Decisión de equipo}No
\fichakey{Relacionados}\hyperref[err:164]{error~164}, \hyperref[err:181]{error~181}, \hyperref[err:182]{error~182}, \hyperref[nota:17]{nota~17}
\end{ficha}
\paragraph{Qué pasa}
En el editor de una relación de tipo árbol (Tree/ADD), cada rama de una variable numérica tiene en su menú contextual «Split interval» (partir el intervalo de la rama por un umbral) y, si la rama no cubre todo el dominio, «Change interval». Las dos abren una ventanilla con campos de texto.

«Split interval» se abre con el campo del umbral vacío. Pulsar «OK» (o Intro) sin escribir nada, o después de escribir algo que no se lee como número ---una letra, o un decimal con coma como «1,5»---, termina en el diálogo de error inesperado que pide enviar el fallo a los desarrolladores. Lo mismo ocurre en «Change interval» si se vacía o se estropea uno de sus dos límites. El árbol no cambia, porque la conversión es lo primero que se hace tras cerrar la ventanilla.

Camino del usuario: abrir \texttt{0\_\allowbreak{}miscellaneous/\allowbreak{}networks/\allowbreak{}mid/\allowbreak{}MID-Chancellor.\allowbreak{}pgmx}, abrir la relación de «State {[}1{]}» o de «Cost lamivudine {[}0{]}» (dos árboles con ramas de «Time in treatment {[}0{]}», una variable numérica; la primera de esas ramas va de 0 a 2), clic derecho en una rama numérica, «Split interval» y «OK».

Por qué pasa: \texttt{Tree\allowbreak{}ADDEditor\allowbreak{}Panel.\allowbreak{}split\allowbreak{}Interval} (módulo \texttt{gui}, línea 643) convierte el texto con \texttt{Float.\allowbreak{}parse\allowbreak{}Float}, y \texttt{change\allowbreak{}Interval} (líneas 449-450) con \texttt{Double.\allowbreak{}parse\allowbreak{}Double}. Ninguno captura \texttt{Number\allowbreak{}Format\allowbreak{}Exception}. \texttt{action\allowbreak{}Performed} solo envuelve en \texttt{Unrecoverable\allowbreak{}Exception} ---que la ventana enseña con su mensaje, como se ve en \hyperref[nota:17]{nota~17}--- las excepciones propias del árbol; \texttt{Number\allowbreak{}Format\allowbreak{}Exception} es de tiempo de ejecución y llega al manejador global como error no previsto. Las ventanillas no validan: \texttt{Split\allowbreak{}Interval\allowbreak{}Dialog.\allowbreak{}do\allowbreak{}Ok\allowbreak{}Click\allowbreak{}Before\allowbreak{}Hide} solo comprueba \texttt{get\allowbreak{}Text(\allowbreak{})\ !=\ null}, que se cumple siempre, y \texttt{Change\allowbreak{}Interval\allowbreak{}Dialog.\allowbreak{}do\allowbreak{}Ok\allowbreak{}Click\allowbreak{}Before\allowbreak{}Hide} devuelve \texttt{true}.

Medido ejecutando por separado lo que no depende de la ventana:

\begin{itemize}
\tightlist
\item
  \texttt{Split\allowbreak{}Interval\allowbreak{}Panel} real: el campo del umbral empieza vacío;
\item
  el \texttt{do\allowbreak{}Ok\allowbreak{}Click\allowbreak{}Before\allowbreak{}Hide} real de \texttt{Split\allowbreak{}Interval\allowbreak{}Dialog} devuelve \texttt{true} (la ventanilla se cerraría) con «», «abc», «1,5» y «2.5»;
\item
  \texttt{Float.\allowbreak{}parse\allowbreak{}Float} lanza \texttt{Number\allowbreak{}Format\allowbreak{}Exception} con los tres primeros («empty String», \texttt{For\ input\ string:\ "abc"}, \texttt{For\ input\ string:\ "1,\allowbreak{}5"}) y acepta «2.5»;
\item
  en la rama de 0 a 2 de «Time in treatment {[}0{]}», \texttt{Change\allowbreak{}Interval\allowbreak{}Panel} rellena los límites con «0.0» (no editable, es el borde del dominio) y «2.0» (editable); \texttt{Double.\allowbreak{}parse\allowbreak{}Double} con el campo vaciado lanza la misma excepción.
\end{itemize}
\paragraph{El código}
gui/src/main/java/org/openmarkov/gui/dialog/treeadd/TreeADDEditorPanel.java:406-420

\begin{Shaded}
\begin{Highlighting}[]
\ControlFlowTok{case}\NormalTok{ ActionCommands}\OperatorTok{.}\FunctionTok{SPLIT\_INTERVAL} \OperatorTok{{-}\textgreater{}} \OperatorTok{\{}
    \CommentTok{// node must be a branch}
    \ControlFlowTok{try} \OperatorTok{\{}
        \FunctionTok{splitInterval}\OperatorTok{(}\NormalTok{ae}\OperatorTok{,} \OperatorTok{(}\NormalTok{TreeADDBranch}\OperatorTok{)}\NormalTok{ node}\OperatorTok{,}\NormalTok{ path}\OperatorTok{);}
    \OperatorTok{\}} \ControlFlowTok{catch} \OperatorTok{(}\NormalTok{InvalidLimitInTreeADDException }\OperatorTok{|}\NormalTok{ TriedToSplitIntervalOutsideBoundsException e}\OperatorTok{)} \OperatorTok{\{}
        \ControlFlowTok{throw} \KeywordTok{new} \FunctionTok{UnrecoverableException}\OperatorTok{(}\NormalTok{e}\OperatorTok{);}
    \OperatorTok{\}}
\OperatorTok{\}}
\ControlFlowTok{case}\NormalTok{ ActionCommands}\OperatorTok{.}\FunctionTok{CHANGE\_INTERVAL} \OperatorTok{{-}\textgreater{}} \OperatorTok{\{}
    \ControlFlowTok{try} \OperatorTok{\{}
        \FunctionTok{changeInterval}\OperatorTok{(}\NormalTok{ae}\OperatorTok{,} \OperatorTok{(}\NormalTok{TreeADDBranch}\OperatorTok{)}\NormalTok{ node}\OperatorTok{,}\NormalTok{ path}\OperatorTok{);}
    \OperatorTok{\}} \ControlFlowTok{catch} \OperatorTok{(}\NormalTok{InvalidArgumentException }\OperatorTok{|}\NormalTok{ ChangeDomainOfTreeADDIsNotAllowedException e}\OperatorTok{)} \OperatorTok{\{}
        \ControlFlowTok{throw} \KeywordTok{new} \FunctionTok{UnrecoverableException}\OperatorTok{(}\NormalTok{e}\OperatorTok{);}
    \OperatorTok{\}}
\OperatorTok{\}}
\end{Highlighting}
\end{Shaded}

gui/src/main/java/org/openmarkov/gui/dialog/treeadd/TreeADDEditorPanel.java:643

\begin{Shaded}
\begin{Highlighting}[]
\DataTypeTok{float}\NormalTok{ introducedLimit }\OperatorTok{=} \BuiltInTok{Float}\OperatorTok{.}\FunctionTok{parseFloat}\OperatorTok{(}\NormalTok{panel}\OperatorTok{.}\FunctionTok{getLimit}\OperatorTok{().}\FunctionTok{getText}\OperatorTok{());}
\end{Highlighting}
\end{Shaded}

gui/src/main/java/org/openmarkov/gui/dialog/treeadd/SplitIntervalDialog.java:79-81

\begin{Shaded}
\begin{Highlighting}[]
\AttributeTok{@Override} \KeywordTok{protected} \DataTypeTok{boolean} \FunctionTok{doOkClickBeforeHide}\OperatorTok{()} \OperatorTok{\{}
    \ControlFlowTok{return} \FunctionTok{getJPanelSplitInterval}\OperatorTok{().}\FunctionTok{getLimit}\OperatorTok{().}\FunctionTok{getText}\OperatorTok{()} \OperatorTok{!=} \KeywordTok{null}\OperatorTok{;}
\OperatorTok{\}}
\end{Highlighting}
\end{Shaded}
\paragraph{Cómo arreglarlo}
Validar en la ventanilla antes de cerrarla: que \texttt{do\allowbreak{}Ok\allowbreak{}Click\allowbreak{}Before\allowbreak{}Hide} de \texttt{Split\allowbreak{}Interval\allowbreak{}Dialog} y de \texttt{Change\allowbreak{}Interval\allowbreak{}Dialog} intente convertir sus textos y, si alguno no es un número, avise y devuelva \texttt{false}, para que el usuario corrija sin perder la ventanilla. La otra opción es capturar \texttt{Number\allowbreak{}Format\allowbreak{}Exception} en \texttt{split\allowbreak{}Interval} y \texttt{change\allowbreak{}Interval} y lanzar una excepción del árbol con mensaje, que \texttt{action\allowbreak{}Performed} ya envuelve en \texttt{Unrecoverable\allowbreak{}Exception}.

Las tres conversiones están en \texttt{Tree\allowbreak{}ADDEditor\allowbreak{}Panel}: línea 643 (\texttt{split\allowbreak{}Interval}) y líneas 449-450 (\texttt{change\allowbreak{}Interval}). Si se quiere aceptar la coma decimal, hay que hacerlo en las tres.

Prueba de regresión: con \texttt{Split\allowbreak{}Interval\allowbreak{}Panel} con el texto vacío y con «abc», aceptar no debe cerrar la ventanilla (o no debe salir ninguna excepción que no sea del árbol); con «1» en la rama de 0 a 2 de «Time in treatment {[}0{]}», la rama debe partirse en dos.

\setcounter{subsection}{161}
\subsection[La pestaña «Parents» ofrece padres que las restricciones rechazan]{La pestaña «Parents» del diálogo de un nodo ofrece como posibles padres nodos que las restricciones de la red rechazan, y elegir uno termina en un diálogo de error técnico}\label{err:162}
\begin{ficha}
\fichakey{Estado}Presente
\fichakey{Módulo}gui
\fichakey{Paquete}org.openmarkov.gui.dialog.common
\fichakey{Ficheros}\path{gui/src/main/java/org/openmarkov/gui/dialog/common/PrefixedDataTablePanel.java:165-190} \newline \path{gui/src/main/java/org/openmarkov/gui/dialog/common/PrefixedDataTablePanel.java:115-143} \newline \path{gui/src/main/java/org/openmarkov/gui/dialog/common/KeyTablePanel.java:362-380} \newline \path{core/src/main/java/org/openmarkov/core/action/base/linkEdits/AddLinkEdit.java:69-187} \newline \path{core/src/main/java/org/openmarkov/core/action/base/PNEdit.java:37-56} \newline \path{gui/src/main/java/org/openmarkov/gui/dialog/common/KeyListSelectionDialog.java:139}
\fichakey{Comprobado}Leyendo el código y ejecutándolo
\fichakey{Esfuerzo}Pequeño (horas)
\fichakey{Decisión de equipo}No
\fichakey{Relacionados}\hyperref[err:23]{error~23}, \hyperref[nota:17]{nota~17}, \hyperref[nota:24]{nota~24}, \hyperref[nota:25]{nota~25}, \hyperref[nota:153]{nota~153}
\end{ficha}
\paragraph{Qué pasa}
El diálogo de propiedades de un nodo (\texttt{Node\allowbreak{}Properties\allowbreak{}Dialog}, módulo \texttt{gui}) tiene una pestaña «Parents» (\texttt{Node\allowbreak{}Parents\allowbreak{}Panel}) con una tabla de padres (\texttt{Prefixed\allowbreak{}Data\allowbreak{}Table\allowbreak{}Panel}) y un botón para añadir. El botón abre una lista con los candidatos, y en ella aparecen nodos que la red no permite enlazar: por ejemplo, a «VisitToAsia» de \texttt{0\_\allowbreak{}miscellaneous/\allowbreak{}networks/\allowbreak{}bn/\allowbreak{}BN-asia.\allowbreak{}pgmx} se le ofrecen «X-ray», «Dyspnea», «Tuberculosis» y «TuberculosisOrCancer», que son descendientes suyos y cerrarían un ciclo.

Si el usuario elige uno, el programa no se cae: sale un diálogo de error con un texto poco legible y el enlace no se añade. Eligiendo «X-ray» para «VisitToAsia», el diálogo tiene el título «There is a cycle» y el mensaje «Could not met the constraint: No cycles paths are allowed. There is a cycle traversing between chance node of Finite states variable X-ray with states {[}state no, state yes{]} and chance node of Finite states variable VisitToAsia with states {[}state no, state yes{]}.»

Además, la lista admite selección múltiple: si se eligen varias filas y una falla, las anteriores ya se han añadido, y la lista de candidatos (\texttt{absent\allowbreak{}Data}) no se recalcula antes de salir por la excepción.

La lista es un \texttt{Key\allowbreak{}List\allowbreak{}Selection\allowbreak{}Dialog} con los candidatos que calcula \texttt{absent\allowbreak{}Prefixed\allowbreak{}Data} (líneas 165-190): todos los nodos de la red que no son ya padres y no son el propio nodo. Por cada uno prepara una edición \texttt{Add\allowbreak{}Link\allowbreak{}Edit} sin ejecutarla. No consulta ninguna restricción de la red; la de ciclos (\texttt{No\allowbreak{}Cycle}), por ejemplo, solo actúa al ejecutar la edición.

Al aceptar, \texttt{action\allowbreak{}Performed\allowbreak{}Add\allowbreak{}Value} (líneas 115-143) ejecuta la edición de cada fila elegida. \texttt{PNEdit.\allowbreak{}execute\allowbreak{}Edit} comprueba entonces las restricciones (\texttt{Add\allowbreak{}Link\allowbreak{}Edit.\allowbreak{}check\allowbreak{}Constraints\allowbreak{}Will\allowbreak{}Be\allowbreak{}Met}, líneas 69-187: enlace repetido, máximo de padres, enlace hacia un corte temporal anterior, ciclo, bucle, padres mezclados, padres de utilidad\ldots) y, si alguna se incumple, lanza \texttt{Do\allowbreak{}Edit\allowbreak{}Exception.\allowbreak{}Cannot\allowbreak{}Do\allowbreak{}Edit\allowbreak{}Exception}. \texttt{Key\allowbreak{}Table\allowbreak{}Panel.\allowbreak{}action\allowbreak{}Performed} (líneas 362-380) la envuelve en \texttt{Unrecoverable\allowbreak{}Exception}, que el manejador global convierte en un diálogo de error (no cierra el programa).

Medido con un pequeño programa de prueba que construye el panel real para cada nodo y pregunta a cada edición preparada si cumpliría las restricciones (\texttt{PNEdit.\allowbreak{}try\allowbreak{}Constraints\allowbreak{}Will\allowbreak{}Be\allowbreak{}Met}):

\begin{itemize}
\tightlist
\item
  \texttt{0\_\allowbreak{}miscellaneous/\allowbreak{}networks/\allowbreak{}bn/\allowbreak{}BN-asia.\allowbreak{}pgmx}: se ofrecen 48 candidatos en total y 18 los rechaza la restricción de ciclos (todos descendientes).
\item
  \texttt{0\_\allowbreak{}miscellaneous/\allowbreak{}networks/\allowbreak{}id/\allowbreak{}ID-decide-test.\allowbreak{}pgmx}: 34 ofrecidos, 17 rechazados (ciclo, padre de utilidad no permitido y combinaciones).
\item
  \texttt{0\_\allowbreak{}miscellaneous/\allowbreak{}networks/\allowbreak{}mid/\allowbreak{}MID-Chancellor-corrected.\allowbreak{}pgmx}: 120 ofrecidos, 52 rechazados (incluye enlaces a un corte temporal anterior).
\end{itemize}

El diálogo de error citado arriba es el que el manejador global enseñaría al ejecutar la edición «X-ray → VisitToAsia».
\paragraph{El código}
gui/src/main/java/org/openmarkov/gui/dialog/common/PrefixedDataTablePanel.java:165-181

\begin{Shaded}
\begin{Highlighting}[]
\KeywordTok{private} \BuiltInTok{Object}\OperatorTok{[][]} \FunctionTok{absentPrefixedData}\OperatorTok{()} \OperatorTok{\{}
    \BuiltInTok{List}\OperatorTok{\textless{}}\BuiltInTok{Node}\OperatorTok{\textgreater{}}\NormalTok{ allNodes }\OperatorTok{=}\NormalTok{ node}\OperatorTok{.}\FunctionTok{getProbNet}\OperatorTok{().}\FunctionTok{getNodes}\OperatorTok{();}
    \BuiltInTok{List}\OperatorTok{\textless{}}\BuiltInTok{Node}\OperatorTok{\textgreater{}}\NormalTok{ nodes }\OperatorTok{=} \KeywordTok{new} \BuiltInTok{ArrayList}\OperatorTok{\textless{}}\BuiltInTok{Node}\OperatorTok{\textgreater{}();}
\NormalTok{    edits}\OperatorTok{.}\FunctionTok{clear}\OperatorTok{();}
    
    \ControlFlowTok{for} \OperatorTok{(}\BuiltInTok{Node}\NormalTok{ otherNode }\OperatorTok{:}\NormalTok{ allNodes}\OperatorTok{)} \OperatorTok{\{}
        \ControlFlowTok{if} \OperatorTok{(!}\NormalTok{node}\OperatorTok{.}\FunctionTok{getParents}\OperatorTok{().}\FunctionTok{contains}\OperatorTok{(}\NormalTok{otherNode}\OperatorTok{)} \OperatorTok{\&\&}\NormalTok{ otherNode }\OperatorTok{!=}\NormalTok{ node}\OperatorTok{)} \OperatorTok{\{}
\NormalTok{            AddLinkEdit linkEdit }\OperatorTok{=} \KeywordTok{new} \FunctionTok{AddLinkEdit}\OperatorTok{(}\NormalTok{node}\OperatorTok{.}\FunctionTok{getProbNet}\OperatorTok{(),}\NormalTok{ otherNode}\OperatorTok{.}\FunctionTok{getVariable}\OperatorTok{(),}\NormalTok{ node}\OperatorTok{.}\FunctionTok{getVariable}\OperatorTok{(),}
                                                   \KeywordTok{true}\OperatorTok{);}
            \KeywordTok{...}
\NormalTok{            edits}\OperatorTok{.}\FunctionTok{add}\OperatorTok{(}\NormalTok{linkEdit}\OperatorTok{);}
\NormalTok{            nodes}\OperatorTok{.}\FunctionTok{add}\OperatorTok{(}\NormalTok{otherNode}\OperatorTok{);}
\end{Highlighting}
\end{Shaded}

gui/src/main/java/org/openmarkov/gui/dialog/common/KeyTablePanel.java:377-379

\begin{Shaded}
\begin{Highlighting}[]
\OperatorTok{\}} \ControlFlowTok{catch} \OperatorTok{(}\NormalTok{DoEditException }\OperatorTok{|}\NormalTok{ ThereIsNoNodeInDataException ex}\OperatorTok{)} \OperatorTok{\{}
    \ControlFlowTok{throw} \KeywordTok{new} \FunctionTok{UnrecoverableException}\OperatorTok{(}\NormalTok{ex}\OperatorTok{);}
\OperatorTok{\}}
\end{Highlighting}
\end{Shaded}
\paragraph{Cómo arreglarlo}
Filtrar la lista con la misma regla que aplicará la edición: añadir el candidato solo si \texttt{link\allowbreak{}Edit.\allowbreak{}constraints\allowbreak{}Will\allowbreak{}Be\allowbreak{}Met(\allowbreak{})} es verdadero (el método ya existe en \texttt{PNEdit}). Así la regla sigue escrita en un solo sitio (\texttt{Add\allowbreak{}Link\allowbreak{}Edit.\allowbreak{}check\allowbreak{}Constraints\allowbreak{}Will\allowbreak{}Be\allowbreak{}Met}) y no se crea otra copia en la interfaz. La lista debe recalcularse tras cada alta o baja (hoy ya se recalcula en los caminos normales).

Sitios que añaden enlaces desde la interfaz y deben comportarse igual:

\begin{itemize}
\tightlist
\item
  la pestaña «Parents» (\texttt{Prefixed\allowbreak{}Data\allowbreak{}Table\allowbreak{}Panel});
\item
  la creación de enlaces con el ratón en el lienzo (\texttt{Visual\allowbreak{}Network}, que ya consulta las restricciones para colorear la flecha; su coste se describe en \hyperref[nota:25]{nota~25}, donde no hay nada que corregir);
\item
  las opciones de crear enlace padre/hijo del menú contextual de un nodo;
\item
  \texttt{Multi\allowbreak{}Add\allowbreak{}Link\allowbreak{}Edit}, que descarta en silencio los pares imposibles (ver \hyperref[nota:24]{nota~24}, donde tampoco hay nada que corregir).
\end{itemize}

Conviene también que, si a pesar de todo la edición falla, se ejecuten todas o ninguna de las filas elegidas, y que el mensaje de \texttt{There\allowbreak{}Is\allowbreak{}ACycle} nombre los nodos por su nombre y no por \texttt{to\allowbreak{}String(\allowbreak{})}.

Prueba: sobre \texttt{0\_\allowbreak{}miscellaneous/\allowbreak{}networks/\allowbreak{}bn/\allowbreak{}BN-asia.\allowbreak{}pgmx}, construir la pestaña de padres de «VisitToAsia» y comprobar que la lista no contiene «X-ray», «Dyspnea», «Tuberculosis» ni «TuberculosisOrCancer», y que todos los candidatos ofrecidos se pueden añadir sin excepción.

\setcounter{subsection}{162}
\subsection[El límite superior del último intervalo de revelación no se puede cambiar]{Escribir el límite superior del último intervalo de las condiciones de revelación de un enlace lanza una excepción inesperada y deja la tabla diciendo algo que el enlace no tiene}\label{err:163}
\begin{ficha}
\fichakey{Estado}Presente
\fichakey{Módulo}gui
\fichakey{Paquete}org.openmarkov.gui.component
\fichakey{Ficheros}\path{gui/src/main/java/org/openmarkov/gui/component/RevelationArcDiscretizeTablePanel.java:97-141} \newline \path{gui/src/main/java/org/openmarkov/gui/action/RevelationIntervalEdit.java:144-150} \newline \path{gui/src/main/java/org/openmarkov/gui/window/MainPanelListenerAssistant.java:259-264} \newline \path{gui/src/main/java/org/openmarkov/gui/validator/RevelationArcValidator.java}
\fichakey{Comprobado}Leyendo el código y ejecutándolo
\fichakey{Esfuerzo}Pequeño (horas)
\fichakey{Decisión de equipo}No
\end{ficha}
\paragraph{Qué pasa}
En una red de análisis de decisiones (DAN), un enlace que sale de un nodo de azar numérico puede llevar condiciones de revelación: los intervalos de valores del padre que revelan el valor del hijo. Se editan con clic derecho sobre el enlace → «Edit revelation conditions», que abre \texttt{Revelation\allowbreak{}Arc\allowbreak{}Edit\allowbreak{}Dialog} con una tabla \texttt{Revelation\allowbreak{}Arc\allowbreak{}Discretize\allowbreak{}Table\allowbreak{}Panel} (una fila por intervalo, con límite inferior y superior). El botón de añadir pone un intervalo nuevo al final (\texttt{Revelation\allowbreak{}Interval\allowbreak{}Edit.\allowbreak{}get\allowbreak{}New\allowbreak{}Partitioned\allowbreak{}Interval}: el primero es (-∞, ∞) y los siguientes empiezan donde acaba el anterior y terminan en ∞).

Si el usuario escribe un número en el límite superior del último intervalo, sale el diálogo de «excepción inesperada» que pide enviarla a los desarrolladores: la excepción no se captura en el panel y llega al gestor global (\texttt{OMException\allowbreak{}Handler}). La celda se queda con el número tecleado, pero el intervalo del enlace no cambia. Es el caso corriente: el último intervalo es el que se acaba de añadir y el que el usuario va a acotar. En la práctica, el límite superior del último intervalo no se puede cambiar desde la tabla.

Al escribir un número en una celda de límite, \texttt{table\allowbreak{}Changed} comprueba que no se monte sobre el intervalo vecino. Para el límite inferior mira el intervalo anterior con \texttt{if\ (\allowbreak{}row\ \textgreater{}\ 0)}, bien. Para el límite superior mira el siguiente con \texttt{if\ (\allowbreak{}row\ \textless{}\ num\allowbreak{}Rows)}, que se cumple también en la última fila (\texttt{row\ =\ num\allowbreak{}Rows\ -\ 1}), y entonces \texttt{table.\allowbreak{}get\allowbreak{}Value\allowbreak{}At(\allowbreak{}row\ +\ 1,\allowbreak{}\ …)} lee una fila que no existe.

Medido con un pequeño programa de prueba que construye el panel real sin pantalla sobre una red de análisis de decisiones con un nodo numérico A y un enlace A → B:

\begin{itemize}
\tightlist
\item
  con un único intervalo (-∞, ∞), escribir 5 en su límite superior lanza \texttt{Array\allowbreak{}Index\allowbreak{}Out\allowbreak{}Of\allowbreak{}Bounds\allowbreak{}Exception:\ 1\ \textgreater{}=\ 1};
\item
  con dos intervalos (-∞, 5) y {[}5, ∞), escribir 3 en el límite superior del primero funciona, y escribir 10 en el del último lanza \texttt{Array\allowbreak{}Index\allowbreak{}Out\allowbreak{}Of\allowbreak{}Bounds\allowbreak{}Exception:\ 2\ \textgreater{}=\ 2}; después de la excepción la celda muestra 10 pero el intervalo del enlace sigue siendo {[}5, ∞), porque la edición no llegó a ejecutarse.
\end{itemize}
\paragraph{El código}
gui/src/main/java/org/openmarkov/gui/component/RevelationArcDiscretizeTablePanel.java:110-132

\begin{Shaded}
\begin{Highlighting}[]
\ControlFlowTok{if} \OperatorTok{(}\NormalTok{lower}\OperatorTok{)} \OperatorTok{\{}
    \DataTypeTok{double}\NormalTok{ upperLimit }\OperatorTok{=} \OperatorTok{(}\BuiltInTok{Double}\OperatorTok{)}\NormalTok{ table}\OperatorTok{.}\FunctionTok{getValueAt}\OperatorTok{(}\NormalTok{row}\OperatorTok{,}\NormalTok{ DiscretizeTablePanel}\OperatorTok{.}\FunctionTok{UPPER\_BOUND\_VALUE\_COLUMN\_INDEX}\OperatorTok{);}
    \KeywordTok{...}
    \ControlFlowTok{if} \OperatorTok{(}\NormalTok{row }\OperatorTok{\textgreater{}} \DecValTok{0}\OperatorTok{)} \OperatorTok{\{}
        \DataTypeTok{double}\NormalTok{ previousLimit }\OperatorTok{=} \OperatorTok{(}\BuiltInTok{Double}\OperatorTok{)}\NormalTok{ table}\OperatorTok{.}\FunctionTok{getValueAt}\OperatorTok{(}\NormalTok{row }\OperatorTok{{-}} \DecValTok{1}\OperatorTok{,}\NormalTok{ DiscretizeTablePanel}\OperatorTok{.}\FunctionTok{UPPER\_BOUND\_VALUE\_COLUMN\_INDEX}\OperatorTok{);}
        \KeywordTok{...}
    \OperatorTok{\}}
\OperatorTok{\}} \ControlFlowTok{else} \OperatorTok{\{}
    \DataTypeTok{double}\NormalTok{ lowerLimit }\OperatorTok{=} \OperatorTok{(}\BuiltInTok{Double}\OperatorTok{)}\NormalTok{ table}\OperatorTok{.}\FunctionTok{getValueAt}\OperatorTok{(}\NormalTok{row}\OperatorTok{,}\NormalTok{ DiscretizeTablePanel}\OperatorTok{.}\FunctionTok{LOWER\_BOUND\_VALUE\_COLUMN\_INDEX}\OperatorTok{);}
    \KeywordTok{...}
    \ControlFlowTok{if} \OperatorTok{(}\NormalTok{row }\OperatorTok{\textless{}}\NormalTok{ numRows}\OperatorTok{)} \OperatorTok{\{}
        \DataTypeTok{double}\NormalTok{ nextLimit }\OperatorTok{=} \OperatorTok{(}\BuiltInTok{Double}\OperatorTok{)}\NormalTok{ table}\OperatorTok{.}\FunctionTok{getValueAt}\OperatorTok{(}\NormalTok{row }\OperatorTok{+} \DecValTok{1}\OperatorTok{,}\NormalTok{ DiscretizeTablePanel}\OperatorTok{.}\FunctionTok{LOWER\_BOUND\_VALUE\_COLUMN\_INDEX}\OperatorTok{);}
        \ControlFlowTok{if} \OperatorTok{(}\NormalTok{nextLimit }\OperatorTok{\textless{}}\NormalTok{ newValue}\OperatorTok{)} \OperatorTok{\{}
            \ControlFlowTok{throw} \KeywordTok{new} \FunctionTok{UnrecoverableException}\OperatorTok{(}\KeywordTok{new} \FunctionTok{InvalidArgumentException}\OperatorTok{(}\NormalTok{newValue}\OperatorTok{,} \StringTok{"New value"}\OperatorTok{,} \StringTok{"it is greater than the next bound\textquotesingle{}s lower limit ("} \OperatorTok{+}\NormalTok{ nextLimit }\OperatorTok{+} \StringTok{")"}\OperatorTok{));}
        \OperatorTok{\}}
    \OperatorTok{\}}
\OperatorTok{\}}
\end{Highlighting}
\end{Shaded}
\paragraph{Cómo arreglarlo}
Cambiar la condición a \texttt{if\ (\allowbreak{}row\ +\ 1\ \textless{}\ num\allowbreak{}Rows)}, simétrica de \texttt{row\ \textgreater{}\ 0}. Es el único sitio con esta comprobación de vecinos para las condiciones de revelación; la tabla de tramos del nodo (\texttt{Discretize\allowbreak{}Table\allowbreak{}Panel.\allowbreak{}table\allowbreak{}Changed}) recoloca los límites de otra manera y no tiene esta lectura.

Conviene además que un valor rechazado devuelva la celda al valor del enlace (hoy la celda se queda con el número tecleado cuando algo lanza una excepción).

Prueba de regresión: construir \texttt{Revelation\allowbreak{}Arc\allowbreak{}Discretize\allowbreak{}Table\allowbreak{}Panel} sin pantalla sobre un enlace con uno y con dos intervalos, escribir el límite superior de la última fila y comprobar que no lanza ninguna excepción y que el intervalo del enlace cambia.

\setcounter{subsection}{163}
\subsection[Un número no válido en el indicador da un error y vacía la pestaña]{Escribir un número no válido en el panel del indicador lanza una excepción no prevista y vacía la pestaña de la relación}\label{err:164}
\begin{ficha}
\fichakey{Estado}Presente
\fichakey{Módulo}gui
\fichakey{Paquete}org.openmarkov.gui.dialog.common
\fichakey{Ficheros}\path{gui/src/main/java/org/openmarkov/gui/dialog/common/IndicatorPotentialPanel.java:80-89} \newline \path{core/src/main/java/org/openmarkov/core/model/network/potential/IndicatorPotential.java:155-178} \newline \path{gui/src/main/java/org/openmarkov/gui/dialog/node/NodePropertiesDialog.java:236-262, 307-324, 423-430} \newline \path{gui/src/main/java/org/openmarkov/gui/dialog/node/PotentialEditDialog.java:110-113} \newline \path{gui/src/main/java/org/openmarkov/gui/dialog/common/OkCancelDialog.java:68-79}
\fichakey{Comprobado}Leyendo el código y ejecutándolo
\fichakey{Esfuerzo}Pequeño (horas)
\fichakey{Decisión de equipo}No
\fichakey{Relacionados}\hyperref[err:109]{error~109}, \hyperref[err:152]{error~152}, \hyperref[err:161]{error~161}, \hyperref[nota:13]{nota~13}
\end{ficha}
\paragraph{Qué pasa}
El panel del indicador tiene dos campos de texto libre, el tiempo hasta el suceso y la probabilidad de que ocurra. Si el usuario escribe algo que no es un número, un tiempo negativo o una probabilidad fuera de {[}0, 1{]}, al guardar sale la ventana de «error no previsto» (\texttt{Unexpected\allowbreak{}Throwable\allowbreak{}Dialog}).

El guardado se ejecuta al pulsar OK del diálogo de propiedades del nodo (\texttt{Node\allowbreak{}Properties\allowbreak{}Dialog.\allowbreak{}do\allowbreak{}Ok\allowbreak{}Click\allowbreak{}Before\allowbreak{}Hide}, línea 416) y también al cambiar de pestaña dejando la de la relación (\texttt{handle\allowbreak{}Change\allowbreak{}Tab}, línea 248). En los dos casos el diálogo no se cierra (OK no llega a \texttt{dispose}, \texttt{Ok\allowbreak{}Cancel\allowbreak{}Dialog}, líneas 67-78), pero la pestaña de la relación se queda sin panel y lo escrito se pierde.

Desde el diálogo propio de editar la relación (\texttt{Potential\allowbreak{}Edit\allowbreak{}Dialog}, líneas 109-111) el panel no se quita: sale la misma ventana de error y el usuario puede corregir y volver a pulsar OK.

No se ha abierto ninguno de los dos diálogos; el camino y el efecto sobre la pestaña salen de la lectura del código.

Al guardar, \texttt{save\allowbreak{}Changes} (líneas 80-89) hace \texttt{Double.\allowbreak{}parse\allowbreak{}Double} sobre cada campo y llama a \texttt{Indicator\allowbreak{}Potential.\allowbreak{}set\allowbreak{}Tte} y \texttt{setp\allowbreak{}Occurrence} (módulo \texttt{core}, líneas 155-178), que lanzan \texttt{Invalid\allowbreak{}Argument\allowbreak{}Exception} (subclase de \texttt{Illegal\allowbreak{}Argument\allowbreak{}Exception}) si el tiempo es negativo o la probabilidad no está entre 0 y 1. No hay ninguna comprobación mientras se escribe, y ninguna de esas excepciones se caza.

Los dos caminos de guardado del diálogo del nodo pasan por \texttt{accept\allowbreak{}Potential\allowbreak{}Edit\allowbreak{}Changes} (líneas 297-314), que \textbf{primero quita el panel de la pestaña} (línea 302) y después guarda, cazando solo tres excepciones comprobadas. Las de arriba son de tiempo de ejecución, así que suben hasta el manejador global.

Medido con un pequeño programa de prueba sobre el panel real:

\begin{Verbatim}[breaklines,breakanywhere,fontsize=\small]
tte=abc p=0.5  -> NumberFormatException: For input string: "abc"
tte=-1  p=0.5  -> InvalidArgumentException: Time to event  out of range 0..
tte=2   p=1.5  -> InvalidArgumentException: Occurrence out of range 0<=pOccurrence <=1
tte=2   p=0.25 -> guardado
\end{Verbatim}
\paragraph{El código}
gui/src/main/java/org/openmarkov/gui/dialog/common/IndicatorPotentialPanel.java:80-89

\begin{Shaded}
\begin{Highlighting}[]
\AttributeTok{@Override} \KeywordTok{public} \DataTypeTok{boolean} \FunctionTok{saveChanges}\OperatorTok{()} \KeywordTok{throws}\NormalTok{ DoEditException}\OperatorTok{,} \KeywordTok{...} \OperatorTok{\{}
    \DataTypeTok{boolean}\NormalTok{ result }\OperatorTok{=} \KeywordTok{super}\OperatorTok{.}\FunctionTok{saveChanges}\OperatorTok{();}
\NormalTok{    IndicatorPotential oldPotential }\OperatorTok{=} \OperatorTok{(}\NormalTok{IndicatorPotential}\OperatorTok{)}\NormalTok{ node}\OperatorTok{.}\FunctionTok{getPotentials}\OperatorTok{().}\FunctionTok{get}\OperatorTok{(}\DecValTok{0}\OperatorTok{);}
\NormalTok{    IndicatorPotential newPotential }\OperatorTok{=}  \KeywordTok{new} \FunctionTok{IndicatorPotential}\OperatorTok{(}\NormalTok{oldPotential}\OperatorTok{);}
\NormalTok{        newPotential}\OperatorTok{.}\FunctionTok{setTte}\OperatorTok{(}\BuiltInTok{Double}\OperatorTok{.}\FunctionTok{parseDouble}\OperatorTok{(}\NormalTok{tteTextField}\OperatorTok{.}\FunctionTok{getText}\OperatorTok{()));}
\NormalTok{        newPotential}\OperatorTok{.}\FunctionTok{setpOccurrence}\OperatorTok{(}\BuiltInTok{Double}\OperatorTok{.}\FunctionTok{parseDouble}\OperatorTok{(}\NormalTok{pOcurrenceTextField}\OperatorTok{.}\FunctionTok{getText}\OperatorTok{()));}
    \KeywordTok{new} \FunctionTok{PotentialChangeEdit}\OperatorTok{(}\NormalTok{node}\OperatorTok{,}\NormalTok{ oldPotential}\OperatorTok{,}\NormalTok{ newPotential}\OperatorTok{).}\FunctionTok{executeEdit}\OperatorTok{();}
    \ControlFlowTok{return}\NormalTok{ result}\OperatorTok{;}
\OperatorTok{\}}
\end{Highlighting}
\end{Shaded}
\paragraph{Cómo arreglarlo}
Dos caminos:

\begin{itemize}
\tightlist
\item
  Pedir los dos números con controles que solo admitan valores válidos (un \texttt{JSpinner} o un campo con formato: tiempo ≥ 0, probabilidad en {[}0, 1{]}).
\item
  Validar en \texttt{save\allowbreak{}Changes} y comunicar el error con una excepción comprobada que el diálogo sepa enseñar sin perder lo escrito.
\end{itemize}

La regla de rango debe seguir en los métodos de \texttt{Indicator\allowbreak{}Potential}, que son los que protegen también la lectura de ficheros.

Aparte, \texttt{accept\allowbreak{}Potential\allowbreak{}Edit\allowbreak{}Changes} quita el panel antes de saber si el guardado ha ido bien; eso afecta a cualquier panel de relación que lance una excepción de tiempo de ejecución, no solo a este.

Prueba: construir el panel sobre un nodo de suceso, escribir «abc» y «-1» y comprobar que el guardado no deja escapar \texttt{Number\allowbreak{}Format\allowbreak{}Exception} ni \texttt{Illegal\allowbreak{}Argument\allowbreak{}Exception} y que la relación del nodo no cambia.

\setcounter{subsection}{164}
\subsection[El doble clic en la ventanilla de expresiones mete dos veces el texto]{La ventanilla de expresiones mete dos veces la variable o la función cuando se hace el doble clic que pide su ayuda}\label{err:165}
\begin{ficha}
\fichakey{Estado}Presente
\fichakey{Módulo}gui
\fichakey{Paquete}org.openmarkov.gui.dialog.common
\fichakey{Ficheros}\path{gui/src/main/java/org/openmarkov/gui/dialog/common/ArithmeticExpressionDialog.java:110-145} \newline \path{gui/src/main/resources/gui/localize/ArithmeticExpressionEvaluator_en.xml:5-6}
\fichakey{Comprobado}Leyendo el código
\fichakey{Esfuerzo}Pequeño (horas)
\fichakey{Decisión de equipo}No
\fichakey{Relacionados}\hyperref[err:158]{error~158}
\end{ficha}
\paragraph{Qué pasa}
\texttt{Arithmetic\allowbreak{}Expression\allowbreak{}Dialog} («Enter an expression») es la ventanilla para escribir una expresión: un campo de texto y dos listas, la de variables y la de funciones. La ayuda emergente de cada lista dice «Double click to introduce the selected variable in the expression» y «Double click to introduce the call to the selected function in the expression».

Si el usuario hace ese doble clic, el texto entra dos veces (\texttt{\{Age\}\{Age\}}, o \texttt{pow(\allowbreak{})pow(\allowbreak{})}), y la expresión que queda no es válida y hay que corregirla a mano. Además, un solo clic para elegir una variable de la lista ya modifica la expresión.

Se usa desde muchos sitios: \texttt{GLMPanel} (añadir o reescribir una covariable), \texttt{Function\allowbreak{}Potential\allowbreak{}Panel}, \texttt{Univariate\allowbreak{}Distr\allowbreak{}Potential\allowbreak{}Panel}, \texttt{Augmented\allowbreak{}Prob\allowbreak{}Table\allowbreak{}Potential\allowbreak{}Panel} y \texttt{Table\allowbreak{}With\allowbreak{}Events\allowbreak{}Panel}. Camino del usuario, por ejemplo: propiedades de un nodo con potencial «Linear combination» → botón de añadir covariable → doble clic en una variable de la lista.

Los dos oyentes de ratón salen solo si \texttt{e.\allowbreak{}get\allowbreak{}Click\allowbreak{}Count(\allowbreak{})\ \textless{}\ 1}. Swing entrega \texttt{mouse\allowbreak{}Clicked} con \texttt{get\allowbreak{}Click\allowbreak{}Count(\allowbreak{})} igual a 1 en el primer clic y 2 en el segundo de un doble clic; nunca vale menos de 1. Por eso \textbf{cada clic} inserta el texto: un clic simple ya mete \texttt{\{variable\}}, y el doble clic lo mete dos veces.

Comprobado leyendo; no se ejecutó porque la ventanilla es un \texttt{JDialog} y no se puede construir sin pantalla.
\paragraph{El código}
gui/src/main/java/org/openmarkov/gui/dialog/common/ArithmeticExpressionDialog.java:110-121

\begin{Shaded}
\begin{Highlighting}[]
\NormalTok{variableList}\OperatorTok{.}\FunctionTok{addMouseListener}\OperatorTok{(}\KeywordTok{new} \BuiltInTok{MouseAdapter}\OperatorTok{()} \OperatorTok{\{}
    \AttributeTok{@Override} \KeywordTok{public} \DataTypeTok{void} \FunctionTok{mouseClicked}\OperatorTok{(}\BuiltInTok{MouseEvent}\NormalTok{ e}\OperatorTok{)} \OperatorTok{\{}
        \KeywordTok{super}\OperatorTok{.}\FunctionTok{mouseClicked}\OperatorTok{(}\NormalTok{e}\OperatorTok{);}
        \ControlFlowTok{if} \OperatorTok{(}\NormalTok{e}\OperatorTok{.}\FunctionTok{getClickCount}\OperatorTok{()} \OperatorTok{\textless{}} \DecValTok{1}\OperatorTok{)} \OperatorTok{\{}
            \ControlFlowTok{return}\OperatorTok{;}
        \OperatorTok{\}}
        \ControlFlowTok{try} \OperatorTok{\{}
            \FunctionTok{insertTextInExpression}\OperatorTok{(}\StringTok{"\{"} \OperatorTok{+}\NormalTok{ variableList}\OperatorTok{.}\FunctionTok{getSelectedValue}\OperatorTok{()} \OperatorTok{+} \StringTok{"\}"}\OperatorTok{);}
        \OperatorTok{\}} \ControlFlowTok{catch} \OperatorTok{(}\BuiltInTok{BadLocationException}\NormalTok{ ex}\OperatorTok{)} \OperatorTok{\{}
            \ControlFlowTok{throw} \KeywordTok{new} \FunctionTok{UnrecoverableException}\OperatorTok{(}\NormalTok{ex}\OperatorTok{);}
        \OperatorTok{\}}
    \OperatorTok{\}}
\OperatorTok{\});}
\end{Highlighting}
\end{Shaded}
\paragraph{Cómo arreglarlo}
Cambiar la condición a \texttt{e.\allowbreak{}get\allowbreak{}Click\allowbreak{}Count(\allowbreak{})\ !=\ 2} (o \texttt{\textless{}\ 2}) en los dos oyentes: el de \texttt{variable\allowbreak{}List} (línea 113) y el de \texttt{function\allowbreak{}List} (línea 135). Con \texttt{!=\ 2} un triple clic no inserta una tercera vez.

Prueba de regresión: si se extrae el cuerpo de los oyentes a un método que reciba el número de clics, se puede probar sin pantalla que con 1 no inserta y con 2 inserta una vez; si no, prueba manual con doble clic en cada lista.

\setcounter{subsection}{165}
\subsection[Tras cambiar el tipo de un nodo, una copia suya no encuentra sus padres]{Tras cambiar el tipo de un nodo, una copia del nodo ya no encuentra sus padres en la red, y la vista de tabla de un modelo canónico se construye sin ellos}\label{err:166}
\begin{ficha}
\fichakey{Estado}Presente
\fichakey{Módulo}core
\fichakey{Paquete}org.openmarkov.core.model.network
\fichakey{Ficheros}\path{core/src/main/java/org/openmarkov/core/model/network/Node.java:80, 133, 141-149, 272-283, 339-341, 413-424} \newline \path{core/src/main/java/org/openmarkov/core/model/graph/Graph.java:49-61, 69-71} \newline \path{gui/src/main/java/org/openmarkov/gui/dialog/node/ICIOptionListenerAssistant.java:62-70, 101-115} \newline \path{gui/src/main/java/org/openmarkov/gui/dialog/common/TablePotentialPanel.java:360-362} \newline \path{gui/src/main/java/org/openmarkov/gui/dialog/node/NodeDefinitionPanel.java:332-341}
\fichakey{Comprobado}Leyendo el código y ejecutándolo
\fichakey{Esfuerzo}Pequeño (horas)
\fichakey{Decisión de equipo}No
\fichakey{Relacionados}\hyperref[err:3]{error~3}, \hyperref[err:50]{error~50}, \hyperref[nota:180]{nota~180}
\end{ficha}
\paragraph{Qué pasa}
Camino del usuario (leído, no ejecutado): el tipo de un nodo se cambia desde el diálogo de propiedades del nodo, con los botones de tipo (\texttt{Node\allowbreak{}Definition\allowbreak{}Panel.\allowbreak{}java:332-341} → \texttt{Change\allowbreak{}Node\allowbreak{}Type\allowbreak{}Edit}, que llama a \texttt{set\allowbreak{}Node\allowbreak{}Type}). El único sitio de la aplicación que construye una copia de un nodo real es la vista «TPC» (tabla de probabilidad condicional) del panel de un modelo canónico: al pulsar ese botón, \texttt{ICIOption\allowbreak{}Listener\allowbreak{}Assistant.\allowbreak{}item\allowbreak{}State\allowbreak{}Changed\allowbreak{}TPC} hace \texttt{new\ Node(\allowbreak{}panel.\allowbreak{}get\allowbreak{}Node(\allowbreak{}))} y construye con esa copia un \texttt{CPTable\allowbreak{}Panel} (líneas 101-115). \texttt{Table\allowbreak{}Potential\allowbreak{}Panel.\allowbreak{}how\allowbreak{}Many\allowbreak{}Rows} calcula las filas de la tabla con \texttt{n.\allowbreak{}get\allowbreak{}Parents(\allowbreak{}).\allowbreak{}size(\allowbreak{})} (líneas 341-343).

Así, si el nodo se creó con otro tipo (decisión, utilidad), luego se pasó a azar y se le puso un modelo canónico (OR/MAX/MIN ruidoso), la vista TPC se construye como si el nodo no tuviera padres: le faltan las filas de cabecera de los padres mientras la tabla expandida sí los tiene. No se ha abierto el diálogo, así que no está medido qué ve exactamente el usuario: una tabla mal dibujada o una excepción.

En Java, dos objetos iguales según \texttt{equals} deben tener el mismo \texttt{hash\allowbreak{}Code}; si no, las tablas de dispersión (\texttt{Hash\allowbreak{}Map}, \texttt{Hash\allowbreak{}Set}) no los encuentran. \texttt{Node.\allowbreak{}equals} (líneas 426-433) considera iguales dos nodos con la misma variable (por identidad), la misma red y \textbf{el mismo tipo actual}. \texttt{Node.\allowbreak{}hash\allowbreak{}Code} devuelve un campo \texttt{final} calculado en el constructor con la variable y el tipo que tenía entonces (líneas 80, 133 y 148). \texttt{set\allowbreak{}Node\allowbreak{}Type} (líneas 272-283) cambia el tipo pero no ese campo.

Resultado: tras cambiar el tipo de un nodo, un nodo igual construido después (por ejemplo con el constructor de copia \texttt{Node(\allowbreak{}Node)}, líneas 141-149, que calcula el código con el tipo nuevo) tiene otro código de dispersión.

Con el propio objeto no pasa nada: su código no cambia y es igual a sí mismo, así que el grafo (que guarda padres e hijos en \texttt{Hash\allowbreak{}Map\textless{}Node,\allowbreak{}\ …\textgreater{}}, \texttt{Graph.\allowbreak{}java:49-61}) lo sigue encontrando. Cambiar el \textbf{nombre} de una variable tampoco afecta, porque \texttt{Variable} no redefine \texttt{equals} ni \texttt{hash\allowbreak{}Code} y se compara por identidad. El problema aparece solo cuando se consulta la red con \textbf{otro objeto} \texttt{Node} igual.

Medido con un pequeño programa de prueba sobre un diagrama de influencia con el enlace P → X, con X creado como decisión. Antes de cambiar nada, \texttt{get\allowbreak{}Parents(\allowbreak{}X)} y \texttt{get\allowbreak{}Parents(\allowbreak{}new\ Node(\allowbreak{}X))} dan \texttt{{[}P{]}}. Tras \texttt{X.\allowbreak{}set\allowbreak{}Node\allowbreak{}Type(\allowbreak{}CHANCE)}: \texttt{new\ Node(\allowbreak{}X).\allowbreak{}equals(\allowbreak{}X)} es verdadero pero los códigos difieren; \texttt{get\allowbreak{}Parents(\allowbreak{}X)} sigue dando \texttt{{[}P{]}} y \texttt{get\allowbreak{}Parents(\allowbreak{}copia)} da \texttt{{[}{]}}; un \texttt{Hash\allowbreak{}Set} que contiene X dice que no contiene la copia.
\paragraph{El código}
core/src/main/java/org/openmarkov/core/model/network/Node.java:272-283 y 426-437

\begin{Shaded}
\begin{Highlighting}[]
\KeywordTok{public} \DataTypeTok{void} \FunctionTok{setNodeType}\OperatorTok{(}\NormalTok{NodeType newNodeType}\OperatorTok{)} \OperatorTok{\{}
\NormalTok{    NodeType oldNodeType }\OperatorTok{=} \KeywordTok{this}\OperatorTok{.}\FunctionTok{nodeType}\OperatorTok{;}
    \ControlFlowTok{if}\OperatorTok{(}\NormalTok{oldNodeType }\OperatorTok{==}\NormalTok{ newNodeType}\OperatorTok{)\{}
        \ControlFlowTok{return}\OperatorTok{;}
    \OperatorTok{\}}
    \KeywordTok{this}\OperatorTok{.}\FunctionTok{probNet}\OperatorTok{.}\FunctionTok{nodeDepot}\OperatorTok{.}\FunctionTok{removeNode}\OperatorTok{(}\KeywordTok{this}\OperatorTok{);}
    \KeywordTok{this}\OperatorTok{.}\FunctionTok{nodeType} \OperatorTok{=}\NormalTok{ newNodeType}\OperatorTok{;}
    \KeywordTok{this}\OperatorTok{.}\FunctionTok{probNet}\OperatorTok{.}\FunctionTok{nodeDepot}\OperatorTok{.}\FunctionTok{addNode}\OperatorTok{(}\KeywordTok{this}\OperatorTok{);}
\OperatorTok{\}}
\KeywordTok{...}
\AttributeTok{@Override} \KeywordTok{public} \DataTypeTok{boolean} \FunctionTok{equals}\OperatorTok{(}\BuiltInTok{Object}\NormalTok{ obj}\OperatorTok{)} \OperatorTok{\{}
    \DataTypeTok{boolean}\NormalTok{ equals }\OperatorTok{=}\NormalTok{ obj }\KeywordTok{instanceof} \BuiltInTok{Node}\OperatorTok{;}
    \ControlFlowTok{if} \OperatorTok{(}\NormalTok{equals}\OperatorTok{)} \OperatorTok{\{}
        \BuiltInTok{Node}\NormalTok{ otherNode }\OperatorTok{=} \OperatorTok{(}\BuiltInTok{Node}\OperatorTok{)}\NormalTok{ obj}\OperatorTok{;}
\NormalTok{        equals }\OperatorTok{=}\NormalTok{ variable}\OperatorTok{.}\FunctionTok{equals}\OperatorTok{(}\NormalTok{otherNode}\OperatorTok{.}\FunctionTok{variable}\OperatorTok{)} \OperatorTok{\&\&}\NormalTok{ probNet}\OperatorTok{.}\FunctionTok{equals}\OperatorTok{(}\NormalTok{otherNode}\OperatorTok{.}\FunctionTok{probNet}\OperatorTok{)} \OperatorTok{\&\&}\NormalTok{ nodeType }\OperatorTok{==}\NormalTok{ otherNode}\OperatorTok{.}\FunctionTok{nodeType}\OperatorTok{;}
    \OperatorTok{\}}
    \ControlFlowTok{return}\NormalTok{ equals}\OperatorTok{;}
\OperatorTok{\}}

\AttributeTok{@Override} \KeywordTok{public} \DataTypeTok{int} \FunctionTok{hashCode}\OperatorTok{()} \OperatorTok{\{}
    \ControlFlowTok{return}\NormalTok{ hashCode}\OperatorTok{;}
\OperatorTok{\}}
\end{Highlighting}
\end{Shaded}
\paragraph{Cómo arreglarlo}
Hacer que \texttt{equals} y \texttt{hash\allowbreak{}Code} usen los mismos datos y que ninguno de ellos cambie mientras el nodo está en el grafo. Opciones:

\begin{itemize}
\tightlist
\item
  \begin{enumerate}
  \def\labelenumi{(\alph{enumi})}
  \tightlist
  \item
    Quitar el tipo de \texttt{equals} y del código de dispersión: un nodo se identifica por su variable y su red, y el tipo es un atributo que cambia.
  \end{enumerate}
\item
  \begin{enumerate}
  \def\labelenumi{(\alph{enumi})}
  \setcounter{enumi}{1}
  \tightlist
  \item
    No redefinir ninguno de los dos y comparar nodos por identidad, que es lo que el grafo necesita.
  \end{enumerate}
\end{itemize}

Con la variable pasa lo mismo: \texttt{set\allowbreak{}Variable} (líneas 166-170) cambia la variable pero no el código, así que la regla afecta también a ese método (llamado desde \texttt{Prob\allowbreak{}Net\allowbreak{}Operations.\allowbreak{}convert\allowbreak{}Numerical\allowbreak{}Variables\allowbreak{}To\allowbreak{}FS}, líneas 451 y 538, y \texttt{DANOperations.\allowbreak{}java:159}). Tras el cambio hay que revisar además los 72 usos de \texttt{Map\textless{}Node,\allowbreak{}…\textgreater{}}/\texttt{Set\textless{}Node\textgreater{}} del código de producción.

Prueba de regresión en \texttt{core}: la de la medida (cambiar el tipo y comprobar que \texttt{get\allowbreak{}Parents(\allowbreak{}new\ Node(\allowbreak{}nodo))} devuelve los padres y que \texttt{new\ Node(\allowbreak{}nodo).\allowbreak{}hash\allowbreak{}Code(\allowbreak{})\ ==\ nodo.\allowbreak{}hash\allowbreak{}Code(\allowbreak{})} cuando son iguales), repetida con \texttt{set\allowbreak{}Variable}.

\setcounter{subsection}{166}
\subsection[Una columna con incertidumbre fuera del primer estado se deja editar]{Una columna con incertidumbre en un estado que no es el primero se deja editar, y una con incertidumbre solo en el primero se bloquea entera}\label{err:167}
\begin{ficha}
\fichakey{Estado}Presente
\fichakey{Módulo}core, gui
\fichakey{Paquete}org.openmarkov.core.model.network.potential; org.openmarkov.gui.dialog.common
\fichakey{Ficheros}\path{core/src/main/java/org/openmarkov/core/model/network/potential/TablePotential.java:413-420} \newline \path{core/src/main/java/org/openmarkov/core/model/network/potential/AbstractIndexedPotential.java:262-283} \newline \path{gui/src/main/java/org/openmarkov/gui/dialog/common/TablePotentialPanel.java:337-354, 616-641, 823-842, 849-866} \newline \path{io/src/main/java/org/openmarkov/io/probmodel/reader/PGMXReader_0_2.java:1324-1355}
\fichakey{Comprobado}Leyendo el código y ejecutándolo
\fichakey{Esfuerzo}Pequeño (horas)
\fichakey{Decisión de equipo}No
\fichakey{Relacionados}\hyperref[err:76]{error~76}
\end{ficha}
\paragraph{Qué pasa}
En una tabla de probabilidad condicionada con incertidumbre (un \texttt{Uncertain\allowbreak{}Table\allowbreak{}Potential}, que guarda para cada casilla una distribución de segundo orden o nada), el panel de la tabla (\texttt{Table\allowbreak{}Potential\allowbreak{}Panel}, módulo \texttt{gui}) decide qué columnas llevan incertidumbre para tres cosas:

\begin{itemize}
\tightlist
\item
  bloquear la edición a mano de esa columna y pintarla con el icono de incertidumbre (\texttt{get\allowbreak{}Uncertainty\allowbreak{}In\allowbreak{}Columns}, líneas 318-335, que alimenta \texttt{get\allowbreak{}Not\allowbreak{}Editable\allowbreak{}Positions}, líneas 614-620);
\item
  habilitar «Assign/Edit/Remove uncertainty» en el menú contextual (\texttt{update\allowbreak{}Contextual\allowbreak{}Menu\allowbreak{}Options}, línea 810);
\item
  abrir el diálogo de incertidumbre con doble clic (\texttt{double\allowbreak{}Click\allowbreak{}Event}, línea 842).
\end{itemize}

Si una columna tiene incertidumbre en la casilla de un estado del nodo que no es el primero, pero no en la del primero, la columna queda editable entera y sin icono, y el menú contextual ofrece «asignar» incertidumbre donde ya la hay. Editarla cambia los números y deja la distribución de segundo orden como estaba, sin aviso, de modo que los valores y la incertidumbre de la columna dejan de corresponderse. Al revés, una columna con incertidumbre solo en el primer estado queda bloqueada entera.

Camino de llegada: el diálogo de incertidumbre (\texttt{Uncertain\allowbreak{}Values\allowbreak{}Dialog} con \texttt{Uncertain\allowbreak{}Values\allowbreak{}Edit}) siempre escribe la columna completa, así que desde la interfaz no se produce. Sí se produce abriendo un \texttt{.\allowbreak{}pgmx} cuyo \texttt{\textless{}Uncertain\allowbreak{}Values\textgreater{}} tenga huecos (\texttt{\textless{}Value\ /\allowbreak{}\textgreater{}}) en algunas casillas de una columna: \texttt{PGMXReader\_\allowbreak{}0\_\allowbreak{}2.\allowbreak{}get\allowbreak{}Uncertain\allowbreak{}Values} lo acepta casilla a casilla.

Las tres llamadas del panel usan \texttt{Table\allowbreak{}Potential.\allowbreak{}has\allowbreak{}Uncertainty(\allowbreak{}configuration)}, donde la configuración tiene un hallazgo por cada padre (el nodo no está). \texttt{has\allowbreak{}Uncertainty} traduce esa configuración a una posición con \texttt{Abstract\allowbreak{}Indexed\allowbreak{}Potential.\allowbreak{}get\allowbreak{}Position(\allowbreak{}Evidence\allowbreak{}Case)}, que, al faltar la variable del nodo, pone su coordenada a 0. Resultado: solo se consulta la casilla del primer estado del nodo en esa columna.

Medido con un pequeño programa de prueba sobre una red A→B (B con estados b0, b1, b2): con incertidumbre solo en la casilla (A=a0, B=b1), \texttt{has\allowbreak{}Uncertainty(\allowbreak{}A=a0)} da \texttt{false} y la columna a0 queda editable entera y sin icono; con incertidumbre solo en (A=a1, B=b0), la columna a1 queda bloqueada entera. Editar la columna a0 cambia los números con \texttt{Table\allowbreak{}Potential\allowbreak{}Value\allowbreak{}Edit} y deja la distribución de b1 como estaba.

También medido: en los 167 ficheros \texttt{.\allowbreak{}pgmx} legibles de \texttt{0\_\allowbreak{}miscellaneous/\allowbreak{}networks} y de los recursos de prueba de \texttt{core}, \texttt{io}, \texttt{gui} e \texttt{inference} hay 102 tablas con incertidumbre y ninguna columna con incertidumbre incompleta.
\paragraph{El código}
core/src/main/java/org/openmarkov/core/model/network/potential/AbstractIndexedPotential.java:262-283

\begin{Shaded}
\begin{Highlighting}[]
    \KeywordTok{public} \DataTypeTok{int} \FunctionTok{getPosition}\OperatorTok{(}\NormalTok{EvidenceCase configuration}\OperatorTok{)} \OperatorTok{\{}
        \KeywordTok{...}
        \ControlFlowTok{if} \OperatorTok{(}\NormalTok{varsTable}\OperatorTok{.}\FunctionTok{size}\OperatorTok{()} \OperatorTok{!=}\NormalTok{ sizeEvi}\OperatorTok{)} \OperatorTok{\{}
\NormalTok{            sizeCoordinates }\OperatorTok{=}\NormalTok{ sizeEvi }\OperatorTok{+} \DecValTok{1}\OperatorTok{;}
\NormalTok{            startLoop       }\OperatorTok{=} \DecValTok{1}\OperatorTok{;}
\NormalTok{            coordinates     }\OperatorTok{=} \KeywordTok{new} \DataTypeTok{int}\OperatorTok{[}\NormalTok{sizeCoordinates}\OperatorTok{];}
\NormalTok{            coordinates}\OperatorTok{[}\DecValTok{0}\OperatorTok{]}  \OperatorTok{=} \DecValTok{0}\OperatorTok{;}
        \OperatorTok{\}} \ControlFlowTok{else} \OperatorTok{\{}
        \KeywordTok{...}
\end{Highlighting}
\end{Shaded}
\paragraph{Cómo arreglarlo}
Hay que fijar el significado y aplicarlo en un solo sitio: que una columna «tiene incertidumbre» si alguna de sus casillas la tiene (lo más prudente para bloquear la edición), recorriendo los estados del nodo en \texttt{has\allowbreak{}Uncertainty} o en un método nuevo de columna. Los tres usos de \texttt{Table\allowbreak{}Potential\allowbreak{}Panel} (líneas 329, 810 y 842) deben pasar por esa misma comprobación.

Alternativa: rechazar al leer el fichero una columna con incertidumbre incompleta; hoy el lector no hace ninguna comprobación de ese tipo (el método \texttt{check\allowbreak{}Uncertain\allowbreak{}Table}, que no comprobaba nada, se quitó en el commit \texttt{13a2b89}).

Prueba: tabla con incertidumbre solo en el segundo estado de una columna; la columna debe quedar no editable y con las opciones «Edit/Remove uncertainty».

\setcounter{subsection}{167}
\subsection[Tras alejar el zoom, los nombres de estado siguen recortados al acercarlo]{Tras alejar el zoom en modo inferencia, los nombres de estado recortados siguen recortados al volver a acercarlo}\label{err:168}
\begin{ficha}
\fichakey{Estado}Presente
\fichakey{Módulo}gui
\fichakey{Paquete}org.openmarkov.gui.graphic
\fichakey{Ficheros}\path{gui/src/main/java/org/openmarkov/gui/graphic/VisualState.java:37, 317-320} \newline \path{gui/src/main/java/org/openmarkov/gui/graphic/VisualElement.java:132-146} \newline \path{gui/src/main/java/org/openmarkov/gui/graphic/FSVariableBox.java:64-82, 100-106} \newline \path{gui/src/main/java/org/openmarkov/gui/graphic/NumericVariableBox.java:134-135} \newline \path{gui/src/main/java/org/openmarkov/gui/window/edition/networkEditorPanel/NetworkEditorPanel.java:241-251}
\fichakey{Comprobado}Leyendo el código y ejecutándolo
\fichakey{Esfuerzo}Pequeño (horas)
\fichakey{Decisión de equipo}No
\end{ficha}
\paragraph{Qué pasa}
En modo inferencia, cada nodo desplegado dibuja sus estados con una barra de probabilidad y, delante, el nombre del estado. Si el usuario aleja el zoom (por ejemplo al 50 \%), algunos nombres que caben enteros al 100 \% se recortan; al volver al 100 \%, siguen dibujándose recortados, como «ne\ldots ve» en lugar de «negative». Vuelven a salir enteros cuando se recrean los estados visuales (\texttt{FSVariable\allowbreak{}Box.\allowbreak{}update\allowbreak{}Num\allowbreak{}Cases}, por ejemplo al cambiar el número de casos de evidencia). Es una consecuencia solo visual y menor.

\texttt{Visual\allowbreak{}State.\allowbreak{}paint} (módulo \texttt{gui}, paquete \texttt{org.\allowbreak{}openmarkov.\allowbreak{}gui.\allowbreak{}graphic}) acorta el nombre con \texttt{Visual\allowbreak{}Element.\allowbreak{}adjust\allowbreak{}Text} para que quepa en 52 unidades (\texttt{Inner\allowbreak{}Box.\allowbreak{}BAR\_\allowbreak{}HORIZONTAL\_\allowbreak{}POSITION}) delante de la barra: si no cabe, deja el principio del nombre, «\ldots» y las dos últimas letras. La línea 319 guarda el resultado en el campo \texttt{state\allowbreak{}Name} en lugar de usarlo solo para dibujar, y el nombre original se pierde en ese objeto.

El ancho disponible es una constante, así que agrandar la ventana o dibujar el nodo más ancho no cambia nada. Pero el ancho medido del texto sí depende del zoom: \texttt{Network\allowbreak{}Editor\allowbreak{}Panel.\allowbreak{}do\allowbreak{}Paint} escala el \texttt{Graphics2D} con el zoom (línea 251) y \texttt{adjust\allowbreak{}Text} mide con ese contexto. Con el zoom alejado, un nombre se recorta, queda guardado recortado y ya no recupera su forma completa.

Otros usos del nombre guardado: \texttt{FSVariable\allowbreak{}Box.\allowbreak{}get\allowbreak{}Visual\allowbreak{}State(\allowbreak{}String)} busca por nombre, pero no tiene llamadores; \texttt{Numeric\allowbreak{}Variable\allowbreak{}Box.\allowbreak{}update\allowbreak{}Num\allowbreak{}Cases} (líneas 134-135) recrea su estado visual con el nombre guardado, que en un nodo numérico es el texto de su único estado.

Medido con la letra de los estados (11 puntos): «negative» mide 48 a escala 1 y 56 a escala 0,5; «Therapy», 46 y 52. Es decir, con el zoom al 50 \% esos nombres se recortan (≥ 52). No se ha visto en la aplicación; se ha deducido de la lectura y de la medida de anchos.
\paragraph{El código}
gui/src/main/java/org/openmarkov/gui/graphic/VisualState.java:317-320

\begin{Shaded}
\begin{Highlighting}[]
\NormalTok{        g}\OperatorTok{.}\FunctionTok{setColor}\OperatorTok{(}\NormalTok{GUIColors}\OperatorTok{.}\FunctionTok{Inference}\OperatorTok{.}\FunctionTok{BOX\_TEXT}\OperatorTok{.}\FunctionTok{getColor}\OperatorTok{());}
\NormalTok{        g}\OperatorTok{.}\FunctionTok{setFont}\OperatorTok{(}\NormalTok{STATES\_FONT}\OperatorTok{);}
        \KeywordTok{this}\OperatorTok{.}\FunctionTok{stateName} \OperatorTok{=} \FunctionTok{adjustText}\OperatorTok{(}\KeywordTok{this}\OperatorTok{.}\FunctionTok{stateName}\OperatorTok{,}\NormalTok{ InnerBox}\OperatorTok{.}\FunctionTok{BAR\_HORIZONTAL\_POSITION}\OperatorTok{,} \DecValTok{2}\OperatorTok{,}\NormalTok{ STATES\_FONT}\OperatorTok{,}\NormalTok{ g}\OperatorTok{);}
\NormalTok{        g}\OperatorTok{.}\FunctionTok{drawString}\OperatorTok{(}\KeywordTok{this}\OperatorTok{.}\FunctionTok{stateName}\OperatorTok{,} \OperatorTok{(}\DataTypeTok{int}\OperatorTok{)}\NormalTok{ xName}\OperatorTok{,} \OperatorTok{(}\DataTypeTok{int}\OperatorTok{)}\NormalTok{ yText}\OperatorTok{);}
\end{Highlighting}
\end{Shaded}
\paragraph{Cómo arreglarlo}
Dibujar el texto acortado en una variable local y no tocar \texttt{state\allowbreak{}Name}. Buscar si hay otros \texttt{adjust\allowbreak{}Text} que asignen a un campo (por ejemplo, en \texttt{Visual\allowbreak{}Node} para el nombre o el título del nodo) y aplicar la misma regla.

Prueba: pintar un \texttt{Visual\allowbreak{}State} con escala 0,5 y luego con escala 1 sobre una imagen en memoria y comprobar que \texttt{get\allowbreak{}State\allowbreak{}Name(\allowbreak{})} sigue devolviendo el nombre completo.

\setcounter{subsection}{168}
\subsection[Con precisión 0.0, 2.0 o 0.5, los límites de los tramos acaban en punto]{Con una precisión que no contiene un «1» (0.0, 2.0, 0.5), los límites de los tramos se muestran con un punto final colgando}\label{err:169}
\begin{ficha}
\fichakey{Estado}Presente
\fichakey{Módulo}gui
\fichakey{Paquete}org.openmarkov.gui.component
\fichakey{Ficheros}\path{gui/src/main/java/org/openmarkov/gui/component/DiscretizeTablePanel.java:596-639} \newline \path{gui/src/main/java/org/openmarkov/gui/component/DiscretizeTablePanel.java:563-568} \newline \path{gui/src/main/java/org/openmarkov/gui/component/DiscretizeTablePanel.java:1011} \newline \path{gui/src/main/java/org/openmarkov/gui/dialog/node/NodeDomainValuesTablePanel.java:300-316} \newline \path{io/src/main/java/org/openmarkov/io/probmodel/reader/PGMXReader_0_2.java:1007-1008} \newline \path{core/src/main/java/org/openmarkov/core/model/network/Util.java:178-202}
\fichakey{Comprobado}Leyendo el código y ejecutándolo
\fichakey{Esfuerzo}Pequeño (horas)
\fichakey{Decisión de equipo}No
\fichakey{Relacionados}\hyperref[err:40]{error~40}
\end{ficha}
\paragraph{Qué pasa}
En la pestaña de dominio del diálogo de propiedades de un nodo numérico o discretizado, la tabla de tramos muestra los límites como texto. Si la precisión de la variable no contiene un «1» (por ejemplo 0.0, 2.0 o 0.5), cada límite finito sale con un punto final colgando: \texttt{5.\allowbreak{}}, \texttt{10.\allowbreak{}}, \texttt{1234.\allowbreak{}}.

Camino del usuario: abrir \texttt{0\_\allowbreak{}miscellaneous/\allowbreak{}networks/\allowbreak{}id/\allowbreak{}ID-arthronet-ce.\allowbreak{}pgmx}, abrir las propiedades del nodo de utilidad «EVAC Implante» (precisión 2.0), pestaña de dominio, y escribir 5 en el límite inferior (que era -∞). La tabla pasa a mostrar \texttt{(\allowbreak{},\allowbreak{}\ 5.\allowbreak{},\allowbreak{}\ ,\allowbreak{},\allowbreak{}\ ∞,\allowbreak{}\ )}. Lo mismo con «Reward for coffee {[}0{]}» de \texttt{0\_\allowbreak{}miscellaneous/\allowbreak{}networks/\allowbreak{}pomdp/\allowbreak{}POMDP-coffee-robot.\allowbreak{}pgmx} (precisión 0.0).

Esas precisiones no son rebuscadas. El lector de \texttt{.\allowbreak{}pgmx} pone precisión 0.0 a una variable numérica cuyo elemento no trae \texttt{\textless{}Precision\textgreater{}} (\texttt{PGMXReader\_\allowbreak{}0\_\allowbreak{}2.\allowbreak{}java:1008}), y \texttt{new\ Variable(\allowbreak{}String)} también la deja en 0.0. En las redes de ejemplo del repositorio (\texttt{0\_\allowbreak{}miscellaneous/\allowbreak{}networks}) hay 1251 nodos de utilidad con \texttt{\textless{}Precision\textgreater{}0.\allowbreak{}0\textless{}/\allowbreak{}Precision\textgreater{}} y 327 con \texttt{\textless{}Precision\textgreater{}2.\allowbreak{}0\textless{}/\allowbreak{}Precision\textgreater{}} (todos los de \texttt{0\_\allowbreak{}miscellaneous/\allowbreak{}networks/\allowbreak{}id/\allowbreak{}ID-arthronet-ce.\allowbreak{}pgmx}, por ejemplo). La lista desplegable de precisión del diálogo solo ofrece 1, 0.1, 0.01, 0.001 y 0.0001, que se muestran bien, así que el defecto aparece con precisiones que vienen de fichero.

\texttt{Discretize\allowbreak{}Table\allowbreak{}Panel.\allowbreak{}set\allowbreak{}Data\allowbreak{}From\allowbreak{}Partitioned\allowbreak{}Interval} convierte los límites con \texttt{convert\allowbreak{}To\allowbreak{}String\allowbreak{}Limit\allowbreak{}Values(\allowbreak{}limites,\allowbreak{}\ Double.\allowbreak{}to\allowbreak{}String(\allowbreak{}precision))}, que decide cuántos decimales mostrar buscando en el texto de la precisión la posición del punto y la del primer «1»: con «0.01» salen dos decimales y con «0.001» tres. Si la precisión no contiene un «1», el número de decimales queda en 0, y 0 no significa «sin decimales» sino «cortar justo después del punto».

Hay casos peores con la misma causa, que solo llegan desde un fichero: con precisión 10.0 el número de decimales sale negativo y el texto se corta antes del punto, de modo que 1234 se muestra como \texttt{123} y 3.14 como una cadena vacía; con precisión 1.0E-10 la rama de notación científica lee solo la primera cifra del exponente y muestra un decimal. Con esas precisiones el daño no se queda en lo que se ve: \texttt{Util.\allowbreak{}round\allowbreak{}With\allowbreak{}Precision}, que se aplica al límite tecleado antes de guardarlo (\texttt{Discretize\allowbreak{}Table\allowbreak{}Panel}, línea 962), convierte con precisión 10.0 el 1234 en 123 y el 12345.678 en 1234, y lanza \texttt{Number\allowbreak{}Format\allowbreak{}Exception} con 5 o con 3.14; con precisión 100.0 convierte 1234 en 12, y \texttt{convert\allowbreak{}To\allowbreak{}String\allowbreak{}Limit\allowbreak{}Values} lanza \texttt{String\allowbreak{}Index\allowbreak{}Out\allowbreak{}Of\allowbreak{}Bounds\allowbreak{}Exception}, de modo que, por lectura, la tabla de tramos no se puede ni rellenar (medido con las dos funciones aisladas). En las redes del repositorio no hay ninguna variable con esas precisiones. \texttt{Util.\allowbreak{}round\allowbreak{}With\allowbreak{}Precision} (el redondeo que se aplica al número tecleado) interpreta el texto de la precisión con la misma lógica.

Medido con un pequeño programa de prueba que construye el panel real sin pantalla (el camino del usuario de arriba) y con la función aislada: precisión 0.0 → \texttt{{[}-∞,\allowbreak{}\ 0.\allowbreak{},\allowbreak{}\ 10.\allowbreak{},\allowbreak{}\ 3.\allowbreak{},\allowbreak{}\ 2.\allowbreak{},\allowbreak{}\ 1234.\allowbreak{},\allowbreak{}\ ∞{]}}; precisión 0.01 → \texttt{{[}-∞,\allowbreak{}\ 0.\allowbreak{}00,\allowbreak{}\ 10.\allowbreak{}00,\allowbreak{}\ 3.\allowbreak{}14,\allowbreak{}\ 2.\allowbreak{}50,\allowbreak{}\ 1234.\allowbreak{}00,\allowbreak{}\ ∞{]}}.
\paragraph{El código}
gui/src/main/java/org/openmarkov/gui/component/DiscretizeTablePanel.java:596-630

\begin{Shaded}
\begin{Highlighting}[]
\KeywordTok{public} \DataTypeTok{static} \BuiltInTok{String}\OperatorTok{[]} \FunctionTok{convertToStringLimitValues}\OperatorTok{(}\DataTypeTok{double}\OperatorTok{[]}\NormalTok{ limits}\OperatorTok{,} \BuiltInTok{String}\NormalTok{ precision}\OperatorTok{)} \OperatorTok{\{}
    \KeywordTok{...}
    \DataTypeTok{int}\NormalTok{ indexE }\OperatorTok{=}\NormalTok{ precision}\OperatorTok{.}\FunctionTok{indexOf}\OperatorTok{(}\CharTok{\textquotesingle{}E\textquotesingle{}}\OperatorTok{);}
    \ControlFlowTok{if} \OperatorTok{(}\NormalTok{indexE }\OperatorTok{!=} \OperatorTok{{-}}\DecValTok{1}\OperatorTok{)} \OperatorTok{\{}
\NormalTok{        numDecimals }\OperatorTok{=} \BuiltInTok{Integer}\OperatorTok{.}\FunctionTok{parseInt}\OperatorTok{(}\NormalTok{precision}\OperatorTok{.}\FunctionTok{substring}\OperatorTok{(}\NormalTok{indexE }\OperatorTok{+} \DecValTok{2}\OperatorTok{,}\NormalTok{ indexE }\OperatorTok{+} \DecValTok{3}\OperatorTok{));}
    \OperatorTok{\}} \ControlFlowTok{else} \OperatorTok{\{}
        \DataTypeTok{int}\NormalTok{ decimalPoint }\OperatorTok{=}\NormalTok{ precision}\OperatorTok{.}\FunctionTok{indexOf}\OperatorTok{(}\CharTok{\textquotesingle{}.\textquotesingle{}}\OperatorTok{);}
        \DataTypeTok{int}\NormalTok{ one }\OperatorTok{=}\NormalTok{ precision}\OperatorTok{.}\FunctionTok{indexOf}\OperatorTok{(}\CharTok{\textquotesingle{}1\textquotesingle{}}\OperatorTok{);}
        \ControlFlowTok{if} \OperatorTok{(}\NormalTok{decimalPoint }\OperatorTok{!=} \OperatorTok{{-}}\DecValTok{1} \OperatorTok{\&\&}\NormalTok{ one }\OperatorTok{!=} \OperatorTok{{-}}\DecValTok{1}\OperatorTok{)} \OperatorTok{\{}
\NormalTok{            numDecimals }\OperatorTok{=}\NormalTok{ one }\OperatorTok{{-}}\NormalTok{ decimalPoint}\OperatorTok{;}
        \OperatorTok{\}} \ControlFlowTok{else} \OperatorTok{\{}
\NormalTok{            numDecimals }\OperatorTok{=} \DecValTok{0}\OperatorTok{;}
        \OperatorTok{\}}
    \OperatorTok{\}}
    \KeywordTok{...}
\NormalTok{            roundedStringDecimalPlace }\OperatorTok{=}\NormalTok{ rounded}\OperatorTok{.}\FunctionTok{indexOf}\OperatorTok{(}\CharTok{\textquotesingle{}.\textquotesingle{}}\OperatorTok{);}
            \DataTypeTok{int}\NormalTok{ finalLength }\OperatorTok{=}\NormalTok{ roundedStringDecimalPlace }\OperatorTok{+}\NormalTok{ numDecimals }\OperatorTok{+} \DecValTok{1}\OperatorTok{;}
            \ControlFlowTok{if} \OperatorTok{(}\NormalTok{finalLength }\OperatorTok{\textless{}=}\NormalTok{ rounded}\OperatorTok{.}\FunctionTok{length}\OperatorTok{())} \OperatorTok{\{}
\NormalTok{                rounded }\OperatorTok{=}\NormalTok{ rounded}\OperatorTok{.}\FunctionTok{substring}\OperatorTok{(}\DecValTok{0}\OperatorTok{,}\NormalTok{ finalLength}\OperatorTok{);}
\end{Highlighting}
\end{Shaded}
\paragraph{Cómo arreglarlo}
Calcular el número de decimales a partir del valor numérico de la precisión y no de su texto (por ejemplo, con \texttt{Big\allowbreak{}Decimal.\allowbreak{}value\allowbreak{}Of(\allowbreak{}precision).\allowbreak{}strip\allowbreak{}Trailing\allowbreak{}Zeros(\allowbreak{}).\allowbreak{}scale(\allowbreak{})}, acotado a 0 por abajo), y cuando sea 0 no dejar el punto (cortar en \texttt{rounded\allowbreak{}String\allowbreak{}Decimal\allowbreak{}Place}, sin el \texttt{+\ 1}). Nunca truncar la parte entera.

La misma lectura del texto de la precisión está en \texttt{Util.\allowbreak{}round\allowbreak{}With\allowbreak{}Precision(\allowbreak{}double,\allowbreak{}\ String)} (módulo \texttt{core}, \texttt{Util.\allowbreak{}java:178-202}), que debería usar la misma regla para que el valor redondeado y el mostrado coincidan.

Prueba de regresión: una prueba unitaria de \texttt{convert\allowbreak{}To\allowbreak{}String\allowbreak{}Limit\allowbreak{}Values} con precisiones 0.0, 2.0, 0.5, 1.0, 10.0, 100.0, 0.01 y 1.0E-10 que exija que ningún texto termine en punto, que la parte entera nunca se recorte y que no se lance ninguna excepción; y otra de \texttt{Util.\allowbreak{}round\allowbreak{}With\allowbreak{}Precision} con las mismas precisiones que exija que no quite cifras enteras (1234 no puede pasar a 123) ni lance.

\setcounter{subsection}{169}
\subsection[Tras A, B y C, el nodo de azar nuevo se llama «E», la letra de los sucesos]{Con los nodos A, B y C ya creados, el siguiente nodo de azar se llama «E», la letra reservada para los nodos de suceso}\label{err:170}
\begin{ficha}
\fichakey{Estado}Presente
\fichakey{Módulo}core
\fichakey{Paquete}org.openmarkov.core.model.network
\fichakey{Ficheros}\path{core/src/main/java/org/openmarkov/core/model/network/Util.java:370-405, 429-445} \newline \path{gui/src/main/java/org/openmarkov/gui/window/edition/networkEditorPanel/EditorInputHandler.java:698}
\fichakey{Comprobado}Leyendo el código y ejecutándolo
\fichakey{Esfuerzo}Pequeño (horas)
\fichakey{Decisión de equipo}No
\end{ficha}
\paragraph{Qué pasa}
Al crear un nodo de azar con doble clic en el lienzo o desde el menú contextual, en una red que ya tiene los nodos A, B y C, el nombre propuesto es «E», y el siguiente nodo de suceso pasaría a llamarse «E1». No hay error, pero rompe la convención de que la E es de los nodos de suceso, que \texttt{get\allowbreak{}Next\allowbreak{}Event\allowbreak{}Node\allowbreak{}Name} sigue.

\texttt{Editor\allowbreak{}Input\allowbreak{}Handler.\allowbreak{}create\allowbreak{}Node} pide el nombre a \texttt{Util.\allowbreak{}get\allowbreak{}Next\allowbreak{}Node\allowbreak{}Name} (módulo \texttt{core}). Para un nodo de azar, \texttt{get\allowbreak{}Next\allowbreak{}Chance\allowbreak{}Node\allowbreak{}Name} recorre las letras desde la A buscando la primera libre y quiere saltarse la D (decisiones), la U (utilidades) y la E (sucesos). Pero el salto es un \texttt{if} que avanza una sola letra: desde la C avanza a la D, la reconoce y avanza a la E, y en la siguiente vuelta prueba la E como candidata sin volver a mirar si está reservada.

La U y la D sí se saltan bien (tras la U viene la V, que no está reservada, y a la D solo se llega desde la C). El javadoc dice que las únicas letras que nunca devuelve son la D y la U, lo que tampoco coincide con la intención de saltar la E.

Medido con un pequeño programa de prueba sobre \texttt{Util.\allowbreak{}get\allowbreak{}Next\allowbreak{}Node\allowbreak{}Name}: con \texttt{\{A,\allowbreak{}\ B,\allowbreak{}\ C\}} el nodo de azar se llama «E»; con \texttt{\{A,\allowbreak{}\ B,\allowbreak{}\ C,\allowbreak{}\ E\}} se llama «F».
\paragraph{El código}
core/src/main/java/org/openmarkov/core/model/network/Util.java:390-400

\begin{Shaded}
\begin{Highlighting}[]
            \ControlFlowTok{while} \OperatorTok{(!}\NormalTok{found }\OperatorTok{\&\&} \OperatorTok{(}\NormalTok{letter }\OperatorTok{\textless{}=} \CharTok{\textquotesingle{}Z\textquotesingle{}}\OperatorTok{))} \OperatorTok{\{}
\NormalTok{                name }\OperatorTok{=}\NormalTok{ letter }\OperatorTok{+} \OperatorTok{((}\NormalTok{index }\OperatorTok{\textgreater{}} \DecValTok{0}\OperatorTok{)} \OperatorTok{?} \BuiltInTok{Integer}\OperatorTok{.}\FunctionTok{toString}\OperatorTok{(}\NormalTok{index}\OperatorTok{)} \OperatorTok{:} \StringTok{""}\OperatorTok{);}
                \ControlFlowTok{if} \OperatorTok{(!}\NormalTok{existingNames}\OperatorTok{.}\FunctionTok{contains}\OperatorTok{(}\NormalTok{name}\OperatorTok{))} \OperatorTok{\{}
\NormalTok{                    found }\OperatorTok{=} \KeywordTok{true}\OperatorTok{;}
                \OperatorTok{\}} \ControlFlowTok{else} \OperatorTok{\{}
\NormalTok{                    letter}\OperatorTok{++;}
                    
                    \ControlFlowTok{if} \OperatorTok{((}\NormalTok{letter }\OperatorTok{==} \CharTok{\textquotesingle{}D\textquotesingle{}}\OperatorTok{)} \OperatorTok{||} \OperatorTok{(}\NormalTok{letter }\OperatorTok{==} \CharTok{\textquotesingle{}U\textquotesingle{}}\OperatorTok{)} \CommentTok{/*Events added*/} \OperatorTok{||} \OperatorTok{(}\NormalTok{letter }\OperatorTok{==} \CharTok{\textquotesingle{}E\textquotesingle{}}\OperatorTok{))} \OperatorTok{\{}
\NormalTok{                        letter}\OperatorTok{++;}
                    \OperatorTok{\}}
                \OperatorTok{\}}
            \OperatorTok{\}}
\end{Highlighting}
\end{Shaded}
\paragraph{Cómo arreglarlo}
Cambiar el \texttt{if} por un \texttt{while} que avance mientras la letra sea D, E o U (o comprobar la reserva antes de probar cada candidata, lo que además cubre el caso de empezar por una letra reservada), y corregir el javadoc para que diga las tres letras.

Prueba: \texttt{get\allowbreak{}Next\allowbreak{}Node\allowbreak{}Name(\allowbreak{}CHANCE,\allowbreak{}\ \{A,\allowbreak{}\ B,\allowbreak{}\ C\})} debe dar «F»; \texttt{\{\}} debe dar «A»; \texttt{\{A…T\}} sin D ni E debe dar «V».

% ══════════════════════════════════════════════════════════════════════
\section{Interfaz: menús, diálogos y herramientas}\label{sec:menus}

Los menús, los diálogos generales y las herramientas del menú Tools.

\setcounter{subsection}{170}
\subsection[Abrir una red desde internet apaga «Guardar» para todas las redes]{Abrir una red desde una dirección de internet apaga «Guardar» para todas las redes durante el resto de la sesión, y a la vez deja esa red apuntando a un fichero local inventado}\label{err:171}
\begin{ficha}
\fichakey{Estado}Presente
\fichakey{Módulo}gui
\fichakey{Paquete}org.openmarkov.gui.window
\fichakey{Ficheros}\path{gui/src/main/java/org/openmarkov/gui/window/NetworkFileHandler.java:179-186, 203, 219-222, 246-254, 368-382, 391-402} \newline \path{gui/src/main/java/org/openmarkov/gui/window/MainPanelMenuAssistant.java:90, 278-285, 413, 853-857} \newline \path{gui/src/main/java/org/openmarkov/gui/window/MainPanel.java:557-565}
\fichakey{Comprobado}Leyendo el código
\fichakey{Esfuerzo}Medio (días)
\fichakey{Decisión de equipo}No
\fichakey{Relacionados}\hyperref[nota:177]{nota~177}
\end{ficha}
\paragraph{Qué pasa}
Menú Archivo → «Open network from URL\ldots» (\texttt{Action\allowbreak{}Commands.\allowbreak{}OPEN\_\allowbreak{}NETWORK\_\allowbreak{}URL}, atajo Ctrl+Alt+O) abre una red escribiendo su dirección de internet (URL), por ejemplo \texttt{https:/\allowbreak{}/\allowbreak{}www.\allowbreak{}cisiad.\allowbreak{}uned.\allowbreak{}es/\allowbreak{}redes/\allowbreak{}asia.\allowbreak{}pgmx}. El código está en \texttt{Network\allowbreak{}File\allowbreak{}Handler.\allowbreak{}open\allowbreak{}Network\allowbreak{}Source} (módulo \texttt{gui}). Hay dos defectos enlazados.

\textbf{1. La red abierta por URL recibe como fichero local el camino de la dirección.} Para una dirección, \texttt{network\allowbreak{}File} es \texttt{url.\allowbreak{}get\allowbreak{}File(\allowbreak{})} (línea 189), es decir \texttt{/\allowbreak{}redes/\allowbreak{}asia.\allowbreak{}pgmx}, y se pasa a \texttt{create\allowbreak{}New\allowbreak{}Frame(\allowbreak{}net\allowbreak{}Read\allowbreak{}From,\allowbreak{}\ network\allowbreak{}File)}, que lo guarda con \texttt{set\allowbreak{}Network\allowbreak{}File}. El lector devuelve además un escritor (\texttt{Prob\allowbreak{}Net\allowbreak{}Info.\allowbreak{}writer(\allowbreak{})}, no nulo para \texttt{.\allowbreak{}pgmx} y Elvira).

Con eso, \texttt{save\allowbreak{}Network} (línea 250) considera que la red ya tiene fichero y formato: intenta primero una copia de seguridad en \texttt{/\allowbreak{}redes/\allowbreak{}asia.\allowbreak{}bak} (si falla, solo escribe un aviso por consola) y después escribe la red en \texttt{/\allowbreak{}redes/\allowbreak{}asia.\allowbreak{}pgmx}, en la raíz del disco local. Si esa carpeta no existe, el guardado falla con un error; si existe y se puede escribir, se sobrescribe lo que haya. Hay caminos que llegan a \texttt{save\allowbreak{}Network} sin pasar por el botón «Guardar» del menú:

\begin{itemize}
\tightlist
\item
  el menú contextual de la pestaña (clic derecho → «Save», \texttt{Main\allowbreak{}Panel}, línea 557), que no mira la bandera del punto 2;
\item
  cerrar la pestaña con cambios y contestar «Sí» a guardar (\texttt{network\allowbreak{}Can\allowbreak{}Be\allowbreak{}Closed}, línea 384);
\item
  «Guardar», cuando vuelve a estar encendido al cambiar de pestaña (ver punto 2).
\end{itemize}

Además, al salir con la opción de reabrir la última sesión, \texttt{close\allowbreak{}Application} guarda \texttt{/\allowbreak{}redes/\allowbreak{}asia.\allowbreak{}pgmx} en la lista de ficheros de la sesión, y el siguiente arranque intenta abrirlo como fichero local.

El filtro de «la red ya está abierta» de \texttt{open\allowbreak{}Network\allowbreak{}Source} compara \texttt{network\allowbreak{}File.\allowbreak{}equals(\allowbreak{}editor.\allowbreak{}get\allowbreak{}Network\allowbreak{}File(\allowbreak{}))}, que ya tolera redes sin fichero, y la pregunta de recargar solo aparece si alguna copia está modificada; si no lo está, la pestaña anterior se quita sin preguntar (commit \texttt{a5da86c}). Por el punto 1, dos direcciones distintas con el mismo camino (\texttt{http:/\allowbreak{}/\allowbreak{}a/\allowbreak{}redes/\allowbreak{}asia.\allowbreak{}pgmx} y \texttt{http:/\allowbreak{}/\allowbreak{}b/\allowbreak{}redes/\allowbreak{}asia.\allowbreak{}pgmx}) se consideran la misma red, y abrir la segunda cierra la primera sin preguntar si no tenía cambios.

\textbf{2. Para tapar lo anterior se apaga «Guardar», pero la bandera es de toda la ventana y no se baja nunca.} Tras abrir por URL se llama a \texttt{update\allowbreak{}Options\allowbreak{}Network\allowbreak{}Opened\allowbreak{}URL(\allowbreak{}true)} (línea 225), que pone \texttt{network\allowbreak{}Opened\allowbreak{}URL\ =\ true} en \texttt{Main\allowbreak{}Panel\allowbreak{}Menu\allowbreak{}Assistant} (el ayudante que enciende y apaga opciones de menú; hay uno para toda la ventana principal) y apaga «Guardar» y «Guardar y reabrir». No existe ninguna llamada con \texttt{false} en todo el proyecto.

Cada vez que se hace, deshace o rehace una edición en \textbf{cualquier} red abierta, \texttt{update\allowbreak{}Options\allowbreak{}Network\allowbreak{}Modified} (líneas 279-286) vuelve a aplicar \texttt{set\allowbreak{}Option\allowbreak{}Enabled(\allowbreak{}SAVE\_\allowbreak{}NETWORK,\allowbreak{}\ !network\allowbreak{}Opened\allowbreak{}URL)}, es decir, apaga «Guardar». Resultado: desde que se abre una red por URL, en ninguna red (ni las abiertas después desde disco, ni las nuevas) se puede guardar con el menú, el botón de la barra o Ctrl+S justo después de modificarla, que es cuando se quiere guardar.

Al cambiar de pestaña, \texttt{update\allowbreak{}Options\allowbreak{}Network\allowbreak{}Dependent} (línea 415) enciende «Guardar» según \texttt{get\allowbreak{}Modified(\allowbreak{})} de la red seleccionada, sin mirar la bandera, y la siguiente edición lo vuelve a apagar. «Guardar como» no está afectado. «Guardar y reabrir» vuelve a encenderse al abrir o crear otra red (\texttt{update\allowbreak{}Options\allowbreak{}New\allowbreak{}Network\allowbreak{}Open} enciende \texttt{FILING\_\allowbreak{}ACTION\_\allowbreak{}COMMANDS}).

Todo esto sale de la lectura del código; no se ha ejecutado.
\paragraph{El código}
gui/src/main/java/org/openmarkov/gui/window/NetworkFileHandler.java:187-190, 207, 223-226

\begin{Shaded}
\begin{Highlighting}[]
\BuiltInTok{String}\NormalTok{ networkFile }\OperatorTok{=} \ControlFlowTok{switch} \OperatorTok{(}\NormalTok{source}\OperatorTok{)} \OperatorTok{\{}
    \ControlFlowTok{case}\NormalTok{ NetworkSource}\OperatorTok{.}\FunctionTok{SourceFile}\NormalTok{ sourceFile }\OperatorTok{{-}\textgreater{}}\NormalTok{ sourceFile}\OperatorTok{.}\FunctionTok{file}\OperatorTok{.}\FunctionTok{getAbsolutePath}\OperatorTok{();}
    \ControlFlowTok{case}\NormalTok{ NetworkSource}\OperatorTok{.}\FunctionTok{SourceURL}\NormalTok{ sourceURL }\OperatorTok{{-}\textgreater{}}\NormalTok{ sourceURL}\OperatorTok{.}\FunctionTok{url}\OperatorTok{.}\FunctionTok{getFile}\OperatorTok{();}
\OperatorTok{\};}
\KeywordTok{...}
\NormalTok{NetworkEditorPanel networkPanel }\OperatorTok{=} \KeywordTok{this}\OperatorTok{.}\FunctionTok{createNewFrame}\OperatorTok{(}\NormalTok{netReadFrom}\OperatorTok{,}\NormalTok{ networkFile}\OperatorTok{);}
\KeywordTok{...}
\ControlFlowTok{case}\NormalTok{ NetworkSource}\OperatorTok{.}\FunctionTok{SourceURL}\NormalTok{ sourceURL }\OperatorTok{{-}\textgreater{}} \OperatorTok{\{}
    \CommentTok{//LastOpenFiles.setLastFileName(sourceURL.url.getFile());}
\NormalTok{    mainPanel}\OperatorTok{.}\FunctionTok{getMainPanelMenuAssistant}\OperatorTok{().}\FunctionTok{updateOptionsNetworkOpenedURL}\OperatorTok{(}\KeywordTok{true}\OperatorTok{);}
\OperatorTok{\}}
\end{Highlighting}
\end{Shaded}

gui/src/main/java/org/openmarkov/gui/window/MainPanelMenuAssistant.java:279-286, 859-863

\begin{Shaded}
\begin{Highlighting}[]
\KeywordTok{public} \DataTypeTok{void} \FunctionTok{updateOptionsNetworkModified}\OperatorTok{(}\DataTypeTok{boolean}\NormalTok{ canUndo}\OperatorTok{,} \DataTypeTok{boolean}\NormalTok{ canRedo}\OperatorTok{)} \OperatorTok{\{}
    \KeywordTok{...}
    \CommentTok{// If the network has been opened from a URL the save button has to remain disabled}
    \FunctionTok{setOptionEnabled}\OperatorTok{(}\NormalTok{ActionCommands}\OperatorTok{.}\FunctionTok{SAVE\_NETWORK}\OperatorTok{,} \OperatorTok{!}\NormalTok{networkOpenedURL}\OperatorTok{);}
\OperatorTok{\}}
\KeywordTok{...}
\KeywordTok{public} \DataTypeTok{void} \FunctionTok{updateOptionsNetworkOpenedURL}\OperatorTok{(}\DataTypeTok{boolean}\NormalTok{ networkOpenedURL}\OperatorTok{)} \OperatorTok{\{}
    \KeywordTok{this}\OperatorTok{.}\FunctionTok{networkOpenedURL} \OperatorTok{=}\NormalTok{ networkOpenedURL}\OperatorTok{;}
    \FunctionTok{setOptionEnabled}\OperatorTok{(}\NormalTok{ActionCommands}\OperatorTok{.}\FunctionTok{SAVE\_NETWORK}\OperatorTok{,} \OperatorTok{!}\NormalTok{networkOpenedURL}\OperatorTok{);}
    \FunctionTok{setOptionEnabled}\OperatorTok{(}\NormalTok{ActionCommands}\OperatorTok{.}\FunctionTok{SAVE\_OPEN\_NETWORK}\OperatorTok{,} \OperatorTok{!}\NormalTok{networkOpenedURL}\OperatorTok{);}
\OperatorTok{\}}
\end{Highlighting}
\end{Shaded}
\paragraph{Cómo arreglarlo}
Una red traída de una dirección no tiene fichero local: \texttt{create\allowbreak{}New\allowbreak{}Frame} debe recibir \texttt{null} como fichero para ella (el nombre de la red puede seguir saliendo del último tramo de la dirección, como hoy en la línea 185). Con eso \texttt{save\allowbreak{}Network} cae en «Guardar como» y pregunta dónde, igual que con una red nueva, por cualquiera de los caminos (menú, barra, Ctrl+S, menú contextual de la pestaña, cerrar con cambios). Después se puede borrar la bandera \texttt{network\allowbreak{}Opened\allowbreak{}URL}, \texttt{update\allowbreak{}Options\allowbreak{}Network\allowbreak{}Opened\allowbreak{}URL} y la línea de \texttt{update\allowbreak{}Options\allowbreak{}Network\allowbreak{}Modified} que la aplica.

Sitios donde la misma regla («una red sin fichero no se guarda en un camino inventado») tiene que cumplirse:

\begin{itemize}
\tightlist
\item
  \texttt{Network\allowbreak{}File\allowbreak{}Handler.\allowbreak{}save\allowbreak{}Network};
\item
  \texttt{save\allowbreak{}Open\allowbreak{}Network} (líneas 283-292, que usa \texttt{get\allowbreak{}Network\allowbreak{}File(\allowbreak{})} para reabrir);
\item
  \texttt{close\allowbreak{}Application}, que no debe guardar en la lista de la sesión un fichero nulo o inventado;
\item
  el filtro de «la red ya está abierta» de \texttt{open\allowbreak{}Network\allowbreak{}Source}, que debería comparar la dirección completa para las redes traídas por URL (o no comparar).
\end{itemize}

Si algún día se separa de la ventana un documento de red con su fichero (ver \hyperref[nota:177]{nota~177}, una propuesta de diseño sin nada que corregir), esta distinción cae ahí de forma natural.

Prueba: abrir una red por URL (se puede servir un \texttt{.\allowbreak{}pgmx} con un servidor local o un \texttt{file:} URL si el lector lo admite), comprobar que \texttt{get\allowbreak{}Network\allowbreak{}File(\allowbreak{})} es nulo, que tras modificar otra red abierta desde disco «Guardar» está encendido, y que guardar la red traída pide un destino.

\setcounter{subsection}{171}
\subsection[«Edits history» deja un núcleo al cien por cien y puede colgar el programa]{La herramienta «Edits history» deja un núcleo del procesador al cien por cien hasta que se cierra el programa, y sus botones de deshacer y rehacer pueden colgar la aplicación}\label{err:172}
\begin{ficha}
\fichakey{Estado}Presente
\fichakey{Módulo}full
\fichakey{Paquete}org.openmarkov.full
\fichakey{Ficheros}\path{full/src/main/java/org/openmarkov/full/EditsHistoryPlugin.java:36-40} \newline \path{full/src/main/java/org/openmarkov/full/EditsHistoryPlugin.java:42-143} \newline \path{core/src/main/java/org/openmarkov/core/action/base/PNESupport.java:104-116} \newline \path{core/src/main/java/org/openmarkov/core/action/base/PNESupport.java:142-154}
\fichakey{Comprobado}Leyendo el código
\fichakey{Esfuerzo}Pequeño (horas)
\fichakey{Decisión de equipo}\textcolor{decisionrojo}{\textbf{Sí.}} Si esta herramienta, que su documentación llama «de desarrollo», debe aparecer en el menú Herramientas para todos los usuarios
\end{ficha}
\paragraph{Qué pasa}
\texttt{Edits\allowbreak{}History\allowbreak{}Plugin} (módulo \texttt{full}) añade al menú Herramientas la entrada «Edits history»: abre una ventana no modal con un botón por cada edición del historial de la red activa (hechas y deshechas); pulsar uno deshace o rehace hasta esa edición. Su javadoc la llama herramienta de desarrollo, pero \texttt{Tool\allowbreak{}Plugin\allowbreak{}Manager} la carga como cualquier otro complemento y aparece para todos.

\textbf{1. Espera activa sin fin.} Desde que se abre la herramienta, un núcleo del procesador queda ocupado entero hasta que se sale del programa, aunque se cierre la ventana; cada vez que se vuelve a abrir, se ocupa otro núcleo más.

\texttt{action(\allowbreak{})} (líneas 42-143) lanza un \texttt{new\ Thread(\allowbreak{}(\allowbreak{})\ -\textgreater{}\ \{\ while\ (\allowbreak{}true)\ \{\ .\allowbreak{}.\allowbreak{}.\allowbreak{}\ \}\ \})}. En cada vuelta pide la red activa, sus listas de ediciones hechas y deshechas, las compara con las de la vuelta anterior y, solo si cambian, rehace los botones. No hay ninguna espera entre vueltas (ni \texttt{sleep}, ni bloqueo, ni aviso por oyente). Tampoco comprueba si la ventana sigue abierta: al cerrarla (se oculta) el hilo continúa, y cada apertura lanza un hilo nuevo. Si no hay red abierta, la excepción se recoge y se vuelve a empezar igual de deprisa. Además, ese hilo modifica componentes Swing (quita y añade botones, \texttt{pack(\allowbreak{})}) fuera del hilo de la interfaz, lo que Swing no admite.

\textbf{2. Botones que pueden no terminar.} Si el botón pulsado ha quedado desfasado respecto al historial, la interfaz se queda congelada y hay que matar el programa; antes de colgarse, además, deshace todo el historial de esa red.

Cada botón de una edición hecha ejecuta, en el hilo de la interfaz, \texttt{while\ (\allowbreak{}!current\allowbreak{}Prob\allowbreak{}Net.\allowbreak{}get\allowbreak{}PNESupport(\allowbreak{}).\allowbreak{}undo(\allowbreak{}).\allowbreak{}contains(\allowbreak{}edit\allowbreak{}And\allowbreak{}Done.\allowbreak{}edit))\ \{\ \}} (líneas 88-92), y los de ediciones deshechas lo mismo con \texttt{redo(\allowbreak{})} (líneas 110-114). \texttt{PNESupport.\allowbreak{}undo(\allowbreak{})} devuelve una lista vacía cuando ya no queda nada que deshacer. Si la edición del botón no está en el historial actual ---porque el historial cambió y el botón es de la lista anterior (el refresco lo hace otro hilo sin coordinación con los clics), o porque la red capturada en \texttt{current\allowbreak{}Prob\allowbreak{}Net} ya no es la del historial que se muestra---, el bucle pide deshacer una y otra vez, recibe listas vacías y no sale nunca.

No se ha ejecutado (la ventana es un \texttt{JDialog} y no se puede construir sin pantalla). El consumo de procesador del punto 1 se sigue directamente del código; el cuelgue del punto 2 necesita que el botón quede desfasado respecto al historial.
\paragraph{El código}
full/src/main/java/org/openmarkov/full/EditsHistoryPlugin.java:56-60, 74, 86-92

\begin{Shaded}
\begin{Highlighting}[]
\KeywordTok{new} \BuiltInTok{Thread}\OperatorTok{(()} \OperatorTok{{-}\textgreater{}} \OperatorTok{\{}
    \ControlFlowTok{while} \OperatorTok{(}\KeywordTok{true}\OperatorTok{)} \OperatorTok{\{}
\NormalTok{        ProbNet currentProbNet }\OperatorTok{=}\NormalTok{ MainPanel}\OperatorTok{.}\FunctionTok{getCurrentProbNet}\OperatorTok{();}
        \ControlFlowTok{try} \OperatorTok{\{}
            \KeywordTok{...}
            \ControlFlowTok{if} \OperatorTok{(!}\NormalTok{newEdits}\OperatorTok{.}\FunctionTok{equals}\OperatorTok{(}\NormalTok{edits}\OperatorTok{))} \OperatorTok{\{}
                \KeywordTok{...}
                \ControlFlowTok{if} \OperatorTok{(}\NormalTok{editAndDone}\OperatorTok{.}\FunctionTok{done}\OperatorTok{)} \OperatorTok{\{}
\NormalTok{                    editButton}\OperatorTok{.}\FunctionTok{setIcon}\OperatorTok{(}\NormalTok{IconBind}\OperatorTok{.}\FunctionTok{UNDO\_ENABLED}\OperatorTok{.}\FunctionTok{icon}\OperatorTok{());}
\NormalTok{                    editButton}\OperatorTok{.}\FunctionTok{addActionListener}\OperatorTok{(}\NormalTok{e }\OperatorTok{{-}\textgreater{}} \OperatorTok{\{}
                        \ControlFlowTok{while} \OperatorTok{(!}\NormalTok{currentProbNet}\OperatorTok{.}\FunctionTok{getPNESupport}\OperatorTok{().}\FunctionTok{undo}\OperatorTok{().}\FunctionTok{contains}\OperatorTok{(}\NormalTok{editAndDone}\OperatorTok{.}\FunctionTok{edit}\OperatorTok{))} \OperatorTok{\{}

                        \OperatorTok{\}}
                    \OperatorTok{\});}
\end{Highlighting}
\end{Shaded}
\paragraph{Cómo arreglarlo}
Antes de arreglarla hay que decidir si la herramienta debe estar en el menú de todos los usuarios o solo en un modo de desarrollo.

Refresco: sustituir el hilo por un oyente de ediciones (\texttt{PNEdit\allowbreak{}Listener} registrado en el \texttt{PNESupport} de la red, más un oyente de cambio de pestaña para la red activa) que rehaga los botones en el hilo de la interfaz (\texttt{Swing\allowbreak{}Utilities.\allowbreak{}invoke\allowbreak{}Later}), y darlo de baja al cerrar la ventana. Si se prefiere un sondeo, un \texttt{javax.\allowbreak{}swing.\allowbreak{}Timer} con un intervalo y parado al cerrar.

Botones: antes de tocar el historial, buscar la posición de la edición en \texttt{get\allowbreak{}Done\allowbreak{}Edits(\allowbreak{})} (o \texttt{get\allowbreak{}Undone\allowbreak{}Edits(\allowbreak{})}) de la red actual; si no está, no hacer nada (o avisar); si está, llamar a \texttt{undo(\allowbreak{})}/\texttt{redo(\allowbreak{})} exactamente ese número de veces. Así el bucle termina siempre.

Prueba: sin ventana, extraer a un método el cálculo de «cuántos pasos hasta la edición X» sobre un \texttt{PNESupport} con tres ediciones y comprobar que deshace hasta la segunda en dos pasos y que con una edición que no está en el historial no hace nada.

\setcounter{subsection}{172}
\subsection[Aceptar los dominios estándar sin marcar ninguno lanza una excepción]{Aceptar el diálogo de dominios estándar sin marcar ninguno lanza una excepción inesperada}\label{err:173}
\begin{ficha}
\fichakey{Estado}Presente
\fichakey{Módulo}gui
\fichakey{Paquete}org.openmarkov.gui.dialog.node
\fichakey{Ficheros}\path{gui/src/main/java/org/openmarkov/gui/dialog/node/StandardDomainPanel.java:33-62} \newline \path{gui/src/main/java/org/openmarkov/gui/dialog/node/StandardDomainsDialog.java:56-76} \newline \path{gui/src/main/java/org/openmarkov/gui/dialog/node/NodeDomainValuesTablePanel.java:1070-1112} \newline \path{gui/src/main/java/org/openmarkov/gui/dialog/common/OkCancelDialog.java:67-80}
\fichakey{Comprobado}Leyendo el código y ejecutándolo
\fichakey{Esfuerzo}Pequeño (horas)
\fichakey{Decisión de equipo}No
\fichakey{Relacionados}\hyperref[err:24]{error~24}, \hyperref[err:174]{error~174}, \hyperref[err:180]{error~180}
\end{ficha}
\paragraph{Qué pasa}
En el diálogo de propiedades de un nodo con variable de estados finitos o discretizada, la pestaña de dominio tiene el botón «Standard domains», que abre una ventana con un botón de radio por cada dominio. Ninguno sale marcado, y el botón «OK» está activo desde el principio (y es el botón por defecto, así que también responde a la tecla Intro).

Si el usuario pulsa «OK» sin marcar nada, sale el diálogo de «excepción inesperada» que pide enviarla a los desarrolladores. La red no cambia y el diálogo del nodo sigue abierto. Cerrar la ventana con la X equivale a «Cancel» y no falla.

\texttt{Node\allowbreak{}Domain\allowbreak{}Values\allowbreak{}Table\allowbreak{}Panel.\allowbreak{}action\allowbreak{}Performed\allowbreak{}Standard\allowbreak{}Domain} abre \texttt{Standard\allowbreak{}Domains\allowbreak{}Dialog}, cuyo \texttt{Standard\allowbreak{}Domain\allowbreak{}Panel} pone un botón de radio por cada dominio que ofrece \texttt{GUIDefault\allowbreak{}States.\allowbreak{}get\allowbreak{}All\allowbreak{}Domains(\allowbreak{})} (los dominios propios del usuario en preferencias, seguidos de los seis de \texttt{Standard\allowbreak{}Domain}). \texttt{chosen\allowbreak{}States} solo se asigna en el oyente de cada botón, así que vale \texttt{null} hasta que el usuario marca uno.

Al pulsar «OK» sin marcar nada, \texttt{request\allowbreak{}Values(\allowbreak{})} devuelve \texttt{Ok} y \texttt{action\allowbreak{}Performed\allowbreak{}Standard\allowbreak{}Domain} llama a \texttt{get\allowbreak{}Chosen\allowbreak{}States(\allowbreak{})}, que ejecuta \texttt{Collections.\allowbreak{}unmodifiable\allowbreak{}List(\allowbreak{}null)} y lanza \texttt{Null\allowbreak{}Pointer\allowbreak{}Exception}. Como no es \texttt{Do\allowbreak{}Edit\allowbreak{}Exception}, \texttt{action\allowbreak{}Performed} no la captura y llega al gestor global.

Medido con un pequeño programa de prueba sobre el panel real: los seis botones (\texttt{absent\ -\ present}, \texttt{no\ -\ yes}, \texttt{negative\ -\ positive}, \texttt{absent\ -\ mild\ -\ moderate\ -\ severe}, \texttt{low\ -\ medium\ -\ high}, \texttt{not\ happened\ -\ happened}) salen con \texttt{selected=false}, y \texttt{get\allowbreak{}Chosen\allowbreak{}States(\allowbreak{})} lanza \texttt{Null\allowbreak{}Pointer\allowbreak{}Exception:\ Cannot\ invoke\ "java.\allowbreak{}util.\allowbreak{}List.\allowbreak{}get\allowbreak{}Class(\allowbreak{})"\ because\ "list"\ is\ null}; tras marcar el segundo devuelve \texttt{{[}no,\allowbreak{}\ yes{]}}.

El diálogo hermano de criterios estándar de la red falla igual al aceptar sin marcar nada, y además borra antes los criterios: está en \hyperref[err:174]{error~174}.
\paragraph{El código}
gui/src/main/java/org/openmarkov/gui/dialog/node/StandardDomainPanel.java:43-62

\begin{Shaded}
\begin{Highlighting}[]
\ControlFlowTok{for} \OperatorTok{(}\BuiltInTok{List}\OperatorTok{\textless{}}\BuiltInTok{String}\OperatorTok{\textgreater{}}\NormalTok{ defaultState }\OperatorTok{:}\NormalTok{ states}\OperatorTok{)} \OperatorTok{\{}
    \BuiltInTok{JRadioButton}\NormalTok{ radioButton }\OperatorTok{=} \KeywordTok{new} \BuiltInTok{JRadioButton}\OperatorTok{(}\NormalTok{GUIDefaultStates}\OperatorTok{.}\FunctionTok{getString}\OperatorTok{(}\NormalTok{defaultState}\OperatorTok{));}
\NormalTok{    radioButtons}\OperatorTok{.}\FunctionTok{add}\OperatorTok{(}\NormalTok{radioButton}\OperatorTok{);}
\NormalTok{    buttonGroup}\OperatorTok{.}\FunctionTok{add}\OperatorTok{(}\NormalTok{radioButton}\OperatorTok{);}
\NormalTok{    radioButton}\OperatorTok{.}\FunctionTok{addItemListener}\OperatorTok{(}\NormalTok{e }\OperatorTok{{-}\textgreater{}}\NormalTok{ chosenStates }\OperatorTok{=}\NormalTok{ defaultState}\OperatorTok{);}
    \FunctionTok{add}\OperatorTok{(}\NormalTok{radioButton}\OperatorTok{,} \BuiltInTok{BorderLayout}\OperatorTok{.}\FunctionTok{CENTER}\OperatorTok{);}
\OperatorTok{\}}
\KeywordTok{...}
\KeywordTok{public} \BuiltInTok{List}\OperatorTok{\textless{}}\BuiltInTok{String}\OperatorTok{\textgreater{}} \FunctionTok{getChosenStates}\OperatorTok{()} \OperatorTok{\{}
    \ControlFlowTok{return} \BuiltInTok{Collections}\OperatorTok{.}\FunctionTok{unmodifiableList}\OperatorTok{(}\KeywordTok{this}\OperatorTok{.}\FunctionTok{chosenStates}\OperatorTok{);}
\OperatorTok{\}}

\KeywordTok{private} \BuiltInTok{List}\OperatorTok{\textless{}}\BuiltInTok{String}\OperatorTok{\textgreater{}}\NormalTok{ chosenStates}\OperatorTok{;}
\end{Highlighting}
\end{Shaded}

gui/src/main/java/org/openmarkov/gui/dialog/node/NodeDomainValuesTablePanel.java:1088-1093

\begin{Shaded}
\begin{Highlighting}[]
\ControlFlowTok{if} \OperatorTok{(}\NormalTok{standardDomainDialog}\OperatorTok{.}\FunctionTok{requestValues}\OperatorTok{()} \OperatorTok{!=}\NormalTok{ OkCancelDialog}\OperatorTok{.}\FunctionTok{ChosenOption}\OperatorTok{.}\FunctionTok{Ok}\OperatorTok{)} \OperatorTok{\{}
    \ControlFlowTok{return}\OperatorTok{;}
\OperatorTok{\}}
\DataTypeTok{var}\NormalTok{ chosenStates }\OperatorTok{=}\NormalTok{ standardDomainDialog}\OperatorTok{.}\FunctionTok{getJPanelStandardDomains}\OperatorTok{().}\FunctionTok{getChosenStates}\OperatorTok{();}

\BuiltInTok{State}\OperatorTok{[]}\NormalTok{ newStates }\OperatorTok{=}\NormalTok{ chosenStates}\OperatorTok{.}\FunctionTok{stream}\OperatorTok{().}\FunctionTok{map}\OperatorTok{(}\BuiltInTok{State}\OperatorTok{::}\KeywordTok{new}\OperatorTok{).}\FunctionTok{toArray}\OperatorTok{(}\BuiltInTok{State}\OperatorTok{[]::}\KeywordTok{new}\OperatorTok{);}
\end{Highlighting}
\end{Shaded}
\paragraph{Cómo arreglarlo}
Dos opciones equivalentes:

\begin{itemize}
\tightlist
\item
  desactivar «OK» hasta que se marque un botón (el oyente de cada botón lo activa);
\item
  hacer que \texttt{Standard\allowbreak{}Domains\allowbreak{}Dialog.\allowbreak{}do\allowbreak{}Ok\allowbreak{}Click\allowbreak{}Before\allowbreak{}Hide} devuelva \texttt{false} si no hay nada marcado, para que el diálogo no se cierre.
\end{itemize}

Conviene que el oyente asigne \texttt{chosen\allowbreak{}States} solo cuando el botón pasa a seleccionado (hoy el \texttt{Item\allowbreak{}Listener} también se dispara al deseleccionar, aunque asigna su propio dominio y el siguiente lo corrige).

Si se cambia el panel, conviene aplicar la misma regla a \texttt{Standard\allowbreak{}Criteria\allowbreak{}Dialog}/\texttt{Standard\allowbreak{}Criteria\allowbreak{}Panel}, que ofrecen la misma elección con el mismo defecto (\hyperref[err:174]{error~174}).

Prueba de regresión: construir \texttt{Standard\allowbreak{}Domains\allowbreak{}Dialog} sin mostrarlo, comprobar que «OK» está desactivado (o que \texttt{do\allowbreak{}Ok\allowbreak{}Click\allowbreak{}Before\allowbreak{}Hide} devuelve \texttt{false}) mientras no hay selección, y que tras marcar un botón \texttt{get\allowbreak{}Chosen\allowbreak{}States(\allowbreak{})} devuelve sus estados.

\setcounter{subsection}{173}
\subsection[Aceptar los criterios estándar sin marcar nada borra los criterios de la red]{Aceptar «Standard criteria» sin marcar ninguna opción borra los criterios de decisión de la red, sale como error inesperado y la red deja de marcarse como modificada}\label{err:174}
\begin{ficha}
\fichakey{Estado}Presente
\fichakey{Módulo}gui
\fichakey{Paquete}org.openmarkov.gui.dialog.network
\fichakey{Ficheros}\path{gui/src/main/java/org/openmarkov/gui/dialog/network/StandardCriteriaDialog.java:82-112} \newline \path{gui/src/main/java/org/openmarkov/gui/dialog/network/StandardCriteriaDialog.java:126-136} \newline \path{gui/src/main/java/org/openmarkov/gui/dialog/network/StandardCriteriaPanel.java:38-57} \newline \path{gui/src/main/java/org/openmarkov/gui/dialog/network/DecisionCriteriaTablePanel.java:178-195} \newline \path{gui/src/main/java/org/openmarkov/gui/dialog/network/NetworkPropertiesDialog.java:80-84} \newline \path{gui/src/main/java/org/openmarkov/gui/dialog/network/NetworkPropertiesDialog.java:269-282} \newline \path{gui/src/main/java/org/openmarkov/gui/dialog/common/OkCancelDialog.java:67-78} \newline \path{core/src/main/java/org/openmarkov/core/action/base/EditsHistoryStacker.java:48-60} \newline \path{core/src/main/java/org/openmarkov/core/action/base/PNESupport.java:286-295} \newline \path{gui/src/main/java/org/openmarkov/gui/window/MainPanelMenuAssistant.java:415} \newline \path{gui/src/main/java/org/openmarkov/gui/window/NetworkFileHandler.java:372-374}
\fichakey{Comprobado}Leyendo el código y ejecutándolo
\fichakey{Esfuerzo}Pequeño (horas)
\fichakey{Decisión de equipo}No
\fichakey{Relacionados}\hyperref[err:35]{error~35}, \hyperref[err:59]{error~59}, \hyperref[err:173]{error~173}, \hyperref[err:180]{error~180}
\end{ficha}
\paragraph{Qué pasa}
En las propiedades de una red que no es solo de azar (un diagrama de influencia, una red de análisis de decisiones\ldots), la pestaña de criterios de decisión tiene el botón «Standard criteria». Abre una ventana con seis botones de radio: «Cost-benefit» en €, £ o \$, y «Cost-effectiveness» en €/QALY, £/QALY o \$/QALY (QALY: años de vida ajustados por calidad). Ninguno sale marcado y «OK» está activo.

Si el usuario pulsa «OK» sin marcar nada, sale el diálogo de error inesperado, pero antes de fallar el diálogo ya ha quitado todos los criterios de la red. Lo que pasa después depende de lo que haga el usuario:

\begin{itemize}
\tightlist
\item
  La ventana de criterios estándar sigue abierta, porque la excepción impide cerrarla, y la tabla de la pestaña sigue enseñando los criterios viejos, porque solo se refresca cuando esa ventana termina con «OK» (por lectura).
\item
  Si cierra esa ventana y acepta las propiedades de la red, la red se queda sin ningún criterio. Si se guarda y se reabre, el lector le pone el criterio por omisión «---» y sus nodos de utilidad se quedan sin criterio. Lo que ocurre al evaluar una red sin criterios está en \hyperref[err:59]{error~59}.
\item
  Si cancela las propiedades de la red, los criterios vuelven.
\item
  En los dos casos la red se queda, para el resto de la sesión, con una sub-historia de ediciones abierta que nadie cierra, y desde ese momento el indicador de «red modificada» deja de actualizarse. Si la red no estaba modificada, ninguna edición posterior la marca: «Save» sigue apagado (\texttt{Main\allowbreak{}Panel\allowbreak{}Menu\allowbreak{}Assistant}, línea 415) y cerrar la red no pregunta si se quiere guardar (\texttt{Network\allowbreak{}File\allowbreak{}Handler.\allowbreak{}network\allowbreak{}Can\allowbreak{}Be\allowbreak{}Closed}, líneas 372-374), así que el trabajo posterior se puede perder sin aviso (por lectura; el indicador congelado sí está medido).
\item
  Si, tras la excepción, marca una opción y vuelve a pulsar «OK», se añaden los criterios elegidos, pero la sub-historia del primer intento sigue abierta.
\end{itemize}

Por qué pasa: \texttt{Standard\allowbreak{}Criteria\allowbreak{}Dialog.\allowbreak{}do\allowbreak{}Ok\allowbreak{}Click\allowbreak{}Before\allowbreak{}Hide} (módulo \texttt{gui}) abre una sub-historia, quita todos los criterios con \texttt{Decision\allowbreak{}Criteria\allowbreak{}Edit(\allowbreak{}REMOVE)} y solo después llama a \texttt{get\allowbreak{}Default\allowbreak{}Criteria}, que busca el botón marcado. Si no hay ninguno, \texttt{selected\allowbreak{}Default\allowbreak{}Criteria} vale \texttt{null} y \texttt{selected\allowbreak{}Default\allowbreak{}Criteria.\allowbreak{}equals(\allowbreak{}.\allowbreak{}.\allowbreak{}.\allowbreak{})} (línea 136) lanza \texttt{Null\allowbreak{}Pointer\allowbreak{}Exception}. \texttt{Ok\allowbreak{}Cancel\allowbreak{}Dialog} vuelve a lanzar las excepciones de tiempo de ejecución, así que la ventana no se cierra y la excepción llega al manejador global. \texttt{close\allowbreak{}Sub\allowbreak{}Edit\allowbreak{}History} (línea 109) no se ejecuta: la sub-historia queda encima de la que abrió el diálogo de propiedades, y al aceptar o cancelar este se cierra la de dentro y la suya queda abierta. \texttt{PNESupport.\allowbreak{}on\allowbreak{}Edit\allowbreak{}Changes} no recalcula el indicador de modificada mientras la historia activa no sea la principal.

Medido invocando el \texttt{do\allowbreak{}Ok\allowbreak{}Click\allowbreak{}Before\allowbreak{}Hide} real con un \texttt{Standard\allowbreak{}Criteria\allowbreak{}Panel} real y sin construir la ventana (un \texttt{JDialog} no se puede construir sin pantalla), sobre \texttt{0\_\allowbreak{}miscellaneous/\allowbreak{}networks/\allowbreak{}id/\allowbreak{}ID-arthronet-ce.\allowbreak{}pgmx} (criterios «coste» y «efectividad») y con la sub-historia que abre el diálogo de propiedades:

\begin{itemize}
\tightlist
\item
  los seis botones salen sin marcar; el método lanza \texttt{Null\allowbreak{}Pointer\allowbreak{}Exception:\ Cannot\ invoke\ "String.\allowbreak{}equals(\allowbreak{}Object)"\ because\ "selected\allowbreak{}Default\allowbreak{}Criteria"\ is\ null} en la línea 136, y la red queda con la lista de criterios vacía y dos sub-historias abiertas;
\item
  tras aceptar las propiedades (\texttt{close\allowbreak{}Sub\allowbreak{}Edit\allowbreak{}History}): sin criterios, una sub-historia abierta y la red sin marcar como modificada; el fichero guardado no lleva \texttt{\textless{}Decision\allowbreak{}Criteria\textgreater{}} y al reabrirlo la red tiene el criterio «---» y sus nodos de utilidad ninguno;
\item
  tras cancelar las propiedades (\texttt{cancel\allowbreak{}Last\allowbreak{}Sub\allowbreak{}Edit\allowbreak{}History}): vuelven «coste» y «efectividad», y sigue una sub-historia abierta;
\item
  con esa sub-historia abierta, añadir después un criterio deja la red sin marcar como modificada; el mismo cambio con el historial bien cerrado la marca.
\end{itemize}
\paragraph{El código}
gui/src/main/java/org/openmarkov/gui/dialog/network/StandardCriteriaDialog.java:82-109

\begin{Shaded}
\begin{Highlighting}[]
\AttributeTok{@Override} \KeywordTok{protected} \DataTypeTok{boolean} \FunctionTok{doOkClickBeforeHide}\OperatorTok{()} \KeywordTok{throws}\NormalTok{ DoEditException }\OperatorTok{\{}
    \CommentTok{// This must be a single operation}
\NormalTok{    probNet}\OperatorTok{.}\FunctionTok{getPNESupport}\OperatorTok{().}\FunctionTok{openNewSubEditHistory}\OperatorTok{();}
    \CommentTok{// Remove all criterion in the probNet with edits}
    \ControlFlowTok{while} \OperatorTok{(!}\NormalTok{probNet}\OperatorTok{.}\FunctionTok{getDecisionCriteria}\OperatorTok{().}\FunctionTok{isEmpty}\OperatorTok{())} \OperatorTok{\{}
\NormalTok{        Criterion criterionToBeDeleted }\OperatorTok{=}\NormalTok{ probNet}\OperatorTok{.}\FunctionTok{getDecisionCriteria}\OperatorTok{().}\FunctionTok{get}\OperatorTok{(}\DecValTok{0}\OperatorTok{);}
\NormalTok{        DecisionCriteriaEdit edit }\OperatorTok{=} \KeywordTok{new} \FunctionTok{DecisionCriteriaEdit}\OperatorTok{(}\NormalTok{probNet}\OperatorTok{,}\NormalTok{ StateAction}\OperatorTok{.}\FunctionTok{REMOVE}\OperatorTok{,}\NormalTok{ criterionToBeDeleted}\OperatorTok{,} \KeywordTok{null}\OperatorTok{);}
\NormalTok{        edit}\OperatorTok{.}\FunctionTok{executeEdit}\OperatorTok{();}
    \OperatorTok{\}}
    \CommentTok{//Get the list with the new criteria}
    \BuiltInTok{List}\OperatorTok{\textless{}}\NormalTok{Criterion}\OperatorTok{\textgreater{}}\NormalTok{ defaultCriteria }\OperatorTok{=} \FunctionTok{getDefaultCriteria}\OperatorTok{();}
    \KeywordTok{...}
\NormalTok{    probNet}\OperatorTok{.}\FunctionTok{getPNESupport}\OperatorTok{().}\FunctionTok{closeSubEditHistory}\OperatorTok{();}
\end{Highlighting}
\end{Shaded}

gui/src/main/java/org/openmarkov/gui/dialog/network/StandardCriteriaDialog.java:129-136

\begin{Shaded}
\begin{Highlighting}[]
\BuiltInTok{String}\NormalTok{ selectedDefaultCriteria }\OperatorTok{=} \KeywordTok{null}\OperatorTok{;}
\ControlFlowTok{for} \OperatorTok{(}\BuiltInTok{JRadioButton}\NormalTok{ radioButton }\OperatorTok{:}\NormalTok{ standardCriteriaPanel}\OperatorTok{.}\FunctionTok{getRadioButtons}\OperatorTok{())} \OperatorTok{\{}
    \ControlFlowTok{if} \OperatorTok{(}\NormalTok{radioButton}\OperatorTok{.}\FunctionTok{isSelected}\OperatorTok{())} \OperatorTok{\{}
\NormalTok{        selectedDefaultCriteria }\OperatorTok{=}\NormalTok{ radioButton}\OperatorTok{.}\FunctionTok{getText}\OperatorTok{();}
    \OperatorTok{\}}
\OperatorTok{\}}

\ControlFlowTok{if} \OperatorTok{(}\NormalTok{selectedDefaultCriteria}\OperatorTok{.}\FunctionTok{equals}\OperatorTok{(}\NormalTok{stringDatabase}\OperatorTok{.}\FunctionTok{getString}\OperatorTok{(}\StringTok{"defaultCriteria.costBenefitEuros.Text"}\OperatorTok{)))} \OperatorTok{\{}
\end{Highlighting}
\end{Shaded}
\paragraph{Cómo arreglarlo}
Comprobar que hay una opción marcada antes de tocar la red. Dos maneras, equivalentes para el usuario:

\begin{itemize}
\tightlist
\item
  desactivar «OK» hasta que se marque un botón;
\item
  que \texttt{do\allowbreak{}Ok\allowbreak{}Click\allowbreak{}Before\allowbreak{}Hide} devuelva \texttt{false}, sin abrir la sub-historia, si no hay nada marcado.
\end{itemize}

Además, conviene que la sub-historia que abre este método se cierre también si algo falla entre \texttt{open\allowbreak{}New\allowbreak{}Sub\allowbreak{}Edit\allowbreak{}History} y \texttt{close\allowbreak{}Sub\allowbreak{}Edit\allowbreak{}History} (con \texttt{try}/\texttt{finally}, cancelándola), para que un fallo del diálogo no deje la red a medias ni congele el indicador de modificada.

El diálogo hermano de dominios estándar tiene el mismo problema de selección (\hyperref[err:173]{error~173}); conviene resolver los dos con la misma regla. Los dos se abren además sin título (\hyperref[err:180]{error~180}).

Prueba de regresión: con el panel sin marcar, aceptar no debe cambiar la lista de criterios, y tras aceptar o cancelar las propiedades de la red la historia activa debe ser la principal; marcando «Cost-effectiveness (€/QALY)», la red debe quedar con los criterios «Cost» y «Effectiveness».

\setcounter{subsection}{174}
\subsection[«OpenMarkov configuration» edita unas preferencias que el programa no usa]{«OpenMarkov configuration» enseña y edita un almacén de preferencias que el programa ya no usa}\label{err:175}
\begin{ficha}
\fichakey{Estado}Presente
\fichakey{Módulo}gui
\fichakey{Paquete}org.openmarkov.gui.dialog.configuration
\fichakey{Ficheros}\path{gui/src/main/java/org/openmarkov/gui/dialog/configuration/PreferencesDialog.java:95-103,161-196,302-384,424} \newline \path{gui/src/main/java/org/openmarkov/gui/dialog/configuration/PreferencesTableModel.java} \newline \path{gui/src/main/java/org/openmarkov/gui/dialog/configuration/PreferenceTreeNode.java} \newline \path{gui/src/main/java/org/openmarkov/gui/configuration/UserPreference.java:70-79} \newline \path{gui/src/main/java/org/openmarkov/gui/configuration/UserPreferenceResolveStrategy.java:10-14,128-138} \newline \path{gui/src/main/java/org/openmarkov/gui/window/EditAndViewHandler.java:106-108} \newline \path{gui/src/main/java/org/openmarkov/gui/menutoolbar/menu/MainMenu.java:250,374} \newline \path{gui/src/main/java/org/openmarkov/gui/window/settings/SettingsDialog.java}
\fichakey{Comprobado}Leyendo el código
\fichakey{Esfuerzo}Medio (días)
\fichakey{Decisión de equipo}\textcolor{decisionrojo}{\textbf{Sí.}} Si la ventana antigua se retira del menú Tools (el diálogo «Settings\ldots» del menú Edit ya hace copia, restauración y valores por omisión sobre el almacén bueno) o se rehace sobre \texttt{User\allowbreak{}Preferences}
\fichakey{Relacionados}\hyperref[nota:68]{nota~68}, \hyperref[nota:90]{nota~90}, \hyperref[nota:98]{nota~98}, \hyperref[nota:162]{nota~162}
\end{ficha}
\paragraph{Qué pasa}
El menú Tools termina en «OpenMarkov configuration» (\texttt{Main\allowbreak{}Menu}, líneas 250 y 374), que abre \texttt{Preferences\allowbreak{}Dialog} (\texttt{Edit\allowbreak{}And\allowbreak{}View\allowbreak{}Handler.\allowbreak{}show\allowbreak{}User\allowbreak{}Configuration\allowbreak{}Dialog}, líneas 106-108): un árbol de familias de preferencias a la izquierda, una tabla clave/valor a la derecha y cinco botones (Save, Reset user preferences, Export preferences, Import preferences, Cancel).

\textbf{Mira donde no están las preferencias.} Lo que enseña la ventana depende de la máquina:

\begin{itemize}
\tightlist
\item
  en una instalación nueva, el árbol sale con la raíz y nada debajo, y la tabla vacía;
\item
  en una máquina donde corrió una versión antigua de OpenMarkov, sale lo que aquella dejó y ya nadie lee. Comprobado por lectura del disco en la máquina de desarrollo: \texttt{\textasciitilde{}/\allowbreak{}.\allowbreak{}java/\allowbreak{}.\allowbreak{}user\allowbreak{}Prefs/\allowbreak{}OPENMARKOV} tiene las familias \texttt{colors}, \texttt{directories}, \texttt{formats}, \texttt{languages} y \texttt{positions}, con fechas de 2025, mientras las preferencias vigentes están en \texttt{\textasciitilde{}/\allowbreak{}.\allowbreak{}openmarkov/\allowbreak{}preferences/\allowbreak{}\{directories,\allowbreak{}formats,\allowbreak{}network,\allowbreak{}user\_\allowbreak{}interface\}}.
\end{itemize}

En los dos casos, editar un valor en la tabla no cambia nada del programa. «Export preferences» (líneas 346-356) escribe ese almacén viejo a un fichero XML, e «Import preferences» (líneas 361-373) lo mete de vuelta en él, donde nadie lo lee. «Reset user preferences» (líneas 378-384) es el único botón que actúa sobre las preferencias reales (llama a \texttt{clear} e \texttt{initialize} de todas las \texttt{User\allowbreak{}Preferences}), y lo hace sin preguntar y sin que la tabla lo refleje.

Las preferencias del programa (\texttt{User\allowbreak{}Preferences}: últimos ficheros, directorios, tema, escala, dominios propios, acciones de arranque\ldots) se guardan con la primera estrategia de \texttt{User\allowbreak{}Preference\allowbreak{}Resolve\allowbreak{}Strategy} que sea capaz de escribir y borrar un valor de prueba (\texttt{User\allowbreak{}Preference}, líneas 70-79). El orden es \texttt{USER\_\allowbreak{}FOLDER} (ficheros bajo \texttt{\textasciitilde{}/\allowbreak{}.\allowbreak{}openmarkov/\allowbreak{}preferences}), \texttt{INSTALLED\_\allowbreak{}LOCATION}, \texttt{BACKING\_\allowbreak{}STORE} (la API \texttt{java.\allowbreak{}util.\allowbreak{}prefs}) y \texttt{SESSION}. En cualquier máquina con un directorio personal escribible gana la primera, así que \texttt{java.\allowbreak{}util.\allowbreak{}prefs} no se usa. El diálogo, en cambio, construye su árbol sobre \texttt{Preferences.\allowbreak{}user\allowbreak{}Root(\allowbreak{}).\allowbreak{}node(\allowbreak{}"OPENMARKOV")} (líneas 191-196), es decir, sobre \texttt{java.\allowbreak{}util.\allowbreak{}prefs}.

\textbf{Cancelar no cancela.} Lo escrito en la tabla se guarda al teclearlo (\texttt{Preferences\allowbreak{}Table\allowbreak{}Model.\allowbreak{}set\allowbreak{}Value\allowbreak{}At} hace \texttt{put} y \texttt{sync}). El «deshacer» que llama Cancelar (líneas 325-330) crea dos modelos nuevos sobre las raíces del sistema y del usuario y llama a \texttt{undo(\allowbreak{})}, que asigna a \texttt{pref} el campo \texttt{pref\allowbreak{}Saved}, que es la misma referencia (\texttt{this.\allowbreak{}pref\allowbreak{}Saved\ =\ pref} en el constructor): no restaura nada.

\textbf{Otros restos en la misma ventana.}

\begin{itemize}
\tightlist
\item
  \texttt{create\allowbreak{}Tree} quita siempre el primer hijo del nodo raíz (líneas 179-180, comentario «not shown color preferences in tree»), suponiendo que es la familia de colores; el orden de \texttt{children\allowbreak{}Names(\allowbreak{})} no está garantizado, así que en una máquina con familias antiguas oculta una cualquiera.
\item
  \texttt{Color\allowbreak{}Table\allowbreak{}Cell\allowbreak{}Editor} y \texttt{Color\allowbreak{}Table\allowbreak{}Cell\allowbreak{}Renderer} (144 líneas entre los dos) no los usa ninguna clase.
\item
  «Export» e «Import» escriben «Export selected» / «Import selected» por la salida estándar (líneas 349 y 362).
\item
  Los textos de la ventana están escritos en el código, pero eso no forma parte de este error.
\end{itemize}

Existe otro diálogo, \texttt{Settings\allowbreak{}Dialog} (menú Edit → «Settings\ldots»), que sí trabaja sobre \texttt{User\allowbreak{}Preferences}: escala, tema, acciones de arranque, dominios propios, crear y restaurar una copia de la configuración y volver a los valores por omisión. La ventana antigua de Tools duplica esas funciones sobre el almacén equivocado.

No se ha abierto la ventana (es un \texttt{JDialog}); lo anterior es lectura del código y del disco.
\paragraph{El código}
gui/src/main/java/org/openmarkov/gui/dialog/configuration/PreferencesDialog.java:191-196,325-330

\begin{Shaded}
\begin{Highlighting}[]
\KeywordTok{private} \BuiltInTok{MutableTreeNode} \FunctionTok{createUserRootNode}\OperatorTok{(}\BuiltInTok{Object}\NormalTok{ obj}\OperatorTok{)} \KeywordTok{throws} \BuiltInTok{BackingStoreException} \OperatorTok{\{}
    \ControlFlowTok{if} \OperatorTok{(}\NormalTok{obj }\OperatorTok{==} \KeywordTok{null}\OperatorTok{)} \OperatorTok{\{}
        \ControlFlowTok{return} \KeywordTok{new} \FunctionTok{PreferenceTreeNode}\OperatorTok{(}\BuiltInTok{Preferences}\OperatorTok{.}\FunctionTok{userRoot}\OperatorTok{());}
    \OperatorTok{\}}
    \ControlFlowTok{return} \KeywordTok{new} \FunctionTok{PreferenceTreeNode}\OperatorTok{(}\BuiltInTok{Preferences}\OperatorTok{.}\FunctionTok{userRoot}\OperatorTok{().}\FunctionTok{node}\OperatorTok{((}\BuiltInTok{String}\OperatorTok{)}\NormalTok{ obj}\OperatorTok{));}   \CommentTok{// "OPENMARKOV"}
\OperatorTok{\}}
\KeywordTok{...}
\KeywordTok{protected} \DataTypeTok{void} \FunctionTok{actionPerformedCancel}\OperatorTok{()} \KeywordTok{throws} \BuiltInTok{BackingStoreException} \OperatorTok{\{}
    \KeywordTok{new} \FunctionTok{PreferencesTableModel}\OperatorTok{(}\BuiltInTok{Preferences}\OperatorTok{.}\FunctionTok{systemRoot}\OperatorTok{()).}\FunctionTok{undo}\OperatorTok{();}
    \KeywordTok{new} \FunctionTok{PreferencesTableModel}\OperatorTok{(}\BuiltInTok{Preferences}\OperatorTok{.}\FunctionTok{userRoot}\OperatorTok{()).}\FunctionTok{undo}\OperatorTok{();}
    \KeywordTok{this}\OperatorTok{.}\FunctionTok{setVisible}\OperatorTok{(}\KeywordTok{false}\OperatorTok{);}
    \KeywordTok{this}\OperatorTok{.}\FunctionTok{dispose}\OperatorTok{();}
\OperatorTok{\}}
\end{Highlighting}
\end{Shaded}

gui/src/main/java/org/openmarkov/gui/configuration/UserPreferenceResolveStrategy.java:10-14

\begin{Shaded}
\begin{Highlighting}[]
\KeywordTok{enum}\NormalTok{ UserPreferenceResolveStrategy }\OperatorTok{\{}
\NormalTok{    USER\_FOLDER}\OperatorTok{,}
\NormalTok{    INSTALLED\_LOCATION}\OperatorTok{,}
\NormalTok{    BACKING\_STORE}\OperatorTok{,}
\NormalTok{    SESSION}\OperatorTok{;}
\end{Highlighting}
\end{Shaded}
\paragraph{Cómo arreglarlo}
Exige una decisión de producto, con dos salidas:

\begin{enumerate}
\def\labelenumi{\arabic{enumi}.}
\tightlist
\item
  Retirar la entrada «OpenMarkov configuration» del menú Tools (\texttt{Main\allowbreak{}Menu}, líneas 250 y 374, y \texttt{Action\allowbreak{}Commands.\allowbreak{}CONFIGURATION} en \texttt{Main\allowbreak{}Panel\allowbreak{}Listener\allowbreak{}Assistant}), junto con \texttt{Preferences\allowbreak{}Dialog}, \texttt{Preferences\allowbreak{}Table\allowbreak{}Model}, \texttt{Preference\allowbreak{}Tree\allowbreak{}Node} y los dos controles de color sin uso, porque \texttt{Settings\allowbreak{}Dialog} ya cubre la función sobre el almacén real.
\item
  Rehacerla sobre \texttt{User\allowbreak{}Preferences.\allowbreak{}get\allowbreak{}All\allowbreak{}Preferences(\allowbreak{})} (el almacén que resuelve \texttt{User\allowbreak{}Preference}), con un Cancelar que restaure los valores leídos al abrir, o sin Cancelar si los cambios se aplican al momento; la exportación e importación deberían usar el mismo formato que la copia de \texttt{Settings\allowbreak{}Dialog}.
\end{enumerate}

En cualquiera de las dos, conviene que «Reset user preferences» pida confirmación, porque borra la lista de ficheros recientes y los dominios propios.

Prueba: con \texttt{User\allowbreak{}Preference.\allowbreak{}IGNORE\_\allowbreak{}STORAGE} o un directorio personal temporal, comprobar que lo que enseña la ventana coincide con \texttt{User\allowbreak{}Preferences.\allowbreak{}get\allowbreak{}All\allowbreak{}Preferences(\allowbreak{})} y que Cancelar deja los valores como estaban.

\setcounter{subsection}{175}
\subsection[Las escalas de criterio se redondean a tres decimales y pueden quedar en 0]{El diálogo de opciones de inferencia guarda la escala de cada criterio redondeada a tres decimales, aunque solo se toque la celda: −1/30 000 queda en 0 y la evaluación deja de contar los costes}\label{err:176}
\begin{ficha}
\fichakey{Estado}Presente
\fichakey{Módulo}gui
\fichakey{Paquete}org.openmarkov.gui.dialog.inference.common
\fichakey{Ficheros}\path{gui/src/main/java/org/openmarkov/gui/dialog/inference/common/InferenceOptionsDialog.java:606-620} \newline \path{gui/src/main/java/org/openmarkov/gui/dialog/inference/common/InferenceOptionsDialog.java:642-651} \newline \path{gui/src/main/java/org/openmarkov/gui/dialog/inference/common/InferenceOptionsDialog.java:682-691} \newline \path{gui/src/main/java/org/openmarkov/gui/dialog/inference/common/InferenceOptionsDialog.java:725-728} \newline \path{core/src/main/java/org/openmarkov/core/model/network/UtilityOperations.java:67-86}
\fichakey{Comprobado}Leyendo el código y ejecutándolo
\fichakey{Esfuerzo}Pequeño (horas)
\fichakey{Decisión de equipo}No
\fichakey{Relacionados}\hyperref[err:11]{error~11}, \hyperref[err:12]{error~12}
\end{ficha}
\paragraph{Qué pasa}
En una red con más de un criterio, el diálogo de opciones de inferencia (el que aparece al entrar en modo inferencia, al pedir el análisis de coste-efectividad, el de sensibilidad o el árbol de decisión, y desde el menú de opciones de inferencia) enseña una tabla con la escala de cada criterio. Lo que el usuario escribe en una escala se guarda redondeado a tres decimales: una escala menor que 0,0005 en valor absoluto queda en 0, y las cercanas pierden precisión (−0,0006 pasa a −0,001). Pasa también sin cambiar nada: basta con hacer clic en la celda y salir de ella, así que una escala más fina que venga del fichero se pierde en cuanto se toca.

Es un uso corriente. Para expresar el resultado en años de vida ajustados por calidad, la escala del coste es −1 dividido por la disposición a pagar: con 30 000 por año, −0,0000333. Las redes del repositorio usan la convención contraria (coste −1, efectividad 30 000), que no se ve afectada. Al salir de la celda, la tabla se vuelve a construir y enseña la escala redondeada (−0.0), sin otro aviso.

La consecuencia es un resultado distinto sin aviso: con la escala del coste en 0, la evaluación de un solo criterio (la del modo inferencia, por ejemplo) ignora los costes. En el análisis de coste-efectividad, una escala 0 hace que \texttt{Utility\allowbreak{}Operations.\allowbreak{}apply\allowbreak{}CEUtility\allowbreak{}Scaling} quite de la red los nodos de ese criterio (líneas 81-84, por lectura).

Por qué pasa: en la tabla de \texttt{Inference\allowbreak{}Options\allowbreak{}Dialog}, \texttt{table\allowbreak{}Changed} toma el texto de la escala, lo formatea con \texttt{Decimal\allowbreak{}Format} y el patrón \texttt{\#.\allowbreak{}\#\#\#} y guarda en la copia del criterio el número que resulta (líneas 647-650 para la escala de coste-efectividad y 687-690 para la de un solo criterio). Al terminar de editar una celda, \texttt{JTable.\allowbreak{}editing\allowbreak{}Stopped} llama siempre a \texttt{set\allowbreak{}Value\allowbreak{}At} con el texto del editor, aunque no haya cambiado, y eso dispara \texttt{table\allowbreak{}Changed}. Las columnas de descuento usan el mismo patrón sobre el porcentaje (líneas 660-663 y 700-703), que ahí solo pierde cifras por debajo de 0,001 \%.

Medido construyendo la tabla del diálogo sin ventana con \texttt{0\_\allowbreak{}miscellaneous/\allowbreak{}networks/\allowbreak{}id/\allowbreak{}ID-CEA-test-2therapies.\allowbreak{}pgmx} (criterios «cost» y «effectiveness») y escribiendo en la escala de «cost»:

\begin{Verbatim}[breaklines,breakanywhere,fontsize=\small]
escrito -0.0000333                                   -> escala guardada -0.0
escrito -3.3333333333333335E-5                       -> -0.0
escrito -0.0005                                      -> -0.001
escrito -0.0006                                      -> -0.001
escrito 0.0004                                       -> 0.0
escrito -1                                           -> -1.0
escala ya en -3.33E-5, se sale de la celda sin tocar -> -0.0
escala de coste-efectividad, escrito 0.000001        -> 0.0
\end{Verbatim}

Con la misma red y la escala de la efectividad en 1, la evaluación da utilidad 8,9856 con la escala del coste en −1/30 000, y con prueba positiva recomienda «therapy 1»; con la escala redondeada (−0,0), 9,3937, y con prueba positiva recomienda «therapy 2», la de 70 000, y sin prueba «therapy 1» en lugar de no tratar.
\paragraph{El código}
gui/src/main/java/org/openmarkov/gui/dialog/inference/common/InferenceOptionsDialog.java:682-691

\begin{Shaded}
\begin{Highlighting}[]
\ControlFlowTok{if} \OperatorTok{(}\NormalTok{column }\OperatorTok{==}\NormalTok{ UNICRITERIA\_SCALE\_COLUMN}\OperatorTok{)} \OperatorTok{\{}
    \BuiltInTok{String}\NormalTok{ scale }\OperatorTok{=}\NormalTok{ table}\OperatorTok{.}\FunctionTok{getValueAt}\OperatorTok{(}\NormalTok{row}\OperatorTok{,}\NormalTok{ UNICRITERIA\_SCALE\_COLUMN}\OperatorTok{).}\FunctionTok{toString}\OperatorTok{();}
    \ControlFlowTok{if} \OperatorTok{(}\NormalTok{scale}\OperatorTok{.}\FunctionTok{indexOf}\OperatorTok{(}\CharTok{\textquotesingle{} \textquotesingle{}}\OperatorTok{)} \OperatorTok{!=} \OperatorTok{{-}}\DecValTok{1}\OperatorTok{)} \OperatorTok{\{}
\NormalTok{        scale }\OperatorTok{=}\NormalTok{ scale}\OperatorTok{.}\FunctionTok{substring}\OperatorTok{(}\DecValTok{0}\OperatorTok{,}\NormalTok{ scale}\OperatorTok{.}\FunctionTok{indexOf}\OperatorTok{(}\CharTok{\textquotesingle{} \textquotesingle{}}\OperatorTok{));}
    \OperatorTok{\}}
    \BuiltInTok{DecimalFormat}\NormalTok{ format }\OperatorTok{=} \OperatorTok{(}\BuiltInTok{DecimalFormat}\OperatorTok{)} \BuiltInTok{NumberFormat}\OperatorTok{.}\FunctionTok{getInstance}\OperatorTok{(}\BuiltInTok{Locale}\OperatorTok{.}\FunctionTok{ENGLISH}\OperatorTok{);}
\NormalTok{    format}\OperatorTok{.}\FunctionTok{applyLocalizedPattern}\OperatorTok{(}\StringTok{"\#.\#\#\#"}\OperatorTok{);}
\NormalTok{    scale }\OperatorTok{=}\NormalTok{ format}\OperatorTok{.}\FunctionTok{format}\OperatorTok{(}\BuiltInTok{Double}\OperatorTok{.}\FunctionTok{parseDouble}\OperatorTok{(}\NormalTok{scale}\OperatorTok{));}
\NormalTok{    decisionCriteria}\OperatorTok{.}\FunctionTok{get}\OperatorTok{(}\NormalTok{row }\OperatorTok{{-}} \DecValTok{1}\OperatorTok{).}\FunctionTok{setUnicriterizationScale}\OperatorTok{(}\BuiltInTok{Double}\OperatorTok{.}\FunctionTok{parseDouble}\OperatorTok{(}\NormalTok{scale}\OperatorTok{));}
\OperatorTok{\}}
\end{Highlighting}
\end{Shaded}
\paragraph{Cómo arreglarlo}
Guardar el número tal como se escribe (\texttt{Double.\allowbreak{}parse\allowbreak{}Double} del texto, sin pasar por \texttt{Decimal\allowbreak{}Format}) y redondear, si se quiere, solo al enseñarlo, en el dibujante de la celda. Los cuatro sitios del mismo método deben cambiar a la vez: las dos escalas (líneas 642-651 y 682-691) y los dos descuentos (líneas 654-665 y 694-705).

Conviene también que un número que no se puede leer se rechace con un mensaje en la celda, en vez de la \texttt{Number\allowbreak{}Format\allowbreak{}Exception} que hoy lanza \texttt{set\allowbreak{}Value\allowbreak{}At} (línea 613).

Prueba de regresión: sin ventana, con la tabla del diálogo o con el cálculo extraído a un método, escribir \texttt{-3.\allowbreak{}3333333333333335E-5} en la escala de un criterio y comprobar que se guarda ese número; y que salir de la celda sin cambiarla no modifica la escala.

% ══════════════════════════════════════════════════════════════════════
\section{Textos y títulos que se ven mal}\label{sec:textos}

Textos que la aplicación, en inglés, enseña equivocados, vacíos o con el nombre de una clave.

\setcounter{subsection}{176}
\subsection[La ayuda de atajos da dos teclas equivocadas, y una cierra el programa]{La ayuda de atajos promete dos teclas equivocadas, y siguiendo una se cierra el programa}\label{err:177}
\begin{ficha}
\fichakey{Estado}Presente
\fichakey{Módulo}gui
\fichakey{Paquete}org.openmarkov.gui.dialog
\fichakey{Ficheros}\path{gui/src/main/resources/gui/localize/Dialogs_en.xml:702-709} \newline \path{gui/src/main/java/org/openmarkov/gui/dialog/ShortcutsBox.java:38-45, 117-120} \newline \path{gui/src/main/java/org/openmarkov/gui/menutoolbar/menu/MainMenu.java:203, 205, 209, 218, 331, 342} \newline \path{gui/src/main/java/org/openmarkov/gui/window/MainPanel.java:149-158} \newline \path{gui/src/main/java/org/openmarkov/gui/window/MainPanelListenerAssistant.java:119} \newline \path{gui/src/main/java/org/openmarkov/gui/window/NetworkFileHandler.java:391-411}
\fichakey{Comprobado}Leyendo el código
\fichakey{Esfuerzo}Pequeño (horas)
\fichakey{Decisión de equipo}No
\fichakey{Relacionados}\hyperref[nota:18]{nota~18}, \hyperref[nota:127]{nota~127}
\end{ficha}
\paragraph{Qué pasa}
Help → «Shortcuts» abre \texttt{Shortcuts\allowbreak{}Box} (módulo \texttt{gui}), dos tablas de «tecla → función» cuyos textos salen de \texttt{Dialogs\_\allowbreak{}en.\allowbreak{}xml} (claves \texttt{Shortcuts\allowbreak{}Box.\allowbreak{}Specific.\allowbreak{}\textless{}entrada\textgreater{}.\allowbreak{}Shortcut} y \texttt{.\allowbreak{}Function}, \texttt{Shortcuts\allowbreak{}Box}, líneas 38-45 y 117-120). Comparadas con los aceleradores que pone la barra de menús (\texttt{Main\allowbreak{}Menu.\allowbreak{}create\allowbreak{}All\allowbreak{}Items}) y con el mapa de teclas de las pestañas (\texttt{Main\allowbreak{}Panel}, líneas 149-158), dos filas de la tabla de atajos específicos no dicen la verdad:

\begin{itemize}
\tightlist
\item
  \textbf{\texttt{Ctrl\ +\ W}}: la ayuda dice «Save and reopen» (\texttt{Dialogs\_\allowbreak{}en.\allowbreak{}xml}, líneas 703-706). Esa tecla cierra la red de la pestaña actual («Close network», \texttt{CLOSE\_\allowbreak{}TAB}, \texttt{Main\allowbreak{}Menu} línea 205). La tecla real de guardar y reabrir es \texttt{Ctrl\ +\ Alt\ +\ S} («Save network and reopen», línea 203).
\item
  \textbf{\texttt{Ctrl\ +\ A/\allowbreak{}Q}}: la ayuda dice «Previous/Next OpenMarkov sub-window» (líneas 699-702). \texttt{Ctrl\ +\ Q} sale del programa («Exit», \texttt{EXIT\_\allowbreak{}APPLICATION}, línea 209), y \texttt{Ctrl\ +\ A} no está atado a nada en la barra de menús ni en el mapa de pestañas. Las teclas reales para pasar a la pestaña anterior o siguiente son \texttt{Ctrl\ +\ Shift\ +\ ←} / \texttt{Ctrl\ +\ Shift\ +\ →} (menú View, líneas 331 y 342, y \texttt{Main\allowbreak{}Panel}, líneas 149-158).
\end{itemize}

Consecuencia: quien quiera pasar a la pestaña siguiente siguiendo la ayuda pulsa \texttt{Ctrl\ +\ Q} y sale del programa (\texttt{Main\allowbreak{}Panel\allowbreak{}Listener\allowbreak{}Assistant}, línea 128 → \texttt{Network\allowbreak{}File\allowbreak{}Handler.\allowbreak{}close\allowbreak{}Application}, líneas 395-405, que cierra todas las pestañas y llama a \texttt{System.\allowbreak{}exit(\allowbreak{}0)}). Por lectura, cada pestaña se cierra por el camino normal (\texttt{Network\allowbreak{}Editor\allowbreak{}Panel.\allowbreak{}close} → \texttt{network\allowbreak{}Can\allowbreak{}Be\allowbreak{}Closed}), así que las redes modificadas deberían pedir confirmación antes de perderse; lo que no se evita es la salida. Y quien quiera guardar y reabrir con \texttt{Ctrl\ +\ W} cierra la red.

El resto de filas sí coincide con el mapa: \texttt{Ctrl\ +\ I} (cambiar modo), \texttt{Ctrl\ +\ N}, \texttt{Ctrl\ +\ O}, \texttt{Ctrl\ +\ S}, \texttt{Ctrl\ +\ E} (seleccionar todo, línea 218), \texttt{Ctrl\ +\ X/\allowbreak{}C/\allowbreak{}V}, \texttt{Ctrl\ +\ Z}, \texttt{Ctrl\ +\ Y}. Establecido leyendo el código; no se ha pulsado ninguna tecla.
\paragraph{El código}
gui/src/main/resources/gui/localize/Dialogs\_en.xml:699-706

\begin{Shaded}
\begin{Highlighting}[]
\NormalTok{\textless{}}\KeywordTok{PreviousNext}\NormalTok{\textgreater{}}
\NormalTok{    \textless{}}\KeywordTok{Shortcut}\OtherTok{ value=}\StringTok{"Ctrl + A/Q"}\NormalTok{/\textgreater{}}
\NormalTok{    \textless{}}\KeywordTok{Function}\OtherTok{ value=}\StringTok{"Previous/Next OpenMarkov sub{-}window"}\NormalTok{/\textgreater{}}
\NormalTok{\textless{}/}\KeywordTok{PreviousNext}\NormalTok{\textgreater{}}
\NormalTok{\textless{}}\KeywordTok{SaveReopen}\NormalTok{\textgreater{}}
\NormalTok{    \textless{}}\KeywordTok{Shortcut}\OtherTok{ value=}\StringTok{"Ctrl + W"}\NormalTok{/\textgreater{}}
\NormalTok{    \textless{}}\KeywordTok{Function}\OtherTok{ value=}\StringTok{"Save and reopen"}\NormalTok{/\textgreater{}}
\NormalTok{\textless{}/}\KeywordTok{SaveReopen}\NormalTok{\textgreater{}}
\end{Highlighting}
\end{Shaded}

gui/src/main/java/org/openmarkov/gui/menutoolbar/menu/MainMenu.java:203-209

\begin{Shaded}
\begin{Highlighting}[]
\FunctionTok{createItem}\OperatorTok{(}\NormalTok{MenuItemNames}\OperatorTok{.}\FunctionTok{FILE\_SAVE\_OPEN\_MENUITEM}\OperatorTok{,}\NormalTok{ ActionCommands}\OperatorTok{.}\FunctionTok{SAVE\_OPEN\_NETWORK}\OperatorTok{,}\NormalTok{ IconBind}\OperatorTok{.}\FunctionTok{SAVE\_ENABLED}\OperatorTok{,}\NormalTok{ MainMenu}\OperatorTok{.}\FunctionTok{ctrlAlt}\OperatorTok{(}\BuiltInTok{KeyEvent}\OperatorTok{.}\FunctionTok{VK\_S}\OperatorTok{));}
\KeywordTok{...}
\FunctionTok{createItem}\OperatorTok{(}\NormalTok{MenuItemNames}\OperatorTok{.}\FunctionTok{FILE\_CLOSE\_MENUITEM}\OperatorTok{,}\NormalTok{ ActionCommands}\OperatorTok{.}\FunctionTok{CLOSE\_TAB}\OperatorTok{,} \KeywordTok{null}\OperatorTok{,}\NormalTok{ MainMenu}\OperatorTok{.}\FunctionTok{ctrl}\OperatorTok{(}\BuiltInTok{KeyEvent}\OperatorTok{.}\FunctionTok{VK\_W}\OperatorTok{));}
\KeywordTok{...}
\FunctionTok{createItem}\OperatorTok{(}\NormalTok{MenuItemNames}\OperatorTok{.}\FunctionTok{FILE\_EXIT\_MENUITEM}\OperatorTok{,}\NormalTok{ ActionCommands}\OperatorTok{.}\FunctionTok{EXIT\_APPLICATION}\OperatorTok{,} \KeywordTok{null}\OperatorTok{,}\NormalTok{ MainMenu}\OperatorTok{.}\FunctionTok{ctrl}\OperatorTok{(}\BuiltInTok{KeyEvent}\OperatorTok{.}\FunctionTok{VK\_Q}\OperatorTok{));}
\end{Highlighting}
\end{Shaded}
\paragraph{Cómo arreglarlo}
Corregir los dos textos en \texttt{Dialogs\_\allowbreak{}en.\allowbreak{}xml}: \texttt{Ctrl\ +\ Alt\ +\ S} para «Save and reopen» y \texttt{Ctrl\ +\ Shift\ +\ ←/\allowbreak{}→} para «Previous/Next tab» (y, si se quiere, añadir filas para \texttt{Ctrl\ +\ W} «Close network» y \texttt{Ctrl\ +\ Q} «Exit»). El fichero \texttt{Dialogs\_\allowbreak{}es.\allowbreak{}xml} tiene las mismas filas, pero no se mantiene.

Para que no vuelva a pasar, la ayuda podría construirse leyendo los aceleradores de los elementos de la barra de menús en vez de repetirlos a mano. Como mínimo, una prueba que compare cada fila de \texttt{Shortcuts\allowbreak{}Box.\allowbreak{}Common.\allowbreak{}*} y \texttt{Shortcuts\allowbreak{}Box.\allowbreak{}Specific.\allowbreak{}*} que nombre una tecla con \texttt{Ctrl} con el acelerador del elemento de menú correspondiente (\texttt{Main\allowbreak{}Menu.\allowbreak{}create\allowbreak{}All\allowbreak{}Items} y \texttt{build\allowbreak{}View\allowbreak{}Menu}).

\setcounter{subsection}{177}
\subsection[Pegar un nodo de sucesos donde no se admite da «\textgreater\textgreater\textgreater{} eventNode is null \textless\textless\textless»]{Al pegar un nodo de sucesos en una red que no los admite, el mensaje de error dice «\textgreater\textgreater\textgreater{} eventNode is null \textless\textless\textless»}\label{err:178}
\begin{ficha}
\fichakey{Estado}Presente
\fichakey{Módulo}core
\fichakey{Paquete}org.openmarkov.core.action.core
\fichakey{Ficheros}\path{core/src/main/java/org/openmarkov/core/action/core/AddNodeEdit.java:88-92, 145-148} \newline \path{gui/src/main/java/org/openmarkov/gui/action/PasteEdit.java:85} \newline \path{gui/src/main/java/org/openmarkov/gui/window/edition/networkEditorPanel/NetworkEditorPanel.java:169} \newline \path{core/src/main/resources/core/localize/core_class_localizations_en.xml:456-457}
\fichakey{Comprobado}Leyendo el código y ejecutándolo
\fichakey{Esfuerzo}Pequeño (horas)
\fichakey{Decisión de equipo}No
\fichakey{Relacionados}\hyperref[nota:153]{nota~153}
\end{ficha}
\paragraph{Qué pasa}
Crear directamente un nodo de sucesos (tipo EVENT, los nodos de las redes de simulación de eventos discretos) en una red que no los admite no es posible: la herramienta y el botón de tipo de nodo se deshabilitan si la red lleva la restricción \texttt{No\allowbreak{}Event\allowbreak{}Nodes}. Pero el portapapeles es único para todas las pestañas (\texttt{CLIPBOARD\_\allowbreak{}ASSISTANT} es estático en \texttt{Network\allowbreak{}Editor\allowbreak{}Panel}), y pegar ejecuta un \texttt{Add\allowbreak{}Node\allowbreak{}Edit} por nodo copiado (\texttt{Paste\allowbreak{}Edit}, línea 85) con el tipo del nodo original.

Copiar un nodo de sucesos en una red de simulación de eventos discretos y pegarlo en una red bayesiana o en un diagrama de influencia lanza la violación, que la ventana muestra en un diálogo de error. El texto de ese diálogo dice «\textgreater\textgreater\textgreater{} eventNode is null \textless\textless\textless{} is of type event.» en lugar del nombre del nodo. El pegado se rechaza entero y la red no cambia; el defecto es solo el texto.

\texttt{Add\allowbreak{}Node\allowbreak{}Edit} añade un nodo a la red. Su comprobación previa (\texttt{check\allowbreak{}Constraints\allowbreak{}Will\allowbreak{}Be\allowbreak{}Met}) rechaza un nodo de tipo EVENT cuando la red lleva \texttt{No\allowbreak{}Event\allowbreak{}Nodes}, y para ello construye \texttt{Cannot\allowbreak{}Have\allowbreak{}Event\allowbreak{}Node\allowbreak{}Exception(\allowbreak{}constraint,\allowbreak{}\ this.\allowbreak{}new\allowbreak{}Node)}. Pero \texttt{new\allowbreak{}Node} solo se asigna en \texttt{do\allowbreak{}Edit(\allowbreak{})}, que todavía no se ha ejecutado: en la comprobación vale \texttt{null}. El texto de esa excepción es «\{eventNode\} is of type event.», y con el nodo a \texttt{null} el sistema de textos lo resuelve como «\textgreater\textgreater\textgreater{} eventNode is null \textless\textless\textless{} is of type event.».

Medido con un pequeño programa de prueba que construye ese \texttt{Paste\allowbreak{}Edit} sin ventana: el mensaje que se entregaría es «Multiple constraints have been violated: - Could not met the constraint: No event nodes allowed --- \textgreater\textgreater\textgreater{} eventNode is null \textless\textless\textless{} is of type event. - Could not met the constraint: Only chance nodes are allowed --- Node for Event variable Ev with states {[}state event{]} must be a chance node.» En la red bayesiana salta además \texttt{Only\allowbreak{}Chance\allowbreak{}Nodes}; en un diagrama de influencia solo saldría la primera línea, con el marcador.
\paragraph{El código}
core/src/main/java/org/openmarkov/core/action/core/AddNodeEdit.java:88-92

\begin{Shaded}
\begin{Highlighting}[]
\ControlFlowTok{if} \OperatorTok{(}\NormalTok{probNet}\OperatorTok{.}\FunctionTok{getConstraintOfClass}\OperatorTok{(}\NormalTok{NoEventNodes}\OperatorTok{.}\FunctionTok{class}\OperatorTok{)} \KeywordTok{instanceof}\NormalTok{ NoEventNodes constraint}\OperatorTok{)} \OperatorTok{\{}
    \ControlFlowTok{if} \OperatorTok{(}\KeywordTok{this}\OperatorTok{.}\FunctionTok{nodeType} \OperatorTok{==}\NormalTok{ NodeType}\OperatorTok{.}\FunctionTok{EVENT}\OperatorTok{)} \OperatorTok{\{}
\NormalTok{        constraintChecker}\OperatorTok{.}\FunctionTok{addException}\OperatorTok{(}\KeywordTok{new}\NormalTok{ ConstraintViolatedException}\OperatorTok{.}\FunctionTok{CannotHaveEventNodeException}\OperatorTok{(}\NormalTok{constraint}\OperatorTok{,} \KeywordTok{this}\OperatorTok{.}\FunctionTok{newNode}\OperatorTok{));}
    \OperatorTok{\}}
\OperatorTok{\}}
\end{Highlighting}
\end{Shaded}
\paragraph{Cómo arreglarlo}
La excepción necesita algo que exista antes de ejecutar: la variable (\texttt{this.\allowbreak{}variable}) en lugar del nodo. Eso obliga a cambiar el tipo del campo \texttt{event\allowbreak{}Node} de \texttt{Cannot\allowbreak{}Have\allowbreak{}Event\allowbreak{}Node\allowbreak{}Exception} (o añadir un constructor con \texttt{Variable}) y el texto de \texttt{core\_\allowbreak{}class\_\allowbreak{}localizations\_\allowbreak{}en.\allowbreak{}xml} (líneas 456-457).

La misma regla vive en \texttt{No\allowbreak{}Event\allowbreak{}Nodes.\allowbreak{}check\allowbreak{}Prob\allowbreak{}Net} (que sí tiene el nodo); conviene que las dos vías produzcan el mismo texto.

Prueba: red bayesiana o diagrama de influencia, \texttt{Add\allowbreak{}Node\allowbreak{}Edit} con tipo EVENT, y afirmar que el mensaje localizado contiene el nombre de la variable y no contiene «\textgreater\textgreater\textgreater».

\setcounter{subsection}{178}
\subsection[El selector del tipo de relación enseña nombres internos como «ProbTable»]{El selector del tipo de relación enseña nombres internos como «ProbTable» o «UnivariateDistr»}\label{err:179}
\begin{ficha}
\fichakey{Estado}Presente
\fichakey{Módulo}gui, core
\fichakey{Paquete}org.openmarkov.gui.dialog.node; org.openmarkov.core.model.network.potential.plugin
\fichakey{Ficheros}\path{gui/src/main/java/org/openmarkov/gui/dialog/node/PotentialEditPanel.java:228-256} \newline \path{core/src/main/java/org/openmarkov/core/model/network/potential/plugin/PotentialUtils.java:34-46} \newline \path{gui/src/main/resources/gui/localize/Selectables_en.xml:3-49} \newline \path{io/src/main/java/org/openmarkov/io/probmodel/writer/PGMXWriter_1_0.java:116, 191}
\fichakey{Comprobado}Leyendo el código y ejecutándolo
\fichakey{Esfuerzo}Pequeño (horas)
\fichakey{Decisión de equipo}\textcolor{decisionrojo}{\textbf{Sí.}} Qué nombre legible se enseña para cada tipo de relación
\end{ficha}
\paragraph{Qué pasa}
En el diálogo de editar el potencial de un nodo, el desplegable «tipo de relación» mezcla nombres legibles («Same as previous», «Linear combination», «OR / MAX») con identificadores pegados («ProbTable», «AugmentedProbTable», «DistributionTable», «UnivariateDistr») o poco claros («Exact»). \texttt{Cycle\allowbreak{}Length\allowbreak{}Shift} también sale con su identificador, aunque no apareció en los nodos probados.

\texttt{Potential\allowbreak{}Edit\allowbreak{}Panel.\allowbreak{}get\allowbreak{}Potential\allowbreak{}Type\allowbreak{}JCombobox} (líneas 227-241) llena el desplegable con las clases de potencial permitidas para el nodo y las pinta con \texttt{Potential\allowbreak{}Utils::get\allowbreak{}Potential\allowbreak{}Name}, que devuelve el primer valor de \texttt{names} de la anotación \texttt{@Potential\allowbreak{}Type}. Ese mismo nombre es el que los escritores de \texttt{.\allowbreak{}pgmx} guardan como tipo del potencial (\texttt{PGMXWriter\_\allowbreak{}1\_\allowbreak{}0}, líneas 116 y 191), así que es un identificador de formato.

\texttt{Selectables\_\allowbreak{}en.\allowbreak{}xml} tiene un bloque \texttt{algorithm\allowbreak{}Relation\allowbreak{}Type} con textos legibles, pero ningún código Java lo lee, y sus claves (\texttt{Table}, \texttt{Same\allowbreak{}As\allowbreak{}Previous}, \texttt{Tree\_\allowbreak{}ADD}, \texttt{ICI-OR}, \texttt{ICI-MAXCausal}\ldots) no coinciden con los nombres actuales de las anotaciones.

Medido con un pequeño programa de prueba que llama a \texttt{Potential\allowbreak{}Utils.\allowbreak{}get\allowbreak{}Filtered\allowbreak{}Potential\allowbreak{}Classes} sobre nodos reales (quitando las clases abstractas, como hace el panel):

\begin{itemize}
\tightlist
\item
  nodo de azar de estados finitos («Disease» de \texttt{0\_\allowbreak{}miscellaneous/\allowbreak{}networks/\allowbreak{}id/\allowbreak{}ID-decide-test.\allowbreak{}pgmx}): AugmentedProbTable, Conditional Gaussian, Delta, Discretized Cauchy, Hazard (Exponential), ProbTable, Uniform;
\item
  nodo de utilidad («Health state», de la misma red): DistributionTable, Exact, Exponential, Function, Linear combination, Tree/ADD, Uniform, UnivariateDistr;
\item
  nodo de azar numérico («Duration of treatment {[}0{]}» de \texttt{0\_\allowbreak{}miscellaneous/\allowbreak{}networks/\allowbreak{}mid/\allowbreak{}DMHEE/\allowbreak{}MID-dmhee-2.\allowbreak{}5.\allowbreak{}pgmx}): Binomial, Delta, DistributionTable, Exact, Exponential, Function, Linear combination, Product, Sum, Uniform, UnivariateDistr.
\end{itemize}
\paragraph{El código}
gui/src/main/java/org/openmarkov/gui/dialog/node/PotentialEditPanel.java:235-237

\begin{Shaded}
\begin{Highlighting}[]
\NormalTok{            filteredPotentialNames}\OperatorTok{.}\FunctionTok{sort}\OperatorTok{(}\BuiltInTok{Comparator}\OperatorTok{.}\FunctionTok{comparing}\OperatorTok{(}\NormalTok{PotentialUtils}\OperatorTok{::}\NormalTok{getPotentialName}\OperatorTok{));}
            \KeywordTok{this}\OperatorTok{.}\FunctionTok{potentialTypeComboBox} \OperatorTok{=} \KeywordTok{new} \BuiltInTok{JComboBox}\OperatorTok{\textless{}}\BuiltInTok{Class}\OperatorTok{\textless{}?} \KeywordTok{extends}\NormalTok{ Potential}\OperatorTok{\textgreater{}\textgreater{}(}\NormalTok{filteredPotentialNames}\OperatorTok{.}\FunctionTok{toArray}\OperatorTok{(}\KeywordTok{new} \BuiltInTok{Class}\OperatorTok{[}\DecValTok{0}\OperatorTok{]));}
            \KeywordTok{this}\OperatorTok{.}\FunctionTok{potentialTypeComboBox}\OperatorTok{.}\FunctionTok{setRenderer}\OperatorTok{(}\KeywordTok{new}\NormalTok{ JComboBoxFunctionRender}\OperatorTok{\textless{}}\BuiltInTok{Class}\OperatorTok{\textless{}?} \KeywordTok{extends}\NormalTok{ Potential}\OperatorTok{\textgreater{}\textgreater{}(}\NormalTok{PotentialUtils}\OperatorTok{::}\NormalTok{getPotentialName}\OperatorTok{));}
\end{Highlighting}
\end{Shaded}
\paragraph{Cómo arreglarlo}
No cambiar \texttt{names} de las anotaciones, porque es lo que se escribe y se lee en los ficheros. Separar el nombre para mostrar:

\begin{itemize}
\tightlist
\item
  una clave de texto por clase de potencial en \texttt{Selectables\_\allowbreak{}en.\allowbreak{}xml} (rehaciendo el bloque \texttt{algorithm\allowbreak{}Relation\allowbreak{}Type} con claves que coincidan, o usando el nombre de la clase como clave);
\item
  un método que devuelva el texto si existe y el identificador si no;
\item
  usarlo en el pintor del desplegable y también en la ordenación de la línea 235, para que el orden sea el que ve el usuario.
\end{itemize}

Hay que decidir los textos.

Prueba: para cada clase concreta anotada con \texttt{@Potential\allowbreak{}Type}, el nombre para mostrar no debe ser la clave rodeada \texttt{\textgreater{}\textgreater{}\textgreater{}\ …\ \textless{}\textless{}\textless{}}.

\setcounter{subsection}{179}
\subsection[El diálogo de dominios estándar se abre sin título]{El diálogo de dominios estándar se abre sin título}\label{err:180}
\begin{ficha}
\fichakey{Estado}Presente
\fichakey{Módulo}gui
\fichakey{Paquete}org.openmarkov.gui.dialog.node
\fichakey{Ficheros}\path{gui/src/main/java/org/openmarkov/gui/dialog/node/StandardDomainsDialog.java:24-37} \newline \path{gui/src/main/java/org/openmarkov/gui/dialog/common/DialogBase.java:49} \newline \path{gui/src/main/resources/gui/localize/Buttons_en.xml:3-7}
\fichakey{Comprobado}Leyendo el código
\fichakey{Esfuerzo}Pequeño (horas)
\fichakey{Decisión de equipo}No
\fichakey{Relacionados}\hyperref[err:173]{error~173}, \hyperref[err:174]{error~174}, \hyperref[err:181]{error~181}
\end{ficha}
\paragraph{Qué pasa}
Al pulsar «Standard domains» en la pestaña de dominio del diálogo de un nodo, la ventana \texttt{Standard\allowbreak{}Domains\allowbreak{}Dialog} se abre con la barra de título vacía.

Su constructor llama a \texttt{super(\allowbreak{}owner)}, y ninguna clase de la cadena (\texttt{Ok\allowbreak{}Cancel\allowbreak{}Dialog}, \texttt{Bottom\allowbreak{}Panel\allowbreak{}Button\allowbreak{}Dialog}, \texttt{Dialog\allowbreak{}Base}, que llama a \texttt{JDialog(\allowbreak{}owner)} sin título) llama a \texttt{set\allowbreak{}Title}. El diálogo hermano \texttt{Standard\allowbreak{}Criteria\allowbreak{}Dialog} tampoco lo hace.

Comprobado leyendo; no se ha abierto la ventana.
\paragraph{Cómo arreglarlo}
Añadir en el constructor \texttt{set\allowbreak{}Title(\allowbreak{}string\allowbreak{}Database.\allowbreak{}get\allowbreak{}String(\allowbreak{}"Standard\allowbreak{}Domain.\allowbreak{}Text"))} (hoy «Standard domains», el texto del botón que la abre) o una clave propia en el fichero de textos. Lo mismo para \texttt{Standard\allowbreak{}Criteria\allowbreak{}Dialog} con \texttt{Standard\allowbreak{}Criteria.\allowbreak{}Text}.

Prueba: construir el diálogo sin mostrarlo y comprobar que \texttt{get\allowbreak{}Title(\allowbreak{})} no está vacío.

\setcounter{subsection}{180}
\subsection[Seis de las ventanillas del editor de árboles se abren sin título]{Seis de las ventanillas del editor de árboles se abren sin título}\label{err:181}
\begin{ficha}
\fichakey{Estado}Presente
\fichakey{Módulo}gui
\fichakey{Paquete}org.openmarkov.gui.dialog.treeadd
\fichakey{Ficheros}\path{gui/src/main/java/org/openmarkov/gui/dialog/treeadd/AddStatesToBranchDialog.java:48-61} \newline \path{gui/src/main/java/org/openmarkov/gui/dialog/treeadd/RemoveStatesDialog.java:48-58} \newline \path{gui/src/main/java/org/openmarkov/gui/dialog/treeadd/AddVariablesDialog.java:47-57} \newline \path{gui/src/main/java/org/openmarkov/gui/dialog/treeadd/RemoveVariablesDialog.java:46-56} \newline \path{gui/src/main/java/org/openmarkov/gui/dialog/treeadd/ChangeIntervalDialog.java:48-62} \newline \path{gui/src/main/java/org/openmarkov/gui/dialog/treeadd/SetReferenceDialog.java} \newline \path{gui/src/main/java/org/openmarkov/gui/dialog/treeadd/SplitIntervalDialog.java:33} \newline \path{gui/src/main/java/org/openmarkov/gui/dialog/treeadd/TreeADDEditorPanel.java:1245}
\fichakey{Comprobado}Leyendo el código
\fichakey{Esfuerzo}Pequeño (horas)
\fichakey{Decisión de equipo}No
\fichakey{Relacionados}\hyperref[err:161]{error~161}, \hyperref[err:180]{error~180}, \hyperref[err:182]{error~182}, \hyperref[nota:125]{nota~125}
\end{ficha}
\paragraph{Qué pasa}
Varias operaciones del menú contextual del editor de árboles (\texttt{Tree\allowbreak{}ADDEditor\allowbreak{}Panel}, módulo \texttt{gui}) abren una ventanilla antes de actuar. En seis de ellas la barra de la ventana aparece vacía, y el usuario no ve qué operación está confirmando:

\begin{itemize}
\tightlist
\item
  \texttt{Add\allowbreak{}States\allowbreak{}To\allowbreak{}Branch\allowbreak{}Dialog}, que abre «Add states to branch»;
\item
  \texttt{Remove\allowbreak{}States\allowbreak{}Dialog}, que abre «Remove states from branch»;
\item
  \texttt{Add\allowbreak{}Variables\allowbreak{}Dialog}, que abre «Add variables to potential»;
\item
  \texttt{Remove\allowbreak{}Variables\allowbreak{}Dialog}, que abre «Remove variables»;
\item
  \texttt{Change\allowbreak{}Interval\allowbreak{}Dialog}, que abre «Change interval»;
\item
  \texttt{Set\allowbreak{}Reference\allowbreak{}Dialog}, que abre «Set reference».
\end{itemize}

Todas heredan de \texttt{Ok\allowbreak{}Cancel\allowbreak{}Dialog} → \texttt{Bottom\allowbreak{}Panel\allowbreak{}Button\allowbreak{}Dialog} → \texttt{Dialog\allowbreak{}Base} → \texttt{JDialog}, y ninguna de esas clases base pone un título. Las cinco primeras tienen su llamada a \texttt{set\allowbreak{}Title} dentro de un bloque de comentario (restos de una API de textos antigua); \texttt{Set\allowbreak{}Reference\allowbreak{}Dialog} no tiene ni siquiera la llamada comentada.

La séptima, \texttt{Split\allowbreak{}Interval\allowbreak{}Dialog}, sí tiene título: \texttt{set\allowbreak{}Title(\allowbreak{}"Split\ Interval")} en \texttt{initialize} (línea 33); el bloque comentado de más abajo no lo quita. La de poner etiqueta es un \texttt{JOption\allowbreak{}Pane.\allowbreak{}show\allowbreak{}Input\allowbreak{}Dialog(\allowbreak{}null,\allowbreak{}\ "Enter\ label\ name:\ ",\allowbreak{}\ "Set\ label",\allowbreak{}\ 1)} (\texttt{Tree\allowbreak{}ADDEditor\allowbreak{}Panel}, línea 1245), también con título.

Los rótulos escritos en inglés dentro del código («Threshold value: », «Included in first interval {]}(», «Split Interval», «Set label», «Enter label name: ») no forman parte de este error: la aplicación está solo en inglés y se leen bien.

No se ha abierto ninguna ventanilla (no se puede sin pantalla); la ausencia de título se deduce de que no hay ningún \texttt{set\allowbreak{}Title} activo en la clase ni en sus bases.
\paragraph{El código}
gui/src/main/java/org/openmarkov/gui/dialog/treeadd/AddVariablesDialog.java:47-57

\begin{Shaded}
\begin{Highlighting}[]
\KeywordTok{private} \DataTypeTok{void} \FunctionTok{configureComponentsPanel}\OperatorTok{()} \OperatorTok{\{}
    \CommentTok{/*dialogStringResource =}
\CommentTok{            StringResourceLoader.getUniqueInstance().getBundleDialogs();}
\CommentTok{    messageStringResource =}
\CommentTok{            StringResourceLoader.getUniqueInstance().getBundleMessages();}
\CommentTok{    setTitle(dialogStringResource}
\CommentTok{            .getValuesInAString("NodePotentialDialog.Title));*/}
    \FunctionTok{getComponentsPanel}\OperatorTok{().}\FunctionTok{setLayout}\OperatorTok{(}\KeywordTok{new} \BuiltInTok{BorderLayout}\OperatorTok{(}\DecValTok{5}\OperatorTok{,} \DecValTok{5}\OperatorTok{));}
    \FunctionTok{getComponentsPanel}\OperatorTok{().}\FunctionTok{add}\OperatorTok{(}\FunctionTok{getJPanelVariables}\OperatorTok{(),} \BuiltInTok{BorderLayout}\OperatorTok{.}\FunctionTok{CENTER}\OperatorTok{);}
\OperatorTok{\}}
\end{Highlighting}
\end{Shaded}
\paragraph{Cómo arreglarlo}
Dar a cada ventanilla un título. Lo natural es el texto de la opción de menú que la abre (claves \texttt{Tree\allowbreak{}ADD.\allowbreak{}Join\allowbreak{}Branches}, \texttt{Tree\allowbreak{}ADD.\allowbreak{}Dissociate\allowbreak{}States}, \texttt{Tree\allowbreak{}ADD.\allowbreak{}Add\allowbreak{}Variables}, \texttt{Tree\allowbreak{}ADD.\allowbreak{}Remove\allowbreak{}Variables}, \texttt{Tree\allowbreak{}ADD.\allowbreak{}Change\allowbreak{}Interval}, \texttt{Tree\allowbreak{}ADD.\allowbreak{}Set\allowbreak{}Reference} de \texttt{Menus\_\allowbreak{}en.\allowbreak{}xml}), que es el criterio que ya siguen otros diálogos. Quitar de paso los bloques comentados. Si se quiere uniformidad, \texttt{Split\allowbreak{}Interval\allowbreak{}Dialog} puede tomar también su título de \texttt{Tree\allowbreak{}ADD.\allowbreak{}Split\allowbreak{}Interval}, aunque no es necesario porque la aplicación está solo en inglés.

Es la misma familia que \hyperref[err:180]{error~180} (diálogo de dominios estándar sin título); conviene revisar a la vez las demás subclases de \texttt{Ok\allowbreak{}Cancel\allowbreak{}Dialog} que no llaman a \texttt{set\allowbreak{}Title}.

Prueba: una prueba de interfaz (AssertJ-Swing, con pantalla) que abra cada operación del menú del árbol y compruebe que el título de la ventana no está vacío (estas ventanas no se pueden construir sin pantalla).

\setcounter{subsection}{181}
\subsection[Varios títulos de error parten las siglas: «tree a d d», «X m l invalid»]{Los títulos de varias ventanas de error parten las siglas letra a letra: «Invalid limit in tree a d d», «X m l invalid»}\label{err:182}
\begin{ficha}
\fichakey{Estado}Presente
\fichakey{Módulo}core
\fichakey{Paquete}org.openmarkov.core.exception
\fichakey{Ficheros}\path{core/src/main/java/org/openmarkov/core/exception/IBundledOpenMarkovException.java:22-43} \newline \path{gui/src/main/java/org/openmarkov/gui/dialog/ExceptionDialog.java:52-69} \newline \path{gui/src/main/java/org/openmarkov/gui/dialog/ExceptionDialog.java:87-103} \newline \path{gui/src/main/java/org/openmarkov/gui/dialog/OMExceptionHandler.java:57} \newline \path{gui/src/main/java/org/openmarkov/gui/exception/InvalidLimitInTreeADDException.java} \newline \path{gui/src/main/java/org/openmarkov/gui/exception/ChangeDomainOfTreeADDIsNotAllowedException.java} \newline \path{core/src/main/java/org/openmarkov/core/exception/ProbNetParserException.java:147-160} \newline \path{io/src/main/java/org/openmarkov/io/probmodel/exception/PGMXParserException.java:92-103} \newline \path{io/src/main/java/org/openmarkov/io/probmodel/reader/PGMXReader_0_2.java:136-141} \newline \path{io/src/main/java/org/openmarkov/io/probmodel/reader/PGMXReader_0_2.java:176-180}
\fichakey{Comprobado}Leyendo el código y ejecutándolo
\fichakey{Esfuerzo}Pequeño (horas)
\fichakey{Decisión de equipo}No
\fichakey{Relacionados}\hyperref[err:161]{error~161}, \hyperref[err:181]{error~181}, \hyperref[err:183]{error~183}, \hyperref[nota:17]{nota~17}
\end{ficha}
\paragraph{Qué pasa}
Cuando una operación se rechaza con una excepción propia de OpenMarkov, la aplicación enseña una ventana de error cuyo mensaje sale del fichero de textos y cuyo título se fabrica a partir del nombre de la clase de la excepción. Si el nombre lleva una sigla, el título la parte en letras sueltas y en minúscula. Casos que se ven en la aplicación:

\begin{itemize}
\tightlist
\item
  en el editor de árboles, partir un intervalo por un umbral que coincide con un límite de la rama y del mismo lado: «Invalid limit in tree a d d»; cambiar el intervalo de una rama saliéndose del dominio: «Change domain of tree a d d is not allowed»;
\item
  abrir un \texttt{.\allowbreak{}pgmx} que no es XML bien formado, por ejemplo un fichero cortado: «X m l invalid»;
\item
  por lectura, con la misma causa: un \texttt{.\allowbreak{}pgmx} que no pasa la comprobación de versión o de estructura, «P g m x invalid»; y uno con una rama de árbol sin potencial ni referencia o sin sus dos umbrales, «Tree a d d without potential or referenfe» (con la errata del nombre de la clase) y «Tree a d d without two thresholds».
\end{itemize}

Por qué pasa: \texttt{IBundled\allowbreak{}Open\allowbreak{}Markov\allowbreak{}Exception.\allowbreak{}get\allowbreak{}Exception\allowbreak{}Title} (módulo \texttt{core}) devuelve \texttt{titlify\allowbreak{}Exception\allowbreak{}Name}, que quita el sufijo «Exception» y cambia cada letra mayúscula por un espacio seguido de esa letra en minúscula, sin tratar de otra forma las mayúsculas seguidas. \texttt{Exception\allowbreak{}Dialog.\allowbreak{}show} usa ese título para toda excepción de OpenMarkov. Le llegan las que el manejador global recibe envueltas en \texttt{Unrecoverable\allowbreak{}Exception} (\texttt{OMException\allowbreak{}Handler}, línea 57): las del editor de árboles (\hyperref[nota:17]{nota~17}) y las de abrir un fichero, que \texttt{GUIUtils.\allowbreak{}execute\allowbreak{}UIAction} envuelve así.

Medido:

\begin{itemize}
\tightlist
\item
  títulos devueltos por las excepciones construidas: \texttt{Invalid\allowbreak{}Limit\allowbreak{}In\allowbreak{}Tree\allowbreak{}ADDException} → «Invalid limit in tree a d d»; \texttt{Change\allowbreak{}Domain\allowbreak{}Of\allowbreak{}Tree\allowbreak{}ADDIs\allowbreak{}Not\allowbreak{}Allowed\allowbreak{}Exception} → «Change domain of tree a d d is not allowed»; \texttt{Tried\allowbreak{}To\allowbreak{}Split\allowbreak{}Interval\allowbreak{}Outside\allowbreak{}Bounds\allowbreak{}Exception} → «Tried to split interval outside bounds», que no tiene siglas y sale bien;
\item
  leer con \texttt{PGMXReader} un \texttt{.\allowbreak{}pgmx} cortado a mitad lanza \texttt{Prob\allowbreak{}Net\allowbreak{}Parser\allowbreak{}Exception\$XMLInvalid} con título «X m l invalid»;
\item
  recorriendo todas las clases de \texttt{org.\allowbreak{}openmarkov} que usan el título fabricado (170), nueve tienen siglas partidas: las seis citadas arriba (\texttt{Invalid\allowbreak{}Limit\allowbreak{}In\allowbreak{}Tree\allowbreak{}ADDException}, \texttt{Change\allowbreak{}Domain\allowbreak{}Of\allowbreak{}Tree\allowbreak{}ADDIs\allowbreak{}Not\allowbreak{}Allowed\allowbreak{}Exception}, \texttt{XMLInvalid}, \texttt{PGMXInvalid}, \texttt{Tree\allowbreak{}ADDWithout\allowbreak{}Potential\allowbreak{}Or\allowbreak{}Referenfe}, \texttt{Tree\allowbreak{}ADDWithout\allowbreak{}Two\allowbreak{}Thresholds}), \texttt{Only\allowbreak{}DESNets\allowbreak{}Allowed\allowbreak{}Exception} («Only d e s nets allowed», de la simulación de sucesos; no se ha buscado cómo llega a la ventana) y dos del formato Elvira (\texttt{Some\allowbreak{}Subpotentials\allowbreak{}Arent\allowbreak{}Linked\allowbreak{}To\allowbreak{}An\allowbreak{}ICIPotential} e \texttt{ICIModel\allowbreak{}Not\allowbreak{}Supported\allowbreak{}By\allowbreak{}Elvira}), que no se tocan porque ese formato está congelado; el arreglo de abajo las corrige de paso.
\end{itemize}
\paragraph{El código}
core/src/main/java/org/openmarkov/core/exception/IBundledOpenMarkovException.java:22-43

\begin{Shaded}
\begin{Highlighting}[]
\AttributeTok{@Override} \AttributeTok{@Nullable} \KeywordTok{default} \BuiltInTok{String} \FunctionTok{getExceptionTitle}\OperatorTok{()} \OperatorTok{\{}
    \ControlFlowTok{return}\NormalTok{ IBundledOpenMarkovException}\OperatorTok{.}\FunctionTok{titlifyExceptionName}\OperatorTok{(}\KeywordTok{this}\OperatorTok{);}
\OperatorTok{\}}
\KeywordTok{...}
\AttributeTok{@NotNull} \DataTypeTok{static} \BuiltInTok{String} \FunctionTok{titlifyExceptionName}\OperatorTok{(}\AttributeTok{@NotNull} \BuiltInTok{Object}\NormalTok{ exception}\OperatorTok{)} \OperatorTok{\{}
    \BuiltInTok{String}\NormalTok{ simpleName }\OperatorTok{=}\NormalTok{ exception}\OperatorTok{.}\FunctionTok{getClass}\OperatorTok{().}\FunctionTok{getSimpleName}\OperatorTok{();}
    \ControlFlowTok{if} \OperatorTok{(}\NormalTok{simpleName}\OperatorTok{.}\FunctionTok{endsWith}\OperatorTok{(}\StringTok{"Exception"}\OperatorTok{))} \OperatorTok{\{}
\NormalTok{        simpleName }\OperatorTok{=}\NormalTok{ simpleName}\OperatorTok{.}\FunctionTok{substring}\OperatorTok{(}\DecValTok{0}\OperatorTok{,}\NormalTok{ simpleName}\OperatorTok{.}\FunctionTok{length}\OperatorTok{()} \OperatorTok{{-}} \StringTok{"Exception"}\OperatorTok{.}\FunctionTok{length}\OperatorTok{());}
    \OperatorTok{\}}
    \ControlFlowTok{for} \OperatorTok{(}\DataTypeTok{var}\NormalTok{ uppercaseLetter }\OperatorTok{:}\NormalTok{ IBundledOpenMarkovException}\OperatorTok{.}\FunctionTok{UPPERCASE\_LETTERS}\OperatorTok{)} \OperatorTok{\{}
\NormalTok{        simpleName }\OperatorTok{=}\NormalTok{ simpleName}\OperatorTok{.}\FunctionTok{replace}\OperatorTok{(}\NormalTok{uppercaseLetter}\OperatorTok{,} \StringTok{" "} \OperatorTok{+}\NormalTok{ uppercaseLetter}\OperatorTok{.}\FunctionTok{toLowerCase}\OperatorTok{());}
    \OperatorTok{\}}
    \KeywordTok{...}
    \ControlFlowTok{return}\NormalTok{ simpleName}\OperatorTok{.}\FunctionTok{replace}\OperatorTok{(}\StringTok{"}\SpecialCharTok{\textbackslash{}\textbackslash{}}\StringTok{n"}\OperatorTok{,} \BuiltInTok{System}\OperatorTok{.}\FunctionTok{lineSeparator}\OperatorTok{());}
\OperatorTok{\}}
\end{Highlighting}
\end{Shaded}
\paragraph{Cómo arreglarlo}
Cambiar \texttt{titlify\allowbreak{}Exception\allowbreak{}Name} para que una serie de mayúsculas se quede junta y en mayúsculas: cortar antes de una mayúscula que sigue a una minúscula, y antes de la última mayúscula de una serie cuando le sigue una minúscula («TreeADDWithout» → «Tree ADD without», «XMLInvalid» → «XML invalid», «OnlyDESNets» → «Only DES nets»). La otra opción es dar a estas clases un título propio en el fichero de textos; arreglar la función cubre también las clases que se creen después.

Aparte, dos nombres de clase llevan erratas que se ven en el título: \texttt{Tree\allowbreak{}ADDWithout\allowbreak{}Potential\allowbreak{}Or\allowbreak{}Referenfe} («referenfe») y \texttt{Some\allowbreak{}Subpotentials\allowbreak{}Arent\allowbreak{}Linked\allowbreak{}To\allowbreak{}An\allowbreak{}ICIPotential} («arent»; es del formato Elvira).

Prueba de regresión: una prueba unitaria de \texttt{titlify\allowbreak{}Exception\allowbreak{}Name} con \texttt{Invalid\allowbreak{}Limit\allowbreak{}In\allowbreak{}Tree\allowbreak{}ADDException}, \texttt{Prob\allowbreak{}Net\allowbreak{}Parser\allowbreak{}Exception.\allowbreak{}XMLInvalid} y \texttt{Only\allowbreak{}DESNets\allowbreak{}Allowed\allowbreak{}Exception} que exija «Invalid limit in tree ADD», «XML invalid» y «Only DES nets allowed», y que un nombre sin siglas como \texttt{Tried\allowbreak{}To\allowbreak{}Split\allowbreak{}Interval\allowbreak{}Outside\allowbreak{}Bounds\allowbreak{}Exception} siga dando «Tried to split interval outside bounds».

\setcounter{subsection}{182}
\subsection[El aviso de zoom fuera de rango dice «Requested zoomManager»]{El aviso de zoom fuera de rango dice «Requested zoomManager» en vez de «Requested zoom»}\label{err:183}
\begin{ficha}
\fichakey{Estado}Presente
\fichakey{Módulo}gui
\fichakey{Paquete}org.openmarkov.gui.menutoolbar.toolbar
\fichakey{Ficheros}\path{gui/src/main/resources/gui/localize/gui_class_localizations_en.xml:102-103} \newline \path{gui/src/main/java/org/openmarkov/gui/menutoolbar/toolbar/ZoomComboBox.java:125-143} \newline \path{gui/src/main/resources/gui/localize/Dialogs_en.xml:907-909}
\fichakey{Comprobado}Leyendo el código y ejecutándolo
\fichakey{Esfuerzo}Pequeño (horas)
\fichakey{Decisión de equipo}No
\fichakey{Relacionados}\hyperref[err:182]{error~182}, \hyperref[nota:70]{nota~70}
\end{ficha}
\paragraph{Qué pasa}
La caja de zoom de la barra de herramientas (\texttt{Zoom\allowbreak{}Combo\allowbreak{}Box}) es editable. Si el usuario escribe un valor fuera del rango 10-500 \% o que no es un número (por ejemplo «600»), sale un aviso que dice «Requested zoomManager (600\%) is out of the range {[}10..500{]}.»: el usuario ve el nombre de un campo interno del programa donde debería leer «zoom».

\texttt{item\allowbreak{}State\allowbreak{}Changed} restaura el valor anterior y lanza \texttt{Zoom\allowbreak{}Out\allowbreak{}Of\allowbreak{}Range\allowbreak{}Exception} envuelta en \texttt{Unrecoverable\allowbreak{}Exception}; el gestor global muestra su mensaje, que sale de \texttt{gui\_\allowbreak{}class\_\allowbreak{}localizations\_\allowbreak{}en.\allowbreak{}xml}. Ese texto dice «Requested zoomManager (\{requestedZoom\}) is out of the range {[}\{minZoom\}..\{maxZoom\}{]}.».

Medido con un pequeño programa de prueba: el mensaje construido es \texttt{Requested\ zoom\allowbreak{}Manager\ (\allowbreak{}600\%)\ is\ out\ of\ the\ range\ {[}10.\allowbreak{}.\allowbreak{}500{]}.\allowbreak{}}

En \texttt{Dialogs\_\allowbreak{}en.\allowbreak{}xml} hay otro resto igual, \texttt{Select\allowbreak{}Zoom.\allowbreak{}Title} = «Select zoomManager», pero ningún código usa esa clave, así que no se ve.

El mismo cambio que dejó \texttt{zoom\allowbreak{}Manager} en estos textos estropeó también varios textos del fichero en castellano; están descritos en \hyperref[nota:123]{nota~123}, donde no hay nada que corregir.
\paragraph{El código}
gui/src/main/resources/gui/localize/gui\_class\_localizations\_en.xml:102-103

\begin{Shaded}
\begin{Highlighting}[]
\NormalTok{\textless{}}\KeywordTok{Localization}\OtherTok{ class=}\StringTok{"org.openmarkov.gui.exception.ZoomOutOfRangeException"}
\OtherTok{              value=}\StringTok{"Requested zoomManager (\{requestedZoom\}) is out of the range [\{minZoom\}..\{maxZoom\}]."}\NormalTok{/\textgreater{}}
\end{Highlighting}
\end{Shaded}
\paragraph{Cómo arreglarlo}
Cambiar el texto a «Requested zoom (\{requestedZoom\}) is out of the range {[}\{minZoom\}..\{maxZoom\}{]}.» y, de paso, «Select zoomManager» por «Select zoom» (o borrar esa clave, que nadie usa). Buscar \texttt{zoom\allowbreak{}Manager} en todos los ficheros \texttt{*\_\allowbreak{}en.\allowbreak{}xml} de todos los módulos: hoy solo aparece en esos dos.

Prueba de regresión: una prueba que recorra los ficheros de textos en inglés y falle si algún valor contiene un identificador con forma de nombre de campo Java como \texttt{zoom\allowbreak{}Manager}, o, más simple, que compruebe el mensaje de \texttt{Zoom\allowbreak{}Out\allowbreak{}Of\allowbreak{}Range\allowbreak{}Exception}.

\setcounter{subsection}{183}
\subsection[El rótulo del panel del indicador dice «Probability of ocurrence»]{El rótulo del panel del indicador dice «Probability of ocurrence», con una errata}\label{err:184}
\begin{ficha}
\fichakey{Estado}Presente
\fichakey{Módulo}gui
\fichakey{Paquete}org.openmarkov.gui.dialog.common
\fichakey{Ficheros}\path{gui/src/main/java/org/openmarkov/gui/dialog/common/IndicatorPotentialPanel.java:60,65}
\fichakey{Comprobado}Leyendo el código
\fichakey{Esfuerzo}Pequeño (horas)
\fichakey{Decisión de equipo}No
\fichakey{Relacionados}\hyperref[err:152]{error~152}, \hyperref[err:186]{error~186}, \hyperref[err:187]{error~187}, \hyperref[err:188]{error~188}, \hyperref[nota:120]{nota~120}, \hyperref[nota:140]{nota~140}
\end{ficha}
\paragraph{Qué pasa}
\texttt{Indicator\allowbreak{}Potential\allowbreak{}Panel} (módulo \texttt{gui}) es el panel con que se edita la relación «indicador» de un nodo de suceso en una red de simulación de eventos discretos (DES): cuándo ocurre el suceso (tiempo hasta el suceso) y con qué probabilidad.

El segundo de sus rótulos dice «Probability of ocurrence:» (línea 65), con una errata que se ve en la aplicación: en inglés se escribe «occurrence», con dos erres y dos ces.

Los dos rótulos, «Time-to-event:» (línea 60) y «Probability of ocurrence:», están escritos en el código. Que no salgan de un fichero de textos no tiene consecuencia en la aplicación en inglés y no forma parte de este error.
\paragraph{El código}
gui/src/main/java/org/openmarkov/gui/dialog/common/IndicatorPotentialPanel.java:65

\begin{Shaded}
\begin{Highlighting}[]
\BuiltInTok{JLabel}\NormalTok{ probabilityTextLabel }\OperatorTok{=} \KeywordTok{new} \BuiltInTok{JLabel}\OperatorTok{(}\StringTok{"Probability of ocurrence:"}\OperatorTok{);}
\end{Highlighting}
\end{Shaded}
\paragraph{Cómo arreglarlo}
Cambiar el texto a «Probability of occurrence:». Si se aprovecha para llevarlo a un fichero de textos, seguir el estilo del resto de paneles de relación. No hace falta prueba.

\setcounter{subsection}{184}
\subsection[La ayuda de «Add → Utility node» dice «Create an utility node»]{La ayuda emergente de «Add → Utility node», en el menú contextual del lienzo, dice «Create an utility node»}\label{err:185}
\begin{ficha}
\fichakey{Estado}Presente
\fichakey{Módulo}gui
\fichakey{Paquete}org.openmarkov.gui.menutoolbar.menu
\fichakey{Ficheros}\path{gui/src/main/java/org/openmarkov/gui/menutoolbar/menu/NetworkContextualMenu.java:106-133} \newline \path{gui/src/main/java/org/openmarkov/gui/componentBuilder/JMenuItemBuilder.java:166-167}
\fichakey{Comprobado}Leyendo el código
\fichakey{Esfuerzo}Pequeño (horas)
\fichakey{Decisión de equipo}No
\fichakey{Relacionados}\hyperref[err:188]{error~188}, \hyperref[nota:120]{nota~120}
\end{ficha}
\paragraph{Qué pasa}
Con una red abierta en modo edición, el clic derecho sobre una zona vacía del lienzo abre un menú con el submenú «Add», que ofrece «Chance node», «Decision node», «Utility node» y «Event node». Al dejar el ratón sobre «Utility node» aparece la ayuda emergente «Create an utility node»; en inglés es «a utility node», porque «utility» empieza por sonido de consonante. Las otras tres ayudas («Create a chance node», «Create a decision node», «Create an event node») están bien.

El texto está escrito en el código, en \texttt{Network\allowbreak{}Contextual\allowbreak{}Menu.\allowbreak{}get\allowbreak{}Create\allowbreak{}Elements\allowbreak{}Menu\allowbreak{}Item} (módulo \texttt{gui}, línea 114), y \texttt{JMenu\allowbreak{}Item\allowbreak{}Builder.\allowbreak{}build} lo pone como ayuda emergente con \texttt{set\allowbreak{}Tool\allowbreak{}Tip\allowbreak{}Text} (líneas 166-167). Este submenú es hoy una de las formas de crear nodos, además del doble clic en el lienzo; los botones de la barra de herramientas no están, por decisión.

Comprobado leyendo; no se ha abierto el menú. Fuera de este texto, «an utility» solo aparece en comentarios del código.
\paragraph{El código}
gui/src/main/java/org/openmarkov/gui/menutoolbar/menu/NetworkContextualMenu.java:114-116

\begin{Shaded}
\begin{Highlighting}[]
\KeywordTok{new} \FunctionTok{NodeMenuGenerator}\OperatorTok{(}\NormalTok{NodeType}\OperatorTok{.}\FunctionTok{UTILITY}\OperatorTok{,} \StringTok{"Utility node"}\OperatorTok{,} \StringTok{"jmenuItemCreateUtilityNode"}\OperatorTok{,}\StringTok{"Create an utility node"}\OperatorTok{,}
\NormalTok{                      IconBind}\OperatorTok{.}\FunctionTok{UTILITY\_ENABLED}\OperatorTok{.}\FunctionTok{icon}\OperatorTok{(),}
\NormalTok{                      ActionCommands}\OperatorTok{.}\FunctionTok{UTILITY\_CREATION}\OperatorTok{,} \OperatorTok{!}\NormalTok{currentNetwork}\OperatorTok{.}\FunctionTok{hasConstraintOfClass}\OperatorTok{(}\NormalTok{OnlyChanceNodes}\OperatorTok{.}\FunctionTok{class}\OperatorTok{)),}
\end{Highlighting}
\end{Shaded}
\paragraph{Cómo arreglarlo}
Cambiar el texto a «Create a utility node». Los rótulos y ayudas de este submenú están escritos en el código; pasarlos al fichero de textos es opcional y no hace falta para corregir la errata (\hyperref[nota:120]{nota~120}).

Comprobación: tras el cambio, «an utility» no debe aparecer en ningún texto visible de la interfaz.

\setcounter{subsection}{185}
\subsection[El mensaje de evidencia incompatible dice «old evidente» y es ilegible]{El mensaje de evidencia incompatible dice «old evidente» y describe cada hallazgo con su representación interna en lugar de nombrar la variable y el estado}\label{err:186}
\begin{ficha}
\fichakey{Estado}Presente
\fichakey{Módulo}core
\fichakey{Paquete}org.openmarkov.core.exception
\fichakey{Ficheros}\path{core/src/main/resources/core/localize/core_class_localizations_en.xml:219-222, 337-338} \newline \path{core/src/main/java/org/openmarkov/core/exception/IncompatibleEvidenceException.java:26-34} \newline \path{core/src/main/java/org/openmarkov/core/model/network/EvidenceCase.java:108-117} \newline \path{inference/src/main/java/org/openmarkov/inference/algorithm/variableElimination/tasks/VEPropagation.java:260-266} \newline \path{gui/src/main/java/org/openmarkov/gui/window/edition/networkEditorPanel/EditorInputHandler.java:260-269} \newline \path{gui/src/main/java/org/openmarkov/gui/dialog/UnexpectedThrowableDialog.java:35-39, 51, 101-110}
\fichakey{Comprobado}Leyendo el código y ejecutándolo
\fichakey{Esfuerzo}Pequeño (horas)
\fichakey{Decisión de equipo}No
\fichakey{Corregir antes}\hyperref[err:48]{error~48}: si se resuelve con un mensaje propio para la evidencia imposible, este texto deja de verse en los casos que cubra, y ese mensaje debe nombrar la variable y el estado.
\fichakey{Relacionados}\hyperref[err:48]{error~48}, \hyperref[err:184]{error~184}, \hyperref[err:187]{error~187}
\end{ficha}
\paragraph{Qué pasa}
Cuando un hallazgo contradice a otro sobre la misma variable, OpenMarkov lanza \texttt{Incompatible\allowbreak{}Evidence\allowbreak{}Exception\$Evidence\allowbreak{}Is\allowbreak{}Incompatible\allowbreak{}With\allowbreak{}Other} y la ventana de error enseña su mensaje, que sale del fichero de textos en inglés (\texttt{core\_\allowbreak{}class\_\allowbreak{}localizations\_\allowbreak{}en.\allowbreak{}xml}, línea 220): «The new evidence \{newFinding\} is not compatible with an old evidente: \{oldFinding\}».

El usuario ve dos fallos en ese texto:

\begin{itemize}
\tightlist
\item
  «evidente» en lugar de «evidence»;
\item
  los dos huecos no se rellenan con la variable y el estado, sino con la descripción genérica de un objeto \texttt{Finding} (línea 338: «Finding of variable \{variable\} with numerical value \{numericalValue\} and state index \{stateIndex\}»), que a su vez describe la variable con la lista entera de sus estados.
\end{itemize}

Camino del usuario: abrir \texttt{0\_\allowbreak{}miscellaneous/\allowbreak{}networks/\allowbreak{}mid/\allowbreak{}MID-hip-Briggs.\allowbreak{}pgmx}, entrar en modo inferencia (funciona) y hacer doble clic en el estado «dead» de «State {[}0{]}». La tabla de «State {[}0{]}» da probabilidad 1 a «primary THR», y antes de propagar \texttt{VEPropagation} añade a la evidencia ese hallazgo forzado por el modelo; al juntarlo con el del usuario (\texttt{get\allowbreak{}All\allowbreak{}Evidence}, líneas 260-266), \texttt{Evidence\allowbreak{}Case.\allowbreak{}add\allowbreak{}Finding} lanza la excepción. El doble clic la envuelve como error no previsto (\texttt{Editor\allowbreak{}Input\allowbreak{}Handler}, líneas 270-277), y la ventana de error no previsto pone el mensaje al principio de su texto; por el diálogo «Add finding» sale en una ventana de error normal titulada «Evidence is incompatible with other». Pasa lo mismo con los otros estados de «State {[}0{]}» y, en \texttt{0\_\allowbreak{}miscellaneous/\allowbreak{}networks/\allowbreak{}mid/\allowbreak{}MID-Chancellor-corrected.\allowbreak{}pgmx}, \texttt{MID-Chancellor-new.\allowbreak{}pgmx}, \texttt{DMHEE/\allowbreak{}MID-dmhee-3.\allowbreak{}5.\allowbreak{}pgmx} y \texttt{DMHEE/\allowbreak{}MID-dmhee-4.\allowbreak{}8.\allowbreak{}pgmx}. Que la evidencia imposible acabe en una ventana de error, y en cuál, es \hyperref[err:48]{error~48}; esta ficha trata solo del texto.

Medido con un pequeño programa de prueba que hace la propagación del modo inferencia con ese hallazgo en \texttt{MID-hip-Briggs.\allowbreak{}pgmx} y pide a la excepción su título y su mensaje:

\begin{Verbatim}[breaklines,breakanywhere,fontsize=\small]
The new evidence Finding of variable Finite states variable State [0] with states [state dead,
state successful revised, state revision THR, state successful primary, state primary THR] with
numerical value NaN and state index 0 is not compatible with an old evidente: Finding of variable
Finite states variable State [0] with states [state dead, state successful revised, state revision
THR, state successful primary, state primary THR] with numerical value NaN and state index 4
\end{Verbatim}

Con dos hallazgos contradictorios puestos directamente en un \texttt{Evidence\allowbreak{}Case} sobre «X-ray» de \texttt{0\_\allowbreak{}miscellaneous/\allowbreak{}networks/\allowbreak{}bn/\allowbreak{}BN-asia.\allowbreak{}pgmx} el formato es el mismo.
\paragraph{El código}
core/src/main/resources/core/localize/core\_class\_localizations\_en.xml:219-222

\begin{Shaded}
\begin{Highlighting}[]
\NormalTok{\textless{}}\KeywordTok{Localization}\OtherTok{ class=}\StringTok{"org.openmarkov.core.exception.IncompatibleEvidenceException$EvidenceIsIncompatibleWithOther"}
\OtherTok{              value=}\StringTok{"The new evidence \{newFinding\} is not compatible with an old evidente: \{oldFinding\}"}\NormalTok{/\textgreater{}}
\NormalTok{\textless{}}\KeywordTok{Localization}\OtherTok{ class=}\StringTok{"org.openmarkov.core.exception.IncompatibleEvidenceException$FindingVariableIsMissingAState"}
\OtherTok{              value=}\StringTok{"The variable \{variable.getName\} has no state \{state\}, so this finding cannot be translated."}\NormalTok{/\textgreater{}}
\end{Highlighting}
\end{Shaded}

core/src/main/resources/core/localize/core\_class\_localizations\_en.xml:337-338

\begin{Shaded}
\begin{Highlighting}[]
\NormalTok{\textless{}}\KeywordTok{Localization}\OtherTok{ class=}\StringTok{"org.openmarkov.core.model.network.Finding"}
\OtherTok{              value=}\StringTok{"Finding of variable \{variable\} with numerical value \{numericalValue\} and state index \{stateIndex\}"}\NormalTok{/\textgreater{}}
\end{Highlighting}
\end{Shaded}
\paragraph{Cómo arreglarlo}
Reescribir el texto de la línea 220 para que nombre la variable y los dos estados, con la misma forma de hueco encadenado que ya usa la línea 222 (\texttt{\{variable.\allowbreak{}get\allowbreak{}Name\}}); por ejemplo, «The finding \{newFinding.getVariable.getName\} = \{newFinding.getState\} contradicts the finding \{oldFinding.getVariable.getName\} = \{oldFinding.getState\} already in the evidence». Para un hallazgo sobre una variable numérica, \texttt{get\allowbreak{}State} no da el valor observado: si ese caso importa, el texto necesita otro hueco para el número (\texttt{get\allowbreak{}Numerical\allowbreak{}Value}).

Es el único texto que usa \texttt{\{new\allowbreak{}Finding\}} y \texttt{\{old\allowbreak{}Finding\}}; la descripción genérica de \texttt{Finding} (línea 338) no hace falta tocarla para esto.

Si antes se resuelve \hyperref[err:48]{error~48} con una excepción nueva y un mensaje propio para la evidencia imposible, ese mensaje debe nombrar también variable y estado, y este texto dejará de verse en los casos que ella cubra.

Prueba: una prueba de \texttt{core} que construya \texttt{Evidence\allowbreak{}Is\allowbreak{}Incompatible\allowbreak{}With\allowbreak{}Other} con dos hallazgos sobre «X-ray» y compruebe que \texttt{get\allowbreak{}Exception\allowbreak{}Message(\allowbreak{})} contiene «X-ray», «yes» y «no», y no contiene «evidente» ni «state index».

\setcounter{subsection}{186}
\subsection[El error de muestras sin peso dice «Samples weigth is 0»]{Cuando ninguna muestra de la exportación de la propagación estocástica es compatible con la evidencia, la ventana de error dice «Samples weigth is 0»}\label{err:187}
\begin{ficha}
\fichakey{Estado}Presente
\fichakey{Módulo}core
\fichakey{Paquete}org.openmarkov.core.exception
\fichakey{Ficheros}\path{core/src/main/resources/core/localize/core_class_localizations_en.xml:223-224, 463-464} \newline \path{core/src/main/java/org/openmarkov/core/exception/IncompatibleEvidenceException.java:57-63} \newline \path{inference/src/main/java/org/openmarkov/inference/algorithm/likelihoodWeighting/StochasticPropagation.java:179-181} \newline \path{stochasticPropagationOutput/src/main/java/org/openmarkov/stochasticPropagationOutput/StochasticPropagationOutputFrame.java:194-202} \newline \path{gui/src/main/resources/gui/localize/Menus_en.xml:269-272}
\fichakey{Comprobado}Leyendo el código y ejecutándolo
\fichakey{Esfuerzo}Pequeño (horas)
\fichakey{Decisión de equipo}No
\fichakey{Relacionados}\hyperref[err:184]{error~184}, \hyperref[err:186]{error~186}
\end{ficha}
\paragraph{Qué pasa}
El menú Tools → «Export stochastic propagation data» abre la ventana «Stochastic propagation output», que muestrea la red bayesiana abierta con la evidencia puesta en el modo inferencia y guarda las muestras en una hoja de cálculo. Si ninguna muestra es compatible con la evidencia, al pulsar «Export to a .xlsx file» sale una ventana de error titulada «Samples weight is zero» con el mensaje «Samples weigth is 0»: «weigth» en lugar de «weight».

Camino del usuario: abrir \texttt{0\_\allowbreak{}miscellaneous/\allowbreak{}networks/\allowbreak{}bn/\allowbreak{}BN-asia.\allowbreak{}pgmx}, entrar en modo inferencia y poner los hallazgos VisitToAsia = yes, Tuberculosis = yes y X-ray = no (la propagación exacta funciona); abrir la herramienta, elegir «Logic sampling», dejar las 10 000 muestras de omisión y pulsar «Export to a .xlsx file». \texttt{Stochastic\allowbreak{}Propagation.\allowbreak{}get\allowbreak{}Posterior\allowbreak{}Values} lanza \texttt{Incompatible\allowbreak{}Evidence\allowbreak{}Exception.\allowbreak{}Samples\allowbreak{}Weight\allowbreak{}Is\allowbreak{}Zero} cuando el peso acumulado de las muestras es cero; \texttt{Stochastic\allowbreak{}Propagation\allowbreak{}Output\allowbreak{}Frame.\allowbreak{}action\allowbreak{}Performed} la envuelve en \texttt{Unrecoverable\allowbreak{}Exception}, y el manejador global la enseña en una ventana de error normal, con el título fabricado a partir del nombre de la clase y el mensaje del fichero de textos.

El fichero \texttt{core\_\allowbreak{}class\_\allowbreak{}localizations\_\allowbreak{}en.\allowbreak{}xml} tiene dos entradas para esa excepción:

\begin{itemize}
\tightlist
\item
  líneas 463-464, con el nombre de clase correcto (\texttt{Samples\allowbreak{}Weight\allowbreak{}Is\allowbreak{}Zero}) y el texto «Samples weigth is 0»: es la que se ve;
\item
  líneas 223-224, con el nombre de clase mal escrito (\texttt{Samples\allowbreak{}Weigth\allowbreak{}Is\allowbreak{}Zero}) y el texto «Samples weigth is zero»: no corresponde a ninguna clase y no se ve nunca.
\end{itemize}

Medido con un pequeño programa de prueba que repite las llamadas de la ventana con esos hallazgos y 10 000 muestras, veinte veces por algoritmo (la ventana no fija la semilla): con «Logic sampling», 17 de las 20 ejecuciones lanzan la excepción, con título «Samples weight is zero» y mensaje «Samples weigth is 0»; con «Likelihood weighting», que es el algoritmo marcado por omisión, ninguna.

En \texttt{gui/\allowbreak{}src/\allowbreak{}main/\allowbreak{}resources/\allowbreak{}gui/\allowbreak{}localize/\allowbreak{}Menus\_\allowbreak{}en.\allowbreak{}xml} (línea 271) hay otra errata, «An error ocurred while trying to process command:» («ocurred» por «occurred»), en la clave \texttt{Tools.\allowbreak{}Plugin.\allowbreak{}Error}. Ningún código lee esa clave ni su vecina \texttt{Tools.\allowbreak{}Plugin.\allowbreak{}Not\allowbreak{}Available} (buscado en las fuentes y en las clases del jar completo), así que no se ve.
\paragraph{El código}
core/src/main/resources/core/localize/core\_class\_localizations\_en.xml:223-224

\begin{Shaded}
\begin{Highlighting}[]
\NormalTok{\textless{}}\KeywordTok{Localization}\OtherTok{ class=}\StringTok{"org.openmarkov.core.exception.IncompatibleEvidenceException$SamplesWeigthIsZero"}
\OtherTok{              value=}\StringTok{"Samples weigth is zero"}\NormalTok{/\textgreater{}}
\end{Highlighting}
\end{Shaded}

core/src/main/resources/core/localize/core\_class\_localizations\_en.xml:463-464

\begin{Shaded}
\begin{Highlighting}[]
\NormalTok{\textless{}}\KeywordTok{Localization}\OtherTok{ class=}\StringTok{"org.openmarkov.core.exception.IncompatibleEvidenceException$SamplesWeightIsZero"}
\OtherTok{              value=}\StringTok{"Samples weigth is 0"}\NormalTok{/\textgreater{}}
\end{Highlighting}
\end{Shaded}
\paragraph{Cómo arreglarlo}
Corregir el texto de las líneas 463-464 («Samples weight is 0») y borrar la entrada de las líneas 223-224, que no corresponde a ninguna clase. Si se aprovecha para mejorar el mensaje, puede decir que ninguna muestra es compatible con la evidencia y que se pruebe con más muestras o con «Likelihood weighting».

En \texttt{Menus\_\allowbreak{}en.\allowbreak{}xml}, corregir «ocurred» o quitar el bloque \texttt{Tools.\allowbreak{}Plugin}, que nadie usa.

Prueba: una prueba de \texttt{core} que construya \texttt{Samples\allowbreak{}Weight\allowbreak{}Is\allowbreak{}Zero} y compruebe que \texttt{get\allowbreak{}Exception\allowbreak{}Message(\allowbreak{})} no contiene «weigth»; y, si se quiere cubrir las demás, una que busque «weigth» y «ocurred» en los ficheros \texttt{*\_\allowbreak{}en.\allowbreak{}xml} de todos los módulos.

\setcounter{subsection}{187}
\subsection[La ayuda de «Unique superparent» dice «commont»]{La ayuda emergente de «Unique superparent», en las opciones del clasificador de superpadre, dice «commont» en vez de «common»}\label{err:188}
\begin{ficha}
\fichakey{Estado}Presente
\fichakey{Módulo}learning.gui
\fichakey{Paquete}org.openmarkov.learning.algorithm.nbderived.spnb.gui
\fichakey{Ficheros}\path{learning.gui/src/main/resources/learning/gui/localize/Learning_en.xml:154-156} \newline \path{learning.gui/src/main/java/org/openmarkov/learning/algorithm/nbderived/spnb/gui/SuperParentNaiveBayesParametersDialog.java:30-31}
\fichakey{Comprobado}Leyendo el código y ejecutándolo
\fichakey{Esfuerzo}Pequeño (horas)
\fichakey{Decisión de equipo}No
\fichakey{Relacionados}\hyperref[err:184]{error~184}, \hyperref[err:185]{error~185}
\end{ficha}
\paragraph{Qué pasa}
En Herramientas («Tools») → «Learning», con el tipo «Discriminative» y el algoritmo «Superparent naive bayes», el botón «Options» abre la ventana «SPNB Algorithm : Options». Al dejar el ratón sobre el rótulo «Unique superparent» aparece la ayuda emergente «If selected only the superparent would be commont to all features»: «commont» en lugar de «common».

\texttt{Super\allowbreak{}Parent\allowbreak{}Naive\allowbreak{}Bayes\allowbreak{}Parameters\allowbreak{}Dialog} (módulo \texttt{learning.\allowbreak{}gui}, líneas 30-31) pone como ayuda del rótulo el texto de la clave \texttt{Learning.\allowbreak{}Super\allowbreak{}Parent\allowbreak{}Naive\allowbreak{}Bayes.\allowbreak{}Same\allowbreak{}SP.\allowbreak{}Tooltip}, que está en \texttt{Learning\_\allowbreak{}en.\allowbreak{}xml}, línea 155.

Comprobado leyendo el código y pidiendo la clave a \texttt{String\allowbreak{}Database}, que devuelve ese texto; la ventana no se ha abierto. «commont» no aparece en ningún otro fichero de textos.
\paragraph{El código}
learning.gui/src/main/resources/learning/gui/localize/Learning\_en.xml:154-156

\begin{Shaded}
\begin{Highlighting}[]
\NormalTok{\textless{}}\KeywordTok{SameSP}\OtherTok{ value=}\StringTok{"Unique superparent "}\NormalTok{\textgreater{}}
\NormalTok{    \textless{}}\KeywordTok{Tooltip}\OtherTok{ value=}\StringTok{"If selected only the superparent would be commont to all features"}\NormalTok{/\textgreater{}}
\NormalTok{\textless{}/}\KeywordTok{SameSP}\NormalTok{\textgreater{}}
\end{Highlighting}
\end{Shaded}
\paragraph{Cómo arreglarlo}
Cambiar «commont» por «common» en la línea 155 de \texttt{Learning\_\allowbreak{}en.\allowbreak{}xml}. No hace falta prueba.

% ══════════════════════════════════════════════════════════════════════
\section{Rendimiento}\label{sec:rendimiento}

Trabajo de más en el camino que más se recorre, y operaciones largas que bloquean la ventana.

\setcounter{subsection}{188}
\subsection[Entrar en modo inferencia con una red grande congela la ventana]{Entrar en modo inferencia con una red grande deja la ventana entera congelada durante decenas de segundos, sin progreso ni forma de cancelar}\label{err:189}
\begin{ficha}
\fichakey{Estado}Presente
\fichakey{Módulo}gui
\fichakey{Paquete}org.openmarkov.gui.window.edition.networkEditorPanel
\fichakey{Ficheros}\path{gui/src/main/java/org/openmarkov/gui/window/MainPanelListenerAssistant.java:146-163, 180-181} \newline \path{gui/src/main/java/org/openmarkov/gui/window/InferenceHandler.java:58-119} \newline \path{gui/src/main/java/org/openmarkov/gui/window/edition/networkEditorPanel/EvidenceManager.java:324-363, 580-599} \newline \path{gui/src/main/java/org/openmarkov/gui/util/GUIUtils.java:57-81} \newline \path{gui/src/main/java/org/openmarkov/gui/window/MainGUI.java:108-118} \newline \path{gui/src/main/java/org/openmarkov/gui/window/edition/networkEditorPanel/NetworkEditorPanel.java:465-472}
\fichakey{Comprobado}Leyendo el código y ejecutándolo
\fichakey{Esfuerzo}Medio (días)
\fichakey{Decisión de equipo}\textcolor{decisionrojo}{\textbf{Sí.}} Qué hacer con \texttt{Main\allowbreak{}GUI.\allowbreak{}freeze(\allowbreak{})/\allowbreak{}unfreeze(\allowbreak{})} (borrarlo o reactivarlo)
\fichakey{Relacionados}\hyperref[nota:105]{nota~105}, \hyperref[nota:167]{nota~167}, \hyperref[nota:176]{nota~176}, \hyperref[nota:177]{nota~177}
\end{ficha}
\paragraph{Qué pasa}
Entrar en modo inferencia con \texttt{integration\allowbreak{}Tests/\allowbreak{}src/\allowbreak{}main/\allowbreak{}resources/\allowbreak{}networks/\allowbreak{}bn/\allowbreak{}BN-cpcs422b.\allowbreak{}pgmx} (422 nodos) deja la ventana sin responder unos 29 segundos en la máquina donde se midió (más en una lenta), y cada hallazgo nuevo otro tanto. No hay barra de progreso ni botón de cancelar, y el usuario no sabe si el programa se ha colgado. Lo mismo ocurre al propagar evidencia (\texttt{PROPAGATE\_\allowbreak{}EVIDENCE}, líneas 188-189 de \texttt{Main\allowbreak{}Panel\allowbreak{}Listener\allowbreak{}Assistant}) y, con propagación automática, al añadir o quitar cada hallazgo, que repite la propagación para cada caso de evidencia en memoria.

Swing tiene un único hilo que pinta la ventana y atiende los clics (el EDT, \emph{Event Dispatch Thread}). Todo lo que se ejecuta dentro de un \texttt{action\allowbreak{}Performed} ocupa ese hilo, y mientras dura la ventana no se repinta ni responde.

En OpenMarkov, pulsar el botón de modo inferencia llega a \texttt{Main\allowbreak{}Panel\allowbreak{}Listener\allowbreak{}Assistant.\allowbreak{}action\allowbreak{}Performed} (líneas 155-171) → \texttt{Inference\allowbreak{}Handler.\allowbreak{}toggle\allowbreak{}Working\allowbreak{}Mode}/\texttt{set\allowbreak{}Working\allowbreak{}Mode} → \texttt{Network\allowbreak{}Editor\allowbreak{}Panel.\allowbreak{}update\allowbreak{}Individual\allowbreak{}Probabilities\allowbreak{}And\allowbreak{}Utilities} → \texttt{Evidence\allowbreak{}Manager.\allowbreak{}update\allowbreak{}Individual\allowbreak{}Probabilities\allowbreak{}And\allowbreak{}Utilities} (líneas 328-360), que llama a \texttt{do\allowbreak{}Propagation} para cada caso de evidencia (líneas 580-599), y esta construye un \texttt{VEPropagation} y le pide todas las probabilidades a posteriori. Todo en el mismo hilo: \texttt{GUIUtils.\allowbreak{}execute\allowbreak{}UIAction} (líneas 57-81) solo envuelve excepciones, no crea ningún hilo. \texttt{do\allowbreak{}Propagation} captura \texttt{Out\allowbreak{}Of\allowbreak{}Memory\allowbreak{}Error}, pero no hay forma de interrumpir un cálculo largo.

Medido con un pequeño programa de prueba que hace la misma llamada que \texttt{do\allowbreak{}Propagation} (un \texttt{VEPropagation} con todas las variables de interés, sin evidencia): tarda 330 ms en \texttt{integration\allowbreak{}Tests/\allowbreak{}src/\allowbreak{}main/\allowbreak{}resources/\allowbreak{}networks/\allowbreak{}bn/\allowbreak{}BN-cpcs179.\allowbreak{}pgmx} (179 nodos) y \textbf{28 688 ms} en \texttt{integration\allowbreak{}Tests/\allowbreak{}src/\allowbreak{}main/\allowbreak{}resources/\allowbreak{}networks/\allowbreak{}bn/\allowbreak{}BN-cpcs422b.\allowbreak{}pgmx} (422 nodos). Leer el fichero tarda 150-500 ms, así que abrir no es el problema. No se ha arrancado la aplicación: el tiempo es el del cálculo que la ventana ejecuta en su hilo, medido fuera de ella.

Aparte, \texttt{Main\allowbreak{}GUI.\allowbreak{}freeze(\allowbreak{})} y \texttt{unfreeze(\allowbreak{})} empiezan con \texttt{if\ (\allowbreak{}true)\ return;} (líneas 108-118), así que no hacen nada. Su único uso es alrededor del diálogo modal de propiedades del nodo (\texttt{Network\allowbreak{}Editor\allowbreak{}Panel.\allowbreak{}java:472-482}); no está relacionado con la propagación y desactivado no produce ningún efecto visible. Hay que decidir si se borra o se reactiva.
\paragraph{El código}
gui/src/main/java/org/openmarkov/gui/window/edition/networkEditorPanel/EvidenceManager.java:580-588

\begin{Shaded}
\begin{Highlighting}[]
\KeywordTok{public} \DataTypeTok{void} \FunctionTok{doPropagation}\OperatorTok{(}\NormalTok{EvidenceCase evidenceCase}\OperatorTok{,} \DataTypeTok{int}\NormalTok{ caseNumber}\OperatorTok{)} \KeywordTok{throws} \KeywordTok{...} \OperatorTok{\{}
    \BuiltInTok{Map}\OperatorTok{\textless{}}\NormalTok{Variable}\OperatorTok{,}\NormalTok{ TablePotential}\OperatorTok{\textgreater{}}\NormalTok{ individualProbabilities }\OperatorTok{=} \KeywordTok{null}\OperatorTok{;}
    \ControlFlowTok{try} \OperatorTok{\{}
        \KeywordTok{this}\OperatorTok{.}\FunctionTok{calculateMinAndMaxUtilityRanges}\OperatorTok{();}
\NormalTok{        Propagation vePosteriorValues }\OperatorTok{=} \KeywordTok{new} \FunctionTok{VEPropagation}\OperatorTok{(}\KeywordTok{this}\OperatorTok{.}\FunctionTok{networkEditorPanel}\OperatorTok{.}\FunctionTok{getVisualNetwork}\OperatorTok{().}\FunctionTok{getProbNet}\OperatorTok{());}
\NormalTok{        vePosteriorValues}\OperatorTok{.}\FunctionTok{setVariablesOfInterest}\OperatorTok{(}\KeywordTok{this}\OperatorTok{.}\FunctionTok{networkEditorPanel}\OperatorTok{.}\FunctionTok{getVisualNetwork}\OperatorTok{().}\FunctionTok{getProbNet}\OperatorTok{().}\FunctionTok{getVariables}\OperatorTok{());}
\NormalTok{        vePosteriorValues}\OperatorTok{.}\FunctionTok{setPreResolutionEvidence}\OperatorTok{(}\KeywordTok{this}\OperatorTok{.}\FunctionTok{preResolutionEvidence}\OperatorTok{);}
\NormalTok{        vePosteriorValues}\OperatorTok{.}\FunctionTok{setPostResolutionEvidence}\OperatorTok{(}\NormalTok{evidenceCase}\OperatorTok{);}
\NormalTok{        individualProbabilities }\OperatorTok{=}\NormalTok{ vePosteriorValues}\OperatorTok{.}\FunctionTok{getPosteriorValues}\OperatorTok{();}
\end{Highlighting}
\end{Shaded}

gui/src/main/java/org/openmarkov/gui/window/MainGUI.java:108-118

\begin{Shaded}
\begin{Highlighting}[]
\KeywordTok{public} \DataTypeTok{void} \FunctionTok{freeze}\OperatorTok{()} \OperatorTok{\{}
    \ControlFlowTok{if} \OperatorTok{(}\KeywordTok{true}\OperatorTok{)} \ControlFlowTok{return}\OperatorTok{;}
    \KeywordTok{this}\OperatorTok{.}\FunctionTok{frameMirror}\OperatorTok{.}\FunctionTok{freeze}\OperatorTok{();}
\OperatorTok{\}}

\CommentTok{//Conditionally disabled}
\KeywordTok{public} \DataTypeTok{void} \FunctionTok{unfreeze}\OperatorTok{()} \OperatorTok{\{}
    \ControlFlowTok{if} \OperatorTok{(}\KeywordTok{true}\OperatorTok{)} \ControlFlowTok{return}\OperatorTok{;}
    \KeywordTok{this}\OperatorTok{.}\FunctionTok{frameMirror}\OperatorTok{.}\FunctionTok{unfreeze}\OperatorTok{();}
\OperatorTok{\}}
\end{Highlighting}
\end{Shaded}
\paragraph{Cómo arreglarlo}
Ejecutar el cálculo fuera del hilo de la interfaz con un \texttt{Swing\allowbreak{}Worker} (el patrón que ya usa \texttt{Splash\allowbreak{}Screen\allowbreak{}Loader}): el trabajador construye la propagación y calcula; al terminar, \texttt{done(\allowbreak{})} pinta los resultados (\texttt{Inference\allowbreak{}Presenter.\allowbreak{}paint\allowbreak{}Inference\allowbreak{}Results}) y actualiza menús en el hilo de la interfaz. Mientras dura, un diálogo de progreso con cancelar.

Conviene un único servicio para todos los caminos que llaman a \texttt{do\allowbreak{}Propagation}:

\begin{itemize}
\tightlist
\item
  entrar en modo inferencia (\texttt{Inference\allowbreak{}Handler.\allowbreak{}set\allowbreak{}Working\allowbreak{}Mode});
\item
  propagar evidencia (\texttt{Network\allowbreak{}Editor\allowbreak{}Panel.\allowbreak{}propagate\allowbreak{}Evidence});
\item
  añadir o quitar hallazgos con propagación automática;
\item
  cambiar de caso de evidencia (\texttt{Inference\allowbreak{}Handler.\allowbreak{}evidence\allowbreak{}Cases\allowbreak{}Navigation\allowbreak{}Option}).
\end{itemize}

Cuidado: la red no debe poder editarse mientras el trabajador lee de ella, y el cálculo trabaja sobre la copia del algoritmo, que comparte potenciales con la red del usuario. La cancelación de verdad necesita que la tarea de inferencia la admita; esa propuesta de diseño está en \hyperref[nota:167]{nota~167}, donde no hay nada que corregir.

Prueba: con \texttt{integration\allowbreak{}Tests/\allowbreak{}src/\allowbreak{}main/\allowbreak{}resources/\allowbreak{}networks/\allowbreak{}bn/\allowbreak{}BN-cpcs422b.\allowbreak{}pgmx}, comprobar (por ejemplo con AssertJ-Swing) que durante la propagación el hilo de la interfaz sigue atendiendo eventos.

\setcounter{subsection}{189}
\subsection[Cada resultado de una operación sobre tablas reserva una tabla que tira]{Cada resultado de una operación sobre tablas reserva y rellena una tabla entera que tira acto seguido}\label{err:190}
\begin{ficha}
\fichakey{Estado}Presente
\fichakey{Módulo}core
\fichakey{Paquete}org.openmarkov.core.model.network.potential
\fichakey{Ficheros}\path{core/src/main/java/org/openmarkov/core/model/network/potential/TablePotential.java:55-75}
\fichakey{Comprobado}Leyendo el código y ejecutándolo
\fichakey{Esfuerzo}Pequeño (horas)
\fichakey{Decisión de equipo}No
\fichakey{Relacionados}\hyperref[err:191]{error~191}, \hyperref[err:192]{error~192}
\end{ficha}
\paragraph{Qué pasa}
Toda operación de la inferencia sobre tablas gasta tiempo en preparar una tabla que no usa. Con tablas grandes, el desperdicio es del orden de una cuarta parte de la operación.

Las operaciones (\texttt{multiply}, \texttt{sum}, \texttt{multiply\allowbreak{}And\allowbreak{}Marginalize}, \texttt{divide}, \texttt{evaluate\allowbreak{}Function\allowbreak{}Potential}, la fusión\ldots) calculan primero el array de resultados y luego construyen el potencial con \texttt{new\ Table\allowbreak{}Potential(\allowbreak{}variables,\allowbreak{}\ papel,\allowbreak{}\ tabla)}. Ese constructor empieza por \texttt{this(\allowbreak{}variables,\allowbreak{}\ role)}, que hace \texttt{values\ =\ new\ double{[}table\allowbreak{}Size{]}} (la máquina virtual pone a cero todas las casillas) y \texttt{set\allowbreak{}Uniform(\allowbreak{})} (dos rellenos más si el papel es de probabilidad, ver \hyperref[err:191]{error~191}). Después comprueba la longitud de la tabla recibida y hace \texttt{this.\allowbreak{}values\ =\ table}, tirando el array recién preparado. El tamaño de esas tablas crece exponencialmente con el número de variables.

Medido con un pequeño programa de prueba (tabla de 4 194 304 casillas): construir con papel de probabilidad condicionada cuesta unos 11,8 ms, de los que 5,8 son la copia del array que el propio programa pasa; con papel de utilidad, 6,7 ms. Es decir, alrededor de 6 ms tirados por tabla de probabilidad (reserva más dos rellenos) y alrededor de 1 ms por tabla de utilidad (reserva). Una multiplicación que produce una tabla de ese tamaño tarda unos 23 ms.
\paragraph{El código}
core/src/main/java/org/openmarkov/core/model/network/potential/TablePotential.java:55-75

\begin{Shaded}
\begin{Highlighting}[]
\KeywordTok{public} \FunctionTok{TablePotential}\OperatorTok{(}\BuiltInTok{List}\OperatorTok{\textless{}}\NormalTok{Variable}\OperatorTok{\textgreater{}}\NormalTok{ variables}\OperatorTok{,}\NormalTok{ PotentialRole role}\OperatorTok{)} \OperatorTok{\{}
    \KeywordTok{super}\OperatorTok{(}\NormalTok{variables}\OperatorTok{,}\NormalTok{ role}\OperatorTok{);} \CommentTok{// AbstractIndexedPotential: computes dimensions, offsets, tableSize}
\NormalTok{    values }\OperatorTok{=} \KeywordTok{new} \DataTypeTok{double}\OperatorTok{[}\NormalTok{tableSize}\OperatorTok{];}
    \FunctionTok{setUniform}\OperatorTok{();} \CommentTok{// Initializes the table as an uniform potential}
\OperatorTok{\}}

\KeywordTok{public} \FunctionTok{TablePotential}\OperatorTok{(}\BuiltInTok{List}\OperatorTok{\textless{}}\NormalTok{Variable}\OperatorTok{\textgreater{}}\NormalTok{ variables}\OperatorTok{,}\NormalTok{ PotentialRole role}\OperatorTok{,} \DataTypeTok{double}\OperatorTok{[]}\NormalTok{ table}\OperatorTok{)} \OperatorTok{\{}
    \KeywordTok{this}\OperatorTok{(}\NormalTok{variables}\OperatorTok{,}\NormalTok{ role}\OperatorTok{);}
    \ControlFlowTok{if} \OperatorTok{(}\NormalTok{table}\OperatorTok{.}\FunctionTok{length} \OperatorTok{!=}\NormalTok{ tableSize}\OperatorTok{)} \OperatorTok{\{}
        \ControlFlowTok{throw} \KeywordTok{new} \FunctionTok{InvalidArgumentException}\OperatorTok{(}\NormalTok{table}\OperatorTok{.}\FunctionTok{length}\OperatorTok{,} \StringTok{"table"}\OperatorTok{,}
                \StringTok{"its length must be "} \OperatorTok{+}\NormalTok{ tableSize }\OperatorTok{+} \StringTok{", the product of the numbers of states of the variables"}\OperatorTok{);}
    \OperatorTok{\}}
    \KeywordTok{this}\OperatorTok{.}\FunctionTok{values} \OperatorTok{=}\NormalTok{ table}\OperatorTok{;}
\OperatorTok{\}}
\end{Highlighting}
\end{Shaded}
\paragraph{Cómo arreglarlo}
Que el constructor de tres argumentos llame a \texttt{super(\allowbreak{}variables,\allowbreak{}\ role)} (que calcula dimensiones, desplazamientos y \texttt{table\allowbreak{}Size}), compruebe la longitud y asigne \texttt{this.\allowbreak{}values\ =\ table}, sin reservar ni llamar a \texttt{set\allowbreak{}Uniform}.

Revisar las tres subclases de \texttt{Table\allowbreak{}Potential} (\texttt{Strategic\allowbreak{}Table\allowbreak{}Potential} y \texttt{Uncertain\allowbreak{}Table\allowbreak{}Potential}, que llaman a \texttt{super(\allowbreak{}variables,\allowbreak{}\ role,\allowbreak{}\ table)}, y \texttt{Table\allowbreak{}With\allowbreak{}Functions}) por si alguna dependía del relleno.

Verificar con el banco de referencia de las operaciones, \texttt{Operations\allowbreak{}Reference\allowbreak{}Times\allowbreak{}Test} en \texttt{integration\allowbreak{}Tests} (red cpcs422b), antes y después; la batería entera de pruebas debe pasar sin cambios de valores.

\setcounter{subsection}{190}
\subsection[Cada tabla de probabilidad nueva se rellena dos veces con el mismo valor]{Cada tabla de probabilidad nueva se rellena dos veces seguidas con el mismo valor}\label{err:191}
\begin{ficha}
\fichakey{Estado}Presente
\fichakey{Módulo}core
\fichakey{Paquete}org.openmarkov.core.model.network.potential
\fichakey{Ficheros}\path{core/src/main/java/org/openmarkov/core/model/network/potential/TablePotential.java:465-513}
\fichakey{Comprobado}Leyendo el código y ejecutándolo
\fichakey{Esfuerzo}Pequeño (horas)
\fichakey{Decisión de equipo}No
\fichakey{Relacionados}\hyperref[err:190]{error~190}
\end{ficha}
\paragraph{Qué pasa}
Cada tabla de probabilidad que se construye, incluidas las que produce cada operación de la inferencia, se rellena dos veces seguidas con el mismo valor uniforme. Con tablas grandes es tiempo perdido en el camino más frecuente de la inferencia.

\texttt{Table\allowbreak{}Potential.\allowbreak{}set\allowbreak{}Uniform(\allowbreak{})} inicializa una tabla como distribución uniforme. Lo llama el constructor de dos argumentos, es decir, toda tabla que se construye (ver \hyperref[err:190]{error~190}). Cuando la tabla tiene variables, ninguna numérica, y el papel es probabilidad condicionada, política, probabilidad conjunta o restricción de enlace, calcula el valor uniforme y hace \texttt{Arrays.\allowbreak{}fill(\allowbreak{}values,\allowbreak{}\ value)} en la línea 501. Como en esa rama también pone \texttt{set\allowbreak{}Value\ =\ true}, al salir del \texttt{if} vuelve a hacer \texttt{Arrays.\allowbreak{}fill(\allowbreak{}values,\allowbreak{}\ value)} en la línea 510 sobre el mismo array con el mismo valor. Para tablas de utilidad (papel sin especificar) no se rellena ninguna vez.

Medido con un pequeño programa de prueba (tabla de 4 194 304 casillas, 22 variables binarias): construir \texttt{new\ Table\allowbreak{}Potential(\allowbreak{}variables,\allowbreak{}\ papel,\allowbreak{}\ tabla)} cuesta unos 6 ms más con papel de probabilidad condicionada que con papel de utilidad, y ese exceso son los dos rellenos de este error y la reserva de \hyperref[err:190]{error~190}, juntos. Como referencia, una multiplicación de dos tablas que produce una tabla de ese tamaño tarda unos 23 ms en total, así que el trabajo tirado ronda la cuarta parte. Quitar solo el relleno duplicado ahorra una de las dos pasadas.
\paragraph{El código}
core/src/main/java/org/openmarkov/core/model/network/potential/TablePotential.java:471-510

\begin{Shaded}
\begin{Highlighting}[]
\ControlFlowTok{if} \OperatorTok{(}\NormalTok{numVariables }\OperatorTok{\textgreater{}} \DecValTok{0} \OperatorTok{\&\&} \FunctionTok{noNumericVariables}\OperatorTok{()} \OperatorTok{\&\&} \OperatorTok{(}
\NormalTok{        role }\OperatorTok{==}\NormalTok{ PotentialRole}\OperatorTok{.}\FunctionTok{CONDITIONAL\_PROBABILITY} \OperatorTok{||}\NormalTok{ role }\OperatorTok{==}\NormalTok{ PotentialRole}\OperatorTok{.}\FunctionTok{POLICY} \OperatorTok{||}
\NormalTok{                role }\OperatorTok{==}\NormalTok{ PotentialRole}\OperatorTok{.}\FunctionTok{JOINT\_PROBABILITY} \OperatorTok{||}\NormalTok{ role }\OperatorTok{==}\NormalTok{ PotentialRole}\OperatorTok{.}\FunctionTok{LINK\_RESTRICTION}\OperatorTok{))} \OperatorTok{\{}
\NormalTok{    setValue }\OperatorTok{=} \KeywordTok{true}\OperatorTok{;}
\NormalTok{    value }\OperatorTok{=} \FloatTok{0.0}\OperatorTok{;}
    \ControlFlowTok{switch} \OperatorTok{(}\NormalTok{role}\OperatorTok{)} \OperatorTok{\{}
        \KeywordTok{...}
    \OperatorTok{\}} \CommentTok{// When role = UTILITY {-}\textgreater{} value = 0.0 (default)}
    \BuiltInTok{Arrays}\OperatorTok{.}\FunctionTok{fill}\OperatorTok{(}\NormalTok{values}\OperatorTok{,}\NormalTok{ value}\OperatorTok{);}
\OperatorTok{\}} \ControlFlowTok{else} \ControlFlowTok{if} \OperatorTok{(}\NormalTok{numVariables }\OperatorTok{==} \DecValTok{0}\OperatorTok{)} \OperatorTok{\{}
\NormalTok{    setValue }\OperatorTok{=} \KeywordTok{true}\OperatorTok{;}
    \KeywordTok{...}
\OperatorTok{\}}
\ControlFlowTok{if} \OperatorTok{(}\NormalTok{setValue}\OperatorTok{)} \BuiltInTok{Arrays}\OperatorTok{.}\FunctionTok{fill}\OperatorTok{(}\NormalTok{values}\OperatorTok{,}\NormalTok{ value}\OperatorTok{);}
\end{Highlighting}
\end{Shaded}
\paragraph{Cómo arreglarlo}
Borrar el \texttt{Arrays.\allowbreak{}fill} de la línea 501 (el de la 510 cubre las dos ramas). No cambia ningún valor.

Medirlo con el banco de referencia de las operaciones, \texttt{Operations\allowbreak{}Reference\allowbreak{}Times\allowbreak{}Test} en \texttt{integration\allowbreak{}Tests}, antes y después. La batería de pruebas existente de \texttt{Table\allowbreak{}Potential} basta como red de seguridad; no hace falta prueba nueva de comportamiento.

\setcounter{subsection}{191}
\subsection[La maximización reserva la tabla de la unión de variables sin usarla]{La maximización reserva una tabla del tamaño de la unión de variables solo para calcular unos desplazamientos}\label{err:192}
\begin{ficha}
\fichakey{Estado}Presente
\fichakey{Módulo}core
\fichakey{Paquete}org.openmarkov.core.model.network.potential.operation
\fichakey{Ficheros}\path{core/src/main/java/org/openmarkov/core/model/network/potential/operation/TablePotentialMaximization.java:86-105} \newline \path{core/src/main/java/org/openmarkov/core/model/network/potential/operation/TablePotentialMaximization.java:206-225} \newline \path{core/src/main/java/org/openmarkov/core/model/network/potential/AbstractIndexedPotential.java:183-186} \newline \path{core/src/main/java/org/openmarkov/core/model/network/potential/AbstractIndexedPotential.java:307-309} \newline \path{core/src/main/java/org/openmarkov/core/model/network/potential/operation/TablePotentialElimination.java:113}
\fichakey{Comprobado}Leyendo el código
\fichakey{Esfuerzo}Pequeño (horas)
\fichakey{Decisión de equipo}No
\fichakey{Relacionados}\hyperref[err:190]{error~190}, \hyperref[nota:90]{nota~90}, \hyperref[nota:93]{nota~93}, \hyperref[nota:98]{nota~98}, \hyperref[nota:162]{nota~162}
\end{ficha}
\paragraph{Qué pasa}
Absorber un nodo de decisión desde el menú (\texttt{Absorb\allowbreak{}Node\allowbreak{}Edit} → \texttt{Node\allowbreak{}Absorption\allowbreak{}Handler}, línea 111) ejecuta \texttt{Table\allowbreak{}Potential\allowbreak{}Maximization.\allowbreak{}multiply\allowbreak{}And\allowbreak{}Maximize} (línea 98), que reserva y pone a cero un array del tamaño de la unión de variables sin usarlo. El tamaño de la unión es el de la utilidad del hijo, así que es un gesto de edición, no el camino caliente de la inferencia. No se ha medido.

El algoritmo de desplazamientos acumulados recorre varias tablas a la vez sin construir la tabla conjunta: solo necesita, para cada operando, cuánto avanza su posición cuando cambia cada variable de la unión. \texttt{multiply\allowbreak{}And\allowbreak{}Maximize} y \texttt{multiply\allowbreak{}And\allowbreak{}Maximize\allowbreak{}Uniformly} obtienen esos desplazamientos construyendo \texttt{new\ Table\allowbreak{}Potential(\allowbreak{}union\allowbreak{}Variables,\allowbreak{}\ null)} y llamando a su método de instancia \texttt{get\allowbreak{}Accumulated\allowbreak{}Offsets(\allowbreak{}potential.\allowbreak{}get\allowbreak{}Variables(\allowbreak{}))}, que se limita a llamar a la versión estática \texttt{Abstract\allowbreak{}Indexed\allowbreak{}Potential.\allowbreak{}get\allowbreak{}Accumulated\allowbreak{}Offsets(\allowbreak{}this.\allowbreak{}variables,\allowbreak{}\ otras)}.

Construir ese potencial reserva un array \texttt{double{[}{]}} del tamaño de la unión (la variable que se maximiza por las que se conservan), que es justo la tabla que el algoritmo existe para no construir. Como el papel es \texttt{null}, \texttt{set\allowbreak{}Uniform} no lo rellena: el coste es la reserva y puesta a cero de ese array. \texttt{Table\allowbreak{}Potential\allowbreak{}Elimination} (línea 113) ya usa la versión estática.

\texttt{multiply\allowbreak{}And\allowbreak{}Maximize\allowbreak{}Uniformly} (línea 218) no tiene llamadores de producción, solo pruebas.
\paragraph{El código}
core/src/main/java/org/openmarkov/core/model/network/potential/operation/TablePotentialMaximization.java:98-105

\begin{Shaded}
\begin{Highlighting}[]
\NormalTok{TablePotential unionPotential }\OperatorTok{=} \KeywordTok{new} \FunctionTok{TablePotential}\OperatorTok{(}\NormalTok{unionVariables}\OperatorTok{,} \KeywordTok{null}\OperatorTok{);}
\ControlFlowTok{for} \OperatorTok{(}\DataTypeTok{int}\NormalTok{ i }\OperatorTok{=} \DecValTok{0}\OperatorTok{;}\NormalTok{ i }\OperatorTok{\textless{}}\NormalTok{ numProperPotentials}\OperatorTok{;}\NormalTok{ i}\OperatorTok{++)} \OperatorTok{\{}
\NormalTok{    TablePotential potential }\OperatorTok{=}\NormalTok{ properPotentials}\OperatorTok{.}\FunctionTok{get}\OperatorTok{(}\NormalTok{i}\OperatorTok{);}
\NormalTok{    tables}\OperatorTok{[}\NormalTok{i}\OperatorTok{]} \OperatorTok{=}\NormalTok{ potential}\OperatorTok{.}\FunctionTok{getValues}\OperatorTok{();}
\NormalTok{    initialPositions}\OperatorTok{[}\NormalTok{i}\OperatorTok{]} \OperatorTok{=}\NormalTok{ potential}\OperatorTok{.}\FunctionTok{getInitialPosition}\OperatorTok{();}
\NormalTok{    currentPositions}\OperatorTok{[}\NormalTok{i}\OperatorTok{]} \OperatorTok{=}\NormalTok{ initialPositions}\OperatorTok{[}\NormalTok{i}\OperatorTok{];}
\NormalTok{    accumulatedOffsets}\OperatorTok{[}\NormalTok{i}\OperatorTok{]} \OperatorTok{=}\NormalTok{ unionPotential}\OperatorTok{.}\FunctionTok{getAccumulatedOffsets}\OperatorTok{(}\NormalTok{potential}\OperatorTok{.}\FunctionTok{getVariables}\OperatorTok{());}
\OperatorTok{\}}
\end{Highlighting}
\end{Shaded}
\paragraph{Cómo arreglarlo}
Sustituir en los dos sitios por \texttt{accumulated\allowbreak{}Offsets{[}i{]}\ =\ Table\allowbreak{}Potential.\allowbreak{}get\allowbreak{}Accumulated\allowbreak{}Offsets(\allowbreak{}union\allowbreak{}Variables,\allowbreak{}\ potential.\allowbreak{}get\allowbreak{}Variables(\allowbreak{}));} y borrar \texttt{union\allowbreak{}Potential}. No cambia valores; la prueba existente de maximización basta.

Si se borra \texttt{multiply\allowbreak{}And\allowbreak{}Maximize\allowbreak{}Uniformly} (el código sin llamadores se recoge en \hyperref[nota:162]{nota~162}, donde no hay nada que corregir), queda un solo sitio.


% ═════════════════════════════════════════════════════════════════════════════
\appendix
\renewcommand{\thesection}{\Alph{section}}

\section{Lo que se comprobó y no entra en la lista}\label{app:notas}

Estas notas recogen todo lo que se comprobó al preparar la lista y quedó fuera de ella. No hay nada
que corregir en ellas; están para que nadie tenga que volver a mirarlo y para que, si algún día
cambia algo (por ejemplo, un código que hoy no llama nadie empieza a usarse), se sepa qué había.


\subsection*{Ya corregidos}\addcontentsline{toc}{subsection}{Ya corregidos}
\begin{longtable}{@{}>{\raggedright\arraybackslash}p{1.1cm}>{\raggedright\arraybackslash}p{\dimexpr\linewidth-1.1cm-2\tabcolsep\relax}@{}}
\textbf{1}\phantomsection\label{nota:1} & \textbf{Reparto de una columna que se queda a cero al restringir un enlace}\newline Corregido en el commit \texttt{fd17916}: \texttt{Link\allowbreak{}Restriction\allowbreak{}Potential\allowbreak{}Operations.\allowbreak{}redistribute\allowbreak{}Probabilities} reparte ahora la columna a partes iguales entre los estados que siguen permitidos, sea cual sea su número. \\[3pt]
\textbf{2}\phantomsection\label{nota:2} & \textbf{Deshacer una absorción de nodo dejaba al hijo sin potenciales}\newline Corregido en el commit \texttt{ca779d4}: \texttt{Absorb\allowbreak{}Node\allowbreak{}Edit.\allowbreak{}do\allowbreak{}Edit} guarda ya los potenciales viejos y nuevos que usan \texttt{undo} y \texttt{redo}. \\[3pt]
\textbf{184}\phantomsection\label{nota:184} & \textbf{Pegar un nodo con un enlace a un nodo no copiado dejaba un nodo a medias}\newline Corregido en el commit \texttt{2c4fca6} (issue 663): \texttt{Visual\allowbreak{}Network.\allowbreak{}export\allowbreak{}To\allowbreak{}Clipboard} copia ya solo los enlaces cuyos dos extremos están entre los nodos seleccionados, y es el único sitio que llena el portapapeles de «Copy» y «Cut». Era el error~32 de la versión anterior de esta lista. \\[3pt]
\end{longtable}

\subsection*{No son errores}\addcontentsline{toc}{subsection}{No son errores}
\begin{longtable}{@{}>{\raggedright\arraybackslash}p{1.1cm}>{\raggedright\arraybackslash}p{\dimexpr\linewidth-1.1cm-2\tabcolsep\relax}@{}}
\textbf{3}\phantomsection\label{nota:3} & \textbf{La poda de nodos estériles declara estéril a un padre con un solo hijo sin comprobarlo}\newline En \texttt{Prob\allowbreak{}Net\allowbreak{}Operations.\allowbreak{}remove\allowbreak{}Barren\allowbreak{}Nodes}, el único hijo de ese padre es, por construcción, el nodo estéril desde el que se llega a él, así que la comprobación que se omite daría siempre «sí». El atajo está dentro del método y ningún otro llamador puede romper esa condición. \\[3pt]
\textbf{4}\phantomsection\label{nota:4} & \textbf{Cambiar una decisión con política por un nodo de azar vacía potenciales que no existen}\newline \texttt{Task\allowbreak{}Utilities.\allowbreak{}replace\allowbreak{}Decisions\allowbreak{}With\allowbreak{}Policies\allowbreak{}By\allowbreak{}Chance\allowbreak{}Nodes} recorre para quitarlos los potenciales de un nodo recién creado, que no tiene ninguno, así que el bucle no hace nada. El nodo queda con la política como único potencial, que es lo que se quiere; solo sobra una línea engañosa. \\[3pt]
\textbf{5}\phantomsection\label{nota:5} & \textbf{Los mapas locales de \texttt{reset\allowbreak{}Link} guardan restricciones y revelaciones que nadie lee}\newline \texttt{Variable\allowbreak{}State\allowbreak{}Operations.\allowbreak{}reset\allowbreak{}Link} copia las restricciones de enlace y las condiciones de revelación en dos mapas locales que mueren al salir del método: es código muerto sin consecuencia propia, porque \texttt{Node\allowbreak{}State\allowbreak{}Edit} ya guarda y repone esos datos (commit \texttt{f7081df}). La pérdida al deshacer un dominio estándar es un defecto de \texttt{Node\allowbreak{}Replace\allowbreak{}States\allowbreak{}Edit}, descrito en \hyperref[err:24]{error~24}, y no se arregla rescatando estos mapas. \\[3pt]
\textbf{6}\phantomsection\label{nota:6} & \textbf{El tipo de análisis guardado en el fichero se lee al abrirlo}\newline \texttt{PGMXReader\_\allowbreak{}0\_\allowbreak{}2.\allowbreak{}get\allowbreak{}Multicriteria\allowbreak{}Options} lee \texttt{\textless{}Selected\allowbreak{}Analysis\allowbreak{}Type\textgreater{}} y fija el tipo en las opciones de la red, y \texttt{PGMXReader\_\allowbreak{}1\_\allowbreak{}0} hereda ese método. Medido: \texttt{integration\allowbreak{}Tests/\allowbreak{}src/\allowbreak{}main/\allowbreak{}resources/\allowbreak{}networks/\allowbreak{}IDCEAn\allowbreak{}Therapies/\allowbreak{}DAN-1tests.\allowbreak{}pgmx}, \texttt{0\_\allowbreak{}miscellaneous/\allowbreak{}networks/\allowbreak{}dan/\allowbreak{}DAN-CE-2-test-problem.\allowbreak{}pgmx} y \texttt{0\_\allowbreak{}miscellaneous/\allowbreak{}networks/\allowbreak{}id/\allowbreak{}ID-CEA-minimal.\allowbreak{}pgmx} se abren como coste-efectividad. \\[3pt]
\textbf{7}\phantomsection\label{nota:7} & \textbf{La lognormal se muestrea con una inversa aproximada de la normal y se acota con la exacta}\newline \texttt{Log\allowbreak{}Normal\allowbreak{}Function.\allowbreak{}get\allowbreak{}Sample} usa la aproximación de Odeh y Evans y \texttt{get\allowbreak{}Interval} la inversa exacta de Apache Commons Math, pero las dos difieren como mucho en unas cienmillonésimas para las probabilidades que da el generador en la práctica. Medido con un millón de extracciones: la fracción de muestras dentro del intervalo del 95 \% es 0,95007. \\[3pt]
\textbf{8}\phantomsection\label{nota:8} & \textbf{La tabla de cuentas de la información mutua condicionada no sigue su orden declarado}\newline \texttt{Conditional\allowbreak{}Mutual\allowbreak{}Information\allowbreak{}Metric.\allowbreak{}build\allowbreak{}Absolute\allowbreak{}Freq\allowbreak{}Cross\allowbreak{}Tab} rellena la tabla con un orden de variables distinto del declarado, pero el único método que la lee, \texttt{score}, la recorre con ese mismo orden. El número resultante es la información mutua condicionada correcta; medido contra un cálculo independiente sobre los casos. \\[3pt]
\textbf{9}\phantomsection\label{nota:9} & \textbf{El algoritmo PC consulta dos veces la caché con la misma llave al orientar colisionadores}\newline \texttt{Collider\allowbreak{}Orientation.\allowbreak{}creates\allowbreak{}Contradictory\allowbreak{}Collider} busca \texttt{new\ Node\allowbreak{}Pair(\allowbreak{}from,\allowbreak{}\ existing\allowbreak{}Parent)} y luego \texttt{new\ Node\allowbreak{}Pair(\allowbreak{}existing\allowbreak{}Parent,\allowbreak{}\ from)}, pero \texttt{Node\allowbreak{}Pair} ordena sus dos nodos por nombre y las dos llaves son iguales. La segunda consulta sobra, pero el resultado del algoritmo PC de aprendizaje de estructura es correcto. \\[3pt]
\textbf{10}\phantomsection\label{nota:10} & \textbf{Pedir la red modelo sin marcar posiciones ni estructura no deja el aprendizaje sin red}\newline \texttt{Learning\allowbreak{}Manager.\allowbreak{}apply\allowbreak{}Model\allowbreak{}Net} devuelve \texttt{null} en ese caso, pero solo se llama cuando \texttt{Model\allowbreak{}Net\allowbreak{}Use.\allowbreak{}is\allowbreak{}Use\allowbreak{}Model\allowbreak{}Net(\allowbreak{})} es verdadero, y el constructor de \texttt{Model\allowbreak{}Net\allowbreak{}Use} lo pone a falso si no se marca ninguna de las dos opciones. \texttt{Learning\allowbreak{}Manager} crea entonces una red con un nodo por variable de la base de casos. \\[3pt]
\textbf{11}\phantomsection\label{nota:11} & \textbf{La letra negrita de 30 puntos de la tabla de probabilidades nunca llega a la pantalla}\newline \texttt{Values\allowbreak{}Table\allowbreak{}Cell\allowbreak{}Renderer.\allowbreak{}set\allowbreak{}Cell\allowbreak{}Fonts} fija esa letra, pero justo después \texttt{Default\allowbreak{}Table\allowbreak{}Cell\allowbreak{}Renderer.\allowbreak{}get\allowbreak{}Table\allowbreak{}Cell\allowbreak{}Renderer\allowbreak{}Component} la sustituye por la de la tabla. Medido: todas las casillas salen con la letra de la tabla y sin texto recortado; \texttt{set\allowbreak{}Cell\allowbreak{}Fonts} es código sin efecto. \\[3pt]
\textbf{12}\phantomsection\label{nota:12} & \textbf{Una columna con incertidumbre no se pinta de rojo en un nodo con restricciones de enlace}\newline \texttt{Values\allowbreak{}Table\allowbreak{}With\allowbreak{}Link\allowbreak{}Restriction\allowbreak{}Cell\allowbreak{}Renderer.\allowbreak{}get\allowbreak{}Cell\allowbreak{}Colors} pone fondo rojo a toda casilla no editable, pero para una columna con incertidumbre \texttt{Values\allowbreak{}Table\allowbreak{}Cell\allowbreak{}Renderer.\allowbreak{}get\allowbreak{}Table\allowbreak{}Cell\allowbreak{}Renderer\allowbreak{}Component} devuelve la etiqueta del icono de incertidumbre, con su propio fondo. Medido: solo sale roja la casilla que el enlace prohíbe. \\[3pt]
\textbf{13}\phantomsection\label{nota:13} & \textbf{Un número mal escrito en la tabla de probabilidades se avisa con una ventana}\newline La excepción que lanza \texttt{Values\allowbreak{}Table.\allowbreak{}set\allowbreak{}Value\allowbreak{}At} llega a \texttt{OMException\allowbreak{}Handler}, que abre la ventana «Mismatched value» con lo tecleado; medido con «0,3», «abc» y «-0.3». Que no se admita la coma decimal queda fuera por la decisión de que la aplicación sea solo en inglés, y \texttt{Open\allowbreak{}Markov.\allowbreak{}main} fija el formato de números inglés. \\[3pt]
\textbf{14}\phantomsection\label{nota:14} & \textbf{Los textos con huecos usan dos convenciones, tilde y llaves numeradas}\newline Las claves con huecos \texttt{\textasciitilde{}} se rellenan con \texttt{String\allowbreak{}Database.\allowbreak{}get\allowbreak{}Formatted\allowbreak{}String} y las de huecos \texttt{\{0\}} con \texttt{Message\allowbreak{}Format}; comprobadas todas las claves inglesas con huecos, cada una va por su mecanismo y ningún texto sale sin rellenar. Tener dos convenciones es una incomodidad de mantenimiento, no un fallo. \\[3pt]
\textbf{15}\phantomsection\label{nota:15} & \textbf{Propagar un dominio a otros ciclos crea varias ediciones, pero se deshacen de una vez}\newline \texttt{Node\allowbreak{}Properties\allowbreak{}Dialog} abre una sub-historia de ediciones al construirse y, al aceptar, la cierra como una sola entrada \texttt{List\allowbreak{}PNEdit}, así que un único «deshacer» devuelve todos los ciclos; comprobado ejecutando con \texttt{0\_\allowbreak{}miscellaneous/\allowbreak{}networks/\allowbreak{}mid/\allowbreak{}MID-mammography.\allowbreak{}pgmx}. Lo que sí se pierde al deshacer es el contenido de las tablas, que es \hyperref[err:24]{error~24}. \\[3pt]
\textbf{16}\phantomsection\label{nota:16} & \textbf{La suma de utilidades conserva los planes de decisión}\newline La suma de potenciales conserva los árboles de estrategia, así que no hay ningún defecto de OpenMarkov que corregir. \\[3pt]
\textbf{17}\phantomsection\label{nota:17} & \textbf{El editor de árboles rechaza las operaciones imposibles con un mensaje claro}\newline \texttt{Tree\allowbreak{}ADDEditor\allowbreak{}Panel.\allowbreak{}action\allowbreak{}Performed} envuelve esas excepciones en \texttt{Unrecoverable\allowbreak{}Exception}, y \texttt{OMException\allowbreak{}Handler} enseña un diálogo con su mensaje sin cerrar la aplicación; las excepciones se lanzan antes de modificar el árbol. Medido: quitar todos los estados de una rama, partir un intervalo por un límite imposible y salir del dominio dan los tres un mensaje claro en inglés, aunque con un título mal formado (\hyperref[err:182]{error~182}). Un umbral que no es un número, en cambio, sí sale como error no previsto (\hyperref[err:161]{error~161}). \\[3pt]
\textbf{18}\phantomsection\label{nota:18} & \textbf{La entrada «Settings\ldots» del menú Edición sí tiene texto}\newline La clave \texttt{Edit.\allowbreak{}Open\allowbreak{}Settings} está en \texttt{gui/\allowbreak{}src/\allowbreak{}main/\allowbreak{}resources/\allowbreak{}gui/\allowbreak{}localize/\allowbreak{}Tool\allowbreak{}Bars\_\allowbreak{}en.\allowbreak{}xml} y \texttt{String\allowbreak{}Database} junta todos los ficheros de textos, así que la opción del menú se lee «Settings\ldots». Solo le falta la clave del mnemónico (la letra subrayada que permite elegirla con el teclado), igual que a otras entradas del mismo menú: no tiene atajo de letra, pero no hay texto ilegible. \\[3pt]
\textbf{19}\phantomsection\label{nota:19} & \textbf{La herramienta de poner sucesos no enseña ninguna clave sin texto}\newline La ayuda del botón de sucesos pide \texttt{Edit.\allowbreak{}Mode.\allowbreak{}Event.\allowbreak{}Tool\allowbreak{}Tip}, que existe en \texttt{Tool\allowbreak{}Bars\_\allowbreak{}en.\allowbreak{}xml}, y la barra de edición y el menú contextual del lienzo escriben «Event» en el código. La clave \texttt{Edit.\allowbreak{}Mode.\allowbreak{}Event}, que no tiene texto, no se pide para enseñarla; la tabla \texttt{Node\allowbreak{}Mode\allowbreak{}Descriptor} que la usaba se retiró en el commit \texttt{8055530}. \\[3pt]
\textbf{20}\phantomsection\label{nota:20} & \textbf{La coherencia «débil» pedida por la validación cruzada no llega a los resultados}\newline \texttt{Cross\allowbreak{}Validation\allowbreak{}Dialog} pide la plantilla de medidas con \texttt{Coherence.\allowbreak{}WEAK}, pero \texttt{Learning\allowbreak{}Evaluator.\allowbreak{}run\allowbreak{}Evaluator} construye el resultado con el constructor de copia de \texttt{Measures\allowbreak{}Set}, que marca todas las variables como usadas. Medido: los resultados no muestran la frase sobre variables sin evidencia, que en la validación cruzada sería falsa; el \texttt{Coherence.\allowbreak{}WEAK} escrito a mano queda como código engañoso sin efecto. \\[3pt]
\textbf{21}\phantomsection\label{nota:21} & \textbf{El paseo guiado «First tour» señala el nodo correcto en el segundo cambio de nombre}\newline Los pasos eligen el nodo con \texttt{get\allowbreak{}Visual\allowbreak{}Network(\allowbreak{}).\allowbreak{}get\allowbreak{}All\allowbreak{}Nodes(\allowbreak{}).\allowbreak{}get(\allowbreak{}1)}, y \texttt{Network\allowbreak{}Editor\allowbreak{}Panel.\allowbreak{}do\allowbreak{}Paint} reordena esa lista en cada repintado poniendo primero los nodos seleccionados: tras renombrar B, que queda seleccionado, la lista es {[}C, A{]} y \texttt{get(\allowbreak{}1)} es A. Medido el mecanismo con un pequeño programa de prueba; elegir el nodo por una posición que depende de la selección es frágil, pero hoy acierta. \\[3pt]
\textbf{22}\phantomsection\label{nota:22} & \textbf{Rehacer un nodo añadido conserva lo que puso \texttt{add\allowbreak{}Node\allowbreak{}Consistently}}\newline \texttt{Add\allowbreak{}Node\allowbreak{}Edit.\allowbreak{}redo} solo vuelve a insertar el nodo, pero el potencial por omisión, la política y la posición viven en el propio objeto \texttt{Node}, que \texttt{undo} saca del grafo sin tocarlo. Medido en un diagrama de influencia: deshacer y rehacer devuelve los nodos idénticos. \\[3pt]
\textbf{23}\phantomsection\label{nota:23} & \textbf{Deshacer un enlace añadido no tiene restricciones de enlace que reponer}\newline Un enlace recién añadido no tiene restricciones ni estados reveladores, y \texttt{Distinct\allowbreak{}Links} y \texttt{No\allowbreak{}Multiple\allowbreak{}Links} impiden añadirlo encima de otro, así que \texttt{Add\allowbreak{}Link\allowbreak{}Edit.\allowbreak{}undo} no tiene nada que reponer. Las restricciones que se pongan después las introducen ediciones propias, con su propio deshacer. \\[3pt]
\textbf{24}\phantomsection\label{nota:24} & \textbf{Enlazar varios nodos a la vez omite los pares ya enlazados y los lazos prohibidos}\newline \texttt{Multi\allowbreak{}Add\allowbreak{}Link\allowbreak{}Edit} descarta los pares que ya tienen ese enlace o que serían un lazo prohibido, que es justo el resultado que el usuario pide; durante el arrastre, \texttt{Visual\allowbreak{}Network.\allowbreak{}update\allowbreak{}Link\allowbreak{}Creation} pinta esas flechas en gris o transparentes. Avisarlo de forma más visible sería un cambio de interfaz, no la corrección de un fallo. \\[3pt]
\textbf{25}\phantomsection\label{nota:25} & \textbf{Comprobar restricciones en cada movimiento del ratón al trazar enlaces cuesta muy poco}\newline \texttt{Visual\allowbreak{}Network.\allowbreak{}update\allowbreak{}Link\allowbreak{}Creation} comprueba las restricciones de cada enlace candidato en cada movimiento del ratón. Medido sobre \texttt{integration\allowbreak{}Tests/\allowbreak{}src/\allowbreak{}main/\allowbreak{}resources/\allowbreak{}networks/\allowbreak{}bn/\allowbreak{}BN-cpcs422b.\allowbreak{}pgmx} (422 nodos, 867 enlaces): 0,14 ms de media con diez enlaces candidatos y 2,6 ms en el peor caso, sin contar el repintado, muy por debajo de lo que se nota. \\[3pt]
\textbf{26}\phantomsection\label{nota:26} & \textbf{El recolector de violaciones de restricciones no quita repetidos, pero el mensaje final sí}\newline \texttt{Constraint\allowbreak{}Checker} guarda las excepciones en un \texttt{Hash\allowbreak{}Set} que no las funde, pero \texttt{Multiple\allowbreak{}Constraints\allowbreak{}Violateds} las agrupa por su texto, así que un mensaje idéntico sale una sola vez (medido con \texttt{No\allowbreak{}Multiple\allowbreak{}Links}). Que \texttt{No\allowbreak{}Cycle} nombre cada enlace de un ciclo es información, no un repetido, y hoy solo lo ve quien abre desde un fichero una red con ciclo (\hyperref[err:37]{error~37}). \\[3pt]
\textbf{27}\phantomsection\label{nota:27} & \textbf{Varias operaciones copian la lista de variables de cada potencial dentro de sus bucles}\newline \texttt{Potential.\allowbreak{}get\allowbreak{}Variables(\allowbreak{})} devuelve una copia en cada llamada y varios bucles la piden una vez por potencial, no por casilla, así que el coste es pequeño y acotado. Es una observación de rendimiento, no un fallo; si algún día se busca velocidad, dar acceso de solo lectura a la lista interna sería una vía barata. \\[3pt]
\textbf{28}\phantomsection\label{nota:28} & \textbf{Tres modelos con «Age {[}0{]}» uniforme no se analizan sin darle un valor: lo exige la red}\newline En \texttt{0\_\allowbreak{}miscellaneous/\allowbreak{}networks/\allowbreak{}mid/\allowbreak{}MID-CHD-Walker.\allowbreak{}pgmx}, \texttt{MID-Colorectal.\allowbreak{}pgmx} y \texttt{MID-mammography.\allowbreak{}pgmx}, «Age {[}0{]}» es un nodo numérico sin distribución que el modelo espera recibir como hallazgo, y sin él no se puede discretizar: \texttt{MID-Colorectal.\allowbreak{}pgmx} guarda ese hallazgo en el fichero y su coste-efectividad global termina (medido: un tramo, coste 1837 y efectividad 23,11). \texttt{MID-mammography.\allowbreak{}pgmx} no lo guarda, y con «Age {[}0{]} = 50» falla más adelante por \hyperref[err:13]{error~13}; \texttt{MID-CHD-Walker.\allowbreak{}pgmx} es un esqueleto cuyos siete nodos de utilidad son uniformes, sin números que calcular. --- cambios --- \\[3pt]
\end{longtable}

\subsection*{Defectos sin consecuencia visible hoy}\addcontentsline{toc}{subsection}{Defectos sin consecuencia visible hoy}
\begin{longtable}{@{}>{\raggedright\arraybackslash}p{1.1cm}>{\raggedright\arraybackslash}p{\dimexpr\linewidth-1.1cm-2\tabcolsep\relax}@{}}
\textbf{29}\phantomsection\label{nota:29} & \textbf{El javadoc de la comparación de motivos del algoritmo PC dice lo contrario que el código}\newline \texttt{PCEdit\allowbreak{}Motivation.\allowbreak{}compare\allowbreak{}To} hace ganar al conjunto de separación más pequeño, que es el orden por tamaño creciente del algoritmo PC de aprendizaje de estructura publicado, mientras que su javadoc dice que gana el más grande. Solo el comentario está mal; el código hace lo correcto. \\[3pt]
\textbf{30}\phantomsection\label{nota:30} & \textbf{Tres búsquedas del editor de árboles comparan nombres con «==» en vez de \texttt{equals}}\newline \texttt{Tree\allowbreak{}ADDEditor\allowbreak{}Panel.\allowbreak{}add\allowbreak{}Subtree}, \texttt{associate\allowbreak{}States} y \texttt{remove\allowbreak{}Variables\allowbreak{}From\allowbreak{}Potential} comparan con \texttt{==} el texto elegido por el usuario con el nombre de la variable o del estado, lo que compara la identidad de los objetos y no su contenido. Hoy funciona porque, medido, ese texto es el mismo objeto \texttt{String} que el nombre; es un defecto latente que aparecería si el texto llegara como otro objeto con el mismo contenido. \\[3pt]
\textbf{31}\phantomsection\label{nota:31} & \textbf{El orden de las entradas empatadas del menú Tools depende del descubrimiento de clases}\newline \texttt{Main\allowbreak{}Menu.\allowbreak{}build\allowbreak{}Tools\allowbreak{}Menu} ordena los complementos solo por prioridad, y las seis entradas empatadas quedan en el orden en que \texttt{Plugin\allowbreak{}Loader} las descubre. En el jar ejecutable ese orden es fijo; solo al ejecutar desde directorios de clases depende de \texttt{File.\allowbreak{}list\allowbreak{}Files}, cuyo orden no está definido, y no se ha observado ningún cambio real. \\[3pt]
\textbf{32}\phantomsection\label{nota:32} & \textbf{La red auxiliar que enseña la política óptima de una decisión duplica cada enlace}\newline \texttt{Network\allowbreak{}Editor\allowbreak{}Panel.\allowbreak{}show\allowbreak{}Optimal\allowbreak{}Policy\allowbreak{}Of\allowbreak{}Node} añade el potencial de la política con \texttt{Prob\allowbreak{}Net.\allowbreak{}add\allowbreak{}Potential}, que ya traza los enlaces, y después los vuelve a trazar con \texttt{add\allowbreak{}Link}, que no comprueba duplicados; medido con \texttt{0\_\allowbreak{}miscellaneous/\allowbreak{}networks/\allowbreak{}dan/\allowbreak{}DAN-dating.\allowbreak{}pgmx} e \texttt{integration\allowbreak{}Tests/\allowbreak{}src/\allowbreak{}main/\allowbreak{}resources/\allowbreak{}networks/\allowbreak{}id/\allowbreak{}ID-decide-test-sv.\allowbreak{}pgmx}. Esa red no se dibuja y el diálogo solo enseña el potencial, así que no hay consecuencia visible, aunque cualquier código que recorriera sus enlaces los contaría de más. \\[3pt]
\textbf{33}\phantomsection\label{nota:33} & \textbf{El generador de bases de casos propone 5000 casos y el comentario del código dice 1000}\newline \texttt{DBGenerator\allowbreak{}GUI} hace \texttt{case\allowbreak{}Number.\allowbreak{}set\allowbreak{}Selected\allowbreak{}Index(\allowbreak{}2);\ /\allowbreak{}/\allowbreak{}\ 1000\ cases}, y el índice 2 de la lista es «5000». El número está a la vista en la ventana, así que el usuario no queda engañado; lo que hay es una contradicción entre comentario y código, y decidir cuál de los dos vale. \\[3pt]
\textbf{185}\phantomsection\label{nota:185} & \textbf{Multiplicar deja fuera las variables de un solo estado de los potenciales constantes}\newline \texttt{Table\allowbreak{}Potential\allowbreak{}Arithmetic.\allowbreak{}multiply} trata como constante cualquier potencial de una sola casilla y deja fuera del producto sus variables, que tienen un estado cada una; por lo mismo, el árbol de estrategia de un potencial de una casilla con variables no llega al resultado. La suma, la marginalización, la maximización y la evaluación de fórmulas solo tratan como constante el potencial sin variables, y conservan esas variables. No cambia ningún número, porque una variable de un estado no añade ninguna dimensión a la tabla. Comprobado el 23 de septiembre en las 165 redes distintas del repositorio que tienen nodos numéricos de azar o son de tipo DAN (propagación, evaluación, políticas, árbol de estrategia, coste-efectividad y evaluación de DAN, con evidencia en los nodos numéricos y sin ella): hoy ningún resultado lleva una variable de un estado. Conservarlas rompe dos cosas que hoy funcionan: «Absorb parents» sobre un producto con dos padres de regresión sin padres propios (\hyperref[err:193]{error~193}) y la utilidad esperada de \texttt{DAN-dating.\allowbreak{}pgmx}, que sale con el mismo valor pero con la variable \texttt{U\ mExp} (\hyperref[nota:186]{nota~186}). Era lo que quedaba del \hyperref[err:1]{error~1}. \\[3pt]
\textbf{186}\phantomsection\label{nota:186} & \textbf{La evaluación de una DAN deja una utilidad a cero sin criterio, que cuenta como probabilidad}\newline Cuando una restricción deja un nodo sin estados posibles y la utilidad que habría que quitar es la última de la red, \texttt{DAN\allowbreak{}Operations.\allowbreak{}replace\allowbreak{}Utility\allowbreak{}Potential\allowbreak{}With\allowbreak{}Zero} (líneas 231-238, llamado en las líneas 136 y 198) la deja con una tabla a cero sobre su propia variable y sin criterio. Sin criterio, la eliminación de variables la cuenta como una probabilidad: \texttt{Variable\allowbreak{}Elimination\allowbreak{}Core.\allowbreak{}get\allowbreak{}Probability} (línea 322) la multiplica con las demás. Pasa al evaluar \texttt{DAN-dating.\allowbreak{}pgmx}, y el resultado es el correcto: \texttt{DAN\allowbreak{}Evaluation\allowbreak{}Test.\allowbreak{}test\allowbreak{}DAN\allowbreak{}Dating\allowbreak{}Ask\allowbreak{}No} fija a mano 9,407584 y la evaluación lo da. No se ha comprobado si ese cero cambia el resultado de otra red. Hoy el producto deja fuera la variable de esa utilidad (\hyperref[nota:185]{nota~185}); si la conservara, la utilidad esperada de \texttt{DAN-dating} saldría con la variable \texttt{U\ mExp}. \\[3pt]
\end{longtable}

\subsection*{Defectos que hoy no alcanza nadie}\addcontentsline{toc}{subsection}{Defectos que hoy no alcanza nadie}
\begin{longtable}{@{}>{\raggedright\arraybackslash}p{1.1cm}>{\raggedright\arraybackslash}p{\dimexpr\linewidth-1.1cm-2\tabcolsep\relax}@{}}
\textbf{34}\phantomsection\label{nota:34} & \textbf{Tras guardar, el indicador de red modificada de \texttt{PNESupport} sigue diciendo que sí}\newline \texttt{PNESupport.\allowbreak{}on\allowbreak{}Save} no recalcula \texttt{network\allowbreak{}Is\allowbreak{}Modified}, que sigue valiendo \texttt{true} hasta la edición siguiente. Su único lector, \texttt{Network\allowbreak{}Editor\allowbreak{}Panel}, lo consulta solo después de ejecutar, deshacer o rehacer una edición, cuando ya se ha recalculado; y el asterisco de la pestaña, el botón de guardar y la pregunta al cerrar dependen del indicador propio del panel, que se pone a falso al guardar. \\[3pt]
\textbf{35}\phantomsection\label{nota:35} & \textbf{\texttt{Remove\allowbreak{}Potentials\allowbreak{}Edit} no se puede deshacer}\newline La edición hereda el \texttt{undo(\allowbreak{})} vacío de \texttt{PNEdit}, así que quitar potenciales con ella no se podría deshacer (el mismo problema en \texttt{Potential\allowbreak{}Change\allowbreak{}Edit} ya se corrigió en el commit \texttt{b81524f}). Ningún menú, diálogo ni algoritmo la construye: solo aparece en su propio fichero, en el fichero de textos y en el catálogo de pruebas de deshacer, que la aparta por no tener llamadores. \\[3pt]
\textbf{36}\phantomsection\label{nota:36} & \textbf{Subir el primer criterio de decisión o bajar el último se sale de la lista}\newline \texttt{Decision\allowbreak{}Criteria\allowbreak{}Edit.\allowbreak{}do\allowbreak{}Edit} no comprueba los extremos en las acciones \texttt{UP} y \texttt{DOWN}, y llamada así lanza \texttt{Index\allowbreak{}Out\allowbreak{}Of\allowbreak{}Bounds\allowbreak{}Exception}. Desde la ventana no se llega: \texttt{Decision\allowbreak{}Criteria\allowbreak{}Table\allowbreak{}Panel} es el único llamador, \texttt{Key\allowbreak{}Table\allowbreak{}Panel.\allowbreak{}value\allowbreak{}Changed} apaga «subir» en la primera fila y «bajar» en la última, y tras cada movimiento se pierde la selección. \\[3pt]
\textbf{37}\phantomsection\label{nota:37} & \textbf{\texttt{Temporal\allowbreak{}Evaluation} no arranca desde el jar completo de la aplicación}\newline \texttt{Temporal\allowbreak{}Evaluation.\allowbreak{}resolve} usa \texttt{Log\allowbreak{}Manager.\allowbreak{}get\allowbreak{}Logger(\allowbreak{})} sin argumentos, que desde \texttt{Open\allowbreak{}Markov.\allowbreak{}jar} lanza \texttt{Unsupported\allowbreak{}Operation\allowbreak{}Exception} porque el manifiesto del jar completo no declara \texttt{Multi-Release:\ true}; con el jar normal de log4j funciona. Hoy nadie llega: \texttt{Temporal\allowbreak{}Evaluation} no tiene llamadores de producción y solo la usan pruebas, que corren con el jar normal. \\[3pt]
\textbf{38}\phantomsection\label{nota:38} & \textbf{La copia superficial de una red recupera las restricciones de su tipo que se quitaron}\newline \texttt{Prob\allowbreak{}Net\allowbreak{}Copier.\allowbreak{}shallow\allowbreak{}Copy} crea la copia con las restricciones del tipo de red y le añade las del original, así que una restricción quitada vuelve (medido con \texttt{Distinct\allowbreak{}Links}: la copia no pasa \texttt{check\allowbreak{}Constraints(\allowbreak{})} y el original sí). En producción solo \texttt{Super\allowbreak{}Parent\allowbreak{}NBAlgorithm} quita una restricción del tipo, y no se ha encontrado nada que lea después las restricciones de una copia de esa red. \\[3pt]
\textbf{39}\phantomsection\label{nota:39} & \textbf{Una propagación de Hugin reutilizada con otra evidencia mezcla la evidencia vieja}\newline \texttt{Cluster\allowbreak{}Propagation.\allowbreak{}introduce\allowbreak{}Evidence} añade los potenciales de la evidencia nueva sin vaciar los anteriores, y \texttt{collect\allowbreak{}Evidence} reutiliza el mensaje guardado de la vez anterior; medido sobre \texttt{0\_\allowbreak{}miscellaneous/\allowbreak{}networks/\allowbreak{}bn/\allowbreak{}BN-asia.\allowbreak{}pgmx}, la segunda pregunta no da la respuesta correcta, {[}0,99, 0,01{]}. Los dos llamadores de producción construyen un objeto nuevo en cada consulta o fijan toda la evidencia antes de la primera pregunta, así que ningún usuario ve un número equivocado. \\[3pt]
\textbf{40}\phantomsection\label{nota:40} & \textbf{Convertir otra camarilla en raíz del árbol de Hugin lanzaría siempre una excepción}\newline \texttt{Hugin\allowbreak{}Forest.\allowbreak{}convert\allowbreak{}To\allowbreak{}Root} espera \texttt{null} de \texttt{Graph.\allowbreak{}get\allowbreak{}Parents} y de \texttt{Stack.\allowbreak{}pop}, que en realidad devuelven una lista vacía y lanzan excepción, así que acabaría siempre en \texttt{Index\allowbreak{}Out\allowbreak{}Of\allowbreak{}Bounds\allowbreak{}Exception} o \texttt{Empty\allowbreak{}Stack\allowbreak{}Exception}; además empieza el recorrido por la raíz vieja y no por \texttt{new\allowbreak{}Root}. El método no tiene ningún llamador en el proyecto, ni en producción ni en pruebas. \\[3pt]
\textbf{41}\phantomsection\label{nota:41} & \textbf{\texttt{Strategy\allowbreak{}Manager} lee la política de una decisión con las elecciones de otra}\newline \texttt{Strategy\allowbreak{}Manager.\allowbreak{}get\allowbreak{}Potential\allowbreak{}Form} calcula mal dónde empiezan las elecciones de cada decisión cuando las decisiones no tienen el mismo número de combinaciones de padres, y \texttt{evaluate} busca la política por número de ciclo y no por decisión. La clase no tiene ningún llamador en el proyecto; su paquete solo aparece exportado en \texttt{module-info.\allowbreak{}java}. \\[3pt]
\textbf{42}\phantomsection\label{nota:42} & \textbf{La búsqueda de estrategias por fuerza bruta divide por cero con límites pequeños}\newline \texttt{Strategy\allowbreak{}Manager.\allowbreak{}brute\allowbreak{}Force} lanza \texttt{Arithmetic\allowbreak{}Exception} con cualquier límite menor que 20, evalúa dos estrategias más de las que promete y lanza \texttt{Null\allowbreak{}Pointer\allowbreak{}Exception} si ninguna estrategia mejora a la base; \texttt{random\allowbreak{}Walk} tiene la misma división. Nadie llama a estos métodos. \\[3pt]
\textbf{43}\phantomsection\label{nota:43} & \textbf{La búsqueda de estrategias recorre elecciones del último ciclo que la evaluación no usa}\newline \texttt{Strategy\allowbreak{}Manager.\allowbreak{}order\allowbreak{}Nodes\allowbreak{}By\allowbreak{}Slice} incluye las decisiones del ciclo \texttt{horizon}, pero \texttt{evaluate} solo recorre hasta \texttt{horizon\ -\ 1}, así que esas elecciones multiplican el número de estrategias sin cambiar ninguna utilidad. La clase no tiene ningún llamador en el proyecto. \\[3pt]
\textbf{44}\phantomsection\label{nota:44} & \textbf{El texto de una partición de coste-efectividad usa la puntuación numérica del sistema}\newline \texttt{CEP} crea sus \texttt{Decimal\allowbreak{}Format} con la puntuación por omisión del proceso, así que con el sistema en castellano escribe \texttt{Eff:\ 8,\allowbreak{}5}. La aplicación de escritorio fija el formato inglés en \texttt{Open\allowbreak{}Markov.\allowbreak{}main} antes de crear ninguna partición, de modo que solo lo vería quien use el núcleo como biblioteca sin fijarlo. \\[3pt]
\textbf{45}\phantomsection\label{nota:45} & \textbf{La «lambda» de una partición de coste-efectividad cierra todos sus tramos por la derecha}\newline \texttt{CEP.\allowbreak{}get\allowbreak{}Lambda} crea la lista de lados con un elemento de más, así que todos los tramos quedan abiertos por la izquierda y cerrados por la derecha, al revés de lo que dicen sus nombres de estado «λ in {[}a , b)». Solo la llama \texttt{CEP.\allowbreak{}get\allowbreak{}Intervention(\allowbreak{})} sin argumentos, que no tiene ningún llamador. \\[3pt]
\textbf{46}\phantomsection\label{nota:46} & \textbf{Copiar la partición de coste-efectividad de probabilidad cero lanza una excepción}\newline \texttt{CEP.\allowbreak{}clone} recorre cuatro vectores que en la partición de probabilidad cero son \texttt{null}, así que \texttt{CEP.\allowbreak{}get\allowbreak{}Zero\allowbreak{}Partition(\allowbreak{}).\allowbreak{}clone(\allowbreak{})} lanza \texttt{Null\allowbreak{}Pointer\allowbreak{}Exception} (medido). La única copia de una partición es la de \texttt{CEBase\allowbreak{}Operations.\allowbreak{}cut\allowbreak{}Partition}, que siempre recibe particiones construidas por \texttt{deterministic\allowbreak{}CEA}, con los vectores llenos. \\[3pt]
\textbf{47}\phantomsection\label{nota:47} & \textbf{La heurística híbrida de eliminación no triangula: repite siempre la misma variable}\newline \texttt{Hybrid\allowbreak{}Elimination.\allowbreak{}get\allowbreak{}Deletion\allowbreak{}Sequence} elige la variable sobre una red y borra enlaces en otra sin quitar nunca la elegida, así que repite siempre la misma y el grafo que devuelve no está triangulado; además \texttt{is\allowbreak{}Super\allowbreak{}Value\allowbreak{}Node} declara supervalor un nodo con un solo padre de utilidad. Ningún código de producción construye un \texttt{Hybrid\allowbreak{}Elimination}, y su prueba está comentada. \\[3pt]
\textbf{48}\phantomsection\label{nota:48} & \textbf{La evolución temporal de varias utilidades falla si no se fuerza un solo criterio}\newline El constructor \texttt{MIDTemporal\allowbreak{}Evolution(\allowbreak{}Prob\allowbreak{}Net,\allowbreak{}\ List\textless{}Node\textgreater{})} deja el objeto en modo de coste-efectividad y \texttt{evaluate\allowbreak{}Slice} lanza \texttt{Class\allowbreak{}Cast\allowbreak{}Exception} al tratar la partición como una tabla (medido sobre \texttt{0\_\allowbreak{}miscellaneous/\allowbreak{}networks/\allowbreak{}mid/\allowbreak{}MID-Chancellor-corrected.\allowbreak{}pgmx}). Su único llamador, \texttt{Trace\allowbreak{}Temporal\allowbreak{}Evolution\allowbreak{}Dialog}, llama a \texttt{force\allowbreak{}Unicriterion(\allowbreak{})} justo después de construirlo, y así el cálculo termina bien. \\[3pt]
\textbf{49}\phantomsection\label{nota:49} & \textbf{La evaluación ciclo a ciclo traga la excepción de un potencial no proyectable}\newline \texttt{Temporal\allowbreak{}Evaluation.\allowbreak{}resolve} captura \texttt{Non\allowbreak{}Projectable\allowbreak{}Potential\allowbreak{}Exception}, la registra y evalúa un modelo al que le faltan ese potencial y los siguientes, sin avisar a quien llama. Solo la construyen pruebas de \texttt{integration\allowbreak{}Tests}, y en la práctica ni siquiera llega a ese punto: falla antes al pedir el registro de mensajes (\hyperref[nota:37]{nota~37}). \\[3pt]
\textbf{50}\phantomsection\label{nota:50} & \textbf{El lector de XMLBIF no lee ningún fichero, ni pone enlaces ni posiciones}\newline \texttt{XMLBIFReader} falla con \texttt{Null\allowbreak{}Pointer\allowbreak{}Exception} incluso con \texttt{io/\allowbreak{}src/\allowbreak{}test/\allowbreak{}resources/\allowbreak{}net\allowbreak{}AB.\allowbreak{}xml}, el fichero de su prueba, porque \texttt{PGMXReader\_\allowbreak{}0\_\allowbreak{}2.\allowbreak{}get\allowbreak{}Variable} es estático y no usa su forma de leer nombres; además \texttt{get\allowbreak{}Links} está vacío y no lee las posiciones de los nodos. Su anotación de formato está comentada y ningún llamador lo construye, así que hoy no se alcanza; lo que sí le pasa a quien abre desde «Abrir» un \texttt{.\allowbreak{}xml} de ese formato es \hyperref[err:117]{error~117}. \\[3pt]
\textbf{51}\phantomsection\label{nota:51} & \textbf{El exportador de árboles de decisión a Amua no lo usa ningún menú ni llamador}\newline El paquete \texttt{org.\allowbreak{}openmarkov.\allowbreak{}io.\allowbreak{}amua} (fachada \texttt{Amua\allowbreak{}Exporter}, 13 clases y 37 pruebas) convierte un árbol de decisión en un fichero del programa Amua, pero fuera del paquete y de sus pruebas solo se nombra en la línea \texttt{exports} del \texttt{module-info.\allowbreak{}java} del módulo \texttt{io}. No hay entrada de menú, acción ni anotación de complemento que lo ofrezca. \\[3pt]
\textbf{52}\phantomsection\label{nota:52} & \textbf{El exportador a Amua escribe los decimales con coma si el sistema está en castellano}\newline \texttt{Amua\allowbreak{}DTWriter} usa \texttt{new\ Decimal\allowbreak{}Format(\allowbreak{}"\#.\allowbreak{}\#\#\#\#")} con los símbolos del idioma del sistema, así que con \texttt{es\_\allowbreak{}ES} escribe \texttt{0,\allowbreak{}8}, que Amua no puede leer, y además redondea a cuatro decimales. El exportador no tiene llamadores (\hyperref[nota:51]{nota~51}), y dentro de la aplicación de escritorio tampoco se vería, porque \texttt{Open\allowbreak{}Markov.\allowbreak{}main} fija el formato numérico inglés al arrancar. \\[3pt]
\textbf{53}\phantomsection\label{nota:53} & \textbf{Pedir al algoritmo PC el motivo de una orientación de una sola flecha lanza una excepción}\newline \texttt{PCAlgorithm.\allowbreak{}get\allowbreak{}Motivation(\allowbreak{}PNEdit)} lee siempre la segunda flecha de un \texttt{COrient\allowbreak{}Links\allowbreak{}Edit}, y con uno de una sola flecha, que \texttt{Collider\allowbreak{}Orientation.\allowbreak{}try\allowbreak{}Orient\allowbreak{}Collider} construye a veces, lanza \texttt{No\allowbreak{}Such\allowbreak{}Element\allowbreak{}Exception}. Solo lo llaman pruebas: el diálogo de aprendizaje interactivo toma el motivo que ya lleva cada propuesta, construido correctamente. \\[3pt]
\textbf{54}\phantomsection\label{nota:54} & \textbf{La discretización por anchura igual lanza una excepción con una columna no numérica}\newline \texttt{Discretization.\allowbreak{}calculate\allowbreak{}Variable\allowbreak{}Min} y \texttt{calculate\allowbreak{}Variable\allowbreak{}Max} convierten los nombres de estado con \texttt{Double.\allowbreak{}parse\allowbreak{}Double} sin protección, y \texttt{Discretization.\allowbreak{}process} no comprueba si la variable es numérica. Los dos diálogos que la llaman, el de aprendizaje y el de preprocesado de \texttt{bn\allowbreak{}Evaluation}, deshabilitan la discretización y la dejan en «ninguna» para las variables no numéricas. \\[3pt]
\textbf{55}\phantomsection\label{nota:55} & \textbf{\texttt{Configuration.\allowbreak{}add\allowbreak{}Finding} descarta sin aviso el valor de una variable discretizada}\newline El \texttt{switch} de \texttt{Configuration.\allowbreak{}add\allowbreak{}Finding(\allowbreak{}Variable,\allowbreak{}\ double)} no contempla las variables discretizadas y vuelve sin hacer nada, como dice su javadoc. Sus únicos llamadores, en la simulación de eventos discretos, solo le pasan variables de sucesos y de nodos de azar, y \texttt{Chance\allowbreak{}Record} rechaza antes de simular los nodos de azar que no son numéricos ni de estados finitos (cómo llega ese rechazo al usuario se describe en \hyperref[err:106]{error~106}). \\[3pt]
\textbf{56}\phantomsection\label{nota:56} & \textbf{Una tabla de funciones sin variables tendría una fórmula nula}\newline El constructor de \texttt{Table\allowbreak{}With\allowbreak{}Functions} deja la única casilla a \texttt{null} cuando la lista de variables está vacía, y evaluarla fallaría dentro del evaluador de expresiones. Su único constructor de producción es el de \texttt{Table\allowbreak{}With\allowbreak{}Events}, que siempre le pasa al menos la variable condicionada. \\[3pt]
\textbf{57}\phantomsection\label{nota:57} & \textbf{\texttt{Table\allowbreak{}With\allowbreak{}Functions.\allowbreak{}equals} no compara las fórmulas}\newline \texttt{Table\allowbreak{}With\allowbreak{}Functions.\allowbreak{}equals} compara los números heredados y no las fórmulas, así que dos tablas que solo difieren en una fórmula salen iguales (medido). Ningún código de producción la compara: la tabla solo vive dentro de \texttt{Table\allowbreak{}With\allowbreak{}Events}, que hereda la igualdad de \texttt{Potential}, y la ventana de potenciales compara por reflexión campo a campo. \\[3pt]
\textbf{58}\phantomsection\label{nota:58} & \textbf{El análisis probabilístico de sensibilidad de la simulación de sucesos no cambia nada}\newline La rama de análisis probabilístico de sensibilidad de \texttt{DESInference} escribe los valores sorteados donde el muestreo no los lee, y sobre tablas cuyo \texttt{get\allowbreak{}Uncertain\allowbreak{}Values(\allowbreak{})} es siempre \texttt{null}, así que no tiene efecto. Solo corre con \texttt{Monte\allowbreak{}Carlo\allowbreak{}Options.\allowbreak{}is\allowbreak{}Psa(\allowbreak{})} verdadero, y nada en producción lo activa: la casilla «Perfom PSA» se crea desactivada y no se añade a la ventana, y el lector de ficheros no fija esa opción. \\[3pt]
\textbf{59}\phantomsection\label{nota:59} & \textbf{La razón coste-efectividad incremental del resumen de la simulación vale siempre cero}\newline \texttt{Simulation\allowbreak{}Summary\allowbreak{}Results} no calcula nunca la razón coste-efectividad incremental, porque el método que lo haría no tiene llamadores, y \texttt{to\allowbreak{}String} escribe \texttt{ICER;0.\allowbreak{}0}. Ese cero no llega al usuario: \texttt{DESResults\allowbreak{}Window} no enseña esa columna y la exportación a CSV (valores separados por punto y coma) la borra. \\[3pt]
\textbf{60}\phantomsection\label{nota:60} & \textbf{\texttt{Stats\allowbreak{}Properties} pide la mediana y los cuartiles con el percentil en fracción}\newline \texttt{Stats\allowbreak{}Properties} llama a \texttt{Stat\allowbreak{}Utils.\allowbreak{}percentile} con 0,5, 0,25 y 0,75, cuando ese método espera un valor de 0 a 100, así que la mediana y los cuartiles salen iguales al mínimo (medido con los valores 1 a 100). La clase no tiene ninguna referencia en el proyecto, ni en producción ni en pruebas. \\[3pt]
\textbf{61}\phantomsection\label{nota:61} & \textbf{El diálogo que ofrece elegir entre inglés y español no se puede abrir desde ningún menú}\newline \texttt{Language\allowbreak{}Dialog} muestra la lista fija «English» y «Spanish» de \texttt{Languages}, pero la opción que lo abre se crea en \texttt{Main\allowbreak{}Menu.\allowbreak{}create\allowbreak{}All\allowbreak{}Items} y \texttt{build\allowbreak{}Help\allowbreak{}Menu} no la añade a ningún menú. Lo único alcanzable es la opción de arranque \texttt{-l}/\texttt{-language}, y \texttt{String\allowbreak{}Database.\allowbreak{}set\allowbreak{}Language} rechaza cualquier idioma distinto del inglés con un aviso en el registro, como pide la decisión de que la aplicación sea solo en inglés. \\[3pt]
\textbf{62}\phantomsection\label{nota:62} & \textbf{El respaldo de \texttt{Languages.\allowbreak{}get\allowbreak{}String} para un idioma sin nombre no se ejecuta nunca}\newline \texttt{Languages.\allowbreak{}get\allowbreak{}String} espera una \texttt{Missing\allowbreak{}Resource\allowbreak{}Exception} que \texttt{String\allowbreak{}Database.\allowbreak{}get\allowbreak{}String} nunca lanza (devuelve la clave rodeada de \texttt{\textgreater{}\textgreater{}\textgreater{}} y \texttt{\textless{}\textless{}\textless{}}), así que su respaldo no funciona; lo mismo pasa en \texttt{get\allowbreak{}Short\allowbreak{}String}. Solo se llama con «English» y «Spanish», cuyos nombres existen, y su único consumidor, \texttt{Language\allowbreak{}Dialog}, no se puede abrir (\hyperref[nota:61]{nota~61}). \\[3pt]
\textbf{63}\phantomsection\label{nota:63} & \textbf{Con una clave repetida en dos ficheros de textos, gana uno según el orden de lectura}\newline \texttt{String\allowbreak{}Database} junta los textos de todos los ficheros en un mapa y el último leído pisa al anterior, en un orden que nadie fija; y dos ficheros con el mismo nombre base en módulos distintos harían desaparecer uno entero. Medido en inglés: la única clave repetida, \texttt{BUNDLEFILE.\allowbreak{}Text}, tiene el mismo texto en todos sus ficheros y ningún nombre base se repite; las colisiones reales están en los ficheros en castellano, fuera por la decisión de que la aplicación sea solo en inglés. \\[3pt]
\textbf{64}\phantomsection\label{nota:64} & \textbf{El constructor \texttt{Partitioned\allowbreak{}Interval(\allowbreak{}Object{[}{]}{[}{]})} no acepta la tabla que produce su pareja}\newline El constructor lee siete columnas por tramo donde \texttt{convert\allowbreak{}To\allowbreak{}Table\allowbreak{}Format} escribe seis, compara símbolos con \texttt{==} e interpreta el corchete al revés en la rama de varios tramos; medido con un intervalo de dos tramos, lanza \texttt{Class\allowbreak{}Cast\allowbreak{}Exception}. Ninguna construcción de \texttt{Partitioned\allowbreak{}Interval} del proyecto, ni en producción ni en pruebas, usa ese constructor. \\[3pt]
\textbf{65}\phantomsection\label{nota:65} & \textbf{\texttt{Discretize\allowbreak{}Table\allowbreak{}Panel.\allowbreak{}set\allowbreak{}Partitioned\allowbreak{}Interval} empareja cada tramo con el estado equivocado}\newline El método pega los nombres de estado invertidos sobre un volcado que no está invertido, así que la fila \emph{i} mostraría el tramo \emph{i} con el nombre del estado \emph{n-1-i}. Su único llamador de producción, \texttt{Revelation\allowbreak{}Arc\allowbreak{}Panel.\allowbreak{}set\allowbreak{}Fields\allowbreak{}From\allowbreak{}Properties}, trabaja con la subclase \texttt{Revelation\allowbreak{}Arc\allowbreak{}Discretize\allowbreak{}Table\allowbreak{}Panel}, que redefine el método, y el panel de dominio del nodo se llena con \texttt{set\allowbreak{}Data\allowbreak{}From\allowbreak{}Partitioned\allowbreak{}Interval}, que empareja bien. \\[3pt]
\textbf{66}\phantomsection\label{nota:66} & \textbf{Cuatro ediciones de opciones de la red ejecutan su cambio dos veces al rehacer}\newline \texttt{Temporal\allowbreak{}Options\allowbreak{}Edit}, \texttt{Multicriteria\allowbreak{}Edit}, \texttt{Monte\allowbreak{}Carlo\allowbreak{}Options\allowbreak{}Edit} y \texttt{Cycle\allowbreak{}Length\allowbreak{}Edit} redefinen \texttt{redo(\allowbreak{})} como \texttt{super.\allowbreak{}redo(\allowbreak{});\ do\allowbreak{}Edit(\allowbreak{});}, y \texttt{super.\allowbreak{}redo(\allowbreak{})} ya ejecuta \texttt{do\allowbreak{}Edit(\allowbreak{})}. Tres de ellas solo se crean en pruebas; \texttt{Multicriteria\allowbreak{}Edit} sí se puede rehacer desde la aplicación, pero su \texttt{do\allowbreak{}Edit(\allowbreak{})} solo asigna valores y repetirlo deja la red igual. Es el mismo patrón que \hyperref[nota:84]{nota~84} y \hyperref[nota:85]{nota~85}. \\[3pt]
\textbf{67}\phantomsection\label{nota:67} & \textbf{El nombre de la fila de un nodo numérico se recorta con la medida de los nodos de estados}\newline \texttt{Visual\allowbreak{}State.\allowbreak{}paint} recorta el nombre de la fila a 52 unidades aunque en una fila numérica la barra empieza en la 32, así que un nombre más largo quedaría debajo de la barra. Las filas numéricas solo dibujan « EU» en los nodos de utilidad, que mide 23 píxeles, o una cadena vacía en los de azar numéricos, y nadie llama a \texttt{Visual\allowbreak{}State.\allowbreak{}set\allowbreak{}State\allowbreak{}Name}. \\[3pt]
\textbf{68}\phantomsection\label{nota:68} & \textbf{La opción «Save evidence\ldots» no está en ningún menú y su código no guarda nada}\newline \texttt{Network\allowbreak{}File\allowbreak{}Handler.\allowbreak{}save\allowbreak{}Evidence} solo imprime la ruta elegida por la salida estándar, y además duplicaría la evidencia previa a la resolución. La entrada se crea en \texttt{Main\allowbreak{}Menu.\allowbreak{}create\allowbreak{}All\allowbreak{}Items}, pero \texttt{recharge\allowbreak{}File\allowbreak{}Menu} no la añade, no tiene atajo de teclado y ningún otro componente emite su mando; la evidencia sí se guarda al guardar la red. \\[3pt]
\textbf{69}\phantomsection\label{nota:69} & \textbf{«Save evidence\ldots», «Link properties» y cambiar idioma se crean y no están en ningún menú}\newline \texttt{Main\allowbreak{}Menu.\allowbreak{}create\allowbreak{}All\allowbreak{}Items} crea esas tres entradas, pero ningún método que construye los menús las coloca y ningún otro componente emite sus mandos. Lo que hay detrás es código muerto sin consecuencia visible: un guardado que no guarda (\hyperref[nota:68]{nota~68}), \texttt{Network\allowbreak{}Editor\allowbreak{}Panel.\allowbreak{}change\allowbreak{}Link\allowbreak{}Properties} con el cuerpo vacío y \texttt{Language\allowbreak{}Dialog} (\hyperref[nota:61]{nota~61}). \\[3pt]
\textbf{70}\phantomsection\label{nota:70} & \textbf{Quedan once constantes y sus textos de un submenú de zoom del menú View que ya no existe}\newline Ningún código usa las constantes \texttt{VIEW\_\allowbreak{}ZOOM\_\allowbreak{}*} de \texttt{Menu\allowbreak{}Item\allowbreak{}Names}: \texttt{Main\allowbreak{}Menu.\allowbreak{}build\allowbreak{}View\allowbreak{}Menu} solo pone las entradas de pestaña siguiente y anterior, y el zoom está en la barra de botones. Las constantes apuntan además a claves \texttt{View.\allowbreak{}Zoom\allowbreak{}Manager.\allowbreak{}*} que no existen (los textos de \texttt{Menus\_\allowbreak{}en.\allowbreak{}xml} están bajo \texttt{View.\allowbreak{}Zoom.\allowbreak{}*}); es limpieza, y el mismo renombrado sí estropea un texto visible, que es \hyperref[err:183]{error~183}. \\[3pt]
\textbf{71}\phantomsection\label{nota:71} & \textbf{Seis botones de opciones del panel de un modelo canónico no tienen escuchador}\newline En \texttt{ICIOptions\allowbreak{}Panel} solo la pareja «Canonical parameters» / «Whole table» tiene escuchador; las otras seis opciones tienen el \texttt{add\allowbreak{}Item\allowbreak{}Listener} comentado. Pero se crean con \texttt{set\allowbreak{}Enabled(\allowbreak{}false)} y nada las habilita (medido), así que nadie puede pulsarlas; lo único visible son seis opciones en gris, una con la errata «Independant parameters» en \texttt{Dialogs\_\allowbreak{}en.\allowbreak{}xml}. \\[3pt]
\textbf{72}\phantomsection\label{nota:72} & \textbf{\texttt{String\allowbreak{}With\allowbreak{}Properties.\allowbreak{}put} falla con un agente leído de fichero sin propiedades}\newline \texttt{PGMXReader\_\allowbreak{}0\_\allowbreak{}2.\allowbreak{}get\allowbreak{}Agents} sustituye la tabla de propiedades del agente por \texttt{null} cuando el \texttt{\textless{}Agent\textgreater{}} no tiene propiedades, y un \texttt{put(\allowbreak{}clave,\allowbreak{}\ valor)} posterior lanzaría \texttt{Null\allowbreak{}Pointer\allowbreak{}Exception}; un agente creado con \texttt{String\allowbreak{}With\allowbreak{}Properties(\allowbreak{}String)} sí tiene tabla. En producción nadie llama a \texttt{put(\allowbreak{}clave,\allowbreak{}\ valor)} sobre un agente, y el escritor comprueba que la tabla no sea nula. \\[3pt]
\textbf{73}\phantomsection\label{nota:73} & \textbf{Elegir el criterio de un nodo de utilidad apila cuatro ediciones iguales}\newline \texttt{Node\allowbreak{}Definition\allowbreak{}Panel.\allowbreak{}get\allowbreak{}JCombo\allowbreak{}Box\allowbreak{}Decision\allowbreak{}Criteria} añade el panel dos veces como escuchador y \texttt{item\allowbreak{}State\allowbreak{}Changed} no descarta la notificación de deselección, así que elegir otro criterio apila cuatro \texttt{Node\allowbreak{}Decision\allowbreak{}Criteria\allowbreak{}Edit}, y elegir otro agente, dos \texttt{Node\allowbreak{}Agent\allowbreak{}Edit} (medido). El diálogo de propiedades agrupa todas sus ediciones en una sola al aceptar y las deshace en orden inverso al cancelar, así que ni Ctrl+Z ni «Cancel» muestran diferencia. \\[3pt]
\textbf{74}\phantomsection\label{nota:74} & \textbf{Muestrear un árbol de relaciones vacía la configuración de padres que recibe}\newline \texttt{Tree\allowbreak{}ADDPotential.\allowbreak{}sample\allowbreak{}Conditioned\allowbreak{}Variable} y \texttt{Tree\allowbreak{}With\allowbreak{}Excluded\allowbreak{}Events\allowbreak{}Potential} quitan hallazgos de la configuración del llamador y no la devuelven como estaba (medido: entra \texttt{{[}E,\allowbreak{}\ Age{]}} y sale vacía). Todos los llamadores de producción, en la simulación de eventos discretos, crean una configuración nueva para cada muestreo y no la vuelven a leer. \\[3pt]
\textbf{75}\phantomsection\label{nota:75} & \textbf{La tabla de la evolución temporal fallaría si se construyera antes que la gráfica}\newline \texttt{Temporal\allowbreak{}Evolution\allowbreak{}Table\allowbreak{}Pane} lee los valores por índice sin repetirlos cuando el resultado de un ciclo no depende de la decisión, y construida sola lanza \texttt{Array\allowbreak{}Index\allowbreak{}Out\allowbreak{}Of\allowbreak{}Bounds\allowbreak{}Exception} (medido). En la ventana, \texttt{Trace\allowbreak{}Temporal\allowbreak{}Evolution\allowbreak{}Dialog.\allowbreak{}create\allowbreak{}Series} construye antes la gráfica y deja guardados los valores ya repetidos, así que la tabla sale bien; queda un acoplamiento oculto entre las dos pestañas. \\[3pt]
\textbf{76}\phantomsection\label{nota:76} & \textbf{El informe de la evolución temporal fallaría con una tabla sin columna de nombres}\newline \texttt{Temporal\allowbreak{}Evolution\allowbreak{}Report.\allowbreak{}write} convierte a \texttt{String} la primera columna de la tabla y lanzaría \texttt{Class\allowbreak{}Cast\allowbreak{}Exception} si esa columna tuviera números, como en la tabla por variable de una utilidad sin decisión. Pero \texttt{Trace\allowbreak{}Temporal\allowbreak{}Evolution\allowbreak{}Dialog.\allowbreak{}get\allowbreak{}Table\allowbreak{}Pane} nunca construye esa tabla: para utilidades usa el constructor por criterio, cuya primera columna son siempre nombres; medido, el informe de una utilidad se escribe bien con y sin decisión. \\[3pt]
\textbf{77}\phantomsection\label{nota:77} & \textbf{El plano de coste-efectividad de la simulación de sucesos fallaría con un solo criterio}\newline \texttt{CEDESDialog.\allowbreak{}get\allowbreak{}CEPlane\allowbreak{}Chart\allowbreak{}Panel} usa \texttt{criteria.\allowbreak{}get(\allowbreak{}1)} sin comprobar, así que con un solo criterio lanzaría \texttt{Index\allowbreak{}Out\allowbreak{}Of\allowbreak{}Bounds\allowbreak{}Exception} y tampoco se vería la tabla de resultados. Pero \texttt{DESInference.\allowbreak{}initialize}, por la que pasan los dos constructores, avisa y no simula con menos de dos criterios, y no hay otro llamador de \texttt{CEDESDialog}. \\[3pt]
\textbf{78}\phantomsection\label{nota:78} & \textbf{La casilla «Always append» lleva el nombre de componente de un campo antiguo}\newline \texttt{Node\allowbreak{}Definition\allowbreak{}Panel.\allowbreak{}get\allowbreak{}JCheckbox\allowbreak{}Always\allowbreak{}Append} crea la casilla con el nombre de componente \texttt{j\allowbreak{}Checkbox\allowbreak{}Override\allowbreak{}Time\allowbreak{}Stamp}, que no dice lo que es; ese nombre solo sirve para encontrarla desde pruebas de interfaz o paseos guiados. Ningún código la busca por ese nombre ni por \texttt{j\allowbreak{}Checkbox\allowbreak{}Always\allowbreak{}Append}; si se corrige, conviene hacerlo junto con \hyperref[err:34]{error~34}, que es el defecto de deshacer de esta misma casilla. \\[3pt]
\textbf{79}\phantomsection\label{nota:79} & \textbf{Invertir un enlace recalculando los potenciales puede cerrar un ciclo}\newline \texttt{Invert\allowbreak{}Link\allowbreak{}And\allowbreak{}Update\allowbreak{}Potentials\allowbreak{}Edit} no comprueba ninguna restricción y, ejecutada directamente, deja una red bayesiana que incumple \texttt{No\allowbreak{}Cycle} (medido), cuando su hermana \texttt{Invert\allowbreak{}Link\allowbreak{}Edit} rechazaría esa inversión. Su único llamador de producción es la opción del menú contextual del enlace, que antes pregunta a \texttt{Link\allowbreak{}Inversion\allowbreak{}With\allowbreak{}Potentials\allowbreak{}Update\allowbreak{}Validator} y sale deshabilitada si la inversión cerraría un ciclo. \\[3pt]
\textbf{80}\phantomsection\label{nota:80} & \textbf{El inverso de orientar un enlace es la propia orientación}\newline \texttt{Orient\allowbreak{}Link\allowbreak{}Edit.\allowbreak{}get\allowbreak{}Undo\allowbreak{}Edit(\allowbreak{})} devuelve la misma edición en vez de una que devuelva el enlace a no dirigido. Su único consumidor, \texttt{Metric.\allowbreak{}after\allowbreak{}Undoing\allowbreak{}Edit}, solo lo usan los algoritmos de puntuación y búsqueda, que nunca crean orientaciones, y el algoritmo PC, que sí las crea, no usa \texttt{Metric}. \\[3pt]
\textbf{81}\phantomsection\label{nota:81} & \textbf{Subir el primer criterio de decisión o bajar el último: la ventana no lo permite}\newline \texttt{Decision\allowbreak{}Criteria\allowbreak{}Edit.\allowbreak{}do\allowbreak{}Edit} no comprueba que haya sitio al subir o bajar un criterio, pero sus únicos llamadores con \texttt{UP} o \texttt{DOWN} son los botones de \texttt{Decision\allowbreak{}Criteria\allowbreak{}Table\allowbreak{}Panel}, que \texttt{Key\allowbreak{}Table\allowbreak{}Panel.\allowbreak{}value\allowbreak{}Changed} deshabilita en la primera y la última fila y sin selección, y el criterio sale siempre de la propia lista. Es el mismo caso que \hyperref[nota:36]{nota~36}; el defecto vecino de deshacer el renombrado de un criterio es \hyperref[err:57]{error~57}. \\[3pt]
\textbf{82}\phantomsection\label{nota:82} & \textbf{Pegar una selección de solo enlaces revienta al calcular el recuadro}\newline \texttt{Paste\allowbreak{}Edit} calcula el máximo de las coordenadas de los nodos pegados, y con una lista vacía lanzaría \texttt{No\allowbreak{}Such\allowbreak{}Element\allowbreak{}Exception}. \texttt{Main\allowbreak{}Panel\allowbreak{}Menu\allowbreak{}Assistant} deshabilita Copiar y Cortar (en el menú, en sus atajos y en la barra) cuando solo hay enlaces seleccionados, así que el portapapeles siempre lleva al menos un nodo. Un nodo acompañado de un enlace cuyo otro extremo no se copió sí llega, y falla de otra manera: \hyperref[nota:184]{nota~184}. \\[3pt]
\textbf{83}\phantomsection\label{nota:83} & \textbf{Editar un parámetro de un modelo canónico deja dos ediciones y un deshacer vacío}\newline \texttt{ICITable\allowbreak{}Potential\allowbreak{}Value\allowbreak{}Edit.\allowbreak{}do\allowbreak{}Edit} ejecuta dentro un \texttt{ICIPotential\allowbreak{}Edit} como edición aparte, y su propio \texttt{undo(\allowbreak{})} no hace nada. Las dos entradas caen en el historial provisional de \texttt{Potential\allowbreak{}Edit\allowbreak{}Panel}: al aceptar se descarta y se sustituye por un único \texttt{Set\allowbreak{}Potential\allowbreak{}Edit} que se deshace bien, y al cancelar se deshace entero y la edición interna repone los parámetros. \\[3pt]
\textbf{84}\phantomsection\label{nota:84} & \textbf{Rehacer «quitar hallazgo» repite la retirada}\newline \texttt{Remove\allowbreak{}Finding\allowbreak{}Edit.\allowbreak{}redo(\allowbreak{})} no apaga el interruptor \texttt{typical\allowbreak{}Redo}, así que la clase base ejecuta \texttt{do\allowbreak{}Edit(\allowbreak{})} y después la clase vuelve a quitar el hallazgo. La segunda retirada no encuentra nada y no cambia nada; el cambio de contrato de las ediciones descrito en \hyperref[nota:154]{nota~154} lo haría desaparecer de raíz. \\[3pt]
\textbf{85}\phantomsection\label{nota:85} & \textbf{Rehacer un cambio de potencial aplica el cambio dos veces}\newline \texttt{Set\allowbreak{}Potential\allowbreak{}Edit.\allowbreak{}redo(\allowbreak{})} llama a la clase base sin apagar el interruptor de rehacer y después repite \texttt{set\allowbreak{}Potential\allowbreak{}With\allowbreak{}Restrictions}. Repetirlo no cambia ningún valor, porque las casillas prohibidas ya valen cero y cada columna ya suma uno (salvo, a lo sumo, redondeo en el último decimal); también desaparecería con el cambio de contrato de \hyperref[nota:154]{nota~154}. \\[3pt]
\textbf{86}\phantomsection\label{nota:86} & \textbf{El constructor de un argumento de \texttt{Set\allowbreak{}Potential\allowbreak{}Edit} no comprueba la lista vacía}\newline \texttt{Set\allowbreak{}Potential\allowbreak{}Edit(\allowbreak{}Node)} toma \texttt{get\allowbreak{}Potentials(\allowbreak{}).\allowbreak{}get\allowbreak{}First(\allowbreak{})} sin comprobar que el nodo tenga potenciales, al contrario que el de dos argumentos. No tiene ningún llamador, ni en producción ni en pruebas; borrarlo o darle la misma guarda es limpieza. \\[3pt]
\textbf{87}\phantomsection\label{nota:87} & \textbf{El conjunto de oyentes de ediciones promete ser seguro entre hilos y no lo es}\newline \texttt{PNESupport} guarda los oyentes en un conjunto sincronizado que \texttt{PNEdit.\allowbreak{}execute\allowbreak{}Edit}, \texttt{undo} y \texttt{redo} recorren sin tomar su cerrojo, y \texttt{set\allowbreak{}Listeners} lo vacía y lo rellena en dos pasos, contra lo que dice su javadoc. No se ha encontrado ningún camino de producción que ejecute ediciones o cambie oyentes desde otro hilo mientras se avisa, ni que añada un oyente dentro de un aviso con otros oyentes detrás. \\[3pt]
\textbf{88}\phantomsection\label{nota:88} & \textbf{Tres restricciones convierten la red a \texttt{Prob\allowbreak{}Net} sin comprobar su clase}\newline \texttt{Proper\allowbreak{}Utility\allowbreak{}Potentials}, \texttt{Utility\allowbreak{}Nodes} y \texttt{Only\allowbreak{}One\allowbreak{}Agent} hacen un molde del \texttt{Graph\allowbreak{}Network} recibido a \texttt{Prob\allowbreak{}Net} o \texttt{Potential\allowbreak{}Network} sin mirar, lo que fallaría con otra implementación de esas interfaces. Hoy \texttt{Prob\allowbreak{}Net} es la única; \texttt{Valid\allowbreak{}Criterion\allowbreak{}Name} muestra la forma segura, con \texttt{instanceof}. \\[3pt]
\textbf{89}\phantomsection\label{nota:89} & \textbf{La sección de restricciones adicionales del fichero escribe frases y lee nombres de clase}\newline \texttt{PGMXWriter\_\allowbreak{}0\_\allowbreak{}2.\allowbreak{}get\allowbreak{}Additional\allowbreak{}Constraints} se salta la primera restricción, solo escribe las que el tipo de red prohíbe y las escribe como texto para el usuario, mientras que el lector espera nombres de clase, así que no podría reabrir ese fichero (medido). Ninguna red construida desde la aplicación lleva restricciones que su tipo prohíba, así que hoy el escritor no escribe nada en esa sección. \\[3pt]
\textbf{90}\phantomsection\label{nota:90} & \textbf{La rama de respaldo de \texttt{Set\allowbreak{}Potential\allowbreak{}Variables\allowbreak{}Edit} desordena la tabla, pero no se alcanza}\newline Si \texttt{reorder(\allowbreak{}List)} devuelve \texttt{null}, \texttt{Set\allowbreak{}Potential\allowbreak{}Variables\allowbreak{}Edit.\allowbreak{}do\allowbreak{}Edit} cambia el orden de las variables del potencial original sin mover sus valores, y la tabla quedaría mal. Su único constructor de producción es el botón de reordenar variables de \texttt{Potential\allowbreak{}Edit\allowbreak{}Panel}, que solo llega con tablas y modelos canónicos, cuyos \texttt{reorder} devuelven siempre un potencial. \\[3pt]
\textbf{91}\phantomsection\label{nota:91} & \textbf{Coste-efectividad declara inalcanzable la evidencia incompatible, y hoy lo es}\newline En el ámbito global, la única vía hacia \texttt{Evidence\allowbreak{}Is\allowbreak{}Incompatible\allowbreak{}With\allowbreak{}Other} era la discretización que pisaba la evidencia del usuario, corregida en el commit \texttt{66d32ab}, y la conversión que queda en \texttt{VECEAnalysis.\allowbreak{}get\allowbreak{}Utility} quita el hallazgo antes de añadir el de la variable convertida. En el ámbito por decisión, \texttt{Scope\allowbreak{}Selector\allowbreak{}Panel} fija y deshabilita el estado de las variables ya observadas. Los dos tratamientos distintos siguen escritos; unificarlos entra en \hyperref[err:48]{error~48}. \\[3pt]
\textbf{92}\phantomsection\label{nota:92} & \textbf{Las operaciones públicas de \texttt{Potential\allowbreak{}Operations} rechazan subclases de \texttt{Table\allowbreak{}Potential}}\newline \texttt{Auxiliary\allowbreak{}Operations.\allowbreak{}check\allowbreak{}Objects\allowbreak{}Collection\allowbreak{}Type} compara la clase exacta, así que \texttt{multiply}, \texttt{multiply\allowbreak{}And\allowbreak{}Maximize} y las dos \texttt{multiply\allowbreak{}And\allowbreak{}Marginalize} de \texttt{Potential\allowbreak{}Operations} rechazan una \texttt{Strategic\allowbreak{}Table\allowbreak{}Potential} o una \texttt{Uncertain\allowbreak{}Table\allowbreak{}Potential} (medido), aunque \texttt{Discrete\allowbreak{}Potential\allowbreak{}Operations} sabe tratarlas. Su único llamador de producción es \texttt{Cluster\allowbreak{}Of\allowbreak{}Variables.\allowbreak{}get\allowbreak{}Downgoing\allowbreak{}Potential}, que nadie llama; la propagación de Hugin usa directamente \texttt{Discrete\allowbreak{}Potential\allowbreak{}Operations}. \\[3pt]
\textbf{93}\phantomsection\label{nota:93} & \textbf{Maximizar operandos que son todos constantes devuelve una tabla de elecciones vacía}\newline El atajo de \texttt{Table\allowbreak{}Potential\allowbreak{}Maximization.\allowbreak{}multiply\allowbreak{}And\allowbreak{}Maximize} para operandos sin variables no añade ninguna elección, cuando debería dejar empatados todos los estados de la decisión, como hace \texttt{multiply\allowbreak{}And\allowbreak{}Maximize\allowbreak{}Uniformly} (medido: tabla de elecciones de tamaño 0). El único lector de esa tabla, \texttt{Strategy.\allowbreak{}Policy}, solo se construye desde \texttt{Strategy(\allowbreak{}Strategy\allowbreak{}Utilities)}, que nadie llama, y la absorción de un nodo de decisión solo usa la tabla de valores. \\[3pt]
\textbf{94}\phantomsection\label{nota:94} & \textbf{Multiplicar una tabla estratégica sin árboles de estrategia lanza \texttt{Null\allowbreak{}Pointer\allowbreak{}Exception}}\newline \texttt{Table\allowbreak{}Potential\allowbreak{}Arithmetic.\allowbreak{}multiply} y \texttt{evaluate\allowbreak{}Function\allowbreak{}Potential} leen \texttt{strategy\allowbreak{}Trees} de una \texttt{Strategic\allowbreak{}Table\allowbreak{}Potential} sin comprobar que no sea nulo, cosa que sí hacen \texttt{sum} y otras operaciones del mismo paquete (medido). Todos los sitios de producción que construyen una \texttt{Strategic\allowbreak{}Table\allowbreak{}Potential} le asignan los árboles antes de soltarla, y el único que no lo hace, \texttt{VEOptimal\allowbreak{}Intervention} sin decisiones, no la multiplica. \\[3pt]
\textbf{95}\phantomsection\label{nota:95} & \textbf{Sumar fuera una variable de una utilidad estratégica sin árboles lanza una excepción}\newline \texttt{Table\allowbreak{}Potential\allowbreak{}Elimination.\allowbreak{}multiply\allowbreak{}And\allowbreak{}Marginalize(\allowbreak{}probabilidad,\allowbreak{}\ utilidad,\allowbreak{}\ variable)} y el constructor de \texttt{Sum\allowbreak{}Out\allowbreak{}Variable} leen \texttt{strategy\allowbreak{}Trees} sin comprobar que no sea nulo (medido en los dos). Como en \hyperref[nota:94]{nota~94}, ninguna ruta de producción entrega a estas operaciones una \texttt{Strategic\allowbreak{}Table\allowbreak{}Potential} sin árboles. \\[3pt]
\textbf{96}\phantomsection\label{nota:96} & \textbf{La clase abstracta \texttt{ICIPotential} figura entre los tipos de potencial aplicables a un nodo}\newline \texttt{Potential\allowbreak{}Utils.\allowbreak{}get\allowbreak{}Filtered\allowbreak{}Potential\allowbreak{}Classes} devuelve \texttt{ICIPotential} para cualquier nodo con padres, porque está anotada con \texttt{@Potential\allowbreak{}Type} y su \texttt{validate} la acepta, aunque construirla lanza \texttt{Unreachable\allowbreak{}Exception} (medido). De sus dos llamadores de producción, \texttt{Potential\allowbreak{}Edit\allowbreak{}Panel.\allowbreak{}get\allowbreak{}Potential\allowbreak{}Type\allowbreak{}JCombobox} quita las clases no concretas y \texttt{Change\allowbreak{}Node\allowbreak{}Type\allowbreak{}Edit.\allowbreak{}do\allowbreak{}Multi\allowbreak{}Step\allowbreak{}Edit} solo pregunta si la lista contiene la clase del potencial actual, que nunca es la abstracta. \\[3pt]
\textbf{97}\phantomsection\label{nota:97} & \textbf{El MAX ruidoso marca su tabla como probabilidad conjunta y normalizarla la estropea}\newline \texttt{Max\allowbreak{}Potential} construye sus tablas intermedias con papel de probabilidad conjunta y \texttt{Min\allowbreak{}Potential} con papel de condicionada, así que la tabla del MAX sale como conjunta y \texttt{Table\allowbreak{}Potential\allowbreak{}Transform.\allowbreak{}normalize} la divide por la suma de toda la tabla (medido: cada columna queda sumando 0,25). La vía que llevaría a normalizarla, la evolución temporal de un diagrama de influencia de Markov, no normaliza nunca (es \hyperref[err:67]{error~67}), y no se encontró otro consumidor de producción cuyo resultado dependa de ese papel. \\[3pt]
\textbf{98}\phantomsection\label{nota:98} & \textbf{Ordenar potenciales por el orden total de las decisiones devuelve cada potencial dos veces}\newline \texttt{Table\allowbreak{}Potential\allowbreak{}Merge.\allowbreak{}order\allowbreak{}Potentials\allowbreak{}By\allowbreak{}Total\allowbreak{}Order} añade al final la lista de entrada entera en vez de solo los potenciales que quedan sin colocar (medido: entran dos potenciales y salen cuatro), así que multiplicar esa lista contaría cada factor dos veces. El método, expuesto también en \texttt{Discrete\allowbreak{}Potential\allowbreak{}Operations}, no tiene ningún llamador, ni en producción ni en pruebas. \\[3pt]
\textbf{99}\phantomsection\label{nota:99} & \textbf{\texttt{matrix\allowbreak{}Potential} con una lista de solo probabilidades da todo ceros}\newline \texttt{Table\allowbreak{}Potential\allowbreak{}Transform.\allowbreak{}matrix\allowbreak{}Potential} multiplica el producto de las probabilidades por la suma de las utilidades, y sin utilidades esa suma es la constante 0 (medido: \texttt{{[}0.\allowbreak{}3,\allowbreak{}\ 0.\allowbreak{}7{]}} da \texttt{{[}0.\allowbreak{}0,\allowbreak{}\ 0.\allowbreak{}0{]}}). Ni el método ni su fachada en \texttt{Discrete\allowbreak{}Potential\allowbreak{}Operations} tienen llamadores, tampoco en pruebas. \\[3pt]
\textbf{100}\phantomsection\label{nota:100} & \textbf{La copia vieja de \texttt{Util} de las operaciones formatea mal negativos y usa coma decimal}\newline \texttt{org.\allowbreak{}openmarkov.\allowbreak{}core.\allowbreak{}model.\allowbreak{}network.\allowbreak{}potential.\allowbreak{}operation.\allowbreak{}Util}, copia antigua de \texttt{org.\allowbreak{}openmarkov.\allowbreak{}core.\allowbreak{}model.\allowbreak{}network.\allowbreak{}Util}, da \texttt{print\allowbreak{}Integer(\allowbreak{}-1510000L)} = \texttt{"0"} y, con la configuración regional española, \texttt{rounded\allowbreak{}String(\allowbreak{}3.\allowbreak{}4,\allowbreak{}\ "0.\allowbreak{}001")} = \texttt{"3,\allowbreak{}400"} (medido). Ningún código de producción la importa; solo la usa \texttt{Util\allowbreak{}Test}, y lo razonable es borrarla con su prueba. \\[3pt]
\textbf{101}\phantomsection\label{nota:101} & \textbf{Reordenar un modelo canónico da a la copia las variables auxiliares del original}\newline \texttt{ICIPotential.\allowbreak{}reorder} copia el potencial y después rellena el mapa de la copia con las variables auxiliares del original, así que los dos comparten esos objetos (medido con un modelo de ajuste); se llega desde el botón de reordenar variables del diálogo del nodo. No da ningún resultado distinto porque la edición pone la copia en el nodo y guarda el original solo para deshacer, de modo que nunca coinciden en una misma inferencia; si coincidieran, la eliminación los trataría como la misma variable. \\[3pt]
\textbf{102}\phantomsection\label{nota:102} & \textbf{La tabla expandida que un modelo canónico guarda en caché no está protegida entre hilos}\newline \texttt{ICIPotential.\allowbreak{}get\allowbreak{}Probability} guarda la tabla completa en \texttt{expanded\allowbreak{}Potential}, un campo sin \texttt{volatile} ni sincronización, así que dos hilos a la vez podrían ver una tabla a medio publicar o anterior a un cambio. Su único llamador de producción, \texttt{Likelihood\allowbreak{}Weighting}, corre en un solo hilo, y el único cálculo con varios hilos, \texttt{VECEPSA}, usa eliminación de variables, que no pasa por aquí. Es el mismo caso que \hyperref[nota:104]{nota~104}. \\[3pt]
\textbf{103}\phantomsection\label{nota:103} & \textbf{El análisis probabilístico de sensibilidad con varios hilos repite muestras}\newline Con \texttt{use\allowbreak{}Multithreading} activado, las tareas de \texttt{VECEPSA} sortean y copian la misma red compartida, y entre el sorteo de un hilo y su copia otro hilo puede volver a sortear (medido sobre \texttt{integration\allowbreak{}Tests/\allowbreak{}src/\allowbreak{}main/\allowbreak{}resources/\allowbreak{}networks/\allowbreak{}mid/\allowbreak{}MID-dmhee-4.\allowbreak{}7.\allowbreak{}pgmx}: 71 resultados repetidos de 300). El único llamador de producción, el diálogo del análisis probabilístico de coste-efectividad, pone \texttt{set\allowbreak{}Use\allowbreak{}Multithreading(\allowbreak{}false)}; la rama de varios hilos solo la activan las pruebas. \\[3pt]
\textbf{104}\phantomsection\label{nota:104} & \textbf{La caché de la tabla expandida de un modelo canónico no se usa desde dos hilos}\newline \texttt{ICIPotential.\allowbreak{}get\allowbreak{}Probability} calcula y guarda \texttt{expanded\allowbreak{}Potential} sin sincronización, y los métodos que cambian parámetros la anulan, así que un hilo podría leer una tabla vieja. Su único llamador de producción, \texttt{Likelihood\allowbreak{}Weighting}, corre en un hilo, y \texttt{VECEPSA} con varios hilos (\hyperref[nota:103]{nota~103}) evalúa con eliminación de variables, que no pasa por esa caché. Es el mismo caso que \hyperref[nota:102]{nota~102}. \\[3pt]
\textbf{105}\phantomsection\label{nota:105} & \textbf{Una propagación reutilizada con otra evidencia contesta con la consulta anterior}\newline \texttt{VEPropagation.\allowbreak{}get\allowbreak{}Posterior\allowbreak{}Values} no vuelve a calcular si ya tiene un resultado, \texttt{Cluster\allowbreak{}Propagation} acumula la evidencia vieja (\hyperref[nota:39]{nota~39}) y \texttt{Stochastic\allowbreak{}Propagation} quita para siempre de las variables de interés las que tuvieron hallazgo; medido sobre \texttt{0\_\allowbreak{}miscellaneous/\allowbreak{}networks/\allowbreak{}bn/\allowbreak{}BN-asia.\allowbreak{}pgmx}, \texttt{VEPropagation} reutilizada devuelve exactamente la respuesta anterior. Todos los llamadores de producción construyen un objeto nuevo por consulta; fijar un ciclo de vida único para las tareas es la propuesta de \hyperref[nota:167]{nota~167}. \\[3pt]
\textbf{106}\phantomsection\label{nota:106} & \textbf{Un complemento de formato o herramienta que no se construye inutiliza su gestor}\newline \texttt{Format\allowbreak{}Manager} y \texttt{Tool\allowbreak{}Plugin\allowbreak{}Manager} construyen todos sus complementos en su inicializador estático, y si uno falla la primera llamada lanza \texttt{Exception\allowbreak{}In\allowbreak{}Initializer\allowbreak{}Error} y las siguientes \texttt{No\allowbreak{}Class\allowbreak{}Def\allowbreak{}Found\allowbreak{}Error}, sin decir cuál; \texttt{Inference\allowbreak{}Manager.\allowbreak{}build}, en cambio, deja fuera solo al culpable. Medido: hoy todos los lectores, escritores y herramientas se construyen sin error, así que solo lo dispararía un complemento defectuoso futuro (la política única para instanciar complementos es \hyperref[nota:171]{nota~171}). \\[3pt]
\textbf{107}\phantomsection\label{nota:107} & \textbf{El preprocesado del aprendizaje no devuelve a «no discretizar» una variable no numérica}\newline En \texttt{Learning\allowbreak{}Dialog.\allowbreak{}discretize\allowbreak{}Combo\allowbreak{}Box\allowbreak{}Action\allowbreak{}Performed} (línea 814) se llama a \texttt{set\allowbreak{}Selected\allowbreak{}Item(\allowbreak{}0)} donde se quería \texttt{set\allowbreak{}Selected\allowbreak{}Index(\allowbreak{}0)}, así que una variable no numérica que la red modelo tiene discretizada se quedaría con «Use model network». Pero la pestaña de preprocesado no se añade a la ventana de aprendizaje (la línea 544 que la añade está comentada), así que nadie puede cambiar esa opción; el panel equivalente de \texttt{bn\allowbreak{}Evaluation} usa \texttt{set\allowbreak{}Selected\allowbreak{}Index} y no tiene el defecto. \\[3pt]
\textbf{108}\phantomsection\label{nota:108} & \textbf{Pedir variables al muestreo en otro orden cruza sus probabilidades, pero nadie lo hace}\newline \texttt{Stochastic\allowbreak{}Propagation.\allowbreak{}get\allowbreak{}Posterior\allowbreak{}Values} (líneas 164-172) suma cada muestra en la fila de la k-ésima variable de interés que encuentra al recorrer las variables en orden de muestreo, pero la guarda a nombre de la k-ésima de la lista que le dio el llamador; medido en \texttt{0\_\allowbreak{}miscellaneous/\allowbreak{}networks/\allowbreak{}bn/\allowbreak{}BN-asia.\allowbreak{}pgmx} con las variables en el orden de \texttt{get\allowbreak{}Variables(\allowbreak{})}, P(Tuberculosis \textbar{} X-ray = yes) sale {[}0,43, 0,57{]} en lugar de {[}0,91, 0,09{]}. Nadie llega: el único llamador de producción, \texttt{Stochastic\allowbreak{}Propagation\allowbreak{}Output\allowbreak{}Frame}, no llama a \texttt{set\allowbreak{}Variables\allowbreak{}Of\allowbreak{}Interest} y usa la lista del constructor, ordenada topológicamente, y la ventana principal propaga siempre con \texttt{VEPropagation}. --- cambios --- \\[3pt]
\textbf{109}\phantomsection\label{nota:109} & \textbf{Hugin usa como probabilidades la tabla de un nodo numérico, escrita a mano en el fichero}\newline En una red bayesiana P → N con N numérico de azar y un \texttt{Table\allowbreak{}Potential} con los valores 3 y 7, \texttt{Hugin\allowbreak{}Propagation} devuelve P = {[}0,3, 0,7{]} (multiplica esos valores como si fueran probabilidades), el muestreo ≈ {[}0,5, 0,5{]} y \texttt{VEPropagation} falla con \texttt{Null\allowbreak{}Pointer\allowbreak{}Exception} (\hyperref[err:42]{error~42}). Esa tabla solo la trae un \texttt{.\allowbreak{}pgmx} escrito a mano, porque el lector acepta \texttt{Table} sobre una variable numérica; el editor no la crea (pone una relación uniforme y \texttt{Table\allowbreak{}Potential.\allowbreak{}validate} excluye las variables numéricas), guardar con \texttt{PGMXWriter\_\allowbreak{}0\_\allowbreak{}2} o \texttt{PGMXWriter\_\allowbreak{}1\_\allowbreak{}0} y reabrir da \texttt{Exact\allowbreak{}Distr\allowbreak{}Potential} o \texttt{Univariate\allowbreak{}Distr\allowbreak{}Potential}, y ninguna red bayesiana del repositorio tiene nodos numéricos. Con las relaciones que ofrece el editor («Exact», «UnivariateDistr») la herramienta de propagación estocástica falla antes, sin dar números distintos ni decir que no admite nodos numéricos (\hyperref[err:108]{error~108}). \\[3pt]
\textbf{110}\phantomsection\label{nota:110} & \textbf{Escalar una utilidad de regresión ya muestreada no la cambia, pero hoy ningún análisis llega}\newline \texttt{scale\allowbreak{}Potential} de \texttt{Exponential\allowbreak{}Potential} y de \texttt{Linear\allowbreak{}Combination\allowbreak{}Potential} solo cambia \texttt{coefficients}, y \texttt{GLMPotential.\allowbreak{}table\allowbreak{}Project} usa \texttt{sampled\allowbreak{}Coefficients} cuando existen; como \texttt{VECEPSA} muestrea antes de que \texttt{VECEAnalysis} descuente y aplique la escala de coste-efectividad, esas dos operaciones se perderían (medido en el potencial: muestreado y escalado por 0,5 sigue valiendo 100). Hoy no se llega: las únicas redes con utilidades de regresión inciertas, \texttt{0\_\allowbreak{}miscellaneous/\allowbreak{}networks/\allowbreak{}mid/\allowbreak{}MID-CHAP-Ryan-Griffin.\allowbreak{}pgmx} y sus copias, no pasan de la discretización, y una utilidad de regresión en un nodo hace fallar el análisis de coste-efectividad con \texttt{Class\allowbreak{}Cast\allowbreak{}Exception} antes de muestrear (medido con \texttt{0\_\allowbreak{}miscellaneous/\allowbreak{}networks/\allowbreak{}id/\allowbreak{}ID-CEA-minimal.\allowbreak{}pgmx} y \texttt{0\_\allowbreak{}miscellaneous/\allowbreak{}networks/\allowbreak{}mid/\allowbreak{}MID-hip-Briggs.\allowbreak{}pgmx}). \\[3pt]
\textbf{111}\phantomsection\label{nota:111} & \textbf{La casilla «Results Per Series (.xlsx file)» no hace nada, pero está siempre desactivada}\newline \texttt{Monte\allowbreak{}Carlo\allowbreak{}Options\allowbreak{}Panel.\allowbreak{}get\allowbreak{}JCheck\allowbreak{}Box\allowbreak{}Excel\allowbreak{}Results\allowbreak{}Per\allowbreak{}Series} desactiva siempre la casilla (línea 333, «Provisional for paper .jar»), y \texttt{PGMXReader\_\allowbreak{}1\_\allowbreak{}0} no lee esa opción, así que en la ventana de la simulación de sucesos sale gris y desmarcada (medido construyendo el panel). Aunque se marcara, nadie lee \texttt{Monte\allowbreak{}Carlo\allowbreak{}Options.\allowbreak{}is\allowbreak{}Results\allowbreak{}To\allowbreak{}Excel(\allowbreak{})} y la clase que escribiría el Excel, \texttt{DESResults\allowbreak{}Excel\allowbreak{}Writer}, no tiene ninguna referencia. Si se decide retirar controles del mismo panel (\hyperref[err:109]{error~109}), esta casilla y esa clase van con ellos. \\[3pt]
\end{longtable}

\subsection*{Fuera de alcance por decisión de equipo}\addcontentsline{toc}{subsection}{Fuera de alcance por decisión de equipo}
\begin{longtable}{@{}>{\raggedright\arraybackslash}p{1.1cm}>{\raggedright\arraybackslash}p{\dimexpr\linewidth-1.1cm-2\tabcolsep\relax}@{}}
\textbf{112}\phantomsection\label{nota:112} & \textbf{Una relación de Elvira que ningún modelo canónico reclama hace perder la red entera}\newline \texttt{Elvira\allowbreak{}Parser.\allowbreak{}add\allowbreak{}Sub\allowbreak{}Potentials} lanza \texttt{Some\allowbreak{}Subpotentials\allowbreak{}Arent\allowbreak{}Linked\allowbreak{}To\allowbreak{}An\allowbreak{}ICIPotential} si queda alguna subrelación sin reclamar, y no se abre nada de la red. El formato Elvira está congelado por decisión de equipo; comprobado solo leyendo el código. \\[3pt]
\textbf{113}\phantomsection\label{nota:113} & \textbf{El escritor de Elvira repite la variable de cada utilidad y crea un enlace a sí misma}\newline \texttt{Elvira\allowbreak{}Writer} escribe la variable de utilidad y después otra vez todas las variables de la relación (\texttt{relation\ U\ U\ A\ D2\ E}), y la red releída tiene un enlace \texttt{U\ -\textgreater{}\ U}. El formato Elvira está congelado por decisión de equipo. \\[3pt]
\textbf{114}\phantomsection\label{nota:114} & \textbf{El escritor de Elvira escribe \texttt{Generalized\allowbreak{}Min}, que su propio analizador no reconoce}\newline \texttt{Elvira\allowbreak{}Writer} escribe el modelo MIN generalizado como \texttt{Generalized\allowbreak{}Min}, pero \texttt{Reserved\allowbreak{}Word\allowbreak{}Tokens} no tiene esa palabra, así que el fichero no se puede releer. El formato Elvira está congelado por decisión de equipo. \\[3pt]
\textbf{115}\phantomsection\label{nota:115} & \textbf{El escritor de Elvira muere al escribir un nodo de supervalor y deja el fichero a medias}\newline \texttt{Elvira\allowbreak{}Writer} convierte las relaciones que no son tablas con \texttt{table\allowbreak{}Project(\allowbreak{}null,\allowbreak{}\ null)}, y una suma o un producto de utilidades lanza porque necesita las tablas de sus hijos; el fichero queda sin cerrar. El formato Elvira está congelado por decisión de equipo. \\[3pt]
\textbf{116}\phantomsection\label{nota:116} & \textbf{Guardar una red en Elvira pierde el título de todos los nodos}\newline \texttt{Elvira\allowbreak{}Writer} vuelve a intercambiar nombre y título antes de escribir, pero no escribe ninguna línea \texttt{title}, así que al releer los nodos se llaman por su nombre corto. El formato Elvira está congelado por decisión de equipo. \\[3pt]
\textbf{117}\phantomsection\label{nota:117} & \textbf{Un nodo continuo de Elvira sin precisión, mínimo o máximo hace fallar la lectura}\newline \texttt{Elvira\allowbreak{}Parser} convierte \texttt{Min}, \texttt{Max} y \texttt{Precision} con \texttt{Double.\allowbreak{}parse\allowbreak{}Double} sin comprobar que existan, y \texttt{Elvira\allowbreak{}Writer} solo escribe la precisión si el nodo la tenía en sus propiedades, así que un fichero que se leía puede dejar de leerse tras guardarlo. El formato Elvira está congelado por decisión de equipo. \\[3pt]
\textbf{118}\phantomsection\label{nota:118} & \textbf{Una tabla generalizada de Elvira con incertidumbre hace fallar la lectura}\newline \texttt{Elvira\allowbreak{}Scanner.\allowbreak{}get\allowbreak{}Next\allowbreak{}Token} convierte los números de la tabla con \texttt{Double.\allowbreak{}parse\allowbreak{}Double} sin contemplar la barra que separa la incertidumbre ni la marca \texttt{\#} de las casillas deducidas, lo que afecta a los ficheros \texttt{.\allowbreak{}elv} de análisis de sensibilidad. El formato Elvira está congelado por decisión de equipo. \\[3pt]
\textbf{119}\phantomsection\label{nota:119} & \textbf{La pestaña de la tabla del nodo solo enseñaría el nombre de la clave en castellano}\newline Las claves \texttt{Node\allowbreak{}Properties\allowbreak{}Dialog.\allowbreak{}Edit\allowbreak{}Potential\allowbreak{}Tab.\allowbreak{}*} que usa \texttt{Node\allowbreak{}Properties\allowbreak{}Dialog.\allowbreak{}update\allowbreak{}Potential\allowbreak{}Tab\allowbreak{}Title\allowbreak{}And\allowbreak{}Button} están todas en \texttt{Dialogs\_\allowbreak{}en.\allowbreak{}xml}, así que en inglés la pestaña y los botones salen bien. Solo faltan en el fichero en castellano, que no se mantiene porque la aplicación es solo en inglés. \\[3pt]
\textbf{120}\phantomsection\label{nota:120} & \textbf{Varios rótulos y menús llevan su texto inglés escrito en el código}\newline «Numeric value:» en \texttt{Add\allowbreak{}Finding\allowbreak{}Dialog}, los submenús «Add» de \texttt{Network\allowbreak{}Contextual\allowbreak{}Menu} y «Alignment» de \texttt{Node\allowbreak{}Contextual\allowbreak{}Menu} y los rótulos de otros paneles siguen escritos en el código, pero están en inglés y la aplicación es solo en inglés, así que no hay nada visible. Aparte, la ayuda emergente «Create an utility node» sí es una errata visible: está en \hyperref[err:185]{error~185}. \\[3pt]
\textbf{121}\phantomsection\label{nota:121} & \textbf{En castellano, la opción «Discretas y continuas» cuelga de una etiqueta mal escrita}\newline En \texttt{Dialogs\_\allowbreak{}es.\allowbreak{}xml} el elemento que pide \texttt{Network\allowbreak{}Variables\allowbreak{}Panel.\allowbreak{}get\allowbreak{}List\allowbreak{}Of\allowbreak{}Types} está bajo \texttt{\textless{}items\textgreater{}} en minúscula, así que en castellano la clave no existiría. En \texttt{Dialogs\_\allowbreak{}en.\allowbreak{}xml} las tres entradas están bien y la aplicación solo sirve inglés (\texttt{String\allowbreak{}Database.\allowbreak{}SUPPORTED\_\allowbreak{}LANGUAGES}). \\[3pt]
\textbf{122}\phantomsection\label{nota:122} & \textbf{Los mensajes asociados a excepciones solo existen en inglés}\newline Los textos de las clases de excepción viven en ocho ficheros \texttt{*\_\allowbreak{}class\_\allowbreak{}localizations\_\allowbreak{}en.\allowbreak{}xml} y no hay versión en castellano. Como \texttt{String\allowbreak{}Database} solo sirve inglés, por decisión de equipo, no hay consecuencia visible. \\[3pt]
\textbf{123}\phantomsection\label{nota:123} & \textbf{Varios textos de los ficheros en castellano están estropeados o faltan}\newline «Ampliar zoomManager», «Reducir zoomManager», «Proposito» y las claves que faltan están solo en los ficheros \texttt{*\_\allowbreak{}es.\allowbreak{}xml}, como \texttt{Menus\_\allowbreak{}es.\allowbreak{}xml}, que la aplicación no carga porque es solo en inglés. La misma sustitución de \texttt{zoom} por \texttt{zoom\allowbreak{}Manager} sí dejó un texto visible en inglés, que es \hyperref[err:183]{error~183}. \\[3pt]
\textbf{124}\phantomsection\label{nota:124} & \textbf{Los estados del dominio «not happened - happened» no tienen texto en el fichero de textos}\newline Sin texto, \texttt{GUIDefault\allowbreak{}States.\allowbreak{}get\allowbreak{}String} devuelve el nombre interno del estado, que en inglés es justo lo que debe verse, como en el resto de dominios de \texttt{Selectables\_\allowbreak{}en.\allowbreak{}xml}. Solo se notaría en castellano, que queda fuera porque la aplicación es solo en inglés. \\[3pt]
\textbf{125}\phantomsection\label{nota:125} & \textbf{El aviso de cambio de dominio al partir un intervalo del árbol no usa su clave de textos}\newline \texttt{Tree\allowbreak{}ADDEditor\allowbreak{}Panel.\allowbreak{}split\allowbreak{}Interval} muestra «Be careful, you changed variable domain» escrito en el código, y la clave \texttt{Tree\allowbreak{}ADD.\allowbreak{}Domain\allowbreak{}Change\allowbreak{}Warning} de \texttt{Menus\_\allowbreak{}en.\allowbreak{}xml} no la usa nadie. En la aplicación, que es solo en inglés, el aviso se lee bien; la clave sin uso es texto muerto. \\[3pt]
\textbf{126}\phantomsection\label{nota:126} & \textbf{El título de la ventana de editar una relación solo tiene texto en el fichero inglés}\newline \texttt{Potential\allowbreak{}Edit\allowbreak{}Dialog.\allowbreak{}get\allowbreak{}Base\allowbreak{}Title} pide la clave \texttt{Node\allowbreak{}Properties\allowbreak{}Dialog.\allowbreak{}Edit\allowbreak{}Potential\allowbreak{}Tab.\allowbreak{}Edit\allowbreak{}Potential\allowbreak{}Title}, que está en \texttt{Dialogs\_\allowbreak{}en.\allowbreak{}xml} y se resuelve como «Edit probability» (medido). Solo falta en \texttt{Dialogs\_\allowbreak{}es.\allowbreak{}xml}, que no se mantiene porque la aplicación es solo en inglés. \\[3pt]
\textbf{127}\phantomsection\label{nota:127} & \textbf{Las dos entradas del menú Ver no tienen texto en castellano}\newline Las claves \texttt{View.\allowbreak{}Go\allowbreak{}Next\allowbreak{}Tab} y \texttt{View.\allowbreak{}Go\allowbreak{}Previous\allowbreak{}Tab} que usa \texttt{Main\allowbreak{}Menu.\allowbreak{}build\allowbreak{}View\allowbreak{}Menu} están en \texttt{Menus\_\allowbreak{}en.\allowbreak{}xml} y dan «Go to next tab» y «Go to previous tab» (medido). Solo faltan en \texttt{Menus\_\allowbreak{}es.\allowbreak{}xml}, que no se mantiene porque la aplicación es solo en inglés. \\[3pt]
\textbf{128}\phantomsection\label{nota:128} & \textbf{El título de la pestaña del árbol de decisión está escrito en el código}\newline \texttt{Inference\allowbreak{}Handler.\allowbreak{}show\allowbreak{}Decision\allowbreak{}Tree} abre la pestaña con «Decision tree for » más el nombre de la red, escrito en el código y no en un fichero de textos. Se lee bien en inglés, y la aplicación es solo en inglés. \\[3pt]
\textbf{129}\phantomsection\label{nota:129} & \textbf{El menú del árbol de decisión dice en castellano «Expandir N niveles» y expande uno}\newline La entrada \texttt{Expand\allowbreak{}Next} de \texttt{Tree\allowbreak{}Contextual\allowbreak{}Menu} ejecuta \texttt{Decision\allowbreak{}Tree\allowbreak{}Editor.\allowbreak{}inference\allowbreak{}Expand\allowbreak{}Next\allowbreak{}Level}, que expande un nivel, y en inglés dice «Expand next level», que es lo que hace. Solo el texto castellano promete otra cosa, y los ficheros \texttt{*\_\allowbreak{}es.\allowbreak{}xml} no se mantienen. \\[3pt]
\textbf{130}\phantomsection\label{nota:130} & \textbf{Dos claves de opciones temporales y de inferencia cuelgan de otro padre en castellano}\newline \texttt{Temporal\allowbreak{}Evolution\allowbreak{}Dialog.\allowbreak{}get\allowbreak{}Slices\allowbreak{}Panel} pide \texttt{Inference.\allowbreak{}Temporal\allowbreak{}Options} y el menú pide \texttt{Inference.\allowbreak{}Inference\allowbreak{}Options}; en \texttt{Menus\_\allowbreak{}en.\allowbreak{}xml} las dos cuelgan de \texttt{\textless{}Inference\textgreater{}} y dan «Temporal options» e «Inference options» (medido). Solo en \texttt{Menus\_\allowbreak{}es.\allowbreak{}xml} cuelgan de \texttt{\textless{}Edit\textgreater{}}, y ese fichero no se mantiene. \\[3pt]
\textbf{131}\phantomsection\label{nota:131} & \textbf{Los rótulos del panel de opciones de Monte Carlo están escritos en el código}\newline \texttt{Monte\allowbreak{}Carlo\allowbreak{}Options\allowbreak{}Panel} crea sus rótulos con textos ingleses en el código, que se leen bien en la aplicación, que es solo en inglés. La errata «Perfom PSA» está en la casilla de \texttt{get\allowbreak{}Psa\allowbreak{}Check\allowbreak{}Box(\allowbreak{})}, que nadie añade al panel (\hyperref[nota:58]{nota~58}). \\[3pt]
\textbf{132}\phantomsection\label{nota:132} & \textbf{Seis rótulos del análisis probabilístico de coste-efectividad están en el código}\newline \texttt{CEProbabilistic\allowbreak{}Dialog} escribe en el código «WTP for evaluation» (WTP: disposición a pagar), «Display:», «Relative to:» y otros, mientras que \texttt{CEDecision\allowbreak{}Results} pide cuatro de ellos por clave a \texttt{costeffectiveness\_\allowbreak{}en.\allowbreak{}xml}, con el mismo texto. En inglés las dos ventanas dicen lo mismo, y la aplicación es solo en inglés. \\[3pt]
\textbf{133}\phantomsection\label{nota:133} & \textbf{Las claves «Decision» e «ICER» del análisis de sensibilidad faltan en castellano}\newline \texttt{Sensitivity\allowbreak{}Analysis.\allowbreak{}General.\allowbreak{}Decision} y \texttt{Sensitivity\allowbreak{}Analysis.\allowbreak{}General.\allowbreak{}ICER} (ICER: razón coste-efectividad incremental) están en \texttt{sensitivityanalysis\_\allowbreak{}en.\allowbreak{}xml}, y \texttt{CESpider\allowbreak{}Dialog} muestra «Decision:» correctamente. Solo faltan en castellano, y la aplicación es solo en inglés. \\[3pt]
\textbf{134}\phantomsection\label{nota:134} & \textbf{Los rótulos de los resultados de la simulación de sucesos están escritos en el código}\newline \texttt{DESResults\allowbreak{}Window}, \texttt{CEDESDialog} y el monitor «Running simulation» de \texttt{Main\allowbreak{}Panel\allowbreak{}Listener\allowbreak{}Assistant} llevan sus textos en el código, con un aviso \texttt{FIXME\ use\ String\allowbreak{}Database}. Están en inglés y se ven bien; solo importaría con otro idioma, y la aplicación es solo en inglés. \\[3pt]
\textbf{135}\phantomsection\label{nota:135} & \textbf{El formato de bases de casos de Elvira (\texttt{.\allowbreak{}dbc}) no está registrado ni se ofrece}\newline \texttt{Elvira\allowbreak{}Data\allowbreak{}Base\allowbreak{}IO} lleva comentada la anotación \texttt{@Case\allowbreak{}Database\allowbreak{}Format}, así que \texttt{Case\allowbreak{}Database\allowbreak{}Manager} no lo encuentra y ningún selector de ficheros lo ofrece. Es un formato de Elvira, congelado por decisión de equipo; además, las bases de aprendizaje de las pruebas ya se convirtieron a CSV (valores separados por comas) en el commit \texttt{9c36f04}. \\[3pt]
\textbf{136}\phantomsection\label{nota:136} & \textbf{El escritor de bases de casos de Elvira (\texttt{.\allowbreak{}dbc}) no escribe los estados}\newline \texttt{Elvira\allowbreak{}Data\allowbreak{}Base\allowbreak{}IO.\allowbreak{}write\allowbreak{}Variables} no escribe \texttt{num-states} ni \texttt{states}, y \texttt{Elvira\allowbreak{}DBParser} falla al leer lo escrito porque esos estados le llegan nulos. Es código exclusivo del formato Elvira, congelado por decisión de equipo, y además no está registrado (\hyperref[nota:135]{nota~135}). \\[3pt]
\textbf{137}\phantomsection\label{nota:137} & \textbf{El lector de bases de casos de Elvira (\texttt{.\allowbreak{}dbc}) se fía del número de casos de la cabecera}\newline \texttt{Elvira\allowbreak{}DBParser} reserva la tabla de casos con el número que declara la cabecera, así que más casos dan \texttt{Array\allowbreak{}Index\allowbreak{}Out\allowbreak{}Of\allowbreak{}Bounds\allowbreak{}Exception} y menos dejan filas a cero que pasan por casos. Es código exclusivo del formato Elvira, congelado por decisión de equipo y no registrado (\hyperref[nota:135]{nota~135}); ningún otro lector comparte ese mecanismo. \\[3pt]
\textbf{138}\phantomsection\label{nota:138} & \textbf{La lista de algoritmos de aprendizaje toma los nombres de las anotaciones}\newline \texttt{General\allowbreak{}Tab\allowbreak{}Panel} muestra el \texttt{name} de \texttt{@Learning\allowbreak{}Algorithm\allowbreak{}Type} y no usa las claves \texttt{Learning.\allowbreak{}\textless{}Algoritmo\textgreater{}.\allowbreak{}Algorithm} de los ficheros de textos. En inglés esos nombres equivalen a los de \texttt{Learning\_\allowbreak{}en.\allowbreak{}xml}; la diferencia solo se vería con otro idioma, y la aplicación es solo en inglés. \\[3pt]
\textbf{139}\phantomsection\label{nota:139} & \textbf{La casilla de prueba causal del algoritmo PC tiene su rótulo y su ayuda en el código}\newline \texttt{PCParameters\allowbreak{}Dialog} crea la casilla «Use ANM causal direction test (phase 3)» (ANM: modelo de ruido aditivo) y su ayuda con textos escritos en el código. En inglés se leen bien y la aplicación es solo en inglés; que la ventana no nombre «phase 3» en otro sitio es cuestión de redacción. \\[3pt]
\textbf{140}\phantomsection\label{nota:140} & \textbf{Los textos de las herramientas de evaluación de redes y de casos están en el código}\newline Las cuatro entradas del menú Herramientas del módulo \texttt{bn\allowbreak{}Evaluation} y los textos de sus ventanas, como \texttt{set\allowbreak{}Title(\allowbreak{}"Bayesian\ network\ evaluation")}, están escritos en inglés en el código, sin fichero de textos. En inglés se leen bien, y la aplicación es solo en inglés; la única diferencia visible es de estilo, las mayúsculas de «Split Dataset» frente a las otras tres entradas. \\[3pt]
\textbf{141}\phantomsection\label{nota:141} & \textbf{Guardar en \texttt{.\allowbreak{}pgmx} un OR ruidoso o un MAX causal lo reabre como MAX general}\newline \texttt{PGMXWriter\_\allowbreak{}0\_\allowbreak{}2.\allowbreak{}get\allowbreak{}ICIPotential} no escribe el tipo de modelo y \texttt{PGMXPotential\allowbreak{}Parsers.\allowbreak{}get\allowbreak{}ICIPotential} construye siempre el tipo general, aunque la tabla de probabilidad es idéntica antes y después (medido). Solo lo alcanzan modelos importados de Elvira y solo se nota al volver a exportar a Elvira, formato congelado por decisión de equipo; el índice de parámetros por padre y el \texttt{equals} de \texttt{ICIPotential} ya se corrigieron en el commit \texttt{6260a4c}. \\[3pt]
\end{longtable}

\subsection*{Propuestas de diseño y limpieza}\addcontentsline{toc}{subsection}{Propuestas de diseño y limpieza}
\begin{longtable}{@{}>{\raggedright\arraybackslash}p{1.1cm}>{\raggedright\arraybackslash}p{\dimexpr\linewidth-1.1cm-2\tabcolsep\relax}@{}}
\textbf{142}\phantomsection\label{nota:142} & \textbf{Una edición que modifica la red y después falla no se revierte}\newline Propuesta: que \texttt{PNEdit.\allowbreak{}execute\allowbreak{}Edit(\allowbreak{})} revierta ante cualquier fallo de \texttt{do\allowbreak{}Edit(\allowbreak{})} y que \texttt{Multi\allowbreak{}Step\allowbreak{}Edit.\allowbreak{}Step\allowbreak{}Executer.\allowbreak{}execute} apunte cada paso antes de ejecutarlo, en vez de después. No es un error por sí mismo, sino una garantía que falta en el diseño: los casos concretos que hoy llegan al usuario tienen su propia ficha: \hyperref[err:28]{error~28} y \hyperref[nota:184]{nota~184}. \\[3pt]
\textbf{143}\phantomsection\label{nota:143} & \textbf{El comentario de \texttt{No\allowbreak{}Backward\allowbreak{}Link} describe una regla más estricta que la del código}\newline Hay que decidir cuál es la regla: el comentario exige que un nodo temporal con un hijo atemporal esté en el ciclo cero, y \texttt{No\allowbreak{}Backward\allowbreak{}Link.\allowbreak{}allowed\allowbreak{}Link} acepta el enlace en cualquier ciclo. No se ha encontrado un fallo observable, y la comprobación de la red entera y la de \texttt{Add\allowbreak{}Link\allowbreak{}Edit} usan el mismo método, así que no divergen entre sí. \\[3pt]
\textbf{144}\phantomsection\label{nota:144} & \textbf{Cada restricción se da de alta como oyente de ediciones sin escuchar nada}\newline Propuesta: quitar el \texttt{implements\ PNEdit\allowbreak{}Listener} de \texttt{PNConstraint} y su alta en \texttt{Prob\allowbreak{}Net.\allowbreak{}add\allowbreak{}Constraint} (ver \hyperref[nota:158]{nota~158}). No es un error visible: ninguna restricción atiende ningún aviso, y aunque \texttt{set\allowbreak{}Network\allowbreak{}Type} deja dados de alta oyentes de restricciones ya quitadas (medido: dos, tras pasar de red bayesiana a diagrama de influencia, a simulación de eventos discretos y otra vez a bayesiana), su único coste es recorrerlos en cada edición. \\[3pt]
\textbf{145}\phantomsection\label{nota:145} & \textbf{El nombre de \texttt{@Constraint} no lo lee nadie y en seis casos no coincide con la clase}\newline Propuesta: si algún día se guardan las restricciones en el fichero por un nombre estable (\hyperref[nota:89]{nota~89}), usar este atributo corrigiendo antes sus dos erratas, «InitialNodeConstrain» y «OnySelfLoops\ldots». No es un error: \texttt{Constraint\allowbreak{}Manager} solo lee \texttt{default\allowbreak{}Behavior(\allowbreak{})}, así que el \texttt{name(\allowbreak{})} es un dato que nadie usa y que no coincida con la clase no tiene efecto observable. \\[3pt]
\textbf{146}\phantomsection\label{nota:146} & \textbf{\texttt{Only\allowbreak{}Unlabeled\allowbreak{}Links} no impone nada y no llega a ninguna red}\newline Hay que decidir si se implementa o se borra: su \texttt{check\allowbreak{}Prob\allowbreak{}Net} es un cuerpo vacío generado automáticamente y ninguna edición la consulta. No hay comportamiento que corregir, porque es opcional, nadie la añade y nunca está en una red; es una de las ocho restricciones opcionales de \hyperref[err:39]{error~39}. \\[3pt]
\textbf{147}\phantomsection\label{nota:147} & \textbf{\texttt{Valid\allowbreak{}Name} está inactiva, y su comprobación ya la hacen otras dos restricciones}\newline Propuesta de limpieza: borrar \texttt{Valid\allowbreak{}Name} dejando su método estático \texttt{name\allowbreak{}Is\allowbreak{}Already\allowbreak{}Present} donde se use, o unificarla con \texttt{No\allowbreak{}Empty\allowbreak{}Name}, cuyo \texttt{check\allowbreak{}Prob\allowbreak{}Net} es el mismo código. No es un agujero: en \texttt{Node\allowbreak{}Base\allowbreak{}Name\allowbreak{}Edit.\allowbreak{}check\allowbreak{}Constraints\allowbreak{}Will\allowbreak{}Be\allowbreak{}Met} las líneas que consultan \texttt{Valid\allowbreak{}Name} son código muerto, pero justo antes \texttt{Distinct\allowbreak{}Variable\allowbreak{}Names} y \texttt{No\allowbreak{}Empty\allowbreak{}Name} ya rechazan el nombre repetido y el vacío. \\[3pt]
\textbf{148}\phantomsection\label{nota:148} & \textbf{\texttt{Model\allowbreak{}Network\allowbreak{}Constraint} solo trabaja por la vía de las ediciones, a propósito}\newline Propuesta: cuando cada restricción sepa juzgar una edición (\hyperref[nota:153]{nota~153}), esta restricción encajaría allí de forma natural. No es un error: su \texttt{check\allowbreak{}Prob\allowbreak{}Net} está vacío a propósito, porque la regla (qué enlaces de la red modelo no se pueden tocar durante el aprendizaje) solo tiene sentido ante una edición, y la consultan \texttt{Add\allowbreak{}Link\allowbreak{}Edit}, \texttt{Remove\allowbreak{}Link\allowbreak{}Edit}, \texttt{Invert\allowbreak{}Link\allowbreak{}Edit} y \texttt{Orient\allowbreak{}Link\allowbreak{}Edit}. \\[3pt]
\textbf{149}\phantomsection\label{nota:149} & \textbf{Deshacer tiene el cuerpo vacío en la clase base, y olvidarlo en una edición no avisa}\newline Propuesta: hacer obligatorio el contrato de edición (\hyperref[nota:154]{nota~154}), para que olvidar \texttt{undo(\allowbreak{})} o \texttt{check\allowbreak{}Constraints\allowbreak{}Will\allowbreak{}Be\allowbreak{}Met(\allowbreak{})} en una subclase de \texttt{PNEdit} no pase inadvertido. No es un error por sí mismo: los casos graves ya se corrigieron, y las ediciones que hoy no redefinen \texttt{undo(\allowbreak{})} no llegan al usuario, como \texttt{Remove\allowbreak{}Potentials\allowbreak{}Edit} (\hyperref[nota:35]{nota~35}). \\[3pt]
\textbf{150}\phantomsection\label{nota:150} & \textbf{No está decidido dónde toma una edición la foto del estado anterior}\newline Propuesta: con el contrato de \hyperref[nota:154]{nota~154}, tomar la foto siempre dentro de \texttt{do\allowbreak{}Edit(\allowbreak{})}, donde se construye la reversión, porque una foto tomada en el constructor caduca si la edición se ejecuta más tarde. No es un error hoy: las ediciones que fotografían en el constructor (\texttt{Node\allowbreak{}State\allowbreak{}Edit}, \texttt{Table\allowbreak{}Potential\allowbreak{}Value\allowbreak{}Edit}) se ejecutan justo después de construirse, y \texttt{Add\allowbreak{}Link\allowbreak{}Edit}, que la ventana construye sin ejecutar para colorear flechas, fotografía en \texttt{do\allowbreak{}Edit(\allowbreak{})}. \\[3pt]
\textbf{151}\phantomsection\label{nota:151} & \textbf{El protocolo de rehacer depende de un interruptor que cada clase usa a su manera}\newline Propuesta: eliminar el interruptor \texttt{typical\allowbreak{}Redo}, de modo que rehacer vuelva a ejecutar \texttt{do\allowbreak{}Edit(\allowbreak{})} y guarde su nueva reversión (\hyperref[nota:154]{nota~154}); hoy conviven al menos seis variantes. No es un error: las variantes comprobadas repiten operaciones que no cambian el resultado (\hyperref[nota:84]{nota~84}, \hyperref[nota:85]{nota~85} y \hyperref[nota:66]{nota~66}). \\[3pt]
\textbf{152}\phantomsection\label{nota:152} & \textbf{Ediciones que ejecutan otras fuera de la composición, y un \texttt{equals} que cambia el estado}\newline Propuesta: prohibir que una edición ejecute otra desde su \texttt{do\allowbreak{}Edit(\allowbreak{})} sin ser una edición compuesta, porque deja dos entradas en el historial por una sola acción. No es un error visible: el caso concreto (\hyperref[nota:83]{nota~83}) no llega al usuario, y \texttt{COrient\allowbreak{}Links\allowbreak{}Edit.\allowbreak{}equals} genera y marca sus sub-ediciones al compararse, pero son las mismas que marcaría al ejecutarse. \\[3pt]
\textbf{153}\phantomsection\label{nota:153} & \textbf{Cada restricción debería saber juzgar una edición, en vez de repetir su regla fuera}\newline Propuesta: un método \texttt{check\allowbreak{}Edit} en \texttt{PNConstraint}, recorrido por \texttt{execute\allowbreak{}Edit(\allowbreak{})}, al que se migre restricción a restricción la lógica que hoy repiten las ediciones (\texttt{Add\allowbreak{}Link\allowbreak{}Edit} nombra quince clases de restricción) y algunos validadores de la ventana, para que una restricción nueva descubierta como complemento se imponga también sobre las ediciones. Es un rediseño, no un error: los defectos concretos que causa la duplicación tienen sus propias fichas, como \hyperref[err:178]{error~178}, \hyperref[err:38]{error~38}, \hyperref[err:39]{error~39} y \hyperref[nota:79]{nota~79}. \\[3pt]
\textbf{154}\phantomsection\label{nota:154} & \textbf{Contrato de edición obligatorio: \texttt{do\allowbreak{}Edit} devuelve su propia reversión}\newline Propuesta: que \texttt{do\allowbreak{}Edit(\allowbreak{})} devuelva un objeto de reversión, que \texttt{undo(\allowbreak{})} y \texttt{redo(\allowbreak{})} sean finales en \texttt{PNEdit}, que desaparezca \texttt{typical\allowbreak{}Redo}, que las modificaciones de \texttt{Prob\allowbreak{}Net} devuelvan su reversión y que \texttt{execute\allowbreak{}Edit(\allowbreak{})} se niegue a ejecutarse dentro del \texttt{do\allowbreak{}Edit(\allowbreak{})} de una edición no compuesta. Es un rediseño de fondo, sin defecto concreto propio. \\[3pt]
\textbf{155}\phantomsection\label{nota:155} & \textbf{Que \texttt{do\allowbreak{}Edit} devuelva cómo deshacerse y que cada restricción sepa juzgar una edición}\newline Son dos propuestas de rediseño ya descritas en \hyperref[nota:154]{nota~154} (la reversión devuelta por \texttt{do\allowbreak{}Edit(\allowbreak{})}) y en \hyperref[nota:153]{nota~153} (un método \texttt{check\allowbreak{}Edit} en la restricción). No tienen defecto concreto propio. \\[3pt]
\textbf{156}\phantomsection\label{nota:156} & \textbf{Que los cambios de \texttt{Prob\allowbreak{}Net} devuelvan su reversión y comprobar la red al abrir}\newline Propuestas: que cada modificación de \texttt{Prob\allowbreak{}Net} devuelva su reversión, que forma parte de \hyperref[nota:154]{nota~154}, y comprobar la red al terminar de leer un fichero y decidir qué hacer con una red inválida. No tienen defecto propio; el defecto concreto de la segunda es \hyperref[err:37]{error~37}. \\[3pt]
\textbf{157}\phantomsection\label{nota:157} & \textbf{Devolver la regla a la restricción y decidir qué significa una restricción opcional}\newline Devolver a la restricción la regla que hoy repiten las ediciones es la misma propuesta que \hyperref[nota:153]{nota~153}, y decidir qué significa \texttt{OPTIONAL} es la decisión pendiente de \hyperref[err:39]{error~39}, donde está su defecto concreto. No hay defecto propio. \\[3pt]
\textbf{158}\phantomsection\label{nota:158} & \textbf{Una sola estrategia para las ediciones de varios pasos y limpieza de las restricciones}\newline Propuestas: quitar la anulación vacía de \texttt{Multi\allowbreak{}Step\allowbreak{}Edit} y la edición verificadora de \texttt{Paste\allowbreak{}Edit}, apuntar cada paso antes de ejecutarlo y revertir en \texttt{execute\allowbreak{}Edit(\allowbreak{})} (\hyperref[nota:142]{nota~142}); y limpiar en las restricciones el registro inerte de oyentes (\hyperref[nota:144]{nota~144}), el \texttt{name(\allowbreak{})} de la anotación (\hyperref[nota:145]{nota~145}), los tres moldes (\hyperref[nota:88]{nota~88}) y las clases sin uso. Son reorganización y limpieza; lo que tienen de defecto concreto está en \hyperref[err:31]{error~31}, \hyperref[err:28]{error~28} y \hyperref[nota:89]{nota~89}. \\[3pt]
\textbf{159}\phantomsection\label{nota:159} & \textbf{Prohibir que una edición ejecute otra fuera de una edición compuesta}\newline Propuesta: que \texttt{execute\allowbreak{}Edit(\allowbreak{})} se niegue a ejecutarse si se llama desde el \texttt{do\allowbreak{}Edit(\allowbreak{})} de una edición que no sea compuesta, para que una acción no pueda dejar dos entradas en el historial; es la misma propuesta que \hyperref[nota:152]{nota~152}. No hay defecto que llegue al usuario: el caso de \texttt{ICITable\allowbreak{}Potential\allowbreak{}Value\allowbreak{}Edit} (\hyperref[nota:83]{nota~83}) queda en un historial provisional que se descarta. \\[3pt]
\textbf{160}\phantomsection\label{nota:160} & \textbf{La suma y la multiplicación localizan potenciales con \texttt{List.\allowbreak{}remove} e \texttt{index\allowbreak{}Of}}\newline Propuesta de limpieza de rendimiento menor: que \texttt{Table\allowbreak{}Potential\allowbreak{}Arithmetic.\allowbreak{}sum} aparte las constantes y \texttt{multiply} localice el factor con árboles de estrategia guardando su posición mientras recorren, en vez de buscarlos por igualdad. No es un error: \texttt{Potential.\allowbreak{}equals} descarta antes por clase y por lista de variables, así que no se comparan tablas enteras, y dos tablas iguales tienen las mismas variables y posiciones, así que el resultado no cambia. \\[3pt]
\textbf{161}\phantomsection\label{nota:161} & \textbf{Código comentado y comentarios en castellano en las operaciones sobre potenciales}\newline Propuesta de limpieza: borrar los dos métodos comentados de \texttt{Link\allowbreak{}Restriction\allowbreak{}Potential\allowbreak{}Operations} (unas 83 líneas que ni siquiera se podrían descomentar) y las sentencias comentadas de \texttt{Potential\allowbreak{}Operations}, y pasar a inglés los cuatro comentarios en castellano de \texttt{Discrete\allowbreak{}Potential\allowbreak{}Operations}, \texttt{Max\allowbreak{}Out\allowbreak{}Variable}, \texttt{Max\allowbreak{}Potential} y \texttt{Min\allowbreak{}Potential}. No tiene consecuencia observable; en los modelos canónicos no hay hoy código comentado, solo dos javadoc huérfanos y comentarios que cuentan la historia del código. \\[3pt]
\textbf{162}\phantomsection\label{nota:162} & \textbf{Código sin llamadores en operaciones, ediciones, restricciones y potenciales}\newline Propuesta de limpieza: borrar lo que sigue sin llamadores de producción, como las dos \texttt{multiply\allowbreak{}And\allowbreak{}Maximize\allowbreak{}Uniformly}, \texttt{maximize(\allowbreak{}Collection)}, \texttt{matrix\allowbreak{}Potential} (\hyperref[nota:99]{nota~99}), \texttt{Remove\allowbreak{}Potentials\allowbreak{}Edit} (\hyperref[nota:35]{nota~35}), \texttt{PNEdit.\allowbreak{}set\allowbreak{}Prob\allowbreak{}Net}, \texttt{Prob\allowbreak{}Net.\allowbreak{}remove\allowbreak{}Constraints}, \texttt{Prob\allowbreak{}Net.\allowbreak{}remove\allowbreak{}All\allowbreak{}Constraints} y \texttt{No\allowbreak{}Mixed\allowbreak{}Parents.\allowbreak{}parent\allowbreak{}Node\allowbreak{}Is\allowbreak{}Not\allowbreak{}Mixed}. No tiene consecuencia observable; varios de esos métodos ya tienen pruebas propias, que habría que borrar con ellos. \\[3pt]
\textbf{163}\phantomsection\label{nota:163} & \textbf{\texttt{initial\allowbreak{}Position} es un concepto a medias, y «constante» se decide de dos formas}\newline Propuesta: retirar o completar la posición inicial de \texttt{Abstract\allowbreak{}Indexed\allowbreak{}Potential}, que leen la eliminación y la maximización pero no \texttt{multiply} ni \texttt{sum}, y unificar el criterio de potencial constante (tabla de longitud 1 frente a lista de variables vacía). La primera parte no tiene defecto alcanzable, porque ningún potencial llega a tener una posición inicial distinta de cero; la consecuencia concreta de la segunda es \hyperref[err:1]{error~1}. \\[3pt]
\textbf{164}\phantomsection\label{nota:164} & \textbf{Contratos sin escribir en las operaciones sobre potenciales}\newline Propuesta: documentar que \texttt{Table\allowbreak{}Potential.\allowbreak{}table\allowbreak{}Project} puede devolver el propio argumento, que \texttt{Discrete\allowbreak{}Potential\allowbreak{}Operations.\allowbreak{}normalize} modifica el potencial que recibe y que \texttt{add\allowbreak{}Variable} y \texttt{remove\allowbreak{}Variable} modifican el receptor en los árboles pero devuelven uno nuevo en las tablas, y unificar las dos definiciones de «tiene intervenciones». No se encontró un defecto alcanzable: los llamadores que ignoran el resultado de \texttt{add\allowbreak{}Variable} trabajan sobre árboles, y las dos definiciones de «tiene intervenciones» solo discrepan en código que nadie llama. \\[3pt]
\textbf{165}\phantomsection\label{nota:165} & \textbf{Las interfaces de capacidad de los potenciales conviven con los métodos que sustituían}\newline Propuesta: completar la sustitución para la que se crearon \texttt{Projectable}, \texttt{Reorderable}, \texttt{Scalable}, \texttt{Uncertainty\allowbreak{}Carrier}, \texttt{Strategy\allowbreak{}Carrier} y \texttt{CEUtility\allowbreak{}Potential}, en lugar de los métodos de \texttt{Potential} que devuelven nulo o lanzan «no soportado». No hay error observable: en producción nadie pregunta por las cuatro primeras y las operaciones siguen preguntando por la clase concreta; el riesgo es que quien lea las interfaces crea que el sistema de tipos protege lo que no protege. \\[3pt]
\textbf{166}\phantomsection\label{nota:166} & \textbf{Decidir si el aprendizaje de parámetros debe aprender modelos canónicos}\newline Hay que decidir si la rama de \texttt{EMAlgorithm.\allowbreak{}adapt\allowbreak{}Network} que aprende modelos canónicos se mantiene como capacidad documentada o se retira; con muchos padres no escala, porque la función que combina crece exponencialmente. No es un error: la rama funciona, sus fallos anteriores ya se corrigieron y la cubre \texttt{EMWith\allowbreak{}Canonical\allowbreak{}Models\allowbreak{}Test}. \\[3pt]
\textbf{167}\phantomsection\label{nota:167} & \textbf{Un único ciclo de vida para las tareas de inferencia, con aviso de progreso y cancelación}\newline Propuesta: documentar en \texttt{Task} un solo ciclo de vida (una consulta por objeto, o invalidar lo calculado cuando cambian la evidencia o las variables de interés) y añadir en \texttt{core} una interfaz pequeña de progreso y cancelación, que hoy se informa de cuatro formas distintas, entre ellas un \texttt{Progress\allowbreak{}Monitor} obligatorio en \texttt{DESInference} que impide simular sin pantalla. Es una propuesta de contrato: los defectos concretos que cuelgan de ella son \hyperref[nota:105]{nota~105} y \hyperref[nota:39]{nota~39}. \\[3pt]
\textbf{168}\phantomsection\label{nota:168} & \textbf{Subir el preprocesamiento de la inferencia al marco común y dejar una sola proyección}\newline Propuesta: pasar el preprocesamiento (expansión temporal, políticas, evidencia, descuentos, discretización y absorción de nodos numéricos) de \texttt{Variable\allowbreak{}Elimination} a \texttt{Inference\allowbreak{}Algorithm}, y dejar una sola de las tres copias de la proyección a red de Markov (\texttt{Task\allowbreak{}Utilities}, \texttt{Cluster\allowbreak{}Propagation}, \texttt{Hybrid\allowbreak{}Elimination}). Es una reorganización: la diferencia de resultados entre algoritmos solo aparece en redes bayesianas con variables numéricas (medido con una red hecha a mano, \hyperref[nota:109]{nota~109}), y no se encontró un camino real de la aplicación donde se note. \\[3pt]
\textbf{169}\phantomsection\label{nota:169} & \textbf{Declarar las versiones de terceros en un solo sitio y que cada módulo declare lo que usa}\newline Propuesta: mover las versiones de las bibliotecas de terceros a \texttt{\textless{}dependency\allowbreak{}Management\textgreater{}} del \texttt{pom.\allowbreak{}xml} raíz (el fichero que configura la construcción con Maven), que cada módulo declare lo que usa y verificarlo con \texttt{mvn\ dependency:analyze}. Hoy todas se declaran en \texttt{\textless{}dependencies\textgreater{}} del raíz y las heredan todos los módulos, incluido \texttt{core}; es una reorganización sin consecuencia visible en la aplicación. \\[3pt]
\textbf{170}\phantomsection\label{nota:170} & \textbf{Podar los \texttt{requires} que no se usan y quitar a \texttt{core} la dependencia del escritorio}\newline Propuesta: quitar de los \texttt{module-info.\allowbreak{}java} las dependencias sin uso (en \texttt{core}, \texttt{org.\allowbreak{}apache.\allowbreak{}poi.\allowbreak{}poi} y \texttt{java.\allowbreak{}desktop}; en \texttt{inference}, \texttt{java.\allowbreak{}desktop} y \texttt{jeval}; en \texttt{gui}, \texttt{jdk.\allowbreak{}compiler}), borrar los restos de escritorio de \texttt{core}, como el import muerto de \texttt{JOption\allowbreak{}Pane} en \texttt{Purpose\allowbreak{}Edit} y la \texttt{Cannot\allowbreak{}Undo\allowbreak{}Exception} de \texttt{Monte\allowbreak{}Carlo\allowbreak{}Options\allowbreak{}Edit.\allowbreak{}undo(\allowbreak{})}, y retirar entonces \texttt{requires\ java.\allowbreak{}desktop} de \texttt{core}. Es limpieza sin consecuencia visible; el único resto de escritorio que sí se nota, el diálogo de \texttt{Tree\allowbreak{}With\allowbreak{}Excluded\allowbreak{}Events\allowbreak{}Potential}, es \hyperref[err:105]{error~105}. \\[3pt]
\textbf{171}\phantomsection\label{nota:171} & \textbf{Una sola política para instanciar complementos y descubrirlos en todo el camino de clases}\newline Propuesta: un ayudante común con la política de \texttt{Inference\allowbreak{}Manager.\allowbreak{}build} (aislar, registrar y seguir), avisos ante nombres repetidos y clases descartadas, descubrimiento por anotación en todo el camino de clases y no solo bajo \texttt{org.\allowbreak{}openmarkov} (\texttt{Plugin\allowbreak{}Loader}), y que el algoritmo preferido lo declare el complemento en vez de nombrarlo \texttt{core} por texto. Es una unificación sin consecuencia hoy: por ejemplo, el nombre del algoritmo preferido coincide con la anotación de \texttt{VEPropagation}; el gestor que se cae entero por un complemento defectuoso es \hyperref[nota:106]{nota~106}. \\[3pt]
\textbf{172}\phantomsection\label{nota:172} & \textbf{Borrar la API de complementos sin uso y reactivar la prueba de requisitos de construcción}\newline Propuesta: borrar \texttt{Filter} y \texttt{Plugin\allowbreak{}Manager}, que nadie importa; reactivar \texttt{Implementation\allowbreak{}Requirements\allowbreak{}Are\allowbreak{}Met\allowbreak{}Test}, hoy con \texttt{@Disabled}, congelando la lista actual de incumplimientos para que no crezca; y decidir el destino del \texttt{name(\allowbreak{})} de \texttt{@Constraint} (\hyperref[nota:145]{nota~145}). Es limpieza y red de seguridad de pruebas, sin consecuencia visible; hoy nada comprueba que una subclase nueva de \texttt{Potential} tenga el constructor que necesita su copia por reflexión. \\[3pt]
\textbf{173}\phantomsection\label{nota:173} & \textbf{Abrir el punto de extensión de los formatos de fichero}\newline Propuesta: que cada complemento de formato declare su extensión y versión (o sepa decir si puede leer un fichero), mover a \texttt{core} el contrato de lectores y escritores, que hoy vive en \texttt{io}, y añadir lectura y escritura sobre flujos. Es un rediseño: hoy \texttt{Format\allowbreak{}Manager.\allowbreak{}get\allowbreak{}Prob\allowbreak{}Net\allowbreak{}Reader} elige el lector solo por extensión y supone un XML con \texttt{format\allowbreak{}Version}, así que un formato nuevo que no sea XML no se podría abrir; el lector de XMLBIF es \hyperref[nota:50]{nota~50}. \\[3pt]
\textbf{174}\phantomsection\label{nota:174} & \textbf{Que cada excepción construya su mensaje y avisar ante claves de texto repetidas}\newline Propuesta: que cada excepción construya su mensaje en inglés en vez de buscarlo por el nombre de su clase en los ficheros de textos (hoy renombrarla rompe el mensaje sin aviso del compilador), sacar de \texttt{Bundle\allowbreak{}Search}, en \texttt{core}, los nombres de paquetes de \texttt{sensitivity\allowbreak{}Analysis} y \texttt{stochastic\allowbreak{}Propagation\allowbreak{}Output}, usar un registrador por clase y avisar cuando \texttt{String\allowbreak{}Database} funda claves repetidas. Es reorganización del diagnóstico: medido, las 172 subclases concretas de \texttt{Open\allowbreak{}Markov\allowbreak{}Exception} tienen su texto, y las claves repetidas no tienen hoy consecuencia en inglés (\hyperref[nota:63]{nota~63}). \\[3pt]
\textbf{175}\phantomsection\label{nota:175} & \textbf{Renombrar el módulo \texttt{inference.\allowbreak{}DES} y separar el motor de simulación de sus ventanas}\newline Propuesta: renombrar el módulo a la convención \texttt{org.\allowbreak{}openmarkov.\allowbreak{}*} y partirlo, con el motor \texttt{DESInference} abajo y sin escritorio, y \texttt{CEDESDialog}, \texttt{DESResults\allowbreak{}Window} y la exportación a Excel en un módulo encima de \texttt{gui}. Es una reorganización sin consecuencia visible; hoy el \texttt{Progress\allowbreak{}Monitor} obligatorio del motor impide simular sin pantalla, y sus ventanas no pueden usar componentes de \texttt{gui} sin crear un ciclo. \\[3pt]
\textbf{176}\phantomsection\label{nota:176} & \textbf{Que la ventana pida los algoritmos al registro de complementos y no los construya a mano}\newline Propuesta: una fábrica de tareas que resuelva la implementación en el registro, en lugar de los \texttt{new\ VEPropagation}, \texttt{new\ VEEvaluation}, \texttt{new\ VECEAnalysis} y demás que hacen la ventana y los módulos de análisis. No cambia ningún resultado hoy, porque cada tarea de la ventana tiene un solo algoritmo aplicable; lo que sí pasa es que registrar un algoritmo nuevo como complemento no cambia lo que ejecuta la aplicación. \\[3pt]
\textbf{177}\phantomsection\label{nota:177} & \textbf{Extraer de la ventana un documento sin Swing con la red, su fichero y su estado}\newline Propuesta: un \texttt{Network\allowbreak{}Document} sin Swing, que el panel observe y que \texttt{Nets\allowbreak{}IO} reciba, con la red, la ruta del fichero, el lector o escritor con que se abrió, el indicador de modificado y la evidencia, que hoy se reparten \texttt{Network\allowbreak{}Editor\allowbreak{}Panel} y \texttt{PNESupport}. Es una reorganización de responsabilidades sin consecuencia visible propia; el defecto que causa que las copias de red hereden ese estado es \hyperref[err:35]{error~35}. \\[3pt]
\textbf{178}\phantomsection\label{nota:178} & \textbf{Mover a \texttt{core} las ediciones del módulo gráfico que no dependen de Swing}\newline Propuesta: mover a \texttt{core/\allowbreak{}action} las ediciones de \texttt{gui/\allowbreak{}src/\allowbreak{}main/\allowbreak{}java/\allowbreak{}org/\allowbreak{}openmarkov/\allowbreak{}gui/\allowbreak{}action/\allowbreak{}} que no tocan clases visuales, como \texttt{Change\allowbreak{}Node\allowbreak{}Type\allowbreak{}Edit}, \texttt{Impose\allowbreak{}Policy\allowbreak{}Edit} o \texttt{Table\allowbreak{}Potential\allowbreak{}Value\allowbreak{}Edit}, junto con la redistribución de probabilidades de \texttt{Potentials\allowbreak{}Table\allowbreak{}Panel\allowbreak{}Operations}, y adaptar las que dependen de \texttt{Visual\allowbreak{}Node} o \texttt{Visual\allowbreak{}Network}. Es un traslado sin consecuencia visible; hoy quien quiera usarlas sin interfaz depende del módulo Swing, como hace \texttt{learning.\allowbreak{}gui} para \texttt{Auto\allowbreak{}Arrange\allowbreak{}Edit}. \\[3pt]
\textbf{179}\phantomsection\label{nota:179} & \textbf{Desmontar el objeto global \texttt{Main\allowbreak{}GUI.\allowbreak{}INSTANCE} pasando los colaboradores por constructor}\newline Propuesta: desmontar \texttt{Main\allowbreak{}GUI.\allowbreak{}INSTANCE} por capas, pasando los colaboradores por constructor y empezando por los cuatro usos del módulo \texttt{learning.\allowbreak{}gui}. Es una reorganización sin consecuencia visible propia; hoy se llega a ese objeto global desde 59 sitios, a menudo con cadenas largas de llamadas. \\[3pt]
\textbf{180}\phantomsection\label{nota:180} & \textbf{Que la red sea la dueña de renombrar variables y cambiar el tipo de los nodos}\newline Propuesta: que esas modificaciones pasen por \texttt{Prob\allowbreak{}Net} (\texttt{rename\allowbreak{}Variable}, \texttt{change\allowbreak{}Node\allowbreak{}Type}) en vez de por \texttt{Node.\allowbreak{}set\allowbreak{}Variable} y \texttt{Node.\allowbreak{}set\allowbreak{}Node\allowbreak{}Type}, que tocan por su cuenta la red, y que \texttt{Potential.\allowbreak{}deep\allowbreak{}Copy} falle cuando la red destino no tiene una variable. Es una reorganización; los defectos concretos tienen sus fichas: \hyperref[err:166]{error~166} para el cambio de tipo de un nodo y \hyperref[err:3]{error~3} para \texttt{deep\allowbreak{}Copy} con nulos. \\[3pt]
\textbf{181}\phantomsection\label{nota:181} & \textbf{Decidir qué es el módulo \texttt{resttemplate}, hoy un ejemplo que responde «Hello, World!»}\newline Hay que decidir el destino de \texttt{resttemplate}, un envoltorio REST (interfaz de programación por HTTP) con Spring Boot cuyo único controlador ofrece \texttt{GET\ /\allowbreak{}hello} y dos \texttt{POST} de juguete sobre \texttt{Example\allowbreak{}Data}. No expone ninguna función de OpenMarkov, así que no hay nada que falle; si se decide que vaya en serio, necesita un \texttt{core} sin escritorio (\hyperref[nota:170]{nota~170}) y leer redes desde un flujo (\hyperref[nota:173]{nota~173}). \\[3pt]
\textbf{182}\phantomsection\label{nota:182} & \textbf{Congelar la biblioteca interna \texttt{org.\allowbreak{}openmarkov.\allowbreak{}java} y revisar qué paquetes exporta \texttt{core}}\newline Propuesta: dejar escrita la regla de mirar la biblioteca estándar de Java y Apache Commons antes de añadir utilidades a \texttt{org.\allowbreak{}openmarkov.\allowbreak{}java} (18 ficheros, con duplicados como \texttt{Null\allowbreak{}Utils.\allowbreak{}equals}, que hace lo mismo que \texttt{Objects.\allowbreak{}equals}), y estudiar cuáles de los 48 paquetes que exporta el \texttt{module-info.\allowbreak{}java} de \texttt{core} son de verdad interfaz pública. Es una norma de mantenimiento sin consecuencia visible. \\[3pt]
\textbf{183}\phantomsection\label{nota:183} & \textbf{Dejar escrito que el sistema de módulos de Java se usa solo para declarar dependencias}\newline Propuesta: documentar que los descriptores son \texttt{open\ module} porque el descubrimiento de complementos por reflexión lo exige, y que la aplicación corre en el camino de clases, así que en ejecución no hay encapsulación; si algún día se quiere, el camino sería \texttt{java.\allowbreak{}util.\allowbreak{}Service\allowbreak{}Loader}. Es una nota de documentación, no trabajo sobre el código. \\[3pt]
\end{longtable}


\section{Procedencia de cada ficha}\label{app:origen}

Cada error y cada nota salen de uno o de varios de estos dos documentos internos: el
\emph{plan de cambios} (con los identificadores de sus puntos y de sus hallazgos) y el
\emph{registro de fallos} (con el número de su entrada). Cuando dos entradas describían el mismo
defecto se han juntado en una sola ficha. Los \emph{hallazgos nuevos} no estaban en ninguno de los
dos: se vieron al comprobar otro error de la lista, que se indica, y se comprobaron después con los
mismos criterios.

{\small
\begin{longtable}{@{}p{1.8cm}p{\dimexpr\linewidth-1.8cm-2\tabcolsep\relax}@{}}
\toprule
\textbf{Ficha} & \textbf{Origen} \\
\midrule
\endhead
\hyperref[err:1]{Error 1} & registro de fallos, nº 49 \\
\hyperref[err:2]{Error 2} & registro de fallos, nº 54 \\
\hyperref[err:3]{Error 3} & registro de fallos, nº 180 \\
\hyperref[err:4]{Error 4} & plan de cambios, RP15; plan de cambios, F2-d; registro de fallos, nº 74 \\
\hyperref[err:5]{Error 5} & plan de cambios, RP20 \\
\hyperref[err:6]{Error 6} & plan de cambios, RP29 \\
\hyperref[err:7]{Error 7} & plan de cambios, RP25; plan de cambios, F8-d \\
\hyperref[err:8]{Error 8} & registro de fallos, nº 39 \\
\hyperref[err:9]{Error 9} & hallazgo nuevo, visto junto al \hyperref[err:8]{error~8} \\
\hyperref[err:10]{Error 10} & registro de fallos, nº 71 \\
\hyperref[err:11]{Error 11} & hallazgo nuevo, visto junto al \hyperref[err:12]{error~12} \\
\hyperref[err:12]{Error 12} & registro de fallos, nº 69 \\
\hyperref[err:13]{Error 13} & registro de fallos, nº 70 \\
\hyperref[err:14]{Error 14} & registro de fallos, nº 73 \\
\hyperref[err:15]{Error 15} & hallazgo nuevo, visto junto al \hyperref[err:10]{error~10}; incluye registro de fallos, nº 32, que se había dado por inalcanzable \\
\hyperref[err:16]{Error 16} & registro de fallos, nº 75 \\
\hyperref[err:17]{Error 17} & registro de fallos, nº 150 \\
\hyperref[err:18]{Error 18} & registro de fallos, nº 152 \\
\hyperref[err:19]{Error 19} & registro de fallos, nº 96 \\
\hyperref[err:20]{Error 20} & registro de fallos, nº 95 \\
\hyperref[err:21]{Error 21} & hallazgo nuevo, visto junto al \hyperref[err:20]{error~20} \\
\hyperref[err:22]{Error 22} & registro de fallos, nº 176 \\
\hyperref[err:23]{Error 23} & registro de fallos, nº 177 \\
\hyperref[err:24]{Error 24} & registro de fallos, nº 145 \\
\hyperref[err:25]{Error 25} & hallazgo nuevo, visto junto a la \hyperref[nota:23]{nota~23} \\
\hyperref[err:26]{Error 26} & registro de fallos, nº 37 \\
\hyperref[err:27]{Error 27} & hallazgo nuevo, visto junto a la \hyperref[nota:23]{nota~23} \\
\hyperref[err:28]{Error 28} & plan de cambios, M11; plan de cambios, D4 \\
\hyperref[err:29]{Error 29} & registro de fallos, nº 29 \\
\hyperref[err:30]{Error 30} & hallazgo nuevo, visto junto al \hyperref[err:29]{error~29} \\
\hyperref[err:31]{Error 31} & plan de cambios, D3 \\
\hyperref[nota:184]{Nota 184} (antes, error~32) & hallazgo nuevo, visto junto a la \hyperref[nota:82]{nota~82} \\
\hyperref[err:33]{Error 33} & registro de fallos, nº 35 \\
\hyperref[err:34]{Error 34} & registro de fallos, nº 166 \\
\hyperref[err:35]{Error 35} & plan de cambios, F5-a \\
\hyperref[err:36]{Error 36} & registro de fallos, nº 112 \\
\hyperref[err:37]{Error 37} & plan de cambios, R1; plan de cambios, F3-b \\
\hyperref[err:38]{Error 38} & registro de fallos, nº 28; plan de cambios, C2; plan de cambios, R2 \\
\hyperref[err:39]{Error 39} & plan de cambios, R3; plan de cambios, F3-c \\
\hyperref[err:40]{Error 40} & registro de fallos, nº 40 \\
\hyperref[err:41]{Error 41} & registro de fallos, nº 120 \\
\hyperref[err:42]{Error 42} & hallazgo nuevo, visto junto al \hyperref[err:44]{error~44} \\
\hyperref[err:43]{Error 43} & registro de fallos, nº 29 \\
\hyperref[err:44]{Error 44} & registro de fallos, nº 41 \\
\hyperref[err:45]{Error 45} & hallazgo nuevo, visto junto al \hyperref[err:44]{error~44} \\
\hyperref[err:46]{Error 46} & registro de fallos, nº 121 \\
\hyperref[err:47]{Error 47} & registro de fallos, nº 123 \\
\hyperref[err:48]{Error 48} & plan de cambios, E1; plan de cambios, F7-g \\
\hyperref[err:49]{Error 49} & plan de cambios, F5-a \\
\hyperref[err:50]{Error 50} & registro de fallos, nº 180 \\
\hyperref[err:51]{Error 51} & registro de fallos, nº 160 \\
\hyperref[err:52]{Error 52} & registro de fallos, nº 184; registro de fallos, nº 42 \\
\hyperref[err:53]{Error 53} & registro de fallos, nº 43 \\
\hyperref[err:54]{Error 54} & hallazgo nuevo, visto junto a la \hyperref[nota:46]{nota~46} \\
\hyperref[err:55]{Error 55} & registro de fallos, nº 185 \\
\hyperref[err:56]{Error 56} & registro de fallos, nº 33 \\
\hyperref[err:57]{Error 57} & registro de fallos, nº 34 \\
\hyperref[err:58]{Error 58} & registro de fallos, nº 36 \\
\hyperref[err:59]{Error 59} & plan de cambios, RP24 \\
\hyperref[err:60]{Error 60} & registro de fallos, nº 167 \\
\hyperref[err:61]{Error 61} & registro de fallos, nº 169 \\
\hyperref[err:62]{Error 62} & registro de fallos, nº 168 \\
\hyperref[err:63]{Error 63} & registro de fallos, nº 48 \\
\hyperref[err:64]{Error 64} & registro de fallos, nº 47 \\
\hyperref[err:65]{Error 65} & registro de fallos, nº 147 \\
\hyperref[err:66]{Error 66} & registro de fallos, nº 182 \\
\hyperref[err:67]{Error 67} & registro de fallos, nº 46; plan de cambios, E5; plan de cambios, F1-s \\
\hyperref[err:68]{Error 68} & registro de fallos, nº 45; registro de fallos, nº 189 \\
\hyperref[err:69]{Error 69} & registro de fallos, nº 188 \\
\hyperref[err:70]{Error 70} & registro de fallos, nº 72 \\
\hyperref[err:71]{Error 71} & hallazgo nuevo, visto junto al \hyperref[err:10]{error~10} \\
\hyperref[err:72]{Error 72} & registro de fallos, nº 38 \\
\hyperref[err:73]{Error 73} & registro de fallos, nº 134 \\
\hyperref[err:74]{Error 74} & hallazgo nuevo, visto junto al \hyperref[err:80]{error~80} \\
\hyperref[err:75]{Error 75} & registro de fallos, nº 65 \\
\hyperref[err:76]{Error 76} & registro de fallos, nº 62 \\
\hyperref[err:77]{Error 77} & registro de fallos, nº 63 \\
\hyperref[err:78]{Error 78} & registro de fallos, nº 137 \\
\hyperref[err:79]{Error 79} & registro de fallos, nº 136 \\
\hyperref[err:80]{Error 80} & registro de fallos, nº 197 \\
\hyperref[err:81]{Error 81} & registro de fallos, nº 198 \\
\hyperref[err:82]{Error 82} & registro de fallos, nº 199 \\
\hyperref[err:83]{Error 83} & hallazgo nuevo, visto junto al \hyperref[err:80]{error~80} \\
\hyperref[err:84]{Error 84} & registro de fallos, nº 206 \\
\hyperref[err:85]{Error 85} & hallazgo nuevo, visto junto a la \hyperref[nota:46]{nota~46} \\
\hyperref[err:86]{Error 86} & hallazgo nuevo, visto junto a la \hyperref[nota:46]{nota~46}; incluye registro de fallos, nº 205, que se había dado por inalcanzable \\
\hyperref[err:87]{Error 87} & hallazgo nuevo, visto junto a la \hyperref[nota:46]{nota~46} \\
\hyperref[err:88]{Error 88} & registro de fallos, nº 94 \\
\hyperref[err:89]{Error 89} & hallazgo nuevo, visto junto al \hyperref[err:88]{error~88} \\
\hyperref[err:90]{Error 90} & hallazgo nuevo, visto junto al \hyperref[err:88]{error~88} \\
\hyperref[err:91]{Error 91} & registro de fallos, nº 179 \\
\hyperref[err:92]{Error 92} & registro de fallos, nº 195 \\
\hyperref[err:93]{Error 93} & registro de fallos, nº 196 \\
\hyperref[err:94]{Error 94} & registro de fallos, nº 207 \\
\hyperref[err:95]{Error 95} & registro de fallos, nº 102 \\
\hyperref[err:96]{Error 96} & registro de fallos, nº 104 \\
\hyperref[err:97]{Error 97} & registro de fallos, nº 105 \\
\hyperref[err:98]{Error 98} & registro de fallos, nº 106 \\
\hyperref[err:99]{Error 99} & hallazgo nuevo, visto junto al \hyperref[err:98]{error~98} \\
\hyperref[err:100]{Error 100} & registro de fallos, nº 174 \\
\hyperref[err:101]{Error 101} & registro de fallos, nº 173 \\
\hyperref[err:102]{Error 102} & registro de fallos, nº 210 \\
\hyperref[err:103]{Error 103} & registro de fallos, nº 213 \\
\hyperref[err:104]{Error 104} & registro de fallos, nº 107 \\
\hyperref[err:105]{Error 105} & registro de fallos, nº 175; plan de cambios, F6-b \\
\hyperref[err:106]{Error 106} & hallazgo nuevo, visto junto a la \hyperref[nota:55]{nota~55} \\
\hyperref[err:107]{Error 107} & registro de fallos, nº 193 \\
\hyperref[err:108]{Error 108} & hallazgo nuevo, visto junto a la \hyperref[nota:168]{nota~168} \\
\hyperref[err:109]{Error 109} & registro de fallos, nº 192 \\
\hyperref[err:110]{Error 110} & registro de fallos, nº 240 \\
\hyperref[err:111]{Error 111} & plan de cambios, F6-e \\
\hyperref[err:112]{Error 112} & hallazgo nuevo, visto junto al \hyperref[err:8]{error~8} \\
\hyperref[err:113]{Error 113} & hallazgo nuevo, visto junto al \hyperref[err:8]{error~8} \\
\hyperref[err:114]{Error 114} & hallazgo nuevo, visto junto al \hyperref[err:8]{error~8} \\
\hyperref[err:115]{Error 115} & registro de fallos, nº 100 \\
\hyperref[err:116]{Error 116} & hallazgo nuevo, visto junto al \hyperref[err:22]{error~22} \\
\hyperref[err:117]{Error 117} & registro de fallos, nº 214 \\
\hyperref[err:118]{Error 118} & registro de fallos, nº 224 \\
\hyperref[err:119]{Error 119} & registro de fallos, nº 225 \\
\hyperref[err:120]{Error 120} & registro de fallos, nº 226 \\
\hyperref[err:121]{Error 121} & registro de fallos, nº 227 \\
\hyperref[err:122]{Error 122} & registro de fallos, nº 228 \\
\hyperref[err:123]{Error 123} & registro de fallos, nº 228 \\
\hyperref[err:124]{Error 124} & registro de fallos, nº 68 \\
\hyperref[err:125]{Error 125} & registro de fallos, nº 85 \\
\hyperref[err:126]{Error 126} & hallazgo nuevo, visto junto al \hyperref[err:125]{error~125} \\
\hyperref[err:127]{Error 127} & registro de fallos, nº 86 \\
\hyperref[err:128]{Error 128} & hallazgo nuevo, visto junto al \hyperref[err:127]{error~127} \\
\hyperref[err:129]{Error 129} & hallazgo nuevo, visto junto al \hyperref[err:127]{error~127} \\
\hyperref[err:130]{Error 130} & hallazgo nuevo, visto junto al \hyperref[err:127]{error~127} \\
\hyperref[err:131]{Error 131} & registro de fallos, nº 84 \\
\hyperref[err:132]{Error 132} & registro de fallos, nº 88 \\
\hyperref[err:133]{Error 133} & hallazgo nuevo, visto junto al \hyperref[err:127]{error~127} \\
\hyperref[err:134]{Error 134} & hallazgo nuevo, visto junto al \hyperref[err:127]{error~127} \\
\hyperref[err:135]{Error 135} & registro de fallos, nº 87 \\
\hyperref[err:136]{Error 136} & registro de fallos, nº 76 \\
\hyperref[err:137]{Error 137} & registro de fallos, nº 83 \\
\hyperref[err:138]{Error 138} & hallazgo nuevo, visto junto al \hyperref[err:139]{error~139} \\
\hyperref[err:139]{Error 139} & registro de fallos, nº 81 \\
\hyperref[err:140]{Error 140} & hallazgo nuevo, visto junto al \hyperref[err:139]{error~139} \\
\hyperref[err:141]{Error 141} & hallazgo nuevo, visto junto al \hyperref[err:139]{error~139} \\
\hyperref[err:142]{Error 142} & hallazgo nuevo, visto junto al \hyperref[err:139]{error~139} \\
\hyperref[err:143]{Error 143} & registro de fallos, nº 89 \\
\hyperref[err:144]{Error 144} & registro de fallos, nº 90 \\
\hyperref[err:145]{Error 145} & registro de fallos, nº 93; registro de fallos, nº 239 \\
\hyperref[err:146]{Error 146} & hallazgo nuevo, visto junto a la \hyperref[nota:20]{nota~20} \\
\hyperref[err:147]{Error 147} & plan de cambios, E8 \\
\hyperref[err:148]{Error 148} & hallazgo nuevo, visto junto al \hyperref[err:147]{error~147} \\
\hyperref[err:149]{Error 149} & hallazgo nuevo, visto junto al \hyperref[err:147]{error~147} \\
\hyperref[err:150]{Error 150} & registro de fallos, nº 204 \\
\hyperref[err:151]{Error 151} & registro de fallos, nº 201 \\
\hyperref[err:152]{Error 152} & registro de fallos, nº 174 \\
\hyperref[err:153]{Error 153} & registro de fallos, nº 126 \\
\hyperref[err:154]{Error 154} & registro de fallos, nº 119 \\
\hyperref[err:155]{Error 155} & hallazgo nuevo, visto junto a la \hyperref[nota:23]{nota~23} \\
\hyperref[err:156]{Error 156} & registro de fallos, nº 118 \\
\hyperref[err:157]{Error 157} & registro de fallos, nº 202 \\
\hyperref[err:158]{Error 158} & registro de fallos, nº 200 \\
\hyperref[err:159]{Error 159} & registro de fallos, nº 151 \\
\hyperref[err:160]{Error 160} & registro de fallos, nº 153 \\
\hyperref[err:161]{Error 161} & hallazgo nuevo, visto junto a la \hyperref[nota:17]{nota~17} \\
\hyperref[err:162]{Error 162} & registro de fallos, nº 162 \\
\hyperref[err:163]{Error 163} & registro de fallos, nº 138 \\
\hyperref[err:164]{Error 164} & registro de fallos, nº 172 \\
\hyperref[err:165]{Error 165} & registro de fallos, nº 203 \\
\hyperref[err:166]{Error 166} & plan de cambios, F7-f \\
\hyperref[err:167]{Error 167} & registro de fallos, nº 115 \\
\hyperref[err:168]{Error 168} & registro de fallos, nº 122 \\
\hyperref[err:169]{Error 169} & registro de fallos, nº 131 \\
\hyperref[err:170]{Error 170} & registro de fallos, nº 128 \\
\hyperref[err:171]{Error 171} & registro de fallos, nº 232 \\
\hyperref[err:172]{Error 172} & registro de fallos, nº 241 \\
\hyperref[err:173]{Error 173} & registro de fallos, nº 142 \\
\hyperref[err:174]{Error 174} & hallazgo nuevo, visto junto al \hyperref[err:173]{error~173} \\
\hyperref[err:175]{Error 175} & registro de fallos, nº 171 \\
\hyperref[err:176]{Error 176} & hallazgo nuevo, visto junto al \hyperref[err:12]{error~12} \\
\hyperref[err:177]{Error 177} & registro de fallos, nº 178 \\
\hyperref[err:178]{Error 178} & plan de cambios, M1 \\
\hyperref[err:179]{Error 179} & registro de fallos, nº 125 \\
\hyperref[err:180]{Error 180} & registro de fallos, nº 144 \\
\hyperref[err:181]{Error 181} & registro de fallos, nº 155 \\
\hyperref[err:182]{Error 182} & hallazgo nuevo, visto junto a la \hyperref[nota:17]{nota~17} \\
\hyperref[err:183]{Error 183} & registro de fallos, nº 139 \\
\hyperref[err:184]{Error 184} & registro de fallos, nº 172 \\
\hyperref[err:185]{Error 185} & hallazgo nuevo, visto junto a la \hyperref[nota:120]{nota~120} \\
\hyperref[err:186]{Error 186} & hallazgo nuevo, visto junto al \hyperref[err:48]{error~48} \\
\hyperref[err:187]{Error 187} & hallazgo nuevo, visto junto al \hyperref[err:48]{error~48} \\
\hyperref[err:188]{Error 188} & hallazgo nuevo, visto junto al \hyperref[err:127]{error~127} \\
\hyperref[err:189]{Error 189} & plan de cambios, F7-b \\
\hyperref[err:190]{Error 190} & plan de cambios, F2-b \\
\hyperref[err:191]{Error 191} & plan de cambios, F2-a \\
\hyperref[err:192]{Error 192} & plan de cambios, RP31 \\
\hyperref[err:193]{Error 193} & hallazgo nuevo, visto junto al \hyperref[err:1]{error~1} \\
\hyperref[nota:1]{Nota 1} & registro de fallos, nº 31; plan de cambios, RP11 \\
\hyperref[nota:2]{Nota 2} & registro de fallos, nº 127 \\
\hyperref[nota:3]{Nota 3} & registro de fallos, nº 24 \\
\hyperref[nota:4]{Nota 4} & registro de fallos, nº 25 \\
\hyperref[nota:5]{Nota 5} & registro de fallos, nº 27 \\
\hyperref[nota:6]{Nota 6} & registro de fallos, nº 51 \\
\hyperref[nota:7]{Nota 7} & registro de fallos, nº 61 \\
\hyperref[nota:8]{Nota 8} & registro de fallos, nº 78 \\
\hyperref[nota:9]{Nota 9} & registro de fallos, nº 82 \\
\hyperref[nota:10]{Nota 10} & registro de fallos, nº 92 \\
\hyperref[nota:11]{Nota 11} & registro de fallos, nº 113 \\
\hyperref[nota:12]{Nota 12} & registro de fallos, nº 116 \\
\hyperref[nota:13]{Nota 13} & registro de fallos, nº 117 \\
\hyperref[nota:14]{Nota 14} & registro de fallos, nº 141 \\
\hyperref[nota:15]{Nota 15} & registro de fallos, nº 143 \\
\hyperref[nota:16]{Nota 16} & registro de fallos, nº 148 \\
\hyperref[nota:17]{Nota 17} & registro de fallos, nº 156 \\
\hyperref[nota:18]{Nota 18} & registro de fallos, nº 164 \\
\hyperref[nota:19]{Nota 19} & registro de fallos, nº 165 \\
\hyperref[nota:20]{Nota 20} & registro de fallos, nº 238 \\
\hyperref[nota:21]{Nota 21} & registro de fallos, nº 242 \\
\hyperref[nota:22]{Nota 22} & plan de cambios, M2 \\
\hyperref[nota:23]{Nota 23} & plan de cambios, M3 \\
\hyperref[nota:24]{Nota 24} & plan de cambios, M12 \\
\hyperref[nota:25]{Nota 25} & plan de cambios, M14 \\
\hyperref[nota:26]{Nota 26} & plan de cambios, M15 \\
\hyperref[nota:27]{Nota 27} & plan de cambios, pot§6.3 \\
\hyperref[nota:28]{Nota 28} & hallazgo nuevo, visto junto al \hyperref[err:10]{error~10} \\
\hyperref[nota:29]{Nota 29} & registro de fallos, nº 79 \\
\hyperref[nota:30]{Nota 30} & registro de fallos, nº 149 \\
\hyperref[nota:31]{Nota 31} & registro de fallos, nº 159 \\
\hyperref[nota:32]{Nota 32} & registro de fallos, nº 186 \\
\hyperref[nota:33]{Nota 33} & registro de fallos, nº 236 \\
\hyperref[nota:34]{Nota 34} & registro de fallos, nº 26 \\
\hyperref[nota:35]{Nota 35} & registro de fallos, nº 30 \\
\hyperref[nota:36]{Nota 36} & registro de fallos, nº 36; plan de cambios, M5 \\
\hyperref[nota:37]{Nota 37} & registro de fallos, nº 44 \\
\hyperref[nota:38]{Nota 38} & registro de fallos, nº 50 \\
\hyperref[nota:39]{Nota 39} & registro de fallos, nº 52 \\
\hyperref[nota:40]{Nota 40} & registro de fallos, nº 53 \\
\hyperref[nota:41]{Nota 41} & registro de fallos, nº 55 \\
\hyperref[nota:42]{Nota 42} & registro de fallos, nº 56 \\
\hyperref[nota:43]{Nota 43} & registro de fallos, nº 57 \\
\hyperref[nota:44]{Nota 44} & registro de fallos, nº 58 \\
\hyperref[nota:45]{Nota 45} & registro de fallos, nº 59 \\
\hyperref[nota:46]{Nota 46} & registro de fallos, nº 60 \\
\hyperref[nota:47]{Nota 47} & registro de fallos, nº 64 \\
\hyperref[nota:48]{Nota 48} & registro de fallos, nº 66 \\
\hyperref[nota:49]{Nota 49} & registro de fallos, nº 67 \\
\hyperref[nota:50]{Nota 50} & registro de fallos, nº 214 \\
\hyperref[nota:51]{Nota 51} & registro de fallos, nº 222 \\
\hyperref[nota:52]{Nota 52} & registro de fallos, nº 223 \\
\hyperref[nota:53]{Nota 53} & registro de fallos, nº 80 \\
\hyperref[nota:54]{Nota 54} & registro de fallos, nº 91 \\
\hyperref[nota:55]{Nota 55} & registro de fallos, nº 97 \\
\hyperref[nota:56]{Nota 56} & registro de fallos, nº 98 \\
\hyperref[nota:57]{Nota 57} & registro de fallos, nº 99 \\
\hyperref[nota:58]{Nota 58} & registro de fallos, nº 101 \\
\hyperref[nota:59]{Nota 59} & registro de fallos, nº 103 \\
\hyperref[nota:60]{Nota 60} & registro de fallos, nº 108 \\
\hyperref[nota:61]{Nota 61} & registro de fallos, nº 109 \\
\hyperref[nota:62]{Nota 62} & registro de fallos, nº 110 \\
\hyperref[nota:63]{Nota 63} & registro de fallos, nº 111 \\
\hyperref[nota:64]{Nota 64} & registro de fallos, nº 129 \\
\hyperref[nota:65]{Nota 65} & registro de fallos, nº 130 \\
\hyperref[nota:66]{Nota 66} & registro de fallos, nº 132 \\
\hyperref[nota:67]{Nota 67} & registro de fallos, nº 140 \\
\hyperref[nota:68]{Nota 68} & registro de fallos, nº 158 \\
\hyperref[nota:69]{Nota 69} & registro de fallos, nº 160 \\
\hyperref[nota:70]{Nota 70} & registro de fallos, nº 161 \\
\hyperref[nota:71]{Nota 71} & registro de fallos, nº 163 \\
\hyperref[nota:72]{Nota 72} & registro de fallos, nº 169 \\
\hyperref[nota:73]{Nota 73} & registro de fallos, nº 170 \\
\hyperref[nota:74]{Nota 74} & registro de fallos, nº 176 \\
\hyperref[nota:75]{Nota 75} & registro de fallos, nº 190 \\
\hyperref[nota:76]{Nota 76} & registro de fallos, nº 194 \\
\hyperref[nota:77]{Nota 77} & registro de fallos, nº 212 \\
\hyperref[nota:78]{Nota 78} & registro de fallos, nº 243 \\
\hyperref[nota:79]{Nota 79} & plan de cambios, G8; plan de cambios, F1-n \\
\hyperref[nota:80]{Nota 80} & plan de cambios, M4 \\
\hyperref[nota:81]{Nota 81} & plan de cambios, M5; registro de fallos, nº 36 \\
\hyperref[nota:82]{Nota 82} & plan de cambios, M6 \\
\hyperref[nota:83]{Nota 83} & plan de cambios, M7 \\
\hyperref[nota:84]{Nota 84} & plan de cambios, M8 \\
\hyperref[nota:85]{Nota 85} & plan de cambios, M9 \\
\hyperref[nota:86]{Nota 86} & plan de cambios, M10 \\
\hyperref[nota:87]{Nota 87} & plan de cambios, M13 \\
\hyperref[nota:88]{Nota 88} & plan de cambios, R6 \\
\hyperref[nota:89]{Nota 89} & plan de cambios, R7; plan de cambios, F3-f \\
\hyperref[nota:90]{Nota 90} & plan de cambios, F8-c \\
\hyperref[nota:91]{Nota 91} & plan de cambios, E10 \\
\hyperref[nota:92]{Nota 92} & plan de cambios, RP22 \\
\hyperref[nota:93]{Nota 93} & plan de cambios, RP23 \\
\hyperref[nota:94]{Nota 94} & plan de cambios, RP26 \\
\hyperref[nota:95]{Nota 95} & plan de cambios, RP27 \\
\hyperref[nota:96]{Nota 96} & plan de cambios, RP28 \\
\hyperref[nota:97]{Nota 97} & plan de cambios, RP3 \\
\hyperref[nota:98]{Nota 98} & plan de cambios, F8-c; plan de cambios, RP-nota1 \\
\hyperref[nota:99]{Nota 99} & plan de cambios, RP-nota2 \\
\hyperref[nota:100]{Nota 100} & plan de cambios, RP-nota3 \\
\hyperref[nota:101]{Nota 101} & plan de cambios, RP-nota4 \\
\hyperref[nota:102]{Nota 102} & plan de cambios, pot§6.4 \\
\hyperref[nota:103]{Nota 103} & plan de cambios, F5-b \\
\hyperref[nota:104]{Nota 104} & plan de cambios, F5-b \\
\hyperref[nota:105]{Nota 105} & plan de cambios, F5-c \\
\hyperref[nota:106]{Nota 106} & plan de cambios, F6-c \\
\hyperref[nota:107]{Nota 107} & registro de fallos, nº 234 \\
\hyperref[nota:108]{Nota 108} & hallazgo nuevo, visto junto a la \hyperref[nota:168]{nota~168} \\
\hyperref[nota:109]{Nota 109} & hallazgo nuevo, visto junto a la \hyperref[nota:168]{nota~168} \\
\hyperref[nota:110]{Nota 110} & hallazgo nuevo, visto junto al \hyperref[err:10]{error~10} \\
\hyperref[nota:111]{Nota 111} & hallazgo nuevo, visto junto a la \hyperref[nota:58]{nota~58} \\
\hyperref[nota:112]{Nota 112} & registro de fallos, nº 215 \\
\hyperref[nota:113]{Nota 113} & registro de fallos, nº 216 \\
\hyperref[nota:114]{Nota 114} & registro de fallos, nº 217 \\
\hyperref[nota:115]{Nota 115} & registro de fallos, nº 218 \\
\hyperref[nota:116]{Nota 116} & registro de fallos, nº 219 \\
\hyperref[nota:117]{Nota 117} & registro de fallos, nº 220 \\
\hyperref[nota:118]{Nota 118} & registro de fallos, nº 221 \\
\hyperref[nota:119]{Nota 119} & registro de fallos, nº 114 \\
\hyperref[nota:120]{Nota 120} & registro de fallos, nº 124 \\
\hyperref[nota:121]{Nota 121} & registro de fallos, nº 133 \\
\hyperref[nota:122]{Nota 122} & registro de fallos, nº 135 \\
\hyperref[nota:123]{Nota 123} & registro de fallos, nº 139 \\
\hyperref[nota:124]{Nota 124} & registro de fallos, nº 146 \\
\hyperref[nota:125]{Nota 125} & registro de fallos, nº 154 \\
\hyperref[nota:126]{Nota 126} & registro de fallos, nº 157 \\
\hyperref[nota:127]{Nota 127} & registro de fallos, nº 164 \\
\hyperref[nota:128]{Nota 128} & registro de fallos, nº 181 \\
\hyperref[nota:129]{Nota 129} & registro de fallos, nº 183 \\
\hyperref[nota:130]{Nota 130} & registro de fallos, nº 187 \\
\hyperref[nota:131]{Nota 131} & registro de fallos, nº 191 \\
\hyperref[nota:132]{Nota 132} & registro de fallos, nº 208 \\
\hyperref[nota:133]{Nota 133} & registro de fallos, nº 209 \\
\hyperref[nota:134]{Nota 134} & registro de fallos, nº 211 \\
\hyperref[nota:135]{Nota 135} & registro de fallos, nº 229 \\
\hyperref[nota:136]{Nota 136} & registro de fallos, nº 230 \\
\hyperref[nota:137]{Nota 137} & registro de fallos, nº 231 \\
\hyperref[nota:138]{Nota 138} & registro de fallos, nº 233 \\
\hyperref[nota:139]{Nota 139} & registro de fallos, nº 235 \\
\hyperref[nota:140]{Nota 140} & registro de fallos, nº 237 \\
\hyperref[nota:141]{Nota 141} & plan de cambios, RP21 \\
\hyperref[nota:142]{Nota 142} & plan de cambios, D4 \\
\hyperref[nota:143]{Nota 143} & plan de cambios, R2 \\
\hyperref[nota:144]{Nota 144} & plan de cambios, R4 \\
\hyperref[nota:145]{Nota 145} & plan de cambios, R5 \\
\hyperref[nota:146]{Nota 146} & plan de cambios, C1 \\
\hyperref[nota:147]{Nota 147} & plan de cambios, C3 \\
\hyperref[nota:148]{Nota 148} & plan de cambios, C4 \\
\hyperref[nota:149]{Nota 149} & plan de cambios, D1 \\
\hyperref[nota:150]{Nota 150} & plan de cambios, D5 \\
\hyperref[nota:151]{Nota 151} & plan de cambios, D6 \\
\hyperref[nota:152]{Nota 152} & plan de cambios, D7 \\
\hyperref[nota:153]{Nota 153} & plan de cambios, F3-a \\
\hyperref[nota:154]{Nota 154} & plan de cambios, F4 \\
\hyperref[nota:155]{Nota 155} & plan de cambios, P1 \\
\hyperref[nota:156]{Nota 156} & plan de cambios, P2 \\
\hyperref[nota:157]{Nota 157} & plan de cambios, P3 \\
\hyperref[nota:158]{Nota 158} & plan de cambios, P4 \\
\hyperref[nota:159]{Nota 159} & plan de cambios, P5 \\
\hyperref[nota:160]{Nota 160} & plan de cambios, F2-c \\
\hyperref[nota:161]{Nota 161} & plan de cambios, RP30 \\
\hyperref[nota:162]{Nota 162} & plan de cambios, F8-c \\
\hyperref[nota:163]{Nota 163} & plan de cambios, F8-b \\
\hyperref[nota:164]{Nota 164} & plan de cambios, F8-f \\
\hyperref[nota:165]{Nota 165} & plan de cambios, F8-a \\
\hyperref[nota:166]{Nota 166} & plan de cambios, F8-h \\
\hyperref[nota:167]{Nota 167} & plan de cambios, F5-c \\
\hyperref[nota:168]{Nota 168} & plan de cambios, F5-d \\
\hyperref[nota:169]{Nota 169} & plan de cambios, F6-a \\
\hyperref[nota:170]{Nota 170} & plan de cambios, F6-b \\
\hyperref[nota:171]{Nota 171} & plan de cambios, F6-c \\
\hyperref[nota:172]{Nota 172} & plan de cambios, F6-d \\
\hyperref[nota:173]{Nota 173} & plan de cambios, F6-e \\
\hyperref[nota:174]{Nota 174} & plan de cambios, F6-f \\
\hyperref[nota:175]{Nota 175} & plan de cambios, F6-g \\
\hyperref[nota:176]{Nota 176} & plan de cambios, F7-a \\
\hyperref[nota:177]{Nota 177} & plan de cambios, F7-c \\
\hyperref[nota:178]{Nota 178} & plan de cambios, F7-d \\
\hyperref[nota:179]{Nota 179} & plan de cambios, F7-e \\
\hyperref[nota:180]{Nota 180} & plan de cambios, F7-f \\
\hyperref[nota:181]{Nota 181} & plan de cambios, F8-e \\
\hyperref[nota:182]{Nota 182} & plan de cambios, F8-f \\
\hyperref[nota:183]{Nota 183} & plan de cambios, F8-g \\
\hyperref[nota:185]{Nota 185} & hallazgo nuevo, visto junto al \hyperref[err:1]{error~1} \\
\hyperref[nota:186]{Nota 186} & hallazgo nuevo, visto junto al \hyperref[err:1]{error~1} \\
\bottomrule
\end{longtable}
}

\end{document}
